\documentclass[11pt,a4paper]{article}
\pdfmapfile{+euler.map}
\pdfmapfile{+rsfs.map}
\usepackage[T1]{fontenc}
\usepackage[utf8]{inputenc}
\usepackage{lmodern,microtype}
\usepackage[margin=25mm]{geometry}
\usepackage{amsmath,amssymb,amsthm,mathtools,mathrsfs}
\usepackage{enumitem,booktabs,array,xurl,needspace,longtable,capt-of}
\usepackage{xcolor}
\usepackage[colorlinks=true,linkcolor=blue!45!black,citecolor=blue!45!black,urlcolor=blue!45!black]{hyperref}
\allowdisplaybreaks[2]
\setlength{\emergencystretch}{2em}
\numberwithin{equation}{section}
\newtheorem{theorem}{Theorem}[section]
\newtheorem{proposition}[theorem]{Proposition}
\newtheorem{lemma}[theorem]{Lemma}
\newtheorem{corollary}[theorem]{Corollary}
\newtheorem{question}[theorem]{Question}
\theoremstyle{definition}
\newtheorem{definition}[theorem]{Definition}
\theoremstyle{remark}
\newtheorem{remark}[theorem]{Remark}
\newcommand{\C}{\mathcal C}
\newcommand{\D}{\mathcal D}
\newcommand{\E}{\mathcal E}
\newcommand{\M}{\mathcal M}
\newcommand{\Z}{\mathcal Z}
\newcommand{\Pcat}{\mathcal P}
\newcommand{\Qcat}{\mathcal Q}
\newcommand{\one}{\mathbf1}
\providecommand{\Vec}{}
\renewcommand{\Vec}{\mathrm{Vec}}
\newcommand{\Rep}{\operatorname{Rep}}
\newcommand{\FPdim}{\operatorname{FPdim}}
\newcommand{\Irr}{\operatorname{Irr}}
\newcommand{\Hom}{\operatorname{Hom}}
\newcommand{\End}{\operatorname{End}}
\newcommand{\Inv}{\operatorname{Inv}}
\newcommand{\rank}{\operatorname{rank}}
\newcommand{\ad}{\mathrm{ad}}
\newcommand{\pt}{\mathrm{pt}}
\newcommand{\rev}{\mathrm{rev}}
\newcommand{\Gal}{\operatorname{Gal}}
\newcommand{\hs}{\widehat s}
\newcommand{\K}{K_0}
\newcommand{\kk}{\mathbb C}
\newcommand{\R}{\mathbb R}
\newcommand{\ZZ}{\mathbb Z}
\newcommand{\QQ}{\mathbb Q}
\newcommand{\CC}{\mathbb C}
\newcommand{\Tr}{\operatorname{Tr}}
\newcommand{\lcm}{\operatorname{lcm}}
\newcommand{\sVec}{\mathrm{sVec}}
\DeclareMathOperator{\ord}{ord}
\newcommand{\OO}{\mathcal O}
\newcommand{\cost}{\operatorname{cost}}
\newcommand{\WGT}{\mathrm{WGT}}
\setcounter{tocdepth}{2}

\theoremstyle{plain}
\newtheorem{maintheorem}{Theorem}
\renewcommand{\themaintheorem}{\Alph{maintheorem}}
\theoremstyle{plain}
\newtheorem{proofclaim}[theorem]{Claim}
\newtheorem{conditionaltheorem}[theorem]{Conditional Theorem}
\newtheorem{conditionalproposition}[theorem]{Conditional Proposition}
\newtheorem{conditionalcorollary}[theorem]{Conditional Corollary}
\newtheorem{conditionalmaintheorem}[maintheorem]{Conditional Theorem}

\hypersetup{pdftitle={Arithmetic restrictions on braided fusion categories of dimension p6 q6 r4},pdfauthor={},bookmarksdepth=2}
\title{Arithmetic restrictions on braided fusion categories\\
of dimension \(p^6q^6r^4\)}
\author{}
\date{}
\begin{document}
\maketitle
\begin{abstract}
Let \(p,q,r\) be distinct primes. We study weak group-theoreticity for
braided fusion categories whose dimensions divide \(p^6q^6r^4\).
A minimal obstruction is integral, perfect, fusion-simple, and modular,
with exact dimension \(p^6q^6r^4\) and the quartic exponent on the largest
prime. The sole remaining even dimension is \(5400^2\). In the odd part
of this family, the smallest parameter not excluded by our uniform
arithmetic criteria is \(5^6 7^6 23^4\).
Our methods combine character conductors, primary modular operators,
and positive cyclotomic words. Complementary residue blocks give a
criterion for coprime braided decomposition. At \(5400^2\), every pair
of dimension primes occurs in an actual twist of such an obstruction,
and simple objects of dimension six are excluded. For a simple integral source of dimension \(5400\) with simple
integral Drinfeld center, we develop canonical-algebra and trace-form
restrictions. The stronger length and source-rank bounds are stated
conditionally on complete exact verification of the specified finite
exclusions; the compact companion supplies bounded checks rather than
the large certificate archives. For a fusion ring with exactly
one nonlinear component, of degree two, its full aligned linear
character table and duality determine the ring up to its opposite.
Cayley--Hamilton identities give local coefficient tests before
reconstruction.
\end{abstract}

\tableofcontents

\section{Introduction}\label{paper:introduction}

Weakly group-theoretical fusion categories are the categories
categorically Morita equivalent to nilpotent fusion categories. Their
stability under extensions, equivariantizations, and Morita equivalence
makes dimension a natural organizing invariant
\cite{ENOsolvable,NataleWGT}. The question of whether weak integrality
forces weak group-theoreticity leads to a more concrete problem:
which arithmetic configurations could support a minimal counterexample?
For braided categories, modular character theory and local congruence
representations impose restrictions that are substantially stronger than
the dimension equation alone.

This paper studies the three-prime family \(p^6q^6r^4\), including all
assignments of the exponents to the primes. Its main conclusions have
three different scopes. First, dimension-family arguments reduce a
possible braided obstruction to an ordered exact modular case and
determine its even and odd arithmetic frontiers. Second, the exceptional
even dimension \(5400^2\) admits strong restrictions on actual twists,
character fields, and small tensor products. Third, if a category at
that dimension is a Drinfeld center, the canonical algebra relates its
modular data to the possibly noncommutative fusion ring of a source of
dimension \(5400\). The center hypothesis is used only in this third
part.

\subsection{The dimension-family theorem}

A fusion category is \emph{perfect} if the unit is its only invertible
simple object, and \emph{fusion-simple} if it has no nontrivial proper
fusion subcategory. The following statement combines the family
reduction with its numerical consequences.

\begin{maintheorem}[The family and its first residual dimensions]
\label{paper:main:family}
Let \(p,q,r\) be distinct primes.
\begin{enumerate}[label=\textup{(\roman*)},leftmargin=*]
\item If a braided fusion category of dimension dividing
\(p^6q^6r^4\) is not weakly group-theoretical, a divisor-minimal
obstruction is integral, perfect, fusion-simple, and modular. Its
dimension is exactly \(p^6q^6r^4\), and after interchanging the two
sextic primes one has \(p<q<r\).
\item Every braided fusion category of even dimension in the family is
weakly group-theoretical, except possibly at dimension
\[
       2^6 3^6 5^4=5400^2=29\,160\,000.
\]
\item Every braided fusion category of odd dimension in the family below
\[
       D_{\mathrm{odd}}
       =5^6 7^6 23^4
       =514\,422\,090\,765\,625
       =22\,680\,875^2
\]
is weakly group-theoretical. The parameter \((5,7,23)\) is not
excluded by the uniform scalar, short-support, and combined symbolic
criteria proved here.
\end{enumerate}
\end{maintheorem}

The ordering in part \textup{(i)} is a conclusion: the exponent four
cannot be assigned to either smaller prime in a minimal obstruction.
The reduction also disposes of the relevant proper divisors and the
slightly degenerate alternative. The only residual \emph{even} dimension
is \(5400^2\); odd parameters remain. The number \(D_{\mathrm{odd}}\)
is the smallest dimension in this family not excluded by those uniform
arithmetic criteria. This conclusion is restricted to the specified
dimension family.

The family proof appears in
Section~\ref{paper:family:consequences}. The even-family exclusion and
the short-support estimates are stated with their attribution to
\cite{BDPfamily}. The numerical odd frontier follows from the uniform
criterion and an explicit cyclic difference-set witness. No existence
claim is attached to either residual dimension.

The odd conclusion is part of a uniform restriction on the prime parameters.
Any remaining ordered odd triple satisfies
\[
 5\le p<q<r,\qquad 3q\le\min(r,p^2),\qquad
 q^{3/2}<r<(p^2/q)^3,\qquad q^3<p^4.
\]
The whole sector with least prime seven is excluded. With least prime five,
the only parameter left by these criteria and the short-support bounds is
$(5,7,23)$; see Corollaries~\ref{paper:family:cor:remaining-region},
\ref{paper:family:family-criterion}, and~\ref{paper:family:first-odd}.

\subsection{Arithmetic and structural mechanisms}

For a modular category of dimension \(D\), let \(d_X\) be a simple
dimension and
\[
             \chi_Y(X)=S_{XY}/d_Y,\qquad f_Y=D/d_Y^2
\]
its character and formal codegree. Here \(S\) is unnormalized,
so \(S_{\mathbf1,Y}=d_Y\). The arithmetic of the full character field
is essential: the field of one entry does not determine the field
of the corresponding character.

The first mechanism is local integrality. For a cyclotomic
character field \(K\) with rational formal codegree \(f\), the
different divides \(f\mathcal O_K\), and the conductor divides
\(2f\). Symmetric integral modular data sharpen the conductor bound
to \(f\). Residue blocks and prime-primary Weyl operators then link
formal-codegree valuations to actual twists and force complementary
principal witnesses.

The second mechanism is positivity. Balancing expresses
\[
          \frac{S_{XY}}{\operatorname{lcm}(d_X,d_Y)}
\]
as a positive rational word in roots of unity of mass
\(\gcd(d_X,d_Y)\). Integrality, prescribed residues, and strict character
bounds constrain its support. Exposed supports propagate to tensor
adjoints; difference sets and Galois translates then convert categorical
positivity into finite scalar and cyclic-mask obstructions. This yields
both uniform prime inequalities and the finite tests at the odd frontier.

The third mechanism concerns actual categorical decomposition.
For a set \(J\) of dimension primes, let \(\mathscr P_J\) be the join
of the full residue-block partitions, and let \(P_J\) be its principal
component. A join is a common coarsening, not an intersection of blocks.

\begin{maintheorem}[A canonical coprime splitting criterion]
\label{paper:main:splitting}
Let \(\mathcal M\) be an integral modular category of dimension \(D\).
Partition its dimension primes into two nonempty sets \(J,K\), and put
\[
 D_J=\prod_{\ell\in J}\ell^{v_\ell(D)},\qquad D_K=D/D_J.
\]
The condition \(P_J\cap P_K=\{\mathbf1\}\) is equivalent to a braided
decomposition
\[
       \mathcal M\simeq\mathcal M_J\boxtimes\mathcal M_K,
       \qquad
       \FPdim(\mathcal M_J)=D_J,\quad
       \FPdim(\mathcal M_K)=D_K.
\]
The factors are the subcategories on the principal joined components
and are each other's M\"uger centralizers.
\end{maintheorem}

The proof uses a stronger linear statement: every complementary
modular corner \(s_{B,C}\), with \(s=S/\sqrt D\), is a partial
isometry of rank \(|B\cap C|\). Integral block masses and positive
products provide additional constraints when the principal
intersection is larger. These results distinguish spectral invariance
from fusion closure and isolate the extra implication needed for a
categorical contradiction.

\subsection{The exceptional even dimension}

The arithmetic at \(5400^2\) is unusually rigid, but the remaining
obstruction is not eliminated.

\begin{maintheorem}[Twists and small dimensions at \(5400^2\)]
\label{paper:main:modular}
Let \(\mathcal M\) be an integral, perfect, fusion-simple modular
category of dimension \(5400^2\).
\begin{enumerate}[label=\textup{(\roman*)},leftmargin=*]
\item Each of \(6,10,15\) divides the order of an actual simple twist.
\item Complementary principal witnesses can be chosen with dimension
pair \((12,20)\) or \((36,20)\). Their normalized crossing is either
\[
 1+3\xi
 \quad\text{or}\quad
 -\sum_{t=1}^{4}\eta_t\xi^t,
 \qquad
 \{\eta_t\}=\{1,1,\omega,\omega^2\},
\]
where \(\xi\) and \(\omega\) have orders five and three.
Every signed crossing forces an actual twist whose order is divisible
by thirty.
\item No simple object has dimension six.
\end{enumerate}
\end{maintheorem}

The three pairwise twist conditions do not assert that a single
twist is divisible by thirty. This leaves the ordinary crossing with
or without such a twist, and the signed tame crossing. The endpoint-field
and short-adjoint restrictions in
Section~\ref{paper:modular:section} constrain
these alternatives without assuming a center realization.

\paragraph{A combined rank consequence.}
Theorem~\ref{paper:main:modular}\textup{(i)} and the full-graph rows of
\cite[Theorem~29.1]{BDPfamily} give $\rank(\mathcal M)\ge106$, or
$\rank(\mathcal M)\ge130$ when no simple object has dimension $900$.
If an actual twist order is divisible by $30$, the respective bounds
are $112$ and $142$. Indeed, the missing-$6$ and missing-$10$
configurations are excluded. For the two remaining configurations, the
degree budgets proved in that reference give
\[
 900^2R\ge A,\qquad
 600^2R+95625k\ge B\quad\text{if }m_{900}=0,
\]
where $R=\rank(\mathcal M)$, $k=m_{675}\le7$, and
\[
 (A,B)=
 \begin{cases}
 (85\,373\,451,47\,123\,451),&\text{no twist order divisible by }30,\\
 (89\,941\,851,51\,691\,851),&\text{some twist order divisible by }30.
 \end{cases}
\]
Here $m_d$ denotes the number of simples of dimension $d$.
Taking ceilings gives the stated bounds. This consequence combines the
earlier full-graph degree budgets with the new missing-edge exclusions;
it does not use the later source enumerations.

\subsection{Verification status of finite exclusions}
\label{paper:verification-status}
This research draft separates structural proofs from large finite
exclusions. The distributed companion is intentionally small: it contains
bounded exact checks, their outputs, a theorem-status map and preserved
source code, but no large field-state or negative-certificate archives.
The bounded checks authenticate only their explicitly stated scopes.
In particular they do not establish complete rank-thirteen or rank-fifteen
exclusions, the length-twelve exclusion, or the rank-sixteen catalogue.

For an enumeration-dependent statement, \emph{verification hypothesis
$\mathsf V$ for its declared scope} means that its complete necessary
input universe has been independently regenerated, every field and
involution alternative and every admissible row in the stated ranges
has been covered, and every claimed exclusion has a witness checked by
exact arithmetic on that same universe. Overlapping exclusions require
exact identifier unions and differences, and all branch certificates
must close every branch. The mathematical necessity of each filter
remains an obligation of its cited proof; a successful program run
alone does not supply it. This is a separate finite-verification
hypothesis, not a claim that the missing checks have been performed here.

Statements labelled \emph{Conditional} assume $\mathsf V$ for the
finite screens, certificate ledgers and prerequisite exclusions specified
in their proofs. Their proofs explain the implication from those checks.
The accompanying numerical tables and their downstream uses have the
same conditional status. The small regression
\path{checks/check_certificate_scope.py} and its record
\path{reference/certificate_scope.json} reject a witness invalidated
by an added row for each certificate type. This checks the corrected
rule on explicit inputs; it does not verify the new fields, involutions,
repeated roots, dimension ranges or quaternion alternatives of a large
historical catalogue. They are retained as precise research targets,
not as independently reproducible global conclusions of this compact
release. The standalone fixed-space exclusion in
Theorem~\ref{paper:symplectic-sixty-exclusion} does not require this
hypothesis; its separate global uniqueness assertion does.

\subsection{Simple sources of dimension \(5400\)}

Now suppose that \(\mathcal M\simeq\mathcal Z(\mathcal C)\), where
\(\mathcal C\) is integral and fusion-simple of dimension
\(n=5400\), and \(\mathcal M\) is integral and fusion-simple.
This is an additional hypothesis. Let \(I\) be the right adjoint
to the forgetful functor, and write
\[
 I(\mathbf1)=\bigoplus_E m_E Z_E,\qquad
 L=\sum_E m_E,\qquad r=\rank(\mathcal C).
\]
The index \(E\) runs through the irreducible representations of
\(K_0(\mathcal C)\otimes_{\mathbb Z}\mathbb C\); one has
\(m_E=\dim E\), \(\FPdim(Z_E)=n/f_E\), and
\(r=\sum_E m_E^2\). Source dimensions, canonical dimensions, and
representation degrees are thus different quantities.

\begin{conditionalmaintheorem}[Canonical length and source ranks]
\label{paper:main:source}
Assume the preceding center hypotheses and verification hypothesis
$\mathsf V$ for the exclusions and catalogues in
Theorem~\ref{paper:source:global-bounds} and
Theorem~\ref{paper:source:length-thirteen}. Then
\[
 L\ge13,\qquad
 \begin{cases}
 14\le r\le462,&K_0(\mathcal C)\text{ commutative},\\
 16\le r\le490,&K_0(\mathcal C)\text{ noncommutative}.
 \end{cases}
\]
If \(L=13\), then
\[
 r\in\{17,19,21,23,25,27,29\},\qquad \dim E\le4.
\]
At noncommutative rank sixteen,
\[
 K_0(\mathcal C)\otimes_{\mathbb Z}\mathbb C
       \simeq\mathbb C^{12}\oplus M_2(\mathbb C),
 \qquad L=14.
\]
\end{conditionalmaintheorem}

The conditional deduction combines the character--center correspondence,
positive trace forms, Galois-stable canonical packets, source-row moment
identities, and the specified complete finite exclusions. The compact
release does not discharge the latter hypothesis. The smaller conclusion
$L\ge11$ is proved with the bounded check in
Theorem~\ref{paper:canonical-length-at-least-eleven}. A single
certificate principle organizes the rank and length exclusions.
The source upper bounds and initial dimension restrictions are used
with their attribution to \cite[Section~30]{BDPfamily}.

The structural part goes beyond these numerical bounds. Integral
commutator tensors bound the noncentral support independently of source
rank. Symplectic duality forces \(f_E/4\) to be an algebraic integer
for each individually preserved degree-two component when all nonlinear
components have that form. If the complexified fusion ring has exactly
one nonlinear component, of degree two, the full aligned linear
character table and basis duality determine the ring up to passage to
its opposite. The Cayley--Hamilton formulas
of Section~\ref{paper:ch:sec:local-coefficients} recover square coefficients
from two rows, before a full table is available.

Under $\mathsf V$, the specified noncommutative rank-sixteen screening
would leave \(600\,558\) necessary field-and-duality states: \(599\,070\) orthogonal and
\(1\,488\) symplectic. These are exact arithmetic inputs to further
reconstruction, not fusion rings. The commutative rank-fourteen and noncommutative rank-sixteen cases
remain open even under this verification hypothesis; without it the
stronger claimed exclusions below those ranks are not established by
the compact release.

\subsection{Organization and verification}

The paper develops character arithmetic and cyclotomic words first,
then proves the family theorem and the general block-decomposition
criterion. It next treats the exceptional modular dimension and the
additional center problem. The latter is organized by its mathematical
mechanisms---canonical spectra, trace forms, commutators, local corners,
and finite certificates.

The accompanying \path{verification/README.md} distinguishes freshly
rerun bounded checks from historical source code and the unavailable
large ledgers. The file \path{verification/theorem_status.json} maps
numerical claims to their verification scope. Complete large coverage is
not claimed for the compact release. An unfinished search supplies no
exclusion.
Verification of a fusion ring's multiplication table does not establish
categorification or simplicity of a Drinfeld center.

\subsection{Provenance and relation to the broader draft}
\label{paper:provenance}
The broader research draft \cite{BDPfamily} is retained separately; this
family-focused manuscript does not supersede it. In particular, its
global dimension and general rank bounds, families involving more
primes, positive reduced degree-six exclusion and arithmetic/action-graph
factorization criteria remain in that work. The versioned locators below
identify shared arguments and prerequisites, not independent validation
of them. ``Additional'' means additional relative to that draft; it is
not a claim of priority over the entire literature.

{\small
\setlength{\tabcolsep}{4pt}
\begin{longtable}{@{}p{.30\textwidth}p{.30\textwidth}p{.34\textwidth}@{}}
\toprule
Results here & Earlier location in \cite{BDPfamily} & Relationship and scope\\
\midrule\endhead
Character conductors and residue-block invariance
(Sections~\ref{paper:family:arithmetic}--\ref{paper:sec:local-action}) &
Sections 2--3 and 6 &
Inherited trace-dual, rotation and local-Weyl mechanisms, with expanded
proofs.\\[3pt]
Theorem~\ref{paper:main:family}: proper-divisor, even and first odd
frontiers & Section 25.8; Theorems 26.2--26.3; Theorem 25.9 and
Corollary 25.10 & Inherited numerical conclusions. The frontiers
$5400^2$ and $5^6 7^6 23^4$ are unchanged.\\[3pt]
Exact compression, phase allocation and dyadic model
(Section~\ref{paper:sec:local-action}) & Proposition 24.20 and
Lemma 26.18 & Strengthened projection and channel formulas; the explicit
level-eight model supplies additional bounds.\\[3pt]
Prime-level and odd-support refinements
(Sections~\ref{paper:sec:words} and~\ref{paper:sec:family-criterion}) &
Propositions 8.40 and 25.14; Section 27 &
The zero-slack mask obstruction and cyclic-quotient tests are inherited;
positive-slack rigidity and fuller phase/field restrictions strengthen them.\\[3pt]
Theorem~\ref{paper:main:splitting} and complementary block corners &
Sections 3, 24 and 31 & Additional exact partial-isometry and canonical
coprime-splitting criterion; complementary to the earlier graph criteria.\\[3pt]
Theorem~\ref{paper:main:modular}(i)--(ii) &
Theorems 26.14--26.15 and 29.1 &
The crossing-word mechanism is inherited. Excluding both missing-edge
cases strengthens the actual-twist alternatives.\\[3pt]
Theorem~\ref{paper:main:modular}(iii), degree six &
Theorem 28.1 & Same ordinary-endpoint conclusion with a different
short-adjoint/Galois-overlap proof; the older reduced-case result has
wider scope.\\[3pt]
Spectral annihilators and finite moment windows & Lemma 16.5 and
Section 30 & The factorial projector method is inherited;
actual annihilators and finite exact integrality tests sharpen it.\\[3pt]
Character--center bridge and initial source bounds & Section 30 &
Inherited, ultimately using Ostrik's correspondence. The upper bounds
462 and 490 are unchanged.\\[3pt]
Canonical algebra, trace forms, quaternion alternatives and local corners
(Sections~\ref{paper:source}--\ref{paper:source:corners}) &
Initial source constraints in Section 30 & Additional systematic
constraints; all field, involution and center hypotheses remain essential.\\[3pt]
Commutator, one-/two-matrix reconstruction and symplectic codegrees
(Section~\ref{paper:source:commutators}) & No corresponding general
package & Additional fusion-ring theory; aligned labelled characters and
the stated component hypotheses are required.\\[3pt]
Conditional Theorem~\ref{paper:main:source} and residual catalogues &
Section 30: $L\ge10$, source ranks at least 11/13 &
Stronger proposed bounds and catalogues. They require $\mathsf V$;
small bounded checks do not certify the complete large exclusions.\\
\bottomrule
\end{longtable}}

\paragraph{Conventions.}
All categories are over \(\mathbb C\). Weakly integral categories carry
their positive spherical structures. The notation \(D\) denotes the
dimension of a modular category; in the center part \(D=n^2\) and
\(n=5400\) is the source dimension. Characters are evaluated on the
distinguished simple basis, and duality is always an involution on
basis labels. An \(\ell\)-active simple has a nontrivial
\(\ell\)-primary component of its actual twist. Actual twists are
distinguished from scalar-normalized diagonal entries of a congruence
representation. A positive cyclotomic word has nonnegative rational
coefficients; the word is called integral when those coefficients are
integers. Its mass is their sum.


\section{Character arithmetic and local orders}
\label{paper:family:arithmetic}\label{paper:sec:character-orders}\label{paper:sec:2}

All categories are defined over $\mathbb C$. We use the positive spherical
structure on a weakly integral fusion category, so categorical and
Frobenius--Perron dimensions agree. A category is \emph{perfect} if the
unit is its only invertible simple, and \emph{fusion-simple} if it has no
nontrivial proper fusion subcategory. Weak group-theoreticity is stable
under the reductions in the next lemma. This allows character arithmetic
to be applied to a category with strict character bounds.
We write $\mathcal C_{\pt},\mathcal C_{\ad},\mathcal C'$ and
$U(\mathcal C)$ for the pointed subcategory, adjoint subcategory,
M\"uger center and universal grading group.

\subsection{Minimal counterexamples and orthogonality}

\begin{lemma}[Minimal counterexample]\label{paper:family:2.1}
Suppose every braided category of dimension a proper divisor of $D$ is WGT. A counterexample $\C$ at $D$ is integral, equals $\C_{\ad}$, and has one of two forms:
\begin{enumerate}[label=\textup{(\roman*)},leftmargin=*]
\item $\C$ is modular, perfect and fusion-simple;
\item $\C'=\C_{\pt}=\sVec$, and the only fusion subcategories are $\Vec,\sVec,\C$. The ring modulo the transparent fermion is perfect and simple.
\end{enumerate}
\end{lemma}
\begin{proof}
A nontrivial universal grading reduces the question to its neutral component. A proper nontrivial nondegenerate subcategory splits off with its centralizer. A nontrivial Tannakian $\Rep(G)$ reduces to the neutral de-equivariantization, of dimension $D/(|G||H|)$ for grading support $H\le G$. In each case the smaller dimensions divide $D$ and WGT passes back by extension, equivariantization or tensor product \cite[Propositions~8.15 and~8.20]{ENOannals}, \cite[Proposition~4.1]{ENOsolvable}, \cite[Theorems~3.13(i) and~4.44, Proposition~4.26]{DGNO}, \cite[Proposition~4.1]{NataleWGT}. Thus $\C=\C_{\ad}$ is integral by the dimension grading \cite[Theorem~3.10]{GNnil}, and its symmetric center is $\Vec$ or $\sVec$. Every proper subcategory is integral and WGT, hence pointed in the absence of Tannakian subcategories \cite[Theorem~3.1]{NataleCore}. It thus lies in $\C_{\pt}$, which is trivial in the modular case by $(\C_{\ad})'=\C_{\pt}$, and is $\sVec$ in the slightly degenerate case by \cite[Lemma~7.1]{NataleWGT}. In the latter case pull back based subrings from the fermion quotient.
\end{proof}

We use \emph{ordinary modular table} for the character table of an integral
modular category, and \emph{positive reduced table} for the table of an
integral slightly degenerate category after identifying its transparent
fermion with the unit. A \emph{minimal table} is one of the two categorical
tables in Lemma~\ref{paper:family:2.1}, with the simplicity and perfectness
conditions stated there. In particular, its twists, traces and indicators
come from the category. Statements formulated for based rings or symmetric
integral character data alone use the algebraic hypotheses specified in
those statements.

Put $\Delta=D$ in case (i) and $\Delta=D/2$ in case (ii). In (ii) use the fusion ring $R=K_0(\C)/(f-\one)$ for the transparent fermion $f$. The twist gives $fX\not\simeq X,X^*$. Duality-compatible orbit representatives containing $\one$ give structure constants $N_{XY}^Z+N_{XY}^{fZ}$ and unit coefficient $\delta_{Y,X^*}$. These two rings are instances of the following definition.

\begin{definition}[Symmetric integral character data]\label{paper:family:def:symmetric-integral-data}
Let $R$ be a commutative fusion ring with basis $[X]$, duality involution and unit $0$. Its trace $\tau$ is the coefficient of the unit; $\tau([X][Y]^*)=\delta_{XY}$ makes $a\mapsto\tau(aa^*)$ positive definite when $*$ is extended conjugate-linearly to the complex algebra. Fix a complete bijection from basis labels $Y$ to the irreducible characters $\chi_Y$ and require $\chi_0(X)=d_X\in\mathbb Z_{>0}$. Put $\Delta=\sum_Xd_X^2$. We require that all $\chi_Y(X)$ lie in the integer ring of a common cyclotomic field, that $S_{XY}=d_Y\chi_Y(X)$ is symmetric, and that
\begin{equation}\label{paper:family:eq:wgt-modular-data}
 \begin{gathered}S_{\one X}=d_X,\quad S^\dagger S=\Delta I,\quad S_{X,Y^*}=\overline{S_{XY}},\\
 \chi_Y(X)=S_{XY}/d_Y,\quad f_Y=\|\chi_Y\|^2=\Delta/d_Y^2\in\ZZ.\end{gathered}
\end{equation}
The involution and orthogonality above give the coefficient trace
\begin{align}
 \tau(B)&=\sum_Y\chi_Y(B)/f_Y,\label{paper:family:eq:fusion-coefficient-trace}\\
 f_Ye_Y&=\sum_X\overline{\chi_Y(X)}[X],\qquad
 \sum_X\chi_Y(X)\overline{\chi_Z(X)}=\delta_{YZ}f_Y.
 \label{paper:family:lw:character-identities}
\end{align}
Galois conjugation acts on the complete set of characters. Since it commutes with duality and fixes the rational $f_Y$, the induced permutation $\pi_\sigma$ preserves $d_Y$. Symmetry therefore gives the unsigned covariance
\begin{align}
 \sigma S_{XW}&=S_{\pi_\sigma X,W}=S_{X,\pi_\sigma W},
                  \label{paper:family:eq:galois-permutations}\\
 \sigma\chi_X(Y)&=\chi_X(\pi_\sigma Y),\qquad
 d_{\pi_\sigma Y}=d_Y.\label{paper:family:lw:argument-galois}
\end{align}
\end{definition}
We use $I_s=\{X:d_X=s\}$ and $n_s=|I_s|$. Twists, the congruence representation and categorical indicators are additional categorical inputs; they are introduced explicitly in the local arguments below.
For the modular table these identities are standard \cite[Sections~8.13--8.14]{EGNO}. For the reduced table, the normalized columns give all characters by \cite[Proposition~3.3]{YuSlightlyDegenerate} (see \cite[Proposition~2.7]{BGNPWSuper} for the unitary case). The ratio $\Delta/d_Y^2$ is an algebraic integer by \cite[Corollary~3.4]{YuSlightlyDegenerate}, hence an integer here.


\begin{proposition}[LCM energy and prime companions]\label{paper:family:2.2}
Let $\C$ be integral braided, with its positive ribbon structure, and let $S$
be its unnormalized matrix. For $X$ simple set
$h_X=\sum_{T\in\Irr(\C')}N_{XX^*}^{T}d_T$. Then
\[
 \sum_{Y:S_{XY}\ne0}\lcm(d_X,d_Y)^2\le D h_X.
\]
If $X$ centralizes $\C_{\pt}$, is not transparent, and $p\mid d_X$ while
$p\nmid\FPdim(\C_{\pt})$, there is a noninvertible $Y$ with
$p\nmid d_Y$, $S_{XY}\ne0$, and
$\FPdim(\C_{\pt})d_X^2+\lcm(d_X,d_Y)^2\le Dh_X$.
\end{proposition}
\begin{proof}
The Hopf-link characters $a_Y(X)=S_{XY}/d_Y$ are ring characters
\cite[Proposition~8.13.11]{EGNO}. B\'ezout applied to $a_Y(X)$ and $a_X(Y)$ shows
$S_{XY}/\lcm(d_X,d_Y)$ is integral, as in the modular criterion of
Dong--Palcoux \cite[Lemma~2.7 and the proof of Theorem~1.1]{DPLCM}. The regular element
$R=\sum_Y d_Y[Y]$ satisfies $[T]R=d_T R$. Accordingly $a_T(R)$ is $D$
when $a_T=d$ and zero otherwise. In the positive structure $a_T=d$
means precisely $T\in\C'$ by the double-braiding criterion
\cite[Proposition~8.20.5(i)]{EGNO}. Symmetry and $a_Y(X^*)=\overline{a_Y(X)}$ give
\[
 \sum_Y |S_{XY}|^2
  =\sum_Y d_Y^2a_Y(XX^*)=\sum_T N_{XX^*}^{T}d_Ta_T(R)=Dh_X.
\]
Average this rational equality over a cyclotomic splitting field. Every
nonzero integral normalized entry has mean squared modulus at least one
by its norm. This proves the first assertion. For the second,
$a_X(R)=0$ and the invertibles contribute $\FPdim(\C_{\pt})$.
Reduction at a prime above $p$ forces a noninvertible $p$-free hit;
retain it and all invertibles in the first inequality.
\end{proof}
In a modular category $h_X=1$. In a slightly degenerate category the same
holds for the full table: $gX\not\simeq X$ excludes the transparent fermion
from $XX^*$. Passage to the reduced table halves the full energy sum
and the number of invertible columns.

\begin{lemma}[Averaging over a character orbit]\label{paper:family:2.3}\label{paper:family:lem:trace-identities}
In ordinary or positive reduced symmetric integral data, let $H$ be a subgroup of the cyclotomic Galois group, $e=|H\cdot Z|$ and $A_H(x)=|H|^{-1}\sum_{\sigma\in H}\sigma(x)$. Then
\[
  f_Z/e=\sum_X |A_H(\chi_Z(X))|^2,\qquad
  0=\sum_X d_X A_H(\chi_Z(X))\quad (Z\ne\one).
\]
For the full Galois group write $K_X=\QQ(\chi_X)$, $F_X=\QQ(\chi_Z(X))$ and $T_X=\Tr_{F_X/\QQ}(\chi_Z(X))$. Grouping columns into orbits gives
\begin{align}
 f_Z&=\sum_{[X]}[K_Z:F_X][K_X:F_X]T_X^2,\label{paper:family:eq:r8-tracenorm}\\
 0&=\sum_{[X]}d_X[K_X:F_X]T_X\qquad(Z\ne\one).\label{paper:family:eq:r8-tracedim}
\end{align}
For any $H$, let $A$ be a set of nonunits whose entries $z_X=S_{XZ}/\lcm(d_X,d_Z)$ are fixed by $H$. With $\mathsf M(z)$ the mean of $|\sigma z|^2$ over all conjugates, we have
\[
 e\left(d_Z^2+\sum_{X\in A}\lcm(d_X,d_Z)^2\mathsf M(z_X)\right)\le\Delta .
\]
\end{lemma}
\begin{proof}
The $e$ conjugates of the row $\chi_Z$ are orthogonal with squared norm $f_Z$, and their average at $X$ is $A_H(\chi_Z(X))$ by \eqref{paper:family:eq:galois-permutations}. Averaging preserves orthogonality to the dimension character. For the full group, the value at each of the $[K_X:\QQ]$ conjugates of $X$ is $T_X/[F_X:\QQ]$. For the inequality multiply the norm identity by $d_Z^2$, retain the vacuum and $A$, then take the full Galois mean. The Galois group is abelian, so the conjugates of these entries are also fixed by $H$.
\end{proof}
Write $GZ$, $I_p Z$, $I_p^{\rm tame}Z$ for the full, inertia and chosen tame inertia orbits in the splitting field; at odd $p$ the tame subgroup is the Teichm\"uller subgroup of the primary factor. Conjugates of an entry $z$ are indexed by $\operatorname{Gal}(\QQ(z)/\QQ)$, independently of the orbit of its row label. Thus $A_H$ is a normalized trace, $T_X$ an unnormalized field trace, and $\mathsf M$ a mean squared modulus. In particular, for pairwise disjoint nonvacuum sets $A_i$ of $H$-fixed entries and lower bounds $c_i\le\lcm(d_X,d_Z)^2\mathsf M(z_X)$ for $X\in A_i$, the last inequality reads
\begin{equation}\label{paper:eq:disjoint-energy}
       |HZ|\left(d_Z^2+\sum_i|A_i|c_i\right)\le\Delta .
\end{equation}
This applies to subsets of columns, even when several columns have the same entry. Two choices of $H$ will recur. For one nonunit $X$, take $H$ fixing $L=\QQ(z_X)$ and $A$ equal to the full Galois orbit of $X$. Then $e=[K_Z:L]$ and the sum is $[K_X:\QQ]\lcm(d_X,d_Z)^2\mathsf M(z_X)$. For inertia at $s$, all columns with $K_X$ unramified at $s$ have fixed entries. If $Z$ is principal at $t\ne s$ and $t\nmid d_X$, then $\chi_Z(X)\equiv d_X$ above $t$ gives $z_X\ne0$, hence $\mathsf M(z_X)\ge1$ by the algebraic norm bound. This applies simultaneously to any set of such columns.

\begin{lemma}[Zeros and square-fullness]\label{paper:family:2.4}
In perfect symmetric integral data, every nonunit row contains a zero and $\Delta$ is square-full. For $s=d_X,t=d_Y$, $S_{XY}/\lcm(s,t)$ is an algebraic integer. In a simple perfect table $S_{XY}=0$ whenever $X,Y$ are nonunit and $(s,t)=1$; in particular no nonunit dimension is a prime power.
If $\Delta$ is odd then duality has no nonunit fixed point.
\end{lemma}
\begin{proof}
Apply the B\'ezout and Galois-mean bounds of Proposition~\ref{paper:family:2.2} to either character table. If a nonunit row of dimension $s$ had no zero, averaging its norm would give
\[
 \Delta/s^2\ \ge\ \sum_Y (d_Y/\gcd(s,d_Y))^2
 \ >\ \sum_Y d_Y^2/s^2=\Delta/s^2;
\]
the strict inequality holds already at $\one$. For ordinary modular categories the row-vanishing theorem is due to
Gelaki, Naidu and Nikshych \cite[Theorem~6.1]{GNNcenters}.
Square-fullness in the perfect integral modular case is
\cite[Corollary~10.10]{BP}, and the coprime-degree and prime-power
obstructions in the integral modular case follow from
Etingof, Nikshych and Ostrik \cite[Lemma~7.1 and Corollary~7.2]{ENOsolvable}. For the fusion-ring vanishing theorem see \cite[Theorem~1.2]{BP}. Now $P=\prod_X XX^*$ vanishes on every non-Perron character and hence, by \eqref{paper:family:eq:fusion-coefficient-trace}, $\tau(P)=\prod_X d_X^2/\Delta$ is an integer. Since $d_X^2\mid\Delta$, no prime can occur only once in $\Delta$.

For a non-Perron character, the labels at which $|\chi(X)|=d_X$ form a based subring: equality in the triangle inequality for a product propagates to each constituent, and the set is closed under duality. Simplicity therefore makes the inequality strict off the unit, since equality everywhere would force codegree $\Delta$. For coprime $s,t>1$, all conjugates of $S_{XY}/(st)$ have modulus strictly below one; its norm forces it to be zero. On the other hand a nonunit row satisfies
\begin{equation}\label{paper:family:eq:reduced-unit-orthogonality}
 1+\sum_{Y\ne\one}d_Y\chi_X(Y)=0.
\end{equation}
Modulo any prime dividing $\Delta$, the row must meet a nonunit label of degree prime to that prime. A prime-power row cannot do so.

In odd dimension all $1/f_X$ reduce to one at two; hence \eqref{paper:family:eq:fusion-coefficient-trace} gives $\tau(B)\equiv\Tr(N_B)\pmod2$ for $B\in R$. Apply this congruence to $B=X^2$ and then to $B=X$, using $\Tr(N_X^2)\equiv\Tr(N_X)\pmod2$. It gives $\delta_{X,X^*}\equiv\delta_{X,\one}\pmod2$.
\end{proof}

\begin{lemma}[Small ordinary codegrees]\label{paper:lem:small-ordinary-codegrees}
In a nontrivial simple perfect integral modular category every formal codegree
is at least five.
\end{lemma}
\begin{proof}
Use the unnormalized modular matrix $S$ with $S_{0Y}=d_Y$.
For a nonunit $X$ put $n=d_X$, $f=D/n^2$ and
$u_Y=S_{XY}/n$. Write $H=\{Y\ne\one:u_Y\ne0\}$.
Orthogonality gives $\sum_{Y\in H} d_Yu_Y=-1$ and
$\sum_{Y\in H}|u_Y|^2=f-1$. Thus the degrees in $H$ have gcd one.
The $(X,0)$ entry of the modular relation $STS=p^+T^{-1}ST^{-1}$,
where $|p^+|=\sqrt D$, gives
\[
             \sqrt f\,n\le 1+\sum_{Y\in H}d_Y|u_Y|.
\]
Symmetry and algebraic integrality of the character values give
$u_Y\in(d_Y/\gcd(d_Y,n))\mathcal O$, where $\mathcal O$ denotes the
algebraic integers. Averaging the squared-norm identity over the
cyclotomic Galois group, each nonzero term contributes at least
$(d_Y/\gcd(d_Y,n))^2$, by the arithmetic--geometric mean inequality
applied to its conjugates. If $f\le4$, we conclude that
$|H|\le f-1$ and every $d_Y$ in $H$ divides $n$.

Perfection and the gcd condition require at least two terms. If there
are two, both degrees are proper divisors of $n$. Cauchy's inequality
then gives $\sqrt f\,n\le1+(n/2)\sqrt{2(f-1)}$, impossible for
$f=3,4$ and $n\ge2$. For three terms we have $f=4$, and all three
entries are roots of unity: each has mean square one, so all its
conjugates have modulus one. At most one of the degrees can be $n$,
by the gcd condition. If one is $n$, the others cannot both equal
$n/2$ since $n/2>1$. Ordering the proper divisors therefore gives
$2n\le1+11n/6$. This forces $n\le6$. A nonunit degree in our category has at least two distinct prime factors
by Lemma~\ref{paper:family:2.4}; hence $n=6$, and its only divisor
admissible as a nonunit degree is $6$, contradicting the gcd condition.
The vacuum codegree is $D$ and is larger than four, since such a
category contains a nonunit of degree at least six.
\end{proof}

\subsection{Fields and residue classes of characters}\label{paper:sec:2.1}
For an irreducible representation $E$ of a fusion ring $R$, put $K=\QQ(\chi_E(X):X)$ and let $e_E$ be its central primitive idempotent. The formal codegree $f_E$ is the Schur element of $\tau$ \cite[Proposition~1.15]{ChlSchur}. Thus $\dim(E) f_E=\sum_X|\chi_E(X)|^2$ and \cite[Lemma~2.3]{OstrikFormalCodegrees}
\begin{equation}\label{paper:family:eq:character-formal-codegree-general}
 \sum_X\chi_E(X^*)[X]=f_Ee_E .
\end{equation}
For a categorified ring $K$ is cyclotomic in the sense of lying in a cyclotomic field \cite[Theorem~8.51]{ENOannals}; the reduced characters also lift.

\begin{theorem}[The different divides a rational formal codegree]\label{paper:family:2.5}
For an irreducible representation of a categorified fusion ring, or of a fusion ring with cyclotomic character values, a rational formal codegree $f$ is a positive integer. The different of its character field satisfies
\begin{equation}\label{paper:family:eq:character-different-divisibility}
 \mathfrak D_K\mid f\mathcal O_K .
\end{equation}
Only primes dividing $f$ ramify.
\end{theorem}
\begin{proof}
For $a\in\mathfrak D_K^{-1}$ the element $\sum_X\Tr_{K/\QQ}(a\chi_E(X^*))[X]$ belongs to $R$. By \eqref{paper:family:eq:character-formal-codegree-general} it equals
$f\sum_{\sigma\in\Gal(K/\QQ)}\sigma(a)e_{\sigma E}$, with distinct conjugate characters. Its eigenvalue $af$ is integral. Thus $f\mathfrak D_K^{-1}\subseteq\mathcal O_K$, and $a=1$ gives the integrality of $f$.

\end{proof}
\begin{corollary}[Conductor of a character field]\label{paper:family:2.6}
If a cyclotomic character field $K$ has rational formal codegree $f$,
its conductor divides $2f$.
\end{corollary}
\begin{proof}
Work after a finite unramified extension of $\QQ_p$, so the ramified
part $E$ of $K$ lies in the totally ramified cyclotomic extension
$L$ of minimal conductor $p^m$. Let $e=[E:F]$, where $F$ is the
unramified base, and let $d_E$ be the different exponent at $E$.
Theorem~\ref{paper:family:2.5} gives $d_E/e\le v_p(f)$.
The derivative of $\Phi_{p^m}$ gives
$d_L/[L:F]=m-1/(p-1)$.
For odd $p$ and $e=e_0p^s>1$, $e_0\mid p-1$, cyclic wild inertia
gives $m=s+1$ and $L/E$ is tame. Thus
$(d_E+1)/e=(d_L+1)/[L:F]$, or
\begin{equation}\label{paper:family:eq:odd-normalized-different}
 \delta_p(e):=d_E/e=s+1-\frac{1-p^{-s}}{p-1}-\frac1e .
\end{equation}
This also gives zero at $e=1$. For $s=0$, $e>1$ it is positive,
and for $s\ge1$ it lies strictly between $s$ and $s+1$.
Hence $p^m\mid f$.
At two, conductor four gives normalized different one.
For $m\ge3$, minimality excludes $1+2^{m-1}$ from
$\Gal(L/E)$. Every nontrivial subgroup of $\langle5\rangle$
contains that element, so $[L:E]\le2$. The possible quadratic
generators are $-1$ and $2^{m-1}-1$. For either generator $\sigma$,
a primitive $2^m$th root $\zeta$ satisfies
$v_L(\zeta-\sigma\zeta)=2$; this is the relative different exponent.
Transitivity therefore gives
$d_E/e\ge m-1-2/2^{m-1}>m-2$, whence $2^{m-1}\mid f$.
\end{proof}
By symmetry the field of $S_{XY}/\lcm(d_X,d_Y)$ lies in both character
fields; its different divides both codegrees by transitivity, and the next theorem sharpens its conductor bound. In particular, the field of a character of $p$-power codegree lies
in a $p$-power cyclotomic field by Kronecker--Weber and the conductor bound
(for odd $p$, adjoining a factor of two to an odd level does not change the field).
Total ramification there gives a unique prime over $p$ in the character field
\cite[Lemma~3.1.13 and Theorem~7.5.5]{Sharifi}.

\begin{theorem}[Conductor at an individual label]\label{paper:family:por:th:exact-conductor}
In symmetric integral character data (Definition~\ref{paper:family:def:symmetric-integral-data}), write $K_Z$ for the field generated by the character $\chi_Z$. Then the conductor of $K_Z$ divides $f_Z$. In particular the conductor of the lcm-normalized $XY$ crossing divides $\gcd(f_X,f_Y)$.
\end{theorem}
\begin{proof}
The different bound (Corollary~\ref{paper:family:2.6}) already puts $K_Z$ in a cyclotomic field of conductor dividing $2f_Z$. It remains to remove its last possible dyadic level. Put $m=2^{v_2(f_Z)}$. For $m=1$, adjoining a factor of two to an odd cyclotomic level does not change the field, so there is nothing to check. Take $L=F(\zeta)$ with $F$ an odd cyclotomic field large enough and $\zeta$ primitive of order $2m$. Suppose the dyadic conductor is $2m$. In $G=(\mathbb Z/2m)^\times$ the stabilizer $H$ of the row does not contain $e=1+m$, since the fixed field of $e$ is $F(\zeta^2)$.

Write $b_Y=\chi_Z(Y)=\sum_{k=0}^{m-1}z_{kY}\zeta^k$. The coefficients belong to $\mathcal O_F$ by coprimality of conductors. The orthogonality of the $m$ odd powers of $\zeta$ gives
\[
 z_{1Y}=m^{-1}\sum_{g\in G}\zeta^{-g}\chi_{gZ}(Y).
\]
Different rows in this sum are orthogonal; the squared norm of each is $f_Z$. Thus
\[
 \sum_Y|z_{1Y}|^2=\frac{f_Z}{m^2}\sum_{h\in H}\sum_{g\in G}\zeta^{g(h-1)}
                 =\frac{f_Z}{m}.
\]
Indeed the inner geometric sum vanishes unless $h=1$ or $h=e$, and only $1$ belongs to $H$. Extend $e$ to the common splitting field, trivially on $F$; the extension need not have order two. Argument-side covariance and uniqueness of the integral expansion give
\[
                         z_{1,\pi_eY}=-z_{1Y}.
\]
Every permutation cycle carrying a nonzero coefficient has even
length. Its squared-norm contribution belongs to $2\mathcal O_F$.
Thus $f_Z/m\in2\mathcal O_F\cap\mathbb Z=2\mathbb Z$,
contrary to its oddness. Symmetry gives the crossing assertion.
\end{proof}

\paragraph{Local projectors.}
Fix a prime $\mathfrak p$ of a splitting field over $p$. The full residue classes (or $p$-blocks) group the characters of the commutative order $R$ with equal reductions on every basis element. Over a complete unramified coefficient ring $\mathcal O$ containing the residues in its residue field, the special-fiber idempotents lift. Hence for a union $L$ of full blocks, $e_L=\sum_{X\in L}e_X$ lies in $R\otimes\mathcal O$, and for locally integral $B$,
\begin{equation}\label{paper:family:eq:local}
 \tau(e_LB)=\sum_{X\in L}\chi_X(B)/f_X\in\mathcal O .
\end{equation}

The character transform $\Phi(B)=(\chi_z(B))_z$ sends $e_z$ to the coordinate vector $v_z$ and intertwines multiplication by $e_L$ with $P_L=\sum_{z\in L}v_zv_z^*$, the coordinate projection. We keep the algebra element $e_L$ in trace identities and the operator $P_L$ in matrix identities.

\begin{lemma}[Principal residues at every prime]\label{paper:family:lem:principal-all-places}
In the data of Definition~\ref{paper:family:def:symmetric-integral-data}, if $Z$ is principal at one place over $p$, then $\chi_Z(X)\equiv d_X$ at every place over $p$ for every $X$. The congruence may be divided by a rational $p$-unit. In particular, a locally integral algebraic number $u$ which vanishes at all places above $p$ belongs to $p\mathcal O_{K,(p)}$ if its field $K$ is unramified at $p$.
\end{lemma}
\begin{proof}
For any $\sigma$, $\sigma\chi_Z(X)=\chi_Z(\pi_\sigma X)$ and $d_{\pi_\sigma X}=d_X$, so the congruence survives transport to each conjugate place. The last assertion follows from $p\mathcal O_{K,(p)}=\prod_{\mathfrak p\mid p}\mathfrak p\mathcal O_{K,(p)}$.
\end{proof}

\begin{theorem}[Support in the principal residue class]\label{paper:family:2.7}
In symmetric integral character data, for any label set $\mathcal A$ put
$L=\{X:\chi_X(W)\equiv d_W\pmod{\mathfrak p}\ \text{for every } W\in\mathcal A\}$.
If $f_Z=p^a$, the support of $\chi_Z$ lies in $L$; it equals $L$ when $\mathcal A=I_{d_Z}$ and $p\nmid d_Z$. For any $Z$ and any Galois-stable $\mathcal A$, $\sum_{X\in L}|\chi_Z(X)|^2$ is an integer divisible by $p^{v_p(f_Z)}$.
\end{theorem}
\begin{proof}
Write $\alpha$ for the sum. Symmetry gives $\alpha/f_Z=\tau(e_LZZ^*)\in\mathcal O$. If $\mathcal A$ is Galois-stable, so is $L$ at the fixed prime $\mathfrak p$: by \eqref{paper:family:eq:galois-permutations}, permuting $X$ permutes the tests on $W$ and preserves $d_W$. Thus $\alpha\in\ZZ$ with the asserted divisibility. Apply this to the set of \emph{all} labels. Its principal block contains the unit, so $1\le\alpha\le f_Z=p^a$ forces $\alpha=f_Z$. The support lies in this block and hence in every $L$. Conversely for $\mathcal A=I_{d_Z}$, $X\in L$ gives $\chi_X(Z)\equiv d_Z$, nonzero if $p\nmid d_Z$.
\end{proof}

\paragraph{Blocks of defect one.}
After a sufficiently large finite unramified base change with fraction
field $F$, integer ring $\mathcal O$ and residue field $k$, the field factors of a block
correspond to its full inertia orbits and are totally ramified.
An orbit is \emph{tame} when wild inertia fixes it pointwise.
The following observation explains why a single character of small
defect can determine the entire block.

\begin{lemma}[A block with a conductor of length one]\label{paper:2.9}
Let $A$ be a commutative local symmetric $\mathcal O$-order with
normalization $M=\prod_i\mathcal O_i$, where all residue fields
are $k$. Suppose its perfect trace on factor $i$ is
$\Tr_{E_i/F}/f_i$, and put $c_i=e_i v_p(f_i)-d_i$, where $e_i,d_i$
are its ramification index and different exponent. A factor with
$c_i=0$ is the whole block and $A=\mathcal O_i$. If $c_i=1$,
there are precisely two factors, both with $c=1$, and
$A=\{(x,y)\in M:\bar x=\bar y\}$.
\end{lemma}
\begin{proof}
Since $A$ is self-dual, the conductor $(A:M)$ is the trace dual
$M^\#=\prod_i f_i\mathfrak D_i^{-1}$. A unit conductor on one
factor puts its idempotent in the local ring $A$. If $c_i=1$,
the conductor contains its maximal ideal $\mathfrak m_i$ on that
factor alone. Writing $B$ for the projection to the other factors,
we have $A=\mathcal O_i\times_k B$: common residues are necessary,
and the copy of $\mathfrak m_i$ lifts every compatible pair.
There must be another factor, since otherwise $A=\mathcal O_i$
has conductor exponent zero. The perfect trace identifies $A$
with $\Hom_{\mathcal O}(A,\mathcal O)$, so $A$ is Gorenstein.
The fiber-product criterion \cite[Fact~2.6 and Corollary~2.7]{NSTV}
makes $B$ a discrete valuation ring. The conductor of this
two-factor order is the product of the two maximal ideals.
\end{proof}

\begin{proposition}[The block of a defect-one character]\label{paper:2.10}
In a local symmetric order as above, with cyclotomic characters and
rational integral codegrees, suppose a character has defect
$v_p(f)=1$. Except for a single ramified quadratic branch at two, the
block has two inertia orbits and is the fiber product of their valuation
rings over $k$. In particular no unramified higher-defect character
belongs to this block.

For odd $p$ both orbits are tame and of defect one. At $p=3$ their
lengths in symmetric integral character data are one and two.
At two, if the given branch is unramified and the other has degree
$e$, the block has rank $1+e$ and the other defect is $1+v_2(e)$.
Consequently uniform defect one at two forces rank two; in an
even-rank dyadic block of rank greater than two all defects are at
least two.
\end{proposition}
\begin{proof}
At odd $p$ the given branch is tame by
Corollary~\ref{paper:family:2.6}, so $e_i-d_i=1$; an unramified
defect-one branch at two has the same conductor length. By
Lemma~\ref{paper:2.9} the other branch satisfies
$v_p(f)=(d_E+1)/e$, which gives defect one if it is unramified.
At odd $p$ the right side for wild degree $p^s>1$ lies strictly
between $s$ and $s+1$ by \eqref{paper:family:eq:odd-normalized-different};
hence both branches are tame of defect one. At three their lengths
are one or two. The integral identity trace
$\sum e_i/f_i=\sum e_i d_{X_i}^2/\Delta$ forces the sum of the
lengths to vanish modulo three, since after canceling the common
powers of three each unit square is one there.

At two a ramified defect-one branch satisfies $2(e_i-1)\le d_i\le e_i$
by the first two ramification groups and the different bound. Hence
$e_i=d_i=2$ and the conductor has length zero. In the two-branch
case the identity trace $1/f_1+e/f_2$ is integral; its first summand
has valuation $-1$, so the other does too. Since $e$ is a power of
two, even rank $1+e$ forces $e=1$. A defect-zero branch is
unramified with conductor length zero, hence has rank one by itself.
\end{proof}

\subsection{Principal heights and mixed defects}

\begin{corollary}[Principal valuation bound]\label{paper:family:4.1}
If $v_p(\Delta)\ge2$, every ordinary or reduced principal-block row
satisfies
\[
  v_p(f_X)\ge2,\qquad
  v_p(d_X)\le\lfloor(v_p(\Delta)-2)/2\rfloor .
\]
In particular the defect is at least three at odd valuation, and
at valuation two or three every principal degree is $p$-free.
\end{corollary}
\begin{proof}
The dimension character is unramified of defect $v_p(\Delta)$.
Proposition~\ref{paper:2.10} excludes defect one; defect zero has
an integral idempotent separating it from the dimension character.
Now use $d_X^2f_X=\Delta$.
\end{proof}

\begin{proposition}[A mixed-defect block bound]\label{paper:4.3}
Let $A$ be a commutative local symmetric order over an unramified
$p$-adic integer ring, with separable cyclotomic generic algebra,
rational positive integral codegrees, and a common character residue.
After an unramified extension write its generic factors as
$E_j/K$, $1\le j\le t$, with totally ramified degrees
$e_j=u_jp^{s_j}$, $u_j\mid p-1$. Put $a_j=v_p(f_j)$.
Then
\[
 \sum_{j=1}^t\bigl(p^{s_j}-(a_j-1)e_j\bigr)\le2.
\]
At two this becomes $\sum_j(2-a_j)e_j\le2$.
\end{proposition}
\begin{proof}
Put $M=\prod_j\OO_{E_j}$. The trace on $E_j$ is $\Tr_{E_j/K}/f_j$,
so the discriminant-index formula and self-duality of $A$ give
$2\operatorname{length}(M/A)=\sum_j(a_je_j-d_j)$.
Common residues impose $t-1$ independent conditions on $M$;
the first two ramification groups give $d_j\ge e_j+p^{s_j}-2$.
Combining yields the bound, with $e_j=2^{s_j}$ at two.
\end{proof}

\subsection{Residue energy and small codegrees}\label{paper:sec:residue-energy}

\label{paper:sec:residue-energy-2}
For a cyclotomic integer $x$ write $\mathsf M(x)$ for the mean of $|\sigma x|^2$ over its conjugates and put $E(x)=[\mathbb Q(x):\mathbb Q]\mathsf M(x)$.  These quantities measure the norm spent by one orbit of entries in a character row.

\begin{theorem}[Energy at a prescribed prime residue]\label{paper:5.3}
If a nonzero cyclotomic integer $x$ vanishes at all primes over $p$,
then $E(x)\ge2p$. For odd $p$ write the ramification index of
$\QQ(x)$ as $e=e_0p^s$, $e_0\mid p-1$. Then
\[
 \mathsf M(x)\ge W_p(e):=
 \begin{cases}
 p^2 & e=1,\\
 2p/e_0 &s=0,\ e_0>1,\\
 1+(p-1)/e_0 & s>0 .
 \end{cases}
\]
In particular $E(x)\ge p^2$ if $e=1$ and
$E(x)\ge p^{s+1}$ if $s>0$.

For odd $p$, in a $p$-power cyclotomic field a prescribed integer residue
$0\le k<p$ gives $\mathsf M(x)\ge k(p-k)/(p-1)$.
\end{theorem}
\begin{proof}
For $p=2$, vanishing at a single prime already gives $|N(x)|\ge2$, whence
$E(x)\ge h\,2^{2/h}\ge4$ for $h=[\QQ(x):\QQ]\ge1$. Now take $p$ odd.
When $e=1$, $x/p$ is integral and its nonzero norm gives
$\mathsf M(x)\ge p^2$.

In the tame case take $x\in F(\zeta_p)$ with $F$ cyclotomic of
conductor prime to $p$. In an integral expression
$x=\sum_{j\bmod p}a_j\zeta_p^j$, the sum of the $a_j$ vanishes at
every prime of $F$ over $p$, hence lies in $p\mathcal O_F$.
Subtract its average from every coefficient. This gives a unique
expression of coefficient sum zero, with
$\mathsf M(x)=\frac p{p-1}\sum_j\mathsf M(a_j)$.
Inertia fixing $x$ has order $g=(p-1)/e_0$ and its nonzero-index
orbits have length $g$. Each nonzero orbit of coefficients costs at
least $g$ by the norm. If there is just one such orbit with common
value $a$, then $a_0=-ga$, so the cost is at least $g+g^2\ge2g$.
This proves the tame bound.

For $s>0$ the inertia conductor is $p^{s+1}$. With
$\zeta=\zeta_{p^{s+1}}$ and $F$ as above, the relative integral
basis over $F(\zeta_p)$ gives a unique expansion
$x=\sum_{j=0}^{p^s-1}\zeta^j\gamma_j$. The terms are orthogonal
for the relative trace, hence $\mathsf M(x)=\sum_j\mathsf M(\gamma_j)$.
Inertia fixing $x$ has order $g=(p-1)/e_0$; it acts freely on
nonzero indices and preserves coefficient means, since the coefficients
are conjugate up to $p$th roots of unity. Some nonzero-index orbit
must occur by wild ramification. Two such orbits cost at least
$2g\ge g+1$, and one together with $\gamma_0\ne0$ costs $g+1$.
For a single orbit with $\gamma_0=0$, all its coefficients have the
same residue at each prime over $p$. Their residues sum to zero,
so each vanishes there since $p\nmid g$. The preceding tame
estimate (or the unramified estimate) gives mean at least
$2p/(p-1)$ for each. Their sum is at least $g+1$.
Finally $[\QQ(x):\QQ]\ge e$ gives the energy bounds.

For the prescribed residue in $\QQ(\zeta_p)$ choose integer
coefficients $c_j$ with $\sum c_j=k$. The same quadratic form gives
$\mathsf M(x)=(p\sum c_j^2-k^2)/(p-1)\ge k(p-k)/(p-1)$.
At higher level the orthogonal expansion above adds the means of its
coefficients. Their least residues sum to $k$ modulo $p$, and
$u(p-u)+v(p-v)\ge w(p-w)$ for $w\equiv u+v\bmod p$ in $[0,p)$.
Iteration gives the claim.
\end{proof}

\begin{corollary}[Conductor and row budget]\label{paper:5.5}
In a symmetric integral table fix an odd prime $p$ and put $c=v_p(\Delta)$ and
$t_X=v_p(f_X)=c-2v_p(d_X)$. A row field has zero $p$-conductor
exponent when $t_X=0$ and conductor exponent at most $t_X$
otherwise. The same holds for an entry field with $t_X$ replaced by
$t=\min(t_z,t_X)$. When $t\ge1$ its wild inertia has order at most
$p^{t-1}$.

Suppose the entries $\chi_z(X)$ on $p$-bearing degrees vanish
modulo $p$ (at every prime above $p$). This holds if the block of
$z$ contains a $p$-unit degree. It also holds if $p\mid d_z$ and
each nonzero entry on such a degree has $X\in P_p$ with
$v_p(d_X)\ge v_p(d_z)$. For each $p$-bearing degree $d$ let $u_d=(d/\gcd(d,d_z))_{p'}$.
Partition the nonzero entries $\chi_z(X)/u_d$ at degree $d$ into
Galois orbits $O$ counted with multiplicity. Write $h_O=|O|$,
$e_O=e_{0,O}p^{s_O}$ for the degree and ramification index of the
entry field. Then
\begin{equation}\label{paper:eq:merged-row-budget}
 \sum_{d:\,p\mid d}\ \sum_{O\text{ at }d}
                 u_d^2 h_O W_p(e_O)\ \le\ f_z-1 .
\end{equation}
\end{corollary}
\begin{proof}
The stated conductor bounds follow from Theorem~\ref{paper:family:por:th:exact-conductor},
applied to either row since their normalized entry lies in both fields.
The bound from the different is attained at conductor $p^t$:
the full cyclotomic field has normalized different $t-1/(p-1)$.

On a $p$-bearing degree $d$, every character of $p$-unit degree
vanishes above $p$ by lcm integrality. The common residue of its
block gives the same vanishing for $\chi_z$. In the second case use symmetry:
$\chi_z(X)=(d_X/d_z)\chi_X(z)$, with $\chi_X(z)\equiv d_z$
and $d_X/d_z$ integral at $p$. In both cases Galois covariance
gives vanishing on all conjugates. Division by the
integral lcm scale $u_d$, prime to $p$, preserves both integrality
and vanishing. Sum the energy bounds over the orbits and reserve
the unit's squared contribution one.
\end{proof}

\begin{corollary}[Small budgets and forced hits]\label{paper:5.6}
For a row satisfying this vanishing condition, if $f_z=p^2m$, $p\nmid m$, at most $m-1$
nonzero orbits on $p$-bearing degrees are unramified or wild in
total; the wild depth is at most one. For $m=1$ all such orbits are
nontrivially tame and number at most $(p-1)/2$.

For any prime $p$, suppose there are just three dimension primes
$p,q,r$, every nonunit degree involves at least two, and the
nonunit row $z$ has zero residue on its supported $p$-bearing degrees
at every prime above $p$. Then
\[
                          f_z\ge [K_z:\QQ]+2p^2 .
\]
\end{corollary}
\begin{proof}
For the first assertion divide the budget $p^2m-1$ by $p^2$ or by
$2p$ and use the conductor bound at $t_z=2$. For the second apply
\eqref{paper:family:eq:r8-tracedim} modulo $q$ and $r$. Its vacuum term is one,
so each reduction supplies an orbit of nonzero trace on a nonunit
degree prime to the respective prime. The two orbits are distinct
and both degrees are $p$-bearing. Each trace is a nonzero integer
divisible by $p$. Their contributions to \eqref{paper:family:eq:r8-tracenorm}
together with the vacuum give the asserted bound.
\end{proof}

\begin{proposition}[Degrees supporting a nonzero trace]\label{paper:5.7}
Let $Z\ne\one$ have formal codegree $p^2$ in a perfect ordinary or
positive reduced integral table. The nonunit degrees at which
$\chi_Z$ has nonzero absolute trace all divide
\[
                         M=d_Z/\gcd(p,d_Z)
\]
and have greatest common divisor one.
If the table is fusion-simple, then $M$ has at least three distinct
prime divisors; when it has exactly three, each pair is the prime
support of some such degree. In particular at a quartic prime of a
three-prime minimal table there is no row of codegree $p^2$, for any $p$.
\end{proposition}
\begin{proof}
By \eqref{paper:family:eq:r8-tracenorm} a nonzero trace $T_X$ at $X\ne\one$
is prime to $p$: a nonzero multiple of $p$ would exhaust the norm
before counting the unit. Now $d_X^2\mid\Delta=p^2d_Z^2$, so
$d_X/\gcd(d_X,d_Z)$ is one or $p$. In the latter case lcm
integrality makes $\chi_Z(X)$, and hence its trace, a multiple of
$p$. Thus $d_X\mid d_Z$. Moreover $X\in P_p$ by
Theorem~\ref{paper:family:2.7}. If
$v_p(d_X)=v_p(d_Z)>0$, then
$\chi_Z(X)=(d_X/d_Z)\chi_X(Z)$ vanishes above $p$ by principality.
Lemma~\ref{paper:family:lem:principal-all-places} gives this vanishing for every conjugate, hence $p\mid T_X$.
Hence $d_X\mid M$. The dimension identity \eqref{paper:family:eq:r8-tracedim}
expresses $-1$ as an integral combination of the indicated degrees.

In the fusion-simple case no nonunit has prime-power degree. For
each prime dividing $M$ some indicated degree omits it; this proves
both assertions about prime supports. Finally $v_p(\Delta)=4$
gives $v_p(d_Z)=1$, so with only two other dimension primes $M$
cannot have three.
\end{proof}


\begin{proposition}[The trace budget of a prime-power codegree]
\label{paper:primepower-trace}
Let $\mathcal C$ be an ordinary perfect fusion-simple integral
modular category with exactly three dimension primes $\ell,u,v$.
If a nonunit $Z$ has $f_Z=\ell^{2a}$, then
\[
 Z\in P_\ell,\qquad
 \ell^{2a}\ge [K_Z:\mathbb Q]+2\ell^2.
\]
In particular, no formal codegree is the square of a prime.

More precisely, put $v_\ell(D)=2m$, $h=m-a=v_\ell(d_Z)$,
and suppose $h>0$. If the character value $\chi_Z(Y)$ has
nonzero absolute trace and $\ell\mid d_Y$, then
\[
 a\ge2,\qquad v_\ell(d_Y)\le m-2.
\]
The nonzero-trace nonunit degrees have greatest common divisor one.
Thus when $m=3$ and $a=2$, there are nonzero-trace labels
of degrees $\ell u^i$ and $\ell v^j$.
\end{proposition}
\begin{proof}
At every dimension prime other than $\ell$, the full residue
block of $Z$ is its singleton, because the codegree is a unit.
Its complementary joined component is consequently $\{Z\}$.
The nonzero vacuum entry and the complementary-corner theorem
force $Z\in P_\ell$. This also follows directly by factoring
$s$ into its $\ell$-local Weyl operator and its complementary
factors: the latter preserve the line of $Z$, while the former
preserves the principal $\ell$-block.

Principality gives $\chi_Z(Y)\equiv d_Y$ at every place over
$\ell$. On every $\ell$-bearing degree the value therefore
vanishes at all these places, and its absolute trace is an
integer divisible by $\ell$. The trace-dimension identity
reduced at $u$ and at $v$ supplies two distinct nonzero-trace
nonunit orbits, respectively on $u$-free and $v$-free degrees.
Each such degree is $\ell$-bearing, since prime-power degrees
are excluded. Each orbit spends at least $\ell^2$ in the
trace-norm identity, and the vacuum spends $[K_Z:\mathbb Q]$.
This proves the budget and excludes $a=1$.

For the sharper assertion, write $j=v_\ell(d_Y)$. The row
support lies in $P_\ell$ by codegree support, so $Y\in P_\ell$
whenever its trace is nonzero. If $j\ge h$, symmetry gives
\[
 \chi_Z(Y)=\ell^{j-h}c\,\chi_Y(Z),
 \qquad v_\ell(c)=0.
\]
Here $c=d_Y/(\ell^{j-h}d_Z)$ is rational and a local unit.
Since $h>0$, principality of $Y$ makes $\chi_Y(Z)$ vanish at
all $\ell$-places. Its trace, after multiplication by $c$,
is in $\ell\mathbb Z_{(\ell)}$. Combining this with the
first principality argument when $j<h$ shows that the integer
absolute trace $T_Y$ satisfies
\[
 v_\ell(T_Y)\ge 1+\max\{j-h,0\}.
\]
The trace-norm identity reserves a positive vacuum term,
so $T_Y\ne0$ implies $|T_Y|<\ell^a$. Consequently $a\ge2$
and $j\le h+a-2=m-2$. The trace-dimension identity expresses
$-1$ as an integral combination of the nonzero-trace nonunit degrees,
so their gcd is one. At $m=3,a=2$, reducing at $u$ and $v$
therefore supplies the claimed degrees, of $\ell$-height one.
\end{proof}

\begin{proposition}[Character-field size at a prime-power codegree]
\label{paper:primepower-field-budget}
In Proposition~\ref{paper:primepower-trace}, suppose $\ell$ is odd
and put $e=[K_Z:\mathbb Q]$. Then
\[
                \ell^{2a}\ge\frac{3\ell}{\ell-1}\,e.
\]
If $D$ is odd and $Z$ is $\ell$-active, its full character field is a full
cyclotomic field $K_Z=\mathbb Q(\zeta_{\ell^b})$, with
\[
                            b\le2a-1.
\]
\end{proposition}
\begin{proof}
Every entry field $F=\mathbb Q(\chi_Z(Y))$ is $\ell$-primary.
Write $[F:\mathbb Q]=e_0\ell^s$, with $e_0\mid\ell-1$.
For $s\ge1$, the trace of every integer of $F$ is divisible
by $\ell^s$. To see this directly, embed $F$ in its minimal
cyclotomic field $L=\mathbb Q(\zeta_{\ell^{s+1}})$.
The traces of its power-basis elements are multiples of
$\ell^s$, while $[L:F]=(\ell-1)/e_0$ is prime to $\ell$.
For $s=0$, a value vanishing at the prime over $\ell$ has
trace divisible by $\ell$, including $F=\mathbb Q$.
Consequently a nonzero trace $T$ of such a value satisfies
\[
 \frac{T^2}{[F:\mathbb Q]}
 \ge \begin{cases}
       \ell^2/(\ell-1),&s=0,\\
       \ell^s/(\ell-1),&s\ge1,
      \end{cases}
 \qquad\text{and hence}\qquad
 \frac{T^2}{[F:\mathbb Q]}\ge\frac{\ell}{\ell-1}.
\]
The two opposite nonzero-trace orbits from the preceding
proposition therefore each cost at least
$e\ell/(\ell-1)$ in the trace norm.

Reduce the trace-dimension identity also modulo $\ell$.
All $\ell$-bearing terms vanish, so it supplies a further,
$\ell$-free nonunit orbit for which both its trace and its
relative column index are $\ell$-units. Its entry field is
therefore tame by the trace divisibility just proved, and has
degree at most $\ell-1$. Its norm contribution is at least
$e/(\ell-1)$. The three selected orbits are distinct.
Adding the vacuum contribution $e$ gives the stated inequality.

If $D$ is odd and $Z$ is $\ell$-active, its full inertia stabilizer is an
$\ell$-group by the active-inertia theorem. In a minimal
$\ell$-power cyclotomic field, however, every nontrivial wild
subgroup lowers the conductor. Its stabilizer is therefore
also a tame group, and hence trivial. Thus $K_Z$ is the full
cyclotomic field at its conductor. The general conductor
bound gives $b\le2a$. If $b=2a$, the displayed inequality
would give $\ell^{2a}\ge3\ell^{2a}$, a contradiction.
\end{proof}


\begin{theorem}[No codegree sixteen at dyadic valuation six]
\label{paper:dyadic-sixteen}
Let $\mathcal C$ be an ordinary perfect fusion-simple integral
modular category with exactly three dimension primes
$2<q<r$ and $v_2(\dim\mathcal C)=6$. No simple has formal
codegree $16$. Consequently every purely dyadic formal
codegree is $64$.
\end{theorem}
\begin{proof}
Suppose $f_Z=16$. Then
\[
 D=64M^2,\qquad M=q^b r^c,\qquad b,c\ge1,
 \qquad d_Z=2M.
\]
The full character field of $Z$ is dyadic, and $Z\in P_2$
by Proposition~\ref{paper:primepower-trace}. Consequently its
row meets every odd-degree label. There are a positive
multiple of four such labels, by the dimension equation
modulo four; their total row-norm contribution is at least
four, by Galois averaging and the algebraic norm bound.

The trace-dimension identity modulo $q$ and modulo $r$,
together with Proposition~\ref{paper:primepower-trace}, supplies
two distinct nonzero-trace even orbits, represented by
\[
                   d_X=2q^j,\qquad d_Y=2r^k.
\]
Their heights are one because the dyadic dimension valuation
is six and their nonzero traces have height at most one.
For either orbit the entry at $Z$ vanishes at the dyadic
place. The residue-energy bound
Theorem~\ref{paper:5.3} therefore gives entry energy at least
four.

Cross $X$ with a $q$-free nonunit of $P_r$. Such a label
has even degree, since nonunit prime-power degrees are
excluded. The normalized crossing has mass two, residue
two at $r$, and strict house below two. The large-residue
word lemma makes it the sum of two distinct $r$-primary
roots. Its tame stabilizer has order at most two, so the
$r$-inertia orbit of $X$ has size at least $(r-1)/2\ge2$.
The entry field $F=\mathbb Q(\chi_Z(X))$ is dyadic and hence
fixed by $r$-inertia. Thus $[K_X:F]\ge2$, and this orbit
spends at least eight in the $Z$-row norm: the exact sum
over the full $X$-orbit is
$[K_X:F]E(\chi_Z(X))$. The $Y$-orbit
spends at least four. Altogether
\[
                            16\ge4+8+4=16.
\]
Every inequality is an equality. There are exactly four
odd labels, every odd entry is a root of unity, and the only
even entries in the support are the two chosen orbits,
with total squared norms eight and four, respectively.
For the first orbit $[K_X:F]=2$ and its entry energy is four;
for the second the relative index is one and its entry
energy is four.

A nonzero dyadic integer $\alpha$ vanishing at the dyadic
place and having energy four has field degree at most two.
Indeed, if $h=[\mathbb Q(\alpha):\mathbb Q]$, its norm has
absolute value at least two, so
\[
               4=E(\alpha)\ge h\,2^{2/h}>4
                    \quad\text{when }h\ge3.
\]
The two even orbits therefore have sizes at most four and
two. Cauchy's inequality bounds their sums of absolute
entry values by $\sqrt{32}$ and $\sqrt8$.

The modular $STS$ relation, as used in the small-codegree
lemma, gives
\[
  \sqrt{f_Z}\,d_Z
       \le 1+\sum_{T\ne\one}d_T|\chi_Z(T)|.
\]
Each odd nonunit has degree at most $M$, the $X$-orbit has
degree at most $2M/r^c$, and the $Y$-orbit has degree at most
$2M/q^b$. Consequently
\[
 8M\le 1+3M+\frac{2M}{r^c}\sqrt{32}
                    +\frac{2M}{q^b}\sqrt8
      <1+\left(3+\frac{12}{5}+2\right)M
      =1+\frac{37}{5}M<8M.
\]
Here $q^b\ge3$, $r^c\ge5$, $M\ge15$, and
$\sqrt2<3/2$. This contradiction proves the theorem.
For the final assertion, a purely dyadic codegree makes
the other dimension valuations even; square divisibility
then leaves only $1,4,16,64$. The small-codegree exclusion
and the result just proved leave $64$.
\end{proof}

\section{Local modular action and principal witnesses}\label{paper:sec:local-action}

\label{paper:sec:local-weyl}
Work first with an ordinary integral modular category, with the positive spherical structure and $n=\sqrt D$.  Write $I$ for its simple labels and $\mathbf e_X$ for coordinate vectors.  The character identities and fields are those established above.

\subsection{Congruence normalization}

Put $\xi=p^+/n$.  The congruence and anomaly theorems of
Dong--Lin--Ng allow a root of unity $\gamma$ of order dividing eight,
with $\gamma^3=\xi$, such that
\begin{equation}\label{paper:lw:congruence}
 \rho\!\begin{pmatrix}0&-1\\1&0\end{pmatrix}=s,
 \qquad
 \rho\!\begin{pmatrix}1&1\\0&1\end{pmatrix}=t:=T/\gamma
\end{equation}
factors through $\operatorname{SL}_2(\mathbb Z/M)$, where
$M=\operatorname{ord}(t)$
\cite[Theorem~II, Proposition~6.7 and Theorem~6.10]{DLNCongruence}.
For $a\in(\mathbb Z/M)^\times$ let $\sigma_a$ act on roots of unity
by the $a$th power.  The same theorem gives
\begin{equation}\label{paper:lw:galois}
 \sigma_a(\rho(g))=
 \rho\!\left(\begin{pmatrix}a&0\\0&1\end{pmatrix}
 g\begin{pmatrix}a^{-1}&0\\0&1\end{pmatrix}\right).
\end{equation}
If $H_a$ denotes the represented diagonal torus, then
$H_a=\epsilon_D(a)P_a$, where $P_a$ is a permutation matrix and
\[
 \epsilon_D(a)=\sigma_a(\sqrt D)/\sqrt D\in\{1,-1\}.
\]
Indeed comparison with $s_{0X}=d_X/\sqrt D$ makes every Galois
sign equal to $\epsilon_D(a)$ and preserves the dimensions.
The unsigned permutations $P_a$ fix $0$.  They are the tori
whenever $D$ is a square, or for any automorphism fixing
$\sqrt D$.  The signs cancel in the unnormalized covariance \eqref{paper:family:eq:galois-permutations}.  We continue to use the conductor bound of Theorem~\ref{paper:family:por:th:exact-conductor}.

\subsection{Rotation traces and full residue blocks}

We record explicitly the published categorical input used below.
Index the central insertion by $\mathcal U_{ij}=U_{i,j^*}$
in Ng--Schauenburg's notation. Its underlying object is
$X_i\otimes X_j^*$ and its central twist is $\theta_i/\theta_j$.
Put $L_m^{ij}(V)=\nu_{m,1}^{\mathcal U_{ij}}(V)$, $m\ge1$.
The dual on the second index gives the unbarred second
$S$-factor in the formula below.
The generalized indicators are additive in $V$, are traces of
finite-order rotations, and satisfy the tensor-grouping identities
\cite[Lemmas~2.5, 2.7, Proposition~2.8 and Corollary~5.5]{NSind}.
All power traces of a rotation lie in $\mathbb Q(\zeta_N)$, where
$N=\operatorname{ord}(T)=\operatorname{FSexp}(\mathcal C)$
\cite[Proposition~5.7]{NSind}\cite[Theorem~7.7]{NSfs}.
Their modular formula is
\begin{equation}\label{paper:lw:indicator-formula}
 L_m^{ij}([V])=
 \sum_{A,B}s_{iA}s_{jB}(\theta_A/\theta_B)^mN_{AB}^{V}
 \quad\cite[Proposition~6.1]{NSind}.
\end{equation}
These statements, rather than an unproved divisibility property of
higher indicators, suffice for the following arguments.

\begin{lemma}[Two consequences of rotation]\label{paper:lw:rotation}
Extend $L_m^{ij}$ linearly over a cyclotomic integer ring.  If
$\ell\mid m$ and $y$ has integral coefficients, then, at every place
above $\ell$,
\begin{equation}\label{paper:lw:rotation-mod}
 (L_m^{ij}(y))^\ell\equiv L_{m/\ell}^{ij}(y^\ell).
\end{equation}
For $A_m=sT^ms^{-1}$ one has
\begin{equation}\label{paper:eq:mixed-rotation}
 L_m^{ij}(f_Ze_Z)=f_Z(A_m)_{iZ}\overline{(A_m)_{jZ}}.
\end{equation}
In particular,
\begin{equation}\label{paper:eq:indicator-entry}
 L_m^{ii}(e_Z)=|(A_m)_{iZ}|^2.
\end{equation}
At $m=1$, the diagonal insertion recovers actual adjoint multiplicity:
\begin{equation}\label{paper:adjoint-rotation}
                 L_1^{ii}([V])=N_{X_iX_i^*}^{V}.
\end{equation}
For a genuine object $V$, put $b\ge\max(1,v_\ell(N))$ with
$\ell^b\equiv1\pmod{N_{\ell'}}$.  If $h\ge2b$, then
\begin{equation}\label{paper:lw:large-power}
 L_{N_{\ell'}}^{ii}(V^{\ell^h})
 \equiv L_{N_{\ell'}\ell^b}^{ii}(V^{\ell^{h-b}})
 \pmod{\ell^{h-2b+1}\overline{\mathbb Z}_\ell}.
\end{equation}
\end{lemma}
\begin{proof}
In residue characteristic $\ell$, the $\ell$th power of the trace
of a rotation is the trace of its $\ell$th power.  Grouping
$\ell$ successive strands gives the right side of
\eqref{paper:lw:rotation-mod} for a genuine object.  Additivity and the
Frobenius map in the commutative residue fusion algebra extend this
to integral virtual coefficients.

Contract \eqref{paper:lw:indicator-formula} with
$f_Ze_Z=\sum_V\overline{\chi_Z(V)}[V]$.
Multiplicativity of $\chi_Z$ separates the two sums and gives
$f_Z(A_m)_{iZ}(A_{-m})_{jZ}$.
Since $s^{-1}=Cs$, where charge conjugation $C$ commutes with $s,T$,
the matrix $A_m$ is symmetric and unitary.  Thus
$(A_{-m})_{jZ}=\overline{(A_m)_{jZ}}$.
At $m=1$, the modular relation gives
$|(A_1)_{iZ}|=|s_{iZ}|$; Verlinde's formula and reciprocity
therefore give
$L_1^{ii}([V])=\sum_Z|s_{iZ}|^2\chi_Z(V)=N_{X_iX_i^*}^{V}$.

For \eqref{paper:lw:large-power}, the adjoint insertion $\mathcal U_{ii}$ has twist one.
Both sides are traces of the powers $\ell^h$ and $\ell^{h-b}$ of
the same rotation on $N_{\ell'}\ell^h$ strands.  Its eigenvalues
have order dividing $N_{\ell'}\ell^h$.  Since all its power traces
belong to $\mathbb Q(\zeta_N)$, its eigenvalue multiplicities are
constant on Galois orbits over that field.  The two powers agree
on the prime-to-$\ell$ components and on components of order at
most $\ell^{h-b}$.  At order $\ell^t>\ell^{h-b}$, the subgroup of
units congruent to one modulo $\ell^b$ has orbit size
$\ell^{t-b}\ge\ell^{h-2b+1}$, and both powers are constant on its
orbits.  Each orbit's contribution to the difference is divisible
by this integer, proving the congruence.
\end{proof}

Fix a place above $\ell$.  A full $\ell$-block is an equivalence
class of characters with the same reduction on every basis element.
Its idempotent $e_B=\sum_{X\in B}e_X$ is locally integral: over a
complete unramified coefficient ring containing the residue fields,
the idempotents of the special fiber lift.  In particular
\begin{equation}\label{paper:lw:block-trace}
 \tau(e_Bb)=\sum_{X\in B}\frac{\chi_X(b)}{f_X}
\end{equation}
is locally integral for every locally integral $b$.
Equation \eqref{paper:family:lw:argument-galois} shows that the full block
partition is independent of the chosen place above $\ell$.
Inertia preserves every block.  If $\ell\nmid f_X$, the expression
for $e_X$ in \eqref{paper:family:lw:character-identities} is locally integral,
so $\{X\}$ is a full block.

For a set of labels $E$, write $P_E$ for its coordinate projector.
The algebra element $e_B$ and the coordinate operator $P_B$ are
different objects; no unproved commutation between coordinate
projectors and fusion multiplication is used below.

\begin{theorem}[Block orthogonality and local invariance]
\label{paper:lw:block}
The assertion and proof hold also for weakly integral modular categories.
For every full $\ell$-block $B$,
\[
 \sum_{Z\in B}\overline{s_{Zi}}s_{Zj}=0
       \quad\text{if }\theta_{i,\ell}\ne\theta_{j,\ell}.
\]
Its coordinate subspace is invariant under the entire $\ell$-local
factor of the congruence representation.  The same invariance
holds after specifying all complementary primary twist components.
In particular, if
$E=P_\ell\cap\{Y:\theta_{Y,\ell'}=1\}$, where $P_\ell$ is the
principal block, then
\begin{equation}\label{paper:lw:principal-fiber-orthogonality}
 \sum_{Y\in E}d_Y\chi_X(Y)=0
 \qquad(\theta_{X,\ell}\ne1).
\end{equation}
\end{theorem}
\begin{proof}
In the finite residue algebra, fix a maximal ideal with residue
field $\mathbb F_{\ell^h}$.  Its geometric residue characters form
Frobenius-conjugate blocks $B_0,\ldots,B_{h-1}$.  For any
$a\in\mathbb F_{\ell^h}^\times$, the Chinese remainder theorem
gives an element with residue $a$ at this maximal ideal and zero
at the others.  Lift it to a genuine object $V$ by choosing
nonnegative integral coefficients in the simple basis.

Choose $b$ as in Lemma~\ref{paper:lw:rotation}, also divisible by $h$,
and let the exponent $k$ tend to infinity through multiples of
$h$.  In the split algebra, both $V^{\ell^k}$ and
$V^{\ell^{k-b}}$ converge $\ell$-adically to
$\sum_j[a]^{\ell^j}e_{B_j}$, with Teichm\"uller lifts $[a]$.
Equation \eqref{paper:lw:large-power} therefore makes
$L_{N_{\ell'}}^{ii}$ and $L_{N_{\ell'}\ell^b}^{ii}$ agree on this
sum.  The distinct characters $a\mapsto[a]^{\ell^j}$ are linearly
independent; hence the two functionals agree on each $e_{B_j}$.
As $N\mid N_{\ell'}\ell^b$, \eqref{paper:eq:mixed-rotation} gives
\[
 \sum_{Z\in B}|(sT^{N_{\ell'}}s^{-1})_{iZ}|^2
       =\mathbf1_{i\in B}.
\]
Unitarity implies that this matrix preserves $B$.  The power
$T^{N_{\ell'}}$ distinguishes the primary $\ell$-twists, so
conjugating back gives block orthogonality.

The projector $P_B$ commutes with the diagonal local upper
unipotent and, by what was just proved, with its lower conjugate
by $s$.  These two elementary unipotents generate the local
$\operatorname{SL}_2$ factor.  Complementary twist projectors
commute with that factor as well, proving the fiber assertion.
For the stated $E$, the operator $sP_Es^{-1}$ commutes with the
actual primary twist $T_\ell$.  Its $(X,0)$ entry is
$d_XD^{-1}\sum_{Y\in E}d_Y\chi_X(Y)$, proving
\eqref{paper:lw:principal-fiber-orthogonality}.
\end{proof}


\subsection{Local Weyl columns}

For a set $J$ of prime divisors of $M$, write $W_J,U_J,H_{a,J}$
for the products of the corresponding Weyl, upper-unipotent and
torus operators.  Write $\gamma_J$ and $T_J$ for the associated
primary parts of $\gamma$ and $T$, so $U_J=T_J/\gamma_J$.
The complementary product to $W_\ell$ is denoted $W_{\ell'}$;
it includes any dyadic anomaly factor.  The local identities are
\begin{equation}\label{paper:lw:local-identities}
 U_J(W_JU_JW_J^{-1})U_J=W_J^{-1},\qquad
 \sigma_a(W_J)=H_{a,J}W_J=W_JH_{a,J}^{-1}.
\end{equation}
They follow from matrix multiplication and \eqref{paper:lw:galois}.
In particular $W_J^2=H_{-1,J}$ is an involutive signed permutation,
$\overline{W_J}=W_J^{-1}$, and $W_J$ is symmetric.
Its entries lie in the cyclotomic field on the levels selected by
$J$ and are unramified at every prime outside $J$.

\begin{theorem}[Integral vacuum coefficients and foreign denominators]
\label{paper:lw:columns}
If $D=n^2$ is a square, the vector $w=W_J\mathbf e_0=W_J^{-1}\mathbf e_0$ is rational, and
\begin{equation}\label{paper:lw:k}
 k_{J,X}:=\frac{n w_X}{d_X}\in\mathbb Z.
\end{equation}
If $\ell\notin J$ and $\ell\mid f_X$, then $\ell\mid k_{J,X}$.
Without a square-dimensional hypothesis, if $a=v_\ell(f_X)$ and $\ell\notin J$, every
nonzero entry of the $X$-column of $W_J^{\pm1}$ satisfies
\begin{equation}\label{paper:lw:foreign-bound}
 v_\ell((W_J^{\pm1})_{YX})\ge
 \begin{cases}0,&a=0,\\-\lfloor(a-1)/2\rfloor,&a>0.
 \end{cases}
\end{equation}
The valuation is normalized by $v_\ell(\ell)=1$ and the bound holds
at every place.  More generally, let $E$ be a coordinate subset invariant under
the local factors in $J$. If every $J$-active label $X\in E$
has $v_\ell(f_X)\le2$, then $W_J^{\pm1}|_E$ are integral at
$\ell$. This includes $E$ equal to the entire label set.
\end{theorem}
\begin{proof}
For the first assertion assume $D$ is square.  All tori then fix $\mathbf e_0$, so \eqref{paper:lw:local-identities} makes $w$ Galois
fixed and gives $W_J^{-1}\mathbf e_0=W_J\mathbf e_0$.  Choose a positive integer
$m$ equal to one on the primary levels in $J$ and zero on the
others.  The Bruhat identity gives
\begin{equation}\label{paper:lw:crt-rotation}
 A_m=sT^ms^{-1}
    =\gamma_J^3T_J^{-1}W_J^{-1}T_J^{-1}.
\end{equation}
Thus \eqref{paper:eq:mixed-rotation} with $i=j=0$ is
$L_m^{00}(f_Xe_X)=f_X|w_X|^2=k_{J,X}^2$.
It is an algebraic integer; as $k_{J,X}$ is rational, it is an
integer.  For $\ell\notin J$, choose the same $m$ divisible by
$\ell$.  Applying \eqref{paper:lw:rotation-mod} to $y=f_Xe_X$, using
$y^\ell=f_X^{\ell-1}y$, shows that $k_{J,X}^2$ vanishes above
$\ell$ when $\ell\mid f_X$, proving the stated divisibility.

For the general bound no square hypothesis is needed.  Put
$w=W_J^{-1}$ and write $H=W_J^2=\varepsilon P$ with $P$ an
involutive permutation and $\varepsilon\in\{1,-1\}$.
At $j=Pi$ one has $\overline{w_{jX}}=\varepsilon w_{iX}$.
Equations \eqref{paper:eq:mixed-rotation} and \eqref{paper:lw:crt-rotation}
therefore show that $f_Xw_{iX}^2$, up to a root of unity, is integral.
When $a>0$ its valuation is strictly positive by
\eqref{paper:lw:rotation-mod}.  The entries of $w$ belong to a field
unramified at $\ell$, so their normalized valuations are integers.
Hence $a+2v_\ell(w_{iX})\ge0$, strictly when $a>0$, which is
\eqref{paper:lw:foreign-bound}.  Multiplication by $H$ proves the same
bound for $W_J$.

Finally, let $u_X=W_J\mathbf e_X$ for the $J$-active
labels in $E$. Their columns and adjoints are integral at every
$\ell$-place. Invariance gives, on the coordinate space of $E$,
\[
 W_JT_JW_J^{-1}
   =I_E+\sum_{\substack{X\in E\\X\ J\text{-active}}}
                     (\theta_{X,J}-1)u_Xu_X^*.
\]
This matrix is integral entrywise. The normalization scalars and
diagonal twist entries are roots of unity. Substitution into the
Bruhat identity therefore proves integrality of $W_J^{-1}|_E$
and then of its adjoint.

\end{proof}

\begin{corollary}[The square-dimensional vacuum entry and support]
\label{paper:lw:vacuum-support}
Assume $D=n^2$ is square and choose $m$ as in
\eqref{paper:lw:crt-rotation}.  Then
\[
 D^{-1}\sum_Xd_X^2\theta_X^m
   =\gamma_J^3(W_J)_{00},\qquad
 nW_J\mathbf e_0\in\mathbb Z^I.
\]
If $\ell\mid n$ and $\ell\notin J$, every coordinate of
$nW_J\mathbf e_0$ is divisible by $\ell$.
For $J=\{\ell\}$ the vacuum column is supported on $P_\ell$.
\end{corollary}
\begin{proof}
Take the vacuum entry in \eqref{paper:lw:crt-rotation}.
The integrality assertion is $n w_X=d_Xk_{J,X}$.
If $\ell\nmid d_X$, then $\ell\mid f_X$, so
Theorem~\ref{paper:lw:columns} supplies the required factor $\ell$;
otherwise it is already in $d_X$.  The support statement is
local block invariance applied to the vacuum.
\end{proof}

\begin{theorem}[Parity of the complementary Weyl coefficient]
\label{paper:global-parity}
Let $\mathcal C$ be an ordinary integral modular category of odd
square dimension $D=n^2$. Let $J$ be a set of its dimension primes
and let $W_{J'}$ be the complementary local Weyl product, including
any dyadic anomaly factor, so that
$s=W_JW_{J'}$. With the integer vacuum coefficients of
Theorem~\ref{paper:lw:columns}, one has
\[
 H_{-1,J}X=X\quad\Longleftrightarrow\quad k_{J',X}\text{ is odd}.
\]
In particular, a label fixed by this local central torus belongs
to the complementary principal join $P_{J'}$
(the component obtained by joining the complementary full-block
partitions; see Section~\ref{paper:sec:joined-calculus}) and is inactive at
every prime in $J$. No label active at an odd prime $p$ is fixed
by $H_{-1,p}$.
\end{theorem}
\begin{proof}
Write $d=(d_X)_X$. All formal codegrees and all simple dimensions are odd. Since
$2\notin J$, Theorem~\ref{paper:lw:columns} makes the entries of $W_J$
integral at every place above two. The local tori fix the vacuum,
and the local factors commute, so
\[
 W_Jd=nW_J^2W_{J'}\mathbf e_0=nW_{J'}\mathbf e_0.
\]
The $X$th coordinate on the right is $d_Xk_{J',X}$.
At a fixed place $\mathfrak l$ above two, all the odd integers $d_U$ reduce
to one. Symmetry therefore gives
\[
 (H_{-1,J})_{XX}
   =\sum_U(W_J)_{XU}^{2}
   \equiv\left(\sum_U(W_J)_{XU}d_U\right)^2
   \equiv k_{J',X}^{2}\pmod{\mathfrak l}.
\]
The torus is an unsigned permutation, so its diagonal entries
are zero or one. This proves the equivalence.
A nonzero complementary vacuum coefficient lies in the
complementary principal join. It is also inactive at each prime
in $J$: the corresponding actual primary twist commutes with
$W_{J'}$ and fixes $\mathbf e_0$.
\end{proof}

\begin{corollary}[Joint inertia at the active primes]
\label{paper:active-inertia}\label{paper:full-primary-orbit}
Under the hypotheses of Theorem~\ref{paper:global-parity}, let $A$ be a set
of primes at which $X$ is active, and write
$\operatorname{ord}(\theta_{X,\ell})=\ell^{a_\ell}$ for $\ell\in A$.
Work in a cyclotomic splitting field containing the full local levels
$\ell^{b_\ell}$, and let
\[
               I_A=\prod_{\ell\in A}(\mathbb Z/\ell^{b_\ell})^\times
\]
act through primary inertia. The stabilizer of $X$ in $I_A$ is
contained in
\[
 \prod_{\ell\in A}
 \ker\!\left((\mathbb Z/\ell^{b_\ell})^\times
                  \longrightarrow(\mathbb Z/\ell^{a_\ell})^\times\right).
\]
Consequently
\[
                  \prod_{\ell\in A}\varphi(\ell^{a_\ell})
                         \mid |I_A\cdot X|.
\]
More generally, let $L\subseteq K_X$ and put
$\alpha_\ell=v_\ell(\operatorname{cond}L)$. Then
\begin{equation}\label{paper:conductor-weight}
 \prod_{\ell\in A}
 \begin{cases}
       \varphi(\ell^{a_\ell}),&\alpha_\ell=0,\\
       \ell^{\max(a_\ell-\alpha_\ell,0)},&\alpha_\ell\ge1
 \end{cases}
                         \ \mid\ [K_X:L].
\end{equation}
In particular, if $L$ is unramified at every prime of $A$, then
\[
                  \prod_{\ell\in A}\varphi(\ell^{a_\ell})
                         \mid [K_X:L].
\]
For a single active prime $\ell$, the full inertia stabilizer is
therefore a subgroup of the principal units and is an $\ell$-group.
\end{corollary}
\begin{proof}
If $u=(u_\ell)\in I_A$ fixes $X$, the square action on actual
twists gives $u_\ell^2\equiv1\pmod{\ell^{a_\ell}}$.
Thus $u_\ell=\varepsilon_\ell v_\ell$, with
$\varepsilon_\ell\in\{1,-1\}$ and
$v_\ell\equiv1\pmod{\ell^{a_\ell}}$. Every $v_\ell$ has odd
order. Raising $u$ to a common odd multiple of these orders leaves
exactly the partial inversion $(\varepsilon_\ell)_\ell$, which
still fixes $X$. If any sign is negative, let $J$ be the nonempty
set of its negative coordinates. The preceding local-parity theorem
says that a label fixed by $H_{-1,J}$ is inactive at every prime of
$J$, contrary to the definition of $A$. All signs are therefore
positive, proving the stabilizer containment and the orbit divisibility.

For the field-degree assertion, use the product of full primary
inertia where $\alpha_\ell=0$ and the principal congruence subgroup
$u_\ell\equiv1\pmod{\ell^{\alpha_\ell}}$ where $\alpha_\ell\ge1$.
This product fixes $L$. At a present conductor prime, a stabilizer
element cannot have negative sign modulo $\ell^{a_\ell}$;
its actual twist image has size
$\ell^{\max(a_\ell-\alpha_\ell,0)}$. At the absent primes, the same
odd-power argument as above excludes every nontrivial partial
inversion. The stabilizer is therefore contained in the product
of the corresponding twist kernels, whose index is exactly the
product in \eqref{paper:conductor-weight}. In the abelian cyclotomic
Galois group, the resulting orbit length divides the orbit length
over $L$, which is $[K_X:L]$.
Finally, each single-prime congruence kernel in the first displayed
containment is an $\ell$-group, proving the last assertion.
\end{proof}

\begin{proposition}[Every primary phase is visible in the principal block]
\label{paper:principal-phase-support}
Assume $D=n^2$, fix a dimension prime $p$, and let $\mathcal F$ be
any set of actual $p$-primary twist values. Put
\[
 F=\{X:\theta_{X,p}\in\mathcal F\},\qquad E_F=\sum_{X\in F}e_X.
\]
Then multiplication by $E_F$ is
\[
                       N_{E_F}=W_p^*P_FW_p,
\]
so it preserves every full $p$-block. If
$k_{p,Z}=n(W_p)_{Z0}/d_Z\in\mathbb Z$ is the vacuum coefficient
of Theorem~\ref{paper:lw:columns}, then
\[
       \sum_{X\in F}d_X^2
           =\sum_{Z\in P_p\cap F}d_Z^2 k_{p,Z}^2.
\]
Consequently every actual $p$-primary twist value occurring anywhere
occurs at a label in $P_p$ with $k_{p,Z}\ne0$. In particular the
principal full $p$-block attains the global $p$-primary twist exponent.
For each prime $t\ne p$ dividing $D$, the common mass in the displayed
identity is divisible by $t^2$.
\end{proposition}
\begin{proof}
The complementary local Weyl operators commute with the actual
$p$-primary twist, hence with $P_F$. Thus $s^*P_Fs=W_p^*P_FW_p$.
Local block invariance proves the first assertion. Taking the vacuum
matrix coefficient gives
\[
 \frac1D\sum_{X\in F}d_X^2
      =\sum_{Z\in F}|(W_p)_{Z0}|^2.
\]
The local vacuum column is supported on $P_p$, is rational, and has
the integral coordinates $n(W_p)_{Z0}=d_Zk_{p,Z}$. This proves the
sum-of-squares identity and its support consequence. Finally, for
$t\ne p$, each $d_Zk_{p,Z}$ is divisible by $t$: this is immediate
if $t\mid d_Z$, and otherwise $t\mid f_Z$ and Theorem~\ref{paper:lw:columns}
gives $t\mid k_{p,Z}$. Every square on the right is therefore
divisible by $t^2$.
\end{proof}

\begin{lemma}[A Weyl line and the support of a twist prime]
\label{paper:family:lw:weyl-line}
The Weyl-line and support assertions below hold also for weakly
integral modular categories.  If $W_\ell$ preserves the coordinate line of $X$, then
$\ell\nmid f_X$.  Consequently
\[
 \ell\nmid f_X\quad\Longrightarrow\quad
 \theta_{X,\ell}=1,\qquad
 S_{XY}=0\ \text{when }\theta_{Y,\ell}\ne1.
\]
Conversely, a row whose codegree is divisible by $\ell$ meets an
$\ell$-active column.  For odd $\ell$, the $\ell$-primary order of
$\theta_X$ divides the $\ell$-primary part of $f_X$.
If no twist contains both distinct primes $\ell,t$, every
$\ell$-active label has $t$-unit codegree.
\end{lemma}
\begin{proof}
Choose $m$ zero on the $\ell$-level and one on all complementary
levels, with $\ell\mid m$.  If $W_\ell$ preserves $\mathbf e_X$ up to a
phase, the identity $W_{\ell'}^{-1}=s^{-1}W_\ell$ gives
$|(W_{\ell'}^{-1})_{0X}|=|(s^{-1})_{0X}|=1/\sqrt{f_X}$.
This identity does not require unsigned tori.  Thus
$L_m^{00}(f_Xe_X)=1$ by \eqref{paper:eq:mixed-rotation} and
\eqref{paper:lw:crt-rotation}.  Congruence \eqref{paper:lw:rotation-mod}
excludes $\ell\mid f_X$.

If $\ell\nmid f_X$, its singleton block and Theorem~\ref{paper:lw:block}
give the asserted zeros, since the vacuum entry is nonzero.
Conversely, if a row misses every active column, the local lower
unipotent $s(T_\ell/\gamma_\ell)s^{-1}$ acts on its line by
$\gamma_\ell^{-1}$.  The diagonal upper unipotent also preserves
that line, so the whole local group does.  These two conjugate
unipotents must act there by the same scalar, giving
$\theta_{X,\ell}=1$; the Weyl-line argument gives $\ell\nmid f_X$.

For the last assertion put $a=v_\ell(f_X)$.  The case $a=0$ was
just proved.  If $a>0$, Theorem~\ref{paper:family:por:th:exact-conductor} makes all local
tori congruent to one modulo $\ell^a$ fix $X$.  The square action
on twists, $H_bT_\ell H_b^{-1}=T_\ell^{b^2}$, and an element
$b=1+\ell^a$ show that the twist order divides $\ell^a$.

For the final assertion, the $t$-local group preserves each
nontrivial $\ell$-twist fiber.  There its actual $t$-twist is the
identity, so its upper unipotent is scalar.  Its lower conjugate
and, by Bruhat, its Weyl are scalar too.  Apply the first assertion
to every coordinate line of that fiber.
\end{proof}

\begin{corollary}[Parity forced by separation]\label{paper:lw:separated-parity}
If two dimension primes $\ell,t$ of an integral modular category
never occur in the same twist, both $v_\ell(D)$ and $v_t(D)$ are
even.  Thus every dimension prime of odd valuation occurs in a
joint twist with every other dimension prime.  With at least two
dimension primes, an isolated prime forces $D$ to be a square.
\end{corollary}
\begin{proof}
Categorical Cauchy supplies active labels at both primes
\cite[Theorem~3.9]{BNRWFinite}.  A $t$-active label $X$ has
$\ell$-unit codegree, so $v_\ell(D)=2v_\ell(d_X)$; interchange
the primes.  For an isolated prime, its active labels similarly
force even valuation at every foreign prime, and any foreign
active label forces even valuation at the isolated prime itself.
\end{proof}

\begin{proposition}[Isolated primes and separated character fields]
\label{paper:lw:isolated}
Assume $D=n^2$ is square.  Let $\ell$ be any one of the three dimension primes, write
$n=\ell^a m$ with $\ell\nmid m$, and suppose no twist contains
$\ell$ together with another prime.  Then $\gamma=\xi=\pm1$.
Set
\[
 \mathcal W_t=\gamma_t^3W_t,\qquad
 b=\ell^a\mathcal W_\ell\mathbf e_0.
\]
Every $\ell$-active label has degree $\ell^j m$ for some
$0\le j<a$, and its character field has $\ell$-power conductor.
The vector $b$ is integral, and, in $\mathbb Z_\ell$,
\begin{equation}\label{paper:lw:isolated-congruence}
 b_X\equiv \xi d_X/m\pmod{\ell^a}.
\end{equation}
Thus every label with $v_\ell(d_X)<a$ is principal at $\ell$,
and $v_\ell(b_X)=v_\ell(d_X)$ for those labels.  Moreover
\begin{equation}\label{paper:lw:isolated-vacuum}
 \ell^a-b_0
 =\ell^{-a}\sum_{X\ \ell\text{-active}}
           \ell^{2j_X}(1-\theta_{X,\ell})\in\mathbb Z.
\end{equation}
For odd $\ell$, $\mathcal W_\ell=W_\ell$.
Separately, in an integral modular category of arbitrary dimension,
if two odd dimension primes $\ell,t$ never occur in the same
twist, both have even valuation in $D$ and no character field
ramifies at both.  In square dimension the same field assertion
allows a dyadic prime whenever $\gamma=\pm1$.
\end{proposition}
\begin{proof}
A local group at $t\ne\ell$ preserves every $\ell$-twist fiber.
On a nontrivial such fiber its actual $t$-twist is the identity,
so its upper unipotent is scalar.  Its lower conjugate is the
same scalar, and Bruhat makes its Weyl scalar too.  By
Lemma~\ref{paper:family:lw:weyl-line}, every $\ell$-active label has $t$-unit
codegree.  Its codegree is consequently an $\ell$-power, proving
the degree and field assertions.  Conversely every label active
at a prime other than $\ell$ has $\ell$-unit codegree.

There is a nonempty exact odd-twist fiber with trivial actual
dyadic twist.  If $\ell$ is odd use an $\ell$-active label; if
$\ell=2$ use a label active at either odd dimension prime.
The required active labels exist by categorical Cauchy
\cite[Theorem~3.9]{BNRWFinite}.  The dyadic group preserves that
fiber.  On it its upper and Weyl are the scalars
$\gamma^{-1}$ and $\gamma^3$.  The Weyl square is an unsigned
torus permutation, so $\gamma^6=1$.  As $\gamma$ has two-power
order, $\gamma=\xi=\pm1$.
For $\gamma=-1$, tensor the dyadic factor with the sign character
of $\operatorname{SL}_2(\mathbb F_2)$.  The corrected local uppers
are the actual primary twists, the corrected Weyls are precisely
$\mathcal W_t$, and the tori are unchanged.  Their product is
$\xi s$.

Let $V_\ell$ be the moving space of this corrected $\ell$-local
group, the orthogonal complement of its fixed space.  It is
generated by the local-group translates of active coordinate
vectors: conjugates of the upper unipotent generate the group.
Isolation makes every complementary local group trivial on these
coordinates, and hence on $V_\ell$.  On its orthogonal complement
$\mathcal W_\ell=I$.  Therefore
\begin{equation}\label{paper:lw:isolated-split}
 \xi s=\mathcal W_{\ell'}+(\mathcal W_\ell-I).
\end{equation}
The last clause of Theorem~\ref{paper:lw:columns} makes $W_\ell$
integral away from $\ell$, and makes $W_{\ell'}$ integral at
$\ell$.  Together with rationality and
$n(W_\ell\mathbf e_0)_X=d_Xk_{\{\ell\},X}$, this proves $b\in\mathbb Z^I$.
Multiply the vacuum column of \eqref{paper:lw:isolated-split} by
$\ell^a$ to obtain \eqref{paper:lw:isolated-congruence}.
The stated exact valuations follow.  Since $W_\ell\mathbf e_0$ is
supported on the principal block, every such nonzero coordinate
is principal.  Finally \eqref{paper:lw:crt-rotation} gives
\[
 D(\mathcal W_\ell)_{00}=\sum_Xd_X^2\theta_{X,\ell},
\]
which is \eqref{paper:lw:isolated-vacuum}.

For the separate assertion at arbitrary integral $D$, an active
label at $t$ has $\ell$-unit codegree by the scalar-line argument.
Thus $v_\ell(D)=2v_\ell(d_X)$ is even, and conversely for $t$.
The local inertia actions at these odd primes therefore fix
$\sqrt D$, so their tori are unsigned permutations.  Use the
uncorrected local groups at odd primes, and the corrected groups
in square dimension if a dyadic
factor is present and $\gamma=\pm1$.  Separation makes the
$t$-local group trivial on $V_\ell$ and the $\ell$-local group
trivial on $V_t$.  These two moving spaces are orthogonal, and
\[
 (H_{a,\ell}-I)(H_{b,t}-I)=0
\]
for all tori.  If both permutations moved a label $X$,
application to $\mathbf e_X$ would contradict linear independence of
the coordinate vectors: the positive term $\mathbf e_X$ cannot cancel
either negative term.  Thus at least one entire inertia group
fixes $X$, proving the unramifiedness assertion.
\end{proof}

\begin{lemma}[The short isolated-quartic argument]\label{paper:lw:isolated-quartic}
Suppose $n=\ell^2m$ and $\ell$ is isolated.  If there are at least
two distinct simple labels of $\ell$-height one, then some active
label has degree $\ell m$ and codegree $\ell^2$.
Consequently an ordinary perfect fusion-simple three-prime table
with two such labels cannot have an isolated quartic prime, once
the prime-square-codegree trace-support obstruction is applied.
\end{lemma}
\begin{proof}
Suppose no active label has degree $\ell m$.  Then all active
degrees equal $m$.  Put $L=\#\{Z:\theta_{Z,\ell}\ne1\}$ and
$\delta=\sum_{Z\ \ell\text{-active}}(1-\theta_{Z,\ell})$.
Use $b=\ell^2\mathcal W_\ell\mathbf e_0$ from
Proposition~\ref{paper:lw:isolated}.  Its vacuum entry is an
$\ell$-unit and
\[
 b_0=\ell^2-\delta/\ell^2.
\]
In particular $L>0$.

Choose a height-one label $X$.  It is principal at $\ell$,
so its entry at each active $Z$ is nonzero: its character value
there reduces to the $\ell$-unit $d_Z=m$.
Since $d_X=\ell u$ with $u\mid m$, LCM integrality gives
$S_{XZ}/(\ell m)$ integral.  The Galois-mean contribution of
these entries to the normalized row norm is at least
$L/\ell^2$.  The vacuum contributes positively, so $L<\ell^2$.
Thus
\[
 1-\delta/\ell^2
 =\left(1-\frac L{\ell^2}\right)
       +\frac1{\ell^2}\sum_{Z\ \ell\text{-active}}\theta_{Z,\ell}
\]
is an algebraic integer all of whose conjugates have modulus
strictly less than one.  It is zero by the norm, giving
$\delta=\ell^2$ and $b_0=\ell^2-1$.
Each of the two height-one labels has $v_\ell(b_X)=1$, hence
$|b_X|\ge\ell$.  Unitarity now gives the contradiction
\[
 \ell^4=\|b\|^2
 \ge(\ell^2-1)^2+2\ell^2=\ell^4+1.
\]
The argument includes $\ell=2$ and uses no residue-energy estimate.
\end{proof}

\begin{lemma}[A height-two label at an isolated sextic prime]
\label{paper:isolated-height-two}
Let $\mathcal C$ be an ordinary perfect fusion-simple modular
category of odd dimension $D=\ell^6m^2$ with exactly three
dimension primes and $\ell\nmid m$.
Suppose $\ell$ is isolated in the actual twist support.
If a label of $\ell$-height two exists, every $\ell$-active
label has degree $\ell m$. Every label of $\ell$-height two
has an $\ell$-primary full character field and complete twist one.
\end{lemma}
\begin{proof}
Isolation and the prime-square-codegree exclusion leave the
active degrees $m,\ell m$.
Let $N_0$ count active labels of degree $m$.
A height-two label is principal at $\ell$ and therefore
meets all these labels, since $m$ is an $\ell$-unit.
Their LCM energy, with the positive vacuum contribution,
gives $N_0<\ell^2$.

At each height and exact twist order $\ell^b$, the Galois
square action makes phase multiplicities constant on each
primitive square class. Each multiplicity is even, by full
active inertia. For $b\ge2$ the sum of roots in either
class is zero. For $b=1$ their difference is a nonzero
irrational quadratic Gauss sum.
Write $a_j,b_j$ for the multiplicities on the two prime-level
classes at height $j=0,1$. Rationality of the weighted vacuum
Gauss sum gives
\[
       (a_0-b_0)+\ell^2(a_1-b_1)=0.
\]
The bound $N_0<\ell^2$ implies
$|a_0-b_0|<2\ell^2/(\ell-1)\le\ell^2$.
Thus $a_j=b_j=n_j$ for $j=0,1$.

The same count bound excludes order $\ell^3$ or higher at
height zero. Put $N_{0,2}=\ell(\ell-1)t$, $t\ge0$,
for its order-$\ell^2$ count. For
$\delta_j=\sum_{Z\text{ active},\ d_Z=\ell^jm}(1-\theta_Z)$,
we have
\[
 N_0=(\ell-1)(n_0+\ell t),\qquad
 \delta_0=\ell(n_0+(\ell-1)t),\qquad
 \delta_1\in\ell\mathbb Z.
\]
The last assertion uses the prime-level equality and the
$\varphi(\ell^b)$ divisibility of the higher-level counts.
The isolated vacuum identity says
\[
 \ell^3-\frac{\delta_0+\ell^2\delta_1}{\ell^3}\in\mathbb Z.
\]
Consequently $\ell^2\mid n_0+(\ell-1)t$, while
\[
 0\le n_0+(\ell-1)t\le n_0+\ell t
             <\frac{\ell^2}{\ell-1}<\ell^2.
\]
Hence $n_0=t=0$ and $N_0=0$.

Let $X$ have height two. If its field ramifies at a foreign
prime, isolation makes it unramified at $\ell$. For any active
$Z$ of degree $\ell m$, the intersection with $K_Z$ is then
$\mathbb Q$, so $\chi_Z(X)$ is rational. Symmetry and
$X\in P_\ell$ make this integer divisible by $\ell^2$.
A nonzero value would exhaust the norm $f_Z=\ell^4$ before
counting the vacuum. Thus $X$ would miss every active label,
contrary to its positive $\ell$-defect and the Weyl-line lemma.
Its field is therefore $\ell$-primary. No foreign prime can
be active, by full active inertia, while $\ell$-activity at
height two would give degree $\ell^2m$ and forbidden codegree
$\ell^2$. Its complete twist is consequently one.
\end{proof}

\begin{corollary}[Principal witnesses under least-prime isolation]
\label{paper:isolated-principal-quartic}
In the exact odd family, if $p$ is isolated, every $q$-free
nonunit in $P_r$ has degree $p^3r$, and every $r$-free
nonunit in $P_q$ has degree $p^3q^j$, $j\in\{1,2\}$.
\end{corollary}
\begin{proof}
The complementary witness heights are two or three.
Suppose first that $X\in P_r$ has degree $p^2r$.
The lemma makes $K_X$ $p$-primary, and isolation puts $X$
in $P_p$. Choose a $p$-free $B\in P_r$, of degree $q^jr$.
Then $S_{XB}\ne0$ by $p$-principality.
The normalized crossing $z=S_{XB}/(p^2q^jr)$ has
strict house below $r$ and zero residue at all $r$-places.
Its field is unramified at $r$, so $z/r$ is a nonzero
algebraic integer of house below one, a contradiction.

Now suppose $A\in P_q$ is $r$-free of $p$-height two.
Its full field is $p$-primary by the lemma.
Choose a $q$-free $Y\in P_r$, now necessarily of degree
$p^3r$. Its $p$-unit codegree makes $K_Y$ unramified at $p$.
The compulsory opposite crossing is consequently a rational
integer $z$, with
\[
              |z|<p^2,\qquad z\equiv p^2\pmod{qr}.
\]
Corollary~\ref{paper:family:cor:remaining-region} already excludes
$m\le2$, where $m=\lfloor\min(r,p^2)/q\rfloor$.
Thus a surviving category has $r>3q$, and
$qr>3q^2>2p^2$. No such integer $z$ exists.
\end{proof}

\begin{proposition}[A bilinear active-mass gap for an isolated sextic prime]
\label{paper:isolated-adjoint-gap}
Let $\mathcal C$ be an ordinary perfect fusion-simple modular
category of odd dimension $D=p^6M^2$, with exactly three
dimension primes, $p\nmid M$, and with $p$ isolated in the
actual twist support. Let $X,Y$ both have $p$-height three.
If $X\otimes Y$ has a $p$-active constituent, then the total
Frobenius--Perron dimension of its $p$-active constituents,
counted with multiplicity, is at least
\[
                         p^2(p-1)M.
\]
When $Y=X^*$, this lower bound improves to $2p^2(p-1)M$.
In the ordered exact family, the product of any two simples
of degree $p^3r$ therefore has only $p$-inactive constituents.
In particular this applies to $X^2$, $X\otimes X^*$, and
all pairwise products of full Galois conjugates of such an $X$.
\end{proposition}
\begin{proof}
The codegrees of $X,Y$ are $p$-units, so their full character
fields and twists are unramified at $p$. Write $n_T=N_{XY}^T$.
Split the balancing sum into labels of height three, inactive
labels of smaller height, and active labels. The first group
has dimensions divisible by $p^3$ and twists unramified at $p$.
Every label in the second group has complete twist one:
a foreign active prime would force $p$-unit codegree by
isolation. Write $I=\sum_{T\text{ in the second group}}n_Td_T$.

The prime-square-codegree exclusion and isolation make every
active dimension $d_T=p^{j_T}M$, with $j_T\in\{0,1\}$.
Put
\[
 L=\sum_{T\text{ active}}n_Tp^{j_T},\qquad
 G=\sum_{T\text{ active}}n_Tp^{j_T}\theta_T.
\]
Balancing and the $p$-unramified field of $S_{XY}$ show
that $G$ is unramified at $p$. It is a $p$-primary cyclotomic
integer, so $G\in\mathbb Z$. LCM integrality makes
$S_{XY}$ divisible by $p^3$ in the algebraic integers.
After subtracting the height-three terms from balancing,
we obtain $p^3\mid I+MG$. On the other hand the dimension
identity gives $p^3\mid I+ML$. Since $p\nmid M$,
\[
                         \delta:=L-G\in p^3\mathbb Z.
\]
If an active constituent occurs, the real contribution of
every summand to $\delta$ is positive, so $\delta\ge p^3$.
The normalized trace of $1-\theta$ for a nontrivial $p$-power
root is at most $p/(p-1)$. Averaging the rational number
$\delta$ therefore gives
\[
                    p^3\le\delta\le\frac p{p-1}L.
\]
The active mass $ML$ is at least $p^2(p-1)M$.

For $Y=X^*$, the product is self-dual. In odd dimension each
nonunit occurs with its distinct equal-twist dual, with equal
multiplicity, so $\delta\in2\mathcal O$ and hence
$\delta\in2\mathbb Z$. Thus $\delta\ge2p^3$, giving the
stated improvement.

For $d_X=d_Y=p^3r$ in the exact family, $M=q^3r^2$ and
\[
 \frac{d_Xd_Y}{M}=\frac{p^6}{q^3}<p^2(p-1).
\]
Indeed $q\ge p+2$, and
\[
 (p-1)(p+2)^3-p^4=5p^3+6p^2-4p-8>0.
\]
The entire product has less than the required active mass,
so its active part is empty. Galois conjugates preserve the
degree, giving the final assertion.
\end{proof}

\begin{remark}[What the product conclusion does and does not give]
Each product in the final assertion has both coefficient
support and character support on the $p$-inactive labels.
The character assertion follows because either input has
$p$-unit codegree. Hence both the local upper and lower
unipotents fix its coefficient vector, and the full $p$-local
group fixes it. For $A=X\otimes X^*$, this is an actual
nonnegative integral vector with unit coefficient one.
Fusion simplicity makes the fusion subcategory generated by
$A$ the entire category, but the invariant subspace of a local
modular group is not known to be closed under fusion
multiplication. The pairwise assertion for degree-$p^3r$
inputs alone does not establish that all powers of $A$ remain
supported on inactive labels.
\end{remark}
\begin{lemma}[A ternary column across a four-element foreign orbit]
\label{paper:ternary-column}
Let $D$ be square and let $X$ be a $3$-inactive label with $v_3(f_X)=2$.
Suppose its full character field is tame at three, the column
$W_3\mathbf e_X$ is integral away from three, and an inertia group at a
foreign prime has an orbit of at least four labels on $X$ while fixing
every $3$-active label. Then two $3$-active labels have degree $4d_X$.
\end{lemma}
\begin{proof}
Write $W=W_3$, $U=T_3$. The local anomaly at three is trivial; use the normalization in
\eqref{paper:lw:congruence}.
Tameness at three and \eqref{paper:lw:local-identities} put
$W\mathbf e_X$ in $\mathbb Q(\omega)^I$, where $\omega$ is a primitive
cube root. The square-integrality bound in the proof of
Theorem~\ref{paper:lw:columns} gives
\[
             \beta=3W\mathbf e_X\in\mathbb Z[\omega]^I,
             \qquad \sum_T|\beta_T|^2=9.
\]
Every squared modulus here is a nonnegative rational integer.
Let $X_1,\ldots,X_t$ be the foreign inertia orbit, with $X_1=X$.
The active coordinates of their orthonormal Weyl columns agree, since
foreign inertia fixes every active label and commutes with $W$.
Projecting $t^{-1/2}\sum_iW\mathbf e_{X_i}$ to the active coordinates
gives
\[
             \sum_{T\;3\text{-active}}|\beta_T|^2\le9/t\le9/4.
\]
This sum is nonzero: otherwise both the upper and lower local
unipotents fix $\mathbf e_X$, making its Weyl line invariant and
contradicting Lemma~\ref{paper:family:lw:weyl-line}.

Wild local tori fix $X$ and therefore fix this column. At an exact
twist order $3^b$ with $b\ge2$, the wild square action has an orbit
of at least three labels. This is incompatible with the displayed
integral bound. Thus only primitive cubic phases occur in the active
support. If $n_1,n_2$ are the respective integer squared masses at
$\omega,\omega^2$, the diagonal local relation
$WUW^{-1}=U^{-1}W^{-1}U^{-1}$ gives
\[
3\overline{\beta_X}
 =9-n_1-n_2+n_1\omega+n_2\omega^2
 =9-n_1-2n_2+(n_1-n_2)\omega.
\]
Hence $3\mid n_1-n_2$. Since $0<n_1+n_2\le2$, one has
$n_1=n_2=1$. Consequently there are precisely two active support
labels $A,B$, of phases $\omega,\omega^2$, and $\beta_A,\beta_B$
are Eisenstein units. It also follows that $t=4$.

For $i,j\in\{1,2,3,4\}$, the same relation gives
\[
\overline{W_{X_iX_j}}
 =\delta_{ij}+\frac{(\omega-1)+(\omega^2-1)}9
 =\delta_{ij}-\frac13.
\]
These four real coordinates spend $7/9$ of each column norm; the
active coordinates spend the remaining $2/9$. Thus each column is
supported on the six labels $X_1,\ldots,X_4,A,B$.
Local inversion sends the $X$-column to its complex conjugate.
They differ only at $A,B$, and an Eisenstein unit $c$ satisfies
$|c-\bar c|^2\le3$. Their squared distance is therefore at most
$6/9<2$. Distinct unitary columns have squared distance two, so
local inversion fixes $X$. The column is consequently real, and
$\beta_A=a$, $\beta_B=b$ with $a,b\in\{1,-1\}$. The same applies
to all four $X_i$.

All $X_i$ are inactive and fixed by the local central torus.
Let $V$ be the span of the six labels $X_1,\ldots,X_4,A,B$.
This space is $W$-invariant. Indeed $W^2\mathbf e_{X_i}$ and
$WU W\mathbf e_{X_i}=U^{-1}W\mathbf e_{X_i}$ both belong to $V$;
subtracting the inactive parts gives $W(a\mathbf e_A+b\mathbf e_B)$
and $W(a\omega\mathbf e_A+b\omega^2\mathbf e_B)$ in $V$, whence
$W\mathbf e_A,W\mathbf e_B\in V$.

The local central torus fixes $V$, so $W^2=I$ there. In the signed
basis $\mathbf e_{X_1},\ldots,\mathbf e_{X_4},a\mathbf e_A,b\mathbf e_B$,
write $J_{h,k}$ for the all-ones matrix. The restriction has form
\[
 W|_V=
 \begin{pmatrix}
 I_4-\frac13J_{4,4}&\frac13J_{4,2}\\
 \frac13J_{2,4}&C
 \end{pmatrix},
 \qquad
 C\binom11=\frac13\binom11,\quad
 C^2=I_2-\frac49J_{2,2}.
\]
In particular $C$ is invertible. The complementary Weyl vacuum
column is supported on $3$-inactive labels. Since
$W_3d=\sqrt D\,W_{3'}\mathbf e_0$, its coordinates at $A,B$ vanish.
Every $X_i$ has degree $d_X$, so
\[
 C\binom{a d_A}{b d_B}=-\frac{4d_X}{3}\binom11.
\]
The displayed eigenvalue of $C$ gives
$a d_A=b d_B=-4d_X$. Hence $d_A=d_B=4d_X$.
\end{proof}
\begin{lemma}[An inactive Gaussian column of norm four is real]
\label{paper:gaussian-column}
Let $W$ be a symmetric unitary matrix, let $U$ be a diagonal
root-of-unity matrix, and suppose
\[
 WUW=U^{-1}WU^{-1},\qquad W^2=H,\qquad \overline W=W^{-1},
\]
where $H$ is an involutive permutation commuting with $U$.
If $U\mathbf e_Y=\mathbf e_Y$ and
$u=2W\mathbf e_Y\in\mathbb Z[i]^I$, then $H Y=Y$ and $u$ is real.
\end{lemma}
\begin{proof}
Suppose $Y'=HY\ne Y$. Since $W\mathbf e_{Y'}=\overline{W\mathbf e_Y}$,
unitarity gives $\sum_Zu_Z^2=0$ and $\sum_Z|u_Z|^2=4$.
Write $u=x+iy$ with integral real vectors $x,y$.
Then $\|x\|^2=\|y\|^2=2$ and $\langle x,y\rangle=0$.
Each support has two positions with coefficients $\pm1$;
the supports are either equal or disjoint.
Also $Hx=x$ and $Hy=-y$.

If the supports are equal, $H$ exchanges their two positions,
which therefore have a common twist $\rho$. Thus
$UW^{-1}\mathbf e_Y=\rho W^{-1}\mathbf e_Y$.
The Bruhat identity
$U(WUW^{-1})U=W^{-1}$, equivalent to the displayed relation,
gives $W^{-1}\mathbf e_Y=\rho\mathbf e_Y$.
This makes $H Y=Y$, a contradiction.

If the supports are disjoint, write $R$ for the sum of the twists
on the two real positions. The two imaginary positions are
exchanged by $H$, so have one common twist $\rho$.
Symmetry gives $u_{Y'}=\overline{u_Y}$.
The $(Y,Y)$ and $(Y,Y')$ entries of the displayed modular relation
are respectively
\[
 R-2\rho=2u_Y,\qquad R+2\rho=2\overline{u_Y}.
\]
Consequently $\rho=-i\operatorname{Im}(u_Y)$.
Since $\rho$ is nonzero, $u_Y=\pm i$, and $Y$ is one of the
imaginary positions. Its twist is therefore $\rho=1$, whereas
the last equation gives $\rho=\mp i$. This contradiction proves
$HY=Y$. Finally $\overline{W\mathbf e_Y}=W\mathbf e_{HY}$
shows that the column is real.
\end{proof}
\begin{lemma}[A dyadic column and a foreign orbit]
\label{paper:dyadic-column}
Let $D$ be square and let $X$ be $2$-inactive, with $v_2(f_X)=2$.
Suppose $K_X$ is unramified at two and the column $W_2\mathbf e_X$
is integral away from two. If inertia at a foreign prime fixes every
$2$-active label, its orbit on $X$ has fewer than four labels.
\end{lemma}
\begin{proof}
Write $W=W_2$, $U=T_2/\gamma$, where the dyadic anomaly scalar
$\gamma$ has order dividing eight. Local Galois covariance makes
$W\mathbf e_X$ rational, and the square-integrality bound of
Theorem~\ref{paper:lw:columns} gives
\[
                  u=2W\mathbf e_X\in\mathbb Z^I,
                  \qquad \sum_Tu_T^2=4.
\]
Suppose the foreign inertia orbit has $t\ge4$ labels $X_i$.
Their Weyl columns agree on active coordinates. Orthogonality of
these columns, applied to their normalized sum, gives
\[
                  \sum_{T\;2\text{-active}}u_T^2\le4/t\le1.
\]
The left side cannot vanish. Otherwise the upper and lower local
unipotents would both preserve the line of $X$, since both $X$ and
its Weyl image would have actual dyadic twist one. Their generation
of the local group would contradict Lemma~\ref{paper:family:lw:weyl-line}.
Thus $t=4$, with precisely one active support coordinate $A$, where
$W_{AX_i}=a/2$ for all four $i$, with $a\in\{1,-1\}$.
These four entries exhaust the norm of the $A$-row, so $W_{AA}=0$.
The four $X_i$ are inactive and therefore have $U$-eigenvalue
$\gamma^{-1}$. Consequently
\[
 (WUW^{-1})_{AA}=\sum_i|W_{AX_i}|^2\gamma^{-1}
                =\gamma^{-1}.
\]
On the other hand the local identity
$WUW^{-1}=U^{-1}W^{-1}U^{-1}$ makes this diagonal entry zero,
because $(W^{-1})_{AA}=\overline{W_{AA}}=0$. This is a contradiction.
\end{proof}


\begin{lemma}[A dimension congruence for an integral local fixed vector]
\label{paper:fixed-vector-dimension-congruence}
Assume $D=n^2$, let $p\mid D$, and let
$v=\sum_Zv_Z\mathbf e_Z$ have integral coordinates and satisfy
$W_pv=v$. Then $\sum_Zv_Zd_Z\equiv0\pmod p$.
\end{lemma}
\begin{proof}
One has $W_pd=nW_{p'}\mathbf e_0$, since the local central torus
fixes the complementary vacuum column. Unitarity and $W_pv=v$
therefore give
\[
 \sum_Zv_Zd_Z=\langle v,W_pd\rangle
             =\langle v,nW_{p'}\mathbf e_0\rangle.
\]
Every coordinate of $nW_{p'}\mathbf e_0$ is an integer divisible
by $p$, by Corollary~\ref{paper:lw:vacuum-support}.
\end{proof}

\subsection{Local compression, integral columns and braid energy}
\label{paper:sec:compression}

The following local identity will supply the Fourier bounds in
Section~\ref{paper:sec:fourier}. It also identifies the equality
case and the invariant term in a trivial torus channel.

\begin{lemma}[Exact local compression on a torus channel]
\label{paper:exact-local-compression}\label{paper:twisted-inactive-bound}\label{paper:exact-local-compression-2}
Let $\rho$ be a unitary representation of
$G=\operatorname{SL}_2(\mathbb Z/\ell^n)$, where $\ell$ is odd and $n\ge1$,
and let $U,L,H,W$ denote its upper and lower unipotents, diagonal
tori and Weyl element.  Write
\[
 K_j=\ker\bigl(G\longrightarrow
                  \operatorname{SL}_2(\mathbb Z/\ell^j)\bigr)
 \quad(1\le j\le n),
\]
and let $P_A$ denote the orthogonal projector onto the $A$-fixed
vectors.  Suppose $U(t)v=v$ and $H(a)v=\psi(a)v$.
If $\psi$ is nontrivial of conductor $\ell^u$, then
\begin{equation}\label{paper:twisted-local-bound}
              \|P_UWv\|^2=\ell^{-u}\|P_{K_u}v\|^2.
\end{equation}
If $\psi=1$, then
\begin{equation}\label{paper:exact-trivial-compression}
 \begin{split}
 P_UWv&=P_Gv-\ell^{-1}(P_{K_1}-P_G)v,\\
 \|P_UWv\|^2
   &=\|P_Gv\|^2+\ell^{-2}\|(P_{K_1}-P_G)v\|^2.
 \end{split}
\end{equation}
In particular the trivial torus channel contracts in squared norm
by at least $\ell^{-2}$ after its $G$-fixed component is removed;
its component orthogonal to the $K_1$-fixed space is killed.
\end{lemma}
\begin{proof}
Put $R=\mathbb Z/\ell^n$ and
$f(t)=\langle v,L(t)v\rangle$, with the inner product linear in
its second argument.  Upper-unipotent averaging gives
\[
                 \|P_UWv\|^2=|R|^{-1}\sum_{t\in R}f(t).
\]
For $u=1$, the unit parameters have zero total contribution.
For a unit $t$, the Bruhat factorization is
\[
 L(-t)=U(-t^{-1})H(-t^{-1})WU(-t^{-1}).
\]
Its coefficient on $v$ is
$\psi(-t)\langle v,Wv\rangle$, so the sum over units vanishes.
For $u\ge2$,
the identity
\[
 U(a)L(t)=L\!\left(\frac{t}{1+at}\right)
            H(1+at)U\!\left(\frac{a}{1+at}\right)
\]
gives $f(t)=\psi(1+at)f(t/(1+at))$ whenever $1+at$ is a unit.
Torus conjugation gives $f(tz^2)=f(t)$.  Choose
$b\in1+\ell^{u-1}R$ with $\psi(b)\ne1$.  For
$v_\ell(t)<u$, the equation $1+at=b$ is soluble; since a
principal-unit group has odd order and hence every element is a square, it follows that
$f(t)=0$.  Consequently, in both cases the remaining average is
\[
 \ell^{-u}\langle v,P_{L(\ell^uR)}v\rangle.
\]
Every element of $K_u$ has a unique lower--diagonal--upper
factorization with all three parameters congruent to their
identity values modulo $\ell^u$.  Thus its uniform average is
\[
 P_{K_u}=P_{L(\ell^uR)}P_{H(1+\ell^uR)}P_{U(\ell^uR)}.
\]
The last two factors fix $v$, because the conductor of $\psi$
is $\ell^u$.  Hence $P_{K_u}v=P_{L(\ell^uR)}v$, proving
\eqref{paper:twisted-local-bound}.

Now let $\psi=1$, and let $B$ be the upper Borel of $G$.
Let $I$ be the inverse image of the upper Borel of
$\operatorname{SL}_2(\mathbb F_\ell)$.  The big cell $BWB$
is precisely $G\setminus I$: its matrices have unit lower-left
entry.  Averaging uniformly over this double coset, and using
$[G:I]=\ell+1$, gives
\[
 P_BWP_B=\frac{\ell+1}{\ell}P_G-\frac1\ell P_I.
\]
Here $v$ is $B$-fixed, $P_UWv$ is also torus-fixed, and
$I=K_1B$, so $P_Iv=P_{K_1}v$.  This gives the first formula
in \eqref{paper:exact-trivial-compression}; the two summands are
orthogonal, proving the second.
\end{proof}

\begin{lemma}[Integral Weyl columns and orbit averages]
\label{paper:integral-orbits}
Let an ordinary integral modular category have square dimension,
and let $t$ be an odd prime. Let $\mathcal O$ be a set of $e\ge2$ labels
with $t$-inactive twists, character fields unramified at $t$, and
$t\mid f_Z$ for every $Z\in\mathcal O$. Suppose all differences
$\mathbf e_Z-\mathbf e_{Z'}$ are fixed by $W_t$, and
$tW_t\mathbf e_Z$ is an integral vector for one (hence every)
$Z\in\mathcal O$. Then there are integers $k>0$ and $m_A$ such that
\begin{equation}\label{paper:coalescence}
 tW_t\mathbf e_Z
 =t\mathbf e_Z-k\sum_{V\in\mathcal O}\mathbf e_V
                  +\sum_{A\notin\mathcal O}m_A\mathbf e_A,
 \qquad
 2tk=e k^2+\sum_{A\notin\mathcal O}m_A^2.
\end{equation}
At least one active $m_A$ is nonzero, and $e\le2t-1$.

If the dimension is odd, the local cyclic module generated by
$\mathcal O$ factors through $\operatorname{SL}_2(\mathbb F_t)$.
For every $b\in(\mathbb Z/t)^\times$ one has
\begin{equation}\label{paper:phase-squares}
 \sum_{\theta_{A,t}=\zeta_t^b}m_A^2=k,
 \qquad k\in2\mathbb Z_{>0},
\end{equation}
and
\begin{equation}\label{paper:inactive-squares}
 e k^2+
 \sum_{\substack{A\notin\mathcal O\\\theta_{A,t}=1}}m_A^2
       =k(t+1).
\end{equation}
In particular $e\le(t+1)/2$. The same nontrivial phase fibers
satisfy the signed degree identities
\begin{equation}\label{paper:linear-phase}
 \sum_{\theta_{A,t}=\zeta_t^b}m_A d_A
 =\sum_{\substack{A\notin\mathcal O\\\theta_{A,t}=1}}m_A d_A
       -k\sum_{Z\in\mathcal O}d_Z.
\end{equation}
\end{lemma}
\begin{proof}
Unramified character fields make every full $t$-inertia permutation
fix each source label. Since the dimension is a square, these local
tori are unsigned. The local Galois identity makes the columns
$W_t\mathbf e_Z$ rational and fixed by those tori.
The vector $w=W_t\mathbf e_Z-\mathbf e_Z$ is independent of $Z$.
Symmetry makes its entries on $\mathcal O$ equal. Norm comparison
therefore gives $w_Z=-\|w\|^2/2$. The Weyl-line lemma excludes
$w=0$. Integrality of the scaled columns now gives
\eqref{paper:coalescence}, with $k>0$.

If the whole Weyl column were inactive, both upper and lower
unipotents would fix its source; generation of the local group
would give a Weyl-fixed coordinate line, again contradicting
$t\mid f_Z$. Thus some active $m_A$ is nonzero. Consequently
\[
 e\le\frac{2tk-1}{k^2}\le2t-1.
\]

Now suppose the dimension is odd. The coefficients $m_A$ are
constant on full $t$-inertia orbits. A nonzero coordinate with
actual $t$-twist of order at least $t^2$ would contribute at least
$t(t-1)$ to their squared norm, by
Corollary~\ref{paper:active-inertia}. But
\[
 \sum m_A^2=2tk-e k^2\le t^2/e\le t^2/2<t(t-1).
\]
Thus every supported actual primary twist has order dividing $t$.
It follows that $U_t(t)$ fixes both $\mathbf e_Z$ and
$W_t\mathbf e_Z$. Hence the lower unipotent with parameter $t$
fixes $\mathbf e_Z$ too. All local tori already fix it. The
lower--diagonal--upper factorization of the first congruence
kernel proves kernel fixedness; normality proves the assertion
on the entire cyclic module.

Write $v=\mathbf e_Z$. In the resulting prime-level representation,
$v$ is Borel-fixed. For $a\ne0$ in $\mathbb F_t$, Bruhat gives
\[
 \langle W_tv,U_t(a)W_tv\rangle
 =\langle v,L_t(-a)v\rangle
 =\langle v,W_tv\rangle=1-k/t.
\]
At $a=0$ the coefficient is one. Fourier inversion on the additive
group $\mathbb F_t$ gives squared norm $k/t^2$ in every nontrivial
upper-unipotent eigenspace. This proves the first identity in
\eqref{paper:phase-squares}. Partial inversion has no fixed active
label in odd dimension. It preserves each such phase fiber, and
the rational column has equal coefficients on its two-cycles.
Each sum in \eqref{paper:phase-squares} is therefore even, so $k$
is positive and even. Subtracting the active norm $k(t-1)$ from
\eqref{paper:coalescence} proves \eqref{paper:inactive-squares}.
Thus $e k\le t+1$, and $k\ge2$ gives the stated orbit bound.

For the signed identities, the dimension functional
$\delta(\mathbf e_A)=d_A$ is fixed by the lower local unipotents:
its vector is $n s\mathbf e_0$, and the upper unipotents fix
the vacuum. For a unit $a$, the factorization
\[
 U_t(a)W_t=L_t(a^{-1})H_t(a)U_t(-a^{-1})
\]
therefore gives $\delta(U_t(a)W_tv)=d_Z$.
Write $C$ for the right side of \eqref{paper:linear-phase} and
$M_b$ for its left side. Substitution of the integral column yields
\[
 C+\sum_{b\ne0}M_b\zeta_t^{ab}=0\qquad(a\ne0).
\]
Fourier inversion makes all $t$ coefficients equal, so $M_b=C$.
\end{proof}

\begin{lemma}[Energy and exact transport under the braid relation]
\label{paper:braid-rank}\label{paper:braid-energy}
Let $W,U$ be unitary operators satisfying
\[
                         WUW=U^{-1}WU^{-1}.
\]
Let $P,Q$ be orthogonal projections with $PQ=0$, commuting
with $U$ and $W^2$, of positive ranks $m,n$. Put
\[
 \begin{gathered}
 R=I-P-Q,\qquad V=UW,\qquad
 Z=\sum_{j=0}^2V^j(nP-mQ)V^{-j},\\
 L_P=\|PWR\|_{\mathrm{HS}}^2,\qquad
 L_Q=\|QWR\|_{\mathrm{HS}}^2.
 \end{gathered}
\]
Then
\begin{equation}\label{paper:braid-energy-identity}
 \|PWQ\|_{\mathrm{HS}}^2+
 \frac{n^2L_P+m^2L_Q+\|Z\|_{\mathrm{HS}}^2/6}{(m+n)^2}
                       =\frac{3mn}{2(m+n)}.
\end{equation}
In particular an equal-rank crossblock has squared
Hilbert--Schmidt norm at most $3m/4$.

The following exact-rank conclusion needs no ambient
$W^2$-invariance. Let $E_A,E_B$ be mutually orthogonal
coordinate spaces preserved by $U$, and suppose
$A\subseteq E_A$, $B\subseteq E_B$ satisfy $WA=B$ and
$W^2A=A$. If
\[
 d=\dim A=\dim B,\qquad
 c_A=\dim E_A-d,\qquad c_B=\dim E_B-d,
\]
then $d\le c_A+c_B$.
\end{lemma}
\begin{proof}
The braid identity gives $V^3=W^2=C$. Since $C$ commutes
with $W$ and $V$, it commutes with $U$ too. Conjugation by
$V$ cyclically permutes the three conjugates of each projection.
Commutation with $U$ shows that all three off-diagonal
$P$-overlaps equal $a=\|PWP\|_{\mathrm{HS}}^2$;
the corresponding $Q$-overlaps equal
$b=\|QWQ\|_{\mathrm{HS}}^2$. All six cross-overlaps equal
$s=\|PWQ\|_{\mathrm{HS}}^2$, and the corresponding
$P,Q$ overlaps are zero. Hence
\[
 \|Z\|_{\mathrm{HS}}^2
 =3(n^2m+m^2n)+6(n^2a+m^2b-2mns).
\]
Unitarity gives $a+s+L_P=m$ and $b+s+L_Q=n$.
Substitution proves \eqref{paper:braid-energy-identity}.

For the rank assertion, the hypotheses imply $WB=A$.
Unitarity gives $P_AW=WP_B$, so compression of $WUW$ to
$A$ is, through the two Weyl identifications, $P_BU|_B$.
Its kernel has dimension at most $c_B$, so its rank is at
least $d-c_B$. On the other side of the braid relation,
$P_{E_A}W|_{E_A}$ kills $A$ and has rank at most $c_A$.
Since $U^{-1}$ preserves $E_A$, compression of
$U^{-1}WU^{-1}$ to $A$ factors through this operator.
Thus $d-c_B\le c_A$.
\end{proof}

\begin{lemma}[Torus conductor and two unipotent depths]
\label{paper:two-depth-compression}
Let $\rho$ be a unitary representation of
$\operatorname{SL}_2(\mathbb Z/\ell^a)$, with upper and lower
unipotents $U(x),L(x)$, torus $H(t)=\operatorname{diag}(t,t^{-1})$
and Weyl element $W$. Write $U_j=\{U(x):\ell^j\mid x\}$,
and let $P_{U_j}$ be its orthogonal averaging projector.
Suppose $v$ is fixed by $U_c$ and is a torus eigenvector
whose character has exact conductor $\ell^\beta$, where
$2\le\beta\le a$ and $0\le c<\beta$.
For every $d\ge0$ with $c+d<\beta$, one has
\[
 \|P_{U_d}Wv\|^2
 =\ell^{-(\beta-c-d)}\|P_{U_{\beta-c}}Wv\|^2
 \le\ell^{-(\beta-c-d)}\|v\|^2.
\]
In particular, a fully upper-fixed vector satisfies
$\|P_{U_{\beta-1}}Wv\|^2\le\|v\|^2/\ell$,
regardless of the ambient depth $a$.
\end{lemma}
\begin{proof}
Put $w=Wv$; it is fixed by $L_c$ and has inverse torus
character $\psi$. Set $f(x)=\langle w,U(x)w\rangle$.
For $t\in\ell^c\mathbb Z/\ell^a$ with $1+tx$ a unit,
Gauss factorization gives
\[
 L(t)U(x)=U\!\left(\frac{x}{1+tx}\right)
 H\!\left(\frac1{1+tx}\right)
 L\!\left(\frac{t}{1+tx}\right),
\]
and hence
\[
 f(x)=\psi\!\left((1+tx)^{-1}\right)
                       f\!\left(\frac{x}{1+tx}\right).
\]
Fix $d\le j<\beta-c$ and $s\in\mathbb F_\ell^\times$.
On the shell $x/\ell^j\equiv s\pmod\ell$, choose
$t=\ell^{\beta-j-1}t_0$ with $t_0$ a unit.
This lies in $\ell^c\mathbb Z/\ell^a$.
The map $x\mapsto x/(1+tx)$ permutes this shell,
while its torus multiplier is constant and nontrivial:
$tx$ has valuation $\beta-1$, and $\psi$ is nontrivial
on the last unit quotient at that level.
The sum of $f$ on this shell is therefore zero.
All shells below depth $\beta-c$ vanish.
The average over $\ell^d\mathbb Z/\ell^a$ thus equals
$\ell^{-(\beta-c-d)}$ times the average over
$\ell^{\beta-c}\mathbb Z/\ell^a$.
For either subgroup this average is
$\langle w,P_Uw\rangle=\|P_Uw\|^2$, proving the identity.
\end{proof}

\begin{lemma}[One prime-level quadratic phase class]
\label{paper:half-channel-local}
Let $\ell$ be an odd prime, let $a\ge1$, and let $\rho$ be a
unitary representation of $G=\operatorname{SL}_2(\mathbb Z/\ell^a)$.
Write $U_j=\{U(x):\ell^j\mid x\}$, let $H$ be the full diagonal
torus, let $W$ be a Weyl operator, and let $K_1$ be the kernel
of reduction modulo $\ell$. Suppose $v$ is fixed by $U_0$ and $H$.
For one quadratic class $C\subset\mathbb F_\ell^\times$, let $P_C$
project onto the upper-unipotent characters
$x\mapsto\zeta_\ell^{kx}$ with $k\in C$. Then
\[
 \|P_CWv\|^2
  =\frac{\ell^2-1}{2\ell^2}
       \|(1-P_G)P_{K_1}v\|^2
  \le\frac{\ell^2-1}{2\ell^2}\|v\|^2.
\]
For $a=1$, interpret $K_1$ as the trivial group.
\end{lemma}
\begin{proof}
Every element of $K_1$ has a unique lower--diagonal--upper
factorization, with parameters in $\ell\mathbb Z/\ell^a$.
Averaging this factorization on $v$ gives
$P_{K_1}v=P_{L_1}v$. Consequently
\[
 P_{U_1}Wv=WP_{L_1}v=WP_{K_1}v.
\]
The last vector lies in the level-$\ell$ subrepresentation,
and $P_C\le P_{U_1}$, so it remains to work with
$\operatorname{SL}_2(\mathbb F_\ell)$ and a Borel-fixed vector.

The permutation representation on $\mathbb P^1(\mathbb F_\ell)$
is the direct sum of the trivial representation and its irreducible
orthogonal complement: double transitivity makes the commutant
have dimension two. The latter is the Steinberg representation.
By Frobenius reciprocity these are precisely the irreducible
representations with a Borel-fixed vector, each with a
one-dimensional fixed space. It therefore suffices to calculate
on their fixed unit vectors.

The trivial vector has no nontrivial upper-unipotent component.
In the permutation model a Steinberg fixed unit vector is
\[
 v_0=\frac{\ell e_\infty-\sum_{x\in\mathbb F_\ell}e_x}
                 {\sqrt{\ell(\ell+1)}}.
\]
After applying $W$, the coefficient at $0$ is
$\ell/\sqrt{\ell(\ell+1)}$ and every other finite coefficient
is $-1/\sqrt{\ell(\ell+1)}$. Projection onto any nontrivial
additive character has squared norm $(\ell+1)/\ell^2$.
There are $(\ell-1)/2$ characters in $C$, giving the stated
factor. The calculation is unchanged on multiple copies, since
it acts as that scalar on the Borel-fixed multiplicity space.
\end{proof}

\begin{corollary}[A half-budget for prime-level activity]
\label{paper:half-channel-crossing}
Use the notation of Theorem~\ref{paper:inactive-fourier-budget}.
Let $\ell$ be an odd prime unramified in $K$. Suppose one input
is $\ell$-inactive and the other has actual $\ell$-primary twist
of exact order $\ell$. Then every Fourier channel satisfies
\[
 ab|\widehat w(\chi)|^2
       \le\frac{\ell^2-1}{2\ell^2}\frac D{c^2},
 \qquad
 abE(w)\le\frac{\ell^2-1}{2\ell^2}\frac D{c^2}.
\]
\end{corollary}
\begin{proof}
Take the inactive input as the source of the normalized orbit
vector. Since $K$ is unramified at $\ell$, every channel has
trivial inflated $\ell$-torus character, so the source is
Borel-fixed at $\ell$. Complementary Weyl operators commute
with this local group and preserve that property.
The full Galois orbit of the other input has its actual
$\ell$-twists in one quadratic class, by the square action on
twists. Its coordinate projection is therefore dominated by the
projection in Lemma~\ref{paper:half-channel-local}. Combine that
lemma with the exact orbit-projection normalization
\eqref{paper:orbit-projection}. Summing the channel inequalities
and using Parseval proves the energy bound.
\end{proof}


\begin{lemma}[A nonquadratic tame torus character]
\label{paper:nonquadratic-half-local}
Use the local notation of Lemma~\ref{paper:half-channel-local}.
Suppose $v$ is fixed by $U_0$ and transforms under the full torus
by a character $\chi$ factoring through $\mathbb F_\ell^\times$,
with $\chi^2\ne1$. Then for either quadratic class $C$,
\[
       \|P_CWv\|^2
        =\frac{\ell-1}{2\ell}\|P_{K_1}v\|^2
        \le\frac{\ell-1}{2\ell}\|v\|^2.
\]
\end{lemma}
\begin{proof}
The Gauss-factorization argument is unchanged, since the torus
character is trivial on the first congruence subgroup. At level
$\ell$, the induced representation from the Borel character
$\chi$ is irreducible: the two Bruhat double cosets contribute
one intertwiner from the identity coset and none from the Weyl
coset, the latter because it would identify $\chi$ with
$\chi^{-1}$. Its dimension is $\ell+1$ and its Borel
$\chi$-space is a line. In the standard coset model, this line
is $\mathbb Ce_\infty$. The Weyl operator takes $e_\infty$
to a unit multiple of $e_0$, and the upper unipotent acts by
translations on the affine coordinates. Each additive character
therefore has squared projection $1/\ell$. Summing over the
$(\ell-1)/2$ elements of $C$ proves the identity. Frobenius
reciprocity applies this scalar calculation to every multiplicity
space of an arbitrary representation.
\end{proof}


\begin{lemma}[A primitive torus character at arbitrary prime-power level]
\label{paper:primitive-torus-half}
Let $\ell$ be odd and let $v$ belong to a unitary representation of
$\operatorname{SL}_2(\mathbb Z/\ell^A)$. Suppose $v$ is fixed by
the full upper unipotent and has torus character $\chi$ such that
both $\chi$ and $\chi^2$ have exact conductor $\ell^a$, $a\le A$.
Let $K_a$ be the kernel of reduction to level $\ell^a$. For
$1\le c\le a$, and one square class $C$ of primitive
$\ell^c$-th additive characters, one has
\[
 \|P_CWv\|^2
 =\frac{\varphi(\ell^c)}{2\ell^a}\|P_{K_a}v\|^2
 \le\frac{\ell-1}{2\ell}\|v\|^2.
\]
The inactive upper-unipotent eigenspace instead has squared
projection $\ell^{-a}\|P_{K_a}v\|^2$.
\end{lemma}
\begin{proof}
Lower--diagonal--upper factorization of $K_a$ gives
$P_{K_a}v=P_{L_a}v$, since its diagonal and upper factors fix
$v$. Consequently $P_{U_a}Wv=WP_{K_a}v$. Every target space
in the assertion is contained in the image of $P_{U_a}$, so
it suffices to work at level $\ell^a$.

Let $B$ be the upper triangular subgroup and extend $\chi$
trivially over its unipotent radical. The induced representation
$\operatorname{Ind}_B^G\chi$ is irreducible. Indeed, Mackey's
intertwining criterion has the identity double coset, the big
Bruhat cell, and intermediate representatives
$g=L(y)$ with $y=\ell^ku$, $1\le k<a$, and $u$ a unit.
The big cell would require $\chi=\chi^{-1}$, which is excluded.
For an intermediate representative put
$b=U(\ell^{a-k-1}u^{-1})$. Then $b\in B$ and
$g^{-1}bg\in B$, with diagonal entry $1+\ell^{a-1}$;
$\chi(b)=1$ while $\chi(g^{-1}bg)\ne1$. Thus that
intertwining space vanishes as well. This argument also applies
when intermediate double cosets split by unit square class.

Frobenius reciprocity now says that this is the only irreducible
representation containing the prescribed $B$-character, whose
multiplicity inside it is one. In the projective-line coset
model its $B$-line is $\mathbb Ce_\infty$, and $W$ maps that
line to $\mathbb Ce_0$. The upper unipotent acts regularly by
translation on the $\ell^a$ affine coordinates. Every additive
character therefore receives squared norm $\ell^{-a}$.
There are $\varphi(\ell^c)/2$ characters in the selected class.
The identity extends unchanged across arbitrary multiplicity
spaces, proving the assertion.
\end{proof}


\begin{lemma}[Dyadic inactive columns and the equality case]
\label{paper:dyadic-eight-bound}
Let $\rho$ be a unitary representation of a dyadic local modular
group, with Weyl operator $W$, full diagonal torus $H$, and
normalized upper unipotent $U=T/\gamma$, where $\gamma^8=1$.
Let $V_0$ be the space fixed by $H$ on which the actual twist
$T$ equals one, and let $P_0$ be its orthogonal projector.
Then, for orthogonal unit vectors $u,v\in V_0$,
\[
                         |\langle u,Wv\rangle|\le\frac34.
\]
For a nonidentity eighth root $\theta$, let $P_\theta$ be the
actual-twist projector. For every $v\in V_0$,
\[
                        \|P_\theta Wv\|^2\le\frac34\|v\|^2.
\]
If equality holds for nonzero $v$, then $\gamma\in\{1,-1\}$,
$\theta=-1$, and
\[
                        P_0Wv=-\frac\gamma2v.
\]
For every other pair $(\gamma,\theta)$ the second bound is
$\|P_\theta Wv\|^2\le\frac12\|v\|^2$.
\end{lemma}
\begin{proof}
Inflate the local level to at least eight if necessary. Let $K_8$
be the kernel of reduction modulo eight. Since the upper
unipotents with parameter divisible by eight and the congruence
torus fix $v$, lower--diagonal--upper factorization gives
$P_{K_8}v=P_{L_8}v$. Consequently
$P_{U_8}Wv=WP_{K_8}v$. Every target in the statement is
$U_8$-fixed. Hence the indicated compressions vanish on the
non-$K_8$ part of $V_0$, and it suffices to work over
$G=\operatorname{SL}_2(\mathbb Z/8)$.

Let $B$ be its upper triangular subgroup. On $V_0$ its character
is
\[
 \tau_\gamma\begin{pmatrix}a&b\\0&a^{-1}\end{pmatrix}
                           =\gamma^{-b/a}.
\]
This is a character because every odd unit squares to one modulo
eight. Frobenius reciprocity allows all the required operator
bounds to be checked on the twelve-dimensional induced module
$\operatorname{Ind}_B^G\tau_\gamma$: it contains every
irreducible with a $B$-vector of this character, with its full
character multiplicity.

The following explicit induced matrices give an exact finite
verification; the computational certificates also record these matrices.
Use coset representatives
\[
 g_x=\begin{pmatrix}x&-1\\1&0\end{pmatrix}\ (x\bmod8),
 \qquad h_y=\begin{pmatrix}1&0\\2y&1\end{pmatrix}\ (y\bmod4).
\]
For $g\in G$, the induced matrix has entry
$\tau_\gamma(g_i^{-1}gg_j)$ in position $(i,j)$ when
$gg_j\in g_iB$, and zero otherwise. Thus all its entries lie
in $\mathbb Q(\zeta_8)$. Set
\[
 P_B=\left(\frac14\sum_{a\in(\mathbb Z/8)^\times}H_a\right)
        \left(\frac18\sum_{j=0}^7\gamma^jU^j\right),
 \quad A=P_BWP_B,\quad C_\theta=P_BWP_\theta WP_B.
\]
Direct matrix multiplication gives
\[
\begin{array}{c|c}
\gamma&\text{annihilating polynomial of }A\\ \hline
1&t(2t+1)(t-1)\\
-1&t(t+1)(2t-1)\\
\text{otherwise}&t(2t+1)(2t-1).
\end{array}
\]
For $(\gamma,\theta)=(\pm1,-1)$ one has
\[
 C_\theta^2=\frac34C_\theta,
 \qquad AC_\theta=-\frac\gamma2C_\theta;
\]
for every other nonidentity $\theta\in\mu_8$ one has
\[
                    C_\theta(4C_\theta-I)(2C_\theta-I)=0.
\]
The supplied check verifies these identities for all eight
values of $\gamma$ and seven values of $\theta$ using exact
arithmetic in $\mathbb Q(\zeta_8)$, with no numerical eigenvalue
tests. The operators $A$ and $C_\theta$ are self-adjoint, and
$C_\theta$ is positive. Their spectral bounds and the equality
assertion follow from these polynomial identities. For the first
bound, subtract the midpoint of the spectral interval of $A$;
its interval has length at most $3/2$, and orthogonality makes
the scalar midpoint contribute zero. Adding a zero operator on
the non-$K_8$ part does not enlarge any displayed spectral
interval, completing the general-level proof.
\end{proof}


\subsection{Fermion orders and reduced local operators}

\begin{lemma}[An integral extension of a minimal fermion category]
\label{paper:integral-extension}
Let $\mathcal F$ be integral and slightly degenerate, with pointed
subcategory $\langle g\rangle=\mathrm{sVec}$ and no other proper
nontrivial fusion subcategory. Suppose $\dim\mathcal F>2$.
It has an integral minimal modular extension $\mathcal M$.
Every minimal extension has pointed subcategory $\mathrm{sVec}$,
adjoint subcategory $\mathcal F$, and no other proper subcategory
besides $\mathrm{Vec},\mathrm{sVec},\mathcal F$.
Noninvertible simples of coprime dimensions in an integral extension
have zero $S$-entry.
\end{lemma}
\begin{proof}
The minimal-extension theorem \cite[Corollary~1.3]{JFR} gives a
minimal modular extension $\mathcal M_0$ of dimension
$2\dim\mathcal F$. Its fermion-monodromy grading has neutral
component $\mathcal F$. An odd invertible would generate with $g$
a nondegenerate pointed subcategory of dimension four: its
bicharacter pairs $g$ with that invertible by $-1$; in the cyclic
case the self-pairing has square $-1$. Splitting this subcategory
off would split $\mathcal F$ as $\mathrm{sVec}\boxtimes\mathcal N$
with $\dim\mathcal N>1$, contrary to its subcategory hypothesis.
Thus $(\mathcal M_0)_{\rm pt}=\mathrm{sVec}$ and
$(\mathcal M_0)_{\rm ad}=\mathcal F$.
Any proper subcategory crossing the grading has dimension twice
that of its intersection with $\mathcal F$, hence dimension two
or four. It is pointed or Ising. The former supplies an odd
invertible; the latter splits off nondegenerately and gives the
same forbidden decomposition of $\mathcal F$.

For completeness, integrality of an extension can be arranged.
At this stage $\mathcal M_0$ is only weakly integral. Its anomaly
$\alpha=p^+/p^-$ satisfies $\alpha^8=1$ by
\cite[Theorem~6.10]{DLNCongruence}. Put
$\xi=p^+/\sqrt{\dim\mathcal M_0}$, so $\xi^2=\alpha$ and
$\xi^{16}=1$. Since cubing is an automorphism of $\mu_{16}$, choose
$\gamma\in\mu_{16}$ with $\gamma^3=\xi$. The linear modular
representation with normalized diagonal $t=T/\gamma$ has the
congruence and Galois properties of \cite[Theorem~II]{DLNCongruence}.
Thus odd inertia fixes the normalized vacuum eigenvalue $\gamma^{-1}$.
The stronger integral bound $\alpha^4=1$ of
\cite[Proposition~6.7(ii)]{DLNCongruence}, and the corresponding choice
$\gamma\in\mu_8$, are not used before integrality has been obtained. The Galois image of
the vacuum is invertible and has actual twist one, so it is the
vacuum rather than $g$. The degrees are consequently unramified
at every odd prime. As their squares are integers, each degree is
an integer or $\sqrt2$ times an integer. If $\mathcal M_0$ is not
integral, its integral subcategory is exactly $\mathcal F$.
Take an Ising category with fermion $\psi$ and noninvertible
$\sigma$, and condense the diagonal boson $g\boxtimes\psi$
in $\mathcal M_0\boxtimes\mathcal I$. The local-module category
has dimension $2\dim\mathcal F$ and contains $\mathcal F$
fully faithfully. An odd object paired with $\sigma$ has dimension
$2n$; a free boson orbit gives a simple of dimension $2n$, while
a fixed orbit splits into two simples of dimension $n$.
All resulting simples are integral.

Finally, for any noninvertible $X$ in this extension,
$\langle XX^*\rangle=\mathcal F$: containment in $\mathrm{sVec}$
would give $d_X^2\le2$. A nonzero entry between coprime-dimensional
noninvertibles implies projective centrality by the algebraic-norm
argument. It would put one object in
$\mathcal F'=\mathrm{sVec}$, a contradiction.
\end{proof}

\begin{lemma}[Positive and signed fermion orders]
\label{paper:fermion-orders}
Let $\mathcal M$ be integral modular with fermion $g$, and put
$\mathcal F=\langle g\rangle'$, $N=\dim\mathcal M=4\Delta$.
Using even fermion-orbit representatives, form
\[
                 R_\pm=K_0(\mathcal F)/(g\mp1).
\]
They are symmetric integral orders with the same reduction modulo
two. A full modular character restricts to $R_+$ or $R_-$
according as its label is even or odd, respectively. Restriction
fibers are exactly fermion orbits, and the restricted codegree is
\begin{equation}\label{paper:fermion-codegrees}
 f_X^{\pm}=\begin{cases}f_X/2,&gX=X,\\ f_X/4,&gX\ne X,
                 \end{cases}\qquad f_X=N/d_X^2.
\end{equation}
At two, a restricted character is principal if and only if its
full character is principal. The corresponding positive and
signed blocks have equal rank.
\end{lemma}
\begin{proof}
The even fermion action is free and does not identify a label with
its dual, by opposite twists. Thus both quotients have basis
representatives with unit-coefficient trace
$\tau(XY^*)=\delta_{XY}$. Their multiplication tables agree
modulo two. With $\varepsilon_Z$ the fermion-monodromy sign,
$S_{gX,Z}=\varepsilon_ZS_{XZ}$. Orthogonality with the even
projector therefore gives
\[
 \sum_{Z\text{ even}}\chi_X(Z)\overline{\chi_Y(Z)}
   =\frac{f_X}{2}(\delta_{XY}+\delta_{X,gY}).
\]
Every character of the even subalgebra extends by semisimplicity;
the displayed identity identifies its fibers and, after dividing
the even sum into pairs, gives the codegree formulas.
If a full character has dimension residue on even objects, each
odd $Y$ satisfies $Y^2$ even. Squaring is injective in the residue
field of characteristic two, so the dimension residue also holds
on $Y$. Finally identical special fibers have the same ranks of
the lifted local blocks.
\end{proof}

\begin{lemma}[Principal representatives at an odd prime]
\label{paper:odd-principal-lift}
At a fixed place over an odd prime, a positive principal fermion
pair has a full principal representative. Both representatives
are principal if all odd degrees vanish in the residue field;
otherwise exactly one is principal.
\end{lemma}
\begin{proof}
For one even representative $X$, write $a_Y$ for the residue of
$\chi_X(Y)$ and $b_Y$ for that of $d_Y$. They agree on even
objects. For odd $Y,Z$, their products are even, so
$a_Y^2=b_Y^2$ and $a_Ya_Z=b_Yb_Z$. If every odd $b_Y$ is zero,
so is every odd $a_Y$. Otherwise choose $Z$ with $b_Z\ne0$.
Then $a_Z/b_Z=\pm1$ and $a_Y$ is that same sign times $b_Y$
for all odd $Y$. Replacing $X$ by $gX$ reverses every odd
evaluation and fixes the even evaluations, proving the assertion.
\end{proof}

\paragraph{Fermion coordinate spaces.}
\label{paper:fermion-dictionary}
Let $Fv_X=v_{gX}$ and $Gv_X=\varepsilon_Xv_X$. Denote their joint
eigenspaces by
\[
 A=(+,+),\quad B=(-,+),\quad C=(+,-),\quad H=(-,-).
\]
The global normalized matrix $s=S/\sqrt N$ exchanges $F,G$.
Write $W=W_2$, $W_o=W_{2'}$, and $U=\lambda T_2$ with
$\lambda^8=1$. Then $W$ preserves $A,H$ and interchanges $B,C$;
$U$ interchanges $A,B$ and preserves $C,H$. Set $K=W|_A$.
The positive normalized matrix is
$\widehat s=s|_A=K W_o|_A$.
For an even pair its basis vectors are
$(v_X\pm v_{gX})/\sqrt2$; a free odd orbit has the analogous
symmetric and antisymmetric vectors. On free pair coordinates the
$B,C$ and $H,H$ entries of $s$ are $2S_{XY}/\sqrt N$.
The local dyadic group preserves every lifted two-block, and its
restriction $K$ preserves positive two-blocks.


\begin{corollary}[Defect one in the fermion orders]
\label{paper:dyadic-node}
The conductor-node and defect-one results
(Lemma~\ref{paper:2.9} and Proposition~\ref{paper:2.10}) apply
to both fermion orders $R_\pm$.  Thus a factor with conductor
exponent $c_i=e_iv_2(f_i)-d_i=0$ constitutes the whole block;
one with $c_i=1$ gives exactly two factors, both of conductor
exponent one, joined over their residue field.
A defect-one character has at most two inertia orbits in its
block.  A ramified branch is quadratic and constitutes the
whole block; otherwise a branch is unramified, and no
unramified higher-defect character belongs to that block.
\end{corollary}
\begin{proof}
Lemma~\ref{paper:fermion-orders} supplies symmetric integral orders
with cyclotomic characters and rational codegrees.  The cited
order-theoretic proofs require these properties, not positivity
of their structure constants, so apply unchanged to $R_-$ as
well as $R_+$.
\end{proof}
\begin{corollary}[Restriction to positive reduced data]
\label{paper:lw:reduced-block}
For a full $\ell$-block $B$ of the positive reduced table
$\widehat s$ of an integral slightly degenerate category,
\[
 \sum_{Z\in B}\overline{\widehat s_{Zi}}\widehat s_{Zj}=0
\]
when the odd primary twists at $i,j$ differ; at two it holds
when their squared dyadic twists differ.  At odd $\ell$ the
reduced block subspace is invariant under the whole local group.
\end{corollary}
\begin{proof}
Use any weakly integral minimal modular extension supplied by
\cite[Corollary~1.3]{JFR}. The ordinary local-block theorem
already applies to weakly integral modular categories, so no
integrality hypothesis on the extension is needed, and the
original fermion category need not satisfy the minimality
assumption of Lemma~\ref{paper:integral-extension}.
Select all its labels whose characters on the even fusion
subring have the prescribed residue of $B$.  Their coordinate
projector $P$ is a sum of full ordinary block projectors, since
equality of characters on the full ring implies equality on the
even subring.  A label and its fermion translate have the same
character on the even subring.  Hence $P$ preserves the positive
fermion subspace and restricts there to $P_B$.  The normalized
ordinary $s$ preserves that subspace and restricts to
$\widehat s$.  Odd primary twists preserve it as well, and the
two members of a fermion pair have equal odd primary twists and
equal squared dyadic twists.  Restrict
$[P,sT_\ell s^{-1}]=0$ at odd $\ell$, and its squared-operator
version at two.  These give the displayed orthogonality.
At odd $\ell$ the upper and lower generators both preserve the
positive subspace, so ordinary local block invariance restricts
to the asserted reduced invariance.
\end{proof}

\subsection{Twists, phase masses and omitted phases}

\begin{proposition}[Cauchy and vanishing at odd primes]\label{paper:3.4}
For an integral slightly degenerate category and any odd $p\mid\Delta$,
some simple has $p$-active twist. All such labels have zero $S$-entry on
rows of $p$-unit codegree.
\end{proposition}
\begin{proof}
If all even labels of a minimal extension were $p$-inactive,
the $p$-local upper unipotent would fix the positive fermion subspace
pointwise (the anomaly is dyadic). The lower unipotent would too,
since $s$ preserves this subspace. The fixed space of the local
factor is invariant under the entire modular representation and
contains $v_0+v_g$, hence also $T(v_0+v_g)=v_0-v_g$ and $sv_0$.
Every coordinate of $sv_0$ is nonzero. Thus the $p$-primary twist
would be trivial everywhere, contrary to Cauchy's theorem in the
extension \cite[Theorem~3.9]{BNRWFinite}. Vanishing follows from
block orthogonality on the singleton block of a $p$-unit codegree.
\end{proof}

\begin{theorem}[Twist orders and formal codegrees]\label{paper:3.5}
In a weakly integral modular category put $m_z=\ord(\theta_z)$ and $f_z=D/d_z^2$. Then
\[
 (m_z)_{\rm odd}\mid f_z,\qquad
 m_z\mid \begin{cases}4f_z&\text{in the integral case},\\8f_z&\text{in general},\end{cases}
 \qquad \ord(\theta_x/\theta_y)\mid4\lcm(f_x,f_y),\qquad N\mid4D.
\]
The odd-part assertion holds also for an integral slightly degenerate category with reduced codegree $(D/2)/d_z^2$.
\end{theorem}
\begin{proof}
The lift $t=T/\gamma$ has $\ord(\gamma)\mid8$ in the integral case and $\ord(\gamma)\mid16$ in general \cite[Proposition~6.7 and Theorem~6.10]{DLNCongruence}. If $p^k$ is the exact $p$-part of $\ord(t_{zz})$, each automorphism fixing $K_z$ fixes $z$, so its square fixes $t_{zz}$ \cite[Theorem~II]{DLNCongruence}. At odd $p$ its exponent is $\pm1\bmod p^k$; at two and $k\ge4$ it is $\pm1\bmod 2^{k-1}$. Thus $K_z$ contains the real subfield $F$ at conductor $p^a$, with $a=k$ for odd $p$ and $a=k-1$ at two.

The conductor of $F$ is $p^a$ unless $F=\QQ$. By
Corollary~\ref{paper:family:2.6}, at odd $p$ we obtain
$p^k\mid f_z$; the exception $p^k=3$ follows from
Lemma~\ref{paper:family:lw:weyl-line}. At two and $k\ge4$ we get
$2^{k-2}\mid f_z$. For an actual twist whose two-part exceeds
$\ord(\gamma)$ this gives the same estimate. Otherwise use
$2\mid f_z$ for every nontrivial actual two-part, again by
Lemma~\ref{paper:family:lw:weyl-line}. This gives the two objectwise bounds.

For a ratio with two-part $2^k$, $k\ge4$, one of $t_{xx},t_{yy}$ has two-part at least $2^k$, so $2^{k-2}\mid\lcm(f_x,f_y)$. For $1\le k\le3$ at least one actual twist has even order. Combined with the odd-part result this proves the ratio bound; taking $y=0$ gives $N\mid4D$ since the codegrees divide $D$. In a minimal modular extension each original label has codegree four times its reduced codegree, giving the last assertion.
\end{proof}

\begin{corollary}[The exact dyadic refinement]
\label{paper:dyadic-twist-codegree}
In ordinary integral data, a dyadic twist of order $2^v$, $v\ge4$,
has $2^{v-1}\mid f_X$ at its label $X$.
\end{corollary}
\begin{proof}
The square Galois action puts $\mathbb Q(\zeta_{2^{v-1}})^+$
in $K_X$; its conductor is $2^{v-1}$.  The normalization scalar
has order at most eight, so does not change this conclusion.
Apply Theorem~\ref{paper:family:por:th:exact-conductor}.
\end{proof}


\begin{proposition}[Odd-valuation slices]\label{paper:3.8}
Suppose $p$ and $v_p(D)$ are odd. Let $B$ be a full $p$-block or the whole label set. Let $x,y$ have trivial $p$-twists, with $y$ fixed by $p$-inertia. For a complementary twist $\eta$ put
\[
 M_j=\sum_{\substack{i\in B,\ \theta_{i,p'}=\eta\\
                 \ord(\theta_{i,p})=p^j}}s_{xi}\overline{s_{yi}}.
\]
Then $M_1=(p-1)M_0$, also after summing over $\eta$. The same formula
holds in integral positive reduced data with $D$ replaced by $\Delta$,
provided the complementary fiber is taken up to the fermion sign. Thus on each such fiber the squared-dimension mass of order-$p$ twists is $p-1$ times the inactive mass, and the active proportion is at least $1-1/p$. For the entire label set, activity at all primes in any set $J$ of odd primes of odd valuation has proportion at least $\prod_{p\in J}(1-1/p)$, even conditional on fixed twist components outside $J$.
\end{proposition}
\begin{proof}
Let $E$ project onto $B\cap\{i:\theta_{i,p'}=\eta\}$. It commutes with the local $p$-factor by Theorem~\ref{paper:lw:block} and commutation of distinct factors; so does $Q=sEs^{-1}$ since $s$ normalizes that factor. Write $u=T_p,w,h_t$ for its upper unipotent, Weyl and torus. For a unit $n$ modulo the local level $p^a$, the $n$th weighted moment of the $p$-twist on the fiber is
\[
 (Qwu^nw^{-1})_{xy}=(Qu^{-1/n}wh_{-n}u^{-1/n})_{xy}
 =\Big(\frac{-n}{p}\Big)(Qw)_{xy}.
\]
The middle identity is Bruhat. The outer unipotents act trivially at $x,y$. Since inertia fixes $y$, the torus acts there by the Galois sign of $d_y/\sqrt D$ \cite[Theorem~II and Remark~4.5]{DLNCongruence}; at odd $v_p(D)$ this is the displayed quadratic character, from $\sqrt{(-1)^{(p-1)/2}p}$. Average over $n$: the right side vanishes, while the left gives $M_0-M_1/(p-1)$ by the means of $p$-power roots. For the reduced formula, use the restricted projector of
Theorem~\ref{paper:lw:block} and the complementary twists modulo sign.
The odd local group acts on the positive subspace, and its torus sign
is the same since the two dimension normalizations differ by two.
At $x=y=0$ the weights are $d_i^2/D$ or $d_i^2/\Delta$. The active bound holds on each complementary fiber and thus on every union of them; iteration over $J$ proves the product bound.
\end{proof}

\begin{corollary}[The two masses in a defect-one block]\label{paper:3.9}
Let $B$ be a full $p$-residue block of an integral modular category,
where $p$ is odd, and assume that $v_p(f_i)=1$ for every $i\in B$.
Then $B$ consists of one inertia orbit with trivial $p$-primary twist
and one inertia orbit with $p$-primary twist of order $p$.
If their respective lengths and dimensions are $a,r$ and $b,s$, then
\[
                        b s^2=(p-1)a r^2.
\]
\end{corollary}
\begin{proof}
The codegree identity implies that $v_p(D)$ is odd. Proposition~\ref{paper:2.10}
gives exactly two inertia orbits; their lengths divide $p-1$ and
the dimension is constant on each. Theorem~\ref{paper:3.5} bounds
all their $p$-primary twist orders by $p$. Inertia preserves these
orders. Proposition~\ref{paper:3.8}, with the
complementary condition omitted, says that both the trivial and the
order-$p$ masses are nonzero and have ratio $p-1$. Thus one orbit
is of each kind and the displayed equation follows.
\end{proof}

\begin{lemma}[The visible class at defect one]\label{paper:lem:quintic-common-class}
In an ordinary integral modular table, let $t,s$ be distinct odd primes
with $s\nmid t-1$. Suppose $\mathcal U$ is a nonempty set of
$s$-free labels of $t$-defect one and $\mathcal V$ a nonempty set of
$t$-defect-one labels in $P_s$. All these rows see only $1$ and a
single common square class of primitive $t$th roots as $t$-phases.
\end{lemma}
\begin{proof}
A defect-one block has an inactive and an active orbit; the latter
has one square class of primitive phases. Their lengths
$a_0,a_1\mid t-1$ and degrees $d_0,d_1$ satisfy
$a_1d_1^2=(t-1)a_0d_0^2$
(Proposition~\ref{paper:2.10}, Corollary~\ref{paper:3.9}).
Thus the two degrees have equal $s$-valuation. In each block of a label
in $\mathcal U$, all active representatives have $s$-unit degree, so
every $V\in\mathcal V$ hits them by principality. Its row can see
only phases in its own block: the local Weyl preserves the block
and the complementary Weyl preserves $t$-phases. The active classes
of every such pair of blocks therefore agree.
\end{proof}
\subsubsection{Indicator packets and omitted phases}
We first record the common indicator computation for the odd and dyadic phase arguments. Fix a prime $p$, write $W_o$ for the complementary Weyl, $w=W_ov_0$, and suppose a Galois-stable label set $J$ satisfies $|w_j|^2=(h/D)d_j^2\mathbf1_J(j)$ for an integer $h$. Choose $m$ zero on the $p$-level and one on the complement. The complementary upper--lower--upper identity gives $|(sT^ms^{-1})_{0j}|=|w_j|$. Thus
\[
 b_X=\nu_m(X)=\frac hD\sum_{j\in J}d_j^2\chi_j(X),
 \qquad b=\sum_X b_X[X^*].
\]
These indicators are algebraic integers by \eqref{paper:eq:indicator-entry} and rational by Galois symmetry. The $b_X$ are therefore integers, constant on character orbits by symmetry of $S$. Orthogonality gives $\chi_j(b)=h\mathbf1_J(j)$, so $b^2=hb$. Also $b=b^*$, $b_0=\|w\|^2=1$ and $\|b\|^2=h$. In each application $0\in J$, giving $\dim b=h$.


\begin{lemma}[An odd-order phase cannot be missing]\label{paper:7.4}
Let an integral perfect modular category have odd dimension $D$, with
$v_p(D)$ odd. If all $p$-primary twists have order at most $p$,
every $p$th root of unity occurs.
\end{lemma}

\begin{proof}
Write $U,W$ for the $p$-local upper unipotent and Weyl, $P$ for the projection onto $\ker(U-I)$. Then $Pw=w$, $Ww=sv_0$, and the local torus acts on $w$ by the Legendre character as in Proposition~\ref{paper:3.8}. For $1\le t<p$, Bruhat gives $PWU^tW^{-1}w=(\frac{-t}{p})PWw$. If $\zeta_p^k$ is missing, its Fourier projection vanishes. Conjugating by $W$ and applying $P$ gives $w=-G Psv_0$ with $G=\sum_{t\ne0}(\frac{-t}{p})\zeta_p^{-kt}$, $|G|=\sqrt p$. Apply the indicator computation with $J$ the $p$-inactive set and $h=p$.

The coefficient vector $p sP s^{-1}v_0$ of $b$ is fixed by the local lower unipotent and transforms under the torus by the Legendre character. It therefore has an active coordinate $Y$, since invariance under both unipotents would give trivial torus action. Twist covariance gives $\mathbb Q(\zeta_p)^+\subseteq K_Y$ and $Y\ne Y^*$ in odd dimension makes $K_Y$ nonreal. Its orbit has at least $p-1$ labels. Their equal nonzero coefficients exhaust the nonvacuum norm $p-1$, so $p=\dim b=1\pm(p-1)d_Y$, contrary to perfectness.
\end{proof}

\begin{lemma}[Twist activity in an invariant fiber]\label{paper:7.8}
Let $s$ be odd, with $v_s(\Delta)$ odd in ordinary or positive reduced
integral modular data. Every nonempty coordinate subspace invariant
under the $s$-local group contains an $s$-active label. This applies
to each fiber of a complementary twist, in either fermion sector
of an integral minimal extension, and to each squared dyadic twist
fiber in the positive quotient.
\end{lemma}
\begin{proof}
On an inactive subspace the $s$-local upper unipotent is the identity,
as are its conjugates there. They generate the local group, yet the
torus of a nonsquare unit at $s$ acts by minus a permutation since
its Galois automorphism negates $\sqrt\Delta$. Complementary twists
and the fermion grading commute with the odd local group.
\end{proof}

\begin{lemma}[Propagation from visible phases]\label{paper:lem:visible-phase-transfer}
Let $a$ be a prime and $\Phi$ a set of primary phases at a different prime. For each $a$-active column in a specified sector suppose there is a nonzero entry on a witness row, all of whose hits in that sector have phase in $\Phi$. Every fiber in the sector outside $\Phi$ is then $a$-inactive. If its coordinate span is invariant under the full $a$-local group of an ordinary modular table, every label on it has $a$-unit codegree. The same conclusion holds on a positive even fiber for odd $a$ in an integral minimal extension. In the latter assertion the fiber is defined on pairs, using the squared dyadic twist if necessary.
\end{lemma}
\begin{proof}
An active column outside $\Phi$ would contradict its witness. On such a fiber $T_a=I$, so the normalized upper unipotent acts by a scalar, possibly nontrivial at two. By invariance every local conjugate is scalar there; the upper and lower elementary subgroups generate the local group over $\mathbb Z/a^e$. Thus $W_a$ preserves each coordinate line, and Lemma~\ref{paper:family:lw:weyl-line} gives the ordinary assertion. For a positive even fiber at odd $a$ the action can first be restricted to $A$. If $W_a$ preserves the symmetrized line of $Y,gY$, it also preserves both full lines: it commutes with the dyadic upper unipotent, whose eigenvalues on those two summands are opposite. Apply the ordinary assertion in the extension and divide the full codegree by four. The invariance required here holds for any fiber of a complementary local twist operator, intersected, at odd $a$, with a fermion sector.
\end{proof}

\begin{lemma}[Tame orbits of a short positive word]\label{paper:lem:short-tame-orbit}
Let $s$ be an odd prime and let $z$ be an unsigned word of $g<s$ roots of $s$-power order, of strict house below $g$. Its tame inertia orbit has size at least $(s-1)/\max\{h:h\mid s-1,\ h\le g\}$. The same lower bound holds for any character label whose field contains $z$.
\end{lemma}
\begin{proof}
The coefficient vector at a common $s$-power level is unique among positive words of length $g$: any difference is constant on every $\mu_s$-coset, and a nonzero constant would require $s$ letters on one side. Strictness gives a nonidentity root. Its orbit under the stabilizer has size equal to the stabilizer order, since every nonidentity tame unit is nontrivial modulo $s$. This orbit fits inside the $g$ letters.
\end{proof}

\subsection{Prime avoidance}\label{paper:sec:principal-family}

\begin{proposition}[Prime avoidance on a principal block]\label{paper:family:4.2}
Let $P$ be a full principal $p$-block of ordinary integral modular data,
or of positive reduced data at an odd prime. Let $E=P$; in ordinary
data one may also intersect $P$ with the fiber of trivial
prime-to-$p$ twist. Assume that the unit is the only invertible in
$E$. Fix a label $X$ and a prime $s\ne p$. If $X$ is $p$-active,
or if $p$ and $v_p(\Delta)$ are odd and $s\nmid p-1$, then at each
prime over $s$ there is a nonunit $Y\in E$ for which
$d_Y\chi_X(Y)$ has nonzero residue. In particular $s\nmid d_Y$.
\end{proposition}
\begin{proof}
The block projector commutes with the local group
(Theorem~\ref{paper:lw:block}); so does the indicated
complementary projector. Conjugation by the global Weyl normalizes
the local group. For an active $X$ the $X,0$ entry of the conjugate
projector is zero, whence $\sum_{Y\in E}d_Y\chi_X(Y)=0$.
For an inactive $X$ in the second case the slice identity
(Proposition~\ref{paper:3.8}) is
\[
 \sum_{\substack{Y\in E\\\ord(\theta_{Y,p})=p}}d_Y\chi_X(Y)
 =(p-1)\sum_{\substack{Y\in E\\\theta_{Y,p}=1}}d_Y\chi_X(Y).
\]
In either identity the vacuum contribution is a unit at $s$, so a
nonunit term must survive modulo each prime over $s$.
\end{proof}

In either setting of the proposition, at valuation two or three the
nonunit degrees in $P_p$ of a perfect table are $p$-free and have
gcd one: apply prime avoidance to an active column at every other prime.

For larger defects an index inequality bounds the total ramification of a block.


\section{Cyclotomic words and phase control}\label{paper:sec:words}\label{paper:family:words}

\label{paper:sec:arithmetic}

Throughout this section, a positive word means a finite sum
$\sum c_\xi\xi$ with $c_\xi\in\mathbb Q_{\ge0}$ and roots of
unity $\xi$. Its mass is $\sum c_\xi$. An integer word has integer
coefficients. A signed integer word is written with nonnegative
integer positive and negative multiplicities; their sums are denoted
$P,N$. Residues at a rational prime always mean residues at
\emph{every} place over that prime. All roots can be put at one finite
common level. The house of an algebraic number is the maximum of the
absolute values of its conjugates.

\subsection{Integral optimization with two retained primes}
The following cyclotomic unimodularity is established in \cite{BH}; we include the configuration argument needed here.


\begin{lemma}[Coset relations and integral optimization]
\label{paper:family:lem:two-prime-lattice}
At level $s^a$, the relations among the indicated roots are exactly
the constant coefficient vectors on the cosets of $\mu_s$.
Omitting one root from each coset gives an integral basis.
These assertions hold over the integer ring of any coprime
cyclotomic field, and after reduction at a different prime.

For roots supported on at most two odd primes, with an arbitrary
dyadic factor also permitted, the configuration of
their columns in an integral basis is unimodular. Consequently an
integral cyclotomic number with a positive rational word of mass
at most $g$ has a positive integer word of mass at most $g$.
The same assertion holds for minimizing any positive linear cost
on signed columns. Every rational optimum is a convex combination
of integer optima.
\end{lemma}
\begin{proof}
The first statements follow from
$\Phi_{s^a}(x)=\sum_{j=0}^{s-1}x^{js^{a-1}}$ and the tensor
product description of coprime cyclotomic integer rings. The
relations have disjoint supports and their number is the difference
between the number of roots and the degree of the field.

Here is a direct two-prime argument. At square-free level $uv$,
index the $uv$ roots by the edges of $K_{u,v}$. Their relation
lattice is generated by the sums on the stars of its vertices,
with one redundancy. Giving the two vertex classes opposite signs
identifies this relation matrix with an oriented incidence matrix.
Such a matrix is totally unimodular: expansion at a column with
at most one nonzero entry reduces any square minor, while if
every column has two nonzero entries the row sum is zero.
Delete a row and choose a spanning tree. Its incidence submatrix
has determinant $\pm1$. Eliminating those tree columns gives a
matrix for the quotient lattice whose maximal minors are,
up to sign, complementary maximal incidence minors, hence are
$0,1,-1$. This is the desired configuration. At prime-power
levels, partitioning both odd coordinates into prime-order cosets
makes a direct sum of copies of this configuration. The
one-prime case is the same argument for a simplex.
For a dyadic factor, the roots of order dividing $2^a$
are signed copies of an integral basis obtained by taking
one root from each $\{\xi,-\xi\}$ pair. Tensoring
therefore gives a direct sum of signed copies of the
odd configuration, which is again unimodular.

Adjoining negative columns preserves unimodularity. Every vertex
of $Ax=y$, $x\ge0$, for an integral $y$ is therefore integral:
solve for its independent support columns and use their unit
maximal minor. Positive costs make a nonempty optimal face
bounded. An optimum exists, and its vertices are integral.
\end{proof}

\begin{lemma}[Contraction without loss of actual phases]
\label{paper:family:lem:contraction}
Retain one or two primes in a cyclotomic level $M$, and put all
other primes in a coprime level $N$. Prescribe rational integer
residues at primes dividing $M$. Suppose every signed integer
word in $\mu_M$ satisfying these residues has cost at least $g$,
where positive and negative letters cost $1,c$, respectively.
If an integral number with the same residues has a positive
rational word of mass at most $g$, then:
\begin{enumerate}
\item for $c=1$, its complementary components are signs;
\item for $1\le c<\ell-1$, where $\ell$ is the least prime
dividing $N$, its complementary components are all $1$;
\item after the permitted signs are absorbed, its actual
coefficient vector is a convex combination of integral optima.
\end{enumerate}
\end{lemma}
\begin{proof}
Write the integer ring as
$\mathbb Z[\mu_M]\otimes\mathbb Z[\mu_N]$.
An integral functional on the second factor fixing $1$ preserves
the stipulated residues: the radical of a retained prime is
generated by $1-\zeta_{s^{v_s(M)}}$, and contraction preserves
this ideal. By Lemma~\ref{paper:family:lem:two-prime-lattice}, the cost
of every contracted rational expression is at least $g$.

For $c=1$, omit one root in each prime-order coset of each
complementary factor, keeping $1$ and a selected nonsign
component. The coefficient-of-$1$ functional in the product
basis takes values $0,\pm1$ on roots and kills the selected
component. Its positive occurrence would strictly lower mass.

For the second assertion, in a complementary $u$-power factor
choose a functional with value $1$ at $1$, value $-1$ at one
uniformly chosen nonidentity $u$th root, and zero elsewhere.
Its sums on prime-order cosets are zero. Tensor these independent
functionals. A term with a nontrivial complementary component
has expected signed cost at most $c/(u-1)<1$ for one such prime
$u$, whereas every term costs at most one in expectation.
Again a positive coefficient would make some contraction cost
less than $g$. The remaining rational word lies on a bounded
optimal face of the retained configuration, so the final claim
follows from Lemma~\ref{paper:family:lem:two-prime-lattice}.
\end{proof}

\begin{corollary}[One positive residue above twice the mass]
\label{paper:cor:one-large-residue}
Let $z$ be an integral cyclotomic number with a positive
rational word of mass at most $g\in\mathbb Z_{>0}$,
allowing even-order phases. If its residue is $g$ at
every place over a prime $r>2g$, then every such word
has only $r$-primary phases, with nonnegative integer
combined coefficients of sum $g$. This unsigned word
is unique at every common $r$-power level.
\end{corollary}
\begin{proof}
An integral signed $r$-word of length at most $g$
has signed augmentation in $[-g,g]$, congruent to
$g$ modulo $r$. Since $r>2g$, the augmentation is
$g$, forcing exactly $g$ positive letters.
Thus the minimum signed cost with the prescribed
residue is $g$. Lemma~\ref{paper:family:lem:contraction}
with both signs of cost one eliminates nonsign
spectators and expresses the actual signed vector
as a convex combination of these unsigned optima.
Their reduced coefficient vector is unique because
its length $g<r$ leaves a vacant position in each
prime-order coset. The actual coefficients equal
this unique integral vector.
\end{proof}

\begin{lemma}[Reciprocity preserves foreign twists at small gcd]
\label{paper:reciprocal-small-gcd}
Let $r$ be a prime, let $Y\in P_r$ have pure $r$-primary
twist and degree divisible by $r$, and let $V$ be a nonunit
constituent of $XY$. If
\[
                         2\gcd(d_V,d_Y)<r,
\]
then $V$ and $X$ have the same actual primary twist at every
prime other than $r$.
\end{lemma}
\begin{proof}
Put $g=\gcd(d_V,d_Y)$ and $L=\operatorname{lcm}(d_V,d_Y)$.
The inequality gives $r\nmid g$, and $r\mid d_Y$ therefore
gives $r\nmid d_V$. The integral entry
\[
                 z=\frac{S_{V,Y^*}}L
                   =\frac{\chi_{Y^*}(V)}{d_V/g}
\]
has an actual positive word of mass $g$ and residue $g$ at
every place above $r$. Here $Y^*\in P_r$, and the denominator
$d_V/g$ is an $r$-unit. Corollary~\ref{paper:cor:one-large-residue}
puts every actual phase of $VY^*$ in the $r$-primary roots.
Frobenius reciprocity gives $X\prec VY^*$. Its balancing phase
is $\theta_X/(\theta_V\theta_{Y^*})$, so its components
away from $r$ are one. Since $\theta_{Y^*}$ is $r$-primary,
the claimed equality of foreign twists follows.
\end{proof}


\begin{theorem}[Two positive residues below the larger prime]
\label{paper:family:thm:two-positive}
Let $q<r$ be primes, let $g<r$ be a positive integer, and let
$z$ be an integral cyclotomic number of positive rational mass
at most $g$, with residue $g$ at $q$ and $r$. Then
\[
 z=h+q\sum_{j=1}^m\xi_j,\qquad h+qm=g,
 \quad h,m\ge0,
\]
where the $\xi_j$ are nonidentity roots of $r$-power order,
unless $r=g+1$ and $q\mid g-1$. In this exceptional situation
the further possibility is
\[
 z=-\sum_{t=1}^{r-1}\eta_t\zeta_r^t,
\]
where the $q$-power roots $\eta_t$ form one copy of $1$ and
$(g-1)/q$ complete $\mu_q$-cosets, with repetitions allowed.
When the exception is unavailable, every actual positive word
of mass at most $g$ has precisely the unsigned primary
multiplicities displayed above. In particular $z=g$ if $g<q$,
and $|z|=g$ if $g=q$.
\end{theorem}
\begin{proof}
First use signed integer words on the $q,r$ alphabet of length
at most $g$. Erase the $r$-parts of the chosen word to obtain
$w$ in the $q$-primary field. This erasure is used only to
compute residues. That field is unramified at $r$, so
$(w-g)/r$ is integral. All conjugates of
$1+(w-g)/r$ lie in the disk of center $1-g/r$ and radius $g/r$.
The disk meets the unit circle only at $1$. The algebraic norm
therefore gives $w=g$ or $w=g-r$.

Erase the $q$-parts instead and denote the signed $r$-coefficients
by $c_i$. On each $\mu_r$-coset the coefficients
$c_i-g\delta_{i0}$ are constant modulo $q$. If every coset has
a vacant position after marking zero, these constants vanish.
The signed augmentation is then $g$ modulo $q$, excluding
$w=g-r$. Equality $w=g$ forces exactly $g$ positive letters
whose $q$-parts are $1$, so each off-identity multiplicity is
divisible by $q$.

Otherwise the length bound forces $g=r-1$ and exactly one letter
at every nonidentity position of $\mu_r$. The signs are all
congruent to $-g$ modulo $q$. A positive sign would force
$q\mid r$; at $q=2$ neither sign has the required parity.
Thus $q$ is odd, all signs are negative, and $q\mid g-1$.
Then $w=-1$ and $\sum_t\eta_t=1$. The prime-power coset
relations give the asserted one identity plus complete cycles.
The same argument shows that no integral signed word with the
prescribed residues has length below $g$.

Apply Lemma~\ref{paper:family:lem:contraction} with cost one. It removes
all spectator components except signs and places the original
word in the convex hull of integral optima. When the exception
is unavailable these optima are unsigned $r$-words of length
$g<r$. Their coefficient vector is unique: each coset has
minimum coefficient zero, and two representatives differ by
a constant on that coset. Thus the actual word has these
same multiplicities. The final assertions follow immediately.
\end{proof}

\begin{corollary}[Full tame orbit in the first positive range]
\label{paper:cor:first-positive-orbit}
Under Theorem~\ref{paper:family:thm:two-positive}, suppose
$q\le g<2q$ and the house of $z$ is strictly below $g$.
The tame $r$-inertia orbit of $z$ has size $r-1$,
including in the exceptional signed branch.
\end{corollary}
\begin{proof}
In the unsigned branch strictness forces $m=1$.
Uniqueness of its coefficient vector then makes every
stabilizing tame multiplier fix its single nonidentity
root, hence be the identity.
In the exceptional branch the nonidentity $r$th roots
are an integral basis over the $q$-primary integer ring.
A stabilizer of order $a$ therefore makes the coefficients
$\eta_t$ constant on its free orbits. Their sum is $1$,
so $1=a\sum_{\text{orbits}}\eta_t$. The orbit sum is
an algebraic integer, forcing $a=1$.
\end{proof}

\begin{proposition}[Dyadic rational extension]
\label{paper:prop:dyadic-three}
Let $2<q<r$ be odd primes. If an integral cyclotomic number
of positive rational mass at most four has residue four
at $2,q,r$, it equals four.
\end{proposition}
\begin{proof}
Theorem~\ref{paper:family:thm:two-positive} applies since $r>4$.
In its ordinary branch, either the value is four, or
$q=3$ and it is $1+3\xi$ with $\xi$ a nonidentity
$r$-primary root. Its residue at two is $1-\xi$, which
is a unit there because its norm is a power of $r$.
In the exceptional branch the nonidentity $r$th roots
form an integral basis over the $q$-primary integer ring,
and their coefficients $-\eta_t$ are units at two.
The value is therefore not in twice that integer ring.
The field has odd conductor and is unramified at two,
so it cannot vanish at every two-place. Both nonconstant
possibilities are excluded.
\end{proof}

\begin{lemma}[One positive residue in an odd alphabet]
\label{paper:lem:one-residue}
Suppose every prime in a positive rational word is at least
$\ell\ge3$. Let its mass be $g\in\mathbb Z_{>0}$, its value
be integral, and its residue at a prime $s>g$ be $g$.
If $(s-g)(\ell-1)>g$, all its actual phases are $s$-primary,
and its combined coefficients are nonnegative integers of sum $g$.
\end{lemma}
\begin{proof}
Choose $\max(1,g/(s-g))<c<\ell-1$. An integral signed
$s$-word has augmentation $g$ modulo $s$. If this augmentation
is negative, its cost is at least $c(s-g)>g$; if nonnegative,
cost at most $g$ forces augmentation $g$ and exactly $g$
positive letters. Lemma~\ref{paper:family:lem:contraction} removes all
spectators, and uniqueness of the reduced unsigned word of
length below $s$ determines the actual coefficients.
\end{proof}

\begin{lemma}[Odd prime mass with a zero residue]
\label{paper:family:lem:prime-mass}
Let $q<r$ be odd primes. An integral number with a positive
odd-order word of mass $r$, strict house below $r$, and residues
$r,0$ at $q,r$, respectively, has the form
\[
 z=\eta\left(h+q\sum_{j=1}^m\xi_j\right),\qquad
 h+qm=r,\quad h,m\ge1,
\]
with $\eta$ $q$-primary and the $\xi_j$ nonidentity
$r$-primary roots. Every actual positive word has the same
phases and combined coefficients.
\end{lemma}
\begin{proof}
Give negative letters cost $3/2$. A signed integer word of cost
below $r$ has length at most $r-1$. Erasing its $r$-parts gives
an integral $q$-word divisible by $r$ with house below $r$,
hence zero. Erasing its $q$-parts gives coefficients whose
deviations from $r\delta_0$ are constant modulo $q$ on
$\mu_r$-cosets. A vacancy after marking zero would make the
augmentation both $r$ and zero modulo $q$. Therefore all
$r-1$ nonidentity positions of the identity coset occur, once
each. Their common sign must be negative, as in the proof of
Theorem~\ref{paper:family:thm:two-positive}, and their cost is
$3(r-1)/2>r$. This contradiction gives the required cost
lower bound. Contraction removes all spectators because every
spectator prime is odd.

An unsigned integral optimum on the retained alphabet has
length $r$. Erase its $r$-parts to obtain $w$. The positive
$q$-word $w$ cannot vanish: its length $r$ is not divisible
by $q$. Thus $w/r$ is a nonzero algebraic integer of house
at most one, hence a root of unity. Equality in the triangle
inequality makes every $q$-part the same root $\eta$.
The remaining $r$-word has no full coset, since that would
give value zero. Nor can all nonidentity positions of the
identity coset occur: the one remaining letter, whether at
zero, at an existing position or at a new coset, contradicts
the coefficient congruences modulo $q$. All coset constants
are therefore zero. This gives the formula, with $h>0$ since
$q\nmid r$ and $m>0$ by strictness.

For two optima, the ratio of their $q$-parts belongs both to a
$q$-primary and an $r$-primary field, so it is rational and
equals $1$. The remaining unsigned coefficient vectors are
the unique reduced vectors in their cosets. Convexity of the
optimal face now proves the assertion for the actual word.
\end{proof}

\begin{proposition}[Three residues at prime mass]
\label{paper:three-residue-prime-mass}
Let $p<q<r$ be odd primes. Suppose an integral cyclotomic number $z$
has an actual positive odd-order word of mass $r$, strict house below
$r$, and residue $r$ at every place over each of $p,q,r$.
Then its actual word is
\[
 z=h+pq\sum_{\xi\ne1}m_\xi\xi,\qquad
 h+pq\sum_{\xi\ne1}m_\xi=r,
\]
where the $\xi$ are $r$-primary roots, $h\ge1$, and the $m_\xi$
are nonnegative integers with positive sum.
\end{proposition}
\begin{proof}
Lemma~\ref{paper:family:lem:prime-mass}, applied at $q,r$, gives the actual word
\[
 z=\eta\left(h+q\sum_{\xi\ne1}n_\xi\xi\right),\qquad
 h+q\sum n_\xi=r,
\]
where $\eta$ is $q$-primary, $h\ge1$, and the $n_\xi$ are nonnegative
integers with positive sum. The number of distinct marked $r$-roots,
including $1$, is at most $1+(r-1)/q<r$. At a common $r$-power level,
omit an unmarked root from each $\mu_r$-coset. The remaining roots form
an integral basis over the ring of integers of the $q$-primary field.
This basis includes every marked root and $1$.

The field on the $q,r$ alphabet is unramified at $p$. The all-places
$p$-residue therefore says $z-r\in p\mathcal O$ in this field. Expand
in the indicated integral basis. Each nonidentity marked coefficient
$q\eta n_\xi$ lies in $p\mathcal O_q$, so $p\mid n_\xi$. Consequently
$h\equiv r\pmod p$. The coefficient at $1$ is $h\eta-r$, and is also
in $p\mathcal O_q$; since $p\nmid r$, this gives
$\eta-1\in p\mathcal O_q$. A nonidentity $q$-primary root has
$\eta-1$ a unit at every prime above $p\ne q$. Thus $\eta=1$.
Writing $n_\xi=pm_\xi$ proves the assertion.
\end{proof}


\subsection{Exposed supports and descent}

\begin{lemma}[Exposure with the actual least prime]
\label{paper:family:lem:exposure}
Fix a primary prime $s$ and a nontrivial complementary
cyclotomic alphabet of coprime conductor whose least prime
is $\ell$, allowing $\ell=2$. A set $\mathscr{D}$ of pure
$s$-primary roots is exposed in the full alphabet if and
only if each $\mu_s$-coset $C$ has
\[
              \ell|\mathscr{D}\cap C|<(\ell-1)s.
\]
Here exposed means that a rational linear functional has value
$1$ precisely on the marked roots and value less than $1$
on every other allowed root. In the odd family this applies
with $(s,\ell)=(r,p)$ and with $(s,\ell)=(q,p)$.
\end{lemma}
\begin{proof}
For $a=|\mathscr{D}\cap C|$, give the $s$-primary functional value $1$
on marked positions and $-a/(s-a)$ on unmarked positions.
Its coset sum is zero. Extend it by normalized trace on the
complementary field. A complementary root has normalized trace
at least $-1/(\ell-1)$, and every nonidentity complementary root
has trace strictly below $1$. Thus the strict condition is
$a/((s-a)(\ell-1))<1$, as claimed.
Conversely average any exposing functional over the complementary
Galois group. At every unmarked position its primary value must
exceed $-(\ell-1)$, since a primitive $\ell$th root has trace
$-1/(\ell-1)$. The zero coset sum then gives the same inequality.
\end{proof}

\begin{lemma}[Common nonunit constituents]
\label{paper:family:lem:common}
Fix a nonunit row $Y$ of a perfect fusion-simple integral
modular category, and let $U,V$ be nonunit simple objects.
If the full positive character words of
$UU^*$ and $VV^*$ are supported on exposed sets $\mathscr{D},F$, their
adjoints have a common nonunit constituent whose actual phases
at $Y$ lie in $\mathscr{D}\cap F$.
\end{lemma}
\begin{proof}
Rigidity and commutativity give
$\langle UU^*,VV^*\rangle=\langle UV,UV\rangle$.
If the unit were their only common constituent, it has
multiplicity one in both adjoints, so $UV$ would be simple.
Its double braiding would then be scalar, contradicting the
strict character bound at two nonunit labels. Apply both
exposing functionals to the same balancing expression of a
common nonunit. Positivity and equality of total masses put
each of its phases in both exposed sets.
\end{proof}

\begin{lemma}[An exposed three-by-five difference set]
\label{paper:wild-difference-exposure}
Let $a\ge2$, let $\lambda$ have order $3^a$, and let $\omega$
and $\xi$ have orders three and five. Let $t_0,t_0',t_1',t_2'$
be the four distinct nonzero elements of $\mathbb F_5$. Put
\[
 S=\{\xi^{t_0},\lambda\xi^{t_0'},
             \lambda\omega\xi^{t_1'},
             \lambda\omega^2\xi^{t_2'}\},\qquad \mathscr{D}=S/S.
\]
Then $\mathscr{D}$ is exposed in every finite root alphabet containing it.
Moreover, if $\sigma$ fixes the
$3$-primary roots and inverts the $5$-primary roots, then
\[
                         \mathscr{D}\cap\sigma(\mathscr{D})=\{1\}.
\]
\end{lemma}
\begin{proof}
Put $F=\mathbb Q(\omega,\xi)$ and $K=\mathbb Q(\lambda,\xi)$.
The $3^{a-1}$ subspaces $\lambda^cF$, indexed by
$c\pmod{3^{a-1}}$, form a direct sum equal to $K$.
The difference set meets only the three distinct subspaces
$F$, $\lambda F$, and $\lambda^{-1}F$.

First consider $\mathscr{D}_0=\mathscr{D}\cap F$. The three cyclic differences
$t'_{i+1}-t'_i$, with indices modulo three, are nonzero and sum to
zero in $\mathbb F_5$. They cannot be three distinct nonzero
residues: the sum of all four nonzero residues is zero, so the sum
of any three is nonzero. They therefore have multiset
$\{d,d,-2d\}$ for a nonzero $d$. Consequently
\[
 \mathscr{D}_0=\{1,\omega\xi^d,\omega\xi^{-2d},
                 \omega^2\xi^{-d},\omega^2\xi^{2d}\}.
\]
After multiplying the column labels by $d^{-1}$, prescribe the
values of a rational linear functional on $\omega^i\xi^t$ by
\[
 \begin{array}{c|rrrrr}
       &0&1&2&3&4\\ \hline
 i=0&1&-1/4&-1/4&-1/4&-1/4\\
 i=1&-1/2&1&-3/4&1&-3/4\\
 i=2&-1/2&-3/4&1&-3/4&1
 \end{array} .
\]
Every row sum and every column sum is zero. These are exactly the
linear relations between the fifteen roots $\omega^i\xi^t$:
the relations have rank $3+5-1=7$, and $[F:\mathbb Q]=8$.
Thus the table defines a rational linear functional on $F$.
Its value is one precisely on $\mathscr{D}_0$ and is at most $3/4$ on
every other signed root of $F$.

The part $\mathscr{D}_+=\mathscr{D}\cap\lambda F$ is
\[
 \mathscr{D}_+=\{\lambda\omega^i\xi^{c_i}:0\le i\le2\},
 \qquad c_i=t_i'-t_0.
\]
The three $c_i$ are distinct and nonzero. Prescribe value one
at these three marked positions, value $-1/2$ at the other two
positions in each of their three columns, and zero in both
remaining columns. Each row and column again sums to zero.
This gives a rational linear functional on $\lambda F$, equal
to one precisely on $\mathscr{D}_+$ among its signed roots and at most
$1/2$ on every unmarked signed root. The same construction,
using the inverse marks, applies to $\mathscr{D}_-=\mathscr{D}_+^{-1}$ in
$\lambda^{-1}F$.

Combine these three functionals using the direct sum and put
zero on all remaining subspaces. The resulting functional
$\ell:K\longrightarrow\mathbb Q$ has value one precisely
on $\mathscr{D}$ and at most $3/4$ on every other root in
$\pm\mu_{3^a5}$; in addition, its absolute value at every
such root is at most one. To pass to an arbitrary larger
alphabet, compose $\ell$ with normalized relative trace to
$K$. Increasing a primary level at $2,3,$ or $5$ either
fixes the root or sends it to zero: a strictly higher
$3$- or $5$-primary level gives a vanishing geometric sum,
and a root of $2$-power order at least four has rational
trace zero. The complementary field has conductor coprime
to $30$. A nonidentity root in it has normalized rational
trace zero or a product of factors $-1/(p-1)$ with
$p\ge7$, whose absolute value is at most $1/6$.
Thus every newly allowed root still has value strictly
below one. This proves exposure in the full alphabet.

Finally $\sigma$ preserves the three primary subspaces.
In the neutral subspace the columns $\{d,-2d\}$ become
$\{-d,2d\}$, which are disjoint in $\mathbb F_5^\times$;
the same holds in the inverse row. Thus the only intersection
in $\mathscr{D}_0$ is the identity. In either cross subspace each row
has just one marked nonzero column $c$, which is different
from $-c$. There is no intersection there. This proves the
last assertion.
\end{proof}

\begin{theorem}[A forbidden four-position character support]
\label{paper:wild-support-forbidden}
In a perfect fusion-simple integral modular category, a
nonunit character at a nonunit row has no full positive
word supported inside $\alpha S$, where $S$ is the set
of Lemma~\ref{paper:wild-difference-exposure} and $\alpha$
is an arbitrary root of unity. The positive coefficients
need not be equal.
\end{theorem}
\begin{proof}
Suppose such a word exists for $X$ at the fixed nonunit
row $Y$.
Multiplying this word by its conjugate gives a full positive
word for $XX^*$ supported on $\mathscr{D}=S/S$; the common scalar cancels.

Choose a cyclotomic Galois automorphism $\sigma$ that fixes
the $3$-primary roots and inverts the $5$-primary roots.
Unsigned Galois covariance and the invariance of integral
dimensions give, at the same row $Y$, a full positive word
for $X^\sigma(X^\sigma)^*$ supported on $\sigma(\mathscr{D})$.
Both supports are exposed by the preceding lemma.
Lemma~\ref{paper:family:lem:common} supplies a common nonunit
constituent whose full actual character word at $Y$ is
supported on $\mathscr{D}\cap\sigma(\mathscr{D})=\{1\}$. Its character value
therefore equals its dimension, contradicting strictness.
\end{proof}

\begin{corollary}[A forbidden three-position support at three and five]
\label{paper:three-five-three-positions}
Let $\omega$ and $\xi$ have orders three and five. In a perfect
fusion-simple integral modular category, no nonunit character at
a nonunit row has a full positive word supported inside
\[
 \alpha\{\xi^{t_0},\omega\xi^{t_1},\omega^2\xi^{t_2}\},
\]
where $\alpha$ is a root of unity and $t_0,t_1,t_2$ are distinct
modulo five.
\end{corollary}
\begin{proof}
The three nonzero cyclic differences of the $t_i$ have multiset
$\{d,d,-2d\}$ for some nonzero $d\pmod5$. Thus the adjoint
difference support is precisely the neutral set $\mathscr{D}_0$ in the
proof of Lemma~\ref{paper:wild-difference-exposure}. The displayed
three-by-five table exposes it in the full root alphabet, and
it meets its pure fifth inverse only at one. The common
constituent argument of Theorem~\ref{paper:wild-support-forbidden}
applies without change.
\end{proof}


\begin{lemma}[The two possible odd-primary orders for two positions]
\label{paper:two-position-orders}
In a perfect fusion-simple integral modular category, suppose a full
positive character word of a nonunit at a nonunit row has exactly two
positions, and their ratio has odd prime-power order. That order is
either three or five. The dimension of the category may be even.
\end{lemma}
\begin{proof}
Write the ratio as $\xi$, of order $s^a$, with $s$ an odd prime.
The full adjoint word has difference support
$\mathscr{D}=\{1,\xi,\xi^{-1}\}$. If $a\ge2$, these three roots lie
in distinct $\mu_s$-cosets, so Lemma~\ref{paper:family:lem:exposure}
exposes $\mathscr{D}$ in the full root alphabet, even with least
complementary prime two. If $a=1$ and $s\ge7$, the same lemma
applies because $2|\mathscr{D}|=6<s$.

In either situation choose a cyclotomic Galois unit whose
$s$-component is two and whose complementary components are one.
The conjugate simple has adjoint support
$\mathscr{D}^2=\{1,\xi^2,\xi^{-2}\}$, which is also exposed, and
$\mathscr{D}\cap \mathscr{D}^2=\{1\}$. Lemma~\ref{paper:family:lem:common} therefore gives
a common nonunit constituent with full character word supported
only at one. Its character value equals its dimension, contrary
to strictness. The only orders left are three and five.
\end{proof}


\begin{theorem}[Uniform difference descent]
\label{paper:family:thm:descent}
Suppose a nonunit character at a fixed nonunit row has, up
to a scalar root, a full positive $r$-word on at most $m+1$
positions. Put
\[
 E=\frac{m(m+1)}2,\qquad H=\frac{r-1}2,\qquad
 u(h)=\min\{k\ge1:k(k+1)\ge h\}.
\]
If $p(2E+1)<(p-1)r$, then some $1\le h\le E$ satisfies
\begin{equation}\label{paper:family:eq:remaining}
                 (H-1)u(h)\le h(h-1).
\end{equation}
\end{theorem}
\begin{proof}
Among all full primary words at this row choose a nonunit
$U$ minimizing the number $h$ of opposite pairs of nonzero
differences of its support. Strictness gives $h>0$ and the
initial word gives $h\le E$. Let $\mathscr{D}$ be this difference set.
It has $2h+1$ elements and is exposed by
Lemma~\ref{paper:family:lem:exposure}.

For every tame unit class $a$ modulo sign set
$k_a=|(\mathscr{D}\cap a\mathscr{D})\setminus\{0\}|/2$.
The Galois conjugate $U^a$ has difference set $a\mathscr{D}$ at the
same row. Lemma~\ref{paper:family:lem:common} gives a nonunit word
on $\mathscr{D}\cap a\mathscr{D}=\{0,\pm x_1,\ldots,\pm x_{k_a}\}$.
Its differences modulo sign lie among $x_i,2x_i,x_i\pm x_j$,
so number at most $k_a(k_a+1)$. Minimality gives
$k_a\ge u(h)$. Tame units modulo sign act freely on nonzero
opposite pairs, even at higher prime-power level. Each ordered
pair of the $h$ pairs therefore contributes to at most one
$k_a$, giving $\sum_a k_a\le h^2$. Since $k_1=h$, this is
exactly \eqref{paper:family:eq:remaining}.
\end{proof}

\begin{proposition}[Exposure is automatic at the family bound]
\label{paper:family:prop:auto-exposure}
For odd primes $p<q<r$ set
\[
 m=\left\lfloor\frac{\min(r,p^2)}q\right\rfloor,
 \qquad E=m(m+1)/2.
\]
If $m\ge1$, then $m\le p-2$ and
$p(2E+1)<(p-1)r$. Thus the only numerical condition left
by Theorem~\ref{paper:family:thm:descent} is
\eqref{paper:family:eq:remaining}; no extra exposure hypothesis is needed.
\end{proposition}
\begin{proof}
Since $q\ge p+2$, $(p-1)q>p^2$, giving $m\le p-2$.
Also $r>mq$. For $1\le m\le p-2$,
\[
 \frac{p-1}{p}m(p+2)-\bigl(m(m+1)+1\bigr)
       =m\left(p-m-\frac2p\right)-1>0.
\]
Together these inequalities give the assertion.
\end{proof}

\begin{lemma}[Differences across primary cosets]
\label{paper:family:coset-differences}
Let $s$ be an odd prime, $a\ge1$, and $S\subseteq\mathbb Z/s^a$ a nonempty
set, and $K$ the subgroup of order $s$. Put $\mathscr{D}=S-S$ and
\[
 h=(|\mathscr{D}|-1)/2,\qquad x=(|\mathscr{D}\cap K|-1)/2,\qquad y=h-x.
\]
If $S$ is not contained in one coset of $K$, then
\begin{equation}\label{paper:family:eq:low-high}
                x\le\binom y2,\qquad y(y+1)\ge2h.
\end{equation}
The assertion is unchanged by translating $S$.
\end{lemma}
\begin{proof}
For a nonzero $u-v\in \mathscr{D}\cap K$, choose $w\in S$ outside the
common $K$-coset of $u,v$. The two higher-level differences
$u-w,v-w$ lie in the same nonzero $K$-coset and have difference
$u-v$. The quotient has odd order, so every nonzero coset is
distinct from its negative. Choose one coset from each opposite
pair meeting $\mathscr{D}$, and let its population be $n_i$. Since $\mathscr{D}=-\mathscr{D}$,
we have $\sum_i n_i=y$. The chosen coset and its negative give
the same opposite differences in $K$, at most $\binom{n_i}{2}$.
Every one of the $x$ low pairs occurs in this way. Therefore
$x\le\sum_i\binom{n_i}{2}\le\binom y2$, and $h=x+y$ gives
the second inequality.
\end{proof}

\begin{theorem}[Intrinsic reduction to prime level]
\label{paper:family:intrinsic-conductor}
Assume the exposure hypothesis $p(2E+1)<(p-1)r$ of
Theorem~\ref{paper:family:thm:descent}, and use its minimizing word. Put
$k=u(h)$, $H=(r-1)/2$ and
$v=v(h)=\min\{j\ge1:j(j+1)\ge2h\}$. Every minimizing
difference mask is at prime level if
\begin{equation}\label{paper:family:eq:conductor-test}
 Hk>\max\left\{\frac{h^2}{r},\ (h-v)^2+\frac{v^2}{r}\right\}.
\end{equation}
At the exact-family bound of
Proposition~\ref{paper:family:prop:auto-exposure}, this condition is
equivalent to
\begin{equation}\label{paper:family:eq:conductor-slack}
 \delta<B(h),\qquad
 \delta=h(h-1)-(H-1)k,\qquad
 B(h)=v(2h-v)-(h-k).
\end{equation}
For $h\ge2$, one has $B(h)>2(h-1)$.
\end{theorem}
\begin{proof}
The exposed common-constituent argument in
Theorem~\ref{paper:family:thm:descent} applies to every primary-unit
multiplier, including wild units. Thus any minimizing mask
meets each such dilate in at least $k$ pairs. If $h_j$ of its
pairs have exact order $r^j$, averaging over the full unit group
at a common conductor gives
\[
 k\le\sum_{j\ge1}\frac{h_j^2}{H r^{j-1}}.
\]
If the support is not in one $\mu_r$-coset, put
$x=h_1$, $y=h-x$. Lemma~\ref{paper:family:coset-differences} gives
$y\ge v$, whence
\[
 Hk\le x^2+\frac{y^2}{r}
       \le\max_{v\le t\le h}\left((h-t)^2+\frac{t^2}{r}\right)
       =\max\left\{(h-v)^2+\frac{v^2}{r},\frac{h^2}{r}\right\}.
\]
This proves the first assertion by contradiction. A support in
one $\mu_r$-coset has prime-level differences, independently
of its common scalar phase.

At the family bound, $v\le m$ and
$r>mq>m(m+1)=2E\ge2h$, while $r>m^2\ge v^2$.
Consequently $Hk\ge h>h^2/r$ and $0<v^2/r<1$.
Since $Hk$ is integral, the other strict inequality in
\eqref{paper:family:eq:conductor-test} is exactly
$Hk\ge(h-v)^2+1$. Substitution of
$Hk=h^2-h+k-\delta$ gives \eqref{paper:family:eq:conductor-slack}.
Finally, $h\ge2$ gives $v\ge2$ and
\[
 B(h)-2(h-1)
  =(v-2)(2h-v-2)+h+k-2>0.
\]
\end{proof}


\begin{lemma}[The exact scalar frontier]
\label{paper:family:scalar-frontier}
For $E\ge1$, put $K=u(E)$. Then
\begin{equation}\label{paper:family:eq:scalar-frontier}
 C(E):=1+\max_{1\le h\le E}\frac{h(h-1)}{u(h)}
 =1+\max\left\{\frac{E(E-1)}K,\ K^3-K^2-K\right\}.
\end{equation}
Thus the scalar descent excludes the seed precisely when
\[
              (H-1)K>E(E-1),\qquad H-1>K^3-K^2-K.
\]
The second comparison is automatic when $K=1$.
\end{lemma}
\begin{proof}
On each interval $u(h)=j$, the ratio $h(h-1)/j$ increases.
The completed-interval maxima are
$(j+1)(j(j+1)-1)=j^3+2j^2-1$, which increase with $j$.
Only the final point $E$ and, if $K\ge2$, the preceding
endpoint $(K-1)K$ need be checked. The latter gives
$K^3-K^2-K$. When $K=1$ this expression is $-1$ and does
not affect the maximum. The two comparisons are exactly
$H>C(E)$.
\end{proof}

This is an exact numerical form of the scalar support criterion.

\subsubsection{Closure and short-support exclusions}

\begin{lemma}[Closure of exposed difference envelopes]
\label{paper:family:forced-envelope}
Fix a nonunit row in a perfect fusion-simple integral modular category,
and an odd prime $s$. A full primary word means a positive word of
total mass $d_X$ for the character of a nonunit $X$, supported, up
to a common scalar root, on $s$-power roots. Put
\[
 u(h)=\min\{j\ge1:j(j+1)\ge h\},\qquad H=(s-1)/2.
\]
Use $\lambda=2$ in arbitrary dimension, or $\lambda=3$ in odd
dimension. Suppose a full primary word has at most $k$ positions,
put $E=\binom{k}{2}$, and assume
$\lambda(1+2E)<(\lambda-1)s$. Let $L$ be the least integer
$1\le h\le E$ for which
\[
                    (H-1)u(h)\le h(h-1).
\]
If no such $L$ exists, the word is impossible. Otherwise every
full primary word at this row has at least $L$ nonzero opposite
difference pairs. With $b=\binom{k-1}{2}$, the condition
\begin{equation}\label{paper:family:one-coset-test}
       \max\{E^2,\,s b^2+(E-b)^2\}<sH u(L)
\end{equation}
forces the initial positions into one coset of $\mu_s$.

At prime level write an \emph{envelope} as a symmetric set
$\mathscr{D}\subseteq\mathbb F_s$ containing zero and containing the
adjoint difference support of some nonunit at the fixed row.
If $\mathscr{D},F$ are exposed envelopes, meaning
$\lambda|\mathscr{D}|,\lambda|F|<(\lambda-1)s$, then for every
$a\in\mathbb F_s^\times$ the set
\[
                   (\mathscr{D}\cap aF)-(\mathscr{D}\cap aF)
\]
is another envelope. Consequently an initial support is excluded
if this closure ever produces fewer than $L$ nonzero opposite
pairs, or an intersection with fewer than $u(L)$ pairs.
One may also use a stronger lower bound proved for the minimum
among prime-level full words. A work limit or an unprocessed
consequence is inconclusive.
\end{lemma}
\begin{proof}
Theorem~\ref{paper:family:thm:descent}, using
Lemma~\ref{paper:family:lem:exposure} with the conservative least-prime
bound $\lambda$, proves the scalar assertion and the lower bound
$L$. If the initial positions meet more than one $\mu_s$-coset,
at most $b$ opposite differences have order $s$. If their number
is $x$, averaging intersections over all units at the full
primary level gives an upper bound
\[
             \frac{x^2}{H}+\frac{(E-x)^2}{sH}.
\]
The order-$s$ orbit has $H$ pairs; every higher-order orbit has
at least $sH$ pairs. Every intersection must contain at least
$u(L)$ pairs, by the common-nonunit argument and the fact that
$\{0,\pm x_1,\ldots,\pm x_j\}$ has at most $j(j+1)$
opposite difference pairs. Convexity on $0\le x\le b$ now
contradicts \eqref{paper:family:one-coset-test}.

For the closure assertion apply Lemma~\ref{paper:family:lem:common} to
the two underlying nonunits, translating the second by the
Galois multiplier $a$. Their adjoints have a common nonunit
whose full positive phases lie in $\mathscr{D}\cap aF$. Its adjoint
differences lie in the displayed set. This remains valid when
an envelope strictly contains the actual support, so supports
may be padded and envelopes may be stored up to dilation.
After the one-coset reduction all these operations remain at
prime level. The two rejection rules follow from the lower
bound $L$ and the bound $j(j+1)$ just used. Omitting a valid
consequence cannot create an exclusion.
\end{proof}

\begin{theorem}[Uniform short-support exclusions]
\label{paper:family:short-support}
In an ordinary perfect fusion-simple integral modular category,
a full positive $s$-primary word of a nonunit at a nonunit row
cannot have at most $k$ positions in the following ranges.
A common scalar root and arbitrary primary conductor are allowed.
\[
\begin{array}{c|c|c|c|l}
 k&\text{Parity}&s\ge T_k&
       \text{Scalar alone}&\text{Checks}\\ \hline
 2&\text{odd}&5&s>3&\text{none}\\
 3&\text{odd}&11&s>9&\text{none}\\
 4&\text{odd}&29&s>33&\mathrm{D20,D21}\\
 5&\text{odd}&47&s>63&\mathrm{D20}\\
 6&\text{none}&83&s>108&\mathrm{D22}\\
 7&\text{none}&127&s>193&\mathrm{D24,D25}
\end{array}
\]
The two-position conclusion also holds in arbitrary dimension
for $s\ge7$.

These support bounds are from \cite{BDPfamily}.
The finite checks concern exponent sets of roots of unity.
Their exact input sets and rejection invariants are given in the proof.
\end{theorem}
\begin{proof}
For $E=\binom{k}{2}$, the maxima of
$1+\max_{1\le h\le E}h(h-1)/u(h)$ give the scalar thresholds
in the table. The exposure bounds hold throughout the stated
ranges, using $\lambda=3$ for $k\le5$ and $\lambda=2$ for
$k=6,7$. The arbitrary-parity two-position case uses
$\lambda=2$ and $6<s$.

For $k=4,5,6$, the only primes not settled by the scalar bound
are respectively
\[
 \{29,31\},\qquad \{47,53,59,61\},\qquad
 \{83,89,97,101,103,107\}.
\]
Condition~\eqref{paper:family:one-coset-test} holds at every one of
these inputs. Translate and dilate the support to contain
$0,1$, pad it to $k$ points, and apply the envelope closure.
The exact checks D20--D22 have no surviving support at these
inputs. Each examines all $\binom{s-2}{k-2}$ normalized
supports; stopping a closure search counts as survival.

For seven positions, the same one-coset test holds at every
prime from $127$ to $193$. We may therefore minimize the
number $h$ of differences among prime-level full words, with
$h\le21$. At $s=127$, the scalar lower bound is $17$.
If $h=17$ or $18$, every nonidentity dilation has overlap
at least four, whereas the mean is at most $306/62<5$.
Thus some intersection is $\{0,\pm x_1,\ldots,\pm x_4\}$.
The four-pair part of D25 rejects all such starting envelopes
at lower bound $17$, raising the prime-level minimum to $19$.
For the primes from $127$ to $167$, D25 then checks the
seven-point supports, with respective lower bounds
\[
\begin{array}{c|rrrrrrrrr}
 s&127&131&137&139&149&151&157&163&167\\ \hline
 L&19&17&17&17&18&18&19&19&19.
\end{array}
\]
All are excluded. The affine normalization chooses a triple
with least normalized third point; its prefix pruning retains
a representative of every support. The additional prefix test
uses only the maximum number of differences that the unchosen
edges could still add. Its work cap again counts as survival.

At the primes from $173$ to $193$, the scalar possibilities
are $h=19,20$ at $173$, and only $h=20$ at the other primes.
The mean nonidentity overlap is at most $380/85<5$, so a
four-pair intersection again exists. Check D24 enumerates all
such four-pair sets after scaling one pair to $\pm1$; none
has a difference envelope of the required size whose every
dilation overlaps it in at least four pairs. This excludes
the last possibilities.
\end{proof}

The finite support checks used in this proof have exact enumeration
certificates. Their algorithms preserve a representative of every
admissible support, and a work limit always retains the unresolved
support. Reproducible commands and integrity hashes accompany the
computational files.

\subsection{Actual phases and their energy}

For an integer $a$ and $c>0$, let $\lambda_c(a)=\lambda_{q,r;c}(a)$
be the minimum of $P+cN$ over signed integer $q,r$-words with residue
$a$ at both primes. Arbitrary primary levels are allowed. Weighted
two-prime optimization makes this a lower bound for the same cost
of every signed rational expression of an integral value with those
residues. The next proof reduces the minimum to prime-order level.

\begin{lemma}[Weighted residue distance for $r\ge3q$]
\label{paper:lem:large-ratio}
Let $5\le p<q<r$ be odd primes and suppose $r\ge3q$.
For $c=p-4/3$ one has
\[
                         \lambda_c(p^2)=p^2.
\]
\end{lemma}
\begin{proof}
Put $g=p^2$. First reduce primary levels without increasing either
sign's mass. In the integral decomposition
\[
 \mathbb Z[\zeta_{q^b}]
 =\bigoplus_{i=0}^{q^{b-1}-1}
       \zeta_{q^b}^{i}\mathbb Z[\zeta_q],
\]
send each displayed basis multiplier to $1$. This integer linear
map sends roots to $q$th roots, preserves reduction at $q$, and,
over the $r$-factor, preserves the radical ideal of $r$. Apply the
same contraction in the $r$-factor. It suffices to consider a
signed integer word on $\mu_q\mu_r$ of cost
$P+cN\le g$.  Since
\[
 cq\ge(p-4/3)(p+2)=g+(2p-8)/3>g,
\]
its negative mass satisfies $N<q$.  Erasing one coordinate
cannot increase either sign's mass.

We first exclude mixed signs in the off-identity coefficients
of either erased margin.  Suppose $n$ coefficients have the
same nonzero residue $a$ modulo $s$, with $1\le a<s$, and
$k$ of them are negative, where $1\le k<n$.  Their cost is at
least
\[
                        (n-k)a+ck(s-a).
\]
For $1<c<n-1$, minimizing this bilinear expression over the
integer rectangle gives the lower bound
\[
          \min\{n-1+c(s-1),\ s-1+c(n-1)\}.
\]
Indeed its minimum in $k$ is at $1$ or $n-1$; the two
remaining linear functions are minimized respectively at
$a=1$ and $a=s-1$.
For $(n,s)=(r-1,q)$ or $(q-1,r)$ this lower bound is at least
\[
 r-1+c(q-2)\ge3p+5+(p-4/3)p>g.
\]
Mixed signs are therefore impossible.

For the $r$-margin, a nonzero off-identity class cannot be
entirely negative, since it has $r-1>N$ positions.
A zero class has no negative coefficient either: a negative
multiple of $q$ would use at least $q$ negative letters.
Hence its off-identity coefficients are nonnegative, with
least class representative
\[
                         0\le v\le g/(r-1).
\]
For the $q$-margin the same argument leaves one additional
possibility: all $q-1$ off-identity coefficients are negative.
Because $N<q$, they must then all equal $-1$.
Thus its common representative can be chosen as
\[
                u=-1\quad\text{or}\quad0\le u\le g/(q-1).
\]
The zero class again has no negative coefficient, since its
modulus is $r>N$.

Let $A=P-N$ be the signed augmentation.  The two residue
conditions give
\[
 A\equiv g+qu\pmod r,\qquad
 A\equiv g+rv\pmod q.
\]
Therefore $A\equiv B:=g+qu+rv\pmod{qr}$.
Both numbers lie strictly between $-q$ and $qr-q$.
For $A$, use $N<q$ and $A\le P\le g<qr-q$.
For $B$, its lower bound is $g-q>-q$, while
\[
 B\le3g+\frac{g}{q-1}+\frac{g}{r-1}<3g+2q<qr-q.
\]
The last inequality follows from $g\le(q-2)^2$, $r\ge3q$,
and $3(q-2)^2+3q<3q^2$.
Consequently $A=B$ literally.

If $u\ge0$, then $A\ge g$ and
$P+cN=A+(c+1)N\ge g$.
If $u=-1$ and $v\ge1$, then $A=g-q+rv>g$, impossible.
The remaining case $u=-1,v=0$ has $A=g-q$ and $N\ge q-1$,
so
\[
 P+cN\ge g-q+(c+1)(q-1)=g+c(q-1)-1>g,
\]
also impossible.  Thus every such word has cost at least $g$,
and the constant word $g$ attains it.
\end{proof}

\begin{theorem}[Weighted concentration of the actual phases]
\label{paper:thm:mass-descent}\label{paper:prop:weighted}
Let $p<q<r$ be primes. Suppose an integral value $z$ has a positive
rational word of mass $M\le p^2$, with primary orders supported on
$p,q,r$, and residue $p^2$ at every place over $q,r$. Let $m_0$ be
the mass on phases whose $p$-primary component is $1$. For every
$0<c<p-1$,
\begin{equation}\label{paper:eq:weighted-concentration}
 \begin{split}
 (p-1-c)m_0+cM&\ge(p-1)\lambda_c(p^2),\\
 M-m_0&\le\frac{p-1}{p-1-c}\bigl(M-\lambda_c(p^2)\bigr).
 \end{split}
\end{equation}
In particular, if $p\ge5$ and $r\ge3q$, then
\[
                         M=m_0=p^2.
\]
Thus every phase of the actual word has trivial $p$-primary part;
no residue condition at $p$ is required.
\end{theorem}
\begin{proof}
Partition the $p$-primary coordinate into $\mu_p$-cosets. In the
identity coset let $\alpha_i$ be the $q,r$-word at position $i$,
of mass $m_i$, with position zero corresponding to $1$.
For $i\ne0$, omit position $i$ in the identity coset and one
arbitrary position in every other coset when choosing an integral
basis of the $p$-power factor. The coefficient of the identity basis
vector is $\alpha_0-\alpha_i$. Integrality of $z$ makes this
difference integral; comparison in the same basis after reduction
at $q$ or $r$ gives residue $p^2$. No individual slice is assumed
integral. Weighted two-prime optimization therefore gives
\[
                         m_0+c m_i\ge\lambda_c(p^2).
\]
Sum over the $p-1$ nonidentity positions and use
$\sum_{i\ne0}m_i\le M-m_0$. Rearrangement gives the two
inequalities. For $p\ge5$ and $r\ge3q$, take $c=p-4/3$ and
apply Lemma~\ref{paper:lem:large-ratio}. Since $M\le p^2$, the
second inequality forces $M=m_0=p^2$.
\end{proof}


\begin{proposition}[Sparse margins and uniqueness at the optimum]
\label{paper:sparse-optimum}
Let $5\le p<q<r$ be odd primes with $r\ge3q$, and put $g=p^2$.
An integral value with residue $g$ at every place over $q,r$ and a
positive rational $q,r$-word of mass at most $g$ has a unique such
coefficient vector. Its coefficients are nonnegative integers and
its mass is exactly $g$.

Write the vector as $(n_{\eta,\xi})$, with $\eta$ $q$-primary and
$\xi$ $r$-primary. Its margins satisfy
\[
 r\mid\sum_\xi n_{\eta,\xi}\quad(\eta\ne1),\qquad
 q\mid\sum_\eta n_{\eta,\xi}\quad(\xi\ne1).
\]
Consequently the numbers of different $q$- and $r$-primary components
in its support are at most
\[
 A=1+\lfloor g/r\rfloor,\qquad B=1+\lfloor g/q\rfloor,
\]
respectively. The conclusion holds at arbitrary primary levels.
\end{proposition}
\begin{proof}
Lemma~\ref{paper:lem:large-ratio} makes $g$ the minimum positive mass.
At one finite common level containing the given word, integral
optimization on the retained $q,r$-alphabet gives an integer optimum
of mass $g$, and expresses every rational optimum as a convex
combination of integer optima. We prove the margins and uniqueness
for those integer optima.

Erase the $r$-parts. On each $\mu_q$-coset $C$, the residue condition
at $r$ says that the coefficients, after subtracting $g$ at the
identity, have a common residue $u_C$ modulo $r$. Choose
$0\le u_C<r$. In the identity coset all $q-1$ nonidentity positions
have coefficient at least $u_C$; in every other coset all $q$
positions do. Thus, with $U=\sum_Cu_C$,
\[
 (q-1)U\le g.
\]
Summing all coefficients modulo $r$ gives
$g\equiv g+qU\pmod r$, so $r\mid U$. Since
$g<(q-1)r$, one has $U=0$. Every $u_C$ vanishes, proving the first
margin assertion, including all higher-level cosets. Interchanging
the primes proves the second assertion: the analogous nonnegative
sum is divisible by $q$ and at most $g/(r-1)<q$.
The support bounds follow by charging at least $r$, respectively
$q$, to each nonidentity marginal position.

Let two integer optima represent the same value at a common level.
The inequalities $g\le(q-2)^2$ and $r\ge3q$ imply
\[
 2A\le2+\frac{2(q-2)^2}{3q}<q,\qquad
 2B\le2+\frac{2(q-2)^2}{q}<r.
\]
Their difference decomposes over products of prime-order cosets.
In each product its matrix is a row-plus-column relation
$a_i+b_j$, by the two-prime relation lattice. The union of the two
supports has fewer than $q$ occupied rows and fewer than $r$ occupied
columns. There is therefore a common vacant row and a common vacant
column in that product. Vanishing on the row makes all $b_j$ equal;
vanishing on the column makes all $a_i$ the opposite constant. The
difference is zero. Hence the integer optimum is unique, and the
convex-combination statement proves integrality and uniqueness of
the original rational vector.
\end{proof}

\begin{corollary}[The actual phase levels are the field conductor]
\label{paper:phase-conductors}
Under Proposition~\ref{paper:sparse-optimum}, let $q^a,r^b$ be the
largest orders of the actual primary components, with exponent zero
when all components at that prime are trivial. For $K=\mathbb Q(z)$,
\[
                       \operatorname{cond}K=q^a r^b.
\]
If $a>0$, the tame $q$-inertia stabilizer of $z$ has order $d_q$
satisfying
\[
 d_q\mid q-1,\qquad d_q\le\lfloor p^2/r\rfloor,\qquad
 e_q(K/\mathbb Q)=\varphi(q^a)/d_q.
\]
If $b>0$, the corresponding assertions are
\[
 d_r\mid r-1,\qquad d_r\le\lfloor p^2/q\rfloor,\qquad
 e_r(K/\mathbb Q)=\varphi(r^b)/d_r.
\]
\end{corollary}
\begin{proof}
Put $F=\mathbb Q(\mu_{q^a r^b})$. Uniqueness of the actual
coefficient vector makes every element of $\operatorname{Gal}(F/K)$
preserve that vector, and hence its primary support.
If $a\ge2$, the last wild subgroup of $q$-inertia has order $q$
and acts freely on every primitive $q^a$ component. It cannot fix
$z$, since there are at most $\lfloor p^2/r\rfloor<q-1$
nonidentity $q$-components. Since $q$-inertia is cyclic, its
stabilizer has no nontrivial $q$-part. In particular the conductor
exponent at $q$ is exactly $a$. For $a=1$, an unramified field would
require invariance under all $q-1$ tame automorphisms, contradicting
the same support bound.

The remaining stabilizer is tame. Its action on every nonidentity
$q$-component is free, so its order divides $q-1$ and is at most
the number of those components. Its index in $q$-inertia is the
ramification index of $K$. This proves the assertions at $q$;
interchange the primes for those at $r$. No other prime ramifies
in $F$, proving the conductor formula.
\end{proof}

\begin{lemma}[Energy at the last primary level]
\label{paper:last-level-energy}
Let $\ell$ be an odd prime, $u\ge2$, $h\in\mathbb Q$ and
$t\in\mathbb Q_{>0}$. Suppose
\[
 z=h+t\sum_{\eta\in\mu_{\ell^u}\setminus\{1\}}k_\eta\eta,
 \qquad k_\eta\in\mathbb Z_{\ge0},\qquad
                    0<\sum_\eta k_\eta<\ell,
\]
has conductor $\ell^u$, its coefficients are invariant under
$H=\operatorname{Gal}(\mathbb Q(\zeta_{\ell^u})/\mathbb Q(z))$,
and some primitive $\ell^u$th root has nonzero coefficient.
Then
\begin{equation}\label{paper:last-level-lower}
                       E(z)\ge t^2\ell^{u-1}(\ell-1).
\end{equation}
\end{lemma}
\begin{proof}
The stabilizer $H$ is tame: a nontrivial subgroup of the cyclic
wild part contains the last order-$\ell$ subgroup, which would
lower the conductor. Put
$\Gamma=(\mathbb Z/\ell^u)^\times/H$, $e=|\Gamma|$, and
let $N\subset\Gamma$ be that last subgroup, of order $\ell$.
Write $m_C$ for the primitive coefficients on $H$-cosets and
$\kappa=\sum_Cm_C$, so $1\le\kappa<\ell$. For an $N$-coset
$C_0$ put $s_{C_0}=\sum_{C\in C_0}m_C$. Partial Parseval and
nonnegative integrality give
\[
 \begin{split}
 \sum_{\chi|_N\ne1}|\widehat m(\chi)|^2
 &=\frac e\ell\left(\ell\sum_Cm_C^2-\sum_{C_0}s_{C_0}^2\right)\\
 &\ge\frac e\ell\kappa(\ell-\kappa)
 \ge\frac e\ell(\ell-1).
 \end{split}
\]
The lower levels of $z$ vanish in these channels. Primitive
Gauss sums have squared modulus $\ell^u$, so Parseval for $z$
gives $eE(z)\ge t^2\ell^u\sum_{\chi|_N\ne1}
|\widehat m(\chi)|^2$, proving the assertion.
\end{proof}

\begin{proposition}[Energy of a sparse residue word]
\label{paper:phase-energy}
Under Proposition~\ref{paper:sparse-optimum}, put $g=p^2$ and
$K=\mathbb Q(z)$. If both primary supports are nontrivial, so that
$\operatorname{cond}K=q^a r^b$ with $a,b\ge1$, then
\begin{equation}\label{paper:mixed-energy}
 E(z)\ \ge\ q^a r^b
 \frac{(qr-r-g)(qr-q-g)}{(g+r)(g+q)}
 \ >\ \frac{qr}{g}\,q^a r^b.
\end{equation}
If just one primary support is nontrivial, write its prime as $s$,
the other prime as $t$, and its conductor as $s^a$. Then the actual
word has the form
\[
 z=h+t\sum_{\eta\ne1} k_\eta\eta,
 \qquad k_\eta\in\mathbb Z_{\ge0},\qquad
 k=\sum_{\eta\ne1}k_\eta\ge1,
 \qquad h=g-tk>0.
\]
For some $d\mid\gcd(k,s-1)$ one has
\begin{equation}\label{paper:one-prime-energy}
 E(z)\ge\frac{s^{a-1}}d
       \bigl(s(h^2+t^2k)-g^2\bigr)
       >\frac23t^2s^a.
\end{equation}
Here $E(z)=[K:\mathbb Q]\mathsf M(z)$ is the energy defined in
Section~\ref{paper:sec:residue-energy-2}.
\end{proposition}
\begin{proof}
Suppose first that both supports are nontrivial. Let $u,v$ count
their nonidentity $q$- and $r$-components. The margin divisibilities
give $ru\le g$ and $qv\le g$. The identity belongs to both primary
supports: its marginal mass is $g$ minus a multiple of the other
prime, and hence is positive. Thus the two supports have sizes
$u+1,v+1$.

Put $F=\mathbb Q(\mu_{q^ar^b})$ and $H=\operatorname{Gal}(F/K)$.
Uniqueness of the actual coefficient vector makes $H$ preserve it.
Its projection onto the $q$-factor acts freely on each occurring
primitive $q^a$ component, so that projection has order at most $u$.
The other projection has order at most $v$. Consequently
$|H|\le uv$ and
\[
 [K:\mathbb Q]\ge
       \frac{q^{a-1}(q-1)r^{b-1}(r-1)}{uv}.
\]
The normalized trace Gram matrix of distinct $q$-primary roots has
diagonal entries one; its off-diagonal entries are zero or
$-1/(q-1)$. On the $u+1$ occurring roots its least eigenvalue is
therefore at least $(q-1-u)/(q-1)$. The analogous bound at $r$ is
$(r-1-v)/(r-1)$. Coprimality makes the joint Gram matrix their
tensor product. Cauchy's inequality for the coefficient array gives
\[
 E(z)\ge q^{a-1}r^{b-1}g^2
       \frac{q-1-u}{u(u+1)}
       \frac{r-1-v}{v(v+1)}.
\]
Each fraction decreases with its positive argument in the range
used here. Substitute $u\le g/r$, $v\le g/q$ to obtain the first
inequality in \eqref{paper:mixed-energy}.

For completeness, the elementary strict inequality needed for the
second one is
\[
 g(qr-r-g)(qr-q-g)>qr(g+r)(g+q),
 \qquad 3q\le g\le(q-2)^2,\quad r\ge3q,\quad q\ge7.
\]
The lower bound $g\ge3q$ follows here from $u\ge1$ and $r\le g$.
Let $F_q(g,r)$ be the difference of the two sides. As a polynomial
in $r$ it is
\[
 q\bigl(g(q-2)-q\bigr)r^2
 -g\bigl(2q^2-q+(3q-1)g\bigr)r+g^2(q+g).
\]
It increases for $r\ge3q$: using $3q\le g\le q^2$, its derivative
there is at least $3q^2(3q^2-15q+1)>0$.
At $r=3q$ it is concave in $g$ on the stated interval. Its values
at the endpoints $g=3q$ and $g=(q-2)^2$ are, respectively,
\[
 9q^3(3q^2-18q+8)
 \quad\text{and}\quad
 q^4(q^2+4q-62)+4q(29q^2-5q-32)+64,
\]
both positive for $q\ge7$. This proves the strict inequality.

Now suppose only the $s$-primary support is nontrivial. The margins
give the stated coefficients, and $h>0$ since $t\nmid g$.
By Corollary~\ref{paper:phase-conductors}, the stabilizer in
$\operatorname{Gal}(\mathbb Q(\mu_{s^a})/\mathbb Q)$ is tame, of
some order $d\mid s-1$. It acts freely on the nonidentity support
and preserves its coefficients. Thus $d\mid k$ and $d\le k$.
The normalized trace of the quotient of two distinct $s$-primary
roots is at least $-1/(s-1)$. Positivity of the coefficients and
$\sum k_\eta^2\ge k$ therefore give the first inequality in
\eqref{paper:one-prime-energy}. Its numerator can be written as
\[
 (s-1)\left(h-\frac{tk}{s-1}\right)^2
       +st^2k\left(1-\frac{k}{s-1}\right).
\]
If $(s,t)=(q,r)$, then
$k\le g/r\le(q-2)^2/(3q)<(q-1)/3$.
If $(s,t)=(r,q)$, then
$k\le g/q\le(q-2)^2/q<q-1<(r-1)/3$.
In either case $k/(s-1)<1/3$. Divide by $d\le k$ to obtain the
strict lower bound.
\end{proof}


\subsection{Complementary principal degrees}

\begin{proposition}[Inertia forced by several positive residues]\label{paper:4.8}
Let $z$ be a cyclotomic integer, all of whose conjugates have
modulus strictly less than a positive integer $g$. Suppose $z\equiv g$
at every prime above each prime in a finite set $P$.
For every nonempty proper subset $R\subset P$ let $H$ be the subgroup
of $\Gal(\QQ(z)/\QQ)$ generated by the inertia groups at $s\in R$.
Then
\[
                         \prod_{s\in R}s < |H|g.
\]
In particular any $s\ge g$ in a pair of prescribed positive residues
ramifies. If $z\ne0$ vanishes at every prime above $s$, then $s\ge g$ forces ramification at $s$.
\end{proposition}
\begin{proof}
Put $e=|H|$, $M=\prod_{s\in R}s$ and $w=\sum_{\sigma\in H}\sigma z$.
Its field is unramified at the primes in $R$, and $w\equiv eg$ at
every prime above $P$. Therefore $(w-eg)/M$ is integral. If $M\ge eg$,
every conjugate of $1+(w-eg)/M$ has modulus strictly less than one,
forcing $w=eg-M$. A prime of $P\setminus R$ would divide $M$.
In the last assertion every prime of $\QQ(z)$ above $s$ divides $(z)$. If $s$ were unramified, their product would be $(s)$, so $z/s$ would be a nonzero algebraic integer of house less than one. Lemma~\ref{paper:family:lem:principal-all-places} supplies the all-places hypothesis when the residue comes from a principal row.
\end{proof}

\begin{lemma}[Complementary principal degrees at three primes]\label{paper:family:4.10}
Let a perfect fusion-simple ordinary or positive reduced table have
three dimension primes $t,u,v$. In the reduced case assume $u$ odd.
The principal block $P_u$ contains $v$-free nonunits. Every such
degree is $t^iu^j$ with $i,j\ge1$ and
$2j+2\le v_u(\Delta)$. Thus at a quartic $u$ both complementary
degrees $ut^i,uv^k$ occur in $P_u$, each of defect two.

When $u<v$ are odd, the $t$-parts of all $v$-free nonunits in $P_u$
and all $u$-free nonunits in $P_v$ are at least $u$.
\end{lemma}
\begin{proof}
Prime avoidance (Proposition~\ref{paper:family:4.2}) for a
$u$-active column produces a $v$-free nonunit. It cannot have
prime-power degree, and Corollary~\ref{paper:family:4.1}
bounds its $u$-height. At $v_u(\Delta)=4$ interchange $t,v$.
For $u<v$ odd pair opposite labels $X,Y$ of degrees
$t^iu^j,t^kv^l$. The lcm-normalized entry has weight
$g=t^{\min(i,k)}$ by balancing and strict house below $g$.
Principality gives residue $g$ over both $u,v$, on all conjugates
by Galois permutation of labels and symmetry. Theorem~\ref{paper:family:thm:two-positive}
therefore gives $g\ge u$.
\end{proof}

\begin{corollary}[The two-adic height in a reduced table]\label{paper:4.11}
For a slightly degenerate minimal counterexample of dimension
$2^a q^b r^c$ with $q<r$ odd, one has
$2^{\lfloor(a-1)/2\rfloor}\ge q$. In particular $a\ge5$.
\end{corollary}
\begin{proof}
The positive table has two-valuation $a-1$. Square-divisibility
bounds all two-parts by $2^{\lfloor(a-1)/2\rfloor}$, and
Lemma~\ref{paper:family:4.10} bounds some two-part below by $q$.
Sources losing a repeated prime are covered by Dong.
\end{proof}


\section{Proper divisors: prime order and odd faces}\label{paper:sec:odd-faces}

\subsection{Small exponents and the order of the primes}

\begin{theorem}[Three primes with a small exponent]\label{paper:4.12}
Let $p,q,r$ be distinct primes and $0\le a,b,c\le6$. Every braided fusion category of
dimension $p^a q^b r^c$ is WGT if $c\le3$.
\end{theorem}
\begin{proof}
A divisor-minimal failure has a perfect simple ordinary or reduced
table of square-full dimension $\Delta=D$ or $D/2$
(Lemma~\ref{paper:family:2.1}); Dong's theorem, including its coprime square-free factor, excludes fewer than three repeated primes \cite{DongWeakIntegral}. Thus all three primes survive in $\Delta$. In ordinary data every prime valuation is at least four by
Lemma~\ref{paper:family:4.10}. In reduced data that lemma
applies at every odd prime, and Corollary~\ref{paper:4.11}
treats two.

\end{proof}

\begin{lemma}[Separation at a quartic prime]\label{paper:7.1}
Let an ordinary perfect fusion-simple integral modular table have dimension $t^a u^4v^b$, $u<v$, with three repeated primes. No twist order is divisible by $uv$. The same holds in a reduced minimal table for odd $u,v$.
\end{lemma}
\begin{proof}
Lemma~\ref{paper:family:4.10} gives $X\in P_u$ of degree $ut^i$. Every nonunit $t$-free $Y\in P_v$ has degree $u^hv^j$, $h=1,2$. If $h=1$, $S_{XY}/\lcm(d_X,d_Y)$ has weight $u$, strict house below $u$ and residue $u$ at both $u,v$ by principality (at $u$ both are zero). This contradicts Theorem~\ref{paper:family:thm:two-positive}, also at $u=2$. Thus all such $Y$ have $u$-unit codegree and vanish on $u$-active columns. For a joint $u,v$-active column, reduction of its orthogonality on $P_v$ modulo $t$ would leave only the vacuum. At odd $u,v$ the same argument uses positive balancing and reduced block orthogonality.
\end{proof}

\begin{theorem}[A quartic below a larger prime]\label{paper:7.3}
Let $p<q<r$ and $0\le a,b,c\le6$. Every braided category of dimension dividing
$p^a q^b r^c$ is WGT when
$\min(a,b)\le4$.
\end{theorem}
\begin{proof}
By descent and Theorem~\ref{paper:4.12}, only an exact quartic $h$
below a larger prime $k$ remains. An ordinary table with $h$ least
is already impossible: Lemma~\ref{paper:7.1} separates $h$
from both larger primes, contrary to Lemma~\ref{paper:lw:isolated-quartic}.
Corollary~\ref{paper:4.11} disposes of a
reduced table at $h=2$. Thus suppose $h,k$ are odd, and denote the
third prime by $t$. In a reduced table $t=2$, since a category containing $\sVec$ has even dimension.

We work in the ordinary table or its integral minimal extension.
In the latter, separation of $h,k$ in the even sector extends to all
labels. Indeed the joint active space is invariant under the dyadic
Weyl. This operator interchanges the fermion permutation $F$ and the
monodromy sign $G$. A space supported on odd labels contains no
$G$-fixed vector, but contains the $F$-fixed vector $v_X+v_{gX}$
for each of its fermion pairs. Thus $h,k$ are separated by
Lemma~\ref{paper:7.1} and no character field ramifies at both
(Proposition~\ref{paper:lw:isolated}).

There is a joint $t,h$-active column $J$: in the ordinary case
otherwise $h$ is isolated; in the extension choose an even $h$-active
pair, whose opposite twists allow a dyadically active member.
Choose $A\in P_h$ of degree $hk^j$ and $B,C\in P_k$ respectively of
degrees $h^v k^l,t^w k^m$, with all powers positive, by Lemma~\ref{paper:family:4.10} and Proposition~\ref{paper:family:4.2}. Choose full principal representatives in the extension by Lemma~\ref{paper:odd-principal-lift}. Orthogonality on $P_t$ modulo $k$ now gives a $k$-free
noninvertible hit $V$ on $J$: the invertibles contribute one, or two
in the extension since $J$ is even. Defect-zero vanishing bounds its
$h$-height by one. It hits both $A$ and $C$ by principality, so
coprime vanishing gives $d_V=t^i h$, $i>0$.

The two normalized entries $z_Y=S_{VY}/\lcm(d_V,d_Y)$,
$Y=A,B$, belong to $K_V$ and have strict house less than $h$.
At $A$ one has $t^i z_A=\chi_A(V)\equiv d_V\equiv0$ at every place above $h$: Galois covariance permutes $V$ through labels of the same degree (Lemma~\ref{paper:family:lem:principal-all-places}), and $t$ is an $h$-unit. Yet
$z_A\ne0$ since $V\in P_t$ and $t\nmid d_A$. At $B$ the
principalities of $V,B$ give residue $h$ at both $t$ and $k$.
The zero-residue and two-residue assertions of
Proposition~\ref{paper:4.8} force ramification of $K_V$
at $h$ and at $k>h$, respectively, contradicting separation.
\end{proof}

\subsection{The quartic alternative}

\begin{proposition}[A primary seed at the quartic prime]
\label{paper:family:prop:quartic-seed}
Let an ordinary perfect fusion-simple integral modular
category have precisely the odd dimension primes $p<q<r$
and $r$-valuation four; the other two valuations may be
arbitrary. Put $m=\lfloor r/q\rfloor$. If $m+1<p$,
some nonunit character at a nonunit row has, up to a
scalar root, a full positive $r$-word on at most $m+1$
positions. In the target family, $r<p^2$ implies
$m+1<p$, since $(p-1)q>p^2$.
\end{proposition}
\begin{proof}
Choose $X\in P_r$ of degree $p^i r$. A $p$-free nonunit
$B\in P_q$ has degree $q^j r^v$, where $v=1$ or $2$.
If one has $v=1$, its crossing with $X$ has normalized
mass $r$, strict house below $r$, and residues $r,0$
at $q,r$. Lemma~\ref{paper:family:lem:prime-mass}, followed by
rescaling to character mass, gives the claimed seed.

Suppose instead all these $B$ have $r$-height two.
Their codegrees are $r$-units, so they miss every
$r$-active column. A joint $q,r$-active column would,
by prime avoidance on $P_q$ modulo $p$, have a nonzero
hit at such a $B$. Thus $q,r$ are separated.
The quartic $r$ is not isolated, so there is a joint
$p,r$-active column $J$. Prime avoidance on $P_p$ in
the trivial prime-to-$p$ twist fiber, modulo $q$, gives
a $q$-free nonunit $V$ hitting $J$, with degree $p^e r$.
Its twist is $p$-primary; its $r$-height is one because
height two has $r$-unit codegree and misses $J$.

Take any $p$-free $A\in P_r$, of degree $q^i r$.
The normalized $V,A$ crossing has mass $r$, nonzero
value by $V\in P_p$, and zero residue at $r$ by
$A\in P_r$. If unramified at $r$, division by $r$
would be a nonzero algebraic integer of house below one.
Thus $K_V$ ramifies at $r$ and, by separation, is
unramified at $q$. For each above $B$, the normalized
$V,B$ crossing $z_B$ lies in $K_V\cap K_B$; the
second field is unramified at $r$. Hence $z_B$ is
$p$-primary. It has mass $r$, strict house below $r$,
and residues $r$ at $p,q$.

Average its actual word over the complementary $q,r$
factor. A negative complementary trace has absolute
value at most $1/(q-1)$; replace a negative coefficient
by positive coefficients on the other $p-1$ positions
of its $\mu_p$-coset. This gives a positive rational
$p$-word of mass at most $r$. Subtract its minimum
coefficient on each coset. The resulting canonical
coefficients $c_i$ are integers by the integral basis
lemma, and their length $L$ satisfies
\[
 0\le L\le r,\quad L\equiv r\pmod p,\qquad
 c_i-r\delta_{i0}\ \text{constant modulo $q$ on each coset}.
\]
Put $\rho=r-mq$ and use the hypothesis $m+1<p$.
Every nonidentity coset has a zero coefficient and
therefore consists of multiples of $q$. If zero were
the sole vacant position of the identity coset, all
its other coefficients would be $q-\rho$ modulo $q$.
Thus $L=(p-1)(q-\rho)+qn$ for $0\le n\le m$.
Reduction modulo $p$ gives $n\equiv m+1\pmod p$,
which is impossible. The identity coset also has a
vacant nonidentity position. All coset constants now
vanish, so $L\equiv r\pmod{pq}$. Since $r<pq$,
this forces $L=r$.

Equality of masses identifies the actual phases.
Every nontrivial complementary component loses mass
under the averaging and conversion: a positive trace
is strictly below one, while a negative trace costs
at most $(p-1)/(q-1)<1$. Thus every actual phase of
$V\otimes B$ is $p$-primary. Since $V$ and $B$ are
$r$-inactive, every constituent of $VB$ is $r$-inactive.
The Fourier transform of its fusion vector is
\[
       (s[VB])_T=d_T\chi_T(V)\chi_T(B)/\sqrt D,
\]
which also vanishes on $r$-active labels, since $B$
has $r$-unit codegree. Both the upper and lower local
unipotents fix $[VB]$, and hence its local torus fixes
it. At the even $r$-valuation this torus is the unsigned
inertia permutation. As $B$ is $r$-unramified,
\begin{equation}\label{paper:family:eq:quartic-inertia}
                  B V^\sigma=BV\qquad(\sigma\in I_r).
\end{equation}

The preceding ramification argument applies to each
$p$-free $A\in P_r$. Some $\sigma\in I_r$ changes
$\chi_A(V)$, so evaluating \eqref{paper:family:eq:quartic-inertia}
at $A$ gives $S_{AB}=0$ for every $p$-free nonunit
$B\in P_q$. Such an $A$ is $q$-inactive: separation
makes each $q$-active label have $r$-unit codegree,
whereas $r\mid f_A$. Choose a $q$-active column $T$.
Principal block orthogonality gives
\[
          0=\sum_{W\in P_q}\chi_A(W)\overline{\chi_T(W)}.
\]
The $p$-free nonunit terms vanish. At $p$-bearing $W$,
LCM integrality and $p\nmid d_A$ make $\chi_A(W)$
divisible by $p$. The unit contributes one, giving a
contradiction after reduction at $p$.
\end{proof}

\begin{corollary}[The initial quartic ratio]
\label{paper:cor:first-ratio}
Under the general ordinary hypotheses of
Proposition~\ref{paper:family:prop:quartic-seed}, $r<2q$
is impossible.
\end{corollary}
\begin{proof}
Here $m=1$, so the proposition gives a full primary
word on at most two positions. The difference set
has at most three positions. Since $r\ge p+4$,
$3p<(p-1)r$ and it is exposed. In
Theorem~\ref{paper:family:thm:descent}, $E=1$ and the
necessary inequality is $H-1\le0$, contrary to
$r>q>p\ge3$.
\end{proof}

\subsection{The two odd quintic faces}

\label{paper:app:odd-faces}
Only two proper-divisor faces of the ordered family remain after the cubic and smaller-quartic exclusions: $p^5q^6r^4$ and $p^6q^5r^4$.  We prove precisely the ordinary odd-table obstructions needed for them.  No assumption about the target corner is used in these proofs.


\begin{lemma}[Gauss sums at a quadratic class]\label{paper:fam:Q:5.1}
Let $D=p^{2k+1}n^2$, $k\ge1$, $p\nmid n$, be the dimension of
an ordinary integral modular category, with $p$ odd. Fix one of the
two square classes of primitive $p$th roots. Suppose every label
whose $p$-part of the twist is neither $1$ nor in that class has
degree $p^i n$ for some $i\ge0$. Let $L_i$ count the degree-$p^i n$
labels at orders $p^a$ with $a\ge2$, and let $M_i$ count such labels
at one fixed primitive root of the other class. Then
\[
 p^{k+1}\ \Bigm|\ \sum_i p^{2i}(L_i+pM_i).
\]
Each $L_i$ is divisible by $p$. In odd dimension $L_i$ and $M_i$
are even.
\end{lemma}
\begin{proof}
Square Galois covariance preserves degrees. A square-unit orbit of
primitive roots of order $p^a$, $a\ge2$, is a union of full
$\mu_p$-cosets, so its sum is zero and its size is divisible by $p$.
At order $p$, multiplicity is constant within each square class.
In odd dimension every nonunit has a distinct dual of the same
twist (Lemma~\ref{paper:family:2.4}).

Write $L=\sum_i p^{2i}L_i$, $M=\sum_i p^{2i}M_i$ and put
$A=(p^{2k+1}-L)/p$. By the odd-valuation slice
(Proposition~\ref{paper:3.8}), the squared-dimension mass at twist
one is $n^2 A$, and the mass at order $p$ is $(p-1)n^2 A$.
In units of $n^2$ the mass per primitive phase is therefore $M$
on the other square class and $2A-M$ on the chosen class.
The sums of the roots in these classes add to $-1$ and differ
by a quadratic Gauss sum of squared modulus $p$. Thus for
$T=\sum_Xd_X^2\theta_{X,p}$ we have $|T|^2=p n^4(A-M)^2$.
The regular-object indicator at an exponent selecting the
$p$-primary twist is $|T|^2/D$, a rational algebraic integer
by \cite[Theorem~7.5]{NSfs} summed over the simples with their dimension multiplicities. It follows that $p^k\mid A-M$, as asserted.
\end{proof}

\begin{corollary}[Exceptional labels and principal inertia]\label{paper:fam:Q:6.6}
In Lemma~\ref{paper:fam:Q:5.1}, suppose either that all exceptional labels
have degree $n$, or that $k\le2$. Then $p^{k+1}\mid L_0+pM_0$.
If $L_0+M_0>0$, the number $N$ of degree-$n$ exceptional labels satisfies
$2N\ge c(p-1)p^k$, where $c=2$ in odd dimension and $c=1$ otherwise.
For $U\in P_p$ of $p$-height $j$, an inertia orbit on $U$ of
length $e$ at a prime $s\ne p$ then satisfies
\[
                       c(p-1)e p^{2j}<2p^{k+1}.
\]
\end{corollary}
\begin{proof}
Higher strata in the congruence vanish under either hypothesis:
$p\mid L_i$ gives $p^3\mid p^{2i}(L_i+pM_i)$ for $i\ge1$.
Thus a nonzero $L_0+pM_0$ is at least $cp^{k+1}$, by parity in
odd dimension. Its contribution to $N$ is
$L_0+(p-1)M_0/2\ge(p-1)(L_0+pM_0)/(2p)$.
Each of these labels has codegree $p^{2k+1}$, so its field is
unramified at $s$. Lemma~\ref{paper:family:2.3} gives $eN(p^jn)^2<D$.
\end{proof}

\begin{lemma}[A trace on the trivial twist fiber]\label{paper:fam:Q:5.2}
Let an ordinary odd minimal table at three primes have dimension
$D=t^5n^2$, with $t$ prime and $t\nmid n$. Suppose the $t$-primary twist level
is $t$ and one quadratic class of primitive $t$th roots occurs
only on rows of codegree $t^3$ and field $\QQ(\zeta_t)$.
If that class occurs, then for each prime $s\mid n$ at least one
$s$-free nonunit in $P_t$ is fixed by $s$-inertia.
\end{lemma}
\begin{proof}
Fix a label $Z$ in that class and an outer prime $s\mid n$.
Its degree is $tn$. The two complementary local groups act by
scalars on each twist fiber of this class. Indeed the fibers are
preserved by each complementary group, their labels have
complementary unit codegrees, and the upper unipotents and their
conjugates generate these groups. Thus the $t$-local Weyl $W$
agrees with $s$ up to a root of unity on each such coordinate vector.
The same holds for the torus translates of $Z$.

Put $U=U_t$. For $a\in\mathbb F_t$ set
\[
 M_a=\sum_V d_V\overline{\chi_Z(V)}\theta_{V,t}^{a}
     =\frac{D}{d_Z}(sU^a s^{-1})_{0Z}.
\]
Orthogonality gives $M_0=0$. For $a\ne0$ use the local Bruhat identity
\[
               WU^aW^{-1}=U^{-1/a}WH_{t,-a}U^{-1/a}.
\]
It gives $M_a$ equal to a root of unity times $\sqrt D$. Indeed
the diagonal factors contribute phases, and the vacuum entry of
$W$ at the torus translate has modulus $d_Z/\sqrt D$ and
agrees with that of $s$ up to a root of unity. Since
$\sqrt D=t^2 n\sqrt t$ and $s\mid n$, summing over $a$ gives
\[
           \sum_{\theta_{V,t}=1} d_V\overline{\chi_Z(V)}
                      \equiv0\pmod {s\overline{\mathbb Z}} .
\]
This set of labels is Galois stable. Its vacuum contribution is
one. Grouping by Galois orbit therefore supplies an $s$-free
nonunit $X$ with trivial $t$-twist and nonzero absolute trace of
$x=\chi_Z(X)$. Its row meets $Z$, so $X\in P_t$ by the
prime-power codegree support theorem (Theorem~\ref{paper:family:2.7}).

Suppose every such $s$-free nonunit is moved by $s$-inertia.
Let $A$ be the span of the $s$-local translates of the
$s$-active coordinates. It is $t$-locally invariant. The chosen
quadratic class is absent from $A$, since it is absent on
$s$-active coordinates and the $s$-group preserves their
$t$-phases. On $A^\perp$ the $s$-upper unipotent and all its
conjugates are the identity. Let $P$ project onto the $U$-fixed
space. For a nonzero torus eigenvector $v\in PA$ with
eigencharacter $\lambda$, the Fourier projector at a missing
phase $\zeta_t^c$ vanishes on $W^{-1}v$. The same Bruhat identity
gives
\[
 0=\sum_{a\bmod t}\zeta_t^{-ca}P WU^aW^{-1}v
   =v+\left(\sum_{a\ne0}\zeta_t^{-ca}\lambda(-a)\right)PWv .
\]
Consequently $PWv\ne0$, and the parenthesized sum has the same
nonzero value for every $c$ in that class. Replacing $c$ by
$cu^2$ multiplies it by $\lambda(u^2)^{-1}$. Thus every square
torus acts trivially on $PA$. On $P A^\perp$ the $s$-group acts
trivially, so on the entire $t$-inactive space
\[
                (H_{t,u^2}-1)(H_{s,b}-1)=0 .
\]
Both tori act by unsigned Galois permutations since their
automorphisms fix $\sqrt D$. Take $b$ moving $X$. If $H_{t,u^2}$
also moved $X$, the coefficient at $X$ in the mixed difference
of $v_X$ would be positive. Hence all square $t$-units fix $X$.
The entry field $F=\QQ(x)\subseteq K_X\cap K_Z$ has degree at
most two, since $K_Z=\QQ(\zeta_t)$. Also $[K_X:F]\ge2$, since
$s$-inertia fixes $K_Z$ and moves $X$.

Write $u'$ for the other outer prime. The $s$-free nonunit $X$
has degree $t(u')^i$: the principal valuation bound gives
$t$-height at most one and a nonunit prime-power degree is
impossible. The $t$-valuation of $d_Z$ is also one.
Since $X\in P_t$, symmetry gives
$x=(d_X/d_Z)\chi_X(Z)\equiv0$ above $t$. Its nonzero absolute
field trace is an integer multiple of $t$. In the trace-norm
identity (Lemma~\ref{paper:family:2.3}) this orbit contributes
\[
 [K_Z:F][K_X:F]\operatorname{Tr}_{F/\QQ}(x)^2
                         \ \ge\ (t-1)t^2 .
\]
Reducing the ordinary dimension--trace identity at $u'$ supplies
a second, distinct orbit, of degree $t s^j$, with a nonzero trace
also divisible by $t$. It contributes at least $t^2$.
Together with the positive vacuum term this exceeds $t^3=f_Z$.
\end{proof}

\begin{proposition}[An odd quintic with two complementary witnesses]\label{paper:fam:Q:5.3}
Let an ordinary odd minimal table have three dimension primes
$t,u,s$, where $s>\max(t,u)$ and $v_t(D)=5$. Suppose the
$s$-free nonunits of $P_u$ and the $u$-free nonunits of $P_s$
all have $t$-height two. Some $s$-free nonunit of $P_t$ is fixed
by $s$-inertia.
\end{proposition}
\begin{proof}
Suppose all $s$-free nonunits of $P_t$ are moved. Call the two
sets in the hypotheses the witnesses. They exist by prime avoidance.
Each witness has $t$-defect one. Lemma~\ref{paper:lem:quintic-common-class}
gives a single visible quadratic class for both sets. Each $s$-active or $u$-active label hits
a corresponding witness by principal-block orthogonality modulo
the omitted prime (Proposition~\ref{paper:family:4.2}). All their
$t$-phases are therefore $1$ or in the chosen class. There is a
label outside these phases: either a deeper twist occurs or
Lemma~\ref{paper:7.4} gives the opposite primitive class.
Lemma~\ref{paper:lem:visible-phase-transfer} applied at $s,u$ gives every label $Z$ of such a fiber pure $t$-power codegree. It follows that $D=t^5n^2$ and the
possible codegrees $t,t^3,t^5$ have degrees $t^2n,tn,n$,
respectively. Each exceptional Galois orbit contains at least
$t-1$ labels: square covariance gives at least $(t-1)/2$
distinct phases, and both duals occur at each phase in odd
dimension.

The row of $Z$ is supported on $P_t$. The dimension--trace
equation reduced at each outer prime gives separate nonunit
orbits of degrees $tu^i$ and $ts^j$ with nonzero entry trace.
Each contributes at least one to the trace norm, while the vacuum
contributes $[K_Z:\QQ]\ge t-1$. Hence $f_Z\ne t$.
If an exceptional label had degree $n$, take an $s$-free nonunit
$U\in P_t$. It has degree $tu^i$ and inertia orbit length $e\ge2$
by assumption. Corollary~\ref{paper:fam:Q:6.6}, with $k=2$, $j=1$ and
$c=2$, would give $(t-1)e<t$. Every exceptional row therefore has
codegree $t^3$.

To determine its field, take the space $A$ spanned by the $s$-local
translates of the $s$-active labels as in Lemma~\ref{paper:fam:Q:5.2}.
On $A$ the only $t$-phases have order at most $t$. Thus $U_t(t)$
and its conjugates are trivial there. The upper--lower
factorization
\[
 U_t(1)L_t(b)=L_t\bigl(b/(1+b)\bigr)H_{t,1+b}U_t\bigl(1/(1+b)\bigr)
\]
for the lower unipotent $L_t(b)$ at $b=tj$ gives $H_{t,1+tj}=1$ there:
$U_t(1/(1+tj))$ agrees with $U_t(1)$ and the lower parameters
are multiples of $t$. On $A^\perp$ the $s$-group is trivial.
The mixed-difference argument of Lemma~\ref{paper:fam:Q:5.2}
shows that every label moved by $s$-inertia is fixed by wild
$t$-inertia; both tori again fix $\sqrt D$. Apply this to
the $s$-free trace-contributing orbit above. Its entry field $F$
is tame at $t$, and its nonzero trace is divisible by $t$ since
$d_Z=tn$. If $K_Z$ were wild at $t$, then $t\mid[K_Z:F]$ and
this orbit alone would cost $t^3$, before the vacuum term.
The codegree makes $K_Z$ a field of $t$-power conductor
(Theorem~\ref{paper:family:2.5} and Corollary~\ref{paper:family:2.6}). Tameness now gives
$K_Z\subseteq\QQ(\zeta_t)$,
with equality by its degree bound. The local Galois unit $1+t$
fixes $Z$, and its square fixes the twist of $Z$. Hence that
twist has $t$-part of order at most $t$. This holds on every
exceptional row, and the other rows already have $t$-phase of
order at most $t$. The global $t$-primary level is thus $t$,
and Lemma~\ref{paper:fam:Q:5.2} gives the required contradiction.
\end{proof}


\begin{proposition}[The least quintic face]\label{paper:family:least-quintic}
There is no ordinary odd minimal table with three ordered primes
\(p<q<r\) and \(v_p(D)=5\).
\end{proposition}
\begin{proof}
Opposite \(r\)-free nonunits in \(P_q\) and \(q\)-free nonunits
in \(P_r\) have common \(p\)-part at least \(q\), by
Lemma~\ref{paper:family:4.10}.  Square divisibility bounds their
\(p\)-heights by two; since \(q>p\), all these heights are two.
Proposition~\ref{paper:fam:Q:5.3}, with \(t=p,u=q,s=r\), gives
an \(r\)-free nonunit \(X\in P_p\) fixed by \(r\)-inertia.
Its degree is \(pq^i\), by the principal height bound.
Pair it with a \(q\)-free witness in \(P_r\).
The normalized entry has mass \(p\), strict house below \(p\),
and residue \(p\) at every place over \(r\).
Since \((r-p)(p-1)>p\), Lemma~\ref{paper:lem:one-residue}
gives an unsigned \(r\)-word of length \(p<r-1\).
Lemma~\ref{paper:lem:short-tame-orbit} makes its \(r\)-inertia orbit
nontrivial.  Its field is contained in \(K_X\), a contradiction.
\end{proof}

\begin{lemma}[The middle quintic when its short words are exposed]
\label{paper:family:middle-small}
An ordinary minimal table of dimension \(p^6q^5r^4\), for
odd primes \(p<q<r\), cannot satisfy
\[
 j=\lfloor q/p\rfloor,\qquad 1+j(j+1)<2q/3.
\]
\end{lemma}
\begin{proof}
Every \(p\)-free nonunit \(A\in P_r\) has \(q\)-height two.
If its height were one, pair it with an \(r\)-free member
of \(P_q\), of degree \(qp^i\).  The common mass is \(q\), and the two
positive residues at \(q,r\) force house \(q\) by
Theorem~\ref{paper:family:thm:two-positive}, contradicting strictness.

Suppose an \(r\)-free nonunit \(U\in P_p\) has \(q\)-height one.
Choose \(A\in P_r\) of degree \(q^2r\) as above and
\(B\in P_q\) of degree \(qr^k\).  On the single row \(U\),
the two normalized crossings have common mass \(q\).
The crossing with \(A\) has positive residues \(q\) at \(p,r\),
and is an \(r\)-word on at most \(1+j\) positions.
The exceptional alternative would require \(r=q+1\), impossible
for odd primes.  The crossing with \(B\) has residues \(q,0\)
at \(p,q\), so Lemma~\ref{paper:family:lem:prime-mass} makes it,
up to a scalar, a \(q\)-word on at most \(1+j\) positions.
After taking adjoints the two full character words are supported
on \(r\)-primary and \(q\)-primary difference sets of at most
\(1+j(j+1)\) positions.  Both are exposed by
Lemma~\ref{paper:family:lem:exposure}; their intersection is \(\{1\}\).
A common nonunit constituent of \(AA^*\) and \(BB^*\)
(Lemma~\ref{paper:family:lem:common}) would therefore have only phase \(1\)
on the row \(U\), contradicting strictness.
Consequently all \(r\)-free nonunits in \(P_p\) have \(q\)-height
two.

Proposition~\ref{paper:fam:Q:5.3}, now with \(t=q,u=p,s=r\),
gives an \(r\)-free nonunit \(X\in P_q\) fixed by \(r\)-inertia.
Here \(d_X=qp^i\).  Pairing it with a \(p\)-free witness in
\(P_r\) gives a mass-\(q\) entry of residue \(q\) at \(r\).
If \((r-q)(p-1)>q\), Lemma~\ref{paper:lem:one-residue} and
Lemma~\ref{paper:lem:short-tame-orbit} contradict that inertia fixation.
Otherwise \(r<2q\), which is excluded by
Corollary~\ref{paper:cor:first-ratio}.
\end{proof}

\begin{proposition}[The middle quintic face]\label{paper:family:middle-quintic}
There is no ordinary odd minimal table of dimension
\(p^6q^5r^4\), for \(p<q<r\).
\end{proposition}
\begin{proof}
Every \(r\)-free nonunit in \(P_q\) and \(q\)-free nonunit in
\(P_r\) has \(p\)-height two or three:
their common \(p\)-part is at least \(q>p\).
These heights cannot all be three.  A nonempty \(p\)-active
fiber contains a \(q\)-active column, because the \(q\)-valuation
is odd (Lemma~\ref{paper:7.8}).
Prime avoidance on \(P_q\) modulo \(r\) forces a complementary
witness hitting that column, whereas a witness of \(p\)-height
three has \(p\)-unit codegree and misses it.

Some opposite crossing therefore has mass \(G=p^2\).
Put \(M=\lfloor p^2/q\rfloor\).  If \(r>p^2\),
Theorem~\ref{paper:family:thm:two-positive} gives an \(r\)-word on
at most \(1+M\) positions.  If \(r<p^2\), choose an
\(r\)-active fiber.  The odd \(q\)-valuation supplies a joint
\(q,r\)-active column, and prime avoidance supplies a
\(p\)-free member of \(P_q\) hitting it.  Its \(r\)-height
is one by defect-zero vanishing.  Pairing it with a
\(q\)-free member of \(P_r\) gives, by
Lemma~\ref{paper:family:lem:prime-mass}, an \(r\)-word on at most
\(1+\lfloor r/q\rfloor\) positions.
In both cases the word has full character mass after rescaling,
and
\[
 m=\left\lfloor\frac{\min(r,p^2)}q\right\rfloor,\qquad
 1\le m\le M\le p-2.
 \tag{Q}
\]
Its adjoint support is exposed by
Proposition~\ref{paper:family:prop:auto-exposure}.

Suppose first that \(q\ge(M+1)^2\).  Then
\(r>mq\ge m(m+1)^2\).  Put \(E=m(m+1)/2\) and
\(g(h)=\min\{k\ge1:k(k+1)\ge h\}\).
For \(1\le h\le E\) one has \(h/g(h)\le m\):
if \(g(h)<m\), use \(h\le g(h)(g(h)+1)\), and if
\(g(h)=m\), use \(h\le E\).
Consequently
\[
 \frac{h(h-1)}{g(h)}\le m(E-1)
                    <\frac{r-3}{2}.
 \tag{Q'}
\]
The uniform difference descent
(Theorem~\ref{paper:family:thm:descent}) contradicts \textup{(Q')}.

It remains that \(q<(M+1)^2\).  As \(q\ge5\), we have
\(M\ge2\), hence \(q<p^2/2\).  Put \(j=\lfloor q/p\rfloor\).
Then \(j\le(p-1)/2\), and \(1+j(j+1)<2q/3\).
For \(j=1\), this follows from \(q\ge5\).
For \(j\ge2\), \(p\ge5\) and
\[
 2p-3(j+1)\ge(p-3)/2\ge1,\qquad
 3+3j(j+1)<2(pj+1)\le2q.
\]
Lemma~\ref{paper:family:middle-small} supplies the contradiction.
\end{proof}


\section{Fermion exclusion and even faces}\label{paper:sec:even-faces}

\subsection{Dyadic phases and signed transfer}

\begin{lemma}[Positive phases and signed pair transfer]
\label{paper:fermion-phase-pair}
Assume $v_2(N)$ is odd.
\begin{enumerate}
\item A positive two-block containing a defect-one character
has squared dyadic twists $+1,-1$, with equal degree-square mass.
Its positive rows vanish on all other squared dyadic phases.
\item If all even actual dyadic twists have order dividing four,
the full dyadic level divides eight. Each block in (1) has two
equal-degree positive labels, both unramified at two. A purely
dyadic full character takes Gaussian-integer values on even objects.
\item Under this bound $U^4$ is scalar on $C$. Above a block in
(1), of common positive degree $d$, there are exactly two odd
fermion orbits. Both are free, their spins have ratio $\pm i$,
and their degrees satisfy
\[
       u^2+v^2=2d^2,\qquad v_2(u)=v_2(v)=v_2(d).
\]
If all odd labels have only two dyadic spins of ratio $\pm i$,
$R_+$ has a nontrivial invertible basis element.
\end{enumerate}
\end{lemma}
\begin{proof}
Let $E$ be the block projector on $A$. It commutes with $K$.
Since $\widehat s^{-1}\mathbb C^E=W_o^{-1}\mathbb C^E$
has nonzero vacuum projection and $W_o$ preserves the squared
dyadic twist, the block contains squared twist one.
Here is the mass computation detecting the other sign. For odd
$k$ put $l=1+4k$ and
\[
 M(k)=\sum_{Y\in E}\frac{d_Y^2}{\Delta}\theta_{Y,2}^{-2k}.
\]
The torus $H_l$ negates the positive vacuum: it fixes its label
and negates $\sqrt N$. The vacuum entry of the Bruhat identity
$U(2)L(2k)=L(2k/l)H_lU(2/l)$, after the conjugate block projector,
is
\[
 \lambda^{2-2k}M(k)=-\lambda^{(2-2k)/l}M(k/l).
\]
The scalar exponents agree modulo eight. The substitution
$k\mapsto k/(1+4k)$ permutes the odd residues. Averaging gives
equal mass at squared phases $+1,-1$, because the average of
$\theta^{-2k}$ is respectively $1,-1,0$ at these two phases and
at higher order. By Corollary~\ref{paper:dyadic-node} there are at most
two inertia orbits. The two signs already occupy distinct orbits,
so no others occur. The equality of subspaces above then gives
the row-support assertion.

Assume now the even-twist bound. Balancing puts even-even entries
in $\mathbb Q(i,\zeta_m)$ for odd $m$. A dyadic Galois unit
$l\equiv1\pmod4$ fixes every even restricted character, hence
every even fermion pair. Its square fixes the actual twists, and
opposite twists distinguish the two representatives. Symmetry
then fixes the values of every character on even objects, so
every odd fermion orbit is fixed as a set. Taking $l=5$ and using
the square action on twists gives $\theta_{Y,2}^{24}=1$ for odd
$Y$. Thus the full level divides eight, and purely dyadic values
on even objects lie in $\mathbb Z[i]$. Positive inertia orbits
have length at most two. One branch in the defect-one block is
unramified; an orbit of length two on the other branch would make
the equal-mass identity an equality of an integer square with
twice an integer square. Hence both branches have length one
and equal degrees.

Put $Z=T_2^2|_A$ and $F_0=KZK^{-1}$. The torus $H_5=-I$ on
$A$. The same Bruhat identity modulo eight gives
$ZF_0=-F_0Z$. Also $H_3=-K^2$ commutes with both. Direct
matrix multiplication modulo eight gives
\[
 (WU)^{-1}U^2(WU)|_A=\lambda^4 F_0H_3Z.
\]
Its square is scalar, and $WU:A\to C$, proving the assertion
about $U^4$. Above a positive block of rank two, $H_3$ is scalar,
so the two eigenvalues of the right side are opposites. Therefore
the two odd spins have ratio $\pm i$.

The two odd characters have equal codegree valuations. If inertia
interchanges them this is immediate. Otherwise their signed local
order is $\{(x,y)\in\mathcal O^2:x\equiv y\pmod{2^b}\}$;
its conductor has length $b$ on both branches, so trace duality
gives equal codegree valuations. The restriction-codegree formula
in Lemma~\ref{paper:fermion-orders} then gives the same fixed/free
type and the same degree height. Block orthogonality of $1,g$
makes the odd mass equal the even mass $4d^2$. Two fixed labels
would give $u^2+v^2=4d^2$ with odd valuation on the left and even
valuation on the right. Thus both are free, giving the displayed
equation and height equality.

Finally suppose only two odd spins occur. On $C$ one has
$U=aI+bU^2$, with $|a|=|b|=1/\sqrt2$. Restricting
$WUW=U^{-1}WU^{-1}$ and using the preceding transfer identity
gives on $A$
\[
       K=a'ZK^2+b'F_0,\qquad |a'|=|b'|=1/\sqrt2.
\]
For $P=(1+Z)/2$, anticommutation gives
$PKP=a'PK^2P$. The complementary vacuum vector
$w=W_ov_{[1]}$ satisfies $Pw=w$, and $K^2w$ is a scalar multiple
of $w$. Since $Kw=\widehat s v_{[1]}$, it follows that
\[
 |w_{[j]}|^2=(2d_j^2/\Delta)\mathbf1_{\theta_{j,2}^2=1}.
\]
In the full vacuum column exactly one even representative of each
such pair occurs, the one of actual dyadic twist one. At a CRT
exponent zero on the dyadic level and one on the odd level, the
indicator formula therefore says that the coefficients of
\[
               1+J,\qquad J=\sum_{[j]}\theta_{j,2}^2e_{[j]},
\]
are algebraic integers. Galois invariance of the signs makes them
rational, hence integers. Now $J=J^*$, $J^2=1$, and
$\operatorname{FPdim}(J)=1$. Its based norm is one, so it is a
positive invertible basis element. Both signs occur, so it is
not the unit.
\end{proof}

\subsection{The fermion sextic exclusion}

\begin{theorem}[No minimal fermion category at dyadic valuation six]
\label{paper:fermion-six}
There is no integral slightly degenerate category $\mathcal F$
with pointed subcategory $\mathrm{sVec}$, no other proper
nontrivial fusion subcategory, and dimension $2^6q^br^c$, where
$q<r$ are distinct odd primes. The exponents $b,c$ are arbitrary.
\end{theorem}
\begin{proof}
Put $\Delta=\dim\mathcal F/2$. Its positive quotient is perfect
and fusion-simple. Opposite principal witnesses at $q,r$ have
degrees $2^i q^j,2^h r^l$. Their common dyadic weight is at least
$q$ by the two positive residues and at most four because
$v_2(\Delta)=5$. Hence $q=3$ and the common weight is four.
If $r>5$, the short-word theorem gives $1+3\xi$, with $\xi$
a nontrivial $r$-power root; the two-position obstruction excludes
it. Consequently $r=5$. This is a restriction inside a family
with both exponents still arbitrary. Every complementary witness
has degree $4\cdot3^i$ or $4\cdot5^j$, and positive dyadic
defect one.

Choose the integral extension of Lemma~\ref{paper:integral-extension}
and use the spaces and matrices of the preceding section.
Its dimension is $N=4\Delta$, of dyadic valuation seven.

\paragraph{The even twists have order at most four.}
By Lemma~\ref{paper:fermion-phase-pair}(1), complementary witness
rows see only squared dyadic phases $+1,-1$. An even pair at
any other squared phase cannot be active at three or five:
principal prime avoidance modulo the opposite prime would force
a complementary hit. On such an extra phase fiber, the odd local
groups preserve the fiber and their upper unipotents act as the
identity. The conjugate unipotents generate, so those groups fix
the full coordinate lines. The complementary-codegree principle
therefore makes the positive codegree of each extra pair purely
dyadic. It follows that $b,c$ are even; write
$n=3^{b/2}5^{c/2}$, so $\Delta=32n^2$.

For an extra pair $V$, $K v_V=\widehat s v_V$ has nonzero
vacuum coefficient, so $V\in P_2$. Its possible positive
codegrees are $32,8,2$. Only $32$ can occur. Indeed its character
vanishes modulo two on every supported even degree. Apply
dimension orthogonality after Galois trace and reduce once modulo
three and once modulo five. The vacuum contribution is one, so
there are two distinct prime-omitting nonunit orbits having
nonzero trace; both have even degree, since a nonunit degree
cannot be a prime power. Each trace is a nonzero even integer.
The Galois-averaged row norm therefore gives
\[
                       f_V\ge [K_V:\mathbb Q]+4+4>8.
\]
Thus $d_V=n$ for every extra pair.

Dimension orthogonality modulo five gives a five-free nonunit hit
$X$ on such a row. Pure-dyadic codegree puts it in $P_2$, whose
positive height bound gives $d_X=2\cdot3^h$. Crossing it with
a complementary $4\cdot5^j$ witness gives a nonscalar unsigned
five-word of length two. Its tame five-inertia orbit has length
$t\ge2$. Every extra pair is hit by this orbit, and its field is
five-unramified. If $L$ is the number of extra pairs, averaging
the $X$-row over this orbit and retaining the vacuum gives
\[
                              4tL<32.
\]
Hence $L\le3$.

If extra pairs exist, let $H_0$ be those whose squared dyadic
phase has maximal order $2^m$, $m\ge2$. On $A$ the corresponding
powered twist is $J=1-2P_{H_0}$. The group commutator
$JKJ^{-1}K^{-1}$ has algebraic-integer trace. Its scalar
normalization cancels, and expansion gives
\[
 \operatorname{Tr}(JKJ^{-1}K^{-1})
 =\dim A-4|H_0|+\frac18
                   \sum_{V,Y\in H_0}|\chi_V(Y)|^2.
\]
The last sum is rational by Galois invariance, so it belongs to
$8\mathbb Z$. Its entries reduce to one at two by principality,
so $|H_0|$ is even. Since $1\le|H_0|\le3$, it equals two.
If both labels are individually Galois fixed, the four terms are
odd integer squares and their sum is four modulo eight.
Otherwise they form a quadratic dyadic orbit. Symmetry makes the
squared-modulus matrix have diagonal entries $y$ and off-diagonal
entries $\sigma y$. If $y$ is rational it is an odd integer;
otherwise $y=a+b'\sqrt2$ with $a$ odd. The sum is respectively
$4y$ or $4a$, again four modulo eight. This contradiction proves
the even-twist bound.

Lemma~\ref{paper:fermion-phase-pair} now gives dyadic level at most
eight, Gaussian values on even objects for purely dyadic
characters, and two equal-degree positive labels in each
complementary block.

\paragraph{The odd spins leave a nonempty exceptional packet.}
Above each complementary block are two free odd orbits of spins
with ratio $\pm i$, and degrees $u,v$ satisfying
$u^2+v^2=2d^2$ and dyadic height two. If $d=4\cdot3^i$,
successive reduction modulo three gives $u=v=d$. If
$d=4\cdot5^j$, reduction modulo three gives $3\nmid uv$.
These odd mates provide a label of unit degree at either odd
prime for the opposite witness.

Each complementary positive witness has a full principal lift
at its odd prime. To see this, evaluate its character on the
square of an odd label of prime-free degree. The residue of its
value on that label is one of the two signs of the dimension.
Tensoring the witness by $g$ chooses the positive sign. Evaluation
on the even product with any other odd label then gives the
dimension residue there as well. Such a lift hits both odd mates
of every opposite block. The dyadic Weyl preserves its full
two-block, while the complementary Weyl preserves actual dyadic
spin. Therefore the same pair of spins, say $\delta,i\delta$,
occurs above every complementary block.

The full principal block at an odd prime contains only even
labels, by evaluation at $g$, and excludes $g$, by evaluation at
an odd label of prime-free degree. Principal prime avoidance
therefore makes every odd label active at three or five hit a
complementary witness. All such labels have spin $\delta$ or
$i\delta$.

There must be another odd spin, since otherwise the last assertion
of Lemma~\ref{paper:fermion-phase-pair} contradicts perfectness of
$R_+$. Since $U^4$ is scalar on $C$, the other spins are
$-\delta,-i\delta$. Denote their odd fermion orbits by $E$.
The odd local groups fix their full lines by the same fiber
argument used above. Thus their full codegrees are purely
dyadic; in particular $b,c$ are even, and we again write
$n=3^{b/2}5^{c/2}$. Each $V\in E$ is full two-principal because
$Wv_V=sv_V$ has nonzero vacuum entry. The full principal height
bound gives $d_V=2^j n$ for $j=0,1,2$.

\paragraph{A signed norm forces every exceptional orbit to be free and odd.}
The positive principal two-block has at least eight odd-degree
labels. Indeed $Z=T_2^2|_A$ and $KZK^{-1}$ anticommute, so its
rank is even. Its degree-square mass is divisible by $32$.
The positive principal height bound permits only odd degrees and
twice-odd degrees. If their counts are $k_0,k_1$, then
$4\mid k_0$, $2\mid k_1$, and
$k_0+4k_1\equiv0\pmod8$. The vacuum therefore gives $k_0\ge8$.
The signed character of $V\in E$ is nonzero on all these
odd-degree labels by principality; its Gaussian values contribute
at least eight to its norm.

Compare this signed character with an odd mate of a
$4\cdot5^j$ witness. Their restricted rows are orthogonal,
because their spins differ. Modulo three the vacuum survives,
and the three-bearing terms vanish on the three-free mate by
LCM integrality. There is therefore a three-free nonunit hit
$X_3$ on the $V$-row. Similarly there is a five-free hit $X_5$.
Pure-codegree support and the positive principal height bound give
\[
                     d_{X_3}=2\cdot5^l,
                     \qquad d_{X_5}=2\cdot3^h.
\]
Their nonzero Gaussian values vanish above two. An orbit of such
values costs at least four per label unless the squared modulus
is two. In the latter case $\chi_V(X)^2=\pm2i$, and
\[
                 \chi_X(V^2)=(d_V/d_X)^2\chi_V(X)^2.
\]
Since $V^2$ is even, $i$ lies in the positive character field of
$[X]$. Five-inertia fixes $i$, so the full orbit then has a second
coset with the same absolute squares. The five-inertia orbit of
$[X_5]$ has size at least two by its length-two crossing, while
the trivial lower bound for $[X_3]$ is one. These two disjoint
prime-omitting supports therefore cost at least eight and four.
The signed codegree is at least
\[
                               8+8+4=20.
\]
By Lemma~\ref{paper:fermion-orders} it is $32/4^j$ for a free
orbit and $64/4^j$ for a fixed one. Thus $j=0$.
A fermion-fixed odd label would have even $d_V^2$, because
$VV^*$ is $g$-invariant and all even constituents occur in free
pairs. Therefore every orbit in $E$ is free and has degree $n$.

\paragraph{The packet has fewer than four orbits.}
Choose $X=X_5$. Every $V\in E$ hits $X$ by two-principality;
LCM integrality makes $\chi_V(X)$ a nonzero multiple of two in
$\mathbb Z[i]$. Its squared modulus is at least four.
Let the five-inertia orbit of the positive label $[X]$ have
size $t\ge2$. It lifts to $t$ full representatives with a common
actual dyadic twist: five-inertia fixes their nonzero entries at
$V$, whereas replacing a representative by its fermion mate
negates those entries.

Write $b_X=(v_X-v_{gX})/\sqrt2$ and let $v$ be the sum over
these $t$ representatives. Then $\|v\|^2=t$. At the symmetric
odd coordinate of any $V\in E$, the entries of $sb_X$ agree and
have squared modulus at least $1/8$. Hence $t|E|\le8$.
If $|E|=4$, then $t=2$ and equality confines $sv$ to that odd
packet. The odd Weyls fix it, so $W_ov=v$. Applying $U$ and using
the common dyadic twist gives $W_oa=a$ for the corresponding
sum of symmetric even vectors. Consequently $sa=Ka$ is supported
on the positive principal two-block. But a complementary
$4\cdot5^j$ witness lies outside that block, and the corresponding
coordinate of $sa$ is a nonzero scalar times a sum whose residue
at five is $2d_X\ne0$. This is impossible. Thus
\[
                              1\le|E|\le3.
\]

\paragraph{A finite-group trace forces the packet size to be a multiple of four.}
Work on the odd antisymmetric space $H$. Set $x=\lambda\delta$
and $T=U/x$. Its eigenvalues lie in $\{1,i,-1,-i\}$. The packet
projection satisfies
\[
 J=1-2P_E=\tfrac12((1-i)T+(1+i)T^{-1}).
\]
Unitarity and cyclicity give
\[
 \operatorname{Tr}_H(JWJW^{-1})
 =\operatorname{Im}\operatorname{Tr}_H(TWTW^{-1})
  +\operatorname{Re}\operatorname{Tr}_H(TWT^{-1}W^{-1}).
\]
The words in $U,W$ underlying these traces are represented modulo
eight by
\[
 h=\begin{pmatrix}0&1\\-1&1\end{pmatrix},\qquad
 k=\begin{pmatrix}2&1\\1&1\end{pmatrix}.
\]
They satisfy $h^3=-I$, $k^3=5I$, and each is conjugate to its
inverse. The Weyl inverts $k$; the coordinate swap times $I+2h$
has determinant one and inverts $h$. Their traces are therefore
real algebraic integers in $\mathbb Q(\zeta_6)$, hence integers.
Replacing $U$ by $T$ multiplies the first trace by
$x^{-2}\in\mu_4$ and leaves the second unchanged. Thus the
displayed trace of $JWJW^{-1}$ is integral.
The odd Weyls fix $E$ pointwise, so projector expansion gives
\[
                   \sum_{V,Y\in E}|\chi_V(Y)|^2\in8\mathbb Z.
\]

Here every entry lies in $L=\mathbb Q(\zeta_8)$ and is a dyadic
unit by principality. Galois symmetry in either index holds up
to a fermion sign, so squared moduli are well defined on the
orbits. Partition $E$ into Galois orbits of sizes
$n_i\in\{1,2,4\}$. The sum on a rectangle of sizes $n_i,n_j$
is $n_in_jm_{ij}$, where $m_{ij}$ is the normalized trace of an
entry's squared modulus. For
$\alpha=\sum_{k=0}^3a_k\zeta_8^k$ one has
$m_{ij}=\sum_ka_k^2$, an odd integer since $\alpha$ is a dyadic
unit. If either orbit is a singleton, $\alpha$ lies on a rational
character line of $L$; integrality and unit valuation restrict it
to an odd integer times $1$ or $i$. Then $m_{ij}\equiv1\pmod8$.
For the other diagonal rectangles $4\mid n_i^2$, and the other
off-diagonal rectangles come in transpose pairs with
$8\mid2n_in_j$. Consequently
\[
 \sum_{V,Y\in E}|\chi_V(Y)|^2
                 \equiv\Bigl(\sum_i n_i\Bigr)^2=|E|^2\pmod8.
\]
It follows that $4\mid|E|$, contradicting $1\le|E|\le3$.
\end{proof}

\paragraph{Application to a minimal category in the exact family.}
Assume the minimal subcategory hypotheses of
Theorem~\ref{paper:fermion-six}; no statement about proper divisor
dimensions is needed for this implication itself.
If two has a sextic exponent, Theorem~\ref{paper:fermion-six} applies.
If two has the quartic exponent, the positive quotient has dyadic
valuation three, while opposite principal witnesses at the odd
primes require a common dyadic factor at least the smaller odd
prime. Square divisibility bounds that common factor by two,
a contradiction. Thus the exact family has no minimal slightly
degenerate branch, independently of a lower-face WGT theorem.

\subsection{The even proper-divisor faces}

The required faces are $2^5q^6r^4$ and $2^6q^5r^4$, with odd
primes $q<r$. Cubic exponents, a quartic below a larger prime,
and dimensions with fewer than three repeated primes were excluded
in Theorems~\ref{paper:4.12} and~\ref{paper:7.3} and by Dong's theorem~\cite{DongWeakIntegral}. Thus a divisor-minimal obstruction on these
faces has the displayed exponents, apart from the lower cases
explicitly eliminated along the way. The arguments below are
uniform in $q,r$.

\subsubsection{The ordinary dyadic quintic}

\begin{lemma}[A dyadic defect-one row and two visible phases]
\label{paper:ef:ordinary-dyadic-phases}
In an integral ordinary modular table with odd $v_2(D)$, a
defect-one row sees only actual dyadic phases $1,\eta$, for one
primitive fourth root $\eta$. Its full two-block has two labels
of equal degree, one of each phase. If an ordinary perfect
modular table has only those two dyadic phases globally, it is
impossible.
\end{lemma}
\begin{proof}
The dyadic conductor-node lemma gives at most two inertia orbits
in the block. Its local Weyl invariance and the complementary
Weyl's preservation of actual dyadic phase show that the block
contains phase one. The dyadic Bruhat mass calculation, with
weights $d_X^2/D$, equates the mass at orders one and two
together with the mass at order four. Explicitly, for odd $k$
and $l=1+4k$, the torus negates the vacuum and gives
\[
 \lambda^{2-2k}M(k)=-\lambda^{(2-2k)/l}M(k/l),\qquad
 M(k)=\sum_{X\in B}\frac{d_X^2}{D}\theta_{X,2}^{-2k}.
\]
Averaging over odd $k$ proves the assertion about the masses.
The two occupied inertia orbits must therefore have phases one
and a single primitive fourth root. On this block $U(4),L(4)$
are scalar. The identity
$U(1)L(4)=L(4/5)H_5U(1/5)$ makes $H_5$ scalar, so units
$1\bmod4$ fix the labels. The unramified orbit is a singleton,
the other has length at most two, and equal square mass rules
out length two. Thus the block has two equal-degree labels.
The factorization $s=W_2W_{2'}$ gives the row-support assertion.

If these two phases are global, set $Z=T_2^2$, $F_0=sZs^{-1}$,
$P=(I+Z)/2$. The level is eight and $H_5=-I$; the dyadic
Bruhat identity gives $ZF_0=-F_0Z$. Write
$T_2=aI+cZ$ with $|a|=|c|=1/\sqrt2$. The Weyl factorization
then shows that $PW_2P$ is scalar on $P$, of modulus
$1/\sqrt2$. The complementary vacuum vector lies in $P$.
The indicator at a CRT exponent zero at two and one on the odd
level gives an integral self-adjoint virtual element $b$ with
$b^2=2b$, unit coefficient one, dimension two and squared norm
two. Its sole nonunit coefficient is therefore $+1$ on an
invertible simple, contradicting perfectness.
\end{proof}

\begin{lemma}[At least four exceptional phases]
\label{paper:ef:dyadic-mass}
Let $D=32n^2$, $n$ odd. Suppose every label whose dyadic phase
is outside $\{1,i\}$ has degree $n$, and at least one exists.
There are at least four such labels.
\end{lemma}
\begin{proof}
If the maximal dyadic twist order is $2^s\ge8$, let $L$ count
labels at that order. The exponent selecting the corresponding
highest-order involution gives the Gauss sum $n^2(32-2L)$.
The regular-object indicator is $n^2(32-2L)^2/32$.
Since $n$ is odd, indicator integrality implies
$(32-2L)^2/32\in\mathbb Z$, so $4\mid L$.
Otherwise let $u,v$ count the exceptional labels at $-1,-i$.
The Bruhat mass equality puts mass $D/2$ at orders one and two
and mass $D/2$ at order four. The dyadic Gauss sum is therefore
$n^2((16-2u)+i(16-2v))$. Its indicator gives
$8\mid u^2+v^2$, hence $u,v$ are even and their halves have
the same parity. Thus $4\mid u+v$, which is positive.
\end{proof}

\begin{proposition}[Ordinary exclusion at dyadic valuation five]
\label{paper:ef:ordinary-five}
There is no ordinary perfect fusion-simple table of dimension
$2^5q^br^c$ for odd primes $q<r$ and arbitrary $b,c$.
\end{proposition}
\begin{proof}
Every pair of opposite principal witnesses at $q,r$ has common
dyadic part at least $q$ and at most four. Thus each witness has
dyadic height two. Its defect is one. By
Lemma~\ref{paper:ef:ordinary-dyadic-phases}, all witness rows see
$1$ and a common primitive fourth root: opposite principal
witnesses hit the equal-degree active representatives of each
other's blocks. Prime avoidance makes every $q$- or $r$-active
column hit a witness. Perfectness and the same lemma give an
occupied dyadic phase outside the common pair. On any such
fiber the complementary local groups act scalarly on each line,
so its labels have pure dyadic codegree. Consequently $b,c$ are
even; write $n=q^{b/2}r^{c/2}$.

An exceptional label has degree $2^j n$ and codegree
$2^{5-2j}$ for $j=0,1,2$. Its row is supported on $P_2$.
Reducing the trace-dimension identity modulo $q$ and $r$ gives
two distinct nonunit trace orbits of degrees $2q^h$ and $2r^l$,
by the principal height bound. This already rules out codegree
two. At codegree eight, the exceptional label and both witnesses
have dyadic height one; symmetry and principality make each
nonzero trace even. Their squared traces cost at least $4+4$
besides the vacuum, impossible. Thus every exceptional label
has degree $n$ and codegree $32$.

Take the $r$-free hit $U$ of degree $2q^h$ and a complementary
principal $r$-witness of degree $4r^l$. Their normalized crossing
is a nonscalar unsigned $r$-word of length two, so the
$r$-inertia orbit of $U$ has size $e\ge(r-1)/2\ge2$.
It hits every exceptional degree-$n$ label by two-principality.
These entries are fixed by $r$-inertia, and their LCM scale is
$2n$. Averaging the row norm and retaining the vacuum gives
$4eN<32$, where $N$ is the number of exceptional labels.
Thus $N\le3$, contrary to Lemma~\ref{paper:ef:dyadic-mass}.
\end{proof}

\subsubsection{The reduced dyadic quintic}

\begin{lemma}[A positive singleton and its odd mate]
\label{paper:ef:singleton-mate}
A positive singleton two-block in an integral fermion extension
lifts to a full block consisting of two even mates and a free
odd pair of the same degree. Its row sees only the odd pair's
actual dyadic phase on odd columns.
For opposite complementary singleton witnesses at odd primes
$q,r$, all their odd mates have one common dyadic phase; every
odd $q$- or $r$-active column has that phase. Their positive
rows see only squared dyadic phase one.
\end{lemma}
\begin{proof}
The positive and signed orders have identical special fibers,
so the signed block has rank one. Hence there is one odd orbit.
Block orthogonality of the full columns $1,g$ equates its
degree-square mass with the even mass $2d_X^2$. A fixed orbit
would have degree $\sqrt2d_X$, impossible. The odd pair is
free of degree $d_X$. Local block invariance and complementary
phase preservation give its row's asserted odd support.
The same factorization on the positive space, whose vacuum
entry is nonzero, gives squared phase one on its positive row.

An opposite odd mate has degree prime to the witness's odd
prime. Choosing the witness's full principal lift makes it hit
that mate. Thus all opposite witnesses have the same visible
odd phase. The principal block $P_q$ contains only even labels,
since an odd label takes value $-1$ on $g$, whereas the vacuum
takes value $1$. Moreover the opposite odd mate $T$ has
$q\nmid d_T$, and $\chi_g(T)=-d_T\not\equiv d_T$
above $q$; hence $g\notin P_q$. The same argument with $q,r$
interchanged gives $g\notin P_r$. Prime avoidance modulo the other
odd prime therefore makes any odd active column hit a witness,
giving the last support assertion.
\end{proof}

\begin{lemma}[A single odd spin is impossible]
\label{paper:ef:one-odd-spin}
In an integral minimal extension, the even dyadic twists cannot
all be signs while all odd dyadic twists have one value.
\end{lemma}
\begin{proof}
Use $A,B,C,W,U$ from the fermion dictionary, and put $J=W|_A$.
On $A,B$, $U^2=\lambda^2$; on $C$, $U$ is a scalar $c$.
Restricting $WUW=U^{-1}WU^{-1}$ to $UA$ gives
$cUJ^2=\lambda^{-2}UJ$, so $J$ is scalar.
The complementary vacuum column is supported on the even
representatives $E$ with actual dyadic twist one and has
squared weights $4d_j^2/N$. The indicator construction gives an
integral self-adjoint element $b$ with $b^2=4b$, dimension four,
and even part $1+g$. Hence $v=b-1-g$ has odd support and
\[
             gv=v,\qquad v^2=2(1+g),\qquad \dim v=2,
             \qquad \|v\|^2=2.
\]
A free fermion pair in its support would give an odd invertible.
Otherwise positivity of dimensions forces $v=Y-Z$ for two
fermion-fixed self-dual odd simples. Now $vY$ has norm four
and coefficients one at $1,g$; its remaining part is an even
free pair $A'+gA'$, with positive sign by dimension. Thus
$d_Y=d_{A'}+1$. Reciprocity and $\langle vA',v\rangle=0$
give $vA'=Y+Z$; the norm of $vA'$ is two, so there are no
other terms. Hence $d_Z=d_{A'}-1$. Coprime-degree vanishing
makes $\chi_{A'}(Y)=\chi_{A'}(Z)=0$, contrary to
$\chi_{A'}(v)^2=4$.
\end{proof}

\begin{lemma}[Sign twists exclude a principal defect-two label]
\label{paper:ef:sign-fs-square}
Let an integral fermion extension have dimension $N=64n^2$,
with $n$ odd. Suppose all even actual dyadic twists are signs.
Its positive principal two-block has no defect-two character.
\end{lemma}
\begin{proof}
Dyadic inertia fixes the even labels individually: balancing
makes their restricted values unramified at two, and opposite
twists distinguish the representatives. Since $N$ is a square,
the local torus sign is positive. Thus $W^2=I$ on $A,B,C$.
Put $J=W|_A$ and $P=WU:A\to C$. Restricting the Weyl relation
gives $UP=\lambda^{-2}PJ$. Since $J^2=I$, the odd dyadic
spins lie in a pair $\{\delta,-\delta\}$.

Choose a CRT exponent $k$ zero on the dyadic level and two
on the odd level. In the modular indicator formula for an even
$V$, the two fusion inputs are in the same fermion sector;
their dyadic twist ratio is a sign. Hence
$\nu_k(V)=\nu_2(V)$. For
$b=\sum_Y\nu_k(Y)[Y^*]$, the indicator diagonalization is
\[
                \chi_x(b)=f_x|(sT^ks^{-1})_{0x}|^2.
\]
Odd-local Bruhat at parameter $-2$ identifies the absolute
entries with a torus permutation of the odd-Weyl vacuum column.
Because $N$ is a square that column is rational, and it has
even support of actual dyadic twist one. Therefore $gb=b$;
for an even pair at most one of its two eigenvalues is nonzero,
and each nonzero eigenvalue is a rational square.

In the positive quotient set
$W^+=\sum_{[V]}\nu_2(V)[V^*]$. Its eigenvalues are
\[
                w'_x=(\chi_x(b)+\chi_{gx}(b))/4.
\]
They are rational squares and algebraic integers, hence ordinary
integer squares. Modulo two, $(W^+)^2$ is the Casimir
$\sum_{[V]}[V][V^*]$: nonselfdual pairs cancel, and a selfdual
positive pair has a selfdual representative by opposite twists.
Thus $4\mid w'_x$ whenever its positive codegree is even.

Suppose $x$ is principal of defect two. Sign-only balancing
makes its positive field unramified at two, so
$v(V)=(\chi_x(V)-d_V)/2$ is locally integral. Orthogonality
gives
\[
 \sum_{[V]}v(V)v(V^*)=(f_x^++\Delta)/4\equiv1\pmod2,
 \qquad \Delta=16n^2.
\]
Cancel nonselfdual pairs. Frobenius in the residue field implies
$\sum_{[V]}\nu_2(V)v(V^*)\not\equiv0\pmod2$.
But this sum is $(w'_x-w'_0)/2$, divisible by two since both
eigenvalues are integer squares divisible by four. This is a
contradiction.
\end{proof}

\begin{proposition}[The reduced quintic face]
\label{paper:ef:reduced-five}
There is no minimal slightly degenerate obstruction of dimension
$2^5q^br^c$ for odd primes $q<r$ and arbitrary $b,c$.
\end{proposition}
\begin{proof}
Initially the positive quotient has dimension $\Delta=16q^br^c$.
Opposite principal witnesses have common dyadic part at least
$q$ and at most four; each therefore has height two and
positive odd codegree. They are singleton two-blocks.
Their positive rows see only squared dyadic phase one, and
their odd rows see one common phase $\delta$, by
Lemma~\ref{paper:ef:singleton-mate}. If one companion valuation were
odd, every occupied complementary dyadic fiber in each fermion
sector would contain an active label at that companion: its
nonsquare local torus acts by minus a permutation. Thus all even
dyadic twists would be signs and all odd spins would equal
$\delta$, contrary to Lemma~\ref{paper:ef:one-odd-spin}. Hence $b,c$
are even. Put $n=q^{b/2}r^{c/2}$, so the integral extension has
dimension $64n^2$.
By Lemma~\ref{paper:ef:singleton-mate}, any positive pair whose
squared dyadic twist is nontrivial is inactive at both $q,r$.
The complementary groups fix its full lines, so its positive
codegree is purely dyadic and its pair is two-principal.
The possibilities are $16,4,1$. Codegree one exhausts the row
norm; codegree four is excluded by the two distinct nonzero
even traces forced modulo $q,r$, each costing at least four.
Thus every such deep pair has degree $n$ and codegree $16$.

If there are $N_0$ deep pairs, dimension orthogonality on a
deep row, reduced modulo $r$, gives an $r$-free nonunit hit.
Pure-codegree support and the principal height bound give its
degree $2q^i$. Its crossing
with a complementary $4r^j$ witness gives an inertia orbit of
length $e\ge2$. The orbit norm gives $4eN_0<16$, so there is
exactly one deep pair $Z$. Its dyadic twist power acts on the
positive space as $R=I+(u-1)P_Z$, $u\ne1$. The trace of the
finite-group commutator $RKR^{-1}K^{-1}$ makes
\[
                    (2-u-u^{-1})|K_{ZZ}|^2
\]
an algebraic integer. Yet $|K_{ZZ}|^2=|\chi_Z(Z)|^2/16$
has dyadic valuation $-4$ by principality, whereas
$v_2(2-u-u^{-1})\le2$. This contradiction eliminates deep pairs.

All even dyadic twists are signs. Lemma~\ref{paper:ef:sign-fs-square}
and the principal height bound now make every positive
two-principal degree odd. All complementary odd mates and all
odd $q$- or $r$-active columns have one common phase $\delta$.
Lemma~\ref{paper:ef:one-odd-spin} gives another odd phase. Its fiber
is inactive at both odd primes, so any label $Y$ on it has
pure dyadic full codegree.

Take an odd mate $T$ of an $r$-free witness. Its degree is
$r$-free and its phase is $\delta$. The signed restrictions
of $T,Y$ are distinct and orthogonal. Reducing their
orthogonality modulo $r$, the vacuum forces a nonunit $r$-free
even pair $V$ hit by $Y$: the $r$-bearing terms vanish on $T$
by LCM integrality. Pure-codegree support places $V$ in the
positive principal two-block. Its degree is odd, hence a power
of $q$, impossible in the perfect fusion-simple positive table.
\end{proof}

\subsubsection{Uniform ternary lemmas for the second face}

\begin{lemma}[Weight three with residues at two and three]
\label{paper:ef:weight-three}
In an ordinary minimal table on dimension primes $2,3,r$, let
an LCM-normalized crossing have common weight three, residue one
at every place above two, and residue zero at every place above
three. Then it has the form $\lambda\sqrt{-3}$ for a dyadic
root of unity $\lambda$.
\end{lemma}
\begin{proof}
The two-odd-prime root configuration is unimodular; adjoining a
dyadic basis gives disjoint signed copies of it. Consequently
the crossing has an integral signed word of length at most three.
Erase its three-parts. The remaining word $w$ is divisible by
three in a three-unramified field and has house at most three.
If $w=0$, its terms are a cancelling pair, since the complementary
alphabet has no vanishing triangle. Thus the original word is
$\lambda(\alpha-\beta)$ for three-power roots $\alpha,\beta$.
If $w\ne0$, the algebraic norm and equality in the triangle
inequality give $\lambda(\alpha+\beta+\gamma)$.
In both cases inversion of the $r$-part, followed by reduction
above two, shows that $\lambda$ has trivial $r$-part: reduction
is faithful on odd-order roots, and both residues are one.
Thus $\lambda$ is dyadic.

Reduce the three-primary factors modulo two. In the pair case,
the three terms $\alpha,\beta,1$ form the identity
$\mu_3$-coset, so $\{\alpha,\beta\}=\{\omega,\omega^{-1}\}$.
In the triple case the parity vector of
$\alpha,\beta,\gamma,1$ is a union of complete triples with
even cardinality at most four; it is therefore zero. The word
is $\lambda(1+2\xi)$ for a three-power root $\xi\ne1$.
If $\xi$ has order at least nine, its adjoint support
$\{1,\xi,\xi^{-1}\}$ has one point in each occupied
$\mu_3$-coset and is exposed. Its dilation by the three-adic
unit four meets it only at one. The common nonunit constituent
of the two adjoints would then have a scalar character value of
full degree, contradicting strictness. Thus $\xi=\omega^{\pm1}$.
Both alternatives equal a dyadic scalar times $\sqrt{-3}$.
\end{proof}

\begin{lemma}[Transfer from a two-principal ternary height-one row]
\label{paper:ef:ternary-transfer}
An ordinary minimal table on primes $2,3,r$ with odd ternary
valuation has no two-principal row of degree $2^i3$.
\end{lemma}
\begin{proof}
Suppose $U$ is such a row. Let $\iota$ invert the three-power
roots, fixing the other primary factors, and put
$H=\chi_U+\chi_{U^\iota}$. These rows have the same ternary
twist. An odd-degree nonunit $B\in P_3$ has degree $3^h r^j$,
with $h,j>0$. Its normalized crossing with $U$ satisfies
Lemma~\ref{paper:ef:weight-three}; it is nonzero and negated by
$\iota$. Thus $H(B)=0$, and $B^\iota\ne B$.
On every even degree, both summands of $H$ vanish above two.

If an odd-degree $J$ had a different ternary twist, principal
block orthogonality against $H$ would have vacuum term two.
Every remaining nonzero term has dyadic valuation greater than
one: its $H$ factor has positive valuation, and its $\chi_J$
factor is divisible by two by LCM integrality. This is impossible.
The vacuum therefore implies that all odd-degree labels are
ternary inactive. The odd members of $P_3$ number one modulo
two, since all but the vacuum pair under $\iota$. At odd ternary
valuation the principal mass identity is
\[
 \sum_{B\in P_3,\,\operatorname{ord}(\theta_{B,3})=3}d_B^2
       =2\sum_{B\in P_3,\,\theta_{B,3}=1}d_B^2.
\]
The left side is zero modulo four and the right side is two,
a contradiction.
\end{proof}

\begin{lemma}[A ternary phase outside a prescribed pair]
\label{paper:ef:ternary-missing-phase}
In an ordinary minimal table of odd ternary valuation, the
ternary phases are not contained in $\{1,\eta\}$ for any
primitive cube root $\eta$.
\end{lemma}
\begin{proof}
Suppose they were. Let $P$ project onto the ternary inactive
coordinates and $w=W_{3'}e_0$. The nonsquare ternary torus
acts on the vacuum by its Legendre sign. The two nonzero
Bruhat identities are
\[
 PW_3U_3^tW_3^{-1}w
       =\left(\frac{-t}{3}\right)Pse_0\quad(t=1,2).
\]
The missing Fourier coefficient gives $w=-G Pse_0$ with
$|G|^2=3$. The indicator at a CRT exponent zero at three and
one on its complement therefore produces
\[
 b=3\sum_{\theta_{X,3}=1}e_X\in R,\qquad
 b=b^*,\quad b^2=3b,\quad b_1=1,\quad\dim b=3,
 \quad\|b\|^2=3.
\]
Here the coefficients are algebraic integers by the indicator
formula and rational by Galois stability of the inactive set.

Set $M=N_b$. It is symmetric integral with $M^2=3M$.
A row with diagonal one or two has exactly two off-diagonal
entries, each $\pm1$; rows with diagonal zero or three are
isolated. The unit component is a cycle. A cycle longer than
four has nonneighbors with a unique length-two path, contrary
to $M^2=3M$. A triangle has trace three, since the unit diagonal
is one and the trace is a multiple of three. It has rank one;
its positive dimension eigenvector forces all three dimensions
to equal one, contradicting perfectness.
On a four-cycle the vertex $A$ opposite the vacuum has diagonal
one and its two neighbors $Y,Z$ have diagonal two. Their
dimension equations give $d_Y,d_Z\in\{d_A-1,d_A+1\}$.
Coprime-degree vanishing gives $\chi_A(Y)=\chi_A(Z)=0$.
But $b=1\pm Y\pm Z$, so $\chi_A(b)=1$, whereas its only
possible eigenvalues are zero and three. This is impossible.
\end{proof}

\begin{proposition}[Exceptional phases at a ternary quintic]
\label{paper:ef:ternary-pure}
In an ordinary minimal table of dimension $2^a3^5r^4$, $r>3$,
the exponent $a$ is even. There is a primitive cube root $\eta$
such that a nonempty set of labels has ternary phase outside
$\{1,\eta\}$, and every such label has codegree $27$ or $243$.
\end{proposition}
\begin{proof}
Prime avoidance provides $r$-free nonunits in $P_2$.
Their ternary height is two by Lemma~\ref{paper:ef:ternary-transfer}
and square divisibility. The odd nonunits of $P_r$ likewise
have degree $9r$: height one would give a common weight three
and residues three at both three and $r$, contrary to the
two-residue theorem. These two sets of witnesses have ternary
defect one.

At an odd defect-one prime, the two local inertia orbits have
one inactive and one primitive quadratic phase class. Their
lengths and degrees satisfy
$a_1d_1^2=(3-1)a_0d_0^2$, where $a_0,a_1$ divide two.
In particular their $r$-valuations agree. An $r$-principal
witness hits the active representatives of every $r$-free
witness block. Their visible primitive cube roots therefore
agree, say at $\eta$. Every two-active or $r$-active column
hits a witness by prime avoidance. All their ternary phases
lie in $\{1,\eta\}$.

Lemma~\ref{paper:ef:ternary-missing-phase} supplies a phase outside
this pair. Its fiber is inactive at two and $r$, so the local
groups preserve its coordinate lines. The codegrees on that
fiber are pure powers of three; hence $a$ is even and the
possibilities are $243,27,3$. Codegree three is impossible:
reducing row orthogonality modulo each of the three dimension
primes forces three distinct nonunit hits, since no nonunit
degree is a prime power. Their Galois mean squares, together
with the vacuum, exceed three.
\end{proof}

\begin{lemma}[The two trace orbits at pure codegree twenty-seven]
\label{paper:ef:ternary-trace}
In the setting of Proposition~\ref{paper:ef:ternary-pure}, suppose
$f_Z=27$. There is a unique $r$-free nonunit Galois orbit
contributing nonzero absolute trace on this row. Its degree is
$d_U=2^i3$, its row field equals $K_Z$ and is purely ternary,
and
\[
                         9\cdot2^i\equiv\pm1\pmod r.
\]
There is also a unique odd nonunit trace-contributing orbit,
of degree $3r^j$ for $j=1$ or $2$, with the same row field.
\end{lemma}
\begin{proof}
Pure-codegree support puts every hit in $P_3$. The trace-dimension
identity modulo $r$ and modulo two forces respectively an
$r$-free and an odd nonunit trace orbit. Their ternary heights
are one by the principal bound and the exclusion of prime-power
degrees. Since $d_Z=2^{a/2}3r^2$, symmetry and principality
make their character entries vanish above three. Each nonzero
absolute trace is therefore a multiple of three.
In the orbit trace norm an entry with field $F$ contributes
\[
 [K_U:F][K_Z:F]\operatorname{Tr}_{F/\mathbb Q}(\chi_Z(U))^2.
\]
Each required orbit costs at least nine. A third such orbit,
or either relative field index exceeding one, would cost at
least $27$ besides the positive vacuum contribution. Thus the
two traces are $\pm3$, their fields equal both row fields,
and the required orbits are unique. Reducing the dimension
trace identity modulo $r$ gives $1\pm3d_U=0$, as claimed.
\end{proof}

\begin{lemma}[Mass at the maximal exceptional ternary codegree]
\label{paper:ef:ternary-mass}
Let an integral modular table have dimension $243n^2$, with
$3\nmid n$.  If every label outside $\{1,\eta\}$, for a
primitive cube root $\eta$, has degree $n$ and the outside set
is nonempty, that set has at least nine labels.
\end{lemma}
\begin{proof}
Apply Corollary~\ref{paper:fam:Q:6.6} with $p=3$, $k=2$ and all
exceptional degrees equal to $n$.  It gives
$2N\ge(p-1)p^k=18$.
\end{proof}

\begin{theorem}[The dyadic sextic, middle quintic face]
\label{paper:ef:sextic-quintic-face}
There is no ordinary minimal table of dimension $2^6q^5r^4$
for odd primes $q<r$. The minimal slightly degenerate branch
at this dimension is also empty.
\end{theorem}
\begin{proof}
The reduced assertion is Theorem~\ref{paper:fermion-six}.
In the ordinary table, a dyadically active fiber contains a
$q$-active label: otherwise the $q$-local group is trivial on
that invariant fiber, contrary to the nonsquare torus acting
there by minus a permutation at odd $q$-valuation.
Prime avoidance on $P_q$ modulo $r$ now gives a nonunit
$r$-free row meeting a joint two--$q$ column. Its degree is
$2^i q$, with $i\le2$ by defect-zero vanishing on the chosen
dyadic activity. Opposite principal witnesses force
$2^i\ge q$. Hence $q=3$, $i=2$.

Call this row $U$, so $d_U=4\cdot3$. Every three-free nonunit
of $P_r$ has degree $4r$ or $8r$. Their crossings with any
such height-two row have common weight four and residues four
above three and $r$. If $r>5$, the word is $1+3\xi$ and the
two-position obstruction applies. Thus the uniform necessary
conditions include $3<r\le5$. At the remaining allowable
prime the two-residue theorem, including its signed alternative,
makes the tame $r$-orbit of the entry have length $r-1$.
In particular every degree-$4\cdot3$ row in $P_3$ is
$r$-ramified, and the orbit of $U$ has size $e\ge r-1$.

Apply Proposition~\ref{paper:ef:ternary-pure}. A codegree-$27$ row
would require the $r$-unramified trace witness of
Lemma~\ref{paper:ef:ternary-trace}, of degree $2^i3$.
Complementary principal weights force $i\ge2$, while the
dyadic sextic cap gives $i\le3$. Height two is excluded by the
ramification just proved. Height three would require
$r\mid (9\cdot2^3)^2-1$, which has no prime divisor in
$3<r\le5$. Hence all exceptional codegrees are $243$.

They all have degree $n=2^3r^2$, and
Lemma~\ref{paper:ef:ternary-mass} gives at least nine of them.
The row $U\in P_3$ hits each, and all those entries are fixed
by its $r$-inertia orbit. Their LCM scale is $3n$.
The orbit-averaged norm, with the vacuum retained, gives
$9eN_0<243$, equivalently $eN_0<27$. But
$eN_0\ge9(r-1)\ge36$, a contradiction.
\end{proof}

\paragraph{Closure of the required even proper faces.}
Theorem~\ref{paper:ef:sextic-quintic-face} eliminates the exact
$(6,5,4)$ even face in both categorical branches.
Propositions~\ref{paper:ef:ordinary-five} and~\ref{paper:ef:reduced-five}
eliminate the exact $(5,6,4)$ face. The already-proved cubic,
smaller-quartic, and fewer-prime reductions remove their proper
divisors. This supplies precisely the even lower-face input
needed for the $(6,6,4)$ family.

\begin{proposition}[The required even divisor faces]
\label{paper:family:even-faces}
Let $q<r$ be odd primes. Every weakly integral braided fusion
category whose dimension divides $2^5q^6r^4$ or $2^6q^5r^4$
is weakly group-theoretical.
\end{proposition}
\begin{proof}
Choose a counterexample of least dimension among divisors of
these two numbers. The divisor collection is closed under taking
divisors, so the generic minimal-counterexample reduction applies.
The fewer-prime, cubic-exponent and
smaller-quartic theorems leave three repeated primes, quartic
valuation at $r$, and dyadic valuation five or six.
At dyadic valuation five,
Propositions~\ref{paper:ef:ordinary-five} and~\ref{paper:ef:reduced-five}
exclude both categorical branches, independently of the two
companion exponents. At dyadic valuation six the second divisor
box is the only possibility; the middle valuation is exactly
five by the smaller-quartic theorem. Both branches are excluded
by Theorem~\ref{paper:ef:sextic-quintic-face}. Thus no minimal
counterexample exists.
\end{proof}


\section{Reduction to the ordered exact family}\label{paper:sec:exact-family}\label{paper:family:consequences}

\subsection{Proof of the family reduction}\label{paper:sec:exact-corner}

\begin{theorem}[Ordered exact reduction]\label{paper:family:family:exact-corner}
Let $p,q,r$ be distinct primes. If a braided fusion category of dimension
dividing $p^6q^6r^4$ is not weakly group-theoretical, then $r$ is the
largest prime. After interchanging $p$ and $q$, a divisor-minimal
counterexample is an ordinary integral perfect fusion-simple modular
category of exact dimension $p^6q^6r^4$, with $p<q<r$.
\end{theorem}
\begin{proof}
Choose a divisor-minimal counterexample. The minimal-counterexample
lemma leaves an ordinary table or the positive quotient of a slightly
degenerate minimal category. Dong's theorem excludes dimensions with
at most two repeated prime factors \cite{DongWeakIntegral}.
The small-exponent result therefore forces all three valuations to be
at least four. The quartic exclusion makes the prime assigned exponent
four the largest prime, with actual valuation four. The two other
valuations are five or six.

For an odd dimension the table is ordinary. The least quintic exclusion
first removes valuation five at $p$; the middle quintic exclusion then
removes it at $q$. In even dimension the two corresponding divisor
faces are WGT. Thus both sextic exponents are attained. The remaining
slightly degenerate possibility is excluded by the fermion theorem at
two-valuation six. Only the stated ordinary exact dimension remains.
\end{proof}
can therefore be applied directly to this ordered corner.
The exponent bookkeeping is particularly small:

\begin{center}
\begin{tabular}{@{}lccc@{}}
\toprule
Prime & Dimension valuation & Maximum simple height
      & Maximum principal height\\
\midrule
\(p\) & \(6\) & \(3\) & \(2\)\\
\(q\) & \(6\) & \(3\) & \(2\)\\
\(r\) & \(4\) & \(2\) & \(1\)\\
\bottomrule
\end{tabular}
\end{center}

The last column follows from the exclusion of principal defects
zero and one.  It is stronger than square divisibility alone
and is responsible for the compulsory complementary crossings.

\subsection{The ordinary dyadic sextic}

\label{paper:sec:dyadic-ordinary}

\begin{theorem}[Two cannot be isolated at sextic valuation]
\label{paper:family:family:dyadic-isolation}
Let an ordinary minimal table have exactly three dimension primes
\(2<q<r\), with \(v_2(D)=6\).  Some actual twist is divisible
by \(2q\) or \(2r\).
\end{theorem}
\begin{proof}
Suppose two is isolated.  By categorical Cauchy and
Lemma~\ref{paper:family:lw:weyl-line}, the dimension is
\(D=64n^2\), with \(n\) odd.  Every dyadically active label
has pure dyadic codegree. Theorem~\ref{paper:dyadic-sixteen}
therefore makes every active degree $n$, of codegree $64$.
The isolated-column formula gives an integral vector
\[
 u=8W_2v_0,\qquad \|u\|^2=64,\qquad
 v_2(u_X)=v_2(d_X)\quad\text{when }v_2(d_X)<3,
 \tag{D1}
\]
and puts all these lower-height labels in \(P_2\).
An irrelevant scalar sign on \(W_2\) has no effect here.
Separation also makes any field ramified at an odd prime
unramified at two.

An \(r\)-free label \(X\) of degree \(2q^j\) has
\(r\)-inertia orbit length \(t\ge(r-1)/2\ge2\).
Indeed, pair it with a \(q\)-free nonunit in \(P_r\),
of degree \(2^h r^l\).  The normalized crossing has
mass two, strict house below two, and residue two at \(r\).
Since \(r>4\), Corollary~\ref{paper:cor:one-large-residue} gives two
distinct \(r\)-power roots.  Their tame stabilizer has order
at most two.  This also proves that \(K_X\) ramifies at \(r\).

Let $N$ count the active labels. Choose the greatest dyadic
twist order among them, and let $L$ be the number at that order.
An exponent can turn these phases into $-1$ and every smaller
phase into $1$, while killing the odd primary factors. The
scaled Gauss sum is $64-2L$. The regular-object indicator
makes $(64-2L)^2/64$ integral, so $4\mid L$ and $N\ge4$.
Prime avoidance gives an \(r\)-free even \(X\in P_2\)
of degree \(2^i q^j\), where \(i=1\) or \(2\).
It hits all \(N\) active labels.  Let \(t\) be the length
of its \(r\)-inertia orbit.  The active entry fields have
dyadic conductor, so the averaged-row energy bound gives
\[
                              4^i tN<64.                    \tag{D2}
\]
This first forces \(i=1\), and then the orbit estimate gives
\(t\ge2\) and \(N\le7\).
The vacuum Gauss formula and \textup{(D1)} now yield
\[
 |u_0|=\frac18\left|64-\sum_{Z\ {\rm active}}
                       (1-\theta_{Z,2})\right|
       \ge\frac{64-2N}{8}>6.
\]
Thus \(u_0^2\ge49\).  The \(N\) odd active coordinates cost
at least four in \(\|u\|^2\); the \(t\ge2\) even
\(r\)-free coordinates cost at least eight; and an even
\(q\)-free member of \(P_2\) costs at least four.
These disjoint contributions give \(64\ge49+4+8+4\),
the final contradiction.
\end{proof}

\begin{corollary}[A crossing of common height two at two]
\label{paper:family:family:dyadic-crossing}
In an ordinary minimal table of dimension \(2^6q^6r^4\),
some opposite \(r\)-free member of \(P_q\) and \(q\)-free
member of \(P_r\) have common dyadic part four.  Necessarily \(q=3\).
\end{corollary}
\begin{proof}
Complementarity and square divisibility make the common
dyadic part four or eight.  If every opposite pair had
part eight, every witness on both sides would have
odd codegree.  Every \(q\)-active column hits an
\(r\)-free witness in \(P_q\), by prime avoidance, and
would therefore be dyadically inactive by defect-zero
vanishing.  The same argument at \(r\) isolates two,
contrary to the theorem.  The complementary-degree bound
(Lemma~\ref{paper:family:4.10}) gives $q\le4$, hence $q=3$.
\end{proof}

\begin{corollary}[The unconditional even-family exclusion]
\label{paper:family:dyadic-exclusion}
In the ordered exact family with $p=2$, every member with $q>3$
or $r>5$ is weakly group-theoretical. Thus a minimal even
obstruction can have only the exceptional assignment $q=3,r=5$.
\end{corollary}
\begin{proof}
Corollary~\ref{paper:family:family:dyadic-crossing} forces $q=3$ and a
compulsory crossing whose common weight is four. For $r\ge7$,
Theorem~\ref{paper:family:thm:two-positive} and strictness give its
actual word $z=1+3\xi$, where $\xi$ is a nonidentity
$r$-primary root. The exceptional signed branch would require
$r=5$. Since
\[
                         \chi_Y(X)=\frac{d_X}{4}z,
\]
this is a full positive two-position character word, of mass
$d_X$. The arbitrary-parity two-position case of
Theorem~\ref{paper:family:short-support} excludes it at arbitrary
primary conductor. This argument uses no homogeneity,
isolated fiber, flatness or input inactivity.
\end{proof}

The same even-family exclusion is obtained in \cite{BDPfamily}.
Thus its only residual even assignment is $(p,q,r)=(2,3,5)$,
of dimension $5400^2$.

\subsection{Compulsory crossings and their heights}

\label{paper:sec:crossing-blocks}

\begin{proposition}[The compulsory complementary crossings]
\label{paper:family:family:crossings}
In an ordinary minimal table of dimension \(p^6q^6r^4\),
with \(p<q<r\), there are nonunits
\[
 X\in P_q,\quad d_X=p^i q^j,\qquad
 Y\in P_r,\quad d_Y=p^h r,
 \quad i,h\in\{2,3\},\ j\in\{1,2\}.
\]
They may be chosen with \(q\)-primary and \(r\)-primary
twists, respectively.  Every opposite pairing has an
integral normalized entry
\[
 z=\frac{S_{XY}}{\lcm(d_X,d_Y)},\qquad
 g=p^{\min(i,h)}\in\{p^2,p^3\},
\]
of positive rational mass \(g\), strict house below \(g\),
and residue \(g\) at every place over both \(q\) and \(r\).
\end{proposition}
\begin{proof}
Use prime avoidance on the principal blocks, restricted to
the trivial complementary twist fibers, to obtain the
prime-omitting nonunits.  The principal height bound gives
\(j\le2\) and \(v_r(d_Y)=1\).  The exclusion of prime-power
degrees makes the remaining heights positive.
Lemma~\ref{paper:family:4.10} gives
\(p^{\min(i,h)}\ge q>p\), so both heights are at least two.
Square divisibility bounds them by three.
LCM integrality, balancing and strictness give the assertions
about \(z\).  Principality gives the two residues, and
argument-side Galois covariance gives them at every place.
\end{proof}

\begin{proposition}[Residue descent at a top-height crossing]
\label{paper:unramified-crossing}
Use the complementary labels \(X,Y\), heights \(i,h\), and
normalized crossing \(z\) of Proposition~\ref{paper:family:family:crossings}.
If \(\max(i,h)=3\), the field \(\mathbb Q(z)\) is unramified
at \(p\).  More precisely,
\begin{enumerate}
\item if \(\{i,h\}=\{2,3\}\), then \(z\) is a unit at some
place over \(p\);
\item if \(i=h=3\), then
\(\min_{\mathfrak p\mid p}v_{\mathfrak p}(z)\le1\), with
\(v_{\mathfrak p}(p)=1\).
\end{enumerate}
These assertions hold also for \(p=2\).
\end{proposition}
\begin{proof}
Let \(T\) be a label of height three among \(X,Y\), and let
\(V\) be the other label.  Its codegree \(f_T=D/d_T^2\)
is prime to \(p\).  The conductor theorem
(Theorem~\ref{paper:family:por:th:exact-conductor}) therefore makes
\(K_T=\mathbb Q(\chi_T(A):A\in\Irr(\C))\) unramified at \(p\).
Since
\[
 z=\frac{d_T}{\lcm(d_X,d_Y)}\chi_T(V),
\]
the same is true of \(\mathbb Q(z)\).  Its normalized
\(p\)-valuations are thus integers.

Put \(k=\min(i,h)\).  Suppose that every valuation of \(z\)
over \(p\) were at least \(k-1\).  Then \(z/p^{k-1}\)
would be an algebraic integer.  Dividing its actual balancing
word gives positive rational mass \(p\), and its residues at
every place over \(q,r\) are both \(p\).  Since \(p<q<r\),
Theorem~\ref{paper:family:thm:two-positive} forces
\(z/p^{k-1}=p\).  This contradicts the strict house bound
\(\lvert z\rvert_{\rm house}<p^k\).  Consequently some
integer \(p\)-valuation is at most \(k-2\), proving both
claims.
\end{proof}

At a mixed-height unit place, the two character values are
asymmetric: the height-three character has a unit value on
the height-two label, whereas the reverse value has valuation
one.  Thus this descent gives the unit crossing without
making the lower-height block homogeneous.  Homogeneity still requires both heights to be two; see Proposition~\ref{paper:family:homogeneous-crossing}.

\begin{corollary}[Two mixed neighbors give height-two witnesses]
\label{paper:two-neighbors}
If an actual twist is divisible by \(pq\), the complementary
witness \(X\in P_q\) may be chosen of \(p\)-height two.
If an actual twist is divisible by \(pr\), the witness
\(Y\in P_r\) may be chosen of \(p\)-height two.
In particular, if both kinds of mixed twist occur, an
opposite pair with both heights two exists.
\end{corollary}
\begin{proof}
For a \(p,q\)-active column, prime avoidance on the
trivial complementary-twist fiber of \(P_q\), at the
omitted prime \(r\), gives an \(r\)-free nonunit
whose row meets that column.  Its \(p\)-height is at
least two by the complementary-degree bound.  Height
three would give \(p\)-unit codegree and, by
Lemma~\ref{paper:family:lw:weyl-line}, a zero on every \(p\)-active
column.  The height is therefore two.  Interchange
\(q,r\) for the second assertion and choose the two
witnesses independently.
\end{proof}


\section{Support and crossing constraints in odd dimension}

\label{paper:sec:family-criterion}

The categorical statements in this section concern an ordinary
perfect fusion-simple integral modular category of exact dimension
$p^6q^6r^4$, with $p<q<r$ odd. They use the principal-block
orthogonality, prime avoidance, LCM integrality, strictness,
codegree conductor bound, and local Weyl facts proved in the
structural sections. In particular: separated activity at two
primes forbids simultaneous ramification of a character field;
an active label has codegree divisible by that prime; a label
of prime-power codegree is supported on the corresponding
principal block; and an isolated quartic prime in three-prime
ordinary data is impossible. No classification at an individual
dimension is used here.

\subsection{The height-three trace budget}

\begin{proposition}[Isolated least sextic]
\label{paper:family:prop:isolated-sextic}
Suppose $p$ is isolated in the exact odd family. Set
\[
 e_s=\frac{s-1}{\max\{a:a\mid s-1,\ a\le p\}},
 \qquad s=q,r.
\]
Then
\begin{equation}\label{paper:family:eq:sextic-budget}
                  (p-1)(e_q+e_r)<p^2.
\end{equation}
In particular an isolated least sextic is impossible
when $r>p^2$.
\end{proposition}
\begin{proof}
Write $D=p^6n^2$, where here $n=q^3r^2$.
The isolated-prime local structure and
Proposition~\ref{paper:primepower-trace} give every active label
$Z$ degree $p^j n$, $j\in\{0,1\}$, and codegree $p^{6-2j}$;
its field is $p$-primary and its row is supported on $P_p$.

First suppose some active label has $j=1$. The same trace
proposition supplies, for each $s=q,r$, an $s$-free nonunit
$X$ of degree $pu^i$, where $u$ is the other companion,
with nonzero trace $T_X$ divisible by $p$. Hence
\begin{equation}\label{paper:family:eq:trace-cost}
                   d_Z^2T_X^2\ge p^4n^2.
\end{equation}

Cross $X$ with a $u$-free nonunit of $P_s$.
The normalized entry has weight $p$ and residue $p$
at $s$. Since $s\ge p+2$ and
$(s-p)(p-1)>p$, Lemma~\ref{paper:lem:one-residue}
makes its actual word a unique positive integral
$s$-word of weight $p$. Strictness gives at least
one nonidentity letter. The tame stabilizer permutes
such letters freely and has order $a\mid s-1$,
$a\le p$. Thus $[K_X:\mathbb Q]\ge e_s$ and
$K_X$ ramifies at $s$, because $s-1>p$.
Separation from $p$ makes $K_X$ unramified at $p$,
so $F_X\subseteq K_X\cap K_Z=\mathbb Q$.
The square Galois action on the active twist gives
at least $(p-1)/2$ distinct twist values in the
$Z$-orbit. In odd dimension no nonunit is self-dual,
and duals have the same twist, so
$[K_Z:\mathbb Q]\ge p-1$.

The two chosen $X$-orbits are disjoint, since one
is $q$-free and the other $r$-free, and nonunits
cannot have prime-power degree. In the trace norm,
\eqref{paper:family:eq:trace-cost} therefore makes their
total cost at least
$p^4n^2(p-1)(e_q+e_r)$. The unit has positive cost.
Comparison with $D=p^6n^2$ gives
\eqref{paper:family:eq:sextic-budget}.

It remains that every active degree equals $n$.
Let $N$ be the number of active labels. The local
vacuum identity is
\[
  p^{-3}\sum_{Z\ \mathrm{active}}(1-\theta_{Z,p})
            =p^3(1-(W_p)_{00})\in\mathbb Z.
\]
Its real part is positive. If this integer is $K$,
pairing distinct equal-twist duals gives $p^3K\in2\mathcal O$.
Since $K$ is rational integral and $p$ is odd, $K$ is even.
The normalized trace of each $1-\theta_{Z,p}$ is
at most $p/(p-1)$, so $N\ge2p^2(p-1)$.

Every $X\in P_p$ hits all these active labels by
principality, since their degree $n$ is a $p$-unit.
A height-two such row would have nonzero integral
entries $S_{XZ}/(p^2n)$. Their Galois mean squared
moduli are at least one by the norm inequality.
The row norm, including the unit, would imply
$Np^4n^2<D$, or $N<p^2$, a contradiction.
For each $s$, prime avoidance gives an $s$-free
nonunit of $P_p$, necessarily of height one.
The previous crossing argument gives
$[K_X:\mathbb Q]\ge e_s$ and makes $K_X$
unramified at $p$. Its orbit therefore agrees on
all active columns, with nonzero integer entries
$S_{XZ}/(pn)$. Averaging that orbit of rows gives
\[
              Np^2 e_s<p^6,
 \qquad (p-1)e_s<p^2/2.
\]
Adding for the two companions again proves
\eqref{paper:family:eq:sextic-budget}.

Finally, if $r>p^2$, the integer $e_r$ exceeds
$p$ and hence is at least $p+1$. Since $e_q\ge1$,
the left side of \eqref{paper:family:eq:sextic-budget} is
at least $(p-1)(p+2)>p^2$ for odd $p\ge3$.
\end{proof}

\begin{proposition}[A half-channel budget for an isolated sextic row]
\label{paper:isolated-half-budget}
In the ordered odd exact family, suppose $p$ is isolated and
some $p$-active label $Z$ has degree $pM$, where $M=q^3r^2$.
With $e_q,e_r$ as in Proposition~\ref{paper:family:prop:isolated-sextic},
one has
\[
 (p-1)\bigl(1+p^2(e_q+e_r)\bigr)
       \le \frac{p^2(p^2-1)}2,
 \qquad e_q+e_r<\frac{p+1}{2}.
\]
In particular this conclusion holds if a label of $p$-height
two exists.
\end{proposition}
\begin{proof}
Retain the exact orbit size $e=[K_Z:\mathbb Q]$ in the
rational-witness trace norm in the isolated-sextic proof.
The two witness orbits then give
\[
                   e\,p^2(e_q+e_r)<p^4.
\]
By Proposition~\ref{paper:primepower-field-budget}, the active
$p$-primary field $K_Z$ is a full cyclotomic field. Since
$e_q+e_r\ge2$, its degree is less than $p^2/2$, excluding
conductor $p^2$ or higher. Thus $K_Z=\mathbb Q(\zeta_p)$;
active inertia also makes the actual twist have order $p$.
Let $u_Z$ be the normalized
sum of its $p-1$ Galois labels. Their twists lie in one
quadratic class $C$ of the nontrivial prime-level phases.
Every complementary local group fixes this whole orbit
coordinatewise, because its labels have complementary unit
codegrees. Including the possible dyadic anomaly scalar,
$s u_Z=\varepsilon W_pu_Z$ with $|\varepsilon|=1$.

Let $P_B$ project onto the vectors fixed by both the full
$p$-torus and the upper $p$-unipotent. The adjoint form of
Lemma~\ref{paper:half-channel-local} gives
\[
 \|P_B W_pu_Z\|^2\le\frac{p^2-1}{2p^2}.
\]
The two nonzero-trace witness orbits in the isolated-sextic
proof have degrees $pq^i$ and $pr^j$, respectively. Their
full fields are unramified at $p$, and their complete twists
are one: a foreign active label would have $p$-height three,
whereas a $p$-active label has both companion prime factors
to their full heights. Consequently each coordinate vector
in the two witness orbits is Borel-fixed at $p$.
The orbits have sizes at least $e_r,e_q$, respectively, and
their values against $Z$ are nonzero rational integers
divisible by $p$. These values are constant across the
$Z$-orbit and across each witness orbit. Since
$s_{XZ}=\chi_Z(X)/p^2$, their combined contribution to
$\|P_B W_pu_Z\|^2$ is at least
\[
                  \frac{(p-1)(e_q+e_r)}{p^2}.
\]
The vacuum is an additional orthogonal Borel-fixed coordinate
and contributes $(p-1)/p^4$. Combining these contributions
proves the first inequality, and dropping its strictly
positive vacuum correction proves the second.
Finally a height-two label forces every active degree to be
$pM$ by Lemma~\ref{paper:isolated-height-two}, while categorical
Cauchy ensures that an active label exists.
\end{proof}

\begin{remark}
The result strengthens the scalar isolated trace budget but
does not exclude all isolated least sextic primes. In
particular, it does not apply by itself to the alternative
where every active label has degree $M$ and no height-two
label occurs. Even in the degree-$pM$ case, the resulting
rough bound $r<p^2/2$ does not contradict the
uniform support inequalities for all $p$.
\end{remark}

\begin{theorem}[The uniform exact-family support criterion]
\label{paper:family:thm:uniform-family}
Let $p<q<r$ be odd primes and put
\[
 m=\left\lfloor\frac{\min(r,p^2)}q\right\rfloor,
 \quad E=m(m+1)/2,\quad H=(r-1)/2.
\]
If a full positive $r$-primary character word on at most $m+1$
positions is impossible at a nonunit row of an ordinary perfect
fusion-simple integral modular category of this dimension,
then every braided fusion category of dimension $p^6q^6r^4$
is weakly group-theoretical. In particular, if $m=0$,
or if
\begin{equation}\label{paper:family:eq:family-test}
          (H-1)u(h)>h(h-1)\qquad(1\le h\le E),
\end{equation}
there is no ordinary perfect fusion-simple integral
modular category of dimension $p^6q^6r^4$.
\end{theorem}
\begin{proof}
The case $r<p^2$ supplies a seed on at most $m+1$
positions by Proposition~\ref{paper:family:prop:quartic-seed}.
The assumed word exclusion gives a contradiction. For the
displayed scalar test, exposure is automatic by
Proposition~\ref{paper:family:prop:auto-exposure}, and
Theorem~\ref{paper:family:thm:descent} gives that exclusion.

If $r>p^2$, pair any $r$-free nonunit of $P_q$
with any $q$-free nonunit of $P_r$. Their degrees
are $p^i q^j,p^k r$, respectively, where
$i,k\in\{2,3\}$; the lower bound follows from
Theorem~\ref{paper:family:thm:two-positive} if their common
weight were below $q$. A weight-$p^2$ crossing
has the two positive residues and, since $r>p^2$,
a full primary word on at most $m+1$ positions.
The exceptional case in
Theorem~\ref{paper:family:thm:two-positive} is impossible
because $p^2+1$ is even and exceeds two.
If $m=0$ its word has one position and contradicts
strictness; otherwise the assumed word exclusion applies.

All these complementary witnesses must consequently
have $p$-height three. On a $q$-active column,
prime avoidance on $P_q$ modulo $r$ gives an
$r$-free hit of this type; its $p$-unit codegree
forces the column to be $p$-inactive. The same
argument with $P_r$ separates $p$ from $r$.
Thus $p$ is isolated, contradicting
Proposition~\ref{paper:family:prop:isolated-sextic}.
\end{proof}

The proof is uniform in the three primes. For the scalar
sufficient test, the remaining numerical possibility is the
existence of some $1\le h\le E$ satisfying
\eqref{paper:family:eq:remaining}. The criterion does not prove that this
possibility is empty for all triples.

\begin{corollary}[A uniform bound on the remaining region]
\label{paper:family:cor:remaining-region}\label{paper:family:uniform-prime-range}
The support criterion excludes the exact odd family
when $m\le2$, and when $r>\max\{33,m^3\}$.
Every triple not excluded by this criterion satisfies
\[
            q^3<p^4,\qquad
            q^{3/2}<r<(p^2/q)^3.
\]
\end{corollary}
\begin{proof}
For $m=1,2,3$, direct substitution gives
$C(E)=1,4,16$, respectively. The first two cases
have $r>3,9$ from $r>mq$ and $q>p\ge3$.
For $m\ge4$, if $u(h)\le m-2$ use
$h\le u(h)(u(h)+1)$; otherwise use
$h-1\le E-1\le(m-1)^2$. Both give
$(h-1)/u(h)\le m-1$. Therefore
\[
 2C(E)+1\le3+2E(m-1)
             =m^3-m+3\le m^3.
\]
This proves the sufficient condition.
If it fails and $r>33$, then $r\le m^3$.
The strict inequalities $m<r/q$ and $m<p^2/q$
give the displayed lower and upper bounds on $r$.

Suppose instead that $r\le33$, and put $g=p^2$.
Since $m\ge3$, both $r$ and $g$ are at least $3q$.
Equality $r=3q$ is impossible, so $r>3q$ and $q<11$.
As $q$ is odd, $q\le9$, whence $r>3q\ge q^{3/2}$.
Also $g>3q$; these are odd integers, so $g\ge3q+2$ and
\[
                (g/q)^3\ge(3+2/9)^3>33\ge r.
\]
Thus the bounds hold throughout the surviving range. Their
conjunction gives $q^3<p^4$.
\end{proof}

\begin{corollary}[Uniform family coverage from short support]
\label{paper:family:family-criterion}
Let $p<q<r$ be odd primes and put
\[
            k=1+\left\lfloor\frac{\min(r,p^2)}q\right\rfloor.
\]
Every braided fusion category of dimension $p^6q^6r^4$ is
weakly group-theoretical if $k=1$, or if $2\le k\le7$ and
$r\ge T_k$ in Theorem~\ref{paper:family:short-support}.
In particular the following complete parameter ranges are covered:
\[
\begin{array}{c|c}
 \text{prime conditions}&\text{covered range}\\ \hline
 7\le p<q<r&r<5q\\
 7\le p<q<r,\ q\ge17&r<6q\\
 11\le p<q<r,\ q\ge19&r<7q.
\end{array}
\]
The full ordered sector with least prime $p=7$ is also covered.
The even-family conclusion is Corollary~\ref{paper:family:dyadic-exclusion}.

\end{corollary}
\begin{proof}
Apply the word-to-family transfer of
Theorem~\ref{paper:family:thm:uniform-family}.
Theorem~\ref{paper:family:short-support} supplies the required word
exclusion, and $k=1$ contradicts strictness directly.
For the three displayed ranges, $k\le5,6,7$, respectively.
When each largest value occurs, the inequalities
$r>4q$, $r>5q$ or $r>6q$ give its stated threshold; in the
last case primality gives $r\ge127$ from $q\ge19$.
The smaller values of $k$ satisfy their thresholds as well.

For $p=7$, one has $q\ge11$ and $m=k-1\le4$.
If $m=4$, then $r>4q\ge44$, so primality gives $r\ge47$.
If $m=3$, then $r>3q\ge33$, exceeding the four-position
threshold. The cases $m\le2$ already follow from the scalar
test. This proves the full least-prime-seven assertion.
\end{proof}

\subsection{Floor bounds and multiplier quotients}

\begin{corollary}[A cubic bound and a uniform floor test]
\label{paper:family:cor:floor-test}
For the exact odd family, the support criterion holds whenever
\[
                  r\ge\frac{(m+1)^3}{\sqrt2}.
\]
For fixed $p<q$, put $J=\lfloor p^2/q\rfloor$.
Every prime $r>q$ is excluded if
\[
 J\le2\qquad\text{or}\qquad (J+1)^3\le\sqrt2\,Jq.
\]
Two convenient sufficient conditions for exclusion are
\[
 (J+2)^2\le\sqrt2\,q,
 \qquad (p^2/q+1)^3\le\sqrt2\,p^2.
\]
In particular $\sqrt2\,q^3\ge(p^2+2q)^2$ suffices.
For any fixed $\alpha>1/\sqrt2$, the condition
$q^3\ge\alpha p^4$ suffices whenever
$p^{-2/3}\le2^{1/6}-\alpha^{-1/3}$.
\end{corollary}
\begin{proof}
The case $m=0$ follows from strictness. For $m\ge1$
and $1\le h\le E=m(m+1)/2$,
\[
 \frac{h}{u(h)}\le\sqrt{h+1/4}+1/2
                 \le\frac{m+1}{\sqrt2}+\frac12.
\]
Consequently $3+2h(h-1)/u(h)<(m+1)^3/\sqrt2$.
Indeed, after replacing $2(h-1)$ by $m(m+1)-2$,
the difference between the claimed upper bound and
the resulting expression is
\[
 \frac{(m+1)(m+3)}{\sqrt2}-\frac{m^2+m}{2}-2>0;
\]
replacing $1/\sqrt2$ by $1/2$ already bounds this
difference below by $(3m-1)/2>0$.
Thus $r\ge(m+1)^3/\sqrt2$ gives $H>C(E)$.

For the floor statement, $m\le J$ and $r>mq$.
The cases $m\le2$ were proved in
Corollary~\ref{paper:family:cor:remaining-region}.
For $m\ge3$, monotonicity of $(x+1)^3/x$ on $x\ge1$
gives
$(m+1)^3\le\sqrt2\,mq<\sqrt2\,r$.
The two sufficient tests follow from
$(J+1)^3/J<(J+2)^2$ and the same monotonicity at
$x=p^2/q$. The first test follows in turn from
$\sqrt2\,q^3\ge(p^2+2q)^2$.
For the last assertion,
\[
 p^{-2/3}(p^2/q+1)
       \le\alpha^{-1/3}+p^{-2/3}\le2^{1/6},
\]
which is the second sufficient test.
\end{proof}

\begin{proposition}[Necessary conditions in a multiplier quotient]
\label{paper:family:prop:multiplier-quotient}
Suppose a seed survives in the exact odd family. With
$m,E,H$ as in Theorem~\ref{paper:family:thm:uniform-family},
some $1\le h\le E$ has
\[
 k=u(h),\qquad \delta=h(h-1)-(H-1)k\ge0.
\]
Put $v=v(h)=\min\{j\ge1:j(j+1)\ge2h\}$.
If $\delta<v(2h-v)-(h-k)$, then for every divisor $b$ of $H$
there are integers $0\le x_i\le H/b$, indexed by
$i\in\mathbb Z/b$, with $\sum_i x_i=h$, and
nonnegative integers $e_j=e_{-j}$, with $\sum_j e_j=\delta$,
satisfying
\begin{equation}\label{paper:family:eq:quotient}
 \sum_i x_i x_{i+j}
       =\frac Hb k+(h-k)\mathbf1_{j=0}+e_j.
\end{equation}
For $b=H$ one additionally has $e_0=0$ and $x_i\in\{0,1\}$.
If $H$ is even, in particular,
$[h-k-\delta,h-k+\delta]$ contains the square of an
integer congruent to $h$ modulo two.
\end{proposition}
\begin{proof}
Theorem~\ref{paper:family:intrinsic-conductor} puts every minimizing
mask $A$ in the cyclic group
$G=\mathbb F_r^\times/\{1,-1\}$ of size $H$.
Every nonidentity overlap is at least $k$ by the common-nonunit
argument. Its excess over $k$ is nonnegative, invariant under
inversion, and has total $\delta$; the identity excess is zero.

Let $x_i$ count the elements of $A$ in class $i$
modulo $b$. The left side of \eqref{paper:family:eq:quotient}
counts ordered pairs with prescribed difference
modulo $b$. Each of its $H/b$ full-group
differences occurs at least $k$ times, with the
zero difference occurring $h$ times. The total
excess is exactly $\delta$, proving the formula. Each quotient
fiber has $H/b$ elements, giving the capacity bound. Inversion
gives $e_j=e_{-j}$, and at $b=H$ the identity excess is zero.
For $b=2$, subtract the odd equation from the
even one to obtain
$(x_0-x_1)^2=h-k+e_0-e_1$; its square root has
the parity of $x_0+x_1=h$.
\end{proof}

\begin{lemma}[An outside-point containment test]
\label{paper:outside-containment}
Let $A,C$ be subsets of a finite abelian group $G$, with
$|G|=H$, $|A|=h$, $|C|=h+t$, $t\ge0$, and $n=h-k>0$.
Suppose all nonidentity autocorrelations of both sets, and all
mixed intersections $|C\cap gA|$, are at least $k\ge1$. Write
\[
 |A\cap gA|=k+n\mathbf1_{g=1}+e_g,\qquad
 |C\cap gC|=k+(n+t)\mathbf1_{g=1}+e'_g,
 \qquad \delta=\sum_g e_g.
\]
The excesses are nonnegative and $e_1=e'_1=0$. For
$y_g=|C\cap gA|-k$, one has
\begin{equation}\label{paper:eq:unequal-moments}
 \sum_g y_g=n+\delta+ht,\qquad
 \sum_g y_g^2=n(n+t)+kt^2+\sum_g e_ge'_g.
\end{equation}
If $C$ contains no translate of $A$, then
\begin{equation}\label{paper:eq:outside-cap}
                         0\le y_g\le kt+\delta.
\end{equation}
Consequently $C$ contains a translate of $A$ whenever
\begin{equation}\label{paper:eq:containment-test}
 n(n+t)+kt^2>(kt+\delta)(n+\delta+ht).
\end{equation}
\end{lemma}
\begin{proof}
Put $c_g=|C\cap gA|$. Double counting gives
\[
 \sum_g c_g=h(h+t),\qquad
 \sum_g c_g^2=h(h+t)+(H-1)k^2
                   +k\left(\delta+\sum_g e'_g\right)+\sum_g e_ge'_g.
\]
Here $\delta=h(h-1)-(H-1)k$ and
$\sum_g e'_g=(h+t)(h+t-1)-(H-1)k$. Subtracting the
baseline $k$ proves \eqref{paper:eq:unequal-moments}.
For each point $x\in G$, a second incidence count gives
\begin{equation}\label{paper:eq:outside-point}
 \sum_{g:\,x\in gA}y_g
       =n\mathbf1_C(x)+kt+\sum_{c\in C}e_{x/c}.
\end{equation}
Indeed, the number of translates of $A$ containing both $x,c$
is $|A\cap(x/c)A|$. If no translate of $A$ lies in $C$, choose
$x\in gA\setminus C$ for each $g$. Positivity in
\eqref{paper:eq:outside-point} bounds $y_g$ by $kt+\delta$.
Now \eqref{paper:eq:outside-cap} bounds the second moment by
$(kt+\delta)\sum_g y_g$, contradicting
\eqref{paper:eq:unequal-moments} under the stated strict inequality.
\end{proof}

\begin{lemma}[Prime level and rigidity in a linear slack range]
\label{paper:prime-level-rigidity}
Use the minimizing full primary word at the fixed nonunit row in
Theorem~\ref{paper:family:thm:descent}, at the family support bound. Put
\[
 k=u(h),\qquad n=h-k,\qquad
 \delta=h(h-1)-(H-1)k.
\]
If $0\le\delta<B(h)$, with $B(h)$ as in
Theorem~\ref{paper:family:intrinsic-conductor}, every minimizing difference
mask is an $h$-element subset of
$G=\mathbb F_r^\times/\{1,-1\}$.
For two such masks $A,A'$, write
\[
 |A\cap gA|=k+n\mathbf1_{g=1}+e_g,\qquad
 |A'\cap gA'|=k+n\mathbf1_{g=1}+e'_g.
\]
Both excess sums equal $\delta$, and
\begin{equation}\label{paper:eq:cross-moments}
 y_g:=|A'\cap gA|-k\ge0,\qquad
 \sum_g y_g=n+\delta,\qquad
 \sum_g y_g^2=n^2+\sum_g e_ge'_g.
\end{equation}
All minimizing masks are translates if
\begin{equation}\label{paper:eq:rigidity-range}
                              3\delta<2n.
\end{equation}
This last inequality itself implies the prime-level hypothesis.
\end{lemma}
\begin{proof}
Exposure is automatic by Proposition~\ref{paper:family:prop:auto-exposure}.
The common-constituent argument applies to any two minimizing
adjoints after any primary-unit dilation. An intersection with
$b<k$ opposite pairs would produce a nonunit word with at most
$b(b+1)<h$ differences, contradicting minimality. Thus all mixed
intersections and nonidentity autocorrelations are at least $k$.
Theorem~\ref{paper:family:intrinsic-conductor} supplies prime level, and
Lemma~\ref{paper:outside-containment} with $t=0$ gives
\eqref{paper:eq:cross-moments}.

If $A,A'$ are not translates, neither contains a translate of the
other, so the same lemma gives $0\le y_g\le\delta$. For
$\delta\le n/2$, this implies
$\sum_g y_g^2\le\delta(n+\delta)\le3n^2/4$.
For $n/2<\delta<2n/3$, concentration under this cap gives
\[
 \sum_g y_g^2\le2\delta^2+(n-\delta)^2
                  =n^2+\delta(3\delta-2n)<n^2.
\]
Both contradict \eqref{paper:eq:cross-moments}. Finally,
$3\delta<2n$ implies $\delta<2(h-1)<B(h)$ for $h\ge2$;
the case $h=1$ does not satisfy the scalar condition here.
\end{proof}

\begin{theorem}[Critical minima and a superlinear slack bound]
\label{paper:positive-slack}\label{paper:triangular-slack}
In the minimizing-mask conditions of
Proposition~\ref{paper:family:prop:multiplier-quotient}, suppose
$h=k(k+1)$, $k\ge2$. If $3\delta<2k^2$, put
\[
 \begin{split}
 H&=k^3+2k^2-\delta/k,\qquad f_* = k+\lfloor\delta/2\rfloor,\\
 N_*&=\left\lfloor
       2+\frac{(k-2)f_*(f_*-1)+\delta f_*}{k(k-1)}
                                      \right\rfloor.
 \end{split}
\]
Then necessarily
\begin{equation}\label{paper:eq:critical-incidence-test}
                         H-1-\delta\le N_*(N_*-1).
\end{equation}
In particular,
\begin{equation}\label{paper:eq:critical-slack-bounds}
 \begin{array}{ll}
 \delta>0 &(k\ge2),\\
 \delta\ge2k &(k\ge3),\\
 \delta\ge3k &(k\ge14),\\
 \delta>k^{5/4} &(k\ge16).
 \end{array}
\end{equation}
The last inequality also gives
$\delta\ge k(\lfloor k^{1/4}\rfloor+1)$.
\end{theorem}
\begin{proof}
The slack is a multiple of $k$, and $n=h-k=k^2$.
Under $3\delta<2k^2$, Lemma~\ref{paper:prime-level-rigidity}
makes all minimizing masks translates of one $A\subset G$.
Its stabilizer is trivial: a nonidentity stabilizer would
contribute excess $k^2>\delta$. Thus there are $H$ distinct
translates. Let $\nu$ be the class of $2$, and put
\[
                       F=A\cap\nu^{-1}A,\qquad f=|F|.
\]
Here $r>5$, so $\nu\ne1$ and $\nu\ne\nu^{-1}$.
Autocorrelation symmetry gives
\begin{equation}\label{paper:eq:doubling-size}
                         k\le f\le k+\lfloor\delta/2\rfloor.
\end{equation}
Call an ordered pair of distinct translates good if its
intersection $B$ has exactly $k$ pairs. A common nonunit word
has phases in $\{0\}\cup\{\pm b:b\in B\}$. The differences
of this envelope lie among
\[
                         b,\quad 2b,\quad b\pm b',
\]
and have at most $k(k+1)=h$ opposite pairs. Minimality forces
its entire difference envelope to equal a minimizing mask,
hence a translate $tA$. In particular it contains $B,2B$, so
\begin{equation}\label{paper:eq:intersection-cover}
                                  B\subseteq tF.
\end{equation}
The multiplier $2$ is realized by a global cyclotomic unit with
trivial complementary components, by the Chinese remainder theorem.

Let $N$ count translates of $A$ meeting $F$ in at least $k$
points, and put $c_t=|tA\cap F|$. An ordered-pair count gives
\[
 \sum_t c_t(c_t-1)
  =k f(f-1)+\sum_{g\ne1}e_g|F\cap gF|
  \le k f(f-1)+\delta f.
\]
The two distinct translates $A,\nu^{-1}A$ contain all of $F$.
Therefore
\begin{equation}\label{paper:eq:incidence-bound}
 2f(f-1)+(N-2)k(k-1)\le k f(f-1)+\delta f.
\end{equation}
The resulting bound on $N$ increases with $f$, giving $N\le N_*$.
At most $\delta$ nonidentity translates overlap $A$ in more
than $k$ points. There are consequently at least
$H(H-1-\delta)$ good ordered pairs. By
\eqref{paper:eq:intersection-cover}, the $H$ translates of $F$
cover all such pairs, with at most $N(N-1)$ per translate.
Repetitions only increase this upper bound. Hence
\begin{equation}\label{paper:eq:pair-cover}
                              H-1-\delta\le N(N-1),
\end{equation}
proving \eqref{paper:eq:critical-incidence-test}.

For $\delta=0$, one has $f=k$ and $N\le k$, whereas
$H-1=h(h-1)/k>k(k-1)$. For $\delta=k$, $k\ge4$,
\eqref{paper:eq:incidence-bound} gives $N\le9k/4$, and
\[
 16\left(H-1-k-\frac{9k}{4}\left(\frac{9k}{4}-1\right)\right)
       =k^2(16k-49)+20k-32>0.
\]
For $k=3$, $f\le4$ gives $N\le6$, but
\eqref{paper:eq:pair-cover} would require $40\le30$.
These arguments recover the first two bounds.

If $\delta=2k$ and $k\ge14$, rigidity holds and $f\le2k$.
The incidence bound gives
$N\le\lfloor4k+2/(k-1)\rfloor=4k$. But
\[
 H-1-2k-4k(4k-1)=k^3-14k^2+2k-3>0,
\]
contradicting \eqref{paper:eq:pair-cover}. This proves the third bound.

Finally suppose $k\ge16$ and $\delta\le k^{5/4}$. Rigidity
again holds. Set $a=1+\delta/(2k)$ and $x=k^{1/4}\ge2$.
Then $f\le ka$, and \eqref{paper:eq:incidence-bound} implies
\[
 N\le2+\frac{k^2a^2-(3k-2)a}{k-1}\le(k+1)a^2.
\]
For the last inequality, $1\le a\le k$ gives
$2+[a^2-(3k-2)a]/(k-1)\le2-2a\le0$.
Since $a\le1+x/2\le x$, one has $N\le(k+1)\sqrt{k}$.
Consequently
\[
 \begin{split}
 H-1-\delta-N(N-1)
 &\ge(k+1)\sqrt{k}-(k+1)-\delta(1+1/k)\\
 &\ge(k+1)(\sqrt{k}-1-k^{1/4})>0,
 \end{split}
\]
again impossible. Divisibility of $\delta$ by $k$ gives the
last integer formulation.
\end{proof}

\begin{theorem}[The minimum one below a critical endpoint]
\label{paper:near-critical-minimum}
Under the same minimizing-word hypotheses, the configuration
\[
                    h=k(k+1)-1,\qquad\delta=2
\]
is impossible for every $k\ge4$.
\end{theorem}
\begin{proof}
Here $n=k^2-1$ and
\[
 H=k^3+2k^2-2k-2,\qquad r=2k^3+4k^2-4k-3.
\]
Lemma~\ref{paper:prime-level-rigidity} makes every minimizing
mask a translate of a fixed $A$, with trivial stabilizer.
Use the good pairs and $F=A\cap2^{-1}A$ from the preceding
proof; now $k\le f:=|F|\le k+1$.
For a good intersection $B$ of $k$ pairs, let $T$ be the
opposite difference envelope of $\{0\}\cup\{\pm b:b\in B\}$.
It contains the actual difference mask $C$ of a common
nonunit word, with
\[
                         h\le|C|\le|T|\le h+1.
\]
These supports are exposed: the sufficient odd-dimensional
bound $3(2(h+1)+1)<2r$ follows because the difference is
$4k^3+2k^2-14k-9>0$ for $k\ge4$. Their autocorrelations
and mixed intersections with $A$ are therefore at least $k$,
by the common-constituent argument and minimality of $h$.

If $|C|=h$, it is a translate of $A$. If $|C|=h+1$,
Lemma~\ref{paper:outside-containment} also forces $C$ to contain
a translate of $A$: its lower second moment and capped upper
second moment differ by
\[
 (k^4-k^2+k)-(k+2)(2k^2+k)
 =k(k^3-2k^2-6k-1)>0,
\]
since the cubic equals $(k-4)(k^2+2k+2)+7$.
Thus $T$ contains a translate $tA$ in either case.

At least $k-1$ points of $B$ have both $b,2b$ in $tA$.
Indeed, if $|T|=h$ then $T=tA$. If $|T|=h+1$, all
$h+1$ formal entries $b,2b,b\pm b'$ are distinct modulo
sign, so $B\cap2B$ is empty. The single point of
$T\setminus tA$ spoils at most one pair $b,2b$.
Consequently $|B\cap tF|\ge k-1$.

Let $N$ count translates of $A$ meeting $F$ in at least
$k-1$ points. The same ordered-pair count, with this new
threshold, gives
\[
 2f(f-1)+(N-2)(k-1)(k-2)\le kf(f-1)+2f.
\]
Using $f\le k+1$ yields
\[
 N\le k+4+\frac{4k-2}{(k-1)(k-2)}<k+7,
 \qquad N\le k+6.
\]
The translates of $F$ cover all good ordered pairs, so
$H-3\le N(N-1)$. For $k\ge5$ this contradicts
\[
 H-3-(k+6)(k+5)=k^3+k^2-13k-35
                 =(k-5)(k^2+6k+17)+50>0.
\]
For $k=4$, the order-two quotient in
Proposition~\ref{paper:family:prop:multiplier-quotient} already gives
a contradiction: $H=86$, $h=19$, and $[n-\delta,n+\delta]
=[13,17]$ contains no odd square.
\end{proof}


\begin{proposition}[The cyclotomic norm restriction at equality]
\label{paper:family:prop:multiplier-equality}
Under the hypotheses of
Proposition~\ref{paper:family:prop:multiplier-quotient}, suppose $\delta=0$.
For every $b>1$ dividing $H$ and every prime $\ell\nmid b$
for which $-1$ is a power of $\ell$ modulo $b$, the valuation
$v_\ell(h-k)$ is even. Moreover $h<k(k+1)$ if $k\ge2$.
\end{proposition}
\begin{proof}
The zero-slack equations make the mask a cyclic $(H,h,k)$
difference set. For a character $\chi$ of order $b$, put
$\alpha=\sum_{a\in A}\chi(a)$. Then
\[
                         \alpha\overline\alpha=h-k.
\]
The prime $\ell$ is unramified in $\mathbb Q(\zeta_b)$, and
its decomposition group contains complex conjugation by the
hypothesis on $-1$. Hence $\alpha$ and $\overline\alpha$
have equal valuations at every prime above $\ell$, proving
the even-valuation assertion. The strict bound on $h$ is
the zero-slack case of Theorem~\ref{paper:triangular-slack}.
\end{proof}


\begin{corollary}[The combined symbolic family criterion]
\label{paper:family:cor:enhanced-criterion}
For odd $p<q<r$, put $m,E,H$ as in
Theorem~\ref{paper:family:thm:uniform-family}. The exact family is WGT
if the short-support criterion of
Corollary~\ref{paper:family:family-criterion} applies, or if no integer
$1\le h\le E$ satisfies all the following necessary conditions:
\begin{enumerate}[leftmargin=*]
\item With $k=u(h)$, the slack
$\delta=h(h-1)-(H-1)k$ is nonnegative.
\item Put $v=v(h)=\min\{j\ge1:j(j+1)\ge2h\}$.
If $\delta<v(2h-v)-(h-k)$, every minimizing mask is at prime
level and all the divisor-quotient conditions, capacities and
symmetries in Proposition~\ref{paper:family:prop:multiplier-quotient} hold.
\item At $\delta=0$, the cyclotomic valuations satisfy
Proposition~\ref{paper:family:prop:multiplier-equality}.
\item If $h=k(k+1)$ with $k\ge2$, then all the slack bounds in
\eqref{paper:eq:critical-slack-bounds} hold. If moreover
$3\delta<2k^2$, the explicit incidence condition
\eqref{paper:eq:critical-incidence-test} holds.
\item If $h=k(k+1)-1$ and $k\ge4$, then $\delta\ne2$.
\end{enumerate}
These are unconditional tests on the primes. They require no
homogeneous block, isolated inactive fiber or input inactivity.
\end{corollary}
\begin{proof}
A seed at the family support bound has exposed adjoints and
therefore a minimizing mask satisfying all these conditions.
Its exclusion implies the result through the common transfer
of Theorem~\ref{paper:family:thm:uniform-family}.
\end{proof}


\begin{corollary}[An unconditional critical-window criterion]
\label{paper:family:symbolic-exclusion}\label{paper:family:terminal-family}
Let $p<q<r$ be odd primes, with $m,E,H$ as in
Theorem~\ref{paper:family:thm:uniform-family}. Suppose integers $k,j$
satisfy
\begin{equation}\label{paper:family:eq:symbolic-window}
 r=2k^3+4k^2+1-2j,\qquad
 k(k+1)\le E,\qquad E(E-1)<(H-1)(k+1),
\end{equation}
and one of the following conditions holds:
\[
 \begin{array}{ll}
 j=1,&k\ge3,\\
 j=2,&k\ge14,\\
 1\le j\le k^{1/4},&k\ge16.
 \end{array}
\]
Then every braided fusion category of dimension $p^6q^6r^4$
is weakly group-theoretical.
\end{corollary}
\begin{proof}
Put $h_0=k(k+1)$. We show that $h_0$ is the only value
$1\le h\le E$ satisfying the scalar condition
\eqref{paper:family:eq:remaining}. The stated ranges give $j\le k$.
If $a=u(h)\le k-1$, then
\[
 \frac{h(h-1)}a\le(a+1)(a(a+1)-1)
                 \le k^3-k^2-k<H-1.
\]
If $u(h)=k$ and $h<h_0$, its slack is at most
$jk-2(h_0-1)<0$. If $u(h)\ge k+1$, the last inequality
in \eqref{paper:family:eq:symbolic-window} excludes it. At $h=h_0$,
the slack is exactly $jk$, which
Theorem~\ref{paper:triangular-slack} excludes in each stated
range. The common transfer of
Theorem~\ref{paper:family:thm:uniform-family} gives the result.
\end{proof}


\subsection{The first odd arithmetic candidate}

\begin{corollary}[The first odd frontier in the family]
\label{paper:family:first-odd}
Every braided fusion category of odd dimension $p^6q^6r^4$ strictly below
\[
 D_{\mathrm{odd}}=5^6\,7^6\,23^4
   =514\,422\,090\,765\,625=(22\,680\,875)^2
\]
is weakly group-theoretical, for every assignment of the quartic prime.
The triple $(5,7,23)$ is the smallest ordered odd parameter triple not
excluded by the uniform scalar, short-support and combined symbolic
criteria above. No category at this dimension is asserted to exist.
\end{corollary}
\begin{proof}
The exact reduction permits $p<q<r$. For $p=3$,
$m\le\lfloor9/5\rfloor=1$, so the scalar criterion applies.
For $p=5$ and $q\ge11$, one has $m\le2$. If $p=5,q=7$,
every prime $r<23$ again has $m\le2$, while every prime $r\ge29$
has $m=3$ and meets the four-position threshold. Thus the only possible
parameter in the least-prime-five sector is $(5,7,23)$.
Every ordered triple with $p\ge7$ has dimension at least
$7^6\,11^6\,13^4=5\,952\,751\,597\,721\,929>D_{\mathrm{odd}}$.
This proves the strict dimension bound and its minimality.

At $(5,7,23)$, the parameters are $m=3,E=6,H=11$. The scalar
condition allows $h=5$, $k=u(h)=2$, with $\delta=0$.
In $\mathbb Z/11$, the set $A=\{1,3,4,5,9\}$ satisfies
\[
 |A\cap(A+t)|=\begin{cases}5,&t=0,\\2,&t\ne0.\end{cases}
\]
It therefore satisfies every divisor-quotient, capacity and symmetry
condition. The equality norm is $h-k=3$; the subgroup generated by
$3$ modulo $11$ has order five and does not contain $-1$, so the
cyclotomic valuation test does not reject it. Also $h<k(k+1)$, and
the near-critical restriction requires $k\ge4$. Finally, $23<29$
does not meet the four-position threshold. Hence this parameter passes
the combined symbolic necessary tests. The accompanying exact audit
also enumerates all forty-one ordered odd triples up to
$D_{\mathrm{odd}}$ and verifies scalar exclusion of the forty smaller
ones.
\end{proof}

This frontier concerns the specified three-prime family; it is distinct
from a bound for all odd-dimensional braided fusion categories.


\subsection{Actual phases and the first residue layer}

\begin{lemma}[Arithmetic of a quartic-seed crossing]
\label{paper:quartic-arithmetic}
Let $\C$ be an ordinary integral perfect fusion-simple modular
category of dimension $D=p^6q^6r^4$, with $p<q<r$ odd primes,
and suppose
\[
 X\in P_r,\quad d_X=p^ir,\quad 1\le i\le3,
 \qquad B\in P_q,\quad d_B=q^jr,\quad 1\le j\le2.
\]
Put $c=p^iq^jr$, $z=S_{XB}/c$, $K=\mathbb Q(z)$,
$a=[K_X:K]$, $b=[K_B:K]$, and
$\mathcal B=D/c^2=p^{6-2i}q^{6-2j}r^2$.
The actual balancing word has the unique form
\begin{equation}\label{paper:eq:quartic-word}
 z=\eta w,\qquad
 w=h+q\sum_{\xi\ne1}n_\xi\xi,\qquad
 h+qt=r,\quad t=\sum_\xi n_\xi\ge2,\quad h\ge1,
\end{equation}
where $h,n_\xi$ are integers, $n_\xi\ge0$, $\eta$ is
$q$-primary and the $\xi$ are $r$-primary.
If $q^\alpha$ is the order of $\eta$ and $r^\beta$ is the
largest occurring $r$-order, then
\[
 K=\mathbb Q(\eta)\mathbb Q(w),\qquad
 \operatorname{cond}K=q^\alpha r^\beta,\qquad
 0\le\alpha\le6-2j,\quad 1\le\beta\le2.
\]
The stabilizer of $w$ in the $r$-primary field is tame,
of order $d\mid\gcd(t,r-1)$. With
$Q=h^2+q^2\sum n_\xi^2$ and
$F=h(h-1)+q(q-1)t$, one has
\begin{equation}\label{paper:eq:prime-mass-energy}
 E(z)=\varphi(q^\alpha)E(w),\qquad
 E(w)\ge\frac{r^{\beta-1}}d(rQ-r^2)
       \ge\frac{r^\beta F}d\ge r^\beta q(q-1),
\end{equation}
where $\varphi(q^0)=1$, and $abE(z)<\mathcal B$.
In addition,
\begin{equation}\label{paper:eq:quartic-mass-gap}
                         F>qr/2>2r.
\end{equation}
There is a Fourier channel, primitive at $r^\beta$ and
also at $q^\alpha$ when $\alpha>0$, such that
\begin{equation}\label{paper:eq:quartic-fourier}
 |\widehat z(\chi)|^2\ge q^{\alpha+2}r^\beta,
 \qquad abq^{\alpha+2}r^\beta\le\mathcal B.
\end{equation}
If both inputs are $r$-inactive, this strengthens to
\begin{equation}\label{paper:eq:quartic-inactive}
                    abq^{\alpha+2}r^{2\beta}\le\mathcal B.
\end{equation}
\end{lemma}
\begin{proof}
LCM integrality and balancing give an integral positive word
of mass $r$ with strict house below $r$. Since
$\chi_B(X)=p^iz$ and $B\in P_q$, division by the $q$-unit
$p^i$ gives residue $r$ at every $q$-place. Similarly
$\chi_X(B)=q^jz$ and $X\in P_r$ give residue $r=0$ at
every $r$-place. Lemma~\ref{paper:family:lem:prime-mass} gives
\eqref{paper:eq:quartic-word}, initially with $t\ge1$.
If $t=1$, the full character word, up to a common scalar,
has two $r$-primary positions. This contradicts
Theorem~\ref{paper:family:short-support}, since $r\ge7$.

If a cyclotomic automorphism fixes $z$, then
$\sigma(\eta)/\eta=w/\sigma(w)$ lies in the intersection
of the coprime primary fields. It is therefore a rational
odd-order root of unity, hence one. This proves field splitting
and the energy factorization. Uniqueness of the actual word
makes the stabilizer of $w$ preserve its coefficients.
The bound $0<t<r/q<r-1$ excludes a nontrivial wild
stabilizer and full tame invariance. Thus its conductor is
exactly $r^\beta$, and $d\mid\gcd(t,r-1)$.
The conductor bounds follow from the two endpoint codegrees
and Theorem~\ref{paper:family:por:th:exact-conductor}.

The normalized trace of the quotient of two distinct
$r$-primary roots is at least $-1/(r-1)$. Applying this
to the mass-$r$ word gives the first energy inequality;
$\sum n_\xi^2\ge t$ and $d\le t$ give the others.
The relative-index budget is
\eqref{paper:two-sided-energy}.

For the Fourier assertion, let $m$ be the integer coefficient
function on the primitive $r^\beta$-root cosets of the tame
stabilizer. It has a nonzero primitive Fourier coefficient:
at $\beta=1$ it is nonconstant, and at $\beta\ge2$ it cannot
be constant on every last-wild orbit, since its total
multiplicity is below $r$. Lemma~\ref{paper:integral-fourier}
gives such a channel with $|\widehat m|\ge1$.
The quotient Gauss-sum calculation of
Corollary~\ref{paper:primary-inactive-cost} gives
$\widehat w=q\tau_{r^\beta}\widehat m$, with
$|\tau_{r^\beta}|^2=r^\beta$; lower levels vanish.
When $\alpha>0$, tensor with a primitive $q^\alpha$-character.
Field splitting makes the scalar $\eta$ contribute an
independent Gauss factor of squared norm $q^\alpha$.
This proves the first part of \eqref{paper:eq:quartic-fourier}.
The ordinary and $r$-inactive channel bounds of
Theorem~\ref{paper:inactive-fourier-budget} prove the rest.

For the final mass estimate, $r=h+qt>2q$ and
\[
 F=(q-1)r+h(h-q)\ge(q-1)r-q^2/4.
\]
Thus $F/r>7q/8-1>q/2>2$, since $q\ge5$.
\end{proof}

\begin{lemma}[Heights and primitive content at a quartic crossing]
\label{paper:quartic-height-content}
In Lemma~\ref{paper:quartic-arithmetic}, one necessarily has $i\ge2$.
The normalized crossing $z$ is a unit at some place over $p$.
Both assertions hold for every $r$.
\end{lemma}
\begin{proof}
Pair $X\in P_r$, of degree $p^ir$, with an $r$-free nonunit
of $P_q$. Lemma~\ref{paper:family:4.10} gives $p^i\ge q>p$,
so $i\ge2$.

Write $z=\eta w$ as in \eqref{paper:eq:quartic-word}.
The integer content of its coefficients $h,qn_\xi$ divides
their sum $r$; because $0<h<r$, this content is one.
The support contains fewer than $r$ roots and extends to an
integral basis of the $r$-primary cyclotomic field: omit an
unmarked root from each $\mu_r$-coset. Thus $w\notin p\mathcal O$.
The field of $z$ is unramified at $p$, and $\eta$ is a unit.
If $z$ vanished at every $p$-place, then $w/p$ would be
integral, contrary to the primitive coefficient vector.
\end{proof}


\begin{theorem}[Joint twist budget at every quartic-seed height]
\label{paper:quartic-joint-twists}
In Lemma~\ref{paper:quartic-arithmetic}, put
\[
 A=\max\{\alpha,v_q(\operatorname{ord}\theta_X),
                      v_q(\operatorname{ord}\theta_B)\},
 \qquad
 R=\max\{\beta,v_r(\operatorname{ord}\theta_X),
                    v_r(\operatorname{ord}\theta_B)\}.
\]
Then $1\le R\le2$, $A\le6$, and
\begin{equation}\label{paper:eq:quartic-joint-twists}
                  A+2R\le13-2i-2j.
\end{equation}
If both inputs are $r$-inactive, the right side improves
to $12-2i-2j$. The $q,r$-primary part of every input twist
and every constituent twist of $X\otimes B$ divides
$q^Ar^R$. The possibilities are bounded by the table:
\[
\begin{array}{c|c|c}
 (i,j)&R=1&R=2\\ \hline
 (2,1)&q^5r&q^3r^2\\
 (2,2)&q^3r&qr^2\\
 (3,1)&q^3r&qr^2\\
 (3,2)&qr&\text{impossible}
\end{array}
\]
For $i=3$ both inputs and all constituents are also
$p$-inactive, so the displayed divisors bound their full
twist orders. In particular, at $(i,j)=(3,2)$ all these
twists have order dividing $qr$.
\end{theorem}
\begin{proof}
The twist/codegree divisibility of Theorem~\ref{paper:3.5}
gives $A\le6$ and $R\le2$. Put $s=6-i-j$; then
$\mathcal B\le q^{2s}r^2$.

Suppose first that an input is $r$-active. The relative
$q$-inertia contribution to $ab$ is
$L_q=\varphi(q^A)/\varphi(q^\alpha)$: choose an input
attaining the maximal $q$-twist if $A>\alpha$, and use
Corollary~\ref{paper:active-inertia}. The tame subgroup of
order $d$ fixing $K$ acts freely on an $r$-active input.
If $R>\beta$, choose an input with $r$-twist order $r^R$;
the subgroup fixing $K$ then contributes
$L_r=dr^{R-\beta}$. If the two contributions use different
inputs, their relative indices multiply. If they use the
same input, apply the joint stabilizer containment of
Corollary~\ref{paper:active-inertia} to the product subgroup
fixing $K=\mathbb Q(\eta)\mathbb Q(w)$. Its image in the
two twist quotients has order $L_qL_r$, so $L_qL_r\mid ab$
also in this case. Consequently
\begin{equation}\label{paper:eq:quartic-active-budget}
 \varphi(q^A)r^RF<\mathcal B\le q^{2s}r^2,
 \qquad F\ge2q(q-1).
\end{equation}
If $A+2R\ge2s+2$ and $R=1$, then $A\ge2s$, and the
left side is strictly greater than
$2q^{2s-1}(q-1)r^2>q^{2s}r^2$, by
\eqref{paper:eq:quartic-mass-gap}. If $R=2$ and $s=1$, already
$F>q^2$ is impossible. If $R=2$ and $s\ge2$, then
$A\ge2s-2$, and the left side is at least
$2q^{2s-2}(q-1)^2r^2>q^{2s}r^2$, since $q\ge5$.
This proves \eqref{paper:eq:quartic-joint-twists}.

If both inputs are $r$-inactive, then $R=\beta$.
The same $q$-index bound gives
$abq^\alpha\ge(1-1/q)q^A$.
Equation~\eqref{paper:eq:quartic-inactive} therefore implies
\[
 (1-1/q)q^{A+2}r^{2R-2}\le q^{2s}.
\]
Since $r>q$, the left side is at least
$(1-1/q)q^{A+2R}$. The inequality
$A+2R\ge2s+1$ would make this at least
$(q-1)q^{2s}>q^{2s}$, proving the stronger bound.

Every actual phase in \eqref{paper:eq:quartic-word} is
$\eta\xi$, so balancing gives the stated constituent
twist divisibility. Finally, if $i=3$, then $p\nmid f_X$.
Lemma~\ref{paper:family:lw:weyl-line} makes $X$ $p$-inactive and makes
its row vanish on all $p$-active columns. The nonzero
$q$-residue of $z$ gives $S_{XB}\ne0$, so $B$ is also
$p$-inactive. The actual phases have trivial $p$-part,
and hence so do all constituent twists.
\end{proof}

\begin{corollary}[Maximal phase and product conductors]
\label{paper:quartic-conductor}
In Lemma~\ref{paper:quartic-arithmetic}, take $(i,j)=(3,2)$.
For every $r$,
\begin{equation}\label{paper:eq:quartic-conductor}
                 \operatorname{cond}K\in\{r,qr\}.
\end{equation}
Both input twists and every constituent twist of $X\otimes B$
have order dividing $qr$.
\end{corollary}
\begin{proof}
The joint theorem gives $A+2R\le3$, with $R\ge1$.
Thus $R=1$ and $A\le1$. The $i=3$ clause controls the
full twists, proving both assertions.
\end{proof}
\begin{lemma}[The hypothetical inactive maximal crossing]
\label{paper:inactive-period}
In Lemma~\ref{paper:quartic-arithmetic}, take $(i,j)=(3,2)$
and suppose both inputs are $r$-inactive, without assuming
$r<p^2$. There are a subgroup $H\le\mathbb F_r^\times$ of
odd order $d\ge3$ and $v\in\mathbb F_r^\times$ such that
\[
 z=(r-qd)+q\sum_{u\in H}\zeta_r^{vu},\qquad
 K_X=K_B=K=\mathbb Q(\zeta_r)^H.
\]
Writing $e=[K:\mathbb Q]=(r-1)/d$, one has
\[
 q+1\le e\le2q-2,\qquad e\text{ even},\qquad r-qd\ge2.
\]
Both input twists are one, and every constituent of
$X\otimes B$ has twist order dividing $r$.
\end{lemma}
\begin{proof}
Here $\mathcal B=q^2r^2$, so
\eqref{paper:eq:quartic-inactive} forces
$\alpha=0$, $\beta=1$ and $a=b=1$. Thus both endpoint
fields equal the $r$-primary field $K$. Their $p$- and
$q$-inertia groups fix the labels, and
Corollary~\ref{paper:active-inertia} makes them inactive
there. Together with the hypothesis, this gives twist one;
the actual phases then give the constituent assertion.

Put $H=\operatorname{Gal}(\mathbb Q(\zeta_r)/K)$,
$d=|H|$, and $\Gamma=\mathbb F_r^\times/H$.
Let $m_C$ be the common coefficient $n_\xi$ on an
$H$-coset. Its integer mass $M=t/d$ satisfies
$1\le M\le e/q$, since $qt\le r-1=de$.
The same quotient Gauss calculation and inactive channel
bound as in Lemma~\ref{paper:quartic-arithmetic} give
$|\widehat m(\chi)|\le1$ for every nontrivial character
of $\Gamma$. Since $1\le M\le e/q<e-1$,
Lemma~\ref{paper:integral-fourier-rigidity} makes $m$ a
singleton indicator, proving the period formula and $t=d$.

By Lemma~\ref{paper:family:2.4}, $X$ is not self-dual, so its
character field $K$ is not real. Hence $-1\notin H$,
making $d$ odd and $e$ even. Since $t=d\ge2$, one has
$d\ge3$; the positive integer $r-qd$ is even.

The trivial Fourier transform is
$\widehat z(1)=e(r-qd)-q=r(e-q)$.
The ordinary orbit-projection bound
\eqref{paper:orbit-projection} is strict in this channel:
the positive orbit average on $B$ has nonzero image at
the vacuum, outside the orbit of $X$. As
$c/\sqrt D=1/(qr)$, this gives $|e-q|<q$.
Also $qt\le r-1$ gives $e\ge q$; its even parity excludes
equality. Thus $q<e<2q$, giving the stated bounds.
\end{proof}

\begin{theorem}[An active input at every maximal quartic crossing]
\label{paper:maximal-r-activity}
In Lemma~\ref{paper:quartic-arithmetic}, take $(i,j)=(3,2)$.
At least one of $X,B$ is $r$-active. This holds without
assuming $r<p^2$.
\end{theorem}
\begin{proof}
Suppose both inputs are $r$-inactive. By
Lemma~\ref{paper:inactive-period}, their complete twists
are one, both character fields equal $K$, and
\[
 z=(r-qd)+q\sum_{u\in H}\zeta_r^{vu},\qquad
 d=|H|\ge3,\qquad e=[K:\mathbb Q]\ge q+1.
\]
In particular $r>3q$, so $qr>2p^2$.

For a nontrivial character $\chi$ of
$\Gamma=\operatorname{Gal}(K/\mathbb Q)$, let $u_{B,\chi}$
be its norm-one Fourier vector on $\operatorname{Orb}(B)$.
The single-period formula gives
$|\widehat z(\chi)|^2=q^2r$. Since the two relative indices
are one and $c^2/D=1/(q^2r^2)$,
\eqref{paper:orbit-projection} gives the exact equality
\begin{equation}\label{paper:eq:maximal-channel-saturation}
             \|P_{\operatorname{Orb}(X)}s u_{B,\chi}\|^2=1/r.
\end{equation}
Let $P_r^0$ project onto all $r$-inactive labels. The local
compression bound is also $1/r$, and therefore
\[
 P_r^0s u_{B,\chi}=P_{\operatorname{Orb}(X)}s u_{B,\chi}.
\]
The symmetric statement holds from the $X$-orbit to the
$B$-orbit. Thus, for any $r$-inactive $Y$ outside
$\operatorname{Orb}(B)$, all nontrivial Fourier coefficients
of $\sigma\mapsto S_{Y,\sigma X}$ vanish. Because
$S_{YX}\in K_X=K$, Galois covariance gives
\begin{equation}\label{paper:eq:maximal-rational-fiber}
 S_{XY}\in\mathbb Q\qquad
 (\theta_{Y,r}=1,\quad Y\notin\operatorname{Orb}(B)).
\end{equation}

The same equality also fixes the complementary local action
on the orbit differences. Let $P_q^0$ be the $q$-inactive
projection and put $u=u_{B,\chi}$, $v=W_{(qr)'}u$.
The vector $v$ is fixed by the upper groups at $q,r$ and
has the nontrivial prime-level $r$-torus character $\chi$.
Commuting the complementary factors and using exact
$r$-compression gives
\[
 \frac1r\le\|P_q^0P_r^0su\|^2
   =\frac1r\|P_{K_{r,1}}P_q^0W_qv\|^2
   \le\frac1r\|P_q^0W_qv\|^2\le\frac1r.
\]
Equality says that $W_qv$ is upper-unipotent fixed at $q$.
The vector $v$ is therefore fixed by both upper and lower
unipotents, hence by the full $q$-local group. The complementary
factor commutes with this group, so it fixes $u$ too.
As $\chi$ varies, these vectors span the orbit augmentation.
Consequently $W_q$ fixes every difference
$\mathbf e_Z-\mathbf e_{Z'}$ with
$Z,Z'\in\operatorname{Orb}(B)$.

We next show that $p$ and $q$ are separated. For any
$q$-active column $J$, Proposition~\ref{paper:family:4.2}, applied
to the trivial prime-to-$q$ twist fiber of $P_q$ modulo $r$,
supplies an $r$-free nonunit $Y\in P_q$ with pure $q$-twist
and $S_{YJ}\ne0$. The degree caps and exclusion of prime-power
degrees give $d_Y=p^i q^j$, with $1\le i\le3$ and
$1\le j\le2$.
The label $Y$ is $r$-inactive and lies outside the $B$-orbit.
If $i\le2$, then
\[
                     z_Y=\frac{S_{XY}}{p^3q^jr}
\]
is an integer by \eqref{paper:eq:maximal-rational-fiber} and
LCM integrality. It has $|z_Y|<p^i$ and residues $p^i$ at
both $q,r$, from the principalities of $Y,X$, respectively.
Thus $z_Y\equiv p^i\pmod{qr}$, which is impossible since
$qr>2p^2\ge2p^i$. Every such witness therefore has $i=3$.
Its $p$-unit codegree makes $J$ $p$-inactive by
Lemma~\ref{paper:family:lw:weyl-line}. Since $J$ was arbitrary, no twist
contains both $p$ and $q$.

The final assertion of Lemma~\ref{paper:family:lw:weyl-line} now makes
every $q$-active label have $p$-unit codegree. The active-column
criterion in Theorem~\ref{paper:lw:columns} makes $W_q$ integral
at $p$. Its $B$-column is rational, since $K_B$ is $r$-primary
and the local Galois identity fixes that column under all
$q$-inertia. Here $f_B=p^6q^2r^2$. The foreign-column bound
gives integrality at $r$; the rotation identity
$f_B(W_q)_{AB}^2\in\mathcal O$ gives integrality at all
other primes and $q$-valuation at least $-1$. Consequently
\[
                            qW_q\mathbf e_B\in\mathbb Z^I.
\]
Apply Lemma~\ref{paper:integral-orbits} to
$\mathcal O=\operatorname{Orb}(B)$ at $q$. Its labels are
inactive with unramified character fields and $q$-divisible
codegrees; their Weyl differences are fixed and the scaled
column is integral. The odd-dimensional conclusion gives
$e\le(q+1)/2$, contradicting $e\ge q+1$.
\end{proof}


\begin{lemma}[An integer obstruction to a half-channel bound]
\label{paper:integer-half-obstruction}
Let $q\ge5$, $M\ge1$, $k\ge0$ be integers with $k^2\le q$,
and put $e=qM+k$. No nonnegative integer function $m$ of
mass $M$ on a cyclic group of order $e$ satisfies
\[
              |\widehat m(\chi)|^2\le\frac e{2q}
                        \quad(\chi^2\ne1).
\]
\end{lemma}
\begin{proof}
There is at most one nontrivial quadratic character, and its
Fourier coefficient has modulus at most $M$. Since
$\sum m^2\ge M$, Parseval would give
\[
 eM-2M^2\le\sum_{\chi^2\ne1}|\widehat m(\chi)|^2
                           \le\frac{e(e-1)}{2q}.
\]
Twice $q$ times the difference of the two outer terms is
\[
             q(q-4)M^2+qM+k-k^2>0,
\]
a contradiction.
\end{proof}

\begin{theorem}[Elimination of the nontrivial scalar branch]
\label{paper:maximal-scalar-trivial}
In the remaining ordered odd sector of the exact family
$p^6q^6r^4$, every maximal quartic crossing has $\eta=1$.
Consequently both inputs are $q$-inactive, at least one input
is $r$-active, the phase field has conductor $r$, and all
input and product twists divide $r$.
\end{theorem}
\begin{proof}
Retain the maximal quartic crossing and its notation. The following
claims analyze a possibly nontrivial scalar phase. Their explicitly
stated broader ranges are preserved; the final contradiction uses
the surviving-family bounds.

\begin{proofclaim}[Activity in the maximal quartic crossing]
\label{paper:quartic-activity}
In the maximal case $(i,j)=(3,2)$ of
Theorem~\ref{paper:quartic-joint-twists}, retain $z=\eta w$
and the relative field indices $a,b$. These conclusions hold
for every $r$.
If $\eta\ne1$, then $abq\le r$ and an input is
$q$-active. If $\eta=1$, the inputs cannot both be
$q$-active. In either case an input is $r$-active by
Theorem~\ref{paper:maximal-r-activity}.
\end{proofclaim}
\noindent\emph{Proof of the claim.}
The joint theorem gives $\alpha\le1$ and $\beta=1$.
The tame stabilizer order $d$ divides $ab$, since an
input is $r$-active. Suppose first that $\eta\ne1$.
The mixed channel in \eqref{paper:eq:quartic-fourier}
has squared modulus at least $q^3r$, giving $abq\le r$.

Suppose both inputs were $q$-inactive. On
$\Gamma_r=\mathbb F_r^\times/H$, of order $e=(r-1)/d$,
let $m$ be the nonnegative integer coefficient vector of
$w$ on the nonidentity $H$-cosets. Put
\[
 M=\sum_Cm_C=t/d\ge1,\qquad
 k=e-qM=(h-1)/d\in\mathbb Z_{\ge0}.
\]
The $r$-trivial trace is
$\operatorname{Tr}_{\mathbb Q(w)/\mathbb Q}(w)=rk$.
Tensoring it with a nontrivial $q$-character gives a
channel of squared modulus $qr^2k^2$; the $q$-inactive
budget yields $abk^2\le1$. Thus $0\le k\le1$.
For every nontrivial $\chi_r$, the mixed channel gives
\[
                  abq^2|\widehat m(\chi_r)|^2\le r.
\]
Summing and applying Parseval and $\sum m_C^2\ge M$
shows $abq^2M(e-M)\le r(e-1)$. Since $d\le ab$ and
$r=de+1$, it follows that
\[
 q^2M(e-M)\le(e+1/d)(e-1)\le(qM+2)qM.
\]
But $e=qM+k$ implies $e-M\ge(q-1)M$, so this would
give $q(q-2)M\le2$, impossible for $q\ge5$.

Finally suppose $\eta=1$ and both inputs are $q$-active.
The two independent $q$-inertia contributions, together
with the tame $r$-contribution at an active input, give
$d(q-1)^2\mid ab$ by the joint stabilizer containment
of Corollary~\ref{paper:active-inertia}. Using $F>2r$
in \eqref{paper:eq:prime-mass-energy} gives
\[
 abE(z)\ge(q-1)^2rF>2(q-1)^2r^2>q^2r^2,
\]
contradicting the two-sided budget.
\par\smallskip

\begin{proofclaim}[Maximal activity detects the scalar phase]
\label{paper:maximal-q-activity-equivalence}
At a maximal quartic crossing of Theorem~\ref{paper:quartic-joint-twists},
write $z=\eta w$ as in Lemma~\ref{paper:quartic-arithmetic}.
Some input is $q$-active if and only if $\eta\ne1$.
\end{proofclaim}
\noindent\emph{Proof of the claim.}
Claim~\ref{paper:quartic-activity} already proves the forward
existence of a $q$-active input when $\eta\ne1$.
Suppose $\eta=1$ and some input is $q$-active. The same
corollary makes the other input $q$-inactive, and the joint
twist bound makes the active order exactly $q$. The phase
field is $r$-primary, so Corollary~\ref{paper:half-channel-crossing}
gives
\[
                abE(z)\le\frac{q^2-1}{2}r^2.
\]
There is an $r$-active input by Theorem~\ref{paper:maximal-r-activity}.
The joint active-inertia bound supplies $d(q-1)\mid ab$, where
$d$ is the tame $r$-stabilizer order. With
$F=h(h-1)+q(q-1)t$ and $r=h+qt$, the word-energy estimate gives
\[
                 abE(z)\ge(q-1)rF.
\]
But $r>2q$ and
$F\ge(q-1)r-q^2/4$ imply
\[
 \frac{(q-1)F}{r}
   >(q-1)(7q/8-1)
   >\frac{q^2-1}{2},
\]
since the difference in the last comparison is
$3(q-1)(q-4)/8>0$ for $q\ge5$. This is a contradiction.
\par\smallskip

\begin{proofclaim}[Synchronized maximal activity and minimal relative indices]
\label{paper:maximal-synchronized}
At the maximal quartic crossing, retain $z=\eta w$, the endpoint
indices $a,b$, and the tame $r$-stabilizer order $d$. Then
\[
 \begin{array}{c|c}
  \eta=1&\text{both inputs are $q$-inactive},\\
  \eta\ne1&\text{both inputs are $q$-active and }ab=d.
 \end{array}
\]
In the second case, if both inputs are $r$-active then
$d=a=b=1$. If exactly one is $r$-active, its relative index is
$d$ and the other relative index is one.
\end{proofclaim}
\noindent\emph{Proof of the claim.}
The first assertion is Claim~\ref{paper:maximal-q-activity-equivalence}.
Suppose $\eta\ne1$ and exactly one input is $q$-active; existence
of at least one is already known. The phase field splits as
$K=\mathbb Q(\zeta_q)\mathbb Q(w)$, with the second factor
$r$-primary, and both actual input $q$-orders divide $q$.
Choose a nonquadratic character $\chi_q$ of
$\mathbb F_q^\times$, possible since $q\ge5$.
For every character $\chi_r$ of $\Gal(\mathbb Q(w)/\mathbb Q)$,
the corresponding source orbit vector is upper-fixed and has
torus character $\chi_q^{\pm1}$ at $q$. Its target orbit lies
in one primitive quadratic phase class. Lemma~\ref{paper:nonquadratic-half-local}
and the crossing normalization give
\[
 abq\,|\widehat w(\chi_r)|^2
   \le\frac{q-1}{2q}\,q^2r^2.
\]
Here the factor $q$ is the squared norm of the nontrivial
prime-level Gauss sum from $\eta$. Summing over $\chi_r$ and
using $E(z)=(q-1)E(w)$ gives
\[
                  abE(z)\le\frac{(q-1)^2}{2}r^2.
\]
An $r$-active input always exists, so $d\mid ab$. The word
energy bound gives $abE(z)\ge(q-1)rF$. Since
$F/r>7q/8-1>(q-1)/2$, this is a contradiction. Both inputs
are therefore $q$-active.

Finally $ab/d$ is a positive integer and
\[
 \frac{ab}{d}(q-1)rF<q^2r^2.
\]
The estimate $F/r>7q/8-1$ gives
$2(q-1)F/r>q^2$ for $q\ge5$: the difference is greater than
$(3q^2-15q+8)/4>0$. Hence $ab/d=1$.
For each $r$-active endpoint, $d$ divides its individual
relative index. The stated alternatives now follow from $ab=d$.
\par\smallskip

\begin{proofclaim}[Nontrivial scalar phase forces both $r$-activities]
\label{paper:maximal-r-synchronized}
At a maximal quartic crossing with $\eta\ne1$, both inputs
are $r$-active, and $a=b=d=1$.
\end{proofclaim}
\noindent\emph{Proof of the claim.}
By Claim~\ref{paper:maximal-synchronized}, $ab=d$, and at
least one input is $r$-active. Suppose exactly one is. Write
$\Gamma=\mathbb F_r^\times/H$, $e=(r-1)/d$, and let
$m_C\in\mathbb Z_{\ge0}$ be the coefficient of the word on
the coset $C$. Set
\[
 M=\sum_Cm_C=t/d\ge1,\qquad k=(h-1)/d\ge0.
\]
The mass identity gives $e=qM+k$. The trivial $r$-channel
of $w$ is $rk$, whereas any nontrivial $q$-channel of
$\eta$ has squared modulus $q$. The ordinary channel bound
therefore gives $dk^2\le q$, in particular $k^2\le q$.

For a character $\chi$ of $\Gamma$ with $\chi^2\ne1$,
\[
 |\widehat w(\chi)|^2=q^2r|\widehat m(\chi)|^2.
\]
The $r$-inactive source is upper-fixed, the $r$-active target
orbit lies in one primitive quadratic class, and the source
torus character is the tame inflation of $\chi^{\pm1}$.
Apply Lemma~\ref{paper:nonquadratic-half-local} at $r$, also
using a nontrivial $q$-channel. Since $ab=d$, this gives
\[
 dq^3r|\widehat m(\chi)|^2
       \le\frac{r-1}{2r}q^2r^2,
 \qquad |\widehat m(\chi)|^2\le\frac e{2q}.
\]
These inequalities contradict
Lemma~\ref{paper:integer-half-obstruction}.
Both inputs are $r$-active; the last assertion
now follows from Claim~\ref{paper:maximal-synchronized}.
\par\smallskip

\begin{proofclaim}[A fourth-moment mass bound]
\label{paper:maximal-fourth-moment}
Let $q\ge5$, $t\ge2$ and $h\ge1$ be integers, put $r=h+qt$,
and suppose $h\le1+\sqrt q$. Let $m$ be a nonnegative
integer-valued function of mass $t$ on a cyclic group of
order $e=r-1$. If
\[
          |\widehat m(\chi)|^2\le r/q\qquad(\chi\ne1),
\]
then $t\ge q$ and $r>q^2$. In particular these conclusions
hold at a maximal quartic crossing with $\eta\ne1$.
\end{proofclaim}
\noindent\emph{Proof of the claim.}
Put $S=\sum_xm(x)^2\ge t$, $L=r/q=t+h/q$, and
$x_\chi=|\widehat m(\chi)|^2$ for $\chi\ne1$.
Let $c(g)=\sum_xm(x)m(gx)$. These are nonnegative integers,
$c(1)=S$, and $\sum_gc(g)=t^2$. Consequently
\[
 \sum_gc(g)^2\ge S^2+t^2-S.
\]
The second and fourth Parseval identities give
\[
 \sum_{\chi\ne1}x_\chi=eS-t^2,\qquad
 \sum_{\chi\ne1}x_\chi^2=e\sum_gc(g)^2-t^4.
\]
Since $x_\chi^2\le Lx_\chi$, it follows that
\[
 e(S^2+t^2-S)-t^4\le L(eS-t^2).
\]
The left side minus the right side is strictly increasing
in $S\ge t$, since $2S-1-L\ge t-1-h/q>0$. Substituting
$S=t$ is therefore necessary. After division by $t$, the
resulting left side minus right side is
\[
 D(t,h,q)=t(t-1)(q-t)-(h+t-1)+\frac{h(t-h+1)}q.
\]
Suppose $2\le t\le q-1$. This expression decreases in $h$
throughout $1\le h\le1+\sqrt q$, because
$\partial D/\partial h=-1+(t-2h+1)/q<0$.
Put $s=\sqrt q$. The identity
\[
 t(t-1)(q-t)-t-(2q-6)
       =(t-2)\{t(q-1-t)+q-3\}\ge0
\]
therefore gives
\[
 \begin{split}
 D(t,h,q)&\ge D(t,s+1,q)\\
 &\ge2q-6-s+\frac{(s+1)(2-s)}q\\
 &=2q-7-\sqrt q+\frac1{\sqrt q}+\frac2q>0.
 \end{split}
\]
The final inequality holds for every $q\ge5$, since
$2q-7>\sqrt q$ there. This contradicts the necessary inequality.
Hence $t\ge q$, and $r=h+qt>q^2$.

For the maximal crossing, Claim~\ref{paper:maximal-r-synchronized}
gives $a=b=d=1$. The coefficient function of its nonconstant
word is defined on $\mathbb F_r^\times$ and has mass $t\ge2$.
The trivial $r$-channel, paired with a nontrivial $q$-channel,
gives $(h-1)^2\le q$. The other channels, using their
prime-level Gauss sums, give $|\widehat m(\chi)|^2\le r/q$.
Thus all the general hypotheses hold.
\par\smallskip

\noindent\emph{The contradiction.}
Corollary~\ref{paper:family:cor:remaining-region} gives $5\le p<q<r$ and $r>3q$.
Suppose $\eta\ne1$. Claim~\ref{paper:maximal-r-synchronized}
gives $a=b=d=1$. The full endpoint Galois orbits are disjoint
because their degrees are different. They are invariant under
$T/\gamma$ and charge conjugation, so the global braid budget
in Lemma~\ref{paper:crossing-energy} gives
\[
                         E(z)<\tfrac34 q^2r^2.
\]
On the other hand, $E(z)\ge(q-1)rF$, and
\[
 \frac Fr\ge q-1-\frac{q^2}{4r}>\frac{11q}{12}-1.
\]
For $q\ge11$ this contradicts the braid budget, since
\[
 (q-1)(11q/12-1)-3q^2/4
                   =\frac{2q^2-23q+12}{12}>0.
\]
For $q<11$, the least-prime bound and primality
imply $p\le5$ and $q\ge7$. Thus
\[
 1+\left\lfloor\frac{\min(r,p^2)}q\right\rfloor\le4.
\]
Claim~\ref{paper:maximal-fourth-moment} gives $r>q^2\ge49$.
The uniform four-position support theorem, equivalently
Corollary~\ref{paper:family:family-criterion}, excludes this entire
parameter range. This uses only a uniform support bound,
not a check of any individual dimension or prime triple.
Hence $\eta=1$ in all cases.

The remaining assertions follow from
Claim~\ref{paper:maximal-synchronized},
Theorem~\ref{paper:maximal-r-activity}, and the maximal
joint-twist theorem. In particular both inputs are $p$- and
$q$-inactive, and the pure $r$ scalar word makes every product
constituent $r$-primary as well.
\end{proof}

\begin{corollary}[Full character fields at maximal crossings]
\label{paper:maximal-fields}\label{paper:quartic-tame-fields}
At a maximal quartic crossing, both full endpoint character
fields have $r$-conductor exponent one. Their actual twists
are $r$-primary, but this does not imply that their full
character fields are $r$-primary.
\end{corollary}
\begin{proof}
Theorem~\ref{paper:maximal-scalar-trivial} gives $\beta=1$
and an $r$-active input. Its relative index contains the
order $d$ of the tame $r$-stabilizer. If an endpoint field
had $r$-conductor exponent at least two, its relative index
would contain the additional factor $r$. These coprime costs
multiply, including when they occur at the same endpoint.
Thus $rd\mid ab$, and the word-energy bound gives
\[
 abE(z)\ge r^2F\ge2q(q-1)r^2>q^2r^2,
\]
contrary to the ordinary relative-index budget.
\end{proof}

\begin{corollary}[No triple-residue prime-mass crossing in the surviving family]
\label{paper:no-triple-residue-crossing}
In the surviving exact odd family, no positive rational multiple of a
nonunit character at a nonunit row can satisfy
Proposition~\ref{paper:three-residue-prime-mass}.
\end{corollary}
\begin{proof}
Put $m=\sum m_\xi\le\lfloor r/(pq)\rfloor$. Strictness gives $m\ge1$,
and the actual $r$-word has at most $m+1$ positions. The
surviving-region bound gives
\[
 r<(p^2/q)^3<(pq)^{3/2}.
\]
Thus $m^3<r$ and $m^2<r/\sqrt{pq}$. Since $p\ge5$,
\[
 p\bigl(m(m+1)+1\bigr)\le3pm^2
 <\frac{3p}{\sqrt{pq}}r<3r<(p-1)r.
\]
Exposure therefore holds. Also $m\ge1$ and $h\ge1$ give
$r>pq\ge35>33$. The scalar estimate in the proof of
Corollary~\ref{paper:family:cor:remaining-region}, applied to the support bound
$m+1$ in Theorem~\ref{paper:family:thm:descent}, now excludes the word because
$r>\max\{33,m^3\}$.
\end{proof}

\begin{corollary}[A forbidden double-principal active label]
\label{paper:no-double-principal-active}
In the surviving exact odd family, any $p$-free nonunit in
$P_q\cap P_r$ forces $p,r$ to be separated and is $r$-inactive.
Symmetrically, any $q$-free nonunit in $P_p\cap P_r$ forces
$q,r$ to be separated and is $r$-inactive.
\end{corollary}
\begin{proof}
Let $B$ be such a $p$-free nonunit. Its degree cannot be a prime
power. Its membership in $P_q$ and $P_r$ also excludes top $q$-height
three and top $r$-height two, since a codegree prime to the corresponding prime gives a singleton
full block. Thus $d_B=q^jr$ with $j=1$ or $2$.
Suppose a $pr$-active column exists. Apply prime avoidance on $P_p$,
omitting $q$, to this column. The resulting
nonunit $U\in P_p$ is $q$-free and has a nonzero hit on the column.
Its degree is $p^a r^b$ with $1\le a\le2$ and $1\le b\le2$.
If $b=2$, its $r$-unit codegree makes it miss every $r$-active column.
Hence $b=1$.

The normalized crossing
\[
 z=\frac{S_{UB}}{p^a q^j r}
\]
is an integral positive rational word of mass $r$ and strict house
below $r$. Since $\chi_B(U)=p^a z$ and $B\in P_q\cap P_r$, its
residue is $r$ at both $q,r$. Since $\chi_U(B)=q^jz$ and $U\in P_p$,
its residue is also $r$ at $p$. This contradicts
Corollary~\ref{paper:no-triple-residue-crossing}. Thus $p,r$ are
separated. Since $B$ has $p$-divisible codegree, the separation
Weyl-line lemma makes it $r$-inactive. The symmetric assertion follows
by interchanging the two sextic principal blocks; the three residues
and the least-prime exposure bound are unchanged.
\end{proof}

\begin{corollary}[No wild $q$-ramification at a maximal crossing]
\label{paper:no-maximal-q-wild}
Both full character fields in a maximal quartic crossing are tame at $q$.
\end{corollary}
\begin{proof}
Suppose an endpoint field is wild at $q$. We determine the exact
hypothetical field pair before contradicting the double-principal
restriction above.

\begin{proofclaim}[Field structure forced by wild ramification]
\label{paper:maximal-q-wild-exact}
In the surviving ordered odd family, consider a maximal quartic
crossing and suppose an endpoint character field is wild at $q$.
Let $E_q$ be the unique degree-$q$ subfield of
$\mathbb Q(\zeta_{q^2})$. Then $\eta=1$, both inputs are
$q$-inactive and $r$-active, and
\[
 d=1,\qquad ab=q,\qquad
 \{K_X,K_B\}=\{\mathbb Q(\zeta_r),E_q\mathbb Q(\zeta_r)\}.
\]
\end{proofclaim}
\noindent\emph{Proof of the claim.}
Theorem~\ref{paper:maximal-scalar-trivial} gives $\eta=1$
and $q$-inactivity of both inputs. The maximal
joint-twist theorem gives $r$-orders at most $r$, and
Theorem~\ref{paper:maximal-r-activity} gives an $r$-active input.

The nontrivial wild $q$-orbit contributes a factor $q$ to the
relative indices. The tame $r$-stabilizer of order $d$ acts
freely at an $r$-active endpoint. These costs multiply even
if assigned to the same input: a stabilizer element has
$r$-component $\pm1$ by the square action on the active
$r$-twist, and the odd-dimensional inversion theorem removes
the nontrivial sign after taking an odd power that kills its
wild $q$-component. Thus $qd\mid ab$. The ordinary budget and
$E(w)\ge rF/d$, with $F>qr/2$, imply
\[
                       \frac{ab}{d}<2q.
\]
Consequently $ab=qd$.

Put $\Gamma=\mathbb F_r^\times/H$, $e=(r-1)/d$, and let
$m_C$ be the common nonnegative integer coefficient on an
$H$-coset. Write
\[
 M=\sum_Cm_C=t/d\ge1,\qquad
 k=(h-1)/d\ge0,\qquad e=qM+k.
\]
The trivial $r$-channel is $rk$, so its ordinary bound gives
$dk^2\le q$. Suppose exactly one input is $r$-active. For
each $\chi\in\widehat\Gamma$ with $\chi^2\ne1$, the
nonquadratic half-channel lemma at $r$ gives
\[
 qd\,q^2r|\widehat m(\chi)|^2
       \le\frac{r-1}{2r}q^2r^2,
 \qquad |\widehat m(\chi)|^2\le\frac e{2q}.
\]
The integer obstruction of
Lemma~\ref{paper:integer-half-obstruction} contradicts these
bounds. Hence both inputs are $r$-active.

The two active $r$-orbits, together with the wild $q$-orbit,
now give $qd^2\mid ab=qd$ by the same joint stabilizer
argument. Thus $d=1$ and $ab=q$. In particular
$K=\mathbb Q(\zeta_r)$ and one endpoint field equals $K$,
while the other is a degree-$q$ extension of $K$ with wild
$q$-ramification. The $r$-conductor exponent is already one.
Its $q$-conductor exponent is exactly two, because a higher
exponent would contribute $q^2$ to the relative degree.
There is no remaining tame $q$-factor. Nor can a $p$-component
occur: the degree of every $p$-primary cyclotomic field has
prime factors among $p$ and divisors of $p-1$, all smaller
than $q$. Since the field already contains
$\mathbb Q(\zeta_r)$, there is no entanglement with that
factor. This proves the asserted pair of full fields.
\par\smallskip

\noindent\emph{The contradiction.}
If either field were $q$-wild, the exact-field theorem gives
$K_B=\mathbb Q(\zeta_r)$ or $K_B=E_q\mathbb Q(\zeta_r)$, and makes
$B$ $r$-active. In either case partial inversion at $p,q$ fixes its
full character field and therefore its label. Global parity puts
$B$ in the individual principal $r$-block. It already lies in $P_q$
and has degree $q^2r$, contradicting
Corollary~\ref{paper:no-double-principal-active}.
\end{proof}

\begin{corollary}[Full-field cap at a maximal quartic crossing]
\label{paper:maximal-full-field-cap}
For a maximal quartic crossing, the full character fields satisfy
\[
 \operatorname{cond}(K_X)\mid qr,\qquad
 \operatorname{cond}(K_B)\mid p^2qr.
\]
Thus the only possible wild ramification in either endpoint field
is a $p^2$-conductor at the degree-$q^2r$ endpoint $B$.
\end{corollary}
\begin{proof}
Both fields are tame at $q$ by Corollary~\ref{paper:no-maximal-q-wild},
and their $r$-conductor exponents are at most one by the maximal
quartic field bound.  The degree-$p^3r$ label $X$ has $p$-unit
codegree, so its field is unramified at $p$. This proves the first
assertion.

For the second, put $K=\mathbb Q(z)\subseteq\mathbb Q(\zeta_r)$,
$a=[K_X:K]$, $b=[K_B:K]$, and
$d=[\mathbb Q(\zeta_r):K]$.  The ordinary norm bound and
$E(z)\ge rF/d$ with $F>qr/2$ give
\[
                         \frac{ab}{d}<2q.
\]
Suppose the $p$-conductor exponent of $K_B$ were at least three.
Its wild $p$-inertia orbit then has size divisible by $p^2$.
An $r$-active endpoint has the free tame stabilizer cost $d$.
These costs multiply: if both occur at $B$, any simultaneous
stabilizer has tame $r$-component $\pm1$ by the square action on
its active twist; raising to an odd power kills the wild
$p$-component, and partial-inversion parity excludes the negative
sign.  The remaining wild component cannot reduce its prescribed
$p^2$ orbit.  If the $r$-active endpoint is $X$, the costs occur
in separate relative indices.  Hence $p^2d\mid ab$ in either case.
But the surviving support range has $3q\le p^2$, so
\[
                    p^2\le\frac{ab}{d}<2q\le\frac23p^2,
\]
a contradiction. Thus the $p$-conductor exponent of $K_B$ is at most
two, proving the result.
\end{proof}

\begin{corollary}[An even index cost for the remaining wild field]
\label{paper:maximal-p-wild-index}
Suppose the degree-$q^2r$ endpoint $B$ of a maximal quartic crossing
has $p$-conductor $p^2$. Put $K=\mathbb Q(z)$,
$a=[K_X:K]$, $b=[K_B:K]$, and
$d=[\mathbb Q(\zeta_r):K]$. Then $p,r$ occur together in an actual
twist, partial inversion at $p,q$ moves $B$, and
\[
                         2pd\mid ab.
\]
Consequently $ab/d$ is a positive multiple of $2p$ satisfying
$2p\le ab/d<2q$.
If $K_B$ is unramified at $q$, its local $p$-inertia field is nonreal.
\end{corollary}
\begin{proof}
The crossing field has conductor $r$, so $K_B$ ramifies at both
$p$ and $r$. The separated-character-fields assertion of
Proposition~\ref{paper:lw:isolated} therefore forces a $pr$-mixed actual
twist. Let $J=H_{-1,\{p,q\}}$. If $J$ fixed $B$, global parity
would put $B$ in $P_r$. Since $B\in P_q$, the double-principal
corollary would force $p,r$ to be separated, a contradiction.
Thus $J$ moves $B$.

Work through the common field $\mathbb Q(\zeta_{p^2qr})$ supplied
by the full-field cap. Let $P$ be its wild $p$-inertia, of order $p$.
It acts faithfully on $B$ because the $p$-conductor is exactly $p^2$.
The product $P\times\langle J\rangle$ therefore has a free orbit
on $B$: raising a stabilizer element to its $p$th power would make
$J$ fix $B$ if its $J$-component were nontrivial. This gives
$2p\mid b$.

If $X$ is $r$-active, $d\mid a$, so the desired divisibility follows.
Otherwise $B$ is $r$-active. Work with
$P\times\langle J\rangle\times H$, where
$H=\operatorname{Gal}(\mathbb Q(\zeta_r)/K)$ has order $d$.
A stabilizer element has $r$-component $1$ or $-1$, by the square
action on the active $r$-twist. Raise it to the odd power $p$
to remove the wild component. The three nonidentity possibilities
are then $J$, inversion at $r$, or full inversion. The first moves
$B$ by the preceding paragraph, the second cannot fix an $r$-active
label, and full inversion is duality and fixes no nonunit label in
odd dimension. Thus the remaining components are trivial, and
faithfulness of $P$ finishes the proof that this orbit is free.
It has size $2pd$ and fixes $K$, so $2pd\mid b$.
The norm bound $ab/d<2q$ supplies the stated interval.

Finally, if $q$-inertia fixes $B$, the action of $J$ on it is exactly
inversion at $p$. Since $J$ moves $B$, its local $p$-inertia field
is nonreal.
\end{proof}


\begin{lemma}[Wild variance in a prime-square cyclotomic subfield]
\label{paper:general-wild-lattice}
Let $\ell$ be an odd prime, let $H$ be a subgroup of order
$h\mid\ell-1$ in $(\mathbb Z/\ell^2)^\times$, and put
\[
 L=\mathbb Q(\zeta_{\ell^2})^H,\qquad
 L_0=L\cap\mathbb Q(\zeta_\ell),\qquad
 e=(\ell-1)/h.
\]
Let $A=\ell^{-1}\operatorname{Tr}_{L/L_0}$ be the orthogonal
projection onto the tame subfield. For every
$\alpha\in\mathcal O_L$ with $\alpha\notin L_0$,
\[
 \frac1{[L:\mathbb Q]}\sum_\sigma
       |\sigma(\alpha-A\alpha)|^2\ge h.
\]
In particular its total mean square is at least $h$.
\end{lemma}
\begin{proof}
Partition the primitive $\ell^2$-th roots by their nonzero
residue exponent modulo $\ell$. In each residue class their
sum is zero, and deleting one root gives an integral basis
for its primitive part. Choose the deletions equivariantly
under $H$: its action on the nonzero residue classes is free.
Together with the tame cyclotomic integer ring, the retained
roots give an $H$-stable direct-sum decomposition of
$\mathcal O_{\mathbb Q(\zeta_{\ell^2})}$. The primitive
summand has a basis freely permuted by $H$.

Its invariant lattice is consequently generated integrally
by the Gaussian periods $\eta_{C,j}$, where $C$ ranges over
the $e$ cosets of $H$ in $\mathbb F_\ell^\times$, and
$j=0,\ldots,\ell-1$ indexes the $\ell$ lifts of that coset.
Their only relations are
$\sum_j\eta_{C,j}=0$, separately for each $C$.
Every period has relative trace zero to $L_0$.
Thus $\alpha-A\alpha$ is an integer linear combination
$\sum_{C,j}c_{C,j}\eta_{C,j}$; normalize one coefficient
in each $C$ to zero.

Prime-square Ramanujan sums give the mean Hermitian Gram matrix
\[
 \langle\eta_{C,j},\eta_{D,k}\rangle=
 \begin{cases}
  0,&C\ne D,\\
  h,&C=D,\ j=k,\\
  -1/e,&C=D,\ j\ne k.
 \end{cases}
\]
Its value on the coefficient vector is therefore
\[
 \frac1e\sum_C\sum_{j<k}(c_{C,j}-c_{C,k})^2.
\]
If the wild component is nonzero, at least one of the integer
vectors $(c_{C,j})_j$ is nonconstant. At least $\ell-1$
integer pair differences in that vector are nonzero, giving
lower bound $(\ell-1)/e=h$. Orthogonality to $L_0$ also
proves the final assertion.
\end{proof}

\begin{proposition}[Low-height detection of the remaining wild field]
\label{paper:maximal-p-wild-detector}
Suppose a maximal quartic crossing has a $p^2$-conductor at
$B$, of degree $q^2r$. The field generated by the entries of
$W_p\mathbf e_B$ is
\[
 L=\mathbb Q(\zeta_{p^2})^H,
 \qquad 2\le h:=|H|\mid p-1.
\]
Let $v$ be a norm-one nontrivial Fourier vector on the wild
$p$-inertia orbit of $B$. Then at least $1-1/p$ of the
squared norm of $W_pv$ is supported on labels $Z$ with
\[
 v_p(d_Z)\le1,\qquad
 \operatorname{ord}(\theta_{Z,p})\ge p^2,\qquad
 \theta_{Z,q}=1,\qquad
 \theta_{Z,r}=\theta_{B,r}.
\]
In particular, if $B$ is $r$-active these labels are jointly
$p^2$- and $r$-active. No assertion about their membership
in $P_q$ is made.
\end{proposition}
\begin{proof}
The full-field cap bounds the $p$-conductor by $p^2$.
The local Galois identity and the invertibility of $W_p$
show that the stabilizer of its $B$-column is exactly the
$p$-inertia stabilizer of $B$. Hence its column field has
conductor $p^2$ and is of the displayed form, initially
with $h\ge1$.

Suppose $h=1$. Then the $p$-inertia action on $B$ has
faithful quotient $(\mathbb Z/p^2)^\times$, of order
$p(p-1)$. Retain $K=\mathbb Q(z)$, the relative indices
$a,b$, and $d=[\mathbb Q(\zeta_r):K]$. There is an
$r$-active endpoint. If it is $X$, the two relative
indices supply $d$ and $p(p-1)$ separately. If it is $B$,
consider the product of its $p$-inertia quotient and the
tame $r$-stabilizer of order $d$. Every stabilizer element
has $r$-component $\pm1$. Its square has trivial
$r$-component, so faithfulness of the $p$-action makes its
$p$-component also $\pm1$. A nonidentity stabilizer would
therefore be a partial inversion involving the active
prime $r$, forbidden by global parity. Thus in either case
\[
                  p(p-1)d\mid ab.
\]
But $ab/d<2q$ and $3q\le p^2$, while
$p(p-1)>2p^2/3$ for $p\ge5$. This contradiction proves
$h\ge2$.

Consider an entry $w=(W_p)_{ZB}$ with $v_p(d_Z)\ge2$.
The rotation bound, applied by symmetry at the $Z$-column,
gives $v_p(w)\ge-1$ because $v_p(f_Z)\le2$.
The foreign-column bound at $q,r$ uses
$v_q(f_B)=v_r(f_B)=2$ and gives integrality there.
At all other primes the ordinary rotation bound gives
integrality. Consequently $pw\in\mathcal O_L$.

Suppose $Z$ is $p$-active and $w$ is not in the tame
subfield $L_0=L\cap\mathbb Q(\zeta_p)$. Its inertia
orbit has at least $p(p-1)$ elements: nontrivial wild
ramification supplies the factor $p$, and active inertia
supplies the full tame factor $p-1$. The fixed $B$-column
has Galois-conjugate entries along this orbit. Lemma~\ref{paper:general-wild-lattice} therefore makes its squared norm at least
\[
             p(p-1)\frac h{p^2}
                       =\frac{h(p-1)}p>1,
\]
contradicting unitarity. Every $p$-active row of $p$-height
at least two consequently has a tame entry in this column.
By the local Galois identity that row is constant on the
wild orbit of $B$. It has zero coordinate in $W_pv$.

All labels in the wild orbit of $B$ are $p$-inactive, so
$v$ is upper-unipotent fixed. Decompose it, if necessary,
into full torus characters. Each character occurring has
conductor $p^2$. Applying two-depth compression with
$\beta=2,c=0,d=1$ to these orthogonal components gives
\[
                      \|P_{U_1}W_pv\|^2\le1/p.
\]
Thus at least $1-1/p$ lies at actual $p$-orders at least
$p^2$, where the preceding exclusion forces $p$-height
at most one. The local $p$-Weyl operator commutes with
the complementary actual twists. Those twists are constant
on the wild $p$-orbit of $B$, proving the assertions at
$q,r$.
\end{proof}

\begin{corollary}[A second ramified prime forces mixed activity]
\label{paper:p-wild-q-ramification}
Under Proposition~\ref{paper:maximal-p-wild-detector}, suppose
$K_B$ is also ramified at $q$. Then some actual twist has
$p$-part of order at least $p^2$ and nontrivial $q$-part,
while its $r$-part equals $\theta_{B,r}$. In particular,
if $B$ is $r$-active then a triple-active twist exists.
\end{corollary}
\begin{proof}
Let $P$ be the orthogonal projection onto the full local
$q$-group invariants and set $\rho=\langle e_B,Pe_B\rangle$.
The $q$-torus orbit has size $s\ge2$, so $\rho\le1/s\le1/2$.
The projector $P$ commutes with wild $p$-inertia; consequently
its trace on the coordinate space of the wild $p$-orbit of
$B$ is $p\rho$. Let $v_1,\ldots,v_{p-1}$ be the norm-one
nontrivial Fourier vectors on that orbit.

Let $Q$ project onto the coordinates of $q$-height three.
These coordinate lines are fixed by the full local $q$-group,
so $Q\le P$. Since $W_p$ commutes with $P$,
\[
 \sum_{j=1}^{p-1}\|QW_pv_j\|^2
 \le\sum_{j=1}^{p-1}\|Pv_j\|^2\le p\rho\le p/2.
\]
The total squared mass of their Weyl images at actual
$p$-orders at least $p^2$ is at least
$(p-1)(1-1/p)>p/2$, because $p\ge5$. Some such coordinate
therefore has $q$-height at most two. It retains
$q$-inactivity and the $r$-twist of $B$.

Its actual $p,r$-twist fiber is invariant under the full
local $q$-group. If the entire fiber were $q$-inactive,
that group would act trivially there, and the Weyl-line
lemma would force $q$-unit codegree at the selected
coordinate. This contradicts its $q$-height at most two.
Thus the same fiber has a $q$-active label, as required.
\end{proof}


These statements concern mass-$r$ quartic seeds. The subsequent
mass-$p^2$ crossings have their own hypotheses; neither construction
forces maximal seed heights.


\begin{theorem}[The lower square-height inactive crossing]
\label{paper:lower-inactive-crossing}
In Lemma~\ref{paper:quartic-arithmetic}, take $(i,j)=(2,2)$
and suppose both inputs are $r$-inactive. Retain
$a=[K_X:K]$, $b=[K_B:K]$ and the scalar exponent $\alpha$.
Then the actual $r$-phases have order $r$, and all input
$q$-twists and actual $q$-phases have order dividing $q$.
More precisely,
\begin{equation}\label{paper:eq:lower-inactive-budget}
                   \beta=1,\qquad \alpha\in\{0,1\},
                   \qquad abq^\alpha\le p^2.
\end{equation}
The following two consequences apply.
\begin{enumerate}[label=\textup{(\roman*)},leftmargin=*]
\item If $\alpha=1$, both inputs are $p$-inactive and at
least one is $q$-active. Both complete endpoint fields have
conductor dividing $pqr$. Every input twist and every
constituent twist of $X\otimes B$ has order dividing $qr$.
The full $p$-block of $X$ is
nonprincipal and homogeneous of height two; its labels
are all $r$-inactive.
\item If both inputs are $p$-active, then $\alpha=0$,
$a=b=p-1$, both $p$-twists have order $p$, and both inputs
are $q$-inactive. Their complete field conductors divide
$pqr$. If instead one input has $p$-twist order $p^2$,
then $\alpha=0$, its relative index is $p(p-1)$ and the
other index is one; both inputs are $q$-inactive. The
active endpoint has conductor dividing $p^2qr$ and the
other endpoint field is exactly $K$.

In both cases of \textup{(ii)}, the actual crossing is a
single period:
\begin{equation}\label{paper:eq:lower-period}
 z=(r-qd)+q\sum_{u\in H}\zeta_r^{vu},
 \qquad H\le\mathbb F_r^\times,\quad |H|=d,
 \quad v\ne0.
\end{equation}
No equality of an endpoint field with a cyclotomic
compositum is implicit in these conductor bounds.
\end{enumerate}
\end{theorem}
\begin{proof}
All surviving odd parameters have $p\ge5$ and $q\ge7$.
The inactive channel inequality gives
\[
                 abq^\alpha r^{2\beta-2}\le p^2.
\]
Thus $\beta=1$. Let $A$ be the largest $q$-exponent
among the input twists and the actual phases. The relative
inertia bound used in Theorem~\ref{paper:quartic-joint-twists}
gives $abq^\alpha\ge(1-1/q)q^A$. If $A\ge2$, this is
at least $q(q-1)>p^2$, a contradiction. This proves
\eqref{paper:eq:lower-inactive-budget} and the twist bounds.

Suppose $\alpha=1$. An active $p$-input would contribute
$p-1$ to $ab$, whereas $q(p-1)>p^2$. Both are therefore
$p$-inactive. If they were also $q$-inactive, apply the
two inactive local compressions to the primitive $q,r$
Fourier channel of \eqref{paper:eq:quartic-fourier}.
Its normalized squared norm is at least $abq/(p^2r)$,
but the simultaneous upper bound is $1/(qr)$. This gives
$abq^2\le p^2$, impossible. Hence an input is $q$-active.

Furthermore $ab\le p^2/q<p$. Wild inertia at any of
$p,q,r$ fixes $K$, whose conductor is $qr$; a nontrivial
image on an endpoint would make that prime divide its
relative index. This is impossible, proving the full-field
conductor bounds. Balancing, the actual phase orders and
input $p$-inactivity give the full twist bound.
A $q$-active input has $p$-divisible
codegree. By Lemma~\ref{paper:family:lw:weyl-line}, $p$ and $q$
cannot therefore be globally separated. An actual $pq$-mixed
twist exists, so Corollary~\ref{paper:two-neighbors} supplies
an $r$-free witness $Y\in P_q$ of $p$-height two.
Pairing $Y$ with $X$ and applying
Proposition~\ref{paper:family:homogeneous-crossing} proves
homogeneity and nonprincipality of the full $p$-block of
$X$. Its common complementary twist has trivial $r$-part.

If both inputs are $p$-active, the relative indices are
each divisible by $p-1$. Since
$p^2/(p-1)^2\le25/16<2$, their product must equal
$(p-1)^2$ and hence $a=b=p-1$. Equation
\eqref{paper:eq:lower-inactive-budget} now forces $\alpha=0$.
Neither index can afford a higher $p$-twist or the extra factor $q-1$ from a
$q$-active input; the joint stabilizer bound applies
when these costs occur at the same endpoint.
If one input instead has twist order $p^2$, its index
is divisible by $p(p-1)$, and the same budget gives
that exact index and index one at the other endpoint,
again with $\alpha=0$.
The same argument excludes $q$-activity and higher
$p$-twist orders.

In either situation let $e=(r-1)/d$, and let $m$ be
the nonnegative integral coefficient vector on the
$H$-cosets of nonzero $r$-exponents. Its mass is
$M=t/d\ge1$, and $r=h+qdM$ gives $e\ge qM$.
Every nontrivial Fourier coefficient satisfies
$|\widehat m|^2\le C$, where
\[
 C=\frac{p^2}{ab}\le\frac{25}{16}.
\]
Parseval gives $M(e-M)\le C(e-1)$. If $M\ge2$,
then $e\ge2q\ge14$ and $M\le e/q<e/2$, whence
\[
 \frac{M(e-M)}{e-1}\ge\frac{2(e-2)}{e-1}
       \ge\frac{24}{13}>\frac{25}{16}.
\]
Thus $M=1$, proving \eqref{paper:eq:lower-period}.

Finally, no wild inertia at $q$ or $r$ can occur in
an extension of $K$ of degree $p-1$ or $p(p-1)$.
At $p$, the full active inertia already has degree
$\varphi(p^k)$, with $k=1$ or $2$ as above. A larger
conductor would require an additional factor $p$.
This proves the remaining conductor statements.
\end{proof}

\begin{corollary}[A quadratic phase grid]
\label{paper:lower-phase-grid}
In case \textup{(i)} of Theorem~\ref{paper:lower-inactive-crossing},
write
\[
 \theta_{X,q}=\zeta_q^{\kappa_X},\qquad
 \theta_{B,q}=\zeta_q^{\kappa_B},\qquad \eta=\zeta_q^{\varepsilon},
 \qquad \varepsilon\ne0.
\]
For $a,b\in\mathbb F_q^\times$, let $X_a,B_b$ denote
the translates under the $q$-inertia automorphisms of
exponents $a,b$. Every constituent of $X_a\otimes B_b$
has $q$-twist
\[
             \zeta_q^{Q(a,b)},\qquad
             Q(a,b)=\kappa_Xa^2+\varepsilon ab+\kappa_Bb^2.
\]
Consequently products with different values of $Q$ have
disjoint simple supports. Here at least one of $\kappa_X,\kappa_B$ is
nonzero.
\end{corollary}
\begin{proof}
Argument-side Galois covariance gives
$S_{X_aB_b}/(p^2q^2r)=\eta^{ab}w$.
Uniqueness of the actual word in
Lemma~\ref{paper:quartic-arithmetic} identifies its
$q$-scalar. The square Galois action on the two input
twists and the balancing relation now give the formula.
Distinct $q$-twists cannot occur at the same simple label.
\end{proof}

\begin{proposition}[Fusion transitions in a quadratic phase grid]
\label{paper:quadratic-grid-transitions}
Let $q\ge5$ be prime, and let $\{X_a\}_{a\in\mathbb F_q^\times}$
and $\{Y_b\}_{b\in\mathbb F_q^\times}$ be two full local
$q$-inertia orbits, indexed so that $\sigma_hX_a=X_{ha}$ and
$\sigma_hY_b=Y_{hb}$. Suppose every simple constituent of
$X_aY_b$ has actual $q$-primary twist
\[
 \zeta_q^{Q(a,b)},\qquad Q(a,b)=A a^2+Eab+B b^2,
 \qquad E\ne0,\quad (A,B)\ne(0,0).
\]
For a simple $Z$ put
\[
 M_Z(a,c)=N_{ZX_c}^{X_a},\qquad
 L_Z(b,d)=N_{Z^*Y_d}^{Y_b},
\]
and let $H_Z\le\mathbb F_q^\times$ be its stabilizer under the
tame local inertia subgroup (using Teichm\"uller lifts at the
ambient $q$-power level). If $|H_Z|\ge3$, then every nonzero product
$M_Z(a,c)L_Z(b,d)$ has
\[
                 (c,d)=(a,b)\quad\text{or}\quad(-a,-b).
\]
If $A=0$ or $B=0$, the same assertion holds when
$|H_Z|\ge2$. Consequently any other common transition forces
a local $q$-inertia orbit of size at least $(q-1)/2$ for $Z$,
or of size $q-1$ when one input orbit is $q$-inactive.

If $Z$ is fixed by tame local $q$-inertia and both rectangles
are nonzero, then they are either $(kI,lI)$ or $(kC,lC)$,
where $C$ is the permutation $a\mapsto-a$ and $k,l$ are positive
integers. In a perfect fusion-simple category, a nonunit $Z$
has $k,l<d_Z$.
\end{proposition}
\begin{proof}
Frobenius reciprocity gives the nonnegative Gram expansion
\[
 \langle X_aY_b,X_cY_d\rangle
       =\sum_Z M_Z(a,c)L_Z(b,d).
\]
The left side is zero when $Q(a,b)\ne Q(c,d)$, so every
summand vanishes separately there. For $h\in H_Z$, fusion
covariance preserves a nonzero $M_Z(a,c)$ after replacing
$(a,c)$ by $(ha,hc)$. Keeping the other rectangle fixed gives
\[
 A(a^2-c^2)h^2+E(ab-cd)h+B(b^2-d^2)=0
                       \qquad(h\in H_Z).
\]
Three distinct roots make this polynomial identically zero.
Since $E\ne0$ and one of $A,B$ is nonzero, this forces the two
displayed possibilities. If $A=0$, the polynomial has degree
at most one; if $B=0$, remove its nonzero factor $h$ first.
This proves the stronger stabilizer assertion and the orbit bounds.

When $H_Z=\mathbb F_q^\times$, any edge of one nonzero
rectangle forces every edge of the other to have the same
diagonal or opposite-index type. A mixture of the two types
is impossible. Galois covariance makes all entries along the
chosen permutation equal. Finally, $k\le d_Z$ follows from
the dimension equation. Equality leaves no constituent of
$ZX_a$ outside the $X$-orbit and makes its positive orbit sum
a Perron eigenvector on a proper support, contrary to fusion
simplicity. Apply the same argument to $Z^*$ and the $Y$-orbit.
\end{proof}

\begin{corollary}[Reciprocity constraints on the phase grid]
\label{paper:quadratic-reciprocity}
Under Corollary~\ref{paper:lower-phase-grid}, put
\[
 M(a,b;c,d)=\langle X_aB_b,X_cB_d\rangle.
\]
If this nonnegative integer is nonzero, then
\[
 \begin{split}
 (a+c)\bigl(2\kappa_X(a-c)+\varepsilon(b-d)\bigr)&=0,\\
 (b+d)\bigl(\varepsilon(a-c)+2\kappa_B(b-d)\bigr)&=0
 \end{split}
 \qquad\text{in }\mathbb F_q.
\]
In particular, if $4\kappa_X\kappa_B-\varepsilon^2\ne0$,
$a+c\ne0$ and $b+d\ne0$, then nonzero overlap forces
$a=c$ and $b=d$.
\end{corollary}
\begin{proof}
The phase grid first requires
$Q(a,b)=Q(c,d)$.
Rigidity and commutativity give
\[
 \langle X_aB_b,X_cB_d\rangle
       =\langle X_aB_d^*,X_cB_b^*\rangle.
\]
The actual word at $X_a,B_d^*$ has $q$-scalar
$\eta^{-ad}$, while the dual has the same twist as the
original label. Thus nonzero overlap also requires
\[
 \kappa_Xa^2-\varepsilon ad+\kappa_Bd^2
       =\kappa_Xc^2-\varepsilon cb+\kappa_Bb^2.
\]
Adding and subtracting these two equalities gives the
factored equations. The final assertion follows by inverting
the indicated two-by-two coefficient matrix.
\end{proof}

\begin{corollary}[Every mass-$p^2$ compulsory crossing has a unit]
\label{paper:cor:proper-crossing}
In an ordinary perfect fusion-simple integral modular category
at the ordered exact family, every compulsory complementary
crossing $z$ of mass $p^2$, as in
Proposition~\ref{paper:family:family:crossings}, is a unit at some place over $p$.
For $p\ge5$, its actual balancing word has no nontrivial
$p$-primary components.  Consequently, for $p\ge5$, every simple constituent $T$
of the corresponding product $X\otimes Y$ satisfies
\[
                    \theta_{T,p}=\theta_{X,p}\theta_{Y,p}.
\]
For $p\ge5$ and pure-$q$ and pure-$r$ witnesses, every such constituent
is $p$-inactive.  This assertion concerns these specified products;
it does not assert fusion closure of all $p$-inactive labels.
\end{corollary}
\begin{proof}
The crossing is integral, has strict house below $p^2$, and
has residue $p^2$ at every place over $q,r$.
The case $p=3$ is already excluded by
Corollary~\ref{paper:family:cor:remaining-region}, since then $m\le1$.
At $p=2$, vanishing at every two-place would give $z=4$ by
Proposition~\ref{paper:prop:dyadic-three}, contrary to strictness.

For $p\ge5$, Corollary~\ref{paper:family:cor:remaining-region} excludes
$m=\lfloor\min(r,p^2)/q\rfloor\le2$.  Thus any surviving
category has $r\ge3q$.
Lemma~\ref{paper:lem:large-ratio} and weighted concentration force
all actual $p$-primary components of this word to be trivial.
If $z$ vanished at every $p$-place, the unramified $q,r$-field
would make $z/p$ integral.  Its positive rational mass is $p$,
and its residues at $q,r$ are $p$.
Theorem~\ref{paper:family:thm:two-positive}, with $p<q<r$, would give
$z/p=p$.  In every case this contradicts strictness.
Finally, every actual constituent has positive coefficient in
the balancing word, with phase $\theta_T/(\theta_X\theta_Y)$.
Its trivial $p$-primary part gives the displayed twist identity.
\end{proof}

\begin{corollary}[Integral phase fibers of the compulsory crossings]
\label{paper:phase-fibers}
Suppose an ordinary perfect fusion-simple integral modular category
lies at the ordered exact family and $p\ge5$. Take a compulsory
opposite pair $X,Y$, put
\[
 L=\lcm(d_X,d_Y),\qquad z=S_{XY}/L,
\]
and suppose either
\begin{enumerate}[label=(\roman*)]
\item $\min(v_p(d_X),v_p(d_Y))=2$, and put $e=0$; or
\item $v_p(d_X)=v_p(d_Y)=3$, and $z$ vanishes at every place over
      $p$, and put $e=1$.
\end{enumerate}
Then the actual balancing word for $z/p^e$ has mass $p^2$, no
nontrivial $p$-primary component, and unique nonnegative integer
combined coefficients $n_{\eta,\xi}$. In particular, if
\[
 F_{\eta,\xi}=
 \sum_{\substack{T:\\
  (\theta_T/(\theta_X\theta_Y))_q=\eta,\\
  (\theta_T/(\theta_X\theta_Y))_r=\xi}}
                   N_{XY}^{T}d_T,
\]
then
\[
 F_{\eta,\xi}=p^eL\,n_{\eta,\xi}.
\]
For every $\eta\ne1$ and $\xi\ne1$, respectively,
\[
 rp^eL\mid\sum_\xi F_{\eta,\xi},\qquad
 qp^eL\mid\sum_\eta F_{\eta,\xi}.
\]
At most $1+\lfloor p^2/r\rfloor$ different $q$-components and
$1+\lfloor p^2/q\rfloor$ different $r$-components occur in the
actual balancing phases. If $d_X=p^i q^j$ and $d_Y=p^h r$ as
in Proposition~\ref{paper:family:family:crossings}, these components have orders
dividing $q^{6-2j}$ and $r^2$, respectively.
\end{corollary}
\begin{proof}
The family support criterion forces $r\ge3q$.
In case (i), Corollary~\ref{paper:cor:proper-crossing} removes every
actual $p$-primary component. In case (ii),
Proposition~\ref{paper:unramified-crossing} places $z$ in a field
unramified at $p$. Its vanishing at all $p$-places makes $z/p$
integral. Dividing the actual balancing expression gives mass
$p^2$ and residues $p^2$ at $q,r$. Weighted concentration
(Theorem~\ref{paper:thm:mass-descent}) again removes every actual
$p$-primary component. Apply
Proposition~\ref{paper:sparse-optimum} to this actual $q,r$-word.
Its coefficient at $(\eta,\xi)$ is precisely
$F_{\eta,\xi}/(p^eL)$, proving the divisibility and support assertions.
Finally $\mathbb Q(z/p^e)=\mathbb Q(z)\subseteq K_X\cap K_Y$.
The codegree conductor bound gives conductor exponents at most
$6-2j$ at $q$ and two at $r$ in this intersection.
Corollary~\ref{paper:phase-conductors} transfers these bounds to
the orders of the actual phases.
\end{proof}

\begin{corollary}[The first residue layer of a height-three crossing]
\label{paper:top-dichotomy}
For a compulsory pair of $p$-height three in the surviving odd
family, exactly one of the following alternatives holds:
\begin{enumerate}[label=(\roman*)]
\item $z$ is a unit at some place over $p$;
\item $z/p$ is integral and a unit at some place over $p$, and
      every constituent $T$ of $X\otimes Y$ is $p$-inactive.
\end{enumerate}
In the second case the integral phase-fiber and sparse-support
conclusions of Corollary~\ref{paper:phase-fibers} hold with $e=1$.
Any $p$-active constituent therefore forces alternative (i).
\end{corollary}
\begin{proof}
The unramification and valuation bound in
Proposition~\ref{paper:unramified-crossing} say that the least
normalized $p$-valuation of $z$ is either zero or one.
In the latter case Corollary~\ref{paper:phase-fibers} applies and
gives the twist assertion.  Both inputs have $p$-unit codegree,
so Lemma~\ref{paper:family:lw:weyl-line} makes their $p$-primary twists trivial.
Thus every constituent is $p$-inactive.  A least valuation of one means exactly
that $z/p$ is integral and is a unit at some $p$-place.
\end{proof}

\subsection{Galois Fourier bounds and sparse phase supports}
\label{paper:sec:fourier}


\begin{lemma}[Field indices and the crossing-energy budget]
\label{paper:crossing-budget}\label{paper:two-sided-budget}
\label{paper:crossing-energy}
Let $X,Y$ be nonunits in an ordinary integral modular category,
let $c>0$ be rational, and suppose $w=S_{XY}/c$ is integral.
Put
\[
 K=\mathbb Q(w),\qquad a=[K_X:K],\qquad b=[K_Y:K].
\]
Then
\begin{equation}\label{paper:two-sided-energy}
                             abE(w)<D/c^2.
\end{equation}
If $D$ is square and the full Galois orbits of $X,Y$ are
disjoint, then also
\begin{equation}\label{paper:indexed-global-energy}
 c^2E(w)+\frac{b d_X^2+a d_Y^2}{(a+b)^2}
                         \le\frac{3D}{2(a+b)}.
\end{equation}
Consequently
\begin{equation}\label{paper:three-quarter-energy}
 E(w)<\min\left\{\frac1{ab},\frac3{2(a+b)}\right\}
                          \frac D{c^2}\le\frac{3D}{4c^2}.
\end{equation}
This improvement concerns total energy. The Fourier-channel
bounds below retain their stated constants.

At either label, the number of actual twist values in the orbit
under the full Galois group fixing $w$ divides the corresponding
field index. In odd dimension, if $w$ is real, that index is at
least twice this number. The same assertions hold after retaining
only a primary component of the twists. In odd square dimension,
the stronger joint divisor in \eqref{paper:conductor-weight},
with $L=K$, divides each corresponding field index; no reality
hypothesis on $w$ is needed for that divisor.

The compulsory endpoints have different degrees and hence
disjoint full orbits. Thus for the words of
Corollary~\ref{paper:phase-fibers}, with division exponent
$\varepsilon\ge0$, one has the strengthened family budget
\begin{equation}\label{paper:scaled-crossing-budget}
 E(z/p^\varepsilon)<\frac{3D}{4p^{2\varepsilon}L^2}
       =\frac34p^{6-2\max(i,h)-2\varepsilon}q^{6-2j}r^2.
\end{equation}
\end{lemma}
\begin{proof}
Work in a common cyclotomic splitting field and let $H$ fix $w$.
Its orbit of $X$ has size $a$. Galois covariance and commutativity
make those rows agree on the full Galois orbit of $Y$.
Their sum divided by $\sqrt a$ has squared norm $D$, so
\[
 a\sum_{Y'\in\operatorname{Orb}(Y)}|S_{XY'}|^2
                     \le D-a d_X^2<D.
\]
The sum equals $c^2bE(w)$, proving
\eqref{paper:two-sided-energy}.

Under the additional hypotheses, let $P,Q$ project onto the
two full input orbits and put $e=[K:\mathbb Q]$.
Their ranks are $ae,be$, and Galois covariance gives
\[
                     \|PsQ\|_{\mathrm{HS}}^2
                          =\frac{abe c^2}{D}E(w).
\]
The projections commute with $T/\gamma$ and with charge
conjugation $s^2$. Apply Lemma~\ref{paper:braid-energy} to
$W=s$, $U=T/\gamma$. The vacuum lies outside both orbits
and contributes at least $ae d_X^2/D$ and $be d_Y^2/D$
to the two leakages. Substitution gives
\eqref{paper:indexed-global-energy}. Its positive correction,
together with \eqref{paper:two-sided-energy}, proves
\eqref{paper:three-quarter-energy}.

Twists, and their primary components, are equivariant images of
the label orbit. Equal fiber sizes give the divisibility assertion.
If $w$ is real, complex conjugation belongs to $H$. In odd
dimension it pairs every nonunit label with its distinct dual,
and the duals have the same actual twist; each twist fiber
therefore has at least two labels. The stronger joint divisor
is Corollary~\ref{paper:active-inertia}. Finally use
$c=p^\varepsilon L$ for the family formula.
\end{proof}

\begin{theorem}[Fourier channels of an inactive crossing]
\label{paper:inactive-fourier-budget}
Let $X,Y$ be nonunits in an ordinary integral modular category
of square dimension $D$. Let $c>0$ be rational and put
\[
\begin{gathered}
 w=S_{XY}/c,\qquad K=\mathbb Q(w),\qquad
 \Gamma=\operatorname{Gal}(K/\mathbb Q),\\
 e=|\Gamma|,\qquad a=[K_X:K],\qquad b=[K_Y:K].
\end{gathered}
\]
For a character $\chi$ of $\Gamma$, define
\[
       \widehat w(\chi)=\sum_{\sigma\in\Gamma}
                                  \chi(\sigma)\sigma(w).
\]
Let $u_\chi$ be the vector on the full Galois orbit of $Y$
whose coefficient at each label over $\tau\in\Gamma$ is
$\chi(\tau)/\sqrt{be}$. It has norm one.
Every channel satisfies
\begin{equation}\label{paper:ordinary-channel}
                         ab|\widehat w(\chi)|^2\le D/c^2.
\end{equation}
Let $A$ be a set of odd primes where both inputs are inactive,
and suppose the inflation of $\chi$ to full $\ell$-inertia
has nontrivial conductor $\ell^{u_\ell}$ for every $\ell\in A$.
Using the local congruence kernels of
Lemma~\ref{paper:exact-local-compression}, put
\[
 \Lambda_A(\chi)=\prod_{\ell\in A}\ell^{u_\ell},
 \qquad Q_A(\chi)=\prod_{\ell\in A}P_{K_{\ell,u_\ell}}.
\]
Then
\begin{equation}\label{paper:inactive-channel}
 ab|\widehat w(\chi)|^2
 \le\frac{D}{c^2\Lambda_A(\chi)}
               \|Q_A(\chi)u_\chi\|^2
 \le\frac{D}{c^2\Lambda_A(\chi)}.
\end{equation}
For $A=\varnothing$, the products are one and the identity,
respectively, and this includes the ordinary channel bound.
\end{theorem}
\begin{proof}
Work in a common cyclotomic splitting field. The full Galois
orbits of $X$ and $Y$ map onto $\Gamma$, with fibers of sizes
$a$ and $b$. Unsigned covariance gives
\[
                        S_{\sigma X,\tau Y}=c\,\sigma\tau(w).
\]
At each label of the $X$-orbit over $\sigma$, the vector
$Su_\chi$ therefore has coordinate
\[
        c\sqrt{b/e}\,\chi(\sigma)^{-1}\widehat w(\chi).
\]
Writing $s=S/\sqrt D$, we obtain the exact normalization
\begin{equation}\label{paper:orbit-projection}
 \|P_{\operatorname{Orb}(X)}s u_\chi\|^2
                      =\frac{ab c^2}{D}|\widehat w(\chi)|^2.
\end{equation}
Unitarity gives \eqref{paper:ordinary-channel}.

For $\ell\in A$, every label in the $Y$-orbit has trivial
actual $\ell$-twist, so the local upper unipotents fix
$u_\chi$. The local torus permutes its $K$-fibers through
its restriction to $\Gamma$ and acts on $u_\chi$ by the
inverse of the inflated character $\chi$. This character has
conductor $\ell^{u_\ell}$, regardless of the fiber size $b$.

Factor $s=(\prod_{\ell\in A}W_\ell)W_{A'}$. The complementary
product includes any dyadic anomaly factor and commutes with
the selected local groups. Thus $v=W_{A'}u_\chi$ retains
their fixedness and torus characters. Every selected upper
operator is the actual primary twist operator, since the anomaly
is dyadic. The $X$-orbit coordinate projection is dominated by
$\prod_{\ell\in A}P_{U_\ell}$. Therefore
\[
 \|P_{\operatorname{Orb}(X)}s u_\chi\|^2
 \le\left\|\prod_{\ell\in A}(P_{U_\ell}W_\ell)v\right\|^2
 =\Lambda_A(\chi)^{-1}\|Q_A(\chi)u_\chi\|^2.
\]
The equality follows by applying the exact local compression
successively. The congruence-kernel projectors commute with
their own groups by normality and with every complementary
local factor. Combining this with
\eqref{paper:orbit-projection} proves the result.
\end{proof}

\subsubsection{Integral coefficients and mixed support}

\begin{lemma}[Integral Fourier nonvanishing]\label{paper:integral-fourier}
Let $\Gamma$ be a finite abelian group, let $f:\Gamma\to\mathbb Z$,
and put
\[
 \widehat f(\chi)=\sum_{\gamma\in\Gamma}f(\gamma)\chi(\gamma)^{-1}.
\]
Let $\mathcal X$ be a set of characters closed under
$\chi\mapsto\chi^k$ whenever $k$ is coprime to the order of $\chi$.
If $\widehat f(\chi)\ne0$ for some $\chi\in\mathcal X$, then
\[
                 |\widehat f(\psi)|\ge1
                 \qquad\hbox{for some }\psi\in\mathcal X.
\]
In particular this applies when $\mathcal X$ is specified by
nontriviality on designated subgroups, or by exact local
conductor exponents.
\end{lemma}
\begin{proof}
If $\chi$ has order $m$, then $\widehat f(\chi)$ is an algebraic
integer in $\mathbb Q(\zeta_m)$.  Its conjugates are
$\widehat f(\chi^k)$ for $(k,m)=1$.  A nonzero algebraic integer
has nonzero integral norm, so at least one conjugate has
absolute value at least one.  Character powers with exponent
coprime to $m$ preserve the character's kernel, hence preserve
the indicated subgroup and conductor conditions.
\end{proof}

\begin{lemma}[Sharp integral Fourier rigidity]
\label{paper:integral-fourier-rigidity}
Let $\Gamma$ be a finite abelian group of order $e\ge2$,
and let $m:\Gamma\to\mathbb Z$. Then
$|\widehat m(\chi)|\le1$ for every nontrivial character
if and only if, for some integer $a$ and $g\in\Gamma$,
\[
                   m=a\mathbf1,\qquad
                   m=a\mathbf1+\delta_g,
             \quad\text{or}\quad m=a\mathbf1-\delta_g.
\]
In particular, if $m\ge0$ and
$1\le\sum_gm(g)<e-1$, then $m=\delta_g$ for some $g$.
\end{lemma}
\begin{proof}
Writing $M=\sum_gm(g)$, Parseval gives
\[
 \sum_{g<g'}(m(g)-m(g'))^2
       =e\sum_gm(g)^2-M^2
       =\sum_{\chi\ne1}|\widehat m(\chi)|^2\le e-1.
\]
If the maximum and minimum differ by at least two, their
pair contributes at least four, and the two pairs joining
any other coordinate to them contribute at least two.
The total is at least $2e$, a contradiction. A nonconstant
vector therefore takes adjacent integer values. If the
higher value occurs $k$ times, its variance is
$k(e-k)\le e-1$, forcing $k=1$ or $k=e-1$.
This proves the classification; its converse follows
because a singleton has Fourier coefficients of modulus one.
The final mass condition excludes a nonzero constant
baseline and the complementary singleton pattern.
\end{proof}

\begin{corollary}[The cost of an inactive primary word]
\label{paper:prime-inactive-cost}\label{paper:primary-inactive-cost}
In Theorem~\ref{paper:inactive-fourier-budget}, suppose
$K$ has conductor exactly $\ell^\beta$, where $\ell$ is odd,
$\theta_{X,\ell}=\theta_{Y,\ell}=1$, and
\[
 w=h+t\sum_{\xi\in\mu_{\ell^\beta}\setminus\{1\}}k_\xi\xi,
 \qquad h\in\mathbb Z,\quad k_\xi\in\mathbb Z_{\ge0},
 \quad t\in\mathbb Z_{>0},\quad 0<\sum_\xi k_\xi<\ell.
\]
Assume that the coefficient vector is fixed by
$H=\operatorname{Gal}(\mathbb Q(\zeta_{\ell^\beta})/K)$.
Then
\begin{equation}\label{paper:prime-inactive-energy}
                         ab\,t^2\ell^{2\beta}\le D/c^2.
\end{equation}
The displayed word and coefficient-invariance hypotheses hold for the actual
one-primary words of
Corollary~\ref{paper:phase-fibers}, with $\ell$ its nontrivial
phase prime and $t$ the other prime among $q,r$.
Endpoint inactivity remains an additional assumption.
\end{corollary}
\begin{proof}
Let $d=|H|$. There is an occurring primitive $\ell^\beta$
root, since the value has conductor $\ell^\beta$.
The coefficient vector is invariant under $H$, whose action
on primitive roots is free. If $H$ had nontrivial wild part,
its last order-$\ell$ subgroup would force at least $\ell$
nonzero primitive coefficients, contrary to the mass bound.
Thus $H$ is tame.

Put $\Gamma=(\mathbb Z/\ell^\beta)^\times/H$, with
$e=|\Gamma|$, and let $m_C$ be the common coefficient on
its primitive-root coset $C$. These are nonnegative integers,
not all zero, with $\sum_Cm_C<\ell$.
For $\beta=1$ the vector $m$ is nonconstant, since a constant
vector would make $w$ rational. For $\beta\ge2$, let $W$
be the image in $\Gamma$ of the last wild subgroup. It has
order $\ell$. The vector $m$ cannot be constant on every
$W$-orbit, since a nonzero such orbit would have coefficient
mass at least $\ell$.

Call a character of $\Gamma$ primitive if its inflation has
conductor $\ell^\beta$. At $\beta=1$ these are the nontrivial
characters. At $\beta\ge2$ they are exactly the characters
nontrivial on $W$. Lower-order roots have zero Fourier
coefficient in these channels because the last wild subgroup
fixes them. For each primitive character $\chi$, the remaining
primitive-root terms give
\[
 \widehat w(\chi)=t\,\tau(\chi)\widehat m(\chi),
 \qquad \widehat m(\chi)=\sum_Cm_C\chi(C)^{-1},
 \qquad |\tau(\chi)|^2=\ell^\beta.
\]
For completeness, the primitive Gauss-sum norm follows by
additive Fourier orthogonality: the transform of a primitive
multiplicative character vanishes on nonunits, has constant
modulus on units, and has total squared norm
$\ell^\beta\varphi(\ell^\beta)$.
The conductor-sensitive channel bound therefore gives
\[
 ab\,t^2\ell^{2\beta}|\widehat m(\chi)|^2\le D/c^2
                  \qquad(\chi\text{ primitive}).
\]

Fourier orthogonality supplies a nonzero primitive coefficient:
for $\beta=1$ the integer vector $m$ is nonconstant, and for
$\beta\ge2$ it is nonconstant on a last-wild orbit.
The primitive characters are closed under powers coprime to
their orders. Lemma~\ref{paper:integral-fourier} therefore gives
one with $|\widehat m(\chi)|\ge1$, proving
\eqref{paper:prime-inactive-energy}.

For the stated family words, phase-fiber uniqueness gives
$H$-invariance. Their coefficients have the displayed form
by the one-prime margin divisibilities, and their nonidentity
coefficient mass is less than $\ell$ by
Proposition~\ref{paper:phase-energy}.
\end{proof}

\begin{lemma}[A mixed word at a squared inactive prime]
\label{paper:mixed-wild-cost}
Let $\ell,t$ be distinct odd primes, let $\alpha\ge2$, and let
\[
 w=\sum_{\eta\in\mu_{\ell^\alpha},\,\xi\in\mu_t}
                         n_{\eta,\xi}\eta\xi,
 \qquad n_{\eta,\xi}\in\mathbb Z_{\ge0},
 \qquad \operatorname{cond}\mathbb Q(w)=\ell^\alpha t.
\]
Put
\[
 G=G_\ell\times G_t
   =\operatorname{Gal}(\mathbb Q(\mu_{\ell^\alpha t})/\mathbb Q),
 \quad H=\operatorname{Gal}(\mathbb Q(\mu_{\ell^\alpha t})/
                                        \mathbb Q(w)),
 \quad \Gamma=G/H.
\]
Assume that $H$ fixes the coefficient vector.  Suppose that
fewer than $\ell$ nonidentity $\ell$-primary components and
fewer than $t-1$ nonidentity $t$-components occur, and that
\[
                     t\mid\sum_\xi n_{\eta,\xi}
                      \qquad(\eta\ne1).
\]
Let $a,b\ge1$ and $\mathcal B>0$.  Assume the effective
inactive-prime budget
\begin{equation}\label{paper:wild-effective-budget}
 ab|\widehat w(\chi)|^2\le\mathcal B/\ell^\alpha
\end{equation}
for every character of $\Gamma$ whose inflation to $G_\ell$
has conductor $\ell^\alpha$.  Here
$\widehat w(\chi)=\sum_{\sigma\in\Gamma}\chi(\sigma)\sigma(w)$.
Then
\begin{equation}\label{paper:mixed-wild-cost-eq}
                         ab\,\ell^{2\alpha}t\le\mathcal B.
\end{equation}
In the ordinary crossing budget one may take
$\mathcal B=D/c^2$ whenever both inputs are $\ell$-inactive.
No inactivity hypothesis at $t$ is needed.
\end{lemma}
\begin{proof}
Identify the primitive $\ell^\alpha$ roots with $G_\ell$,
and the nonidentity $t$-roots with $G_t$.  Let $A$ be the
coefficient matrix on this product.

First suppose $A\ne0$.  Let $W\le G_\ell$ be the last wild
subgroup, of order $\ell$.  Some $W$-orbit of rows contains a
nonzero entry.  That orbit has a zero row because fewer than
$\ell$ primitive $\ell^\alpha$ components occur.  The matrix
also has a zero column because fewer than $t-1$ nonidentity
$t$-components occur.

There is consequently a character of $G$ which is nontrivial
on both $W$ and $G_t$ and for which $\widehat A$ is nonzero.
Indeed, if all these Fourier coefficients vanished, then on
each product of a $W$-orbit with $G_t$ the matrix would be a
sum of a row function and a column function.  A zero row and
a zero column make such a matrix zero, contradicting the
chosen nonzero entry.

The matrix $A$ is $H$-invariant, so its Fourier coefficients
vanish on characters not trivial on $H$.  Regard $A$ as the
integer function $m$ on $\Gamma$; then
$\widehat A=|H|\widehat m$ on characters of $\Gamma$.
By Lemma~\ref{paper:integral-fourier}, some character $\chi$ of
$\Gamma$, primitive at $\ell$ and nontrivial at $t$, satisfies
$|\widehat m(\chi)|\ge1$.  Primitive Gauss summation gives
\[
 \widehat w(\chi)
       =\tau_\ell(\chi_\ell)\tau_t(\chi_t)\widehat m(\chi),
 \qquad
 |\tau_\ell(\chi_\ell)\tau_t(\chi_t)|^2=\ell^\alpha t.
\]
Terms with lower $\ell$-order vanish in this channel, as do
the terms with $\xi=1$.  Substitution in
\eqref{paper:wild-effective-budget} proves
\eqref{paper:mixed-wild-cost-eq}.

Now suppose $A=0$.  Every occurring primitive
$\ell^\alpha$ row is concentrated at $\xi=1$, so the margin
condition writes
\[
                         n_{\eta,1}=t k_\eta
             \qquad(\eta\hbox{ primitive of order }\ell^\alpha).
\]
There is such a row because the field has conductor exponent
$\alpha$ at $\ell$.  Write
$H_\ell=\operatorname{pr}_\ell H$ and $d_t=|H\cap G_t|$.
The integer function $k$ on the primitive rows is
$H_\ell$-invariant.  It is nonzero and cannot be constant on
every $W$-orbit: any nonzero constant orbit would occupy
$\ell$ primitive rows.  Hence it has a nonzero primitive
Fourier coefficient on $G_\ell/H_\ell$.  Lemma
\ref{paper:integral-fourier} makes one such coefficient have
modulus at least one.

Inflate the corresponding character to $G$ by making it
trivial on $G_t$.  It is trivial on $H$.  Summing first over
$G_t$ and then over the $H_\ell$-cosets gives
\[
 \widehat w(\chi)
       =t R\,\tau_\ell(\chi_\ell)\widehat k(\chi_\ell),
 \qquad
 R=\frac{t-1}{d_t}\in\mathbb Z_{>0}.
\]
All terms of lower $\ell$-order again vanish.  Consequently
\[
                       ab\,t^2R^2\ell^{2\alpha}\le\mathcal B.
\]
This implies the claimed bound.  The projection $H_\ell$
is automatically tame, since it preserves a nonempty set of
fewer than $\ell$ primitive rows.  In neither case has $H$
been assumed to be a product of its two projections.
\end{proof}

\begin{lemma}[Sparse Fourier estimates]\label{paper:sparse-fourier}
Let $G_1,G_2$ be finite abelian groups of orders $m,n$.
Use the unnormalized Fourier transform.
\begin{enumerate}[label=\textup{(\roman*)},leftmargin=*]
\item If a real function $b$ on $G_1$ is supported on at most
$u$ positions, where $1\le u<m$, and $B=\sum b$, then
\[
 \sum_{\chi\ne1}|\widehat b(\chi)|^2
      =m\sum b^2-B^2
      \ge \frac{m}{u}\left(1-\frac um\right)B^2.
\]
\item If a real matrix $A$ on $G_1\times G_2$ is supported on
at most $u$ rows and $v$ columns, where $1\le u<m$ and
$1\le v<n$. If $C=\sum A$, then
\[
 \sum_{\chi_1\ne1,\,\chi_2\ne1}
          |\widehat A(\chi_1,\chi_2)|^2
 \ge \frac{mn}{uv}
             \left(1-\frac um\right)
             \left(1-\frac vn\right)C^2.
\]
\end{enumerate}
\end{lemma}
\begin{proof}
The first identity is Parseval with the trivial character
removed; Cauchy's inequality gives $\sum b^2\ge B^2/u$.
For the second assertion, let $P_m,P_n$ be the orthogonal
projections away from the constant vectors.  Parseval gives
\[
 \sum_{\chi_1\ne1,\,\chi_2\ne1}|\widehat A|^2
                       =mn\|P_m A P_n\|_{\mathrm F}^2.
\]
For a vector supported on at most $u$ positions,
$\|P_m x\|^2\ge(1-u/m)\|x\|^2$.  Apply this to the columns
of $AP_n$, whose row support is unchanged, and then apply
the analogous estimate to the rows of $A$.  Finally
$\|A\|_{\mathrm F}^2\ge C^2/(uv)$.  These arguments allow
signed entries and do not require separated stabilizers.
\end{proof}

\begin{proposition}[Prime-level mixed support exceeds the inactive budget]
\label{paper:squarefree-mixed-exclusion}
Let $p<q<r$ be odd primes and put $g=p^2$, with
\[
                 r\ge3q,\qquad g\le(q-2)^2.
\]
Suppose
\[
 w=\sum_{\eta\in\mu_q,\,\xi\in\mu_r}n_{\eta,\xi}\eta\xi,
 \qquad n_{\eta,\xi}\in\mathbb Z_{\ge0},\qquad
 \sum_{\eta,\xi}n_{\eta,\xi}=g,
\]
has conductor $qr$, has coefficient vector invariant over
$K=\mathbb Q(w)$, and satisfies
\[
 r\mid\sum_\xi n_{\eta,\xi}\quad(\eta\ne1),\qquad
 q\mid\sum_\eta n_{\eta,\xi}\quad(\xi\ne1).
\]
Let $a,b\ge1$ and let $\mathcal B\le q^4r^2/g$.  It is
impossible that all its nontrivial Fourier channels satisfy
the effective inactive-prime bounds
\begin{equation}\label{paper:mixed-effective-budget}
 ab|\widehat w(\chi)|^2
 \le\frac{\mathcal B}
 {q^{\epsilon_q(\chi)}r^{\epsilon_r(\chi)}},
 \qquad
 \epsilon_\ell(\chi)=
 \begin{cases}
 1,&\chi_\ell\ne1,\\
 0,&\chi_\ell=1.
 \end{cases}
\end{equation}
It is enough to assume these bounds for the joint nontrivial
channels and the $q$-only nontrivial channels.
\end{proposition}
\begin{proof}
Put
\[
 \begin{gathered}
 G=G_q\times G_r,\quad |G_q|=m=q-1,\quad |G_r|=n=r-1,\\
 H=\operatorname{Gal}(\mathbb Q(\mu_{qr})/K),\quad
 h=|H|,\quad
 H_q=\operatorname{pr}_qH,\quad
 h_q=|H_q|,\quad d_r=|H\cap G_r|.
 \end{gathered}
\]
Thus $h=h_qd_r$.  Let $u,v$ be the numbers of nonidentity
$q$- and $r$-components in the actual support.  They are
positive, and the margins imply
\[
 u\le g/r,\qquad v\le g/q.
\]
The two projections of $H$ preserve their respective
nonidentity supports and act freely on them.  Therefore
\begin{equation}\label{paper:mixed-stabilizer-bounds}
                       h\le uv,\qquad d_r\le v.
\end{equation}
No assertion that $H=H_q\times\operatorname{pr}_rH$ is used.
Set
\[
 \kappa_q=1-\frac{u}{q-1},\qquad
 \kappa_r=1-\frac{v}{r-1},\qquad
 T=\sqrt{\frac{\mathcal B}{abq^2r^2}}
                   \le\frac{q}{\sqrt{abg}}.
\]
The numerical assumptions give
\[
 \frac{u}{q-1}
 \le\frac{(q-2)^2}{3q(q-1)}<\frac13,\qquad
 \frac{v}{r-1}
 \le\frac{(q-2)^2}{q(3q-1)}<\frac13.
\]
In particular $\kappa_q,\kappa_r>2/3$.

Let $A$ be the interior coefficient matrix, with
$\eta\ne1,\xi\ne1$, and put
\[
 C=\sum_{\eta\ne1,\,\xi\ne1}n_{\eta,\xi}\ge0.
\]
For a character of $G$ trivial on $H$ and nontrivial on
both factors, the exact Gauss-sum normalization is
\[
 \widehat w(\chi)
       =\frac{\tau_q(\chi_q)\tau_r(\chi_r)}h\,
                            \widehat A(\chi_q,\chi_r),
 \qquad |\tau_q\tau_r|^2=qr.
\]
The joint channel bound in
\eqref{paper:mixed-effective-budget} consequently says
\[
                 |\widehat A(\chi_q,\chi_r)|\le hT.
\]
All other joint coefficients vanish unless the character is
trivial on $H$, because $A$ is $H$-invariant.  There are at
most $mn/h$ such characters.  Lemma
\ref{paper:sparse-fourier} therefore gives
\[
 \frac{mn\kappa_q\kappa_r}{uv}C^2
 \le\sum_{\chi_q\ne1,\,\chi_r\ne1}|\widehat A|^2
 \le\frac{mn}{h}h^2T^2.
\]
Hence
\begin{equation}\label{paper:interior-mass-bound}
 C\le\sqrt{\frac{uvh}{\kappa_q\kappa_r}}\,T
   \le\frac{uvT}{\sqrt{\kappa_q\kappa_r}}
   <\frac32uvT.
\end{equation}

For each $\eta\ne1$, write the row margin as
\[
 \sum_\xi n_{\eta,\xi}=rk_\eta,\qquad
 k_\eta\in\mathbb Z_{\ge0},
 \qquad
 b_\eta=n_{\eta,1}-k_\eta.
\]
Put $k=\sum_{\eta\ne1}k_\eta$ and
$B=\sum_{\eta\ne1}b_\eta$.  Every occupied row has
$k_\eta\ge1$, so $k\ge u$.  Summing the definitions gives the
essential margin identity
\begin{equation}\label{paper:row-margin-identity}
                              (r-1)k=C+B.
\end{equation}
The signed function $b$ is supported on at most $u$
positions and is invariant under $H_q$.

For a nontrivial character $\chi_q$ trivial on $H_q$,
summing the $r$-roots first gives
\[
 \begin{split}
 \widehat w(\chi_q,1)
  &=\frac{\tau_q(\chi_q)}h
        \sum_{\eta\ne1}\chi_q(\eta)^{-1}
       \left((r-1)n_{\eta,1}
                        -\sum_{\xi\ne1}n_{\eta,\xi}\right)\\
  &=\frac{r\tau_q(\chi_q)}h\,\widehat b(\chi_q).
 \end{split}
\]
Here the notation $\chi_q(\eta)$ identifies the primitive
$q$-roots with the corresponding elements of $G_q$.
Since $|\tau_q|^2=q$, the $q$-only channel bound gives
$|\widehat b(\chi_q)|\le hT$.  Its other nontrivial
coefficients vanish unless $\chi_q|_{H_q}=1$; there are at
most $m/h_q=md_r/h$ such characters.  The one-dimensional
estimate of Lemma~\ref{paper:sparse-fourier} gives
\[
 \frac{m\kappa_q}{u}B^2
 \le\sum_{\chi_q\ne1}|\widehat b(\chi_q)|^2
 \le \frac{md_r}{h}h^2T^2.
\]
Consequently
\begin{equation}\label{paper:margin-correction-bound}
 |B|\le\sqrt{\frac{uhd_r}{\kappa_q}}\,T
      \le\frac{uvT}{\sqrt{\kappa_q}}
      <\frac32uvT.
\end{equation}

Combining \eqref{paper:row-margin-identity} with the two
bounds yields
\[
 (r-1)u\le(r-1)k=C+B
                \le C+|B|<3uvT.
\]
Cancel $u>0$ and use $v\le g/q$:
\[
 r-1<3vT
       \le3\sqrt{\frac{g}{ab}}\le3p.
\]
This contradicts $r\ge3q$ and $q>p$.  The proof permits
$C=0$; mixed primary support does not have to contain an
interior term.
\end{proof}

\begin{corollary}[Prime-level mixed support forces activity at the smaller prime]
\label{paper:squarefree-q-activity}
In the word and numerical setting of
Proposition~\ref{paper:squarefree-mixed-exclusion}, suppose
only the $q$-inactive channel bounds are available:
\[
 ab|\widehat w(\chi)|^2\le\mathcal B/q
       \qquad(\chi_q\ne1),\qquad
 \mathcal B\le q^4r^2/g.
\]
Write $d_r=|H\cap G_r|$ as in that proposition.
If $d_r\le ab$, these hypotheses are already impossible.

Consequently, for a crossing in an ordinary integral modular
category of odd square dimension, a mixed word of conductor
$qr$ forces at least one input to have nontrivial actual
$q$-twist, provided the effective budget supplies the
conductor factor at every selected inactive prime.
\end{corollary}
\begin{proof}
Use all notation from the preceding proof.  The $q$-only
channel bound is unchanged.  The joint channels now give
$|\widehat A|\le h\sqrt r\,T$, instead of $hT$.  Thus
\[
 C\le
      \sqrt{\frac{uvhr}{\kappa_q\kappa_r}}\,T,
 \qquad
 |B|\le
      \sqrt{\frac{uhd_r}{\kappa_q}}\,T.
\]
Since $h=h_qd_r$, $h_q\le u$, and $d_r\le ab$, one has
$h\le uab$.  Also $d_r\le v\le g/q$.  Substitute
$T\le q/\sqrt{abg}$ and $v\le g/q$ to obtain
\[
 \frac Cu\le\sqrt{\frac{qr}{\kappa_q\kappa_r}},
 \qquad
 \frac{|B|}u\le\sqrt{\frac q{\kappa_q}}.
\]
The margin identity and $k\ge u$ would imply
\begin{equation}\label{paper:q-inactive-impossible}
 r-1\le\sqrt{\frac{qr}{\kappa_q\kappa_r}}
                      +\sqrt{\frac q{\kappa_q}}.
\end{equation}

The nonempty $q$-support gives $3q\le r\le p^2$,
so $p\ge5$ and $q\ge7$.  Both sparsity constants satisfy
the useful uniform bound
\[
 \kappa_q,\kappa_r\ge\kappa_0
                 :=\frac{2(q+1)}{3q}.
\]
Indeed, both $u/(q-1)$ and $v/(r-1)$ are at most
$(q-2)/(3q)$, by $g\le(q-2)^2$ and $r\ge3q$.
After division by $r$, the required inequality would be
\[
 1-\frac1r\le\frac{\sqrt{q/r}}{\kappa_0}
                  +\frac1r\sqrt{\frac q{\kappa_0}}.
\]
Its left side increases with $r$ and its right side decreases.

It therefore suffices to contradict this inequality at
$r=3q$.  After division by $q$, its right side there is
\[
 \frac{3\sqrt3}{2}\frac q{q+1}
              +\sqrt{\frac{3}{2(q+1)}}.
\]
Put $x=(q+1)^{-1/2}$, $A=3\sqrt3/2$ and $B_0=\sqrt{3/2}$.
Completing the square gives
\[
 A-Ax^2+B_0x
       \le A+\frac{B_0^2}{4A}
       =\frac{19\sqrt3}{12}
       <\frac{20}{7}\le3-\frac1q
                      \qquad(q\ge7),
\]
contradicting \eqref{paper:q-inactive-impossible}.

For the categorical assertion, suppose both inputs were
$q$-inactive.  If both were also $r$-inactive, the preceding
proposition applies.  Otherwise an $r$-active input has
an $r$-group as its full $r$-inertia stabilizer, by
Corollary~\ref{paper:active-inertia}.  The tame subgroup of
order $d_r$ fixing $K$ therefore acts freely on that input.
Its relative field index is at least $d_r$, so $ab\ge d_r$.
The argument just proved again gives a contradiction.
\end{proof}

\begin{corollary}[A mixed wild larger prime still forces smaller-prime activity]
\label{paper:wild-q-activity}
Retain the mass, margin, coefficient-invariance and numerical
hypotheses of Proposition~\ref{paper:squarefree-mixed-exclusion},
but allow $r$-components in $\mu_{r^2}$ and suppose
$\operatorname{cond}K=qr^2$.  Put
\[
 G=G_q\times G_{r^2},\quad
 H=\operatorname{Gal}(\mathbb Q(\mu_{qr^2})/K),\quad
 h=|H|,\quad d_r=|H\cap G_{r^2}|.
\]
Suppose the $q$-inactive effective budget holds:
\[
 ab|\widehat w(\chi)|^2\le\mathcal B/q
       \qquad(\chi_q\ne1),\qquad
 \mathcal B\le q^4r^2/g.
\]
If $d_r\le ab$, these hypotheses are impossible.
Consequently a crossing in an ordinary integral modular
category of odd square dimension with this mixed conductor
forces an input active at $q$, as well as an input active
at $r$, provided the effective budget supplies the conductor
factor at every selected inactive prime.
\end{corollary}
\begin{proof}
Let $u$ count nonidentity $q$-components, and let $v_1,v_2$
count $r$-components of exact orders $r,r^2$, respectively.
The margins give
\[
                u\le g/r,\qquad v_1+v_2\le g/q.
\]
Both $u$ and $v_2$ are positive.  In particular
$3q\le r\le g$, so $p\ge5$ and $q\ge7$.
Write $H_q=\operatorname{pr}_qH$, $h_q=|H_q|$.  The projection
of $H$ to $G_{r^2}$ preserves the nonempty primitive
$r^2$-support and acts freely on it.  Its order is therefore
at most $v_2<r$ and has no $r$-part.  Reduction modulo $r$ is injective on this projection.
Consequently the map $G_q\times G_{r^2}\to G_q\times G_r$
is injective on $H$, and its image still has order $h$.
Also
\[
           h=h_qd_r\le uab,\qquad d_r\le v_2\le g/q.
\]

Let $A_1,A_2$ be the interior matrices for $r$-components
of exact orders $r,r^2$, with total masses $C_1,C_2$.
They are indexed by $G_q\times G_r$ and
$G_q\times G_{r^2}$, respectively.  Put
\[
 \begin{gathered}
 \kappa_q=1-\frac{u}{q-1},\qquad
 \kappa_1=1-\frac{v_1}{r-1},\qquad
 \kappa_2=1-\frac{v_2}{r},\\
 T=\sqrt{\frac{\mathcal B}{abq^2r^2}}
                    \le\frac{q}{\sqrt{abg}}.
 \end{gathered}
\]
All three $\kappa$'s are greater than $2/3$.

If $\chi_q\ne1$ and the $r$-character has exact conductor
$r$, Gauss summation gives
\[
 \widehat w(\chi)
       =\frac{r\tau_q(\chi_q)\tau_r(\chi_r)}h
                                            \widehat A_1(\chi),
 \qquad |r\tau_q\tau_r|^2=qr^3.
\]
Primitive $r^2$ terms vanish in these channels.  Thus
$|\widehat A_1|\le hT/\sqrt r$.  The number of potentially
nonzero characters is at most $(q-1)(r-1)/h$, since the
reduced image of $H$ has order $h$.  Sparse Parseval gives
\begin{equation}\label{paper:r-prime-interior}
 C_1\le
       \sqrt{\frac{uv_1h}{\kappa_q\kappa_1}}\,
                         \frac{T}{\sqrt r}.
\end{equation}
If $v_1=0$, this is interpreted as $C_1=0$.

If $\chi_q\ne1$ and the $r$-character has exact conductor
$r^2$, only $A_2$ contributes, and
\[
 \widehat w(\chi)
        =\frac{\tau_q(\chi_q)\tau_{r^2}(\chi_r)}h
                                            \widehat A_2(\chi),
 \qquad |\tau_q\tau_{r^2}|^2=qr^2.
\]
Thus $|\widehat A_2|\le hT$.  Let $W$ be the last wild
subgroup of $G_{r^2}$, of order $r$, and let $P_W^\perp$
remove the constant function on each $W$-coset.  Since
at most $v_2$ columns occur, every such coset has at most
$v_2$ occupied positions.  Therefore
\[
 \|P_q^\perp A_2P_W^\perp\|_{\mathrm F}^2
                 \ge\kappa_q\kappa_2\|A_2\|_{\mathrm F}^2
                 \ge\frac{\kappa_q\kappa_2}{uv_2}C_2^2.
\]
Parseval on $G_q\times G_{r^2}$, and the upper bound
$(q-1)r(r-1)/h$ for the number of surviving characters,
give
\begin{equation}\label{paper:r-wild-interior}
 C_2\le\sqrt{\frac{uv_2h}{\kappa_q\kappa_2}}\,T.
\end{equation}
No product decomposition of $H$ has been used.

For the $q$-only channels define the signed row function
\[
 F_\eta=(r-1)n_{\eta,1}
          -\sum_{\operatorname{ord}(\xi)=r}n_{\eta,\xi},
 \qquad F=\sum_{\eta\ne1}F_\eta.
\]
The Ramanujan sums over $G_{r^2}$ are $r(r-1)$ at $1$,
$-r$ at roots of exact order $r$, and zero at roots of
exact order $r^2$.  Hence the exact normalization is
\[
              \widehat w(\chi_q,1)
                     =\frac{r\tau_q(\chi_q)}h
                                            \widehat F(\chi_q).
\]
It follows that $|\widehat F|\le hT$.  The function is
$H_q$-invariant and supported on at most $u$ positions;
at most $(q-1)d_r/h$ characters can contribute.
The one-dimensional sparse estimate yields
\begin{equation}\label{paper:r-wild-margin}
                  |F|\le\sqrt{\frac{uhd_r}{\kappa_q}}\,T.
\end{equation}

Use $h\le uab$, $v_1,v_2,d_r\le g/q$ and
$T\le q/\sqrt{abg}$ in the three estimates.  They imply
\[
 \frac{C_1}{u}
       \le\sqrt{\frac{q}{r\kappa_q\kappa_1}}<1,\qquad
 \frac{C_2}{u}
       \le\sqrt{\frac{q}{\kappa_q\kappa_2}}<q,\qquad
 \frac{|F|}{u}
       \le\sqrt{\frac q{\kappa_q}}<q.
\]
Here $r\ge3q$, $q\ge7$ and $\kappa_q,\kappa_1,\kappa_2>2/3$
make the displayed strict inequalities immediate.

Let $M_0=\sum_{\eta\ne1}n_{\eta,1}$ and let
$rk=M_0+C_1+C_2$ be the total nonidentity $q$-margin.
Since every occupied row has positive margin divisible by
$r$, one has $k\ge u$.  The definition of $F$ gives
$F=(r-1)M_0-C_1$, and hence
\[
                    rk=\frac{F}{r-1}
                         +\frac{rC_1}{r-1}+C_2.
\]
The preceding bounds would force
\[
 r\le\frac{rk}{u}
    <\frac{q+r}{r-1}+q
     =q+1+\frac{q+1}{r-1}
     <q+2<3q,
\]
contrary to $r\ge3q$.

For the categorical conclusion, inactivity at $r$ would
give $abqr^4\le\mathcal B$ by
Lemma~\ref{paper:mixed-wild-cost}, or $abgr^2\le q^3$,
which is impossible since $g\ge3q$ and $r>q$.
Thus an input is $r$-active.  The tame subgroup of order
$d_r$ fixing $K$ acts freely on that input by
Corollary~\ref{paper:active-inertia}, so $d_r\le ab$.
If both inputs were $q$-inactive, the result just proved
would give a contradiction.
\end{proof}

\begin{corollary}[The activity forced by mixed phase support]
\label{paper:mixed-inactive-exclusion}
In the surviving odd ordered family, let $w=S_{XY}/c$ be
an actual mass-$p^2$ word of
Corollary~\ref{paper:phase-fibers}. Put
\[
 K=\mathbb Q(w),\qquad a=[K_X:K],\qquad b=[K_Y:K].
\]
Suppose both phase supports occur and
\[
                   \operatorname{cond}K\in\{qr,q^2r,qr^2\}.
\]
Assume an effective Fourier budget with
$\mathcal B\le q^4r^2/p^2$: for every choice of odd phase
primes where both inputs are inactive, their nontrivial local
character conductors contribute their product $\Lambda$ in
\[
                  ab|\widehat w(\chi)|^2\le\mathcal B/\Lambda.
\]
Then an input is $q$-active. If the conductor is $qr^2$,
an input is also $r$-active. In particular, a squared phase
prime requires activity there, and the two inputs cannot be
inactive at every phase prime.

If the same conductor-sensitive bounds hold with an
amplification factor $A_0\ge1$,
\[
 A_0ab|\widehat w(\chi)|^2\le D/(c^2\Lambda),
\]
then $\mathcal B=D/(A_0c^2)$ may be used. This includes
the flat amplification proved in
Corollary~\ref{paper:flat-energy}.
\end{corollary}
\begin{proof}
Put $g=p^2$. Uniqueness makes the coefficient vector invariant
over $K$. The margin bounds give at most $g/r$ nonidentity
$q$-components and $g/q$ nonidentity $r$-components, both
smaller than one third of the corresponding prime-level
unit-group sizes. For conductors $qr$ and $qr^2$, apply
Corollaries~\ref{paper:squarefree-q-activity} and
\ref{paper:wild-q-activity}, respectively.

For conductor $q^2r$, suppose both inputs were $q$-inactive.
Lemma~\ref{paper:mixed-wild-cost} would give
\[
                  abq^4r\le\mathcal B\le q^4r^2/g,
                  \qquad\text{hence }abg\le r.
\]
A nonidentity $q$-row has positive margin divisible by $r$,
so $r\le g$. Equality is impossible because $r$ is prime
and $g=p^2$; hence $r<g$, a contradiction.
\end{proof}

\begin{proposition}[Wild crossing with an inactive $r$-input]
\label{paper:wild-mass-crossing}
In the surviving odd ordered exact family, suppose $r>p^2$ and
take a compulsory pair of degrees
\[
 d_X=p^2q^2,\qquad d_Y=p^2r.
\]
Put $g=p^2$, $w=S_{XY}/(p^2q^2r)$ and
$K=\mathbb Q(w)$. Suppose $K$ has conductor $r^2$ and at least one input is
$r$-inactive. Exactly one input is then $r$-active, its $r$-primary
twist has order $r^2$, and both inputs are $q$-inactive. At most
one input is $p$-active; if one is, its $p$-primary twist has
order $p$ and $(p-1)h^2<q$, with $h$ as below. All full input
and product-constituent twists divide $pr^2$.
Write its unique actual word as
\[
 w=h+q\sum_{\xi\ne1}k_\xi\xi,\qquad
 h+qk=g,\qquad k=\sum k_\xi.
\]
Every occurring nonidentity phase has exact order $r^2$.
Let $d=[\mathbb Q(\zeta_{r^2}):K]$ and
$a=[K_X:K]$, $b=[K_Y:K]$ and $u=ab/d$. Then $u$ is a positive integer, and
\[
 d\mid k,\qquad
 u(r-1)^2h^2\le gq^2d\le gq(g-h),\qquad
 uh^2<q.
\]
In particular $2\le h<\sqrt q$, so $h$ is the least positive
residue of $p^2$ modulo $q$ and $k=\lfloor p^2/q\rfloor$.
For pure-$q$ and pure-$r$ witnesses, respectively,
one has $\theta_X=1$ and $\operatorname{ord}(\theta_Y)=r^2$.
A $q$-active pure-$q$ witness therefore forces prime phase conductor
$r$. Both complete endpoint fields have $r$-conductor
exponent exactly two and are tame at $q$. Wild ramification at $p$
in either endpoint field requires $ph^2<q$.
\end{proposition}
\begin{proof}
The phase-fiber theorem makes the word $r$-primary because
$r>g$, with unique nonnegative integral coefficients and
$h+qk=g$. Strictness gives $k>0$; one has $h>0$, and $h=1$
would give $q\mid(p-1)(p+1)$, impossible for $q>p+1$.
Thus $h\ge2$. The ordinary channel bound is
\begin{equation}\label{paper:wild-mass-budget}
 ab|\widehat w(\chi)|^2\le gq^2r^2.
\end{equation}
The stabilizer $H$ of the unique coefficient vector is tame,
since its nonidentity coefficient mass is $k<g/q<r$.
Its order is $d$, and its action on nonidentity roots is free;
hence $d\mid k$ and $d\le(g-h)/q$.

There is a primitive torus channel at conductor $r^2$ with
$|\widehat w(\chi)|^2\ge q^2r^2$, by the integral
Fourier/Gauss argument in Corollary~\ref{paper:primary-inactive-cost}.
Orient the channel so that its source is an $r$-inactive input.
If the other input had $r$-primary twist order dividing $r$, Lemma~\ref{paper:two-depth-compression}
would improve the right side of \eqref{paper:wild-mass-budget}
by the factor $r$, forcing $abr\le g$. This contradicts
$r>g$. Thus exactly one input is $r$-active, with $r$-primary twist
order at least $r^2$. Its tame stabilizer orbit makes $d$ divide
its relative field index, and hence $u=ab/d$ is a positive integer.

We next remove the order-$r$ letters. Reduction identifies
$H$ with a subgroup of $\mathbb F_r^\times$ of order $d$.
If such letters occurred, their coefficients would give a
nonzero nonnegative integral vector on
$\Gamma_1=\mathbb F_r^\times/H$, whose mass is less than
$g/(qd)<|\Gamma_1|$. The vector is therefore nonconstant.
Some nontrivial character of $\Gamma_1$ has integral Fourier
coefficient of modulus at least one. Inflated to
$\operatorname{Gal}(K/\mathbb Q)$, its channel kills the
primitive $r^2$ letters, while each order-$r$ Gauss sum gains
an extra factor $r$. Consequently
\[
 |\widehat w(\chi)|^2\ge q^2r^3,
\]
contradicting \eqref{paper:wild-mass-budget}, since $abr\le g$
would follow. Thus only primitive $r^2$ letters occur.

Their full traces vanish. The trivial Fourier channel is
therefore
\[
 \widehat w(1)=\frac{r(r-1)}d h.
\]
Use $d\mid ab$ in \eqref{paper:wild-mass-budget} and then the
bound on $d$ to obtain
\[
 u(r-1)^2h^2\le gq^2d\le gq(g-h).
\]
Because $r$ and $g$ are odd and $r>g$, one has $r-1>g$;
hence $uh^2<q$. A higher $r$-conductor at either endpoint
would contribute a factor $r$ to $u$, since $(d,r)=1$, which
is impossible as $u<q<r$. The active $r$-twist therefore has
exact order $r^2$.

If either input were $q$-active, the joint active-inertia
argument would give $d(q-1)\mid ab$: use the product of the
tame $r$-stabilizer and full $q$-inertia at one endpoint when
both activities occur there, and multiply the separate endpoint
costs otherwise. Partial inversion excludes simultaneous sign
stabilizers, exactly as in Corollary~\ref{paper:active-inertia}.
Thus $u\ge q-1$, contradicting $h\ge2$. The same argument
at $p$ gives $\varphi(p^t)\mid u$ for a $p$-active input
of twist order $p^t$. Since $uh^2<q\le p^2/3$ and $h\ge2$,
one must have $t=1$, and $(p-1)h^2<q$. If both inputs were
$p$-active, the two independent costs would give
$(p-1)^2\mid u$, again impossible. Finally $h<\sqrt q<q$
and $h\equiv g\pmod q$ give the stated residue and mass.

The preceding $r$-conductor argument and containment of $K$
make both full $r$-conductor exponents exactly two. Since
$d<p<q$, wild $q$-ramification in either field would force
$q\mid u$, impossible. Wild $p$-ramification would similarly
force $p\mid u$, and hence $ph^2<q$. Finally the actual
phases have order dividing $r^2$, while the input twist bounds
just proved divide $pr^2$; balancing gives the same bound
for every constituent.
\end{proof}

\begin{corollary}[The least-index wild alternative]
\label{paper:wild-mass-index-one}
In Proposition~\ref{paper:wild-mass-crossing}, suppose $u=1$.
Write $V$ for the inactive input and $A$ for the $r$-active one.
Then
\[
 d\text{ is odd},\qquad K_V=K,\qquad
 K_A=\mathbb Q(\zeta_{r^2}),\qquad
 \theta_V=1,\quad\operatorname{ord}(\theta_A)=r^2.
\]
In particular this applies when $q\le2h^2$.
\end{corollary}
\begin{proof}
The tame stabilizer order $d$ divides the active endpoint's
relative field index. The inactive endpoint's index and the
active endpoint's index divided by $d$ are positive integers
with product one. Since $V$ is a nonunit in odd dimension,
it is not self-dual; its full character field $K_V=K$ is
nonreal. Thus $-1\notin H$ and $d$ is odd. The actual
order-$r^2$ primary twist of $A$ makes $K_A$ contain the
maximal real subfield of $\mathbb Q(\zeta_{r^2})$. Together
with $K$, this generates the full cyclotomic field: the two
stabilizers have orders $2$ and odd $d$, so their intersection
is trivial. Finally $[K_A:\mathbb Q]=d[K:\mathbb Q]=\varphi(r^2)$,
which proves equality. No $p$-active input is possible because
$p-1$ would divide $u=1$, so the full twist assertions follow.
The last assertion follows from $1\le u<q/h^2\le2$.
\end{proof}

\begin{corollary}[A twist cap for an active pure-$q$ crossing]
\label{paper:active-pure-crossing-cap}
In Proposition~\ref{paper:wild-mass-crossing}, retain the degree
hypotheses $d_X=p^2q^2$, $d_Y=p^2r$ and $r>g=p^2$, but assume
instead that $X,Y$ have pure-$q$ and pure-$r$ twists, respectively,
with $X$ $q$-active. Then the phase field has conductor $r$,
$\operatorname{ord}(\theta_Y)\mid r$, and every constituent of
$XY$ has twist order dividing $q^2r$. Moreover
\[
 \operatorname{ord}(\theta_X)=q^2
 \quad\Longrightarrow\quad
 \operatorname{ord}(\theta_Y)=r.
\]
\end{corollary}
\begin{proof}
The one-primary phase bound leaves conductors $r$ and $r^2$;
Proposition~\ref{paper:wild-mass-crossing} excludes the latter.
Write $K=\mathbb Q(w)$, $d=[\mathbb Q(\zeta_r):K]$,
$a=[K_X:K]$, $b=[K_Y:K]$, and
\[
 w=h+q\sum_{\xi\ne1}k_\xi\xi,\qquad
 h+q\sum k_\xi=g,\qquad
 Q=h^2+q^2\sum k_\xi^2.
\]
The coefficients are nonnegative integers, and
\[
 Q\ge h^2+q(g-h)\ge qg-q^2/4,\qquad
 E(w)=\frac{rQ-g^2}{d}.
\]
The ordinary energy budget is $abE(w)\le B=gq^2r^2$.
Suppose that $Y$ has twist of order $r^2$. Joint inertia costs
then give
\[
 ab\ge rd\,\varphi(q^t),\qquad
 \operatorname{ord}(\theta_X)=q^t,\quad t\in\{1,2\}.
\]
Here the active $r$-inertia at $Y$ contributes its wild factor
$r$ together with the tame stabilizer of size $d$, while the
$q$-active endpoint contributes $\varphi(q^t)$ independently.
Since $r>g$, one obtains
\[
 abE(w)>r^2\varphi(q^t)\bigl((q-1)g-q^2/4\bigr).
\]
For $q\ge7$ and $g\ge3q$,
\[
 2(q-1)\bigl((q-1)g-q^2/4\bigr)>gq^2.
\]
Indeed the difference is
$g(q^2-4q+2)-(q-1)q^2/2$, which is positive already at
$g=3q$. If $t=1$, the trivial-$q$-torus half-channel bound
of Corollary~\ref{paper:half-channel-crossing} gives
\[
 abE(w)\le\frac{q^2-1}{2q^2}B<B/2,
\]
contradicting the displayed lower bound. If $t=2$, the lower
bound gains a factor $q$ and already contradicts the ordinary
budget $B$. The codegree bound at $Y$ therefore gives
$\operatorname{ord}(\theta_Y)\mid r$.

Finally, if $X$ has twist order $q^2$ and $Y$ were $r$-inactive,
both endpoints would be inactive at the sole phase prime $r$.
Corollary~\ref{paper:primary-inactive-cost}, with $t=q$ and
$\beta=1$, would give $ab\le g$. Active $q$-inertia at $X$
gives $a\ge q(q-1)>p^2=g$, a contradiction. Balancing and
the phase conductor $r$ now give the asserted constituent cap.
\end{proof}

\subsection{The compulsory crossing alternative}

\begin{corollary}[Conductor loss at a top-height crossing]
\label{paper:top-conductor-loss}
In Corollary~\ref{paper:phase-fibers}, suppose $\max(i,h)=3$.
If both primary supports of $z/p^e$ occur, then
\[
 \operatorname{cond}\mathbb Q(z)
    <\frac{p^2}{qr}\frac{D}{p^{2e}L^2}.
\]
In particular its conductor cannot attain $q^{6-2j}r^2$.
\end{corollary}
\begin{proof}
Combine \eqref{paper:mixed-energy} and
\eqref{paper:scaled-crossing-budget}. The last assertion uses
$qr>p^2$ and $e\ge0$.
\end{proof}

\begin{lemma}[The uncancelled energy of a prime-order word]
\label{paper:prime-word-energy}
In the surviving odd ordered family, put $g=p^2$. Suppose an
actual mass-$g$ word from Corollary~\ref{paper:phase-fibers}
has conductor $r$, and write it as
\[
 w=h+q\sum_{\xi\in\mu_r\setminus\{1\}}k_\xi\xi,
 \qquad h=g-qk,\qquad k=\sum_\xi k_\xi\ge1.
\]
Put $K=\mathbb Q(w)$ and $d=[\mathbb Q(\zeta_r):K]$, and let
\[
 Q=h^2+q^2\sum_\xi k_\xi^2,
 \qquad F=rQ-g^2,
 \qquad B=q^2r^2/g.
\]
Then $h\ge2$ and
\begin{equation}\label{paper:exact-period-energy}
 E(w)=F/d,\qquad d\mid k,\quad d\mid r-1,\quad d\le k<g/q<p.
\end{equation}
If $r>g/2$, then
\begin{equation}\label{paper:uncancelled-energy}
                              F>B.
\end{equation}
\end{lemma}
\begin{proof}
The phase margins and uniqueness give the displayed word and
$h>0$. The equality $h=1$ would imply $q\mid p^2-1$, impossible
because $q>p+1$. The stabilizer of the coefficient vector has
order $d$ and acts freely on nonidentity roots, so $d\mid k$.
The prime-order trace formula gives the displayed energy.

Since $\sum k_\xi^2\ge k$ and $h+qk=g$,
\[
 Q\ge h^2+q^2k=qg-q^2/4+(h-q/2)^2\ge qg-q^2/4.
\]
For fixed $g=p^2$, the function
$A(q)=qg-q^2/4-g^2/q$ increases on $p+2\le q\le g/3$,
because $A'(q)/g=1+g/q^2-q/(2g)>0$. At the left endpoint,
\[
 \frac{A(p+2)}g
   =\frac{4(p+1)}{p+2}-\frac{(p+2)^2}{4p^2}>2.
\]
Hence $Q>g^2/q+2g$. With $r>g/2$ and
Corollary~\ref{paper:family:uniform-prime-range}, this gives
\[
             F=rQ-g^2>rg^2/q>q^2r^2/g=B.
\]
\end{proof}

\begin{theorem}[The compulsory crossing alternative]
\label{paper:top-second-height}\label{paper:top-boson}
\label{paper:top-crossing}\label{paper:top-second-unit}
\label{paper:divided-inactive-rigidity}
In the surviving odd ordered family, take a compulsory pair
\[
 d_X=p^iq^j,\qquad d_Y=p^hr,\qquad
 i,h\in\{2,3\},\quad j\in\{1,2\}.
\]
Put $L=p^{\max(i,h)}q^jr$ and $z=S_{XY}/L$.
The assertions below hold for every such pair, without a
preliminary choice of pure primary witnesses.
\begin{enumerate}[label=\textup{(\roman*)},leftmargin=*]
\item If $\min(i,h)=2$ or $j=2$, then $z$ is a unit at
some place over $p$.
\item If $z$ vanishes at every place over $p$, then
$i=h=3$ and $j=1$. Put
\[
 w=z/p,\quad K=\mathbb Q(w),\quad N=\operatorname{cond}K,
 \quad M=\lcm\bigl(N,\operatorname{ord}\theta_X,
                         \operatorname{ord}\theta_Y\bigr).
\]
Then $w$ is an integral actual mass-$p^2$ word and a unit at
some $p$-place, and
\[
                         M\mid q^2r\quad\text{or}\quad M\mid qr^2.
\]
Every constituent of $X\otimes Y$ has actual twist order
dividing the same $M$. If $\ell^2\mid N$, an input is
$\ell$-active. If both phase supports occur, an input is
$q$-active; for $N=qr^2$ an input is also $r$-active.
\item\label{paper:top-first-inactive}
In \textup{(ii)}, suppose both inputs are inactive at every
prime dividing $N$. Then for one $\ell\in\{q,r\}$,
\[
 N=\ell,\qquad \theta_X=\theta_Y=1,\qquad
 [K_X:K][K_Y:K]p^2\le q^2.
\]
Every product twist then has order dividing $\ell$.
In particular this conclusion holds whenever both complete
input twists are one.
\end{enumerate}
\end{theorem}
\begin{proof}
Corollary~\ref{paper:cor:proper-crossing} proves the unit assertion
when $\min(i,h)=2$. Hence suppose $i=h=3$ and $z$ vanishes
at every $p$-place. Put $g=p^2$ and use the notation in
\textup{(ii)} for either value of $j$.
The two input codegrees are $p$-units, so their $p$-parts are
trivial. Corollaries~\ref{paper:top-dichotomy} and
\ref{paper:phase-fibers} give the unique actual mass-$g$ word,
with no nontrivial $p$-part. Its strict house bound excludes
the all-identity word. Write $N=q^\alpha r^\beta$, where
$0\le\alpha\le6-2j$, $0\le\beta\le2$, and
$\alpha+\beta>0$. Put
\[
 a=[K_X:K],\quad b=[K_Y:K],\quad B=q^2r^2/g,
\]
and let $t$ and $s$ be the maximum $q$- and $r$-twist exponents,
respectively, among the two inputs. Then
\[
 A=\max(\alpha,t),\qquad R=\max(\beta,s),\qquad M=q^A r^R.
\]
The two-sided budget is
\begin{equation}\label{paper:divided-two-sided}
                         abE(w)<q^{4-2j}B.
\end{equation}

We use the following orbit counts. When $\alpha\ge1$, the
wild $q$-inertia fixing the phase field gives at least
$q^{\max(t-\alpha,0)}$ different primary twists; the analogous
factor at $r$ is $r^{\max(s-\beta,0)}$ when $\beta\ge1$.
If a phase prime is absent, the full inertia twist image has
size $\varphi(q^t)/2$ or $\varphi(r^s)/2$ when active.
Counts at the two primes multiply: if their maxima occur at
one endpoint, independent inertia factors give the Cartesian
product of its two primary twist images; if at different
endpoints, multiply the two field indices.
Separately, Corollary~\ref{paper:active-inertia} gives the
stronger full label-orbit divisor $\varphi(q^t)$ whenever
$\alpha=0$, and similarly at an absent $r$-phase prime.

First suppose $j=1$. If both phase supports occur, the mixed
energy bound and the two wild counts give
\[
                  abE(w)>\frac{qr}{g}q^A r^R.
\]
Hence $M<q^3r$. As $A,R\ge1$ and $r>q$, this forces
$A,R\le2$ and excludes $A=R=2$.

If only the $q$-support occurs, the one-prime bound gives
\[
 abE(w)>\tfrac23r^2q^A,
 \qquad q^A<\frac{3q^4}{2g}\le\frac{q^3}{2}.
\]
Thus $A\le2$. If either input is $r$-active, multiplication
by its full inertia twist-image count gives
\[
                    gq^A r^{s-1}(r-1)<3q^4.
\]
Since $g\ge3q$ and $r>3q$, this excludes $s\ge2$ and also
excludes $A=2$. The desired cap follows in this case.

If only the $r$-support occurs, the one-prime bound and wild
$r$-inertia give
$abE(w)>(2/3)q^2r^R$. If $R\ge3$, the budget would imply
$r<3q^2/(2g)\le q/2$, a contradiction. Thus $R\le2$.
Suppose $t\ge3$. Since $K$ is unramified at $q$, the active
endpoint has field index at least
\[
                         \varphi(q^t)\ge q^2(q-1).
\]
The one-prime bound excludes $\beta=2$, so $N=r$ and the
budget gives $(q-1)E(w)<B$. In particular
$r>(2/3)(q-1)g>g/2$.
Lemma~\ref{paper:prime-word-energy} now gives $E(w)=F/d$,
$F>B$, and $d<p$. Consequently
\[
                   ab<q^2d<q^2(q-1),
\]
contrary to the displayed full inertia bound. Hence $t\le2$.
The only possible violation of the desired cap is now
$t=R=2$. But the $q$-twist image has size $q(q-1)/2$;
combine it with the wild $r$-factor if $\beta=1$. This gives
\[
              abE(w)>\tfrac13q^3(q-1)r^2>q^4r^2/g,
\]
again a contradiction. This proves the order bound in \textup{(ii)}.

The phase-conductor identity identifies $N$ with the least
common multiple of the actual balancing-phase orders. Each
constituent twist is its phase multiplied by $\theta_X\theta_Y$,
so its order divides $M$. This proves the joint order bound.

Now suppose $j=2$. Mixed support would give $E(w)>B$.
The $q$-only case would give $E(w)>(2/3)r^2q>B$.
Thus only the $r$-support occurs, and the one-prime bound
excludes conductor $r^2$. Hence $N=r$ and $r>2g/3$.
Lemma~\ref{paper:prime-word-energy} gives the actual word
\begin{equation}\label{paper:prime-level-word}
 w=h_0+q\sum_{\xi\in\mu_r\setminus\{1\}}k_\xi\xi,
 \qquad h_0=g-qk\ge2,\quad k=\sum_\xi k_\xi\ge1,
 \quad k_\xi\in\mathbb Z_{\ge0}.
\end{equation}
With $d=[\mathbb Q(\zeta_r):K]$, it also gives $d\mid k$,
$d<p$, and $E(w)=F/d$. Since $r>2g/3>g/2$, the strict
numerator estimate is $F>B$. The two-sided budget therefore gives

\begin{equation}\label{paper:field-fiber-bound}
                              ab<d<p.
\end{equation}

If either input were $q$-active, its field index would be
divisible by $\varphi(q^t)\ge q-1>p$, by
Corollary~\ref{paper:active-inertia}, contradicting
\eqref{paper:field-fiber-bound}.
If either input were $r$-active, the same corollary says
its stabilizer in full $r$-inertia is an $r$-group.
The tame subgroup of order $d$ fixing $w$ would therefore
act freely on it, making the corresponding field index
at least $d$, again impossible. Thus both complete twists
are one. The coefficient vector is invariant over $K$, and
$\sum k_\xi<p<r$. Apply
Corollary~\ref{paper:primary-inactive-cost} with $\ell=r$,
$\beta=1$, $t=q$, and $c=p^4q^2r$. Since
$D/c^2=q^2r^2/p^2$, it would require
\[
                         ab\,q^2r^2\le q^2r^2/p^2,
\]
or $abp^2\le1$, an impossibility. Thus the divided branch
does not occur at $j=2$, proving \textup{(i)}.


We finish the activity assertion of \textup{(ii)} and
\textup{(iii)} together. If both phase supports occur, the
effective constant is $D/c^2=q^4r^2/g$, so
Corollary~\ref{paper:mixed-inactive-exclusion} forces a
$q$-active input and, for conductor $qr^2$, an $r$-active
input. It also excludes the hypothesis of \textup{(iii)}.

It remains to consider one phase prime $s\in\{q,r\}$.
Assume both inputs are $s$-inactive and let $t$ be the other
prime. This will exclude a squared $s$-conductor and prove
\textup{(iii)} in the remaining case.
The actual phase-fiber and one-prime margin statements give
\[
 w=h_0+t\sum_{\xi\ne1}k_\xi\xi,\qquad
 k_\xi\in\mathbb Z_{\ge0},\qquad
 0<\sum k_\xi<s.
\]
The coefficients are invariant over $K$.
Part \textup{(ii)} bounds the phase conductor by $s^\beta$
with $\beta\in\{1,2\}$.
Here $c=p^4qr$ and $D/c^2=q^4r^2/g$.
If $s=q$ and $\beta=2$, Corollary~\ref{paper:primary-inactive-cost}
would give
\[
                      ab\,r^2q^4\le q^4r^2/g,
\]
which is impossible. If $s=r$ and $\beta=2$, it would give
\[
                      ab\,q^2r^4\le q^4r^2/g,
\]
or $abgr^2\le q^2$, also impossible. Thus $\beta=1$.
The same bound now gives $abg\le q^2$ in either case.

If an input were active at $t$, its field index over $K$
would be divisible by the full inertia factor
$\varphi(t^u)\ge t-1\ge q-1$, by
Corollary~\ref{paper:active-inertia}, since $K$ is unramified
at $t$. This contradicts
\[
                   abg\le q^2<g(q-1),
\]
where $g\ge3q$ in the surviving region. Both inputs are
already $p$-inactive by their codegrees and are $s$-inactive
by hypothesis. Thus their complete twists are one.
The actual phases have orders dividing $s$, so the same is
true of all constituent twists.
\end{proof}


\subsection{Degree restrictions inside an exposed adjoint}

\begin{proposition}[The two companion squares]
\label{paper:prop:companion-squares}
In an ordinary perfect fusion-simple integral modular category
of dimension $p^6q^6r^4$, with $p<q<r$ odd, fix
$Y\in P_q$ of degree $p^2q^j$, $j\in\{1,2\}$.
Suppose a nonunit $U$ has, up to a scalar root, a full
positive $r$-primary character word at $Y$, whose adjoint
difference set $\mathscr{D}$ is exposed. Then every nonunit
constituent $Z$ of $UU^*$ has $r$-height at most one.
Its $p$-free possibilities are exactly among
$q^2r,q^3r$, and its $q$-free possibilities exactly
among $p^2r,p^3r$.

If $p\mid d_U$, at least one of the $p$-free types
occurs. If $q\mid d_U$, at least one of the $q$-free
types occurs. In particular a seed of degree $p^2r$
forces an adjoint constituent of degree at least $q^2r$.
\end{proposition}
\begin{proof}
Exposure and positivity put the actual phases of each
constituent at $Y$ in $\mathscr{D}$. The set contains $1$ and
misses a position in every $\mu_r$-coset; hence it
extends to an integral cyclotomic basis. An integral
value expressed on $\mathscr{D}$ therefore has integral
coefficients, also after division whenever the divided
value is integral.

If $v_r(d_Z)=2$, its codegree is $r$-free. Thus $K_Z$
is unramified at $r$, while $\chi_Y(Z)\in K_Z$ up
to a rational factor and is $r$-primary. It is rational.
The basis containing $\mathscr{D}$ and $1$ then makes its full
word the constant $d_Z$, contradicting strictness.

For a $p$-free constituent of degree $qr^c$, LCM
integrality makes $\beta=\chi_Y(Z)/r^c$ integral.
It is $r$-primary of mass $q$. Principality at $q$
makes its residue zero. Its field is unramified at
$q$, so $\beta/q$ is integral. Its coefficients on
$\mathscr{D}$ are nonnegative integers of sum one, forcing
a single phase and contradicting strictness.
Prime-power degrees are excluded, and the degree
caps now leave only $q^2r,q^3r$ for a $p$-free
nonunit.

A $q$-free nonunit has degree $p^a r^c$ with
$a,c>0$. At $a=1$, its normalized crossing with
$Y$ is integral, supported on $\mathscr{D}$, of mass $p$,
and has residue $p$ at $q$. Comparing in a basis
containing $\mathscr{D}$ and $1$, every off-identity
coefficient is zero modulo $q$. Since the mass
$p$ is below $q$, only the constant word $p$
is possible, again contradicting strictness.
The remaining degree caps give $p^2r,p^3r$.

Finally, the unit occurs once in $UU^*$. If all
nonunit constituents had degree divisible by $p$,
its dimension would be one modulo $p$, impossible
when $p\mid d_U$. The argument at $q$ is identical.
\end{proof}

\begin{corollary}[No degree $pq$ on a proper primary word]
\label{paper:cor:no-pq}
Under the same assumptions on the category and row, suppose
$r>p^2$. No nonunit of degree $pq$ has a full positive
$r$-primary character word whose support misses a position
in every $\mu_r$-coset. In particular such a degree does
not occur in the exposed adjoint of
Proposition~\ref{paper:prop:companion-squares}.
\end{corollary}
\begin{proof}
For such a label $Z$, principality at $q$ and
unramifiedness make $\chi_Y(Z)/q$ integral.
Its coefficients on the proper support are
nonnegative integers of sum $p$. It has at most
$p$ positions. Its adjoint difference set has
at most $p(p-1)+1$ positions and is exposed, since
\[
 (p-1)r\ge(p-1)(p^2+2)
       >p\bigl(p(p-1)+1\bigr).
\]
Apply Proposition~\ref{paper:prop:companion-squares}
to $ZZ^*$. It has a $p$-free constituent of degree
at least $q^2r$ and a $q$-free constituent of
degree at least $p^2r$. They are distinct, since
nonunits cannot have prime-power degree. Including
the unit gives
$p^2q^2\ge1+(p^2+q^2)r$, impossible for $r>p^2$.
\end{proof}

\begin{corollary}[An active row removes the full-height option]
\label{paper:cor:active-companion}
In Proposition~\ref{paper:prop:companion-squares}, if $Y$
is $q$-active, every $p$-free nonunit constituent has
degree $q^2r$. If $Y$ is $p$-active, every $q$-free
nonunit constituent has degree $p^2r$.
\end{corollary}
\begin{proof}
Height three at the indicated sextic prime has
unit codegree there and therefore vanishes on an
active row. But a nonzero positive word on a set
extending to a basis cannot vanish. The height-three
option is excluded in each case.
\end{proof}


\begin{theorem}[Sharp companions in a pure wild crossing]
\label{paper:wild-adjoint}\label{paper:wild-adjoint-sharp-companion}
\label{paper:sharp-companion-pure-twists}
Consider the pure-$q$/pure-$r$ crossing of
Proposition~\ref{paper:wild-mass-crossing}: $X\in P_q$ has degree
$p^2q^2$ and twist one, and $Y$ has degree $p^2r$ and twist order
$r^2$. Put $g=p^2$, $K=\mathbb Q(w)$,
$d=[\mathbb Q(\zeta_{r^2}):K]$, and $u=ab/d$.
Every nonunit $p$-free constituent $Z$ of $YY^*$ has degree $q^2r$.
Its normalized character
\[
 v_Z=\frac{S_{XZ}}{p^2q^3r},\qquad F_Z=\mathbb Q(v_Z),
\]
is a positive integral $r$-primary word of mass $q$, and for some
$\beta\in\{1,2\}$,
\[
 \operatorname{ord}(\theta_Z)=\operatorname{cond}(F_Z)=r^\beta,
 \qquad \operatorname{cond}(K_X)\mid pqr^2,
 \qquad \operatorname{cond}(K_Z)\mid pqr^\beta.
\]
If $\beta=2$, then $F_Z=\mathbb Q(\zeta_{r^2})\subseteq K_X$,
every phase of $v_Z$ is primitive of order $r^2$, and $u\ge d$.
Thus $u=1$ forces $d=1$ in this case, while $u<d$ forces prime
conductor for every such adjoint word. At least one such $Z$ occurs.
No assertion of individual principal membership or of a purely
$r$-primary full character field is made.
\end{theorem}
\begin{proof}
\noindent\emph{The adjoint word and its wild field.}
Let $S$ be the phase support of the original mass-$p^2$ word,
and let $\mathscr{D}=S-S$ be its adjoint difference support, written
additively. The family support bound gives
\[
 |S|\le k+1,\quad k=\lfloor p^2/q\rfloor\le p-2,
 \qquad |\mathscr{D}|\le k(k+1)+1<p^2<r.
\]
The set $\mathscr{D}$ is exposed by Proposition~\ref{paper:family:prop:auto-exposure}
and extends to an integral cyclotomic basis. Every actual
word at $X$ of a constituent of $YY^*$ is supported on $\mathscr{D}$.
Proposition~\ref{paper:prop:companion-squares} gives the two
$p$-free degrees.

For degree $q^3r$, the stated normalization is ordinary LCM
integrality and has mass $q^2$. For degree $q^2r$, first put
$z=S_{XZ}/(p^2q^2r)$, of mass $q^2$. Since
$\chi_X(Z)=rz$, principality of $X$ at $q$ gives zero
residue at every $q$-place. Its field is $r$-primary and
therefore unramified at $q$, so $z/q$ is integral. The basis
containing $\mathscr{D}$ makes its actual coefficients nonnegative
integers. This is $v_Z$, of mass $q$.

Assume that $F=F_Z$ has conductor $r^2$, and write
\[
 \delta_Z=[\mathbb Q(\zeta_{r^2}):F],\quad
 a_Z=[K_X:F],\quad b_Z=[K_Z:F].
\]
The ordinary Fourier budget is
\[
 a_Zb_Z|\widehat v_Z(\chi)|^2\le gr^2.
\]
The stabilizer of the coefficient vector is tame: its wild
subgroup would require at least $r$ primitive positions,
whereas $|\mathscr{D}|<r$. Some primitive channel has squared modulus
at least $r^2$, by integral Fourier nonvanishing. Since $X$
is $r$-inactive, Lemma~\ref{paper:two-depth-compression} shows
that $Z$ must have $r$-twist order $r^2$; order at most $r$
would give $a_Zb_Zr\le g<r$. The codegree of $Z$ bounds its
$r$-twist order by $r^2$. Active inertia gives $\delta_Z\mid b_Z$.

If order-$r$ letters occurred, their coefficient vector on
$\mathbb F_r^\times/H_Z$ would be nonconstant: a nonzero
constant vector occupies all $r-1$ nonidentity order-$r$
roots, more than $|\mathscr{D}|$. A nontrivial tame channel would then
have squared modulus at least $r^3$, by the inflated Gauss
sum, contradicting $a_Zb_Zr\le g$. Thus only primitive
$r^2$ letters and possibly the identity remain.

Write $h_Z$ for the constant coefficient. The trivial
channel is $r(r-1)h_Z/\delta_Z$. With $u_Z=a_Zb_Z/\delta_Z\ge1$,
it gives
\[
 u_Z(r-1)^2h_Z^2\le g\delta_Z.
\]
A primitive stabilizer orbit occurs in $\mathscr{D}\setminus\{0\}$,
so $\delta_Z\le |\mathscr{D}|-1<g$. Since $r>g$, this forces $h_Z=0$.
The stabilizer acts freely on every remaining root, so
$\delta_Z$ divides the mass, either $q$ or $q^2$.

To compare fields without assuming $u=1$, put
\[
 L=K_X\cap\mathbb Q(\zeta_{r^2}),\qquad
 \delta_X=[\mathbb Q(\zeta_{r^2}):L],\qquad A=[K_X:L].
\]
One has $K\subseteq L$, so $\delta_X\mid d<p$. Also $F\subseteq L$,
so $\delta_X\mid \delta_Z$. Since $\delta_Z$ divides a power of $q$, one has $\delta_X=1$. Now $a_Z=A \delta_Z$, and $\delta_Z\mid b_Z$;
the primitive channel $a_Zb_Z\le g$ excludes $\delta_Z\ge q$.
Hence $\delta_Z=1$ and $L=F=\mathbb Q(\zeta_{r^2})$.
For the original crossing, $a=A d$ and $d\mid b$, so
$u=Ab\ge A d\ge d$. The specialization at $u=1$ follows.

Finally a rational word is impossible: coefficient uniqueness
on the basis containing $\mathscr{D}$ would make it the constant word
of full mass, contrary to strictness. A non-wild $F_Z$
therefore has conductor exactly $r$.

\noindent\emph{The exact companion degree.}
Retain the support and field notation just introduced. Two bounds
exclude the remaining degree-$q^3r$ possibility. Whenever $d=1$, the
original trivial-channel inequality gives
\[
 (r-1)h\le pq,\qquad q>ph\ge2p,\qquad k<p/2.
\]
Consequently
\begin{equation}\label{paper:small-wild-adjoint-support}
 |\mathscr{D}|\le k(k+1)+1\le(p^2+3)/4<g/2.
\end{equation}
If $d\ge2$, every nonidentity position of $S$ is primitive,
and its coefficient vector is invariant under the tame
stabilizer $H$ of order $d$. This group acts freely on
nonzero residue classes modulo $r$. For each orbit of
such classes, let $n_i$ be the number of positions in each
class. Then $d\sum_i n_i\le k$. Nonzero differences of
order $r$ arise only from two positions in one such class.
Therefore
\begin{equation}\label{paper:small-prime-adjoint-support}
 |\mathscr{D}\cap\mu_r|
 \le1+d\sum_i n_i(n_i-1)
 \le1+k(k/d-1)<g/2.
\end{equation}
The left side identifies additive exponents with roots;
the last inequality follows from $d\ge2$ and $k\le p-2$.

Suppose a $p$-free constituent had degree $q^3r$.
Its word $v_Z$ has mass $M=q^2$, and its effective budget
is $gr^2$. The first part excludes a rational phase field.

If $F_Z$ has conductor $r^2$, the first part gives
$L=F_Z=\mathbb Q(\zeta_{r^2})$ and $u\ge A d$.
The original trivial-channel inequality consequently yields
\[
 A(r-1)^2h^2\le gq^2,
 \qquad q>ph\sqrt A\ge2p.
\]
Thus the bound $|\mathscr{D}|<g/2$ in
\eqref{paper:small-wild-adjoint-support} applies here too,
even when $d>1$. Write $Q$ for the sum of squared
coefficients and $T$ for the number of occupied positions.
The prime-square Ramanujan formula and positivity give
\[
 E(v_Z)\ge r^2Q-rM^2
 \ge r^2M^2\left(\frac1T-\frac1r\right)
 >r^2M^2/g>gr^2,
\]
contradicting the ordinary energy budget.

It remains that $F_Z$ has conductor $r$. Its actual word
is then supported on $\mathscr{D}\cap\mu_r$: wild inertia fixes its
unique coefficient vector, while any primitive position
would require $r$ positions, contrary to $|\mathscr{D}|<r$.
Put $\delta_1=[\mathbb Q(\zeta_r):F_Z]$. Then
\[
 \delta_X\mid \delta_1,\qquad a_Z=A r \delta_1/\delta_X\ge r.
\]
The word is nonconstant and has support smaller than $r-1$,
so integral Fourier nonvanishing supplies a nontrivial
prime-level channel of squared modulus at least $r$.
If $Z$ were $r$-inactive, its inactive channel bound would
give $a_Zb_Z\le g$, impossible since $a_Z\ge r>g$.
If its $r$-twist had order $r^2$, active inertia would give
$b_Z\ge r\delta_1$. The ordinary primitive channel gives
$a_Zb_Z\le gr$, implying
$A r \delta_1^2\le g\delta_X$, again impossible.
Thus $Z$ has exact $r$-twist order $r$, and $b_Z\ge \delta_1$.
The prime-level energy identity now yields
\[
 A\frac{\delta_1}{\delta_X}\left(Q-\frac{M^2}{r}\right)\le g.
\]
For $d=1$ and $d\ge2$, respectively, the applicable support
bound above gives $T<g/2$. Hence
\[
 Q-M^2/r\ge M^2(1/T-1/r)>M^2/g=q^4/g>g,
\]
a contradiction. This excludes degree $q^3r$.

For a degree-$q^2r$ constituent with prime phase conductor,
the same index and primitive-channel arguments force exact
$r$-twist order $r$; they do not depend on the mass $M$.
The wild case is covered by the wild-field argument above.
Finally a nonunit $p$-free constituent exists by
Proposition~\ref{paper:prop:companion-squares}.

\noindent\emph{Foreign activity and the full fields.}
Retain $g=p^2$, $v_Z=S_{XZ}/(p^2q^3r)$, of mass $q$,
and write $Q$ for the sum of its squared integer coefficients.
Then $Q\ge q$. Put
\[
 L=K_X\cap\mathbb Q(\zeta_{r^2}),\qquad
 \delta_X=[\mathbb Q(\zeta_{r^2}):L],\qquad A=[K_X:L].
\]
If $\beta=1$, put $\delta_1=[\mathbb Q(\zeta_r):F_Z]$,
$a_Z=[K_X:F_Z]=Ar\delta_1/\delta_X$ and
$b_Z=[K_Z:F_Z]$. The matching-conductor conclusion makes $Z$
$r$-active of order $r$, so $\delta_1\mid b_Z$.
Since $\delta_X\mid\delta_1$, the integer
\[
 J=A\frac{b_Z}{\delta_1}\frac{\delta_1}{\delta_X}
\]
is positive. The exact prime-level energy identity gives
\begin{equation}\label{paper:sharp-index-budget}
                    J\left(Q-\frac{q^2}{r}\right)\le g.
\end{equation}
If $\beta=2$, The wild-field argument gives
$L=F_Z=\mathbb Q(\zeta_{r^2})$. Put $J=A b_Z$.
The prime-square Ramanujan lower bound
$E(v_Z)\ge r^2Q-rq^2$ gives the same inequality.

For all surviving parameters,
\[
 q\left(1-\frac qr\right)>p.
\]
Indeed $r>p^2$ and $p+2\le q\le p^2/3$; the function
$t(p^2-t)$ increases on that interval, and at $t=p+2$
its difference from $p^3$ is $p^2-4p-4>0$ for $p\ge5$.
Consequently \eqref{paper:sharp-index-budget} implies $J<p$.

Joint inertia at the active $r$-input makes a $q$-active
$Z$ contribute $q-1$ to $J$, which is impossible. The same
joint orbit calculation makes wild ramification at $s=p$ or
$q$ in $K_Z$ contribute $s$ to $J$: at prime $r$-conductor
its wild inertia cost multiplies the free tame $r$-stabilizer
of size $\delta_1$, and at conductor $r^2$ the entire
$r$-primary field is already contained in $K_Z$. Wild
ramification at $s$ in $K_X$ contributes $s$ to $A$.
Thus both complete fields are tame at $p,q$. The codegree of
$Z$ bounds its full $r$-conductor exponent by two. An extra wild
$r$-factor in $K_Z$ when $\beta=1$ would similarly make
$r\mid J$, impossible. The original crossing field already has conductor $r^2$.

It remains to exclude $p$-activity of $Z$. Its contribution
$p-1\mid J$, together with $J<p$, would force $J=p-1$;
a higher $p$-twist order is already impossible. If $q\ge p+4$,
then
\[
 (p-1)q\left(1-\frac qr\right)>p^2.
\]
To verify this at the smaller bound $r>p^2$, use monotonicity
of $q(p^2-q)$ and evaluate at $q=p+4$. The numerator of the
difference is $2p^3-11p^2-8p+16>0$ for $p\ge7$.
The case $p=5,q\ge p+4$ is absent since $q\le p^2/3$.
This contradicts \eqref{paper:sharp-index-budget}.

The only remaining possibility is $q=p+2$ and $\beta=1$:
a wild companion already forces $q>2p$. For the original
crossing, $h\equiv p^2\equiv4\pmod q$ and $0<h<q$;
its constant coefficient and nonidentity mass are therefore
$h=4$, $k=p-2$. Write its original relative indices as $a,b$.
One has $a=Ad/\delta_X$, $d\mid b$, and hence
$Ad/\delta_X\le u$. Its original trivial-channel bound gives
\[
 \frac{A}{\delta_X}\le\frac ud
 <\frac{q^2}{16p^2},\qquad
 \delta_X>\frac{16Ap^2}{(p+2)^2}>3.
\]
Also $\delta_X\mid d\mid k=q-4$, and
$\delta_X\mid\delta_1$.
Let $h_Z$ be the constant coefficient of the prime-level
word $v_Z$. Tame stabilizer invariance gives
$\delta_1\mid q-h_Z$. If $h_Z=0$, then
$\delta_X\mid q$, impossible since $1<\delta_X<p<q$.
If $h_Z=1$, then $\delta_X$ divides both $q-4$ and $q-1$,
so $\delta_X\mid3$, also impossible. Thus $h_Z\ge2$,
and integer positivity gives
\[
                   Q\ge h_Z^2+q-h_Z\ge q+2.
\]
Finally, at $q=p+2$ and $r>p^2$,
\[
 (p-1)\left(q+2-\frac{q^2}{r}\right)>p^2:
\]
after multiplying by $p^2$, the smaller bound differs from
$p^4$ by $2p^3-7p^2+4>0$. This again contradicts
\eqref{paper:sharp-index-budget}. Therefore $Z$ is inactive
at both foreign primes, proving the full twist assertion
and the stated conductor caps.
\end{proof}


\subsection{Limits of support criteria}
The support mechanism has a genuine dense regime in which its
intersection--difference operation supplies no further exposed set.
Here exposure uses the actual least prime $p$, namely
$p|\mathscr{D}|<(p-1)r$, rather than the sufficient bound $3|\mathscr{D}|<2r$.

\begin{proposition}[Dense envelopes and a conditional support model]
\label{paper:dense-envelope-barrier}
Let $r$ be an odd prime and let $\mathscr{D}\subseteq\mathbb F_r$ satisfy
$|\mathscr{D}|>3r/4$. Then, for every $a\ne0$,
\begin{equation}\label{paper:eq:dense-closure}
       (\mathscr{D}\cap a\mathscr{D})-(\mathscr{D}\cap a\mathscr{D})=\mathbb F_r.
\end{equation}
Thus, if $\mathscr{D}$ is an exposed difference envelope, every first
consequence of the closure operation in
Lemma~\ref{paper:family:forced-envelope} is full and unexposed.

There is a uniform conditional construction compatible with the
exact-family initial support bound. Suppose $t\ge17$ and
$t,t+2,t^2+t+1$ are prime, and put
\[
 p=t,\qquad q=t+2,\qquad r=t^2+t+1,\qquad m=t-2.
\]
There is a set $S\subseteq\mathbb F_r$ with $|S|=m+1$ such that
$\mathscr{D}=S-S$ is exposed at least prime $p$, has $|\mathscr{D}|>3r/4$, and has
exactly $E=\binom{m+1}{2}$ nonzero opposite difference pairs.
After translating $S$ to contain zero, the word
\[
                 z=4+q\sum_{s\in S\setminus\{0\}}\zeta_r^s
\]
has mass $p^2$, strict house below $p^2$, and residue $p^2$
at every place over $q$ and $r$.
These are statements about sets and cyclotomic words, not about
fusion categories. No existence or infinitude assertion concerning
the displayed simultaneous primality condition is needed.
\end{proposition}
\begin{proof}
One has $|\mathscr{D}\cap a\mathscr{D}|\ge2|\mathscr{D}|-r>r/2$. A subset $T$ of
$\mathbb F_r$ of size greater than $r/2$ meets $x+T$ for every
$x$, and hence $T-T=\mathbb F_r$. This proves
\eqref{paper:eq:dense-closure}.

For the construction, the cyclic quotient
$\mathbb F_{t^3}^{\times}/\mathbb F_t^{\times}$ has order $r$.
The nonzero vectors in the trace-zero plane, modulo
$\mathbb F_t^{\times}$, give a set $S_0$ of $t+1$ points.
For a nonidentity quotient element represented by $a$, the
linear forms $\operatorname{Tr}(x)$ and $\operatorname{Tr}(ax)$
are independent, by nondegeneracy of the trace pairing.
Their common kernel is therefore a line. Consequently each
nonidentity quotient element occurs exactly once as an ordered
difference of $S_0$. Identify this cyclic group with
$(\mathbb F_r,+)$ and delete two points. The remaining set $S$
has $t-1$ points and still has distinct nonzero ordered differences,
so
\[
 |\mathscr{D}|=1+(t-1)(t-2)=t^2-3t+3.
\]
Now $m=\lfloor\min(r,p^2)/q\rfloor=t-2$, and
\[
 4|\mathscr{D}|-3r=t^2-15t+9>0,\qquad
 (p-1)r-p|\mathscr{D}|=3t^2-3t-1>0.
\]
This proves density, exact exposure and the difference count.
Finally $4+q(t-2)=p^2$; reduction at $q$ leaves $4\equiv p^2$,
and reduction at $r$ sends every $\zeta_r^s$ to one.
Distinct occurring roots and positive coefficients give strictness
at every conjugate.
\end{proof}

The associated set system explains the obstruction precisely.
Take supports $aS$ and $a\mathscr{D}\cap b\mathscr{D}$ for all nonzero $a,b$.
It is stable under dilation. The first supports have difference
sets $a\mathscr{D}$; every support of the second kind has full difference
set. Each pair of exposed adjoint masks therefore has the required
common-word support, but its new adjoint is unexposed. Moreover,
for $A=(\mathscr{D}\setminus\{0\})/\{\pm1\}$, $h=E$ and $H=(r-1)/2$,
\[
 |A\cap aA|\ge2h-H=(t^2-7t+4)/2\ge t-2\ge u(h).
\]
Thus this actual mask satisfies the lower intersection bound and
its consequent quotient and autocorrelation identities.
The set system supplies none of the fusion coefficients, modular
orthogonality or local-group compatibility of a category, and is
not a WGT counterexample. It isolates the need for additional
categorical information in this regime.

An unrestricted small-word bound for common constituents is also
unavailable from the present degree restrictions. At a $q$-principal row of
degree $p^2q^2$, the formally allowed constituent degree
$p^2q^2r$ has neither maximal sextic height and has $r$-height one.
Its LCM-normalized word has mass $p^2q^2$; division by $q$ from
principality still leaves mass $p^2q$, rather than $p^2$.
At a $q$-principal quartic row of degree $q^2r$, the corresponding masses are
$q^2r$ and $qr$. These are allowed degree patterns, not asserted
occurrences. They show why the initial mass bound cannot simply
be imposed on every common word, even at active rows.


\paragraph{Spectral positivity is not coefficient positivity.}
There is also a simple limitation to an abstract projector
argument. For an odd prime $\ell$, take the pointed metric
category on $G=C_{\ell^2}$ with quadratic form
$x\mapsto\zeta_{\ell^2}^{x^2}$. Put
$e_L=|L|^{-1}\sum_{g\in L}g$ for its subgroup $L=C_\ell$,
and define $e_G$ similarly. The rational self-adjoint element
\[
 E=1-e_L+e_G,\qquad A=\ell^2E
\]
is a spectral projection, has augmentation one before scaling,
and satisfies
\[
 A^2=\ell^2A,\quad \FPdim(A)=\ell^2,\quad
 [1]A=\ell^2-\ell+1.
\]
Nevertheless its coefficient on every $g\in L\setminus\{1\}$
is $-(\ell-1)$. Indeed $e_Le_G=e_G$, so the idempotence is
immediate, and the coefficients follow directly from the two
group averages. All twists are $\ell$-primary. Thus these
spectral identities and primary support alone do not give an
object with nonnegative coefficients. This pointed example
does not meet the perfect exact-corner hypotheses; any argument
for the desired positive projector must use more of those
hypotheses, in particular its prescribed principal-block origin.


\section{Residue blocks and categorical decomposition}\label{paper:blocks:section}\label{paper:sec:blocks}
Residue information has a global interpretation that does not depend on
the exponents $(6,6,4)$. Let $\mathcal M$ be an integral modular category
of dimension $D$, with fusion ring $R=K_0(\mathcal M)$. Let $J$ be a
nonempty set of dimension primes. At each prime choose a place in the
cyclotomic splitting field. Let $\mathscr P_J$ be the join of the full residue-block
partitions at primes in $J$: two labels are joined when they can be
connected by a chain of residue-block relations. Write
\[
 D_J=\prod_{\ell\in J}\ell^{v_\ell(D)},\qquad Q_J=D/D_J.
\]
The component containing the unit is denoted $P_J$. When $J$ is proper,
let $K$ be its nonempty complement, so $D_K=Q_J$ and $D_JD_K=D$. For a union
$B$ of joined components put
\[
 \Omega_B=\sum_{X\in B}d_X[X],\qquad D_B=\sum_{X\in B}d_X^2,
 \qquad E_B=\sum_{X\in B}e_X.
\]
The empty complementary prime set has the discrete partition and
principal component $\{\mathbf1\}$.
Coordinate projection onto a label set and multiplication by its
spectral idempotent are distinct operators. Their relation produces
the splitting theorem below.

\subsection{Integral masses and positive products}\label{paper:sec:joined-calculus}\label{paper:sec:joined}
If $E_B=\sum_Z b_{B,Z}[Z]$, character inversion gives
\begin{equation}\label{paper:block:master}
 \chi_Z(\Omega_B)=\frac{D\,\overline{b_{B,Z}}}{d_Z},
 \qquad
 [T](A\Omega_B)=d_T\sum_Z\chi_Z(A)
                         \overline{b_{B,Z}\chi_T(Z)}.
\end{equation}
The coefficients of $E_B$ are integral at every prime in $J$.
For $\ell\in J$ put $a_\ell=v_\ell(D)$ and
$m_\ell=\lfloor a_\ell/2\rfloor$.

\begin{theorem}[Simultaneous block products and masses]
\label{paper:block:joined-products}\label{paper:block:block-products}
For nonempty unions $B_1,\ldots,B_k$ of joined components, let
\[
 h_{\ell,i}=\min_{X\in B_i}v_\ell(d_X),\qquad
 g_i=\prod_{\ell\in J}\ell^{h_{\ell,i}},\qquad
 U_i=\Omega_{B_i}/g_i\in R_{\ge0}.
\]
For $k\ge2$ their product belongs to
\begin{equation}\label{paper:block:multi-product}
 \left(\prod_{\ell\in J}
 \ell^{(k-1)a_\ell-\sum_i h_{\ell,i}-(k-2)m_\ell}\right)R_{\ge0}.
\end{equation}
In particular,
\begin{equation}\label{paper:block:joined-pair}
 K_B:=D_B/D_J\in\mathbb Z_{>0},\qquad
 W_{BC}:=\Omega_B\Omega_C^*/D_J\in R_{\ge0}.
\end{equation}
If every $a_\ell$ in $J$ is even and $n_J=\sqrt{D_J}$, then
\[
              x_B=\Omega_B/n_J
\]
has the property that every product of two or more such vectors,
with duals and repetitions allowed, is an actual Grothendieck
object.  The vector $x_B$ itself need not be integral.
\end{theorem}
\begin{proof}
Fix $\ell\in J$ and put $h_Z=v_\ell(d_Z)\le m_\ell$.
The first formula in \eqref{paper:block:master} gives
\[
 v_\ell\bigl(\chi_Z(U_i)\bigr)
             \ge a_\ell-h_{\ell,i}-h_Z.
\]
In the coefficient trace formula for $U_1\cdots U_k$, each
summand therefore has valuation at least
\[
 (k-1)a_\ell-\sum_i h_{\ell,i}-(k-2)h_Z
 \ge (k-1)a_\ell-\sum_i h_{\ell,i}-(k-2)m_\ell.
\]
The coefficients are nonnegative rational integers, so the
simultaneous divisibilities give \eqref{paper:block:multi-product}.
All displayed exponents are nonnegative.

At $k=2$, restoring $g_Bg_C$ gives
$\Omega_B\Omega_C^*\in D_JR_{\ge0}$.  Its unit coefficient at
$B=C$ is $D_B$, proving \eqref{paper:block:joined-pair}.
Equivalently, mass integrality follows directly from
$\tau(E_B)=D_B/D$, which is integral at every prime in $J$.
When $a_\ell=2m_\ell$, the exponent in
\eqref{paper:block:multi-product} is $km_\ell-\sum_i h_{\ell,i}$.
Canceling the removed degree powers gives integrality of all
products of the $x_B$.  Duality preserves each full-block
partition and hence its join.
\end{proof}

\begin{corollary}[The simultaneous mass budget and connectedness]
\label{paper:block:joined-connected}
The joined components satisfy
\[
                    \sum_{B\in\mathscr P_J}K_B=Q_J.
\]
If $J$ is the full set of dimension primes, the joined partition
has one component.  Thus no nonempty proper label set is a union
of full residue blocks at every dimension prime.
A pair join may already have one component; no properness of the
three complementary-pair partitions is asserted.  Their family
budgets and height consequences are recorded together below.
\end{corollary}
\begin{proof}
Sum $D_B/D_J$ over the components.  For the full prime set,
$Q_J=1$, so its positive integral summands consist of exactly one
component.
\end{proof}


\subsubsection{Complementary modular corners}
\label{paper:sec:corners}

The joined partitions interact particularly rigidly when their prime
sets are complementary.  Let \(J,K\) be nonempty disjoint sets whose
union is the set of dimension primes.  Thus \(D_JD_K=D\).
For a label set \(B\), \(P_B\) denotes its coordinate projection,
whereas \(E_B\) denotes its spectral idempotent in the fusion algebra.
Multiplication by the latter is
\[
                         N_{E_B}=s^*P_Bs.
\]
These two projections have different roles in the next theorem.

\begin{theorem}[Complementary-join corners]\label{paper:block:complementary-corners}
Let \(B\in\mathscr P_J\) and \(C\in\mathscr P_K\).
The rectangular matrix \(s_{B,C}\) is a partial isometry:
all its nonzero singular values are one, and its rank is
\(|B\cap C|\).  In particular,
\begin{equation}\label{paper:block:corner-energy}
 \frac1D\sum_{X\in B,\,Y\in C}|S_{XY}|^2=|B\cap C|.
\end{equation}
Consequently \(S_{B,C}=0\) exactly when \(B\cap C\) is empty.

Equivalently, the normalized characters satisfy
\begin{equation}\label{paper:block:restriction-rank}
 |B\cap C|
  =\dim_{\mathbb C}\operatorname{span}
                    \{\chi_X|_C:X\in B\}.
\end{equation}
For $C=P_K$, rank one is therefore equivalent to exact agreement
of all characters $\chi_X$, $X\in B$, on every label of $P_K$.

The operator \(N_{E_B}\) preserves every complementary component
\(\mathbb C^C\), and its restriction there is an orthogonal
projection of rank \(|B\cap C|\).  Its action on the regular
component vectors and the support of its vacuum column are
\begin{equation}\label{paper:block:corner-regular}
 N_{E_B}\Omega_C=
    \begin{cases}\Omega_C,&B=P_J,\\0,&B\ne P_J,\end{cases}
 \qquad
 \operatorname{supp}(E_B)\subseteq P_K.
\end{equation}
The coefficients of \(E_B\) also vanish on labels active at any
prime in \(J\).  Every component of either partition meets the
principal component of the other.

Writing \(K_B=D_B/D_J\) and \(K_C=D_C/D_K\), one has
\[
                           |B\cap C|\le K_BK_C.
\]
If the category is perfect and fusion-simple, the stronger bound is
\begin{equation}\label{paper:block:strict-corner}
                           |B\cap C|\le K_BK_C-1.
\end{equation}
\end{theorem}
\begin{proof}
The local Weyl factors at primes in \(J\) preserve \(B\), and those
at primes in \(K\) preserve \(C\).  Any congruence factor outside the
dimension primes is scalar: its actual upper unipotent is scalar,
as are its conjugate lower unipotents, and these generate the
local group.  Thus, for a scalar \(|\kappa|=1\),
\[
 s=\kappa W_JW_K,\qquad
 P_BsP_C=\kappa W_JP_{B\cap C}W_K.
\]
The two outer factors restrict to unitaries on their respective
components.  This proves the singular-value and energy assertions.
Dividing row $X$ by $d_X/\sqrt D$ gives the restriction of
$\chi_X$, proving \eqref{paper:block:restriction-rank}.  When $C=P_K$,
these restricted characters all take value one at the unit;
rank one therefore means equality, not merely proportionality.
The same factorization also gives
\[
 N_{E_B}=W_K^*P_BW_K,\qquad
 P_CN_{E_B}P_C=W_K^*P_{B\cap C}W_K.
\]
Hence \(N_{E_B}\) commutes with every \(P_C\), and its restriction
has the stated rank.

The global regular vector \(\Omega=\sum_Xd_X[X]\) is an eigenvector
of multiplication by \(E_B\), with eigenvalue
\(\chi_0(E_B)=\mathbf1_{B=P_J}\).  Commuting with \(P_C\) gives
\eqref{paper:block:corner-regular}.  Taking its vacuum column gives
\(E_B=N_{E_B}[\mathbf1]\in\mathbb C^{P_K}\).
Vanishing at the \(J\)-active labels is the block-orthogonality
statement for each constituent full block.

For \(B=P_J\), the nonzero vector \(\Omega_C\) belongs to the range
of the restricted projection.  Its rank \(|P_J\cap C|\) is
therefore positive.  Interchange \(J,K\) for the other incidence
assertion.

Finally, \(|S_{XY}|\le d_Xd_Y\) gives
\[
 |B\cap C|\le\frac{D_BD_C}{D}=K_BK_C.
\]
In a perfect fusion-simple category every pair of nonunit labels
has strict inequality.  Each component contains a nonunit, because
its mass is divisible by \(D_J>1\), respectively \(D_K>1\).
At least one summand is therefore strict.  The masses and the
intersection cardinality are integers, giving
\eqref{paper:block:strict-corner}.
\end{proof}

\begin{corollary}[Codegree support and complementary height]
\label{paper:joined-support}
If every prime divisor of \(f_Z\) belongs to \(J\), then
\[
                   Z\in P_J,\qquad
                   S_{ZX}=0\quad(X\notin P_J).
\]
For a complementary-pair join in \(p^6q^6r^4\), every label of top
height at the omitted prime belongs to its principal component.
Nonprincipal components have omitted-prime height at most two
when that prime is sextic, or at most one when it is quartic.
\end{corollary}
\begin{proof}
At each prime in \(K\), the codegree of \(Z\) is a unit, so its
full residue block is a singleton.  Hence its \(K\)-join component
is \(\{Z\}\).  The nonzero vacuum entry and
\eqref{paper:block:corner-energy} force \(Z\in P_J\); the same identity
makes its row zero on every other \(J\)-component.
The height conclusions follow from \(f_Z=D/d_Z^2\).
\end{proof}

For the same complementary prime sets put
\[
 L=P_J\cap P_K\cap\{Z:\theta_Z=1\}.
\]
Thus $L$ contains the unit.  The next identity distinguishes
degree-square mass from the cardinality measured by the modular
corner theorem.

\begin{proposition}[Intersection masses and common labels of twist one]
\label{paper:intersection-mass}
Write $E_B=\sum_Z b_{B,Z}[Z]$ for a component of either joined
partition.  For $B\in\mathscr P_J$ and $C\in\mathscr P_K$,
\begin{equation}\label{paper:mass-factorization}
 \frac{D_{B\cap C}}{D}
   =\sum_{Z\in L}b_{B,Z}\overline{b_{C,Z}}.
\end{equation}
Consequently the matrix
\[
 H=\bigl(D_{B\cap C}/D\bigr)_{B,C}
\]
has rank at most $|L|$.  More precisely, on putting
$r_B=D_B/D$ and $c_C=D_C/D$, one has
\begin{equation}\label{paper:centered-mass-factorization}
 H-rc^{\mathsf T}
  =\left(\sum_{Z\in L\setminus\{\mathbf1\}}
                 b_{B,Z}\overline{b_{C,Z}}\right)_{B,C},
 \qquad
 \operatorname{rank}(H-rc^{\mathsf T})\le |L|-1.
\end{equation}
If $L=\{\mathbf1\}$, every complementary pair of components
meets, and its mass is exactly
\begin{equation}\label{paper:independent-mass}
 D_{B\cap C}=\frac{D_BD_C}{D}=K_BK_C,
 \qquad K_B=D_B/D_J,\quad K_C=D_C/D_K.
\end{equation}
Thus any failure of \eqref{paper:independent-mass} forces a
nonunit twist-one label in the common principal component.
Under the same hypothesis $L=\{\mathbf1\}$, a singleton component
in either partition forces the other partition to be connected.
\end{proposition}
\begin{proof}
Primitive spectral idempotents are orthogonal, so
$E_BE_C=E_{B\cap C}$.  Taking the unit coefficient gives
\[
 \tau(E_BE_C)=\frac{D_{B\cap C}}D.
\]
Character inversion gives
$b_{C,Z^*}=\overline{b_{C,Z}}$, whether or not $C$ is stable
under duality.  Therefore the same coefficient equals
$\sum_Zb_{B,Z}\overline{b_{C,Z}}$.
The complementary-corner theorem places the first coefficient
support in $P_K$ and makes it inactive at every prime of $J$;
the second lies in $P_J$ and is inactive at every prime of $K$.
Only labels in $L$ can contribute.  This proves
\eqref{paper:mass-factorization}.

The unit coordinates are $b_{B,\mathbf1}=r_B$ and
$b_{C,\mathbf1}=c_C$.  Removing their outer product proves
\eqref{paper:centered-mass-factorization}; both rank bounds follow
from the displayed matrix factorizations.  When no other coordinate
remains, multiply by $D$ and use $D_JD_K=D$ to obtain
\eqref{paper:independent-mass}.  Both masses are positive, so every
intersection is nonempty.
If one component is a singleton, it can meet every
component of the opposite partition only when that partition
has one component.
\end{proof}

\begin{corollary}[When the only common twist-one label is the unit]
\label{paper:independent-height}
Suppose $J=\{\ell\}$, $v_\ell(D)=2m$, and
$L=\{\mathbf1\}$.  If the complementary joined partition is
nontrivial, then every simple has $\ell$-height at most $m-1$,
and $\ell$ occurs in a mixed twist with every other dimension prime.
For $\ell=p$ in the exact family, every compulsory opposite
witness consequently has $p$-height two, and the full $p$-blocks
containing any opposite pair are nonprincipal and homogeneous
of height two.  No additional unit hypothesis is required.
For $p\ge5$, witnesses with pure $q$- and $r$-primary twists
also have only $p$-inactive constituents in their tensor product.

More generally, if all degrees in a full $\ell$-block $B$
are divisible by $\ell^h$, every complementary mass satisfies
\[
 v_\ell(K_C)\ge\max\{0,2h-v_\ell(K_B)\}.
\]
In particular, a homogeneous height-$m-1$ block with
$\ell$-unit $K_B$ forces every complementary mass to be
divisible by $\ell^{2m-2}$, and permits at most $\ell^2$
complementary components.  This applies to a homogeneous
dyadic defect-two block and to every equal-degree block at the least prime
of Theorem~\ref{paper:singleton}.
\end{corollary}
\begin{proof}
A label of height $m$ has $\ell$-unit codegree and hence a
singleton full $\ell$-block.  The preceding proposition would
make the complementary partition connected.  Thus no such
label exists in the nontrivial case.  If $\ell$ and another
dimension prime $t$ never occurred in the same twist,
Lemma~\ref{paper:family:lw:weyl-line} would make every $t$-active label
have $\ell$-unit codegree.  Categorical Cauchy supplies such
a label \cite[Theorem~3.9]{BNRWFinite}, a contradiction.
The compulsory witness heights
lie in $\{2,3\}$, so the sextic height bound makes them both two.
Corollary~\ref{paper:cor:proper-crossing} supplies the unit crossing,
and Proposition~\ref{paper:family:homogeneous-crossing} gives homogeneity
and nonprincipality of their blocks.  The same crossing corollary gives the
product-support assertion when $p\ge5$.

For the valuation assertion, all summands of $D_{B\cap C}$
are divisible by $\ell^{2h}$; apply
\eqref{paper:independent-mass}.  The number of components is
then bounded using $\sum_CK_C=\ell^{2m}$.
The cited packet results give unit mass respectively by the
dyadic parity identity and by $K_B=u_Z^2$ in the equal-degree case.
\end{proof}


\subsubsection{When the principal corner gives a categorical factor}

\begin{theorem}[A canonical coprime splitting criterion]
\label{paper:block:prime-splitting}
Let $\mathcal C$ be an integral modular category, and partition its
nonempty set of dimension primes into nonempty sets $J,K$.
Use the joined residue partitions $\mathscr P_J,\mathscr P_K$,
their principal components $P_J,P_K$, and the mass factors
$D_JD_K=D$.  The following conditions are equivalent:
\begin{enumerate}[label=(\roman*)]
\item $P_J\cap P_K=\{\mathbf1\}$;
\item $D_{P_J}=D_J$ and $D_{P_K}=D_K$;
\item there is a braided decomposition
\[
 \mathcal C\simeq\mathcal C_J\boxtimes\mathcal C_K,
 \qquad \operatorname{FPdim}(\mathcal C_J)=D_J,
 \quad \operatorname{FPdim}(\mathcal C_K)=D_K.
\]
\end{enumerate}
When these conditions hold, the label sets $P_J$ and $P_K$
are the two nondegenerate factors, and the factors are each
other's M\"uger centralizers.  In particular the factors are
determined by the joined partitions.
\end{theorem}
\begin{proof}
Assume (i).  Theorem~\ref{paper:block:complementary-corners} makes
$s_{P_J,P_K}$ rank one.  Its vacuum row and column are
$d_Y/\sqrt D$ and $d_X/\sqrt D$, so
\[
                   S_{XY}=d_Xd_Y
                   \qquad(X\in P_J, Y\in P_K).
\]
The two principal sets therefore centralize one another.
The corner energy identity now gives
$D_{P_J}D_{P_K}=D$.  Since $D_J\mid D_{P_J}$ and
$D_K\mid D_{P_K}$ with $D_JD_K=D$, this proves (ii).

Let $\mathcal D_J,\mathcal D_K$ be the fusion subcategories
generated by the two sets.  They still centralize one another,
so the centralizer dimension formula
\cite[Theorem~3.14]{DGNO} gives
\[
 D_{P_J}D_{P_K}
 \le\operatorname{FPdim}(\mathcal D_J)
        \operatorname{FPdim}(\mathcal D_K)
 \le D.
\]
Equality at the ends forces their label sets to be exactly
$P_J,P_K$ and makes them mutual centralizers.  Their intersection
is the unit, so they are nondegenerate; the nondegenerate
subcategory decomposition \cite[Theorem~3.13(i)]{DGNO}
proves (iii).  Equivalently, the rank-one principal spectral
projection has vacuum column
\begin{equation}\label{paper:block:positive-corner}
                  E_{P_J}=\frac{\Omega_{P_K}}{D_{P_K}}.
\end{equation}

Conversely, (ii) and the corner energy bound give
$1\le |P_J\cap P_K|\le D_{P_J}D_{P_K}/D=1$, proving (i).
Under (iii), labels and characters are products.  At a prime
$\ell\in J$, the second factor has only $\ell$-unit formal
codegrees, hence singleton residue blocks.  An $\ell$-block
of the product is consequently an $\ell$-block in the first
factor times a fixed second label.  The full-prime join in
the first factor is connected.  Thus
\[
 P_J=\operatorname{Irr}(\mathcal C_J)\times\{\mathbf1\},
 \qquad
 P_K=\{\mathbf1\}\times\operatorname{Irr}(\mathcal C_K),
\]
which proves (i) and identifies the factors canonically.
\end{proof}

\begin{corollary}[A nontrivial principal intersection is compulsory]
\label{paper:block:principal-core}
In a perfect fusion-simple integral modular category, put
\[
 c_J=|P_J\cap P_K|,
 \qquad k_J=D_{P_J}/D_J,\qquad k_K=D_{P_K}/D_K.
\]
Then
\[
                 2\le c_J\le k_Jk_K-1.
\]
If $D$ is odd, then $c_J$ is odd and
\[
                 3\le c_J\le k_Jk_K-1.
\]
Thus $k_Jk_K\ge3$ in general and $k_Jk_K\ge4$ in odd
dimension.  Every common label other than the unit has formal
codegree divisible by a prime of $J$ and by a prime of $K$.
\end{corollary}
\begin{proof}
The splitting theorem excludes $c_J=1$, since $D_J,D_K>1$.
The strict upper bound is the complementary-corner energy bound:
each principal component contains a nonunit, and strictness for
that pair makes $c_J<k_Jk_K$.  Both principal components are
stable under duality.  In odd dimension the unit is the only
self-dual simple, so their intersection has odd cardinality.

If the codegree of a common nonunit $Z$ were supported on $J$,
every full residue block at a prime in $K$ containing $Z$ would
be its singleton.  Thus its $K$-join component would be the
singleton $\{Z\}$, contradicting $Z\in P_K$.  Interchange
$J,K$ for the other assertion.
\end{proof}

For the exact family, take respectively
$J=\{p\},\{q\},\{r\}$ and $K$ the complementary pair.
Each of the three intersections therefore contains a nonunit
whose codegree meets both sides of that prime partition.
The conclusion is a uniform categorical alternative: an
intersection containing only the unit constructs the forbidden
coprime factor directly.  A larger intersection supplies no
positivity assertion for its projector.

\subsubsection{The principal joined projector}

\begin{proposition}[A virtual object at the complementary prime powers]
\label{paper:block:principal-projector}
Let $P_J$ be the component of $\mathscr P_J$ containing the unit,
and set $k=D_{P_J}/D_J$ and $Q=Q_J$.  Then
\[
                    A_J=Q E_{P_J}
\]
is an integral self-dual virtual element of $R$, satisfying
\begin{equation}\label{paper:block:projector-identities}
 A_J^2=Q A_J,\qquad
 \operatorname{FPdim}(A_J)=Q,\qquad
 [\mathbf1]A_J=k,\qquad
 \sum_{X\ne\mathbf1}([X]A_J)^2=k(Q-k).
\end{equation}
For proper $J$, write $K$ for the complementary set of dimension
primes.  The coefficient support satisfies
\[
 \operatorname{supp}(A_J)
 \subseteq P_K\cap\{X:\theta_{X,\ell}=1\text{ for every }\ell\in J\}.
\]
In particular every nonunit whose codegree is supported on $J$
has zero coefficient.  For a complementary-pair join in the exact
family, the nonunit support lies in the principal block at the
omitted prime, so its height there is below the top height.
If $J$ is the full prime set, $A_J=1$.
If the category is fusion-simple and $P_J$ is proper, at least one
nonunit coefficient of $A_J$ is negative.

For a complementary pair in $p^6q^6r^4$, $Q$ is the full fourth
or sixth power of the omitted prime.  Hence either that pair's
join is already connected, or the category has the mixed-sign
virtual element in \eqref{paper:block:projector-identities}, supported on
pure twists of the omitted prime, allowing twist one.
\end{proposition}
\begin{proof}
We first justify rationality at the fixed chosen places.
Unsigned Galois covariance and preservation of dimensions give
\[
 \chi_{\pi_\sigma X}(T)
      =\sigma\chi_X(T)=\chi_X(\pi_\sigma T).
\]
If $X$ and $Y$ have equal reductions on every basis element at the
chosen $\ell$-place, the right-hand side shows that
$\pi_\sigma X$ and $\pi_\sigma Y$ do too, at that same place.
Thus $\pi_\sigma$ permutes each full-block partition and therefore
its join.  It fixes the unit label because the dimension character
is rational.  Consequently $P_J$ is Galois stable and its
idempotent has rational coefficients.  This uses the symmetric
integral modular character table; it is not a general assertion
about arbitrary commutative orders.

The primitive-idempotent formula shows that $D E_{P_J}$ has
algebraic-integer coefficients, now rational and hence integral.
Local integrality at every prime in $J$ removes the entire $D_J$
part of the possible denominator.  Thus $Q E_{P_J}\in R$.
Complex conjugation also preserves $P_J$, giving self-duality.
Idempotency and the dimension character give the first two
identities in \eqref{paper:block:projector-identities}.  Moreover
\[
 \tau(A_J)=Q\frac{D_{P_J}}D=k,\qquad
 \|A_J\|^2=Q^2\tau(E_{P_J})=Qk.
\]
Subtract its unit square $k^2$ to obtain the last identity.
For proper $J$, Theorem~\ref{paper:block:complementary-corners}
gives the support assertion.  A nonunit of codegree supported on
$J$ has a singleton complementary component and cannot be in $P_K$.
The principal-height bound gives the stated family height cap.
The full-prime case follows from connectedness of its join.

Suppose $P_J$ is proper and $A_J$ has no negative coefficient.
The nonnegative self-dual idempotent $E_{P_J}$ has fusion-closed
support: every constituent of a product of supported simples
has a positive coefficient in its square.  Fusion simplicity
makes this support either the unit alone or all simples.
The first case gives $E_{P_J}=1$, contrary to properness.
In the second case its multiplication matrix is strictly positive:
for any pair of simples some supported simple has a nonzero
fusion entry between them.  Perron--Frobenius then makes the
idempotent matrix rank one.  Since its character eigenvalues
are the indicator of $P_J$, this forces $P_J=\{\mathbf1\}$.
But $D_J\mid D_{P_J}$ and $D_J>1$ exclude a unit-only joined
component.  Thus a negative nonunit coefficient is necessary.
\end{proof}
For the exact family, the joint mass budget and the omitted-prime
height constraints take the following common form:
\[
 \begin{array}{c|c|c|c|c}
 \ell&J&Q_J&h_{\rm top}&h_{\rm remaining}\\ \hline
 p&\{q,r\}&p^6&3&\le2\\
 q&\{p,r\}&q^6&3&\le2\\
 r&\{p,q\}&r^4&2&\le1
 \end{array}
\]
Here $\ell$ is the omitted prime, $Q_J=\sum_BK_B$, and every
label of height $h_{\rm top}$ at $\ell$ belongs to $P_J$.
The last bound holds for every label in a nonprincipal joined
component.  It also holds for the nonunit coefficient support of
$A_J$, which lies more precisely in $P_\ell$ and has pure
$\ell$-primary twist, allowing twist one.

\begin{corollary}[Principal coefficient congruences]
\label{paper:principal-coefficients}
In Proposition~\ref{paper:block:principal-projector}, let $J$ contain
all dimension primes except $\ell$, and write
$A_J=\sum_Xa_X[X]$, $h_X=v_\ell(d_X)$ and
$u_X=d_X/\ell^{h_X}$. Then
\[
 \ell^{h_X}\mid a_X\quad\text{for every }X,\qquad
 \frac{a_X}{\ell^{h_X}}\equiv ku_X\pmod\ell
                  \quad(X\in P_\ell).
\]
If $\ell\nmid k$, the coefficient support is exactly $P_\ell$,
and every label in that principal block has pure $\ell$-primary
twist, allowing twist one. If $\ell\mid k$, then
$\ell^{h_X+1}\mid a_X$ for every $X$.
In particular the latter alternative is forced if $P_\ell$
contains a label with a nontrivial complementary twist component.
\end{corollary}
\begin{proof}
Symmetry and the primitive-idempotent formula give
\[
 \frac{a_X}{d_X}
   =\frac1{D_J}\sum_{Y\in P_J}d_Y\chi_X(Y^*).
\]
This rational number is $\ell$-integral because $D_J$ is an
$\ell$-unit and the characters are algebraic integers. This
proves the first divisibility. For $X\in P_\ell$, reducing
at the chosen $\ell$-place replaces $\chi_X(Y^*)$ by $d_Y$,
so the displayed rational number is congruent to
$D_{P_J}/D_J=k$. Multiplication by $u_X$ gives the stated
congruence. Outside $P_\ell$ the coefficient is zero by the
support assertion of Proposition~\ref{paper:block:principal-projector};
inside it, the congruence gives both alternatives. The same
support assertion supplies the twist statement.
\end{proof}
The extra factor when $\ell\mid k$ does not automatically
iterate: the coefficient formula for $A_J/\ell$ has denominator
$\ell D_J$, which consumes the principal congruence just used.

\subsection{Fusion transitions, expansion and integral lattices}\label{paper:sec:fusion}

For the joined partition $\mathscr P_J$, with its
normalization $D_J$, components $B$, and masses $K_B$, the actual
block products give integer fusion transitions.

\begin{proposition}[Integral transitions and their stabilizer]
\label{paper:block-transitions}\label{paper:stabilizer}
For $B,C\in\mathscr P_J$ define $M_Y(B,C)=[Y]W_{BC}$.  Then
\[
 M_Y(B,C)=D_J^{-1}\!\sum_{X\in B,\,X'\in C}
 d_Xd_{X'}N_{YX'}^X\in\mathbb Z_{\ge0}.
\]
Its row sums are $d_YK_B$, its column sums are $d_YK_C$,
$M_Y(B,C)=M_{Y^*}(C,B)$, and
$M_{\mathbf1}(B,C)=\delta_{BC}K_B$.  An entry is zero exactly
when $Y$ sends no member of $C$ to a constituent in $B$.

For a proper block, set $t_B=D_B/D$ and
\[
 F_B=\frac1{t_B}\sum_Z\frac{|b_{B,Z}|^2}{d_Z}[Z].
\]
This positive formal element has dimension one and
\[
 \chi_Y(F_B)=\frac{M_Y(B,B)}{d_YK_B}\in[0,1].
\]
Equality to one means that $Y$ preserves $B$, equivalently
that it centralizes every coefficient-support label of $E_B$.
The equality labels form the proper fusion subcategory
$\mathcal C_{\mathcal C}(\langle Z:b_{B,Z}\ne0\rangle)$.
\end{proposition}
\begin{proof}
The product theorem gives integrality; reciprocity gives the
formula for $M_Y$, duality and the zero criterion.  With
$\Omega=\sum_Xd_X[X]$ one has
$\sum_CW_{BC}=K_B\Omega$, giving the row sums; duality gives
the column sums.

Since $E_B$ is a self-adjoint idempotent,
$\sum_Z|b_{B,Z}|^2=t_B$.  Verlinde's formula gives
\[
 \chi_Y(F_B)=\sum_Z\frac{|b_{B,Z}|^2}{t_B}
                  \frac{S_{YZ}}{d_Yd_Z}
             =\frac{M_Y(B,B)}{d_YK_B}.
\]
The internal fusion mass proves the interval bound and the
preservation criterion.  In the convex sum, equality to one
is precisely centrality on the support.  If every simple
preserved $B$, tensoring a member by its dual would put the
unit in $B$ and then put every simple there, contrary to
propriety.
\end{proof}

\subsubsection{Strict expansion}
\label{paper:sec:block-expansion}

\begin{theorem}[Strict expansion in a minimal table]
\label{paper:principal-proper}\label{paper:strict-contraction}
\label{paper:weighted-expansion}\label{paper:deterministic-transition}
Assume the category is perfect and fusion-simple.  When
$J=\{p\}$, every principal block is proper.  For a nonunit
$Y$ and a nonempty proper set $\mathcal S\subset\mathscr P_J$, let
\[
 \mathcal N_Y(\mathcal S)=
       \{B:M_Y(B,C)>0\text{ for some }C\in\mathcal S\}.
\]
Then
\begin{equation}\label{paper:expansion}
 \sum_{B\in\mathcal N_Y(\mathcal S)}K_B
                  \ge\sum_{C\in\mathcal S}K_C+1.
\end{equation}
In particular a deterministic transition from a proper component $C$
to $B$
requires $K_B>K_C$.  If the partition is nontrivial, every
maximal-mass component splits under $Y$, and no nonunit
permutes the partition.  For every proper component the
exact integer gap is
\begin{equation}\label{paper:escape}
 E_Y(B):=d_YK_B-M_Y(B,B)\ge1,
 \qquad \chi_Y(F_B)\le1-\frac1{d_YK_B}.
\end{equation}
\end{theorem}
\begin{proof}
For $J=\{p\}$, the dimension equation supplies a $p$-free nonunit $Z$.
If all characters were principal, every $\chi_X(Z)$ would
have nonzero residue $d_Z$; symmetry would contradict the
nonunit-row zero theorem.

By Lemma~\ref{paper:family:2.4}, $|\chi_Z(Y)|<d_Y$ for $Z\ne0$.
Unitary diagonalization of $N_Y$ therefore gives
$\|N_Yv\|<d_Y\|v\|$ off the regular line.  Apply this to
$v=\Omega_{\mathcal S}$, whose squared norm is
$D_{\mathcal S}=\sum_{C\in\mathcal S}D_C$.  Its image is
supported on $\mathcal T=\mathcal N_Y(\mathcal S)$ and has
dimension $d_YD_{\mathcal S}$, so Cauchy--Schwarz gives
\[
 d_Y^2D_{\mathcal S}^2
 \le D_{\mathcal T}\|N_Y\Omega_{\mathcal S}\|^2
 <D_{\mathcal T}d_Y^2D_{\mathcal S}.
\]
Thus $D_{\mathcal T}>D_{\mathcal S}$; divisibility by $D_J$
gives \eqref{paper:expansion}.  Its singleton case proves the
deterministic-transition statements and excludes preservation
of a proper block.  The outward mass is therefore a positive
integer, proving \eqref{paper:escape}.
\end{proof}

\begin{proposition}[One projector and averaging classes of different sizes]
\label{paper:averaging-products}\label{paper:averaging}
Let $E=E^*=E^2$ lie in the complexified commutative fusion ring.
Suppose its multiplication operator preserves the coordinate spaces
of disjoint label sets $B_i$, and restricts on $B_i$ to $J_{k_i}/k_i$,
where $k_i=|B_i|$. Assume that all labels in $B_i$ have degree $d_i$.
Put $B=\bigcup_iB_i$, $R_i=\sum_{X\in B_i}[X]$,
$g_{ij}=\gcd(k_i,k_j)$, and $l_{ij}=\operatorname{lcm}(k_i,k_j)$.
Then
\begin{align}
 H_{ij}:=\frac{R_iR_j^*}{l_{ij}}&\in R_{\ge0},
 &R_iX^*&=\frac{k_i}{g_{ij}}H_{ij}\quad(X\in B_j),
 \label{paper:individual-average-product}\\
 [\mathbf1]H_{ij}&=\delta_{ij},
 &\operatorname{FPdim}(H_{ij})&=g_{ij}d_id_j.\notag
\end{align}
Every rectangle $(N_Y)_{B_i,B_j}$ has constant row sums
$\alpha_Y(i,j)$ and constant column sums. More precisely,
\[
 \alpha_Y(i,j)=\frac{k_j}{g_{ij}}[Y]H_{ij},\qquad
 \text{column sum}=\frac{k_i}{g_{ij}}[Y]H_{ij}.
\]
Write $A_Y=(\alpha_Y(i,j))$, extended linearly to the fusion ring.
For positive objects $U,V$,
\begin{equation}\label{paper:averaging-submultiplicativity}
                 A_{UV}-A_UA_V\in\operatorname{Mat}(\mathbb Z_{\ge0}).
\end{equation}
For $U=Y^*$ and $V=Y$, conjugation of this difference by
$\operatorname{diag}(\sqrt{k_i})$ makes it positive semidefinite;
its $j$th diagonal entry is
\begin{equation}\label{paper:averaging-leakage}
 \bigl(A_{Y^*Y}-A_{Y^*}A_Y\bigr)_{jj}
       =\frac1{k_j}\sum_{T\notin B}([T]YR_j)^2.
\end{equation}
An escape from $B$ makes this diagonal a positive integer.
If the category is perfect and fusion-simple, every proper class
$B_i$ satisfies $[Y]H_{ii}\le d_Y-1$ for nonunit simple $Y$.
If every $k_i$ is a square, the vectors $R_i/\sqrt{k_i}$ generate
an integral fusion-stable lattice in their rational module span,
and these vectors are orthonormal.

For representatives $X_i\in B_i$, the positive products also obey
\[
 H_{ij}=g_{ij}E[X_i][X_j]^*,\qquad
 H_{ij}H_{ab}=g_{ij}[X_i][X_j]^*H_{ab}.
\]
In particular, when all class sizes equal $k$, every product of
two averaging objects is coefficientwise divisible by $k$.
\end{proposition}
\begin{proof}
Multiplication by $E$ commutes with every $N_Y$. Its self-adjointness
and coordinate invariance give, on each rectangle $M$,
\[
                 \frac1{k_i}J_{k_i}M=M\frac1{k_j}J_{k_j}.
\]
Thus every row has one common sum $r$ and every column has one common
sum $c$, with $k_ir=k_jc$. Since these sums are nonnegative integers,
$r=(k_j/g_{ij})a$ and $c=(k_i/g_{ij})a$ for a nonnegative integer $a$.
Frobenius reciprocity identifies $c$ with $[Y]R_iX^*$ for
$X\in B_j$. Summing over $X$ gives $[Y]R_iR_j^*=l_{ij}a$.
This proves the object, divisibility, and row-sum assertions.
The unit coefficient follows from the disjoint supports; the degree
formula follows by applying $\operatorname{FPdim}$.

Let $v_i=R_i/\sqrt{k_i}$ and let $Q$ have these orthonormal columns.
Writing $P_B$ for the coordinate projection onto $B$, constant row
sums give
\[
 P_BN_YR_j=\sum_i\alpha_Y(i,j)R_i,\qquad
 T_Y:=Q^*N_YQ=D_k^{1/2}A_YD_k^{-1/2},
 \qquad D_k=\operatorname{diag}(k_i).
\]
Inserting $1=P_B+(1-P_B)$ shows that the difference for $U,V$
has nonnegative entries after this positive diagonal conjugation,
and hence before it. It is integral because every $A_Y$ is integral.
For $U=Y^*$ the difference is the Gram matrix
\[
 Q^*N_{Y^*}(1-P_B)N_YQ.
\]
This proves the stated diagonal formula, positive semidefiniteness,
and integer escape gap. The common degree on $B_i$ gives $[Y]H_{ii}\le d_Y$.
Equality would leave no tensor-product mass outside $B_i$ and
give $N_YR_i=d_YR_i$. The strict Perron bound in a perfect
fusion-simple category excludes this eigenvector on a proper support.

If $k_i=a_i^2$ and $k_j=a_j^2$, then
$l_{ij}/\sqrt{k_ik_j}=a_ia_j/\gcd(a_i,a_j)^2$ is an integer.
Consequently $R_iR_j^*/\sqrt{k_ik_j}$ is integral, and
\[
 \langle Av_i,Cv_j\rangle
      =\tau\left(AC^*\frac{R_iR_j^*}{\sqrt{k_ik_j}}\right)\in\mathbb Z
      \qquad(A,C\in R).
\]
The generated module lies in a fixed rational dilation of $R$, so is
an integral lattice in its rational span.

Finally $E[X_i]=R_i/k_i$ and $E[X_j]^*=R_j^*/k_j$.
Idempotence gives $H_{ij}=g_{ij}E[X_i][X_j]^*$, and every
$H_{ab}$ is fixed by $E$. This proves the product identity.
\end{proof}

\begin{lemma}[A fixed constituent in a positive integral object]
\label{paper:averaging-witness}\label{paper:fixed-adjoint}
Let $R$ be a commutative integral fusion ring, let $\ell$ be a
prime, and let a finite $\ell$-group $G$ act by
dimension-preserving based-ring automorphisms. Suppose
$H\in R_{\ge0}$ is $G$-invariant and
\[
             [\mathbf1]H=1,\qquad
             \ell\mid\operatorname{FPdim}(H).
\]
Writing $h_Y=[Y]H$, there is a nonunit simple constituent $Y$
such that
\[
                         GY=Y,\qquad \ell\nmid d_Yh_Y.
\]
At $\ell=2$, if $H^*=H$, one may also require $Y=Y^*$.
If $R$ is the Grothendieck ring of a braided pivotal fusion
category and $H=XX^*$ for a simple $X$, this self-dual choice
has second Frobenius--Schur indicator $\nu_2(Y)=+1$.

Suppose specifically that a self-adjoint idempotent $E$ averages
a class $Q$ of $k$ equal-degree labels as in
Proposition~\ref{paper:averaging}, with $\ell\mid k$, and
\[
 R_Q=\sum_{X\in Q}[X],\qquad H=H_Q=R_QR_Q^*/k.
\]
If $H_Q$ is $G$-invariant, as holds when $Q$ is $G$-stable,
the selected constituent additionally satisfies $E[Y]\ne[Y]$.
For every $X\in Q$, the total multiplicity of $Y\otimes X$
returning to $Q$ is the same integer $h_Y$.
\end{lemma}
\begin{proof}
Dimensions and coefficients of $H$ are constant on $G$-orbits.
Reduce $\operatorname{FPdim}(H)=\sum_T h_Td_T$ modulo
$\ell$. Every nontrivial orbit contributes zero, while the
unit contributes one and the left side is zero. Some nonunit
fixed label therefore has $\ell\nmid d_Yh_Y$.

At two, duality commutes with the based-ring automorphisms.
If $H^*=H$, adjoining duality gives a two-group fixing $H$,
so the same selection makes $Y$ self-dual. For $H=XX^*$,
the selected multiplicity $N_{XX^*}^{Y}$ is odd. The parity
theorem of Fuchs--Runkel--Schweigert \cite{FRSEvenFusion},
applied to the duality-stable list $(X,X^*,Y)$, makes
$\dim\operatorname{Hom}(\mathbf1,X\otimes X^*\otimes Y)$
even if $\nu_2(Y)=-1$. This dimension is $N_{XX^*}^{Y}$,
so $\nu_2(Y)=+1$.

For the averaging application, the averaging proposition gives
\[
 H_Q=kE[X][X]^*,\qquad [\mathbf1]H_Q=1,
 \qquad \operatorname{FPdim}(H_Q)=kd_X^2\quad(X\in Q).
\]
If $E[T]=[T]$, self-adjointness gives
\begin{equation}\label{paper:tame-coefficient}
                   h_T=k[T]XX^*\in\ell\mathbb Z.
\end{equation}
Such a label cannot be the selected $Y$. The return multiplicity
is the common column-sum identity of
Proposition~\ref{paper:averaging}.
\end{proof}

In the modular applications, square dimension makes the local
Galois tori unsigned. Their permutations preserve dimensions and
fusion coefficients:
\begin{equation}\label{paper:galois-fusion}
                       N_{\sigma A,\sigma B}^{\sigma C}=N_{AB}^{C}.
\end{equation}
Indeed, apply $\sigma$ to Verlinde's formula and use unsigned
covariance, the rational vacuum row and the rational dimension.
Thus these permutations act by based-ring automorphisms; no tensor
autoequivalence is required.

\begin{remark}[A size obstruction for connected \'etale algebras]
\label{paper:etale-size}
An algebra is connected \'etale if it is commutative and separable
and $\dim\operatorname{Hom}(\mathbf1,A)=1$.
In a nondegenerate integral braided category, such an algebra satisfies
\[
 \operatorname{FPdim}(\mathcal C_A^0)
      =\frac{D}{\operatorname{FPdim}(A)^2}\in\mathbb Z_{>0},
\]
where $\mathcal C_A^0$ is the braided fusion category of local
$A$-modules \cite[Corollary~3.32]{DMNO}. The displayed rational
number is an algebraic integer. Hence
$2v_\ell(\operatorname{FPdim}A)\le v_\ell(D)$.
For an equal-degree averaging class of size $k=\ell$ or
$\ell^2$, with $v_\ell(D)=2m$, $m\ge2$, and degrees of
$\ell$-height $m-1$, the associated object $H_{ii}$ has
dimension valuation $2m-2+v_\ell(k)>m$. These particular objects therefore admit
no connected \'etale algebra structure. This obstruction does not
exclude a noncommutative separable algebra structure or a smaller
algebra obtained by a different construction.
\end{remark}

\subsubsection{The defect of compressed fusion multiplication}

Let $V$ have orthonormal columns
$v_B=\Omega_B/\sqrt{D_B}$, put $P=VV^*$ and
$K=\operatorname{diag}(K_B)$.  The compressed fusion operators are
\[
              T_Y=V^*N_YV=K^{-1/2}M_YK^{-1/2},
\]
extended linearly to all fusion-ring elements.

\begin{proposition}[Positive defect and its principal-block gap]
\label{paper:compression-defect}
One has
\begin{equation}\label{paper:compression-gram}
 H_Y:=T_{YY^*}-T_YT_Y^*
       =V^*N_Y(I-P)N_Y^*V\succeq0.
\end{equation}
If $B_0$ is the principal block, define
\[
 a_Y=\frac{d_Ye_{B(Y)}}{\sqrt{D_{B(Y)}}}
                  -\frac{T_Ye_{B_0}}{\sqrt{D_{B_0}}}.
\]
Then $H_Y\succeq a_Ya_Y^*/(1-D_{B_0}^{-1})$.  In a
perfect fusion-simple table with proper $B_0$, every nonunit $Y\in B_0$ satisfies
\[
 (H_Y)_{B_0B_0}\ge
 \frac{E_Y(B_0)^2}{K_{B_0}^2(D_{B_0}-1)}>0.
\]
\end{proposition}
\begin{proof}
Insert $I=P+(I-P)$ between $N_Y$ and $N_Y^*=N_{Y^*}$.
For $w=(I-P)[\mathbf1]$ one has
$\|w\|^2=1-D_{B_0}^{-1}>0$ and $V^*N_Yw=a_Y$.
The projection onto $\mathbb Cw$ is bounded above by $I-P$;
applying $V^*N_Y$ and its adjoint proves the rank-one bound.
For $Y\in B_0$, the $B_0$ coordinate of $a_Y$ is
$E_Y(B_0)/(K_{B_0}\sqrt{D_{B_0}})$, giving the last inequality.
\end{proof}
The equality $H_Y=0$ says that the block-vector span reduces
the one normal operator $N_Y$.  Tensor generation by $Y$
does not imply invariance under its separate simple constituents.
For the single-prime partitions and the compulsory witnesses
in $P_q$ and $P_r$, the respective principal-block gaps are
strictly positive.  A joined partition may already have one
component, so that conclusion is not automatic for a join.

\subsubsection{Integral lattices from divided block sums}

\begin{proposition}[Fusion-stable lattices and unimodular gluing]
\label{paper:integral-lattice}\label{paper:gluing}
Suppose $D_J=n_J^2$.  For any collection $\mathcal A\subseteq\mathscr P_J$,
put
\[
 x_B=\Omega_B/n_J,\quad L=\sum_{B\in\mathcal A}Rx_B,
 \quad L^\#=\{v\in R_{\mathbb Q}:\langle v,L\rangle\subseteq\mathbb Z\},
 \quad R_0=R\cap L^\#,\quad\Lambda=R_0+L.
\]
Then $L$ is an integral fusion-stable lattice in its rational
span, with $\|x_B\|^2=K_B$.  Its character values vanish on
labels active at any prime in $J$.  The lattice $\Lambda$ is full-rank,
integral, unimodular and stable under $R$; it is duality-stable
when $\mathcal A$ is.  If $L\not\subset R$, then
$\Lambda\ne R$ and $[\mathbf1]\notin\Lambda$.
If $\mathcal A=\mathscr P_J$ and $\Lambda$ has an orthonormal
$\mathbb Z$-basis in which every simple fusion operator is
entrywise nonnegative, then
\begin{equation}\label{paper:gluing-rank}
                         \operatorname{rank}R\le Q_J.
\end{equation}
For the full prime set this is impossible unless $R$ has rank one.
\end{proposition}
\begin{proof}
The module lies in $n_J^{-1}R$, is finitely generated, and
\[
 \langle Ax_B,Cx_E\rangle=\tau(AC^*W_{BE})\in\mathbb Z.
\]
Its spectral support follows from the block-transform formula
and $b_{B,Z}=0$ at a label active at any prime in $J$.
Adjointness of fusion multiplication makes $R_0$, and hence
$\Lambda$, stable under $R$.  All pairings in $R_0+L$ are
integral.  Let $F=R+L$, a full lattice.  Since the simple
lattice $R$ is unimodular, $R_0=F^\#$; taking pairing duals
and using integrality of $L$ gives
\[
 \Lambda^\#=F\cap L^\#=(R+L)\cap L^\#=R_0+L=\Lambda.
\]
This proves full rank and unimodularity; the duality assertion
is immediate.  Finally a unit in $\Lambda$ would put all of
$R$ in it by fusion stability, and equal covolumes would give
$\Lambda=R$, impossible if it contains a nonintegral member of $L$.

For the last assertion, using all components puts
$v=\Omega/n_J=\sum_Bx_B$ in $\Lambda$, with
$\|v\|^2=Q_J$. In the proposed basis its coordinates $a_i$
are integers, and regular multiplication has matrix
\[
             N_\Omega=\sum_Xd_XN_X=n_J^2(a_i a_j)_{i,j}.
\]
Its diagonal entries are at least one, since every summand is
nonnegative and $N_{\mathbf1}=I$. Thus every $a_i$ is nonzero,
and $\operatorname{rank}R\le\sum_i a_i^2=Q_J$.
The full prime set has $Q_J=1$.
\end{proof}
The bound concerns the canonical inner product of this gluing
construction. It does not exclude a different lattice or metric,
and a nonnegative integral representation would still require
module-category coherence.

\subsection{Homogeneous blocks, height congruences and packets}\label{paper:sec:homogeneous}

\begin{theorem}[Dense local Weyl matrix and an inactive integer packet]
\label{paper:lw:homogeneous}
Assume $D=n^2$ is square.  Let $\ell$ be any prime, $a=v_\ell(n)>0$, and let $B$ be a full
$\ell$-block with $v_\ell(f_X)=2$ for every $X\in B$.  Write
\[
 n=\ell^a n',\qquad d_X=\ell^{a-1}u_X\quad(X\in B),
 \qquad
 v=\frac n{\ell^{a-1}}P_BW_{\ell'}\mathbf e_0.
\]
Then there are integers $c_X$, supported on the labels with
$\theta_{X,\ell}=1$, such that
\begin{equation}\label{paper:lw:packet}
 v_X=\ell u_Xc_X,\qquad W_\ell v=u,\qquad
 \sum_{X\in B}u_X^2=\ell^2\sum_{X\in B}u_X^2c_X^2,
 \qquad \ell\nmid S_B:=\sum_{X\in B}u_X^2c_X.
\end{equation}
At every $\ell$-place,
\begin{equation}\label{paper:lw:rank-one}
 \ell(W_\ell)_{XY}\equiv S_B^{-1}u_Xu_Y\pmod{\mathfrak l}
 \qquad(X,Y\in B).
\end{equation}
Every entry of $W_\ell|_B$ is nonzero and has normalized valuation
exactly $-1$.  All labels in $B$ have the same prime-to-$\ell$
twist.  When $\ell=2$,
\begin{equation}\label{paper:lw:dyadic-mass}
 v_2\!\left(\sum_{X\in B}u_X^2\right)=2,
 \qquad |B|\equiv4\pmod8.
\end{equation}
\end{theorem}
\begin{proof}
Theorem~\ref{paper:lw:columns} gives
$k_{\ell',X}=n(W_{\ell'}\mathbf e_0)_X/d_X\in\ell\mathbb Z$.
Set $c_X=k_{\ell',X}/\ell$.
The complementary Weyl commutes with the actual $\ell$-twist
and that twist fixes $\mathbf e_0$; hence $c_X$ vanishes at active labels.
This uses the actual twist, so is valid without a dyadic anomaly
restriction.  Block invariance and $W_\ell W_{\ell'}=s$ give
$W_\ell v=u$.  Unitarity proves the sum-of-squares identity.

For $X\in B$ and any label $Z$,
\[
 \ell s_{XZ}=\frac{u_X}{n'}\chi_X(Z).
\]
These rows are integral at $\ell$ and their reductions are
proportional to $u_X$, by the definition of a full block.
The columns indexed by $B$ of $W_{\ell'}^{-1}$ are integral at
$\ell$ by \eqref{paper:lw:foreign-bound}.  Thus
$\ell W_\ell|_B=(\ell sW_{\ell'}^{-1})|_B$ is integral and has
reduction of the form $u b^T$.  Symmetry of $W_\ell$ makes this
$\lambda u u^T$ for a scalar $\lambda$ in the residue field.
Reducing $(\ell W_\ell)(v/\ell)=u$ now gives
\[
 \lambda\sum_{X\in B}u_X^2c_X=1.
\]
Since the sum is a rational integer, it is an $\ell$-unit and
$\lambda=S_B^{-1}\in\mathbb F_\ell^\times$.  This proves the exact
valuation and nonvanishing claims.  Commutation of $W_\ell$ with
every complementary diagonal twist then forces equality of those
twists throughout $B$.

At two, $\sum u_X^2c_X^2\equiv\sum u_X^2c_X\equiv1\pmod2$.
Equation \eqref{paper:lw:packet} gives exact valuation two, and
$u_X^2\equiv1\pmod8$ gives the rank congruence.
\end{proof}

\begin{proposition}[Height-two crossings force homogeneous blocks]
\label{paper:family:homogeneous-crossing}
Suppose \(X,Y\) in Proposition~\ref{paper:family:family:crossings} both
have \(p\)-height two.  The full \(p\)-blocks containing
\(X\) and \(Y\) are
nonprincipal and homogeneous of height two.
Each such block has one common prime-to-\(p\) actual twist
and satisfies the local Weyl conclusions of
Theorem~\ref{paper:lw:homogeneous}.
\end{proposition}
\begin{proof}
Corollary~\ref{paper:cor:proper-crossing} supplies a place over
$p$ at which $z$ is a unit.  At that place,
\(\chi_X(Y)=(\lcm(d_X,d_Y)/d_X)z\) is a unit.
For a label \(T\) of smaller \(p\)-height, LCM divisibility
would make \(\chi_T(Y)\) vanish modulo that place.
Such a label cannot belong to the full block of \(X\).
A label of height three has unit \(p\)-codegree and its
block is a singleton, so it cannot lie in the block of \(X\)
either.  This proves homogeneity; interchange \(X,Y\) for
the other block.  Since \(d_Y\) is divisible by \(p\),
the unit value \(\chi_X(Y)\) excludes principality.
The dense local Weyl matrix in Theorem~\ref{paper:lw:homogeneous}
commutes with complementary twists.  Its nonzero entries
force those twists to agree throughout the block.
\end{proof}

Corollary~\ref{paper:two-neighbors} gives a sufficient condition
for choosing two height-two witnesses.  Mixed-height \((2,3)\)
crossings also have a unit, but the preceding two-block conclusion
requires both heights to be two.

\begin{corollary}[An inactive label with a tame local field]
\label{paper:lw:inactive-fixed}
Every block in Theorem~\ref{paper:lw:homogeneous} contains an inactive
label $Z\in B\cap P_{\ell'}$ with $\ell\nmid c_Z$ fixed by wild
$\ell$-inertia.  Here $P_{\ell'}$ is the principal join at the
other dimension primes.
For odd $\ell$, its local Weyl column belongs to
$\mathbb Q(\zeta_\ell)$.  At two, the whole inertia group fixes
such a $Z$, its character field is unramified at two, and
$W_2\mathbf e_Z$ is rational.
\end{corollary}
\begin{proof}
The vector $P_BW_{\ell'}\mathbf e_0$ is supported on
$B\cap P_{\ell'}$ and fixed by the $\ell$-local tori.
The dimensions, and hence the $c_X$, are constant on their label
orbits.  Nontrivial wild-inertia orbits have size divisible by
$\ell$; the unit sum $S_B$ therefore has a fixed-label contribution
with $\ell\nmid c_Z$.  At two all local inertia is a two-group,
so the same argument uses the entire group.  Finally
$\sigma_a(W_\ell)\mathbf e_Z=W_\ell H_{a,\ell}^{-1}\mathbf e_Z$ and the entries
of $W_\ell$ lie in the local cyclotomic field.  Taking invariants
gives the asserted column fields.
\end{proof}

\begin{corollary}[An actual adjoint witness from every homogeneous block]
\label{paper:homogeneous-adjoint}
In Theorem~\ref{paper:lw:homogeneous}, assume $v_\ell(n)\ge2$.
Some inactive $X\in B\cap P_{\ell'}$ has a nonunit simple
constituent $Y$ in $XX^*$ such that
\[
               \ell\nmid d_YN_{XX^*}^{Y},\qquad
               Y\text{ is fixed by wild }\ell\text{-inertia}.
\]
For odd $\ell$, $K_Y$ is tame at $\ell$ and
$\operatorname{ord}(\theta_{Y,\ell})\mid\ell$.
At two, one may choose
\[
 Y=Y^*,\qquad \nu_2(Y)=1,\qquad
 K_Y\text{ real and unramified at two},\qquad
 \operatorname{ord}(\theta_{Y,2})\mid8.
\]
The conclusion allows every sign pattern and every complementary
principal intersection. In a perfect fusion-simple category,
$N_{XX^*}^{Y}<d_Y$, and at two the gap is at least two.
In the ordered exact family with $\ell=p$, the degree is
$d_Y=q^a r^b$, $1\le a\le3$, $1\le b\le2$.
\end{corollary}
\begin{proof}
Corollary~\ref{paper:lw:inactive-fixed} supplies such an inactive
$X$ fixed by wild inertia, and by full inertia at two.
Its positive height makes $\ell\mid d_X$. Apply
Lemma~\ref{paper:averaging-witness} to the actual object $XX^*$
and this inertia group, using \eqref{paper:galois-fusion}.
At two use also its self-dual adjoint clause.

The field assertions follow from inertia fixedness and duality.
The square action of $1+\ell$ at odd $\ell$, and of $5$ at
two, gives the primary twist bounds. In the perfect fusion-simple case,
$N_{XX^*}^{Y}=N_{YX}^{X}\le d_Y$; equality would give
$N_Y\mathbf e_X=d_Y\mathbf e_X$, contrary to strict Perron
expansion. At two both integers are odd. The degree range
follows from square divisibility and prime-power-degree exclusion.
\end{proof}

\subsubsection{Height congruences and the defect-two layer}

For this subsection take $J=\{p\}$, put $a=v_p(D)$ and
$m=\lfloor a/2\rfloor$, and let $B$ be a full $p$-block.
The joined-block product theorem already includes every
one-prime product formula.  Its second character identity
also controls individual fusion coefficients.

\begin{proposition}[All-height coefficient congruences]
\label{paper:height}
Put $h=\min_{X\in B}v_p(d_X)$ and $U_B=p^{-h}\Omega_B$.
For $A\in R$, the rational number
\[
 q_{B,A}(T)=\frac{p^h}{d_T}[T](AU_B)
\]
is integral at $p$ and has constant residue in $\mathbb F_p$
on each full target block.  Consequently, for
$k=v_p(d_T)\ge h$,
\[
 [T](AU_B)\in p^{k-h}\mathbb Z;
\]
the exponent improves to $k-h+1$ when the target block
contains a label of height less than $h$.
\end{proposition}
\begin{proof}
The master identity gives
\[
 q_{B,A}(T)=\sum_Z\chi_Z(A)\,
                       \overline{b_{B,Z}\chi_T(Z)}.
\]
Its factors are locally integral, and blockwise character
congruence makes the residue constant.  The rational left
side has residue in $\mathbb F_p$.  A target label of height
below $h$ makes this common residue zero.  Multiplication
by $d_T/p^h$ proves both coefficient bounds.
\end{proof}

\begin{corollary}[The defect-two layer modulo a prime]
\label{paper:layer}\label{paper:lower-support}
Suppose $a=2m$, $m\ge2$, and let $\mathcal H$ be the full
blocks homogeneous of height $m-1$.  Write
$d_X=p^{m-1}u_X$ on these blocks and put
\[
 \mathcal J=\sum_{B\in\mathcal H}RU_B,\qquad
 \mathfrak i=(\mathcal J+pR)/pR.
\]
Then $\mathfrak i$ is a square-zero ideal, totally isotropic
for the dot product in the simple basis.  Its height-$m$
coordinates vanish; its height-$(m-1)$ part is supported on
$\mathcal H$ and is constant times $u$ on each block.  Thus
\begin{equation}\label{paper:linear-splitting}
 \mathfrak i=
 \bigoplus_{B\in\mathcal H}\mathbb F_p\overline{U_B}
       \oplus\mathfrak i_<,
 \qquad
 \mathfrak i_<=\mathfrak i\cap
   \operatorname{span}_{\mathbb F_p}\{[X]:v_p(d_X)\le m-2\}.
\end{equation}
If $\mathcal H\ne\varnothing$, $\mathfrak i_<$ contains a
vector with nonzero unit coordinate.  At two, $\mathfrak i$
is a doubly even binary code, so at least four labels have
height at most $m-2$.  In odd dimension at least three do;
at least five do if $p\equiv5,7\pmod8$.
\end{corollary}
\begin{proof}
The coefficient congruences give the stated coordinates:
a height-$m$ label is a defect-zero singleton, and a block
meeting height $m-1$ is either homogeneous there or contains
a lower label.  Subtract the resulting multiples of
$\overline{U_B}$ to obtain the direct sum; their disjoint
supports give independence.

The product theorem gives $\mathcal J^2,\mathcal J\mathcal J^*\subseteq p^2R$, proving
square-zero multiplication and isotropy.  At two,
$\|z\|^2\equiv0\pmod4$ for $z\in\mathcal J$, so its odd-coordinate
support has size divisible by four.  For $X\in B\in\mathcal H$,
$[X]^*U_B$ has unit coefficient $u_X$, a unit modulo $p$;
removing its homogeneous-layer part preserves that coordinate.
This proves the assertion about $\mathfrak i_<$ and the binary
bound.  In odd dimension, average the lower vector with its
dual and normalize its unit coordinate to one.  Nonunit
duals are distinct, so isotropy reads
\[
                         1+2\sum_j a_j^2=0.
\]
At least one dual pair is necessary; one pair is impossible
when $-2$ is a nonsquare, namely for $p\equiv5,7\pmod8$.
\end{proof}
The splitting is one of vector spaces; no fusion invariance
of $\mathfrak i_<$ is asserted.  In the exact family its
lower heights are $0,1$ at a sextic prime and $0$ at the
quartic prime.

\subsubsection{Phase masses and conductor layers}

The following calculation retains the useful family-level content
of the conductor filtration without assuming positivity of a
local invariant projection.  Only norms of orthogonal projections
are used.

\begin{proposition}[Masses in an odd local fiber]\label{paper:lw:filtration}
Assume $D=n^2$ is square.  Fix an odd prime $\ell$ and a nonempty coordinate set $E$ obtained
by specifying a full $\ell$-block and complementary twist
components.  Put
\[
 m(\alpha)=\sum_{X\in E:\theta_{X,\ell}=\alpha}d_X^2,
 \quad
 J_E=\{j\ge2:m(\alpha)>0\text{ for some }\alpha
                       \text{ of order }\ell^j\}.
\]
All nonidentity $\ell$th roots have the same mass $m_*$.  One has
\[
 m_*\le(\ell+1)m(1).
\]
Let $b$ be the number of orbits on the inactive labels under the
diagonal tori preserving $E$.  Then
\[
 b\ge |J_E|+\mathbf1_{m_*>0}
              +\mathbf1_{m_*<(\ell+1)m(1)}\ge |J_E|+1.
\]
For the number of inactive labels, the contribution $|J_E|$ can
be replaced by
$\sum_{j\in J_E}\varphi(\ell^{\lfloor j/2\rfloor})/2$.
\end{proposition}
\begin{proof}
Let $v=nP_EW_{\ell'}\mathbf e_0$ and $W=W_\ell$.  This vector is fixed
by the local upper unipotents and all tori preserving $E$, and
$Wv=(d_X)_{X\in E}$.  Work at local level $\ell^a$.
For $1\le j\le a$, let $P_j$ average the kernel of reduction
modulo $\ell^j$.  Its lower--diagonal--upper factorization shows
that $P_jv$ is the average of the lower unipotents with parameters
divisible by $\ell^j$.  Conjugating by $W$ shows that $WP_jv$ is
the projection of $Wv$ onto twists of order dividing $\ell^j$.
Thus $(P_j-P_{j-1})v$, $j\ge2$, are nonzero exactly for
$j\in J_E$.  They are orthogonal, upper-fixed and fixed by all
tori preserving $E$.

On the quotient at level $\ell$, write $u=P_1v=u_0+u_1$, where
$u_0$ is group-fixed and $u_1$ is its orthogonal complement.
Both are upper-Borel-fixed.  If $P_{\mathcal B}$ averages that
Borel, the two Bruhat double cosets have relative sizes $1:\ell$,
so, on Borel-fixed vectors,
\[
 P_{\mathcal B}WP_{\mathcal B}
      =P_G-\ell^{-1}(I-P_G).
\]
The lower unipotents of nonzero parameter belong to the other
double coset.  Their matrix coefficients are consequently equal;
Fourier inversion gives
\[
 m(1)=\|u_0\|^2+\ell^{-2}\|u_1\|^2,
 \qquad
 m_*=(\ell+1)\ell^{-2}\|u_1\|^2.
\]
The nonzero vectors $u_0,u_1$ and the higher layers give the stated
orbit bound.  There is an inactive label because $v\ne0$ is
upper-fixed; hence $m(1)>0$.

For the stronger label count on a layer $j\ge2$, put
$l=\lceil j/2\rceil$ and $d=\lfloor j/2\rfloor$.
The kernel modulo $\ell^l$ is the abelian group
$I+\ell^l\mathfrak{sl}_2(\mathbb Z/\ell^d)$.
Upper and torus invariance restrict the characters occurring in
the layer vector to $A\mapsto\exp(2\pi i cA_{21}/\ell^d)$.
Some $c$ is a unit, since otherwise the lower unipotent of
parameter $\ell^{j-1}$ would fix this vector, contrary to its
Weyl image having exact twist order $\ell^j$.
The square-unit multiples of $c$ all occur by torus invariance.
Their mutually orthogonal components remain upper-fixed, and
number $\varphi(\ell^d)/2$.  Different layers are orthogonal,
proving the stronger bound.
\end{proof}

\begin{corollary}[Exact mass gap in a homogeneous block]
\label{paper:lw:exact-gap}
For a homogeneous defect-two block $B$ at an odd prime $\ell$,
use the scaled degrees $u_X$ and packet coefficients $c_X$ from
Theorem~\ref{paper:lw:homogeneous}.  Let $m_0$ be the sum of $u_X^2$
over inactive labels, and $m_*$ the common sum over a specified
nonidentity $\ell$th phase.  Then the only possible primary twist
orders are $1,\ell,\ell^2$, and
\begin{equation}\label{paper:lw:exact-gap-eq}
 m_0-m_*=\ell S_B,\qquad v_\ell(m_0-m_*)=1.
\end{equation}
If all inactive labels of $B$ form one orbit under the tori
preserving $B$, then only orders $1,\ell$ occur,
\[
 m_*=(\ell+1)m_0,\qquad
 \sum_{X\in B}u_X^2=\ell^2m_0,\qquad \ell\nmid m_0,
\]
and $c_X=-1$ on the inactive labels.
\end{corollary}
\begin{proof}
The twist-order bound follows from Lemma~\ref{paper:family:lw:weyl-line}.
Use the scaled packet $v=\ell(u_Xc_X)_X$, so $W_\ell v=u$.
It is upper-Borel-fixed.  For a unit $t$ in the local ring,
$W_\ell^{-1}U(t)W_\ell$ lies in the double coset $\mathcal BW_\ell\mathcal B$ of the upper Borel,
so its coefficient at $v$ equals $\langle v,W_\ell v\rangle$.
Averaging $t$ over the units averages a nonidentity $\ell$th root
to $-1/(\ell-1)$ and kills roots of order at least $\ell^2$.
It follows that
$m_0-m_*=\langle v,W_\ell v\rangle=\ell S_B$.
The unit assertion in \eqref{paper:lw:packet} gives the exact valuation.

A homogeneous block is active: otherwise its local upper operator
would be scalar on the whole invariant block, hence so would its
Weyl, contradicting \eqref{paper:lw:rank-one} and its exact valuation.
When there is only one inactive orbit, Proposition~\ref{paper:lw:filtration}
therefore gives no higher layer and $m_*=(\ell+1)m_0$.
Equation \eqref{paper:lw:exact-gap-eq} gives $S_B=-m_0$, which is an
$\ell$-unit.  The packet coefficients are constant on that orbit,
so $S_B=c_Xm_0$ gives $c_X=-1$.
\end{proof}

\subsubsection{Packet norms and complementary principal support}

\begin{proposition}[Packet norm and its minimal foreign factor]
\label{paper:packet-norm}\label{paper:packet-location}
Assume $D=n^2$, write $n=\ell^m n'$ with $m\ge2$ and
$\ell\nmid n'$, and let $B$ be a full $\ell$-block homogeneous
of height $m-1$.  Use $u_X,c_X,S_B$ from the homogeneous Weyl
theorem, put $J=\{\text{dimension primes other than }\ell\}$,
and write $I_B=B\cap P_J$.  Since $B$ is nonprincipal,
connectedness of the full-prime join ensures $J\ne\varnothing$.
Put $x_B=\Omega_B/\ell^m$ and $K_B=D_B/\ell^{2m}$.
Define the integral virtual element
\[
                 C_B=\sum_{X\in B}u_Xc_X[X].
\]
Then
\begin{equation}\label{paper:packet-location-eq}
 \operatorname{supp}(C_B)
     \subseteq I_B\cap\{X:\theta_{X,\ell}=1\}.
\end{equation}
\begin{equation}\label{paper:packet-inner}
 \begin{aligned}
 W_\ell x_B&=C_B,
 &\langle x_B,C_B\rangle&=S_B/\ell\notin\mathbb Z_{(\ell)},\\
 \|C_B\|^2&=K_B
   =\sum_{\substack{X\in I_B\\\theta_{X,\ell}=1}}u_X^2c_X^2.
 \end{aligned}
\end{equation}
For any collection of full $\ell$-blocks containing $B$, let
$\Lambda$ be the lattice of Proposition~\ref{paper:gluing}
constructed for the single-prime set $\{\ell\}$.  Then
$\Lambda$ contains $x_B$ but not $C_B$ and is not preserved by
$W_\ell$.

In a perfect fusion-simple category let $s$ be the least
dimension prime other than $\ell$.  Then
\[
                              K_B\ge s^2.
\]
For the ordered target family this gives $K_B\ge q^2$ at
its least sextic prime, and $K_B\ge p^2$ at either other prime.
If all labels with $c_X\ne0$ have $u_X$ divisible by an integer
$g$, then $g^2\mid K_B$; this is a conditional statement about
support degrees, not a consequence of twist activity alone.
If $D$ is odd and $B=B^*$, then $K_B$ is even and $K_B\ge2s^2$.
If $B\ne B^*$, the two blocks have equal $K_B$ and their total
mass parameter is at least $2s^2$.
\end{proposition}
\begin{proof}
Write $u=(u_X)_{X\in B}$ and $v=\ell(u_Xc_X)_{X\in B}$.
The homogeneous theorem gives $W_\ell v=u$.  The square $W_\ell^2$
is a local torus.  It fixes $v$, because the complementary vacuum
column and the block are torus invariant in square dimension.
Thus $W_\ell u=v$.  Since $x_B=u/\ell$, this proves the first identity.
The defining vacuum vector also gives
$C_B=n'P_BW_{\ell'}\mathbf e_0$.  The complementary Weyl
vacuum column is supported on $P_J$ by joined block invariance;
any anomaly factor outside the dimension primes is scalar.
This proves \eqref{paper:packet-location-eq}.  The norm and pairing
are the packet identities, with the norm sum restricted to this
support.

Every member of $\Lambda=R_0+L$ pairs integrally with $L$:
this holds by definition for $R_0$ and by integrality for $L$.
As $x_B\in L$, the nonintegral pairing in
\eqref{paper:packet-inner} excludes $C_B$ from $\Lambda$.

A label in $B$ has degree $\ell^{m-1}u_X$ with $\ell\nmid u_X$.
A nonunit prime-power degree is impossible in the simple perfect
category.  Hence $u_X$ contains another dimension prime and
$u_X\ge s$.  At least one $c_X$ is nonzero, giving the lower
bound, and the divisibility assertion follows term by term.
The complementary vacuum vector is fixed by full charge
conjugation, so duality identifies $C_B^*$ with $C_{B^*}$.
In odd dimension all nonunits have distinct duals.  If $B=B^*$,
the sum of squares therefore consists of equal pairs; otherwise
the same identity equates the two block norms.
\end{proof}

\begin{corollary}[One inactive orbit gives a unit block trace]
\label{paper:one-orbit}
Under the hypotheses of Proposition~\ref{paper:packet-norm},
let $\ell$ be odd and suppose that the inactive labels of $B$
form one orbit under the tori preserving $B$.  Then the actual
object
\[
                V_B=\sum_{\substack{X\in B\\\theta_{X,\ell}=1}}u_X[X]
\]
satisfies
\[
 C_B=-V_B,\qquad \|V_B\|^2=K_B,\qquad
 \operatorname{FPdim}(V_B)=\ell^{m-1}K_B,\qquad \ell\nmid K_B.
\]
Its constituents have one common full twist.  Moreover the block
trace $t_B=D_B/D$ is a $\ell$-unit.  The coefficient-support consequence is given below.
For a homogeneous dyadic block the unit-trace conclusion holds
without the one-orbit hypothesis.
\end{corollary}
\begin{proof}
The exact odd mass-gap theorem gives $c_X=-1$ on the inactive
labels and $K_B=m_0$ with $\ell\nmid m_0$.  This proves the norm and
dimension formulas.  The full twist is common because the labels
are $\ell$-inactive and have the common complementary twist of $B$.
Writing $n=\ell^m n'$ gives $t_B=K_B/(n')^2$, which is a $\ell$-unit.
At two, the packet norm already gives $K_B$ odd.
\end{proof}

\begin{proposition}[Inactive complementary support]
\label{paper:family:unit-block-trace}
Let \(B\) be a nonprincipal full \(p\)-block in the
ordered exact family.  Write
\[
 E_B=\sum_{X\in B}e_X=\sum_Z b_{B,Z}[Z],\qquad
 t_B=\tau(E_B)=\frac{\sum_{X\in B}d_X^2}{D}.
\]
If \(t_B\) is a unit at the chosen place, there is a
nonunit \(Z\in P_{\{q,r\}}\), inactive at \(p\), such that
\(p\nmid d_Z\) and \(b_{B,Z}\) is a local unit.
Its degree lies in
\[
                    \{qr,\ qr^2,\ q^2r,\ q^2r^2,\ q^3r\}.
\]
If all labels of \(B\) are \(q\)-active, the \(q\)-height
of \(Z\) is at most two.  If they are \(r\)-active, its
\(r\)-height is one.  If both conditions hold,
\(d_Z\in\{qr,q^2r\}\).

For a homogeneous defect-two block at \(p=2\), the unit
trace hypothesis holds automatically.
\end{proposition}
\begin{proof}
Theorem~\ref{paper:block:complementary-corners} places the coefficient
support in $P_{\{q,r\}}$.  Block orthogonality against the vacuum gives
\(b_{B,Z}=0\) at every \(p\)-active \(Z\).
Since \(B\) is nonprincipal, \(\FPdim(E_B)=0\), so
\[
                   0=t_B+\sum_{Z\ne0}d_Zb_{B,Z}.
\]
All coefficients are locally integral.  If \(t_B\) is
a unit, some nonunit summand has both factors units.
Square divisibility and the exclusion of prime-power
degrees leave \(d_Z=q^a r^b\), with
\(1\le a\le3\), \(1\le b\le2\).
The degree \(q^3r^2\) would have codegree prime to both
\(q,r\), giving a singleton complementary join component,
contrary to \(Z\in P_{\{q,r\}}\).
This gives the stated five possibilities.
Full height at \(q\) or \(r\) has unit codegree there
and misses an active block; this proves the refinements.

Finally write \(d_X=2^2u_X\) in a homogeneous dyadic
defect-two block.  The packet identity gives
\(v_2(\sum_Xu_X^2)=2\).
Thus
\[
 v_2(t_B)=v_2\!\left(2^4\sum_Xu_X^2/D\right)=0.
\]
\end{proof}

\paragraph{The distinction between support and residue evaluation.}
The unit obtained here is a coefficient \(b_{B,Z}\) of
the block idempotent.  It need not be a unit value of the
block's residue character on \(Z\).  A trace on a
nonsemisimple residue algebra can be nonzero on its radical.
The proposition therefore identifies an actual support
label; it does not identify that label with a missing
fusion subcategory.


\subsubsection{Integral complementary projections}

Retain the hypotheses and notation of
Proposition~\ref{paper:packet-norm}: thus $D=n^2$, $n=\ell^m n'$,
$m\ge2$, and the full $\ell$-block $B$ is homogeneous of height
$m-1$.  Write $d_X=\ell^{m-1}u_X$, put
$I=B\cap P_{\ell'}$, and set $b=|B|$, $h=|I|$.
Here $P_{\ell'}$ is the principal join at the other dimension
primes.  The complementary-corner theorem gives $h\ge1$.

\begin{theorem}[An integral projection and its intersection bounds]
\label{paper:integral-projection}
Let $W=W_\ell|_B$ and define
\[
 G=\ell^2W^*P_IW
   =\ell^{-(2m-2)}
       \bigl(N_{\ell^{2m}E_{P_{\ell'}}}\bigr)_{B,B}.
\]
Then $G$ is an integral symmetric positive semidefinite matrix,
\[
 G^2=\ell^2G,\qquad \operatorname{rank}G=h,\qquad
 Gu=\ell^2u,\qquad \operatorname{Tr}G=\ell^2h,
\]
and every diagonal entry is positive.  At the residue prime,
\begin{equation}\label{paper:projection-residue}
 G_{XY}\equiv S_B^{-2}u_Xu_Y
                  \sum_{Z\in I}u_Z^2\pmod\ell,
 \qquad S_B=\sum_{X\in B}u_X^2c_X\in\mathbb Z_{(\ell)}^\times.
\end{equation}
Consequently, if $\ell\nmid\sum_{Z\in I}u_Z^2$, then
\begin{equation}\label{paper:unit-intersection-bound}
 b^2(b-1)\le\ell^4h(b-h).
\end{equation}
In particular $b\le\ell^2\sqrt h$ and
$b\le1+\ell^4/4$.
If $\ell\mid\sum_{Z\in I}u_Z^2$, then $G/\ell$ is integral and
\begin{equation}\label{paper:divisible-intersection-bound}
                         b\le\ell h.
\end{equation}
Equality in \eqref{paper:divisible-intersection-bound} partitions
$B$ into $h$ sets of $\ell$ equal-degree labels, on which
$G$ is the direct sum of the matrices $\ell J_\ell$.
Here $J_\ell$ is the all-ones matrix of size $\ell$.
\end{theorem}
\begin{proof}
The signed principal projector
$\ell^{2m}E_{P_{\ell'}}$ is integral by
Proposition~\ref{paper:block:principal-projector}.  Its multiplication
matrix preserves $B$, and the complementary-corner theorem identifies
the indicated restriction with $\ell^{2m}W^*P_IW$.
Every entry of $W$ has $\ell$-valuation $-1$; hence this integral
matrix is divisible by $\ell^{2m-2}$.  This proves integrality of
$G$.  It is real symmetric because the signed projector is
self-dual.  The displayed projection formula gives positivity,
the quadratic identity and rank; the regular-vector identity gives
$Gu=\ell^2u$.  Density of $W$ and $I\ne\varnothing$ make every
diagonal positive.

The homogeneous residue formula gives
$\ell W_{ZX}\equiv S_B^{-1}u_Zu_X$ at every $\ell$-place.
This formula holds at every $\ell$-place with the same residue
in $\mathbb F_\ell$, so it also holds for the conjugate entries.  Summing the
products over $Z\in I$ proves \eqref{paper:projection-residue}.

In the unit case every entry of $G$ is a nonzero integer.
Using its positive diagonal and the trace identity, one obtains
\[
 \ell^4h=\operatorname{Tr}(G^2)
    =\sum_{X,Y}G_{XY}^2
    \ge b(b-1)+\frac{\ell^4h^2}{b}.
\]
This is \eqref{paper:unit-intersection-bound}.  Since $1\le h\le b$,
it implies $b^2\le\ell^4h$; the bound
$h(b-h)\le b^2/4$ gives the second consequence.
In the divisible case, $G/\ell$ is an integral positive
semidefinite matrix with positive diagonal and trace $\ell h$,
which proves \eqref{paper:divisible-intersection-bound}.

At equality every diagonal of $G/\ell$ is one.  Choose real
Gram vectors for this matrix.  They are unit vectors with pairwise inner products in
$\{0,1,-1\}$.  They therefore form mutually orthogonal signed
classes.  The nonzero eigenvalue $\ell$ makes each class have
$\ell$ members.  The positive eigenvector $u$ forces all signs
within a class to agree and the corresponding degrees to be equal.
Thus its blocks are $J_\ell$, as claimed.
\end{proof}

At two these projection identities have a complete elementary
classification.  It concerns the complementary projection, not a
fusion multiplication on $B$.


\begin{corollary}[When the complementary projection is nonnegative]
\label{paper:nonnegative-projection}
In Theorem~\ref{paper:integral-projection}, suppose that $p=\ell$
is the least dimension prime.  The matrix $G$ is entrywise
nonnegative if and only if one of the following occurs:
\begin{enumerate}[label=(\roman*)]
\item $I=B$ and $G=p^2I_b$;
\item $b=p^2$, $h=1$, and $G=J_{p^2}$, with equal degrees on $B$;
\item $b=ph$ and $G=\bigoplus_{i=1}^h pJ_p$, with equal degrees
within each set of $p$ labels.
\end{enumerate}
Thus every other proper principal intersection forces a negative
off-diagonal entry in its complementary projection.
\end{corollary}
\begin{proof}
Suppose $G$ is nonnegative.  Each connected matrix block has
rank one by Perron--Frobenius and $G^2=p^2G$.  Write it
$k vv^{\mathsf T}$ with $v$ a primitive positive integer vector.
Integrality of the entries gives $k\in\mathbb Z_{>0}$, and
$k\sum v_X^2=p^2$.  The positive eigenvector $u$ is proportional
to $v$ on that block, so $v$ is obtained from the corresponding
scaled degrees by dividing their gcd.  Every prime factor of a
$v_X>1$ would exceed $p$, contrary to $v_X^2\le p^2$.
Hence $v=\mathbf1$, and the block is $kJ_s$ with $ks=p^2$.

If $I\ne B$, density of $W$ makes every diagonal entry of
$G/p^2$ strictly less than one.  Thus no $s=1$ block occurs.
The remaining sizes are $p$ and $p^2$.  A $p^2$ block has
nonzero reduction modulo $p$; the rank-one residue formula,
whose scaled degree entries are all units, then makes every
entry of $G$ nonzero.  Such a block is therefore all of $B$.
Otherwise every block has size $p$.  This gives (ii) and (iii);
$I=B$ gives (i).  The converse is immediate.
\end{proof}


\begin{theorem}[The dyadic complementary projection]
\label{paper:dyadic-projection}
Suppose $\ell=2$.  Then precisely one of the following possibilities
occurs, up to reordering the labels of $B$:
\begin{enumerate}[label=(\roman*)]
\item $I=B$ and $G=4I_b$;
\item $b=4$, $h=1$, $G=J_4$, and all four degrees are equal;
\item $b=4$, $h=3$, and
\[
           G=4I_4-\sigma\sigma^{\mathsf T},\qquad
           \sigma\in\{1,-1\}^4,\qquad \sigma^{\mathsf T}u=0;
\]
\item $h=b/2$, $h\equiv2\pmod4$, and $G$ is the direct sum of
$h$ matrices $2J_2$.  The two degrees in each pair are equal.
\end{enumerate}
In (iii), the unique label $T\in B\setminus I$ is fixed by all
local tori and, after absorbing an overall sign in $\sigma$,
\[
                         W_2\mathbf e_T=\frac12\sigma.
\]
\end{theorem}
\begin{proof}
All $u_X$ are odd, so \eqref{paper:projection-residue} says that
every entry of $G$ has parity $h$.  If $h$ is odd, $G-2I_b$ is
an integral matrix all of whose entries are odd, and
\[
                         (G-2I_b)^2=4I_b.
\]
Each row consequently has at most four entries.  The homogeneous
dyadic rank congruence $b\equiv4\pmod8$ gives $b=4$, so $h=1$ or
$3$.  When $h=1$, the integral rank-one positive matrix $G$ has
positive diagonal and trace four.  Its diagonal entries are all one;
the positive eigenvector $u$ gives $G=J_4$ and equal degrees.
When $h=3$, the same argument applied to $4I_4-G$ gives
$4I_4-G=\sigma\sigma^{\mathsf T}$.  The equality $Gu=4u$
gives the stated signed degree relation.

The complement $B\setminus I=\{T\}$ is preserved by local tori,
so they fix $T$.  Moreover
$4W_2^*P_{\{T\}}W_2=\sigma\sigma^{\mathsf T}$.
Thus $W_2^*\mathbf e_T$ is a scalar multiple of $\sigma/2$.
The local torus $W_2^2$ fixes $T$, so $W_2^*\mathbf e_T=W_2\mathbf e_T$.
Symmetry and $W_2^2\mathbf e_T=\mathbf e_T$ give
$\sum_X(W_2)_{XT}^2=1$, so that scalar has square one.
This proves the column assertion.

If $h$ is even, put $H=G/2$.  The matrix $H-I_b$ is integral,
symmetric and orthogonal, hence a signed involution permutation:
each row has one entry equal to $1$ or $-1$.
Positive diagonals of $H$ exclude negative fixed points.
The equation $Hu=2u$, with every $u_X>0$, excludes negative
transpositions and forces equal degrees at each positive
transposition.  A positive fixed point would make $G/4$ fix a
coordinate vector.  If $I\ne B$, this is impossible, since its
image under the dense matrix $W_2$ would have to be supported on
$I$.  Therefore either $I=B$, or every cycle is a positive
transposition.  These alternatives give (i) and (iv), respectively.
Finally $b\equiv4\pmod8$ gives $h\equiv2\pmod4$ in (iv).
\end{proof}

\begin{corollary}[The complementary components meeting a large dyadic block]
\label{paper:dyadic-orbit-concentration}
Let $\ell=2$ and suppose $|B|>4$ and $B\not\subseteq P_{2'}$.  Then
$|B\cap P_{2'}|=|B|/2$, and all nonprincipal complementary
components meeting $B$ belong to one Galois orbit of components.
\end{corollary}
\begin{proof}
The dyadic classification gives the first assertion.  If $C$ is
any Galois-stable union of complementary components containing
$P_{2'}$, its idempotent $E_C$ is rational and integral at the
complementary primes.  Hence $2^{2m}E_C$ is integral, by the same
denominator argument as for the principal projector.
Its restriction fixes the positive regular vector of $B$.
The proof of the dyadic classification therefore applies unchanged
with $B\cap C$ in place of $I$.
Choose a Galois orbit of nonprincipal components meeting $B$ and
adjoin its union to $P_{2'}$.  The resulting intersection has size
strictly greater than $|B|/2$.  Since $|B|>4$, the classification
forces it to be all of $B$.  Thus this one orbit contains every
nonprincipal complementary component meeting $B$.
\end{proof}

\begin{corollary}[Actual products from the dyadic pairs]
\label{paper:paired-fusion}
In case \textup{(iv)} of Theorem~\ref{paper:dyadic-projection}, write the pairs as $\{X_i,X_i'\}$, with common
degree $d_i$.  Then for every $i,j$ one has the identities of
actual Grothendieck objects
\[
 X_iX_j^*=X_i'X_j'^*,\qquad
 X_iX_j'^*=X_i'X_j^*.
\]
Consequently
\[
 H_{ij}:=\frac{(X_i+X_i')(X_j+X_j')^*}{2}
       =X_iX_j^*+X_iX_j'^*
       \in R_{\ge0},
\]
with unit coefficient $\delta_{ij}$ and dimension $2d_id_j$.
For a perfect fusion-simple category and nonunit simple $Y$,
\[
                           [Y]H_{ii}\le d_Y-1.
\]
Equivalently, the rational involution
$2E_{P_{2'}}-1$ interchanges the two labels in every such pair.
This assertion does not identify that involution with an actual
invertible object.
\end{corollary}
\begin{proof}
On $B$, multiplication by $E_{P_{2'}}$ is the direct sum of
$J_2/2$.  Hence every compressed fusion matrix commutes with the
pair permutation, and each $2\times2$ rectangle between pairs has
the form $\left(\begin{smallmatrix}a&b\\b&a\end{smallmatrix}\right)$.
Frobenius reciprocity gives the two product identities.
Proposition~\ref{paper:averaging-products}, with $k=2$, supplies
the remaining integrality, dimension and coefficient assertions.
Finally $2N_{E_{P_{2'}}}-I$ is exactly the pair permutation on $B$.
\end{proof}


\begin{proposition}[Uniform membership and twists of dyadic pairs]
\label{paper:paired-block-uniform}
Let $B$ be a homogeneous dyadic defect-two block in the paired
case of Theorem~\ref{paper:dyadic-projection}. Put $C=P_{2'}$,
$I=B\cap C$, and $r=2E_C-1$. Then exactly one of the following
pair-membership types occurs.
\begin{enumerate}
\item Every canonical pair is mixed, with one label in $C$.
There is one $t\in\{i,-i\}$ such that
$\theta_{X',2}/\theta_{X,2}=t$ for every pair ordered with
$X\in C$. If $h=|B|/2$, there is a symmetric unitary $h\times h$
matrix $V$ such that
\[
 W_2|_B=\frac1{\sqrt2}\begin{pmatrix}V&V\\V&-V\end{pmatrix},
 \qquad T_2|_B=\begin{pmatrix}D&0\\0&tD\end{pmatrix}.
\]
\item No canonical pair is mixed. There is one $t\in\{1,-1\}$
such that the twist ratio is $t$ in every whole principal pair
and $-t$ in every wholly nonprincipal pair.
\end{enumerate}
\end{proposition}
\begin{proof}
Write $W=W_2|_B$, $Q=P_I$, and $P=W^*QW$ for the
complementary principal multiplication projection. Its form
is a direct sum of $J_2/2$. The square $W^2$ is a local torus
permutation preserving $I$, so
\[
                         PW=WQ,\qquad QW=WP.
\]
Therefore $QWQ$ commutes with $K=QPQ$ on principal coordinates.
A mixed pair contributes the eigenvalue $1/2$ to $K$; a whole
principal pair contributes $J_2/2$, with eigenvalues zero and one.
An entry of $W$ between these two types would vanish, contrary
to the dense local Weyl theorem. If any pair is mixed there
is thus no whole principal pair. Since $|I|=|B|/2$, the numbers
of wholly principal and wholly nonprincipal pairs are equal;
neither occurs. This proves uniform membership.

Let $R=2P-I$ interchange canonical pairs and let
$\varepsilon=2Q-I$. The two projection identities give
$RW=W\varepsilon$ and $WR=\varepsilon W$. For every pair
write $X_i'=RX_i$ and
$t_i=\theta_{X_i',2}/\theta_{X_i,2}$.
Since $\varepsilon$ commutes with $T=T_2|_B$,
\[
                  (WTW)_{X_i',X_j'}=(WTW)_{X_i,X_j}.
\]
In the mixed case choose $X_i\in I$; then
$W_{X_i',X_j'}=-W_{X_i,X_j}$. The modular relation
$WTW=\gamma_2^3T^{-1}WT^{-1}$ and density give
$t_it_j=-1$ for all $i,j$. Hence all ratios are one common
$t=\pm i$. Ordering principal members before their partners,
the equations with $R,\varepsilon$ give
$W=\left(\begin{smallmatrix}A&A\\A&-A\end{smallmatrix}\right)$.
Symmetry and unitarity give $A=A^{\mathsf T}$ and $2AA^*=I$;
put $V=\sqrt2A$.

In the unmixed case put $\varepsilon_i=\varepsilon_{X_i}$.
Now $W_{X_i',X_j'}=\varepsilon_i\varepsilon_jW_{X_i,X_j}$.
The same modular comparison gives
$t_it_j=\varepsilon_i\varepsilon_j$, so
$t_i=t\varepsilon_i$ with one $t\in\{1,-1\}$.
\end{proof}

\begin{theorem}[Fourfold energy from a paired inactive fiber]
\label{paper:mixed-inactive-four}
Put $J=B\cap C\cap\{X:\theta_{X,2}=1\}$. In the mixed
alternative, and in the unmixed alternative with principal
twist ratio $-1$, one has
\[
                             \|P_JW_2P_J\|\le\frac12.
\]
For every dyadically inactive $Y\in C$, consequently,
\[
                       \|s_{B,Y}\|^2
                            \ge4\|P_Js_{B,Y}\|^2.
\]
The second inequality also holds for any linear combination
of columns indexed by dyadically inactive labels in $C$.
\end{theorem}
\begin{proof}
First suppose all pairs are mixed. Each label in $B$ has
dyadic formal codegree four. The exact dyadic twist bound
therefore gives $T_2^8=1$ on $B$. The normalizing scalar
$\gamma_2$ has order dividing eight, so the normalized
upper operator $u=T_2/\gamma_2$ also satisfies $u^8=1$.
The individual conductor theorem puts the local character
fields in $\mathbb Q(i)$, so every torus congruent to one
modulo four acts trivially; the local tori are unsigned.

In the local matrix group, for
$w=\left(\begin{smallmatrix}0&-1\\1&0\end{smallmatrix}\right)$,
\[
 [u(2),wu(2)w^{-1}]
   =\begin{pmatrix}13&8\\8&5\end{pmatrix}
   =u(8/5)h_{1/5}\ell(8/5).
\]
The right side acts trivially: both unipotent parameters
are divisible by eight, and $1/5\equiv1\pmod4$.
Scalar normalizations cancel in the commutator. Thus
\[
                        [T_2^2,W_2T_2^2W_2^{-1}]=1.
\]
In the decomposition of Proposition~\ref{paper:paired-block-uniform}, write
$H=2^{-1/2}\left(\begin{smallmatrix}1&1\\1&-1\end{smallmatrix}\right)$
and $Z=\operatorname{diag}(1,-1)$. The semion factors
$Z$ and $HZH^{-1}$ anticommute. It follows that the residual
unitary matrices $D^2$ and $VD^2V^{-1}$ anticommute.

Let $P_+$ project to the eigenvalue one of $D^2$.
Anticommutation makes $D^2$ exchange the eigenvalues one
and minus one of $VD^2V^{-1}$. A vector fixed by $D^2$
has equal squared projections onto these two eigenspaces,
each at most half its squared norm. Hence
$\|P_+VP_+\|\le1/\sqrt2$. The projection $P_1$ onto
$D=1$ is a subprojection of $P_+$, and the nonzero corner
of $P_JW_2P_J$ is $2^{-1/2}P_1VP_1$. This proves the
claimed bound in the mixed case.

Now suppose the pairs are unmixed with principal ratio
$-1$. If $m=|B|/4$, the principal set consists of $m$
whole pairs. The corner $A=QWQ$ has equal rows and columns
within each such pair, and
\[
                         AA^*=QPQ=\bigoplus_{1}^{m}J_2/2.
\]
Thus $A=J_2\otimes V/2$ for a unitary $m\times m$ matrix
$V$, under the pair ordering. At most one member of each
principal pair has twist one. Selecting inactive principal
coordinates compresses $A$ to a submatrix of $V/2$, proving
the same bound.

Finally factor $s=W_2W_{2'}$. The complementary operator
preserves $C$ and commutes with the actual dyadic twist.
For $Y\in C$ with twist one at two,
$v=P_BW_{2'}\mathbf e_Y$ is supported on $J$ and
$s_{B,Y}=W_2v$. Unitarity and the compression bound give
$\|s_{B,Y}\|=\|v\|$ and
$\|P_Js_{B,Y}\|\le\|v\|/2$. The same reasoning applies
to every linear combination of these source columns.
\end{proof}

\begin{corollary}[Orbit energy in the amplified paired types]
\label{paper:mixed-orbit-energy}
Let $X,Y$ be compulsory pure-primary witnesses at $5400^2$,
and suppose $X$ belongs to one of the paired types in the
preceding theorem. Put $z=S_{XY}/L$,
$L=\operatorname{lcm}(d_X,d_Y)$, $K=\mathbb Q(z)$,
$\alpha=[K_X:K]$ and $\beta=[K_Y:K]$. Then
\[
                  \beta(d_Y^2+4\alpha L^2E(z))\le D.
\]
The same factor four applies to the retained energy of
legitimate odd-prime Fourier channels.
\end{corollary}
\begin{proof}
Average the $Y$-column over the subgroup fixing $K$.
The resulting vector has norm squared $D/\beta$,
retains its vacuum coordinate $d_Y$, and retains all
orbit-$X$ entries. The dyadically inactive condition and
membership in $C$ are Galois invariant. Apply the preceding
theorem to the full dyadic blocks meeting the orbit of $X$;
the retained orbit energy is $\alpha L^2E(z)$. Add the
vacuum coordinate for the stated bound.

For Fourier channels the same argument is linear in the
incoming inactive principal column combination. Every
transported homogeneous dyadic block has the same odd
actual twists as its orbit label. Their union is therefore
contained in the odd phase projections used for the channel
upper bounds. This gives one dyadic factor four combined
with the commuting odd local estimates, without multiplying
independent endpoint gains.
\end{proof}

The factor-four estimate leaves the unmixed type with equal twists
inside principal pairs and opposite twists inside nonprincipal pairs.
Its principal inactive corner can have norm one. The local semion
factor is an operator factorization; its virtual class is not known
to be a nonnegative integral object.

\begin{corollary}[An odd-degree witness forced by dyadic averaging]
\label{paper:dyadic-positive-witness}
Work in an ordinary perfect fusion-simple integral modular category
of square dimension. Let $B$ satisfy the homogeneous dyadic
defect-two hypotheses of Theorem~\ref{paper:dyadic-projection},
with complementary principal projection proper and
entrywise nonnegative. Thus it is either the flat four-label form
or the paired form of Theorem~\ref{paper:dyadic-projection}.
Put $E=E_{P_{2'}}$.

In the flat case take $Q=B$ and $k=4$. In the paired case there
is a canonical pair $Q$ fixed pointwise by all dyadic inertia;
take this pair and $k=2$. In both cases the actual object
$H_Q=R_QR_Q^*/k$ has a nonunit self-dual simple constituent $Y$ satisfying
\[
 d_Y\text{ and }[Y]H_Q\text{ are odd},\qquad
 K_Y\text{ is real and unramified at two},\qquad E[Y]\ne[Y].
\]
Its actual dyadic twist has order dividing eight. Its full
two-block is not wholly contained in $P_{2'}$, and it lies
outside every homogeneous block of positive dyadic height.
For each $X\in Q$, the total multiplicity of $Y\otimes X$
returning to $Q$ is the odd integer $[Y]H_Q\le d_Y-2$.
In the paired case $(N_Y)_{Q,Q}$ is invertible modulo two.

In the ordered exact family, this witness has degree
\[
                       d_Y=q^a r^b,
              \qquad 1\le a\le3,\quad 1\le b\le2.
\]
\end{corollary}
\begin{proof}
Dyadic inertia preserves $B$ and acts by dimension-preserving
based-ring automorphisms. The flat averaging class is therefore
stable. In the paired case the invariant rational projector
permutes its canonical pairs. Corollary~\ref{paper:lw:inactive-fixed}
supplies an inertia-fixed label in $B$; its pair is then fixed
pointwise.

Thus $H_Q$ is inertia invariant and satisfies $H_Q^*=H_Q$.
The averaging and self-dual clauses of
Lemma~\ref{paper:averaging-witness} give the odd-degree,
self-dual, inertia-fixed constituent with $E[Y]\ne[Y]$ and
the stated return multiplicity. Its field is real and unramified
at two. The square action of $5$ gives
$\theta_{Y,2}^{24}=1$, hence primary order dividing eight.
If the full two-block of $Y$ lay in $P_{2'}$, the complementary
principal projector would fix its coordinate vector, a
contradiction. Odd degree excludes every homogeneous block of
positive dyadic height.

Strict expansion on the proper averaging class gives
$[Y]H_Q\le d_Y-1$. Oddness improves the gap to two.
In the paired case the compressed matrix is
$\left(\begin{smallmatrix}a&b\\b&a\end{smallmatrix}\right)$
with $a+b=[Y]H_Q$ odd, so its determinant $(a+b)(a-b)$ is odd.
Square divisibility and prime-power-degree exclusion give the
degree range.
\end{proof}

The adjoint witness of Corollary~\ref{paper:homogeneous-adjoint}
has indicator $+1$ without a positivity assumption on the
projection. The present averaging witness additionally escapes
the projector-fixed coordinate set. These need not be the same
constituent: odd multiplicity in $H_Q$ becomes even after
multiplication by $k$, so it supplies no indicator sign by itself.


\begin{proposition}[Averaging classes and the field-index energy budget]
\label{paper:class-energy}
Let $E$ be a rational self-adjoint spectral idempotent, with spectral
support $C$. Suppose multiplication by $E$ averages a Galois-stable
collection of disjoint equal-degree classes of size $k$, restricting
to $J_k/k$ on each class. Let $X$ belong to one of these classes,
let $Y\in C$ be a nonunit, and suppose the union of classes meeting
the Galois orbit $\mathcal O_X$ does not contain the unit label.
For $c\in\mathbb Q_{>0}$ put
\[
 w=S_{XY}/c,\qquad K=\mathbb Q(w),\qquad
 a=[K_X:K],\quad b=[K_Y:K].
\]
Let $e$ be the number of labels of $\mathcal O_X$ in the averaging
class containing $X$. Then $e$ is the same for every class meeting
$\mathcal O_X$, with $1\le e\le k$ and $e\mid a$. One has
\begin{equation}\label{paper:class-energy-bound}
       b\left(d_Y^2+\frac{k a}{e}\,c^2E(w)\right)\le D.
\end{equation}
In particular $k bE(w)<D/c^2$. If each averaging class meets
$\mathcal O_X$ in one label, the full stronger bound is
$kabE(w)<D/c^2$.
\end{proposition}
\begin{proof}
The $Y$th character vector is fixed by $N_E$. On each averaging
class this says that the entries $S_{UY}$ are all equal.
Consequently labels in the same class have the same value of $w$.
Galois covariance and invariance of the class system make its action
on the classes meeting $\mathcal O_X$ transitive, so the intersection
size $e$ is constant, with $1\le e\le k$. The orbit has $a[K:\mathbb Q]$
labels and each value of $w$ occurs $a$ times. Each value fiber is a
disjoint union of class intersections of size $e$, proving $e\mid a$.

Let $\mathcal U$ be the union of these classes and let $H$ be the
subgroup of the cyclotomic Galois group fixing $K$. For every
$U\in\mathcal U$, class constancy identifies $S_{UY}$ with an entry
indexed by a label of $\mathcal O_X$. The field $K$ is abelian
and hence Galois, so all these conjugate entries lie in $K$.
The normalized average of the $H$-orbit of the $Y$th row has squared
norm $D/b$. Its vacuum entry is $d_Y$, and its entries on
$\mathcal U$ are unchanged. Thus
\[
       d_Y^2+\sum_{U\in\mathcal U}|S_{UY}|^2\le D/b.
\]
Replacing the $e$ orbit labels in each class by its full $k$ labels
multiplies their total energy by $k/e$. Galois covariance gives
\[
 \sum_{U\in\mathcal U}|S_{UY}|^2
       =\frac{k}{e}\sum_{Z\in\mathcal O_X}|S_{ZY}|^2
       =\frac{ka}{e}c^2E(w).
\]
This proves the bound. Since $d_Y>0$ and $a/e\ge1$, its strict
consequences follow.
\end{proof}

\begin{theorem}[A singleton inactive principal fiber]
\label{paper:isolated-fiber}\label{paper:inactive-projector}
Assume $D=n^2$, write $n=p^m n'$ with $m\ge2$ and
$p\nmid n'$, and let $B$ be a homogeneous full $p$-block of
height $m-1$. Write $d_U=p^{m-1}u_U$ on $B$. Suppose
\begin{equation}\label{paper:isolated-fiber-hypothesis}
 B\cap P_{p'}\cap\{U:\theta_{U,p}=1\}=\{F\}.
\end{equation}
The product $W_{p'}$ includes every complementary congruence
factor, including the dyadic anomaly when $p$ is odd.
There is a nonzero integer $\kappa$,
prime to $p$, such that
\begin{equation}\label{paper:isolated-column}
 W_p\mathbf e_F=\frac{1}{p\kappa u_F}(u_U)_{U\in B},
 \qquad (W_p)_{FF}=\frac1{p\kappa},
 \qquad\sum_{U\in B}u_U^2=p^2\kappa^2u_F^2.
\end{equation}
The label $F$ is fixed by full local $p$-inertia. The same
integer $\kappa$ occurs on every full Galois translate.

Put
\[
 \mathcal I_p^0=P_{p'}\cap\{A:\theta_{A,p}=1\},\qquad
 E_p^0=\sum_{A\in\mathcal I_p^0}e_A.
\]
This is a rational self-adjoint spectral idempotent. Its
multiplication operator preserves every full $p$-block and
\begin{equation}\label{paper:inactive-projector-formula}
 N_{E_p^0}=W_p^*P_{\mathcal I_p^0}W_p,\qquad
 (N_{E_p^0})_{B,B}
       =\left(\frac{u_Uu_V}{p^2\kappa^2u_F^2}\right)_{U,V\in B}.
\end{equation}
No global entrywise positivity of $N_{E_p^0}$ is asserted.

For any $p$-inactive $Y\in P_{p'}$,
\begin{equation}\label{paper:isolated-factorization}
 s_{FY}=\frac{(W_{p'})_{FY}}{p\kappa},\qquad
 S_{UY}=\frac{u_U}{u_F}S_{FY}\quad(U\in B).
\end{equation}
If $S_{FY}\ne0$, the local cyclic module generated by
$\mathbf e_F$ is a quotient of that generated by
$\mathbf e_Y$. In particular, for every label $Y$ with
$p\nmid f_Y$,
\begin{equation}\label{paper:isolated-top-vanishing}
                              S_{B,Y}=0.
\end{equation}
\end{theorem}
\begin{proof}
The packet of Theorem~\ref{paper:lw:homogeneous} is supported on the
left side of \eqref{paper:isolated-fiber-hypothesis}. It is
therefore $p u_F\kappa\mathbf e_F$ with
$\kappa=c_F\in\mathbb Z\setminus\{0\}$. Its unit sum is
$S_B=u_F^2\kappa$, so $p\nmid\kappa$. The packet equation
and its norm give \eqref{paper:isolated-column}. Local inertia
preserves the full block, the complementary principal component,
and the inactive fiber; their unique common label is fixed.
The complementary vacuum column is fixed by all permutation
tori: its own local tori fix it by rationality, and the other
tori commute with $W_{p'}$ and fix the vacuum. Since dimensions
are Galois invariant, its packet coefficients are unchanged on
transported blocks.

The spectral set $\mathcal I_p^0$ is Galois- and duality-stable.
The full complementary operator preserves both its principal
component and its actual $p$-twist condition. Thus it commutes
with $P_{\mathcal I_p^0}$, and character diagonalization gives
\[
 N_{E_p^0}=s^*P_{\mathcal I_p^0}s
          =W_p^*P_{\mathcal I_p^0}W_p.
\]
Full-block invariance gives its block restrictions. Since
$W_p^2$ fixes $F$, its $F$-column and inverse $F$-column agree;
the displayed packet column gives the outer product in
\eqref{paper:inactive-projector-formula}.

If $Y\in\mathcal I_p^0$, then $W_{p'}\mathbf e_Y$ stays in
$\mathcal I_p^0$. In $s=W_pW_{p'}$, its part on $B$ is
therefore a multiple of $W_p\mathbf e_F$. This proves
\eqref{paper:isolated-factorization}. The map $P_BW_{p'}$
intertwines the entire local $p$-group and sends
$\mathbf e_Y$ to a nonzero multiple of $\mathbf e_F$ when
$S_{FY}\ne0$, proving the cyclic quotient.

If $p\nmid f_Y$, codegree support puts $Y\in P_{p'}$;
the Weyl-line lemma makes it $p$-inactive, and its full $p$-block
is a singleton. A nonzero image under the preceding intertwiner
would make $\mathbb C\mathbf e_F$ locally invariant. This is
impossible: its packet column is nonzero at every label of $B$,
and the norm equation prevents $B$ from being a singleton.
The intertwiner vanishes on $Y$, and so does $S_{B,Y}$.
\end{proof}

\begin{corollary}[Energy and channels from the inactive fiber]
\label{paper:isolated-amplification}
Use Theorem~\ref{paper:isolated-fiber}, let
$Y\in\mathcal I_p^0$ be a nonunit, and suppose
$w=S_{FY}/c$ is integral, where $c>0$ is rational. Put
\[
 K=\mathbb Q(w),\quad a=[K_F:K],\quad b=[K_Y:K],\qquad
 \mathcal B=\frac{D}{p^2\kappa^2c^2}.
\]
Then
\begin{equation}\label{paper:isolated-ordinary-energy}
 b\bigl(d_Y^2+p^2\kappa^2a c^2E(w)\bigr)\le D,
 \qquad abE(w)<\mathcal B.
\end{equation}
If the full input orbits are disjoint, then
\begin{equation}\label{paper:isolated-vacuum-energy}
 E(w)+\frac{b d_F^2}{c^2(a+b)^2}
                        \le\frac{3\mathcal B}{2(a+b)},
\end{equation}
and
\begin{equation}\label{paper:isolated-total-energy}
 E(w)<\min\left\{\frac1{ab},\frac3{2(a+b)}\right\}
                        \mathcal B\le\frac34\mathcal B.
\end{equation}
If $Y$ also satisfies the singleton inactive-fiber hypothesis
and $S_{FY}\ne0$, its packet integer is the same $\kappa$;
one may then retain $a d_Y^2/[c^2(a+b)^2]$ on the left of
\eqref{paper:isolated-vacuum-energy}.

For an odd-prime set $A$ not containing $p$, with both endpoints
inactive at its primes, the notation and hypotheses of
Theorem~\ref{paper:inactive-fourier-budget} give
\begin{equation}\label{paper:isolated-channel}
 p^2\kappa^2ab|\widehat w(\chi)|^2
 \le\frac{D}{c^2\Lambda_A(\chi)}\|Q_A(\chi)u_\chi\|^2
 \le\frac{D}{c^2\Lambda_A(\chi)}.
\end{equation}
The ordinary channel is obtained with $A$ empty. Under the
one-primary-word hypotheses of
Corollary~\ref{paper:primary-inactive-cost}, with phase prime
$\ell\ne p$, this gives
\begin{equation}\label{paper:isolated-primary-cost}
                  p^2\kappa^2ab\,t^2\ell^{2\beta}\le D/c^2.
\end{equation}
The factor $3/4$ improves total energy only.
\end{corollary}
\begin{proof}
Every full Galois translate of $F$ is the unique inactive
principal label in its transported block. Consequently these
blocks are distinct. Write $\mathcal V$ for their union.
For $Y'$ in the full orbit of $Y$, all rows on a transported
block are proportional to its distinguished row, with ratio
$d_U/d_F$. Their squared ratios sum to $p^2\kappa^2$.
Average the $Y$-row over the Galois group fixing $K$. The
averaged row has squared norm $D/b$, retains vacuum entry
$d_Y$, and is unchanged on $\mathcal V$, since the latter
entries are rational multiples of conjugates of $w$.
Their total squared norm is
$p^2\kappa^2a c^2E(w)$, proving
\eqref{paper:isolated-ordinary-energy}.

For the three-quarter bound use the full products
$W_{p'},U_{p'}$, including any dyadic anomaly. They satisfy
the braid relation, and $W_{p'}^2$ is a permutation torus
preserving each full input orbit. If $e=[K:\mathbb Q]$,
the two orbit projections have ranks $ae,be$ and their
crossblock has squared Hilbert--Schmidt norm
$p^2\kappa^2abe c^2E(w)/D$. Taking $Y=0$ in
\eqref{paper:isolated-factorization} gives
$|(W_{p'})_{F'0}|=p|\kappa|d_F/\sqrt D$ on the first
orbit. Its vacuum leakage is therefore at least
$ae p^2\kappa^2d_F^2/D$. Lemma~\ref{paper:braid-energy}
gives \eqref{paper:isolated-vacuum-energy}; its positive
correction and \eqref{paper:isolated-ordinary-energy} give
\eqref{paper:isolated-total-energy}. If both endpoints are
isolated, symmetry and the nonzero crossing equate their
packet integers, and the second vacuum leakage may be retained.

The same row proportionality gives the exact channel norm
\[
 \|P_{\mathcal V}s u_\chi\|^2
        =\frac{p^2\kappa^2ab c^2}{D}|\widehat w(\chi)|^2.
\]
The complementary actual twists are constant on every
homogeneous block. Thus $\mathcal V$ is inactive at any
selected prime where $F$ is inactive. The exact-compression
argument of Theorem~\ref{paper:inactive-fourier-budget}
proves \eqref{paper:isolated-channel}. The primitive Gauss
sum and integral Fourier nonvanishing give
\eqref{paper:isolated-primary-cost}.
\end{proof}

\begin{corollary}[Normalization and uniform local columns]
\label{paper:isolated-normalization}\label{paper:uniform-column}
In the surviving odd ordered exact family, suppose a compulsory
$q$-witness $F$ of degree $p^2q^2$ satisfies
\eqref{paper:isolated-fiber-hypothesis} in a homogeneous
full $p$-block. Then
\begin{equation}\label{paper:uniform-column-conclusion}
 \kappa=\varepsilon\in\{1,-1\},\qquad
 |B|=p^2,\qquad d_U=d_F\ (U\in B),\qquad
 W_p\mathbf e_F=\frac\varepsilon p\mathbf1_B,
 \quad W_p\mathbf1_B=\varepsilon p\mathbf e_F.
\end{equation}
Every opposite compulsory witness has $p$-height two;
inactivity of all those opposite witnesses is not asserted.

More generally, in either parity, if an isolated endpoint of
degree $p^2q^2$ or $p^2r$ has $\kappa=\pm1$, then the uniform
column conclusion holds. On every such block,
\begin{equation}\label{paper:uniform-inactive-projection}
                         (N_{E_p^0})_{B,B}=J_{p^2}/p^2.
\end{equation}
Consequently the common-projector averaging identities,
positive integral products, constant fusion row sums, and
integral lattices of Proposition~\ref{paper:averaging-products}
apply simultaneously to these blocks, using $E_p^0$.
Explicitly, for normalized blocks $B_i,B_j$, their sums
$R_i=\sum_{U\in B_i}[U]$, and representatives $X_i\in B_i$,
$X_j\in B_j$,
\begin{equation}\label{paper:inactive-averaging-products}
 H_{ij}:=\frac{R_iR_j^*}{p^2}
       =R_iX_j^*=X_iR_j^*\in R_{\ge0}.
\end{equation}
In particular $H_{ii}$ contains the actual adjoint object
$X_iX_i^*$ as a nonnegative summand.
An escape conclusion refers to the coordinate lines fixed by
$N_{E_p^0}$; it does not assert escape from the original
complementary principal component.
\end{corollary}
\begin{proof}
The foreign bounds on $(W_p)_{FF}=1/(p\kappa)$, with
$f_F=p^2q^2r^4$, show
$\kappa\in\{1,-1,r,-r\}$. No other prime can divide
$\kappa$: the $q$-valuation is nonnegative, the $r$-valuation
is at least $-1$, and all other foreign valuations are
nonnegative. The packet already makes $\kappa$ a $p$-unit.

Choose a pure-$r$ opposite witness, which is $p$-inactive.
Every opposite crossing is nonzero. The top-height vanishing
in Theorem~\ref{paper:isolated-fiber} excludes height three,
so this witness has degree $p^2r$. Put
$z=S_{FY}/(p^2q^2r)$. If $|\kappa|=r$, its effective
budget is $\mathcal B=q^2$.
In odd dimension a mixed word has
$E(z)>(qr/p^2)\operatorname{cond}\mathbb Q(z)
\ge q^2r^2/p^2>q^2$. A one-primary word has
$E(z)>\frac23t^2\ell^a>q^2$ by
Proposition~\ref{paper:phase-energy}, where
$\{\ell,t\}=\{q,r\}$. Both contradict the amplified budget.
Hence $\kappa=\pm1$ in the stated odd scope. The same
top-height argument excludes every height-three opposite
witness, since such a label is automatically $p$-inactive.

For $d_F=p^2q^2$, the foreign $q$-bound makes every entry
of the $F$-column of $W_p$ integral at $q$. Formula
\eqref{paper:isolated-column} therefore gives $q^2\mid u_U$.
Its norm reads $\sum_U(u_U/q^2)^2=p^2$. These ratios are
positive integers with prime factors only $q,r$, both
greater than $p$. Each ratio must be one. For $d_F=p^2r$
and $\kappa=\pm1$, the identical argument uses the foreign
$r$-bound and $\sum_U(u_U/r)^2=p^2$. This proves the
uniform degree and column assertions. The second column
identity follows because $W_p^2$ fixes $F$.

Substitute the uniform degrees and $\kappa^2=1$ into
\eqref{paper:inactive-projector-formula}. The averaging
consequences now follow from one common spectral idempotent,
exactly as stated in Proposition~\ref{paper:averaging-products}.
Its equal-size identity is
\eqref{paper:inactive-averaging-products}; the summand assertion
follows by expanding $R_iX_i^*$.
This uses no positivity assumption on the original
complementary principal projection.
\end{proof}

\begin{corollary}[The flat specialization]
\label{paper:flat-energy}\label{paper:flat-Fourier}
Suppose $X$ is distinguished in a homogeneous singleton-flat
full $p$-block and $Y\in P_{p'}$ is a nonunit. For integral
$w=S_{XY}/c$, all energy, vacuum, and complementary-prime
channel bounds of Corollary~\ref{paper:isolated-amplification}
hold with $\kappa^2=1$ and $\mathcal B=D/(p^2c^2)$, retaining
their stated orbit and inactivity hypotheses. Its exact
orbit-projection identity also holds with $\kappa^2=1$.
For compulsory degrees $p^2q^j,p^2r$, the effective constant
is $\mathcal B=q^{6-2j}r^2$.
\end{corollary}
\begin{proof}
The singleton-flat packet has $\kappa=\pm1$, and its full
complementary principal intersection is the distinguished
label. Thus the row proportionality in
Theorem~\ref{paper:isolated-fiber} holds for every
$Y\in P_{p'}$. If $Y$ is $p$-active, both this crossing
and the complementary Weyl crossing vanish. The orbit,
vacuum-leakage and compression proof of
Corollary~\ref{paper:isolated-amplification} therefore applies
also in that case. The effective constant follows from the
degrees.
\end{proof}

\subsubsection{Primary levels and saturated flat crossings}

The total-energy improvement lowers the maximal one-primary levels
without input inactivity. Inactive prime phases then give the
exact period and orbit equalities used below.

\begin{proposition}[The maximal one-primary level is absent]
\label{paper:primary-level}
\label{paper:top-conductor-loss-2}
In the surviving odd ordered exact family, take a compulsory crossing
\[
 d_X=p^iq^j,\qquad d_Y=p^hr,\qquad
 i,h\in\{2,3\},\quad\min(i,h)=2,\quad j\in\{1,2\}.
\]
Suppose the $p$-heights are mixed or an endpoint is distinguished
in a singleton-flat full $p$-block. At second $q$-height, any
one-primary actual word has prime conductor. At first $q$-height
a $q$-only word has conductor at most $q^3$.
These two alternatives require no input inactivity or homogeneity
of the lower block. The same conclusions hold if an endpoint lies
in a homogeneous full $p$-block satisfying
\eqref{paper:isolated-fiber-hypothesis} and both inputs are
$p$-inactive.

\end{proposition}
\begin{proof}
For mixed heights the ordinary effective budget is
$\mathcal B=D/(p^3q^jr)^2=q^{6-2j}r^2$.
For a flat endpoint, Corollary~\ref{paper:flat-witness-height}
makes both heights two, and flat amplification gives the same
budget. For the added inactive-fiber alternative,
\eqref{paper:isolated-top-vanishing} again forces both heights
to be two, and Corollary~\ref{paper:isolated-amplification}
gives effective constant $\mathcal B/\kappa^2\le\mathcal B$.
Lemma~\ref{paper:crossing-energy} and the two amplification
corollaries therefore give $E(z)<3\mathcal B/4$ in every case.

Write $\ell$ for the phase prime, $t$ for the other complementary
prime, and $\ell^u$ for the conductor.
Corollaries~\ref{paper:phase-fibers} and
\ref{paper:phase-conductors}, with
Proposition~\ref{paper:phase-energy}, supply invariant
nonnegative integer one-primary coefficients of total mass
less than $(\ell-1)/3$. The basic bound $E(z)>\frac23t^2\ell^u$ already makes
the $q$-conductor at most $q^{6-2j}$, and the $r$-conductor
at most $r^2$ when $j=2$. At that maximal exponent,
Lemma~\ref{paper:last-level-energy} would give
\[
 E(z)\ge r^2q^{5-2j}(q-1)
             =\left(1-\frac1q\right)\mathcal B
             >\frac34\mathcal B,
\]
a contradiction. Thus it is at most $q^{5-2j}$.
At $j=2$, an $r^2$-word similarly costs
$(1-1/r)\mathcal B>3\mathcal B/4$ and is excluded.

\end{proof}

\begin{corollary}[Sharper conductor bounds]
\label{paper:conductor-consequences}
In the odd setting of the preceding proposition, suppose both
primary supports of the actual word occur. Then
\[
 N:=\operatorname{cond}\mathbb Q(z)
                         <\frac34p^2q^{5-2j}r.
\]
In particular, conductor $qr^2$ at second $q$-height, or
$q^3r^2$ at first $q$-height, requires $4r<3p^2$.

Separately, for the divided first-height word of
Theorem~\ref{paper:top-crossing}, mixed support instead gives
$N<3q^3r/4$; its conductor $qr^2$ requires $4r<3q^2$.
If that divided word is one-primary, then
\[
 \begin{aligned}
 N=q^2&\quad\Longrightarrow\quad4p^2(q-1)<3q^3,\\
 N=r^2&\quad\Longrightarrow\quad4p^2(r-1)<3q^2r.
 \end{aligned}
\]
Consequently $4p^2(q-1)\ge3q^3$ forces prime conductor
in the one-primary divided case, without input inactivity.
\end{corollary}
\begin{proof}
For a mixed mass-$p^2$ word,
Proposition~\ref{paper:phase-energy} gives
$E(z)>(qr/p^2)N$. Compare this with $3\mathcal B/4$,
using $\mathcal B=q^{6-2j}r^2$ in the preceding proposition
and $\mathcal B=q^4r^2/p^2$ for the divided first-height word.

In the latter case a $q^2$-word has energy at least
$r^2q(q-1)$ by Lemma~\ref{paper:last-level-energy}; an
$r^2$-word has energy at least $q^2r(r-1)$. Comparison with
$3q^4r^2/(4p^2)$ gives the two necessary inequalities.
Since $r>q$, the displayed sufficient condition contradicts
both. The previous conductor theorem already bounds the
one-primary conductor by $q^2$ or $r^2$.
\end{proof}

\begin{corollary}[Prime-period rigidity from an inactive principal fiber]
\label{paper:flat-period-rigidity}\label{paper:flat-period-rigidity-2}
In the surviving odd ordered family, take a compulsory crossing
\[
 d_X=p^2q^2,\qquad d_Y=p^2r,\qquad
 z=\frac{S_{XY}}{p^2q^2r}.
\]
Suppose one endpoint $F$ has a homogeneous full $p$-block
satisfying \eqref{paper:isolated-fiber-hypothesis}, and both
inputs are $p$-inactive. If both complete input twists are one,
then the actual word is
\begin{equation}\label{paper:flat-period-word}
 z=h+t\sum_{a\in C}\zeta_\ell^a,
 \qquad\{\ell,t\}=\{q,r\},\qquad h=p^2-td\ge2,
\end{equation}
where $C$ is one coset of a subgroup
$H\le(\mathbb Z/\ell)^\times$ of odd order $d=|H|<p$.
Moreover
\begin{equation}\label{paper:flat-period-field}
 K_X=K_Y=K=\mathbb Q(z),\qquad \operatorname{cond}K=\ell,
 \qquad H=\operatorname{Gal}(\mathbb Q(\zeta_\ell)/K).
\end{equation}
The packet integer of $F$ is $\kappa=\pm1$, its block has
$p^2$ labels of degree $d_F$, and its local column is uniform
as in \eqref{paper:uniform-column-conclusion}. The constant $h$
is even, and every constituent of $X\otimes Y$ has twist order
dividing $\ell$.

The same conclusions follow if the actual word is initially
one-primary and both endpoints are inactive at its phase prime,
still retaining the stated $p$-inactivity. Complete twist one
is then a conclusion. A singleton-flat endpoint supplies the
inactive-fiber hypothesis and makes every opposite compulsory
witness $p$-inactive, so the original flat formulation is included.
\end{corollary}
\begin{proof}
Orient the Fourier formulas so that the isolated endpoint is
the target. The actual word has mass $g=p^2$ by
Corollary~\ref{paper:phase-fibers}, with both phase exponents
at most two. Put $K=\mathbb Q(z)$, $a=[K_X:K]$,
$b=[K_Y:K]$. Corollary~\ref{paper:isolated-amplification}
gives effective energy and channel constant
\[
                         B=\frac{q^2r^2}{\kappa^2}\le q^2r^2.
\]
If both complete twists are one, mixed conductor $q^2r^2$
is excluded by \eqref{paper:mixed-energy}, since $qr>g$.
The other mixed conductors $qr,q^2r,qr^2$ are excluded by
Corollary~\ref{paper:mixed-inactive-exclusion}, because
$B\le q^2r^2\le q^4r^2/g$. Thus only one primary support
occurs; strict house excludes the all-identity word.

For either formulation, write the remaining actual word as
\[
 z=h+t\sum_{\xi\ne1}k_\xi\xi,\qquad
 k_\xi\in\mathbb Z_{\ge0},\quad k=\sum k_\xi\ge1,
 \quad h=g-tk>0,\quad\{\ell,t\}=\{q,r\}.
\]
Let its conductor be $\ell^\beta$. Coefficient uniqueness
makes the vector invariant under
$H=\operatorname{Gal}(\mathbb Q(\zeta_{\ell^\beta})/K)$.
The amplified primary cost is
\[
               \kappa^2ab\,t^2\ell^{2\beta}
                          \le q^2r^2=t^2\ell^2.
\]
Hence $\beta=1$, $a=b=1$, and $\kappa=\pm1$.
Corollary~\ref{paper:isolated-normalization} now gives the
uniform local column at either possible endpoint degree.
If only phase-prime inactivity was assumed, equality of the
complete character fields makes full inertia at the other
dimension primes fix both labels. The local inversion parity
theorem then makes their twists trivial there as well.
Thus both complete twists are one in either formulation.

Put $d=|H|$, $e=(\ell-1)/d$, and let $m_C$ be the common
coefficient on a primitive-root $H$-coset $C$. For every
nontrivial $\chi\in\widehat\Gamma$, where
$\Gamma=(\mathbb Z/\ell)^\times/H$, the primitive Gauss
formula and \eqref{paper:isolated-channel} give
\begin{equation}\label{paper:flat-period-saturation}
 |\widehat m(\chi)|^2
              \le\|P_{K_{\ell,1}}u_\chi\|^2\le1.
\end{equation}
The nonnegative integer vector $m$ has mass
$1\le\sum_Cm_C=k/d<e/3$ by sparsity.
Lemma~\ref{paper:integral-fourier-rigidity} therefore makes
$m$ a singleton indicator, giving $k=d$ and the period word.

The field $K_X=K$ is nonreal, since in odd dimension a
nonunit has a distinct dual. Thus $-1\notin H$, and $d$
is odd. Since $td<g$ and $t>p$, one has $d<p$.
The positive integer $h=g-td$ is even. The actual phases
and both input twists now give the assertion about
constituent twists.
\end{proof}


Work in the surviving odd ordered exact family, where $qr>2p^2$,
with a compulsory crossing $d_X=p^2q^2$, $d_Y=p^2r$.
One endpoint $F$, with opposite endpoint $G$, has a homogeneous
full $p$-block satisfying
\eqref{paper:isolated-fiber-hypothesis} and the uniform column
\eqref{paper:uniform-column-conclusion}. Both original inputs
are $p$-inactive. Suppose that, for $\{\ell,t\}=\{q,r\}$,
both are $\ell$-inactive and
\[
 z=\frac{S_{XY}}{p^2q^2r}
   =h+t\sum_{a\in a_0H}\zeta_\ell^a,
 \qquad K_X=K_Y=K=\mathbb Q(z),\qquad\operatorname{cond}K=\ell,
\]
where $H=\operatorname{Gal}(\mathbb Q(\zeta_\ell)/K)$ and
$e=[K:\mathbb Q]>1$. The odd twist-one case is supplied by
Corollary~\ref{paper:flat-period-rigidity}. Under the original
singleton-flat hypothesis, nonzero crossing already supplies
the other input's $p$-inactivity.

Let $\Gamma=\operatorname{Gal}(K/\mathbb Q)$. For
$\sigma\in\Gamma$, let $B_\sigma$ be the full $p$-block of
$\sigma F$. They are distinct, and put
$\mathcal A=\bigcup_\sigma B_\sigma$. Define
\[
 f_\chi=e^{-1/2}\sum_\sigma\chi(\sigma)\mathbf e_{\sigma F},
 \qquad
 u_\chi=e^{-1/2}\sum_\sigma\chi(\sigma)\mathbf e_{\sigma G},
\]
\[
 v_\chi=(p\sqrt e)^{-1}\sum_\sigma\chi(\sigma)
                                      \sum_{T\in B_\sigma}\mathbf e_T.
\]
Let $\Pi_\ell$ project onto the actual $\ell$-inactive labels.
Every varying compulsory witness paired with either fixed
height-two endpoint has a nonrational normalized crossing:
a rational value would be an integer $z_0$ with
\begin{equation}\label{paper:nonrational-crossing}
                  z_0\equiv p^2\pmod{qr},\qquad |z_0|<p^2,
\end{equation}
which is impossible because $qr>2p^2$ in the surviving region.

\begin{theorem}[Equality in the amplified Fourier bound]
\label{paper:saturated-orbits}\label{paper:isolated-saturation}
For every nontrivial $\chi$,
\begin{equation}\label{paper:unified-saturation}
 \Pi_\ell s u_\chi=\ell^{-1/2}\lambda_\chi v_{\chi^{-1}},
 \qquad
 \Pi_\ell s^*v_{\chi^{-1}}
       =\ell^{-1/2}\overline{\lambda_\chi}u_\chi,
\end{equation}
where
\[
 \lambda_\chi=\frac{\chi(a_0)^{-1}\tau_\ell(\chi)}{\sqrt\ell},
 \qquad
 \tau_\ell(\chi)=\sum_{a\in(\mathbb Z/\ell)^\times}
                       \chi(a)\zeta_\ell^a,
 \qquad |\lambda_\chi|=1.
\]
Both complete input twists are one, and both full input orbits
lie in the principal full $\ell$-block. The zero-sum spaces
on those orbits and the corresponding block averages are
fixed by the first $\ell$-congruence kernel and by the full
local $t$-group. The same holds on their local cyclic spans.
Each input orbit lies in one full $t$-block. Moreover:
\begin{enumerate}
\item Both $\chi_X(T)$ and $\chi_Y(T)$ are rational integers
whenever $T$ is $t$-active or its actual $\ell$-twist has order
greater than $\ell$.
\item If $T$ is $\ell$-inactive and outside $\mathcal A$, then
$\chi_G(T)\in\mathbb Z$.
\item If $T$ is both $p$- and $\ell$-inactive, belongs to
$P_{p'}$, and lies outside $\operatorname{Orb}(G)$, then
$\chi_F(T)\in\mathbb Z$.
\end{enumerate}
If $F$ is distinguished in a singleton-flat block, the
$p$-inactivity condition in \textup{(iii)} may be omitted.
\end{theorem}
\begin{proof}
Uniform row proportionality against the $p$-inactive source
gives the exact normalization
\[
 P_{\mathcal A}s u_\chi
        =\frac{\widehat z(\chi)}{qr}v_{\chi^{-1}}
        =\frac{\lambda_\chi}{\sqrt\ell}v_{\chi^{-1}}.
\]
Every label in $\mathcal A$ has the complementary twists
of its distinguished label, so is $\ell$-inactive. Exact
compression bounds $\|\Pi_\ell s u_\chi\|$ by
$\ell^{-1/2}$. The displayed projection attains this norm;
it is the whole inactive projection, and the first
$\ell$-congruence kernel fixes $u_\chi$.
Apply inverse exact compression to $v_{\chi^{-1}}$. Its
inner product with the unit vector $u_\chi$ already has
modulus $\ell^{-1/2}$. Equality in Cauchy--Schwarz gives
the reverse identity and first-kernel fixedness of
$v_{\chi^{-1}}$.

The uniform-column sign is unchanged on the transported
blocks, and $W_pf_\chi=\varepsilon v_\chi$. Commutation
with the $\ell$-local group transfers kernel fixedness to
$f_\chi$. Normality extends fixedness to local cyclic spans.

Put $\alpha=\theta_{G,t}$ and $\beta=\theta_{F,t}$.
The full fields have conductor $\ell$, so pure-$\ell$
Galois representatives show that these components are
constant on each input orbit, and the block averages preserve
$\beta$. Factor $s=W_tW_{t'}$ with the full complementary
product, and put $A=\Pi_\ell W_{t'}$. It commutes with
the local $t$-group. Thus $Au_\chi$ is an upper-unipotent
eigenvector of eigenvalue $\alpha$, whose Weyl image is
a nonzero multiple of $v_{\chi^{-1}}$ with upper eigenvalue
$\beta$. It is simultaneously an eigenvector for upper and
lower unipotents; these generate the local group, so its
line affords a one-dimensional character. Its Weyl image
lies on the same line, forcing $\alpha=\beta$.
The reverse equality transfers this group character to
$u_\chi$, and commutation with $W_p$ transfers it to
$f_\chi$. Their $s$-images have the same character.

Consequently both endpoint character values are rational
integers at any label whose $t$-twist differs from $\alpha$.
Any such compulsory witness would contradict
\eqref{paper:nonrational-crossing}. All compulsory witnesses
therefore have $t$-twist $\alpha$. A pure-$\ell$ witness
exists and has $t$-twist one. Hence $\alpha=1$, and the
one-dimensional group character is trivial by upper and
lower generation. This also applies at $t=3$, without a
perfectness assumption on the local group. The original
$p$-inactivity, assumed $\ell$-inactivity, and proved
$t$-inactivity give both complete twists one.

If two labels of an input orbit belonged to different full
$t$-blocks, project
$W_t(\mathbf e_Z-\mathbf e_{Z'})=\mathbf e_Z-\mathbf e_{Z'}$
onto one of those blocks. This gives $W_t\mathbf e_Z=
\mathbf e_Z$, contrary to $t\mid f_Z$ and the Weyl-line
lemma. Each orbit therefore lies in one such block.
Nontrivial characters span the zero-sum spaces. Their
$s$-images are fixed by the full $t$-group and by the first
$\ell$-kernel, so have zero coordinates at $t$-active
labels and at twists of $\ell$-order greater than $\ell$.
Fourier inversion, field equality and character integrality
give \textup{(i)}. The forward saturated equality gives
\textup{(ii)} in the same way.

For \textup{(iii)}, use the common inactive projector from
Theorem~\ref{paper:isolated-fiber}. Its uniform restriction gives
\[
                    v_\chi=pN_{E_p^0}f_\chi,
 \qquad s^*N_{E_p^0}=P_{\mathcal I_p^0}s^*.
\]
The second identity uses duality stability of the spectral
set. The reverse saturated equality therefore gives the
required Fourier vanishing on its inactive principal fiber,
outside the other orbit. If the full principal intersection
is already a singleton, replace $E_p^0$ by $E_{P_{p'}}$;
the original argument removes the $p$-inactivity condition.

Finally cancel $W_p$ using its uniform column, and cancel
$W_t$ using the proved fixedness. The equalities give
nonzero local $W_\ell$ transport between the inactive
channels of the two endpoint orbits. Each orbit is contained
in one full $\ell$-block because it is represented by
pure-$\ell$ inertia. Local block preservation makes those
two blocks equal. One endpoint is principal at $\ell$,
so both full orbits lie in $P_\ell$.
\end{proof}

\begin{corollary}[Compulsory witnesses in the saturated configuration]
\label{paper:witness-pattern}\label{paper:inactive-witness-pattern}
Let $\mathcal W_q$ be the nonunit labels in $P_q$ of degrees
$p^i q^j$, $i\in\{2,3\}$, $j\in\{1,2\}$, and let
$\mathcal W_r$ be those in $P_r$ of degrees $p^h r$,
$h\in\{2,3\}$. Every such witness is $t$-inactive and
has $\ell$-primary twist order dividing $\ell$.
Write $\mathcal W_q^0,\mathcal W_r^0$ for their $p$-inactive
parts. Then
\[
 \{A\in\mathcal W_q^0:\theta_{A,\ell}=1\}
           =\operatorname{Orb}(X),\qquad
 \{B\in\mathcal W_r^0:\theta_{B,\ell}=1\}
           =\operatorname{Orb}(Y).
\]
These orbits consist of complete twist-one labels. Every
$p$-inactive witness outside them has complete twist order
$\ell$. In particular no $p$- and $\ell$-inactive $q$-witness
has $q$-height one. If $\ell=q$, every pure-$r$ witness lies
in the single $r$-orbit; if $\ell=r$, every pure-$q$ witness
lies in the single $q$-orbit and has $q$-height two.

Under the original singleton-flat hypothesis, the superscript
$0$ may be omitted in both displayed equalities. Every witness
opposite a singleton-flat endpoint is $p$-inactive, so its
complete twist order is $1$ or $\ell$ as above. If both
endpoints are singleton-flat, this holds on both sides.
\end{corollary}
\begin{proof}
Pair every varying witness with the fixed opposite endpoint
and use \eqref{paper:nonrational-crossing}. Part \textup{(i)}
of Theorem~\ref{paper:saturated-orbits} excludes $t$-activity
and $\ell$-twist order greater than $\ell$ for all witnesses.
For a $p$-inactive witness opposite $F$, part \textup{(iii)}
forces membership in $\operatorname{Orb}(G)$. On $F$'s
side, part \textup{(ii)} forces membership in $\mathcal A$;
the witness is in $\mathcal I_p^0$, and
$\mathcal A\cap\mathcal I_p^0=\operatorname{Orb}(F)$.
The original full orbits consist of witnesses with both
inactivities, giving equality. Their Galois representatives
may be chosen purely at $\ell$, so their complete twist-one
property persists, including at two. The other claims for
$p$-inactive and pure-primary witnesses follow immediately.

If $F$ is singleton-flat, its full principal intersection
is just its distinguished label, and the stronger version
of part \textup{(iii)} applies. The same argument gives
the two equalities without restricting to $p$-inactive
witnesses. Corollary~\ref{paper:flat-witness-height} supplies
the stated $p$-inactivity opposite that endpoint. If the
other endpoint is the singleton-flat one, interchange the
two roles in this argument.
\end{proof}

The equalities alone control the zero-sum orbit spaces.
The global witness restriction now supplies the additional
integrality needed to control one orbit average.

\begin{proposition}[Integral columns at the quartic phase prime]
\label{paper:quartic-phase}
Under the hypotheses of Theorem~\ref{paper:saturated-orbits},
suppose $\ell=r$ and $t=q$. Then
\[
                         qW_q\mathbf e_X\in\mathbb Z^I.
\]
In odd dimension, put $d=|H|$, $e=(r-1)/d$ and
$\mathcal O=\operatorname{Orb}(X)$. The local $q$-cyclic module
generated by this full orbit, including its average, factors
through $\operatorname{SL}_2(\mathbb F_q)$, and the column
identities of Lemma~\ref{paper:integral-orbits} hold.
Its vacuum coefficient $m_0=q(W_q)_{0X}$ is an even integer
with $0<|m_0|<p$, and
\begin{equation}\label{paper:vacuum-congruence}
 m_0\equiv e k p^2q^2\pmod r,
 \qquad d m_0+k p^2q^2\equiv0\pmod r.
\end{equation}
Consequently
\[
 e k^2+8\le k(q+1),\qquad
 e\le\frac{q-3}{2},\qquad
 2(r-1)\le d(q-3).
\]
Here $qd=p^2-h\le p^2-2$, $d$ is odd, $d\ge7$, $p\ge11$, and
\begin{equation}\label{paper:quartic-range}
             2q(r-1)\le(p^2-2)(q-3),\qquad r<\frac{p^2}{2}.
\end{equation}
Every odd isolated-fiber second-$q$-height crossing with two
twist-one inputs has $r<p^2$, whichever phase prime occurs; the
stronger half-square bound here concerns phase prime $r$.

\end{proposition}
\begin{proof}
Every nonunit $r$-free label in $P_q$ is a compulsory $q$-witness:
the principal height bound gives $q$-height at most two;
prime-power-degree exclusion forces both $p,q$ to occur;
and the complementary-degree bound makes its $p$-height at least
two. Corollary~\ref{paper:witness-pattern} makes it $q$-inactive.
Thus every $q$-active label in $P_q$ has $r\mid d_A$, hence
$v_r(f_A)\le2$. The block form of Theorem~\ref{paper:lw:columns}
makes $W_q^{\pm1}|_{P_q}$ integral at $r$.

For $X$, the codegree is $f_X=p^2q^2r^4$. The foreign-column
bound already gives integrality at $p$. The indicator formula
gives integrality away from $p,q,r$ and $q$-valuation at least
$-1$. The column is rational because $K_X$ is unramified at $q$.
Its $q$-multiple is consequently an integer vector.
Theorem~\ref{paper:saturated-orbits} fixes all orbit differences;
Lemma~\ref{paper:integral-orbits} now applies.

Suppose the dimension is odd. The column preserves complementary
twists, so it is supported on the complete complementary
twist-one fiber of $P_q$. In particular it is $p$-inactive.
Among its $r$-free nonunits, the restricted witness equalities
of Corollary~\ref{paper:witness-pattern} therefore leave
precisely $\mathcal O$. Thus every supported
label outside $\mathcal O\cup\{0\}$ has degree divisible by $r$.
Reducing \eqref{paper:linear-phase} modulo $r$ gives
$m_0\equiv e k d_X\pmod r$, since its active side is divisible
by $r$. This is the first congruence in
\eqref{paper:vacuum-congruence}; multiply by $d$ and use $ed=r-1$
for the second. Here $0<e,k<r$, and $d_X=p^2q^2$ is an $r$-unit,
so $m_0\ne0$.

The outside-orbit coefficients are constant on full Galois
orbits: the $q$-tori fix the column, pure $p$-inertia fixes $X$,
and pure $r$-inertia permutes the source orbit whose columns
coalesce. They are therefore constant on dual pairs. Only the
vacuum is self-dual in odd dimension, so the inactive-square
identity, with even right side $k(q+1-e k)$, makes $m_0$ even.
Its square costs at least four, proving $e k^2+4\le k(q+1)$.
At $k=2$ this gives $e\le(q-1)/2$; for $k\ge4$ the earlier
bound $e k\le q+1$ gives at least this much.

Together with $e=(r-1)/d$ and $qd=p^2-h\le p^2-2$, this
preliminary bound already gives
\[
 r\le\frac{p^2}{2}-\frac{p^2-2}{2q}<\frac{p^2}{2}.
\]
Since $r\ge3q+2$, it also gives $d\ge2(r-1)/(q-1)>6$;
oddness gives $d\ge7$, and $d<p$ gives $p\ge11$.
Completing the square in the inactive norm now yields
\[
 m_0^2\le\frac{(q+1)^2}{4e}
    =\frac{d(q+1)^2}{4(r-1)}<d\sqrt q<p^2.
\]
Here $r>q^{3/2}$, $d<p$, $q<p^2$, and
$(q+1)^2<4\sqrt q\,(q^{3/2}-1)$ for $q\ge3$.
Thus $m_0$ is a unit at both $p$ and $q$.

Let $C$ be the common active phase-degree sum in
\eqref{paper:linear-phase}. Summing all phases gives
$\delta(W_q\mathbf e_X)=d_X+C$, where $\delta$ is the dimension
functional. Symmetry and $\boldsymbol d=n s\mathbf e_0$ give
\[
 W_q\boldsymbol d=nW_q^2W_{q'}\mathbf e_0
                =nW_{q'}\mathbf e_0.
\]
The full complement includes any dyadic anomaly, and $W_q^2$
fixes the vacuum. Theorem~\ref{paper:lw:columns} gives
$C=(k_{q',X}-1)d_X$, so
\begin{equation}\label{paper:inactive-degree-sum}
 \sum_{\substack{A\notin\mathcal O\\\theta_{A,q}=1}}m_A d_A
          =(k_{q',X}-1+e k)d_X\in p^2q^2\mathbb Z.
\end{equation}
Reduction modulo $p$ and $q$ therefore forces both a $p$-free
and a $q$-free nonunit in this inactive outside support.
Their distinct dual pairs cannot coincide: a label omitting
both primes would have prime-power degree. Each pair costs
at least two squares, in addition to the vacuum's four.
Thus $e k^2+8\le k(q+1)$.
At $k=2$ this gives $e\le(q-3)/2$; for $k\ge4$, the bound
$e\le(q+1)/4$ suffices because $q\ge7$.
Using $e=(r-1)/d$ and $qd\le p^2-2$ gives the strengthened
range in \eqref{paper:quartic-range}.

For the other phase assignment, Corollary~\ref{paper:flat-period-rigidity}
already gives $rd=p^2-h<p^2$. This proves the assertion
independent of the phase assignment.

\end{proof}

\begin{corollary}[The coefficient and support of the quartic-phase column]
\label{paper:column-support}
In the odd-dimensional quartic-phase setting of
Proposition~\ref{paper:quartic-phase}, one has
\begin{equation}\label{paper:vacuum-size}
 \begin{aligned}
 0<|m_0|&<p,\qquad k_{\{q\},X}=p r^2m_0,\\
 v_p(k_{\{q\},X})&=1,\qquad v_r(k_{\{q\},X})=2.
 \end{aligned}
\end{equation}
The support of $qW_q\mathbf e_X$ contains two nonunit twist-one
labels $A,B\in P_q$, with
\[
 \begin{gathered}
 d_A=q^a r^b\quad(1\le a,b\le2),\qquad p\nmid d_A m_A,\\
 d_B=p^i r^j\quad(1\le i\le3,\ 1\le j\le2),\qquad q\nmid d_B m_B.
 \end{gathered}
\]
They may be chosen fixed by wild $p$-inertia and wild
$q$-inertia, respectively.
Every individual label in the column support has character field
tame at $r$. Each $q$-active supported label has character-field
conductor dividing $pqr$. In particular $K_A$ is tame at $p,r$,
and $K_B$ is tame at $q,r$.
\end{corollary}
\begin{proof}
The proposition's size bound makes $m_0$ a unit at $p$ and $r$.
Symmetry and
Theorem~\ref{paper:lw:columns} give
$m_0=q(W_q)_{0X}=k_{\{q\},X}/(p r^2)$, proving
\eqref{paper:vacuum-size}.

For $\ell=p$ or $q$, reduce \eqref{paper:inactive-degree-sum}
modulo $\ell$ and group its nonunit terms into wild
$\ell$-inertia orbits. Their dimensions and coefficients are
constant: pure $p$-inertia fixes $X$ and commutes with $W_q$,
while the $q$-tori fix its column. Every nonfixed orbit has
size divisible by $\ell$. Since the vacuum contributes the
nonzero residue $m_0$, some term is fixed and has
$\ell\nmid d_A m_A$, or $\ell\nmid d_B m_B$, respectively.
Both labels are $q$-inactive by construction and have
complementary twists one. The column lies in $P_q$; the principal
$q$-height bound, square divisibility and prime-power-degree
exclusion give the asserted degree ranges. The $q$-free label
is not asserted to lie in $P_r$.

For the fields at $r$, use
\[
 \sum_{A\notin\mathcal O}m_A^2
    =2qk-e k^2\le\frac{q^2}{e}
    =\frac{d q^2}{r-1}<r.
\]
Indeed $r(r-1)>q^3-r>q^3-q^2>d q^2$, since
$r>q^{3/2}$, $r<p^2<q^2$ and $d<q-1$.
Coalescence makes the outside-orbit coefficients constant on
full $r$-inertia orbits. A nontrivial wild orbit would contribute
at least $r$ squares. The source orbit already has prime-$r$
conductor, so every supported field is tame at $r$.

Each nontrivial $q$-phase has the sharper square mass
\[
 k\le\frac{d(q+1)}{r-1}<p.
\]
Its coefficient vector is fixed by complementary inertia,
and wild $q$-inertia preserves the phase fiber as well.
A nontrivial wild orbit at $p$, $q$ or $r$ would cost at
least that prime many squares. Thus every active supported
label is fixed by all three wild inertia groups. Only
$p,q,r$ can ramify in its character field, giving the claimed
conductor. These individual-field conclusions do not make the
full local cyclic module of each supported label tame.
\end{proof}

\begin{proposition}[Support from the inactive uniform endpoint]
\label{paper:flat-column-support}
Work in the odd-dimensional quartic-phase setting of
Proposition~\ref{paper:quartic-phase}.  Suppose that the full
$p$-block $B$ of the $q$-side endpoint $X$ satisfies
\begin{equation}\label{paper:inactive-uniform}
 W_p\mathbf e_X=\frac{\sigma}{p}\mathbf1_B,\qquad
 B\cap P_{p'}\cap\{A:\theta_{A,p}=1\}=\{X\},
 \qquad \sigma\in\{1,-1\}.
\end{equation}
In particular this holds when $X$ is the distinguished label
of a singleton-flat block.  Then the following conclusions hold.
\begin{enumerate}[label=(\roman*)]
\item Every label of $p$-height three contains both $q$ and $r$
in its degree.  Every compulsory witness on either side has
$p$-height two.
\item Every label in the support of $qW_q\mathbf e_X$ has
$p$-height at most two.  If a nonzero outside-orbit coefficient
$m_A$ has $r$-height two, then
\[
                 d_A=p^i r^2,\qquad 1\le i\le2.
\]
There is a common sign $\omega\in\{1,-1\}$ such that, for all
these labels,
\begin{equation}\label{paper:defect-zero-reciprocity}
             \chi_X(A)=\omega r^2m_A,\qquad
             m_A\equiv\omega p^i\pmod q.
\end{equation}
The $p$-free twist-one support witness of
Corollary~\ref{paper:column-support} can consequently be chosen
with degree $q^a r$, where $1\le a\le2$.
\item Each nontrivial $q$-phase contains a supported label
of $r$-height one.
\end{enumerate}
\end{proposition}
\begin{proof}
Both $X\in P_q$ and $X\in P_r$ by the saturated-orbit theorem.
A label $A$ with $p\nmid f_A$ has a singleton full $p$-block
and is $p$-inactive.  It also belongs to $P_{p'}$ by codegree
support.  The map $P_BW_{p'}$ intertwines the full local
$p$-group and sends $\mathbf e_A$ into $\mathbb C\mathbf e_X$,
by \eqref{paper:inactive-uniform}.  The cyclic module of
$\mathbf e_A$ is one-dimensional, whereas that of
$\mathbf e_X$ is not: its displayed Weyl image has $p^2$
nonzero coordinates.  Therefore this image is zero.
Factoring $s=W_pW_{p'}$ gives $S_{XA}=0$.
Principality of $X$ at both $q$ and $r$ then forces both
primes to divide $d_A$.  In particular a compulsory witness,
which omits one of them, cannot have $p$-height three.

For the stronger assertion about the $W_q$ column, let
$A\in P_q$ have $p$-height three.  The map $P_BW_q$ is also
a local $p$-intertwiner.  Since $W_q$ preserves $P_q$ and
$p$-inactivity, its value on $\mathbf e_A$ is
$(W_q)_{XA}\mathbf e_X$.  The same one-dimensional quotient
argument gives $(W_q)_{XA}=0$.
Symmetry and the support of the source column on $P_q$
prove the first assertion of (ii).  This local argument
is needed: vanishing of $S_{XA}$ alone would not imply
vanishing of $(W_q)_{XA}$.

Now let $m_A\ne0$ and $v_r(d_A)=2$.
Then $r\nmid f_A$ and the full $r$-block of $A$ is a
singleton.  The column preserves complementary twists,
so $\theta_{A,r}=1$.  The upper and lower local unipotents
therefore fix its coordinate line pointwise, and
$W_r\mathbf e_A=\mathbf e_A$.
In odd dimension the possible dyadic anomaly factor is a
scalar sign: its upper unipotent is scalar, and its Weyl
square is an unsigned torus.  Absorb that sign together
with $\sigma$ into $\omega$.  The factorization
$s=W_pW_{p'}$ and the inactive intersection in
\eqref{paper:inactive-uniform} now give
\[
 \frac{S_{XA}}{\sqrt D}
     =\frac{\omega}{p}(W_qW_r)_{XA}
     =\frac{\omega}{p}(W_q)_{XA}
     =\frac{\omega m_A}{pq}.
\]
Since $d_X=p^2q^2$ and $\sqrt D=p^3q^3r^2$, this is the
first identity in \eqref{paper:defect-zero-reciprocity}.
The outside squared norm is less than $r<p^2/2$, by
Corollary~\ref{paper:column-support}; hence every nonzero
outside coefficient satisfies $0<|m_A|<p<q$.
If $q\mid d_A$, principality of $X$ at $q$ would give
$q\mid m_A$, which is impossible.  Thus $A$ is $q$-free.
Prime-power-degree exclusion and the already proved
$p$-height bound give $d_A=p^i r^2$, $1\le i\le2$.
Reducing the same identity modulo $q$ gives its stated
congruence.  The $p$-free witness is $q$-bearing, so its
$r$-height is consequently one.

Finally suppose that a nontrivial $q$-phase were supported
entirely at $r$-height two.  Write $C$ for its signed degree
sum.  For each of its terms, (ii) gives
$m_A d_A\equiv\omega r^2m_A^2\pmod q$.
The phase-square identity and
$C=(k_{q',X}-1)d_X$ would therefore imply
\[
                   0\equiv C\equiv\omega r^2 k\pmod q.
\]
But $0<k<p<q$, a contradiction.
\end{proof}

\begin{proposition}[An integral transform on the complementary principal fiber]
\label{paper:join-fiber}
Under the hypotheses of
Proposition~\ref{paper:flat-column-support}, put
\[
 E=P_{p'}\cap\{A:\theta_{A,p}=1\},\qquad
 M_b=\sum_{\theta_{A,q}=\zeta_q^b}m_A\mathbf e_A
          \quad(b\in\mathbb F_q^\times).
\]
Then
\begin{equation}\label{paper:transformed-phase-integral}
                       rW_rM_b\in\mathbb Z^I.
\end{equation}
Every $q$-active label in $E$ has degree divisible by $r$.
Consequently $W_q^{\pm1}|_E$ are integral at $r$.
\end{proposition}
\begin{proof}
Write $P_b$ for the coordinate projection onto the indicated
$q$-phase.  The $r$-local action commutes with $P_b$ and $W_q$,
so
\[
                    W_rM_b=qP_bW_qW_r\mathbf e_X.
\]
The joined principal component $P_{p'}$ and the $p$-inactive
fiber are preserved by $W_q$ and $W_r$.  Thus this vector is
supported on $E$ and on the chosen nontrivial $q$-phase.
For $A$ in that support fiber, the inactive uniform cancellation
used in the preceding proposition gives
\[
                  \chi_X(A)=\omega r^2(W_rM_b)_A.
\]
The character value is an integer by
Theorem~\ref{paper:saturated-orbits}(i).
All other coordinates of $W_rM_b$ vanish, so
$r^2W_rM_b\in\mathbb Z^I$.

Every label occurring in $M_b$ has degree divisible by $r$,
by Proposition~\ref{paper:quartic-phase}; hence its formal
codegree has $r$-valuation at most two.  The indicator
identity used in the proof of Theorem~\ref{paper:lw:columns}
makes $f_A(W_r)_{TA}^2$, up to a root of unity, an
algebraic integer.  At every place above $r$ it follows
that $v_r((W_r)_{TA})\ge-1$ for these source columns.
The same lower bound holds for their integral linear
combination $W_rM_b$.  Combined with
$r^2W_rM_b\in\mathbb Z^I$, this removes one power of $r$
from its denominator and proves
\eqref{paper:transformed-phase-integral}.

Now let $A$ be any $q$-active label of $E$.
If its $q$-twist has order $q$, the displayed character
identity and \eqref{paper:transformed-phase-integral} give
$r\mid\chi_X(A)$, whether or not its transformed
coordinate is zero.  If that order exceeds $q$, the
prime-level support of $W_q\mathbf e_X$ and preservation
of $q$-twists by $W_r$ give
$(W_qW_r)_{AX}=0$.  Inactive uniform cancellation then
gives $\chi_X(A)=0$.  Since $X\in P_r$,
$\chi_X(A)\equiv d_A\pmod r$, so $r\mid d_A$.
The coordinate space of $E$ is invariant under the full
local $q$-group.  Its $q$-active labels consequently all
have $v_r(f_A)\le2$, and the last clause of
Theorem~\ref{paper:lw:columns} proves the final assertion.
\end{proof}

\begin{theorem}[Exclusion of the minimal phase mass]
\label{paper:phase-mass-exclusion}
Under the hypotheses of
Proposition~\ref{paper:flat-column-support}, one has $k\ge4$.
Consequently
\begin{equation}\label{paper:quarter-square}
 \begin{gathered}
 e\le\frac{q-1}{4},\qquad
 4(r-1)\le d(q-1),\\
 4q(r-1)\le(p^2-2)(q-1),\qquad
 r<\frac{p^2}{4}.
 \end{gathered}
\end{equation}
In particular $d\ge13$ and $p\ge17$.
\end{theorem}
\begin{proof}
Suppose $k=2$.  Each nontrivial $q$-phase has exactly two
coefficients of modulus one.  The local central inversion
$H_{-1,q}$ has no fixed active label, so it interchanges
them.  Their coefficients agree and their dimensions agree.
The common signed phase-degree sum therefore gives one
sign $\varepsilon$ and one degree $\delta$ for all these
pairs, with
\[
                  2\varepsilon\delta
                         =(k_{q',X}-1)p^2q^2.
\]
The full $q$-inertia fixes $X$, so
Theorem~\ref{paper:global-parity} makes $k_{q',X}$ odd.
Thus $p^2q^2\mid\delta$.
The principal $q$-height bound and
Proposition~\ref{paper:flat-column-support} give
$\delta=p^2q^2r^b$, $b=1$ or $2$.
The same proposition excludes $r$-height two in
$q$-bearing support, hence
\[
                              \delta=p^2q^2r.
\]

Fix one phase pair and put
\[
 v=M_b=\varepsilon(\mathbf e_A+\mathbf e_{A'}),\qquad
 w=rW_rv,\qquad x=\langle v,w\rangle.
\]
Proposition~\ref{paper:join-fiber} gives
$w\in\mathbb Z^I$ and $\|w\|^2=2r^2$.
Both vectors are supported in
$E\cap\{\theta_q=\zeta_q^b\}$, where
$E=P_{p'}\cap\{\theta_p=1\}$.
The vector $v$ is upper-$r$-Borel fixed:
its labels are $r$-inactive, and its coefficients are
invariant under full $r$-inertia.
Thus $w$ is fixed by the full $r$-torus.
It is also fixed by $H_{-1,q}$, which fixes $v$ and
commutes with the $r$-local action.
The join-fiber character identity gives
$\chi_X(A)=\omega r w_A$.
Since $X\in P_q$ and $q\mid d_A$, one has $q\mid w_A$.
The two source coefficients of $w$ agree, so
\begin{equation}\label{paper:paired-coefficient}
                         x=2\varepsilon w_A\in2q\mathbb Z.
\end{equation}

Every target in this fiber is $r$-bearing by
Proposition~\ref{paper:join-fiber}.  Its $r$-codegree
valuation is at most two, so
Lemma~\ref{paper:family:lw:weyl-line} bounds its actual $r$-twist
order by $r^2$.
At any label active at both $q$ and $r$, the involutions
$H_{-1,q}$ and $H_{-1,r}$ generate a free four-element
orbit: neither inversion nor their product can fix
such a label, by Theorem~\ref{paper:global-parity}.
They preserve each exact primary-twist fiber and fix $w$.
The squared mass of its integer coefficients at every
nontrivial $r$-phase is consequently divisible by four.

Let $I$ be its squared mass at the inactive $r$-phase.
For a local unit $a$, Bruhat gives
\[
                   \langle w,U_r(a)w\rangle=rx.
\]
Primitive $r^2$-phases contribute zero here: torus
invariance makes their squared masses constant on
the $r$ lifts of a nonzero residue modulo $r$, whose
roots sum to zero.
Fourier inversion modulo $r$ therefore makes the squared
mass at every nontrivial order-$r$ phase a common
$M\in4\mathbb Z_{\ge0}$, with $I-M=rx$.
Let $P_+,P_-\in4\mathbb Z_{\ge0}$ be the squared masses
per primitive phase in the two square classes modulo $r$,
and put $L=(P_++P_-)/2\in2\mathbb Z_{\ge0}$.
The norm identity yields
\[
              M=2r-x-(r-1)L,\qquad
              I=2r+(r-1)(x-L).
\]
If $L\ge4$, then $x\le4-2r$ and
$I\le2r(2-r)<0$.
If $L=2$, positivity gives $0\le x\le2$,
and $4\mid M=2-x$ forces $x=2$,
contrary to \eqref{paper:paired-coefficient}.
Hence $L=0$: the vector $w$ has only $r$-twists of
order dividing $r$.
Now $I=2r+(r-1)x\ge0$ and
\eqref{paper:paired-coefficient} imply $x\ge0$, so
\begin{equation}\label{paper:small-transformed-mass}
                           M=2r-x\le2r<p^2.
\end{equation}

Every nonzero coefficient at a nontrivial $r$-phase
therefore has $|w_T|<p<q$.
The identity $\chi_X(T)=\omega r w_T$ and principality
of $X$ at $q$ force $q\nmid d_T$.
Such a target is $r$-bearing, while $r$-activity
excludes $r$-height two.
The global height-three exclusion in
Proposition~\ref{paper:flat-column-support} and
prime-power-degree exclusion leave
\[
                             d_T=pr\ \text{or}\ p^2r.
\]

It remains to use the signed degree sums.
Let $\Delta(\mathbf e_T)=d_T$ be the dimension functional.
It is fixed by lower $r$-unipotents, so unit Bruhat gives
\[
                   \Delta(U_r(a)w)=r\Delta(v)
                                      =2\varepsilon r\delta
                         \qquad(a\not\equiv0\pmod r).
\]
Symmetry and $W_r\boldsymbol d=nW_{r'}\mathbf e_0$ give
\[
                \Delta(w)=2\varepsilon r\delta K,\qquad
                K=k_{r',A}=k_{r',A'}\in r\mathbb Z.
\]
Here the complementary vacuum column is fixed by the
$q$-torus, which interchanges $A,A'$; its two coefficients
therefore agree.  The integrality and divisibility follow
from Theorem~\ref{paper:lw:columns}, because
$f_A=p^2q^2r^2$ and $r$ is foreign to $r'$.
Fourier inversion on $\mathbb F_r$ now gives, for every
$c\ne0$,
\[
 J_c:=\sum_{\theta_{T,r}=\zeta_r^c}w_Td_T
                         =2\varepsilon\delta(K-1).
\]
Since $K\in r\mathbb Z$, this has $|J_c|\ge2\delta$.
But the target degrees and integral coefficients give
\[
 |J_c|
   \le p^2r\sum_{\theta_{T,r}=\zeta_r^c}|w_T|
   \le p^2r M
   \le2p^2r^2
   <2p^2q^2r=2\delta,
\]
where $r<p^2/2<q^2$.
This contradiction excludes $k=2$.

The phase mass is positive and even, hence $k\ge4$.
From $e k^2+8\le k(q+1)$ one obtains
\[
 e\le\frac{q+1}{k}-\frac8{k^2}\le\frac{q-1}{4};
\]
after multiplication by $4k^2$, the second inequality
is $(k-4)((q-1)k-8)\ge0$.
Substituting $e=(r-1)/d$ and $qd=p^2-h\le p^2-2$
gives the first three bounds in
\eqref{paper:quarter-square}.
The bound $r\ge3q+2$ implies
$d\ge4(r-1)/(q-1)>12$; thus $d\ge13$ and $p>d$
gives $p\ge17$.
Finally $h\ge2$ and $d\ge13$ give
\[
             r\le1+\frac{d(q-1)}4
                =\frac{p^2+4-h-d}{4}<\frac{p^2}{4}.
\]
\end{proof}

The integrality argument depends on the quartic exponent at $r$.
With phase prime $q$, an $r$-active label in $P_r$ is only known
to be $q$-bearing, which permits $v_q(f_A)=4$; the Bruhat products
can still retain foreign denominators. Neither the opposite
orbit average nor an entire residue block is made tame here.

\begin{proposition}[Projection amplification of orbit energy]
\label{paper:projection-energy}
Let $B_1,\ldots,B_s$ be distinct homogeneous defect-two full
$\ell$-blocks, under the hypotheses of
Theorem~\ref{paper:integral-projection}. Choose $X_i\in B_i$, put
$g_i=(G_{B_i})_{X_iX_i}$, and let $Y\in P_{\ell'}$.
Then
\[
 \sum_{i=1}^s\frac{\ell^2}{g_i}|S_{X_iY}|^2
       \le D-d_Y^2<D.
\]
In particular, this applies to a Galois orbit whenever its chosen
labels belong to pairwise distinct full $\ell$-blocks.
\end{proposition}
\begin{proof}
Multiplication by $E_{P_{\ell'}}$ fixes $s^*\mathbf e_Y$ and
preserves every full $\ell$-block. Its restriction to $B_i$ is
the orthogonal projection $P_i=G_{B_i}/\ell^2$. Thus the vector
$v_i=(\overline{S_{ZY}})_{Z\in B_i}$ belongs to its range. The
coordinate Cauchy--Schwarz inequality gives
\[
 |S_{X_iY}|^2
   =|\langle P_i\mathbf e_{X_i},v_i\rangle|^2
   \le\frac{g_i}{\ell^2}\sum_{Z\in B_i}|S_{ZY}|^2.
\]
Sum over the disjoint blocks. They are nonprincipal because
their homogeneous height is $m-1>0$, so none contains the unit.
Column orthogonality and $S_{\mathbf1Y}=d_Y$ give the result.
\end{proof}


\begin{remark}[Scope of the energy bounds]
The projection-specific estimates retain their stated two-height-two
or averaging hypotheses. The ordinary braid bound applies to every
compulsory pair of distinct full orbits. Its flat specialization
requires no inactivity assumption on the partner. The broader
inactive-fiber amplification requires a $p$-inactive partner;
with that condition it gives the same primary-level exclusions.
The top-height vanishing holds for every isolated inactive endpoint,
but $p$-inactivity of all opposite witnesses is proved only under
the stronger singleton-flat hypothesis.
\end{remark}


\subsubsection{Small intersections and their local action}
\label{paper:sec:flat}

\begin{theorem}[A singleton intersection forces an equal-degree block
at the least prime]
\label{paper:singleton}
Under the hypotheses of Proposition~\ref{paper:packet-norm},
suppose $I_B=\{Z\}$.
Put $A_J=\ell^{2m}E_{P_J}$, the integral signed projector of
Proposition~\ref{paper:block:principal-projector}.  Then
\begin{equation}\label{paper:singleton-data}
 c_Z=\varepsilon\in\{1,-1\},\qquad
 K_B=u_Z^2,\qquad
 t_X:=u_X/u_Z\in\mathbb Z_{>0},\qquad
 \sum_{X\in B}t_X^2=\ell^2.
\end{equation}
The label $Z$ is $\ell$-inactive and fixed by all $\ell$-inertia,
and
\begin{equation}\label{paper:singleton-column}
 W_\ell\mathbf e_Z=\frac{\varepsilon}{\ell}
                         \sum_{X\in B}t_X\mathbf e_X,
 \qquad
 A_J[X]=\ell^{2m-2}t_X\sum_{Y\in B}t_Y[Y]
                       \quad(X\in B).
\end{equation}
In particular $\tau(E_B)=u_Z^2/(n')^2$ is an $\ell$-unit.

If $\ell$ is the least dimension prime, every $t_X$ equals one.
Thus $|B|=\ell^2$ and every label in $B$ has the same degree
$d_Z$.  In the ordered target family, a homogeneous block at the
least sextic prime $p$ can consequently have a singleton
complementary principal intersection only if it consists of
exactly $p^2$ labels of equal degree.  If $D$ is odd, a block
with a singleton intersection cannot be self-dual.
\end{theorem}
\begin{proof}
Packet localization gives $C_B=c_Zu_Z[Z]$ and
$K_B=u_Z^2c_Z^2$, with $c_Z\ne0$.  The integral projection
$G$ of Theorem~\ref{paper:integral-projection} has rank one and
fixes $u$ with eigenvalue $\ell^2$.  Thus
\[
 G_{XY}=\frac{u_Xu_Y}{K_B},
 \qquad
 G_{ZZ}=\frac1{c_Z^2}\in\mathbb Z_{>0}.
\]
Consequently $c_Z=\varepsilon=\pm1$ and
$t_X=G_{XZ}=u_X/u_Z$ is a positive integer.  The packet norm
gives $\sum_Xt_X^2=\ell^2$.  Also
\begin{equation}\label{paper:singleton-integrality}
 (N_{A_J})_{XY}=\ell^{2m-2}G_{XY}
                         =\ell^{2m-2}t_Xt_Y.
\end{equation}
This proves the multiplication formula.  The identity
$W_\ell C_B=x_B$ proves the Weyl-column formula.

Both $B$ and its complementary principal intersection are preserved
by local inertia.  Their unique common label $Z$ is therefore fixed.
Its nonzero packet coefficient makes it inactive, and
$D_B=\ell^{2m}K_B$ gives the stated trace.

If $\ell$ is the least dimension prime, every prime factor of an
integer $t_X>1$ is greater than $\ell$, contradicting
$t_X^2\le\sum_Yt_Y^2=\ell^2$.  Thus all $t_X=1$.
If $D$ is odd and $B=B^*$, duality would fix the unique common
label $Z$.  Its degree has positive $\ell$-height, so it is a
nonunit, contradicting the absence of self-dual nonunits in odd
dimension.
\end{proof}


Assume the exceptional configuration of
Theorem~\ref{paper:singleton} at the least dimension prime $p$.
Thus a homogeneous defect-two block $B$ has $p^2$ labels of one
degree, $B\cap P_{p'}=\{Z\}$, and, on $\mathbb C^B$,
\begin{equation}\label{paper:flat-setup}
 W\mathbf e_Z=\frac{\varepsilon}{p}\mathbf1_B,
 \qquad W\mathbf1_B=\varepsilon p\mathbf e_Z,
 \qquad \varepsilon\in\{1,-1\},\qquad W=W_p|_B.
\end{equation}
The label $Z$ is inactive and fixed by all local tori.  Write
$T=\operatorname{diag}(\theta_{X,p})_{X\in B}$ and let
$\mathcal V_Z$ be the cyclic subspace generated by $\mathbf e_Z$
under the local modular group.  The following restrictions use
the local group and Galois relations; no fusion multiplication
on the label set $B$ is assumed.

\begin{corollary}[Integer row sums and the fusion escape bound]
\label{paper:flat-fusion}
Put $R_B=\sum_{X\in B}[X]$, $v_B=R_B/p$ and
$H_B=v_Bv_B^*$.  The object $H_B$ is integral and nonnegative,
with unit multiplicity one and dimension $D_B$.
For every simple $Y$, the actual fusion submatrix
$(N_Y)_{B,B}$ has constant row and column sums
\[
 a_Y=[Y]H_B\in\mathbb Z_{\ge0},\qquad
 M_Y(B,B)=K_Ba_Y.
\]
In a perfect fusion-simple table, every nonunit $Y$ therefore has
\[
 a_Y\le d_Y-1,\qquad
 E_Y(B)=K_B(d_Y-a_Y)\ge K_B,\qquad
 a_{Y^*Y}-a_Y^2\ge1.
\]
Moreover $v_B$ has norm one and $Rv_B$ is an integral
fusion-stable lattice in its rational span.
\end{corollary}
\begin{proof}
The complementary principal projection restricts to $J_{p^2}/p^2$
on $B$.  Apply Proposition~\ref{paper:averaging-products} with
$k=p^2$.  It gives the object, row sums, integer gap and lattice.
The common degree and the definition of $M_Y(B,B)$ give
$M_Y(B,B)=K_Ba_Y$, hence the displayed escape bound.
The last inequality follows from
\eqref{paper:averaging-leakage}: strict expansion forces
$YR_B$ to have a constituent outside the proper block $B$.
\end{proof}

\subsection{Odd flat blocks: phases and local inertia}


\begin{corollary}[Local inversion in an odd flat block]
\label{paper:flat-charge}
In the odd-dimensional singleton-flat setup \eqref{paper:flat-setup},
the local central torus $H_{-1,p}$ fixes exactly $Z$.
Every local-inertia-stable subset of $B\setminus\{Z\}$ has even
cardinality, and every primary phase has even multiplicity on
such a subset. Moreover $\dim\mathcal V_Z\le(p^2+1)/2$.
\end{corollary}
\begin{proof}
All degrees on $B$ agree. The identity $W_pd=nW_{p'}\mathbf e_0$
and \eqref{paper:flat-setup} give
$k_{p',X}=\varepsilon p\,\delta_{XZ}$ for $X\in B$.
Theorem~\ref{paper:global-parity} therefore gives precisely the
one fixed label. The involution preserves every primary phase;
its other cycles have length two. Its positive eigenspace has
dimension $(p^2+1)/2$, and contains $\mathcal V_Z$ because the
central involution fixes its generating vector.
\end{proof}

\begin{theorem}[Three quadratic phase laws]
\label{paper:odd-patterns}
Assume the odd-dimensional singleton-flat setup
\eqref{paper:flat-setup}, and fix a primitive $p^2$th root $\zeta$.
Let $N_1$ be the number of inactive labels and $N_p$ the number
at each specified nonidentity $p$th root. For $a\not\equiv0
\pmod p$, the multiplicity of $\zeta^a$ is a number $b_+$ or
$b_-$ according as $a\bmod p$ is a square or a nonsquare.
Up to interchanging the two square classes, the only possibilities are
\[
\begin{array}{c|c|c|c|c|c}
 \text{law}&\varepsilon&N_1&N_p&(b_+,b_-)&\dim\mathcal V_Z\\ \hline
 \text{anisotropic}&-1&1&p+1&(0,0)&p\\
 \text{split}&+1&2p-1&p-1&(0,0)&p+1\\
 \text{cyclic}&+1&p&0&(2,0)&(p^2+1)/2
\end{array}
\]
For every integer $t$ prime to $p$,
\begin{equation}\label{paper:unit-gauss}
                    \sum_{X\in B}\theta_{X,p}^{\,t}=\varepsilon p.
\end{equation}
In the first two laws the entire local representation factors
through $\operatorname{SL}_2(\mathbb F_p)$; in particular wild
$p$-inertia fixes every label of $B$. The cyclic law has primitive
twists of order $p^2$ and will be called wild.
\end{theorem}
\begin{proof}
The codegree bound gives twist orders dividing $p^2$ and
character-field conductor at $p$ dividing $p^2$. The representation
on $B$ consequently factors through
$G=\operatorname{SL}_2(\mathbb Z/p^2)$: an element congruent
to the identity modulo $p^2$ has a lower--diagonal--upper
decomposition, whose three factors act trivially by these
two bounds. Inflate from level $p$ if necessary.

Let $\mathcal B$ be the upper Borel. Since $p$ is odd, its
upper unipotent is the actual primary twist, so $\mathbf e_Z$
is $\mathcal B$-fixed. A lower unipotent with unit parameter
belongs to $\mathcal B W\mathcal B$. Its matrix coefficient
at $\mathbf e_Z$ is therefore $\varepsilon/p$.
Using \eqref{paper:flat-setup} to express the same coefficient as
a mean of primary twists proves \eqref{paper:unit-gauss}.

Write $N_a=\#\{X:\theta_{X,p}=\zeta^a\}$ for
$0\le a<p^2$. The equality at $t=1$ and the relation
$\Phi_{p^2}(x)=1+x^p+\cdots+x^{(p-1)p}$ give integers $b_j$ with
\begin{equation}\label{paper:phase-coefficients}
 N_a-\varepsilon p\,\delta_{a0}=b_{a\bmod p},
 \qquad
 b_0+\sum_{j=1}^{p-1}b_j=p-\varepsilon.
\end{equation}
All $b_j$ are nonnegative, as $b_0=N_p$.
If $\varepsilon=-1$, then $b_0=N_1+p\ge p+1$, forcing
$b_0=p+1$ and all other $b_j=0$.
For $\varepsilon=1$, the square action of the tori makes
$b_j$ constant on the two nonzero square classes, and
\[
               b_0+\frac{p-1}{2}(b_++b_-)=p-1.
\]
The local central involution pairs every active phase fiber by
Theorem~\ref{paper:global-parity}. Hence $b_+$ and $b_-$ are even.
Their sum is at most two, leaving exactly $(0,0)$ and, up to
interchange, $(2,0)$. This proves the three phase distributions
directly.

For the cyclic ranks, form the Gram matrix of the vectors
$g\mathbf e_Z$, $g\mathcal B\in G/\mathcal B$.
The projective line $G/\mathcal B$ has $p+1$ reduction fibers,
each of size $p$. Representatives can be chosen as
$L(-x)$, $x\in\mathbb Z/p^2$, and $U(pt)W$, $t\in\mathbb F_p$.
Vectors from distinct fibers have inner product $\varepsilon/p$.
Within a fiber the Gram matrix is circulant, with entry at
difference $k$ equal to
\[
 h(k)=\frac1{p^2}\sum_X\theta_{X,p}^{pk}
     =\frac\varepsilon p+\frac1p
                             \sum_{j=0}^{p-1}b_j\zeta_p^{jk}.
\]
Its Fourier eigenvalues are the $b_j$ for $j\ne0$ and
$\varepsilon+b_0$ on the constant vector. Thus the full Gram
matrix, on the fiber-constant subspace, is
$b_0I+\varepsilon\mathbf1\mathbf1^{\mathsf T}$; each nonzero
$b_j$ also contributes one dimension in each of the $p+1$ fibers.
Its rank is
\[
 p\,\mathbf1_{b_0>0}
 +\mathbf1_{b_0+\varepsilon(p+1)>0}
 +(p+1)\#\{j\ne0:b_j>0\},
\]
giving the last column. These dimensions also respect the bound
$\dim\mathcal V_Z\le(p^2+1)/2$: the central involution fixes
exactly $Z$, has a positive eigenspace of that dimension, and
is the identity on the cyclic module generated by $\mathbf e_Z$.

In the first two laws $U(p)$ acts trivially. The level-$p$
congruence kernel is the additive group
$I+p\mathfrak{sl}_2(\mathbb F_p)$, and the conjugates of $U(p)$
span it. The representation therefore factors through level $p$.
\end{proof}

The names describe the phase distributions. For a nonsquare
$\delta\in\mathbb F_p$, the form $x^2-\delta y^2$ on
$\mathbb F_p^2$ has one zero and $p+1$ representations of each
nonzero value, by the norm map from $\mathbb F_{p^2}$.
The split form $xy$ has $2p-1$ zeros and $p-1$ representations
of each nonzero value. On $\mathbb Z/p^2$, the form $ax^2$ for
a unit $a$ has $p$ zeros and two representations of each unit
value in one square class. These counts do not identify the
full local Weyl matrix or the ambient tensor multiplication.

\subsubsection{Complementary intersections}

\begin{theorem}[Complementary orbits in an odd flat block]
\label{paper:flat-orbits}
Suppose that $p$ is odd and that a homogeneous defect-two block $B$
has $p^2$ labels of one degree and
$B\cap P_{p'}=\{Z\}$, as in \eqref{paper:flat-setup}.
Consider the orbits of the full Galois group on the nonprincipal
components of the complementary joined partition.  Retain those
orbits whose unions $C_1,\ldots,C_t$ meet $B$, and put
\[
 I_a=B\cap C_a,\qquad h_a=|I_a|,\qquad
 G_a=p^2W^*P_{I_a}W,
 \qquad W=W_p|_B.
\]
Then
\begin{equation}\label{paper:flat-orbit-projections}
 \begin{gathered}
 G_a\in\operatorname{Mat}_{p^2}(\mathbb Z),\quad
 G_a^{\mathsf T}=G_a\succeq0,\quad
 G_a^2=p^2G_a,\quad \operatorname{rank}G_a=h_a,\\
 G_a\mathbf1=0,\qquad
 G_aG_b=0\quad(a\ne b),\qquad
 J_{p^2}+\sum_{a=1}^tG_a=p^2I_{p^2},\qquad
 (G_a)_{XY}\equiv h_a\pmod p.
 \end{gathered}
\end{equation}
Every $h_a$ is even, every diagonal entry of every $G_a$ is positive, and
\begin{equation}\label{paper:flat-orbit-gap}
                  h_a=p-1\quad\text{or}\quad h_a\ge p+1.
\end{equation}
The same gap holds for the intersection size of any Galois-stable
union of nonprincipal complementary components that meets $B$.

If $h_a=p-1$, there is a partition $\mathcal Q_a$ of $B$ into
$p$ sets of $p$ labels such that
\begin{equation}\label{paper:parallel-projector}
 G_a=p\!\bigoplus_{Q\in\mathcal Q_a}J_p-J_{p^2}.
\end{equation}
For distinct minimal orbits, the corresponding partitions are
orthogonal: each set of one meets each set of the other in exactly
one label.
In particular $t\le p+1$; equality forces every $h_a=p-1$,
and these $p+1$ partitions are the parallel classes of an affine
plane of order $p$ on the label set $B$.
\end{theorem}
\begin{proof}
Each $C_a$ is a Galois-stable union of complementary components.
At the fixed chosen residue places, Galois covariance permutes
the full-block partitions, as in the proof of
Proposition~\ref{paper:block:principal-projector}.  Thus $E_{C_a}$ is
rational and self-dual.  Local integrality at all complementary
primes removes those prime powers from its denominator; hence
$p^{2m}E_{C_a}\in R$, where $v_p(D)=2m$.
The proof of Theorem~\ref{paper:integral-projection} therefore gives
the integrality and symmetry of $G_a$, without requiring $C_a$
to be principal.  Its projection formula gives positivity,
rank and the quadratic identity.  Density of $W$ gives positive
diagonals.  Since $C_a$ omits the unit label, multiplication by $E_{C_a}$
annihilates the global degree vector. It preserves $B$, where all
degrees agree, so $G_a\mathbf1=0$.

The intersections $\{Z\},I_1,\ldots,I_t$ partition $B$.
Their coordinate projections are orthogonal.  The flat column
$W\mathbf e_Z=(\varepsilon/p)\mathbf1$ gives
$p^2W^*P_{\{Z\}}W=J_{p^2}$, proving the sum and orthogonality
in \eqref{paper:flat-orbit-projections}.
All scaled degrees are one common number $u$, and the packet has
the single coefficient $c_Z=\varepsilon$.  Thus $S_B=\varepsilon u^2$;
the residue calculation in \eqref{paper:projection-residue}, now
summed over $I_a$, gives $(G_a)_{XY}\equiv h_a\pmod p$.
Exactly the same argument applies to a Galois-stable union.

Write $G=G_a$ and $h=h_a$.  Suppose first that $0<h<p$.
Every positive diagonal entry is congruent to $h$ modulo $p$,
so is at least $h$.  Since $\operatorname{Tr}G=p^2h$ and $G$
has $p^2$ rows, all diagonal entries equal $h$.
The Gram bound gives $|G_{XY}|\le h$.  Together with the residue
formula, this leaves just the values $h$ and $h-p$.
In a real Gram realization, inner product $h$ means that the two
vectors coincide.  It therefore partitions the labels into classes,
with inner product $h-p$ between distinct classes.
If such a class has size $s$, the zero row sum gives
\[
             sh+(p^2-s)(h-p)=0,
             \qquad s=p(p-h).
\]
All classes have this size, so their number is $p/(p-h)$.
Since $0<p-h<p$ and $p$ is prime, necessarily $h=p-1$.
There are then $p$ classes of size $p$, and their Gram matrix is
\eqref{paper:parallel-projector}.

Every intersection under consideration is stable under local
inertia and omits $Z$, so Corollary~\ref{paper:flat-charge}
makes its size even. In particular $h=p$ is impossible.
This proves \eqref{paper:flat-orbit-gap}.

For two minimal intersections, let $A_a$ be the block matrix
$\bigoplus_{Q\in\mathcal Q_a}J_p$.  Then $A_aJ_{p^2}=pJ_{p^2}$,
and the orthogonality $G_aG_b=0$ becomes
\[
                         A_aA_b=J_{p^2}.
\]
Its entries count intersections of the two classes through the
corresponding labels.  Hence all these intersections have size one.
Finally $\sum_a h_a=p^2-1$ gives $t\le p+1$.
At equality every rank is $p-1$.  The resulting $p+1$ orthogonal
parallel classes give an affine plane: two points lie together
in at most one line, while the $p+1$ lines through a point contain
all $p^2-1$ other points.
\end{proof}

\begin{corollary}[Fusion products and the wild-inertia alternative]
\label{paper:parallel-fusion}
For each minimal orbit in Theorem~\ref{paper:flat-orbits},
multiplication by $E_{P_{p'}}+E_{C_a}$ restricts on $B$ to
$\bigoplus_{Q\in\mathcal Q_a}J_p/p$.
Write $\mathcal Q_a=\{Q_1,\ldots,Q_p\}$ and
$R_i=\sum_{X\in Q_i}[X]$.  One consequently has
\[
 H_{ij}:=\frac{R_iR_j^*}{p}\in R_{\ge0},\qquad
 R_iX^*=H_{ij}\quad(X\in Q_j),\qquad
 [\mathbf1]H_{ij}=\delta_{ij},\qquad
 \operatorname{FPdim}(H_{ij})=p d_B^2.
\]
Here $d_B$ is the common simple degree in $B$.

If there are two minimal orbits, wild $p$-inertia fixes every label
of $B$, and every actual primary twist on $B$ has order dividing
$p$.  Thus only the first two laws of
Theorem~\ref{paper:odd-patterns} remain.
\end{corollary}
\begin{proof}
Add the principal projection $J_{p^2}/p^2$ to
\eqref{paper:parallel-projector} divided by $p^2$.
Proposition~\ref{paper:averaging-products} applies with $k=p$.

Local tori preserve $B$ and each $I_a$, because $C_a$ is a full
Galois-stable union.  The relation $H_{s,p}W=W H_{s,p}^{-1}$ shows that
they commute with $G_a$, so they permute the classes of every
$\mathcal Q_a$.  The distinguished label $Z$ is fixed by all
local tori.  Work at level $p^2$, as allowed by the codegree and
conductor bounds for this block.  A wild torus has order dividing
$p$.  Its permutation of the $p$ classes of a minimal partition
fixes the class containing $Z$, and therefore fixes every class.
With two orthogonal partitions it fixes their singleton
intersections, hence every label of $B$.
The torus at $1+p$ then acts trivially.  Its conjugation of the
actual primary twist is $T\mapsto T^{(1+p)^2}$; since $T^{p^2}=1$
and $p$ is odd, this implies $T^p=1$.

\end{proof}


\subsubsection{The tame laws}

\begin{theorem}[Tame inertia and the minimal complementary classes]
\label{paper:tame-orbits}
Assume the odd-dimensional singleton-flat setup
\eqref{paper:flat-setup}, and suppose its phase law is tame.
Every active label has a regular orbit under the local torus
group $(\mathbb Z/p)^\times$.
\begin{enumerate}[label=(\roman*)]
\item In the anisotropic law the inertia orbits are $\{Z\}$ and
exactly $p+1$ regular orbits of size $p-1$. Hence every nonprincipal
complementary component-orbit intersection has size $n_a(p-1)$,
where $n_a\ge1$ and $\sum_a n_a=p+1$.
\item In the split law there are exactly $p-1$ active regular
orbits. The other $2p-2$ inactive labels, besides $Z$, form
orbits of even lengths dividing $p-1$.
\end{enumerate}

For a minimal nonprincipal complementary intersection $I_a$ of
size $p-1$, let $\mathcal Q_a$ be its averaging partition from
Theorem~\ref{paper:flat-orbits}, let $Q_a$ be the class containing
$Z$, and put
\[
             A_a=\{Z\}\cup I_a,
             \qquad J_a=Q_a\setminus\{Z\}.
\]
Then the following alternatives hold.
\begin{enumerate}[label=(\alph*)]
\item If $I_a$ contains an active label, it is one regular active
orbit. Then $J_a$ is a different regular active orbit, and
$A_a\cap Q_a=\{Z\}$.
\item If $I_a$ is entirely inactive, then the law is split and
$Q_a=A_a$. Its indicator is fixed by the entire local group.
\end{enumerate}
The sets $J_a$ for distinct minimal intersections are disjoint.
If $J_a=I_b$ for another minimal intersection, then $J_b=I_a$.
In particular, in the anisotropic saturated case of $p+1$ minimal
intersections, the parallel classes of the resulting affine plane
are paired by a fixed-point-free involution.
\end{theorem}
\begin{proof}
The tame representation factors through $\operatorname{SL}_2(\mathbb F_p)$.
If a torus fixes an active label, its square action on the nontrivial
$p$th-root phase gives $s^2=1$. The nonidentity possibility $s=-1$
is excluded by Corollary~\ref{paper:flat-charge}. Thus every active
orbit has size $p-1$. The phase counts in
Theorem~\ref{paper:odd-patterns} give the asserted orbit numbers.
Every other nontrivial orbit has even length, since an odd-length
orbit would have a stabilizer containing $-1$.
Complementary component-orbit intersections are local-inertia
stable, giving the anisotropic divisibility and sum.

The local tori permute the classes of $\mathcal Q_a$ and fix
$Q_a$, because they fix $Z$. Thus both $I_a$ and $J_a$ are
inertia-stable sets of size $p-1$. In a tame block each is either
entirely inactive or one regular active orbit.
Write $v_A=\mathbf1_{A_a}$ and $v_Q=\mathbf1_{Q_a}$.
The projection formula, evaluated at $Z$, gives
\[
 W\mathbf1_{I_a}
      =\varepsilon\left(\mathbf1_{Q_a}-p^{-1}\mathbf1_B\right),
 \qquad
 Wv_A=\varepsilon v_Q,
 \qquad Wv_Q=\varepsilon v_A.
\]
Here $W^*=W$ on inertia-invariant indicator vectors. Also
\[
               W^*P_{A_a}W
                    =\bigoplus_{Q\in\mathcal Q_a}J_p/p.
\]
Hence the columns of $W^*$ indexed by $A_a$ are constant on
$Q_a$. Their entries there equal their entries at $Z$, namely
$\varepsilon/p$.

If $A_a=Q_a$, put $\tau=\sum_{X\in A_a}\theta_{X,p}$.
The local relation $WTW=T^{-1}WT^{-1}$, evaluated between $v_A$
and itself, gives
\[
                         \tau=\frac{\varepsilon}{p}\overline\tau^{2}.
\]
Thus $\tau=0$ or $|\tau|=p$. A zero sum of $p$ many $p$th roots
has each phase once, contradicting the even multiplicities away
from $Z$ in Corollary~\ref{paper:flat-charge}. In the second case
equality in the triangle inequality makes all phases equal, hence
all one because $Z\in A_a$. Thus a self-paired $A_a$ is inactive.

Suppose $A_a$ is inactive and $Q_a$ contains active labels.
The law must be split, so $\varepsilon=1$. Its active part $J_a$
is a regular orbit, with two labels at each phase in one square
class. Therefore
\[
 \left|\sum_{X\in Q_a}\theta_{X,p}\right|=\sqrt p,
\]
by the elementary quadratic Gauss sum. On the other hand,
$Tv_A=v_A$ and the same local relation give
\[
 \sum_{X\in Q_a}\theta_{X,p}=\langle v_A,Wv_A\rangle
                          =|A_a\cap Q_a|=1,
\]
a contradiction. The same argument with $A_a,Q_a$ interchanged
excludes an active $A_a$ and inactive $Q_a$.
When both are inactive, the displayed relation gives
$p=|A_a\cap Q_a|$, hence they coincide. Their indicator is then
fixed by $T$, by the tori, and by $W$, so by the local group.
When both contain active labels, equality was already excluded;
their distinct regular orbits are disjoint. This proves (a)--(b).

The distinguished classes for distinct minimal intersections meet
only at $Z$, by orthogonality of the corresponding partitions.
Thus their $J_a$ are disjoint. Finally $Wv_A=\varepsilon v_Q$
and its reverse show that whenever $J_a=I_b$, the distinguished
class attached to $I_b$ is $A_a$. In the saturated anisotropic
case every regular orbit is one minimal intersection, giving the
stated fixed-point-free involution. No coordinatization of the
affine plane is used.
\end{proof}

\begin{remark}[The inactive vectors in the split law]
\label{paper:split-invariant}
Let $I$ be the inactive set of a split flat block. Averaging the
distinguished vector over the local group, using the $p+1$ Borel
cosets, gives
\[
 P_G\mathbf e_Z=\frac{\mathbf e_Z+\mathbf1_I}{p+1}.
\]
In particular $\mathbf e_Z+\mathbf1_I$ is group-fixed. If a
minimal inactive intersection gives the $p$-set $A_a$ above, then
the complementary $p$-set
$A'_a=\{Z\}\cup(I\setminus A_a)$ also has a group-fixed indicator.
This does not assert that $A'_a\setminus\{Z\}$ is a complementary
component-orbit intersection or that either indicator defines an
algebra object.
\end{remark}

\subsubsection{The wild law}

\begin{lemma}[The congruence spectrum of a Borel-fixed vector]
\label{paper:cyclic-nilpotent}
Let $p$ be odd and let $V$ be a unitary representation of
$G=\operatorname{SL}_2(\mathbb Z/p^2)$. Let $K$ be the kernel
of reduction modulo $p$, and suppose $v$ is fixed by every upper
unipotent and diagonal torus. Identify the characters of
$K=1+p\mathfrak{sl}_2(\mathbb F_p)$ by
\[
 \lambda_A(1+pM)=\zeta_p^{\operatorname{Tr}(AM)},
 \qquad A\in\mathfrak{sl}_2(\mathbb F_p).
\]
The cyclic module $V(v)$ is the orthogonal sum of a $K$-trivial
module and at most two further modules, one for each nonzero
nilpotent conjugacy orbit. Each latter module has dimension
$M=(p^2-1)/2$, with each character of its orbit occurring once.
The central torus at $-1$ is the identity on $V(v)$.

Every $G$-module with nontrivial $K$-action has dimension at least
$M$. At equality it is one irreducible nilpotent-orbit module,
with each character once. On any such one-orbit module, each
occurring primitive order-$p^2$ eigenvalue of the upper unipotent
$T=\rho(U(1))$ has multiplicity one.
\end{lemma}
\begin{proof}
Write $A=\bigl(\begin{smallmatrix}x&y\\z&-x\end{smallmatrix}\bigr)$.
The upper and diagonal elements of $K$ fix $v$, so its nonzero
character components satisfy $z=0$ and $2x=0$; their indices
are $A=yE_{12}$. The zero component generates a $K$-trivial
module. For $A=cE_{12}\ne0$, its stabilizer modulo $p$ is
\[
 \left\{\begin{pmatrix}\epsilon&t\\0&\epsilon\end{pmatrix}:
                   \epsilon\in\{1,-1\},\ t\in\mathbb F_p\right\}.
\]
Its inverse image is generated by $K$, upper unipotents and the
central torus at $-1$. They act by scalars on $v_A$, so its
translates give one orthogonal line for each character of an
orbit of size $|\,\operatorname{SL}_2(\mathbb F_p)\,|/(2p)=M$.
Diagonal tori carry the components $v_{cE_{12}}$ through each
nonzero square class. This proves the cyclic decomposition.
Centrality and the fact that $-1$ fixes $v$ prove its central sign.

A nonzero traceless matrix at odd $p$ is nonscalar, and its
commuting matrices are $aI+bA$. There are at most $2p$ of
determinant one, since for each $b$ the equation
$a^2+b^2\det A=1$ has at most two solutions. A nonzero character
therefore brings at least $M$ distinct characters into a
$G$-module. For a nonnilpotent traceless matrix the centralizer
has $p-1$ or $p+1$ elements, by the split or nonsplit quadratic
norm equation, giving the larger orbit sizes $p(p+1)$ or
$p(p-1)$. Equality can consequently consist only of one
nilpotent orbit, with one-dimensional character spaces.
Such a module is irreducible because every nonzero submodule
must contain that full character orbit.

Finally, for a permitted nonzero lower-left entry $z$, the
nilpotent character indices are
\[
           \begin{pmatrix}x&-x^2/z\\z&-x\end{pmatrix},
           \qquad x\in\mathbb F_p.
\]
The upper unipotent cycles these $p$ lines, while $T^p$ acts
as $\zeta_p^z$. Its $p$ eigenvalues solving
$\lambda^p=\zeta_p^z$ therefore occur once each.
\end{proof}

\begin{theorem}[The canonical classes in a wild flat block]
\label{paper:wild-flat}
Suppose an odd flat block in \eqref{paper:flat-setup} has the wild
phase law: $p$ inactive labels, no twist of exact order $p$, and
multiplicity two at each primitive order-$p^2$ phase in one square
class. Write $I$ for its inactive set and $A=B\setminus I$.
Then the central local torus $H_{-1,p}$ fixes exactly $Z$.
The full local $p$-inertia orbits on $B$ are exactly
\[
                      \{Z\},\qquad I\setminus\{Z\},\qquad A,
\]
of sizes $1,p-1,p(p-1)$.

Let $\mathcal Q$ consist of the set $I$ and the $p-1$ wild-inertia
orbits in $A$. Each of its $p$ sets has $p$ labels.
Define the global spectral idempotent
\[
 \mathcal F_p=\{X:\operatorname{ord}(\theta_{X,p})\mid p\},
 \qquad E_p=\sum_{X\in\mathcal F_p}e_X.
\]
Multiplication by $E_p$ preserves every full $p$-block and restricts
on the present block to
\begin{equation}\label{paper:wild-averaging}
                  (N_{E_p})_{B,B}
                   =\bigoplus_{Q\in\mathcal Q}\frac{J_p}{p}.
\end{equation}
Consequently, for $R_Q=\sum_{X\in Q}[X]$, every $Q,Q'\in\mathcal Q$
and $X\in Q'$ give actual Grothendieck objects
\[
              R_QX^*=\frac{R_QR_{Q'}^*}{p}\in R_{\ge0}.
\]
The same single idempotent $E_p$ works on every wild flat $p$-block.
On any full $p$-block all of whose actual primary twists divide
$p$, multiplication by $E_p$ is the identity.
\end{theorem}
\begin{proof}
Put $M=(p^2-1)/2$. The cyclic module $\mathcal V_Z$ has dimension
$M+1$. Lemma~\ref{paper:cyclic-nilpotent} and the single permitted
primitive square class show that it consists of one nilpotent
orbit module $V_+$ of dimension $M$ and a $K$-trivial line.
The latter line is generated by a Borel-fixed vector; since
upper and lower unipotents generate the group, it is the trivial
representation. The nilpotent module contributes one copy of
each occurring primitive phase. The full block has multiplicity
two, so its $M$-dimensional orthogonal complement is not
$K$-trivial. The minimum orbit bound makes this complement one
nilpotent-orbit module $V'$ with each character occurring once.
The permitted square class forces it to have the same orbit
as $V_+$.

The central torus is positive on $\mathcal V_Z$, whose
dimension is $(p^2+1)/2$. Corollary~\ref{paper:flat-charge}
makes this its entire positive eigenspace, so its action on
the complementary module $V'$ is negative. The same corollary
has already identified $Z$ as its only fixed label.

The square action of all local tori is transitive on the
$p(p-1)/2$ primitive phases. The torus at $-1$ interchanges the
two labels at each such phase, so all of $A$ is one orbit.
Wild inertia fixes $I$ pointwise: it has order $p$ on this block,
permutes the $p$ inactive labels, and fixes $Z$.
For a tame torus whose residue $s$ has $s^2\ne1$, no nonzero
nilpotent character is fixed by conjugation. Its trace on
$V_+$ and $V'$ is therefore zero, and its trace on $B$ is one.
The tame torus at $-1$ also has trace one. Thus every nonidentity
tame torus acts freely on $I\setminus\{Z\}$, making this set one
orbit of size $p-1$. The primitive orbit is regular for the full
local torus group of order $p(p-1)$, so its wild orbits have
size $p$.

Let $P_U,P_L,P_H$ be the orthogonal projections onto the fixed
spaces of the upper, lower, and diagonal wild subgroups, and
let $P_K$ project onto the $K$-fixed line. The character support
is zero together with one nilpotent orbit, twice. On such a
character, diagonal invariance means $x=0$; nilpotence then gives
$yz=0$, so it means upper or lower invariance. Both at once mean
the zero character. Consequently
\begin{equation}\label{paper:wild-projection-identity}
                         P_H=P_U+P_L-P_K.
\end{equation}
Here $P_U=P_I$ and $P_L=W^*P_IW$.
The $K$-fixed unit vector is supported on $I$ and constant on
$I\setminus\{Z\}$. Write its coordinates there as $\alpha$ at
$Z$ and $\beta$ elsewhere. It is fixed by $W$, and the flat
$Z$-row gives $\alpha=(\alpha+(p-1)\beta)/p$, hence
$\alpha=\beta$. Thus $P_K=J_I/p$.
The diagonal wild subgroup acts as the identity on $I$ and by
$p$-cycles on $A$. Substitute its permutation average into
\eqref{paper:wild-projection-identity}: this gives
\[
                   W^*P_IW=\bigoplus_{Q\in\mathcal Q}J_p/p.
\]

Proposition~\ref{paper:principal-phase-support} gives globally
\[
       N_{E_p}=s^*P_{\mathcal F_p}s
              =W_p^*P_{\mathcal F_p}W_p.
\]
This preserves every full $p$-block. On the present block
$P_{\mathcal F_p}=P_I$, proving
\eqref{paper:wild-averaging}. The actual product identities follow
from Proposition~\ref{paper:averaging-products}. If every primary
twist on a full block divides $p$, the same formula restricts to
the identity there.
\end{proof}

\begin{corollary}[Complementary components in the wild law]
\label{paper:wild-complementary}
In Theorem~\ref{paper:wild-flat}, the full Galois orbits of
nonprincipal complementary components meeting $B$ have exactly
one of the following intersection patterns:
\begin{enumerate}[label=(\roman*)]
\item one intersection $B\setminus\{Z\}$ of size $p^2-1$;
\item two intersections $I\setminus\{Z\}$ and $A$, of sizes
$p-1$ and $p(p-1)$.
\end{enumerate}
In the second case the minimal-orbit partition from
Theorem~\ref{paper:flat-orbits} is precisely the canonical
partition $\mathcal Q$ of Theorem~\ref{paper:wild-flat}.
If $J=J_{p^2}$ and $A_{\mathcal Q}=\bigoplus_{Q\in\mathcal Q}J_p$,
the nonprincipal integral projections are respectively
\[
 p^2I-J,
 \qquad\text{or}\qquad
 pA_{\mathcal Q}-J\quad\text{and}\quad p^2I-pA_{\mathcal Q}.
\]
\end{corollary}
\begin{proof}
Each fully Galois-stable union has an intersection with $B$
stable under local $p$-inertia. The principal intersection is
$\{Z\}$; the two remaining local orbits are exactly
$I\setminus\{Z\}$ and $A$. The component-orbit intersections
partition their union, giving precisely the two alternatives.
The projection formulas follow from
$p^2W^*P_{\{Z\}}W=J$ and
$p^2W^*P_IW=pA_{\mathcal Q}$.
\end{proof}

\subsection{Fusion support and crossings between full blocks}

\begin{corollary}[Tensor collisions across the tame and wild blocks]
\label{paper:tame-wild-products}
Fix an odd prime $p$ and let
$F_p=\{X:\operatorname{ord}(\theta_{X,p})\mid p\}$,
$E_p=\sum_{X\in F_p}e_X$. Use the canonical $p$-tuple classes
in the wild flat blocks of Theorem~\ref{paper:wild-flat}.
For any simple $Y$ in a full $p$-block all of whose actual primary
twists have order dividing $p$, and any such class $Q$, one has
\[
        XY^*=X'Y^*=\frac1p\left(\sum_{T\in Q}T\right)Y^*
        \in R_{\ge0}\qquad(X,X'\in Q).
\]
The same common idempotent simultaneously supplies averaging products
between classes in different wild flat blocks. Thus those products
are coefficientwise divisible by $p$ after multiplication by any
other averaging object.
\end{corollary}
\begin{proof}
The canonical projector restricts to $J_p/p$ on $Q$, and to the
identity on the coordinate space of the stated tame block.
Treat each of its labels as a class of size one in
Proposition~\ref{paper:averaging}. Equivalently $E_pY^*=Y^*$ and
\[
 XY^*=E_pX\,Y^*=\frac1pR_QY^*.
\]
The last assertion is the common-size product identity.
\end{proof}

\begin{proposition}[A positive witness in a wild block]
\label{paper:wild-support}
Let $p$ be odd and let $B$ be a wild singleton-flat block as in
Theorem~\ref{paper:wild-flat}, and let $E_p$ be the global cutoff
there. For any one of its canonical classes $Q$, put
\[
 R_Q=\sum_{X\in Q}[X],\qquad
 H_Q=\frac{R_QR_Q^*}{p},\qquad h_Y=[Y]H_Q.
\]
Assume the ambient integral modular category is perfect and
fusion-simple. Then $H_Q$ has a nonunit simple constituent $Y$
such that
\[
 p\nmid d_Yh_Y,\qquad
 Y\text{ is fixed by wild }p\text{-inertia},\qquad
 \operatorname{ord}(\theta_{Y,p})\mid p,\qquad E_p[Y]\ne[Y],
\]
and hence the full $p$-block containing $Y$ has a primary twist whose
order is greater than $p$. More precisely the level-$p$ congruence
kernel acts nontrivially on the cyclic local module generated by
$\mathbf e_Y$: its upper wild unipotent fixes this vector, while
its lower wild unipotent does not. The $p$-part of the conductor of $K_Y$
therefore divides $p$, although its full block is not wholly tame.
The label $Y$ sends every member of $Q$ to a product with at least
one constituent in $Q$; the total multiplicity into $Q$ is the same
integer $h_Y$, prime to $p$, for every member.

If $Q$ is a primitive wild-inertia class, the compressed matrix
$(N_Y)_{Q,Q}$ is circulant and invertible modulo $p$.

In the ordered exact family, this witness has degree
\[
                    d_Y=q^a r^b,\qquad 1\le a\le3,
                    \quad1\le b\le2.
\]
In particular its full $p$-block is outside the homogeneous
height-two blocks classified above.
\end{proposition}
\begin{proof}
Every canonical class $Q$ is preserved by full wild
$p$-inertia: the primitive classes are its $p$-orbits, the
higher wild kernel is trivial on the level-$p^2$ block, and
the inactive class is preserved as well. By
\eqref{paper:galois-fusion}, this is a finite $p$-group of
dimension-preserving based-ring automorphisms. Apply
Lemma~\ref{paper:averaging-witness} with $k=p$ and the canonical
cutoff $E_p$. It supplies a nonunit $Y$ with $p\nmid d_Yh_Y$,
fixed by wild inertia and not fixed by $E_p$.

Wild-fixedness means the wild Galois subgroup acts trivially on
$K_Y$, proving the conductor assertion. The square action of the
torus at $1+p$ gives
$\theta_{Y,p}^{(1+p)^2-1}=1$; because $p$ is odd, this implies
$\theta_{Y,p}^p=1$. The operator $P_{\mathcal F_p}$ projects onto
the fixed space of the upper wild unipotent $U(p)$, and
$N_{E_p}=W_p^*P_{\mathcal F_p}W_p$ projects onto the fixed space
of its lower conjugate. Thus the former fixes $\mathbf e_Y$ and
the latter does not. The cyclic module has nontrivial congruence
kernel action, and the full block cannot be wholly tame.

The return assertion is the last part of
Lemma~\ref{paper:averaging-witness}.

For a primitive class, a generator
of wild inertia acts as a $p$-cycle on $Q$. It fixes $Y$, so fusion
covariance makes $(N_Y)_{Q,Q}$ commute with that cycle, hence be
circulant. Over $\mathbb F_p$ its representing polynomial $f(t)$
has $f(1)=h_Y\ne0$ in
$\mathbb F_p[t]/(t^p-1)=\mathbb F_p[t]/((t-1)^p)$.
It is therefore invertible.

Finally square divisibility gives $d_Y=q^ar^b$ in the exact family;
the exclusion of nonunit prime-power degrees forces both exponents
to be positive and the global exponents give $a\le3$, $b\le2$.
Its $p$-height is zero, so it cannot lie in a homogeneous height-two
block.
\end{proof}

\begin{proposition}[Wild-fixed fusion and the primary phase classes]
\label{paper:wild-transfer}
Work in an ordinary integral modular category of odd square dimension,
and fix an odd prime $p$. Let $B$ and $C$ be wild singleton-flat
full $p$-blocks as in Theorem~\ref{paper:wild-flat}; they may coincide.
Let $\mathcal Q_B$ and $\mathcal Q_C$ be their canonical $p$-class
partitions. For a canonical class $Q$, write
\[
 R_Q=\sum_{X\in Q}[X],\qquad
 \lambda_Q=\theta_{X,p}^{p}\quad(X\in Q).
\]
This scalar is independent of $X\in Q$. It equals $1$ precisely
for the inactive class. Within a wild block every other occurring
value belongs to exactly two classes, interchanged by $H_{-1,p}$.

Suppose the simple label $Y$ is fixed by wild $p$-inertia.
If $\lambda_Q\ne\lambda_{Q'}$, for $Q\in\mathcal Q_B$ and
$Q'\in\mathcal Q_C$, then the entire fusion rectangle
$(N_Y)_{Q,Q'}$ is a constant nonnegative integer matrix.
In particular
\begin{equation}\label{paper:off-phase-transfer}
 p\mid [Y]\frac{R_QR_{Q'}^*}{p}.
\end{equation}
Equivalently, after reduction modulo $p$, the compressed
class-transfer matrices have nonzero entries only between
classes with the same $p$-th primary twist power.
On one wild block their only possible nonzero blocks are the
inactive singleton and the $(p-1)/2$ paired primitive classes.

If $B$ and $C$ have opposite primitive square classes, every
canonical-class rectangle $(N_Y)_{Q,Q'}$ is constant, including
the rectangle between their inactive classes. Thus
\eqref{paper:off-phase-transfer} holds for all their class pairs.
\end{proposition}

\begin{proof}
Use one common local representation and abbreviate the upper,
lower and diagonal wild operators on a block by
\[
 U=\rho\begin{pmatrix}1&p\\0&1\end{pmatrix},\qquad
 L=\rho\begin{pmatrix}1&0\\p&1\end{pmatrix},\qquad
 H=\rho\begin{pmatrix}1+p&0\\0&(1+p)^{-1}\end{pmatrix}.
\]
On each of the two blocks these operators factor through the
abelian level-$p$ congruence kernel in
$\operatorname{SL}_2(\mathbb Z/p^2)$; there
$(1+p)^{-1}=1-p$.
By the wild-block decomposition, its character support consists
of the zero character and one nonzero nilpotent orbit, the latter
with multiplicity two. For
$A=\left(\begin{smallmatrix}x&y\\z&-x\end{smallmatrix}\right)$,
the simultaneous eigenvalues of $U,L,H$ are
$\zeta_p^z,\zeta_p^y,\zeta_p^{2x}$, respectively, and
\[
                             x^2+yz=0.
\]
The cutoff projection on each wild block is the fixed-space
projection of $L$; denote these restrictions by $E_B,E_C$.
Consequently, on $\ker E_B$ or $\ker E_C$ one has $y\ne0$,
and the upper-wild eigenvalue is determined by the lower and
diagonal eigenvalues through
\[
                             z=-x^2/y.
\]
It follows that every map intertwining $L,H$ also intertwines
$U$ after restriction to these two kernels.

Now put $M=(N_Y)_{B,C}$. Every fusion multiplication commutes
with the lower local unipotents: the latter are diagonalized by
the modular matrix, since
$s^*T_p^p s=W_p^*T_p^pW_p$ is $L$ or $L^{-1}$ according to
the Weyl-sign convention. Invariance of $Y$ under wild
inertia and fusion covariance likewise imply that $N_Y$ commutes
with $H$. The full-block restrictions therefore give
\[
 ME_C=E_BM,\qquad
 U_BM(1-E_C)=M U_C(1-E_C).
\]
For each canonical class $Q$, the range of $1-E_B$ in its
coordinate space is the zero-sum subspace $\mathbb C^Q_0$.
Moreover $U_B=T_p^p$ acts there by the scalar $\lambda_Q$.
Thus if $\lambda_Q\ne\lambda_{Q'}$, the component of $M$
from $\mathbb C^{Q'}_0$ to $\mathbb C^Q_0$ is zero.
The identity $ME_C=E_BM$ eliminates the two mixed components
between constant and zero-sum vectors. The whole rectangle
therefore maps only the constant line to the constant line,
so all its entries are equal. They are nonnegative integers
because they are fusion coefficients. The common row sum is
$[Y](R_QR_{Q'}^*/p)$ by Proposition~\ref{paper:averaging},
and is $p$ times that entry, proving the divisibility.

For opposite primitive square classes, the nonzero nilpotent
orbits differ. At $y\ne0$, the pair $(x,y)$ determines the
entire nilpotent matrix $A$, so the two blocks have no common
$(L,H)$-character off their cutoff ranges. Hence
$(1-E_B)M(1-E_C)=0$. The same constant-line argument applies
to every class rectangle, including the two inactive classes.

Finally the primitive phases in a wild block occur twice, in
one square class, and wild inertia runs through their $p$
lifts with a fixed nonzero $p$-th power. The local central
inversion exchanges their two labels at every phase. Thus
each nontrivial value $\lambda_Q$ occurs on exactly two
canonical classes exchanged by $H_{-1,p}$, as stated.
\end{proof}

\begin{remark}[What the divisibility controls]
For the positive witness of Proposition~\ref{paper:wild-support},
the return coefficient to its chosen class is prime to $p$,
whereas its transfer coefficient to every class with a different
$p$-th twist power is divisible by $p$. At the level of individual
fusion products, every such off-phase projection is independent
of the chosen input label within the class. This is a restriction
on actual fusion multiplicities. It does not assert closure of
the wild blocks, or of their phase classes, under constituents
of tensor products.
\end{remark}


\paragraph{Crossings between full blocks.}

\begin{proposition}[Local cyclic modules across a flat crossing]
\label{paper:flat-crossing}
Let $B$ be a singleton-flat block at a prime $p$, and write $X$
for its distinguished label, so that
\[
 B\cap P_{p'}=\{X\},\qquad
 W_p\mathbf e_X=(\varepsilon/p)\mathbf1_B.
\]
For $Y\in P_{p'}$ with $S_{XY}\ne0$, the cyclic local $p$-module
$\mathcal V_X$ is a quotient of $\mathcal V_Y$.
In particular the full $p$-block containing $Y$ is not a singleton.

If $p$ is odd and $Y$ is also distinguished in a singleton-flat
block, then $\mathcal V_X\simeq\mathcal V_Y$. The two blocks have
the same one of the three quadratic phase laws. In the wild law
they have the same primitive square class.
More generally, if $B$ has the wild law, the full block of $Y$
contains at least $(p+1)/2$ inactive labels and at least one label
at each of the $p(p-1)/2$ primitive phases occurring on $B$.
\end{proposition}
\begin{proof}
Since $W_{p'}$ preserves $P_{p'}$ and the actual $p$-twist,
the singleton intersection makes the crossing zero when
$Y$ is $p$-active. Hence a nonzero crossing has
$Y\in\mathcal I_p^0$, and the cyclic quotient assertion,
including exclusion of a singleton block, follows from
Theorem~\ref{paper:isolated-fiber}.

If both labels are distinguished, symmetry gives the reverse
surjective intertwiner, so their finite-dimensional cyclic
modules are isomorphic. Their dimensions $p$, $p+1$ and
$(p^2+1)/2$ distinguish the three odd laws; the primitive
$T$-spectrum distinguishes the wild square classes.
Eigenvalue multiplicities in a quotient do not exceed those
in the source. In the wild case,
Lemma~\ref{paper:cyclic-nilpotent} gives $(p+1)/2$ inactive
eigenvectors and one at each occurring primitive phase,
proving the final assertion.
\end{proof}

\begin{proposition}[A flat block misses the maximal-height stratum]
\label{paper:flat-top-support}
Let $B$ be a singleton-flat full $p$-block, with distinguished
$Z\in B\cap P_{p'}$. If $p\nmid f_Y$, then
\[
                         S_{XY}=0\qquad(X\in B).
\]
Moreover tensoring by $Y^*$ identifies all labels of $B$:
\[
                XY^*=X'Y^*=\frac1{p^2}R_BY^*
                         \qquad(X,X'\in B),
                \qquad R_B=\sum_{X\in B}[X].
\]
The right-hand side is an actual positive integral fusion object.
In the exact family this applies to every simple of $p$-height three.
\end{proposition}
\begin{proof}
The vanishing is \eqref{paper:isolated-top-vanishing}.
The common complementary principal idempotent restricts
to $J_{p^2}/p^2$ on $B$ and to the identity on the
singleton full $p$-block $\{Y\}$, since codegree support
puts $Y\in P_{p'}$. Apply
Proposition~\ref{paper:averaging-products} to the two
classes $B$ and $\{Y\}$. Its individual-product identity
gives the stated positive integral tensor collisions.
\end{proof}

\begin{theorem}[Whole rectangles adjacent to a wild flat block]
\label{paper:wild-rectangles}\label{paper:wild-tame-orthogonality}
Assume the odd-dimensional singleton-flat setup at $p$.
Let $B$ be a wild flat block, let $I_B$ be its inactive set, and
let $C$ be another full $p$-block.
\begin{enumerate}[label=(\roman*)]
\item If all actual $p$-primary twists on $C$ have order dividing
$p$, then
\[
                              S_{B,C}=0.
\]
In particular this covers both tame singleton-flat laws and every
singleton block of $p$-unit codegree.
\item If $C$ is wild singleton-flat with the opposite primitive
square class, then $S_{B,C}=0$.
\item If $C$ is wild singleton-flat with the same primitive
square class, put $M=(p^2-1)/2$. The only possible ranks of
$S_{B,C}$ are
\[
                         0,\quad M,\quad M+1,\quad p^2.
\]
Writing $Z_B,Z_C$ for the distinguished labels, the first two
ranks occur when $S_{Z_BZ_C}=0$ and the latter two occur when
$S_{Z_BZ_C}\ne0$.
\end{enumerate}
\end{theorem}
\begin{proof}
Put $F=P_BW_{p'}P_C$. This is an intertwiner for the full local
$p$-group, and
\[
                         D^{-1/2}S_{B,C}=W_BF.
\]
By Theorem~\ref{paper:wild-flat}, the wild module is the direct
sum of its trivial line, generated by
$u_B=p^{-1/2}\mathbf1_{I_B}$, and the two irreducible modules
on one nonzero nilpotent congruence-character orbit, distinguished
by central signs $+$ and $-$. Both have dimension $M$.

Work in a common local group $\operatorname{SL}_2(\mathbb Z/p^n)$.
If $T_p^p=1$ on $C$, then $U(p)$ acts trivially there. Its normal
closure is the kernel of reduction modulo $p$: it contains upper
and lower unipotents with parameters divisible by $p$, and Gaussian
elimination reduces the remaining assertion to diagonal tori
$H(a)$ with $a\equiv1\pmod p$. The identity
\[
 U(a-1)L(1)=L(a^{-1})H(a)U((a-1)/a)
\]
puts those tori in the same normal closure. Thus $C$ factors
through level $p$, and the image of $F$ lies in the
congruence-fixed line $\mathbb C u_B$.
For $Y\in C\cap P_{p'}$, the column $F\mathbf e_Y$ is supported
on $B\cap P_{p'}=\{Z_B\}$, which has zero intersection with
$\mathbb C u_B$. For $Y\in C\setminus P_{p'}$, complementary
block invariance makes the $Z_B$ coordinate of that column zero,
again impossible for a nonzero multiple of $u_B$. Every column
therefore vanishes, proving (i).

For opposite wild square classes, the only common congruence
character is zero, so the image of $F$ again lies in the trivial
line. The same column argument gives $F=0$, proving (ii).

For equal wild square classes the three irreducible summands
within each block are pairwise inequivalent, each with multiplicity
one. The distinguished vector
generates the sum of the trivial and central-positive summands.
Complementary-principal support gives
$F\mathbf e_{Z_C}=c\mathbf e_{Z_B}$, with $c$ a nonzero scalar
multiple of $S_{Z_BZ_C}$. If $c=0$, the intertwiner kills that
entire cyclic module, leaving only the central-negative summand.
If $c\ne0$, it is an isomorphism on the cyclic module of dimension
$M+1$. Its restriction on the remaining irreducible summand is
either zero or an isomorphism. Multiplication by $W_B$ preserves
rank, proving the list.
\end{proof}

\begin{corollary}[A flat witness forces the opposite height]
\label{paper:flat-witness-height}
In the ordered exact family, suppose the full $p$-block of one
compulsory witness is singleton-flat. Every opposite compulsory
witness then has $p$-height two and trivial actual $p$-twist;
its full $p$-block is homogeneous and nonprincipal. If the first block is odd and wild,
the opposite block contains actual twists of order $p^2$ and has
at least $(p^2+1)/2$ labels. If it too is singleton-flat, both
blocks have the same wild square class.
\end{corollary}
\begin{proof}
Both compulsory witnesses belong to $P_{p'}$, since each belongs
to a principal component at one of the other primes. Their
crossing is nonzero by its residue at those primes. The witness
in the singleton-flat block is its distinguished label and has
height two. An opposite witness of height three would have
$p$-unit codegree and a singleton full $p$-block, contrary to
Proposition~\ref{paper:flat-crossing}. Its height is therefore two. The intertwiner
$P_BW_{p'}$ sends its coordinate vector to a nonzero multiple
of the distinguished flat vector. It commutes with the actual
$p$-twist; that distinguished vector is inactive, so the
opposite witness is inactive as well.
Proposition~\ref{paper:family:homogeneous-crossing} gives homogeneity
and nonprincipality. The remaining assertions follow from the
same cyclic-module quotient.
\end{proof}

\subsection{Four-label dyadic local matrices}

\begin{theorem}[Complete local matrix classification at two]
\label{paper:dyadic-flat}
Suppose $p=2$ in \eqref{paper:flat-setup}, and put
$\Sigma=\varepsilon W$ and $\alpha=\varepsilon\gamma_2^3$.
After relabeling the three labels other than $Z$, one has
\begin{equation}\label{paper:fourier-four}
 \Sigma=\Sigma_\eta:=\frac12
 \begin{pmatrix}
 1&1&1&1\\
 1&1&-1&-1\\
 1&-1&\eta&-\eta\\
 1&-1&-\eta&\eta
 \end{pmatrix}.
\end{equation}
The complete list of actual primary twists and parameters is
\[
\begin{array}{c|c|c}
 \eta&\operatorname{diag}(T)&\alpha\\ \hline
 1&(1,-1,1,1)&1\\
 1&(1,-1,-1,-1)&-1\\
 -1&(1,1,i,-i)&1\\
 -1&(1,-1,i,i)&i\\
 -1&(1,-1,-i,-i)&-i\\
 \pm i&(1,-1,t,t),\quad t^2=-\eta&t
\end{array}
\]
In the last row $t$ is primitive of order eight.
Here $H_{-1,2}=I$ in the first five rows and is one transposition
in the last row.  The latter case occurs exactly when
$\gamma_2$ has order eight.

In particular every entry of $2W$ is a root of unity, so this
local block has no foreign denominators.  The auxiliary Verlinde
multiplication defined by $\Sigma$ is the group-ring multiplication
of $C_2^2$ in the first five rows and of $C_4$ in the last row.
\end{theorem}
\begin{proof}
The dyadic conductor bound makes every torus congruent to one
modulo four fix each label of $B$.  Galois covariance therefore
gives $W\in\operatorname{Mat}_4(\mathbb Q(i))$.
The square action of the torus at $5$ gives primary twist order
dividing eight; the anomaly has the same bound.
The matrix $H=H_{-1,2}=W^2$ is an involutive permutation fixing
$Z$, hence is either the identity or one transposition.
Moreover
\begin{equation}\label{paper:dyadic-relations}
 \Sigma\mathbf e_Z=\tfrac12\mathbf1,
 \quad\Sigma\mathbf1=2\mathbf e_Z,
 \quad\overline\Sigma=H\Sigma,
 \quad\Sigma T\Sigma=\alpha T^{-1}\Sigma T^{-1}.
\end{equation}
The vacuum entry of the last identity gives
$\sum_X\theta_{X,2}=2\alpha$.

Comparison of coefficients in $\mathbb Q(\zeta_8)$ says that
this four-term multiset is two copies of $\alpha$ and one
opposite pair.  Since it contains $1$, it has one of the forms
\begin{equation}\label{paper:dyadic-multisets}
 \{1,1,t,-t\},\quad\alpha=1,
 \qquad\text{or}\qquad
 \{1,-1,t,t\},\quad\alpha=t,
 \qquad t\in\mu_8.
\end{equation}

First suppose $H=I$.  Then $\Sigma$ is rational, real,
symmetric and orthogonal.  In either form in
\eqref{paper:dyadic-multisets}, when $t$ is nonreal the vacuum-row
equation in \eqref{paper:dyadic-relations}, together with rationality,
row sums and orthogonality, forces the row at the second label to be
$(1,1,-1,-1)/2$.  Row sums and orthogonality now give
$\Sigma=\Sigma_\eta$ with $\eta=\pm1$.
Substitution in the same modular relation gives $t^2=\eta$.
Thus $t=\pm i$ and $\eta=-1$, giving the three nonreal cases
in the table.  The real endpoints are the multisets with
one negative sign or with three negative signs.
Here $T^2=I$ and $\operatorname{tr}(T)=2\alpha$, so taking
traces of the modular relation gives $\operatorname{tr}(\Sigma)=2$.
A real symmetric orthogonal matrix with this trace is a rank-one
reflection $\Sigma=I-2vv^T$, $\|v\|=1$.
Its prescribed first row forces
$v=\pm(1,-1,-1,-1)^T/2$, hence $\Sigma=\Sigma_1$.
These are the first two rows.

Now suppose $H$ interchanges the last two labels.  Since $T$
commutes with $H$, those labels have equal twists, so
\eqref{paper:dyadic-multisets} gives $T=(1,-1,t,t)$ and
$\alpha=t$.  Symmetry, conjugation in
\eqref{paper:dyadic-relations}, and the row sums give
\[
 \Sigma=\begin{pmatrix}
 1/2&1/2&1/2&1/2\\
 1/2&a&b&b\\
 1/2&b&c&\overline c\\
 1/2&b&\overline c&c
 \end{pmatrix},\qquad a,b\in\mathbb Q.
\]
The second row has $a+2b=-1/2$ and $a^2+2b^2=3/4$,
so $b=-1/2$ or $1/6$.  On the antisymmetric line, $\Sigma$
has eigenvalue $\eta=c-\overline c=\pm i$.
The modular relation on that line gives $t^2=-\eta$.
In particular $t\ne-1$.  Its second-row vacuum equation is
$(1+2b)(1+t)=0$, forcing $b=-1/2$.  Thus $a=1/2$ and
$c+\overline c=0$, proving $\Sigma=\Sigma_\eta$ and the last row.

Every listed pair satisfies the displayed modular relations.
Since $\alpha=\varepsilon\gamma_2^3$, its order is eight
exactly in the last row.  Finally the rows of $2\Sigma_\eta$
are the characters of $C_2^2$ when $\eta=\pm1$ and of $C_4$
when $\eta=\pm i$.  Their character orthogonality gives the
stated auxiliary Verlinde multiplication.
\end{proof}

\begin{proposition}[Four-label dyadic blocks with proper principal support]
\label{paper:dyadic-four}\label{paper:dyadic-three}
Under the standing hypotheses, let $B$ be a homogeneous dyadic
defect-two block with four labels,
and suppose $I=B\cap P_{2'}$ is proper.  Every entry of $2W_2|_B$
is a root of unity.  Its scaled degrees have one of the forms
\[
           (a,a,b,b),\qquad a/b\in\mathbb Z_{>0}\text{ odd},
           \qquad\text{or}\qquad (3b,b,b,b).
\]
If $|I|=1$, all four degrees are equal.  If $|I|=2$, the second
form is impossible; in the unequal case $I$ is one of the two
pairs of Theorem~\ref{paper:dyadic-projection}, and its two scaled
degrees equal $b$.  A transversal of the two pairs forces equal
degrees.  If $|I|=3$, the unique label outside $I$ has scaled
degree $a$, or $3b$, respectively.

The $3:1$ form has actual primary twists $(-1,1,1,1)$, with
the degree-$3b$ label first, and equal packet coefficients
$\pm1$ on the three other labels.  In the ordered target family
it requires $q=3$.  A nontrivial local central torus forces all
four degrees to be equal.
\end{proposition}
\begin{proof}
Write $W=W_2|_B$.  We first prove flatness.  The singleton case
is Theorems~\ref{paper:singleton} and~\ref{paper:dyadic-flat}.

Suppose $|I|=2$.  The integral projection theorem partitions
$B$ into two pairs of equal scaled degrees; its range consists
of the vectors constant on each pair.  As this range is
$W^*\mathbb C^I$, each row of $W$ indexed by $I$ is constant
on both pairs.  If $I$ is a transversal, order the labels as
the transversal followed by their partners.  Symmetry and
unitarity give
\[
       W=\begin{pmatrix}A&A\\A&-A\end{pmatrix},
       \qquad A=A^{\mathsf T},\qquad 2AA^*=I_2.
\]
Let the two scaled degrees be $a,b$.  The packet norm says
$a^2c_1^2+b^2c_2^2=(a^2+b^2)/2$, so precisely one packet
coefficient is nonzero.  After interchanging the pairs,
$W\mathbf e_1=u/(2ac)$, with $c\in\mathbb Z\setminus\{0\}$.
Consequently
\[
       A=\frac1{2c}\begin{pmatrix}1&k\\k&-1\end{pmatrix},
       \qquad k=b/a>0,\qquad k^2=2c^2-1.
\]
The rational number $k$ is therefore an odd integer, and $c$
is odd.  Write $T=\operatorname{diag}(t_1,t_2,t_3,t_4)$.
The $(1,2)$ entry of $WTW=\gamma_2^3T^{-1}WT^{-1}$ gives
\[
                 \frac{|t_1-t_2+t_3-t_4|}{2|c|}=1.
\]
Thus $|c|\le2$, whence $|c|=k=1$.  This proves flatness and
equality of all degrees in the transversal case.

If $I$ is one whole pair, order it first.  Symmetry,
orthogonality and the constant-pair rows give
\[
 W=\begin{pmatrix}
 \alpha&\alpha&\beta&\beta\\
 \alpha&\alpha&-\beta&-\beta\\
 \beta&-\beta&\zeta&-\zeta\\
 \beta&-\beta&-\zeta&\zeta
 \end{pmatrix},
 \qquad |\alpha|=|\beta|=|\zeta|=\tfrac12.
\]
The central torus preserves $I$, and $\overline W=H_{-1,2}W$
makes $\alpha$ real.  The packet equations, for
$u=(b,b,a,a)$, give $c_1=\alpha+\beta a/b$ and
$c_2=\alpha-\beta a/b$.  Their integrality makes $\beta$
real and $a/b$ an odd positive integer.  The central-torus
relation makes $\zeta$ real or purely imaginary.  Thus all
entries of $2W$ are roots of unity.

For $|I|=3$, Theorem~\ref{paper:dyadic-projection} supplies a
rational signed flat column at the unique label $X\notin I$,
which is fixed by every local torus.  Write
$W\mathbf e_X=\delta s/2$, where $s_X=1$, $s_i=\pm1$ and
$\delta=\pm1$.  The Weyl--torus relation shows that $s$ is
torus fixed, so $D_s=\operatorname{diag}(s)$ commutes with all
permutation tori.  The pair
\[
        \Sigma=\delta D_sWD_s,
        \qquad T'=\theta_{X,2}^{-1}T
\]
has flat distinguished column, $\Sigma\mathbf1=2\mathbf e_X$,
entries in $\mathbb Q(i)$, twists in $\mu_8$, and modular scalar
$\delta\gamma_2^3\theta_{X,2}^{-3}$.  These are precisely the
matrix hypotheses used in the proof of
Theorem~\ref{paper:dyadic-flat}.  Its table proves flatness here
as well.  This shift is only a matrix normalization; the
packet is still supported at the actual twist one.

We finish with the degree restrictions, using one common
calculation for the flat cases.  Corollary~\ref{paper:lw:inactive-fixed}
supplies an inactive packet label fixed by the local tori.
Its now flat rational column permits the same sign normalization
without shifting $T$.  Thus the actual twist table of
Theorem~\ref{paper:dyadic-flat} applies.  Put $v=D_su$ and absorb
the common sign into $c$.  For $M=2\Sigma$,
\[
                 Mv=4(v_ic_i)_i,
                 \qquad c_i=0\text{ if }\theta_{i,2}\ne1.
\]
A twist row with only one inactive label forces a single packet
position; its flat column then makes every $|v_i|$ equal.
This includes all nontrivial-central-torus cases.

There are only two rows with more than one inactive label.
For $\eta=-1$, $T=(1,1,i,-i)$ gives
$v=(b,b,a,a)$, up to signs, and $a/b$ is an odd integer
from the first two packet equations.  When $a>b$, exactly
the two degree-$b$ labels carry the packet.
For $\eta=1$, the row has one active label and three inactive
labels.  Normalize its active column as above, making it
position $0$.  Then $c_0=0$, and the other equations are
\[
             2v_ic_i=v_0+v_i\quad(1\le i\le3).
\]
Thus $k_i=v_0/v_i=2c_i-1$ are nonzero odd integers with
$\sum_{i=1}^3k_i^{-1}=-1$.  If all $|k_i|\ge3$, equality
forces all three to be $-3$.  Otherwise a value $-1$ forces
the other two to be opposites, whereas a value $1$ forces
the other two to equal $-1$.  These are exactly the displayed
paired and $3:1$ forms.  In an unequal paired form the packet
is supported precisely on the two smaller labels; in the
$3:1$ form it has equal coefficients $\pm1$ on all three
smaller labels.  Since the packet is supported in $I$, the
claims about $|I|$ and its complement follow.  Finally a
scaled-degree ratio of three can use only dimension primes
other than two, so it forces $q=3$ in the ordered family.
\end{proof}

\subsection{Fourier transport and its height limitation}\label{paper:sec:transport}

\begin{proposition}[Transport to a common complementary block]
\label{paper:lw:carry}
Assume $D=n^2$ is square.  Let $u,v,w$ be the three dimension primes, and let
$X\in P_u$, $Y\in P_v$ have pure $u$-power and pure $v$-power
twists respectively, allowing twist one.  Let $G$ be the Galois
group of the splitting field, acting by the unsigned tori.
Suppose a character $\psi$ of $G$ has the same kernel after
inverting its $u$-component, and also after inverting its
$v$-component.  If
\[
 \sum_{a\in G}\psi(a)\,a(S_{XY})\ne0,
\]
then $P_u$ and $P_v$ contain twist-one labels in a common full
$w$-block.  Both character fields contain the cyclic field cut
out by $\psi$.  The kernel condition holds in particular for
$\psi$ of prime-power conductor.
\end{proposition}
\begin{proof}
Galois covariance and Fourier projection on the two label orbits
give a $\overline\psi$-eigenvector $x$ supported on the orbit of
$X$ and a $\psi$-eigenvector $y$ supported on the orbit of $Y$,
with $\langle x,sy\rangle\ne0$.
The Weyl factors outside $u,v,w$, if any, are scalar: their actual
twists are trivial, and only a dyadic anomaly factor can remain.
Thus $s=\kappa W_uW_wW_v$ for a nonzero scalar $\kappa$.
Move $W_u^*$ to the first vector and $W_v$ to the second.
The new vectors are supported respectively on $P_u$ with pure
$u$-twists and $P_v$ with pure $v$-twists.  Conjugation by a local
Weyl inverts the corresponding torus component, so their characters
still have kernel $\ker\psi$.

A nonzero matrix coefficient of $W_w$ between these supports
supplies a pair of labels in the same full $w$-block, by
Theorem~\ref{paper:lw:block}.  This entry also preserves both complementary
twists.  A pure $u$-twist and a pure $v$-twist can agree in those
components only when both are one.  Finally, the stabilizer of a
label at a nonzero coordinate of a torus eigenvector lies in the
kernel of its character.  The fixed-field correspondence gives
the two field inclusions.
\end{proof}

\paragraph{Scope of the block conclusion.}
In the target family a homogeneous defect-two block at a sextic
prime has degrees $\ell^2u_X$; at the quartic prime it has degrees
$\ell u_X$.  The preceding theorems apply once that homogeneity
is proved.  Fourier transport does not itself preserve these
degree heights.  Its common block can contain labels of lower
height, and the foreign-column theorem gives integrality, not
vanishing, of the carried coordinates there.  Nor does a common
complementary twist identify a complementary residue block.
Finally, \eqref{paper:lw:rank-one} is a local integrality statement.
Proposition~\ref{paper:dyadic-four} removes foreign denominators
in every proper four-label dyadic case.  Other configurations
need the integral projection bounds; unitarity alone does not
give the global estimate $|B|\le\ell^2$.


The rank-one principal corner in Theorem~\ref{paper:block:prime-splitting}
constructs a categorical factor because it makes an entire spectral
idempotent positive.  A singleton intersection for a nonprincipal
homogeneous block gives exact local data and integer restrictions on
ambient fusion matrices.  The auxiliary multiplication in
Theorem~\ref{paper:dyadic-flat} is not thereby identified with ambient
tensor products, nor do the positive products acquire an algebra
structure from these identities alone.


\section{The exceptional modular dimension}\label{paper:modular:section}\label{paper:even-target}
Assume in this section that $\mathcal M$ is an integral, perfect,
fusion-simple modular category with
\[
 \FPdim(\mathcal M)=5400^2=2^6\,3^6\,5^4.
\]
We continue to write $d_X$ for a simple dimension, $P_\ell$ for the
principal residue block, and $K_X$ for the full character field of $X$.
The character field of a single matrix entry can be strictly smaller
than $K_X$. The conclusions below impose restrictions on an arbitrary
modular category at this dimension; they do not require a realization
as a Drinfeld center.

\subsection{Compulsory crossings and actual twists}
The family reduction supplies complementary witnesses
\[
 X\in P_3,\quad d_X=2^i3^j,\qquad
 Y\in P_5,\quad d_Y=2^h5,
 \qquad i,h\in\{2,3\},\quad j\in\{1,2\},
\]
with pure $3$-primary and $5$-primary twists and $\min(i,h)=2$.
Their normalized crossing has mass four. The positive-word method gives
an exact classification, including the actual balancing phases.

\begin{proposition}[Exact weight-four alternatives at three and five]
\label{paper:modular:weight-four-alternatives}
Let $X\in P_3$ and $Y\in P_5$ be complementary witnesses of degrees
$d_X=2^i3^j$, $d_Y=2^h5$, where $i,h\in\{2,3\}$,
$j\in\{1,2\}$, and $\min(i,h)=2$. Choose the witnesses so that their
actual twists are $3$-primary and $5$-primary, respectively. Put
$L=\operatorname{lcm}(d_X,d_Y)$ and $z=S_{XY}/L$.
Then precisely one of the following alternatives holds:
\begin{enumerate}
\item $z=1+3\xi$, with $\xi$ a primitive fifth root. Every actual
balancing word for this entry has coefficient $1$ at $1$ and
coefficient $3$ at $\xi$. Its field is $\mathbb Q(\zeta_5)$,
its mean squared modulus is $17/2$, and its energy is $34$.
\item $z=-\sum_{t=1}^4\eta_t\zeta_5^t$, where
$\{\eta_t\}=\{1,1,\omega,\omega^2\}$ and $\omega$ has
order three. Every actual balancing word has exactly these
four phases, each of coefficient $1$. Its mean squared modulus
is $19/4$.
\end{enumerate}
In both cases the normalized field trace
$[\mathbb Q(z):\mathbb Q]^{-1}\operatorname{Tr}_{\mathbb Q(z)/\mathbb Q}(z)$
is $1/4$.
In the signed case the field is $\mathbb Q(\zeta_{15})^+$,
of energy $19$, when the two positions whose cubic components
are $1$ are opposite modulo five. Otherwise it is
$\mathbb Q(\zeta_{15})$, of energy $38$.
In particular every signed coefficient is tame at three;
no signed wild crossing occurs. The full endpoint fields
$K_X,K_Y$ can be larger than this entry field.
\end{proposition}
\begin{proof}
Theorem~\ref{paper:family:thm:two-positive} gives the two forms,
initially allowing a nonidentity $5$-primary root in the first
and the multiset $\{1\}\sqcup\lambda\mu_3$ in the second,
with $\lambda$ a $3$-primary root.
Its contraction assertion places every actual vector in the convex
hull of the integral optimal vectors representing the same value.
We show that there is just one such vector.

The ordinary vectors $1+3\xi$ are distinct as their $\xi$ vary.
An exceptional vector lies in the field generated by $\mu_5$ and
$3$-power roots, whereas an ordinary vector with a root of order
at least $25$ has wild $5$-conductor; they cannot coincide. At
order $5$, use the integral basis $\zeta_5,\ldots,\zeta_5^4$
over the $3$-primary integer ring. The ordinary vector has
coefficient $2$ at its selected nonidentity root and coefficient
$-1$ at each other position; every exceptional coefficient is
a root of unity with a minus sign. These are distinct. Two
exceptional vectors representing the same value have identical
coefficients in this basis, hence identical actual phase vectors.
This argument works after enlarging either primary level, so
also proves uniqueness among all the integral optima from the
contraction theorem. The convex hull is therefore a point.

For the ordinary alternative, Lemma~\ref{paper:two-position-orders}
applied to the two actual positions forces $\xi$ to have order five.

If $\operatorname{lcm}(3,\operatorname{ord}\lambda)\ge9$,
the full positive word for $\chi_Y(X)=(d_X/4)z$ has, up to
its common sign, the forbidden support of
Theorem~\ref{paper:wild-support-forbidden}. Thus the coefficient
depth is one and its multiset is $\{1,1,\omega,\omega^2\}$.

For an exceptional vector, relative averaging at $5$ gives
$A_5(z)=1/4$, since $\sum_t\eta_t=1$, and
$A_5(|z|^2)=(5\sum_t|\eta_t|^2-|\sum_t\eta_t|^2)/4=19/4$.
Both right sides are rational, giving the asserted full means.
The ordinary means follow from $A_5(\xi)=-1/4$.

For the signed tame vector, a nontrivial stabilizer must preserve
the two identity positions. Its only possible $5$-component is
$-1$, which preserves those positions exactly when they are
opposite. In that case the other positions are opposite too,
and the simultaneous $3$-inversion swaps $\omega,\omega^2$;
the stabilizer is exactly full complex conjugation. The field
and energy statements follow.
\end{proof}

\begin{proposition}[A signed crossing forces triple activity]
\label{paper:modular:signed-triple-activity}\label{paper:modular:signed-mixed-activity}
Suppose a complementary weight-four crossing is signed. Then some
actual simple twist is divisible by $30$. More precisely, write its coefficient multiset as
$\{1,1,\omega,\omega^2\}$, let $\alpha=\theta_X$ be the
$3$-primary input twist and $\beta=\theta_Y$ the $5$-primary
input twist, and put $L=\operatorname{lcm}(d_X,d_Y)$.
The tensor product $X\otimes Y$ has a direct summand of
Frobenius--Perron dimension at least $L$ supported on simples
whose actual twists are divisible by $30$. The lower bound
improves to $2L$ if $\alpha\ne1$, and to $3L$ if the order
of $\alpha$ is at least nine.
Consequently complementary witnesses can be chosen with both
dyadic heights equal to two. Their full dyadic blocks are
nonprincipal and homogeneous of height two.
\end{proposition}
\begin{proof}
The four exact phase classes have monodromy phases
$-\eta_t\xi^t$, $1\le t\le4$, each carrying constituent
dimension $L$ when counted with fusion multiplicity. Every
constituent is dyadically active. A class is $3$-inactive
precisely when $\alpha\eta_t=1$, and $5$-inactive precisely
when $\beta\xi^t=1$. There are at most two classes of the
first kind and at most one of the second. At least one of
the four classes is therefore active at all three primes.
Each such class has total dimension $L$.

If $\alpha\ne1$, at most one class is $3$-inactive. If
$\operatorname{ord}(\alpha)\ge9$, none is $3$-inactive.
These observations give the two stronger bounds. The triple-active
constituent gives both edges incident to two; apply
Corollary~\ref{paper:two-neighbors} and
Proposition~\ref{paper:family:homogeneous-crossing} for the final assertion.
\end{proof}


\subsection{The compulsory fifteen-edge}

\begin{theorem}[Three and five must occur together]
\label{paper:three-five-mixing}
An ordinary perfect fusion-simple integral modular category of dimension
$2^6 3^6 5^4$ has an actual twist divisible by fifteen.
\end{theorem}
\begin{proof}
Suppose three and five are separated. Lemmas~\ref{paper:family:4.10}
and~\ref{paper:lw:isolated-quartic} show that five cannot be isolated, so
an actual twist divisible by ten exists. Corollary~\ref{paper:two-neighbors}
gives a $3$-free $Y\in P_5$ of degree $20$, with pure $5$-primary twist.
By Proposition~\ref{paper:family:family:crossings}, choose a $5$-free $X\in P_3$, with pure $3$-primary twist and degree
$2^i3^j$, where $i\in\{2,3\}$ and $j\in\{1,2\}$.
Separation makes every $3$-active degree divisible by $25$ and every
$5$-active degree divisible by $27$. Consequently both $X$ and $Y$
have complete twist one.

Their normalized crossing has mass four, positive residues four at
three and five, and a full tame $5$-inertia orbit. By Theorem~\ref{paper:family:thm:two-positive}, the exceptional
signed branch is ramified at both three and five: in its unique
expansion $-\sum_{t=1}^4\eta_t\zeta_5^t$ over the $3$-primary field,
the coefficient multiset is one copy of $1$ and one complete
$\mu_3$-coset, so some coefficient is nonrational. The independence
of the four nonidentity fifth roots detects nontrivial inertia at
three; Corollary~\ref{paper:cor:first-positive-orbit} detects inertia at five.
This contradicts Proposition~\ref{paper:lw:isolated}. Therefore
\[
 S_{XY}=L(1+3\xi),\qquad
 L=\operatorname{lcm}(d_X,20)=2^i3^j5,
\]
where $\xi$ is a nonidentity $5$-power root.

Apply a $5$-inertia automorphism $\sigma$ moving $\xi$. Balancing and
$\theta_X=\theta_Y=1$ show that $S_{XY}-\sigma(S_{XY})$ is divisible
by $27$: only $5$-active constituents contribute and their degrees
are all divisible by $27$. But the same difference equals
$3L(\xi-\sigma\xi)$, whose valuation at three is $j+1$. Hence $j=2$,
so $d_X$ is $36$ or $72$.

The full field $K_X$ is ramified at five and therefore unramified at
three by separation. Its $5$-inertia orbit has at least four labels.
All $3$-active labels have $5$-unit codegree, so $5$-inertia fixes
them, and Theorem~\ref{paper:lw:columns} makes $W_3$ integral
at five. Since $v_2(f_X)$ is zero or two, its $X$-column is integral
at two by \eqref{paper:lw:foreign-bound}; at every other prime outside
three and five the same assertion is automatic. Finally $v_3(f_X)=2$.
Lemma~\ref{paper:ternary-column} therefore supplies $3$-active labels of degree
$4d_X=144$ or $288$. Neither degree is divisible by $25$, contrary
to separation. This contradiction proves the theorem.
\end{proof}

\begin{proposition}[Ternary inactivity at the largest ordinary mixed crossing]
\label{paper:tame-largest-ternary-inactive}
In an ordinary minimal table of dimension $2^6 3^6 5^4$, let
$X\in P_3$, $Y\in P_5$ have pure $3$-primary and pure
$5$-primary twists, respectively. Suppose
\[
 (d_X,d_Y)\in\{(36,40),(72,20)\},\qquad
 \frac{S_{XY}}{\operatorname{lcm}(d_X,d_Y)}=1+3\zeta_5,
\]
where $\zeta_5$ is primitive. Then $\theta_{X,3}=1$.
\end{proposition}
\begin{proof}
Put $c=360$, $K=\mathbb Q(\zeta_5)$,
$a=[K_X:K]$ and $b=[K_Y:K]$. Thus $D/c^2=225$.
For a faithful character $\chi$ of
$\operatorname{Gal}(K/\mathbb Q)\cong C_4$, the Gauss-sum
formula gives
\[
 |\widehat{(1+3\zeta_5)}(\chi)|^2=45.
\]
Consequently the exact orbit-projection identity
\eqref{paper:orbit-projection} has squared norm $ab/5$.

Write $\operatorname{ord}(\theta_{X,3})=3^\alpha$ and
$\operatorname{ord}(\theta_{Y,5})=5^\beta$. The twist cap
in Lemma~\ref{paper:family:lw:weyl-line} gives $\alpha,\beta\le2$.
The square action on twists gives $3\mid a$ when $\alpha=2$:
the ternary inertia orbit of a primitive ninth-root twist has
three members, and this inertia fixes $K$. If $\beta=2$, its
twist orbit has ten members; hence $10\mid[K_Y:\mathbb Q]=4b$
and $5\mid b$. These are square-orbit divisibilities and do
not use the stronger odd-dimensional active-inertia theorem.

Since $X$ is $5$-inactive, the faithful tame torus channel has
the following bounds:
\[
\begin{array}{c|c}
 \beta&45ab\text{ is at most}\\ \hline
 0&225/5=45\\
 1&(2/5)225=90\\
 2&225.
\end{array}
\]
The first is the exact inactive compression of
Theorem~\ref{paper:inactive-fourier-budget}; the second is
Lemma~\ref{paper:nonquadratic-half-local}; the third is ordinary
unitarity. If $\alpha=2$, the condition $3\mid a$ contradicts
the first two bounds. In the last case $5\mid b$, so $ab\ge15$,
again a contradiction.

Suppose $\alpha=1$. At three, $Y$ is inactive and the inflated
channel has trivial torus character. The half-channel factor
in Lemma~\ref{paper:half-channel-local} is $4/9$. This factor
multiplies each of the displayed bounds. Indeed the local
operators, phase projections and torus-isotypic projections
at three and five commute. After replacing the endpoint orbit
projections by their containing phase and torus projections,
the matrix-coefficient estimate is the norm estimate for the
product of the compressed local Weyl operators. The bound at
three may be applied in the reversed direction by taking
adjoints. Thus
\[
 45ab\le
 \begin{cases}
 20,&\beta=0,\\
 40,&\beta=1,\\
 100,&\beta=2.
 \end{cases}
\]
The first two are impossible because $ab\ge1$; the third is
impossible because $5\mid b$. Hence $\alpha=0$.
\end{proof}

\subsection{When no twist order is divisible by ten}

\begin{lemma}[A dyadic unit diagonal excludes lower heights]
\label{paper:nt-unit-diagonal-homogeneity}
Let $X$ have degree $12$ at $n=5400$, let
$W=\gamma^3W_2$ be the corrected dyadic Weyl matrix, and suppose
$W_{XX}=5/6$.  Then the full two-block of $X$ is homogeneous
of two-height two.
\end{lemma}
\begin{proof}
Suppose a label $L$ of smaller two-height belongs to the same
full block.  Their characters have equal reductions at a place
over two.  The $X$-column of the complementary inverse Weyl is
integral there by the foreign denominator bound, because
$v_2(f_X)=2$.  Multiplying the character congruence by that
column gives
\[
 2W_{XX}\equiv \frac{2d_X}{d_L}W_{LX}
       =\frac{d_X}{d_L}(2W_{LX})\equiv0\pmod{\mathfrak p}.
\]
Here $2d_X/n=1/225$ is a two-unit, $2W_{LX}$ is integral
at two by square-integrality, and $d_X/d_L$ is divisible by
two.  This contradicts $2W_{XX}=5/3$.  A label of two-height
three is a defect-zero singleton and cannot belong to this
block.  These are all alternatives to height two.
\end{proof}

\begin{lemma}[A cubic twist-moment estimate]
\label{paper:real-cubic-moment}
Let $v$ be a vector of integers in $F=\mathbb Q(\zeta_9)^+$.
Suppose three-primary Galois automorphisms permute its coordinates,
preserving positive integer weights $d_V$ and acting on
three-primary roots $\vartheta_V$ by the square action. Partition
the supported coordinates into low weights $d_V\le100$ and high
weights $d_V\le300$. At each high coordinate assume that $v_V$
has residue zero at three. Then
\begin{equation}\label{paper:eq:cubic-moment}
\left|\sum_Vd_V\vartheta_Vv_V\right|
                         \le100\|v\|^2.
\end{equation}
If equality holds, $\sum_Vd_Vv_V/100$ is an integer congruent to
$\|v\|^2$ modulo two. Furthermore, any nonzero coordinate orbit
whose coefficients all belong to $2\mathcal O_F$ makes \eqref{paper:eq:cubic-moment} strict.
\end{lemma}
\begin{proof}
Put $z=\zeta_9+\zeta_9^{-1}$ and
$q=\zeta_9^2+\zeta_9^{-2}$. The integers $1,z,q$ form an integral
basis. For
\[
 \alpha=a+bz+cq,\qquad Q=b^2-bc+c^2,
\]
the normalized mean square and three-residue are
\begin{equation}\label{paper:eq:cubic-norm-residue}
M(\alpha)=a^2+2Q,\qquad
       \alpha\equiv a-b-c\pmod{\mathfrak p_3}.
\end{equation}
The nonzero values of $Q$ are $1$ or at least $3$.

Average first over the Galois kernel fixing $F$. An orbit with
$\vartheta_V$ of order at least $27$ contributes zero to the
twist moment: a kernel element multiplies this twist by a
primitive third root while preserving the coefficient and weight.
For the remaining orbits, a nonconstant coefficient or a primitive
ninth-root twist gives, with possible repetitions, a triple of
coefficient--twist pairs. Its energy is
\[
                         E=3a^2+6Q.
\]
If the twist has order at most three, its moment has absolute
value $3d|a|$. If its order is nine, that absolute value is
$3d\sqrt Q$. The latter formula is the three-term Fourier
identity: the constant coefficient cancels and
\[
 |\alpha+\omega\sigma\alpha+\omega^2\sigma^2\alpha|^2=9Q.
\]
Changing the cyclic orientation does not change this identity.
Rational coefficients with twists of order at most three may be
handled one coordinate at a time.

For a high triple, \eqref{paper:eq:cubic-norm-residue} implies
\begin{equation}\label{paper:eq:cubic-high-bound}
3|a|\le a^2+2Q,
          \qquad 3\sqrt Q\le a^2+2Q.
\end{equation}
Indeed, $Q=0$ gives $a\in3\mathbb Z$; $Q=1$ forces
$a\not\equiv0\pmod3$; and both inequalities are strict for
$Q\ge3$. Thus its moment is at most $100E$.
For a low triple, the stronger elementary inequalities
$|a|\le a^2+2Q$ and $\sqrt Q\le a^2+2Q$ suffice.
The analogous rational-coordinate bounds are immediate.
Summing over coordinate orbits proves \eqref{paper:eq:cubic-moment}.

The nonzero equality cases for a high triple have $d=300$.
For a twist of order at most three they are
\[
             (|a|,Q)=(1,1),(2,1),(3,0);
\]
for a primitive ninth-root twist only $(|a|,Q)=(1,1)$ remains.
Its untwisted weighted sum divided by $100$ is $9a$, whose
parity equals that of $E=3a^2+6Q$. A low equality case has
$d=100$ and the rational coefficient $1$ or $-1$; it has the
same parity property. Rational high equality cases have
$d=300$ and coefficient $3$ or $-3$, and also have that
property. Repetitions preserve it, proving the asserted global
congruence. None of the listed equality cases has all three
integral-basis coefficients even. A nonzero orbit in
$2\mathcal O_F$ therefore makes \eqref{paper:eq:cubic-moment} strict.
\end{proof}

\begin{lemma}[No odd ternary packet of norm at most twenty-seven]
\label{paper:odd-ternary-packet}
Let $u$ be a vector of three-primary cyclotomic integers. Suppose
its distinct Galois conjugates are orthogonal and are also coordinate
permutations preserving positive integer weights $d_V$ and acting
on three-primary twists $\vartheta_V$ by the square action.
Partition its support into low weights $d_V\le100$ and high
weights $d_V\le300$, with zero three-residue at every high
coordinate. Write $\delta u=\sum_V d_Vu_V$.
There is no such vector satisfying
\begin{equation}\label{paper:eq:odd-packet-data}
 \|u\|^2=b,\qquad b\le27\text{ odd},\qquad
 \delta u=0,\qquad |\delta(\vartheta u)|=100b.
\end{equation}
\end{lemma}
\begin{proof}
For a three-primary cyclotomic integer, the normalized mean
square is
\[
                         M(\alpha)=\sum(a^2+c^2-ac)
\]
in the two-root integral basis of each $\mu_3$-coset. Thus $M$
is an integer, and its nonzero minimum one occurs exactly at
roots of unity, including their negatives. A nonzero coefficient
with zero three-residue has $M\ge2$.

Put $F=\mathbb Q(u)$ and $t=[F:\mathbb Q]$. The sum of the
coefficient mean squares is $b$, so the support has at most $b$
coordinates. All Galois conjugates have the same zero coordinates;
therefore orthogonality gives $t\le b\le27$. The possibilities are
\begin{equation}\label{paper:eq:odd-packet-field-degrees}
                       t\in\{1,2,3,6,9,18,27\}.
\end{equation}
If a coefficient $\alpha$ generates a subfield of degree $e$,
average $u$ over the Galois kernel fixing it. Orthogonality makes
this projection have squared norm $be/t$. Its unchanged coordinate
orbit contains at least $e$ conjugates of $\alpha$, so
\begin{equation}\label{paper:eq:odd-packet-coefficient-bound}
                              M(\alpha)\le b/t.
\end{equation}
Repetitions of coefficient conjugates only strengthen this bound.

\emph{Rational and quadratic fields.}
For $t=1$, nonzero low coefficients are integers and high ones
are divisible by three. Thus $d_V|u_V|\le100u_V^2$ everywhere.
Equality in \eqref{paper:eq:odd-packet-data} forces low equality
coordinates of degree $100$ with coefficient $\pm1$, and high
ones of degree $300$ with coefficient $\pm3$. Their weighted
sum divided by $100$ is an integer congruent to their energy
$b$ modulo two, contradicting $\delta u=0$.

For $t=2$, put $x=\operatorname{Re}u$. Central three-inertia
inversion preserves actual twists and conjugates coefficients.
Consequently
\[
          \delta(\vartheta x)=\delta(\vartheta u),
                       \qquad\|x\|^2=b/2,
\]
where orthogonality of $u$ and $\overline u$ gives the second
identity. A low real part is a half-integer; a high real part
is in $(3/2)\mathbb Z$, because $a+c\omega\equiv0$ implies
$2a-c\equiv0\pmod3$. Hence $d_V|x_V|\le200x_V^2$, and
\eqref{paper:eq:odd-packet-data} forces equality. Every nonfixed
pair under central inversion costs the even integer $2M(u_V)$
in the norm. Since $b$ is odd, there is a nonzero fixed coordinate,
whose coefficient is rational. The stronger bound
$d_V|x_V|\le100x_V^2$ holds there, making the estimate strict.

\emph{The two ninth-root fields.}
For $t=3$, apply Lemma~\ref{paper:real-cubic-moment} directly.
Its bound must be an equality, and its parity conclusion gives
$\delta u/100\equiv b\pmod2$, a contradiction.
For $t=6$, put $v=u+\overline u\in\mathcal O_{\mathbb Q(\zeta_9)^+}^I$.
Then
\[
                  \|v\|^2=2b,\qquad
                       |\delta(\vartheta v)|=200b.
\]
The cubic estimate must again be an equality. A coordinate orbit
not fixed by central inversion has even size $r$ and costs the
even integer $rM(\alpha)$ in the norm of $u$. Oddness of $b$
therefore supplies a nonzero central-fixed orbit. On it $u$ is
real and $v=2u$, contradicting the strict clause of the cubic
estimate.

\emph{Larger coefficient fields.}
For $t=18$ or $27$, \eqref{paper:eq:odd-packet-coefficient-bound}
and integrality of $M$ make every nonzero coefficient a root
of unity. Hence no high coefficient is possible.

For $t=9$, use $F=\mathbb Q(\zeta_{27})^+$, with integral
basis $1,c_1,\ldots,c_8$, where
$c_j=\zeta_{27}^j+\zeta_{27}^{-j}$. The usual monic recurrence
for the $c_j$ proves the basis assertion. All $c_j$ have trace
zero, whereas $\operatorname{Tr}(1)=9$. The rational projection
of $u$ has norm squared $b/9\le3$. A high coefficient can have
subfield degree $e=1,3$ or $9$; a nonzero trace has absolute
value at least $3,3$ or $9$, respectively. Its orbit contributes
at least $(\operatorname{Tr}\alpha)^2/e$ to that projection.
Degrees one and nine exceed the budget. Degree three could only
exhaust it, leaving one orbit with a nonzero weighted trace,
contrary to $\delta u=0$. All high coefficients therefore have
trace zero.

For a trace-zero integer $\alpha=\sum_{j=1}^8a_jc_j$, the
exact trace Gram matrix gives
\begin{equation}\label{paper:eq:odd-packet-real-twenty-seven-norm}
 M(\alpha)=2\sum_{j=1}^4
          (a_j^2+a_{9-j}^2-a_ja_{9-j}).
\end{equation}
Indeed the diagonal inner products are two and the only
nonzero off-diagonal ones are $-1$ when the indices sum to nine.
A nonzero mean square at most three must be two, represented
by one Eisenstein norm-one pair. Its coefficient is $\pm c_k$
for some $1\le k\le13$, using
$c_j+c_{9-j}=-c_{9+j}$. It has three-residue $\pm2$, not zero.
This contradicts the bound $M(\alpha)\le b/9\le3$ for a high
coefficient. Thus none occurs.

In all three larger-field cases only low weights remain. If
$s$ is the support size, then
\[
 100b=|\delta(\vartheta u)|
        \le100\sum_V|u_V|
        \le100\sqrt{s}\sqrt b\le100b.
\]
Equality gives $b$ supported labels, all of degree $100$, all
coefficients of modulus one, and a common value of
$\vartheta_Vu_V$. Central inversion preserves the twist while
conjugating the coefficient, so every coefficient is real,
hence $1$ or $-1$. Odd order of the twists forces all signs to
agree. This contradicts $\delta u=0$ and exhausts
\eqref{paper:eq:odd-packet-field-degrees}.
\end{proof}

\begin{theorem}[Exclusion of the no-ten graph]
\label{paper:no-ten-excluded}
The activity graph $\{6,15\}$ cannot occur at $n=5400$.
\end{theorem}
\begin{proof}
Assume that no actual twist is divisible by ten. The following
local calculations reduce this hypothesis to the odd ternary packet
excluded above. All operators in these claims have the corrected
dyadic normalization $W=\gamma^3W_2$, $U=T_2$.

\begin{proofclaim}[The remaining witness when two and five are separated]
\label{paper:two-five-height}
Suppose an ordinary perfect fusion-simple integral modular category of
dimension $2^6 3^6 5^4$ has no twist divisible by ten. A compulsory
weight-four crossing cannot have its $5$-free endpoint of degree $36$.
In particular a $5$-free witness of dyadic height two in $P_3$ must
have degree $12$.
\end{proofclaim}
\noindent\emph{Proof of the claim.}
A nonempty $5$-active fiber has actual dyadic twist one. Its dyadic
upper unipotent is the scalar $\gamma^{-1}$, hence its Weyl is the
scalar $\gamma^3$ by Bruhat. Since the Weyl square is an unsigned
torus permutation, $\gamma^6=1$; thus $\gamma=\pm1$.
Proposition~\ref{paper:lw:isolated} now applies to the separated primes two
and five and separates their full character fields.

Suppose $X$ has degree $36$ and pure $3$-primary twist. Pair it with
any opposite compulsory $3$-free witness in $P_5$. The common dyadic
weight is four, and Corollary~\ref{paper:cor:first-positive-orbit}
makes the full $5$-inertia orbit of $X$ have at least four labels.
It also makes $K_X$ ramified at five, hence unramified at two.
Here $v_2(f_X)=v_3(f_X)=2$ and $v_5(f_X)=4$.
The foreign-denominator bound makes the $X$-column of $W_2$ integral
at three. Every $2$-active label has $5$-unit codegree, so the final
assertion of Theorem~\ref{paper:lw:columns} makes $W_2$ integral at five;
these active labels are also fixed by $5$-inertia. Integrality at
all other foreign primes is immediate. Lemma~\ref{paper:dyadic-column}
therefore contradicts the orbit of length at least four.
\par\smallskip

\begin{proofclaim}[An integral column of norm eight outside the foreign orbit]
\label{paper:two-five-norm-eight}
Under the same separation hypothesis, let $X$ be a compulsory
$5$-free witness of degree $12$, with pure $3$-primary twist.
Correct the dyadic Weyl by $W=\gamma^3W_2$, so its upper unipotent
is the actual primary twist $U=T_2$. Its full $5$-inertia orbit has four labels $X_1,\ldots,X_4$.
For one integer
$l\in\{1,2\}$, the integer vector $u=6W\mathbf e_{X_1}$ satisfies
\[
 u_{X_i}=6\delta_{i1}-l,
 \qquad \sum_{T\notin\{X_1,\ldots,X_4\}}u_T^2=8,
 \qquad \sum_{T\;2\text{-active}}u_T^2(1-\theta_{T,2})=6l.
\]
All support labels outside the orbit are $5$-inactive, and their
coefficients are constant on $5$-inertia orbits.
\end{proofclaim}
\noindent\emph{Proof of the claim.}
The preceding proof gives $\gamma=\pm1$ and $K_X$ unramified at two.
The foreign-denominator bound at three, together with $v_3(f_X)=4$,
and integrality at five from separation give
$u=6W\mathbf e_X\in\mathbb Z^I$ and $\|u\|^2=36$.

Let $V_2$ be the moving space of the corrected dyadic group.
Separation makes the full $5$-local group trivial on $V_2$, as in
the proof of Proposition~\ref{paper:lw:isolated}. Thus the projections of
all labels $X_i$ in the $5$-inertia orbit to $V_2$ agree, and
\[
                    W\mathbf e_{X_i}=\mathbf e_{X_i}+v
\]
for a common vector $v$, fixed by the full $5$-local group. In particular
it is supported on $5$-inactive labels. Write $t$ for the orbit size.
Symmetry makes all $t$ coordinates $v_{X_i}$ equal. Unitarity gives
$2v_{X_i}+\|v\|^2=0$, so write $v_{X_i}=-l/6$ with integer $l>0$.
The case $l=0$ would make the Weyl line of $X_i$ invariant.
The $t$ orbit coordinates give $t(l/6)^2\le l/3$, hence $tl\le12$.
Corollary~\ref{paper:cor:first-positive-orbit} makes $t=4\cdot5^s$
for some $s\ge0$. Thus $t=4$ and $l\le3$.
If $l=3$, the four coordinates exhaust the column norm and the
column has no active coordinate, again contradicting the Weyl-line
lemma. Hence $l=1$ or $2$, and the outside norm is
\[
36-(6-l)^2-3l^2=4l(3-l)=8.
\]
Finally the local relation, on the upper-fixed and central-fixed
source $X_i$, gives $\sum_Tu_T^2\theta_{T,2}=6u_{X_i}$.
Subtract from $\sum_Tu_T^2=36$ to obtain the last identity.
\par\smallskip

\begin{proofclaim}[The degree pair when two and five are separated]
\label{paper:two-five-pair}
If no actual twist is divisible by ten, a compulsory crossing can be
chosen with degrees $(12,40)$. It has the ordinary value
$1+3\xi$, with $\xi$ a primitive fifth root.
\end{proofclaim}
\noindent\emph{Proof of the claim.}
Two is not isolated by Theorem~\ref{paper:family:family:dyadic-isolation}.
Thus a twist divisible by six exists, and
Corollary~\ref{paper:two-neighbors} supplies a $5$-free principal-three
witness $X$ of degree $12$ or $36$. The signed crossing is excluded
by Proposition~\ref{paper:modular:signed-triple-activity}.
If an opposite witness $Y$ had degree $20$, then
$v_2(L)=2$. The two actual phase buckets of its ordinary crossing
have total dimensions $L$ and $3L$. At least one bucket is
$5$-active, but separation makes every $5$-active degree divisible
by eight. Neither bucket dimension is divisible by eight, a
contradiction. Thus $d_Y=40$. Claim~\ref{paper:two-five-height}
excludes $d_X=36$, proving the assertion.
\par\smallskip

\begin{proofclaim}[A signed dimension identity and its fixed-space consequence]
\label{paper:column-dimension-identity}
Retain the corrected operators $W,U$, the degree-twelve source
$X$ and the integer $l\in\{1,2\}$ of
Claim~\ref{paper:two-five-norm-eight}. Put
\[
 R_X=\sum_{i=1}^4\mathbf e_{X_i},\qquad
 w=6(W\mathbf e_X-\mathbf e_X),
 \qquad A=P_{\{\theta_2=1\}}w.
\]
Let $I,M,H$ be the squared masses of $w$ outside the source orbit
on actual dyadic twists $1,-1$, and orders at least four,
respectively. Then
\begin{equation}\label{paper:dyadic-signed-dimension}
 -48l+\sum_{\substack{Z\notin\operatorname{supp}R_X\\
                       \theta_{Z,2}=1}}d_Zw_Z
   =\sum_{\theta_{Z,2}=-1}d_Zw_Z\in25\mathbb Z,
\end{equation}
and
\[
                 I+M+H=8,\qquad I-M=6l-4l^2.
\]
Moreover
\begin{equation}\label{paper:inactive-defect-zero}
              v_2(d_Y)=3\quad\Longrightarrow\quad\chi_Y(A)=0.
\end{equation}
\end{proofclaim}
\noindent\emph{Proof of the claim.}
The source coordinates of $w$ are all $-l$. The source is
upper-fixed and central-fixed, so the local relations give
\[
                    Ww=-w,\qquad WUw=(U^{-1}-1)w.
\]
The vector $W^{-1}d$ is supported on the dyadically inactive
coordinates. Pairing the second identity with $d$, then using
the first, gives
\[
                        \sum_Z d_Zw_Z\theta_{Z,2}^{-1}=0.
\]
Normalized trace in a common dyadic cyclotomic field is $1$ on
$1$, $-1$ on $-1$, and zero on roots of order at least four.
It gives \eqref{paper:dyadic-signed-dimension}; separation makes
every active degree divisible by $25$. The four source
coordinates of $u=6W\mathbf e_X$ spend
$(6-l)^2+3l^2=36-12l+4l^2$ of its norm. Apply the same trace
to $\sum_Zu_Z^2\theta_{Z,2}=6u_X=36-6l$ and subtract that
source contribution. This proves the squared-mass identities.

Let $V_2$ be the orthogonal complement of the fixed space of the
corrected dyadic local group. The vector $w=6(W-I)\mathbf e_X$
belongs to $V_2$. The upper-unipotent averaging projector
preserves $V_2$, so $A\in V_2$ as well. The global matrix $s$
preserves this space: its dyadic factor belongs to the group,
and its complementary factors commute with the group. If
$v_2(d_Y)=3$, its codegree is odd and
Lemma~\ref{paper:family:lw:weyl-line} makes $\mathbf e_Y$ fixed by the
corrected group. Hence $(sA)_Y=0$, proving
\eqref{paper:inactive-defect-zero}.
\par\smallskip

\begin{proofclaim}[Exclusion of the second norm-eight alternative]
\label{paper:second-norm-eight-excluded}
In Claim~\ref{paper:two-five-norm-eight}, necessarily $l=1$.
\end{proofclaim}
\noindent\emph{Proof of the claim.}
Suppose $l=2$. Use $R_X,w,A$ and $I,M,H$ from
Claim~\ref{paper:column-dimension-identity}. Its signed
dimension identity has source term $-96$, and its norm identities give
\[
                  I+M+H=8,\qquad I-M=-4,
                  \qquad 2I+H=4.
\]
Thus at most two outside inactive labels occur, each with
coefficient $1$ or $-1$.

Every supported outside label has ternary height at most one.
Indeed a label $Z$ of height at least two has $v_3(f_Z)\le2$;
symmetry and Theorem~\ref{paper:lw:columns} make $W_{ZX}$ integral
at three, so $3\mid u_Z$, contrary to the outside norm eight.
A label of dyadic height three has a singleton full dyadic
block and cannot occur in this column. Prime-power degrees
are excluded. Therefore a $5$-free outside label is the vacuum,
or has degree $6$ or $12$.

A degree-six label has a nontrivial $5$-inertia orbit. Pair it
with an opposite witness of degree $20$ or $40$ in $P_5$. Its normalized
crossing has mass two and residue two at five, so
Corollary~\ref{paper:cor:one-large-residue} and strictness make
it a sum of two distinct $5$-primary roots. Its tame stabilizer
has order at most two. Since the coefficients of $w$ are
$5$-inertia invariant, a degree-six outside coefficient consumes
at least two slots of the same degree and sign.

The inactive source contribution is $-96$. If all outside
inactive slots are $5$-free, their signed degree sum cannot be
$96$ modulo $25$: for nonunit slots it is a multiple of six
of absolute value at most $24$, and a vacuum slot gives only
the additional values $\pm1$, $\pm1\pm6$, $\pm1\pm12$.
If both are $5$-bearing there is already a contradiction modulo
five. Thus precisely one is $5$-free. Modulo five its signed
degree must be one. This leaves a vacuum coefficient $+1$ or
a degree-six coefficient $+1$; the latter is excluded by the
orbit bound. Hence the vacuum occurs with coefficient $+1$.

The other label has $5$-bearing degree
\[
      10,15,20,30,60,50,75,100,150,\text{ or }300.
\]
The congruence $1\pm d_T\equiv96\pmod{25}$ leaves exactly
$+20$ and $-30$. These are the only possible inactive supports. They have
$I=2$, hence $H=0$ and $M=6$, so higher dyadic phases do not
occur. The vacuum coefficient forces the common complementary
ternary twist of the column to be one. Finally $T$ is the only
outside label of its degree, so its nonzero coefficient and
$5$-inertia invariance force it to be fixed.

Thus the inactive part of $w$ is $-q$, where
\[
 q=2R_X-\mathbf e_0-\mathbf e_T\quad(d_T=20),
 \qquad\hbox{or}\qquad
 q=2R_X-\mathbf e_0+\mathbf e_T\quad(d_T=30).
\]
Write $w=-q+a$ for its inactive and active parts. These are
integral vectors with
\[
 \|q\|^2=18,\qquad \|a\|^2=6,\qquad Uq=q,\quad Ua=-a.
\]
Both $q$ and $a$ are fixed by the full five-local group,
since $w$ is fixed and that group preserves dyadic twist fibers.
All supported complementary twists are trivial.

We first locate the full residue blocks. The nonzero vacuum
coordinate of the Weyl column puts $X$ in $P_2$; every other
supported coordinate, including $T$, is in this same full
dyadic block. Suppose $X\notin P_5$. Projection of $q$ to
the principal five-block is $-\mathbf e_0$, possibly together
with its coefficient $\pm\mathbf e_T$. This projected vector
is fixed by the full five-local group. Its dimension is $-1$ modulo five, contradicting Lemma~\ref{paper:fixed-vector-dimension-congruence}.
Hence $X\in P_5$. If $T\notin P_5$, projection to its full
five-block is the singleton vector $\pm\mathbf e_T$, which
would make its Weyl line invariant. But $5\mid f_T$, again
impossible. We have proved
\[
                X\in P_2\cap P_3\cap P_5,
                \qquad T\in P_2\cap P_5.
\]

The local Bruhat identity already obtained above gives
$WUw=(U^{-1}-1)w=-2a$. Together with
$w=-q+a$, $Uw=-q-a$, and $Ww=-w=q-a$, this yields
\[
 -Wq+Wa=q-a,\qquad Wq+Wa=2a,
 \qquad Wq=-\frac12q+\frac32a.
\]

Let $\sigma$ be any element of three-inertia, acting by its
unsigned torus permutation. It commutes with $W$ and preserves
the inactive and active dyadic subspaces. Unitarity consequently
gives
\[
                  \langle q,\sigma q\rangle
                         =3\langle a,\sigma a\rangle\in3\mathbb Z.
\]
The set underlying $R_X$ is one full five-inertia orbit, and
its image under $\sigma$ is either the same orbit or disjoint
from it. The degree of $T$ differs from twelve and one.
It follows from the displayed formula for $q$ that
\[
               \langle q,\sigma q\rangle\in\{1,2,17,18\}.
\]
Only eighteen is divisible by three. Thus every three-inertia
element preserves the four-element source orbit and fixes $T$.

The source is fixed by all dyadic tori, and the relation
$H_2W=WH_2^{-1}$ therefore fixes the Weyl column under those
tori. Its inactive part $q$ is fixed by them as well. The unique
degree-$d_T$ coordinate of $q$ makes $T$ fixed by dyadic inertia.
It is also fixed by three-inertia, as just proved, and by
five-inertia from the vacuum enumeration. The conductor theorem
therefore gives $K_T=\mathbb Q$.

Finally put $\lambda=\chi_T(X)\in\mathbb Z$. Strictness gives
$|\lambda|<12$, and $T\in P_2\cap P_5$ gives
\[
                        \lambda\equiv12\pmod{10},
                        \qquad\lambda\in\{-8,2\}.
\]
If $d_T=20$, LCM integrality at the degree pair twelve, twenty
makes $\lambda$ divisible by three. If $d_T=30$, principality
of $X$ at three gives $\chi_X(T)\equiv30\equiv0\pmod3$;
symmetry and $\chi_T(X)=(2/5)\chi_X(T)$ again give
$3\mid\lambda$. Neither possible value of $\lambda$ is
divisible by three. This contradiction excludes $l=2$.
\par\smallskip

\begin{proofclaim}[The four dyadic mass profiles]
\label{paper:dyadic-mass-profiles}
When no twist order is divisible by ten, the column of
Claim~\ref{paper:two-five-norm-eight} has $l=1$ and
\[
 (I,M,H)\in\{(2,0,6),(3,1,4),(4,2,2),(5,3,0)\}.
\]
\end{proofclaim}
\noindent\emph{Proof of the claim.}
Claim~\ref{paper:second-norm-eight-excluded} gives $l=1$.
The preceding signed dimension identities then give
$I-M=2$ and $2I+H=10$. Nonnegative integral masses give the list.
\par\smallskip

\begin{proofclaim}[Homogeneity and maximal companions]
\label{paper:no-ten-maximal-companions}
In the no-ten branch, the full two-block of the compulsory
degree-twelve source has twelve labels. For one sign
$\epsilon\in\{1,-1\}$, its corrected Weyl column satisfies
\[
 w=6(W-I)\mathbf e_X
   =-\sum_{i=1}^4\mathbf e_{X_i}-\mathbf e_T
      +\mathbf e_Z+\epsilon\sum_{j=1}^6\mathbf e_{U_j},
 \qquad
 d_T=12,\quad d_Z=60,\quad d_{U_j}=300.
\]
All twelve labels are distinct. All outside coefficients have
absolute value one, and the block is homogeneous of two-height
two. If $(I,M,H)$ is an initial mass profile, precisely $M$
of the $U_j$ are dyadically inactive, $M$ have twist $-1$,
and $H=6-2M$ have higher dyadic order.
\end{proofclaim}
\noindent\emph{Proof of the claim.}
The exclusion of $l=2$ gives $W_{XX}=5/6$.
Lemma~\ref{paper:nt-unit-diagonal-homogeneity} makes its full
two-block homogeneous. The dense homogeneous-block theorem
then makes every entry of this Weyl column nonzero with exact
two-valuation $-1$. Since $6W\mathbf e_X$ is integral,
each of its nonzero entries is odd. The outside norm is eight,
so it consists of eight coefficients $1$ or $-1$. Density
identifies these eight labels and the four source labels with
the whole block.

Every supported outside label has three-height at most one:
otherwise the foreign-column bound makes its Weyl entry
three-integral, so its nonzero coefficient in $6W\mathbf e_X$
is divisible by three, exceeding the outside norm. Exact
two-height two, prime-power exclusion and $d_V\mid5400$
therefore leave just the degrees
\[
                         12,20,60,100,300.
\]
Separation forces every dyadically active one to have degree
$100$ or $300$.

Write $\delta(\mathbf e_V)=d_V$ and
$k=(Wd)_X/d_X\in2\mathbb Z$, using the complementary
vacuum coefficient and its foreign two-divisibility. Then
\[
                  \delta w=72(k-1)\in72(2\mathbb Z+1).
\]
Let $A_-$ and $A_{\rm hi}$ be the twist-$-1$ and higher-order
parts of $w$. The signed dimension identity gives
$\delta A=\delta A_-$.
Both $A_-$ and $A_{\rm hi}$ have dimensions
divisible by $25$. Thus $25\mid\delta w$ and
\[
 \delta w\in1800(2\mathbb Z+1),\qquad
 |\delta w|=|2\delta A_-+\delta A_{\rm hi}|
       \le300(2M+H)=1800.
\]
Equality is forced: $\delta w=\epsilon1800$.
Every active coefficient therefore has sign $\epsilon$
and degree $300$.

It remains to determine the outside inactive coordinates.
There are $I=M+2$ of them, and their signed dimensions sum to
$48+\epsilon300M$. After division by four, they are signed
values in $\{3,5,15,25,75\}$ with sum
\[
                           12+\epsilon75M.
\]
Modulo five, the number of $+3$ terms minus the number of
$-3$ terms is $-1$ or $4$, since $I\le5$.
The value $4$ uses at least four terms and leaves too little
capacity to supply $\epsilon75M$. In the value $-1$ case,
three or more terms of absolute value three would leave at
most $M-1$ other terms; these could not supply
$15+\epsilon75M$, whose absolute value exceeds $75(M-1)$.
There is consequently exactly one $-3$ term.

The remaining $M+1$ signed values, from $\{5,15,25,75\}$,
sum to $15+\epsilon75M$. For $\epsilon=1$, their deficit
from the maximum $75(M+1)$ is $60$; the possible positive
individual deficits are $50,60,70,80,90,100,150$.
Thus exactly one term is $15$ and the others are $75$.
For $\epsilon=-1$, the increment above the minimum
$-75(M+1)$ is $90$, and the same list of positive increments
forces one term $15$ and all others $-75$.
Restoring the factor four gives the claimed $T,Z$ and the
remaining $M$ degree-three-hundred coordinates.
\par\smallskip

\noindent\emph{The contradiction.}
Claim~\ref{paper:no-ten-maximal-companions} gives a homogeneous
full two-block $\mathcal B$ and the corrected degree-$12$ column
\[
 w=6(W-I)\mathbf e_X
   =-\sum_{i=1}^4\mathbf e_{X_i}-\mathbf e_T
      +\mathbf e_Z+\epsilon\sum_{j=1}^6\mathbf e_{U_j},
 \quad d_T=12,\quad d_Z=60,\quad d_{U_j}=300.
\]
The first six labels are dyadically inactive. The diagonal
local modular relation at $X$ gives
\[
             \sum_{j=1}^6\theta_{U_j,2}=0:
\]
the squared coefficients from the four source labels are
$(25+1+1+1)/36$, and those from $T,Z$ are $2/36$; together
they already equal $W_{XX}=5/6$. A vanishing sum of two-power
roots splits into opposite pairs with equal multiplicities.
The six high labels therefore have an active phase $\phi$
of odd multiplicity $m=1$ or $3$.

Let $A$ be the sum of the $m$ coordinate vectors of this phase,
and put $u=3W_3A$. On the exact dyadic phase fiber $\phi$,
commutation with the complementary operators and the trivial
five-local lines of the degree-$300$ labels give
\begin{equation}\label{paper:eq:odd-phase-character}
                  \chi_X(V)=\epsilon'25u_V,
                         \qquad\epsilon'\in\{1,-1\}.
\end{equation}
Outside the fiber $u=0$. Character integrality gives integrality
at three; the foreign-denominator bound gives integrality at
two, because every $f_{U_j}=324$ has two-valuation two. There
are no other denominators. The local Weyl field assertion in
equation~\eqref{paper:lw:local-identities} places $W_3$, and hence $u$,
in a three-primary cyclotomic field. Principality of $X$ gives
$u_V\equiv\epsilon'd_V$ at every three-place on the support.

Three-inertia preserves the exact dyadic phase and transports
full two-blocks to full two-blocks. If a transported source
packet meets $A$, its full block equals $\mathcal B$, and its
entire phase-$\phi$ subset equals $A$. Distinct transported
packets are thus disjoint. The local Galois identities identify
the conjugates of $u$ with their images under $3W_3$, and also
with coordinate permutations preserving degrees. These
conjugates are orthogonal, and $\|u\|^2=9m$ is odd and at
most $27$.

Each source $U_j$ has a trivial five-local line and belongs to
the nonprincipal full two-block $\mathcal B$. Hence
\[
 (W_3d)_{U_j}=n(W_2W_5\mathbf e_0)_{U_j}=0,
\]
up to an immaterial normalization phase, and $\delta u=0$.
The inverse ternary Weyl transform of $d$ is supported on
ternarily inactive labels. The local Bruhat identity and the
common actual ternary twist on $\mathcal B$ from
Theorem~\ref{paper:lw:homogeneous} therefore give
\[
               |\delta(T_3u)|=3\delta A=900m=100\|u\|^2.
\]

Separation and divisibility of simple degrees initially leave
$50,75,100,150,225,300,450,900$ on this fiber. Ternary height
three is absent because it has a singleton three-block,
which cannot meet a source of degree $300$.
Suppose a supported $V$ has degree $225$, $450$ or $900$.
LCM integrality of $S_{XV}=\epsilon'300u_V$ gives
$u_V=3\beta$ with $\beta$ integral. Some source $U_j$ shares
its full three-block. At that source,
$\chi_{U_j}(X)=\epsilon'u_{U_j}$ has zero three-residue.
Equality of full-block character reductions consequently gives
\[
       \chi_V(X)=\epsilon'\frac{900}{d_V}\beta
                              \equiv0\pmod{\mathfrak p_3}.
\]
The displayed factor is $4$, $2$ or $1$, so $\beta$ has zero
three-residue. If rational, its coordinate costs at least
$81$ in $\|u\|^2$. If nonrational, its mean square is at
least two and its orbit contains at least two conjugates,
costing at least $9\cdot2\cdot2=36$. Both exceed $27$.
Only low degrees $50,100$ and high degrees $75,150,300$
remain. All hypotheses of Lemma~\ref{paper:odd-ternary-packet}
hold, giving the required contradiction.
\end{proof}


\subsection{When no twist order is divisible by six}

\begin{lemma}[A saturated foreign-fixed coordinate cannot occur]
\label{paper:ns-foreign-fixed-zero}
Work under two--three separation at dimension $2^6 3^6 5^4$,
with corrected dyadic operators $W=\gamma^3W_2$, $U=T_2$.
Let $Y$ be a degree-twenty label whose five-inertia orbit has
at least four members. If $2W\mathbf e_Y$ is integral and
has four nonzero coefficients, each $1$ or $-1$, then its
support contains no label fixed by five-inertia.
In particular it contains neither the vacuum nor a label of
five-height two.
\end{lemma}
\begin{proof}
Suppose a fixed label $Z$ occurs. Commutation with five-inertia
makes $W_{Z,Y'}=W_{Z,Y}=\pm1/2$ on the full orbit of $Y$.
Row norm one forces the orbit to have exactly four members,
and exhausts the row on that orbit. All of those labels are
dyadically inactive. Hence $UW^{-1}\mathbf e_Z=W^{-1}\mathbf e_Z$,
so $(WUW^{-1})\mathbf e_Z=\mathbf e_Z$.
The local Bruhat relation gives
$W^{-1}\mathbf e_Z=\theta_{Z,2}^2\mathbf e_Z$.
This is incompatible with its support on four distinct orbit labels.
The vacuum is fixed by inertia. Five-unit codegree makes every
five-height-two label fixed by five-inertia as well.
\end{proof}

\begin{lemma}[Five-prime lattice energy]
\label{paper:ns-five-lattice}
Let $F_b=\mathbb Q(\zeta_{5^b})$, for $b=1,2$, and let
$\mathfrak p_b$ be its unique prime above five. For nonzero
$t\in\mathfrak p_b$, write $E_{F_b}(t)=\operatorname{Tr}_{F_b/\mathbb Q}(t\bar t)$.
Then $E_{F_1}(t)\ge10$ and $E_{F_2}(t)\ge40$.
\end{lemma}
\begin{proof}
For $t=\sum_{j=0}^4c_j\zeta_5^j$ one has
$E_{F_1}(t)=5\sum c_j^2-(\sum c_j)^2$.
Membership in $\mathfrak p_1$ makes the coefficient sum a
multiple of five. Subtracting a common integer makes that sum
zero, so a nonzero coefficient vector has squared norm at least
two. This proves the first assertion.

The unique expansion $t=\sum_{j=0}^4\zeta_{25}^j a_j$, with
$a_j\in\mathbb Z[\zeta_5]$, gives
$E_{F_2}(t)=5\sum_jE_{F_1}(a_j)$ by relative trace orthogonality.
Each nonzero coefficient has energy at least four. Two such
coefficients therefore cost forty; a single nonzero coefficient
must be in $\mathfrak p_1$ and costs at least fifty.
\end{proof}

\begin{proposition}[The quintic packet obstruction]
\label{paper:quintic-packet-obstruction}
Let $u$ be a vector of algebraic integers whose coefficient field
is either $\mathbb Q(\zeta_5)$ or $\mathbb Q(\zeta_{25})$.
Suppose its distinct Galois conjugates are orthogonal, and the
five-primary Galois group also acts by coordinate permutations.
Let these permutations preserve positive integer weights $d_V$
and act on five-primary roots $\vartheta_V$ by the square action.
Partition the support into low coordinates, with $d_V\le108$,
and high coordinates, with $d_V\le540$. Suppose every high
coefficient has residue zero at five. Then the following
conditions are incompatible:
\[
 \|u\|^2=25,\qquad \sum_Vd_Vu_V=0,\qquad
 \sum_Vd_V\vartheta_V^a u_V=2700\tau^{-a^{-1}}
                  \quad(5\nmid a),
\]
where $\tau$ is a five-primary root and inverses are taken at
a common five-power level. An overall common sign on the last
right side makes no difference.
\end{proposition}
\begin{proof}
Write $M(\alpha)$ for normalized mean squared modulus and
$t=[\mathbb Q(u):\mathbb Q]\in\{4,20\}$.
If a coefficient $\alpha$ has field degree $e$, averaging the
orthogonal conjugate columns over the kernel fixing that field
has squared norm $25e/t$. A coordinate orbit of size $r$
containing $\alpha$ retains all its entries under this average.
Its energy is $rM(\alpha)$, so
\begin{equation}\label{paper:quintic-packet-projection}
                      rM(\alpha)\le25e/t.
\end{equation}
This argument permits any repetitions of coefficient conjugates.
Every nonzero algebraic integer has $M\ge1$, so each
Galois-stable part of the support has cardinality at most its
squared norm.

If $t=20$, a nonzero coefficient of residue zero at five has
$M\ge2$, by Lemma~\ref{paper:ns-five-lattice}.
But $r\ge e$, whereas \eqref{paper:quintic-packet-projection}
gives $rM\le5e/4$. Thus no high coefficient occurs.

Suppose $t=4$. A nonzero high coefficient has $M\ge5/2$.
If its coordinate were moved by wild five-inertia, its orbit
would have $r\ge5e$, contradicting
\eqref{paper:quintic-packet-projection}. Consequently every high
coordinate orbit has size $r=1,2$ or $4$, and its actual
five-primary twist has order at most five. Indeed an order
at least twenty-five has square-action orbit at least ten.

Let $T_O=\sum_{V\in O}u_V$ be a high orbit trace. It is a
rational integer divisible by five: its full cyclotomic trace
has that property, and the relative index is prime to five.
The rational projection of $u$ has squared norm $25/4$.
If $T_O\ne0$, its contribution $T_O^2/r$ is at least $25/4$.
Equality exhausts this projection, leaving $O$ as the only
orbit with nonzero trace. Since its degree is constant, this
contradicts the zero weighted sum. Thus every high orbit has
trace zero.

Every nonzero high orbit now has energy at least ten.
For an orbit of size four this is the five-prime lattice bound;
for one of size two its coefficient is a nonzero integer
multiple of $\sqrt5$, which gives the same bound. An orbit
of size one has zero coefficient. Its twist moment is purely
imaginary for every $a$. For an orbit of size four, central
inversion pairs conjugate coefficients at the same twist;
the two pair sums are real and opposite, and their twists
are inverse under the torus at the unit two. The size-two
case is the same identity $c(\vartheta^a-\vartheta^{-a})$.

If any high coordinate occurs, the low part has energy at
most fifteen and at most fifteen coordinates. Hence
\[
 \left|\operatorname{Re}\sum_Vd_V\vartheta_V^a u_V\right|
 \le108\sqrt{15\cdot15}=1620.
\]
For $\tau=1$ the required real part has absolute value $2700$.
For any other five-primary $\tau$, a unit exponent can be
chosen nearest the negative real axis, giving real part of
absolute value at least $2700\cos(\pi/5)>1620$.
This contradiction removes the high part also when $t=4$.

Only low coordinates remain. Their support size is at most
twenty-five, so the moment of absolute value $2700$ forces
equality in
\[
 2700\le108\sum_V|u_V|
       \le108\sqrt{|\operatorname{supp}u|}\,\|u\|
       \le2700.
\]
There are therefore twenty-five coordinates, all of weight
$108$, all with coefficients of modulus one, and all products
$\vartheta_Vu_V$ equal. Central inversion preserves each twist
and conjugates each coefficient while permuting the support.
Thus every coefficient is real, hence $1$ or $-1$.
The twists have odd order, so the equality of their products
forces all these signs to agree. Their weighted sum cannot
vanish, the final contradiction.
\end{proof}

\begin{corollary}[Exclusion of a forced degree-twenty piece]
\label{paper:quintic-four-piece-excluded}
Suppose a forced piece of
Claim~\ref{paper:ns-degree-twenty-four-piece} occurs,
with a principal-five label $Y$ of degree twenty, a common
five-primary twist, and full field $\mathbb Q(\zeta_5)$ or
$\mathbb Q(\zeta_{25})$. Its dyadically active degree-$540$
label $H$ has three-unit codegree. These data are contradictory,
without a global two--three separation assumption.
\end{corollary}
\begin{proof}
Put $u=5W_5\mathbf e_H$. Its codegree $100$ gives global
integrality and norm twenty-five. Its coefficient field is
the five-local character field of $H$, and its distinct
conjugates are orthogonal columns. On the actual dyadic
phase of $H$, the unique matching high coordinate in the
displayed four-piece matrix gives
\[
                       \chi_Y(V)=\pm27u_V.
\]
Here the complementary ternary Weyl fixes $H$ because its
codegree is a three-unit. Principality of $Y$ implies zero
five-residue at all five-divisible degrees.

The column is supported on dyadically active labels. A
five-height-two label has a singleton five-block and cannot
occur. Since $d_V\mid5400$, five-free supported degrees
are at most $108$ and five-divisible ones at most $540$.
The invariant four-piece excludes the vacuum from the
dyadic column at $H$. The ternary Weyl fixes $H$, so
\[
 (W_5d)_H=n(W_2W_3\mathbf e_0)_H=0.
\]
The local Bruhat identity supplies the remaining moment
identity in Proposition~\ref{paper:quintic-packet-obstruction},
with $5d_H=2700$. That proposition yields a contradiction.
\end{proof}

\begin{theorem}[Exclusion of the no-six graph]
\label{paper:no-six-excluded}
The activity graph $\{10,15\}$ cannot occur at $n=5400$.
\end{theorem}
\begin{proof}
Assume two and three are separated. We derive a four-label local
piece and apply the quintic packet obstruction. The piece is a
coordinate invariant space; no fusion closure is used.

\begin{proofclaim}[The degree-twenty field under two--three separation]
\label{paper:no-six-degree-twenty-full-field}\label{paper:no-six-degree-twenty-field}
In an ordinary perfect fusion-simple integral modular category
of dimension $2^6 3^6 5^4$, suppose no actual twist is divisible
by six. Every degree-twenty label has full character field
contained in $\mathbb Q(\zeta_{25})$.
For a compulsory degree-twenty principal-five witness,
\[
             K_Y=\mathbb Q(\zeta_5)\quad\text{or}\quad
             K_Y=\mathbb Q(\zeta_{25}).
\]
\end{proofclaim}
\noindent\emph{Proof of the claim.}
Every dyadically active degree is divisible by $27$, so $Y$ is
inactive. As in Proposition~\ref{paper:lw:isolated}, the dyadic anomaly
has $\gamma=\pm1$. Put $W=\gamma^3W_2$ and $U=T_2$.
Then $WUW=U^{-1}WU^{-1}$, $W^2=H_{-1,2}$ and
$\overline W=W^{-1}$.

The two-part of the conductor of $K_Y$ divides four because
$v_2(f_Y)=2$. Local Galois covariance therefore places the
$Y$-column of $W$ in $\mathbb Q(i)$.
Square-integrality gives $2W\mathbf e_Y$ integral above two.
The foreign-denominator theorem gives integrality at five,
since $v_5(f_Y)=2$, and gives global integrality at three,
since all dyadically active labels have three-unit codegree.
All other foreign primes have unit codegree at $Y$.
Thus $u=2W\mathbf e_Y\in\mathbb Z[i]^I$.
Lemma~\ref{paper:gaussian-column} makes this column rational;
local Galois covariance then gives dyadic unramifiedness and
\[
                         u\in\mathbb Z^I,\qquad\|u\|^2=4.
\]

Suppose now that $K_Y$ ramifies at three. The full three-local
group is trivial on the moving space of the corrected dyadic
group. If $Y_1,\ldots,Y_t$ is the three-inertia orbit of $Y$,
then $t\ge2$ and
\[
                  W\mathbf e_{Y_i}=\mathbf e_{Y_i}+v
\]
for one vector $v$ fixed by the full three-local group.
Symmetry and unitarity give $v_{Y_i}=-l/2$ with an integer
$l>0$ and $tl\le4$. Equality would leave the Weyl column
entirely in the inactive space, giving an invariant Weyl line,
which is impossible because $2\mid f_Y$.
Thus $t=2$ or $3$ and $l=1$.
The source-orbit coordinates of $u$ are $2\delta_{i1}-1$;
the outside squared norm is $4-t$.

Put $w=2v$. The local relations and the fixed source give
\[
 Ww=-w,\qquad WUw=(U^{-1}-1)w,\qquad
 \sum_Z d_Zw_Z\theta_{Z,2}^{-1}=0,\qquad
 \sum_Z u_Z^2\theta_{Z,2}=2.
\]
The weighted identity follows by pairing the middle relation
with $d$, since $W^{-1}d$ is supported on inactive labels;
the last identity is the source diagonal modular relation.

If $t=3$, the unique outside coefficient is $\varepsilon=\pm1$.
Its twist is $-1$, and the weighted identity gives
$d_A=3d_Y=60$, contradicting $27\mid d_A$.
If $t=2$, the two outside coefficients are $\varepsilon,\eta\in\{1,-1\}$
and their twists are opposite. If both are active, they are
$\rho,-\rho$ with $\rho\notin\{1,-1\}$; the weighted identity
\[
                   -40+\rho^{-1}(\varepsilon d_A-\eta d_B)=0
\]
is impossible by its imaginary and real parts.
Otherwise one outside label $Z$ is inactive and the other has
twist $-1$. Hence
\[
                   \varepsilon d_Z-40\in27\mathbb Z.
\]
This makes $d_Z$ prime to three. A label of dyadic height three
has an invariant dyadic Weyl line and cannot occur in this column.
Square divisibility and exclusion of nonunit prime-power degrees
therefore give $d_Z\in\{1,10,20,50,100\}$.
Neither sign of any of these integers is $40$ modulo $27$.
Thus $K_Y$ is unramified at three as well.

The exact conductor bound now gives
$K_Y\subseteq\mathbb Q(\zeta_{25})$.
For the compulsory principal-five witness, the complementary
weight-four crossing is ordinary under two--three separation.
Its value $1+3\xi$, with $\xi$ primitive of order five, gives
$\mathbb Q(\zeta_5)\subseteq K_Y$.
Since $\mathbb Q(\zeta_{25})/\mathbb Q(\zeta_5)$ has prime degree
five, precisely the two stated possibilities remain.
\par\smallskip

\begin{proofclaim}[A forced four-label dyadic piece]
\label{paper:ns-degree-twenty-four-piece}
Suppose no actual twist is divisible by six, and let $Y$ be a
compulsory degree-twenty principal-five witness. There are distinct
labels $Y,Z,A,B$ of respective degrees
\[
                           (20,20,540,540)
\]
whose coordinate space is invariant under the full corrected
dyadic local group. For signs $a,\eta\in\{1,-1\}$, its Weyl
matrix, in this order, is
\[
 W=\frac12\begin{pmatrix}
 1&1&a&a\\
 1&1&-a&-a\\
 a&-a&\eta&-\eta\\
 a&-a&-\eta&\eta
 \end{pmatrix}.
\]
For $\eta=-1$, the actual dyadic twists are
$(1,1,i,-i)$ up to interchanging $A,B$; for $\eta=1$ they
are $(1,1,-1,1)$ up to that interchange.
All four labels have the same full character field, equal to
$\mathbb Q(\zeta_5)$ or $\mathbb Q(\zeta_{25})$.
The corrected complementary vacuum coefficients $(Wd)_V/d_V$
at $Y,Z$ are $1+27a$ and $1-27a$, respectively; in particular
both labels belong to $P_{2'}$.
\end{proofclaim}
\noindent\emph{Proof of the claim.}
Claim~\ref{paper:no-six-degree-twenty-field} gives $K_Y=\mathbb Q(\zeta_5)$ or
$\mathbb Q(\zeta_{25})$ and $u=2W\mathbf e_Y\in\mathbb Z^I$
with squared norm four. A single coefficient of absolute value
two would put $W\mathbf e_Y$ in one twist fiber; Bruhat would
then make its Weyl line invariant, which is impossible.
Thus $u$ has four coefficients $\pm1$.
Lemma~\ref{paper:ns-foreign-fixed-zero} applies.

Every supported label has dyadic height at most two, because a
label of dyadic-unit codegree has an invariant Weyl line.
Every supported active degree is divisible by $27$, by separation.
The lemma excludes five-height two. Together with prime-power
exclusion, this gives the active degree list
\[
                        54,108,135,270,540.                 \tag{*}
\]
The remaining supported degrees are among
\[
 6,10,12,15,18,20,30,36,45,54,60,90,108,135,180,270,540.     \tag{**}
\]

For any dyadically inactive degree-twenty label $V$, its corrected
Weyl column satisfies
\[
 \sum_T (2W_{TV})^2\theta_{T,2}=4W_{VV},\qquad
 (Wd)_V/d_V=\kappa_V\in2\mathbb Z.                         \tag{***}
\]
The first right side is $2(2W_{VV})$; the second identity uses
$\kappa_V=\gamma^3k_{2',V}$, including the common correcting sign.
For $w=2(W-I)\mathbf e_V$, one also has
\[
                 \sum_T d_Tw_T\theta_{T,2}^{-1}=0.          \tag{****}
\]
Write $c=u_Y\in\{-1,0,1\}$. If $c=-1$, the other three twists
are all $-1$. The inactive part of $w$ is then $-3\mathbf e_Y$.
But the moving dyadic space, and hence this inactive part, is
fixed by the full three-local group. This makes $Y$ a
three-local invariant Weyl line, contrary to $3\mid f_Y$.

Suppose $c=0$. The four twists sum to zero, so are two opposite
pairs of dyadic roots. If none is inactive, normalized dyadic
trace in (****) gives $-40=0$. If exactly one is inactive,
the same identity gives $\varepsilon d_T\equiv40\pmod{27}$.
This makes $d_T$ prime to three, and (**) leaves $10$ or $20$,
neither of whose signs has that residue.
Thus there are two inactive coordinates and two of twist $-1$.
Let $b$ be the signed degree sum on the two active coordinates.
Then (****) gives inactive signed sum $40+b$, while (***) gives
\[
                         \kappa_Y=1+b/20\in2\mathbb Z.
\]
Since $27\mid b$ and $|b|\le1080$, this forces $b=\pm540$.
List (*) makes both active degrees $270$, with the same sign.
The inactive signed sum is therefore $580$ or $-500$.
Neither occurs as a signed sum of two degrees in (**): a term
$540$ would require a term $40$, which is absent; without $540$
the only sum of absolute value at least $500$ is $270+270=540$.
This excludes $c=0$.

Hence $c=1$. The other three twists sum to one, so they are
$1,\rho,-\rho$. First let $\rho$ be nonreal. Write their labels
as $Z,A,B$, respectively. The imaginary and real parts of (****)
force $d_Z=20$, its coefficient $+1$, and equal active dimensions
and equal active signs $a$. Thus
\[
                    W\mathbf e_Y=(Y+Z+aA+aB)/2.
\]
The even complementary coefficient gives
$1+a d_A/20\in2\mathbb Z$. List (*) forces $d_A=d_B=540$.
The central dyadic torus fixes all four labels: it fixes $Y$,
then $Z$, and the two distinct active twist fibers separately.
Thus their Weyl columns are real. Put $p=a(A+B)$ and $r=a(A-B)$.
The modular relation gives
\[
                  W(Y+Z)+\rho Wr=Y+Z+\rho^{-1}r.
\]
Its imaginary part gives $Wr=-r$. Orthogonality then forces
$\operatorname{Re}\rho=0$. Consequently $\rho=\pm i$,
$W(Y+Z)=Y+Z$, and $Wp=Y-Z$, proving the invariant matrix
with $\eta=-1$.

Now suppose $\rho$ is real. There is one active label $A$ of
twist $-1$ and two inactive outside labels. Equation (****)
and the even complementary coefficient give $d_A=540$ and
inactive signed sum $20\pm540$. List (**) leaves only a degree
$20$ label $Z$ with coefficient $+1$ and a degree $540$ label
$B$ with the same sign $a$ as $A$. Hence
\[
 WY=(Y+Z+aB+aA)/2,\qquad WA=(aY-aZ-B+A)/2.                 \tag{*****}
\]
The second identity follows from Bruhat, using the inactive
vector $-Y+Z+aB$ and the active vector $aA$.

It remains to prove invariance, which is not assumed. Dyadic
inertia fixes $B$, since it fixes $Y$ and preserves the coefficients,
degrees and twist fibers of $WY$. As $f_B=100$, the same column
integrality bounds give $2WB\in\mathbb Z^I$ of squared norm four.
Its known coefficients at $Y,A$ are $a,-1$.
Its diagonal coefficient cannot be $-1$, because that would make
all three outside twists $-1$, whereas $Y$ is inactive.
If its diagonal is $+1$, the remaining coordinate is inactive;
(****), with source $B$, gives its degree $20$ and coefficient
$-a$. Orthogonality with $WY$ identifies it as $Z$.
Thus $WB=(aY-aZ+B-A)/2$.

If instead its diagonal is zero, the remaining two twists are
opposite. The weighted identity excludes a nonreal opposite pair,
so write $WB=(aY-A+\varepsilon C+\delta D)/2$, with $C$
inactive and $D$ of twist $-1$. Orthogonality with $WY$ gives
$C\ne Z$. Applying $W^2Y=Y$ to (*****) yields
\[
                 WZ=(Y-aA-a\varepsilon C-a\delta D)/2.
\]
This diagonal-zero degree-twenty column has an active coordinate
of degree $540$. The integral complementary coefficient forces
$20\mid d_D$. Separation and the height bound leave
$d_D=540$ or $2700$; the latter has the excluded codegree four.
But $d_D=540$ makes the complementary coefficient
$1-27a(1+\delta)$ odd, a contradiction. Thus zero diagonal is
excluded. The obtained $WB$ and $W^2Y=Y$ now give
$WZ=(Y+Z-aA-aB)/2$, proving the invariant matrix with $\eta=1$.

All four fields are unramified at two. The degree-twenty fields
are unramified at three by Claim~\ref{paper:no-six-degree-twenty-field}; the degree-$540$ fields are
unramified there because their codegree is $100$.
The five-inertia stabilizer of $Y$ fixes all four support labels,
using their degrees and distinct dyadic twist fibers. Conversely,
the column of either high label has its two degree-twenty
coordinates with opposite signs, so its stabilizer fixes $Y$.
The column of $Z$ similarly fixes the two high labels.
Their stabilizers, and hence full fields, coincide.
Finally the matrix gives $Wd=(20+540a,20-540a,0,0)$ on this
piece, proving the complementary coefficients and the stated
membership in $P_{2'}$.
\par\smallskip

\noindent\emph{The contradiction.}
Claim~\ref{paper:ns-degree-twenty-four-piece} supplies a
forced piece with degrees $(20,20,540,540)$. Its dense dyadic
matrix makes the actual five-primary twist common to the
four labels. Corollary~\ref{paper:quintic-four-piece-excluded}
contradicts this piece in either of its possible coefficient
fields.
\end{proof}


\subsection{The full graph: inactive witnesses and reciprocity}


\begin{lemma}[At least four dyadic blocks in a compulsory fifth orbit]
\label{paper:compulsory-dyadic-block-separation}\label{paper:compulsory-fifth-block-separation}
Let $X\in P_3$ and $Y\in P_5$ be a compulsory crossing at
$D=5400^2$, with $d_X=4\cdot3^j$, $d_Y=20$, and
$j\in\{1,2\}$.  Write $L=20\cdot3^j$ and $z=S_{XY}/L$.
Assume one of the following alternatives
\[
 z=1+3\xi,
 \qquad\hbox{or}\qquad
 z=-\sum_{t=1}^4\eta_t\xi^t,
 \quad\{\eta_t\}=\{1,1,\omega,\omega^2\},
\]
where $\xi$ and $\omega$ have orders five and three.
For either endpoint, the stabilizer of its full geometric
dyadic block inside fifth inertia acts trivially on $\mu_5$.
Consequently the fifth-inertia orbit of the endpoint meets at
least four distinct full dyadic blocks. More precisely, its
orbit on blocks has size $4\cdot5^r$ for some $r\ge0$.

In particular, when the endpoint field is
$\mathbb Q(\xi)$, $\mathbb Q(\xi,i)$, or
$\mathbb Q(\zeta_{15})$, its four-label fifth-inertia orbit
meets four different full dyadic blocks.
\end{lemma}
\begin{proof}
First recall the all-places consequence of full-block equality.
Suppose labels $A,B$ have equal character reductions at one
place $\mathfrak p$ above two, on every basis element. For any
Galois automorphism $\tau$, argument-side covariance gives
\[
 \tau\bigl(\chi_A(V)-\chi_B(V)\bigr)
       =\chi_A(\tau V)-\chi_B(\tau V).
\]
The right side vanishes at $\mathfrak p$ for every $V$.
Thus the original difference vanishes at $\tau^{-1}\mathfrak p$.
All places above two are Galois conjugate, so full-block
equality holds at every such place. This also shows that
Galois permutations preserve the block partition.

Let $\sigma$ be a fifth-inertia element stabilizing the
dyadic block of $X$, and write $\sigma(\xi)=\xi^a$,
$a\in(\mathbb Z/5)^\times$. Then $X$ and $\sigma X$
have the same reductions at every two-place. Evaluation at
$Y$ gives
\[
 \chi_{\sigma X}(Y)-\chi_X(Y)
                =5\bigl(\sigma z-z\bigr).
\]
The scalar five is a two-unit. The number $\sigma z-z$
belongs to $E=\mathbb Q(\zeta_{15})$, which is unramified
at two. Vanishing at all its two-places therefore implies
\[
                         \sigma z-z\in2\mathcal O_E.
\]
In the ordinary case this difference is
$3(\xi^a-\xi)$. Reduction on roots of order five is
injective in characteristic two, so this is impossible
unless $a=1$.

In the signed case, fifth inertia fixes every $\eta_t$.
The four elements $\xi,\xi^2,\xi^3,\xi^4$ form an
integral basis of $\mathcal O_E$ over
$\mathcal O_{\mathbb Q(\omega)}$. Indeed
$\mathcal O_E=\mathcal O_{\mathbb Q(\omega)}[\xi]$,
and the indicated basis differs unimodularly from the usual
power basis by $1+\xi+\xi^2+\xi^3+\xi^4=0$.
The displayed divisibility thus gives
\[
                  \eta_{a^{-1}t}-\eta_t
                      \in2\mathcal O_{\mathbb Q(\omega)}
                  \qquad(t=1,2,3,4).
\]
Distinct roots in $\mu_3$ remain distinct modulo two.
Consequently $a$ preserves the coefficient pattern exactly.
The unique position with coefficient $\omega$ must be fixed,
so $a=1$.

For the $Y$ block use instead
\[
 \chi_{\sigma Y}(X)-\chi_Y(X)
            =3^j\bigl(\sigma z-z\bigr),
\]
whose rational prefactor is again odd. The same argument applies.
Fifth inertia is cyclic of order $4\cdot5^h$ at a sufficiently
large splitting level. Its block stabilizer lies in the
kernel of reduction to $(\mathbb Z/5)^\times$, so its index
has the asserted form. The three listed endpoint fields have
fifth-inertia orbits of length exactly four.
\end{proof}

Full-block equality compares every basis element. Argument-side Galois
covariance then gives equality at every two-place, which is the input
used in the integral-basis argument above.


\begin{lemma}[Separation on an exact odd-twist fiber]
\label{paper:no-thirty-local-separation}
Suppose no actual simple twist is divisible by thirty.
Then $\gamma=\pm1$. On an exact nontrivial three-primary
twist fiber, the corrected dyadic and five-local moving spaces
are orthogonal; every dyadically active coordinate there has
five-unit codegree. On an exact nontrivial five-primary twist
fiber, the corresponding assertions hold for two and three.
The separated-field and local-integrality arguments apply
within each such fiber.
\end{lemma}
\begin{proof}
The compulsory fifteen-edge supplies a nonempty exact fiber
with nontrivial three- and five-primary twists. Its actual
dyadic twist is identically one. On this invariant coordinate
space the dyadic upper and Weyl operators are the scalars
$\gamma^{-1}$ and $\gamma^3$. Since the Weyl square is an
unsigned torus permutation, $\gamma^6=1$; its two-power order
therefore gives $\gamma=\pm1$.

Fix a nontrivial three-primary twist. On its fiber a nontrivial
actual dyadic twist forces the actual five-primary twist to be
one. The five-local upper is the identity on every such exact
dyadic fiber, so Bruhat makes the entire five-local group
trivial there. The Weyl-line criterion gives five-unit codegree
at each of its coordinate labels. Commutativity propagates
this trivial action to the moving dyadic space generated by
those coordinates. The same argument with the primes reversed
gives orthogonality of the two moving spaces. This is precisely
the operator argument for separated primes, restricted to an
invariant coordinate space. The foreign-integrality theorem
explicitly permits such a restriction. Interchanging three and
five proves the other assertion.

For clarity, the field argument also localizes: orthogonality
gives $(H_2-I)(H_5-I)=0$ on the first fiber. Because the tori
are unsigned permutations, if any five-torus moves a label,
every dyadic torus fixes it. The analogous statement holds
on the second fiber with three in place of five.
\end{proof}

\begin{proposition}[Inactive compulsory witnesses when thirty is absent]
\label{paper:no-thirty-inactive-witnesses}
Suppose no actual simple twist is divisible by thirty.
For every complementary pair of dyadic-height-two witnesses
\[
 X\in P_3,\quad d_X\in\{12,36\},\quad\theta_X\text{ three-primary},
 \qquad
 Y\in P_5,\quad d_Y=20,\quad\theta_Y\text{ five-primary},
\]
one has $\theta_X=\theta_Y=1$.
\end{proposition}
\begin{proof}
Both missing-edge graphs have already been excluded, so the
witnesses can be chosen at dyadic height two. Their full
dyadic blocks are nonprincipal and homogeneous. A signed
crossing would force a thirty-active constituent; thus the
crossing is ordinary and contains a primitive fifth root.
In particular the five-inertia orbit of $X$ has at least four
members.

Suppose first that $\theta_X\ne1$ and restrict to its exact
three-primary twist fiber. Lemma~\ref{paper:no-thirty-local-separation}
makes $X$ unramified at two and its dyadic Weyl column integral
at five. At degree thirty-six the foreign bound at three
makes $2W\mathbf e_X$ an integral column of norm four; the
foreign-orbit argument then exhausts that column on inactive
labels, a contradiction. Hence $d_X=12$.

At degree twelve, $6W\mathbf e_X$ is an integer column of
norm thirty-six. Its five-inertia orbit has size $t=4\cdot5^s$
because its field contains $\mathbb Q(\zeta_5)$, and its
conjugates coalesce on the moving dyadic space. If $-l$ is
their common correction coefficient, the orbit norm inequality
gives $tl\le12$, hence $t=4$. Write their common correction as
$w/6$. Homogeneous density makes each coefficient of the
column odd. The usual orbit norm inequality allows $l\le3$;
oddness excludes $l=2$, and $l=3$ would exhaust the block
on inactive labels. Therefore $l=1$. Its outside norm is
eight, so the full block consists of twelve labels and all
eight outside coefficients are $1$ or $-1$.

The signed dimension argument is unchanged on this fiber.
Every active companion has degree divisible by twenty-five,
and is at most three hundred. It forces the support and signs
\[
 w=-\sum_{i=1}^4\mathbf e_{X_i}-\mathbf e_T
       +\mathbf e_Z+\epsilon\sum_{j=1}^6\mathbf e_{U_j},
 \qquad (d_T,d_Z,d_{U_j})=(12,60,300),
\]
where $T,Z$ are dyadically inactive. The diagonal modular
relation makes the six high dyadic twists sum to zero.
Their antipodal multiplicities have sum three, so some active
phase occurs with odd multiplicity $m=1$ or $3$.
Let $A$ be that phase packet and put $u=3W_3\sum_{U\in A}\mathbf e_U$.

All needed packet identities remain global. The correction
$w$ is fixed by the full five-local group, and the source term
has dyadic twist one; hence on the selected active phase
$\chi_X(V)=\pm25u_V$. Thus $u$ is integral, has norm $9m$,
zero weighted sum and Bruhat moment of absolute value $900m$.
Three-inertia transports full two-blocks and preserves that
phase, so the transported packets are equal or disjoint and
the conjugate columns are orthogonal.

The higher-height exclusion uses no global separation:
if $v_3(d_V)=2$, lcm-integrality gives $u_V=3\beta$;
character congruence between $V$ and a source in its full
three-block, evaluated at $X$, makes $\beta$ have zero
three-residue. Its orbit costs at least thirty-six, too much
for $\|u\|^2\le27$. Ternary height three is excluded by
the defect-zero singleton property. Thus every supported
degree prime to three is at most $100$, and every other
supported degree is at most $300$.
Lemma~\ref{paper:odd-ternary-packet} uses exactly these bounds,
the prime-three residue and the identities just established;
it gives a contradiction for both possible values of $m$.
This proves $\theta_X=1$.

Suppose next that $\theta_Y\ne1$. On its exact five-primary
twist fiber the local dyadic and ternary actions are separated.
The Gaussian norm-four column argument and the foreign-orbit
argument of Claim~\ref{paper:no-six-degree-twenty-full-field}
therefore apply without change inside this invariant fiber.
They give
\[
 K_Y=\mathbb Q(\zeta_5)\quad\text{or}\quad
 K_Y=\mathbb Q(\zeta_{25}),\qquad
 2W\mathbf e_Y\in\mathbb Z^I.
\]
Every active dyadic companion on the fiber has degree divisible
by twenty-seven. The forced-four-piece proof consequently
applies there and produces degrees $(20,20,540,540)$, with
the same field at all four labels. Its dyadically active high
label has three-unit codegree. Corollary~\ref{paper:quintic-four-piece-excluded}
now gives a contradiction. Hence $\theta_Y=1$ as well.
\end{proof}


\begin{corollary}[A reciprocal degree restriction for signed buckets]
\label{paper:signed-bucket-gcd}
In a signed compulsory crossing with degrees $(12,20)$ or
$(36,20)$, every simple constituent $V$ satisfies
\[
                       4\mid d_V\quad\hbox{or}\quad5\mid d_V.
\]
In particular, degrees $6$, $18$, and $54$ cannot occur in its
actual phase buckets.
\end{corollary}
\begin{proof}
Every constituent has actual dyadic twist $-1$, whereas the
pure-three input $X$ has actual dyadic twist one. The lemma
at the pure-five principal input $Y$ excludes
$\gcd(d_V,20)\in\{1,2\}$. Among divisors of twenty, the
remaining possibilities are $4,5,10,20$, giving the assertion.
\end{proof}

\begin{proposition}[An odd witness for a thirty-active label]
\label{paper:triple-odd-witness}
Let $A$ have an actual twist divisible by thirty. There is a
nonunit $Z\in P_2$, with pure dyadic actual twist, such that
\[
 d_Z\in\{15,45\},\qquad S_{AZ}\ne0,
 \qquad K_Z\text{ is unramified at two}.
\]
Its actual twist has order dividing eight, and
$W_2\mathbf e_Z$ is rational. More precisely, if
$F=\mathbb Q(\chi_A(Z))$, one may choose $Z$ so that both
\[
 [K_Z:F]\quad\text{and}\quad
 \operatorname{Tr}_{F/\mathbb Q}(\chi_A(Z))
\]
are odd integers.
For $d_Z=45$ the vector $8W_2\mathbf e_Z$ is integral with
squared norm $64$; for $d_Z=15$ the vector
$24W_2\mathbf e_Z$ is integral with squared norm $576$.
\end{proposition}
\begin{proof}
Use principal-two orthogonality on its trivial complementary
twist fiber, as in \eqref{paper:lw:principal-fiber-orthogonality}:
\[
 1+\sum_{\substack{Z\ne\mathbf1,\ Z\in P_2\\
                         \theta_{Z,3}=\theta_{Z,5}=1}}
                    d_Z\chi_A(Z)=0.
\]
This set is stable under the full Galois group. For one full
label orbit, argument-side Galois covariance identifies its
character sum with
\[
 [K_Z:F]\operatorname{Tr}_{F/\mathbb Q}(\chi_A(Z)),
       \qquad F=\mathbb Q(\chi_A(Z)).
\]
Reducing the grouped integer identity modulo two supplies an
odd degree with an odd character sum. In particular the entry
is nonzero and both displayed factors are odd.

By square divisibility, an odd degree is $3^j5^k$.
Lemma~\ref{paper:family:2.4} excludes a nonunit prime-power degree,
so $j,k>0$. If $j=3$ or $k=2$, the corresponding formal
codegree is respectively a three-unit or a five-unit.
Lemma~\ref{paper:family:lw:weyl-line} would then make its entry against
the three- and five-active row $A$ vanish. Thus
$j\in\{1,2\}$ and $k=1$, giving fifteen or forty-five.

All character fields are abelian cyclotomic fields. If $F$
ramified at two, its dyadic inertia group would have even
order. At a prime above two, inertia acts trivially on the
residue field. Grouping an integral absolute trace into inertia
cosets would therefore make that trace even. This contradicts
the odd trace just obtained, so $F$ is unramified at two.
The dyadic inertia group of $K_Z$ then lies in
$\operatorname{Gal}(K_Z/F)$, whose order is odd, but inertia
has two-power order. It is trivial, proving the field claim.

All dyadic tori fix $Z$, so local Galois covariance makes its
dyadic Weyl column rational. If $\gamma$ is the dyadic anomaly,
the torus at the unit three fixes the upper eigenvalue
$\theta_Z/\gamma$ after raising it to the ninth power.
Hence $(\theta_Z/\gamma)^8=1$, and $\gamma^8=1$ gives
$\theta_Z^8=1$. This does not assert that the twist is nontrivial.

Finally the formal codegrees are respectively
$2^6 3^2 5^2$ at degree forty-five and $2^6 3^4 5^2$
at degree fifteen. Square-integrality at two and the foreign
denominator bounds of Theorem~\ref{paper:lw:columns} give the two
integral multiples of the rational columns. Their norms follow
from unitarity.
\end{proof}

\begin{remark}
Together with the compulsory witnesses of degrees $12$ or
$36$ in $P_3$ and degree $20$ in $P_5$, this gives three
pure-primary endpoints. Their additional pairings do not
inherit the two nonzero foreign residues of the original
weight-four crossing: normalization at the dyadic-principal
endpoint can divide a power of two. No additional residue
condition is asserted here.
\end{remark}


\begin{proposition}[A signed common witness has positive fifth defect]
\label{paper:signed-common-positive-fifth-defect}
Let $X\in P_3$ and $Y\in P_5$ be compulsory pure-primary
witnesses of degrees $4\cdot3^j$ and twenty, where
$j\in\{1,2\}$, and suppose their crossing is signed tame.
If a label $Z$ satisfies $S_{XZ}S_{YZ}\ne0$, then
\[
                              5\mid f_Z.
\]
This assertion does not require either endpoint twist to be
nontrivial and does not require $Z\in P_2$.
\end{proposition}
\begin{proof}
Suppose $5\nmid f_Z$. By the Weyl-line vanishing criterion,
the $Z$ row vanishes on every five-active label. Its nonzero
entry at $Y$ therefore forces $\theta_Y=1$, because the
twist of $Y$ is pure five-primary. The exact signed crossing
has the four relative phases
\[
                         -\eta_t\xi^t,
                         \qquad t=1,2,3,4,
\]
with $\xi$ primitive fifth and $\eta_t$ three-primary.
Since $X$ is five-inactive and $\theta_Y=1$, every actual
simple constituent $A$ of $XY$ has five-primary twist
$\xi^t\ne1$. Thus $\chi_Z(A)=0$ for every constituent.
Multiplicativity now gives
\[
             0=\chi_Z(XY)=\chi_Z(X)\chi_Z(Y),
\]
contradicting the two nonzero entries.
\end{proof}


\begin{proposition}[A joint dyadic witness for a signed crossing]
\label{paper:signed-joint-witness}
Let $X\in P_3$ and $Y\in P_5$ be complementary pure-primary
witnesses of degrees $d_X=4\cdot3^j$ and $d_Y=20$, where
$j\in\{1,2\}$, and suppose their compulsory crossing is signed.
There is a nonunit $Z\in P_2$ with pure dyadic actual twist such that
\[
 S_{XZ}S_{YZ}\ne0,\qquad v_2(d_Z)\le1,\qquad
 K_Z\text{ is unramified at two},\qquad \theta_Z^8=1.
\]
Put $b_Z=S_{XZ}S_{YZ}/16$ and $F_Z=\mathbb Q(b_Z)$.
The witness can be chosen so that both
\[
 [K_Z:F_Z]\quad\text{and}\quad
 \operatorname{Tr}_{F_Z/\mathbb Q}(b_Z)
\]
are odd. In particular $W_2\mathbf e_Z$ is rational.
If $\theta_X\ne1$, then
\[
                         d_Z\in\{10,15,18,30,45,90\}.
\]
\end{proposition}
\begin{proof}
Let $E=P_2\cap\{Z:\theta_{Z,3}=\theta_{Z,5}=1\}$.
Every simple constituent $A$ of the signed product $XY$ has
$\theta_{A,2}=-1$. Principal-fiber orthogonality and symmetry give
\[
 \sum_{Z\in E}d_Z^2\chi_Z(A)
 =d_A\sum_{Z\in E}d_Z\chi_A(Z)=0.
\]
Sum with the multiplicities of $A$ in $XY$ and use multiplicativity:
\[
 \sum_{Z\in E}S_{XZ}S_{YZ}=0,
 \qquad \sum_{Z\in E}b_Z=0.
\]
LCM integrality makes each $S_{XZ}/4$ and $S_{YZ}/4$ integral,
so every $b_Z$ is integral. The vacuum contribution is
$b_{\mathbf1}=d_Xd_Y/16=5\cdot3^j$, an odd integer.

The set $E$ is Galois-stable. Argument-side covariance groups
its $Z$-orbit contribution as
$[K_Z:F_Z]\operatorname{Tr}_{F_Z/\mathbb Q}(b_Z)$.
For $4\mid d_Z$ one has
\[
 b_Z=\frac{d_Z^2}{16}\chi_Z(X)\chi_Z(Y).
\]
The rational prefactor is an integer. Principality at two
makes both character factors vanish at every place above two,
because $d_X$ and $d_Y$ are even. The integral orbit sum
therefore vanishes at every such place and is an even rational
integer. The odd vacuum term consequently requires a nonunit
orbit with $v_2(d_Z)\le1$ and odd displayed orbit contribution.
Both its relative index and its absolute trace are odd, and its
entry is nonzero.

If $F_Z$ ramified at two, its dyadic inertia group would have
even order and act trivially on the residue field. Grouping its
absolute integral trace into inertia cosets would make the trace
even. Thus $F_Z$ is unramified at two. The dyadic inertia group
of $K_Z$ lies in $\operatorname{Gal}(K_Z/F_Z)$, of odd order;
being a two-group, it is trivial. This proves the field assertion.
The rationality of the Weyl column and the twist bound now follow
from the same local Galois covariance and square-action argument
as in Proposition~\ref{paper:triple-odd-witness}.

If $\theta_X\ne1$, nonzero $S_{XZ}$ and the Weyl-line vanishing
criterion force $v_3(d_Z)\le2$. Proposition~\ref{paper:signed-common-positive-fifth-defect}
gives $v_5(d_Z)\le1$ without an activity assumption on $Y$. Combine these bounds with
$v_2(d_Z)\le1$, square divisibility, the absence of nonunit
prime-power degrees, and Theorem~\ref{paper:modular:no-degree-six}.
This leaves exactly the six displayed degrees.
\end{proof}

\begin{remark}[The remaining parity obstruction]
The result gives a common witness for the two endpoints, which the
single-label odd-trace argument does not do. It does not show that
the common witness has odd degree. At dyadic height one,
$\chi_Z(X)$ and $\chi_Z(Y)$ are divisible by two, but after
dividing $S_{XZ}S_{YZ}$ by sixteen the two powers of two in
$d_Z^2$ exactly cancel this divisibility. Such orbit contributions
need not be even. Therefore the degree-fifteen-or-forty-five
conclusion of Proposition~\ref{paper:triple-odd-witness} cannot
be imposed on this common witness without a further argument.
\end{remark}

\begin{corollary}[Nine degrees for the signed joint witness]
\label{paper:signed-joint-all-degrees}\label{paper:signed-common-nine-degrees}
The witness in Proposition~\ref{paper:signed-joint-witness} can
always be chosen with
\[
                 d_Z\in\{10,15,18,30,45,54,90,135,270\},
\]
even when both endpoint twists are trivial. If
$d_Z=2^a3^b5^c$, then
\[
 m_ZW_2\mathbf e_Z\in\mathbb Z^I,
 \qquad m_Z=2^{3-a}3^{\max(2-b,0)}5^{\max(1-c,0)},
\]
with squared norm $m_Z^2$. Degrees ninety and two hundred
seventy have multiplier four and squared norm sixteen.
\end{corollary}
\begin{proof}
Combine $v_2(d_Z)\le1$, the positive fifth defect,
$d_Z\mid5400$, prime-power exclusion, and the absence of
degree six. These give exactly the displayed list.
The column is rational. Rotation integrality makes
$f_Z(W_2)_{VZ}^2$ an algebraic integer, so its rational
entries have denominators dividing $\sqrt{f_Z}$.
At each odd prime the strict foreign bound improves its
denominator exponent by one whenever that prime divides
$f_Z$. This gives $m_Z$; unitarity gives the norm.
\end{proof}

\begin{remark}[The surviving three-unit alternatives]
If the common witness has degree $54$, $135$, or $270$,
nonzero $S_{XZ}$ forces $\theta_X=1$. Its row vanishes
on the two signed phase buckets with nontrivial cubic
component. Only the two buckets with $\eta_t=1$ can
contribute to $\chi_Z(XY)$; their cancellation is not forced.
\end{remark}

\subsection{Indicator parity and principal dyadic columns}


\begin{proposition}[Principal dyadic parity from rotation indicators]
\label{paper:principal-dyadic-parity}
Assume $D=n^2$ and put $W=W_2$,
\[
 k_A=\frac{nW_{A0}}{d_A}\in\mathbb Z,
 \qquad \kappa=k_0=nW_{00}.
\]
Then
\[
                    k_A\equiv\kappa\pmod2
                    \qquad(A\in P_2).
\]
If $Z\in P_2$ has a character field unramified at two, its
column $W\mathbf e_Z$ is rational.  Suppose also that
$q_i=(n/d_Z)W_{iZ}$ is an integer for every $i$, as holds
when all the simple dimensions divide $n$.  Then
\[
 q_i\equiv d_i k_i\pmod2.
\]
In particular,
\[
 q_i\equiv\kappa d_i\pmod2\quad(i\in P_2),
 \qquad q_i=0\quad(i\notin P_2).
\]
The assertion $q_i\equiv\kappa d_i$ is only made on $P_2$.
\end{proposition}
\begin{proof}
Choose a positive CRT exponent $m$ which is one on the dyadic
level of the congruence representation and zero on the odd levels.
Write $y_Z=f_Ze_Z=\sum_V\overline{\chi_Z(V)}[V]$ and
$\Omega=De_0=\sum_Vd_V[V]$.  Principality of $Z$ means that
the coefficients of $y_Z-\Omega$ vanish at every place above two.
Since each $L_m^{ii}([V])$ is an algebraic integer, the same is
true of
\[
                     L_m^{ii}(y_Z)-L_m^{ii}(\Omega).
\]
The mixed rotation formula and the local CRT identity cancel
all root-of-unity factors for a diagonal insertion.  They give
\[
 f_Z|W_{iZ}|^2\equiv D|W_{i0}|^2\pmod{\mathfrak p_2}.
\]
At $i=0$ the two sides are $k_Z^2$ and $\kappa^2$.
These are rational integers, so the congruence implies
$k_Z\equiv\kappa\pmod2$.

Dyadic inertia fixes $Z$ under the unramifiedness assumption.
Local Galois covariance therefore makes $W\mathbf e_Z$ rational.
If $n/d_Z$ is an integer, rotation integrality makes $q_i^2$
an integer; a rational number with integral square is an integer.
The displayed indicator congruence becomes
\[
                     q_i^2\equiv(d_i k_i)^2\pmod2,
\]
and hence $q_i\equiv d_i k_i\pmod2$.  For $i\in P_2$
use the first assertion.  Outside $P_2$, local block invariance
makes both $W_{iZ}$ and $W_{i0}$ zero.  No assertion that all
odd-dimensional labels belong to $P_2$ is needed or made.
\end{proof}

\begin{corollary}[A parity dichotomy for the principal block]
\label{paper:principal-dyadic-dichotomy}
In the preceding setting the following alternatives hold.
\begin{enumerate}
\item If $\kappa$ is odd, every label in $P_2$ has trivial
      complementary actual twist.  Consequently
      $P_2=P_2\cap\{Z:\theta_{Z,3}=\theta_{Z,5}=1\}$
      in the three-prime problem.
\item If $\kappa$ is even, $Z\in P_2$ is unramified at two,
      and $m_ZW\mathbf e_Z$ is integral,
      with $(n/d_Z)/m_Z$ an odd integer, then
      $(m_Z/2)W\mathbf e_Z$ is integral and has squared norm
      $m_Z^2/4$.
\end{enumerate}
\end{corollary}
\begin{proof}
If $\kappa$ is odd, the first proposition makes every $k_A$
with $A\in P_2$ odd and nonzero.  Thus every such label occurs
in the vacuum column of $W_2$.  This column has trivial
complementary actual twists because all complementary twist
operators commute with $W_2$ and fix the vacuum.

For the second assertion, all coordinates $q_i$ in the
proposition are even: on $P_2$ this follows from evenness of
$\kappa$, and outside $P_2$ they are zero.  Division by the
odd integer $(n/d_Z)/m_Z$ preserves parity.  Unitarity gives
the squared norm after halving.
\end{proof}

\begin{lemma}[An integral Weyl column cannot be monomial at positive defect]
\label{paper:no-monomial-weyl-column}
Assume $p\mid f_Z$.  A rational integral column
$W_p\mathbf e_Z$ is impossible.  More generally, a monomial
column $W_p\mathbf e_Z=c\mathbf e_A$ is impossible.
\end{lemma}
\begin{proof}
If the column is monomial, the line $\mathbb C\mathbf e_A$
is invariant under the diagonal upper unipotent and under the
conjugate lower unipotent $W_pU_pW_p^{-1}$, because
$W_p^{-1}\mathbf e_A=c^{-1}\mathbf e_Z$.  Those unipotents
generate the local group, so that line is $W_p$-invariant.
Unitarity and $W_p\mathbf e_Z=c\mathbf e_A$ then force $A=Z$.
The resulting invariant coordinate line contradicts the
Weyl-line criterion at $p\mid f_Z$.  A rational integral
unit column has exactly one nonzero coordinate, equal to
$\pm1$, and is therefore monomial.
\end{proof}

\begin{corollary}[The even alternative gives four-point columns]
\label{paper:four-point-common-column}
For the signed common witnesses, write
$d_Z=2^a3^b5^c$ and
\[
 m_Z=2^{3-a}3^{\max(2-b,0)}5^{\max(1-c,0)}.
\]
The ratio $(5400/d_Z)/m_Z$ is odd.  If $\kappa$ is even,
all the joint-witness integral multipliers $m_Z$ can consequently
be halved.  In particular, for
$d_Z\in\{90,270\}$ the column has exactly four nonzero
entries, each equal to $1/2$ or $-1/2$.
Degree $450$ is excluded by
Proposition~\ref{paper:signed-common-positive-fifth-defect}.

The single odd witness for a thirty-active label also has its
multiplier halved: degree fifteen gives an integral
$12W_2\mathbf e_Z$ of squared norm $144$, and degree
forty-five gives an integral $4W_2\mathbf e_Z$ of squared
norm $16$.
\end{corollary}
\begin{proof}
The odd ratio is the product of a factor three when $b\le2$
and a factor five when $c\le1$.  Corollary~\ref{paper:principal-dyadic-dichotomy}
therefore applies. At degrees ninety and two hundred seventy,
$m_Z=4$, so the integer vector
$2W_2\mathbf e_Z$ has squared norm four.  Such a vector is
either a single $\pm2$ or four entries $\pm1$.  The first
case is excluded by Lemma~\ref{paper:no-monomial-weyl-column}, since the formal
codegree is even.  The final assertions apply the same parity
calculation to the odd witness multipliers $24$ and $8$.
\end{proof}

\begin{proposition}[The reciprocal parity bound on a rational column]
\label{paper:reciprocal-dyadic-parity}
Continue at $D=5400^2$.  Let $Z\in P_2$ have a rational
dyadic Weyl column, and put $a_A=v_2(d_A)$ for $A\in P_2$.
Every nonzero entry of this column satisfies
\[
 v_2(W_{AZ})\ge a_A-2
       \quad\hbox{if }\kappa d_Z\text{ is even}.
\]
If $\kappa d_Z$ is odd, every label of $P_2$ occurs in the
column, and one has the exact valuation
\[
                    v_2(W_{AZ})=a_A-3
                    \qquad(A\in P_2).
\]
In the even case, if the column has only dyadic denominators,
it has no nonzero coordinate of dyadic height two.
\end{proposition}
\begin{proof}
Apply the proof of Proposition~\ref{paper:principal-dyadic-parity}
with source $A$ and diagonal insertion indexed by $Z$.
Symmetry and rationality of $W_{AZ}$ make its integral
indicator value a rational integer, and give
\[
               f_A W_{AZ}^2\equiv d_Z^2k_Z^2
                              \equiv\kappa d_Z\pmod2.
\]
Since $v_2(f_A)=6-2a_A$ and a rational entry has integral
valuation, strict positivity of the left valuation in the even
case gives $v_2(W_{AZ})\ge a_A-2$.  In the odd case the
entry cannot vanish, and its valuation is determined by
$6-2a_A+2v_2(W_{AZ})=0$.

In the even case, a height-two coordinate is therefore
two-integral.  With only dyadic denominators it is an integer.
A nonzero such coordinate in a unit column would exhaust
its norm and make the column monomial, excluded by
Lemma~\ref{paper:no-monomial-weyl-column}, since $Z\in P_2$
has positive dyadic defect.
\end{proof}

\begin{corollary}[Bounds from the degree-forty-five odd witness]
\label{paper:forty-five-principal-support}
Suppose the unramified odd witness has degree forty-five.
If $\kappa$ is even, its integral vector
$4W_2\mathbf e_Z$ has squared norm sixteen and is supported
only on degrees of dyadic height at most one.

If $\kappa$ is odd, its $8W_2\mathbf e_Z$ column has
support exactly $P_2$, and
\[
              \sum_{A\in P_2}4^{v_2(d_A)}\le64.
\]
In particular $|P_2|\le64$.  If $e=[K_Z:\mathbb Q]$, the
existence in $P_2$ of a height-two label requires $e\le3$,
and the existence of a height-one label requires $e\le15$.
Prime avoidance consequently gives $e\le15$ unconditionally
in this odd-$\kappa$ degree-forty-five case.  If $e\ge4$,
the five-free nonunits in $P_2$ have degrees in $\{18,54\}$,
and the three-free nonunits have degrees in $\{10,50\}$;
both such sets are nonempty.
\end{corollary}
\begin{proof}
The even assertion combines the halved multiplier with the
preceding proposition.  In the odd case, each integral entry
$8W_{AZ}$ has exact two-valuation $v_2(d_A)$, so its square
is at least $4^{v_2(d_A)}$.  Sum the column norm.

For the final assertions let $\mathcal O$ be the full Galois
orbit of $Z$, of size $e$.  Its labels have the same degree,
are unramified at two, and satisfy the same column assertions.
For any $A\in P_2$ of height $a=1$ or $2$, each entry
$W_{AZ'}$, $Z'\in\mathcal O$, has absolute value at least
$2^{a-3}$.  The vacuum is outside $\mathcal O$, and its entry
$W_{A0}$ is nonzero because $k_A$ is odd.  The norm of the
$A$ row therefore gives
\[
                         e\,4^{a-3}<1.
\]
Thus $e<16$ at height one and $e<4$ at height two.
Prime avoidance in $P_2$ supplies both five-free and
three-free nonunits.  Prime-power exclusion makes their
dyadic height positive, and the principal height bound makes
it at most two.  This proves $e\le15$.  If $e\ge4$,
height two is impossible, and the two displayed degree lists
follow from square divisibility and the absence of degree six.
\end{proof}

A four-point Weyl column need not be the whole residue block or have
equal-degree support. Its integral halving uses the principal
congruence modulo two.


\begin{lemma}[Even principal parity excludes purely dyadic codegrees]
\label{paper:even-parity-no-pure-dyadic-codegree}
At $D=5400^2$, if $\kappa=5400(W_2)_{00}$ is even,
no simple label has purely dyadic formal codegree.
In particular degrees $675$, $1350$, and $2700$ are absent.
The remaining formal possibility, degree $5400$, is already
excluded by the positive unit contribution to global dimension.
\end{lemma}
\begin{proof}
If $f_A$ is a power of two, every odd local Weyl preserves
the coordinate line of $A$, by the Weyl-line criterion.
Thus $s\mathbf e_A$ is a root-of-unity multiple of
$W_2\mathbf e_A$. Taking the vacuum coordinate shows that
\[
                  k_{2,A}=\frac{5400(W_2)_{0A}}{d_A}
\]
is a root of unity. It is a rational integer by the integral
vacuum-coefficient theorem, so it is $1$ or $-1$. Its
nonzero vacuum entry puts $A$ in $P_2$. The uniform principal
parity $k_{2,A}\equiv\kappa\pmod2$ is a contradiction.
Since dimensions divide $5400$, the only degrees with purely
dyadic codegree are $5400$, $2700$, $1350$, and $675$.
\end{proof}

\begin{proposition}[The sparse common columns avoid the vacuum]
\label{paper:sparse-common-no-vacuum}
Suppose $\kappa$ is even. If the signed common witness has
degree $90$ or $270$, then its rational dyadic Weyl column
has exactly four nonzero entries, each $\pm1/2$, and its
vacuum coordinate is zero. Equivalently, $k_{2,Z}=0$.
\end{proposition}
\begin{proof}
The four-point assertion is
Corollary~\ref{paper:four-point-common-column}. Its entries have
only dyadic denominators, so reciprocal principal parity
excludes all supported degrees of dyadic height two or more.
Write its support as $\mathcal S$ and
\[
                W_2\mathbf e_Z
                     =\frac12\sum_{A\in\mathcal S}
                                         \varepsilon_A\mathbf e_A,
           \qquad\varepsilon_A\in\{1,-1\}.
\]
Suppose the vacuum belongs to $\mathcal S$.

Let $d=(d_A)_A$ and write $W_{2'}=W_3W_5$.
The local factors commute and the central torus fixes $d$,
so $W_2d=5400W_{2'}\mathbf e_0$. The foreign-prime
divisibility theorem gives $k_{2',Z}\in2\mathbb Z$.
Consequently
\begin{equation}\label{paper:sparse-dimension-sum}
        \sum_{A\in\mathcal S}\varepsilon_A d_A
             =2d_Zk_{2',Z}\in4d_Z\mathbb Z.
\end{equation}

The right side is even. Hence the four supported coordinates
have an even number of odd degrees. If all four were odd,
reduction modulo three would leave just the signed vacuum
term: every nonunit odd degree is divisible by fifteen.
This is impossible because $3\mid d_Z$. Thus there are
exactly two odd degrees and two even degrees in the support.

Write the nonunit odd degree as $d_A$. Reduction of
\eqref{paper:sparse-dimension-sum} modulo three and modulo
five forces one of the even degrees to be three-free and
the other to be five-free. A single even nonunit cannot be
free of both primes, since prime-power degrees are absent.
All supported even degrees have dyadic height one, so the
two possibilities are
\[
                      d_U\in\{10,50\},
               \qquad d_V\in\{18,54\},
\]
using degree-six exclusion. Modulo five, a signed degree
eighteen cannot cancel the signed unit. Thus $d_V=54$
and $\varepsilon_V=\varepsilon_0$. Modulo three gives
\[
 \varepsilon_U=-\varepsilon_0\quad(d_U=10),
 \qquad
 \varepsilon_U=\varepsilon_0\quad(d_U=50).
\]

Lemma~\ref{paper:even-parity-no-pure-dyadic-codegree} excludes degree $675$. The remaining
possible odd degrees are $15,45,75,135,225$, so $d_A\le225$.
The absolute value of the sum in
\eqref{paper:sparse-dimension-sum} is at most
$1+225+50+54=330<4d_Z$. It is therefore zero. In the two
cases for $U$ it has the form
\[
                  45\varepsilon_0+\varepsilon_A d_A,
        \qquad\hbox{or}\qquad
                 105\varepsilon_0+\varepsilon_A d_A.
\]
Only the first can vanish, with $d_A=45$ and
$\varepsilon_A=-\varepsilon_0$. The column is forced to be
\[
       W_2\mathbf e_Z
          =\frac{\varepsilon_0}{2}
              (\mathbf e_0-\mathbf e_{A_{45}}
                           +\mathbf e_{V_{54}}-\mathbf e_{U_{10}}).
\]

For the final contradiction use the twisted dimension moment.
Put $T=T_2$. The complementary vacuum vector $W_2d$ has
actual dyadic twist one. Applying
$W_2TW_2=\gamma^3T^{-1}W_2T^{-1}$ to $d$ therefore gives
\[
                         W_2T^{-1}d=\gamma^{-3}Td.
\]
The $Z$ coordinate of this identity has absolute value $d_Z$.
For the forced column, its left side has absolute value at most
\[
                              \frac{1+45+54+10}{2}=55,
\]
whereas $d_Z\ge90$. This contradiction excludes the vacuum
from the four-point support.
\end{proof}

The common witness has nonzero pairings with $X,Y$ but need not have
a nonzero dyadic vacuum coefficient; the sparse-column conclusion
therefore imposes no further degree exclusion by itself.

\subsection{Small signed columns and their full fields}

\begin{proposition}[The odd-parity shape at degrees ninety and two hundred seventy]
\label{paper:odd-parity-small-column-shape}
Suppose $\kappa$ is odd and a signed common witness $Z$ has
degree $90$ or $270$. Let $O$ be the set of odd-degree
labels in $P_2$, and put $u=4W_2\mathbf e_Z$.
Then exactly one of the following shapes occurs:
\[
\begin{array}{c|c|c}
 |O|&\text{coefficients on }O&\text{other nonzero coefficients}\\\hline
 4&\text{four entries }\pm1&\text{three height-one entries }\pm2\\
 8&\text{eight entries }\pm1&\text{two height-one entries }\pm2.
\end{array}
\]
All coordinates outside $P_2$ and all height-two coordinates
are zero. In the four-odd-label case every odd-degree label
of $P_2$ has actual dyadic twist one.

In the eight-odd-label case the two even support labels have
degrees $10$ or $50$, and $18$, respectively. Each of their
full character fields is contained in $K_Z$.
\end{proposition}
\begin{proof}
Corollary~\ref{paper:signed-joint-all-degrees} gives an integral $u$ of squared norm
sixteen. Principal parity makes all odd-degree $P_2$ coordinates
odd and nonzero, while the even-degree coordinates are even.
Since $d_Z$ is even, reciprocal parity makes every rational
entry at height two two-integral; its dyadic-only denominator
and nonmonomiality make that entry zero. Each nonzero
height-one coefficient is therefore $\pm2$, since $\pm4$
would exhaust the norm despite the nonzero vacuum coefficient.
Each odd coefficient is $\pm1$ or $\pm3$, with at most one
of absolute value three.

The full principal block mass is divisible by $64$.
Reduction modulo four gives $|O|\equiv0\pmod4$.
The signed dimension identity is
\[
                    \sum_Au_A d_A=4d_Zk_{2',Z}
                                      \in8d_Z\mathbb Z.
\]
Every nonunit odd degree is divisible by fifteen. If the
vacuum coefficient is $\pm1$, reduction modulo three and
five therefore requires at least two even support labels:
one three-free and one five-free, and no nonunit can have
both properties. The norm equation then leaves only
$(|O|,n_2)=(4,3)$ or $(8,2)$, with no coefficient of absolute
value three.

For completeness, the only alternatives with a coefficient
of absolute value three have $(|O|,n_2)=(4,1)$ or $(8,0)$.
If that coefficient is not at the vacuum, the preceding
two-label requirement excludes them. If it is at the vacuum,
the $(8,0)$ case is excluded modulo five. In the $(4,1)$
case the sole even support label must be five-free, so it is
different from $Z$. Thus $u_Z=0$. The diagonal modular
relation would then give $\sum_Au_A^2\theta_{A,2}=0$,
which is impossible because the vacuum term has weight nine
and all remaining terms together have weight seven.
The alternatives with no coefficient of absolute value three
but fewer than two even support labels are likewise excluded
by the two-label requirement. This proves the table.

Suppose $|O|=4$. The diagonal modular identity gives
\[
 \sum_{A\in O}\theta_{A,2}
        +4\sum_{A\in\operatorname{supp}(u)\setminus O}
                       \theta_{A,2}
        =16\gamma^3\theta_{Z,2}^{-2}(W_2)_{ZZ}.
\]
The diagonal is zero or $\pm1/2$, so the sum of the four
odd-label phases is divisible by four in the algebraic
integer ring. Its quotient has house at most one. It is
therefore zero or a root of unity. If it were zero, the sum
of the three even-label phases would be zero or twice a
dyadic root. Reduction at a two-place is impossible: the
left side reduces to one and the right to zero. Hence
the quotient is a root of unity. Equality in the triangle
inequality, including the vacuum phase one, makes all four
odd-label phases one.

Now suppose $|O|=8$. Exactly two even labels occur. The
mod-three and mod-five dimension congruences force their
degrees to be
$d_U\in\{10,50\}$ and $d_V\in\{18,54\}$.
Modulo five, the signed unit must be canceled by
$\pm2d_V$. Degree fifty-four gives the residues $\pm3$
and cannot do this; hence $d_V=18$.

Every Galois automorphism fixing $Z$ fixes its rational
column under the torus action. To see this separately at two,
dyadic inertia fixes $Z$ and local Galois covariance fixes
its rational column; at the odd primes use commutation of
the tori with $W_2$. The labels $U,V$ are the unique support
labels of their respective degrees, so each is fixed by
every stabilizer element of $Z$. This is exactly
$K_U,K_V\subseteq K_Z$.
\end{proof}

\begin{lemma}[Fourier norm of an equivariant orbit column]
\label{paper:orbit-column-fourier-norm}
Let a finite cyclic group generated by $g$ permute an
orthonormal basis, and let a unitary operator $W$ commute
with it. Suppose the orbit of $Z$ has length $e$.
On a target orbit $(A,gA,\ldots,g^{r-1}A)$ write the
coordinates of $W\mathbf e_Z$ as $c_0,\ldots,c_{r-1}$.
For $\lambda^e=1$, the norm squared of the target-orbit
projection of
\[
             W\left(e^{-1/2}\sum_{k=0}^{e-1}
                        \lambda^{-k}\mathbf e_{g^kZ}\right)
\]
is zero if $\lambda^r\ne1$, and otherwise is
\[
                  \frac er\left|\sum_{j=0}^{r-1}
                                  c_j\lambda^j\right|^2.
\]
The column coefficients satisfy $c_{j+e}=c_j$ modulo $r$.
In particular, an orbit meeting the support in one position
has $r\mid e$, and one meeting it in at most two positions
has $r/\gcd(r,e)\le2$.
\end{lemma}
\begin{proof}
The stabilizer of $Z$ fixes its column, giving the periodicity
and support assertions. Equivariance identifies the target
projection with the corresponding cyclic character vector.
Summing its coefficients over one period gives the displayed
norm formula. This is the ordinary finite Fourier transform
with the source and target orbit vectors normalized to norm one.
\end{proof}

\begin{lemma}[Large three-unit common witnesses are not rational]
\label{paper:no-rational-degree270-common}\label{paper:no-rational-degree135-common}
A signed common witness $Z$ of degree $135$ or $270$ cannot
have $K_Z=\mathbb Q$, in either principal-parity alternative.
\end{lemma}
\begin{proof}
In either case LCM integrality gives $S_{YZ}=540w$.
Rationality makes $w\in\mathbb Z$, while $Y\in P_5$ gives
$\chi_Y(Z)=27w\equiv d_Z\equiv0\pmod5$.
Thus $5\mid w$. The common-witness entry is nonzero, so
\[
          |\chi_Z(Y)|=\frac{540|w|}{d_Z}
                                  \ge\frac{2700}{d_Z}.
\]
The compulsory signed entry generates a field of degree at
least four, so the full orbit of $Y$ has at least four
members. Every one has this same rational value at $Z$.
Their contribution to the formal codegree is at least
\[
        4\left(\frac{2700}{d_Z}\right)^2
                  =\frac{5400^2}{d_Z^2}=f_Z,
\]
and the unit contributes an additional one. This contradicts
row orthogonality.
\end{proof}

\begin{corollary}[Phasewise divisibility of a three-unit column]
\label{paper:invariant-vector-phase-divisibility}
At $D=5400^2$, suppose $f_Z$ is a three-unit and
$u=mW_2\mathbf e_Z\in\mathbb Z^I$. On every actual dyadic
phase fiber $F$ one has
\[
                            \sum_{A\in F}u_A d_A\in3\mathbb Z.
\]
\end{corollary}
\begin{proof}
The three-local coordinate line of $Z$ is invariant. Its
actual three-primary twist and its odd local anomaly are
both one, so the three-local group acts trivially on that
line and $W_3\mathbf e_Z=\mathbf e_Z$. Commutation with
$W_2$ and with the actual dyadic twist operator makes the
integral phase projection $P_Fu$ a $W_3$-fixed vector.
Apply Lemma~\ref{paper:fixed-vector-dimension-congruence}.
\end{proof}

\begin{theorem}[Odd principal parity excludes a degree-two-hundred-seventy common witness]
\label{paper:odd-degree270-excluded}
If the signed principal dyadic parity $\kappa$ is odd,
a signed common witness cannot have degree $270$.
\end{theorem}
\begin{proof}
Assume that $\kappa$ is odd and a signed common witness has degree
$270$. The general column and phase lemmas above apply. We first
reduce its field and support, and then use both dimension moments.

\begin{proofclaim}[Field choices in the eight-odd-label shape]
\label{paper:odd-parity-two-seventy-fields}
In the eight-odd-label case of the proposition, if $d_Z=270$,
then
\[
                      K_Z=\mathbb Q(\sqrt5)
                   \quad\hbox{or}\quad
                      K_Z=\mathbb Q(\zeta_5).
\]
\end{proofclaim}
\noindent\emph{Proof of the claim.}
The formal codegree is $400=2^4 5^2$, and $K_Z$ is
unramified at two, so $K_Z\subseteq\mathbb Q(\zeta_{25})$.
The degree-eighteen label $V$ has a mass-two quintic
character at the compulsory degree-twenty endpoint. Its
squared absolute value contains $\sqrt5$. Thus
$\mathbb Q(\sqrt5)\subseteq K_V\subseteq K_Z$.
The vacuum coefficient of every full Galois translate of
$Z$ has absolute value $1/4$. The additional nonzero
vacuum diagonal entry, since $\kappa$ is odd, makes its
full orbit size $e$ strictly below sixteen. The subfields
of the cyclic degree-twenty field now leave $e=2,4,10$.

Suppose $e=10$. Restrict to fifth inertia. The ten-label
orbit of $Z$ is its transitive permutation quotient of
order ten, so it has a one-dimensional character of exact
order ten. On the eight odd labels $O$, no orbit has size
divisible by ten. Thus this character is absent from
$\mathbb C O$.

The other support of all these ten columns consists of the
orbits of $U$ and $V$. Their fields are subfields of $K_Z$,
so these orbit lengths divide ten. In each source column
there is exactly one $U$ coordinate and exactly one $V$
coordinate, each of absolute value $1/2$. Commutation with
fifth tori makes each such map equivariant with constant
sign. On a character of exact order ten, a target orbit
of length less than ten contributes zero. A target orbit
of length ten contributes an isometry multiplied by
$1/2$, since the unique supported-coordinate map is then
a bijection. The squared norm of the whole image of a
unit order-ten character vector is therefore at most
$1/4+1/4=1/2$. This contradicts unitarity of $W_2$.
Thus $e\ne10$, leaving exactly the two stated fields.
\par\smallskip

\begin{proofclaim}[The odd-parity degree-two-hundred-seventy field is quadratic]
\label{paper:odd-parity-degree270-quadratic}
Suppose $\kappa$ is odd and a signed common witness $Z$
has degree $270$. Then
\[
                           K_Z=\mathbb Q(\sqrt5).
\]
The two support shapes from
Proposition~\ref{paper:odd-parity-small-column-shape} are not
otherwise excluded by this proposition.
\end{proofclaim}
\noindent\emph{Proof of the claim.}
Write $W=W_2$ and $u=4W\mathbf e_Z$. The field is
unramified at two, and $f_Z=400$, so exact conductor
divisibility gives $K_Z\subseteq\mathbb Q(\zeta_{25})$.
Let $e=[K_Z:\mathbb Q]$. Every Galois conjugate of $Z$
has vacuum column coordinate of absolute value $1/4$.
The vacuum diagonal is also nonzero since $\kappa$ is
odd. Its row norm therefore gives $e<16$.
Consequently
\[
                            e\in\{1,2,4,5,10\}.
\]
The preceding lemma excludes $e=1$. We exclude the other
three nonquadratic possibilities by the cyclic fifth-torus
action. The full orbit of $Z$ has length $e$ under that
action. All its odd support coordinates are $\pm1/4$,
and its even support coordinates are $\pm1/2$.
The odd support set $\mathcal O$ is the full set of
odd-degree labels in $P_2$, by reciprocal parity. It is
therefore stable under the fifth torus. In particular every
odd target orbit has length at most $|\mathcal O|$.

First suppose the support has eight odd labels. There are
exactly two even support labels $U,V$, with
$d_U\in\{10,50\}$ and $d_V=18$, and
$K_U,K_V\subseteq K_Z$. Claim~\ref{paper:odd-parity-two-seventy-fields} leaves only
$e=2,4$ in this support shape. Suppose $e=4$.
For a primitive fourth cyclic character, the two even target
orbits contribute at most $1/4$ each to the image norm.
Among the eight odd labels, including the fixed vacuum,
at most one fifth-torus orbit can have length divisible by
four; it must have length four. Its contribution is at most
$1/2$. Equality of the total norm to one forces both even
orbits to have length four and forces the odd four-orbit
coefficients, in cyclic order, to have the form
\[
                      (a,b,-a,-b)/4,\qquad a,b\in\{1,-1\}.
\]
Its contribution at the order-two character is zero, whereas
the two even orbits contribute $1/2$ in total. The remaining
four odd labels must therefore consist of one two-orbit
with opposite coefficients and two fixed labels, one of which
is the vacuum. Denote the other fixed odd label by $F$.

The twisted dimension identity established in
Proposition~\ref{paper:sparse-common-no-vacuum} gives
\begin{equation}\label{paper:degree270-twisted-dimension}
 \left|\sum_Au_A d_A\theta_{A,2}^{-1}\right|=4d_Z=1080.
\end{equation}
Degrees and dyadic twists are constant on fifth-torus orbits.
The odd four-orbit and two-orbit therefore cancel in this
sum. Its absolute value is at most
\[
                       1+d_F+2d_U+2d_V
                          \le1+675+100+36=812,
\]
a contradiction. This excludes $e=4$ for the eight-odd shape.

Now suppose the support has four odd and three even labels.
The odd support contains the vacuum, so no odd target orbit
can have length four, five, or ten. Reduction of the signed
dimension identity modulo three and five requires a
three-free and a five-free even support label. Their possible
degrees are respectively $10,50$ and $18,54$.
In particular the three even degrees cannot all coincide.

If all three are distinct, each target orbit meets the source
column in one position. On a character of exact order $e$,
for $e\in\{4,5,10\}$, the odd coordinates contribute zero
and each even orbit contributes at most $1/4$. The norm is
at most $3/4$, a contradiction.

Otherwise exactly two even support coordinates share a
degree. Call them the paired coordinates, and call the third
one $V$. If the pair lies in different target orbits, the
same $3/4$ bound applies. Suppose it lies in one orbit of
length $r$. Its stabilizer-periodicity gives
$r/\gcd(r,e)\le2$. Since fifth-torus orbit lengths divide
$4\cdot5^b$, any target orbit that can carry a character
of exact order $e$ has $r=e$ or $r=2e$; the latter is
possible only for $e=5,10$. In the latter case the pair is
a stabilizer orbit, its two coefficients agree, and its
Fourier norm contribution is $1/2$. Together with the
single coordinate's contribution at most $1/4$, this is
again insufficient.

We are left with a paired orbit of length exactly $e$.
At a primitive $e$th character its squared contribution is
\[
                        \frac{2+\varepsilon
                                  (\lambda+\lambda^{-1})}{4},
       \qquad\varepsilon\in\{1,-1\},
\]
where $\lambda\ne1$ is the ratio of the two character
positions. The single coordinate contributes either zero
or $1/4$. For $e=5$, none of the nonidentity fifth-root
traces is $\pm1$ or $\pm2$, so the total cannot equal one.
For $e=10$, equality forces the pair to occupy opposite
positions with opposite coefficients, and forces the single
coordinate's orbit to have length less than ten. At a
primitive fifth character the pair then contributes zero;
the single coordinate contributes at most $1/2$, and the
odd support contributes zero. This is a contradiction.

Finally let $e=4$. Equality similarly forces opposite
positions with opposite coefficients in the paired four-orbit,
and a single-coordinate target orbit of length at most two.
At the order-two character the paired contribution is zero.
Norm one forces the single-coordinate orbit to have length
two and contribute $1/2$, and the odd support to contain
one two-orbit with opposite coefficients contributing the
remaining $1/2$. The other two odd labels are fixed,
including the vacuum; write $F$ for the other one.
The paired even coordinates and the odd two-orbit cancel
in \eqref{paper:degree270-twisted-dimension}. The single
even label has degree at most fifty-four, since the paired
and single degree types must supply the two distinct
prime-avoidance requirements. The remaining sum therefore
has absolute value at most
\[
                          1+675+2\cdot54=784<1080,
\]
the last contradiction. Thus only $e=2$ remains; the unique
quadratic subfield of $\mathbb Q(\zeta_{25})$ is
$\mathbb Q(\sqrt5)$.
\par\smallskip

\begin{proofclaim}[The eight-odd-label shape forces a quadratic degree-ten simple]
\label{paper:odd-parity-degree270-forces-ten}
Under the hypotheses of Proposition~
\ref{paper:odd-parity-degree270-quadratic}, suppose the column
has eight odd support labels. Its two even support labels
then have degrees ten and eighteen. Both are self-dual,
belong to $P_2$, have pure dyadic twists, and have full
character field $\mathbb Q(\sqrt5)$.
\end{proofclaim}
\noindent\emph{Proof of the claim.}
The reviewed support-shape proposition gives even labels
$U,V$ of degrees $10$ or $50$, and $18$, respectively,
with $K_U,K_V\subseteq K_Z=\mathbb Q(\sqrt5)$.
The mass-two quintic character at $V$ already makes
$\mathbb Q(\sqrt5)\subseteq K_V$, so equality holds.

Suppose $K_U=\mathbb Q$. Choose a fifth-torus generator
$g$ that exchanges $Z$ and its conjugate. It also exchanges
$V$ and its conjugate, and fixes $U$. Let $w=W_2\mathbf e_Z$.
The unit vector $(\mathbf e_Z-\mathbf e_{gZ})/\sqrt2$
has image $(w-gw)/\sqrt2$. Its even coordinates contribute
squared norm $1/4$: the $U$ contributions cancel, and the
two $V$-orbit coordinates have absolute value $1/(2\sqrt2)$.
On the eight odd coordinates, each coefficient of $w$ is
$\pm1/4$. Since $g^2w=w$, the odd coordinates whose signs
change form a union of even $g$-orbits with alternating signs.
Each such coordinate contributes $1/8$ to the image norm.
Exactly six must therefore change sign. Their weighted
dyadic twist moment vanishes orbit by orbit.

Only the vacuum and one other odd label $F$ remain in
the unsigned part of the original column. Consequently
the twisted dimension identity gives the same contradiction
\[
        1080\le1+d_F+2d_U+2d_V
                 \le1+675+100+36=812.
\]
Thus $K_U=\mathbb Q(\sqrt5)$. Degree fifty has five-unit
formal codegree and cannot have that field, so $d_U=10$.
Both full fields are real, which makes both labels self-dual.
They lie in $P_2$ by local block invariance, and odd principal
parity makes their actual twists pure dyadic.
\par\smallskip

\begin{proofclaim}[Degree two hundred seventy requires eight odd labels]
\label{paper:odd-degree270-no-four-odd}
Suppose $\kappa$ is odd and a signed common witness $Z$ has
degree $270$. The four-odd-label support shape is impossible.
Consequently the odd-degree part of $P_2$ has exactly eight
labels. The two even labels supported by $W_2\mathbf e_Z$
are self-dual simples of degrees ten and eighteen, both
with full character field $\mathbb Q(\sqrt5)$.
\end{proofclaim}
\noindent\emph{Proof of the claim.}
The field is $K_Z=\mathbb Q(\sqrt5)$ by
Claim~\ref{paper:odd-parity-degree270-quadratic}.
Suppose the four-odd-label shape occurs, put
$w=W_2\mathbf e_Z$, $u=4w$, and let $g$ generate the
fifth torus. Its source orbit has length two, and $g^2w=w$.
The four odd coordinates form the full odd-degree part
$O$ of $P_2$, so $O$ is $g$-stable and contains the fixed
vacuum. Thus its only possible nontrivial $g$-orbit is one
pair. All four odd-label twists equal one. Their coefficients
in $w$ are $\pm1/4$; the three even supported coordinates
have coefficients $\pm1/2$.

The image $W_2(\mathbf e_Z-\mathbf e_{gZ})/\sqrt2$
has norm one. An odd pair contributes $1/4$ to its squared
norm when its signs are opposite, and zero otherwise.
An even target orbit meeting the support in one coordinate
has length one or two; its contribution is respectively
zero or $1/4$.

Reduction of the untwisted dimension moment modulo three
and five requires among the even supported degrees a
three-free degree in $\{10,50\}$ and a five-free degree
in $\{18,54\}$. The three even degrees therefore cannot
all coincide. If two supported even coordinates lie in the
same $g$-orbit, the third has a different degree. The pair's
orbit has length two or four by $g^2w=w$. At length two,
equal coefficients contribute zero and opposite coefficients
contribute one. At length four the supported positions are
opposite, with equal coefficients, and contribute $1/2$.
These are all possibilities by
Lemma~\ref{paper:orbit-column-fourier-norm}.

First suppose the paired even orbit has length two and
opposite coefficients. It exhausts the norm, so the single
even coordinate is fixed and the odd coordinates have no
sign changes. The paired coordinates cancel in both dimension
moments. If $A$ is the untwisted signed odd-degree sum,
and $d$ is the remaining even degree, then $d\le54$ and
\[
             A+2\varepsilon d\in2160\mathbb Z,
   \qquad |A|\le1+3\cdot675=2026.
\]
The absolute value of the left side is at most $2134$, so
it is zero. If $t$ is the remaining even twist, the twisted
moment would consequently have absolute value
\[
          |A+2\varepsilon dt^{-1}|
                =2d|t^{-1}-1|\le4d\le216<1080,
\]
a contradiction. The equal-coefficient length-two case has
total squared norm at most $1/4+1/4<1$ and is impossible.

Next suppose the paired orbit has length four. Norm one
forces an opposite-sign odd pair and a length-two single
even orbit. The odd pair cancels in the twisted dimension
moment. Both even degree types are at most fifty-four,
since between them they supply the three-free and five-free
requirements. With $F$ the other nonunit odd label,
the absolute value of the twisted moment is therefore at most
\[
                  1+d_F+4\cdot54+2\cdot54
                          \le1000<1080.
\]
This is again impossible.

The three even coordinates must consequently lie in three
separate $g$-orbits. To reach norm one, each orbit has length
two and there is an opposite-sign odd pair. Let $F$ be the
remaining nonunit odd label. The odd pair cancels in both
moments. Each supported even label is moved by fifth inertia,
so its formal codegree is divisible by five. With dyadic
height one this leaves the even-degree palette
\[
                         \{10,18,30,54,90,270\}.
\]
The three-free degree must be ten and one degree must be
$18$ or $54$. If $d_F\le225$, the remaining twisted moment
has absolute value at most
\[
                         1+225+2(10+54+270)=894<1080.
\]
Thus $d_F=675$. If the third even degree were at most ninety,
the analogous bound would be
$1+675+2(10+54+90)=984<1080$.
It must therefore be $270$.

The untwisted moment now reads
\[
 \varepsilon_0+675\varepsilon_F+20\varepsilon_U
                 +2d\varepsilon_V+540\varepsilon_R
                           \in2160\mathbb Z,
       \qquad d\in\{18,54\}.
\]
Its absolute value is at most $1344$, so it must vanish.
But
\[
 |675\varepsilon_F+540\varepsilon_R|\ge135,
 \qquad
 |\varepsilon_0+20\varepsilon_U+2d\varepsilon_V|
                                      \le129,
\]
which is impossible. This excludes the four-odd-label shape.
The final assertions follow from
Claim~\ref{paper:odd-parity-degree270-forces-ten}.
\par\smallskip

\noindent\emph{The contradiction.}
Suppose it does. The preceding results give
$K_Z=\mathbb Q(\sqrt5)$ and exactly eight odd-degree
labels in $P_2$. Put $W=W_2$, $w=W\mathbf e_Z$,
and $u=4w$. The odd coordinates of $w$ are $\pm1/4$;
its only even support labels are $U,V$, of degrees ten and
eighteen, with coefficients $\pm1/2$. Both have full field
$\mathbb Q(\sqrt5)$. All supported actual twists are purely
dyadic. Denote the signs in $u$ by
$u_0=\varepsilon_0$, $u_U=2\varepsilon_U$,
$u_V=2\varepsilon_V$.

Let $g$ generate the fifth torus. Its orbits on $Z,U,V$
all have length two, and $g^2w=w$. The image of the
unit vector $(\mathbf e_Z-\mathbf e_{gZ})/\sqrt2$
under $W$ has squared norm one. Its $U$-orbit and $V$-orbit
contributions are $1/4$ each. On an odd coordinate where
$g$ changes the sign of $w$, the contribution is $1/8$;
on every other odd coordinate it is zero. Hence precisely
four odd coordinates change sign. They form complete even
$g$-orbits with alternating coefficients. Their contributions
to both dimension moments cancel, since degrees and dyadic
twists are constant on each $g$-orbit.

Let $R$ be the remaining four odd coordinates. It contains
the vacuum, and $gw=w$ on $R$. Write $\varepsilon_A=u_A$
for $A\in R$. The untwisted dimension moment becomes
\[
           \sum_{A\in R}\varepsilon_A d_A
                   +20\varepsilon_U+36\varepsilon_V
                                      \in2160\mathbb Z.
\]
Its absolute value is at most
$1+3\cdot675+20+36=2082$, so it is zero.
Reduction modulo three and five gives
\[
             \varepsilon_U=\varepsilon_0,
          \qquad\varepsilon_V=-\varepsilon_0,
          \qquad
             \sum_{A\in R}\varepsilon_A d_A
                                         =16\varepsilon_0.
\tag{*}
\]

The formal codegree of $Z$ is $400$, a three-unit. Apply
Corollary~\ref{paper:invariant-vector-phase-divisibility} to each
dyadic phase component of $u$. Every nonunit odd supported
degree is divisible by three, and $d_V=18$ is divisible by
three. Thus the only three-free contributions are those of
the vacuum and $U$. Each has nonzero residue modulo three,
so they must belong to the same phase fiber. Therefore
\[
                              \theta_U=1.
\]

Neither $Z$ nor $gZ$ belongs to the support of $w$, since
both have degree $270$. On the symmetric source vector
$z_+=(\mathbf e_Z+\mathbf e_{gZ})/\sqrt2$, therefore,
$z_+^{\mathsf t}Wz_+=0$. Its actual dyadic twist is the
same as that of $Z$. The local modular identity
$WTW=\gamma^3T^{-1}WT^{-1}$ consequently gives
\[
                (Wz_+)^{\mathsf t}T(Wz_+)=0.
\]
The four sign-changing odd coordinates have canceled from
$Wz_+$; each coordinate in $R$ has square $1/8$, and each
of the two even target orbits has total squared weight $1/4$.
It follows that
\[
                     \sum_{A\in R}\theta_A
                                  =-2(1+\theta_V).
\tag{**}
\]

A vanishing positive word of dyadic roots is a disjoint sum
of opposite pairs, because one representative from each
opposite pair at a fixed dyadic level is linearly independent
over $\mathbb Q$. In \((**)\), the vacuum and the two
copies of $\theta_U=1$ give three copies of one. If
$\theta_V\ne-1$, the remaining three odd letters cannot
both cancel these three copies and cancel the two copies
of $\theta_V$. Hence
\[
             \theta_V=-1,\qquad
             \{\theta_A:A\in R\}=\{1,-1,r,-r\}
\]
for a dyadic root $r$.

Normalize the signs by multiplying the dimension moment
by $\varepsilon_0$. Its two even contributions to the
inverse-twist moment are now $20+36=56$. If all three
nonunit degrees in $R$ were at most $225$, that moment
would have absolute value at most $1+3\cdot225+56=732$,
whereas its required absolute value is $4d_Z=1080$.
Thus one of these degrees is $675$.
The nonunit part of \((*)\) is a signed sum equal to fifteen.
After division by fifteen, its possible degrees belong to
$\{1,3,5,9,15,45\}$. A term of absolute value forty-five
can occur in such a three-term signed sum equal to one only
with a second forty-five of opposite sign and a final one
of positive sign. Thus the three degrees and their normalized
coefficients are
\[
                         (+15,+675,-675).
\]
Write their twists as $t,a,b$, respectively. The required
inverse-twist moment is
\[
                    \left|57+15t^{-1}
                                  +675a^{-1}-675b^{-1}\right|
                                      =1080,
       \qquad \{t,a,b\}=\{-1,r,-r\}.
\]
If $t=-1$, the expression has form $42\pm1350r^{-1}$,
whose absolute value is at least $1308$.
If $a=-1$, then $b=-t$, and the expression is
$-618+690t^{-1}$. Its squared absolute value would force
\[
                    t+t^{-1}=-\frac{8566}{11845}.
\]
If $b=-1$, then $a=-t$, and the expression is
$732-660t^{-1}$, forcing
\[
                    t+t^{-1}=-\frac{1354}{3355}.
\]
Both are nonintegral rational numbers, whereas the trace of
a root of unity is an algebraic integer. These contradictions
exhaust the possible twist assignments and exclude $Z$.
\end{proof}

\begin{proposition}[A diagonal-zero degree-two-hundred-seventy column]
\label{paper:even-degree270-zero-diagonal}
Suppose $\kappa$ is even and a signed common witness $Z$
has degree $270$. Write $W=W_2$, $T=T_2$, and let
$\gamma$ be the dyadic anomaly. If $W_{ZZ}=0$, then there
are four distinct degree-$135$ labels $A,B,C,D$ such that
\[
       W\mathbf e_Z=\tfrac12(\mathbf e_A+\mathbf e_B
                                  -\mathbf e_C-\mathbf e_D).
\]
For $\rho\in\{1,-1\}$, their actual twists are
\[
           \theta_A=\theta_B=\rho,\qquad
           \theta_C=\theta_D=-\rho,\qquad
           \theta_Z=\rho\gamma^3.
\]
Put
\[
 w=\tfrac12(\mathbf e_A+\mathbf e_B-\mathbf e_C-\mathbf e_D),
 \qquad
 v=\tfrac12(\mathbf e_A+\mathbf e_B+\mathbf e_C+\mathbf e_D).
\]
The orthonormal span of $\mathbf e_Z,w,v$ is invariant
under the full dyadic local group, with
\[
 W\mathbf e_Z=w,\quad Ww=\mathbf e_Z,\quad Wv=\rho v,
 \qquad
 T\mathbf e_Z=\rho\gamma^3\mathbf e_Z,
 \quad Tw=\rho v,\quad Tv=\rho w.
\]
At least one of the two twist-one degree-$135$ labels belongs
to $P_2\cap P_5$.
\end{proposition}
\begin{proof}
Proposition~\ref{paper:sparse-common-no-vacuum} gives four entries
$\varepsilon_i/2$, no vacuum coordinate, and no supported
degree of dyadic height at least two. Commutation with the
odd twist operators makes every supported twist purely
dyadic. Even parity excludes
purely dyadic formal codegrees, so the possible supported
degrees are exactly among
\[
 \mathcal D=\{10,15,18,30,45,50,54,75,90,135,150,225,270,450\}.
\]
Write their dyadic twists as $t_i$. The diagonal modular
identity, the dimension sum, and the twisted dimension moment
give respectively
\[
 \sum_{i=1}^4t_i=0,\qquad
 \sum_i\varepsilon_i d_i\in1080\mathbb Z,\qquad
 \left|\sum_i\varepsilon_i d_i t_i^{-1}\right|=540.
\]
A zero sum of four dyadic roots consists of two opposite
pairs. Write them as $\rho,-\rho,\sigma,-\sigma$ and put
\[
       a=\varepsilon_1d_1-\varepsilon_2d_2,
       \qquad b=\varepsilon_3d_3-\varepsilon_4d_4.
\]
Then $|a\rho^{-1}+b\sigma^{-1}|=540$.

Neither $a$ nor $b$ can be zero. If $a=0$, the first
pair has equal degrees and equal coefficients, while
$|b|=540$. The latter pair must have degrees $(450,90)$
or $(270,270)$ with opposite coefficients. Its untwisted
signed sum is therefore $\pm360$ or zero. Combining with
the first pair's signed sum $\pm2d$, the required divisibility
by $1080$ would force $d=180$ or $360$, or would require
a nonzero $2d$ to be divisible by $1080$. None is possible
for $d\in\mathcal D$. The argument for $b=0$ is identical.

The rational number
\[
           \rho/\sigma+\sigma/\rho
                  =\frac{540^2-a^2-b^2}{ab}
\]
is an algebraic integer. A rational trace of a dyadic root
is $-2,0$, or $2$. If the trace is zero, then
$a^2+b^2=540^2$ with both terms nonzero. The only positive
integer leg pair is $(324,432)$: the primitive hypotenuse
must be five, so the triple is $108(3,4,5)$.
Inspection of $\mathcal D$ shows that an absolute signed
difference $324$ requires $(270,54)$ with opposite
coefficients, and $432$ requires $(450,18)$ with equal
coefficients. Their untwisted sums are $\pm216$ and
$\pm468$, never a multiple of $1080$. Thus the trace is
$\pm2$ and all four twists lie in one opposite pair.

Let $P,M$ be the signed degree sums on the two twist classes.
Each has two labels, and
$P+M\in1080\mathbb Z$, $|P-M|=540$.
The possible pairs would require signed two-degree sums
of absolute value $270$ or $810$. The value $810$ does
not occur in $\mathcal D$. The only pairs giving absolute
sum $270$ are $(45,225)$ and $(135,135)$, with equal
coefficient signs within each pair. Thus the column signs
are constant on each twist class and opposite between them.

The field $K_Z$ is a nontrivial subfield of
$\mathbb Q(\zeta_{25})$, by
Lemma~\ref{paper:no-rational-degree270-common}. Some fifth
torus therefore moves $Z$. Its column must be orthogonal
to the column of $Z$. Fifth tori preserve dyadic twists,
so every overlapping coordinate contributes nonnegatively
to this inner product. A degree-$225$ coordinate would be
fixed by every fifth torus, since its formal codegree is a
five-unit, and would contribute $1/4$. This is impossible.
All four degrees must therefore be $135$. Choose the names
so that the positive coefficients occur at $A,B$, and denote
their common twist by $\rho$.

The exact twisted dimension identity, including its scalar,
now gives
\[
                   270\rho^{-1}
                         =270\gamma^{-3}\theta_Z,
       \qquad \rho=\gamma^3\theta_Z^{-1}.
\]
Dyadic inertia fixes $Z$, so $W^2\mathbf e_Z=\mathbf e_Z$
and $Ww=\mathbf e_Z$. The modular relation
$WTW=\gamma^3T^{-1}WT^{-1}$ applied to $\mathbf e_Z$
then gives $Wv=\rho^{-1}v$. Since $W^2$ is an unsigned
permutation and $v$ is a positive vector, the equality
$W^2v=\rho^{-2}v$ forces $\rho^2=1$.
All the stated actions and invariant-span assertions follow.

Finally write $d=(d_E)_E$. Symmetry gives
\[
 \langle Wd,v\rangle=270\rho,\qquad
 \langle Wd,w\rangle=270.
\]
The sum of the coordinates of $Wd$ on the twist-one
pair is consequently $540$ or $-540$, and at least one
such coordinate is nonzero. The identity
$Wd=5400W_3W_5\mathbf e_0$ applies. Each degree-$135$
label has three-unit formal codegree, so its coordinate
line is fixed by the three-local factor. Thus the same
coordinate of $W_5\mathbf e_0$ is nonzero, placing it
in $P_5$. It already lies in $P_2$ by the source column's
block support. This proves the final assertion.
\end{proof}

\begin{remark}
The three-dimensional span above is a local representation
subspace. No fusion closure, algebra object, or categorical
factor is claimed. The proposition leaves the case
$W_{ZZ}=\pm1/2$ and does not exclude the displayed
diagonal-zero configuration.
\end{remark}

\subsection{Actual signed phase buckets}

\begin{proposition}[Degree-sixty signed buckets]
\label{paper:sixty-bucket-palette}
For a signed compulsory crossing with $(d_X,d_Y)=(12,20)$,
each of the four exact phase buckets has total degree sixty.
Every degree in a bucket belongs to
\[
                   \{10,12,15,20,30,36,45,50,60\}.
\]
If the bucket is active at five, degree fifty is absent.
If a bucket has a constituent whose degree is prime to five,
its entire degree partition is one of
\[
                         12+12+12+12+12,
                 \qquad 12+12+36.
\]
Thus each bucket is either supported entirely on degrees
divisible by five or entirely on degrees prime to five.
There are sixteen possible degree partitions in total, fifteen
for a bucket active at five. These are degree partitions with
fusion multiplicity included; equal summands need not denote
the same simple object.
\end{proposition}
\begin{proof}
Every signed constituent has actual dyadic twist minus one,
so the Weyl-line criterion gives $v_2(d_V)\le2$.
The reciprocal signed-bucket restriction gives
$4\mid d_V$ or $5\mid d_V$.  Intersect these conditions with
$d_V\mid5400$, $d_V\le60$, and prime-power and degree-six
exclusion. This gives the displayed palette. An active fifth
twist further gives $v_5(d_V)\le1$, removing fifty.

The only degrees prime to five in the palette are twelve and
thirty-six. Their total contribution is a positive multiple of
five, since all other terms are divisible by five. A positive
multiple of five of the form $12a+36b$, not exceeding sixty,
must equal sixty. Solving $a+3b=5$ gives the two stated
partitions.

The remaining fourteen partitions are
\[
\begin{gathered}
6(10),\quad4(10)+20,\quad3(10)+2(15),\quad3(10)+30,\\
2(10)+2(20),\quad10+2(15)+20,\quad10+20+30,\quad10+50,\\
4(15),\quad2(15)+30,\quad15+45,\quad3(20),\quad2(30),\quad60,
\end{gathered}
\]
where $k(d)$ denotes $k$ repeated summands of degree $d$.
Only $10+50$ is removed by fifth activity.
\end{proof}

\begin{remark}[Why direct exposed-adjoint transfer stops]
The difference support of every signed tame four-position
word contains $\{1,\xi^a,\xi^{-a}\}$ for some primitive fifth
root $\xi$ and $a\not\equiv0\pmod5$: take the two positions
whose ternary coefficients are both one. Such a set cannot be
exposed in the full even root alphabet. Indeed, a linear
functional bounded above by one on every root would take the
value three on their sum if it exposed all three positions.
But the fifth-root relation writes that same sum as the sum
of two negative fifth roots, on which the functional is at
most two. Therefore the exposed-adjoint method used to exclude
the signed wild support cannot be applied directly to either
signed tame type. This is a failure of that proposed method,
not evidence that a signed tame crossing exists.
\end{remark}

\subsection{Mass-two characters and the degree-six exclusion}

\begin{lemma}[The ternary character at common weight two]
\label{paper:modular:degree-ten-ternary-character}
\label{paper:modular:mass-two-ternary-character}
Let $X\in P_3$ be a compulsory witness of degree $12$ or $36$.
If $W$ is simple, $3\nmid d_W$, and $\gcd(d_X,d_W)=2$,
then a $3$-primary root of unity $\eta$ satisfies
\[
                    \chi_X(W)=-\frac{d_W}{2}\eta.
\]
If $W$ is self-dual, its value is $-d_W/2$.
In particular this applies to degrees ten and fifty.
This numerical assertion does not prescribe the actual
balancing phases of $XW$.
\end{lemma}
\begin{proof}
Put $z=S_{XW}/\operatorname{lcm}(d_X,d_W)$.
It is an algebraic integer with an actual positive rational word
of mass two and strict house below two.  Since
$\chi_X(W)=(d_W/2)z$, with $d_W/2$ a three-unit, and $X\in P_3$, its residue is two at every
place over three.

Lemma~\ref{paper:family:lem:two-prime-lattice} supplies a positive integer
word for $z$ with at most two letters.  All roots involved are
supported on the primes $2,3,5$.  Erase the $3$-primary parts
of the letters in this chosen integer word and call the resulting
sum $w$.  Reduction modulo a three-place erases exactly these
parts.  The field containing $w$ is unramified at three, so the
all-places residue gives $w-2\in3\mathcal O$.
Consequently $(w+1)/3$ is an algebraic integer whose conjugates
lie in the closed disk with center $1/3$ and radius $2/3$.
The disk meets the unit circle only at one.  Taking the algebraic
norm therefore gives $w=-1$ or $w=2$.

If $w=-1$, the chosen word has one letter: two roots summing to
$-1$ must be the two primitive cube roots, whereas the letters
of $w$ have orders prime to three.  Thus $z=-\eta$ with
$\eta$ $3$-primary.  If $w=2$, equality in the triangle
inequality gives two letters with trivial prime-to-three parts,
so $z=\eta_1+\eta_2$ for $3$-primary roots $\eta_1,\eta_2$.
Strictness excludes $\eta_1=\eta_2$.

It remains to exclude a ratio $\lambda=\eta_1/\eta_2$ of
order at least nine.  In that case
$S=\{\eta_1,\eta_2\}$ meets each cube-root coset in at most
one position. Lemma~\ref{paper:family:lem:exposure} therefore exposes it
in the entire root alphabet:
the full even-alphabet exposure inequality on every occupied
coset is $2\cdot1<3$.  Since an exposing functional has
value two on $z$, \emph{every actual positive word of mass two}
for $z$ has support in $S$.  In particular the full character
word of $W$ at $X$ has those two phases.

Its adjoint difference set is
$D=\{1,\lambda,\lambda^{-1}\}$.  This set also meets each
cube-root coset in at most one position and is exposed.  Choose
a Galois automorphism acting by the second power on the
$3$-primary roots and trivially at the other primes.  Argument-side
Galois covariance gives the corresponding full character word
of the conjugate simple $W^\sigma$, with adjoint difference
set $D^2$.  Since the order of $\lambda$ is at least nine,
\[
                            D\cap D^2=\{1\}.
\]
Lemma~\ref{paper:family:lem:common} now contradicts
strictness at the nonunit row $X$.

The remaining ratio has order three, and hence
$\eta_1+\eta_2$ itself is the negative of a $3$-primary root.
This proves the numerical assertion in all cases.
Finally, a self-dual $W$ has real character value at $X$.
The only real $3$-primary root is one, so its value is $-d_W/2$.
\end{proof}

\begin{corollary}[No degree-six adjoint with self-dual degree-ten remainder]
\label{paper:modular:degree-six-selfdual-ten-remainder}
There is no degree-six simple object $V$ with
\[
             VV^*=1+Z_{15}+W_{10}+W'_{10},
\]
where $W,W'$ are self-dual; repetitions are allowed.
\end{corollary}
\begin{proof}
The degree-fifteen summand $Z$ is self-dual because it occurs
with multiplicity one in the self-dual object $VV^*$.
Choose a compulsory witness $X$ of degree $12$ or $36$, and
put $\alpha=\chi_X(V)$ and $a=\alpha\overline\alpha$.
The lemma gives
\[
                    a=\chi_X(Z)-9.
\]
Every conjugate of $a$ is nonnegative, and strictness at $Z$
makes it strictly less than six.  LCM integrality gives
$\chi_X(Z)\in5\mathcal O$, so
\[
                  \frac{a-1}{5}=\frac{\chi_X(Z)}5-2
\]
is an algebraic integer.  Its conjugates lie in $[-1/5,1)$,
so its house is strictly below one and it is zero.  Thus $a=1$.
On the other hand, $X\in P_3$ gives
$\alpha\equiv d_V=6\equiv0$ at every place over three,
contradicting $\alpha\overline\alpha=1$.

\end{proof}

\begin{lemma}[The quintic character at common weight two]
\label{paper:modular:degree-six-quintic-character}
\label{paper:modular:mass-two-quintic-character}
Let $Y\in P_5$ be a compulsory witness of degree twenty.
For every simple $V$ with $5\nmid d_V$ and $\gcd(d_V,20)=2$,
there are a five-primary root $\alpha$ and a primitive fifth
root $\xi$ such that
\[
                 \chi_Y(V)=\frac{d_V}{2}\alpha(1+\xi).
\]
The two displayed phases, each with coefficient $d_V/2$,
are the full actual balancing word. In particular
\[
       \chi_Y(VV^*)=\frac{d_V^2}{4}(2+\xi+\xi^{-1}).
\]
For self-dual $V$, one has $\chi_Y(V)=(d_V/2)s$ with
$s=\zeta+\zeta^{-1}$ a primitive fifth-root real trace.
This applies in particular to degrees six, eighteen and fifty-four.
\end{lemma}
\begin{proof}
The normalized crossing $z=S_{VY}/\operatorname{lcm}(d_V,20)$ is integral and has an
actual positive word of mass two.  Since $\chi_Y(V)=(d_V/2)z$, with $d_V/2$ a five-unit,
principality at five gives residue two for $z$ at every five-place.
Corollary~\ref{paper:cor:one-large-residue} applies because $5>2\cdot2$.
Thus the actual word consists of two $5$-primary roots, each
with coefficient one.  They are distinct by strictness.

Let $\lambda$ be their ratio.  If it has order at least $25$,
the actual adjoint difference set
$D=\{1,\lambda,\lambda^{-1}\}$ has one marked position in
each occupied fifth-root coset.  The full even-alphabet exposure
inequality in Lemma~\ref{paper:family:lem:exposure} is $2\cdot1<5$,
so $D$ is exposed.
The automorphism taking fifth-primary roots to their second powers
gives the exposed difference set $D^2$, and
$D\cap D^2=\{1\}$.  Argument-side Galois covariance and the
Lemma~\ref{paper:family:lem:common} contradict strictness.
Consequently $\lambda$ has order five.  Absorb one of the
original phases into $\alpha$ and take $\xi=\lambda$.
Multiplication by the dual character gives the adjoint identity.
If $V$ is self-dual, reality forces $\alpha^2\xi=1$.
Squaring is invertible on five-primary roots, so the two phases
form a nonidentity inverse pair of fifth roots, as asserted.
\end{proof}


\begin{theorem}[No simple object of degree six]
\label{paper:modular:no-degree-six}
\label{paper:modular:no-selfdual-six}
An ordinary perfect fusion-simple integral modular category
of dimension $5400^2$ has no simple object of degree six.
\end{theorem}
\begin{proof}
Suppose that $V$ is such an object. The unit has multiplicity
one in $VV^*$. Simple degrees divide $5400$, nontrivial
invertibles are absent, and nonunit prime-power degrees have
already been excluded. The nonunit degrees at most $35$ are
\[
                         6,10,12,15,18,20,24,30.
\]
Their only odd member is fifteen. The residual dimension $35$
therefore gives
\begin{equation}\label{paper:modular:eq:six-adjoint-profiles}
                  VV^*=\mathbf1+Z+R,
                  \qquad d_Z=15,\quad d_R=20,
\end{equation}
where $Z$ is simple and self-dual, and $R$ is either one
self-dual degree-twenty simple or the sum of two degree-ten
simples. In the latter case duality either fixes both terms
or interchanges them. Corollary~\ref{paper:modular:degree-six-selfdual-ten-remainder}
excludes the former, including a repeated term. Thus the only
degree-ten possibility is $R=W+W^*$ with $W\not\simeq W^*$.

\paragraph{One adjoint component is fixed by Galois.}
The rational vacuum row and unsigned Galois covariance make
the Galois permutations automorphisms of the fusion ring: apply
a cyclotomic automorphism to the Verlinde formula, whose vacuum
denominators are rational. In particular each conjugate $V^\sigma$
has adjoint components $Z^\sigma,R^\sigma$ of the same type
as those in \eqref{paper:modular:eq:six-adjoint-profiles}.

Any two conjugate adjoints have a common nonunit constituent.
Indeed, if the unit were their only common constituent, rigidity
and commutativity would give
\[
 \big\langle V^\sigma(V^\sigma)^*,
                 V^\tau(V^\tau)^*\big\rangle
             =\|V^\sigma V^\tau\|^2=1.
\]
Then $V^\sigma V^\tau$ would be simple, its double braiding
would be scalar, and strictness at two nonunit labels would fail.
Consequently, for every $\sigma,\tau$,
\[
                 Z^\sigma=Z^\tau
                  \quad\hbox{or}\quad R^\sigma=R^\tau.
\]
For the second conclusion, two degree-ten dual pairs sharing
one constituent coincide; the degree-twenty case is immediate.

A family of ordered pairs with this property has a common
first or second coordinate. To see this, if two first coordinates
differ, their second coordinates coincide; any third pair must
then share that second coordinate, since it cannot share both
different first coordinates. Hence either $Z$ or $R$ is fixed
by every Galois permutation. Its character values are therefore
rational algebraic integers.

\paragraph{The quintic character excludes both possibilities.}
Fix a compulsory $Y\in P_5$ of degree twenty and put
$s=\xi+\xi^{-1}$ for the primitive fifth root supplied by
Lemma~\ref{paper:modular:degree-six-quintic-character}. Applying the
character at $Y$ to \eqref{paper:modular:eq:six-adjoint-profiles} gives
\begin{equation}\label{paper:modular:eq:six-quintic-adjoint-sum}
                 \chi_Y(Z)+\chi_Y(R)=17+9s.
\end{equation}
Every conjugate of $\chi_Y(Z)$ has absolute value strictly
below fifteen. Every conjugate of $\chi_Y(R)$ has absolute
value strictly below twenty, by strictness for its one or two
nonunit simple constituents.

If $Z$ is Galois fixed, $\chi_Y(Z)$ is an integer.
LCM integrality gives $\chi_Y(Z)\in3\mathcal O$, while
$Y\in P_5$ gives $\chi_Y(Z)\equiv15\equiv0$ at five.
Thus $\chi_Y(Z)$ is a multiple of fifteen; strictness makes
it zero. But one conjugate of $s$ is
$s_+=(\sqrt5-1)/2>1/3$, so
\eqref{paper:modular:eq:six-quintic-adjoint-sum} would give a conjugate
of $\chi_Y(R)$ equal to $17+9s_+>20$, a contradiction.

If $R$ is Galois fixed, its character is an integer with
principal-five residue $d_R=20\equiv0$. Write
$\chi_Y(R)=5k$, $k\in\mathbb Z$.
Since $\chi_Y(Z)/3$ and $s$ are algebraic integers,
\eqref{paper:modular:eq:six-quintic-adjoint-sum} gives
\[
                         k\equiv1\pmod3.
\]
On the other hand, the two conjugates of $\chi_Y(Z)$ are
\[
                     \frac{25-10k\pm9\sqrt5}{2}.
\]
Strictness at both forces
\[
                    |25-10k|+9\sqrt5<30.
\]
Here $|25-10k|$ is a positive odd multiple of five, and
$5<30-9\sqrt5<15$. Hence $|25-10k|=5$, so $k=2$ or $3$,
contrary to $k\equiv1\pmod3$.
Both Galois alternatives are impossible, proving the theorem.
\end{proof}

\subsection{An integral-spectrum obstruction}

\begin{lemma}[Strict spectral radius for a positive central element]
\label{paper:strict-positive-central}
Let $\mathcal C$ be a nonpointed fusion-simple fusion category,
with $\mathcal C_{\ad}=\mathcal C$. Let $P$ be a nonzero element
of the positive cone of $K_0(\mathcal C)$, supported on nonunit
simples, and assume $P$ is central. If $P$ acts on an irreducible
representation $E$ by the scalar $\lambda_E(P)$, then
\[
       |\lambda_E(P)|<\FPdim(P)
       \qquad\text{unless }E=\FPdim.
\]
\end{lemma}
\begin{proof}
The coefficient-of-unit trace makes the complexified based ring
a finite dimensional $*$-algebra with positive inner product.
Its irreducible representations may be taken to preserve $*$.
For a simple class $X$, left multiplication satisfies
$L_X^*=L_{X^*}$. The nonnegative matrix $L_{X^*X}$ has the
positive dimension vector as an eigenvector of eigenvalue
$d_X^2$. Hence its spectral radius is $d_X^2$, and
$\|E(X)\|\le\|L_X\|=d_X$ in every irreducible representation.

Suppose equality held for $P=\sum_Xc_XX$. Since
$E(P)=\lambda_E(P)\operatorname{id}$, equality in the norm
triangle inequality on every unit vector forces
\[
 E(X)=\zeta d_X\operatorname{id}
 \quad(c_X>0),\qquad |\zeta|=1.
\]
Choose one such nonunit $X$. It generates $\mathcal C$. Products
of $X$ and $X^*$ act as scalars of full possible norm. Applying
the same equality argument to their nonnegative simple
decompositions makes every $E(Y)$ a scalar
$\alpha_Yd_Y\operatorname{id}$ with $|\alpha_Y|=1$.
Irreducibility gives $\dim E=1$. For every constituent $Z$ of
$UV$, equality in the dimension triangle inequality gives
$\alpha_Z=\alpha_U\alpha_V$. Also
$\alpha_{U^*}=\overline{\alpha_U}$, so all constituents of
$UU^*$ have phase one. Since $\mathcal C_{\ad}=\mathcal C$,
all phases are one and $E=\FPdim$, as asserted.
\end{proof}

\begin{theorem}[A factorial divisor from an integral central spectrum]
\label{paper:integral-central-spectrum-factorial}
Let $\mathcal D$ be a nonpointed integral fusion-simple
fusion category with $\mathcal D_{\ad}=\mathcal D$, and put
$n=\FPdim(\mathcal D)$. Let $P$ be a nonzero positive
central element of $K_0(\mathcal D)$ supported on nonunit
simples, of dimension $d$. Suppose its scalar in every
irreducible representation of the complexified fusion ring
is an integer. Then
\[
                              n\mid(2d-1)!.
\]
\end{theorem}
\begin{proof}
Lemma~\ref{paper:strict-positive-central} puts every scalar
other than the dimension character in
$\{-d+1,-d+2,\ldots,d-1\}$. Thus, for
\[
                 p(t)=\prod_{j=-d+1}^{d-1}(t-j)\in\mathbb Z[t],
\]
semisimplicity gives
\[
                         p(P)=(2d-1)!\,e_{\FPdim}.
\]
The left side belongs to the integral based ring. The
coefficient-of-unit trace of the primitive idempotent on
the right is $1/n$, because the dimension character has
formal codegree $n$. Therefore
$\tau(p(P))=(2d-1)!/n$ is an integer, proving divisibility.
\end{proof}

\begin{corollary}[A fully rational commutative ring at dimension 5400]
\label{paper:rational-commutative-source-rank}
Let $\mathcal C$ be a simple integral fusion category of dimension
$5400$. If its Grothendieck ring is commutative and every
irreducible character is rational valued, then every nonunit
simple of $\mathcal C$ has dimension at least six and
\[
                         \operatorname{rank}(\mathcal C)\le149.
\]
\end{corollary}
\begin{proof}
Simplicity at this composite dimension makes $\mathcal C$
nonpointed and gives $\mathcal C_{\ad}=\mathcal C$:
otherwise the adjoint subcategory is trivial, all simples
are invertible, and a simple pointed category has prime dimension.
Every simple class is central, and its eigenvalues are
rational algebraic integers. Apply the factorial theorem
to that class. A dimension at most five would require
$5400$ to divide a factorial with five-valuation at most one,
which is impossible. Hence all nonunit dimensions are at
least six, initially giving rank at most $150$.
At rank $150$, their squared dimensions would satisfy
\[
             \sum_{X\ne\mathbf1}(d_X^2-36)=35.
\]
The only positive increments at most thirty-five are
$7^2-6^2=13$ and $8^2-6^2=28$. No nonnegative combination
of thirteen and twenty-eight equals thirty-five. Thus rank
$150$ is impossible.
\end{proof}

\begin{remark}
The dimension identity
$1+145\cdot6^2+2\cdot7^2+9^2=5400$ has rank $149$.
It only shows that this last numerical bound cannot be
improved using that sum-of-squares relaxation alone; it
does not assert a fusion ring or a categorification of the type.
\end{remark}

\begin{corollary}[No rational degree-ten row in the ordinary modular reduction]
\label{paper:no-rational-degree-ten}
In the ordinary perfect fusion-simple integral modular category
of dimension $5400^2$, no degree-ten simple has rational full
character field. In particular every degree-ten full Galois
orbit has at least two labels. This is a global conclusion;
it does not assume a center realization or the no-thirty branch.
\end{corollary}
\begin{proof}
The modular fusion ring is commutative. If $K_Z=\mathbb Q$,
symmetry of $S$ gives
\[
                  \chi_X(Z)=\frac{d_Z}{d_X}\chi_Z(X)\in\mathbb Q.
\]
These are eigenvalues of an integral fusion matrix, so all
are integers. Apply the theorem with $P=Z$, $d=10$, and
$n=5400^2$. It would require $5400^2\mid19!$, whereas
$v_5(19!)=3<4=v_5(5400^2)$.
\end{proof}
\subsection{Spectral annihilators and exact gap divisibility}
\label{paper:sec:spectral-annihilators}

\begin{theorem}[The non-Perron annihilator divisor]
\label{paper:spectral-annihilator-divisor}
Let $\mathcal D$ be a nonpointed integral fusion-simple category,
put $R=K_0(\mathcal D)$ and $n=\FPdim(\mathcal D)$, and let
$P\in R$ be nonzero and positive, supported on nonunit simples.
Write $d=\FPdim(P)$. Suppose $q(t)\in\mathbb Z[t]$ annihilates
$E(P)$ for every irreducible representation $E$ of $R_{\mathbb C}$
other than the dimension representation. Then
\begin{equation}\label{paper:annihilator-idempotent}
 q(P)=q(d)e_{\FPdim},\qquad n\mid q(d).
\end{equation}
There is a unique monic polynomial of least degree with this
annihilation property; it belongs to $\mathbb Z[t]$ and does not
vanish at $d$. If $P$ is central, or more generally if $PP^*=P^*P$,
this polynomial is the product of the distinct monic minimal
polynomials over $\mathbb Q$ of the non-Perron scalar eigenvalues.
In particular it is squarefree in this case.
\end{theorem}
\begin{proof}
For $P=\sum_Xc_XX$, Lemma~\ref{paper:source:strict-character-operator-norm}
gives
\[
 \|E(P)\|\le\sum_Xc_X\|E(X)\|<\sum_Xc_Xd_X=d
 \qquad(E\ne\FPdim).
\]
The regular representation of $R$ is faithful and its matrices are
integral. Its minimal polynomial at $P$ is therefore monic in
$\mathbb Z[t]$. Semisimplicity of $R_{\mathbb C}$ decomposes the
regular representation into the irreducible matrix blocks. The
strict inequality shows that the dimension block supplies the
factor $t-d$ exactly once and that no other block has eigenvalue
$d$. Dividing by $t-d$ gives the asserted least polynomial
$q_0\in\mathbb Z[t]$ with $q_0(d)\ne0$.
For any annihilator $q$ as in the statement, its value is zero on
all the other blocks and $q(d)$ on the dimension block, proving
the first identity. The coefficient-of-unit trace satisfies
$\tau(e_{\FPdim})=1/n$. Since $q(P)\in R$, its unit coefficient
$q(d)/n$ is an integer, proving divisibility.
Finally a central element acts by scalars. A normal element acts
by normal matrices in the regular $*$-realization. In either case
all the relevant matrices are diagonalizable, so their minimal
polynomials have no repeated roots. Their Galois-stable union of
roots yields the stated product of distinct irreducible factors.
\end{proof}

\begin{corollary}[The actual integral-spectrum divisor]
\label{paper:actual-integer-gap-divisor}
Under the preceding hypotheses, suppose $P$ is central and every
scalar eigenvalue is integral. Let $\Lambda$ be the set of its
\emph{distinct} non-Perron eigenvalues. Then
\[
 n\mid\prod_{\lambda\in\Lambda}(d-\lambda),
 \qquad \Lambda\subset\{-d+1,\ldots,d-1\}.
\]
This strengthens the factorial bound
$n\mid(2d-1)!$: omitting an actual scalar value can remove a prime
factor that the factorial relaxation had supplied.
\end{corollary}
\begin{proof}
Apply Theorem~\ref{paper:spectral-annihilator-divisor} to
$q(t)=\prod_{\lambda\in\Lambda}(t-\lambda)$.
\end{proof}

\begin{corollary}[Forced values at the smallest rational remainders]
\label{paper:forced-small-integral-spectra}
For $n=5400$ and $d=6$ in the preceding corollary, the non-Perron
spectrum contains $1,-4,-3$, and at least one of $0,3$.
It contains at least five distinct values.
For $n=5400^2$ and $d=12$, it contains at least seven distinct
values. If there are exactly seven, they are
\[
                         7,3,2,0,-3,-6,-8.
\]
These statements concern integer spectra; they do not assert
existence of an element with any of the displayed spectra.
\end{corollary}
\begin{proof}
At $d=6$, the available positive gaps are $1,\ldots,11$.
Two factors of five require both gaps $5,10$, giving the values
$1,-4$. Three factors of three require gap $9$ together with at
least one of $3,6$, giving $-3$ and at least one of $3,0$.
These four gaps supply at most two factors of two, whereas
$v_2(5400)=3$, so at least one further gap is required.

At $d=12$, four factors of five require all four gaps
$5,10,15,20$. They supply only one factor of three; to reach
$v_3(5400^2)=6$ requires at least three further gaps. If exactly
three are used, two of them must be $9,18$. The remaining gap
must be one of $3,6,12,21$. The four forced gaps and $9,18$
supply only four factors of two, so the final gap must be $12$.
The resulting seven gaps are $5,9,10,12,15,18,20$, whose product
is exactly $5400^2$. Subtracting them from twelve gives the list.
\end{proof}

\begin{remark}[Exact use in character-table searches]
\label{paper:annihilator-row-implementation}
For a commutative source ring and a simple object of dimension
$d$, a rational character value $a$ contributes $t-a$ once,
regardless of how many characters have that value. A real
quadratic character pair with trace $s$ and squared-norm trace
$u$ contributes
\[
                         t^2-st+\frac{s^2-u}{2};
\]
an imaginary quadratic pair contributes $t^2-st+u/2$.
If the quadratic entries are rational, the repeated linear
factor is retained only once. Equal irreducible factors from
different character orbits are also retained only once.
The resulting product evaluated at $d$ must be divisible by
$5400$. A grouped set of rational columns is retained only if
\emph{some} integer multiset with the specified trace and squared
norm passes this condition. Thus grouping is a relaxation and
cannot create an exclusion by choosing one convenient multiset.

For a noncommutative ring, a characteristic polynomial of a
matrix block is an admissible annihilator, but its repeated
roots must not be discarded without proving diagonalizability.
The theorem does not require that every available annihilator
be minimal; a larger annihilator gives a weaker valid test.
\end{remark}

\subsection{Restrictions on a degree-ten adjoint}

\begin{lemma}[A quintic relative trace at mass five]
\label{paper:quintic-real-mass-five-trace}
Let $z$ be a real integral cyclotomic number supported on the
primes $2,3,5$. Suppose it has a positive rational word of mass
at most five, strict house below five, and residue zero at
every five-place. In a common splitting field let $A$ be
normalized averaging over the automorphisms fixing
$\mathbb Q(\sqrt5)$. Then every conjugate of $A(z)$ has
absolute value at most
\[
                          B=\frac{5+\sqrt5}{2}.
\]
\end{lemma}
\begin{proof}
Integral optimization supplies a word
$z=\sum_{j=1}^n\xi_j\eta_j$, $n\le5$, where
$\xi_j$ has order prime to five and $\eta_j$ is five-primary.
Put $w=\sum_j\xi_j$. Erasing the five-primary factors
preserves the residue at five. The field of $w$ is unramified
there, so $w\in5\mathcal O$. Since $w$ has house at most
five, Kronecker's theorem gives $w=0$ or $w=5\varepsilon$
with $\varepsilon$ a root of unity.

The averaging factors into full averaging on the prime-to-five
field and averaging on the five-primary field over
$\mathbb Q(\sqrt5)$. Write
\[
 t_j=A_{5'}(\xi_j),\qquad u_j=A_5(\eta_j).
\]
Here $t_j$ is real and lies in $[-1,1]$, while
\[
 u_j\in\left\{1,0,\frac{\sqrt5-1}{4},
                         -\frac{\sqrt5+1}{4}\right\}.
\]
The value zero covers every root of order at least twenty-five.
If $w=0$, then $\sum_jt_j=0$. The sum $P$ of the positive
$t_j$ equals the absolute sum of the negative ones. Since
there are at most five terms, $P\le2$. Consequently
\[
 \left|A(z)\right|=\left|\sum_jt_ju_j\right|
 \le P\left(1+\frac{\sqrt5+1}{4}\right)\le B.
\]

If $w=5\varepsilon$, equality in the triangle inequality
forces $n=5$ and $\xi_j=\varepsilon$ for every $j$.
Write $z=\varepsilon v$ with $v=\sum_j\eta_j$.
We may assume $v\ne0$. Reality of $z$, together with the
trivial intersection of the two coprime cyclotomic fields,
gives $\varepsilon^2\in\mathbb Q$, hence
$\varepsilon\in\{1,-1,i,-i\}$. For $\varepsilon=\pm i$
the prime-to-five average is zero.

For $\varepsilon=\pm1$, $v$ is real. Its five-letter
coefficient vector is invariant under inversion unless $v=0$:
indeed the difference from its inverse is constant on every
fifth-root coset. A nonzero difference requires all five
letters in one non-self-inverse coset, whose sum is zero.
Thus nonidentity letters occur in inverse pairs. Strictness
excludes five copies of one, so at most three letters are
one and at least one letter is one. For $\varepsilon=1$,
\[
 A_5(v)\le3+2\frac{\sqrt5-1}{4}=B.
\]
For $\varepsilon=-1$,
\[
 -A_5(v)\le-1+4\frac{\sqrt5+1}{4}=\sqrt5<B.
\]
The same argument for $-z$ bounds the opposite sign.
All hypotheses are preserved by Galois conjugation, so the
bound holds at either embedding of $\mathbb Q(\sqrt5)$.
\end{proof}

\begin{proposition}[No degree-fifty-four constituent]
\label{paper:degree-ten-no-fifty-four}
If $W$ is simple of degree ten, then no degree-fifty-four
simple occurs in $WW^*$.
\end{proposition}
\begin{proof}
Suppose $WW^*=\mathbf1+U+R$ with $d_U=54$ and $d_R=45$.
The constituent $U$ is unique and self-dual. At a compulsory
$Y\in P_5$ of degree twenty,
Lemma~\ref{paper:modular:mass-two-quintic-character} gives
$\chi_Y(U)=27s$, where $s$ is a primitive fifth-root real
trace. Choose an embedding with $s=(-1-\sqrt5)/2$ and average
over the automorphisms fixing $\mathbb Q(\sqrt5)$.

Since $R$ is self-dual of odd dimension, it has a self-dual
odd-degree constituent $Z$. The only possible degrees are
$d_Z=15$ or $45$. LCM integrality writes
$\chi_Y(Z)=(d_Z/5)z$, where $z$ is real, integral, has a
positive word of mass five and strict house below five.
Principality at five makes $z$ vanish at every five-place,
because $d_Z/5$ is a five-unit. By the preceding lemma,
$A(\chi_Y(Z))\le(d_Z/5)B$.
The remaining constituents have character real part at most
their dimensions. Since $B<5$ and $d_Z\ge15$, this gives
\[
 A(\chi_Y(R))\le45-d_Z+\frac{d_Z}{5}B\le30+3B.
\]
Every conjugate of $\chi_Y(WW^*)=|\chi_Y(W)|^2$ is
nonnegative. However,
\[
 0\le A(\chi_Y(WW^*))
 \le1-\frac{27(1+\sqrt5)}2+30+\frac{3(5+\sqrt5)}2
 =25-12\sqrt5<0,
\]
a contradiction.
\end{proof}

\begin{proposition}[Large constituents of a degree-ten adjoint]
\label{paper:degree-ten-large-adjoint}
Let $W$ be simple of degree ten. None of the degrees fifty,
seventy-five, or ninety occurs in $WW^*$. At most four
self-dual degree-ten constituents occur, counted with multiplicity.
\end{proposition}
\begin{proof}
At a compulsory $X\in P_3$ of degree twelve or thirty-six,
Lemma~\ref{paper:modular:mass-two-ternary-character} gives
\[
                         \chi_X(WW^*)=25.
\]
The unit occurs once, and the least nonunit degree is ten.

A degree-fifty constituent $U$ would occur once and be
self-dual, since its multiplicity or its distinct dual would
already use dimension one hundred. Write
$WW^*=\mathbf1+U+R$, where $d_R=49$.
The lemma gives $\chi_X(U)=-25$. Strictness at every
nonunit constituent of $R$ gives
\[
             25=1-25+\operatorname{Re}\chi_X(R)
                    <1-25+49=25,
\]
a contradiction.

A degree-seventy-five constituent is similarly unique and
self-dual. Write $WW^*=\mathbf1+U+R$ with $d_R=24$.
LCM integrality gives $\chi_X(U)=25z$ with $z$ integral.
Since $X\in P_3$ and $25$ is a three-unit, $z$ vanishes
at every three-place. On the other hand
\[
                         \chi_X(R)=24-25z.
\]
The self-duality of $R$ and strictness give
$0<\sigma(z)<48/25<2$ for every embedding $\sigma$.
Thus $z-1$ is an algebraic integer of strict house below one,
so $z=1$, contrary to its three-residue.

A degree-ninety constituent would leave nonunit dimension nine,
which is impossible. Finally, if $m$ self-dual degree-ten
constituents occur, their total character is $-5m$. The
remaining nonunit part has dimension $99-10m>0$, so
\[
                     25<100-15m.
\]
Hence $m<5$.
\end{proof}

\begin{remark}[Exact scope]
The bound counts only self-dual degree-ten constituents.
For a non-self-dual pair the character contribution is
$-5(\eta+\eta^{-1})$, which can be positive. Dimension
alone still permits six degree-ten terms together with degrees
$(15,24)$ or $(15,12,12)$; the preceding argument does not
exclude those alternatives.
\end{remark}

\begin{corollary}[The remaining degree-ten adjoint palette]
\label{paper:degree-ten-adjoint-palette}
If $d_W=10$, every nonunit constituent of $WW^*$ has degree in
\[
           \{10,12,15,18,20,24,30,36,40,45,60,72\}.
\]
The odd degrees are fifteen and forty-five. A degree-forty-five
constituent is unique and self-dual, because two such terms
would leave nonunit dimension nine. At least three nonunit
constituents occur, counted with multiplicity.
\end{corollary}
\begin{proof}
Enumerate the divisors of $5400$ below one hundred, omit prime
powers and degree six, and apply the preceding exclusions.
Neither one nor two terms from the displayed palette can sum
to ninety-nine.
\end{proof}

\begin{lemma}[Averaged common support of conjugate adjoints]
\label{paper:adjoint-orbit-cover}
For a nonunit simple $W$, write
$WW^*-\mathbf1=\sum_U m_U[U]$. For every full Galois orbit
$\mathcal O$ of nonunit labels put
$M_{\mathcal O}=\sum_{U\in\mathcal O}m_U$.
Then
\[
                 \sum_{\mathcal O}
                    \frac{M_{\mathcal O}^2}{|\mathcal O|}\ge1.
\]
For degree ten one also has
$\sum_{\mathcal O}d_{\mathcal O}M_{\mathcal O}=99$.
\end{lemma}
\begin{proof}
Unsigned Galois covariance makes the label permutations based-ring
automorphisms. Two conjugate adjoints must share a nonunit
constituent: otherwise rigidity gives
$\|W^\sigma W^\tau\|^2=1$, making their tensor product
simple and contradicting strictness. Consequently
$\langle WW^*-1,\sigma(WW^*-1)\rangle\ge1$ for every
$\sigma$. Averaging over the finite Galois group gives the
displayed orbit sum. The dimension identity is immediate.
\end{proof}

\begin{lemma}[Two short real ternary mean bounds]
\label{paper:short-real-ternary-means}
Let $\mathsf A$ denote full normalized Galois averaging. Suppose
$z$ is a real integral cyclotomic number supported on $2,3,5$.
\begin{enumerate}
\item If $z$ has a positive rational word of mass at most three,
strict house below three, and residue zero at every three-place,
then $|\mathsf A(z)|\le1$.
\item If $z$ has a positive rational word of mass at most four,
strict house below four, and residue one at every three-place,
then $\mathsf A(z)\le2$.
\end{enumerate}
\end{lemma}
\begin{proof}
Integral optimization gives integer words of the same mass bounds.
For the first assertion write $z=\sum_{j=1}^n\xi_j\eta_j$,
$n\le3$, with $\xi_j$ prime to three and $\eta_j$
three-primary. Put $w=\sum_j\xi_j$. The all-places residue
and unramifiedness give $w\in3\mathcal O$. By its house
bound and Kronecker's theorem, $w=0$ or $3\varepsilon$.

If $w=0$ and the word is nonempty, it has exactly two opposite
letters: three roots sum to zero only in an equilateral triple,
which the prime-to-three alphabet excludes. Thus
$z=\xi(\eta_1-\eta_2)$. If this is nonzero and real,
\[
                  \xi^2=-\frac1{\eta_1\eta_2}.
\]
The two coprime cyclotomic fields intersect in $\mathbb Q$.
Consequently $\xi=\pm i$, $\eta_2=\eta_1^{-1}$, and the
full mean is zero. If $w=3\varepsilon$, all three
prime-to-three letters equal $\varepsilon$. Reality again
gives $\varepsilon\in\{\pm1,\pm i\}$; the imaginary
possibilities have mean zero. A nonzero real three-letter
word on the three-primary roots is invariant under inversion:
otherwise the coset relation consumes all three letters and
their sum is zero. Strictness then leaves
$1+\eta+\eta^{-1}$, with $\eta\ne1$, whose full mean is
zero or one. Multiplication by $\varepsilon=\pm1$ proves
the first bound.

For the second assertion, every nonidentity allowed root has
full mean at most $1/2$. Equality occurs only for primitive
sixth roots; all other nonidentity roots have mean at most
$1/4$. These values follow from
$\mathsf A(\rho)=\mu(n)/\varphi(n)$ for a primitive
$n$th root. A real integer word of length at most four and
mean greater than two must therefore be either $z=3$ or
\[
                         z=2+\rho+\rho^{-1}.
\]
Indeed, with no identity letter its mean is at most two.
With exactly one identity letter, exceeding two requires two
primitive sixth-root letters. If they are equal, their imaginary
parts cannot be cancelled by the one remaining root; if they
are conjugate, the remaining root must be real, and the mean
does not exceed two. With two identity letters the other two
must be inverse or opposite; the latter case gives two. Three
identity letters give three or, after adding $-1$, two; four
identity letters violate strictness.

The value three has the wrong three-residue. In the displayed
form erase the three-primary part of $\rho$, obtaining a
prime-to-three root $\xi$. The residue implies that
$(1+\xi+\xi^{-1})/3$ is integral. Its real conjugates lie
in $[-1/3,1]$, so the norm forces it to be zero or one.
The zero case would make $\xi$ a primitive third root;
the other gives $\xi=1$. Thus $\rho$ is three-primary,
whose full mean is nonpositive unless $\rho=1$. The latter
violates strictness, and the former gives mean at most two.
This proves the second assertion.
\end{proof}

\begin{corollary}[At most three self-dual degree-ten terms]
\label{paper:degree-ten-three-selfdual}
At most three self-dual degree-ten constituents occur in the
adjoint of a degree-ten simple, counted with multiplicity.
\end{corollary}
\begin{proof}
The preceding bound already permits at most four. Suppose four
occur and remove them from $WW^*-\mathbf1$. The remaining
dimension is fifty-nine. Its only possible profiles in the
degree palette have the form
\[
                 Z_{15}+R_{20}+R_{24},
\]
where $Z$ is simple and self-dual, $R_{20}$ is one self-dual
degree-twenty simple or two degree-ten simples, and $R_{24}$
is one degree-twenty-four simple or two degree-twelve simples.
All three parts are self-dual. The two degree-ten terms in
$R_{20}$ cannot both be self-dual by the earlier bound.

At a compulsory ternary witness $X$, LCM integrality and
principality give $\chi_X(Z)=5z$, where $z$ has mass three,
strict house below three, and residue zero at three. Its full
mean is at most five by the lemma. For a self-dual degree-twenty
simple, $\chi_X(U)=5u$, where $u$ has mass four, strict
house below four, and residue one at three. Its mean is at most
ten. If $R_{20}$ is a degree-ten dual pair, its character is
$-5(\eta+\eta^{-1})$ for a three-primary root $\eta$;
the full mean is at most five, hence still at most ten.
Finally strictness gives $\mathsf A(\chi_X(R_{24}))<24$.
The four self-dual degree-ten terms contribute $-20$, whereas
$\chi_X(WW^*)=25$. Averaging this identity now gives
\[
                   25<1-20+5+10+24=20,
\]
a contradiction.
\end{proof}

\subsection{Actual twists, Galois overlap and degree-ten profiles}


\begin{lemma}[Three distinct-degree adjoint constituents]
\label{paper:three-adjoint-fields}
Suppose a nonunit simple $W$ has
\[
             WW^*=\mathbf1+Z_1+Z_2+Z_3,
\]
where the $Z_i$ are simple and have pairwise distinct dimensions.
Then either some $Z_i$ is fixed by every Galois permutation, or
the three full character fields $K_{Z_i}$ are the three distinct
quadratic subfields of a biquadratic field contained in $K_W$.
\end{lemma}
\begin{proof}
The $Z_i$ are self-dual, since their dimensions are distinct.
Let $G$ be the Galois group of a common cyclotomic splitting
field and let $H_i$ be the stabilizer of $Z_i$. Any two conjugate
adjoints have a common nonunit constituent, by the strictness
argument of Lemma~\ref{paper:adjoint-orbit-cover}. Applied to
$WW^*$ and $W^\sigma(W^\sigma)^*$, and using the distinct
dimensions, this says
\[
                         G=H_1\cup H_2\cup H_3.
\]
The quotient by $H_1\cap H_2\cap H_3$ embeds in
$\prod_iG/H_i$, and each of its elements has an identity
coordinate. Lemma~\ref{paper:three-coordinate-support} gives
either a trivial coordinate projection, or a Klein four group
whose three coordinate projections have order two. These are
exactly the stated field alternatives. Finally, any element
fixing $W$ fixes its adjoint and hence each $Z_i$, again by
distinct dimensions. Thus each $K_{Z_i}$, and their compositum,
is contained in $K_W$.
\end{proof}

\begin{lemma}[Two full quintic mean bounds]
\label{paper:short-real-quintic-means}
Let $\mathsf A$ denote full normalized Galois averaging. Let
$z$ be a real integral cyclotomic number supported on $2,3,5$.
\begin{enumerate}
\item If $z$ has a positive rational word of mass at most five,
strict house below five, and residue zero at every five-place,
then $|\mathsf A(z)|\le3$.
\item If $z$ has a positive rational word of mass at most four,
strict house below four, and residue four at every five-place,
then $\mathsf A(z)\le2$.
\end{enumerate}
\end{lemma}
\begin{proof}
Use integral optimization and write each letter as
$\xi_j\eta_j$, with $\xi_j$ prime to five and $\eta_j$
five-primary. Put $w=\sum_j\xi_j$,
$t_j=\mathsf A(\xi_j)$, and $u_j=\mathsf A(\eta_j)$.
The means factor over the coprime primary fields, and
\[
                 -1\le t_j\le1,
       \qquad u_j\in\{1,-1/4,0\}.
\]

For the first assertion, the proof of
Lemma~\ref{paper:quintic-real-mass-five-trace} gives $w=0$ or
$w=5\varepsilon$, with $\varepsilon$ a root of unity.
If $w=0$, the positive and negative parts of $\sum_jt_j=0$
have common mass at most two, since there are at most five
terms. Hence $|\mathsf A(z)|\le2(1+1/4)=5/2$.
If $w=5\varepsilon$, all five prime-to-five factors equal
$\varepsilon$. If $z\ne0$, reality forces
$\varepsilon\in\{1,-1,i,-i\}$. The imaginary cases have
mean zero. In the real cases the five-letter five-primary
word is invariant under inversion unless it is zero, as proved
in that lemma. Strictness excludes five identity letters, so
there are one or three identity letters, with the other letters
in inverse pairs. Its full mean lies in $[0,3]$.
Multiplication by $\varepsilon=\pm1$ proves the bound.

For the second assertion, unramifiedness at five and the
all-places residue imply $w-4\in5\mathcal O$. The algebraic
integer $(w+1)/5$ has its conjugates in the disk with center
$1/5$ and radius $4/5$. The norm argument gives $w=-1$ or
$w=4$. In the first case $\sum_jt_j=-1$. If $P$ is the sum
of the positive $t_j$, the negative part has mass $P+1$.
At most four terms, each of absolute value at most one, imply
$P\le1$: with at most one positive term this is immediate;
with at least two there are at most two negative terms, giving
$P+1\le2$. Therefore
\[
       \mathsf A(z)=\sum_jt_ju_j
             \le P+\frac{P+1}{4}\le\frac32.
\]
If $w=4$, there are four letters and all $\xi_j=1$.
A real four-letter word on the five-primary roots is invariant
under inversion: a nonzero difference from its inverse is
constant on a fifth-root coset and would require at least five
original letters. By strictness, there are at most two identity
letters. Every other letter has nonpositive full mean, so the
mean is at most two.
\end{proof}


\begin{lemma}[A character determines a two-summand twist product]
\label{paper:two-summand-phase-product}
Let $W,B$ be degree-ten simples, and suppose both $WB$ and
$WB^*$ are sums of two distinct nonunit simples, each with
multiplicity one. If $\chi_W(B)\ne0$, the products of the
two relative balancing phases in these two decompositions
are complex conjugates.
\end{lemma}
\begin{proof}
The possible degree pairs are $(10,90)$, $(40,60)$ and
$(50,50)$. Their character squared moduli have the respective
forms
\[
             82+9x,\qquad52+24y,\qquad50+25z,
\]
where $x,y,z$ are traces $\rho+\rho^{-1}$ of roots of unity.
These three sets of values are disjoint below one hundred.
For the first two, equality gives $3x-8y=-10$, hence
$(y+1)/3\in\mathcal O$. All conjugates of this algebraic
integer lie in $[-1/3,1]$, so the norm argument makes it zero
or one. Thus $y=-1$ or $2$. The first gives squared modulus
twenty-eight, below the first pair's lower bound sixty-four; the
second gives the forbidden value one hundred.
Comparing the first and third similarly gives
$(z+1)/3\in\mathcal O$, with $z=-1$ too small and $z=2$
forbidden. Comparing the second and third gives the same
alternatives; $z=-1$ would require $y=-9/8$, not an algebraic
integer. Since the characters at $B$ and $B^*$ are conjugate,
their degree pairs therefore agree.

For unequal degrees, divide the balancing coefficients by
their common divisor. They become $(a,b)=(1,9)$ or $(2,3)$.
A number $C=a\alpha+b\beta$, with $\alpha,\beta$ roots
of unity, has a unique such ordered root decomposition.
Indeed, a second decomposition gives
$a(\alpha-\alpha')=b(\beta'-\beta)$. Coprimality of
$a,b$ makes $(\alpha-\alpha')/b$ an algebraic integer,
whose house is at most $2/b<1$. Hence $\alpha=\alpha'$
and then $\beta=\beta'$. Applying uniqueness
to $C$ and $\overline C$ proves the assertion in these cases.
For equal degrees, if $C=\alpha+\beta\ne0$, then directly
$\alpha\beta=C/\overline C$. This proves the assertion
for the degree-fifty pair as well.
\end{proof}

\begin{proposition}[Vafa rigidity of a uniquely fixed adjoint constituent]
\label{paper:vafa-unique-fixed-adjoint}
Let $d_W=10$ and suppose
\[
             WW^*=\mathbf1+Z_1+Z_2+Z_3
\]
with pairwise distinct nonunit degrees. Assume
$\chi_W(W^*)\ne0$. If a Galois permutation fixes exactly
one of the three constituents, that constituent has twist one.
Consequently at least one constituent is fixed by every
Galois permutation. If exactly one is fixed by every Galois
permutation, its twist is one.
\end{proposition}
\begin{proof}
Let $B=gW$ and let $Z$ be the uniquely fixed constituent.
Common-adjoint overlap gives
\[
 \langle WW^*,BB^*\rangle=2,
 \qquad \|WB\|^2=\|WB^*\|^2=2.
\]
Here $B$ is neither $W$ nor $W^*$, since either equality
would fix all three adjoint constituents. Both products
therefore have two distinct nonunit summands with multiplicity
one. Their characters at $W$ are nonzero by argument-side
Galois covariance.

Apply the four-point Vafa identity
\cite[Theorem~8.18.1]{EGNO} with $(i,j,m,r)=(W,W^*,B,B)$.
The common-adjoint factor is $\theta_Z$; the other two
factors are the products of the two twists in $WB$ and
$W^*B$. Dividing each of these twists by
$\theta_W\theta_B$ cancels the endpoint factor on the
left, because each product has squared multiplicity sum two.
The identity becomes
\[
              \theta_Z\,P(WB)\,P(WB^*)=1,
\]
where $P$ denotes the product of relative phases; replacing
$W^*B$ by its dual $WB^*$ preserves all twists.
The preceding lemma gives $P(WB)P(WB^*)=1$.
Thus $\theta_Z=1$.

If none of the three constituents were globally fixed,
Lemma~\ref{paper:three-adjoint-fields} would give the three
index-two stabilizers of a Klein four quotient. Each
constituent is then the unique fixed one for some Galois
permutation. All three twists would be one, and balancing
would give $|\chi_W(W^*)|=10$, contrary to strictness.
Thus at least one is globally fixed. If exactly one is,
choose a permutation outside the two other proper stabilizers;
it fixes only that constituent, so its twist is one.
\end{proof}

\begin{proposition}[Twists and overlap in the sixty profile]
\label{paper:fifteen-twentyfour-sixty}
Suppose $d_W=10$ and
\[
                   WW^*=\mathbf1+Z_{15}+U_{24}+V_{60}.
\]
Then $\theta_U=1$, $\theta_Z\in\{1,-1\}$, and
\[
 \chi_W(W^*)=\theta_W^{-2}
 \begin{cases}
       4+6\theta_V,&\theta_Z=1,\\
       1+6\theta_V,&\theta_Z=-1.
 \end{cases}
\]
At least one adjoint constituent is fixed by Galois.
If $\theta_Z=-1$ and at least two constituents are not fixed
by Galois, then $\theta_V=-1$, and every permutation moving
two constituents fixes $U$ and gives two degree-fifty
summands in each of $W(gW)$ and $W(gW)^*$.
In this latter case $U$ is fixed by all Galois permutations.
\end{proposition}
\begin{proof}
Balancing gives
$1+15\theta_Z+24\theta_U+60\theta_V\in10\mathcal O$.
Modulo two, the unit-disk norm argument gives
$\theta_Z=\pm1$. Modulo five,
$(1+4\theta_U)/5$ is integral and has house at most one.
It cannot vanish, so equality in the triangle inequality
gives $\theta_U=1$. The character formula follows and is
nonzero in either case. The preceding proposition supplies
a fixed constituent.

Now assume $\theta_Z=-1$ and take a permutation moving
two constituents. Put $s=g(\theta_V)+g(\theta_V)^{-1}$.
The corresponding character squared modulus is $37+6s$.
The degree pair $(10,90)$ is impossible by its lower bound
sixty-four, since $37+6s\le49$.
The pair $(40,60)$ would give
$52+24t=37+6s$, or $4t-s=-5/2$, impossible for algebraic
integers. The pair $(50,50)$ gives
$25t-6s=-13$, hence $(s+2)/5\in\mathcal O$.
Its house is at most $4/5$, so $s=-2$ and $\theta_V=-1$;
then $t=-1$. This proves the degree-pair assertion.
The only twist-one adjoint constituent is now $U$, so every
permutation moving two must fix $U$.
If $U$ were not globally fixed and either of the other two
were not, a permutation outside their proper stabilizers
would uniquely fix the remaining twist-minus-one constituent,
which is impossible. Thus either $U$ is fixed or both $Z,V$
are fixed. Under the present hypothesis of two nonfixed
constituents, the latter alternative cannot occur.
\end{proof}


\begin{lemma}[Twists forced by the forty-five adjoint profile]
\label{paper:forty-five-profile-twists}
In the ordinary perfect fusion-simple reduction at $D=5400^2$,
suppose $d_W=10$ and
\[
                   WW^*=\mathbf1+Z_{45}+U_{36}+V_{18}.
\]
Then
\[
        \theta_Z=-1,\qquad \theta_U=\theta_V=1,
        \qquad \chi_W(W^*)=\theta_W^{-2}.
\]
\end{lemma}
\begin{proof}
Write $t=\theta_Z$, $u=\theta_U$, and $v=\theta_V$.
Balancing and character integrality give
\[
    T:=1+45t+36u+18v
       =\theta_W^2 S_{W,W^*}\in10\mathcal O.
\]
Reducing this divisibility modulo two gives
$1+t\in2\mathcal O$. The algebraic integer $(1+t)/2$
has all conjugates in the unit disk, so it is zero or a root
of unity. Equality in the triangle inequality in the nonzero
case forces $t=1$; the zero case gives $t=-1$.

Reducing $T\in5\mathcal O$ gives
$1+18(2u+v)\in5\mathcal O$, hence
\[
                      h:=\frac{2+2u+v}{5}\in\mathcal O.
\]
All conjugates of $h$ lie in the unit disk. If $h$ is a root
of unity, equality in the triangle inequality, with the two
identity letters present, forces $u=v=1$. The other possibility
$h=0$ would give $v=-2(1+u)$ and therefore
\[
                      u+u^{-1}=-\frac74,
\]
which is impossible for an algebraic integer. Thus $u=v=1$.

If $t=1$, then $T=100$ and $|\chi_W(W^*)|=10$, contrary
to strictness at two nonunit labels. Thus $t=-1$, giving
$T=10$ and the stated character value.
\end{proof}

\begin{theorem}[No adjoint profile $1+45+36+18$]
\label{paper:forty-five-profile-excluded}
An ordinary perfect fusion-simple integral modular category of
dimension $5400^2$ has no degree-ten simple with
\[
                   WW^*=\mathbf1+Z_{45}+U_{36}+V_{18}.
\]
This conclusion does not require the absence of thirty-active
twists, and it excludes the entire degree profile, for every permitted character assignment.
\end{theorem}
\begin{proof}
All three nonunit constituents are self-dual, by their distinct
degrees. At a compulsory $Y\in P_5$ of degree twenty,
Lemma~\ref{paper:modular:mass-two-quintic-character} gives
\[
                           \chi_Y(V)=9s,
\]
where $s$ is a primitive fifth-root real trace. Hence $V$ is
not fixed by all Galois permutations. If both $Z$ and $U$
were fixed, their character values would be rational algebraic
integers. LCM integrality, the principal-five residues, and
strictness would give
\[
                        \chi_Y(Z)=0,
                   \qquad \chi_Y(U)=-9.
\]
Indeed, $\chi_Y(Z)/9$ is an integer of absolute value below
five and divisible by five, whereas $\chi_Y(U)/9$ is an
integer of absolute value below four congruent to four modulo
five. It would follow that
\[
                       |\chi_Y(W)|^2=-8+9s<0,
\]
since even the larger value $s=(\sqrt5-1)/2$ is less than
$8/9$. This is impossible. At least two of $Z,U,V$ are
therefore nonfixed.

Two proper subgroups do not cover a finite group. Choose a
Galois automorphism $g$ moving two nonfixed constituents.
Any two conjugate adjoints share a nonunit constituent, by
Lemma~\ref{paper:adjoint-orbit-cover}; the distinct degrees
therefore force $g$ to fix the third constituent. Consequently
\[
       \left\langle WW^*,(gW)(gW)^*\right\rangle=2.
\]
Rigidity and commutativity identify this with
$\|W(gW)^*\|^2$. Also $gW\ne W$, so the unit does not occur.
Thus $W(gW)^*$ is the sum of exactly two distinct nonunit
simples, each with multiplicity one and with total dimension
one hundred. Simple dimensions divide $5400$, nonunit
prime-power degrees and degree six are absent, so the only
possible pairs of dimensions are
\[
                          (10,90),\qquad(40,60),\qquad(50,50).
\]

By Lemma~\ref{paper:forty-five-profile-twists},
$\chi_W(W^*)$ is a root of unity.
Argument-side Galois covariance now shows that
\[
                   \chi_W((gW)^*)=g(\chi_W(W^*))
\]
is also a root of unity. If the two summands of $W(gW)^*$
are $A,B$, balancing gives
\[
       \chi_W((gW)^*)
           =\frac{d_A\theta_A+d_B\theta_B}
                         {10\theta_W\theta_{gW}}.
\]
For $(d_A,d_B)=(10,90)$, the reverse triangle inequality
makes its absolute value at least eight; for $(40,60)$ the
lower bound is two. Both contradict its absolute value one.
For $(50,50)$ the displayed value is five times an algebraic
integer, multiplied by a root of unity. A root of unity cannot
be divisible by five in the ring of algebraic integers: its
quotient by five would have all conjugates of absolute value
$1/5$. This excludes the final pair and proves the theorem.
\end{proof}


\begin{lemma}[Integer radii of a two-summand degree-ten product]
\label{paper:two-summand-integer-radii}
Let $W,B$ be degree-ten simples in the ordinary reduction at
$D=5400^2$, and suppose $WB$ is the sum of exactly two distinct
nonunit simples, each with multiplicity one. If
$c=|\chi_W(B)|$ is an integer, then
\[
                            c\in\{0,2,5,8\}.
\]
More precisely, the possibilities are
\[
\begin{array}{c|c|c}
\text{summand degrees}&c&\text{ratio of summand twists}\\\hline
(10,90)&8&-1\\
(40,60)&2&-1\\
(50,50)&0&-1\\
(50,50)&5&\text{a primitive cube root}.
\end{array}
\]
\end{lemma}
\begin{proof}
The only allowed dimension pairs summing to one hundred are
$(10,90)$, $(40,60)$, and $(50,50)$. Balancing, after removal
of a common phase, writes the character as respectively
$1+9\rho$, $4+6\rho$, or $5+5\rho$, with $\rho$ a root
of unity. Thus
\[
 c^2=82+9(\rho+\rho^{-1}),\quad
 c^2=52+24(\rho+\rho^{-1}),\quad\text{or}\quad
 c^2=50+25(\rho+\rho^{-1}).
\]
When $c$ is an integer, the trace $\rho+\rho^{-1}$ is a
rational algebraic integer in $[-2,2]$, hence belongs to
$\{-2,-1,0,1,2\}$. Substitution gives the table; the value
$c=10$ is excluded by strictness.
\end{proof}

\begin{lemma}[An integer self-radius forces fixed adjoint constituents]
\label{paper:integer-self-radius-rigidity}
Suppose $d_W=10$ and
$WW^*=\mathbf1+Z_1+Z_2+Z_3$, with distinct nonunit degrees.
If $\chi_W(W^*)=c\varepsilon$, where $c$ is a nonnegative
integer outside $\{0,2,5,8\}$ and $\varepsilon$ a root of
unity, then at least two of the $Z_i$ are fixed by every
Galois permutation.
\end{lemma}
\begin{proof}
If two constituents are nonfixed, their two proper stabilizers
do not cover the Galois group. Choose $g$ moving both.
The common-adjoint argument forces $g$ to fix the third,
so $\|W(gW)^*\|^2=2$. The unit is absent because $gW\ne W$.
Thus $W(gW)^*$ has exactly two distinct nonunit summands,
each of multiplicity one. Argument-side covariance gives
$\chi_W((gW)^*)=c\,g(\varepsilon)$, contradicting the
preceding lemma.
\end{proof}

\begin{proposition}[Global confinement of a degree-seventy-two constituent]
\label{paper:seventy-two-global-confinement}
If a degree-seventy-two simple occurs in $WW^*$ for a
degree-ten simple $W$, then necessarily
\[
                    WW^*=\mathbf1+Z_{15}+V_{12}+U_{72}.
\]
Both $Z$ and $V$ are fixed by Galois, whereas $U$ is not.
Their twists and the self-character satisfy
\[
       \theta_Z=-1,\qquad\theta_V=\theta_U=1,
       \qquad \chi_W(W^*)=7\theta_W^{-2}.
\]
At every compulsory degree-twenty witness $Y\in P_5$,
\[
                       |\chi_Y(W)|^2=25\pm9\sqrt5.
\]
\end{proposition}
\begin{proof}
The degree-ten adjoint palette leaves dimension twenty-seven
besides the degree-seventy-two term. Its only partition is
$15+12$, proving the displayed profile. All three constituents
are self-dual.

Put $t=\theta_Z$, $u=\theta_V$, and $v=\theta_U$.
Balancing and integrality give
\[
       T:=1+15t+12u+72v
                      =\theta_W^2S_{W,W^*}\in10\mathcal O.
\]
Modulo two this gives $1+t\in2\mathcal O$, so the unit-disk
norm argument forces $t=\pm1$. Modulo five it gives
\[
                        \frac{1+2u+2v}{5}\in\mathcal O.
\]
This algebraic integer has house at most one. The zero case
would make $u+v=-1/2$, impossible for an algebraic integer.
In the root-of-unity case, equality in the triangle inequality
forces $u=v=1$. If $t=1$, then $T=100$ contradicts strictness.
Therefore $t=-1$, $T=70$, and
$\chi_W(W^*)=7\theta_W^{-2}$.
Lemma~\ref{paper:integer-self-radius-rigidity} now makes at least
two of the three constituents fixed by Galois.

A fixed $U$ is impossible without any no-thirty hypothesis.
Indeed, at a compulsory $Y$, LCM integrality and the
principal-five residue would give $\chi_Y(U)=-18$.
Writing $\chi_Y(Z)=3z$ and $\chi_Y(V)=3v$, the full mean
bounds in Lemma~\ref{paper:short-real-quintic-means} give
\[
          0\le\mathsf A(|\chi_Y(W)|^2)
             \le1-18+9+6=-2.
\]
Thus it is exactly $Z,V$ that are fixed. Their characters
are consequently $\chi_Y(Z)=0$ and $\chi_Y(V)=-3$.
Write $\chi_Y(U)=18a$. Nonnegativity and strictness imply
\[
             |\chi_Y(W)|^2=18a-2,
          \qquad 1/9\le\sigma(a)<4\quad\text{for all }\sigma.
\]
The real algebraic integer $a-2$ has strict house below two,
so Kronecker's theorem gives $a=2+\eta+\eta^{-1}$ for a
root of unity $\eta$. Its residue $a\equiv4$ at five forces
$\eta$ to be five-primary. Strictness excludes order one.
The lower bound $a\ge1/9$ excludes every order $5^b\ge25$,
since then a conjugate is $2-2\cos(\pi/5^b)<1/9$.
Hence $\eta$ is primitive of order five, yielding
$18a-2=25\pm9\sqrt5$.
\end{proof}

\begin{corollary}[No degree-seventy-two adjoint constituent without thirty]
\label{paper:no-thirty-no-seventy-two}
In the no-thirty branch, the adjoint of a degree-ten simple
has no degree-seventy-two constituent. Equivalently, the full
profile $1+15+12+72$ is excluded there.
\end{corollary}
\begin{proof}
Apply the preceding proposition and
Corollary~\ref{paper:no-thirty-quadratic-norm-obstruction}.
\end{proof}


\begin{lemma}[A multiplicity-one bound in a tensor square]
\label{paper:braiding-square-bound}
Let $W$ be simple in a pseudo-unitary ribbon fusion category
whose twists are roots of unity. If a simple $A$ occurs with
multiplicity one in $W^2$, then
\[
                         d_A\le\frac{d_W^2+d_W}{2}.
\]
In particular, a self-dual degree-ten simple in the present
reduction has neither a degree-sixty nor a degree-seventy-two
constituent in its adjoint.
\end{lemma}
\begin{proof}
The ribbon partial-trace identity
$\operatorname{ptr}_R(c_{W,W})=\theta_W\operatorname{id}_W$
gives
\[
                  \operatorname{Tr}(c_{W,W})=d_W\theta_W;
\]
see \cite[Section~1.1]{GeerPatureau2018}. On the multiplicity
space of a simple summand $B$, balancing makes the square of
the braiding the scalar $\theta_B/\theta_W^2$. Its eigenvalues
are therefore roots of unity, and the absolute value of its
ordinary trace is at most $N_{WW}^B$. The multiplicity-one
summand $A$ contributes a term of absolute value exactly $d_A$.
The reverse triangle inequality gives
\[
             d_W=|\operatorname{Tr}(c_{W,W})|
                    \ge d_A-(d_W^2-d_A).
\]
This is the asserted bound. A constituent of dimension sixty
or seventy-two in a degree-ten square necessarily has
multiplicity one, and both dimensions exceed fifty-five.
\end{proof}

\begin{proposition}[Self-dual degree ten needs at least four adjoint terms]
\label{paper:selfdual-ten-four-terms}
The adjoint of a self-dual degree-ten simple has at least four
nonunit constituents, counted with fusion multiplicity.
This assertion does not assume the no-thirty branch.
\end{proposition}
\begin{proof}
Corollary~\ref{paper:degree-ten-adjoint-palette} requires at least three.
The only three-term
partitions of ninety-nine in its remaining palette are
\[
              (15,12,72),\quad(15,24,60),\quad
              (18,36,45),\quad(24,30,45).
\]
The first two are excluded for a self-dual $W$ by the preceding
lemma. The third is excluded for every $W$ by
Theorem~\ref{paper:forty-five-profile-excluded}.

For the last profile write
$W^2=\mathbf1+Z_{45}+A_{30}+B_{24}$. Balancing integrality
gives
\[
                 1+45t+30v+24u\in10\mathcal O,
       \qquad (t,u,v)=(\theta_Z,\theta_B,\theta_A).
\]
Modulo two, $(1+t)/2$ is integral of house at most one,
forcing $t=\pm1$. Modulo five, $(1+4u)/5$ is integral of
house at most one. The zero case would give $u=-1/4$; the
root-of-unity case forces $u=1$. Thus $t=\pm1$ and $u=1$.

The square is multiplicity-free. Multiply its braiding-trace
identity by $\theta_W$. Each summand then contributes its
dimension times a square root of its twist, and the right
side has absolute value ten. If $t=-1$, the degree-forty-five
contribution is $\pm45i$, while the unit and degree-twenty-four
contributions are real. The degree-thirty contribution cannot
cancel its imaginary part below fifteen, a contradiction.

If $t=1$, the identity has the form
\[
                 |a+30\rho|=10,
      \qquad |a|\in\{20,22,68,70\},
\]
where $a$ is an integer and $\rho$ a root of unity with
$\rho^2=v$. The values sixty-eight and seventy fail the
reverse triangle inequality. After changing signs, $a=22$
would force $\rho+\rho^{-1}=-107/55$, which is not an
algebraic integer. The case $a=20$ forces $\rho=-1$, hence
$v=1$. All four square constituents would then have twist
one, making $|\chi_W(W)|=10$ and contradicting strictness.
This excludes the last three-term profile.
\end{proof}

\begin{lemma}[A four-letter first-and-second-moment constraint]
\label{paper:four-root-moments}
Let $q_1,q_2,q_3,q_4$ be roots of unity and $d_i$ positive
integers satisfying
\[
                  \sum_i d_i=100,\qquad
                  \sum_i d_iq_i=10,\qquad
                  \sum_i d_iq_i^2=70.
\]
Up to permutation, the weighted phases are exactly one of
\[
\begin{array}{c|rrrr}
 &q_1&q_2&q_3&q_4\\\hline
\text{first type}&1&-1&\omega&\omega^{-1}\\
\text{dimensions}&50&30&10&10\\\hline
\text{second type}&1&-1&\xi_6&\xi_6^{-1}\\
\text{dimensions}&40&40&10&10,
\end{array}
\]
where $\omega$ and $\xi_6$ are primitive roots of orders
three and six, respectively.
\end{lemma}
\begin{proof}
For every nonreal root $q$, full normalized Galois averaging
satisfies
\[
             \mathsf A(q^2+q)\le\tfrac12,
         \qquad \mathsf A(q^2-q)\le\tfrac12.
\]
Indeed, for odd order $n>1$ the two means of $q,q^2$ coincide
with $\mu(n)/\varphi(n)$; when this is positive,
$\varphi(n)\ge8$. For order $2m$ with $m>1$ odd, the two
means are opposite, with the same bound. For order $4m$,
$m$ odd, the first mean is zero and the second is
$-\mu(m)/\varphi(m)\le1/2$ (the case $m=1$ gives $-1$).
For order divisible by eight both means vanish.

If no $q_i$ were one, averaging the sum of the two moment
identities would give $80\le50$. If none were minus one,
averaging their difference would give $60\le50$. Hence both
$1$ and $-1$ occur. All four phases cannot be real because
their second moment is seventy, and exactly one nonreal phase
cannot give a real first moment. Thus there is one phase of
each real sign and exactly two nonreal phases. Let their
weights be $a,b,c,d$ and their phases $1,-1,u,v$.

Write
\[
                    cu+dv=R:=10-a+b\in\mathbb Z.
\]
If $R=0$, then $c=d$ and $v=-u$. The second moment forces
$u^2$ real, hence $u^2=-1$. Combining the second moment and
$a+b+2c=100$ then gives $c=15/2$, impossible.
If $R\ne0$, taking absolute values of $R-cu=dv$ gives
\[
                   u+u^{-1}=\frac{R^2+c^2-d^2}{Rc}.
\]
This rational algebraic integer is $-1,0$, or $1$, since
$u$ is nonreal. The same holds for $v$. Thus their orders
belong to $\{3,4,6\}$. The cancellation of imaginary parts
excludes one order four and one order three or six, because
the required ratio of their positive integer weights would
be irrational. Two order-four roots give $R=0$, already
excluded. The remaining possibilities have $c=d$, opposite
imaginary parts, and, since $R\ne0$, equal real parts.
They are therefore conjugate primitive cube roots or conjugate
primitive sixth roots. Their squared contribution is $-c$.
The dimension and second-moment equations give $c=10$ and
$a+b=80$. The first moment then gives $(a,b)=(50,30)$ in
the cubic case and $(40,40)$ in the sixth-root case.
\end{proof}

\begin{proposition}[The tensor square in the surviving seventy-two case]
\label{paper:seventy-two-square}
Suppose a degree-ten simple has a degree-seventy-two adjoint
constituent. Then it is non-self-dual, and its tensor square
has exactly four distinct nonunit constituents, each with
multiplicity one. Their dimension multiset is either
\[
                         (10,10,30,50)
              \qquad\text{or}\qquad(10,10,40,40).
\]
In the first case the two degree-ten constituents have twists
$\theta_W^4\omega$ and $\theta_W^4\omega^{-1}$, and the
degree-thirty and degree-fifty constituents have twist
$\theta_W^4$. Moreover $\theta_{W,5}=1$.
In the second case the two degree-ten twists are again
$\theta_W^4\omega^{\pm1}$, the two degree-forty twists are
$\theta_W^4$, and $\theta_{W,2}=1$.
\end{proposition}
\begin{proof}
Self-duality is excluded by Lemma~\ref{paper:braiding-square-bound}.
The adjoint is multiplicity-free with four terms, so
\[
                         \|W^2\|^2=\|WW^*\|^2=4.
\]
The unit is absent from $W^2$. Its only possible multiplicity
patterns are four distinct summands or twice one degree-fifty
simple. The latter is impossible: the two eigenvalues of
$c_{W,W}$ on that multiplicity space are signs of the same
root, so its categorical trace would have absolute value zero
or one hundred, whereas the partial-trace identity requires ten.

Let $R_i$ be the scalar braidings on the four distinct square
summands, and put $q_i=R_i/\theta_W$. The partial trace gives
$\sum d_iq_i=10$. By Proposition~\ref{paper:seventy-two-global-confinement},
\[
                 S_{W,W^*}=70\theta_W^{-2},
             \qquad S_{W,W}=70\theta_W^2.
\]
Balancing gives $R_i^2=\theta_i/\theta_W^2$ and
$\sum_i d_iR_i^2=S_{W,W}$. Consequently
$\sum_i d_iq_i^2=70$. The preceding moment lemma yields the
two dimension patterns and the twists, since
$\theta_i=\theta_W^4q_i^2$.

In the first case the degree-fifty label has five-unit formal
codegree. Its fifth twist is one, so $\theta_{W,5}^4=1$ and
$\theta_{W,5}=1$.
In the second case both degree-forty labels have two-unit
formal codegree. If $W$ were dyadically active, their entries
against $W$ would vanish by Lemma~\ref{paper:family:lw:weyl-line}. Applying
the character at $W$ to $W^2$ would then express a number of
absolute value $|\chi_W(W)|^2=49$ as the sum of two
degree-ten character values, each of absolute value below
ten. This is impossible. Hence $\theta_{W,2}=1$.
\end{proof}


\begin{lemma}[Four-root moments with negative second moment]
\label{paper:negative-four-root-moments}
Let $q_i$ be four roots of unity and $d_i$ four positive
integers satisfying
\[
             \sum_i d_i=100,\qquad
             \sum_i d_iq_i=10,\qquad
             \sum_i d_iq_i^2=-50.
\]
Then, up to permutation,
\[
          (d_i,q_i)=(20,\omega),(20,\omega^{-1}),
                     (30,\xi_6),(30,\xi_6^{-1}),
\]
where $\omega$ and $\xi_6$ are primitive cube and sixth
roots, respectively.
\end{lemma}
\begin{proof}
We first exclude a phase $q_i=\pm i$. Work over the odd
cyclotomic field and expand the dyadic parts in the usual
basis of a two-power cyclotomic field. The coefficient of the
basis vector $i$ receives exactly the terms whose dyadic parts
are $\pm i$. Their weighted sum is zero, since the first
moment is rational. At most three of the four terms can lie
in that sector; all four would give first moment zero.

A positive rational relation of at most three roots in this
sector, containing $i$ or $-i$, is either a cancellation
between the two phases $i,-i$, or an equal-weight equilateral
triple $i\{1,\omega,\omega^2\}$, up to a common sign.
Here is an elementary justification of this short-relation
claim. After division by one phase, the root ratios have
dyadic part $\pm1$. In a three-term relation
$a+b\rho+c\sigma=0$, taking absolute values makes
$\rho+\rho^{-1}=(c^2-a^2-b^2)/(ab)$ a rational algebraic
integer. Apart from the collinear case, $\rho$ therefore has
order three, four, or six; order four is excluded by its
dyadic part. The other ratio belongs to the same quadratic
field. A positive relation among three sixth roots with no
opposite pair is an equal-weight equilateral triple. Two
terms must simply be opposite with equal weights.

If exactly two terms lie in the $i$ sector, they are $i,-i$
with equal weight $a$. The remaining two weighted roots have
rational first moment ten. The same rational-trace argument
shows that they are either both real, or a conjugate primitive
sixth-root pair of equal weight ten. The alternative primitive
cube-root pair would have negative first moment, and mixing
orders four and three or six would require an irrational ratio
of the integer weights. In the real case the dimension and
second-moment equations give $a=75/2$, impossible. In the
sixth-root case $a=40$ and the second moment is $-90$, again
impossible.

If three terms lie in the sector, the remaining term must be
the identity with weight ten. A collinear cancellation sector
then has total weight ninety and contributes $-90$ to the
second moment, giving total $-80$. An equilateral sector
contributes zero to that moment, giving total ten. Both
contradict the prescribed value $-50$. Thus no $q_i$ is $\pm i$.

For any root $q$ other than $\pm i$, the full Galois mean
of $q^2$ is at least $-1/2$. Equality occurs exactly when
$q^2$ is a primitive cube root: the root mean is
$\mu(n)/\varphi(n)$, whose only value below $-1/2$ is
$-1$ at the root $-1$. Averaging the second-moment identity
now gives equality in the bound
\[
                     -50\ge-\tfrac12\sum_i d_i=-50.
\]
Every $q_i$ is consequently primitive of order three or six.
Let $a,b$ be the total weights at the two primitive cube roots
and $c,d$ those at the two primitive sixth roots, with their
upper-half-plane phases listed first. The real first moment
and total mass give $a+b=40$, $c+d=60$. The two imaginary
moment equations are
\[
                      a-b+c-d=0,
                  \qquad -a+b+c-d=0.
\]
Thus $a=b=20$ and $c=d=30$. All four phase classes have
positive weight and there are exactly four terms, so there
is one term in each class, as asserted.
\end{proof}

\begin{proposition}[The square of a negative three-term profile]
\label{paper:negative-three-term-square}
Suppose a degree-ten simple has a multiplicity-free adjoint
with exactly three nonunit constituents of distinct degrees,
and
\[
                     \chi_W(W^*)=-5\theta_W^{-2}.
\]
Then $W$ is non-self-dual and
\[
                          W^2=A_{20}+A'_{20}+B_{30}+B'_{30},
\]
with four distinct simple objects, each occurring once.
The two actual twist classes are
$\theta_W^4\omega$ and $\theta_W^4\omega^{-1}$;
each class contains one degree-twenty and one degree-thirty
constituent.

In the no-thirty branch,
\[
                    \theta_{W,5}=1,
          \qquad \operatorname{ord}(\theta_{W,2})\mid4.
\]
Hence its complete twist is either three-primary or dyadic
of order at most four. If its dyadic order is four, then
$\chi_X(W)=-5$ at every compulsory ternary witness $X$.
\end{proposition}
\begin{proof}
Self-duality is excluded by
Proposition~\ref{paper:selfdual-ten-four-terms}. Rigidity gives
$\|W^2\|^2=\|WW^*\|^2=4$. If $W^2=2C_{50}$, the
partial trace of its braiding would make two normalized root-of-unity
eigenvalues sum to $1/5$, contrary to algebraic integrality.
Thus there are four distinct nonunit summands. Put $q_i=R_i/\theta_W$
for their scalar braidings. The partial trace gives
$\sum d_iq_i=10$, whereas the assumed self-character and
balancing give
\[
              \sum d_iq_i^2
                   =S_{W,W}/\theta_W^2=-50.
\]
Lemma~\ref{paper:negative-four-root-moments} gives the dimensions
and the twist classes.

In the no-thirty branch a degree-twenty simple cannot be
simultaneously active at three and five, since that would
force degree divisible by eight; nor can it be active at two
and three, since that would force degree divisible by
twenty-five. At least one of the two displayed twist classes
has nontrivial ternary component. Their common fifth component
is $\theta_{W,5}^4$, so it must be one. Their common dyadic
component is $\theta_{W,2}^4$, so it too must be one.
Since a degree-ten simple cannot be active at both two and
three in the no-thirty branch, the complete-twist assertion
follows. The final scalar assertion is
Corollary~\ref{paper:degree-ten-higher-dyadic-scalar}.
\end{proof}


\begin{proposition}[The degree-twenty-four, thirty, forty-five profile]
\label{paper:twentyfour-thirty-fortyfive}
Suppose $d_W=10$ and
\[
                 WW^*=\mathbf1+B_{24}+A_{30}+Z_{45}.
\]
Then $W$ is not self-dual, $\theta_B=1$, and
$\theta_Z\in\{1,-1\}$. Its self-character is
\[
 \chi_W(W^*)=\theta_W^{-2}
 \begin{cases}
       7+3\theta_A,&\theta_Z=1,\\
       -2+3\theta_A,&\theta_Z=-1.
 \end{cases}
\]
At least one of $B,A,Z$ is fixed by all Galois permutations.
If $\theta_Z=1$, at least two are fixed by all Galois
permutations. If $\theta_Z=-1$ and at least two are not
globally fixed, then $\theta_A=-1$, $B$ is globally fixed,
and every permutation moving two adjoint constituents gives
two distinct degree-fifty summands in each of $W(gW)$ and
$W(gW)^*$.
\end{proposition}
\begin{proof}
Non-self-duality is Proposition~\ref{paper:selfdual-ten-four-terms}.
Balancing gives
\[
         1+24\theta_B+30\theta_A+45\theta_Z\in10\mathcal O.
\]
Modulo two, $(1+\theta_Z)/2$ is an algebraic integer with
house at most one; the norm argument gives $\theta_Z=\pm1$.
Modulo five, $(1+4\theta_B)/5$ is integral and has house at
most one. It cannot vanish, and equality in the triangle
inequality gives $\theta_B=1$. This proves the character
formula, which is nonzero in both cases. Proposition~
\ref{paper:vafa-unique-fixed-adjoint} supplies a globally fixed
constituent.

If at least two constituents are not globally fixed, choose
a Galois permutation $g$ moving both. The common-adjoint
overlap argument gives exactly two distinct nonunit summands
in each of $W(gW)$ and $W(gW)^*$. Their possible degree
pairs are $(10,90)$, $(40,60)$ and $(50,50)$, whose squared
character moduli at $W$ have the respective forms
\[
                  82+9r,\qquad52+24r,\qquad50+25r,
\]
where $r=\rho+\rho^{-1}$ is a root-of-unity trace.

First let $\theta_Z=1$ and put
$s=g(\theta_A)+g(\theta_A)^{-1}$, so that the same squared
modulus is $58+21s$. Equality with the first degree pair
gives $7s-3r=8$. Hence $(s+1)/3$ is integral and all its
conjugates lie in $[-1/3,1]$; it is zero or one. The value
$s=-1$ would give $r=-5$, and $s=2$ gives the forbidden
strictness equality $\theta_A=1$.
Equality with the second pair gives $7s-8r=-2$, so
$(r-2)/7=s-r$ is an algebraic integer of house at most
$4/7$. Thus $r=s=2$, again impossible. Equality with the
third pair gives $21s-25r=-8$, making $(s-2)/5$ integral
of house at most $4/5$; again $s=2$. All three pairs are
impossible. Therefore at least two adjoint constituents
are globally fixed in this case.

Now let $\theta_Z=-1$. The squared modulus is $13-6s\le25$,
so $(10,90)$ is impossible by its lower bound sixty-four.
The pair $(40,60)$ would give
$s+4r=-13/2$, contrary to integrality. The pair $(50,50)$
gives $6s+25r=-37$. Consequently $(s+2)/5$ is integral
of house at most $4/5$, giving $s=-2$, $\theta_A=-1$
and $r=-1$. Thus every such cross-product uses the
degree-fifty pair. Its uniquely fixed adjoint constituent
has twist one by Proposition~\ref{paper:vafa-unique-fixed-adjoint},
so it must be $B$. If $B$ were not globally fixed and one
of $A,Z$ were not, a permutation outside their two proper
stabilizers would uniquely fix the other twist-minus-one
constituent, a contradiction. Since at least two constituents
are nonfixed, $B$ must therefore be globally fixed.
\end{proof}

\begin{corollary}[Rational degree-twenty-four endpoint values]
\label{paper:rational-twentyfour-values}
In the profile of Proposition~\ref{paper:twentyfour-thirty-fortyfive},
if $B$ is globally fixed, then
for every compulsory $X\in P_3$ of degree twelve or thirty-six
and $Y\in P_5$ of degree twenty,
\[
                    \chi_X(B)=\chi_Y(B)=-6.
\]
The same conclusion holds for a globally fixed degree-twenty-four
constituent in the profile $(15,24,60)$.
If all three constituents in the $(24,30,45)$ profile are
globally fixed, then
\[
 \begin{split}
       &\chi_Y(Z)=0,\qquad\chi_Y(A)=15,
         \qquad |\chi_Y(W)|^2=10,\\
       &(\chi_X(Z),\chi_X(A))\in\{(15,15),(30,0)\}.
 \end{split}
\]
In the no-thirty branch, this last case forces
$\mathbb Q(\zeta_{15})\subseteq K_Y$.
\end{corollary}
\begin{proof}
Global fixedness makes the displayed character values rational.
LCM integrality and principality at three make $\chi_X(B)$
a multiple of six; strictness gives $|\chi_X(B)|<24$.
The mass-two ternary lemma gives $\chi_X(WW^*)=25$.
The other two nonunit constituents have character values
in $5\mathcal O$, so $\chi_X(B)\equiv24\pmod5$.
The unique multiple of six in $(-24,24)$ with that residue
is $-6$. The same argument uses the degrees fifteen and
sixty in the other profile. At $Y$, LCM integrality gives
$\chi_Y(B)\in6\mathbb Z$, and principality at five gives
the same residue twenty-four, hence again $-6$.

When $Z,A$ are also globally fixed, LCM integrality,
principality and strictness give $\chi_Y(Z)=0$ and
$\chi_Y(A)\in\{-15,0,15\}$. The norm identity
$|\chi_Y(W)|^2=-5+\chi_Y(A)\ge0$ forces the value fifteen.
At $X$ both remaining values are multiples of fifteen,
strictly smaller in modulus than their degrees, and sum
to thirty. The two asserted possibilities follow.
The last assertion is Theorem~\ref{paper:norm-ten-endpoint-field}.
\end{proof}

\subsection{Four-term squares and Vafa twist rigidity}

\begin{lemma}[A Vafa product constraint on a four-term square]
\label{paper:four-square-vafa-product}
Let $W$ be a degree-ten simple with
$WW^*=\mathbf1+Z_1+Z_2+Z_3$, with three distinct-degree
nonunit constituents, and suppose $W$ is not self-dual.
Then $W^2$ has four distinct nonunit summands $A_i$.
Write $R_i$ for the scalar braiding on $A_i$ and put
$q_i=R_i/\theta_W$. Then
\[
 \sum_i d_{A_i}=100,\qquad \sum_i d_{A_i}q_i=10,
 \qquad
 \left(\prod_i q_i\right)^2
       =\left(\prod_{j=1}^3\theta_{Z_j}\right)^{-2}.
\]
Every $d_{A_i}$ belongs to
$\{10,12,15,18,20,24,30,36,40,45,50,54\}$.
\end{lemma}
\begin{proof}
The common-adjoint identity gives $\|W^2\|^2=4$.
The alternative $W^2=2A_{50}$ is excluded by the
braiding partial trace: two normalized eigenvalues would sum to the
nonintegral rational number $1/5$. Thus there are
four distinct summands, and their dimensions are at most
fifty-five by the multiplicity-one bound. The degree palette
follows from divisibility by $5400$ and the absence of
prime-power nonunit degrees and degree six. The partial
trace supplies the first moment.

Apply the four-point Vafa identity with
$(i,j,m,r)=(W,W^*,W,W)$. The fusion-space dimension
is four. The two adjoint factors have exponent two at each
$Z_j$, and the square factor has exponent one at each $A_i$.
Consequently
\[
 \theta_W^{16}
      =\left(\prod_j\theta_{Z_j}\right)^2
                                      \prod_i\theta_{A_i}.
\]
Balancing gives $\theta_{A_i}=\theta_W^4q_i^2$.
Canceling $\theta_W^{16}$ proves the product relation.
\end{proof}

\begin{proposition}[The full square table for the forty-five profile]
\label{paper:global-fortyfive-square-table}
\label{paper:fortyfive-square-table}
Suppose
\[
 WW^*=\mathbf1+B_{24}+A_{30}+Z_{45},\qquad
 \theta_B=1,\quad t=\theta_Z\in\{1,-1\},\quad v=\theta_A.
\]
There is no further assumption on the root of unity $v$.
If $t=-1$, then $v=1$ or $-1$. If $t=1$, then $v\ne1$.
The possible dimension and normalized braiding-phase multisets
of $W^2$ are precisely the following necessary alternatives,
up to permutation and complex conjugation:
\[
\begin{array}{c|c|c|c}
t&v&(d_{A_i})&(q_i)\\\hline
1&v\ne1&(40,30,15,15)&(1,-1,u,-u),\quad u^2=\overline v\\
1&-1&(45,15,20,20)&(1,-1,\omega,\omega^{-1})\\
-1&1&(10,30,30,30)&(1,1,\omega,\omega^{-1})\\
-1&1&(10,30,30,30)&(1,-1,\xi_6,\xi_6^{-1})\\
-1&1&(20,20,30,30)&(1,1,\omega,\omega^{-1})\\
-1&-1&(20,20,30,30)&(\omega,\omega^{-1},\xi_6,\xi_6^{-1}).
\end{array}
\]
Here $\omega$ and $\xi_6$ are primitive cube and sixth roots.
The two degree-fifteen constituents in the first row have
actual twist $\theta_W^4\overline v$; the two real-braiding
constituents have actual twist $\theta_W^4$.
No occurrence of a listed alternative is asserted. The theorem
holds throughout the ordinary reduction, including thirty activity.
\end{proposition}
\begin{proof}
The degree-ten short-adjoint theorem excludes self-duality of
$W$ and gives four distinct nonunit square constituents. The
braiding partial trace, self-character formula and Vafa product
give, with $a=7$ for $t=1$ and $a=-2$ for $t=-1$,
\begin{equation}\label{paper:global-fortyfive-moments}
 \sum_i d_i=100,\qquad \sum_i d_iq_i=10,\qquad
 \sum_i d_iq_i^2=L:=10a+30\overline v,
 \qquad \left(\prod_iq_i\right)^2=v^{-2}.
\end{equation}
The multiplicity-one square bound puts each $d_i$ in
\[
 \{10,12,15,18,20,24,30,36,40,45,50,54\}.
\]
Every $q_i$ has only the primes two, three and five in its order.
The case $t=v=1$ gives the strictness equality and is excluded.

We give the five-term first-moment reduction in full. Regard
$\sum_i d_iq_i-10=0$ as a rational root relation with five
terms. If it is irredundant, choose the least conductor of its
root ratios; the constant term makes this also a conductor
for its roots. A prime-power exponent at least two gives a
degree-$p$ extension on lowering that exponent. Splitting by
its power-basis sectors makes each occupied sector a vanishing
subrelation. Irredundancy therefore removes every such exponent.
If the fifth prime remains, expansion by the five powers of
a fifth root over $\mathbb Q(\mu_6)$ makes the five coefficients
equal. An empty sector contradicts irredundancy. Otherwise
the five terms occupy five different sectors, one each, and
have equal absolute coefficient ten. This contradicts
$\sum_i d_i=100$. Thus all roots in an irredundant relation
belong to $\mu_6$.

If the relation is reducible, its five nonzero terms partition
into a vanishing pair and a vanishing triple. Suppose first
that the pair uses two square phases. They are opposite with
equal weight $c$. The remaining two weighted phases have
rational sum ten. The elementary two-root rational-trace
argument shows that either they are real, with weights
$b+10,b$ and phases $1,-1$, or they are the conjugate
primitive sixth roots, both with weight ten. To recall the
argument, $hr+ks=10$ gives
$r+r^{-1}=(100+h^2-k^2)/(10h)$. A rational root trace
is one of $-2,-1,0,1,2$; imaginary-part cancellation leaves
exactly the two stated positive-sum alternatives.

In the real case $b+c=45$, and the degree palette gives
only $(b+10,b,c,c)=(40,30,15,15)$. Its second moment is
$70+30u^2$. For $t=1$ this gives $u^2=\overline v$ and
the first table row. Its Vafa product relation holds because
the squared unweighted phase product is $u^4=v^{-2}$.
For $t=-1$ it would give $u^2=-3+\overline v$, whose
modulus exceeds one. In the nonreal case $c=40$ and the
second moment is $-10+80u^2$. If $t=-1$, then
$|80u^2|=|-10+30\overline v|\le40$, impossible. If
$t=1$, then $u^2=1+3\overline v/8$; its modulus-one
condition would give $v+v^{-1}=-3/8$, impossible for an
algebraic integer.

If the vanishing pair uses the constant term, there is a
weight-ten phase one and the other three terms sum to zero
with mass ninety. A three-term rational root relation is
collinear or an equal-weight rotated equilateral triple;
this follows by the same conductor reduction, now with three
terms. In the collinear case the second moment is
$10+90u^2$. For $t=1$ one has
$|L-10|=|60+30\overline v|<90$, since $v\ne1$;
for $t=-1$ it is at most sixty. Thus this case is impossible.
In the equilateral case the second moment is ten. This forces
$t=-1,v=1$, and the Vafa relation gives $u^6=1$ for the
common rotation of the three weight-thirty phases. These
are the two table rows of dimensions $(10,30,30,30)$.

It remains to solve the case $q_i\in\mu_6$. Equation
\eqref{paper:global-fortyfive-moments} puts $v$ in
$\mathbb Q(\omega)$, so $v\in\mu_6$. Let
$a_0,b_0,c,d,e,f$ be the total weights at
$1,-1,\omega,\omega^{-1},\xi_6,\xi_6^{-1}$.
Put $S=a_0+b_0$, $C=c+d$, and $E=e+f$. Then
\[
 S+C+E=100,\qquad S=\frac{100+2\Re L}{3},\qquad
 a_0-b_0+\frac{E-C}{2}=10,
\]
\[
 c-d+e-f=0,\qquad
 -c+d+e-f=\frac{2\Im L}{\sqrt3}.
\]
There are at most four occupied classes, each having weight
at least ten.

For $t=1,v=-1$, the equations give either real weights
$(45,15)$ and a weight-twenty conjugate cubic pair, or real
weights $(25,35)$ and a weight-twenty conjugate sixth-root
pair. The latter degrees are forbidden. This gives the second
table row. For $t=1,v=\omega$, the class weights are
$(a_0,b_0,c,d,e,f)=(40,30,15,0,0,15)$, already in the
first row. Conjugation treats $v=\omega^{-1}$. For a
primitive sixth-root $v$, one gets $C+E=10$, while the
imaginary equations require $|c-d|=|e-f|=15$, impossible.

For $t=-1,v=\omega$, the equations give $S=10$,
$C+E=90$, $c-d=15$ and $e-f=-15$. The single real
class has weight ten; either sign for that class leaves
four occupied nonreal classes, contrary to the four-summand
bound. Conjugation again treats $v=\omega^{-1}$.
For $t=-1$ and primitive sixth-root $v$, the corresponding
data are $S=30$, $C+E=70$, with the same absolute
differences fifteen. If both real classes occur, there is
room for at most two nonreal classes; that would require
$C=E=15$, contrary to their sum seventy. With a single
real class, the negative real class gives a negative total
weight. The positive real class forces occupied weights
$(30,35,20,15)$, including the forbidden degree thirty-five.
Thus only $v=1,-1$ remain in the negative branch.

For $t=-1,v=1$, the occupied-class condition gives either
$(a_0,b_0,C,E)=(10,30,0,60)$, already an equilateral
row, or $(40,0,60,0)$. In the latter case one of the three
classes of weights $40,30,30$ must split into two distinct
simple constituents. Splitting a cubic class violates the
Vafa product relation. Splitting the real class leaves
$40=10+30$ or $20+20$, giving the other two rows with
$v=1$. Finally $t=v=-1$ is exactly
Lemma~\ref{paper:negative-four-root-moments}, which gives
the last table row. Balancing
$\theta_{A_i}=\theta_W^4q_i^2$ proves the actual-twist
statement and completes the proof.
\end{proof}


\begin{remark}
The finite $\mu_6$ part admits an exact arithmetic enumeration. The
proof above also covers the opposite-pair family at unrestricted
conductor. The listed rows are necessary arithmetic possibilities;
no categorification is asserted.
\end{remark}

\begin{theorem}[The twists in a fifteen, twenty-four, sixty adjoint]
\label{paper:fifteen-twentyfour-sixty-twists}
Suppose a degree-ten simple satisfies
\[
                        WW^*=\mathbf1+Z_{15}+U_{24}+V_{60}.
\]
Then
\[
                 \theta_Z=-1,\qquad \theta_U=1,
                         \qquad\theta_V\in\{1,-1\}.
\]
This conclusion holds throughout the ordinary reduction,
including the thirty-active branch. In particular the
positive $\theta_Z=1$ branch is entirely excluded.

If $\theta_V=1$, then
$\chi_W(W^*)=7\theta_W^{-2}$ and at least two of the
three adjoint constituents are globally Galois fixed.
The square has degree multiset $(10,10,30,50)$ or
$(10,10,40,40)$, with the two degree-ten twists
$\theta_W^4\omega$ and $\theta_W^4\omega^{-1}$.
If $\theta_V=-1$, the self-character and square are those
of Proposition~\ref{paper:negative-three-term-square}: the self-character is
$-5\theta_W^{-2}$ and the square degrees are
$(20,20,30,30)$.
\end{theorem}
\begin{proof}
Proposition~\ref{paper:fifteen-twentyfour-sixty} gives
$\theta_U=1$, $\theta_Z=t\in\{1,-1\}$,
and excludes self-duality of $W$. Put $v=\theta_V$.
The four square summands and their normalized scalar
braidings satisfy
\[
 \sum_i d_i=100,\quad \sum_i d_iq_i=10,\quad
 \sum_i d_iq_i^2=L:=10a+60\overline v,
 \quad \left(\prod_iq_i\right)^2=v^{-2},
\]
where $a=4$ if $t=1$, and $a=1$ if $t=-1$.
The last identity is Lemma~\ref{paper:four-square-vafa-product}.
The case $t=v=1$ violates strictness and is omitted.

Apply the five-term first-moment reduction from the proof
of Proposition~\ref{paper:fortyfive-square-table}. Either
all $q_i$ lie in $\mu_6$, or a vanishing pair and triple
give one of the alternatives analyzed next. This reduction
uses only the first moment, total dimension and degree palette.

An opposite pair not involving the constant term leaves
either the real pair of weights forty and thirty, with
the opposite pair having weights fifteen, or a conjugate
sixth-root pair of weights ten, with the opposite pair
having weights forty. In the first case
$L=70+30u^2$. If $t=1$, this gives
$u^2=-1+2\overline v$, whose modulus is one only if
$v=1$, already excluded. If $t=-1$, it gives
$u^2=-2+2\overline v$, whose modulus-one condition
would require $v+v^{-1}=7/4$, impossible for an algebraic
integer. In the second case $L=-10+80u^2$. For $t=1$
the modulus-one condition gives $v+v^{-1}=1/10$,
again impossible. For $t=-1$ it forces $v=1$ and
$u^2=1$, so all four phases are already in $\mu_6$.

If the vanishing pair uses the constant term, there is a
weight-ten phase one, and the other three terms sum to
zero with mass ninety. A collinear triple has second
moment $90u^2$, but $|L-10|<90$: for $t=1$ this is
the strict triangle inequality $|30+60\overline v|<90$,
and for $t=-1$ it equals sixty. An equilateral triple
has second moment zero and would give $L=10$, which
would require $v=-1/2$ or $v=0$. Neither is a root of
unity. This exhausts the pair/triple alternatives.

Thus all four $q_i$ belong to $\mu_6$. Their second
moment belongs to $\mathbb Q(\omega)$, so $v$ belongs
to that field and is itself in $\mu_6$. We solve the
six-class moment equations, writing $a_0,b_0,c,d,e,f$
for the total weights at
$1,-1,\omega,\omega^{-1},\xi_6,\xi_6^{-1}$.
Put $S=a_0+b_0$, $C=c+d$, $E=e+f$.

If $t=1,v=-1$, then $L=-20$, so $S=20$ and
$C+E=80$. Real first moment ten gives
$C=a_0-b_0+30$. Since every occupied real class has
weight at least ten, $(a_0,b_0)$ is $(20,0)$,
$(0,20)$, or $(10,10)$. Each choice has $C,E>0$.
Both nonreal pairs have equal weights because $L$ is
real, so there are at least four nonreal classes as well
as a real class. This contradicts four square summands.

If $t=1,v=\omega$, then $L=40+60\omega^{-1}$.
The equations give $S=40$, $C+E=60$,
$c-d=30$, and $e-f=-30$. Consequently the four
occupied classes have weights
\[
                (a_0,b_0,c,f)=(25,15,30,30).
\]
Each must be one simple summand, but degree twenty-five
is excluded. Conjugation deals with $v=\omega^{-1}$.
If $t=1$ and $v$ is primitive sixth, the real second
moment gives $S=80$, hence $C+E=20$, whereas its
imaginary part requires both $C,E\ge30$. This is
impossible. These cases exclude the whole positive branch.

Now let $t=-1$. For $v=\omega$ the equations give
$S=20$, $C+E=80$, $c-d=30$, $e-f=-30$, and
$C=a_0-b_0+30$. The only solution with at most four
occupied classes has weights and phases
\[
                (20,40,10,30),\qquad
                (1,\omega,\omega^{-1},\xi_6^{-1}).
\]
Its squared unweighted phase product is $\omega^{-1}$,
but the Vafa relation requires $v^{-2}=\omega$.
This contradiction excludes $v=\omega$, and conjugation
excludes $v=\omega^{-1}$. For primitive sixth $v$,
one gets $S=60$ and $C+E=40$ while still requiring
$C,E\ge30$, again impossible. Thus only $v=1,-1$
remain, proving the twist statement.

When $v=1$, the self-character has integer modulus seven.
Lemma~\ref{paper:two-summand-integer-radii}, for two-summand
degree-ten products, rules out any Galois permutation moving two
adjoint constituents. The union of two proper stabilizers
cannot cover the Galois group, so at least two constituents
are globally fixed. Lemma~\ref{paper:four-root-moments}, with
second moment seventy, gives the two square patterns.
When $v=-1$, its negative moment lemma gives the final
assertions.
\end{proof}

\begin{corollary}[The quintic norms in an all-fixed sixty profile]
\label{paper:all-fixed-sixty-quintic-norm}
If all three adjoint constituents in the theorem are
globally Galois fixed, then at a compulsory degree-twenty
endpoint $Y$,
\[
 \chi_Y(Z)=0,\qquad\chi_Y(U)=-6,\qquad
               |\chi_Y(W)|^2\in\{10,25,40\}.
\]
In the no-thirty branch the norm-forty theorem removes
forty, leaving only ten or twenty-five. The norm-ten case
has $\theta_W=1$ by Theorem~\ref{paper:norm-ten-untwisted};
the norm-twenty-five case has $\chi_Y(W)=5\rho$ for a
root of unity by Proposition~\ref{paper:norm-twentyfive-boundary}.
If a compulsory degree-thirty-six witness exists in the
no-thirty branch, the norm-ten value is excluded as well;
only $|\chi_Y(W)|^2=25$ remains.
\end{corollary}
\begin{proof}
Rationality, LCM integrality, principality at five and
strictness force $\chi_Y(Z)=0$; Corollary~\ref{paper:rational-twentyfour-values} gives
$\chi_Y(U)=-6$. They also make $\chi_Y(V)$ an integer
multiple of fifteen, strictly between minus sixty and
sixty. Nonnegativity of
$|\chi_Y(W)|^2=\chi_Y(V)-5$ leaves precisely
$\chi_Y(V)=15,30,45$. The additional conclusions are
the stated no-thirty norm results.
The final claim is
Corollary~\ref{paper:thirtysix-forbids-norm-ten}.
\end{proof}

\begin{corollary}[The all-fixed forty-five profile without thirty]
\label{paper:all-fixed-twentyfour-thirty-fortyfive}
\label{paper:all-fixed-fortyfive-twists}
Assume the no-thirty branch and suppose all three nonunit
constituents in the $(24,30,45)$ adjoint profile are globally
Galois fixed. Then $\theta_W=1$, and
\[
 \begin{array}{c|c}
 \theta_{45}&\text{possible }\theta_{30}\\\hline
 1&-1,\omega,\omega^{-1}\\
 -1&1,-1.
 \end{array}
\]
The tensor square belongs to the necessary table in
Proposition~\ref{paper:fortyfive-square-table}. If a compulsory
degree-thirty-six witness exists, this entire all-fixed
profile is impossible.
\end{corollary}
\begin{proof}
Corollary~\ref{paper:rational-twentyfour-values} gives
$|\chi_Y(W)|^2=10$. Theorem~\ref{paper:norm-ten-untwisted}
therefore gives $\theta_W=1$ and
$\mathbb Q(\sqrt{-15})\subseteq K_W$.

Every globally fixed adjoint label $A$ has rational character
value at $W$: symmetry expresses it as a rational multiple of
the rational value $\chi_A(W)$. Hence
$|\chi_W(W)|^2=\chi_W(WW^*)$ is rational. The self-character
formula makes $\theta_{30}+\theta_{30}^{-1}$ a rational
algebraic integer in $[-2,2]$, so the order of $\theta_{30}$
is one of $1,2,3,4,6$. Order six is impossible: a degree-thirty
label active at both two and three would have degree divisible
by twenty-five.

If the order were four, the self-character, now without an
endpoint scalar, would be $a\pm3i$ with $a=7$ or $-2$.
Thus $\mathbb Q(i)\subseteq K_W$. Together with
$\mathbb Q(\sqrt{-15})\subseteq K_W$, this gives ramification
at all three dimension primes and a pure dyadic quadratic
subfield. It contradicts the full-field alternatives in
Proposition~\ref{paper:no-thirty-field-alternatives}, since
none of their three exceptional fields contains $\mathbb Q(i)$.
Order four is excluded. Strictness excludes
$\theta_{45}=\theta_{30}=1$, and Proposition~
\ref{paper:fortyfive-square-table} excludes the negative
degree-forty-five cases with cubic degree-thirty twist.
This proves the table. The degree-thirty-six exclusion is
Corollary~\ref{paper:thirtysix-forbids-norm-ten}.
\end{proof}


\subsection{Further restrictions on the full graph}

\begin{proposition}[Torus fields and block motion when thirty is absent]
\label{paper:no-thirty-block-stability}
Suppose that no actual simple twist is divisible by thirty.  Write
$W=\gamma^3W_2$ for the corrected dyadic Weyl operator, and write
$H_p$ for the unsigned $p$-local torus action.  Then:
\begin{enumerate}
\item For every dyadic group element $g_2$, ternary group element
$g_3$, and five-local group element $g_5$, in the corrected local
representations,
\begin{equation}\label{paper:eq:no-triple-operator}
                (g_2-I)(g_3-I)(g_5-I)=0.
\end{equation}
\item If a full character field $K_X$ contains
$\mathbb Q(\zeta_{15})$, then $K_X$ is unramified at two.
\item Let $\mathcal B$ be a full dyadic residue block containing a
label $X$ with $2\mid f_X$.  Either every ternary torus preserves
$\mathcal B$ or every five-local torus preserves $\mathcal B$.
If some ternary torus moves $\mathcal B$, then
$(W-I)\mathbf e_X$ is fixed by all five-local tori; the assertion
with three and five interchanged also holds.
\end{enumerate}
\end{proposition}
\begin{proof}
Lemma~\ref{paper:no-thirty-local-separation} gives $\gamma=\pm1$,
so the corrected dyadic representation has upper operator equal
to the actual dyadic twist. If $\gamma=-1$, this correction is
tensoring by the sign character inflated from
$\operatorname{SL}_2(\mathbb F_2)$, so it remains a representation.
Decompose the representation of the
product of the three finite local groups into irreducible tensor
products.  A nontrivial irreducible factor at any prime has a
nonidentity upper eigenvalue: otherwise the upper unipotents and
their conjugates would act trivially and the factor would be
trivial.  A tensor product with all three factors nontrivial would
therefore have an actual twist divisible by thirty.  Every
irreducible tensor product has a trivial factor, proving
\eqref{paper:eq:no-triple-operator}.

The torus orbit of $X$ is a transitive permutation representation
of an abelian group.  Its occurring characters are exactly the
characters of the quotient corresponding to $K_X$.  If
$\mathbb Q(\zeta_{15})\subseteq K_X$, this character group contains
the pure ternary sign and the pure fifth characters of order four.
Were $K_X$ ramified at two, an occurring character would have a
nontrivial dyadic component.  Multiplying it by a suitable pure
ternary sign and by a suitable pure fifth character makes all
three components nontrivial.  Such a torus character can occur
only in an irreducible tensor product with three nontrivial local
factors, contrary to the first part.  Thus $K_X$ is unramified at
two.

For the block assertion put $v=(W-I)\mathbf e_X$ and
$\delta(\mathbf e_V)=d_V$.  Local block invariance gives
$\operatorname{supp}(v)\subseteq\mathcal B$.  Integral
complementary vacuum coefficients, with the immaterial correction
sign absorbed, give
\[
       \delta v=(k-1)d_X=:c,\qquad k\in2\mathbb Z.
\]
In particular $c\ne0$.  Foreign tori preserve dimensions and
transport full dyadic blocks, so $\delta(gv)=c$ and
$\operatorname{supp}(gv)\subseteq g\mathcal B$.

For a ternary torus $g$ and a five-local torus $h$, the first
part gives
\[
                       ghv-gv-hv+v=0.
\]
Apply $\delta P_{\mathcal C}$ for any full dyadic block
$\mathcal C$, and divide by $c$.  The result is the unsigned
block identity
\[
 \mathbf1_{gh\mathcal B=\mathcal C}
 -\mathbf1_{g\mathcal B=\mathcal C}
 -\mathbf1_{h\mathcal B=\mathcal C}
 +\mathbf1_{\mathcal B=\mathcal C}=0.
\]
If both $g$ and $h$ moved $\mathcal B$, its specialization to
$\mathcal C=\mathcal B$ would read $1=0$ or $2=0$.
Thus one of the two complete torus groups preserves the block.
If $g\mathcal B\ne\mathcal B$, every $h$ preserves it, and
the vector identity reads $g(hv-v)=hv-v$.  The two sides are
supported in disjoint full blocks, so $hv=v$.  This proves the
last assertion, and the other case is symmetric.
\end{proof}


\begin{lemma}[Three-coordinate support lemma]
\label{paper:three-coordinate-support}
Let $A_1,A_2,A_3$ be abelian groups and let
$H\le A_1\times A_2\times A_3$. Suppose that every element of
$H$ has an identity coordinate. Then either one coordinate
projection of $H$ is trivial, or there are nonidentity elements
$a_i\in A_i$ of order two such that
\[
 H=\{(1,1,1),(a_1,a_2,1),(a_1,1,a_3),(1,a_2,a_3)\}.
\]
\end{lemma}
\begin{proof}
Suppose all three projections are nontrivial. No element may
have precisely one nonidentity coordinate: multiplying such an
element by an element nontrivial at a second coordinate, and
then one nontrivial at the third if needed, would produce an
element nontrivial in all three coordinates. More explicitly,
if $(a,1,1)\in H$ with $a\ne1$, then every element with nontrivial
third coordinate has trivial second coordinate (otherwise it,
or its product with $(a,1,1)$, has full support); similarly every
element with nontrivial second coordinate has trivial third
coordinate. The product of one of each kind then has both
coordinates nontrivial, a contradiction.

Thus every nonidentity element has support of size two. After permuting the coordinates, choose
$x=(a,b,1)$ and $y=(c,1,d)$ with $a,b,c,d\ne1$; these exist
because all projections are nontrivial and singleton support is
excluded. Both $xy$ and $xy^{-1}$ have nontrivial second and
third coordinates, so $ac=ac^{-1}=1$. Hence $a=c$ has order two.
Applying the same argument to $x,xy$ and to $y,xy$ shows that
$b,d$ have order two. If another element has the same support
as $x$, multiplying it by $y$ forces its first coordinate to
equal $a$, and multiplying it by $xy$ forces its second
coordinate to equal $b$. The other two supports are identical
by symmetry. These four elements exhaust $H$.
\end{proof}

\begin{proposition}[Complete character-field alternatives without thirty]
\label{paper:no-thirty-field-alternatives}
Assume the hypotheses at $D=5400^2$ and suppose no actual simple
twist has order divisible by thirty. For every simple label $Z$,
either $K_Z$ is unramified at at least one of $2,3,5$, or
\[
 K_Z=\mathbb Q(\sqrt{-3a},\sqrt{5a}),
 \qquad a\in\{-1,2,-2\}.
\]
In the second alternative $K_Z$ is biquadratic, has no wild
ramification at three or five, and contains none of the three
individual primary quadratic fields
$\mathbb Q(\sqrt a)$, $\mathbb Q(\sqrt{-3})$,
$\mathbb Q(\sqrt5)$.
\end{proposition}
\begin{proof}
Take a splitting field of primary level
$2^{b_2}3^{b_3}5^{b_5}$ large enough for the local action.
The torus orbit of $Z$ is the transitive permutation
representation corresponding to $K_Z$. Its occurring characters
form the subgroup $H_Z$ of Dirichlet characters trivial on the
stabilizer of $Z$. The no-thirty operator identity implies that
no occurring character is nontrivial on all three local torus
groups. Lemma~\ref{paper:three-coordinate-support} therefore
applies to $H_Z$.

A trivial local projection means that the corresponding inertia
group acts trivially on $K_Z$, which is the first alternative.
Otherwise $H_Z$ has four elements with three nontrivial local
quadratic characters. At the odd primes these are uniquely the
characters of $\mathbb Q(\sqrt{-3})$ and
$\mathbb Q(\sqrt5)$. At two the nontrivial quadratic characters
of two-power conductor correspond to
$\mathbb Q(\sqrt a)$ for $a=-1,2,-2$. The three nontrivial
characters of $H_Z$ are the pairwise products of these primary
characters. Their fields generate exactly the displayed
biquadratic field. No individual primary character belongs to
$H_Z$, proving the last assertion.
\end{proof}

\begin{corollary}[Reality and activity in the exceptional field case]
\label{paper:no-thirty-exceptional-field-scope}
Every self-dual simple has a full character field unramified at
at least one of $2,3,5$. Every label in the exceptional
biquadratic case of Proposition~\ref{paper:no-thirty-field-alternatives}
has complete twist one.
\end{corollary}
\begin{proof}
Each displayed biquadratic field contains
$\mathbb Q(\sqrt{-15})$, whereas the field of a self-dual
label is real. This proves the first assertion.
Fix a prime $p$ at which $Z$ is active and restrict to its exact
nontrivial $p$-primary twist fiber. The other two local groups
act there with orthogonal moving spaces, by the no-thirty
hypothesis. Their unsigned torus permutations consequently
satisfy $(H_q-I)(H_r-I)=0$ on that fiber. If some $q$-torus
and some $r$-torus both moved $Z$, the coefficient of
$\mathbf e_Z$ in this identity applied to $\mathbf e_Z$
would be one or two, a contradiction. Thus one of these two
whole inertia groups fixes $Z$. Its field is unramified at
that complementary prime, contradicting the exceptional case.
So $Z$ is inactive at every prime.
\end{proof}

\begin{corollary}[Degrees in the exceptional field case]
\label{paper:exceptional-field-degrees}
In the exceptional field case, $d_Z\mid180$ if $a=-1$, and
$d_Z\mid90$ if $a=2$ or $-2$.
\end{corollary}
\begin{proof}
The three nontrivial quadratic subfields have discriminants
whose least common multiple is $60$ for $a=-1$ and $120$
for $a=\pm2$. These are the respective field conductors.
The exact conductor theorem makes them divide
$f_Z=2^6 3^6 5^4/d_Z^2$. The prime valuations of $f_Z$ are
even, so the first conductor gives heights at most $(2,2,1)$,
and the second at most $(1,2,1)$. This proves the assertions.
\end{proof}

\begin{corollary}[A field dichotomy at every compulsory witness]
\label{paper:compulsory-field-dichotomy}
In the no-thirty branch, let $X\in P_3$ have degree twelve or
thirty-six and let $Y\in P_5$ have degree twenty, with the
compulsory pure-primary twists. Then each of $K_X,K_Y$ is
unramified at two or at three. More generally this conclusion
holds for every full character field containing
$\mathbb Q(\zeta_5)$.
\end{corollary}
\begin{proof}
The compulsory crossing is ordinary and its entry field is
$\mathbb Q(\zeta_5)$, which lies in both endpoint fields.
The fifth coordinate of its Dirichlet-character group therefore
contains a character of order four. This excludes the
biquadratic alternative of the proposition. Its fifth
projection is nontrivial, so one of its two other projections
must vanish.
\end{proof}

\begin{corollary}[A four-field bound at a degree-thirty-six crossing]
\label{paper:ordinary36-field-list}
In the no-thirty branch, suppose a compulsory crossing has
degrees $(36,20)$. Set $K=\mathbb Q(\zeta_5)$,
$a=[K_X:K]$, and $b=[K_Y:K]$. Then $ab\le4$, and each
endpoint field is one of
\[
 K,\qquad K(i),\qquad \mathbb Q(\zeta_{15}),\qquad
 K\,\mathbb Q(\zeta_9)^+.
\]
Their respective relative degrees over $K$ are $1,2,2,3$.
If either endpoint has relative degree three, the other has
field exactly $K$.
\end{corollary}
\begin{proof}
Both twists are one by Proposition~\ref{paper:no-thirty-inactive-witnesses}.
The ordinary crossing is $z=1+3\xi$ and has a nontrivial
fifth-torus Fourier channel of squared modulus $45$. The
inactive fifth compression from
Theorem~\ref{paper:inactive-fourier-budget} gives
\[
 45ab\le\frac{D}{5\operatorname{lcm}(36,20)^2}=180.
\]
Thus $ab\le4$.

The exact conductor bound puts $K_X$ inside
$\mathbb Q(\zeta_{2^2 3^2 5^4})$ and $K_Y$ inside
$\mathbb Q(\zeta_{2^2 3^6 5^2})$.
They contain $K$, so any further wild fifth ramification has
degree divisible by five and is impossible. By
Corollary~\ref{paper:compulsory-field-dichotomy}, each field
omits either two or three. The dyadic factor of the remaining
extension has degree at most two and is $\mathbb Q(i)$ if
nontrivial. The ternary factor is cyclic; among its subfields
of degree at most four the only nontrivial possibilities are
$\mathbb Q(\zeta_3)$ and the cubic
$\mathbb Q(\zeta_9)^+$. This gives the list. The final assertion
is immediate from $ab\le4$.
\end{proof}


\begin{lemma}[Two separated trivial Borel channels]
\label{paper:separated-borel-compression}
Let a unitary representation of $G_2\times G_3$ have no
irreducible tensor summand nontrivial at both factors, where
$G_p=\operatorname{SL}_2(\mathbb Z/p^{n_p})$.
Let $E=P_{B_2}P_{B_3}$, with $B_p$ the upper Borel, and let
$K_{p,1}$ be the kernel of reduction modulo $p$. On $E V$,
\[
 C=E W_2W_3 E
 =P_{G_2}P_{G_3}
 -\frac12(P_{K_{2,1}}-P_{G_2})P_{G_3}
 -\frac13P_{G_2}(P_{K_{3,1}}-P_{G_3}).
\]
Consequently $C$ is self-adjoint, its spectrum is contained in
$\{1,-1/2,-1/3,0\}$, and $\|Cv\|=\|v\|$ for $v\in EV$
only when both full groups fix $v$.
\end{lemma}
\begin{proof}
The big-cell proof of the trivial-channel identity
\eqref{paper:exact-trivial-compression} is valid also for $p=2$:
the preimage of the upper Borel modulo $p$ has index $p+1$,
and its complement is the double coset of the Weyl element.
Thus, on $B_p$-fixed vectors,
\[
 P_{B_p}W_p=P_{G_p}-p^{-1}(P_{K_{p,1}}-P_{G_p}).
\]
Decompose into the mutually orthogonal parts where both groups
are trivial, only $G_2$ may be nontrivial, and only $G_3$ may
be nontrivial. The excluded fourth part is zero by hypothesis.
Multiplying the two displayed compression identities gives the
formula. The three displayed projectors are mutually orthogonal,
which proves the spectral and equality assertions.
\end{proof}

\begin{proposition}[The saturated degree-thirty-six channels]
\label{paper:ordinary36-saturated-channels}
In the no-thirty branch, consider a compulsory ordinary crossing
with degrees $(36,20)$ and suppose
\[
 [K_X:\mathbb Q(\zeta_5)]
 [K_Y:\mathbb Q(\zeta_5)]=4.
\]
For every nontrivial character $\chi$ of
$\operatorname{Gal}(\mathbb Q(\zeta_5)/\mathbb Q)$, the normalized
full-orbit Fourier vectors of both endpoints are fixed by the
entire corrected dyadic group, the entire ternary group, and
the congruence kernel $K_{5,1}$.
Their five-inactive Weyl images are supported precisely on the
opposite endpoint orbit, with squared norm $1/5$.
Moreover $X$ belongs to the individual principal fifth block,
so that $X\in P_3\cap P_5$.
\end{proposition}
\begin{proof}
Let $u_\chi$ be the normalized full-orbit Fourier vector of $Y$
from Theorem~\ref{paper:inactive-fourier-budget}. Both endpoint
twists are one. Since the defining character factors through
$\mathbb Q(\zeta_5)$, $u_\chi$ is fixed by the dyadic and
ternary upper Borels and is a nontrivial fifth-torus eigenvector
fixed by the fifth upper unipotents. The no-thirty identity
excludes simultaneous nontrivial dyadic and ternary factors
on this fifth-torus channel.

Let $E=P_{B_2}P_{B_3}$ and use the corrected local Weyls;
their product differs from $s$ by the immaterial sign $\xi$.
The projection of $s u_\chi$ to the $X$-orbit is still fixed
by both complementary tori. Hence its norm is bounded by the
projection to $EV\cap V^{U_5}$, and
\[
 \begin{split}
 \|P_{\operatorname{Orb}(X)}s u_\chi\|^2
 &\le \|P_{U_5}W_5 C u_\chi\|^2\\
 &=\frac15\|P_{K_{5,1}}C u_\chi\|^2
 \le\frac15\|C u_\chi\|^2\le\frac15.
 \end{split}
\]
The equality is exact fifth-channel compression; all selected
operators at distinct primes commute.
Every nontrivial transform of $1+3\xi$ has squared modulus
$45$. The exact orbit-projection identity makes the left side
$4\cdot45/900=1/5$. Equality throughout, followed by
Lemma~\ref{paper:separated-borel-compression}, gives full dyadic
and ternary fixedness, and then $P_{K_{5,1}}u_\chi=u_\chi$.
Equality in the coordinate projection gives the support
assertion. Interchanging $X$ and $Y$ gives the same conclusions
for the opposite Fourier vector.

The $Y$ orbit lies in $P_5$, because the principal fifth block
is Galois stable. The operators $W_5$ and $P_{U_5}$ preserve
this coordinate subspace. Its exact inactive Weyl image is a
nonzero scalar multiple of an $X$-orbit Fourier vector, whose
coefficient at every label of that orbit is nonzero. Hence
the entire $X$ orbit lies in $P_5$, as asserted.
\end{proof}

\begin{proposition}[Fifth invariance at an omitted prime]
\label{paper:ordinary36-omitted-prime-packet}
Assume the no-thirty degree-$(36,20)$ crossing has $ab=4$.
For either endpoint $Z$, its full field is
$K(i)$ or $\mathbb Q(\zeta_{15})$, where
$K=\mathbb Q(\zeta_5)$. Let $p\in\{2,3\}$ be the
prime absent from $K_Z$ and $q$ the other prime. For every
element $g_p$ of the corrected $p$-local group,
\[
                         (g_p-I)\mathbf e_Z
\]
is fixed by every fifth-local torus.
\end{proposition}
\begin{proof}
The field list and $ab=4$ give $a=b=2$. The endpoint orbit
has torus action $C_4\times C_2$, with first factor
fifth-local and second factor $q$-local. Fix a nontrivial
character of $C_4$. Its Fourier projection of $\mathbf e_Z$
splits into the even and odd parts for the relative $C_2$.
The even part is a multiple of the full-orbit Fourier vector,
fixed by the full $p$-local group by
Proposition~\ref{paper:ordinary36-saturated-channels}.
The odd part is orthogonal to both the $q$- and fifth-local
fixed spaces; the no-thirty operator identity makes it
$p$-locally fixed too. Thus every nontrivial fifth-torus
component of $\mathbf e_Z$ is fixed by $g_p$.
Commutation with the fifth tori proves the assertion.
\end{proof}

\begin{lemma}[Rational principal cancellation at defect two]
\label{paper:principal-defect-two-diagonal}
Let $D=n^2$, let $p^m\Vert n$ with $m\ge2$, and suppose
$X\in P_p$ has $v_p(d_X)=m-1$. Assume also that the
$X$ column of $W_p$ is rational, as holds when $K_X$
is unramified at $p$. Then
\[
                         (W_p)_{XX}\text{ is integral at }p.
\]
More generally, if $A\in P_p$ has $v_p(d_A)=m-1$, then
$(W_p)_{AX}$ is integral at $p$.
\end{lemma}
\begin{proof}
Write $n=p^m n'$ and $d_A=p^{m-1}u_A$. The $X$ column
of the complementary Weyl inverse $V=W_{p'}^{-1}$ is
integral at $p$, by the strict foreign bound and
$v_p(f_X)=2$. Principality gives, at every $p$-place,
\[
                 (ps)_{AZ}=\frac{u_A}{n'}\chi_A(Z)
                         \equiv\frac{u_A}{n'}d_Z.
\]
Since $sV=W_p$, multiplying by that integral column yields
\[
 (pW_p)_{AX}
   \equiv\frac{u_A}{n'}\sum_Zd_ZV_{ZX}
   =u_Ap^m(W_p)_{0X}\pmod{\mathfrak p}.
\]
The primary square bound gives
$v_p((W_p)_{0X})\ge-1$, so the right side vanishes.
Thus $v_p((W_p)_{AX})>-1$. Rationality of the column makes
this valuation an integer, proving the assertion.
\end{proof}

\begin{proposition}[The rational principal column in the omitted-three case]
\label{paper:omitted-three-saturated-column}
Assume the no-thirty degree-$(36,20)$ crossing has $ab=4$ and
$K_X=\mathbb Q(\zeta_5,i)$. Let $\mathcal O=I_5X$ be its
four-label fifth-inertia orbit. Then the ternary Weyl column
is rational, $15W_3\mathbf e_X$ is integral, and
\[
          (W_3)_{XX}=\frac45,\qquad
          (W_3)_{AX}=-\frac15
                   \quad(A\in\mathcal O\setminus\{X\}).
\]
Writing $k_A=15(W_3)_{AX}$ outside $\mathcal O$, one has
\[
                       \sum_{A\notin\mathcal O}k_A^2=54.
\]
All its nonzero coordinates lie in $P_3$ and have trivial
second and fifth primary twists. For any such coordinate of
ternary height two, $3\mid k_A$; if $5\mid d_A$, then
$5\mid k_A$ and consequently $k_A=\pm5$. No nonzero outside
coordinate has both ternary height two and degree divisible
by five.
\end{proposition}
\begin{proof}
The omitted-prime invariant-packet result makes
$(W_3-I)\mathbf e_X$ fifth-torus fixed. Unramifiedness of
$K_X$ at three makes the column rational. Its primary and
strict foreign bounds give the integral multiple fifteen.
The principal cancellation calculation above, now in this
rational column, shows that $(W_3)_{XX}$ is integral at three.
Thus the four entries are $1+c,c,c,c$, with $c\in\frac15\mathbb Z$.
The column norm gives $-1/2\le c\le0$; the Weyl-line
criterion excludes $c=0$. Hence $c=-1/5$ or $-2/5$.

The orthogonal complement of the fifth-torus fixed space in
the four-coordinate orbit has dimension three and is fixed
by the full ternary group. Its projection of $\mathbf e_X$
has squared norm $3/4$. Since $X$ is ternary-inactive and
fixed by ternary tori, $\mathbf e_X$ is upper-Borel-fixed.
Writing $a=\|P_{G_3}\mathbf e_X\|^2$ and
$b=\|(P_{K_{3,1}}-P_{G_3})\mathbf e_X\|^2$, exact trivial
Borel compression gives
\[
 (W_3)_{XX}=a-\frac b3
       \ge a-\frac{1-a}3\ge\frac23,
                     \qquad a\ge\frac34.
\]
This excludes $c=-2/5$ and proves the displayed entries.
Their squared norm is $19/25$, leaving $6/25=54/225$.

Local block invariance and commutation with the other primary
twists give the support assertions. The same principal
cancellation, applied to any height-two row in $P_3$, makes
that rational entry integral at three, hence $3\mid k_A$.
If $5\mid d_A$, symmetry and the strict foreign bound give
five-integrality of the entry. Thus $5\mid k_A$; since
$0<k_A^2\le54$, its only possibilities are $\pm5$.
Divisibility by both three and five would be impossible.
\end{proof}

\begin{lemma}[Coefficients in a rational primary word]
\label{paper:rational-primary-word-coefficients}
Let $p$ be prime and let a nonnegative integral word on
$p$-primary roots have rational value $r\in\mathbb Z_{\ge0}$
and mass $M$. Every nonidentity phase has coefficient at most
$(M-r)/p$, and the coefficient at one is at most
$r+(M-r)/p$.
\end{lemma}
\begin{proof}
Place all roots in $\mu_{p^n}$. The relation lattice consists
of the constant vectors on its $\mu_p$ cosets. In the coset
containing one, the coefficient at one is $r+t$ and all
other coefficients are $t\ge0$. Every other coset has one
common nonnegative coefficient. The sum of all these common
coefficients is $(M-r)/p$, proving both bounds.
\end{proof}

\begin{theorem}[No Gaussian degree-thirty-six endpoint at saturation]
\label{paper:no-gaussian-saturated-endpoint}
In the no-thirty branch, a compulsory degree-$(36,20)$
crossing with $ab=4$ cannot have
$K_X=\mathbb Q(\zeta_5,i)$.
\end{theorem}
\begin{proof}
Suppose $K_X=\mathbb Q(\zeta_5,i)$ and put $W=W_3$.
Proposition~\ref{paper:omitted-three-saturated-column} gives
the entries $4/5,-1/5,-1/5,-1/5$ on
$\mathcal O=I_5X$, and integral outside numerators
$k_A=15W_{AX}$ satisfying
\[
                      \sum_{A\notin\mathcal O}k_A^2=54.
\]
The source and its fifth orbit have complete twist one.
The diagonal relation $WT_3W=T_3^{-1}WT_3^{-1}$ therefore
gives the additional identity
\[
             \sum_{A\notin\mathcal O}k_A^2\theta_{A,3}=9.
\]
Lemma~\ref{paper:rational-primary-word-coefficients}, with
$(p,r,M)=(3,9,54)$, bounds every coefficient by twenty-four
(and every nonidentity coefficient by fifteen). No individual
outside numerator can consequently have absolute value five
or greater. If $5\mid d_A$, the foreign-column bound and
symmetry would force $5\mid k_A$, so such a coordinate is
zero. The four source degrees are thirty-six as well. Thus
the whole column has five-free support.

Every nonzero coordinate lies in $P_3$ and has trivial second
and fifth twists. Apart from the unit, the possible degrees
are
\[
                              12,24,36,72.
\]
Indeed, principal membership excludes ternary height three,
prime-power degrees are absent, and the complementary-degree
bound of Lemma~\ref{paper:family:4.10}, applied with the compulsory
degree-twenty member of $P_5$, excludes dyadic height one.
Explicitly, a degree-eighteen row in $P_3$ paired with $Y$
would have common weight $\gcd(18,20)=2$, whereas the two
principal residues require this weight to be at least three.
The same argument excludes degree six independently of its
previous global exclusion.
The degree-twelve and degree-thirty-six coordinates are
three-inactive by Proposition~\ref{paper:no-thirty-inactive-witnesses}.
The degree-seventy-two coordinates are three-inactive by
Proposition~\ref{paper:tame-largest-ternary-inactive}: their
crossing with $Y$ is ordinary, since a signed crossing would
force thirty activity. Consequently every three-active
coordinate in this column has degree twenty-four.

Such a coordinate has two-unit codegree, so every dyadic torus
fixes it. Let $P_{\mathrm{act}}$ project onto the three-active
coordinates and put $v=P_{\mathrm{act}}W\mathbf e_X$.
The omitted-prime packet theorem makes
$(W-I)\mathbf e_X$ fifth-torus fixed; hence all four columns
on $\mathcal O$ have the same active projection $v$.
The dyadic relative involution on $K_X$ transports
$\mathcal O$ to a disjoint four-label orbit. It fixes every
coordinate in the support of $v$, since those coordinates
have degree twenty-four. Commutation with $W$ therefore
makes the active projections of these four columns equal to
$v$ too. The full $X$ orbit has eight labels, and unitarity
gives
\[
                              8\|v\|^2\le1.
\]

On the other hand, $\mathbf e_X$ is fixed by the ternary
upper Borel. Put
\[
 \alpha=\|P_{G_3}\mathbf e_X\|^2,\qquad
 \beta=\|(P_{K_{3,1}}-P_{G_3})\mathbf e_X\|^2.
\]
Exact trivial Borel compression gives
\[
 \alpha-\frac\beta3=W_{XX}=\frac45,\qquad
 \|P_{U_3}W\mathbf e_X\|^2=\alpha+\frac\beta9,
                       \qquad\alpha+\beta\le1.
\]
Thus $\beta\le3/20$ and the inactive squared norm is at most
$13/15$. It follows that
\[
                         \|v\|^2\ge\frac2{15}>\frac18,
\]
contradicting the eight-column bound. The field list of
Corollary~\ref{paper:ordinary36-field-list} leaves only
$\mathbb Q(\zeta_{15})$ for $K_X$ at saturation.
\end{proof}

\begin{theorem}[No saturated compulsory degree-thirty-six crossing]
\label{paper:no-saturated-compulsory-crossing}
In the no-thirty branch, let a compulsory crossing have
degrees $(36,20)$, and put
\[
 K=\mathbb Q(\zeta_5),\qquad
 a=[K_X:K],\qquad b=[K_Y:K].
\]
Then $ab\le3$. More precisely, the only remaining pairs are
\[
                   (a,b)\in
                   \{(1,1),(1,2),(2,1),(1,3),(3,1)\}.
\]
In particular at least one endpoint has full character field
exactly $K$.
\end{theorem}
\begin{proof}
The four-field list gives $ab\le4$. Suppose $ab=4$.
Then $a=b=2$, and each endpoint field is either $K(i)$ or
$\mathbb Q(\zeta_{15})$.

First suppose $K_X=\mathbb Q(\zeta_{15})$. Write
$W=\gamma^3W_2$ for the corrected dyadic Weyl operator,
and let $B$ be the full dyadic block containing $X$.
The omitted-prime conclusion of
Proposition~\ref{paper:ordinary36-omitted-prime-packet} gives
\[
                          w=(W-I)\mathbf e_X
                 \quad\hbox{fixed by all fifth tori}.
\]
Local block invariance puts the support of $w$ in $B$.
Lemma~\ref{paper:compulsory-fifth-block-separation} supplies a fifth torus $h$ with
$hB\ne B$. Full blocks are disjoint, whereas $hw=w$;
hence $w$ is supported in $B\cap hB=\varnothing$ and is
zero. Thus $W\mathbf e_X=\mathbf e_X$, which makes the
uncorrected Weyl operator preserve the same coordinate line.
The Weyl-line criterion contradicts $2\mid f_X$.

The only remaining saturated possibility has $K_X=K(i)$.
It is excluded by
Theorem~\ref{paper:no-gaussian-saturated-endpoint}.
Therefore $ab=4$ is impossible. Combining $ab\le3$ with
the relative degrees $1,2,3$ in the four-field list
gives exactly the displayed pairs.
\end{proof}

\begin{corollary}[Ternary stability before saturation]
\label{paper:compulsory-ternary-block-stability}
In the no-thirty branch, under
Lemma~\ref{paper:compulsory-fifth-block-separation}, let $Z$ be either compulsory endpoint and let $B_Z$ be its full
dyadic block. Every ternary torus preserves $B_Z$, and
\[
             (W-I)\mathbf e_A
       \text{ is fixed by every ternary torus}
                     \qquad(A\in B_Z),
\]
where $W=\gamma^3W_2$. Every dyadically active label in
$B_Z$ is fixed by all ternary tori, and hence has full
character field unramified at three.
\end{corollary}
\begin{proof}
Some fifth torus moves $B_Z$, by the lemma. The no-thirty
block-stability proposition therefore makes every ternary
torus preserve $B_Z$. Since $2\mid f_Z$, this block is not
a singleton of dyadic-unit codegree; no label $A\in B_Z$
has dyadic-unit codegree, since such a label has a singleton
full block. Apply the same proposition at each $A\in B_Z$,
with the roles of three and five interchanged. This gives
the displayed fixedness.

For completeness, let $V$ be the orthogonal complement of
the ternary-torus fixed space in the coordinate space of
$B_Z$. The displayed identity makes $W$ act as the identity
on $V$. Both $W$ and the corrected dyadic upper unipotent
$U_2=T_2$ preserve $V$, since they commute with the ternary
tori. The local relation
$U_2(WU_2W^{-1})U_2=W^{-1}$ then gives $U_2^3=I$ on
$V$. As $U_2$ has two-power order, it is the identity on
$V$. Hence $V$ is supported on dyadically inactive labels.
Any nontrivial ternary orbit on a dyadically active label
would supply a nonzero vector in $V$ with that same active
dyadic phase, which is impossible. Thus every such label
is ternary-torus fixed, proving the final assertion.
\end{proof}


\begin{lemma}[A quintic transport lattice on an active dyadic fiber]
\label{paper:quintic-transport-phase-lattice}
\label{paper:degree-twenty-transport-scope}
Let $Y$ be any compulsory degree-twenty witness in the no-thirty
branch. Suppose $K_Y$ is unramified at two.
The opposite compulsory degree may be twelve or thirty-six.
For the corrected dyadic
Weyl $W=\gamma^3W_2$, the column $18W\mathbf e_Y$ is
integral. For a nonidentity dyadic phase $\rho$ occurring
in this column put
\[
 v_\rho=P_\rho W\mathbf e_Y,\qquad
 M_\rho=324\|v_\rho\|^2\in\mathbb Z_{>0}.
\]
There is a positive integer $n_\rho$ such that
\begin{equation}\label{paper:quintic-phase-lattice-interval}
              \frac{4M_\rho}{81}\le n_\rho
                         \le\frac{5M_\rho}{81}.
\end{equation}
\end{lemma}
\begin{proof}
The full field contains $K=\mathbb Q(\zeta_5)$.
The conductor bound puts it in
$\mathbb Q(\zeta_{25})\mathbb Q(\zeta_{3^b})$ for some $b\ge0$.
Local Galois covariance makes the dyadic column rational.
At its source,
\[
             f_Y=2^2 3^6 5^2.
\]
The primary square bound at two and strict foreign bounds
at three and five therefore give $18W\mathbf e_Y\in\mathbb Z^I$.

Let $H=H_{-1,5}$. The fifth-block separation lemma makes
$H$ move the full dyadic block of $Y$ to a disjoint block.
Thus, for
\[
                         u_\rho=v_\rho-Hv_\rho,
\]
one has
\[
 Hu_\rho=-u_\rho,
 \qquad \|u_\rho\|^2=\frac{M_\rho}{162}.
\]
The source $Y$ and its fifth translates have twist one.
Commutation of the dyadic Weyl and phase projector with
the fifth upper group consequently gives $U_5u_\rho=u_\rho$.
The central fifth involution acts by minus one on $u_\rho$.
Consequently all of its torus-character components are nontrivial.
For a component of conductor $5^c$, exact local compression
contracts its squared norm by at most $5^{-c}\le1/5$.
Summing over the orthogonal components gives
\begin{equation}\label{paper:quintic-transport-compression}
 \|P_{5,\mathrm{inact}}W_5u_\rho\|^2
                     \le\frac15\|u_\rho\|^2.
\end{equation}
This argument allows wild fifth inertia at the source.
In particular $P_{5,\mathrm{act}}W_5u_\rho$ is nonzero and has squared
norm at least $4\|u_\rho\|^2/5$.

Consider a coordinate $A$ active at both two and five, with
dyadic phase $\rho$. No-thirty separation gives three-unit
codegree, hence
\[
               d_A=27\,2^a5^c,
                       \qquad 0\le a\le2,\quad0\le c\le1.
\]
Its full ternary local line is trivial. Since
$s=W_2W_3W_5$ and $W_5$ preserves the exact dyadic phase,
\[
                \chi_Y(A)=\pm270(W_5v_\rho)_A.
\]
Moreover $\operatorname{lcm}(20,d_A)=540$, so LCM
integrality gives $\chi_Y(A)/27\in\mathcal O$.
The entry belongs to $K_Y\cap K_A$, and $K_A$ is
unramified at three. The displayed containment for $K_Y$ thus gives
\[
          z_A:=10(W_5v_\rho)_A\in\mathbb Z[\zeta_{25}].
\]
All divisibility statements here hold in the full ring of
algebraic integers; no reduction at only one chosen place
is being used.

The vector $v_\rho$ is rational. Local Galois covariance gives
\[
       10(W_5u_\rho)_A=z_A-\overline{z_A}=:\alpha_A.
\]
A nonzero anti-real element of $\mathbb Q(\zeta_{25})$ has
actual field degree four or twenty: these are its only nonreal
subfields. In degree four it lies in $K$, and the anti-real
integer lattice has basis
$\zeta_5-\zeta_5^{-1},\zeta_5^2-\zeta_5^{-2}$. Thus
\begin{equation}\label{paper:quintic-anti-real-trace}
 \operatorname{Tr}_{K/\mathbb Q}(\alpha_A\overline{\alpha_A})
                              =10(r_A^2+s_A^2)
\end{equation}
for integers $r_A,s_A$. In degree twenty, put
$L=\mathbb Q(\zeta_{25})$. Every integer of $L$ has trace
divisible by five, as follows from the power-basis traces
$20,-5,0$. Since $\alpha_A=z_A-\overline{z_A}$,
\[
 \operatorname{Tr}_{L/\mathbb Q}(\alpha_A\overline{\alpha_A})
       =2\operatorname{Tr}_{L/\mathbb Q}(z_A\overline{z_A}-z_A^2)
                                                    \in10\mathbb Z.
\]
This trace is positive when $\alpha_A\ne0$.

Fifth Galois covariance identifies the entries on each target
orbit with all conjugates of $\alpha_A$, with equal repetitions
if its field is smaller than that orbit. Exact dyadic phase and
fifth activity are preserved. Summing the two trace calculations
over complete target orbits gives
\[
 q_\rho:=100\|P_{5,\mathrm{act}}W_5u_\rho\|^2
                                  \in10\mathbb Z_{>0}.
\]
Compression and unitarity give
\[
             \frac{40M_\rho}{81}\le q_\rho
                             \le\frac{50M_\rho}{81}.
\]
Put $n_\rho=q_\rho/10$ to obtain the assertion.
\end{proof}

\begin{corollary}[An integral interval for the dyadic diagonal]
\label{paper:degree-twenty-diagonal-interval}
Under the hypotheses of Lemma~\ref{paper:quintic-transport-phase-lattice}, write
\[
                         W_{YY}=1-\frac l{18}.
\]
Then $l$ is a positive odd integer, and some positive integer
$N$ satisfies
\begin{equation}\label{paper:degree-twenty-l-interval}
                         \frac{4l}{9}\le N\le\frac{5l}{9}.
\end{equation}
In particular $l\notin\{1,3,5,7\}$.
\end{corollary}
\begin{proof}
Homogeneity of the compulsory dyadic block makes the nonzero
diagonal entry have exact two-valuation minus one. Its
integral multiplier eighteen thus makes $l$ odd. Unitarity
gives $l\ge0$, and equality would make the column monomial,
contrary to positive dyadic defect. Hence $l>0$.

Put $k_A=18W_{AY}$ for $A\ne Y$. The norm and the diagonal
corrected modular relation give the positive dyadic word
\[
 \sum_{A\ne Y}k_A^2=36l-l^2,
 \qquad
 \sum_{A\ne Y}k_A^2\theta_{A,2}=18l-l^2.
\]
In any positive integral word on two-primary roots with
rational value, the coefficients at $\eta$ and $-\eta$
are equal for every nonreal pair. If $t$ is the coefficient
at minus one and $m_j$ are the common coefficients at the
nonreal pairs, subtraction of the value from the mass gives
\[
                             t+\sum_jm_j=9l.
\]
Terms of coefficient zero are omitted. For minus one use the
integer supplied by the phase-lattice lemma when $t>0$;
for each nonreal pair choose either representative and use
its integer. The sum $N$ of these positive integers satisfies
\eqref{paper:degree-twenty-l-interval}. The four excluded
values of $l$ give intervals containing no integer.
\end{proof}


\begin{lemma}[A coarse separated Borel bound and its equality case]
\label{paper:separated-borel-coarse-bound}
On a nontrivial fifth-torus channel in the no-thirty setting,
let $u$ be a unit vector fixed by the dyadic and ternary
upper Borels. Write
\[
 E=P_{B_2}P_{B_3},\qquad C=EW W_3E,
                   \qquad\lambda=\langle u,Wu\rangle,
\]
where $W$ is the corrected dyadic Weyl. Then
\[
                       \|Cu\|^2\le\frac{1+\lambda}{2}.
\]
If equality holds, $u$ is fixed by the whole ternary group
and by $K_{2,1}$. On the cyclic span relevant to $u$, $C$
has only the eigenvalues $1,-1/2$. If also $\lambda=1/2$,
the squared norms of these two components are $2/3,1/3$,
respectively.
\end{lemma}
\begin{proof}
Set
\[
 \alpha=\|P_{G_2}u\|^2,
 \qquad\beta=\|(P_{K_{2,1}}-P_{G_2})u\|^2.
\]
Trivial Borel compression gives
$\lambda=\alpha-\beta/2$, with $\alpha+\beta\le1$.
The separated compression formula gives
\[
                \|Cu\|^2\le\alpha+\frac\beta4
                     =\lambda+\frac{3\beta}{4}
                     \le\frac{1+\lambda}{2}.
\]
Equality in the last inequality gives
$\alpha+\beta=1$, hence $K_{2,1}$-fixedness. Equality
in the first means that the $G_2$-fixed component has no
nontrivial ternary part. The $G_2$-moving component already
has trivial ternary factor by the no-thirty operator identity
on a nontrivial fifth channel. Thus $G_3$ fixes $u$.
The stated eigenvalues and, when $\lambda=1/2$, their weights
follow by solving $\alpha-\beta/2=1/2$, $\alpha+\beta=1$.
\end{proof}

\begin{theorem}[No relative cubic endpoint at a degree-thirty-six crossing]
\label{paper:no-cubic-relative-endpoint}
In the no-thirty compulsory degree-$(36,20)$ crossing, put
\[
 K=\mathbb Q(\zeta_5),\qquad
 a=[K_X:K],\qquad b=[K_Y:K].
\]
Then $ab\ne3$. Combined with
Theorem~\ref{paper:no-saturated-compulsory-crossing}, this gives
\[
                  ab\le2,\qquad
                  (a,b)\in\{(1,1),(1,2),(2,1)\}.
\]
In particular the relative cubic field
$K\mathbb Q(\zeta_9)^+$ occurs at neither endpoint.
\end{theorem}
\begin{proof}
Suppose $ab=3$. The field list makes $(a,b)=(3,1)$ or
$(1,3)$; both fields are unramified at two, and their fifth
inertia quotients have order four.

Fix a nontrivial fifth-torus character $\chi$, and let
$u_Y$ be the unit full-orbit Fourier vector at $Y$ with
that character and trivial relative ternary character.
Define $u_X$ with the opposite fifth character, so that
the exact orbit-projection formula gives, for a phase
$\epsilon$,
\begin{equation}\label{paper:cubic-orbit-crossing}
 P_{\operatorname{Orb}(X)}s u_Y
          =\epsilon\sqrt{\frac3{20}}\,u_X.
\end{equation}
Indeed its squared norm is $ab\cdot45/900=3/20$.
Both vectors are fixed by the dyadic and ternary upper
Borels. Put $E=P_{B_2}P_{B_3}$, $C=EW W_3E$, and
$A_5=P_{U_5}W_5$. As in the saturated-channel calculation,
all complementary projections commute, and
\[
 \frac3{20}
 \le\|A_5Cu_Y\|^2
 =\frac15\|P_{K_{5,1}}Cu_Y\|^2
 \le\frac15\|Cu_Y\|^2.
\]
Thus $\|Cu_Y\|^2\ge3/4$.

Write $W_{YY}=1-l/18$ as in
Corollary~\ref{paper:degree-twenty-diagonal-interval}.
The nonsaturated block-stability corollary makes
$(W-I)\mathbf e_Y$ fixed by ternary tori. On the
$b$-label ternary orbit of $Y$, the restricted matrix is
therefore
\[
                         I_b-\frac l{18}J_b.
\]
The four distinct fifth bins lie in distinct full dyadic
blocks. There are no Weyl entries between them; their
within-bin matrices are conjugate by fifth tori. Consequently
the full-orbit Fourier vector has
\begin{equation}\label{paper:cubic-y-dyadic-expectation}
                \lambda_Y=\langle u_Y,Wu_Y\rangle
                             =1-\frac{bl}{18}.
\end{equation}
The coarse Borel bound now gives
\[
              \frac34\le\|Cu_Y\|^2
                          \le1-\frac{bl}{36},
                          \qquad bl\le9.
\]
If $b=3$, the positive odd $l$ is one or three, both excluded
by the integral interval for $l$. Hence $b=1$ and $a=3$.
The same interval excludes $l=1,3,5,7$, leaving $l=9$.
Thus $\lambda_Y=1/2$, and equality holds throughout the
preceding chain of inequalities.

The equality case of the coarse Borel bound gives
$G_3$-fixedness of $u_Y$ and its decomposition into the
$C$-eigenvalues $1,-1/2$ with squared norms $2/3,1/3$.
In particular
\begin{equation}\label{paper:cubic-third-compression-moment}
                    \langle u_Y,C^3u_Y\rangle
                      =\frac23-\frac1{24}=\frac58.
\end{equation}
Equality in fifth compression gives
$P_{K_{5,1}}Cu_Y=Cu_Y$. Equality in the orbit projection
gives the exact identity, after changing the phase if needed,
\[
                  A_5Cu_Y
                    =\epsilon\sqrt{\frac3{20}}\,u_X.
\]
The operators $A_5$ and $C$ commute. Polarization of exact
nontrivial fifth-channel compression gives
$A_5^*A_5=P_{K_{5,1}}/5$ on its upper-unipotent fixed
torus channel. Since $Cu_Y$ and $C^2u_Y$ are fixed by
$K_{5,1}$, it follows that
\[
 \begin{split}
 \langle u_X,Cu_X\rangle
  &=\frac{20}{3}
       \langle A_5Cu_Y,A_5C^2u_Y\rangle\\
  &=\frac43\langle Cu_Y,C^2u_Y\rangle
    =\frac43\cdot\frac58=\frac56.
 \end{split}
\]
The exact image identity also makes $u_X$ fixed by $G_3$.
As $u_X$ is Borel-fixed at both complementary primes,
\[
                       \langle u_X,Cu_X\rangle
                          =\langle u_X,Wu_X\rangle.
\]

Finally, $K_X$ is unramified at two and
$f_X=2^2 3^2 5^4$. The primary and strict foreign bounds
give $10W\mathbf e_X\in\mathbb Z^I$.
Put $c_X=W_{XX}-1\in\frac1{10}\mathbb Z$.
The ternary orbit has three labels and the nonsaturated
fixedness makes its within-orbit matrix $I_3+c_XJ_3$.
Its four fifth translates again occupy distinct dyadic
blocks, so the same full-orbit calculation as above gives
\[
                   \langle u_X,Wu_X\rangle
                      =1+3c_X\in\frac1{10}\mathbb Z.
\]
This contradicts its value $5/6$. Both product-three
possibilities are excluded, proving the theorem.
\end{proof}

\subsection{Exact quintic endpoint fields and the degree-twelve branch}


\begin{lemma}[A ternary-fixed fifth orbit has no mixed dyadic--ternary component]
\label{paper:thirtysix-no-mixed-two-three}
Assume the no-thirty branch. Fix a compulsory degree-twenty
witness $Y$, and let $X\in P_3$ have degree thirty-six and
pure ternary twist. Suppose $K_X=K:=\mathbb Q(\zeta_5)$
and every nontrivial fifth Fourier vector of its four-label
orbit is fixed by $G_3$. Then
\[
 (I-P_{G_2})(I-P_{G_3})\mathbf e_X=0,
 \qquad
 S_{XZ}=0\quad(\theta_{Z,2}\ne1,\ \theta_{Z,3}\ne1).
\]
\end{lemma}
\begin{proof}
The inactive-witness theorem gives $\theta_X=\theta_Y=1$.
Let $\mathcal O=I_5X$. The Fourier hypothesis makes
$(W_3-I)\mathbf e_X$ fifth-torus fixed. Its column is rational,
with $15W_3\mathbf e_X$ integral. Rational principal defect-two
cancellation gives the diagonal three-integrality. Thus its
four orbit entries are $1+c,c,c,c$, with $c\in\frac15\mathbb Z$.
The norm and Weyl-line criterion give $c=-1/5$ or $-2/5$.
The nontrivial Fourier subspace supplies a $G_3$-fixed projection
of $\mathbf e_X$ of squared norm $3/4$. Trivial Borel compression
therefore gives $(W_3)_{XX}\ge2/3$, excluding $c=-2/5$.
The orbit entries are $4/5,-1/5,-1/5,-1/5$.

The outside numerators $k_A=15(W_3)_{AX}$ satisfy
\[
 \sum_{A\notin\mathcal O}k_A^2=54,
 \qquad \sum_{A\notin\mathcal O}k_A^2\theta_{A,3}=9.
\]
The rational-primary-word bound makes every coefficient at most
24, hence $|k_A|<5$. If $5\mid d_A$, the strict foreign bound
and symmetry give $5\mid k_A$. Thus the whole column is five-free.

Every supported coordinate lies in $P_3$ and has trivial second
and fifth twists. A five-free nonunit in $P_3$ has degree
$2^i3^j$, with $i,j\ge1$. The principal height bound gives
$j\le2$. Complementary principal degrees with the fixed $Y$
give $2^i\ge3$, hence $i\ge2$; divisibility gives $i\le3$.
The complete palette is $12,24,36,72$. Degrees twelve and
thirty-six are ternary inactive by
Proposition~\ref{paper:no-thirty-inactive-witnesses}, applied
with the same $Y$. Degree seventy-two is inactive by
Proposition~\ref{paper:tame-largest-ternary-inactive}: its
mass-four crossing with $Y$ is ordinary, since a signed
crossing would force thirty activity. Hence every ternary-active
coordinate has degree twenty-four and a trivial full dyadic line.

Put $w=(I-P_{G_2})\mathbf e_X$. It is fixed by $U_3$, and
commutation of the local groups gives
\[
 P_{3,\mathrm{act}}W_3w
   =(I-P_{G_2})P_{3,\mathrm{act}}W_3\mathbf e_X=0.
\]
Therefore $U_3W_3w=W_3w$. Upper and lower elementary
ternary groups fix $w$, so $G_3$ fixes it. This proves the
projector identity. The global $s$ preserves both local fixed
spaces, whereas a coordinate active at both primes is orthogonal
to each. The asserted zeros follow.
\end{proof}

\begin{lemma}[A tame degree-twelve witness from the mixed-support obstruction]
\label{paper:mixed-support-forces-tame-twelve}
Fix a compulsory degree-twenty witness $Y$ in the no-thirty
branch. Suppose every degree-thirty-six label in
$P_3\cap\{\theta_2=\theta_5=1\}$ has zero entry at every
jointly dyadic--ternary active label. This includes the case
that no such degree-thirty-six label exists. Then a compulsory
$T\in P_3$ can be chosen with $d_T=12$, $\theta_T=1$, and
$|I_5T|=4$. For some jointly dyadic--ternary active $Z$,
$\chi_Z(T)$ is nonzero modulo a chosen place above five.
\end{lemma}
\begin{proof}
Choose $Z$ from edge six of the complete activity graph and put
$E=P_3\cap\{A:\theta_{A,2}=\theta_{A,5}=1\}$.
Principal-fiber orthogonality gives
\[
                         \sum_{A\in E}d_A\chi_Z(A)=0.
\]
Modulo five, the five-divisible degrees vanish and the vacuum
contributes one. The five-free nonunit palette is $12,24,36,72$,
by the principal-height and complementary-degree argument above.
Degrees twenty-four and seventy-two have dyadic-unit codegree
and miss $Z$. Degree thirty-six misses $Z$ by hypothesis.
Thus degree twelve makes a nonzero contribution. Every such
label in $E$ has twist one by the inactive-witness theorem.

Moreover $Z$ has five-unit codegree, so fifth inertia fixes it.
The set $E$ is stable under fifth inertia, and $\chi_Z(A)$ is
constant on each such orbit. The ordinary crossing with $Y$
makes the full field of every selected degree-twelve label
contain $K=\mathbb Q(\zeta_5)$, so its fifth orbit has size
$4\cdot5^r$. Orbits with $r>0$ contribute zero modulo five.
An orbit of size four must therefore contribute nontrivially;
choose $T$ in that orbit.
\end{proof}

\begin{lemma}[Exact equality support excludes an index-two quintic endpoint]
\label{paper:index-two-exact-support-obstruction}
Assume a no-thirty compulsory degree-$(36,20)$ crossing with
$[K_Y:K]=2$ and $K_X=K=\mathbb Q(\zeta_5)$. For a nontrivial
fifth character let $u_Y,u_X$ be the normalized full-orbit Fourier
vectors, with trivial complementary characters. Put
\[
 E=P_{B_2}P_{B_3},\quad C=EWW_3E,\quad
 A_5=P_{U_5}W_5,\quad Q=P_{B_2}P_{B_3}P_{U_5}.
\]
Equality throughout
\[
 \frac12\le\|Cu_Y\|^2
       \le\frac{1+\langle u_Y,Wu_Y\rangle}{2}
\]
with common value $1/2$ is impossible.
\end{lemma}
\begin{proof}
The four fifth translates of the relative pair of $Y$ lie in
distinct full dyadic blocks. Commutation with fifth tori makes
their within-pair Weyl matrices conjugate. Thus
$\langle u_Y,Wu_Y\rangle$ is independent of the chosen
nontrivial fifth character. Every such crossing has squared
orbit norm $1/10$, so the assumed equality propagates to all
three nontrivial fifth characters.
The coarse equality case gives $G_3$-fixedness of $u_Y$.
The orbit crossing has squared norm $1/10$; equality in fifth
compression and the orbit projection therefore gives
\[
                   A_5Cu_Y=\varepsilon\sqrt{1/10}\,u_X.
\]
Consequently $u_X$ is $G_3$-fixed. The same reasoning applies
to every nontrivial fifth character. Furthermore local commutation,
$Eu_Y=u_Y$, and the corrected local factorization of $s$ give
\[
                         Qsu_Y=\varepsilon' A_5Cu_Y.
\]
Thus this genuine orthogonal projection is supported exactly
on the full degree-thirty-six orbit of $X$.

Any other compulsory degree-thirty-six orbit would have a
nonzero ordinary fifth Fourier coefficient against this same
$u_Y$. Its matching orbit vector is in $\operatorname{im}Q$,
since Theorem~\ref{paper:no-cubic-relative-endpoint} makes its full field $K$ and its twist
is one. Exact support rules out such another orbit. Hence every
compulsory degree-thirty-six label is a conjugate of $X$ and
satisfies the preceding mixed-support lemma. The tame-witness
lemma now forces a compulsory degree-twelve label $T$.

Its matching full-orbit Fourier vector $u_T$ is fixed by both
complementary Borels and by the fifth upper group, so
$Qu_T=u_T$. It is orthogonal to the degree-thirty-six orbit.
Exact support gives $\langle u_T,su_Y\rangle=0$. But every
nontrivial fifth transform of $1+3\xi_5$ has squared modulus
45. The exact orbit formula instead gives
\[
 |\langle u_T,su_Y\rangle|^2
 =\frac{[K_T:K]\cdot2\cdot45}{D/60^2}
 =\frac{[K_T:K]}{90}>0.
\]
This is the contradiction.
\end{proof}

\begin{theorem}[No ternary degree-twenty field opposite degree thirty-six]
\label{paper:no-ternary-quintic-endpoint-at-thirtysix}
In the no-thirty branch, a compulsory degree-$(36,20)$
crossing cannot have $K_Y=\mathbb Q(\zeta_{15})$.
\end{theorem}
\begin{proof}
Theorem~\ref{paper:no-cubic-relative-endpoint} makes $K_X=K$.
Write $W_{YY}=1-l/18$.
The rational degree-twenty transport interval makes $l$ positive
odd and excludes $1,3,5,7$. Uniform ternary block stability
gives the within-orbit matrix $I_2-(l/18)J_2$ on the relative
ternary pair of $Y$. Its four fifth translates lie in distinct
full dyadic blocks. Thus $\langle u_Y,Wu_Y\rangle=1-l/9$,
and the Fourier crossing and coarse Borel bound give
\[
                        \frac12\le\|Cu_Y\|^2\le1-l/18.
\]
Therefore $l=9$ and all estimates are equalities. The exact-support
obstruction excludes this case.
\end{proof}

\begin{lemma}[Gaussian column parity at a compulsory degree-twenty endpoint]
\label{paper:gaussian-quintic-column-parity}
Assume the no-thirty branch, and let $Y$ be a compulsory
degree-twenty witness with $K_Y=\mathbb Q(\zeta_{5^b},i)$,
where $b=1$ or $2$. Let
$W=\gamma^3W_2$ be the corrected dyadic Weyl. Write
\[
             18W\mathbf e_Y=(m_A+in_A)_A,
                         \qquad m_A,n_A\in\mathbb Z.
\]
Then
\[
 \sum_A m_A^2=\sum_A n_A^2=162,\qquad
 \sum_A m_A^2\theta_{A,2}=18m_Y,\qquad
 \sum_A n_A^2\theta_{A,2}=-18in_Y.
\]
Every exact dyadic-phase sum of the $n_A^2$ is even.
Consequently $n_Y$ is odd and $m_Y$ is even.
\end{lemma}
\begin{proof}
Dyadic Galois covariance puts the column in $\mathbb Q(i)$.
The primary square bound and strict foreign bounds for
$f_Y=2^2 3^6 5^2$ give the stated integral multiplier eighteen.
Let $H=H_{-1,2}$ and $Y'=HY$. The relative Gaussian
involution is nontrivial on the full field, so $Y'\ne Y$.
Moreover $W^2=H$ and $\overline W=W^{-1}$, whence
\[
                  W\mathbf e_{Y'}=\overline{W\mathbf e_Y}.
\]
The normalized vectors
$u_+=(\mathbf e_Y+\mathbf e_{Y'})/\sqrt2$ and
$u_-=(\mathbf e_Y-\mathbf e_{Y'})/\sqrt2$ have Weyl images
\[
              Wu_+=\frac{\sqrt2}{18}(m_A)_A,
              \qquad Wu_-=\frac{\sqrt2 i}{18}(n_A)_A.
\]
Unitarity proves the two norm identities.

Both source labels have twist one. On the $H=1$ space,
$W^*=W$, while on the $H=-1$ space $W^*=-W$.
Using $WT_2W=T_2^{-1}WT_2^{-1}$ gives
\[
 \langle Wu_+,T_2Wu_+\rangle=\langle u_+,Wu_+\rangle
                                        =m_Y/9,
\]
\[
 \langle Wu_-,T_2Wu_-\rangle=-\langle u_-,Wu_-\rangle
                                        =-in_Y/9.
\]
Substituting the two Weyl images proves the moment identities.

The torus $H$ is an unsigned permutation, sends the vector
$(n_A)_A$ to its negative, and preserves each exact dyadic
phase. Its fixed coordinates have $n_A=0$, and every other
orbit contributes $2n_A^2$ to that phase's mass. Thus each
phase mass is even. Divide the last moment identity by two.
It is a positive integral word on dyadic roots of mass
eighty-one and value $-9in_Y$. Reducing at a place above
two, every dyadic root reduces to one and $i$ also reduces
to one. Hence $81\equiv n_Y\pmod2$, proving oddness.

The compulsory dyadic block is homogeneous of defect two.
Its diagonal Weyl entry has exact two-valuation minus one.
Therefore $m_Y+in_Y=18W_{YY}$ is a unit at $1+i$;
equivalently $m_Y+n_Y$ is odd. The asserted evenness of
$m_Y$ follows.
\end{proof}

\begin{lemma}[Quintic transport of the real Gaussian column]
\label{paper:gaussian-real-column-transport}
Under the preceding hypotheses, set
$v=\operatorname{Re}(W\mathbf e_Y)=(m_A/18)_A$.
For each nonidentity dyadic phase $\rho$ with
$v_\rho=P_\rho v\ne0$, put
$M_\rho=324\|v_\rho\|^2\in\mathbb Z_{>0}$.
There is an integer $n_\rho>0$ with
\[
                   \frac{4M_\rho}{81}\le n_\rho
                                    \le\frac{5M_\rho}{81}.
\]
If $l=9-m_Y>0$, then some positive integer
$N$ satisfies
\[
                           \frac{4l}{9}\le N\le\frac{5l}{9}.
\]
In particular $l\notin\{1,3,5,7\}$.
\end{lemma}
\begin{proof}
Let $H_5=H_{-1,5}$ and put $w_\rho=v_\rho-H_5v_\rho$.
The fifth torus moves the full dyadic block of $Y$ to a
disjoint block. Thus
\[
                         \|w_\rho\|^2=M_\rho/162.
\]
The vector is rational, fifth-upper-fixed, fixed by the
tori congruent to one modulo $5^b$, and odd under $H_5$.
Its torus characters are nontrivial, of conductor $5^u$
with $1\le u\le b$. The exact compression factor is
$5^{-u}\le1/5$, so fifth compression gives
\[
 \|P_{5,\mathrm{act}}W_5w_\rho\|^2
                  \ge\frac45\|w_\rho\|^2.
\]

For a target $A$ active at two and five, no-thirty separation
gives
\[
              d_A=27\,2^a5^c,
                    \qquad 0\le a\le2,\quad0\le c\le1.
\]
Its full ternary local line is trivial. Since
$\operatorname{lcm}(20,d_A)=540$, the same local product
and LCM calculation as in the rational transport lemma gives
\[
  \chi_Y(A)/27=\pm10(W_5P_\rho W\mathbf e_Y)_A
                         \in\mathcal O_{\mathbb Q(\zeta_{5^b},i)}.
\]
Here the exact dyadic phase of $A$ is $\rho$; all operators
at five preserve that phase. The integer ring of
$\mathbb Q(\zeta_{5^b},i)$ is $\mathbb Z[\zeta_{5^b}][i]$, since
the two cyclotomic discriminants are coprime. Taking the
coefficient of $1$ in the basis $1,i$ therefore preserves
integrality. The fifth Weyl is fixed by the dyadic field
involution, and the $Y'$ column is the dyadic conjugate of
the $Y$ column. Consequently
\[
                   z_A:=10(W_5v_\rho)_A\in\mathbb Z[\zeta_{5^b}].
\]
This is the point where a real Gaussian column replaces a
rational coordinate column; no coordinate-only integrality
theorem is applied to an unverified average.

Fifth Galois covariance now gives
\[
          10(W_5w_\rho)_A=z_A-\overline{z_A}.
\]
The actual field of a nonzero such anti-real integer has
degree four or twenty. In degree four its squared-modulus
trace is $10(r_A^2+s_A^2)$ in the usual anti-real integral
basis of $\mathbb Q(\zeta_5)$. In degree twenty it is
$2\operatorname{Tr}(z_A\overline{z_A}-z_A^2)\in10\mathbb Z$,
as proved in the rational scope lemma. Summing over the complete
target inertia orbits, with their possible repetitions, gives
\[
 q_\rho:=100\|P_{5,\mathrm{act}}W_5w_\rho\|^2
                                      \in10\mathbb Z_{>0}.
\]
The lower compression bound and unitarity give
$40M_\rho/81\le q_\rho\le50M_\rho/81$.
Taking $n_\rho=q_\rho/10$ proves the first interval.

The real-column norm and moment have mass $162$ and value
$18m_Y$. In a rational positive dyadic word, each nonreal
opposite pair has equal coefficients. If $t$ is the coefficient
at minus one and $b_j$ the common coefficients on the other
nonreal opposite pairs, subtraction of value from mass gives
\[
                         t+\sum_j b_j=9(9-m_Y)=9l.
\]
Apply the preceding interval to the phase minus one when
present and to one member of every present nonreal pair,
then add. This gives the second interval and its exclusions.
\end{proof}


\begin{theorem}[No Gaussian degree-twenty field opposite degree thirty-six]
\label{paper:no-gaussian-quintic-endpoint}
In the no-thirty branch, a compulsory degree-$(36,20)$
crossing cannot have $K_Y=\mathbb Q(\zeta_5,i)$. Consequently
\[
 K_Y=\mathbb Q(\zeta_5),\qquad
 K_X\in\{\mathbb Q(\zeta_5),\mathbb Q(\zeta_5,i),
                         \mathbb Q(\zeta_{15})\}.
\]
\end{theorem}
\begin{proof}
Suppose $K_Y=\mathbb Q(\zeta_5,i)$;
Theorem~\ref{paper:no-cubic-relative-endpoint} gives $K_X=K$.
Within the relative dyadic pair $Y,Y'$, the Weyl matrix is
\[
 \begin{pmatrix}z&\overline z\\\overline z&z\end{pmatrix},
                   \qquad z=(m_Y+in_Y)/18.
\]
The four fifth translates lie in distinct full dyadic blocks,
so the normalized full-orbit vector has
$\langle u_Y,Wu_Y\rangle=m_Y/9$. Hence
\[
             \frac12\le\|Cu_Y\|^2\le\frac{1+m_Y/9}{2}.
\]
This gives $m_Y\ge0$, while the paired-coordinate norm gives
$m_Y\le9$. Gaussian parity makes $m_Y$ even. Thus
$l=9-m_Y$ is one of $1,3,5,7,9$; the real-column transport
interval leaves only $l=9$ and $m_Y=0$. All estimates are
equalities, contradicting the exact-support obstruction.
The final field list combines both endpoint exclusions with
Theorem~\ref{paper:no-cubic-relative-endpoint}, which gives $ab\le2$.
\end{proof}

\begin{corollary}[No norm ten when a degree-thirty-six witness exists]
\label{paper:thirtysix-forbids-norm-ten}
In the no-thirty branch, suppose a compulsory degree-thirty-six
witness exists. At every compulsory degree-twenty witness $Y$,
every degree-ten simple $V$ satisfies $|\chi_Y(V)|^2\ne10$.
\end{corollary}
\begin{proof}
Theorem~\ref{paper:norm-ten-endpoint-field} puts
$\mathbb Q(\zeta_{15})$ inside $K_Y$, whereas the endpoint
field theorem just proved makes $K_Y=\mathbb Q(\zeta_5)$.
\end{proof}

\begin{theorem}[A forced tame degree-twelve witness at an index-two quintic endpoint]
\label{paper:forced-tame-degree-twelve}
Assume the no-thirty branch and fix a compulsory degree-twenty
witness with
\[
 K_Y=\mathbb Q(\zeta_{15})\quad\hbox{or}\quad
 K_Y=\mathbb Q(\zeta_5,i).
\]
There is a compulsory $X\in P_3$ with $d_X=12$, $\theta_X=1$,
$|I_5X|=4$, and full field in the list
\[
 \boxed{\quad
 \mathbb Q(\zeta_5),\quad\mathbb Q(\zeta_5,i),\quad
 \mathbb Q(\zeta_{15}),\quad
 \mathbb Q(\zeta_5)\mathbb Q(\zeta_9)^+,\quad
 \mathbb Q(\zeta_{45}).\quad}
\]
It can be chosen so that $\chi_Z(X)$ is nonzero modulo a
chosen five-place for a jointly dyadic--ternary active label $Z$.
\end{theorem}
\begin{proof}
Both index-two fields forbid every compulsory degree-thirty-six
endpoint, by the preceding theorems. The mixed-support tame-witness
lemma therefore supplies a degree-twelve endpoint with fifth
orbit four and the stated nonzero residue.

Its conductor divides $f_X=2^2 3^4 5^4$. Tame fifth inertia
and the compulsory field dichotomy put $K_X$ either in $K(i)$,
or in $K\mathbb Q(\zeta_{81})$, where $K=\mathbb Q(\zeta_5)$.
Since $K\subseteq K_X$, the latter field has the form $KF$
with $F$ a subfield of the cyclic field $\mathbb Q(\zeta_{81})$.
Put $a=[K_X:K]$.

In the ternary case write $W_{YY}=1-l/18$; in the Gaussian
case put $l=9-m_Y$. In both cases $l$ is positive odd and the
appropriate transport interval excludes $1,3,5,7$. For the
Gaussian case, $|m_Y|\le9$ and its evenness ensure $l>0$;
the transport proof does not require $m_Y\ge0$.
For either field the full-orbit dyadic expectation is
$\langle u_Y,Wu_Y\rangle=1-l/9$.

The degree-$(12,20)$ orbit crossing has squared norm $a/90$.
Therefore the same coarse bound gives
\[
                    \frac a{18}\le\|Cu_Y\|^2\le1-l/18,
                           \qquad a+l\le18.
\]
Thus $a\le9$. The cyclic ternary degrees, divisors of
54, leave $1,2,3,6,9$.

If $a=9$, then $l=9$ and all estimates are equalities.
The vector $u_Y$ has $C$-eigenvalues $1,-1/2$ with squared
component norms $1/3,2/3$. Equality gives
$A_5Cu_Y=\varepsilon\sqrt{1/10}\,u_X$ and $G_3$-fixedness
of $u_X$. Polarized fifth compression yields
\[
 \langle u_X,Cu_X\rangle
 =2\langle u_Y,C^3u_Y\rangle
 =2\left(\frac13-\frac{2/3}{8}\right)=\frac12.
\]
This is also $\langle u_X,Wu_X\rangle$.

The relative field is then the unique degree-nine field
$F=\mathbb Q(\zeta_{27})^+$. Its ternary orbit has nine
labels. The dyadic column is rational and $30W\mathbf e_X$
is integral, by the primary and strict foreign bounds for
$f_X=2^2 3^4 5^4$. Put $c=W_{XX}-1\in\frac1{30}\mathbb Z$.
Uniform ternary block stability gives the orbit matrix
$I_9+cJ_9$; its four fifth translates occupy distinct blocks.
Consequently $1/2=1+9c$, requiring $c=-1/18$, a contradiction.
Thus $a=9$ is impossible. The remaining relative fields are
$\mathbb Q$, $\mathbb Q(i)$, $\mathbb Q(\zeta_3)$,
$\mathbb Q(\zeta_9)^+$ and $\mathbb Q(\zeta_9)$, as claimed.
\end{proof}

\begin{lemma}[A compulsory degree-twenty witness tame at three]
\label{paper:choose-ternary-tame-twenty}
In the no-thirty branch a compulsory degree-twenty witness
$Y\in P_5$ can be chosen whose ternary inertia orbit has
size one or two. Its full field belongs to
\[
 \mathbb Q(\zeta_5),\quad\mathbb Q(\zeta_{25}),\quad
 \mathbb Q(\zeta_{15}),\quad\mathbb Q(\zeta_{75}),\quad
 \mathbb Q(\zeta_5,i),\quad\mathbb Q(\zeta_{25},i).
\]
\end{lemma}
\begin{proof}
Choose a jointly dyadic--quintic active label $Z$ from edge
ten of the activity graph. No-thirty separation makes its
codegree a three-unit, so ternary inertia fixes it. On
\[
                    E=P_5\cap\{A:\theta_{A,2}=\theta_{A,3}=1\},
\]
principal-fiber orthogonality gives
$\sum_{A\in E}d_A\chi_Z(A)=0$.
Reduce at a place above three. The vacuum contributes one;
three-divisible degrees vanish. The three-free nonunit degrees
in $P_5$ are exactly twenty and forty: the principal height
bound gives fifth height one, the complementary principal
degree bound gives dyadic height at least two, and dimension
divisibility gives dyadic height at most three. The degree-forty
rows have dyadic-unit codegree and miss $Z$.

Group the degree-twenty terms into ternary-inertia orbits.
The set $E$ is stable under this inertia, and $\chi_Z(A)$
is constant on each orbit. An orbit with nontrivial wild
ternary part has size divisible by three and contributes zero
to the reduction. Therefore an orbit of size one or two
contributes nontrivially. Choose $Y$ on that orbit.
It has pure fifth twist and is paired with any existing
compulsory ternary witness. The inactive-witness theorem gives
$\theta_Y=1$, and their ordinary crossing puts
$K=\mathbb Q(\zeta_5)$ inside $K_Y$.

The field conductor divides $f_Y=2^2 3^6 5^2$. Tame ternary
inertia and the compulsory field dichotomy put $K_Y$ in
$\mathbb Q(\zeta_{75})$ or $\mathbb Q(\zeta_{100})$.
Each extension over $K$ is cyclic of degree ten, so its
intermediate fields are exactly the displayed alternatives.
\end{proof}

\begin{proposition}[Finite field choices in the degree-twelve-only branch]
\label{paper:degree-twelve-only-field-choices}
Assume the no-thirty branch and suppose no compulsory
degree-thirty-six witness exists. A compulsory pair $(X,Y)$
can then be chosen with $d_X=12$, $d_Y=20$, tame fifth
inertia at $X$, and tame ternary inertia at $Y$.
Put $K=\mathbb Q(\zeta_5)$, $a=[K_X:K]$, and $b=[K_Y:K]$.
The possible relative degrees obey
\[
\begin{array}{c|c|c}
b&K_Y&\text{possible }a\\\hline
1&K&1,2,3,6,9,18\\
2&\mathbb Q(\zeta_{15})\text{ or }K(i)&1,2,3,6\\
5&\mathbb Q(\zeta_{25})&1,2,3,6\\
10&\mathbb Q(\zeta_{75})\text{ or }\mathbb Q(\zeta_{25},i)&1,2,3.
\end{array}
\]
At $a=2$ the two choices are $K(i)$ and $\mathbb Q(\zeta_{15})$.
All other listed relative degrees give the unique subfield
$KF$ with $F\subseteq\mathbb Q(\zeta_{81})$ of that degree.
\end{proposition}
\begin{proof}
Choose $Y$ by the preceding lemma. With this fixed $Y$,
Lemma~\ref{paper:mixed-support-forces-tame-twelve} chooses
a fifth-tame $X$; its mixed-support hypothesis is vacuous
because the degree-thirty-six witnesses are absent.
The conductor and field dichotomy put $K_X$ in $K(i)$ or
$K\mathbb Q(\zeta_{81})$. Hence its possible relative
degrees before the crossing bound are
$1,2,3,6,9,18,27,54$.

Both twists are one. The exact inactive fifth Fourier bound
at the ordinary degree-$(12,20)$ crossing gives
\[
                     45ab\le\frac{D}{5\cdot60^2},
                                  \qquad ab\le36.
\]
For $b=5,10$ this gives the stated rows immediately. For
$b=2$, Theorem~\ref{paper:forced-tame-degree-twelve} gives
the stronger stated list, choosing its witness against this
same $Y$ if necessary.

It remains to exclude $a=27$ when $b=1$. Write
$W_{YY}=1-l/18$. The rational degree-twenty transport interval
gives positive odd $l\ge9$. On each nontrivial fifth orbit
vector $u_Y$, the dyadic expectation is $1-l/18$. The exact
crossing and coarse Borel bound give
\[
                 \frac a{36}\le\|Cu_Y\|^2\le1-l/36.
\]
Thus $a+l\le36$. If $a=27$, then $l=9$ and all estimates
are equalities. The $C$-eigenvalue weights at $1,-1/2$ are
$2/3,1/3$, and the exact opposite image has squared norm
$3/20$. Polarized fifth compression, as in the proof of
Theorem~\ref{paper:no-cubic-relative-endpoint}, gives
\[
                    \langle u_X,Wu_X\rangle=5/6.
\]
Here $u_X$ is $G_3$-fixed by the equality case.

The ternary orbit of $X$ has twenty-seven labels, and its
dyadic column is rational with integral multiplier thirty.
Uniform ternary block stability and the four disjoint fifth
blocks make the last expectation $1+27c$, with
$c=W_{XX}-1\in\frac1{30}\mathbb Z$. Equality to $5/6$
would give $c=-1/162$, impossible. This removes $a=27$
and proves the table.
\end{proof}

The field table concerns a jointly chosen compulsory pair. Its wild
fifth alternatives remain arithmetic possibilities.

\subsection{Short adjoints at a degree-twelve endpoint}

\begin{lemma}[Integer radii of two-summand degree-twelve products]
\label{paper:degree-twelve-two-summand-radii}
Let $X,B$ be degree-twelve simples in the ordinary perfect
fusion-simple reduction at $D=5400^2$. Suppose $XB$ is the
sum of two distinct nonunit simples, each with multiplicity one.
If $c=|\chi_X(B)|$ is an integer, then
\[
                            c\in\{0,3,6,8\}.
\]
More precisely the possibilities are
\[
\begin{array}{c|c|c}
\text{summand degrees}&c&\text{relative twist ratio}\\\hline
(24,120)&8&-1\\
(36,108)&6&-1\\
(54,90)&3&-1\\
(72,72)&0&-1\\
(72,72)&6&\text{a primitive cube root}.
\end{array}
\]
\end{lemma}
\begin{proof}
The displayed four degree pairs are the only non-prime-power
divisors of $5400$, with degree six omitted, that sum to $144$.
Balancing writes the squared character modulus as, respectively,
\[
 104+20r,\qquad 90+27r,\qquad
 \frac{306+135r}{4},\qquad 72+36r,
\]
where $r=\rho+\rho^{-1}$ for a root of unity $\rho$.
If $c$ is an integer, $r$ is a rational algebraic integer in
$[-2,2]$, hence is one of $-2,-1,0,1,2$. Direct substitution
gives the table; the value $c=12$ is forbidden by strictness.
\end{proof}

\begin{proposition}[Exclusion of the degree-twelve adjoint $1+18+50+75$]
\label{paper:degree-twelve-185075-excluded}
Let $X\in P_3$, $d_X=12$, and let $Y\in P_5$, $d_Y=20$,
be a compulsory pair with an ordinary crossing
\[
                     S_{XY}/60=1+3\xi_5.
\]
Then
\[
                  XX^*\ne\mathbf1+V_{18}+U_{50}+Z_{75}.
\]
The proposition itself does not require the no-thirty hypothesis.
\end{proposition}
\begin{proof}
Suppose the displayed adjoint occurs. Its three constituents
are self-dual, because their degrees are distinct. Put
$t=\theta_Z$, $u=\theta_V$, and $v=\theta_U$.
Balancing and character integrality give
\[
        T:=1+75t+18u+50v\in12\mathcal O.
\]
Modulo three, this makes $v-1$ divisible by three. Since
all conjugates of $v-1$ have modulus at most two, the
algebraic-integer norm argument forces $v=1$.
Modulo two, $(1+t)/2$ is integral, whence $t=1$ or $-1$
by the unit-disk argument. If $t=-1$, then
$T=-24+18u\in12\mathcal O$ would make $3u/2$ integral,
which is impossible. Thus $t=1$. The same divisibility
now gives $1+u\in2\mathcal O$, so $u=1$ or $-1$.
The first choice gives the forbidden self-character modulus
twelve. Consequently
\[
 \theta_Z=\theta_U=1,\quad \theta_V=-1,\qquad
                  \chi_X(X^*)=9\theta_X^{-2}.
\]

At least two adjoint constituents must be fixed by every
Galois permutation. Otherwise choose a permutation $g$
outside the two proper stabilizers of two nonfixed constituents.
The adjoints $XX^*$ and $(gX)(gX)^*$ must have a common
nonunit constituent: if their only common constituent were
the unit, the product $X(gX)^*$ would be simple, contradicting
strictness of the character bound. Distinct degrees therefore
make $g$ fix the remaining constituent and give
$\|X(gX)^*\|^2=2$. Its unit is absent, since $gX\ne X$.
It is thus the sum of two distinct nonunits. Argument-side
Galois covariance gives its character at $X$ modulus nine,
contrary to Lemma~\ref{paper:degree-twelve-two-summand-radii}.

The degree-eighteen constituent cannot be globally fixed:
Lemma~\ref{paper:modular:mass-two-quintic-character} and self-duality
give
\[
                     \chi_Y(V)=9s,
\]
where $s$ is one of the two primitive fifth-root real traces.
Thus $Z$ and $U$ are both globally fixed, and their values
at $Y$ are rational integers. On the other hand the ordinary
crossing gives
\[
       |\chi_Y(X)|^2=90+27s_0
                 =1+\chi_Y(Z)+\chi_Y(U)+9s,
\]
where $s_0$ is also a primitive fifth-root real trace.
The two possibilities are $(-1\pm\sqrt5)/2$. The irrational
coefficient of $27s_0-9s$ cannot vanish. This contradicts
the rationality of the remaining terms.
\end{proof}

\begin{lemma}[A cyclic real field at a compulsory degree-twelve endpoint]
\label{paper:degree-twelve-real-field-cyclic}
In the no-thirty branch, the maximal real subfield of the
full character field of a compulsory degree-twelve endpoint
is cyclic over $\mathbb Q$. In particular it contains no
real biquadratic subfield.
\end{lemma}
\begin{proof}
Its full field contains $\mathbb Q(\zeta_5)$ and has
conductor dividing $f_X=2^2 3^4 5^4$. The compulsory
field dichotomy makes it unramified at two or at three.
It is therefore contained in one of
\[
                   \mathbb Q(\zeta_{81\cdot625}),
                   \qquad \mathbb Q(\zeta_{4\cdot625}).
\]
Both maximal real cyclotomic subfields are cyclic. For the
first, the full Galois group is $C_{54}\times C_{500}$;
quotienting by simultaneous complex conjugation makes its
two-primary part cyclic of order four, while its odd part
is $C_{27}\times C_{125}$, hence cyclic. For the second,
the quotient of $C_2\times C_{500}$ by simultaneous
conjugation is $C_{500}$. Every subfield of a cyclic
extension is cyclic, proving the assertion.
\end{proof}

\begin{proposition}[Exclusion of the degree-twelve adjoint $1+15+20+108$]
\label{paper:degree-twelve-1520108-excluded}
Assume the no-thirty branch. Let $X\in P_3$ have degree
twelve and pure ternary twist, and let $Y\in P_5$ be a
compulsory degree-twenty witness with
\[
 K_Y\in\{\mathbb Q(\zeta_5),\mathbb Q(\zeta_5,i),
                                      \mathbb Q(\zeta_{15})\}.
\]
Then
\[
                     XX^*\ne\mathbf1+Z_{15}+U_{20}+V_{108}.
\]
\end{proposition}
\begin{proof}
All three constituents would be self-dual. The three-adjoint
field dichotomy, Lemma~\ref{paper:three-adjoint-fields}, and
the preceding cyclic-real-field lemma force at least one of
them to be globally Galois fixed.

Write $r_A=\chi_Y(A)$. If $V$ were fixed, LCM integrality
would give $r_V=27n$ for an integer $n$ with $|n|<4$.
Principality at five gives $27n\equiv108\pmod5$, hence
$n\equiv4\pmod5$ and $n=-1$. The adjoint identity would
then imply
\[
 |\chi_Y(X)|^2=1+r_Z+r_U-27<1+15+20-27=9.
\]
But the ordinary compulsory crossing gives
\[
 |\chi_Y(X)|^2=90+27s_0
       \ge\frac{153-27\sqrt5}{2}>9,
\]
a contradiction. Thus a globally fixed constituent must be
$Z$ or $U$.

The full conductor of $V$ divides
$f_V=D/108^2=2500=2^2 5^4$. Put $K=\mathbb Q(\zeta_5)$.
If $K_Y=K$ or $\mathbb Q(\zeta_{15})$, its character at
$Y$ lies in $K_V\cap K_Y\subseteq K$, and is real by
self-duality. LCM integrality consequently gives
\[
                r_V=27z,\qquad
                z\in\mathcal O_{\mathbb Q(\sqrt5)}.
\]
For the Gaussian alternative, LCM integrality first gives
$r_V=27z$ with $z\in\mathcal O_{K(i)}$ real. Write
$z=a+ib$, with $a,b\in\mathcal O_K$; this is an integral
decomposition because $\mathcal O_{K(i)}=\mathcal O_K[i]$.
One of $r_Z,r_U$ is rational, and the adjoint identity makes
the other equal to a number in $K$ minus $27z$. Applying the
involution $i\mapsto-i$ to this last identity gives a difference
of magnitude $54|b|$. All conjugates of each small character
have modulus strictly below its dimension, at most twenty.
Hence every conjugate of $b$ has modulus below $40/54<1$.
The algebraic-integer norm argument gives $b=0$. Since $z$
is real, the same displayed conclusion holds in this case too.
The adjoint identity is
\[
                r_Z+r_U=89+27(s_0-z).
\]
One of $r_Z,r_U$ is rational, since its label is globally
fixed. If the other were irrational, its two conjugates in
$\mathbb Q(\sqrt5)$ would differ in absolute value by at
least $27\sqrt5$: the ring of integers is $\mathbb Z[s_0]$,
and the coefficient of $s_0$ on the right is a multiple of
twenty-seven. Strictness bounds the two absolute values of
that other character by its dimension, at most twenty, so
their difference is strictly below forty. This is impossible.
Therefore both $r_Z,r_U$ are rational integers.

LCM integrality and the principal-five residue give
$r_Z\in3\mathbb Z\cap5\mathbb Z$. Since $|r_Z|<15$,
one has $r_Z=0$. Likewise $r_U=t\in5\mathbb Z$ with
$|t|<20$. The preceding identity now gives
$t\equiv89\equiv8\pmod{27}$, because $s_0-z$ is a
rational algebraic integer. None of
$-15,-10,-5,0,5,10,15$ satisfies this congruence. This is
the final contradiction.
\end{proof}

\begin{corollary}[At least four nonunit occurrences in a compulsory degree-twelve adjoint]
\label{paper:degree-twelve-adjoint-at-least-four}
Under the preceding hypotheses, $XX^*$ contains at least
four nonunit simple occurrences, counted with multiplicity.
\end{corollary}
\begin{proof}
The nonunit part has dimension $143$. The allowed nonunit
dimensions at most $143$ are
\[
 10,12,15,18,20,24,30,36,40,45,50,54,60,72,75,90,
 100,108,120,135.
\]
There is no one- or two-term partition of $143$ from this
list, and its only three-term partitions, allowing repetition,
are $15+20+108$ and $18+50+75$. The two propositions exclude
both, leaving at least four nonunit occurrences.
\end{proof}


\subsection{Small character norms in the no-thirty branch}

\begin{proposition}[A numerical quintic norm with no actual no-thirty word]
\label{paper:beta-face-obstruction}
Assume the no-thirty branch at $D=5400^2$. Let $Y$ be a
compulsory degree-twenty witness, so $\theta_Y=1$, and let
$W$ have degree ten. Put
\[
 s=\zeta+\zeta^{-1},\qquad
 \beta_0=(\omega-s)(1-\zeta^2)(1+2\zeta),
\]
where $\omega,\zeta$ have orders three and five. Then
$\chi_Y(W)$ cannot be a root-of-unity multiple of a Galois
conjugate of $\beta_0$.
\end{proposition}
\begin{proof}
Every dyadically active simple of degree at most forty is
inactive at both odd primes. Indeed, activity at two and three
would force five-unit codegree by
Lemma~\ref{paper:no-thirty-local-separation}, hence degree
divisible by twenty-five and at least fifty. Activity at two
and five similarly forces degree divisible by twenty-seven
and at least fifty-four. Prime-power nonunit degrees are absent.

Define a rational linear functional $\lambda$ on
$\mathbb Q(\zeta_{15})$ by the following values on
$\omega^a\zeta^b$:
\[
 \begin{array}{c|rrrrr}
 a\backslash b&0&1&2&3&4\\\hline
 0&-1&-1&0&1&1\\
 1&1&1&-1&-1&0\\
 2&0&0&1&0&-1.
 \end{array}
\]
The row and column sums vanish, so these values respect all
cyclotomic relations. One has $|\lambda(\rho)|\le1$ for
$\rho\in\mu_{30}$ and $\lambda(\beta_0)=10$.
Extend $\lambda$ to a common cyclotomic field by normalized
relative trace. Every root outside $\mu_{30}$ then has value
zero: its higher primary component is killed by the corresponding
relative trace. Thus every positive word of mass ten for
$\beta_0$ is supported on the ten roots where $\lambda=1$.

Comparing coefficients in the basis
$\omega^i\zeta^j$, $0\le i<2$, $0\le j<4$, gives all
such words:
\begin{align*}
 \beta_0={}&(a+b)\zeta^3+a\zeta^4
 +(1+b)\omega+(2+b)\omega\zeta
 +(1-a-b)\omega^2\zeta^2\\
 &-(3-a-b)-(1-a-b)\zeta
       -a\omega\zeta^2-(2-b)\omega\zeta^3
       -b\omega^2\zeta^4,
 \qquad a,b\ge0,\quad a+b\le1.
\end{align*}
The minus signs here belong to the roots, so all displayed
multiplicities are nonnegative. There are exactly two free
coefficients: the two relations among these ten roots are
obtained by eliminating $\omega^2$ and $\zeta^4$ in the
indicated basis. The parameters may be rational, as in an
actual balancing word.

Suppose such a word were actual. A phase of coefficient $c$
corresponds to a constituent bucket of dimension $20c$.
Write $\vartheta=\theta_W$.
If $\vartheta_2\ne-1$, the phase $-\omega\zeta^3$, whose
coefficient $2-b$ lies in $[1,2]$, gives a dyadically active
bucket of dimension at most forty. It must be pure dyadic,
so $\vartheta_3=\omega^2$ and $\vartheta_5=\zeta^2$.
But the phase $-1$ has coefficient at least two and would
then give a thirty-active constituent, a contradiction.

If $\vartheta_2=-1$, use the phase $\omega$, whose coefficient
$1+b$ also lies in $[1,2]$. Its bucket must be pure dyadic,
giving $\vartheta_3=\omega^2$ and $\vartheta_5=1$.
The two positive phases $\zeta^3$ and $\omega^2\zeta^2$
have coefficients $a+b$ and $1-a-b$, respectively. At least
one occurs, and either would give a thirty-active constituent.
This is again impossible.

The argument allows an arbitrary common twist shift. It therefore
also excludes root-of-unity multiples of $\beta_0$, and
Galois conjugation merely replaces the chosen primitive roots.
\end{proof}

\begin{lemma}[The odd cyclotomic realizations of a particular quadratic norm]
\label{paper:odd-field-norm-realizations}
Let $\omega,\zeta$ be primitive roots of orders three and
five, choose $s=\zeta+\zeta^{-1}=(\sqrt5-1)/2$, and put
\[
 A=25+9\sqrt5,\qquad
 \beta_0=(\omega-s)(1-\zeta^2)(1+2\zeta).
\]
If an algebraic integer $\beta\in
L=\mathbb Q(\zeta_{3^a5^b})$, $a,b\ge1$, satisfies
$\beta\overline\beta=A$, then
\[
            \beta=\rho\,\sigma(\beta_0)
\]
for a root of unity $\rho\in L$ and
$\sigma\in\operatorname{Gal}(\mathbb Q(\zeta_{15})/
\mathbb Q(\sqrt5))$. For the conjugate norm $25-9\sqrt5$,
the analogous assertion uses conjugates of $\beta_0$ carrying
$\sqrt5$ to $-\sqrt5$.
\end{lemma}
\begin{proof}
Put $E=\mathbb Q(\sqrt5)$ and $K=\mathbb Q(\zeta_{15})$.
Directly, $\beta_0\overline\beta_0=A$ and
$N_{E/\mathbb Q}(A)=220=2^2\cdot5\cdot11$.
At two, $K$ has two primes, interchanged by complex conjugation,
and $A$ has valuation one at each. At eleven, only the two
primes over one of the two primes of $E$ occur in $(A)$;
they are also interchanged by conjugation and have valuation
one. At five, $K$ has a unique prime, fixed by conjugation,
and $A$ has valuation two. These assertions follow from
the cyclotomic decomposition data
\[
 \begin{array}{c|ccc}
 p&e(K/\mathbb Q)&f(K/\mathbb Q)&\text{number of primes}\\\hline
 2&1&4&2\\
 5&4&2&1\\
 11&1&2&4.
 \end{array}
\]
Thus a norm realization has two choices for its ideal at two,
two choices at eleven, and exponent one forced at five.

Identify $\operatorname{Gal}(K/\mathbb Q)$ with
$(\mathbb Z/15)^\times$. The subgroup fixing $E$ is
$\{1,4,11,14\}$. Its kernels on the two ideal choices are
$\{1,4\}$ at two and $\{1,11\}$ at eleven. These are
distinct subgroups of index two, so it acts independently
and transitively on all four possible ideal patterns.
Consequently the conjugates of $(\beta_0)$ realize all four.

The same patterns are the only ones in $L$. Indeed,
\[
\begin{array}{ll}
 \operatorname{ord}_{3^a}(2)=2\cdot3^{a-1},&
 \operatorname{ord}_{5^b}(2)=4\cdot5^{b-1},\\
 \operatorname{ord}_{3^a}(11)=2\cdot3^{a-1},&
 \operatorname{ord}_{5^b}(11)=5^{b-1}.
\end{array}
\]
Therefore every prime of $K$ over two or eleven is inert
in $L/K$. At five, the ramification index in $L/K$ is
$5^{b-1}$ and the relative residue degree is $3^{a-1}$;
there is still one prime, since
$\operatorname{ord}_{3^a}(5)=2\cdot3^{a-1}$.
The valuation of $A$ there is $2\cdot5^{b-1}$, so the
valuation of $\beta$ is forced to be $5^{b-1}$, exactly
the extension of the base ideal pattern.

Hence $(\beta)=(\sigma\beta_0)\mathcal O_L$ for a
conjugate fixing $E$. Their quotient $u$ is a unit and
$u\overline u=1$. All conjugates of $u$ have modulus one,
so Kronecker's theorem makes $u$ a root of unity. No
class-number hypothesis is used. Conjugating the entire
argument proves the last assertion.
\end{proof}

\begin{corollary}[An odd-field norm obstruction in the no-thirty branch]
\label{paper:no-thirty-odd-norm-obstruction}
In the no-thirty branch, let $Y$ be a compulsory degree-twenty
witness and let $W$ have degree ten. If $\chi_Y(W)$ lies in
an odd-conductor cyclotomic field, then
\[
                |\chi_Y(W)|^2\ne25\pm9\sqrt5.
\]
\end{corollary}
\begin{proof}
All field primes divide thirty. Enlarge the odd cyclotomic
field to $\mathbb Q(\zeta_{3^a5^b})$ and apply the lemma
and Proposition~\ref{paper:beta-face-obstruction}.
\end{proof}


\begin{lemma}[A low-house mixed-inertia identity]
\label{paper:low-house-mixed-inertia}
Assume the no-thirty branch, let $Y$ be a compulsory
degree-twenty witness, and let $W$ be any simple. Put
$\beta=\chi_Y(W)$ and $\gamma=\theta_W\beta$.
If $\operatorname{house}(\beta)<27/4$, then for every dyadic
inertia automorphism $\sigma_2$ and every fifth-inertia
automorphism $\sigma_5$ in a common cyclotomic splitting field,
\[
              (\sigma_2-1)(\sigma_5-1)\gamma=0.
\]
\end{lemma}
\begin{proof}
Here $\theta_Y=1$. The actual balancing formula is
\[
            20\gamma=\sum_Z N_{YW}^Z d_Z\theta_Z.
\]
The mixed difference annihilates every summand inactive at
two or at five. Each remaining summand is active at both
primes. The no-thirty local separation lemma gives three-unit
formal codegree for that label, hence $27\mid d_Z$.
Consequently, writing $\Delta=(\sigma_2-1)(\sigma_5-1)$,
\[
                       20\Delta\gamma\in27\mathcal O.
\]
Since $20$ is prime to $27$, one has
$\Delta\gamma\in27\mathcal O$. On the other hand,
\[
       \operatorname{house}(\Delta\gamma)
           \le4\operatorname{house}(\beta)<27.
\]
The algebraic integer $\Delta\gamma/27$ has strict house
below one and must be zero.
\end{proof}

\begin{proposition}[A quintic norm obstruction without ternary ramification]
\label{paper:quintic-small-norm-obstruction}
Retain the no-thirty hypotheses. Let $\beta=\chi_Y(W)$ be
contained in $\mathbb Q(i,\zeta_{5^b})$ for some $b\ge1$,
and suppose $5\mid d_W$. Put $F=\mathbb Q(\zeta_{5^b})$.
Assume that
\begin{enumerate}
\item $A=\beta\overline\beta$ belongs to $F$;
\item $\operatorname{house}(\beta)<27/4$, and some real
embedding of $A$ is strictly below twenty-five;
\item the valuation of $A$ at the unique prime of $F$ above
two is odd.
\end{enumerate}
Then these conditions are contradictory.
\end{proposition}
\begin{proof}
The norm $A$ is nonzero by its odd valuation. Since $Y\in P_5$
and $5\mid d_W$, $\beta$ vanishes at every five-place.
Factor the root of unity $\theta_W=\tau_2\tau_3\tau_5$
into primary parts. Increase $b$ if necessary to include
$\tau_5$, and set $\eta=\tau_2\tau_5\beta$.
Increasing $b$ is unramified at two and preserves the odd
valuation. Lemma~\ref{paper:low-house-mixed-inertia}, after
removing the fixed three-primary scalar, gives
\[
                    (\sigma_2-1)(\sigma_5-1)\eta=0.
\]
Also $\eta\overline\eta=A$, and $\eta$ still vanishes at
every five-place.

Suppose first that $\tau_2\notin\mu_4$. Writing
$\beta=a+ib$ with $a,b\in\mathcal O_F$, both dyadic
normalized traces of $\tau_2$ and $i\tau_2$ are zero.
Thus the full dyadic average $E_2(\eta)$ is zero.
Averaging the mixed identity in the dyadic variable gives
$\sigma_5\eta=\eta$ for every fifth-inertia automorphism.
Hence $\eta$ lies in a two-primary cyclotomic field. It
follows that $A\in F\cap\mathbb Q(\mu_{2^\infty})
=\mathbb Q$. By the assumed small embedding, $0<A<25$.
That two-primary field is unramified at five, so vanishing at
all five-places gives $\eta\in5\mathcal O$. Every conjugate
of $\eta/5$ then has absolute value $\sqrt A/5<1$, which is
impossible for a nonzero algebraic integer.

Now suppose $\tau_2\in\mu_4$. Absorb it and $\tau_5$
into the coefficients and write
\[
                         \eta=a+ib,
                       \qquad a,b\in\mathcal O_F.
\]
Apply the mixed identity to the automorphism $i\mapsto-i$.
It says that $b$ is fixed by every fifth-inertia automorphism,
so $b\in\mathbb Z$. Since both $\eta$ and its dyadic
conjugate vanish at every five-place, their difference $2ib$
does too. Therefore $5\mid b$.

The equality $\eta\overline\eta=A\in F$ is fixed by
dyadic conjugation. Comparing it with the corresponding norm
of $a-ib$ gives $b(a-\overline a)=0$. If $b\ne0$, then
$a$ is real and
\[
                         A=a^2+b^2.
\]
Every real embedding of $A$ would consequently be at least
$b^2\ge25$, contrary to the hypothesis. Thus $b=0$, so
$\eta\in F$.

The prime two is inert in $F$: its order modulo $5^b$ is
$4\cdot5^{b-1}$. In particular complex conjugation fixes the
unique prime over two. Hence
\[
                 v_2(A)=v_2(\eta\overline\eta)
                        =2v_2(\eta)
\]
is even, the final contradiction.
\end{proof}

\begin{lemma}[Odd cyclotomic integers of norm ten]
\label{paper:odd-norm-ten-field}\label{paper:odd-norm-ten-realizations}
Let $L=\mathbb Q(\zeta_{3^a5^b})$, $a,b\ge1$, and put
$\alpha=(5+\sqrt{-15})/2$. Every $\beta\in\mathcal O_L$ with
$\beta\overline\beta=10$ is a root-of-unity multiple of
$\alpha$ or $\overline\alpha$, and
$\mathbb Q(\sqrt{-15})\subseteq\mathbb Q(\beta)$.
If $g$ is a fifth-inertia automorphism that is a nonsquare modulo five,
then $\beta\,\overline{g\beta}$ is ramified at five.
\end{lemma}
\begin{proof}
The order of two modulo $3^a5^b$ is
$4\cdot3^{a-1}5^{b-1}=[L:\mathbb Q]/2$.
Thus $L$ has two primes above two, exchanged by conjugation. Their
decomposition subgroup is the kernel of the product of the quadratic
characters modulo three and five, with fixed field
$\mathbb Q(\sqrt{-15})$. The valuations of $\beta$ at these two
unramified primes are $(0,1)$ or $(1,0)$. Any automorphism fixing
$\beta$ preserves the prime with positive valuation, proving the field
inclusion. A nonsquare fifth automorphism and conjugation each exchange
the primes, so $\beta\,\overline{g\beta}$ has valuations $(0,2)$ or
$(2,0)$. It cannot lie in the three-primary subfield, where two is inert;
hence it ramifies at five.

There is a unique prime above five because five has order
$2\cdot3^{a-1}$ modulo $3^a$. Its ramification index is
$4\cdot5^{b-1}$, and the norm equation fixes the valuation of $\beta$
there at $2\cdot5^{b-1}$. No other prime divides its ideal.
The integers $\alpha,\overline\alpha$ realize the two possible ideal
patterns. Dividing $\beta$ by the matching one gives a unit $u$ with
$u\overline u=1$. All conjugates of $u$ have modulus one, so
Kronecker's theorem makes $u$ a root of unity.
\end{proof}

\begin{theorem}[Norm-ten character rigidity]
\label{paper:norm-ten-needs-three}\label{paper:norm-ten-dyadic-twist}
\label{paper:norm-ten-actual-phase}\label{paper:norm-ten-untwisted}
\label{paper:norm-ten-endpoint-field}
\label{paper:norm-ten-negative-profile-fixedness}
Assume no actual twist is divisible by thirty. Let $Y\in P_5$ be a
compulsory degree-$20$ witness, and let $W$ be simple of degree ten with
$\beta=\chi_Y(W)$ and $\beta\overline\beta=10$. Then
\[
 \theta_W=1,\qquad
 \beta\in\left\{\pm\frac{5+\sqrt{-15}}2,
                  \pm\frac{5-\sqrt{-15}}2\right\},\qquad
 \mathbb Q(\zeta_{15})\subseteq K_Y,
\]
and $K_Y$ is unramified at two. In particular no compulsory degree-$36$
witness exists. For every Galois automorphism $g$, a decomposition of
$W(gW)^*$ into two distinct nonunit simples cannot have degree pair
$(50,50)$.
If the adjoint of $W$ has profile $(15,24,60)$ or $(24,30,45)$ and its
odd-degree constituent has twist $-1$, at least two of its three
nonunit constituents are globally Galois fixed.
\end{theorem}
\begin{proof}
The entry lies in $K_W\cap K_Y$, so its conductor divides
$\gcd(D/10^2,D/20^2)=2^2 3^6 5^2$. If it were unramified at three,
it would belong to $\mathbb Q(i,\zeta_{25})$, contradicting
Proposition~\ref{paper:quintic-small-norm-obstruction} at $A=10$:
the two-valuation is odd, $\sqrt{10}<27/4$, and $10<25$.
Since $K_Y$ also contains $\mathbb Q(\zeta_5)$, the compulsory field
dichotomy makes it unramified at two. The lemma therefore applies to
$\beta$ and puts $\mathbb Q(\sqrt{-15})$ in both $K_Y$ and $K_W$.

Both odd inertia groups move $W$. If $\theta_{W,2}\ne1$, the no-thirty
operator identity on that exact dyadic twist fiber gives
$(g_3-I)(g_5-I)\mathbf e_W=0$. Choose $g_3,g_5$ moving $W$.
Their unsigned permutation matrices make the coefficient at $W$ equal
to $1+\mathbf1_{g_3g_5W=W}>0$, a contradiction. Thus
$\theta_{W,2}=1$, including exclusion of the rational phase $-1$.

Write $\gamma=\theta_W\beta$. In a sufficiently large odd cyclotomic
field, the lemma gives $\gamma=\varepsilon\lambda\alpha$ or
$\varepsilon\lambda\overline\alpha$, where $\varepsilon=\pm1$ and
$\lambda$ has odd order. Let $\sigma_3,\sigma_5$ square the roots at
their respective primes and fix the other primary factors, and put
$\Delta=(\sigma_3-1)(\sigma_5-1)$. Balancing gives
\[
 20\gamma=\sum_A N_{YW}^A d_A\theta_A,
 \qquad 20\Delta\gamma\in8\mathcal O,
 \qquad \Delta\gamma\in2\mathcal O.
\]
Indeed only constituents active at both odd primes survive the mixed
difference, and no-thirty separation makes their codegrees odd.
At the prime above two where $\alpha$ is a unit, its reduction is one
and that of $\overline\alpha$ is zero, since
$\alpha^2-5\alpha+10=0$. The two automorphisms each interchange
$\alpha,\overline\alpha$, while their product fixes them. Thus
$0\equiv\Delta(\lambda\alpha)\equiv\lambda^2+\lambda$.
Reduction at two is injective on odd-order roots, so $\lambda=1$.
The conjugate case is identical. Hence
$\theta_W\beta\in\{\pm\alpha,\pm\overline\alpha\}$.

Apply this conclusion also to $W^*$, which has the same twist and
conjugate character value. Their product is $10\theta_W^2$.
After division by ten, the possible products are
$\pm1$, $\pm\alpha/\overline\alpha$, and their inverses.
The latter ratios have unequal two-valuations and are not roots of
unity. Thus $\theta_W^2=\pm1$; its odd order forces $\theta_W=1$.
This proves the stated values. Their quadratic field, together with
$\mathbb Q(\zeta_5)\subseteq K_Y$, gives
$\mathbb Q(\zeta_{15})\subseteq K_Y$. A compulsory degree-$36$
witness would instead force $K_Y=\mathbb Q(\zeta_5)$ by
Theorem~\ref{paper:no-gaussian-quintic-endpoint}, so none exists.

Suppose $B=gW$ and $WB^*$ consists of two distinct degree-$50$ simples.
Rigidity gives $\|WB\|^2=\|WB^*\|^2=2$. The product $WB$ contains
no unit, since otherwise $B=W^*$ and the adjoint of $W$ would have
only one nonunit occurrence. Thus $WB$ also has two distinct nonunit
summands. The conjugate balancing character values at $W$, and
Lemma~\ref{paper:two-summand-phase-product}, force its degree pair to
be $(50,50)$ too; the other possible pairs have positive minimum
modulus, so this also covers a zero character value.
Put $\beta'=g\beta$. One of $\beta\beta'$ and
$\beta\overline{\beta'}$ is $\pm\alpha^2$ or
$\pm\overline{\alpha}^{\,2}$, hence generates
$\mathbb Q(\sqrt{-15})$ and ramifies at five.
But these are $\chi_Y(WB)$ and $\chi_Y(WB^*)$.
Both sums must be unramified at five because a degree-$50$ simple has
formal codegree $11664$, a five-unit. This contradiction excludes the
pair. Finally, in either stated negative adjoint profile, two nonfixed
constituents would produce such a pair by
Propositions~\ref{paper:fifteen-twentyfour-sixty} and
\ref{paper:twentyfour-thirty-fortyfive}. At least two constituents are
therefore globally fixed.
\end{proof}

\begin{corollary}[A higher dyadic twist fixes the ternary scalar]
\label{paper:degree-ten-higher-dyadic-scalar}
In the no-thirty branch, write
$\chi_X(W)= -5\eta$ for a degree-ten $W$ and a compulsory
ternary witness $X$, with $\eta$ three-primary. If
$\operatorname{ord}(\theta_{W,2})\ge4$, then $\eta=1$.
\end{corollary}
\begin{proof}
Apply a mixed dyadic/ternary inertia difference to the actual
balancing formula for $XW$. Every output active at both two
and three has degree divisible by twenty-five. Since $d_X$
is prime to five and $\theta_X=1$, this gives
\[
          (\sigma_2-1)(\sigma_3-1)(-5\theta_W\eta)
                               \in25\mathcal O.
\]
Its house is at most twenty, so it vanishes. The nontrivial
dyadic inertia on $\theta_{W,2}$ consequently forces
$\theta_{W,3}\eta=1$. A degree-ten object active at two
cannot also be active at three in the no-thirty branch,
so $\theta_{W,3}=1$ and $\eta=1$.
The order assumption cannot be weakened to dyadic activity:
the rational phase $-1$ is fixed by all dyadic inertia.
\end{proof}


\begin{corollary}[The quadratic norm is impossible without thirty]
\label{paper:no-thirty-quadratic-norm-obstruction}
In the no-thirty branch, let $Y$ be a compulsory degree-twenty
witness and let $W$ have degree ten. Then
\[
                   |\chi_Y(W)|^2\ne25\pm9\sqrt5.
\]
\end{corollary}
\begin{proof}
Suppose one equality holds and put $\beta=\chi_Y(W)$.
Its full label field $K_W$ contains $\mathbb Q(\sqrt5)$.
The no-thirty field classification therefore makes $K_W$
unramified at two or at three: the exceptional biquadratic
case contains no pure-primary quadratic field, and ramification
at five is already present.

In the first case $\beta$ lies in an odd cyclotomic field,
contradicting Corollary~\ref{paper:no-thirty-odd-norm-obstruction}.
In the second case its conductor divides
that of both full endpoint fields, since
$\beta=S_{WY}/20\in K_W\cap K_Y$.
The exact conductor theorem therefore gives
\[
 \operatorname{cond}(\mathbb Q(\beta))
        \mid\gcd(f_W,f_Y)=72900=2^2 3^6 5^2.
\]
This conductor is prime to three, so $\beta$ belongs to
$\mathbb Q(\zeta_{4\cdot5^2})$. Apply
Proposition~\ref{paper:quintic-small-norm-obstruction}.
Here $4\sqrt{25+9\sqrt5}<27$ and one norm embedding is
$25-9\sqrt5<25$. At the inert prime two of
$\mathbb Q(\sqrt5)$, the valuation of $25+9\sqrt5$ is
one, since its rational norm is $220$. The cyclotomic
extension at two is unramified, so this odd valuation persists.
All hypotheses of the proposition hold, a contradiction.
Galois conjugation handles the other sign.
\end{proof}

\begin{proposition}[The norm-twenty-five boundary]
\label{paper:norm-twentyfive-boundary}
Assume the no-thirty branch. Let $Y\in P_5$ be a compulsory
simple of degree twenty and let $W$ have degree ten. If
$|\chi_Y(W)|^2=25$, then
\[
                         \chi_Y(W)=5\rho
\]
for a root of unity $\rho$.
\end{proposition}
\begin{proof}
Put $\beta=\chi_Y(W)$. By the compulsory field dichotomy,
$K_Y$ is unramified at two or at three. In the first case
$\beta$ belongs to an odd field
$L=\mathbb Q(\zeta_{3^a5^b})$, allowing $a=0$.
There is a unique prime above five in $L$, fixed by complex
conjugation. The norm equation forces $(\beta)=(5)$.
Thus $\beta/5$ is a unit all of whose conjugates have modulus
one, and is a root of unity.

In the second case the conductor bound for the entry gives
$\beta\in\mathbb Q(i,\zeta_{25})$. Principality of $Y$
makes $\beta$ vanish at every place above five. Write
$\theta_W=\tau_2\tau_3\tau_5$ in primary parts, enlarge
$F=\mathbb Q(\zeta_{5^b})$ to include $\tau_5$, and put
$\eta=\tau_2\tau_5\beta$. The low-house mixed-inertia lemma
applies since $5<27/4$. After removing the fixed ternary
scalar, it gives
\[
                   (\sigma_2-1)(\sigma_5-1)\eta=0.
\]
Also $\eta\overline\eta=25$ and $\eta$ vanishes at every
five-place.

If $\tau_2\notin\mu_4$, the full dyadic average of $\eta$
is zero, because $\beta\in F(i)$. Averaging the displayed
identity in the dyadic variable makes $\eta$ fixed by every
fifth-inertia automorphism. Thus $\eta$ lies in a dyadic
cyclotomic field, unramified at five. Vanishing at all its
five-places gives $\eta\in5\mathcal O$. The integral
number $\eta/5$ has every conjugate of modulus one, so it
is a root of unity.

If $\tau_2\in\mu_4$, write
$\eta=a+ib$ with $a,b\in\mathcal O_F$. The mixed identity
for $i\mapsto-i$ makes $b$ rational, hence integral.
The two dyadic conjugates vanish at every five-place, so
$5\mid b$. Comparing their norm equations gives
$b(a-\overline a)=0$. If $b\ne0$, then $a$ is real and
$25=a^2+b^2$ at every embedding. Consequently $b=\pm5$
and $a=0$. If $b=0$, then $\eta\in F$; its norm equation
and the unique prime above five again give $\eta/5$ a unit
of absolute value one at every embedding. In either case
$\eta$, and hence $\beta$, is five times a root of unity.
\end{proof}

\begin{lemma}[Odd cyclotomic integers of norm forty]
\label{paper:odd-norm-forty-classification}
Put $\alpha=(5+\sqrt{-15})/2$. If an algebraic integer
$\beta$ in $L=\mathbb Q(\zeta_{3^a5^b})$, with $a,b\ge1$,
satisfies $\beta\overline\beta=40$, then $\beta$ is a
root-of-unity multiple of one of
\[
              2\alpha,\quad2\overline\alpha,\quad
              \alpha^3/5,\quad\overline\alpha^{\,3}/5.
\]
Moreover none of these root multiples has a positive
rational root word of mass at most ten, allowing arbitrary
higher primary orders at two, three and five.
\end{lemma}
\begin{proof}
As in the norm-ten calculation, $L$ has two primes over
two, exchanged by complex conjugation. The two valuations
of $\beta$ therefore form one of
$(0,3),(1,2),(2,1),(3,0)$. There is a unique prime over
five: its residue degree is $2\cdot3^{a-1}$ and its
ramification index is $4\cdot5^{b-1}$. Conjugation fixes
it, so the valuation there is forced to be half the
valuation of $40$. There are no other prime factors.

The algebraic integers $2\alpha$ and
\[
                 \alpha^3/5=(-5+3\sqrt{-15})/2
\]
and their conjugates realize precisely those four ideal
patterns. The norm is forty in each case. Thus the ratio
of $\beta$ to one of them is a unit $u$ with
$u\overline u=1$. Every conjugate of $u$ has modulus one,
so $u$ is a root of unity.

Here is an exact positive-word certificate. Put
$z=\zeta_{30}$, take
\[
 \sqrt{-15}=(1+2z^{10})(1+2z^6+2z^{24}),
\]
and define rational linear functionals $\lambda_{12}$ and
$\lambda_{13}$ on the power basis $1,z,\ldots,z^7$ by
\[
 \lambda_{12}=(1,-1,1,0,1,0,0,0),\qquad
 \lambda_{13}=(-1,-1,1,0,0,0,0,1).
\]
For $j=0,\ldots,14$, their values on $z^j$ are
\[
\begin{array}{c|rrrrrrrrrrrrrrr}
j&0&1&2&3&4&5&6&7&8&9&10&11&12&13&14\\\hline
\lambda_{12}&1&-1&1&0&1&0&0&0&1&0&-1&1&-1&1&-1\\
\lambda_{13}&-1&-1&1&0&0&0&0&1&1&-1&1&1&0&1&-1.
\end{array}
\]
The remaining fifteen values are the negatives, since
$z^{j+15}=-z^j$. In particular both functionals are at
most one on every root in $\mu_{30}$. Direct reduction
modulo $\Phi_{30}$ gives
\[
       \lambda_{12}(2\alpha)=12,\qquad
       \lambda_{13}(\alpha^3/5)=13.
\]

Extend each functional to a larger cyclotomic field by
normalized relative trace. A root outside $\mu_{30}$
whose order uses only two, three and five has relative
trace zero, because at least one higher primary component
has zero relative trace. Thus the bound one remains valid
on every allowed root. Composing with a Galois conjugation
and multiplying the root word by the inverse of its common
root phase gives the corresponding bound for every listed
root multiple. A positive word of mass at most ten could
not have functional value twelve or thirteen. This proves
the assertion.
\end{proof}

\begin{theorem}[Norm forty is impossible at the compulsory quintic endpoint]
\label{paper:no-thirty-norm-forty}
In the no-thirty branch, let $Y\in P_5$ be compulsory
of degree twenty and let $W$ be simple of degree ten.
Then
\[
                         |\chi_Y(W)|^2\ne40.
\]
\end{theorem}
\begin{proof}
Suppose $\beta=\chi_Y(W)$ has norm forty. Its house
is $\sqrt{40}<27/4$, so Lemma~\ref{paper:low-house-mixed-inertia}
applies.

First suppose the entry field is unramified at three.
The entry conductor then puts $\beta$ in
$\mathbb Q(i,\zeta_{25})$. Write the primary parts of
$\theta_W$ as $\tau_2,\tau_3,\tau_5$, enlarge the fifth
field to $F=\mathbb Q(\zeta_{5^c})$ as needed, and put
$\eta=\tau_2\tau_5\beta$. It satisfies
\[
       (\sigma_2-1)(\sigma_5-1)\eta=0,
       \qquad \eta\overline\eta=40.
\]
Principality of $Y$ makes $\eta$ vanish at every five-place.

If $\tau_2\notin\mu_4$, its dyadic average is zero,
including after multiplication by either coefficient of
$\beta\in F+iF$. Averaging the mixed identity therefore
makes $\eta$ fifth-invariant, so it lies in a two-primary
cyclotomic field. That field is unramified at five;
vanishing at all five-places gives $\eta\in5\mathcal O$.
The algebraic integer $\eta/5$ would have rational norm
$40/25=8/5$, impossible.

If $\tau_2\in\mu_4$, write $\eta=a+ib$ with
$a,b\in\mathcal O_F$. The mixed identity makes
$b\in\mathbb Z$, and vanishing at every five-place
gives $5\mid b$. Comparing the norm with its dyadic
conjugate gives $b(a-\overline a)=0$.
If $b=0$, the inertness of two in $F$ forces the
two-valuation of $\eta\overline\eta$ to be even,
contrary to $v_2(40)=3$. Otherwise $a$ is real and
$a^2+b^2=40$. Thus $b=\pm5$ and $a^2=15$.
This would put $\mathbb Q(\sqrt{15})$ in a field
unramified at three, another contradiction.

Hence $\mathbb Q(\beta)$ is ramified at three. The
compulsory endpoint field contains $\mathbb Q(\zeta_5)$,
so the no-thirty field dichotomy makes $K_Y$ unramified
at two. Consequently $\beta$ belongs to an odd
cyclotomic field, and the preceding lemma applies.
The actual balancing equation for $YW$, divided by
$20\theta_W$ (recall $\theta_Y=1$), is a positive
rational root word for $\beta$ of mass ten. The lemma
excludes every possible norm-forty realization, completing
the contradiction.
\end{proof}

\begin{proposition}[A real norm-twenty-five value in the negative sixty profile]
\label{paper:negative-sixty-norm25-real}
Assume the no-thirty branch and the existence of a compulsory
degree-thirty-six witness. Let $Y$ be any compulsory degree-twenty
witness, and suppose a degree-ten simple $W$ has
\[
 WW^*=\mathbf1+Z_{15}+U_{24}+V_{60},
 \qquad (\theta_Z,\theta_U,\theta_V)=(-1,1,-1).
\]
If $|\chi_Y(W)|^2=25$, then
\[
                            \chi_Y(W)\in\{5,-5\}.
\]
In particular this holds whenever all three adjoint constituents
are globally Galois fixed. No condition on the Galois fixedness
of the adjoint is needed beyond the stated norm hypothesis.
\end{proposition}
\begin{proof}
The endpoint field theorem gives $K_Y=K:=\mathbb Q(\zeta_5)$,
and the norm-twenty-five boundary theorem writes
\[
                   \beta:=\chi_Y(W)=5\rho,
                               \qquad \rho\in\mu_{10}.
\]
Suppose $\beta$ is nonreal, so the fifth-primary part of $\rho$
is nontrivial.

Proposition~\ref{paper:negative-three-term-square} makes
$\theta_W$ either three-primary or an element of $\mu_4$.
It also gives four square constituents of degrees
$(20,20,30,30)$, in the two actual twist classes
\[
                       \theta_W^4\omega,
                       \theta_W^4\omega^{-1}.
\]
Each class contains one degree-twenty and one degree-thirty
simple.

For use at the odd primes, fix a common cyclotomic field and
let $\sigma_p$ square its $p$-primary roots and fix the other
primary roots, for $p=3,5$. Put
$\Delta=(\sigma_3-1)(\sigma_5-1)$. Balancing any simple $B$
against the twist-one $Y$ gives
\[
 20\theta_B\chi_Y(B)=\sum_C N_{YB}^C d_C\theta_C,
 \qquad
              20\Delta(\theta_B\chi_Y(B))\in8\mathcal O.
\tag{*}
\]
Indeed the mixed difference kills every summand inactive at
three or at five, while every remaining label has degree
divisible by eight by no-thirty separation.

If $\theta_W$ is a nontrivial three-primary root, then
\[
 \Delta(\theta_W\beta)
          =5(\theta_W^2-\theta_W)(\sigma_5\rho-\rho).
\]
Both root differences are units at every place above two:
reduction at two is injective on roots of odd order.
This contradicts \((*)\), which requires that displayed mixed
difference be divisible by two at every such place. We have
therefore shown the needed conclusion
\[
                              \theta_W^4=1.
\]

Write $W^2=A_1+A_2+A_3+A_4$ as in the negative-square theorem.
Now every $\theta_{A_j}$ is $\omega$ or $\omega^{-1}$.
Put $b_j=\chi_Y(A_j)\in\mathcal O_K$. Applying \((*)\) to
$A_j$ gives
\[
        20(\sigma_3\theta_{A_j}-\theta_{A_j})
                      (\sigma_5b_j-b_j)\in8\mathcal O.
\]
The first root difference is a two-unit. It follows that
$\sigma_5b_j-b_j$ is divisible by two at every two-place.
Character multiplicativity gives $\sum_j b_j=\beta^2$, and hence
\[
                         \sigma_5(\beta^2)-\beta^2\equiv0
                                                   \pmod{2}.
\]
But $\beta=5\rho$ with nontrivial fifth-primary part makes this
difference a two-unit: writing $\rho=\varepsilon\zeta_5^k$,
$\varepsilon=\pm1$ and $k\not\equiv0\pmod5$, it equals
\[
                        25\zeta_5^{2k}(\zeta_5^{2k}-1).
\]
This contradiction proves that $\rho=\pm1$ and $\beta=\pm5$.

For an all-fixed sixty profile, the norm-forty theorem and the
degree-thirty-six norm-ten exclusion leave exactly norm twenty-five,
so the specialization follows.
\end{proof}

This character-value conclusion imposes no additional twist identity
on $W$; its final congruence uses the negative sixty profile.

\subsection{Amplified crossings and the activity graph}

\begin{proposition}[An amplified ordinary degree-thirty-six crossing]
\label{paper:amplified-ordinary-36}
Let $(d_X,d_Y)=(36,20)$ be a compulsory pure-three/pure-five
crossing. Suppose the full dyadic block of $X$ is paired,
and either all its pairs are mixed or its whole principal
pairs have twist ratio $-1$. If $z=1+3\xi$, then
$\theta_{X,3}=1$.
\end{proposition}
\begin{proof}
Theorem~\ref{paper:mixed-inactive-four} supplies factor four. Since $L=180$
and $D/L^2=900$, the effective Fourier budget is $225$.
A faithful fifth-torus character has energy $45$. Write
$a=[K_X:\mathbb Q(\zeta_5)]$ and
$b=[K_Y:\mathbb Q(\zeta_5)]$. For actual fifth twist
orders $1,5,25$ at $Y$, the faithful-channel bounds are,
respectively,
\[
                         45ab\le45,\quad90,\quad225.
\]
In the last case $5\mid b$. A ternary twist of order nine
at $X$ gives $3\mid a$ and contradicts every row. A
ternary twist of order three supplies the trivial-torus
half-channel factor $4/9$, reducing the right sides to
$20,40,100$, again impossible (use $5\mid b$ in the last
case). These are all possible active ternary orders.
\end{proof}

\begin{proposition}[The amplified nonreal signed equality]
\label{paper:amplified-nonreal-36}
Let $(d_X,d_Y)=(36,20)$ be a compulsory pure-three/pure-five
crossing. Suppose the full dyadic block of $X$ is paired, with
all pairs mixed or all whole principal pairs of twist ratio $-1$.
Suppose $z$ is signed tame and nonreal. Then
\[
 \theta_X=1,\qquad\operatorname{ord}(\theta_Y)=5,
 \qquad K_X=K_Y=\mathbb Q(\zeta_{15}).
\]
The two channel bounds in the proof are both equalities.
These conclusions do not exclude the crossing.
\end{proposition}
\begin{proof}
Here $K=\mathbb Q(\zeta_{15})$ has group $C_2\times C_4$.
The two identity coefficient positions are nonopposite.
For either faithful fifth character $\chi_5$, Gauss
summation gives
\[
 |\widehat z(1,\chi_5)|^2=90,
 \qquad |\widehat z(\operatorname{sgn}_3,\chi_5)|^2=30.
\]
The ternary sum has coefficients $2,2,-1,-1$ on the
four fifth positions; its faithful transform has squared
modulus $9\cdot2\cdot5=90$. The ternary difference has
two coefficients of magnitude $\sqrt3$ in nonopposite
positions and gives $3\cdot2\cdot5=30$.

Put $a=[K_X:K]$, $b=[K_Y:K]$. Including the dyadic gain,
the first channel gives
\[
 360ab\le
 \begin{cases}
 180,&\operatorname{ord}(\theta_Y)=1,\\
 360,&\operatorname{ord}(\theta_Y)=5,\\
 900,&\operatorname{ord}(\theta_Y)=25.
 \end{cases}
\]
A ternary twist of order nine at $X$ gives $3\mid a$;
a fifth twist of order twenty-five gives $5\mid b$.
These divisibilities exclude ternary order nine. At ternary
order three, the first torus character is trivial, so
the factor $4/9$ changes the right sides to $80,160,400$,
also impossible. Thus $\theta_X=1$.

The displayed bounds now exclude fifth orders one and
twenty-five. Fifth order five forces $ab=1$, proving
the full-field assertions. The first channel saturates
$4\cdot90=900(2/5)$. The second has both inputs inactive
at three and is faithful at five, so it saturates
$4\cdot30=900(1/3)(2/5)$.
Both are nontrivial Fourier channels and have zero vacuum
coordinate; equality alone gives no contradiction.
\end{proof}


\begin{proposition}[Mixed ramification at codegree thirty-six]
\label{paper:codegree-thirty-six-field}
Suppose that an ordinary perfect fusion-simple integral modular category
has dimension $2^6 3^6 5^4$. If a simple $Z$ has degree $900$ and is
active at both $2$ and $3$, then its full character field satisfies
\[
 \operatorname{cond}K_Z\mid12\qquad\hbox{or}\qquad
 \operatorname{cond}K_Z\mid9.
\]
Equivalently, its full field cannot simultaneously ramify at two and
have conductor exponent two at three.
\end{proposition}
\begin{proof}
Here $f_Z=36$, so the exact conductor theorem gives
$K_Z\subseteq\mathbb Q(\zeta_{36})$. Write $e=[K_Z:\mathbb Q]$.
Principal-three orthogonality on the pure-three twist fiber gives
\[
 1+\sum_{X\ne\mathbf1,\ X\in P_3,\ \theta_{X,3'}=1}
             d_X\chi_Z(X)=0.
\]
Group this sum into full Galois orbits, and reduce the resulting
integer identity modulo five. Some five-free nonunit orbit has a
nonzero absolute entry trace. Since $Z$ is two-active, a row of
two-height three cannot meet it. Complementary principal degrees
and the principal height bound therefore make the degree of this
chosen $X$ either $12$ or $36$.

Put $F=\mathbb Q(\chi_Z(X))$ and
$T=\operatorname{Tr}_{F/\mathbb Q}\chi_Z(X)\ne0$.
Pair $X$ with a three-free principal-five witness, of degree $20$
or $40$. This crossing has common weight four. The full-tame-orbit
corollary of the two-positive-residue theorem shows that five-inertia
has order at least four on $X$. Since $K_Z$ and hence $F$ are
unramified at five, $[K_X:F]\ge4$. The contribution of this orbit
to the exact trace norm is thus at least
\[
             4[K_Z:F]T^2.
\]
If $d_X=36$, principality at three and symmetry imply that
$\chi_Z(X)=(36/900)\chi_X(Z)$ vanishes at every three-place:
$36/900=1/25$ is a three-unit. Thus $3\mid T$, and the orbit alone
costs at least $36$, contradicting the positive vacuum contribution.
Consequently $d_X=12$. The same strict budget gives $|T|\le2$.
Every integer of a subfield of $\mathbb Q(\zeta_{36})$ that is wild
at three has absolute trace divisible by three. Indeed, the full
cyclotomic trace is divisible by three and the relative index is
prime to three. Hence $F$ is tame at three and lies in
$\mathbb Q(\zeta_{12})$.

Suppose $K_Z$ ramifies at two and is wild at three. Its degree is
either $12$, in which case $K_Z=\mathbb Q(\zeta_{36})$, or $6$.
For degree $12$, a tame entry field $F$ is rational,
$\mathbb Q(\zeta_3)$, $\mathbb Q(i)$, $\mathbb Q(\sqrt3)$, or
$\mathbb Q(\zeta_{12})$. The last three fields have even integral
absolute traces. In every case the chosen orbit costs at least
$24$. The vacuum costs $12$, already exhausting the codegree;
the trace-dimension identity necessarily has further nonunit
orbits, a contradiction.

In degree $6$, the two possible fields with this ramification are
the real subfield of $\mathbb Q(\zeta_{36})$ and
$\mathbb Q(i)\mathbb Q(\zeta_9)^+$. Neither contains
$\mathbb Q(\zeta_3)$. Its tame subfields are rational or the
quadratic field $\mathbb Q(\sqrt3)$ or $\mathbb Q(i)$,
respectively. Their integral absolute traces are even in the
quadratic case. The chosen orbit therefore costs at least $24$,
and the vacuum costs $6$.

Every additional nonzero-trace orbit costs at least $6$ in this
degree-six case. To check this, if its entry field is wild at three,
its nonzero trace is divisible by three and its degree is at most
six; if tame, the preceding rational/quadratic description gives
the same bound. The trace-dimension identity modulo two and modulo
three supplies two further distinct nonzero-trace orbits on
two-free and three-free degrees. Neither is the selected five-free
orbit, since nonunit prime-power degrees are impossible. Thus the
trace norm is at least $24+6+6+6>36$, the final contradiction.
\end{proof}

\begin{theorem}[Complete prime activity]\label{paper:modular:activity}
Every pair among $2,3,5$ divides the order of an actual simple twist.
Consequently the complementary witnesses can be chosen with
\[
                 (d_X,d_Y)=(12,20)\quad\hbox{or}\quad(36,20).
\]
If no actual simple twist has order divisible by $30$, both witnesses
have twist one and their crossing is $1+3\xi$.
\end{theorem}
\begin{proof}
The dyadic sextic argument gives an edge incident to two.
The quartic principal witness and the two-residue word classification
force an edge between three and five
(Theorem~\ref{paper:three-five-mixing}). It remains to rule out the two
connected graphs missing one dyadic edge. If the $2$--$3$ edge is
absent, the local dyadic column forces a four-label piece of dimensions
$(20,20,540,540)$. The quintic packet obstruction excludes this piece
(Theorem~\ref{paper:no-six-excluded}). If the $2$--$5$ edge is absent,
a degree-twelve column forces an active ternary packet with odd squared
norm at most $27$. Its integral trace and twist moment are incompatible
(Theorem~\ref{paper:no-ten-excluded}). These packet arguments use local invariant coordinate spaces; they do
not assume that a residue block is a fusion subcategory.

An edge at each dyadic neighbor forces both principal witnesses to
have dyadic height two, giving the two degree pairs. In the absence
of a twist divisible by thirty, separation on the exact complementary
twist fibers makes both witnesses inactive at their remaining primes
(Proposition~\ref{paper:no-thirty-inactive-witnesses}). A signed crossing
would force triple activity, so only the ordinary word remains.
\end{proof}

Both missing-edge graphs are excluded in full. If no simple
twist is divisible by thirty, Proposition~\ref{paper:no-thirty-inactive-witnesses}
makes every compulsory height-two principal witness inactive,
and the crossing is ordinary. If a thirty-active simple exists,
ordinary or signed tame crossings remain. In both cases
Theorem~\ref{paper:modular:no-degree-six} excludes all degree-six
simples. A thirty-active label must meet an unramified
principal-two witness of degree fifteen or forty-five.
The paired-block bounds address additional local
configurations but do not exhaust either case.


\section{Simple sources with simple center}
\label{paper:source}

We now impose the additional hypothesis
\begin{equation}\label{paper:source:assumption}
 \mathcal M\simeq\mathcal Z(\mathcal C),\qquad
 \operatorname{FPdim}(\mathcal C)=n=5400,
 \qquad \mathcal C\text{ integral and fusion-simple},
\end{equation}
where $\mathcal M$ is integral and fusion-simple of dimension $D=n^2$.
Both categories have their positive spherical structures.  The results
of this section constrain this center case; the existence of such a
source for an arbitrary modular category is not assumed elsewhere.

Write $R=K_0(\mathcal C)$, $r=\operatorname{rank}R$, and
$\tau(x)=[\mathbf1]x$.  The distinguished basis of $R$ is
orthonormal for $\langle x,y\rangle=\tau(xy^*)$.
Let $F:\mathcal M\to\mathcal C$ be the forgetful functor and $I$ its
right adjoint.  The canonical algebra and its length are
\[
 A=I(\mathbf1)=\bigoplus_Zm_ZZ,\qquad L=\sum_Zm_Z.
\]
Canonical multiplicities, dimensions of central simples, and dimensions
of source simples will be kept distinct.  We write $m_E$, $f_E$, and
$K_E$ for the degree, formal codegree, and full character field of an
irreducible representation $E$ of $R_{\mathbb C}$.

\subsection{The canonical character correspondence}

\begin{proposition}[Perfection, absence of braidings, and small generators]
\label{paper:source:source-structure}
The category $\mathcal C$ is perfect, $\mathcal C_{\ad}=\mathcal C$,
and its universal grading group is trivial. It admits no braiding.
Every nonunit simple of $\mathcal M$ forgets to an object with a
nonunit constituent. In particular, $\mathcal C$ has a simple
tensor generator of dimension at most twenty. If $\mathcal M$
has a degree-twelve, respectively degree-ten, simple, this bound
improves to twelve, respectively ten.
\end{proposition}

\begin{proof}
A pointed simple fusion category has prime dimension: its fusion
subcategories are indexed by subgroups of its grading group,
and a group of composite order has a proper nontrivial subgroup.
Thus $\mathcal C$ is not pointed. Simplicity now gives
$\mathcal C_{\pt}=\Vec$ and $\mathcal C_{\ad}=\mathcal C$.
The latter equality makes the universal grading trivial.

A braiding on $\mathcal C$ would give the fully faithful tensor
section $X\mapsto(X,c_{X,-})$ of $F$. Its image in
$\mathcal Z(\mathcal C)$ has dimension $5400$, contrary to
fusion-simplicity of $\mathcal M$.

The objects of $\mathcal M$ that forget to sums of units form
the fusion subcategory $\ker F$. It is proper, since $F$ is
dominant and $\mathcal C\ne\Vec$. Simplicity of $\mathcal M$
therefore gives $\ker F=\Vec$. A nonunit central simple of
dimension $a$ has a nonunit constituent of $F(Z)$ of dimension
at most $a$. Any such constituent generates $\mathcal C$ by
simplicity. The compulsory degree-twenty witness in
Corollary~\ref{paper:two-neighbors} proves the unconditional bound;
the other two bounds follow in the same way.
\end{proof}

\begin{proposition}[The exact character-theoretic bridge]
\label{paper:source:canonical-character-bridge}
For each irreducible representation $E$ of
$K_0(\mathcal C)\otimes_{\mathbb Z}\mathbb C$, let $m_E=\dim E$,
let $f_E$ be its formal codegree, and let
$K_E=\mathbb Q(\chi_E(X):X\in\operatorname{Irr}(\mathcal C))$.
There is a bijection between these representations and the simple
summands $Z_E$ of $A$, with
\[
 m_{Z_E}=m_E,\qquad d_{Z_E}=\frac n{f_E},\qquad
 \theta_{Z_E}=1,\qquad K_{Z_E}=K_E.
\]
In particular $f_E$ is a positive integer divisor of $n$, and
\begin{equation}\label{paper:source:canonical-character-budget}
 \sum_E\frac{m_E}{f_E}=1,\qquad
 \operatorname{rank}(\mathcal C)=\sum_E m_E^2.
\end{equation}
The dimension representation corresponds to the unit and has
$(m_E,f_E)=(1,n)$. Moreover
\begin{equation}\label{paper:source:canonical-source-conductor}
 \mathfrak D_{K_{Z_E}}\mid f_E\mathcal O_{K_{Z_E}},
 \qquad \operatorname{cond}(K_{Z_E})\mid 2f_E.
\end{equation}
\end{proposition}

\begin{proof}
The multiplicity, dimension and rank assertions are Ostrik's
induction theorem \cite[Theorem~2.13]{OstrikRankThree}; the trivial
twist assertion is \cite[Theorem~2.5]{OstrikRankThree}.
Since $F$ preserves dimension, integrality of $\mathcal C$ makes
$d_{Z_E}$ integral. Thus $f_E=n/d_{Z_E}$ is rational, and its
algebraic integrality makes it an integer.

Here is the field assertion, including its normalization. The
central action of $K_0(\mathcal M)$ on $E$ is
\[
 \rho_E(F(Y))=\chi_{Z_E}(Y)\operatorname{id}_E.
\]
Taking traces gives
$\chi_E(F(Y))=m_E\chi_{Z_E}(Y)$ and hence
$K_{Z_E}\subseteq K_E$.
For the reverse inclusion, the induction formula gives
\[
 [F(I(X))]=\sum_{U\in\operatorname{Irr}(\mathcal C)}[U][X][U^*].
\]
For a symmetric semisimple algebra with symmetrizing trace
$\tau=\sum_E\chi_E/f_E$, the Casimir averaging map on the right
acts on $E$ by $f_E\chi_E(X)\operatorname{id}_E$.
Indeed, on the matrix block of $E$, the dual of the matrix unit
$e_{ij}$ is $f_Ee_{ji}$, and
$\sum_{i,j}e_{ij}xe_{ji}=\operatorname{Tr}(x)\operatorname{id}$.
Consequently
\[
                  \chi_{Z_E}(I(X))=f_E\chi_E(X).
\]
As $f_E$ is rational and nonzero, $K_E\subseteq K_{Z_E}$.
Finally apply Theorem~\ref{paper:family:2.5} and
Corollary~\ref{paper:family:2.6} to the representation $E$ itself.
\end{proof}

The character correspondence gives two different dimension budgets:
$L=\sum_E m_E$ and $r=\sum_E m_E^2$.  Consequently
\begin{equation}\label{paper:source:nonlinear-defect}
 r-L=\sum_E m_E(m_E-1)
\end{equation}
is a nonnegative even integer.  This identity will turn a canonical
length bound into a restriction on all possible complex algebra shapes.

\begin{proposition}[The canonical palette and the source type]
\label{paper:source:palette-and-type}
For every non-Perron representation $E$, put $a_E=n/f_E$.  Then
\begin{equation}\label{paper:source:canonical-palette}
 a_E\in\mathcal A
 :=\{a:a\mid5400,\ 10\le a\le900,
                  \ a\text{ has at least two distinct prime factors}\},
\end{equation}
and
\begin{equation}\label{paper:source:codegree-palette}
 f_E\in\{5400/a:a\in\mathcal A\},\qquad
 K_E\subseteq\mathbb Q(\zeta_{10800}).
\end{equation}
A full Galois orbit of degree $e$, representation degree $m$, and
canonical degree $a$ contributes $ema$ to $\dim A$ and $em^2$ to $r$.
For $n_d=|\{X\in\operatorname{Irr}(\mathcal C):d_X=d\}|$, one has
\begin{equation}\label{paper:source:dimension-type}
 n_1=1,\quad n_2=0,\quad
 \sum_{d=3}^{72}d^2n_d=5399,\quad
 r=1+\sum_{d=3}^{72}n_d,
 \quad\sum_{d=3}^{20}n_d\ge1.
\end{equation}
At least three nonunit dimensions occur, their gcd is one, and an
odd-dimensional nonunit simple is selfdual.  The source estimates of
\cite[Section~30]{BDPfamily} give
$r\le462$ in the commutative case and $r\le490$ otherwise, and
\begin{equation}\label{paper:source:dimension-three-density}
 n_3\ge2\quad\Longrightarrow\quad
 n_3\le2\bigl(\min(4n_4,5n_5)+6n_6+9n_9\bigr).
\end{equation}
In the commutative case at most $359$ nonunit labels have odd dimension.
\end{proposition}
\begin{proof}
The canonical character correspondence gives $a_E\mid5400$.
The central degree restrictions exclude prime powers and degree six
(Lemma~\ref{paper:family:2.4} and Theorem~\ref{paper:modular:no-degree-six}).
If $h=5400/a_E$,
the central formal codegree is $h^2$.  Lemma~\ref{paper:lem:small-ordinary-codegrees},
Proposition~\ref{paper:primepower-trace}, and Theorem~\ref{paper:dyadic-sixteen}
exclude $h=1,2,3,4,5$, so $a_E\le900$; the smallest remaining
divisor is ten.  The source conductor bound gives the displayed
cyclotomic containment, and the orbit costs follow from the two
character budgets.

Fusion-simplicity and nonpointedness give $n_1=1$.  A dimension-two
simple would, by \cite[Theorem~3.4]{DNVSmall}, force either a nonunit
invertible or a proper fusion subcategory of dimension $60$.
Thus $n_2=0$.  The starting estimate $L\ge10$ from
\cite[Section~30]{BDPfamily} implies $r\ge10$.  For a source label
of maximal dimension $d$, the other $r-2$ nonunit labels contribute
at least $8\cdot9$ to the squared-dimension budget; hence
$d^2\le5399-72<73^2$.  The small generator gives the occurrence
of a dimension at most twenty.  The assertion about three distinct
dimensions is \cite[Proposition~3.7(v)]{DNVSmall}.  Their gcd $g$
satisfies $g^2\mid5399$; since $5399$ is prime, $g=1$.
Dual pairs contribute an even squared mass, whereas the nonunit mass
is odd, proving the selfdual assertion.  The remaining estimates are
the cited source bounds, whose hypotheses follow from
Proposition~\ref{paper:source:source-structure}.
\end{proof}

Only the canonical dimensions $a_E$ have the displayed divisor
restriction.  No condition $d_X\mid5400$ is imposed on source dimensions.
A commutative source Grothendieck ring also need not admit a braiding.

\begin{proposition}[Invariant multiplicities and exact block masses]
\label{paper:source:canonical-invariant-block-masses}
Let $m=(m_Z)_Z$, let $S$ be the unnormalized modular matrix of
$\mathcal M$, and put $s=S/n$. Then
\[
 m_0=1,\qquad Sm=nm,\qquad Tm=m,
 \qquad m_{gZ}=m_Z
\]
for every Galois permutation $g$. In the central-charge-one
normalization, $m$ is fixed by every local modular group.

If $L$ is a union of full residue blocks at $p\mid n$, then
\begin{equation}\label{paper:source:canonical-block-mass-divisibility}
 p^{v_p(n)}\ \bigm|\ M_L,
 \qquad M_L:=\sum_{Z\in L}m_Zd_Z.
\end{equation}
In particular the moduli at $2,3,5$ are $8,27,25$.
If $L$ is a union of full blocks at every prime in $J$, its mass
is divisible by $\prod_{p\in J}p^{v_p(n)}$.
For every $p\mid n$, the principal full $p$-block contains a
nonunit canonical summand $Z$ with $p\nmid m_Zd_Z$.
\end{proposition}

\begin{proof}
Connectedness gives $m_0=1$. The identities for $S$ and $T$ are
\cite[Theorem~2.5 and Proposition~2.8]{OstrikRankThree}.
The normalized modular representation of a Drinfeld center has
central charge one, so these identities make $m$ fixed by its
global group and hence by each local factor.
The multiplicity and source codegree in
Proposition~\ref{paper:source:canonical-character-bridge} are preserved by
Galois conjugation of $E$; the field equality identifies this with
the modular Galois permutation. This proves the orbit assertion.

For the divisibility, let $e_L$ be the locally integral idempotent
of the full-block calculus in \eqref{paper:family:eq:local}. The fusion-ring trace is
\[
 \tau_{\mathcal M}(e_L[A])
   =\sum_{Z\in L}\frac{\chi_Z(A)}{f_Z}.
\]
The identity $Sm=nm$ gives $\chi_Z(A)=nm_Z/d_Z$, and
$f_Z=n^2/d_Z^2$. Thus the displayed trace is $M_L/n$.
It is integral at $p$ by \eqref{paper:family:eq:local}, while
$M_L$ is an integer. This proves
\eqref{paper:source:canonical-block-mass-divisibility}. The assertion for
several primes follows by coprimality. In the principal block the
unit contributes one; reducing its mass modulo $p$ proves the
last assertion.
\end{proof}

\begin{proposition}[Forced canonical orbits of restricted fields]
\label{paper:source:forced-canonical-tame-orbits}
For each $p\mid n$, the principal full $p$-block contains a
nonunit canonical Galois orbit $\mathcal O$ such that
\[
 p\nmid e\,m\,d,\qquad
 e=|\mathcal O|=[K_Z:\mathbb Q],\quad
 m=m_Z,\quad d=d_Z\quad(Z\in\mathcal O).
\]
Its full character field is tamely ramified at $p$.
At $p=2$ one can choose this orbit with
\[
 d\in\{15,45,75,135,225,675\},\qquad m\text{ odd},
 \qquad e\in\{1,3\}.
\]
Every object in this orbit is self-dual. If $e=3$, then
$d=15$ or $75$ and $K_Z=\mathbb Q(\zeta_9)^+$; otherwise
$K_Z=\mathbb Q$.
At $p=3$ and $p=5$, respectively, the possible dimensions
in the asserted orbit are
\[
 \begin{array}{c|l}
 p&d\\\hline
 3&10,20,40,50,100,200\\
 5&12,18,24,36,54,72,108,216.
 \end{array}
\]
\end{proposition}

\begin{proof}
The principal full block is Galois-stable, by
Lemma~\ref{paper:family:lem:principal-all-places} and Galois covariance.
Its mass is a sum of contributions $e m d$ from canonical
Galois orbits. The unit contributes one, whereas the total
is divisible by $p$ by
Proposition~\ref{paper:source:canonical-invariant-block-masses}.
Thus some nonunit orbit has $p\nmid e m d$.
The character field is abelian, and its wild inertia group
at $p$ is a $p$-group. Since its order divides $e$, this
wild inertia group is trivial.

The central degree palette gives the displayed lists.
For an odd $d$ in that palette, the source conductor bound
places $K_Z$ in $\mathbb Q(\zeta_{2n/d})$.
For $d=45,135,225,675$ these cyclotomic Galois groups have
power-of-two order, so odd $e$ forces $e=1$.
For $d=15,75$ their odd part has order three, and their unique
cubic subfield is $\mathbb Q(\zeta_9)^+$.
These fields are totally real. Complex conjugation of a
modular character corresponds to duality of its label,
so the labels are self-dual.
\end{proof}

\begin{proposition}[The integral Gram factorization]
\label{paper:source:canonical-gram-factorization}
Put $b_{ZX}=[F(Z):X]=[I(X):Z]$, and let $B=(b_{ZX})$ have
central simples as rows and source simples as columns, with the
unit column first. Let $N_A$ denote the matrix with
$(Z,W)$-entry $[A\otimes Z:W]$. Then
\begin{equation}\label{paper:source:canonical-gram-identities}
 N_A=BB^{\mathsf t},\qquad
 B\,d_{\mathcal C}=d_{\mathcal M},\qquad
 B^{\mathsf t}d_{\mathcal M}=n\,d_{\mathcal C}.
\end{equation}
Its nonzero eigenvalues, indexed by the distinct irreducible
representations $E$, are $m_Ef_E$; all other eigenvalues are zero.
In particular its rank is the number of irreducible representations
of $K_0(\mathcal C)\otimes\mathbb C$, not in general the rank of
$\mathcal C$.

If $B'$ is obtained by removing the unit column, then
\begin{equation}\label{paper:source:canonical-reduced-gram}
 N_A-mm^{\mathsf t}=B'(B')^{\mathsf t}.
\end{equation}
This is a nonnegative integral Gram factorization, stronger than
positive semidefiniteness alone. In particular
\[
 \sum_Tm_TN_{ZZ^*}^T-m_Z^2
       =\sum_{X\ne\mathbf1}b_{ZX}^2\ge0.
\]
\end{proposition}

\begin{proof}
The projection formula for the central functor gives
$I(F(Z))\simeq A\otimes Z$. Taking the multiplicity of $W$
and using adjunction gives
$[A\otimes Z:W]=\sum_Xb_{ZX}b_{WX}$.
Dimensions of $F(Z)$ and $I(X)$ give the next two identities.
The modular characters diagonalize $N_A$, with eigenvalue
$\chi_Z(A)=nm_Z/d_Z$. By
Proposition~\ref{paper:source:canonical-character-bridge} this is zero
off $A$ and is $m_Ef_E$ at $Z_E$.
The unit column is $m$, giving
\eqref{paper:source:canonical-reduced-gram}; its diagonal is the last
identity, by Frobenius reciprocity.
\end{proof}

\subsection{The canonical algebra under the center hypothesis}
\label{paper:sec:canonical-structure}

Recall the canonical connected \emph{\'etale} algebra
$A=I(\mathbf1)=\bigoplus_Zm_ZZ$ of dimension $n=5400$
under the additional hypothesis \eqref{paper:source:assumption}.

\begin{proposition}[Every nonunit canonical summand generates the algebra]
\label{paper:canonical-algebra-generation}
Every unit-preserving subalgebra of $A$ is either $\mathbf1$ or
$A$. Consequently, for any simple $W\ne\mathbf1$ and any
nonzero morphism $j:W\to A$, the unit-preserving subalgebra
generated by $j(W)$ equals $A$.

More explicitly, let $\mu_k:W^{\otimes k}\to A$ be the iterated
product of $j$, with $\mu_0$ the unit. Set
\[
 B_k=\sum_{i=0}^k\operatorname{im}(\mu_i)\subseteq A.
\]
The sequence is strictly increasing until it reaches $A$, and
it reaches $A$ by $k\le\operatorname{length}(A)-1$.
\end{proposition}

\begin{proof}
Burciu's subalgebra theorem says that every unit-preserving
subalgebra of the canonical algebra in the center of a
pseudo-unitary fusion category is \emph{\'etale}
\cite[Theorem~1.1]{BurciuSubalgebras2025}.
Here ``unit-preserving'' is the meaning of ``unitary subalgebra''
in that theorem; no $C^*$-structure is being assumed.
The subalgebra is connected because it is a subobject of $A$.
The antiisomorphism between \emph{\'etale} subalgebras of $A$
and fusion subcategories of $\mathcal C$
\cite[Theorem~4.10]{DMNO}, together with fusion-simplicity of
$\mathcal C$, therefore leaves only $\mathbf1$ and $A$.

For completeness the generated subalgebra exists in the given
finite semisimple category. The images defining $B_k$ exist, and
$B_0\subseteq B_1\subseteq\cdots$ is a chain of subobjects of
the finite-length object $A$. If $B_{k+1}=B_k$, multiplication
by $j(W)$ preserves $B_k$. Associativity then gives
$B_{k+i}=B_k$ for every $i\ge0$ and
$B_kB_k\subseteq B_{2k}=B_k$. Thus $B_k$ is a unit-preserving
subalgebra. Any such subalgebra containing $j(W)$ contains all
the $B_i$, so the stabilized object is the generated subalgebra.
It is nontrivial and therefore equals $A$. Each strict inclusion
increases composition length, proving the stated bound.
\end{proof}

\begin{lemma}[Nonvanishing of braided symmetric powers]
\label{paper:canonical-symmetric-powers-nonzero}
Let $U$ be a nonzero subobject of a connected \emph{\'etale}
algebra $A$ in a braided fusion category. Every iterated
multiplication $U^{\otimes k}\to A$ is nonzero.
It factors through the quotient
\[
 Q_k(U)=\operatorname{coker}
 \left(\bigoplus_{i=1}^{k-1}U^{\otimes k}
 \xrightarrow{\ (c_i-\operatorname{id})_i\ }U^{\otimes k}\right),
\]
where $c_i$ braids the adjacent $i$th and $(i+1)$st factors.
In particular $Q_k(U)\ne0$ for every $k\ge1$.
No semisimplicity or finite-image assertion for braid-group
representations is required.
\end{lemma}

\begin{proof}
Let $J$ be the image of $A\otimes U\to A$. It is a nonzero
ideal: it contains $U$, and braided commutativity makes it
two-sided. The regular $A$-module is simple, since the category
of $A$-modules is semisimple and
$\operatorname{End}_A(A)=\operatorname{Hom}(\mathbf1,A)=\mathbb C$.
Hence $J=A$. If the $k$fold product on $U$ vanished,
associativity and braided commutativity would give $J^k=0$,
contradicting $J=A$ and the unit of $A$.
The factorization follows from the adjacent braid invariance
of iterated multiplication. Cokernels exist in the ambient
semisimple category, so a nonzero multiplication forces a
nonzero quotient.
\end{proof}

\begin{proposition}[Canonical capacity of higher braided symmetric powers]
\label{paper:canonical-higher-symmetric-capacity}
Let $W$ be simple with $m=m_W>0$, and use the quotient
$Q_k(W)$ from Lemma~\ref{paper:canonical-symmetric-powers-nonzero}.
Then
\begin{equation}\label{paper:canonical-higher-symmetric-bound}
 \dim\operatorname{Hom}_{\mathcal M}(Q_k(W),A)
 =\sum_Zm_Z[Q_k(W):Z]
 \ge\binom{m+k-1}{k}
 \qquad(k\ge1).
\end{equation}
The common $-1$ eigenspace of the adjacent braids on
$\operatorname{Hom}_{\mathcal M}(W^{\otimes k},A)$
has dimension at least $\binom mk$.
In particular, a canonical degree-ten summand of multiplicity
four requires $\dim\operatorname{Hom}(Q_3(W),A)\ge20$.
\end{proposition}

\begin{proof}
Put $V=\operatorname{Hom}_{\mathcal M}(W,A)$.
Canonical multiplication gives a linear map
\[
 \Phi_k:V^{\otimes k}\longrightarrow
              \operatorname{Hom}_{\mathcal M}(W^{\otimes k},A).
\]
Under adjunction this map tensors $k$ morphisms
$F(W)\to\mathbf1$. It is injective: split
$F(W)=\mathbf1^{\oplus m}\oplus R$ and restrict to
$(\mathbf1^{\oplus m})^{\otimes k}$.
Braided commutativity of $A$ and naturality identify
precomposition by each adjacent braid on the image with
the ordinary adjacent flip on $V^{\otimes k}$.
Consequently $\operatorname{Sym}^kV$ injects into the
common $+1$ eigenspace, which is
$\operatorname{Hom}(Q_k(W),A)$.
Similarly $\bigwedge^kV$ injects into the common $-1$
eigenspace. Their dimensions give the claims.
No decomposition or semisimplicity of the ambient
braid-group representation is used.
\end{proof}

\begin{lemma}[Blockwise canonical orthogonality and conditional prime avoidance]
\label{paper:canonical-prime-avoidance}
If $B$ is a full $p$-block and $X$ is $p$-active, then
\begin{equation}\label{paper:canonical-blockwise-orthogonality}
              \sum_{Y\in B}m_Y\chi_X(Y)=0.
\end{equation}
In particular, let $B=P_p$ be principal and let $s\ne p$
be a prime. At every place over $s$ there is a nonunit
$Y\in P_p$ with $m_Y>0$ such that
\[
 m_Y\chi_X(Y)\not\equiv0,\qquad
 s\nmid m_Y,\qquad v_s(d_Y)\le v_s(d_X).
\]
Thus if $s\nmid d_X$, a canonical $s$-free nonunit can be
chosen in $P_p$. Every such label has actual twist one.
\end{lemma}

\begin{proof}
The full-block coordinate projector $P_B$ commutes with the
local $p$-group by Theorem~\ref{paper:lw:block}.
Since $m$ is invariant, $P_Bm$ is locally invariant.
The global Weyl element normalizes the local $p$-group, so
$(S/n)P_Bm$ is locally invariant as well. A locally invariant
vector has zero coordinate at every $p$-active label.
Its $X$-coordinate is
$n^{-1}d_X\sum_{Y\in B}m_Y\chi_X(Y)$, proving
\eqref{paper:canonical-blockwise-orthogonality}.

For the principal block the unit contributes one.
Reducing the identity at a place over $s$ therefore gives
a nonunit term with nonzero residue. Both $m_Y$ and
$\chi_X(Y)$ are integral, so $s\nmid m_Y$.
LCM integrality gives
\[
 \chi_X(Y)=\frac{d_Y}{\gcd(d_X,d_Y)}\,\beta,
 \qquad\beta\text{ an algebraic integer}.
\]
A nonzero residue forces the stated height inequality.
Finally canonical support has actual twist one.
\end{proof}

\begin{remark}\label{paper:canonical-prime-avoidance-scope}
The foreign-prime-free condition on the active label in
Lemma~\ref{paper:canonical-prime-avoidance} is substantive.
The unconditional principal-block prime-avoidance theorem
weights its sum by $d_Y$; here the weights are $m_Y$.
Consequently this lemma does not, without additional input,
put the previously selected compulsory endpoints in $A$.
\end{remark}

\begin{corollary}[A unique symmetric-square constituent is canonical]
\label{paper:unique-symmetric-square-constituent}
Suppose $W\subset A$ is simple and the eigenvalue-one quotient
of $c_{W,W}$ on $W^{\otimes2}$ is one simple object $V$.
Then $V\subset A$, and the restriction of multiplication
$W^{\otimes2}\to A$ has image exactly one copy of $V$.
If $V\ne\mathbf1$ and $\mathcal C$ is fusion-simple, this copy of $V$ also generates
$A$ as an algebra.
\end{corollary}

\begin{proof}
The product is nonzero by
Lemma~\ref{paper:canonical-symmetric-powers-nonzero}, factors
through $V$, and is therefore injective on that simple quotient.
When $V\ne\mathbf1$, apply Proposition~\ref{paper:canonical-algebra-generation} to its image.
\end{proof}

\begin{remark}\label{paper:canonical-generation-scope}
These statements concern the actual multiplication of $A$.
They do not say that its simple support is fusion closed.
For the surviving positive-sixty degree-ten square profiles,
the unique eigenvalue-one constituent has degree $50$ or $40$;
if the degree-ten object lies in $A$, that constituent lies in
$A$ and generates it. This does not yet exclude either profile:
the higher braided symmetric quotients and their multiplication
images are not determined by the four-term square alone.
\end{remark}
\subsection{Canonical multiplication and short tensor squares}

\begin{lemma}[Canonical multiplicities under tensor products]
\label{paper:canonical-product-multiplicity}
Write $A=I(\mathbf1)=\bigoplus_Z m_ZZ$ in
$\mathcal M=\mathcal Z(\mathcal C)$. Then
\[
 m_Um_V\le\sum_ZN_{UV}^Zm_Z
\]
for any simples $U,V$. More generally, if $m_U>0$, every
positive tensor power of $U$ has a constituent in the support
of $A$.
\end{lemma}
\begin{proof}
Adjunction gives $m_U=[F(U):\mathbf1]$, and the forgetful
functor is tensor. Tensoring the unit summands of $F(U)$ and
$F(V)$ proves the inequality. Its iterated form proves the
last assertion.
\end{proof}

\begin{lemma}[Symmetric and alternating canonical multiplicities]
\label{paper:canonical-square-braiding}
Suppose $m_W=m>0$ and $W^2=\bigoplus_Z Z$ is multiplicity
free. Write $R_Z$ for the scalar braiding on $Z\subset W^2$.
Then $\theta_W=1$, and
\[
 \sum_{R_Z=1}m_Z\ge\binom{m+1}{2},\qquad
 \sum_{R_Z=-1}m_Z\ge\binom m2.
\]
In particular, absence of a $+1$ channel is impossible, and
absence of a $-1$ channel forces $m=1$.
\end{lemma}
\begin{proof}
Put $V=\operatorname{Hom}_{\mathcal M}(W,A)$. The canonical
commutative algebra multiplication induces
\[
 \Phi:V\otimes V\longrightarrow
              \operatorname{Hom}_{\mathcal M}(W^2,A).
\]
Under the adjunction
$\operatorname{Hom}_{\mathcal M}(W,A)
 =\operatorname{Hom}_{\mathcal C}(F(W),\mathbf1)$,
this map tensors two morphisms to the unit. It is injective:
choose a splitting $F(W)=\mathbf1^{\oplus m}\oplus R$ and
restrict to the tensor square of the displayed unit summand.
Commutativity and naturality of braiding give
$\Phi(v\otimes u)=\Phi(u\otimes v)\circ c_{W,W}$.
Thus $\operatorname{Sym}^2V$ and $\bigwedge^2V$ inject into
the $+1$ and $-1$ eigenspaces of precomposition by $c_{W,W}$.
Multiplicity freeness makes their ambient dimensions the two
displayed sums. Finally $\theta_A=\operatorname{id}_A$ gives
$\theta_W=1$.
\end{proof}

\begin{proposition}[Negative degree-ten profiles miss the canonical algebra]
\label{paper:negative-ten-not-canonical}
Suppose a degree-ten simple $W$ has a multiplicity-free adjoint
with three nonunit constituents of distinct degrees and
$\chi_W(W^*)=-5\theta_W^{-2}$. Then $m_W=0$.
This includes the negative $(15,24,60)$ profile and the
$(24,30,45)$ profile with $(\theta_{45},\theta_{30})=(-1,-1)$.
The assertion does not require the no-thirty hypothesis.
\end{proposition}
\begin{proof}
If $m_W>0$, then $\theta_W=1$. Proposition
\ref{paper:negative-three-term-square} gives four square
constituents of dimensions $20,20,30,30$, all with actual
twists $\omega$ or $\omega^{-1}$. None occurs in $A$,
contradicting Lemma~\ref{paper:canonical-product-multiplicity}.
\end{proof}

\begin{proposition}[Positive sixty profiles in the canonical algebra]
\label{paper:canonical-positive-sixty}
Suppose $m_W=m>0$, $d_W=10$, and
$WW^*=\mathbf1+Z_{15}+U_{24}+V_{60}$.
Then $\theta_Z=-1$, $\theta_U=\theta_V=1$.
The square is one of the two patterns
\[
\begin{array}{c|rrrr}
d&50&30&10&10\\
R&1&-1&\omega&\omega^{-1}\\\hline
d&40&40&10&10\\
R&1&-1&\xi_6&\xi_6^{-1}.
\end{array}
\]
If $B_+$ and $B_-$ denote the first two square constituents,
then
\[
 m_{B_+}\ge\binom{m+1}{2},\qquad
 m_{B_-}\ge\binom m2,
 \qquad m_U+m_V\ge m^2-1.
\]
If the already established bound
$\operatorname{rank}(\mathcal C)\le490$ is used, then $m\le4$.
\end{proposition}
\begin{proof}
The twist assertion follows from Theorem
\ref{paper:fifteen-twentyfour-sixty-twists} and Proposition
\ref{paper:negative-ten-not-canonical}. The square phases are
Lemma~\ref{paper:four-root-moments}, with $\theta_W=1$.
The square inequalities follow from
Lemma~\ref{paper:canonical-square-braiding}. The adjoint
inequality follows from Lemma~\ref{paper:canonical-product-multiplicity}
and $m_Z=0$.

The objects $W,W^*,B_+,B_-,U,V,\mathbf1$ are distinct;
$W\ne W^*$ by the short-adjoint self-duality exclusion.
Since $m_{W^*}=m$ and
$\operatorname{rank}(\mathcal C)=\sum_Tm_T^2$, one obtains
\[
 \operatorname{rank}(\mathcal C)\ge
 1+2m^2+\binom{m+1}{2}^{\!2}+\binom m2^{\!2}
       +\left\lceil\frac{(m^2-1)^2}{2}\right\rceil.
\]
At $m=5$ the right side is $664$ and it increases thereafter.
Thus $m\le4$.
\end{proof}

\begin{corollary}[Two forty-five square rows have multiplicity one]
\label{paper:canonical-fortyfive-plus-only}
For an $A$-supported degree-ten object in either row of
Proposition~\ref{paper:fortyfive-square-table} with phases
$(1,1,\omega,\omega^{-1})$, its multiplicity in $A$ is one.
\end{corollary}
\begin{proof}
There is no $-1$ braiding channel; apply
Lemma~\ref{paper:canonical-square-braiding}.
\end{proof}

\begin{corollary}[Supported complementary endpoints have an ordinary crossing]
\label{paper:canonical-endpoints-not-signed}
If both endpoints of a complementary mass-four crossing
occur in $A$, the crossing is ordinary.
\end{corollary}
\begin{proof}
Both endpoint twists are one. In the signed alternative of
Proposition~\ref{paper:modular:weight-four-alternatives}, every actual
constituent twist is $-\eta_t\zeta_5^t$, hence has nontrivial
dyadic component. Thus the tensor product has no constituent
in $A$, contrary to Lemma~\ref{paper:canonical-product-multiplicity}.
\end{proof}

\begin{remark}[Scope]
None of these statements forces a compulsory endpoint to lie
in $A$, and none asserts that the support of $A$ is fusion
closed. Under the standard correspondence $E\leftrightarrow
Z_E\subset A$, the multiplicity $m_{Z_E}=\dim E$ and central
dimension $d_{Z_E}=5400/f_E$ are not the simple dimensions of
$\mathcal C$. The degree-ten conclusions concern formal
codegree $f_E=540$, without excluding all representations of
that formal codegree.
\end{remark}

\subsection{Exact forgetful types and orbit costs}

\begin{lemma}[A content divisibility test for a canonical constituent]
\label{paper:canonical-content-test}
Let $W=Z_E$ occur in $I(\mathbf1)$ with multiplicity
$m=\dim E$, and suppose
$F(W)=m\mathbf1+gP$ for a genuine object $P$ and a positive
integer $g$. Then
\[
                  \frac{\chi_W(W)-m}{g}
\]
is an algebraic integer. In particular, if $\chi_W(W)$ is
an integer, then $g$ divides $\chi_W(W)-m$.
\end{lemma}
\begin{proof}
The image of the forgetful map is central in the fusion ring.
The canonical correspondence identifies its scalar action on
$E$ by
$E(F(Y))=\chi_W(Y)\operatorname{id}_E$.
Consequently $E(P)=((\chi_W(W)-m)/g)\operatorname{id}_E$.
Every eigenvalue of a based-ring element with integer
coefficients is an algebraic integer: it occurs among the
eigenvalues of its integral regular-representation matrix.
This proves the assertion.
\end{proof}

\begin{proposition}[Exact Gram possibilities for a positive seven profile]
\label{paper:positive-seven-gram}
Suppose $m_W=m>0$, $d_W=10$, and
\[
 WW^*=\mathbf1+Z+U+V,
 \qquad \theta_Z=-1,\quad\theta_U=\theta_V=1,
 \qquad \chi_W(W)=7,
\]
where the three nonunit constituents have distinct degrees.
Assume its square has the positive-seventy moment pattern,
with real-braiding constituents $B_+,B_-$, and that these
two constituents are distinct and have degrees different from
$10,d_U,d_V$.
This applies both to the positive $(15,24,60)$ profile and
to the globally confined $(15,12,72)$ profile.

Then $m\le4$. Write
\[
 F(W)=m\mathbf1+P,\qquad
 q=\langle P,P\rangle,\qquad t=\langle P,P^*\rangle.
\]
One has the exact identities
\[
 m_U+m_V=m^2-1+q,\qquad
 m_{B_+}+m_{B_-}=m^2+t.
\]
Moreover, for some integer $0\le x\le t$,
\[
 m_{B_+}=\binom{m+1}{2}+x,
 \qquad
 m_{B_-}=\binom m2+t-x.
\]
The complete remaining possibilities for $(q,t)$ are
\[
\begin{array}{c|c}
m& (q,t)\\\hline
1&(1,0\text{ or }1),\ (2,0,1,2),\ (3,0,1,2,3),
       \ (5,0,1,4,5),\ (9,0\text{ or }9)\\
2,3,4&(1,0\text{ or }1),\ (2,0,1,2).
\end{array}
\]
In this display $(q,0,1,\ldots)$ abbreviates the listed
separate choices of the second coordinate, not a tuple of
more than two entries.
\end{proposition}
\begin{proof}
The proof of $m\le4$ is that of
Proposition~\ref{paper:canonical-positive-sixty}; its rank
estimate uses only the displayed assumptions. Adjunction
and tensoriality of $F$ give
\[
 [F(WW^*):\mathbf1]=\|F(W)\|^2=m^2+q,
 \qquad [F(W^2):\mathbf1]=m^2+t.
\]
Since the nonreal-braiding square constituents have nontrivial
cubic twists and $\theta_Z=-1$, they miss $A$. This proves
the two exact identities. Lemma~\ref{paper:canonical-square-braiding}
then gives the formula with $x$.

Every nonunit simple of $\mathcal C$ has dimension at least
three. Partitioning $\dim P=10-m$ into these dimensions
initially gives
\[
\begin{array}{c|c}
m&q\\\hline
1&1,2,3,5,9\\
2&1,2,4\\
3&1,2\\
4&1,2,4.
\end{array}
\]
For $m=2$ or $4$, the value $q=4$ is exactly the case
$P=2R$, with $d_R=4$ or $3$, respectively. The content
test would make $(7-m)/2$ integral, whereas it is $5/2$
or $3/2$. Thus these two cases are impossible.

For $q=1$, $P$ is simple, so $t=0$ or $1$. For $q=2$,
$P$ consists of two distinct simples; their self-duality or
duality pairing gives $t=0,1,2$. The remaining cases occur
only at $m=1$ and total dimension nine. For $q=3$ there
are three distinct degree-three simples, giving $t=0,1,2,3$.
For $q=5$ one has $P=2R+S$ with distinct degree-three
simples. If both are self-dual then $t=5$; if only $R$
is self-dual then $t=4$; if only $S$ is self-dual then
$t=1$; if $R^*=S$ then $t=4$; and otherwise $t=0$.
For $q=9$, $P=3R$ and $t=9$ or zero. These exhaust the
possibilities and prove the table. No existence is asserted.
\end{proof}

\begin{corollary}[Small actual dimensions below a canonical degree ten]
\label{paper:positive-seven-forgetful-types}
Under the preceding hypotheses, the following are the possible
forgetful types when $m\ge2$:
\[
\begin{array}{c|l}
m&F(W)\\\hline
2&2\mathbf1+R_8,\quad
  2\mathbf1+R_3+S_5,\quad
  2\mathbf1+R_4+S_4\ (R_4\ne S_4)\\
3&3\mathbf1+R_7,\quad 3\mathbf1+R_3+S_4\\
4&4\mathbf1+R_6,\quad
  4\mathbf1+R_3+S_3\ (R_3\ne S_3).
\end{array}
\]
If $K_0(\mathcal C)$ is commutative, then $m=1$ and
\[
 F(W)=\mathbf1+R_9,\qquad
 \mathbf1+R_3+S_6,\qquad\text{or}\qquad
 \mathbf1+R_4+S_5.
\]
Thus in the commutative case the existence of this canonical
profile forces a degree-nine simple, or simples of degrees
three and six, or simples of degrees four and five in
$\mathcal C$. In either two-term remainder, at least one of
the two constituents is not self-dual.
\end{corollary}
\begin{proof}
For $m\ge2$, Proposition~\ref{paper:positive-seven-gram} makes $P$
multiplicity free, of length at most two. Its dimension is
$10-m$, giving the displayed list.
If the fusion ring is commutative, all canonical
multiplicities are zero or one. Thus $m=1$ and
$q=m_U+m_V\le2$, so $P$ again has at most two distinct
summands. The partitions of nine with every part at least
three are the three listed possibilities. The square identity
also gives $1+t=m_{B_+}+m_{B_-}\le2$, hence $t\le1$.
The two remainder degrees differ, so their being both
self-dual would give $t=2$; this is impossible.
\end{proof}

\begin{corollary}[Rank and field costs of a positive seven profile]
\label{paper:positive-seven-orbit-cost}
Under the same hypotheses, let $e=[K_W:\mathbb Q]$.
For $m=2,3,4$ one has, respectively,
\[
 \operatorname{rank}(\mathcal C)\ge 4e+19,\qquad
 \operatorname{rank}(\mathcal C)\ge 9e+87,\qquad
 \operatorname{rank}(\mathcal C)\ge16e+265.
\]
In particular these cases require ranks at least $27,105,297$.
Using $\operatorname{cond}(K_W)\mid1080$ and rank at most
$490$ gives respectively $e\le96,36,12$.
At $m=4$, the full field of $B_+$ has degree at most two,
and that of $B_-$ has degree at most six.
\end{corollary}
\begin{proof}
The full orbit of $W$ contributes $em^2$ to the canonical
rank. Here $e\ge2$ because $W$ is not self-dual. The lower
square multiplicities contribute
$\binom{m+1}{2}^2+\binom m2^2$, and $q\ge1$ gives
$m_U+m_V\ge m^2$, contributing at least
$\lceil m^4/2\rceil$. Including the unit gives the three
displayed bounds. The canonical conductor bound puts
$K_W\subseteq\mathbb Q(\zeta_{1080})$, of degree $288$.
Thus $e$ is an even divisor of $288$, and combining with
rank at most $490$ gives the stated limits.

At $m=4$, write $e_\pm=[K_{B_\pm}:\mathbb Q]$.
The bounds $m_{B_+}\ge10$, $m_{B_-}\ge6$ and
$m_U+m_V\ge16$ give
\[
 490\ge1+32+100e_++36e_-+128.
\]
The two square orbits are disjoint if their degrees differ.
If both square degrees are forty, they are still disjoint:
the exact formulas with $t\le2$ give $m_{B_+}>m_{B_-}$,
whereas canonical multiplicity is constant on an orbit.
Thus the displayed orbit count is valid. It implies
$e_+\le2$ and $e_-\le6$.
\end{proof}

\begin{corollary}[A stronger rank cost in the seventy-two profile]
\label{paper:positive-seventy-two-rank}
If the positive $(15,12,72)$ profile has $m_W=4$, then
\[
                 \operatorname{rank}(\mathcal C)\ge305.
\]
If its degree-twelve adjoint constituent has canonical
multiplicity at least seven, then the rank is at least $348$.
The general positive-sixty bound remains $297$.
\end{corollary}
\begin{proof}
Put $c=m_{12}$ and $h=m_{72}$. The Gram identity gives
$c+h\ge16$. If $c\le6$, then $c^2+h^2\ge136$.
If $c\ge7$, Corollary~\ref{paper:rational-canonical-remainder}
excludes a rational full field for the degree-twelve constituent.
Its full orbit therefore contributes at least $2c^2$, and
$2c^2+h^2\ge179$, with the minimum at $(c,h)=(7,9)$.
The unit, the full orbit of $W$, and the two real-braiding
square targets have disjoint orbits and contribute at least
$1+2\cdot4^2+10^2+6^2=169$. Adding the corresponding
adjoint contribution proves both bounds.
\end{proof}

\begin{proposition}[Canonical multiplicities in the forty-five profile]
\label{paper:canonical-fortyfive-multiplicities}
Suppose $m_W=m>0$, $d_W=10$, and
\[
 WW^*=\mathbf1+U_{24}+V_{30}+Z_{45}.
\]
Then the following necessary alternatives hold.
\begin{enumerate}
\item If $\theta_Z=-1$, then $\theta_V=1$ and $m\le3$.
The two square rows with no negative real braiding have
$m=1$. In the other row there is a degree-ten square
constituent $B$ with
$m_B\ge\binom{m+1}{2}$. In particular $m=3$ forces
another canonical degree-ten constituent of multiplicity
at least six.
\item If $\theta_Z=1$ and the square has degrees
$(40,30,15,15)$, then $m\le4$. If $\theta_V\ne-1$
and $m\ge2$, the remainder $F(W)-m\mathbf1$ is
multiplicity free and has at most two simple constituents.
\item If the square has degrees $(45,15,20,20)$, then
$\theta_Z=1$, $\theta_V=-1$, and $m\le4$.
\end{enumerate}
\end{proposition}
\begin{proof}
The twists and square alternatives are supplied by
Proposition~\ref{paper:global-fortyfive-square-table}.
When $\theta_Z=-1$, its $\theta_V=-1$ possibility is
excluded by Proposition~\ref{paper:negative-ten-not-canonical}.
The two rows without a $-1$ braiding channel have $m=1$.
In the remaining row the unique $+1$ channel is a
degree-ten simple $B$. Thus $m_B\ge\binom{m+1}{2}$.
Every nonunit central simple of degree ten has canonical
multiplicity at most seven: its forgetful image is not
supported entirely on the unit, and its remaining simples
have dimension at least three. Hence $m\le3$.

In the $(40,30,15,15)$ square, the degree-fifteen twists
are $\overline{\theta_V}\ne1$. The degree-forty and
degree-thirty real channels are distinct from the adjoint
degree-twenty-four and degree-forty-five constituents.
The latter two have twists one, whereas $V$ does not.
The same rank estimate as in
Proposition~\ref{paper:canonical-positive-sixty} therefore
gives $m\le4$.
For $m=2$ or $4$, a repeated remainder would be
$F(W)=m\mathbf1+2R$. The content test makes
$(7-m+3\overline{\theta_V})/2$ integral. Since $7-m$
is odd, this makes $(1+\theta_V)/2$ integral. This number
has house at most one, so it is zero or a root of unity;
the nonzero case saturates the triangle inequality and
gives $\theta_V=1$. Thus only $\theta_V=-1$ can permit
the repeated remainder. The other dimension partitions
are exactly the multiplicity-free one- or two-term
remainders listed in Corollary~\ref{paper:positive-seven-forgetful-types}.

For the exceptional square, let $P_{45}$ and $Q_{15}$ be
the real $+1$ and $-1$ channels. If $P\ne Z$, the earlier
disjoint-degree rank argument again gives $m\le4$.
Otherwise write $a=m_Z$, $b=m_Q$, and $c=m_U$.
The square bounds give
$a\ge A_m:=m(m+1)/2$ and $b\ge B_m:=m(m-1)/2$.
The forgetful remainder is nonzero, so the adjoint Gram
identity gives $a+c\ge m^2$.
Consequently, for $m\ge6$,
\[
 \operatorname{rank}(\mathcal C)
 \ge1+2m^2+a^2+b^2+c^2
 \ge1+2m^2+A_m^2+2B_m^2\ge964.
\]
Here the convex quadratic $a^2+(m^2-a)^2$ is increasing
for $a\ge A_m>m^2/2$; allowing $c=0$ beyond $a=m^2$
does not lower the minimum. This contradicts rank at most
$490$, so initially $m\le5$. If $m=5$, the same estimate
gives rank at least $476$. The degree-fifteen square constituent
$Q$ has $m_Q\ge10$. If its full field were nonrational,
its orbit would add at least $100$ to the displayed rank
contribution, contradicting the upper bound $490$. Hence
$Q$ is rational. Corollary~\ref{paper:rational-canonical-remainder}
then gives $15-m_Q\ge6$, contrary to $m_Q\ge10$. Thus
$m\le4$.
\end{proof}

\begin{corollary}[A large multiplicity requires a longer adjoint]
\label{paper:large-ten-canonical-multiplicity}
If a canonical degree-ten constituent has multiplicity at
least five, its adjoint has at least four nonunit occurrences.
In particular the multiplicity-three negative-forty-five
case forces another canonical degree-ten constituent with
at least four nonunit adjoint occurrences.
\end{corollary}
\begin{proof}
The earlier degree-ten classification leaves the positive
sixty, globally confined seventy-two, and forty-five profiles
when there are exactly three nonunit occurrences; the
other short profiles are excluded globally. The positive
seven profiles have multiplicity at most four, and the
preceding proposition bounds every remaining forty-five
case by four. The final assertion follows from
$\binom{3+1}{2}=6$.
\end{proof}

\subsection{A strict character disk after removing the unit}
\label{paper:subsec:shifted-disk}

\begin{theorem}[The canonical shifted disk]
\label{paper:canonical-shifted-disk}
Under \eqref{paper:source:assumption}, if $W\ne\one$ and
$m_W>0$, then for every nonunit canonical constituent
$Z_E\subset A$,
\begin{equation}\label{paper:shifted-disk}
       |\chi_{Z_E}(W)-m_W|<d_W-m_W.
\end{equation}
The same inequality holds after every Galois conjugation.
\end{theorem}
\begin{proof}
The class $F(W)$ is central in $K_0(\mathcal C)$, and
\[
             F(W)=m_W\one+P,
             \qquad\FPdim(P)=d_W-m_W.
\]
Proposition~\ref{paper:source:source-structure} gives $P\ne0$.
The scalar by which $F(W)$ acts on $E$ is
$\chi_{Z_E}(W)$, by the canonical character correspondence.
Write \(P=\sum_{X\ne\mathbf1}b_XX\).  The strict bound of
Lemma~\ref{paper:source:strict-character-operator-norm} gives
\(\|E(P)\|\le\sum_Xb_X\|E(X)\|<\sum_Xb_Xd_X=d_W-m_W\).
Since \(E(P)\) is the scalar \(\chi_{Z_E}(W)-m_W\), this proves
\eqref{paper:shifted-disk}.
The Galois group permutes the nonunit constituents of $A$ and
preserves their multiplicities, giving all conjugate inequalities.
\end{proof}

\begin{corollary}[Multiplicities at supported ordinary endpoints]
\label{paper:canonical-endpoint-multiplicity}
Suppose both endpoints $X,Y$ of a compulsory mass-four crossing
occur in $A$. Then their crossing is ordinary, and
\[
                     m_X<\frac{d_X}{4},\qquad
                     m_Y<\frac{d_Y}{4}.
\]
In particular,
\begin{equation}\label{paper:canonical-endpoint-numeric}
 \begin{array}{c|cc}
 (d_X,d_Y)&m_X&m_Y\\\hline
 (12,20)&1\le m_X\le2&1\le m_Y\le4\\
 (36,20)&1\le m_X\le8&1\le m_Y\le4.
 \end{array}
\end{equation}
\end{corollary}
\begin{proof}
The ordinary assertion is
Corollary~\ref{paper:canonical-endpoints-not-signed}.
Write $z=1+3\xi$ with $\xi$ primitive of order five and
$L=\operatorname{lcm}(d_X,d_Y)$. In the compulsory mass-four
case, $d_Xd_Y=4L$, so
\[
       \chi_X(Y)=\frac{d_Y}{4}(1+3\xi).
\]
Put $a=d_Y/4$, $m=m_Y$, and $c=\operatorname{Re}\xi<1$.
The square of \eqref{paper:shifted-disk} gives
\[
 (a-m)^2+9a^2+6a(a-m)c<(4a-m)^2,
\]
which is equivalent to $6a(m-a)(1-c)<0$.
Thus $m_Y<a$. Interchanging the two endpoints proves the
other bound.
\end{proof}

\begin{corollary}[The sign of a supported norm-ten character matters]
\label{paper:norm-ten-supported-multiplicity}
In the no-thirty branch let $Y$ be a compulsory degree-twenty
witness, let $d_W=10$, and assume $Y,W\subset A$ and
$|\chi_Y(W)|^2=10$. With
$\alpha=(5+\sqrt{-15})/2$, one has
\[
 \begin{array}{c|c}
 \chi_Y(W)&m_W\\\hline
 \alpha\text{ or }\overline\alpha&m_W\le5\\
 -\alpha\text{ or }-\overline\alpha&m_W\le3.
 \end{array}
\]
\end{corollary}
\begin{proof}
The four character values are
Theorem~\ref{paper:norm-ten-untwisted}.
For $\beta=\chi_Y(W)$, the shifted disk reads
$10-2m_W\operatorname{Re}\beta<(10-m_W)^2-m_W^2$.
Equivalently,
$m_W(20-2\operatorname{Re}\beta)<90$.
The real parts are $5/2$ and $-5/2$, giving the two integer
bounds.
\end{proof}

\begin{remark}[A membership condition remains]
These bounds require the named labels to occur in $A$.
Trivial twist is necessary for that membership but does not
imply it. In particular, the compulsory endpoint fields and
twists established earlier do not by themselves place either
endpoint in the canonical algebra.
\end{remark}

\subsection{Positive source rows and canonical remainders}
The coefficient trace and the ordinary regular trace have different
weights:
\[
 \tau(x)=\sum_E\frac{\chi_E(x)}{f_E},\qquad
 \lambda(x):=\operatorname{Tr}L_x=\sum_E m_E\chi_E(x).
\]
Their simultaneous use retains information about the actual forgetful
images, beyond the canonical dimension equation.

\begin{lemma}[Strict operator norm in a simple fusion category]
\label{paper:source:strict-character-operator-norm}
Let $\mathcal C$ be a nonpointed fusion-simple integral fusion
category. If $X$ is a nonunit simple, $d=d_X$, and $E$ is an
irreducible representation of $K_0(\mathcal C)_{\mathbb C}$
different from the dimension representation, then, in a
unitary $*$-realization,
\[
                         \|E(X)\|<d.
\]
Consequently, for a Galois-stable collection $\mathcal B$ of
non-dimension irreducible representations, put
\[
 k=\sum_{E\in\mathcal B}\dim E,\qquad
 s=\sum_{E\in\mathcal B}\chi_E(X),\qquad
 t=\sum_{E\in\mathcal B}\chi_E(XX^*).
\]
Then $s,t$ are integers and
\[
                  0\le t\le kd^2-1,\qquad s^2\le kt.
\]
\end{lemma}

\begin{proof}
The regular $*$-representation first gives the weak bound
$\|E(Y)\|\le d_Y$ for every simple $Y$: the positive
operator $E(YY^*)$ occurs in the integral regular
representation, whose largest eigenvalue is $d_Y^2$.
If equality holds for $X$, choose a unit vector $v$ with
$E(XX^*)v=d^2v$. Write
$XX^*=\mathbf1+\sum_{T\ne\mathbf1}n_TT$.
Equality in
\[
 d^2=1+\sum_T n_T\operatorname{Re}\langle v,E(T)v\rangle
       \le1+\sum_Tn_Td_T=d^2
\]
forces $E(T)v=d_Tv$ for every occurring $T$, and the same
holds for $T^*$ by the weak operator bound. Since $X$ is
not invertible, some such $T$ is nonunit. Simplicity makes
$T$ and $T^*$ generate $\mathcal C$. Applying equality
in the same estimate to their words shows
$E(Y)v=d_Yv$ for every simple $Y$. Thus $\mathbb Cv$ is a
dimension-character subrepresentation, contradicting the
choice of $E$.

Every summand of $t$ is the squared Hilbert--Schmidt norm
of $E(X)$, and is strictly below $(\dim E)d^2$.
Galois stability and integrality of character values make
$s,t$ rational algebraic integers. Cauchy--Schwarz on the
direct sum of the representation matrices gives $s^2\le kt$.
\end{proof}

\begin{proposition}[Positive regular-word traces]
\label{paper:source:positive-regular-traces}
Let the fusion ring have rank $r$. For every nonnegative integral
combination $W$ of basis elements, in particular for every positive
word in $X$ and $X^*$,
\begin{equation}\label{paper:source:regular-trace-domination}
                  \lambda(W)\ge r\,\tau(W)\ge0,
             \qquad \lambda(W)\in\mathbb Z.
\end{equation}
For a nonunit source simple $X$ this gives
\[
       d_X+\sum_{E\ne{\rm FP}}m_E\chi_E(X)\ge0.
\]
The same scalar-value and matrix-determinant choices must be used
simultaneously when comparing the two traces of a word.
\end{proposition}

\begin{proof}
In the distinguished basis,
$L_W=\tau(W)I_r+\sum_{X\ne\mathbf1}n_XL_X$ with $n_X\ge0$.
All matrices in the sum have nonnegative integral entries. Taking
ordinary traces proves \eqref{paper:source:regular-trace-domination}.
The character expression follows from the decomposition of the
left regular module into $m_E$ copies of $E$.
\end{proof}

\begin{proposition}[Canonical support and regular-trace capacity]
\label{paper:source:canonical-regular-trace-identity}
Let $Z_E$ be a simple canonical constituent of $I(\mathbf1)$ and write
\[
                 F(Z_E)=m_E\mathbf1+
                     \sum_{X\ne\mathbf1}b_{X,E}X.
\]
Then the integers $b_{X,E}\ge0$ satisfy
\begin{align}
 \sum_{X\ne\mathbf1}d_Xb_{X,E}&=a_E-m_E,\label{paper:source:canonical-remainder-mass}\\
 \lambda(X)&=\sum_E m_E b_{X,E}
                \quad (X\ne\mathbf1).\label{paper:source:canonical-regular-capacity}
\end{align}
Thus a row used in the remainder of this constituent must have
$\lambda(X)\ge m_E$.

If the full character field $K_E$ is real, $Z_E$ and $F(Z_E)$ are
self-dual. Its remainder is therefore a sum of self-dual source
dimensions $d_X$ and doubled non-self-dual dimensions $2d_X$, with
nonnegative integer multiplicities. In particular the unique rational
degree-two constituent in the one-matrix shape has this stronger
self-dual remainder condition, with $\lambda(X)\ge2$.
\end{proposition}

\begin{proof}
The first identity is the dimension of the forgetful image, with
$b_{\mathbf1,E}=m_E$ by adjunction. The induction formula gives
\[
 F(I(\mathbf1))=\sum_Y YY^*=\sum_E m_EF(Z_E).
\]
By Frobenius reciprocity, the coefficient of $X$ in the middle
expression is
$\sum_Y N_{YY^*}^{X}=\sum_Y N_{XY}^{Y}=\Tr(L_X)$.
This proves \eqref{paper:source:canonical-regular-capacity}.

Proposition~\ref{paper:source:canonical-character-bridge} identifies the full
source character field $K_E$ with the full modular character field
$K_{Z_E}$. In a modular category complex conjugation sends a column
to its dual. A real full field therefore fixes that column, so
$Z_E=Z_E^*$. The forgetful functor preserves duals, and coefficients
of each non-self-dual source pair in its image must agree.
\end{proof}

A source row with $\lambda(X)=0$ is absent from every nonunit canonical
remainder, but may still occur in the source dimension equation.  The
selfduality of a canonical constituent is determined by its full field,
not by the value of a single character entry.

\begin{theorem}[Canonical capacity and an exact row knapsack]
\label{paper:source:canonical-capacity-knapsack}
Under the standing integral center hypotheses, partition the non-Perron
irreducible representations into sets $\mathcal G_\alpha$ closed under
complex conjugation, with common representation degree $m_\alpha$ and
canonical dimension $a_\alpha$ within each set. Write
$c_\alpha=|\mathcal G_\alpha|$ and $R_\alpha=a_\alpha-m_\alpha$.
For a nonunit source simple $X$, put $d=d_X$ and
$\lambda_X=\operatorname{Tr}L_X$, and define
\[
 U_\alpha(X)=
 \begin{cases}
 c_\alpha\lfloor R_\alpha/d\rfloor,&X=X^*,\\[2pt]
 \left\lfloor c_\alpha\lfloor R_\alpha/d\rfloor/2\right\rfloor,
       &X\ne X^*.
 \end{cases}
\]
Then there are integers $B_\alpha(X)$ satisfying
\[
 0\le B_\alpha(X)\le U_\alpha(X),\qquad
 \lambda_X=\sum_\alpha m_\alpha B_\alpha(X).
\]
In particular a proposed source row is impossible if this bounded
integer equation has no solution. A row with $\lambda_X=0$ is retained:
all $B_\alpha(X)=0$ satisfy this necessary condition.
\end{theorem}

\begin{proof}
For the canonical simple $Z_E$ corresponding to $E$, write
\[
 F(Z_E)=m_E\mathbf1+\sum_{Y\ne\mathbf1}b_{Y,E}Y,
 \qquad b_{Y,E}\in\mathbb Z_{\ge0}.
\]
The positive remainder has dimension $a_E-m_E$, so
$b_{X,E}\le\lfloor R_\alpha/d\rfloor$ for
$E\in\mathcal G_\alpha$. Put
$B_\alpha(X)=\sum_{E\in\mathcal G_\alpha}b_{X,E}$.
The regular-forgetful identity gives
$\lambda_X=\sum_E m_Eb_{X,E}=\sum_\alpha m_\alpha B_\alpha(X)$.

If $X=X^*$, summing the individual coefficient bounds proves the
first capacity. If $X\ne X^*$, the remainder identity gives
\[
 b_{X,E}+b_{X^*,E}\le\lfloor R_\alpha/d\rfloor.
\]
Duality of the canonical correspondence and closure of
$\mathcal G_\alpha$ under conjugation give
\[
 \sum_{E\in\mathcal G_\alpha}b_{X^*,E}
  =\sum_{E\in\mathcal G_\alpha}b_{X,E}
  =B_\alpha(X).
\]
Summing the preceding inequalities and using integrality gives
$2B_\alpha(X)\le c_\alpha\lfloor R_\alpha/d\rfloor$,
which is the asserted pair capacity.
\end{proof}

\begin{corollary}[A finite lift by canonical coefficient allocations]
\label{paper:source:canonical-allocation-lift}
For each eligible source-duality row $x$, let $d_x$ be its source
dimension, let $\lambda_x$ be its regular trace, and let
$\epsilon_x=1$ for a selfdual label and $\epsilon_x=2$ for one
representative of a dual pair. Define the finite set
\[
 \mathcal A_x=\left\{(B_\alpha)_\alpha:
       B_\alpha\in\mathbb Z_{\ge0},\quad B_\alpha\le U_\alpha(x),\quad
       \sum_\alpha m_\alpha B_\alpha=\lambda_x\right\}.
\]
If $v_x$ is the ordinary source-count, dimension and character
moment vector, replace it by all the lifted vectors
\[
 \left(v_x,\ (\epsilon_xd_xB_\alpha)_\alpha\right),
                   \qquad B\in\mathcal A_x.
\]
Append the coordinates $(c_\alpha R_\alpha)_\alpha$ to the
ordinary moment target. An actual fusion ring gives a
nonnegative \emph{integer} combination of the lifted vectors
equal to that target. In particular, an exact separating
covector for their nonnegative real cone excludes the branch.
\end{corollary}

\begin{proof}
For each actual nonunit source label choose the vector
$B_\alpha(X)=\sum_{E\in\mathcal G_\alpha}b_{X,E}$ from
Theorem~\ref{paper:source:canonical-capacity-knapsack}. It belongs to
$\mathcal A_x$. The dimension equation for every canonical
column gives
\[
 \sum_{X\ne\mathbf1}d_XB_\alpha(X)
   =\sum_{E\in\mathcal G_\alpha}(a_\alpha-m_\alpha)
   =c_\alpha R_\alpha.
\]
Duality stability gives $B_\alpha(X^*)=B_\alpha(X)$, which
accounts for the factor two on a pair row. Grouping labels
with the same row and allocation gives the required
nonnegative integer coefficients. A separating covector is
nonnegative on any such combination and therefore cannot be
negative on the target.
\end{proof}

\begin{theorem}[Bounded allocation separation]
\label{paper:source:bounded-allocation-separation}
Use the notation of Corollary~\ref{paper:source:canonical-allocation-lift}.
For a row $x$, an integer coefficient vector $z=(z_\alpha)$, and
regular trace $\lambda_x$, put
\[
 M_x(z)=\min\left\{\sum_\alpha z_\alpha B_\alpha:
 0\le B_\alpha\le U_\alpha(x),\quad
 \sum_\alpha m_\alpha B_\alpha=\lambda_x,\quad
 B_\alpha\in\mathbb Z\right\}.
\]
Suppose all row allocation sets are nonempty.  An integer pair
$(y,z)$ is nonnegative on every lifted source row if and only if
\[
 y\cdot v_x+\epsilon_xd_xM_x(z)\ge0
 \quad\text{for every eligible row }x.
\]
If additionally
\[
 y\cdot T+\sum_\alpha z_\alpha c_\alpha R_\alpha<0,
\]
then the proposed canonical data are impossible.
\end{theorem}

\begin{proof}
For a fixed row $x$, minimize the value of the covector on
\[
 (v_x,(\epsilon_xd_xB_\alpha)_\alpha).
\]
Since $\epsilon_xd_x>0$, its minimum is precisely the expression
in the assertion.  Nonnegativity on every lifted generator then
implies nonnegativity on every nonnegative integral combination,
contradicting the strictly negative value on the required target.
\end{proof}

\begin{theorem}[Packet scalar targets]
\label{paper:source:galois-packet-central-remainder}
Under the standing center hypotheses, let $\mathcal P$ be a nonempty
Galois-stable set of nonunit simple constituents of $I(\mathbf1)$.
Their dimensions and canonical multiplicities are denoted by $a_Z$
and $m_Z$. Then
\[
 P_{\mathcal P}=\sum_{Z\in\mathcal P}\bigl(F(Z)-m_Z\mathbf1\bigr)
\]
is a nonzero positive selfdual central integral source-ring element
with unit coefficient zero and dimension
$R=\sum_{Z\in\mathcal P}(a_Z-m_Z)>0$. In every irreducible source
representation it acts by an integer scalar. The Perron scalar is
$R$, and all other scalars have absolute value strictly less than $R$.
Individual constituents in the packet need not have rational full
character fields.
\end{theorem}

\begin{proof}
Adjunction gives unit multiplicity $m_Z$ in $F(Z)$, so the indicated
subtraction leaves a positive integral remainder. Centrality follows
from the central forgetful action. Complex conjugation is one of the
Galois permutations and takes a modular column to its dual, so the
packet and its forgetful sum are selfdual.

Fix a source representation $V$, with associated canonical column
$Z_V$ and dimension $a_V$. By
Proposition~\ref{paper:source:canonical-character-bridge}, the scalar of $F(Z)$
on $V$ is $S_{Z,Z_V}/a_V$. Symmetry of the unnormalized modular
matrix rewrites this as
\[
       \frac{a_Z}{a_V}\,\chi_Z(Z_V).
\]
On every full orbit in the packet, $a_Z$ is a constant rational
integer, and Galois conjugation permutes the characters $\chi_Z$.
Thus the sum of these scalars is rational. It is also an algebraic
integer, since it is an eigenvalue of multiplication by the integral
element $\sum F(Z)$. Subtracting the integer $\sum m_Z$ proves the
integer-scalar assertion.

If the remainder vanished, every $Z$ in the packet would have
forgetful image supported on the unit. The central objects with that
property form a proper fusion subcategory, because forgetting is
dominant onto the nontrivial source. Simplicity of the center rules
out its containing a nonunit simple. The strict non-Perron norm bound
now applies to the positive nonunit support, exactly as in
Theorem~\ref{paper:rational-canonical-scalar-targets}.
\end{proof}

For a rational canonical constituent, the packet can consist of one
label.  If its remainder is $P=B_1+\cdots+B_h$, with repeated selfdual
source atoms listed according to multiplicity, each Galois orbit of
source representations gives a useful further constraint.  If its
total representation degree is $k$, write
$s(B)=\sum_E\chi_E(B)$ and $t(B)=\sum_E\chi_E(B^2)$.
The positive quantities $q(B)=kt(B)-s(B)^2$ are squared norms of the
traceless Hermitian images, and these images sum to zero.  Moreover
$k\mid\sum_\alpha s(B_\alpha)$, because the remainder acts by one
integer scalar on the orbit.  The resulting two- and three-atom
tests and the complete bounded central-spectrum enumeration are
given in the following results
(Theorem~\ref{paper:rational-canonical-scalar-targets} and
Propositions~\ref{paper:small-central-spectrum-problem}--\ref{paper:small-central-atom-completion}).
\subsection{Canonical support covers and forgetful capacities}
\label{paper:canonical-support-subsection}

The Boolean canonical-remainder test can be coupled to the global
character moments. This retains information about which eligible source
rows must actually occur, rather than only which dimensions occur
somewhere in the candidate row list.

\begin{lemma}[Canonical support-cover inequalities]
\label{paper:canonical-support-cover}
Let $Z$ be a nonunit simple constituent of $I(\mathbf1)$, put
$a=\FPdim Z$ and $m=[I(\mathbf1):Z]$, and write $R=a-m$.
Fix any necessary finite list of candidate nonunit source rows, including
source duality and the regular trace $T_X=\Tr(L_X)$. Only candidate rows
with $T_X\ge m$ need be considered as possible constituents of
$F(Z)-m\mathbf1$.

For an arbitrary $Z$, take the eligible dimensions as coin values.
If $Z\simeq Z^*$, take instead the dimensions of eligible selfdual
source rows and twice the dimensions of eligible nonselfdual source rows.
Let $H$ be any set of source dimensions such that $R$ is not in the
additive semigroup generated by the coins belonging to dimensions
outside $H$. Then at least one actual nonunit source simple has dimension
in $H$ and regular trace at least $m$.

Consequently, in a source moment model whose variables count selfdual
objects and nonselfdual pairs, there is the valid inequality
\[
 \sum_{\rho}
  \mathbf1_{\{d_\rho\in H,\ T_\rho\ge m\}}\,n_\rho\ge1,
 \qquad
 \gamma_\rho=\begin{cases}1,&\rho\text{ selfdual},\\
                         2,&\rho\text{ nonselfdual}.
             \end{cases}
\]
This conclusion uses no divisibility assumption on source dimensions.
\end{lemma}
\begin{proof}
Adjunction gives $[F(Z):\mathbf1]=m$. Hence
\[
 F(Z)=m\mathbf1\oplus\bigoplus_{X\ne\mathbf1}b_{X,Z}X,
 \qquad R=\sum_{X\ne\mathbf1}d_Xb_{X,Z},
 \qquad b_{X,Z}\in\mathbb Z_{\ge0}.
\]
Moreover
\[
 T_X=\Tr(L_X)=[F(I(\mathbf1)):X]
     =\sum_Y [I(\mathbf1):Y]\,b_{X,Y}.
\]
Thus $b_{X,Z}>0$ implies $T_X\ge m$. If $Z$ is selfdual,
$F(Z)$ is selfdual, so the coefficients of $X$ and $X^*$ agree and a
nonselfdual pair contributes a coin $2d_X$. Were no actual constituent
of dimension in $H$ present, the displayed positive decomposition of
$R$ would use only the excluded complement coins. This proves existence
of an actual object in $H$. Its source duality orbit has a nonnegative integer multiplicity
at least one, giving the stated linear inequality. The factor
$\gamma_\rho$ counts individual objects when needed for capacities,
but is not needed in this orbit-existence inequality.
\end{proof}

The stated character--center field bridge identifies the full source
character field with that of its canonical central constituent.
Complex conjugation of modular characters is duality. In particular a
real full character field supplies a selfdual canonical constituent, so
the doubled-coin refinement above applies to it. An individual real
entry in a CM field does not make the full field real and does not justify
this refinement.

Dimension-dependent nonnegative costs are a special case. Indicator
costs recover capacity constraints for sets of dimensions or row types.
Unlike the support-cover inequality, these constraints add the demands
of all canonical constituents, with their representation multiplicities.
Constant cost and square-dimension cost give other inexpensive exact
recurrences. Each minimum is found by a finite coin-cost dynamic program.

\paragraph{Exact certificate format.}
For a moment system $\sum_\rho n_\rho v_\rho=t$, append the relevant
cover or capacity coordinates to each $v_\rho$ and append the lower
bounds to $t$. For each appended inequality, add one nonnegative slack
variable with vector minus the corresponding new coordinate vector.
Any integer functional nonnegative on every extended source vector and
every slack vector but negative on the extended target is an exact
certificate. The checker uses integer arithmetic only; a numerical
linear-program solver is merely a proposer.

\begin{lemma}[Global forgetful-capacity inequalities]
\label{paper:forgetful-capacity}
Keep a necessary source-row list and let $c(\rho)$ be any signed
integer cost on source row types $\rho=(d,\text{duality type},T)$.
The cost is constant on the two members of each nonselfdual source pair.
For an actual source simple $X$, write $c_X=c(\rho_X)$. For each canonical constituent
$Z$, let $\mu_Z$ be the minimum cost of a coin decomposition of
$a_Z-m_Z$, using only candidate rows of regular trace at least $m_Z$.
A selfdual coin of row type $\rho$ and dimension $d$ has cost
$c(\rho)$; when $Z$ is selfdual, a paired coin has dimension $2d$
and cost $2c(\rho)$. For a canonical constituent without imposed
selfduality, its ordinary dimension-$d$ coin has cost $c(\rho)$.
All coin dimensions are positive and their total is fixed. Consequently
there are only finitely many decompositions and every nonempty minimum
is finite, even when some costs are negative.
Then
\[
 \sum_{X\ne\mathbf1}c_X\,T_X
       \ \ge\ \sum_{Z\ne\mathbf1}m_Z\,\mu_Z.
\]
In source-row variables the left side is
$\sum_\rho\gamma_\rho c(\rho)T_\rho n_\rho$.
An impossible coin decomposition itself excludes the assignment.
\end{lemma}
\begin{proof}
For each actual forgetful column,
$\sum_{X\ne\mathbf1}c_Xb_{X,Z}\ge\mu_Z$ by the definition of the
minimum. Multiply by $m_Z$, sum over the canonical constituents, and
use $T_X=\sum_Zm_Zb_{X,Z}$. The unit canonical constituent contributes
zero to every nonunit row. All sums are finite; multiplying by the nonnegative representation
multiplicities preserves every inequality.
\end{proof}
\subsection{Simultaneous capacities of canonical columns}
\label{paper:sec:canonical-capacities}

Separate canonical-remainder tests ignore the fact that the same source
simple can contribute to several canonical columns only within its
regular-character budget. The following integral bounds couple those
columns without choosing their individual coefficients.

\begin{corollary}[Simultaneous canonical mass inequalities]
\label{paper:canonical-subset-capacities}
For any set $S$ of the groups in
Theorem~\ref{paper:source:canonical-capacity-knapsack}, put
\[
 T_S=\sum_{\alpha\in S}m_\alpha c_\alpha R_\alpha,
 \qquad
 C_S(X)=\min\left\{\lambda_X,
             \sum_{\alpha\in S}m_\alpha U_\alpha(X)\right\}.
\]
Then
\[
                     \sum_{X\ne\mathbf1}d_X C_S(X)\ge T_S.
\]
Equivalently, a sum over selfdual nonunit labels and one representative
of each nonselfdual pair uses contributions $d_XC_S(X)$ and
$2d_XC_S(X)$, respectively.
\end{corollary}

\begin{proof}
At each source label,
$\sum_{\alpha\in S}m_\alpha B_\alpha(X)\le C_S(X)$.
Multiplying by $d_X$, summing over nonunit labels, and using the
canonical remainder dimensions gives
\[
 \sum_{X\ne\mathbf1}d_X
       \sum_{\alpha\in S}m_\alpha B_\alpha(X)
   =\sum_{\alpha\in S}m_\alpha c_\alpha R_\alpha=T_S.
\]
All selected groups are dual-stable, so the summand is the same on
$X$ and $X^*$, proving the representative formulation.
\end{proof}

\paragraph{Exact implementation.}
The row equation is checked by a finite bitset computation: the
reachable values are the exponents in
\[
          \prod_\alpha
          \left(1+z^{m_\alpha}+\cdots+
                     z^{m_\alpha U_\alpha(X)}\right).
\]
Only support is needed, so binary splitting of each bounded interval
computes it using integer shifts and bitwise unions. The mass
inequalities can be adjoined to the source moment cone. A row receives
its capacity contributions, the target receives $T_S$, and a negative
unit slack vector is included for each additional coordinate. An
integer separating covector is accepted only after checking every
source vector, every slack vector, and the target exactly.

The hypotheses require complete dual-stable groups. An individual
nonreal character is not, by itself, such a group. No claim is made
that the unknown coefficients $b_{X,E}$ are independent, or that these
necessary inequalities suffice to construct any canonical column.

For a subset $S$ of groups, every lifted allocation satisfies
\[
 \sum_{\alpha\in S}m_\alpha B_\alpha
 \le\min\!\left\{\lambda_x,
                 \sum_{\alpha\in S}m_\alpha U_\alpha(x)\right\}.
\]
Thus the subset-capacity inequalities are projections of this
allocation system. The lift retains more of the integral
canonical coefficient data; it is a necessary reconstruction
condition, and is not asserted to construct a fusion ring.
The length-twelve certificates use the knapsack and subset
inequalities above, without requiring enumeration of these
larger lifted row sets.

\section{Character fields, positive trace forms and integral coordinates}

\label{paper:source:character-arithmetic}
\subsection{Rational forms and integral coordinates}
For a nonnegative integer $s$, write $I_s=\langle1,\ldots,1\rangle$.
The quadratic form is considered on the actual rational algebra;
neither a rational matrix component nor a full character field is
replaced by a split component or a discriminant alone.

\begin{theorem}[The positive rational trace form]
\label{paper:source:positive-rational-trace-form}
The rational quadratic form
\[
                         Q_R(x)=\tau(xx^*)
\]
is isometric over $\mathbb Q$ to $I_r$.  If $b$ is the number of
non-self-dual pairs of source simples, its restrictions to the
duality-fixed and duality-skew subspaces are respectively
\begin{equation}\label{paper:source:duality-trace-forms}
                 Q_R^+\simeq I_{r-2b}\perp2I_b,
                 \qquad Q_R^-\simeq2I_b.
\end{equation}
If $R$ is commutative and $K_j$ is the full field of a Galois
orbit of linear characters, of common formal codegree $f_j$, then
\begin{equation}\label{paper:source:commutative-field-trace-form}
 Q_R\simeq\mathop{\perp}_j
       \frac1{f_j}\operatorname{Tr}_{K_j/\mathbb Q}(x\overline x).
\end{equation}
The summand for the dimension character is $\langle1/5400\rangle$.
A real field contributes its entire trace form to $Q_R^+$ and
nothing to $Q_R^-$.  A CM field contributes
$(2/f_j)\operatorname{Tr}_{K_j^+/\mathbb Q}(x^2)$ to $Q_R^+$;
its skew subspace supplies the other half of its dimension.
\end{theorem}

\begin{proof}
For the simple basis one has
$\tau(XY^*)=\delta_{X,Y}$, so $Q_R=I_r$ in that basis.
For each non-self-dual pair replace $X,X^*$ by $X+X^*,X-X^*$.
These vectors are orthogonal of norm two; self-dual basis
elements retain norm one.  This proves
\eqref{paper:source:duality-trace-forms}.

The commutative rational algebra is the product of the character
fields.  The symmetrizing trace is $\sum_E\chi_E/f_E$.
Combining the embeddings in one Galois orbit therefore gives
$\operatorname{Tr}_{K_j/\mathbb Q}/f_j$ on its component.
Duality acts there as complex conjugation.  This proves
\eqref{paper:source:commutative-field-trace-form}.  The field is abelian,
so it is either totally real or CM.  On its real subfield in
the latter case, the trace from $K_j$ is twice the trace from
$K_j^+$, proving the final assertion.
\end{proof}

\begin{corollary}[A necessary local test beyond the squareclass]
\label{paper:source:local-trace-form-test}
For a positive rational diagonal form $Q=\langle a_1,\ldots,a_s\rangle$,
put
\[
  \det Q=\prod_i a_i\pmod{\mathbb Q^{\times2}},\qquad
  \epsilon_p(Q)=\prod_{i<j}(a_i,a_j)_p,
\]
where $(\ ,\ )_p$ denotes the Hilbert symbol.
Every proposed character-field assignment must have the same
dimension, determinant squareclass and local invariants as
the appropriate form in Theorem~\ref{paper:source:positive-rational-trace-form}.
For the total commutative form this requires
$\det Q=1$ and $\epsilon_p(Q)=1$ for every prime $p$.
The fixed and skew forms must be checked separately as well.
\end{corollary}

\begin{proof}
These are invariants of rational isometry.  Their computation
uses only diagonalization over $\mathbb Q$ and Hilbert symbols;
see \cite[Chapters~12 and~14]{VoightQuaternion} for the local
and global quadratic-form framework.  Only necessity is used
for an exclusion.  In particular no rational isometry passing
this test is claimed to preserve an integral fusion basis or
its multiplication.
\end{proof}

\begin{theorem}[Even representation degrees disappear from the positive determinant]
\label{paper:source:even-degree-trace-determinant}
Retain the center hypotheses at dimension $5400$.
Let $\mathcal O$ be a Galois orbit of irreducible source-ring
representations, of common dimension $m$, rational formal
codegree $f$, and full character field $K$ of degree $e$.
Its rational simple-algebra component contributes to the positive
form $\tau(xx^*)$ the determinant squareclass
\begin{equation}\label{paper:source:component-determinant-parity}
             \left(\frac{|\operatorname{disc}K|}{f^e}\right)^m.
\end{equation}
In particular every orbit of even-dimensional representations
contributes a square, regardless of its Schur index or whether
its involution is of the first or second kind. Consequently
\begin{equation}\label{paper:source:odd-component-determinant-condition}
 \prod_{\mathcal O:\,m_{\mathcal O}\text{ odd}}
        \frac{|\operatorname{disc}K_{\mathcal O}|}
             {f_{\mathcal O}^{[K_{\mathcal O}:\mathbb Q]}}
                  \in\mathbb Q^{\times2}.
\end{equation}
\end{theorem}

\begin{proof}
The rational component is a central simple algebra $B$ of degree
$m$ over $K$, possibly nonsplit. The symmetrizing trace on it is
\[
               \frac1f\operatorname{Tr}_{K/\mathbb Q}
                                      \operatorname{Trd}_{B/K}.
\]
Duality induces complex conjugation on its center. The character
field is either totally real or CM.

If $K$ is totally real, the involution is of the first kind.
Lemma~\ref{paper:reduced-trace-discriminant} makes the determinant
over $K$ a square before scaling. For a quadratic form of
$K$-dimension $m^2$, restriction of scalars multiplies its
rational determinant by $(\operatorname{disc}K)^{m^2}$ and
by the norm of its $K$-determinant. Scaling by $1/f$ then
contributes $f^{-em^2}$. Its squareclass is therefore
$(\operatorname{disc}K/f^e)^{m^2}$, which equals
\eqref{paper:source:component-determinant-parity} because $m^2\equiv m\pmod2$.

If $K$ is CM, put $F=K^+$ and write $K=F(\sqrt{-\delta})$
with $\delta$ totally positive. The involution is unitary, and
\[
                 H(x,y)=\operatorname{Trd}_{B/K}(x\sigma(y))
\]
is a nondegenerate Hermitian form on the $m^2$-dimensional
$K$-vector space $B$. Hermitian diagonalization is valid even
when $B$ is a division algebra: it concerns the vector space
and its form, not a matrix presentation of the algebra.
Write its diagonal entries as $h_1,\ldots,h_{m^2}\in F$.
The underlying $F$-trace form is the orthogonal sum of
$2h_i\langle1,\delta\rangle$. Its determinant squareclass
over $F$ is $\delta^{m^2}$. Restriction to $\mathbb Q$
contributes a square power of $\operatorname{disc}F$ and
$\operatorname{N}_{F/\mathbb Q}(\delta)^{m^2}$.
The latter has squareclass $|\operatorname{disc}K|^{m^2}$,
as follows from the relative discriminant of the basis
$1,\sqrt{-\delta}$ and the discriminant formula in a tower.
The scaling contribution is again $f^{-em^2}$, proving
\eqref{paper:source:component-determinant-parity}.

Distinct rational components are orthogonal for $\tau(xx^*)$,
whose determinant is one in the simple basis. Multiplying their
squareclasses and discarding the even $m$ factors proves
\eqref{paper:source:odd-component-determinant-condition}.
\end{proof}

\begin{theorem}[Character-coordinate covariance]
\label{paper:source:character-coordinate-covariance}
Let $K$ be the full character field of a Galois orbit of irreducible
representations, all of common degree $m$ and rational formal codegree
$f$. Assume that this orbit is disjoint from the Perron character.
Let $B=(b_0,\ldots,b_{e-1})$ be a rational basis of $K$, with $b_0=1$,
such that the coordinate vector $z(\alpha)$ is integral whenever
$\alpha\in\mathcal O_K$. It is not required that every integral
coordinate vector represent an algebraic integer. Put
\[
 G_{ij}=\Tr_{K/\mathbb Q}(b_i\overline{b_j}),\qquad
 H=mfG^{-1},\qquad z_0=z(1).
\]
For a representative character $\chi$ and each source simple $X$, write
$z_X=z(\chi(X))$. Then
\[
 z_X\in\mathbb Z^e,\qquad
 z_{\mathbf1}=m z_0,\qquad
 \sum_X z_Xz_X^{\mathsf t}=H,\qquad
 \sum_X d_Xz_X=0.
\]
If $C$ is the coordinate matrix of complex conjugation, then
$z_{X^*}=Cz_X$ as well. No Schur splitting assumption is needed.
\end{theorem}

\begin{proof}
Character values are algebraic integers. Each conjugate character
satisfies $\chi^\sigma(X^*)=\overline{\chi^\sigma(X)}$.
Since character values generate $K$, every embedding intertwines
field conjugation with complex conjugation. Let $E$ be the embedding
matrix with entries $E_{\sigma j}=\sigma(b_j)$. The vector of conjugate
characters on $X$ is $Ez_X$. Character orthogonality on a degree-$m$
orbit gives
\[
 E\left(\sum_Xz_Xz_X^{\mathsf t}\right)E^*=mfI_e.
\]
Since $E^*E=G$, inversion proves the covariance formula. Orthogonality
to the Perron character gives $E\sum_Xd_Xz_X=0$, hence the asserted
weighted sum. The unit character value is $m$, and the positive
involution gives $\chi(X^*)=\overline{\chi(X)}$. These statements
establish the remaining identities.
\end{proof}

\begin{theorem}[The integral skew projection]
\label{paper:source:skew-projection-lattice}
Let $\mathcal O$ be a full Galois orbit of irreducible representations of
a fusion ring, with common rational integral formal codegree $f$, and put
$\chi_{\mathcal O}=\sum_{E\in\mathcal O}\chi_E$. Choose one basis element
$X_i$ from each non-self-dual pair, $1\le i\le b$, and set
$a_i=X_i-X_i^*$. Then
\begin{equation}\label{paper:source:skew-projection-matrix}
 K_{ij}=-\frac12\chi_{\mathcal O}(a_i a_j)
       =\chi_{\mathcal O}(X_iX_j^*)-
          \chi_{\mathcal O}(X_iX_j)
       \quad\text{belongs to }\mathbb Z.
\end{equation}
The matrix $K$ is symmetric positive semidefinite and satisfies
\[
                K^2=fK,\qquad \Tr K=f b_{\mathcal O}.
\]
Here $K/f$ is the orthogonal projection, in the orthonormal coordinates
$a_i/\sqrt2$, onto the skew-adjoint part of the rational simple component.
Projections attached to different components are orthogonal and sum to
$I_b$. If the component has representation degree $m$ and full field
degree $e$, then
\[
 b_{\mathcal O}=
 \begin{cases}
 e m(m-1)/2,&\text{real center, orthogonal involution},\\
 e m(m+1)/2,&\text{real center, symplectic involution},\\
 e m^2/2,&\text{CM center, unitary involution}.
 \end{cases}
\]
No splitting assumption is made on the rational simple algebra.
\end{theorem}

\begin{proof}
The full orbit character is rational-valued and algebraically integral
on the fusion ring, hence integer-valued. Cyclicity and
$\chi_{\mathcal O}(x^*)=\chi_{\mathcal O}(x)$ give the second expression
in \eqref{paper:source:skew-projection-matrix}. The positive coefficient form is
$\langle x,y\rangle=\tau(xy^*)$. On the component belonging to
$\mathcal O$ it is $\chi_{\mathcal O}(xy^*)/f$. Consequently its matrix
on the normalized source skew coordinates is exactly $K/f$. Orthogonal
projection gives all the displayed identities. The dimensions follow
after extending the component and its involution to $\mathbb R$; this
does not require its splitting over the full character field.
\end{proof}

\begin{corollary}[A finite skew parameter for a rational orthogonal block]
\label{paper:source:orthogonal-skew-divisor}
For a rational orthogonal degree-two component, there are a positive
squarefree divisor $\delta$ of $f$ and an integral vector $z$ such that
\begin{equation}\label{paper:source:rank-one-skew-factorization}
                 K=\delta zz^{\mathsf t},\qquad
                 \sum_{i=1}^b z_i^2=f/\delta.
\end{equation}
The vector $z$ need not be primitive. A generator $U$ of the skew line
can be chosen with $U^2=-\delta$, up to rational rescaling. Thus both
the skew squareclass and all possible skew diagonal entries are finite
at a specified $f$ and source pair count $b$.
\end{corollary}

\begin{proof}
An integral positive semidefinite matrix of rank one has the form
$c vv^{\mathsf t}$, with $v$ a primitive integral vector and $c$ a
positive integer. Write $c=\delta q^2$ and $z=qv$. The trace identity
gives \eqref{paper:source:rank-one-skew-factorization}. If initially
$U^2=-\epsilon$, its component norm is $2\epsilon/f$, whereas the source
line has norm squareclass $2\sum v_i^2$. Since
$f=c\sum v_i^2$, it follows that $\epsilon$ has squareclass $c$, hence
$\delta$.
\end{proof}

\begin{corollary}[A bounded discriminant group]
\label{paper:source:skew-discriminant-exponent}
Let $P=K/f$, $d=\operatorname{rank}K$, and
$\Lambda=\operatorname{im}(P)\cap\mathbb Z^b$, with its ordinary
Euclidean integral form. Then
\[
 P(\mathbb Z^b)=\Lambda^*,\qquad f\Lambda^*\subseteq\Lambda,
 \qquad \det\Lambda\mid f^{\min(d,b-d)}.
\]
The positive coefficient form on this component's skew space has
determinant squareclass $2^d\det\Lambda$.
\end{corollary}

\begin{proof}
The lattice $\Lambda$ is primitive. Every integral functional on it
extends to $\mathbb Z^b$, proving the first equality. Integrality of
$K$ proves the inclusion. Restriction of the Euclidean pairing
identifies
$\mathbb Z^b/(\Lambda+\Lambda^\perp)$ with both discriminant groups
$\Lambda^*/\Lambda$ and
$(\Lambda^\perp)^*/\Lambda^\perp$. This group has exponent dividing
$f$ and can be generated by at most $\min(d,b-d)$ elements, giving the
divisibility. Finally the source skew form in the basis $a_i$ is
$2I_b$.
\end{proof}

\subsection{Commutative character-field discriminants}
\begin{proposition}[Discriminant squareclasses of commutative character fields]
\label{paper:commutative-character-discriminant}
Suppose $R=K_0(\mathcal C)$ is commutative, and let
$K_1,\ldots,K_s$ be the full character fields of its distinct
Galois orbits. In positive rational numbers modulo squares,
\[
                \prod_E f_E
                  \equiv\prod_{j=1}^s|\operatorname{disc}(K_j)|.
\]
If an abelian Galois field $K/\mathbb Q$ has exactly one
quadratic subfield $\mathbb Q(\sqrt\Delta)$, where $\Delta$
is a fundamental discriminant, then the absolute discriminant
of $K$ has squareclass $|\Delta|$. If it has no quadratic
subfield or more than one, its absolute discriminant is a square.
\end{proposition}
\begin{proof}
The coefficient-of-unit form $\tau(xy)$ has determinant
$1$ or $-1$ in the simple basis. The regular trace form
is obtained by multiplication by the central element acting
by $f_E$ on character $E$. Thus
$|\operatorname{disc}(R)|=\prod_E f_E$.
The rational algebra is $\prod_j K_j$. Comparing its
component field-trace forms with the order $R$ changes a
discriminant by the square of a rational determinant,
proving the first assertion.

The square root of a field discriminant transforms by the
sign of the permutation of its embeddings. For a Galois
field this is the sign of the regular permutation action
of $G=\operatorname{Gal}(K/\mathbb Q)$.
For an abelian group that sign character is nontrivial
exactly when its Sylow two subgroup is nontrivial and cyclic.
Indeed an element $g$ has sign
$(-1)^{|G|-|G|/\operatorname{ord}(g)}$.
The stated condition is equivalent to $G$ having exactly
one subgroup of index two. Its fixed quadratic field
therefore gives the discriminant squareclass. Taking
absolute values gives the assertion in the imaginary case
as well.
\end{proof}

\subsection{Determinants of nonlinear rational components}
\label{paper:sec:even-component-determinants}

\begin{lemma}[The reduced-trace discriminant]
\label{paper:reduced-trace-discriminant}
Let $B$ be a central simple algebra of degree $m$ over a field $K$
of characteristic zero. The symmetric bilinear form
$\operatorname{Trd}_{B/K}(xy)$ has determinant squareclass
\[
                        (-1)^{m(m-1)/2}
                     \quad\text{in }K^\times/K^{\times2}.
\]
If $\sigma$ is an involution of the first kind on $B$, then
$\operatorname{Trd}_{B/K}(x\sigma(y))$ has square determinant.
The assertions include nonsplit algebras.
\end{lemma}
\begin{proof}
Choose a finite Galois splitting field $L/K$ and an isomorphism
$\phi:B\otimes_KL\simeq M_m(L)$. Let $P$ be its matrix from a
$K$-basis of $B$ to the matrix-unit basis. For every
$g\in\operatorname{Gal}(L/K)$, the two splitting isomorphisms
$g\phi$ and $\phi$ differ by an inner automorphism of $M_m(L)$, by Skolem--Noether
\cite[Main Theorem~7.7.1 and Corollary~7.7.4]{VoightQuaternion}.
Conjugation by an invertible matrix $U$ has linear determinant
$(\det U)^m(\det U^{-1})^m=1$. Hence $g(\det P)=\det P$, so
$\det P\in K^\times$. The reduced-trace Gram determinant therefore
differs from that in the split algebra by a square in $K$.
The matrix-unit basis has $m$ singleton diagonal entries one and
$m(m-1)/2$ off-diagonal pairs with Gram matrix
$\left(\begin{smallmatrix}0&1\\1&0\end{smallmatrix}\right)$.
This proves the first assertion.

Composing the second argument with $\sigma$ multiplies the
Gram determinant by $\det(\sigma)$. Over a splitting field an
orthogonal involution has skew-space dimension $m(m-1)/2$;
a symplectic involution has skew-space dimension $m(m+1)/2$,
and in that case $m$ is even. Thus the resulting squareclass
is one in the orthogonal case and $(-1)^{m^2}=1$ in the
symplectic case.
\end{proof}

\subsection{The noncommutative rank-thirteen algebra}
\label{paper:sec:noncommutative-thirteen}

Write $R=K_0(\mathcal C)$ and suppose its complexified algebra has
\begin{equation}\label{paper:rank-thirteen-block-decomposition}
                  R_{\mathbb C}\simeq\mathbb C^9\oplus M_2(\mathbb C).
\end{equation}
Let $E$ be its unique two-dimensional irreducible representation,
let $f=f_E$, and let $q$ count the nonreal conjugate pairs among
its nine linear characters. The following statements separate
rational algebra constraints from the two possible actions of
duality on the matrix block.

\begin{proposition}[Regular-trace discriminants with a degree-two block]
\label{paper:rank-thirteen-discriminant}
The character of $E$ is rational, and the corresponding simple
summand of $R_{\mathbb Q}$ is a central simple algebra of degree
two over $\mathbb Q$. It may be split or quaternionic.
If $\operatorname{disc}(R)$ denotes the discriminant of the
regular-trace form in the simple-object basis, then
\begin{equation}\label{paper:regular-trace-discriminant}
 |\operatorname{disc}(R)|
   =\prod_H(m_Hf_H)^{m_H^2}
   =(2f)^4\prod_{h\text{ linear}}f_h.
\end{equation}
Let $K_1,\ldots,K_s$ be the fields of the distinct Galois orbits
of linear characters. Then, in positive rational numbers modulo
squares,
\begin{equation}\label{paper:rank-thirteen-field-squareclass}
       \prod_{h\text{ linear}}f_h
       \equiv\prod_{j=1}^s|\operatorname{disc}(K_j)|.
\end{equation}
In a weighted canonical degree profile, removing the two copies
of the degree attached to $E$ therefore leaves the same target
squareclass as the full profile.
\end{proposition}

\begin{proof}
Galois conjugation preserves representation degree and $E$ is
the only representation of degree two. Its character is therefore
rational. The associated rational simple algebra has center
$\mathbb Q$ and becomes $M_2(\mathbb C)$ after complexification,
so it is a central simple algebra of degree two.

Let $\tau$ be the coefficient-of-unit trace, and let
$\lambda(x)=\operatorname{Tr}(L_x)$ be the regular trace.
The character decompositions give
\[
 \tau=\sum_H\frac{\chi_H}{f_H},\qquad
 \lambda=\sum_Hm_H\chi_H=\tau(z\,{-}),\qquad
 z=\sum_Hm_Hf_He_H.
\]
The Gram matrix of $\tau$ in the simple basis is the duality
permutation matrix, of determinant $1$ or $-1$.
On the $m_H^2$-dimensional matrix block, $L_z$ acts by
$m_Hf_H$. Taking determinants proves
\eqref{paper:regular-trace-discriminant}.

Each linear Galois orbit gives the number field $K_j$, with
its usual field-trace form. On a degree-two central simple
algebra the regular trace is twice the reduced trace.
For a quaternion presentation with basis $1,i,j,ij$ and
$i^2=a$, $j^2=b$, $ij=-ji$, its Gram matrix is
\[
           \operatorname{diag}(4,4a,4b,-4ab),
\]
with determinant $-256a^2b^2$. This also covers the split
algebra, and its absolute squareclass is one.
Changes of rational bases, including passage from these
component orders to $R$, change a discriminant by a square.
This proves \eqref{paper:rank-thirteen-field-squareclass}.
Finally the two weighted copies contribute $f^2$, itself a
square, explaining the last assertion.
\end{proof}

\begin{proposition}[Duality type and exact matrix-character moments]
\label{paper:rank-thirteen-duality}
The complex-linear duality anti-involution on the unique
matrix block is either orthogonal or symplectic.
The number $b$ of non-self-dual simple pairs in $\mathcal C$ is
\[
              b=q+1\quad\text{in the orthogonal case},
 \qquad       b=q+3\quad\text{in the symplectic case}.
\]
Consequently the disjoint-support theorem requires,
respectively,
\[
                         q+1\ge2q_6,
 \qquad                   q+3\ge2q_6,
\]
where $q_6$ counts the nonreal linear codegree-six pairs.

In the symplectic case, for a simple $X$ put
$s_X=\chi_E(X)$ and $t_X=\chi_E(XX^*)$. Then
\begin{equation}\label{paper:symplectic-character-moments}
 s_X\in\mathbb Z,\qquad t_X\in2\mathbb Z_{\ge0},
 \qquad 2t_X-s_X^2\ge0.
\end{equation}
If $X=X^*$, the sharper identities are
\[
                       s_X\in2\mathbb Z,
 \qquad                 t_X=s_X^2/2.
\]
In this case $q\le2$.
\end{proposition}

\begin{proof}
The trace of duality on the linear-character summand is
$9-2q$. Its trace on $M_2(\mathbb C)$ is $2$ for an
orthogonal anti-involution and $-2$ for a symplectic one.
The same operator in the simple basis fixes exactly the
self-dual labels. Thus their number is $11-2q$ or $7-2q$,
respectively, giving the two formulas for $b$.
The inequalities follow from
Theorem~\ref{paper:codegree-six-disjoint-supports}.
Since the unit is self-dual, initially $q\le3$ in the
symplectic case. If $q=3$, the unit would be the only
self-dual source label, making
$\operatorname{FPdim}(\mathcal C)=1+2\sum d_X^2$ odd.
This contradicts its even value $5400$, so $q\le2$.

A symplectic anti-involution on $M_2$ is its adjugate map
$A\mapsto\operatorname{tr}(A)I-A$.
Take $E$ to preserve the positive $*$-structure and write
$A=E(X)$. On a basis element the complex-linear duality and
the positive $*$-operation agree, so
\[
 A^\dagger=s_XI-A,\qquad AA^\dagger=\det(A)I.
\]
The rational character is integral, giving $s_X\in\mathbb Z$.
The determinant is rational by the trace formula for a
$2$ by $2$ matrix, and is an algebraic integer because the
eigenvalues of $E(X)$ occur in an integral regular
multiplication matrix. Thus $\det(A)\in\mathbb Z$.
It is nonnegative by the displayed identity, and
$t_X=2\det(A)$. Writing
$A=(s_X/2)I+B$ with $B^\dagger=-B$ and
$\operatorname{tr}(B)=0$ gives $2t_X-s_X^2\ge0$.
If $X$ is self-dual, then $A=A^\dagger$ and consequently
$A=(s_X/2)I$. This rational scalar is an algebraic integer,
so $s_X$ is even and $t_X=s_X^2/2$.
\end{proof}

\begin{proposition}[Exactly one nonreal linear pair, of codegree six]
\label{paper:rank-thirteen-one-six-pair}
Suppose $q=q_6=1$. Then $f$ is even in either duality type.
In the symplectic case one has the sharper restriction
\[
                         f\in\{8,72\},
\]
so the canonical multiplicity-two degree is $675$ or $75$.
\end{proposition}

\begin{proof}
Let $h,\overline h$ be the nonreal linear pair, and choose
one of its two supported source dual pairs. Its difference
$a=X-X^*$ satisfies $h(a)=\pm\sqrt{-3}$.
The two linear characters therefore contribute
$2\cdot3/6=1$ to $\tau(a^*a)=2$.
All real linear characters vanish on $a$, so
\[
              \operatorname{tr}(E(a)^\dagger E(a))=f.
\]
The matrix $E(a)$ is trace zero and skew-Hermitian, hence
$E(a)^2=-\det(E(a))I$ and $\det(E(a))=f/2$.
This determinant is a rational algebraic integer, proving
that $f$ is even.

In the symplectic case there are four source dual pairs.
The two unsupported differences $u,v$ vanish on all linear
blocks. Their based-ring squared norms are two, so the reduced
trace formula gives quaternion norms $f,f$. The matrix-block projection $b=e_Ea$ of a supported
difference has norm $f/2$. Based-ring orthogonality makes
$u,v,b$ mutually orthogonal pure quaternion elements.
The original integral product $uva$ equals $uvb$, since
$u,v$ vanish on every linear block. Consequently
\[
 uva=\pm f\sqrt{f/2}\,e_E,
 \qquad \tau(uva)=\pm2\sqrt{f/2}\in\mathbb Z.
\]
Thus $f/2$ is an integer square. On either unsupported pair,
$f=2t_X-s_X^2$ with $t_X$ even, so $f\equiv0$ or $3\pmod4$.
Writing $f=2k^2$ shows that $k$ is even. Among the available
source codegrees, the numbers of the form $8\ell^2$ are
precisely $8$ and $72$.
\end{proof}

\begin{corollary}[Symplectic states whose nonreal linear pairs all have codegree six]
\label{paper:symplectic-only-six-pairs}
Suppose the matrix-block duality is symplectic and $q=q_6$.
Then the only possible cases are
\[
\begin{array}{c|c|c}
 q=q_6&f&\text{canonical multiplicity-two degree}\\\hline
 0\text{ or }2&36,100&150,54\\
 1&8,72&675,75.
\end{array}
\]
The orthogonal case with $q=q_6=2$ is impossible.
\end{corollary}

\begin{proof}
The canonical budget gives $q_6\le2$.
The cases $q=0,1$ follow from
Propositions~\ref{paper:symplectic-real-linear-codegree} and
\ref{paper:rank-thirteen-one-six-pair}, respectively.
For $q=q_6=2$, the disjoint-support theorem gives four
supported source dual pairs. Symplectic duality gives five
source pairs in total, leaving one unsupported pair.
The matrix-block projections of differences from the two
disjoint supports have quaternion norms $f/2,f/2$ and are
orthogonal. The unsupported difference lies entirely in the
matrix block, has norm $f$, and is orthogonal to both.
The original three differences are integral based-ring
elements; their product lies in the matrix block because
one factor vanishes on all linear characters. Its scalar
value is $\pm f\sqrt f/2$, so its unit coefficient is
$\pm\sqrt f\in\mathbb Z$.
Thus $f$ is a square. The unsupported difference gives
$f=2t_X-s_X^2$ with $t_X$ even, so $4\mid f$ as before.
The allowed codegrees are therefore $36,100$.
For orthogonal duality the source would have $q+1=3$
dual pairs, fewer than the four required by disjoint support.
\end{proof}

\begin{proposition}[The symplectic case with real linear characters]
\label{paper:symplectic-real-linear-codegree}
Assume the symplectic case and $q=0$, so every linear
character is real valued. Then
\[
                         f\in\{36,100\}.
\]
Equivalently the canonical constituent of multiplicity two
has degree $150$ or $54$.
Moreover $s_X$ is even for every simple $X$.
The linear characters are not required to be rational.
\end{proposition}

\begin{proof}
Real linear characters give \(b=3\).  The normalized integral
symplectic tensor has a single coordinate, whose square is \(f/4\),
by Proposition~\ref{paper:source:reduced-symplectic-tensor}.  Thus
\(f=4z^2\); the available formal-codegree palette gives exactly
\(36\) and \(100\).  The same proposition gives even matrix
character values on every source label.
\end{proof}

\subsection{Rational quaternion alternatives and rank-fourteen parity}
\label{paper:sec:rational-matrix-blocks}

For the rank-fourteen algebra shape
\[
                    R_{\mathbb C}\simeq\mathbb C^{10}\oplus M_2(\mathbb C),
\]
the degree-two irreducible representation $E$ is again unique and has
rational character. The weighted canonical length is twelve. Thus its
eleven nonunit weighted canonical degrees form a nonincreasing list
from the established palette, with sum $5399$; one repeated degree
is selected for the unique multiplicity-two constituent.
There are exactly $6\,547$ such lists and $12\,921$ possible selections.
These are necessary numerical data, not fusion rings.
If $q$ counts the nonreal pairs among the ten linear characters,
then $b=q+1$ source dual pairs occur for orthogonal duality and
$b=q+3$ for symplectic duality. Consequently
\[
 0\le q\le4\quad\text{(orthogonal)},\qquad
 0\le q\le3\quad\text{(symplectic)}.
\]
In the latter case the rank-fourteen source has $8-2q$ self-dual
labels. In particular, the rank-thirteen restriction $q\le2$
must not be reused at rank fourteen.

\begin{theorem}[Ramification of a rational degree-two component]
\label{paper:rational-matrix-ramification}
Let $E$ be any rational irreducible degree-two character of an
integral fusion ring, of integer formal codegree $f$.
Its rational central simple algebra $B$ is unramified at every
finite prime not dividing $f$.
Under a positive orthogonal involution $B$ splits at the real
place; under a positive symplectic involution it is ramified
there. In particular, if $f$ is a prime power and the involution
is orthogonal, $B\simeq M_2(\mathbb Q)$.
\end{theorem}

\begin{proof}
The rational primitive central idempotent is
\[
                 e_E=\frac1f\sum_X\chi_E(X^*)X.
\]
For $p\nmid f$ it belongs to $R\otimes\mathbb Z_p$ and gives a
direct-sum decomposition into orders. The coefficient-of-unit
pairing $\tau(xy)$ is unimodular in the simple basis, and the
two central summands are orthogonal for this pairing. Its
restriction to $e_E(R\otimes\mathbb Z_p)$ is therefore
unimodular. On this component the reduced trace is exactly
$f\tau$, so its reduced-trace pairing is unimodular as well.

A maximal order in a ramified quaternion algebra over
$\mathbb Q_p$ has reduced discriminant $p$, hence its trace
pairing has discriminant valuation two, including when
$p=2$; see \cite[Theorem~15.5.5 and Section~23.2.5]{VoightQuaternion}.
Every order is contained in a maximal order, and passing to a
suborder multiplies the trace discriminant by the square of
its index. This excludes ramification at $p$.

A positive involution on the split real degree-two algebra is
orthogonal, whereas on the real Hamilton algebra it is
symplectic. Equivalently, the standard symplectic involution
has positive reduced norm exactly in the definite case; an
orthogonal involution on the Hamilton algebra has a nonzero
pure fixed element whose square is negative, violating
positivity. Finally a rational quaternion ramification set
has even cardinality. An orthogonal block with finite
ramification contained in a singleton therefore has empty
ramification set and is split.
\end{proof}

\begin{corollary}[Finite quaternion choices in the dimension-$5400$ problem]
\label{paper:finite-quaternion-trace-forms}
For a rational degree-two component in the present problem,
all finite ramification primes belong to $\{2,3,5\}$ and must
divide its formal codegree $f$.
In the symplectic case the complete possible positive block
forms are the applicable rows of
\[
\begin{array}{c|c|c}
\text{finite ramification}&B& (f/2)\,\tau(xx^*)\\\hline
\{2\}&(-1,-1)&\langle1,1,1,1\rangle\\
\{3\}&(-1,-3)&\langle1,1,3,3\rangle\\
\{5\}&(-2,-5)&\langle1,2,5,10\rangle\\
\{2,3,5\}&(-3,-10)&\langle1,3,10,30\rangle.
\end{array}
\]
A row is applicable only if every listed prime divides $f$.
In the orthogonal case the finite ramification set is one of
$\varnothing,\{2,3\},\{2,5\},\{3,5\}$, again contained in
the prime divisors of $f$.
If the matrix skew line has positive norm squareclass $\delta$,
any prime in the ramification set must be nonsplit in
$\mathbb Q(\sqrt{-\delta})$.
For each remaining quaternion class represented by
$(-\delta,c)$ with $c>0$, its positive block form is
\[
                 \frac2f\langle1,\delta,c,\delta c\rangle.
\]
Changing $c$ within this quaternion class does not change the
rational isometry class of the displayed form.
\end{corollary}

\begin{proof}
The source codegree palette has prime divisors only two,
three and five. Apply
Theorem~\ref{paper:rational-matrix-ramification}, positivity at
infinity, and the even-cardinality rule for quaternion
ramification. The listed presentations have the stated
ramification sets by the local Hilbert-symbol formula.
Their norm forms, scaled by $2/f$, give the symplectic table.

In the orthogonal case a rational skew generator of square
$-\delta$ embeds the indicated imaginary quadratic field.
It cannot split at a ramified quaternion place, because a
division algebra contains no nontrivial idempotent.
A rational trace-zero fixed element gives a presentation
$(-\delta,c)$, with $c>0$ by positivity. The basis
$1,u,v,uv$ gives the displayed positive form.
Two parameters defining the same quaternion algebra with
this fixed quadratic subfield differ by a norm from
$\mathbb Q(\sqrt{-\delta})$; multiplication by an element
of that field identifies the corresponding binary fixed
planes. This proves independence of the parameter choice.
\end{proof}

\begin{proposition}[Rank-independent restrictions with only one matrix block]
\label{paper:one-matrix-rank-independent-restrictions}
Suppose more generally that
$R_{\mathbb C}\simeq\mathbb C^a\oplus M_2(\mathbb C)$
in the present dimension-$5400$ simple-source problem.
The following restrictions depend on the nonreal linear
characters, not on the total source rank.
\begin{enumerate}
\item If all nonreal linear pairs have formal codegree six,
so $q=q_6\le2$, symplectic duality requires
\[
 f_E\in\{36,100\}\quad(q=0),\qquad
 f_E\in\{8,72\}\quad(q=1).
\]
More generally $q=q_6=2$ would force $16\mid f_E$;
it is impossible here because $f_E\mid5400$.
\item Suppose duality on the matrix block is orthogonal and
the nonreal linear characters are either a single pair
$h,\bar h$ of formal codegree six, or this pair together
with one formal-codegree-eight pair of full field
$\mathbb Q(\sqrt{-2})$.
Then $f_E$ is even and the two source pairs supported by $h$
can be oriented and represented by $T,U$ so that
\[
 (T-T^*)(U-U^*)=\frac{f_E}{2}e_E-3(e_h+e_{\bar h}).
\]
In particular
\[
 \chi_E(X)\equiv h(X)+\bar h(X)\pmod2
 \qquad\text{for every source simple }X.
\]
Neither configuration is possible if $f_E=6$.
\end{enumerate}
\end{proposition}

\begin{proof}
The trace of duality gives $b=q+3$ in the symplectic case
and $b=q+1$ in the orthogonal case, independently of $a$.
This is the only source-pair count used in the relevant
arguments of Proposition~\ref{paper:rank-thirteen-discriminant} and Corollary~\ref{paper:orthogonal-six-finite-forms}, with their rank bounds
made explicit.

For symplectic duality, the skew matrix space is the
three-dimensional pure quaternion space. If $q=q_6=0$,
the three source-pair differences lie entirely in that
space, are orthogonal, and have quaternion norm $f_E$.
Their integral product has unit coefficient
$\pm2\sqrt{f_E}$, forcing $f_E$ to be a square.
On any of these pairs, the identity
$f_E=2t_X-s_X^2$ and the evenness of $t_X$ then force
$4\mid f_E$. If $q=q_6=1$, the two unsupported pair
differences have norms $f_E,f_E$, while the matrix
projection of a supported difference has norm $f_E/2$.
The unit coefficient of their integral product is
$\pm2\sqrt{f_E/2}$. Hence $f_E/2$ is an integer square,
and the unsupported-pair identity further makes
$f_E=8\ell^2$ for an integer $\ell$.
If $q=q_6=2$, use one difference from each of the two
disjoint supported sets and the one unsupported pair.
Their matrix norms are $f_E/2,f_E/2,f_E$, respectively;
the integral triple product has unit coefficient
$\pm\sqrt{f_E}$. Thus $f_E$ is a square divisible by four.
Inspecting the already established formal-codegree palette
gives exactly the displayed possibilities. No count of
self-dual labels or total-rank parity enters these proofs.
There is a further restriction in the two-pair case,
which does not require inspection of the codegree palette:
on a supported pair the skew norm is $f_E/2$, hence
$f_E/2=2t_X-s_X^2\equiv0$ or $3\pmod4$.
We already know that $f_E$ is a square divisible by four.
Consequently $f_E/2$ is even, and the displayed congruence
forces $v_2(f_E)\ge3$. Its valuation is even, because
$f_E$ is a square, so in fact $16\mid f_E$.
This contradicts $v_2(5400)=3$ and $f_E\mid5400$.
Thus $q=q_6=2$ is excluded altogether. The finite scan
retains these states until its later row gates, which
only weakens its preliminary restrictions.

For the orthogonal assertion, orient $T,U$ so that
$h(T-T^*)=h(U-U^*)=\sqrt{-3}$.
If $h,\bar h$ is the only nonreal pair, there are precisely
two source pairs. The $h,\bar h$ skew projection on their
normalized differences is
$\frac12\left(\begin{smallmatrix}1&1\\1&1\end{smallmatrix}\right)$.
The matrix skew projection is its complement, so the two
matrix images are opposite, square to $-(f_E/2)I$, and
have product $(f_E/2)I$.
Their reduced norms $f_E/2$ are rational algebraic integers,
so $f_E$ is even.

With the additional $\mathbb Q(\sqrt{-2})$ pair, its
integral values $u+v\sqrt{-2}$ and codegree eight force
exactly one source pair with nonzero imaginary part.
It fills the entire one-coordinate skew projection there.
Orthogonality makes this pair disjoint from $T,U$.
Since now $b=3$, its matrix skew image is zero and the
same two-coordinate calculation applies. All real linear
characters vanish on pair differences. Therefore the
displayed product has the stated values on every simple
algebra summand, proving the identity. Its coefficient
at $X$ is
$\frac12(\chi_E(X^*)-h(X^*)-\bar h(X^*))$; integrality
proves the parity assertion.

Finally set $f_E=6$. The parity assertion and
$\sum_X\chi_E(X)^2=12$ force trace $\pm1$ on the two
supported pairs, trace $2\eta$ on exactly one off-unit
self-dual label $R$, and zero on all other nonunit labels.
Let $S$ be the real off-unit label of $h$, with value
$\varepsilon$. If $S=R$, orthogonality again forces a
nonunit dimension to be $1$ or $-1$, impossible here.
Otherwise character and dimension orthogonality give
\[
 3(e_E+e_h+e_{\bar h})=2\mathbf1+\eta R+\varepsilon S,
 \qquad 2+\eta d_R+\varepsilon d_S=0.
\]
Thus, naming the smaller dimension $A$ and the larger $B$,
$D=B-A=2\mathbf1-3(e_E+e_h+e_{\bar h})$ satisfies
$D^2=D+2\mathbf1$ and $d_B-d_A=2$.
Lemma~\ref{paper:quadratic-simple-difference} supplies a new
self-dual $Q$ with $DQ=A+B$, forcing
$\chi_E(Q)=\pm2$ beyond an already exhausted character
square budget. This contradiction proves the exclusion
at every source rank with either listed nonreal configuration.
\end{proof}

\begin{theorem}[The universal symplectic codegree-six support]
\label{paper:universal-symplectic-six-support}
Suppose a rational degree-two irreducible representation $E$
of $R$ has formal codegree six and positive symplectic
involution. Then exactly twelve source simple labels have
nonzero image under $E$:
\begin{itemize}
\item the unit and one off-unit self-dual label $R$, with
$E(R)=\eta I$, $\eta\in\{1,-1\}$;
\item two source dual pairs with traces $\pm1$;
\item three further source dual pairs with trace zero.
\end{itemize}
Every nonzero basis image is unitary of determinant one.
All other source simples have image zero.
The rational component is $(-3,-1)_{\mathbb Q}$, with
positive trace form $\frac13\langle1,3,1,3\rangle$.
The result does not require that $E$ be the only nonlinear
representation or that the source rank equal thirteen.
\end{theorem}

\begin{proof}
Write $s_X=\chi_E(X)$ and
$t_X=\chi_E(XX^*)$. The same coefficient trace and
Schur-orthogonality identities as before give
\[
          \sum_Xs_X^2=12,\qquad \sum_Xt_X=24,
          \qquad 0\le t_X<6.
\]
Symplectic positivity gives $t_X\in2\mathbb Z$.
For a source dual-pair difference $a_i=X_i-X_i^*$,
put $K_{ij}=-\frac12\operatorname{tr}(E(a_i)E(a_j))$.
Theorem~\ref{paper:source:skew-projection-lattice} applies to this
individual rational block even if other nonlinear components
are present: $K$ is integral and
positive semidefinite, $K^2=6K$, and its rank is three.
Its diagonal entries belong to $\{0,3,4\}$ and sum to
eighteen. There are only two possible multisets of
nonzero entries: six threes, or two threes and three
fours. A diagonal-three pair has trace $\pm1$, so six
such pairs would already contribute twelve to the
character-square budget in addition to the unit's four.
The six-three alternative is therefore impossible.

A zero diagonal entry of a positive semidefinite matrix
has zero row and column. Restrict $K$ to its five
nonzero positions. Its complementary rank-two matrix
has diagonal $(3,3,2,2,2)$. Call this complement $H=6I-K$. Its nonzero principal
two-minors have a common rational squareclass: they are
the squared two-volumes in one rational plane. A minor
on diagonal entries three and two is $6-H_{ij}^2$,
hence is one of $6,5,2$, all nonzero. A nonzero minor
on two diagonal entries two would be $4$ or $3$, of
different squareclass. Thus the three norm-two vectors
are collinear. Their mutual entries have absolute value
two and exhaust their off-diagonal square budgets in
$H^2=6H$, so they are orthogonal to the other two
vectors. The two remaining vectors are also collinear.
Consequently $H$ has blocks $3vv^{\mathsf t}$ and
$2ww^{\mathsf t}$, with $v\in\{1,-1\}^2$ and
$w\in\{1,-1\}^3$.
The two diagonal-three pairs have traces $\pm1$ and
$t_X=2$. A diagonal-four pair could have trace $\pm2$
only at cost eight in the character-square sum, which
exceeds the remaining budget four; hence these three
pairs have trace zero and $t_X=2$.

If a further source pair has zero skew image, its two
images are equal self-adjoint scalars. Such a rational
scalar is an algebraic integer and hence an integer.
A nonzero scalar pair would cost at least eight in
the remaining square budget four, so it must vanish.
Self-dual simples also have integral scalar images.
Exactly one off-unit self-dual label has scalar $\pm1$,
and every other self-dual image is zero. This proves
the asserted twelve-label support. The five-pair Gram block supplies a pure quaternion
generator of square $-3$ from the first block. After
changing coordinate signs, half the sum of two
generators from the second block has square $-1$.
The two blocks are orthogonal, so these generators
anticommute. They give the stated quaternion algebra
and positive form.
\end{proof}

\begin{corollary}[A nonreal linear codegree-six pair needs rank at least sixteen]
\label{paper:symplectic-six-linear-six-threshold}
In the present simple-source problem, under
Theorem~\ref{paper:universal-symplectic-six-support},
at most one nonreal linear codegree-six pair can occur.
If one occurs, its two imaginary-supported source pairs
have zero image under $E$. Its real off-unit label is
$R$, with value $-\eta$. Consequently the source rank
is at least sixteen. In particular there is no such
linear pair at source rank fourteen.
\end{corollary}

\begin{proof}
The imaginary-part vector of a linear character is
orthogonal to every matrix coefficient of the distinct
irreducible $E$. The two-position support vector of
an Eisenstein codegree-six character must therefore
lie in the kernel of the skew Gram matrix $K$.
On the five active positions this kernel is the direct
sum of the line supported on the two diagonal-three
positions and the line supported on all three
diagonal-four positions. Outside those positions
$K$ is zero. A vector with exactly two nonzero
positions can consequently use either precisely the
two diagonal-three positions or two zero-image pairs.

The first alternative is impossible. If the character's
real extra label is $R$, character and dimension
orthogonality force a supported nonunit dimension to
be $1$ or $-1$, contrary to the established lower bound
$d_X\ge3$ for nonunit simples. Otherwise call the real label $S$.
The same orthogonality identities make the two real
pair numerators opposite to the matrix traces, and
\[
 2+\eta d_R+\varepsilon d_S=0,\qquad
 3(e_E+e_h+e_{\bar h})=2+\eta R+\varepsilon S.
\]
Naming the smaller dimension $A$ and the larger $B$
gives the central difference $D=B-A=2-3(e_E+e_h+e_{\bar h})$.
Lemma~\ref{paper:quadratic-simple-difference} then forces
a new label $Q$ with $\chi_E(Q)=\pm2$, contradicting
the exhausted character-square budget twelve.

Thus the two supported pairs are zero-image pairs.
Character orthogonality reduces to
$2+2\eta\varepsilon\delta_{R,S}=0$, forcing
$S=R$ and $\varepsilon=-\eta$.
These four zero-image labels are in addition to the
twelve nonzero-image labels. Two distinct nonreal
codegree-six pairs would have disjoint pair supports
by Theorem~\ref{paper:codegree-six-disjoint-supports},
but would share $R$ with the same real sign. Their
real-part inner product would be two, impossible.
\end{proof}

\begin{corollary}[Two symplectic codegree-six blocks need rank at least eighteen]
\label{paper:two-symplectic-six-rank-eighteen}
If two distinct rational degree-two irreducible components
have symplectic duality and formal codegree six, the source
rank is at least eighteen. This conclusion allows other
nonlinear components.
\end{corollary}

\begin{proof}
Let $K_1,K_2$ be their skew Gram matrices. The normalized
matrices $K_1/6,K_2/6$ are orthogonal projections onto
orthogonal simple-algebra summands. Consequently
$(K_1)_{ii}+(K_2)_{ii}\le6$ at every source pair.
Each active diagonal multiset is $(3,3,4,4,4)$ by
Theorem~\ref{paper:universal-symplectic-six-support}.
A position with diagonal four for one component must
therefore have diagonal zero for the other. Their two
five-position supports can intersect in at most two
positions, and at least eight source dual pairs occur.
The unit and the off-unit self-dual scalar label supplied
by either component are two further labels. Hence the
rank is at least $2\cdot8+2=18$.
\end{proof}

\subsection{Local forms of a unique rational matrix component}
\label{paper:sec:trace-forms}

\begin{proposition}[Trace forms at the unique matrix block]
\label{paper:matrix-block-trace-forms}
Suppose $R_{\mathbb C}\simeq\mathbb C^{r-4}\oplus M_2(\mathbb C)$.
Let $L,L^+,L^-$ be the positive rational trace form and its fixed and
skew restrictions on the linear-character algebra. Let $q$ count the
nonreal linear-character pairs and let $f$ be the matrix formal codegree.
In the symplectic case $b=q+3$ and
\begin{equation}\label{paper:symplectic-fixed-trace-test}
 L^+\perp\langle2/f\rangle\simeq I_{r-2b}\perp2I_b.
\end{equation}
In the orthogonal case $b=q+1$. The skew squareclass has positive
squarefree representative
\begin{equation}\label{paper:orthogonal-delta}
 \delta=2^{b-1}f/\det L^-\quad\text{in }\mathbb Q^\times/(\mathbb Q^\times)^2,
\end{equation}
and
\begin{equation}\label{paper:orthogonal-skew-test}
 L^-\perp\langle2\delta/f\rangle\simeq2I_b.
\end{equation}
At every prime $p$ for which $-\delta$ is a square in $\mathbb Q_p$,
\begin{equation}\label{paper:orthogonal-local-matrix-test}
 \epsilon_p(L)(\delta,-1)_p=1.
\end{equation}
Every full candidate matrix form $B$ must also satisfy
$L\perp B\simeq I_r$.
\end{proposition}
\begin{proof}
The matrix contribution to the coefficient trace is its reduced trace
divided by $f$. Symplectic duality has scalar fixed space, whose form is
$\langle2/f\rangle$. The positive source form and its duality eigenspaces
give the first identity. In the orthogonal case choose a rational skew
generator $u$ with $u^2=-\delta$. Positivity gives $\delta>0$ and its
norm is $2\delta/f$. Comparing the skew form with $2I_b$, and then
their determinants, gives the next two identities.

Choose a rational fixed traceless vector $v$ with $v^2=c>0$.
It anticommutes with $u$, so the positive matrix form is
$\frac2f\langle1,\delta,c,\delta c\rangle$. Its determinant is a
square and its local Hasse invariant is
$(\delta,-1)_p(-\delta,c)_p$. The full form is $I_r$, so the
determinant of $L$ is also a square and
$(-\delta,c)_p=\epsilon_p(L)(\delta,-1)_p$.
If $-\delta$ is a local square the left side is one, proving the
asserted necessary local condition. No splitting assumption is used.
\end{proof}

\subsection{Integral skew projections and a strict trace--norm defect}
\label{paper:sec:skew-and-trace}

The coefficient trace on a fusion ring is denoted by $\tau$; thus
$\tau(XY^*)=\delta_{X,Y}$ on the distinguished basis. All statements in
this subsection concern this actual based-ring trace.

\begin{theorem}[An integral trace--norm bound and its equality case]
\label{paper:trace-norm-defect}
Let $\mathcal B$ be a Galois-stable collection of irreducible
representations and let $A_E=E(X)$ be the images of a fusion-basis
element. With adjoints taken in positive realizations, put
\[
 s=\sum_{E\in\mathcal B}\Tr A_E,\qquad
 t=\sum_{E\in\mathcal B}\Tr(A_EA_E^*).
\]
Then $s\in\mathbb Z$, $t\in\mathbb Z_{\ge0}$, and
\begin{equation}\label{paper:trace-norm-bound}
                            |s|\le t.
\end{equation}
If equality holds, either all images vanish, or all of them are
orthogonal projections multiplied by the same sign. In particular,
$A_E^3=A_E$ for every $E$ in that block.

Under the center hypotheses at $5400^2$, every nonunit source row has
at least one non-Perron Galois-stable block for which $t>|s|$.
\end{theorem}

\begin{proof}
The product of the characteristic polynomials over $\mathcal B$ is
monic integral. For a nonzero irreducible root orbit of degree $h$,
write $S$ for the sum of its roots and $T$ for the sum of their squared
absolute values. Its nonzero algebraic-integer norm gives $T\ge h$ by
arithmetic--geometric mean, whence
$|S|\le\sqrt{hT}\le T$. Summing over root orbits, with multiplicity,
and using the Schur triangular form of each $A_E$ gives $|s|\le t$.
Equality forces equality at every step. Each nonzero root orbit then
consists entirely of one sign $1$ or $-1$, the signs agree across the
block, and each matrix is normal. This proves the equality assertion.

If equality held in every non-Perron block, the semisimple algebra
identity would be
$X^3-X=(d_X^3-d_X)e_{\rm FP}$. Applying $\tau$ would give
$5400\mid d_X^3-d_X$. The stated source bounds are
$3\le d_X\le72$. Divisibility by $25$ and $27$ leaves only $d_X=26$
in this interval, and $26^3-26$ is not divisible by $8$. Hence no such
row exists. This uses equality in every block, not equality in an
individual block alone.
\end{proof}

\begin{lemma}[Coupled degree-two word traces]
\label{paper:matrix-commutator-traces}
For a complex $2\times2$ matrix put
$\alpha=\Tr A$, $\beta=\det A$ and $\gamma=\Tr(AA^*)$. Then
\begin{align}
 G(A)&=\Tr((AA^*)^2)-\Tr(A^2(A^*)^2)\notag\\
 &=\gamma^2-|\alpha|^2\gamma
       +2\operatorname{Re}(\alpha^2\overline\beta)-4|\beta|^2
       =\tfrac12\Tr([A,A^*]^2)\ge0.
 \label{paper:general-two-matrix-commutator}
\end{align}
If these are all the nonlinear representation blocks of a fusion
ring, the necessary condition is
\[
                     \sum_E G(E(X))/f_E\in\mathbb Z_{\ge0}.
\]
In particular, for a quadratic orbit with common codegree six it is
$\Tr_{K/\mathbb Q}G/6\in\mathbb Z_{\ge0}$; in an imaginary quadratic
field this is $G/3\in\mathbb Z_{\ge0}$. Contributions from distinct
nonlinear blocks must be coupled.
\end{lemma}

\begin{proof}
Cayley--Hamilton gives
$A^2=\alpha A-\beta I$ and the first two expressions in
\eqref{paper:general-two-matrix-commutator}. Expanding the commutator
square proves the last expression. Apply the coefficient trace to
$(XX^*)^2-X^2(X^*)^2$; scalar blocks contribute zero. Its coefficient
of the unit is integral, and every displayed matrix contribution is
nonnegative.
\end{proof}

\begin{remark}\label{paper:trace-norm-nilpotent-caution}
The assertion $t=1\Rightarrow |s|=1$ is valid for scalar field blocks
but not for general matrix blocks: a nonzero nilpotent matrix can have
$(s,t)=(0,1)$. The matrix searches retain this possibility.
\end{remark}

\subsection{Two representation-theoretic bounds at arbitrary rank}
\label{paper:sec:odd-matrix-forms}

\begin{theorem}[The whole positive form of an odd real component]
\label{paper:odd-real-whole-form}
Let $K$ be a totally real number field, and let $B$ be a central
simple algebra of odd degree $m$ over $K$, with a positive
involution $\sigma$ of the first kind. Then $B\cong M_m(K)$ and
the quadratic form
\[
                       \operatorname{Trd}_{B/K}(x\sigma(x))
\]
is isometric over $K$ to $I_{m^2}$. Consequently, a full real
Galois orbit of odd-degree source representations of common
rational formal codegree $f$ contributes the rational positive
form
\[
       \mathop{\perp}\limits_{j=1}^{m^2}
             \left(\frac1f\operatorname{Tr}_{K/\mathbb Q}(x^2)\right).
\]
In particular, for a rational character field the whole form is
$(1/f)I_{m^2}$, independent of the positive adjoint form defining
the involution.
\end{theorem}
\begin{proof}
An involution of the first kind makes the Brauer class have order
dividing two, whereas its order divides the odd degree $m$.
Thus the algebra splits. The involution is orthogonal, and after
diagonalization is adjoint to $H=\operatorname{diag}(h_1,\ldots,h_m)$.
The trace form consists of $m$ entries one and, for $i<j$, the
pair $\langle h_i/h_j,h_j/h_i\rangle$. Each such pair has square
determinant and is isometric to $\langle c_{ij},c_{ij}\rangle$,
where $c_{ij}=h_i/h_j$. Its local Hasse invariant is
$(c_{ij},c_{ij})_v=(c_{ij},-1)_v$.

All cross terms between these pairs have trivial determinant
symbol. Their combined Hasse invariant is
\[
 \left(\prod_{i<j}h_i/h_j,-1\right)_v
       =\left(\prod_{i=1}^m h_i^{m+1-2i},-1\right)_v=1,
\]
since every exponent is even. The determinant is square and all
real signatures are positive. Hasse--Minkowski therefore gives
the claimed isometry with $I_{m^2}$ over $K$. Transfer to
$\mathbb Q$ and scaling by $1/f$ give the last assertions.
\end{proof}

\begin{proposition}[A rank bound from the integral norm budget]
\label{paper:integral-norm-rank-bound}
Let a fusion ring of rank $r$ have a Galois-stable collection
$\mathcal B$ of non-Perron irreducible representations of common
rational integral formal codegree $f$. Put
$k=\sum_{E\in\mathcal B}m_E$ and
$q=\sum_{E\in\mathcal B}m_E^2$. Then
\[
 (r-1)(f-1)\ge fq-k,\qquad
 r\ge q+1+\left\lceil\frac{q-k}{f-1}\right\rceil.
\]
For $c$ representations of common degree $m$, take $k=mc$ and
$q=m^2c$.
\end{proposition}
\begin{proof}
For every nonunit basis element $X$, the aggregate norm trace
$t_X=\sum_E\chi_E(XX^*)$ is a nonnegative integer. The separate
positive Perron contribution to $\tau(XX^*)=1$ gives $t_X/f<1$,
hence $t_X\le f-1$. Matrix-coefficient orthogonality gives
$\sum_Xt_X=fq$, including the unit contribution $k$.
Summing the pointwise bounds over the $r-1$ nonunit labels proves
the result.
\end{proof}

\subsection{A skew determinant for odd matrix components}
\label{paper:sec:odd-skew-determinant}

\begin{lemma}[Odd-degree rational skew determinant]
\label{paper:odd-rational-skew-determinant}
Let a rational irreducible source representation have odd degree $m$
and rational formal codegree $f$. Put $b=m(m-1)/2$. Its positive
coefficient-trace form has whole part isometric to $(1/f)I_{m^2}$, and the determinant
squareclass of its duality-skew part is
\[
                             (2/f)^b.
\]
The assertion does not prescribe the Hasse invariants of that skew
part separately.
\end{lemma}

\begin{proof}
The whole-form assertion is
Theorem~\ref{paper:odd-real-whole-form}. The rational simple algebra is
$M_m(\mathbb Q)$, and its positive involution is adjoint to a positive
diagonal form $\operatorname{diag}(h_1,\ldots,h_m)$.
The skew subspace has the positive diagonal form
\[
                  \frac2f\left\langle h_i/h_j:i<j\right\rangle.
\]
Its determinant is
\[
 (2/f)^b\prod_{i=1}^m h_i^{m+1-2i}.
\]
Every exponent in the product is even because $m$ is odd. The product
therefore has trivial rational squareclass, proving the result.
\end{proof}

Suppose every nonlinear character in an assignment is rational and
odd-dimensional. After choosing full fields for its linear characters,
the entire positive trace form is known, while the preceding lemma
also fixes the determinant squareclass of the whole skew form.
The positive form must be isometric to $I_r$ and the skew determinant
must equal the squareclass of $2^b$, where $b$ is the source-pair count.
These tests preserve every positive adjoint form on each odd matrix
component; no diagonal defining form has been fixed in advance.

\begin{remark}[Totally real odd-degree orbits]
\label{paper:totally-real-odd-skew-determinant}
More generally, suppose a full orbit of odd-degree $m$
representations has totally real character field $K$ of degree
$e$ and common rational formal codegree $f$. Put
$b=m(m-1)/2$. Its rational skew determinant squareclass is
\[
                \operatorname{disc}(K)^b(2/f)^{eb}.
\]
Indeed, the positive involution fixes the totally real center
$K$, so is of the first kind. The Brauer class of the associated
central simple algebra has order dividing two. Its index divides
the odd degree $m$, and its period divides its index; consequently
the Brauer class is trivial and the algebra is $M_m(K)$.
After diagonalizing the positive adjoint form over $K$, the skew
form is $(2/f)\langle h_i/h_j:i<j\rangle$. Its determinant over
$K$ is $(2/f)^b\prod_i h_i^{m+1-2i}$, where every exponent in
the product is even. Determinant transfer to $\mathbb Q$ gives
$\operatorname{disc}(K)^b$ times the field norm of this
determinant. The norm of the displayed product is a rational
square, and the norm of $(2/f)^b$ is $(2/f)^{eb}$, proving the
formula. This generalization is not used by the certified
rational-only odd-matrix computation.
\end{remark}

\subsection{Integral character-coordinate covariances}
\label{paper:character-coordinate-lattices-section}

Full character fields impose integral lattice constraints in addition to
rational trace-form isometry. These constraints can exclude a field
assignment before any candidate source row is generated.

\begin{corollary}[Small integral functional obstructions]
\label{paper:character-coordinate-short-functionals}
Use the hypotheses and notation of the character-coordinate covariance
theorem, Theorem~\ref{paper:source:character-coordinate-covariance}: the full
orbit of degree $m$ and codegree $f$ is disjoint from the Perron orbit,
$G$ is the positive field-trace Gram matrix, $H=mfG^{-1}$,
$z_X=z(\chi(X))$, $z_0=z(1)$, and $C$ represents conjugation.
Then $H$ must be an integral matrix. For
$u\in\mathbb Z^e$, put
\[
 b_X=u^{\mathsf t}z_X,\qquad
 N=u^{\mathsf t}Hu,\qquad b_{\mathbf1}=m u^{\mathsf t}z_0.
\]
The following are impossible:
\begin{enumerate}
\item a nonintegral entry of $H$;
\item a nonzero functional with $N=1$;
\item $N=2$ and $b_{\mathbf1}\ne0$, provided the source has no
nonunit invertible simple when $m=1$;
\item $C^{\mathsf t}u=-u$ and $N$ an odd integer.
\end{enumerate}
For $m\ge2$, the third obstruction needs no hypothesis about invertible
simples: its nonzero unit coordinate already has absolute value at
least two.
\end{corollary}
\begin{proof}
Every covariance entry is a sum of integer products. Moreover
\[
 \sum_X b_X^2=N,\qquad \sum_Xd_Xb_X=0.
\]
If $N=1$, exactly one $b_X$ is $1$ or $-1$, contradicting positivity
of $d_X$. If $N=2$ and the unit coordinate is nonzero, integrality
forces that coordinate to be $1$ or $-1$ and leaves exactly one other
nonzero coefficient of absolute value one. For $m=1$, the weighted
sum forces the corresponding nonunit dimension to be one. For
$m\ge2$, a nonzero unit coefficient is already incompatible with
$N=2$. Finally, an antisymmetric functional vanishes on every
selfdual label, while the two members of each nonselfdual pair
contribute $2b_X^2$. Its total squared norm is therefore even.
\end{proof}

The orbit-disjointness condition is essential. In the present integral
source problem the Perron character is rational, so every non-Perron
canonical orbit satisfies it. One cannot apply the norm-one conclusion
to a nonintegral Perron orbit, such as the real quadratic Perron orbit
of the Ising fusion ring.

The covariance identity also excludes the full fields
$\mathbb Q(\zeta_8)$ at $f=8$, $\mathbb Q(\zeta_9)$ at $f=9$,
$\mathbb Q(\zeta_{20})$ at $f=10$, and
$\mathbb Q(\zeta_{24})$ at $f=12$. The parity refinement also
excludes $\mathbb Q(\zeta_8)$ at $f=12$. These are rank-independent
consequences of the character-coordinate identities.

\paragraph{Safe coordinates from cyclotomic fields.}
For $K=\mathbb Q(\zeta_c)^H$, write its fixed space in the ambient
integral power basis of $\mathbb Q(\zeta_c)$. Rational row reduction
expresses its vectors by free ambient coordinates. Every algebraic
integer of $K$ has integral ambient coordinates, so these free
coordinates satisfy the theorem's integrality requirement. Pivot
coordinates need not be integral for arbitrary integer free vectors;
no assertion that the resulting rational basis is an integral basis
is made. The field trace Gram matrix is the ambient trace Gram
matrix divided by $|H|$. The verifier computes its inverse and the
full conjugation matrix exactly.

\paragraph{Explicit linear-field examples.}
The endpoint implementation uses only $m=1$. In the free coordinates
just described it obtains the following witnesses; $z_0=(1,0,\ldots,0)$.
\[
\begin{array}{c|c|c|c}
 f&K&H=fG^{-1}&\text{obstruction}\\\hline
 8&\mathbb Q(\zeta_{16})^{\langle7\rangle}
    &\operatorname{diag}(2,1,1,1)&N=1\\
 8&\mathbb Q(\zeta_{16})^{\langle15\rangle}
    &\operatorname{diag}(2,1,1,1)&N=1\\
 8&\mathbb Q(\zeta_{16})^{\langle9\rangle}
    &2I_4&N=2,\ b_{\mathbf1}\ne0\\
 9&\mathbb Q(\zeta_9)
    &\begin{pmatrix}2I_3&I_3\\ I_3&2I_3\end{pmatrix}
    &N=2,\ b_{\mathbf1}\ne0.
\end{array}
\]
For the first two rows use the fourth coordinate functional. For the
third use the first. For the last, either the first coordinate alone
or the difference of the first and fourth is a suitable functional.
Thus every codegree-eight quartic field is excluded. The full octic
fields $\mathbb Q(\zeta_{20})$ at $f=10$ and
$\mathbb Q(\zeta_{24})$ at $f=12$ likewise have $H_{00}=2$ and are
excluded by the third obstruction; they were correctly included in
the initial field catalogue before any such exclusions.

\paragraph{Complete finite gate at the endpoint.}
Applying the explicit-witness tests to all $40\,476$ trace-retained
commutative rank-thirteen assignments excludes $13\,235$: $8789$ by a
norm-one witness, $3825$ by an odd antisymmetric norm, and $621$ by a
norm-two witness with nonzero unit coordinate. The gate leaves
$27241$ assignments and removes all assignments on $1301$ complete
profiles. It contains the $11057$-assignment union of the two
quadratic gates and adds $2178$ further exclusions.

Every accepted exclusion stores its exact field and functional (or
nonintegral matrix entry). An independent audit checks all $277$
actual $(f,K)$ pairs, including $47$ cyclotomic fixed-field pairs by
direct polynomial multiplication, then rechecks every one of the
$13\,235$ witnesses. The short-vector search is a means of finding
witnesses; no claim of a complete classification of all short
lattices or all possible integral embeddings is needed.

\begin{corollary}[A full cyclotomic field bound]
\label{paper:full-cyclotomic-unit-coordinate}
Suppose a non-Perron linear character orbit has full field
$K=\mathbb Q(\zeta_n)$ and rational formal codegree $f$, and the
source has no nonunit invertible simple. If $\omega(n)$ is the number
of distinct prime divisors of $n$, then
\[
                    N:=\frac{2^{\omega(n)}f}{n}\in\mathbb Z_{\ge3}.
\]
If $n>2$, then $N$ is even and consequently $N\ge4$.
\end{corollary}
\begin{proof}
For each prime power $p^k\Vert n$, use the integral power basis of
$\mathbb Q(\zeta_{p^k})$. After grouping exponents modulo $p^{k-1}$,
its positive trace Gram matrix has blocks
$p^{k-1}(pI_{p-1}-J_{p-1})$, where $J$ is the all-ones matrix.
The inverse of each block is $p^{-k}(I_{p-1}+J_{p-1})$, so its
unit-coordinate diagonal entry is $2/p^k$.

The tensor product of these prime-power bases is an integral basis of
$\mathbb Q(\zeta_n)$. Indeed, the prime-power roots generate
$\mathbb Z[\zeta_n]$ by the Chinese remainder theorem, and the product
of their degrees is $\varphi(n)$. Its Gram matrix is the tensor
product of their Gram matrices. The integer functional extracting the
unit coordinate therefore has value one at the unit and squared norm
$2^{\omega(n)}f/n$. Covariance integrality makes this an integer;
the norm-one and norm-two-unit obstructions force it to be at least
three.

For the parity refinement, suppose first that $n$ has an odd prime
factor $p$. In the corresponding prime-power inverse Gram matrix,
the unit diagonal entry is $2/p^k$ and an off-diagonal entry in the
same block is $1/p^k$. Tensoring this off-diagonal entry with unit
diagonal entries in all other factors gives a covariance entry
$N/2$. Covariance integrality therefore makes $N$ even. If instead
$n=2^k>2$, the coordinate functional at exponent $2^{k-2}$ is
antisymmetric under conjugation and has covariance norm $N$.
The odd-antisymmetric-norm obstruction again makes $N$ even.
Excluding $N=2$ proves $N\ge4$.
\end{proof}

\begin{proposition}[Coordinate roots have type $A$]
\label{paper:coordinate-roots-type-a}
Let a full non-Perron Galois orbit of irreducible characters have common
degree $m$, common rational formal codegree $f$, and full character
field $K$ of degree $k$.  Choose a rational basis $B=(B_1,\ldots,B_k)$
in which every algebraic integer of $K$ has integral coordinates.
Write $z_X\in\mathbb Z^k$ for the coordinates of $\chi(X)$, and put
\[
 G_{ij}=\Tr_{K/\mathbb Q}(B_i\overline{B_j}),\qquad
 H=mfG^{-1},\qquad
 \Phi(u)=(u^{\mathsf T}z_X)_X.
\]
Then $\Phi$ embeds the integral lattice with Gram matrix $H$ into
the standard Euclidean lattice on the source labels, orthogonally
to the positive Perron vector.  Every vector of squared norm two
has image
\[
                 \pm(e_X-e_Y),\qquad d_X=d_Y.
\]
The sublattice generated by all such vectors is an orthogonal sum
of root lattices $A_s$.  An $A_s$ component uses $s+1$ distinct
source labels of the same dimension; different components have
disjoint supports.

If $K$ is totally real, every label in these supports is selfdual.
In general, let $C$ be the coordinate matrix of complex conjugation.
A root functional satisfying $C^{\mathsf T}u=u$ has two individually
selfdual endpoints, whereas one satisfying $C^{\mathsf T}u=-u$
has endpoints exchanged by source duality.  No selfduality assertion
is made for an arbitrary root over a nonreal field.
\end{proposition}

\begin{proof}
Schur orthogonality and the full embedding matrix of $B$ give
\[
 \sum_Xz_Xz_X^{\mathsf T}=mfG^{-1},\qquad
 \sum_Xd_Xz_X=0.
\]
The second identity uses the non-Perron hypothesis.  Thus $\Phi$
is an isometry into $\mathbb Z^r$, with image perpendicular to the
Perron vector.  An integral vector of squared norm two has two
entries $\pm1$.  Positivity of the Perron dimensions forces opposite
signs and equal dimensions at their positions.

Draw an edge between $X,Y$ when $e_X-e_Y$ lies in the image lattice.
Edges $XY$ and $YZ$ imply the edge $XZ$, because their appropriately
oriented root vectors add to $e_X-e_Z$.  Each nontrivial connected
component is therefore a complete graph, whose roots generate exactly
$A_s$ on its $s+1$ vertices.  Distinct components are supported on
disjoint coordinates and are orthogonal.  This is a statement about
the root-generated sublattice, which need not equal the full image.

Finally $z_{X^*}=Cz_X$, so source duality sends $\Phi(u)$ to
$\Phi(C^{\mathsf T}u)$.  The two assertions follow by inspecting the
unique positive and negative entries of a root.  In a totally real
field $C=I$.  The argument over a nonreal field requires no assertion
that $C$ preserves the chosen ambient coordinate lattice.
\end{proof}

\subsection{Integral quadratic linear-character columns}
\label{paper:sec:quadratic-linear-fields}

\begin{theorem}[Integral norms of a quadratic linear-character column]
\label{paper:quadratic-linear-norms}
\label{paper:imaginary-quadratic-skew-norm}
Let $h$ be a linear irreducible character of a fusion ring whose full
field is quadratic, with fundamental discriminant $\Delta$.  Suppose
its formal codegree $f$ is an integer.  For each source simple $X$,
write
\[
                h(X)=\frac{a_X+c_X\sqrt\Delta}{2},
                \qquad a_X,c_X\in\mathbb Z.
\]
The unit has $c_{\mathbf1}=0$.
\begin{enumerate}
\item If $\Delta<0$ and there are $b$ source dual pairs, then
\[
                   f=|\Delta|\sum_{i=1}^{b}c_i^2,
\]
where one representative of each pair is used.  Thus $|\Delta|\mid f$
and $f/|\Delta|$ is a sum of at most $b$ integer squares.  At a fixed
pair position the squared skew coordinate is $|\Delta|c_i^2$, and
$f/|\Delta|-c_i^2$ must be a sum of at most $b-1$ integer squares.
\item If $\Delta>0$, suppose in addition that the quadratic Galois
orbit is disjoint from the Perron character.  This is automatic for
an integral fusion ring, whose Perron character is rational.  Then
\[
 \sum_Xc_X^2=\frac{2f}{\Delta},\qquad
 \sum_X d_Xc_X=0,\qquad \sum_Xa_Xc_X=0,
\]
where $d_X>0$ is the Frobenius--Perron dimension.  The integer
$2f/\Delta$ is at least two.  In particular a real quadratic
linear-character field with $\Delta=2f$ is impossible.
\end{enumerate}
These statements do not require a source-rank bound.
\end{theorem}

\begin{proof}
Character values are algebraic integers.  For
$\Delta\equiv1\pmod4$, the standard quadratic integer lattice gives
$a_X\equiv c_X\pmod2$.  For $\Delta\equiv0\pmod4$, it gives even
$a_X$, with $c_X$ an arbitrary integer.

First let $\Delta<0$.  Since
$h(X^*)=\overline{h(X)}$, self-dual labels have $c_X=0$, and the two
members of a source pair have opposite coefficients.  The distinct
characters $h$ and $\overline h$ obey
\[
              \sum_X|h(X)|^2=f,\qquad \sum_Xh(X)^2=0.
\]
Subtracting the real part of the second identity from the first gives
\[
 f=2\sum_X\operatorname{Im}(h(X))^2
   =\frac{|\Delta|}{2}\sum_Xc_X^2
   =|\Delta|\sum_{i=1}^{b}c_i^2.
\]
All assertions in the imaginary case follow.

Now let $\Delta>0$, and let $h'$ be the other Galois-conjugate
character.  Its formal codegree is also $f$, because $f$ is an
integer.  Both characters are real on the fusion basis, and
column orthogonality gives
\[
 \sum_Xh(X)^2=\sum_Xh'(X)^2=f,
 \qquad\sum_Xh(X)h'(X)=0.
\]
The squared norm of $h-h'=\sqrt\Delta\,c$ is therefore $2f$.
The assumed orthogonality of both characters to the Perron character gives
$\sum_Xd_Xc_X=0$.  Subtracting their two norm identities gives
$\sum_Xa_Xc_X=0$.  If $\sum_Xc_X^2=1$, exactly one nonunit
coefficient is $1$ or $-1$, contradicting the positivity of its
dimension in $\sum_Xd_Xc_X=0$.  This proves the real-field assertion.
Every character used here has degree one; no higher-dimensional
normalization is being substituted.
\end{proof}

\begin{proposition}[Disjoint supports at real quadratic norm two]
\label{paper:real-quadratic-paired-support}
In the real quadratic case of the theorem, suppose $\Delta=f$.
Then the coefficient vector $c$ has support on exactly two distinct
nonunit self-dual source labels $X,Y$, with opposite coefficients
$1,-1$.  These labels satisfy
\[
                         d_X=d_Y,\qquad a_X=a_Y.
\]
The two-label supports belonging to distinct such quadratic Galois
orbits are disjoint.

Consequently, if there are $t$ such real quadratic orbits and $s$
nonunit self-dual source labels, then $2t\le s$.  If the source ring is
integral, its dimension is even, and its rank is odd, then
\[
                              2t\le s-2.
\]
\end{proposition}

\begin{proof}
The norm identity is $\sum_Xc_X^2=2$.  Values of a real character
agree on $X$ and $X^*$, so $c_X=c_{X^*}$.  One nonselfdual pair
with coefficient $1$ or $-1$ would give the entire norm, but would
make $\sum_Xd_Xc_X=\pm2d_X\ne0$.  The support therefore consists of
two self-dual labels.  Positive dimensions force opposite signs and
equal dimensions.  The identity $\sum_Xa_Xc_X=0$ then gives equal
trace coordinates.

For two distinct quadratic Galois orbits, all four constituent
linear characters are mutually orthogonal across the orbits.  Their
difference vectors are therefore orthogonal.  Two vectors of the form
$\pm(e_X-e_Y)$ cannot have exactly one common support position, and
cannot have the same two-position support.  Their supports are disjoint.

At odd rank the number $s$ is even.  If $2t=s$, all nonunit
self-dual dimensions can be paired into equal pairs.  Their squared
dimension sum is then even; the contribution of every source dual
pair is also even.  The remaining unit contributes one, making the
total dimension odd.  This contradicts the stated even dimension.
Thus $2t<s$, and the parity of both sides gives $2t\le s-2$.
\end{proof}

\begin{corollary}[Quadratic transpositions of the fusion basis]
\label{paper:quadratic-transposition-automorphisms}
Let $R$ be a fusion ring, and let a real quadratic Galois orbit of
linear characters satisfy the hypotheses of
Proposition~\ref{paper:real-quadratic-paired-support}, with fundamental
discriminant $\Delta=f$.  Let $X,Y$ be its two supported source labels.
Then exchanging $X$ and $Y$ and fixing every other source label is a
based-ring automorphism of $R$.

Outside this quadratic Wedderburn factor, the representation images
of $X$ and $Y$ coincide.  In particular $d_X=d_Y$, and their aggregate
row data, consisting of the trace and squared-norm trace on every
full-field representation block, coincide.  If there are $t$ distinct
quadratic orbits of this kind, their transpositions generate a faithful
group of based-ring automorphisms isomorphic to $(\mathbb Z/2\mathbb Z)^t$.
No lift of these automorphisms to tensor autoequivalences is asserted.
\end{corollary}

\begin{proof}
Write $h,h'$ for the two characters, choosing their order so that the
integer difference coefficients satisfy $c_X=1$ and $c_Y=-1$.  For a
linear character the central primitive idempotent in $R_{\mathbb C}$ is
\[
 e_h=\frac1f\sum_Z h(Z^*)Z.
\]
The characters are real on the fusion basis.  Their difference is
supported exactly on $X,Y$, so
\[
 e_h-e_{h'}=\frac{\sqrt\Delta}{f}(X-Y),
 \qquad
 X-Y=\sqrt\Delta\,(e_h-e_{h'}).
\]
Thus $X-Y$ is central and vanishes in every Wedderburn factor outside
the given quadratic factor.  All representation images of $X,Y$
outside that factor are consequently equal.

In the rational semisimple algebra $R_{\mathbb Q}$, apply the nontrivial
field automorphism to this quadratic factor and the identity to every
other factor.  This algebra automorphism sends $X-Y$ to its negative.
It fixes $X+Y$, since the two quadratic difference coordinates cancel,
and fixes every other source label, whose difference coordinate is zero.
It therefore exchanges $X,Y$ and preserves the fusion basis and its
unit, proving the based-ring assertion.  The two affected labels are
self-dual, as established in the cited proposition.

On the affected full quadratic block, conjugation preserves both the
trace and the squared-norm trace; on every other block the images are
equal.  This proves equality of the aggregate row data.  The Perron
character is outside the quadratic orbit by hypothesis, so it also
gives $d_X=d_Y$.  Finally the supports of distinct norm-two quadratic
orbits are disjoint by the same proposition.  Their transpositions
therefore commute and act independently on these disjoint pairs,
giving the claimed faithful elementary abelian group.
\end{proof}

\begin{corollary}[An integral orthogonal frame in the quadratic case]
\label{paper:quadratic-skew-frame}
Suppose the unique degree-two component has orthogonal involution,
and every nonreal linear-character pair has a full imaginary
quadratic field, of discriminant $\Delta_j$ and codegree $f_j$.
Let $q$ be the number of these pairs, let $b=q+1$, and let
$\delta$ be the matrix skew squareclass.  Then an integer
$b\times b$ matrix $M$ satisfies
\[
 M^{\mathsf t}M=\operatorname{diag}
 \left(\frac{f_1}{|\Delta_1|},\ldots,
       \frac{f_q}{|\Delta_q|},\frac{f_E}{\delta}\right).
\]
\end{corollary}

\begin{proof}
Choose one character from each conjugate pair.  The integer vectors
from Theorem~\ref{paper:quadratic-linear-norms} are pairwise orthogonal: apply character
orthogonality both to $h_j,\overline{h_k}$ and to $h_j,h_k$, and
take imaginary parts.  The matrix skew line is orthogonal to these
linear-character summands.  Its integral Gram factorization supplies
an integer vector of squared length $f_E/\delta$.  There are exactly
$q+1=b$ orthogonal nonzero columns, so adjoining these vectors gives
the asserted square matrix.  This is a necessary integral-lattice
condition; no integral-frame enumeration is required by the endpoint
certificates below.
\end{proof}

\subsection{Quadratic degree-two orbits of codegree six}
\label{paper:sec:quadratic-six}

The next results concern a sole nonlinear Galois orbit: two conjugate
irreducible representations of degree two and formal codegree six.
All other irreducible representations are linear. The source is
integral and has no nonunit invertible basis element. These hypotheses
hold in the quadratic two-matrix branch of the conditional $5400$
source problem. They are stated explicitly because a further nonlinear
component would add a commutator contribution.

\begin{lemma}[A small integral projection graph]
\label{paper:projection-cycle-lemma}
Let $H$ be an integral positive semidefinite matrix with $H^2=3H$ and
diagonal entries in $\{0,1,2\}$. Its nonzero part is an orthogonal sum,
up to simultaneous permutations and sign changes of rows and columns,
of the following components:
\begin{itemize}
\item a three-cycle with all diagonal entries one, of rank one;
\item a three-cycle with all diagonal entries two, of rank two;
\item a four-cycle with alternating diagonal entries one and two and
      negative edge-sign product, of rank two.
\end{itemize}
Consequently, rank three in at most six coordinates requires one
three-cycle of each of the first two kinds.
\end{lemma}

\begin{proof}
A diagonal entry one or two leaves off-diagonal squared sum two in
$H^2=3H$, hence has exactly two incident entries $\pm1$. Thus the
support graph is a union of cycles. A cycle of length at least five
has two nonadjacent vertices with exactly one common neighbor,
contradicting $(H^2)_{ij}=0$. On a triangle the edge equation forces
all diagonal entries to be one, with positive edge product, or all to
be two, with negative edge product. On a four-cycle adjacent
diagonals sum to three, and the two length-two paths between opposite
vertices have opposite signs. The ranks follow from $\rank H=\Tr H/3$.
The final assertion follows by counting the coordinates required by
components of total rank three.
\end{proof}

\begin{theorem}[The Gaussian orbit needs rank at least seventeen]
\label{paper:gaussian-matrix-six}
Under the hypotheses above, a nonlinear full character field
$K=\mathbb Q(i)$ requires source rank at least $17$.
\end{theorem}

\begin{proof}
Write $\chi(X)=u_X+i v_X$ with $u_X,v_X\in\mathbb Z$.
Character orthogonality gives
\[
 \sum_X u_X^2=\sum_X v_X^2=6,\qquad
 u_{\mathbf1}=2,\qquad v_{\mathbf1}=0.
\]
The two nonzero off-unit $u$-coordinates are $1$ and $-1$:
Perron orthogonality says $\sum_{X\ne\mathbf1}d_Xu_X=-2$, and two
negative signs would force two nonunit dimensions to sum to two.
Each sign occurs once, so both elements are self-dual. On every
non-self-dual pair $u=0$. Choosing pair representatives, the vector
$v=(v_i)$ therefore has exactly three nonzero entries $\pm1$ and
$\sum v_i^2=3$. In rank at most $16$ the three already forced
self-dual elements imply $b\le6$.

For this orbit the matrix $K$ of
Theorem~\ref{paper:source:skew-projection-lattice} is even, has rank four, and
satisfies $K^2=6K$. Its central imaginary line has orthogonal
projection $vv^{\mathsf t}/3$. Hence
\[
 H=K/2-vv^{\mathsf t}\in\operatorname{Mat}_b(\mathbb Z),\qquad
 H^2=3H,\qquad \rank H=3,\qquad Hv=0.
\]
For a pair row write $\alpha=i v_i$, $\gamma=\Tr(AA^*)$, and
$\beta=\det A$. Its positive contribution to $\tau(XX^*)=1$, with
the Perron contribution still present, gives
$\gamma\in\{0,1,2\}$. The determinant inequality
$|\beta|\le\gamma/2$ leaves zero and Gaussian units.
Using $G\ge0$, $G\in3\mathbb Z$, and $0\le K_{ii}\le6$ gives
\[
 H_{ii}\in\{0,1,2\},\qquad
 H_{ii}=1\Longrightarrow v_i\ne0,\qquad
 H_{ii}=2\Longrightarrow v_i=0.
\]
For completeness $H_{ii}=\gamma+2\operatorname{Re}\beta$;
the entry one comes only from $(\gamma,\beta)=(1,0)$ with
$v_i=\pm1$, and the entry two only from $(2,\pm i)$ with $v_i=0$.
These are finite exact substitutions in
\eqref{paper:general-two-matrix-commutator}.

Lemma~\ref{paper:projection-cycle-lemma} forces one triangle of each
kind. All three nonzero entries of $v$ lie on the rank-one triangle,
which has matrix $zz^{\mathsf t}$ with $z_i\in\{1,-1\}$.
The identity $Hv=0$ would require a sum of three signs to vanish.
This contradiction proves the theorem.
\end{proof}

\begin{lemma}[Two isolated trace-zero rows are impossible]
\label{paper:two-isolated-unitaries}
For a sole imaginary quadratic degree-two orbit of codegree six,
there cannot be two distinct non-self-dual source-pair coordinates
whose orbit skew-projection diagonals equal six, whose matrix traces
are zero, and whose matrix norms are two with determinant of modulus
one and positive real part.
\end{lemma}

\begin{proof}
Write $A,B$ for the images of the chosen elements at one complex
embedding. Norm two and determinant modulus one imply that both
matrices are unitary. As their traces vanish,
$A^*=-A/\det A$ and $B^*=-B/\det B$. The positive-real-part condition
ensures that $A-A^*$ and $B-B^*$ are nonzero scalar multiples of $A,B$.
An entry $K_{ii}=6$ in $K^2=6K$ has all its off-diagonal entries zero.
Thus $\Tr((A-A^*)(B-B^*))=0$: this trace is real because both factors
are skew-Hermitian, so the two conjugate embeddings do not cancel.
It follows that $\Tr(AB)=0$, and the degree-two trace identity gives
$AB=-BA$. Therefore
\[
 \tau([X,Y][X,Y]^*)=
 \frac{\Tr([A,B][A,B]^*)+
       \Tr([\overline A,\overline B][\overline A,\overline B]^*)}{6}
 =\frac{16}{6}.
\]
Scalar representations annihilate the commutator. The left side is
the sum of the squares of its integral fusion-basis coefficients,
and must be an integer. This is a contradiction.
\end{proof}

\begin{theorem}[The Eisenstein orbit needs rank at least seventeen]
\label{paper:eisenstein-matrix-six}
Under the same hypotheses, a nonlinear full field
$K=\mathbb Q(\sqrt{-3})$ requires source rank at least $17$.
\end{theorem}

\begin{proof}
Write $\chi(X)=(s_X+b_X\sqrt{-3})/2$, with integral $s_X,b_X$ of
the same parity. Orthogonality gives
\[
 \sum_Xs_X^2=24,\qquad \sum_Xb_X^2=8,
 \qquad s_{\mathbf1}=4.
\]
On pair representatives $\sum b_i^2=4$. There are two cases.

\emph{One exceptional pair has $b_i=\pm2$.}
Its $s_i$ is even. A value $\pm2$ would exhaust the off-unit square
budget eight and Perron orthogonality would force its dimension to be
one. Thus $s_i=0$. The remaining real character values are two
opposite signs $1,-1$ on self-dual elements, whose dimensions differ
by two. Together with the unit these imply $b\le6$ at rank at most
$16$. The exceptional row has $\chi=\pm\sqrt{-3}$; the norm and
commutator inequalities force $(\gamma,\beta)=(2,-1)$ and $K_{ii}=6$.
It is precisely the central imaginary line. All other pair traces
are zero.

For a trace-zero pair, the same finite inequalities leave
$K_{jj}\in\{0,2,6\}$. Subtracting the exceptional coordinate gives
an integral matrix $H$ on at most five coordinates with
$H^2=6H$, rank three, and diagonal entries $0,2,6$.
A diagonal six is isolated. At a diagonal two, the off-diagonal
squared sum is eight; every possible neighbor also has diagonal two
and the entry has absolute value at most two. With at most four
neighbors the only possibility is two entries of absolute value two.
Thus $H/2$ is integral, squares to three times itself, and decomposes
into rank-one all-one triangles and isolated entries three.
Rank three on at most five coordinates needs at least two isolated
entries. These give two trace-zero diagonals six.

\emph{Four pairs have $b_i=\pm1$.}
Their $s_i=\pm1$ already exhaust the off-unit $s$-square budget.
Their character values are sixth roots of unity. Each such row has
either $(\gamma,\beta)=(1,0)$ or
$(\gamma,\beta)=(2,\chi^2)$, and in both cases $K_{ii}=3$.
The four coordinates contribute twelve to $\Tr K=24$.
At rank at most $16$ there are at most seven pairs, so at most three
other coordinates remain. Their traces are zero and their diagonals
belong to $\{0,2,6\}$. A sum twelve in at most three such coordinates
requires at least two entries six.

In both cases a trace-zero diagonal six has $\gamma=2$ and determinant
an Eisenstein unit with real part $1/2$. The two forced rows contradict
Lemma~\ref{paper:two-isolated-unitaries}.
\end{proof}

\begin{theorem}[The real quadratic orbit is impossible at every rank]
\label{paper:real-quadratic-matrix-six}
Under the same hypotheses, the nonlinear full character field cannot
be $\mathbb Q(\sqrt3)$, at any source rank. Both orthogonal and
symplectic involutions, and nonsplit simple components, are covered.
\end{theorem}

\begin{proof}
Write $\chi(X)=u_X+v_X\sqrt3$ with integral $u_X,v_X$.
Orthogonality gives $\sum u_X^2=6$, $\sum v_X^2=2$ and
$u_{\mathbf1}=2$. As in the Gaussian proof, the two off-unit
nonzero $u$ values are opposite signs, each on a self-dual element.
Perron orthogonality gives $\sum d_Xv_X=0$, so the two nonzero
$v$ values are also opposite signs, each on a self-dual element.
Consequently $\chi(X)=0$ on every non-self-dual pair.

For such a pair put $\gamma=\Tr(AA^*)\in\mathbb Z[\sqrt3]$ and
$\beta=\det A\in\mathbb Z[\sqrt3]$. Total positivity and the strict
bound $\Tr_{K/\mathbb Q}\gamma<6$ give
\[
                 \gamma\in\{0,1,2,2+\sqrt3,2-\sqrt3\}.
\]
At either embedding $|\beta|\le\gamma/2$.
For $\gamma=1$ or $2\pm\sqrt3$ the integral norm of $\beta$ has
absolute value at most $1/4$, hence $\beta=0$; the same is immediate
at $\gamma=0$. At $\gamma=2$ both real conjugates of $\beta$ lie
in $[-1,1]$, so $\beta\in\{-1,0,1\}$.

Since $\chi=0$, the coupled commutator test is
\[
       \Tr_{K/\mathbb Q}(\gamma^2-4\beta^2)/6
                                  \in\mathbb Z_{\ge0}.
\]
It eliminates $\gamma=1$, $\gamma=2\pm\sqrt3$, and
$(\gamma,\beta)=(2,0)$, with respective numerators $2,14,8$.
For $(2,1)$ the skew diagonal
$K_{ii}=\Tr_{K/\mathbb Q}(\gamma+2\beta)=8$ exceeds six.
The only retained rows are $(0,0)$ and $(2,-1)$, both with $K_{ii}=0$.
Hence the whole skew projection vanishes. Its rank must instead be
two in the orthogonal case or six in the symplectic case, giving
the contradiction.
\end{proof}

\begin{corollary}[No quadratic codegree-six matrix orbit at low rank]
\label{paper:quadratic-six-orbit-exclusion}
For a sole nonlinear quadratic degree-two orbit of codegree six,
source rank is at least seventeen. In particular, at rank fifteen
or sixteen two codegree-six matrix irreducibles must be individually
rational if they are the only nonlinear irreducibles.
\end{corollary}

\begin{proof}
The full-field conductor bound is
$c\mid\gcd(2f,f^2)=12$. The quadratic possibilities are exactly
$\mathbb Q(i)$, $\mathbb Q(\sqrt{-3})$, and $\mathbb Q(\sqrt3)$.
Apply Theorems~\ref{paper:gaussian-matrix-six},
\ref{paper:eisenstein-matrix-six} and~\ref{paper:real-quadratic-matrix-six}.
\end{proof}

\subsection{An upper bound for the regular trace}
\label{paper:sec:regular-trace-upper-bound}

The following structural bound applies at every source rank. The
endpoint certificates do not need it as an additional row filter.

\begin{lemma}[Regular trace in a fusion-simple category]
\label{paper:regular-trace-upper-bound}
Let $\mathcal C$ be a fusion category with no proper nontrivial
fusion subcategory.  For every nonunit simple $X$, its left fusion
matrix $N_X$ is irreducible.  If $\mathcal C$ is integral and has
rank $r$, then
\[
           0\le\operatorname{Tr}N_X\le(r-1)(d_X-1),
           \qquad d_X=\operatorname{FPdim}(X).
\]
\end{lemma}

\begin{proof}
Index source simples by $i$, put $d_i=\operatorname{FPdim}(i)>0$,
and use the convention $(N_X)_{ij}=N_{Xj}^{i}$.  Dimension
multiplicativity and Frobenius reciprocity give
\[
 N_X\mathbf d=d_X\mathbf d,\qquad
 \mathbf d^{\mathsf T}N_X=d_X\mathbf d^{\mathsf T},\qquad
 N_{X^*}=N_X^{\mathsf T}.
\]
With $D_{\mathbf d}=\operatorname{diag}(d_i)$, the matrix
\[
 P=\frac1{d_X}D_{\mathbf d}^{-1}N_XD_{\mathbf d}
\]
is row-stochastic and has strictly positive stationary distribution
$\pi_i=d_i^2/\sum_jd_j^2$.

For completeness, this implies that every strong component is
isolated.  Choose a sink strong component $S$ of the finite directed
graph of $P$.  There are no transitions from $S$ to its complement.
Summing stationarity over $S$ gives
\[
 \sum_{i\in S}\pi_i
   =\sum_{i\in S}\pi_i
      +\sum_{i\notin S}\pi_i\sum_{j\in S}P_{ij}.
\]
Since all $\pi_i$ are positive, there are no incoming transitions
either.  Delete this isolated block and repeat.  Thus $N_X$ is
permutation-similar to a block diagonal matrix whose blocks are
its strong components.

The block containing the unit is invariant under both $N_X$ and
$N_X^{\mathsf T}=N_{X^*}$.  Applying them to the unit shows that
the block contains $X$ and $X^*$; repeated application contains
every simple summand of every word in $X,X^*$.  It therefore
contains the fusion subcategory generated by $X$.  Since $X$ is
nonunit, fusion simplicity makes this the whole category.  Hence
$N_X$ has only one strong component and is irreducible.

There is a nonunit simple, so $r>1$.  Every row of this irreducible
matrix has a positive off-diagonal entry.  Consequently
\[
 (d_X-(N_X)_{ii})d_i
       =\sum_{j\ne i}(N_X)_{ij}d_j>0,
\]
and each diagonal entry is strictly less than $d_X$.  In the
integral case it is therefore at most $d_X-1$.  The unit diagonal
entry is zero, since $N_{X\mathbf1}^{\mathbf1}=0$ for nonunit $X$.
Summing the remaining $r-1$ nonnegative diagonal entries proves
the bound.
\end{proof}

\section{Source rows, positive words and integral central spectra}

\subsection{A common finite row model}
Fix all full character fields, Galois orbits, and admissible rational
matrix alternatives.  A row may have values in proper subfields;
fixing a character's full field does not require every entry to generate
that field.  Aggregate rows give necessary supersets when the individual
matrix images have not been constructed.

\begin{lemma}[Row equations and global moments]
\label{paper:source:row-moments}
Partition the non-Perron characters into Galois-stable blocks
$\mathcal B_i$ of common codegree $f_i$.  Set
\[
 a_i=n/f_i,\quad k_i=\sum_{E\in\mathcal B_i}m_E,
 \quad h_i=\sum_{E\in\mathcal B_i}m_E^2,
\]
and, for a source simple $X$, put
\[
 s_i(X)=\sum_{E\in\mathcal B_i}\chi_E(X),\quad
 t_i(X)=\sum_{E\in\mathcal B_i}\chi_E(XX^*),\quad
 u_i(X)=\sum_{E\in\mathcal B_i}\chi_E(X^2).
\]
For a nonunit row of dimension $d$,
\begin{align}
 d+\sum_i a_i s_i&=0,&
 d^2+\sum_i a_i t_i&=n,& 0\le t_i&<f_i,
 \label{paper:source:row-equations}\\
 d^2+\sum_i a_i u_i&=n\delta_{X,X^*}.&&&
 \label{paper:source:row-square-equation}
\end{align}
If a block consists of real-valued linear characters, $u_i=t_i$.
For a single imaginary quadratic linear orbit, $u_i=s_i^2-t_i$.
The matrix blocks retain their separate $u_i$ and $t_i$ data.
Summing over nonunit source simples gives
\begin{align}
 \sum_X d_X^2&=n-1,& \sum_Xd_Xs_i(X)&=-k_i,\\
 \sum_Xs_i(X)^2&=k_if_i-k_i^2,&
 \sum_Xs_i(X)s_j(X)&=-k_ik_j\quad(i\ne j),\\
 \sum_Xt_i(X)&=h_if_i-k_i.&&
\end{align}
For $b$ source dual pairs, these sums contain $r-1-2b$ selfdual
nonunit atoms and $b$ pair atoms, each pair contributing twice.
\end{lemma}
\begin{proof}
Apply the decomposition of $\tau$ to $X$, $XX^*$, and $X^2$.
The unit coefficients are $0$, $1$, and $\delta_{X,X^*}$.
The Perron terms are $d/n$, $d^2/n$, and $d^2/n$;
positivity of each $t_i$ gives its strict upper bound.
The two scalar-block formulas follow by squaring real values or
the conjugate pair $\alpha,\bar\alpha$.

Character orthogonality gives
$\sum_X|\chi_E(X)|^2=m_Ef_E$ and orthogonality between distinct
characters, including the Perron character.  A full block trace is
real and invariant under source duality.  Removing its unit value
$k_i$ gives the stated trace moments.  Finally
$\sum_X\chi_E(XX^*)=m_E^2f_E$ on the whole basis;
subtracting $\chi_E(\mathbf1)=m_E$ proves the last formula.
\end{proof}

These equations give a finite row set.  A positive trace-norm form
bounds every integral coordinate in each specified field; for aggregate
matrix spectra the bounded monic integral-polynomial superset in
Lemma~\ref{paper:aggregate-polynomial-superset} supplies a finite
necessary enumeration without a Schur-index assumption.
For a fixed rank $r$, all source dimensions lie in the sharper ranges
\[
 3\le d\le\left\lfloor\sqrt{5399-9(r-2)}\right\rfloor
 \quad\text{for a selfdual nonunit,}
\]
\[
 3\le d\le
 \left\lfloor\sqrt{\frac{5399-9(r-3)}2}\right\rfloor
 \quad\text{for a nonselfdual pair.}
\]
These bounds use only the dimensions of the remaining labels.

\begin{theorem}[An integral recurrence of degree at most the canonical length]
\label{paper:source:finite-power-trace-test}
For $x\in R$, put
\begin{equation}\label{paper:source:unweighted-character-polynomial}
 p_x(T)=\prod_{E\in\operatorname{Irr}(R_{\mathbb C})}
                  \det(TI_{\dim E}-E(x)).
\end{equation}
Then $p_x\in\mathbb Z[T]$ is monic of degree $L$,
$p_x(x)=0$, and $p_{x^*}=p_x$.
Consequently:
\begin{enumerate}
\item the sequence $c_k=\tau(x^k)$ obeys an integer recurrence
of length $L$. Its terms are all integers if and only if
$c_0,\ldots,c_{L-1}$ are integers;
\item the array $c_{k,\ell}=\tau(x^k(x^*)^\ell)$ obeys the
same recurrence in each variable. Its entries are all integers
if and only if the $L\times L$ initial rectangle
$0\le k,\ell<L$ consists of integers;
\item for a positive integral combination $x$ of simple basis
elements, every $c_k$ and $c_{k,\ell}$ is nonnegative.
Checking finitely many of these inequalities is a necessary
test; the integer recurrence alone does not propagate positivity.
\end{enumerate}
In spectral candidate calculations, the same conclusions about
propagation of integrality hold whenever the chosen spectral data
give the displayed monic integral polynomial and its recurrences.
The choices of character values and matrix invariants must be kept
consistent across all powers being tested.
\end{theorem}

\begin{proof}
The left regular characteristic polynomial of $x$ is
\[
       \prod_E\det(TI-E(x))^{\dim E}\in\mathbb Z[T].
\]
Thus every eigenvalue of every $E(x)$ is an algebraic integer.
The absolute Galois group permutes the irreducible characters
and their characteristic polynomials, so the product in
\eqref{paper:source:unweighted-character-polynomial} has rational
coefficients. These coefficients are also algebraic integers,
and hence are integers. In each simple matrix component the
corresponding factor annihilates $E(x)$ by Cayley--Hamilton.
All factors are polynomials in that same matrix, so the product
annihilates it as well. Semisimplicity gives $p_x(x)=0$.

The coefficient-of-unit form makes every complex simple
component a positive $*$-algebra. In a unitary realization,
$E(x^*)=E(x)^{\dagger}$. Its characteristic polynomial is
the complex conjugate of that of $E(x)$. Therefore
$p_{x^*}=\overline{p_x}=p_x$.
Applying $\tau$ after multiplying the monic relation by powers
of $x$ gives the first recurrence. Applying it independently
to $x$ on the left and $x^*$ on the right gives the two
recurrences for the array. Integer initial data propagate by
induction. Nonnegative fusion coefficients give the final
assertion directly.
\end{proof}

\begin{lemma}[A short positive power]
\label{paper:source:short-positive-power}
For every simple basis element $X$ in a fusion ring of rank $r$,
there is some $1\le k\le L\le r$ for which $\tau(X^k)>0$.
For a nonunit simple, necessarily $k\ge2$.
This assertion does not require that the multiplication matrix
of $X$ be irreducible. More generally $L$ may be replaced by
the degree of any monic polynomial annihilating $X$ itself.
\end{lemma}

\begin{proof}
Let $N$ be left multiplication by $X$, and let $d$ be the
strictly positive vector of simple Frobenius--Perron dimensions.
The dimension homomorphism and Frobenius reciprocity give
$Nd=d_Xd$ and $N^{\mathsf t}d=d_Xd$.
The matrix
\[
              P_{ij}=\frac{N_{ij}d_j}{d_Xd_i}
\]
is stochastic, and has a strictly positive stationary measure
$\pi_i=d_i^2/\sum_jd_j^2$.
Every positive directed edge is contained in a directed cycle:
otherwise the set reachable from its terminal vertex would
have a positive incoming stationary flow and zero outgoing
flow, contradicting stationarity. Each vertex has an outgoing
edge, so the unit vertex lies on a simple directed cycle of
length at most $r$. The corresponding diagonal entry of
$N^k$ is positive and equals $\tau(X^k)$.
The unit coefficient of a nonunit $X$ is zero.

Finally, if all $\tau(X^j)$ for $1\le j\le L$ vanished,
the recurrence from Theorem~\ref{paper:source:finite-power-trace-test},
started at $j=L+1$, would force every later positive-power
coefficient to vanish. This contradicts the cycle just obtained.
The same argument applies to any genuine monic annihilator
of $X$. The inequality $L\le r$ follows from
$\sum_E\dim E\le\sum_E(\dim E)^2$.
\end{proof}

\subsection{Finite power traces and positive word coefficients}
\label{paper:sec:power-word-traces}

The coefficient-of-unit trace contains more information than its
linear and quadratic moments. If $X$ is a simple basis element,
every coefficient $\tau(X^k)$ is a nonnegative integer. This gives
an inexpensive obstruction before one attempts to reconstruct the
fusion multiplication. The following statements apply to any fusion
ring, independently of the dimension-$5400$ center hypothesis.
Write $r=\operatorname{rank}(R)$ and
$L=\sum_{E\in\operatorname{Irr}(R_{\mathbb C})}\dim E$.

\begin{proposition}[Exact degree-two power and word traces]
\label{paper:matrix-power-word-traces}
Let $E$ be a rational degree-two irreducible character and let
$X$ be a source simple. In a unitary complex realization write
\[
 A=E(X),\qquad s=\operatorname{tr}A,\qquad
 n=\det A,\qquad t=\operatorname{tr}(AA^{\dagger}).
\]
Then $s,n,t$ are integers. Define integer sequences by
\[
 u_0=0,\quad v_0=1,\qquad
 u_{k+1}=su_k+v_k,\quad v_{k+1}=-nu_k.
\]
For all $k,\ell\ge0$,
\begin{equation}\label{paper:ordered-matrix-word-trace}
 \operatorname{tr}\bigl(A^k(A^{\dagger})^\ell\bigr)
 =u_ku_\ell t+(u_kv_\ell+v_ku_\ell)s+2v_kv_\ell.
\end{equation}
The traces of $(AA^{\dagger})^k$ satisfy
\begin{equation}\label{paper:alternating-matrix-word-trace}
 b_0=2,\qquad b_1=t,\qquad
 b_k=t b_{k-1}-n^2b_{k-2}\quad(k\ge2).
\end{equation}
After combining these values with the linear-character
contributions, every coefficient
$\tau(X^k(X^*)^\ell)$ and $\tau((XX^*)^k)$ must be a
nonnegative integer. These two families need not coincide
in a noncommutative ring.
\end{proposition}

\begin{proof}
Rational character values of integral ring elements are
algebraic integers, so $s$ and $t$ are integers.
The characteristic polynomial of $A$ is Galois fixed because
the character is rational, and its roots are algebraic integers;
thus $n\in\mathbb Z$. Cayley--Hamilton gives
$A^2-sA+nI=0$, and induction gives $A^k=u_kA+v_kI$.
The adjoint satisfies the same relation, since $s,n$ are real.
Multiplication and tracing prove
\eqref{paper:ordered-matrix-word-trace}.
The matrix $AA^{\dagger}$ has trace $t$ and determinant $n^2$,
so its Cayley--Hamilton recurrence gives
\eqref{paper:alternating-matrix-word-trace}.
All the underlying ring words are positive integral
combinations of simples, proving their coefficient conditions.
\end{proof}

\begin{proposition}[A commutator congruence for a unique degree-two block]
\label{paper:unique-matrix-commutator-congruence}
Suppose $R_{\mathbb C}=\mathbb C^a\oplus M_2(\mathbb C)$,
and let $f$ be the formal codegree of the unique degree-two
character. For the preceding integers $s,n,t$, put
\begin{equation}\label{paper:commutator-polynomial}
              F=(t-2n)(t+2n-s^2).
\end{equation}
Then
\[
        F\ge0,\qquad F/f\in\mathbb Z_{\ge0},\qquad
        \tau\bigl((XX^*-X^*X)^2\bigr)=2F/f.
\]
In particular $f\mid F$. Moreover $F=0$ if and only if
$XX^*=X^*X$ in the fusion ring.
\end{proposition}

\begin{proof}
The unique nonlinear character is rational. In its matrix
component, the preceding recurrences give
\[
 \operatorname{tr}((AA^{\dagger})^2)=t^2-2n^2,
 \qquad
 \operatorname{tr}(A^2(A^{\dagger})^2)=s^2t-2s^2n+2n^2.
\]
Their difference is $F$. Every linear component makes equal
contributions to the two words, and the symmetrizing trace
on the matrix component is $\operatorname{tr}/f$. Hence
\[
   \tau((XX^*)^2)-\tau(X^2(X^*)^2)=F/f\in\mathbb Z.
\]
The self-adjoint commutator $C=AA^{\dagger}-A^{\dagger}A$
satisfies $\operatorname{tr}(C^2)=2F$. Positivity gives
$F\ge0$. The same calculation in $R$, whose coefficient
form is positive definite, gives the stated norm identity
and the equivalence at $F=0$.
\end{proof}

\subsection{Integral entries, matrix rows and central atoms}
\label{paper:source:matrix-row-domains}
The full character fields remain fixed across a table.  The next
enumeration works in their complete rings of integers, including entries
in proper subfields; the matrix tests retain nonsplit rational algebras.
\begin{lemma}[Exact integral-entry power enumeration]
\label{paper:comm12-integral-entry-powers}
For a fixed full field $K=\mathbb Q(\zeta_c)^H$, every possible integral
character entry of norm trace at most $f-1$ can be enumerated exactly in a
rational nullspace of the integral cyclotomic power basis. This enumeration
also includes entries in proper subfields of $K$.
For each entry $\alpha$, its full-field polynomial
\[
       P_{K,\alpha}(T)=\prod_{\sigma:K\hookrightarrow\mathbb C}
                                      (T-\sigma(\alpha))\in\mathbb Z[T]
\]
retains repetitions when $\alpha$ lies in a proper subfield. Its power
traces are determined by Newton identities and the resulting integral
recurrence.
\end{lemma}
\begin{proof}
Write $\alpha=\sum_{j=0}^{\varphi(c)-1}x_j\zeta_c^j$, with $x_j\in\mathbb Z$.
The equations $h(\alpha)=\alpha$ for $h\in H$ are integral linear equations
in these coordinates. Solve them over $\mathbb Q$, choosing free variables
among the original coordinates. An actual integral entry has integer free
coordinates. Conversely, an integer free-coordinate vector gives an
integral entry exactly when the reconstructed pivot coordinates are
integers. Thus enumeration of integer free coordinates, followed by this
pivot check, is exhaustive; no assertion of an unproved integral basis for
$K$ is needed.

The positive norm matrix in the ambient power basis is
$(c_c(i-j))_{i,j}$, with $c_c$ the Ramanujan sum. Restrict it to the rational
nullspace and enumerate its integer vectors of norm at most $|H|(f-1)$.
For a selfdual source row add the conjugation-fixed equations. Taking
ambient traces and dividing by $|H|$ gives the full-field traces.
Direct multiplication modulo the cyclotomic polynomial yields
$\operatorname{Tr}_{K/\mathbb Q}(\alpha^j)$ for $1\le j\le[K:\mathbb Q]$.
Newton identities recover $P_{K,\alpha}$; its integral coefficients are
also checked exactly. The same polynomial then generates all the power
traces needed in the finite source-row test.
\end{proof}

\begin{lemma}[Annihilator tests from a row]
\label{paper:row-annihilator-test}
Suppose first that the fusion ring is commutative and that the fixed full
fields of a row are rational or quadratic. A block of $k$ rational
characters with trace $s$ and norm trace $t$ allows only integer multisets
\[
 (x_1,\ldots,x_k),\qquad \sum_jx_j=s,\quad\sum_jx_j^2=t.
\]
Its distinct value polynomials are the distinct factors $T-x_j$.
For a quadratic orbit its value polynomial is
\[
 T^2-sT+\frac{s^2-t}{2}\quad\text{in the real case},\qquad
 T^2-sT+\frac t2\quad\text{in the imaginary case}.
\]
A repeated rational root contributes only its linear factor. If $p_X(T)$
is the product of the distinct irreducible factors obtained from all blocks,
then a necessary condition is
\begin{equation}
                         5400\mid p_X(d_X).
 \label{paper:row-minimal-gap}
\end{equation}
A row is retained if at least one compatible choice of rational multisets
satisfies this condition.

For a unique rational degree-two component, let $s,t$ be the trace and norm trace of
the separate matrix block. Its determinant is an integer $n$ satisfying
$2|n|\le t$. For a selfdual source object $2n=s^2-t$; for symplectic matrix
involution $2n=t$. In the orthogonal case $2n\ge s^2-t$.
When all linear square traces are known, equation
\eqref{paper:source:row-square-equation} determines
\begin{equation}
 u=\frac{5400\delta_{X,X^*}-d_X^2-
                   \sum_{i\ne M}a_i u_i}{a_M},\qquad
                         n=\frac{s^2-u}{2}.
 \label{paper:row-matrix-determinant}
\end{equation}
Reject a nonintegral or incompatible value. Otherwise use the least common
multiple of the linear-block value polynomials and $T^2-sT+n$ in
\eqref{paper:row-minimal-gap}. A repeated root of the matrix characteristic
polynomial retains multiplicity two, except when selfadjointness proves
that the matrix is diagonalizable.
\end{lemma}
\begin{proof}
The quadratic formulas follow directly from the sum and product of the two
conjugate values. Their coefficients are integers because the values are
algebraic integers and the orbit is Galois stable. A repeated root is
rational and therefore integral. The polynomial in the commutative case
annihilates every non-dimension character. Thus its value at $X$ is
$p_X(d_X)e_{\mathrm{FP}}$, and its coefficient at the unit is
$p_X(d_X)/5400$. Integrality of $p_X(X)$ proves the divisibility.

For the matrix block its rational character makes its determinant a rational
algebraic integer. The product of its singular values gives $2|n|\le t$.
For a selfadjoint matrix, $s^2-2n=\operatorname{tr}(E(X)^2)=t$.
For a symplectic positive involution, $E(X)E(X)^*=nI$.
In the orthogonal case choose its real positive realization: the squared
norm of its skew part gives $t-s^2+2n\ge0$.
Formula~\eqref{paper:row-matrix-determinant} is then the third row identity.
Cayley--Hamilton annihilates the entire matrix, including any Jordan part.
The resulting integral least common multiple annihilates every
non-dimension simple algebra. It therefore has the same unit-coefficient
argument. Strict operator norm excludes $d_X$ as a non-dimension eigenvalue,
so the scalar at the dimension representation is nonzero.
\end{proof}

The orthogonal and symplectic alternatives below refine the same row
data.  They use the complete quaternion alternatives of
Corollary~\ref{paper:finite-quaternion-trace-forms}.
\begin{theorem}[A source row determines an orthogonal quaternion class]
\label{paper:orthogonal-row-quaternion-class}
Let $B$ be a rational central simple algebra of degree two equipped with
a positive orthogonal involution $\dagger$.  Choose a nonzero rational
skew generator $U$ and normalize it so that
$U^2=-\delta$, where $\delta>0$ is squarefree.  For $A\in B$, put
\[
 s=\operatorname{Trd}(A),\qquad n=\operatorname{Nrd}(A),\qquad
 t=\operatorname{Trd}(AA^\dagger),
\]
and define
\[
 J=t-2n,\qquad K=t-s^2+2n.
\]
Then $J,K\ge0$.  If $J>0$, necessarily
\[
                         B\simeq(-\delta,J)_{\mathbb Q}.
\]
If $J=0$, then $A\in\mathbb Q(U)$.

For the image of a source simple, $s,n,t$ are integers.  If the source
label is self-dual, then $K=0$.  If it belongs to a source dual pair,
the integral skew-lattice theorem further requires
\[
 K=\delta z_i^2,\qquad
 \frac f\delta-z_i^2\text{ is a sum of at most }b-1
 \text{ integer squares},
\]
where $f$ is the matrix formal codegree and $b$ is the number of source
dual pairs.  The latter condition is imposed only at pair positions;
a self-dual label is not one of these $b$ positions.
\end{theorem}

\begin{proof}
Decompose
\[
 A=\frac s2\mathbf1+V+W,\qquad
 V^\dagger=V,\quad\operatorname{Trd}(V)=0,\quad W^\dagger=-W.
\]
The skew space has dimension one, so $W\in\mathbb Q U$.  For pure
quaternions $V,W$, the anticommutator is the scalar
$\operatorname{Trd}(VW)$.  Involution invariance and cyclicity of
reduced trace give
\[
 \operatorname{Trd}(VW)=\operatorname{Trd}((VW)^\dagger)
 =-\operatorname{Trd}(WV)=-\operatorname{Trd}(VW),
\]
so $V$ and $W$ anticommute; the same argument applies to $V$ and $U$.
Write $V^2=v\mathbf1$ and $W^2=w\mathbf1$.  Positivity gives $v\ge0$
and $w\le0$, and direct multiplication yields
\[
 n=\frac{s^2}{4}-v-w,\qquad t=\frac{s^2}{2}+2v-2w.
\]
Thus $J=4v$ and $K=-4w$.  If $J>0$, the generators $U$ and $2V$
anticommute and have squares $-\delta$ and $J$.  They generate $B$,
giving the asserted presentation.  If $J=0$, positivity forces $V=0$.

The image of a fusion-basis element satisfies a monic integral
polynomial.  Its rational characteristic polynomial has integral
coefficients, proving $s,n\in\mathbb Z$; $t$ is the integer character
value of $XX^*$.  For a self-dual label $W=0$.  For a pair difference
$A-A^\dagger=2W$, the relevant skew Gram diagonal is
$-\tfrac12\operatorname{Trd}((A-A^\dagger)^2)=K$.  Applying the
rank-one integral Gram factorization $K_{ij}=\delta z_i z_j$ and its
trace identity $\delta\sum z_i^2=f$ proves the final restrictions.
\end{proof}

\begin{proposition}[The analogous symplectic row obstruction]
\label{paper:symplectic-row-embedding}
For a positive symplectic rational degree-two component, retain the
notation $s,n,t$.  Then $t=2n$ and
$\Delta=s^2-4n\le0$.  If $\Delta=0$, the image is the scalar
$(s/2)\mathbf1$, so $s$ is even and $t=s^2/2$.  If $\Delta<0$, then
at every finite ramified prime $p$ of the component,
$\Delta$ is not a square in $\mathbb Q_p$.
\end{proposition}

\begin{proof}
The involution is quaternion conjugation on a definite quaternion
algebra, so $AA^\dagger=n\mathbf1$.  The norm of the pure part of
$A$ is $-\Delta/4$, proving the first assertions by positivity.
A nonscalar element generates the quadratic field
$\mathbb Q(\sqrt\Delta)$.  If this field split at a ramified prime,
its local algebra would contain a nontrivial idempotent embedded in
a division quaternion algebra, which is impossible.
\end{proof}

\paragraph{An exact extra coordinate for row enumeration.}
For a Galois-stable block set
$u_j=\operatorname{Tr}_j(X^2)$ and $w_j=t_j-u_j$.  Positivity gives
$w_j\ge0$, and the coefficient trace identities imply
\[
 \sum_j a_js_j=-d,\qquad
 \sum_j a_jt_j=5400-d^2,\qquad
 \sum_j a_jw_j=5400(1-\varepsilon),
 \quad\varepsilon=\delta_{X,X^*}.
\]
For the rational matrix block $w_E=K$.  For real linear blocks it
vanishes.  For quadratic nonreal blocks it is $2t_j-s_j^2$.
Where a higher complex field leaves $u_j$ unspecified, the third
coordinate is omitted, a safe weakening.
For any remaining collection of blocks with weighted degree mass
$M=\sum a_jk_j$, weighted trace $S$, weighted norm $T$, and weighted
skew coordinate $W$, Cauchy--Schwarz on the fixed parts gives
\[
                         2S^2\le M(2T-W).
\]
This allows exact pruning before the rows are assembled; it does not
add a hypothesis on the source dimensions.

\begin{theorem}[Centrality of a forced single remainder atom]
\label{paper:central-canonical-atom}
Let $A_E$ be a canonical constituent of degree $a$ and multiplicity
$m$, and write $F(A_E)=m\mathbf1+P_E$, with $R=a-m>0$.
If $A_E$ is self-dual, decompose $P_E$ into source duality atoms: a
self-dual simple $X$ of dimension $d_X$, or the sum $X+X^*$ of a
nonselfdual pair, of dimension $2d_X$.  Without self-duality, use
individual source simples as atoms.

Suppose $D$ contains every eligible positive atom dimension at most
$R$, and $R<2\min D$.  Then $P_E$ consists of exactly one atom of
dimension $R$, with multiplicity one, and its class is central in
$K_0(\mathcal C)$.  For the unique rational matrix component:
\begin{enumerate}
\item If the atom is an individual source simple $X$, its image
$A=E(X)$ is scalar.  In the row notation above this requires
$s\in2\mathbb Z$ and $t=s^2/2$.
\item If the atom is $X+X^*$, its image satisfies
$A+A^\dagger=sI$, equivalently $J=t-2n=0$.  Its skew coordinate $K$
need not vanish.
\end{enumerate}
Thus absence of an eligible central atom of the prescribed dimension
is an exact obstruction even when the ordinary remainder semigroup
contains $R$.
\end{theorem}

\begin{proof}
The half-braiding of $A_E$ provides
$F(A_E)\otimes Y\simeq Y\otimes F(A_E)$ for every source object $Y$.
Consequently its Grothendieck class, and hence the class of $P_E$, is
central.  All coefficients in the indicated atom decomposition are
nonnegative integers.  Two atoms, or two copies of one atom, would
have dimension at least $2\min D>R$.  Exactly one atom therefore occurs.

For a unique degree-two matrix component, centrality is
equivalent to a scalar image in the matrix block.  A rational scalar
image of an integral basis element is an integer multiple of $I$,
which gives the first condition.  Conversely, equality in the positive
trace inequality $s^2\le2t$ forces $A=(s/2)I$.
For a pair atom, write $A=(s/2)I+V+W$ as in
Theorem~\ref{paper:orthogonal-row-quaternion-class}.  Its sum with its
dual is $sI+2V$, so it is scalar exactly when $V=0$.  Positivity makes
this equivalent to $J=0$; the same conclusion holds automatically for
a symplectic involution, whose fixed subspace is scalar.
\end{proof}

Only centrality of the Grothendieck class is required: a summand of a
forgetful image is not asserted to inherit a half-braiding.  Eligible
atoms may also be restricted by the regular-character support condition
\(\lambda(X)\ge m\) of Lemma~\ref{paper:canonical-support-cover}.

For several degree-two components, centrality requires the same scalarity
condition in every component.  For a quadratic component put
\(\alpha=\operatorname{Tr}A\), \(\beta=\det A\), and
\(\gamma=\operatorname{Tr}(AA^*)\), with the bar denoting complex
conjugation on its field \(K\).  An individual atom requires
\[
 \alpha/2\in\mathcal O_K,\qquad 2\gamma=\alpha\overline\alpha.
\]
For a pair atom put
\[
 S=\alpha+\overline\alpha,\qquad
 T=\alpha^2+\overline\alpha^2-2\beta-2\overline\beta+2\gamma.
\]
It requires \(S/2\in\mathcal O_K\) and \(2T=S^2\).
Indeed \(S\) and \(T\) are the trace and squared norm of the
selfadjoint sum \(A+A^*\); equality in the positive trace inequality
forces that sum to be scalar.  The individual formula is the same
equality for \(A\).  Integral scalarity gives the stated half-trace
conditions.  The tests are componentwise; distinct rational linear
characters are not required to have a common scalar.

\subsection{Integral aggregate rows and finite polynomial supersets}
\label{paper:sec:aggregate-polynomials}

Fix a Galois-stable collection of $c$ irreducible representations
of common degree $m$, canonical degree $a$, and formal codegree
$f=5400/a$. Put $k=mc$. For a source simple $X$ define
\[
 s=\sum_E\chi_E(X),\qquad t=\sum_E\chi_E(XX^*),\qquad
 w=\sum_E\bigl(\chi_E(XX^*)-\chi_E(X^2)\bigr).
\]
All three numbers are integers. For a nonselfdual pair, $w$ is the
diagonal of the integral skew matrix obtained by summing the
orthogonal projections at this common formal codegree. Thus
\begin{equation}\label{paper:aggregate-skew-row}
 0\le w\le f,\qquad w\equiv t-s\pmod2,\qquad
 2s^2\le k(2t-w).
\end{equation}
For a selfdual label, $w=0$. The parity condition follows from the
first two Newton sums of the integral aggregate characteristic
polynomial. For the last inequality, the Hermitian parts of the
representation matrices have squared Frobenius norm $t-w/2$;
their aggregate trace is $s$, so Cauchy--Schwarz applies in
dimension $k$. For a nonselfdual row the additional global equation is
\[
                            \sum_{(a,m)}a w=5400.
\]
It is the difference between the unit coefficients of $XX^*$ and
$X^2$. For a selfdual row this sum is zero. These assertions apply
to repeated odd-dimensional characters as a Galois-stable
collection, without assuming that the individual characters are
rational.

\begin{lemma}[A complete finite integral polynomial superset]
\label{paper:aggregate-polynomial-superset}
For a Galois-stable aggregate as above, its degree-$k$
characteristic polynomial is monic integral. Write it as
\[
 P(T)=T^k+c_1T^{k-1}+\cdots+c_k.
\]
Then
\[
 c_1=-s,\qquad c_2=(s^2-t+w)/2,\qquad
 k^j c_j^2\le\binom{k}{j}^2t^j\quad(2\le j\le k),
\]
and $c_j=0$ whenever $j>t$. If $p_j$ is its $j$th Newton sum,
then $p_j^2\le t^j$ for every $j\ge2$. For a selfdual source
label, the Hankel matrix $(p_{i+j})_{0\le i,j<k}$ is positive
semidefinite.
\end{lemma}
\begin{proof}
Every characteristic root is an algebraic integer, because it
occurs in the integral multiplication matrix of $X$. Galois
stability makes the product polynomial rational, hence integral.
The first two formulas are Newton's identities. If its roots are
$\lambda_1,\ldots,\lambda_k$, the squared-modulus sum is at
most $t$, by Schur triangularization of the representation
matrices. Cauchy--Schwarz and the elementary symmetric mean
inequality give
\[
 |c_j|\le\binom{k}{j}(t/k)^{j/2}.
\]
Every nonzero irreducible factor of degree $e$ has nonzero integral
norm, so its roots contribute at least $e$ to the squared-modulus
sum. Therefore at most $t$ roots are nonzero, with multiplicity,
proving the asserted vanishing coefficients. For $j\ge2$,
\[
 |p_j|\le\sum_i|\lambda_i|^j
                \le\left(\sum_i|\lambda_i|^2\right)^{j/2}.
\]
Finally, selfduality makes the representation matrices Hermitian.
The Hankel matrix is then the Gram matrix of the real vectors
$(\lambda_h^i)_h$, proving its positivity.
\end{proof}

The lemma gives a finite, completely integral coefficient search.
For each admissible $w$, the first two coefficients are fixed;
the remaining coefficients range through the displayed integer
bounds. Newton recurrence supplies all requested power traces.
Polynomials satisfying these necessary tests need not be realized
by a representation matrix or by the specified full character
field; retaining them is a safe enlargement of the possible rows.

For every positive fusion word $W$, the regular trace also obeys
\[
                    \operatorname{Tr}L_W\ge r\,\tau(W).
\]
Indeed the unit term of $W$ contributes $r\tau(W)$ and all other
basis terms have nonnegative multiplication-matrix diagonals.
For $W=X^j$, the two traces in terms of an aggregate polynomial
choice are
\[
 5400\tau(X^j)=d^j+\sum_{(a,m)}a p_j,\qquad
 \operatorname{Tr}L_{X^j}=d^j+\sum_{(a,m)}m p_j.
\]
One common polynomial choice is required to give nonnegative
integral coefficient traces and the displayed regular-trace
inequalities for all tested powers. It must also give
$\tau(X^2)=\mathbf1_{X=X^*}$ and a positive return by power
eleven when the canonical length is eleven. This method includes
all Schur indices; it uses characteristic polynomials of the
complex irreducible representations, not a rational matrix
presentation of a simple algebra.

\subsection{A third-power residue test before polynomial completion}
\label{paper:sec:cubic-residue-filter}

The positive-word tests can be applied before all characteristic
polynomials have been completed. This is useful for low canonical
degrees, where materializing every later coefficient is expensive.

\begin{lemma}[Complete third-power residue supersets]
\label{paper:cubic-residue-superset}
Let a Galois-stable collection of representations of common formal
codegree $f=D/a$ have total matrix degree $k$. For a source simple
$X$, let $A$ be the direct sum of its matrices in this collection and
put
\[
 s=\operatorname{Tr}A,\qquad
 t=\operatorname{Tr}(AA^*),\qquad
 w=t-\operatorname{Tr}(A^2).
\]
These aggregate traces are integral. With $p_2=t-w$, every possible
third power sum has the form
\[
 p_3=\frac{3sp_2-s^3}{2}+3e_3,
\]
where $e_3=0$ for $k<3$ and, otherwise,
\[
 |e_3|\le
 \left\lfloor\sqrt{\frac{\binom{k}{3}^{\!2}t^3}{k^3}}\right\rfloor,
 \qquad |p_3|\le\lfloor\sqrt{t^3}\rfloor.
\]
Thus the allowed $w$ values and these integer intervals give a finite
necessary set $S\subseteq\mathbb Z/f\mathbb Z$ of residues for $p_3$.
Now let the collections indexed by $j$ partition all non-Perron
irreducible source representations, with each representation occurring
once and with common codegree $f_j=D/a_j$ within each collection.
Then every nonunit source row of dimension $d$ satisfies
\[
                 -d^3\in\sum_j a_jS_j
                    \quad\text{in }\mathbb Z/D\mathbb Z.
\]
Failure of this residue condition excludes the source row.
\end{lemma}

\begin{proof}
Galois stability makes the characteristic polynomial of $A$ integral;
this uses the direct sum of the complex representations and assumes
no splitting of their rational simple algebras. Its first elementary
symmetric coefficients are $e_1=s$ and
$e_2=(s^2-p_2)/2$. The third Newton identity gives
$p_3=sp_2-e_2s+3e_3$, which is the displayed expression.

If $\alpha_1,\ldots,\alpha_k$ are the eigenvalues, Schur triangularization
gives $\sum_i|\alpha_i|^2\le t$. Cauchy--Schwarz on the elementary
symmetric sum, followed by the elementary symmetric mean inequality,
gives
\[
 |e_3|^2\le\binom{k}{3}
          \sum_{i<j<\ell}|\alpha_i\alpha_j\alpha_\ell|^2
       \le\binom{k}{3}^{\!2}(t/k)^3.
\]
Also $|p_3|\le\sum_i|\alpha_i|^3\le t^{3/2}$.
Taking integer bounds proves the finite residue superset. For
$k<3$ the third elementary symmetric coefficient is zero.

The coefficient of the unit in $X^3$ is
\[
                   \tau(X^3)=\frac1D
                       \left(d^3+\sum_j a_jp_{3,j}\right)\in\mathbb Z.
\]
Since $a_jf_j=D$, each contribution depends only on the residue of
$p_{3,j}$ modulo $f_j$. This proves the last assertion.
\end{proof}

For a selfdual row, $w=0$. For a nonselfdual row, the computation
retains every $w$ admitted by its necessary skew conditions. The
residue test may allow different choices from another necessary
row test; this only enlarges the retained set. No completion of
$e_4,\ldots,e_k$, no irreducibility of a proposed characteristic
polynomial, and no existence of a source simple is inferred.

For a fixed $(k,s,t)$, all admitted $w$ have the same parity.
Changing $w$ by two changes the Newton base by $-3s$, so every
admitted $w$ contributes to the same coset modulo $\gcd(3,f)$.
The residue search may therefore stop once it has found
$f/\gcd(3,f)$ distinct residues.

The integer residue sets are combined by circular shifts and bitwise
unions in $\mathbb Z/5400\mathbb Z$. Their exactness is independently
checked on integer matrices by direct multiplication; the proof above
establishes the general necessity of the test, including nonsplit
components and nonrational individual characters.

\subsection{Orthogonal frames of character rows}
\label{paper:sec:orthogonal-frames}

Aggregate moments do not impose pairwise orthogonality between distinct
source rows. At a fixed full-field assignment this supplies a further
finite obstruction, before any fusion coefficients need be reconstructed.

\begin{proposition}[The character-frame obstruction]
\label{paper:character-frame-obstruction}
Suppose the source Grothendieck ring is commutative of rank $r$.
Fix its full character fields $K_1,\ldots,K_h$, one for each
non-Perron Galois orbit, and the corresponding canonical degrees $a_i$.
For every simple $X$ put
\[
 v_X=(d_X,\alpha_{1,X},\ldots,\alpha_{h,X})
         \in\mathbb Q\oplus\bigoplus_{i=1}^hK_i,
 \qquad \alpha_{i,X}=\chi_i(X)\in\mathcal O_{K_i}.
\]
The ambient space has rational dimension
$1+\sum_i[K_i:\mathbb Q]=r$. With the positive rational form
\[
 \langle u,v\rangle_a=u_0v_0+
       \sum_i a_i\Tr_{K_i/\mathbb Q}(u_i\overline{v_i}),
\]
one has
\[
 \langle v_X,v_Y\rangle_a=5400\,\delta_{X,Y},
       \qquad v_{\mathbf1}=(1,\ldots,1).
\]
Let $\mathcal V$ be any finite set containing every possible nonunit
row, including all conjugate orientations and all orderings of
repeated equal-field factors.
Join two distinct rows by an edge when their weighted inner product
is zero. If this graph has no clique of size $r-1$, the field
assignment is impossible.
\end{proposition}

\begin{proof}
Character values are algebraically integral, and
$\chi_i(X^*)=\overline{\chi_i(X)}$. The coefficient-trace formula yields
\[
 \tau(XY^*)=\frac1{5400}
   \left(d_Xd_Y+\sum_i a_i
       \Tr_{K_i/\mathbb Q}(\alpha_{i,X}\overline{\alpha_{i,Y}})\right).
\]
Distinct basis rows must therefore be distinct orthogonal vectors.
All $r-1$ nonunit basis elements would give the specified clique.
Column permutations may be fixed globally, but cannot be chosen
independently between rows when testing their inner products.
\end{proof}

When every $K_i=\mathbb Q$, all coordinates are integers and all source
labels are self-dual, since the faithful character embedding identifies
$X$ and $X^*$. The form then has diagonal
$\langle1,a_1,\ldots,a_{r-1}\rangle$. This is the case used in the
following exact calculation; no nonrational frame calculation is claimed
for that certificate.

\begin{conditionalproposition}[An exact obstruction for a surviving rational assignment]
\label{paper:rational-8370-frame}
Assume verification hypothesis $\mathsf V$ of
Section~\ref{paper:verification-status} for the finite scopes and
prerequisite exclusions specified in the proof below.
The all-rational commutative rank-thirteen assignment with weighted
nonunit canonical profile
\[
                    (900^4,675^2,150^2,75,50,12^2)
\]
is impossible.
\end{conditionalproposition}

\begin{proof}
The exact upstream row enumeration retains $410$ aggregate rows.
In each row every rational block has exactly one value multiset
compatible with its sum and squared norm. Expanding every within-block
permutation gives $648$ distinct ordered character rows. The verifier
checks their weighted squared norm $5400$ and orthogonality to the unit.

The resulting integer orthogonality graph has $13169$ edges. A compact
certificate successively removes eight vertices. In the remaining
graph at each removal, its neighborhood has a proper coloring with
respectively
\[
                         5,5,5,5,3,3,5,5
\]
colors. Any twelve-clique containing that vertex would require an
eleven-clique in its neighborhood, so the vertex cannot belong to one.
The residual $640$ vertices admit a proper coloring with eleven
colors. Hence no twelve-clique exists.

The certificate checker directly verifies that each displayed color
class has no pair with weighted inner product zero, that the classes
cover exactly the relevant neighborhood or residual vertex set,
and that there are no repeated vertices. It does not invoke a
maximum-clique solver or accept an optimization status as a proof.
The upstream checker regenerates the complete scalar multisets and
their expansions, so the conclusion is not based on an unaudited
candidate-row file.
\end{proof}

This argument is deliberately a necessary-table obstruction. Even an
orthogonal frame would still have to satisfy nonnegative integral
fusion coefficients and the categorical requirements. The remaining
full-field assignments are tested separately.

\subsection{Galois traces of small canonical remainders}
\label{paper:sec:trace-remainders}

Centrality of a source object in this subsection refers to its class
in $K_0(\mathcal C)$. For a number field $K$, write
\[
 \operatorname{Tr}_{K/\mathbb Q}(\mathcal O_K)=t_K\mathbb Z,
 \qquad t_K>0.
\]
The integer $t_K$ divides $[K:\mathbb Q]$, since the trace of
$1$ is the field degree. This trace ideal can contain more
information than the field degree alone.

\begin{theorem}[A trace-ideal factorial restriction]
\label{paper:canonical-trace-ideal-factorial}
Under the center hypothesis at dimension $n=5400$, let
$W\ne\mathbf1$ occur in $A$ with multiplicity $m$, put
$r=d_W-m>0$, and write $K=K_W$, $e=[K:\mathbb Q]$, and
$t=t_K$. Then, with the integer $\delta=er/t$,
\begin{equation}\label{paper:trace-factorial-divisor}
                 5400\mid t^{\,2\delta-1}(2\delta-1)!.
\end{equation}
More precisely, the sum
\[
 P_{\mathcal O}=\sum_{V\in\mathcal O_W}
                         \bigl(F(V)-m\mathbf1\bigr)
\]
is a positive central integral element of dimension $er$,
supported on nonunits. Its scalar eigenvalues are all in
$t\mathbb Z$, and the product of its distinct non-Perron
spectral gaps is divisible by $5400$.
In particular, if $5\nmid t$, then
\[
                            \frac{e}{t}\,r\ge6.
\]
\end{theorem}
\begin{proof}
Canonical multiplicities and dimensions are Galois invariant,
so each displayed remainder is a nonzero genuine object
supported on nonunit source simples. The forgetful image is
central. In the irreducible representation indexed by $Z_E$,
its scalar is
\[
 \sum_{V\in\mathcal O_W}\bigl(\chi_{Z_E}(V)-m\bigr)
 =\operatorname{Tr}_{K/\mathbb Q}
                           \bigl(\chi_{Z_E}(W)-m\bigr).
\]
Indeed symmetry of $S$ places $\chi_{Z_E}(W)$ in $K_W$, and
the unsigned Galois symmetry in the integral modular category
identifies the summands with its $e$ conjugates, including
repetitions when this entry generates a smaller field.
The scalar is therefore a multiple of $t$.
The strict norm bound places every non-Perron scalar in
\[
                  t\{-\delta+1,\ldots,\delta-1\}.
\]
The polynomial with these roots belongs to $\mathbb Z[x]$,
and at the Perron scalar $er=t\delta$ its value is
$t^{2\delta-1}(2\delta-1)!$. Apply
Theorem~\ref{paper:spectral-annihilator-divisor}. Its least
annihilator gives the more precise distinct-gap statement.
If $5\nmid t$ and $\delta\le5$, the factorial on the right
has five-valuation at most one, contradicting
$v_5(5400)=2$.
\end{proof}

\begin{proposition}[A residue obstruction at a rational canonical constituent]
\label{paper:rational-canonical-residue-obstruction}
Let $W$ be a nonunit rational-field canonical constituent,
put $d=d_W$, $m=m_W$, and $r=d-m$. Let $g$ be any positive
integer dividing $m_Ed_{Z_E}$ for every non-Perron canonical
constituent $Z_E\ne W$. Then an integer $\mu$ must exist with
\begin{equation}\label{paper:rational-canonical-residue-condition}
                 -r<\mu<r,\qquad r+md\,\mu\equiv0\pmod g.
\end{equation}
Thus the absence of such an integer excludes the canonical
profile. This applies without commutativity of the source ring.
\end{proposition}
\begin{proof}
The positive central remainder $P=F(W)-m\mathbf1$ has dimension
$r$, zero unit coefficient, and integer scalar eigenvalues
$\mu_E$, by rationality of the full field and $S$-symmetry.
The strict spectral bound gives $|\mu_E|<r$ away from the
Perron representation. The symmetrizing trace is
$\tau=\sum_E\chi_E/f_E$, so
\[
 0=5400\tau(P)=r+
       \sum_{E\ne\FPdim}m_Ed_{Z_E}\mu_E.
\]
Reducing modulo $g$ leaves the displayed condition for the
representation corresponding to $W$.
\end{proof}

\begin{proposition}[High degree-ten multiplicity forces several source simples]
\label{paper:high-ten-multiplicity-source-orbit}
Suppose $d_W=10$ and $m=m_W\ge5$. Then $m\in\{5,6,7\}$,
and
\[
                           F(W)=m\mathbf1+R
\]
for one central source simple $R$ of dimension $r=10-m$.
There are at least two distinct source simples of dimension
$r$. More precisely, all the distinct remainders
$F(gW)-m\mathbf1$ are central source simples of that dimension;
their number $\ell$ is a divisor of $[K_W:\mathbb Q]$ and
satisfies $\ell\ge2$. In particular
\[
 m=5\Longrightarrow n_5\ge2,\qquad
 m=6\Longrightarrow n_4\ge2,\qquad
 m=7\Longrightarrow n_3\ge2.
\]
\end{proposition}
\begin{proof}
Every nonunit source simple has dimension at least three.
The forgetful kernel is trivial, so $10-m\ge3$ and $m\le7$.
For $m\ge5$, the positive remainder has dimension at most
five, hence consists of a single simple. It is central because
$F(W)$ is central.
Galois symmetry acts on the central scalar functions, so the
distinct remainders form an orbit whose size divides
$[K_W:\mathbb Q]$. If they all coincided with $R$, every
scalar of $R$ would be fixed by $\operatorname{Gal}(K_W/\mathbb Q)$:
the scalars lie in $K_W$ by $S$-symmetry. They would therefore
be rational algebraic integers. The integral-spectrum factorial
bound would require $5400\mid(2r-1)!$, impossible for $r\le5$.
Thus $\ell\ge2$ and the claims follow.
\end{proof}

\begin{corollary}[Quadratic fields at high degree-ten multiplicity]
\label{paper:high-ten-quadratic-fields}
Under the preceding hypotheses, if $[K_W:\mathbb Q]=2$, then
\[
 K_W\in\bigl\{\mathbb Q(\sqrt5),\;
              \mathbb Q(\sqrt{-3}),\;
              \mathbb Q(\sqrt{-15})\bigr\}.
\]
In the real case the two distinct degree-$r$ remainders are
self-dual. In either imaginary case they form a nontrivial
dual pair. Consequently $m=7$ forces either two self-dual
source simples of degree three or a dual pair of degree-three
source simples in this quadratic case.
\end{corollary}
\begin{proof}
For a quadratic field of fundamental discriminant $\Delta$,
its integer ring has trace ideal $\mathbb Z$ if $\Delta$ is
odd and $2\mathbb Z$ if $\Delta$ is even. In the latter case
$e/t=1$, and Theorem~\ref{paper:canonical-trace-ideal-factorial}
would force $r\ge6$. Thus $\Delta$ is odd.
The source codegree of $W$ is $540$, so
$\operatorname{cond}(K_W)\mid1080$. The only nontrivial odd
fundamental discriminants whose absolute values divide $1080$
are $5,-3,-15$.
The two remainders are distinct by
Proposition~\ref{paper:high-ten-multiplicity-source-orbit}.
Real modular character fields give self-dual labels, so their
forgetful remainders are self-dual. For an imaginary quadratic
field its nontrivial automorphism is complex conjugation,
which sends $W$ to $W^*$ and $R$ to $R^*$. Distinctness then
makes this a nontrivial dual pair.
\end{proof}

\begin{corollary}[No odd field degree at high degree-ten multiplicity]
\label{paper:high-ten-even-field-degree}
If $d_W=10$ and $m_W\ge5$, then $[K_W:\mathbb Q]$ is even.
The stated rank bound at most $490$ consequently leaves
only the following possible field degrees:
\[
\begin{array}{c|l}
 m_W &[K_W:\mathbb Q]\\\hline
 5&2,4,6,8,12,16,18\\
 6&2,4,6,8,12\\
 7&2,4,6,8.
\end{array}
\]
These are upper-level necessary field degrees, not asserted
fields or categories.
\end{corollary}
\begin{proof}
The conductor bound embeds $K_W$ into
$\mathbb Q(\zeta_{1080})$, of degree $288=2^5 3^2$.
Its odd-degree subfields have degrees $1,3,9$. The rational
case is already excluded globally for central degree ten.
The other two are uniquely
$\mathbb Q(\zeta_9)^+$ and $\mathbb Q(\zeta_{27})^+$.
For $a=2,3$, put $L=\mathbb Q(\zeta_{3^a})$ and
$K=L^+$, so $[K:\mathbb Q]=e=3^{a-1}$.
The traces to $\mathbb Q$ of all integral powers of
$\zeta_{3^a}$ are multiples of $e$; explicitly they are
$0,-e,2e$. Since $\mathcal O_L=\mathbb Z[\zeta_{3^a}]$,
every trace from $\mathcal O_L$ is divisible by $e$.
For $x\in\mathcal O_K$ one has
$\operatorname{Tr}_{L/\mathbb Q}(x)=2\operatorname{Tr}_{K/\mathbb Q}(x)$.
The integer $e$ is odd, so $e$ divides every trace from
$\mathcal O_K$. The trace of $1$ is $e$, proving $t_K=e$.
The trace-ideal theorem again contradicts $r\le5$.
Thus the field degree is even.
Its full orbit contributes $e m_W^2$ to the canonical rank,
besides the unit. Hence $1+e m_W^2\le490$.
The even divisors of $288$ satisfying this inequality are
exactly those displayed.
\end{proof}

\subsection{A forbidden canonical residue pattern}
\label{paper:sec:forbidden-24-200-pattern}

The congruence $75\mid(d-h(X))$ associated with canonical degrees
$24$ and $200$ already appears in the canonical-length argument of
\cite[Section~30.1]{BDPfamily}. We first give an exact general
two-character modulus, then isolate that particular arithmetic as a
rank-independent forbidden pattern for the present simple source.

\begin{proposition}[A modulus from two rational canonical characters]
\label{paper:two-rational-canonical-modulus}
Let $\mathcal C$ be an integral nonpointed fusion-simple category with
$\mathcal C_{\mathrm{ad}}=\mathcal C$, of source dimension $n_0$. Suppose two
nonunit canonical constituents have degrees $a,b>1$, multiplicity one,
and rational corresponding source characters. Suppose $M\mid n_0$ and
every other nonunit canonical constituent has degree divisible by $M$.
For $p^k\parallel M$, put $\alpha_p=v_p(a)$ and $\beta_p=v_p(b)$, and
define
\begin{equation}\label{paper:two-character-forced-modulus}
 Q(a,b;M)=\prod_{p^k\parallel M}
 p^{\min\{k,\,\beta_p+
       \lceil\max\{0,k-\alpha_p-\beta_p\}/2\rceil\}}.
\end{equation}
Then every nonunit source simple has dimension strictly greater than
$Q(a,b;M)/2$. In particular these canonical data are impossible whenever
\[
                         Q(a,b;M)\ge2(a-1).
\]
The same test can be applied after interchanging $a$ and $b$.
\end{proposition}
\begin{proof}
Write $h,g$ for the rational linear characters corresponding to the
degrees $a,b$. For a nonunit simple $X$, let $d=d_X$, $h=h(X)$ and
$g=g(X)$. These character values are integers. The canonical dimension
identity gives $1+a+b\equiv0\pmod M$. The unit traces of $X$ and
$XX^*$ give
\[
 d+ah+bg\equiv0\pmod M,\qquad
 d^2+ah^2+bg^2\equiv0\pmod M.
\]
Every remaining degree block is Galois stable and has integral
aggregate character trace, so its contribution vanishes modulo $M$.

Set $t=d-h$ and $u=g-h$. The two congruences become
\[
                       t+bu\equiv0,\qquad
                       t^2+bu^2\equiv0\pmod M.
\]
Using $1+a+b\equiv0$ once more, these are equivalent to
$t\equiv-bu\pmod M$ and $M\mid ab u^2$. Hence, for each
$p^k\parallel M$,
\[
 v_p(u)\ge\left\lceil\max\{0,k-\alpha_p-\beta_p\}/2\right\rceil,
 \qquad
 v_p(t)\ge\min\left\{k,\beta_p+
       \left\lceil\max\{0,k-\alpha_p-\beta_p\}/2\right\rceil\right\}.
\]
Thus $Q(a,b;M)\mid(d-h)$. The strict character bound gives
$|h|<d$, so $0<d-h<2d$ and $d>Q(a,b;M)/2$.

The canonical constituent of degree $a$ forgets to
$\mathbf1+P$, where $P$ is a nonzero positive object supported on
nonunit simples and has dimension $a-1$. Some such simple therefore
has dimension at most $a-1$. This contradicts the lower bound when
$Q(a,b;M)\ge2(a-1)$.
\end{proof}

\begin{proposition}[The degree-twenty-four and degree-two-hundred pattern]
\label{paper:forbidden-canonical-24-200-pattern}
Under the center hypotheses at dimension $5400$, suppose that
$A=I(\mathbf1)$ has two nonunit constituents of degrees $24$ and
$200$, each of multiplicity one, and every other nonunit
canonical constituent has degree divisible by $225$.
Then no such category exists.
This obstruction does not require a specified source rank or
commutativity of its Grothendieck ring.
\end{proposition}
\begin{proof}
The general modulus in
Proposition~\ref{paper:two-rational-canonical-modulus} is
$Q(24,200;225)=75>2(24-1)$, so that proposition already applies.
The following direct congruence calculation records this particular
obstruction explicitly.
The degree-$24$ and degree-$200$ constituents are individually
Galois fixed, since their canonical degrees occur nowhere else.
Their source representations are therefore rational linear
characters; call them $h$ and $g$, respectively.
For any nonunit source simple $X$, put $d=d_X$, $h=h(X)$ and
$g=g(X)$. These values are integers.
The two unit-trace identities
$\tau(X)=0$ and $\tau(XX^*)=1$ give
\begin{equation}\label{paper:24-200-two-congruences}
 d+24h+200g\equiv0\pmod{225},\qquad
 d^2+24h^2+200g^2\equiv0\pmod{225}.
\end{equation}
Indeed every other nonunit canonical degree is divisible by
$225$; summing character traces over its Galois-stable degree
block gives integers. The Perron contribution is $d$ or $d^2$,
and $5400$ is divisible by $225$.

Modulo $25$ the first congruence gives $h\equiv d$.
Modulo $9$, eliminate $g$ between the two congruences.
This gives $6(d^2+dh+h^2)\equiv0\pmod9$, hence
$3\mid(d-h)$ because $d^2+dh+h^2\equiv(d-h)^2\pmod3$.
It follows that
\[
                              75\mid(d-h).
\]
The strict operator-norm theorem gives $|h|<d$ for this
non-Perron character. Thus $0<d-h<2d$, so every nonunit
source simple has dimension at least $38$.

On the other hand, the canonical degree-$24$ constituent $W$
has $F(W)=\mathbf1+P$ with $P$ a nonzero genuine object,
supported on nonunits, of dimension $23$. Every simple
constituent of $P$ would have dimension at most $23$.
This contradicts the bound just obtained.
\end{proof}

\subsection{Small rational canonical remainders as vector sums}
\label{paper:sec:central-remainder-geometry}

Canonical remainder tests can use centrality and rationality before
individual multiplication coefficients have been reconstructed.  The
following formulation treats the linear and matrix components together.
A \emph{source-duality atom} is either one selfdual nonunit simple or
the sum $X+X^*$ of a nonselfdual pair.

\begin{theorem}[Integral scalar targets for a rational canonical column]
\label{paper:rational-canonical-scalar-targets}
Under the standing center hypotheses at source dimension $n=5400$, suppose
that a nonunit constituent $Z$ of $I(\mathbf1)$ has full character
field $\mathbb Q$, dimension $a$, and canonical multiplicity $m$.
Then
\[
 P=F(Z)-m\mathbf1
\]
is a nonzero positive selfdual integral element supported on nonunits,
central in the source ring, of dimension $R=a-m$.  In every irreducible
source representation $E$, its image is an integer scalar
$\lambda_E I_{m_E}$, with $|\lambda_E|<R$ except for the Perron
representation, where $\lambda_{\rm FP}=R$.

Partition the irreducible source representations into their full
Galois orbits $j$.  Let $k_j$ be the size of an orbit and $m_j$ its
common representation degree, and put $n_j=k_jm_j$.  For a
source-duality atom $B$, define
\[
 s_j(B)=\sum_{E\in j}\chi_E(B),\qquad
 t_j(B)=\sum_{E\in j}\chi_E(B^2),\qquad
 q_j(B)=n_jt_j(B)-s_j(B)^2.
\]
These are integers, with $q_j(B)\ge0$.  If
$P=B_1+\cdots+B_h$ is its atom decomposition, listing repeated atoms
with their multiplicities, then, for every orbit $j$,
\begin{align}
 n_j&\mid\sum_{\alpha=1}^h s_j(B_\alpha),
 &\lambda_j&=\frac1{n_j}\sum_{\alpha=1}^h s_j(B_\alpha),
 \label{paper:central-atom-trace-divisibility}\\
 \sum_{\alpha=1}^h U_{j,\alpha}&=0,
 &\|U_{j,\alpha}\|^2&=q_j(B_\alpha),
 \label{paper:central-atom-vector-sum}
\end{align}
for vectors in a positive real inner product space.
\end{theorem}

\begin{proof}
The forgetful-image and canonical-character identities of
Proposition~\ref{paper:source:canonical-character-bridge} and
Proposition~\ref{paper:source:canonical-regular-trace-identity} give positivity,
centrality, selfduality and the stated dimension.  Simplicity of the
modular center also ensures $P\ne0$: if $F(Z)=m\mathbf1$, the
nonunit $Z$ belongs to the fusion subcategory of central objects whose
forgetful images are sums of the unit.  This subcategory is proper,
because the forgetful functor is dominant onto the nontrivial source,
contradicting simplicity of the center.  Symmetry of the
modular $S$-matrix places every scalar eigenvalue of $F(Z)$ in the
full field of $Z$: the dimensions occurring in the two column
normalizations are rational in the integral setting.  That field is
$\mathbb Q$, and the scalar is an algebraic integer, so it is an
integer.  Subtracting $m$ proves the same assertion for $P$.
Lemma~\ref{paper:source:strict-character-operator-norm}, applied to its
positive nonunit support, gives the strict bound.

Fix an orbit $j$ and choose unitary realizations of its complex
representations.  Since $B=B^*$, the matrix
\[
 A_j(B)=\bigoplus_{E\in j}E(B)
\]
is Hermitian of size $n_j$. Its ordinary trace is $s_j(B)$ and
its squared Hilbert--Schmidt norm is $t_j(B)$. Both are rational
Galois-invariant algebraic integers. Project off the common scalar
line and set
\[
 U_j(B)=\sqrt{n_j}\left(A_j(B)-\frac{s_j(B)}{n_j}I_{n_j}\right).
\]
Its squared norm is $q_j(B)$. Every conjugate representation has
the same integer scalar on $P$, so summing these projections gives
zero and summing their traces gives
\eqref{paper:central-atom-trace-divisibility}.
\end{proof}

A sufficient way to know that a canonical field is rational is that
its pair $(a,m)$ occurs only once: Galois symmetry preserves both
entries.  In applying the theorem, each index $j$ must be a
\emph{full Galois orbit}. A group of several inequivalent orbits
sharing the same $(a,m)$ need not have one common scalar on $P$.

\begin{corollary}[One-, two-, and three-atom certificates]
\label{paper:central-remainder-three-atoms}
In the preceding setting suppose every nonunit source dimension is
at least $d_{\min}$ and $R<4d_{\min}$.  A canonical remainder has
at most three source-duality atoms.  For each proposed atom multiset
of total dimension $R$, the trace divisibilities above and the
following conditions are necessary, separately for every full orbit:
\[
\begin{array}{c|l}
\text{number of atoms}&\text{necessary condition}\\\hline
1&q_1=0,\\
2&q_1=q_2,\\
3&2(q_1q_2+q_2q_3+q_3q_1)-(q_1^2+q_2^2+q_3^2)\ge0.
\end{array}
\]
Here $q_\alpha=q_j(B_\alpha)$. In particular, at dimension $5400$
these tests apply whenever $0<R\le11$.

The proposed integers $\lambda_j$ also satisfy
\begin{equation}\label{paper:central-remainder-distinct-gap}
 5400\ \Bigm|\ \prod_{\lambda\in\{\lambda_j:j\ne{\rm FP}\}}
                         (R-\lambda),
\end{equation}
where repeated scalar values occur only once in the product.
Failure of every eligible atom multiset gives an exact exclusion.
\end{corollary}

\begin{proof}
Each atom has dimension at least $d_{\min}$, so there are at most
three. For one or two vectors summing to zero, the first two
conditions are immediate. For three, write $U_3=-U_1-U_2$ and
apply Cauchy--Schwarz to
\[
 \langle U_1,U_2\rangle=(q_3-q_1-q_2)/2.
\]
This gives $(q_3-q_1-q_2)^2\le4q_1q_2$, equivalent to the table.
The tests are only necessary. Repeated occurrences of one actual
atom must still represent the same vector; forgetting this can
only enlarge the set of candidates.

For the last assertion put
$p(T)=\prod_{\lambda\in\{\lambda_j:j\ne{\rm FP}\}}(T-\lambda)$.
The strict bound ensures $p(R)\ne0$, and centrality gives
$p(P)=p(R)e_{\rm FP}$. The unit coefficient of $e_{\rm FP}$ is
$1/5400$, whereas $p(P)$ is an integral ring element. Hence
$5400\mid p(R)$, as also follows from the spectral annihilator
principle of Theorem~\ref{paper:spectral-annihilator-divisor}.
\end{proof}

\paragraph{Computing pair-atom data from a source row.}
For a source simple $X$ write
\[
 S_j=\sum_{E\in j}\chi_E(X),\qquad
 T_j=\sum_{E\in j}\chi_E(XX^*),\qquad
 P_{2,j}=\sum_{E\in j}\chi_E(X^2).
\]
For a selfdual $X$, the atom data are $(s,t)=(S_j,T_j)$. For a
nonselfdual pair they are
\[
 s_j(X+X^*)=2S_j,\qquad
 t_j(X+X^*)=2(T_j+P_{2,j}).
\]
The full Galois sums are real integers, and cyclicity of trace gives
the formula for the cross terms. In a rational degree-two component,
if $s=\Tr E(X)$, $t=\Tr(E(X)E(X)^*)$ and $n=\det E(X)$, this becomes
\[
 q(X+X^*)=4(t-2n).
\]
Thus a forced central pair has $t=2n$, recovering the earlier
central-atom condition while imposing it simultaneously in every
component. For a linear full orbit the quantity $kt-s^2$ also
checks that the sum becomes rational, which centrality alone would
not impose.

Every atom used in a column of multiplicity $m$ must also satisfy
the regular-character support threshold $\lambda(X)\ge m$.
No such threshold is imposed on unrelated source basis rows.

\begin{corollary}[A rational canonical constituent has remainder at least six]
\label{paper:rational-canonical-remainder}
Under the center assumption at dimension $5400$, let a
nonunit $Z$ occur in $A=I(\mathbf1)$ with multiplicity $m_Z$.
If $K_Z=\mathbb Q$, then
\[
                               d_Z-m_Z\ge6.
\]
More precisely, $5400$ divides $(2(d_Z-m_Z)-1)!$.
\end{corollary}
\begin{proof}
Write $F(Z)=m_Z\mathbf1+P$. The kernel statement for the
forgetful functor makes $P\ne0$; it is positive, central,
and supported on nonunits. In the irreducible representation
corresponding to a canonical constituent $Z_E$, the scalar
of $F(Z)$ is $\chi_{Z_E}(Z)$. Symmetry gives
\[
        \chi_{Z_E}(Z)=\frac{d_Z}{d_{Z_E}}\chi_Z(Z_E)\in\mathbb Q.
\]
It is an algebraic integer as a fusion-matrix eigenvalue,
so the scalar of $P$ is an integer. Apply
Theorem~\ref{paper:integral-central-spectrum-factorial}.
For $d_Z-m_Z\le5$, its factorial has five-valuation at most
one, while $v_5(5400)=2$. Thus the remainder is at least six.
\end{proof}

\subsection{Integer spectra of central remainders}
For the packet remainders constructed above, every scalar eigenvalue
is an integer.  This additional fact gives a smaller exact moment test
than the general annihilator recurrence.

\begin{theorem}[A finite prime-primary criterion]
\label{paper:source:primary-integral-moments}
Let $N$ be a positive integer and let $(w_i,\lambda_i)$ be a finite
collection of pairs of integers; the weights may be signed.  The
following conditions are equivalent:
\begin{enumerate}
\item $N\mid\sum_i w_i\lambda_i^k$ for every integer $k\ge0$;
\item for each prime power $p^v\Vert N$, each $t\in\{0,\ldots,p-1\}$,
and each $0\le j<v$, one has
\[
 p^{v-j}\ \Bigm|\
 \sum_{\lambda_i\equiv t\pmod p}
       w_i\left(\frac{\lambda_i-t}{p}\right)^j.
\]
\end{enumerate}
\end{theorem}

\begin{proof}
For necessity, every integer polynomial has a weighted sum divisible
by $p^v$.  The polynomial
\[
 Q_t(T)=1-(T-t)^{\varphi(p^v)}
\]
is congruent to one modulo $p^v$ on the residue class $t$ modulo $p$,
and to zero on the other classes.  Here
$\varphi(p^v)\ge v$, so the first assertion holds also when
$p\mid T-t$. Apply the polynomial condition to
$Q_t(T)(T-t)^j$ and divide the resulting divisibility by $p^j$.
For sufficiency, expand $(t+pu)^k$ in each residue class.
The terms of binomial degree $j<v$ are divisible by $p^v$ by the
asserted conditions, and the terms of degree at least $v$ already
contain $p^v$.  This holds for every prime power in $N$.
\end{proof}

\begin{theorem}[The optimal universal factorial window]
\label{paper:source:integer-spectrum-factorial-window}
For $N\ge2$, define
\[
 h(N)=\min\{j\ge1:N\mid j!\}.
\]
For arbitrary finite signed integer spectra, integrality of the
weighted moments of degrees $0,\ldots,h(N)-1$ implies integrality
of every weighted moment.  No shorter consecutive initial window
works uniformly for all signed integer spectra.
\end{theorem}

\begin{proof}
The monic polynomial
\[
 A_h(T)=T(T-1)\cdots(T-h+1)
\]
takes values divisible by $h!$, hence by $N$, at every integer,
including negative integers.  Multiplying $A_h(\lambda_i)$ by
$w_i\lambda_i^s$ and summing gives a recurrence modulo $N$
expressing the moment of degree $h+s$ in lower moments.
Induction proves sufficiency.

For optimality take $\lambda_i=i$, $0\le i<h$, and
\[
 w_i=(-1)^{h-1-i}\binom{h-1}{i}.
\]
Finite differences give zero moments in degrees below $h-1$,
whereas the moment of degree $h-1$ is $(h-1)!$, not divisible by
$N$ by the definition of $h(N)$.
\end{proof}

For the source denominator $n=5400$, the prime-primary criterion consists of the residue-class
congruences
\[
\begin{array}{c|c|c}
p&t& (j,\text{modulus})\\\hline
2&0,1&(0,8),\ (1,4)\\
3&0,1,2&(0,27),\ (1,9),\ (2,3)\\
5&0,1,2,3,4&(0,25),\ (1,5).
\end{array}
\]
The binary quadratic congruences are redundant because $u^2\equiv u$
modulo two.  Equivalently, $h(5400)=10$, so degrees zero through nine
test integrality of every weighted moment of an integer spectrum.
The weights for a central packet are $1$ at the Perron character and
$m_Ea_E$ elsewhere; their sum is $5400$, so degree zero is automatic.
This reduction requires integer scalar eigenvalues.  General source
rows must still satisfy the full positive-return test through $L$.

\subsection{Prime-primary integral spectral moments}
\label{paper:sec:integer-spectral-moments}

The next arithmetic criteria apply whenever a central element has an
integer spectrum.  They do not require its individual eigenspaces or the
source multiplication table to have been constructed.

\begin{corollary}[Ten initial moments at $5400$]
\label{paper:ten-integer-spectral-moments}
For $h(N)=\min\{j\ge1:N\mid j!\}$, as in
Theorem~\ref{paper:source:integer-spectrum-factorial-window},
one has $h(5400)=10$.  For a rational canonical column, the
zeroth weighted moment is already $5400$, so degrees $1$ through
$9$ are enough to test the integrality of every central spectral
moment.
\end{corollary}

\begin{proof}
The factorial $10!$ is divisible by $2^3\cdot3^3\cdot5^2$,
whereas $9!$ has only one factor of $5$.
\end{proof}

\paragraph{Scope of the window.}
This is a theorem about integer eigenvalues.  It does not replace
the full canonical-length return test for a general source simple,
whose matrix eigenvalues may be nonrational algebraic integers.
In the length-thirteen word verifier a positive return is allowed
through degree $13$, including the Perron contribution.  The
degree-$12$ cutoff of an early exploratory script is not used.

\begin{corollary}[The canonical masses at $5400$]
\label{paper:canonical-primary-masses}
Suppose a rational canonical constituent has dimension $a$ and
multiplicity $m$, and put $R=a-m$.  Its positive central remainder
$P=F(Z)-m\mathbf1$ has integer scalar eigenvalues
$\lambda_E$, with Perron value $R$.  The weights in Theorem~\ref{paper:source:primary-integral-moments} are
\[
 w_{\rm FP}=1,\qquad w_E=m_Ea_E,
 \qquad 1+\sum_{E\ne{\rm FP}}m_Ea_E=5400.
\]
At $5400=2^3\cdot3^3\cdot5^2$, the theorem gives $25$ congruences.
The two binary quadratic congruences are implied by the corresponding
linear ones, because $u^2\equiv u\pmod2$.  Thus the following
$23$ congruences suffice:
\[
\begin{array}{c|c|c}
p&t&\text{pairs }(j,\text{modulus})\\\hline
2&0,1&(0,8),\ (1,4)\\
3&0,1,2&(0,27),\ (1,9),\ (2,3)\\
5&0,1,2,3,4&(0,25),\ (1,5)
\end{array}
\]
The modulus applies to the weighted residue-class sum.  No minimality
of the number $23$ is asserted.
\end{corollary}

\begin{proof}
The rational-column integer-spectrum statement is
Theorem~\ref{paper:rational-canonical-scalar-targets}.  The coefficient
trace satisfies
\[
 [\mathbf1]P^k
 =\frac{R^k+\sum_{E\ne{\rm FP}}m_Ea_E\lambda_E^k}{5400}
 \in\mathbb Z.
\]
Apply the prime-primary moment theorem and the binary redundancy just noted.
\end{proof}

\subsection{Complete small central spectra and atom completions}
\label{paper:sec:small-central-spectra}

\begin{proposition}[A necessary bounded central-spectrum problem]
\label{paper:small-central-spectrum-problem}
Suppose a canonical group $(a,m,c)$ has $c=1$, and put $R=a-m$.
Then every non-Perron irreducible source representation $E$ has
an integer $\lambda_E$ such that
\[
 1-R\le\lambda_E\le R-1,\qquad
 R+\sum_{E\ne{\rm FP}}m_Ea_E\lambda_E=0,
\]
and
\[
 R^k+\sum_{E\ne{\rm FP}}m_Ea_E\lambda_E^k
       \equiv0\pmod{5400}\qquad(1\le k\le9).
\]
Writing each distinct non-Perron scalar once, one also has
\[
 5400\mid\prod_{\lambda\in\{\lambda_E:E\ne{\rm FP}\}}
                     (R-\lambda).
\]
Failure of this finite integer problem excludes the canonical input.
\end{proposition}

\begin{proof}
Galois symmetry preserves $(a,m)$, so $c=1$ makes the full
canonical field rational.  Apply
Theorem~\ref{paper:rational-canonical-scalar-targets} to the
unit-free positive central remainder.  The first displayed equality
is its vanishing unit coefficient.  The moment congruences follow
from Corollary~\ref{paper:ten-integer-spectral-moments}, and the
distinct-gap divisor is
Corollary~\ref{paper:central-remainder-three-atoms}.
\end{proof}

\paragraph{A complete one-sided enumeration.}
The program expands every group $(a_E,m_E,c_E)$ into $c_E$
separate scalar variables, each of weight $a_Em_E$.  In particular,
sharing $(a_E,m_E)$ never forces their scalars to be equal.  This
even permits distinct scalars within one full Galois orbit, thereby
enlarging the necessary search universe. A negative certificate checked
on this entire enlarged universe excludes the original cases. A witness
checked only before the enlargement must be rechecked on every added
assignment.

At a partial assignment let $g$ be the greatest common divisor of
$5400$ and the weights of the remaining variables.  Each residual
moment must vanish modulo $g$.  These congruences index all possible
values of the next scalar.  The exact residual first moment is also
bounded by $(R-1)$ times the sum of the remaining weights; the last
scalar is obtained from that exact first moment.  Equal-weight
variables may be ordered, since all conditions in this relaxed
problem are symmetric in them.  The resulting recursion therefore
loses no admissible spectrum.  A deterministic search budget may
stop a column calculation, but such a stop is recorded only as
``no exclusion.''  Every accepted negative certificate comes from
an exhausted finite search and is exhausted again in exact replay.

\begin{proposition}[Integral spectrum completion of small atom multisets]
\label{paper:small-central-atom-completion}
Let $P$ be the rational canonical remainder in the preceding
proposition.  Enumerate all necessary source-duality atoms that
can occur in $P$, including the regular-trace support threshold
of the canonical column.  If the smallest eligible atom dimension
is $d_0$ and $R<4d_0$, then every actual decomposition of $P$
occurs among the multisets of at most three such atoms of total
dimension $R$.

For every full rational group $(a_j,m_j,1)$, write an atom's
trace and square trace as $s_j,t_j$, and put
$q_j=m_jt_j-s_j^2$.  For a proposed multiset the zero-sum tests
of Corollary~\ref{paper:central-remainder-three-atoms} are necessary.
Moreover, if $S_j$ is its total trace in any aggregate group
$(a_j,m_j,c_j)$, there must exist $c_j$ integers
\[
 \lambda_{j,1},\ldots,\lambda_{j,c_j}\in[1-R,R-1],\qquad
 m_j\sum_{u=1}^{c_j}\lambda_{j,u}=S_j,
\]
whose complete weighted spectrum satisfies the preceding
proposition.  If every eligible atom multiset fails a necessary
test, the canonical input is impossible.
\end{proposition}

\begin{proof}
Every atom in an actual remainder is an eligible source row and
has dimension at least $d_0$.  The strict inequality on $R$ bounds
the number of atoms by three, counting repeated copies.  The
zero-sum geometry applies to each full rational group separately.
Centrality and rationality of the column make its scalar on each
individual representation an integer.  Adding traces over an
aggregate group proves the displayed sum equation, whether or
not the group splits into several Galois orbits.  The completed
spectrum is subject to the first proposition.
\end{proof}

\paragraph{Pair atoms and relaxed choices.}
For a nonselfdual source row write $s=\operatorname{Tr}E(X)$,
$t=\operatorname{Tr}(E(X)E(X)^*)$, and
$w=t-\operatorname{Tr}(E(X)^2)$.  The pair atom $X+X^*$ has
trace $2s$ and square trace $4t-2w$.  The program permits every
row-admissible value of $w$, independently between rational groups,
and permits repeated copies to make independent relaxed choices.
Both permissions enlarge the atom set.  The zero-sum geometry is
never imposed on a group with $c_j>1$ merely because its entries
have the same $(a_j,m_j)$.  For such a group the completion retains
separate integer scalars with their required sum.  Larger or
unfinished completion problems are retained without exclusion.

\subsection{A degree-five packet at common codegree six}
\label{paper:sec:degree-five-projection}

\begin{lemma}[A complement Cauchy inequality]
\label{paper:common-codegree-complement}
Let $\mathcal B$ be a Galois-stable collection of non-Perron
representations with common rational formal codegree $f$, and put
$m=\sum_{E\in\mathcal B}\dim E<f$. For a nonunit basis element $X$
write
\[
 s=\sum_{E\in\mathcal B}\chi_E(X),\qquad
 t=\sum_{E\in\mathcal B}\chi_E(XX^*).
\]
Then $s,t$ are integers, $t\ge|s|$, and
\[
                    s^2\le(f-m)(f-t).
\]
If equality holds, then $d_X=-s/(f-m)$.
\end{lemma}

\begin{proof}
The trace--norm inequality is
Theorem~\ref{paper:trace-norm-defect}. Let $e$ be the central projection
onto the indicated representations. The complementary projections
of $\mathbf1$ and $X$ have Gram matrix
\[
 \frac1f\begin{pmatrix}f-m&-s\\-s&f-t\end{pmatrix}.
\]
It is positive semidefinite, proving the inequality. Equality makes
the projection of $X$ equal to $-s/(f-m)$ times that of the unit.
The Perron component lies in the complement; evaluating it gives the
asserted dimension.
\end{proof}

\begin{theorem}[Five signed support labels]
\label{paper:degree-five-codegree-six-support}
Suppose every nonunit source dimension is at least three and a
Galois-stable non-Perron packet has common codegree six and total
ordinary representation degree five. Its aggregate character $\chi$
has exactly five nonzero nonunit values, all equal to $1$ or $-1$.
Both signs occur and this support is stable under source duality.
If $e$ is its central projection, then
\[
 W=6(1-e)=\mathbf1-\sum_{X\ne\mathbf1}\chi(X^*)X
\]
is an integral selfdual central element satisfying
\[
 W^2=6W,\qquad \tau(W)=1,\qquad \tau(WW^*)=6,\qquad
 \operatorname{FPdim}(W)=6.
\]
The coefficients of $W$ may have either sign.
\end{theorem}

\begin{proof}
The preceding inequality and $t\ge|s|$ imply $|s|\le2$. If $|s|=2$,
then $t=2$ and equality holds, giving $d_X=-s$, contrary to the lower
bound on nonunit dimensions. Therefore $s\in\{0,\pm1\}$.
Character orthogonality gives
\[
 \sum_X\chi(X)^2=6\cdot5=30,
 \qquad \chi(\mathbf1)=5,
 \qquad \sum_Xd_X\chi(X)=0.
\]
Exactly five nonunit terms remain. Their weighted sum is $-5$, which
is impossible if all signs are positive or if all are negative and
all five dimensions are at least three. Rationality gives
$\chi(X^*)=\chi(X)$, so the support is duality stable.
The central idempotent formula is
$e=\frac16\sum_X\chi(X^*)X$. It gives the displayed expression for
$W$ and all its identities; the Perron projection is in $1-e$.
\end{proof}

\begin{theorem}[A rank bound from the signed complement]
\label{paper:degree-five-complement-rank}
In the preceding setting put $q=\sum_{E\in\mathcal B}(\dim E)^2$
and let $r$ be the source rank. Then
\[
                             5r\ge6q+6.
\]
For the partitions of five the resulting bounds are
\[
\begin{array}{c|rrrrrrr}
\text{degrees}&5&4+1&3+2&3+1+1&2+2+1&2+1+1+1&1+1+1+1+1\\\hline
q&25&17&13&11&9&7&5\\
r\text{ at least}&32&22&17&15&12&10&8.
\end{array}
\]
These bounds require this packet. They are not uniform source-rank
bounds in its absence.
\end{theorem}

\begin{proof}
Let $N$ be left multiplication by $W$ in the orthonormal source
basis. It is an integral symmetric positive semidefinite matrix,
$N^2=6N$, with rank $r-q$ and trace $6(r-q)$. The Perron vector
$\mathbf d>0$ satisfies $N\mathbf d=6\mathbf d$. No diagonal entry
can be zero: positivity would make its row zero. Thus every diagonal
entry is at least one. The unit diagonal is exactly one, and there
are five entries $\pm1$ elsewhere in its row. At any such position,
a diagonal entry of one would give equality in the Gram
Cauchy inequality and make that row equal up to sign to the unit
row. Applying both rows to $\mathbf d$ would force dimension one or
negative dimension. Hence those five diagonals are at least two,
and $6(r-q)\ge r+5$.

Equality is impossible. To see this, realize $N$ as a Gram matrix
of vectors. Let $u$ be the unit vector for the unit label. The five
special vectors are then
$v_i=\varepsilon_i u+w_i$, where $\varepsilon_i=\pm1$,
$w_i\perp u$, and $\|w_i\|=1$. All remaining vectors are orthogonal
to $u$. Since the nonzero eigenvalues are six, the frame operator
on the span is $6I$. Applying it to $u$ gives
\[
                          \sum_{i=1}^5\varepsilon_iw_i=0.
\]
Each inner product of two $w_i$ is an integer of absolute value at
most one. Thus any two are equal up to sign or orthogonal. Their
classes are orthogonal unit axes, and cancellation on each axis
requires an even number of the five vectors. That is impossible.
The integral strict improvement is $6(r-q)\ge r+6$, as claimed.
\end{proof}

\section{Integral commutators and reconstruction}

\label{paper:source:commutators}
The following results apply to fusion rings themselves.  A center
hypothesis is needed only when the dimension-$5400$ canonical
restrictions are subsequently imposed.  Order the distinguished basis
once and for all and retain both forms
\[
 \beta(x,y)=\tau(xy),\qquad
 \langle x,y\rangle=\tau(xy^*).
\]
The first is a symmetric associative bilinear form and the second is
positive definite; on basis elements $\beta(X,Y)=\delta_{X,Y^*}$.

\subsection{The commutator tensor and its integral support}

\begin{theorem}[The commutator three-form]
\label{paper:source:commutator-three-form-audit}
For any fusion ring, set
\[
 c_{XYZ}=\tau([X,Y]Z)=N_{XY}^{Z^*}-N_{YX}^{Z^*}.
\]
These integers form an alternating three-form.  If the complex simple
components have degrees $m_E$ and formal codegrees $f_E$, then
\begin{equation}\label{paper:source:general-commutator-square-sum-audit}
 \sum_{X<Y<Z}c_{XYZ}^{\,2}
   =\frac13\sum_E f_E m_E(m_E^2-1).
\end{equation}
In particular the expression on the right is a nonnegative integer.
\end{theorem}

\begin{proof}
The coefficient expression proves integrality.  Antisymmetry in $X,Y$
is immediate, and cyclicity of $\tau$ gives
\[
 \tau([X,Y]Z)=\tau(XYZ-XZY)=\tau([Y,Z]X).
\]
This proves full alternation.

Consider the squared Hilbert--Schmidt norm of the bilinear commutator
map, using the positive coefficient-trace inner product.  In the source
orthonormal basis it equals
\[
 \sum_{X,Y}\|[X,Y]\|^2
  =\sum_{X,Y,Z}\bigl|\tau([X,Y]Z^*)\bigr|^2
  =6\sum_{X<Y<Z}c_{XYZ}^{\,2}.
\]
Changing to any complex orthonormal basis preserves this norm.  The
coefficient trace on the component $M_{m_E}(\mathbb C)$ is
$\Tr/f_E$, so the elements $\sqrt{f_E}\,e_{ab}$ are orthonormal.
Different simple components commute.  Within one component,
\[
 [\sqrt f e_{ab},\sqrt f e_{cd}]
     =f(\delta_{bc}e_{ad}-\delta_{da}e_{cb}).
\]
Summing its squared norms over $a,b,c,d$ gives
$2f(m^3-m)$: each of the two square terms sums to $fm^3$,
and their combined cross term is $-2fm$.  Summing over components
and dividing by six proves the formula.
\end{proof}

\begin{theorem}[The commutator Gram operator]
\label{paper:source:commutator-gram-spectrum}
For a fusion ring with ordered distinguished basis, put
\[
 c_{XYZ}=\tau([X,Y]Z),\qquad
 G_{XY}=\sum_{A<B}c_{XAB}c_{YAB}.
\]
Then $G$ is an integral positive semidefinite matrix. Its kernel is
the center of the complexified ring, and its image is the commutator
subspace $[R_{\mathbb C},R_{\mathbb C}]$. On the trace-zero part of a
simple component of degree $m_E$ and formal codegree $f_E$, it has
one eigenvalue $m_Ef_E$, of multiplicity $m_E^2-1$. Equivalently,
\[
 G_{XY}=\sum_E\left(
 m_E\chi_E(XY^*)-\chi_E(X)\overline{\chi_E(Y)}\right).
\]
In particular, with exactly two degree-two components its nonzero
spectrum is
\[
 (2f_1)^{(3)},\qquad(2f_2)^{(3)}.
\]
If $f_1=f_2=f$, then $G^2=2fG$ and $\operatorname{rank}G=6$.
For distinct $f_1,f_2$, the projectors onto their pure components are
\[
 P_i=\frac{G(G-2f_j I)}{4f_i(f_i-f_j)}\qquad(\{i,j\}=\{1,2\}).
\]
The tensor also obeys
\[
 c_{\mathbf1,Y,Z}=0,\qquad
 c_{X^*,Y^*,Z^*}=-c_{XYZ},\qquad
 \sum_Xd_Xc_{XYZ}=0.
\]
\end{theorem}

\begin{proof}
The displayed Gram definition proves integrality and positivity.
Its kernel consists exactly of the $x$ for which
$\tau([x,A]B)=0$ for all basis elements $A,B$. Nondegeneracy of
$\beta$ makes this equivalent to $[x,A]=0$ for every $A$.
The center is stable under the positive involution. Its orthogonal
complement for $\tau(xy^*)$ is consequently the commutator subspace,
proving the assertion about the image.

The Gram operator is half the Hilbert--Schmidt Gram operator of the
maps $\operatorname{ad}X$: summing over ordered pairs $A,B$ counts
each unordered pair twice, while replacing $B$ by $B^*$ merely
permutes the orthonormal basis. On $M_m(\mathbb C)$ the trace is
$\Tr/f$, so $\sqrt f\,e_{ab}$ is an orthonormal basis. A direct
matrix-unit calculation gives
\[
 \frac12\sum_{a,b}
 \langle[A,\sqrt f e_{ab}],[B,\sqrt f e_{ab}]\rangle
   =m\Tr(AB^*)-\Tr(A)\overline{\Tr(B)}.
\]
This is zero on scalars and equals $mf$ times the original positive
form on the traceless subspace. Summing over components proves the
formula and spectrum. The projector formulas follow by evaluating
their polynomials at $0,2f_1,2f_2$.

The unit commutes with every element. Taking duals reverses a
commutator; cyclicity of the coefficient trace proves the second
identity. Finally $\sum_Xd_XX$ is the Perron regular element and is
central, proving the weighted identity.
\end{proof}

The diagonal of $G$ is a row constraint, but its off-diagonal entries,
rank, and spectrum require one common alternating tensor.  This is the
additional compatibility needed to pass from independently admissible
rows to a multiplication table.

\begin{theorem}[The support cost of a one-matrix commutator]
\label{paper:source:noncentral-support-bound}
Suppose $R_{\mathbb C}\simeq\mathbb C^\ell\oplus M_2(\mathbb C)$,
with nonlinear formal codegree $f$. Let $\nu$ be the number of
noncentral source labels. Then the integral commutator three-form
has at least $4,6,8$ nonzero coordinates when $\nu=4,5,6$, respectively,
and at least $3(\nu-3)$ when $\nu\ge7$. In particular, for $f\ge6$,
\[
                         \nu\le\left\lfloor\frac{2f}{3}\right\rfloor+3.
\]
The assertion is independent of source rank and of orthogonal versus
symplectic matrix duality.
\end{theorem}

\begin{proof}
Project each source label to its traceless matrix part. The nonzero
vectors span a three-dimensional space. Their positive Perron
weighted sum is zero, since the regular Perron element is central.
Consequently their represented rank-three matroid has neither loops
nor coloops. Its bases are precisely the nonzero coordinates of the
commutator volume form.

For a connected rank-three matroid on $\nu$ elements the minimal-base
theorem gives at least $3(\nu-3)+1$ bases; see
\cite[Theorem~2.1]{FerroniMinimal}. If it is disconnected, the ranks
of its components partition three. In the $1+2$ case their sizes
are $a\ge2$, $b\ge3$, and $a+b=\nu$. Their numbers of bases are at
least $a$ and $2b-3$, so the product is at least $a(2b-3)$.
As a concave quadratic in $a$, its minimum is at an endpoint,
giving at least $\min\{4\nu-14,3(\nu-3)\}$. For $\nu\ge5$ this is
$3(\nu-3)$. In the $1+1+1$ case all component sizes are at least two;
their product is at least $4(\nu-4)$, which is at least $3(\nu-3)$
for $\nu\ge7$. At $\nu=6$ it gives eight. The cases $\nu=4,5$ give
four and six by the same connected/disconnected alternatives.
The bound for $\nu\ge7$ is sharp for
$U_{1,\nu-3}\oplus U_{2,3}$.

The squared norm of the integral three-form is $2f$ by
Theorem~\ref{paper:source:commutator-gram-spectrum}. Each nonzero integral
coordinate costs at least one. For $\nu\ge7$ this gives
$3(\nu-3)\le2f$; for $\nu\le6$ the displayed conclusion follows
directly from $f\ge6$.
\end{proof}

\subsection{One matrix component}
An aligned character table means that all character values refer to the
same ordered source basis.  Choosing a separate Galois conjugation for
each row would lose this alignment and does not supply the following
reconstruction data.

\begin{theorem}[Linear characters determine the one-matrix ring up to opposite]
\label{paper:source:one-matrix-reconstruction}
Let a fusion ring $R$ have complexified algebra
\[
                  R_{\mathbb C}\simeq\mathbb C^\ell\oplus M_2(\mathbb C).
\]
The full aligned table of its linear characters, their formal codegrees,
the formal codegree $f$ of the unique degree-two character, and source
duality determine all multiplication coefficients up to passage to the
opposite ring. This includes a nonsplit rational quaternion component.

More explicitly, let $\mathcal L$ be the set of linear characters,
including the Perron character, and define
\begin{align}
 K_{XY}&=f\left(\delta_{X,Y}-
       \sum_{\lambda\in\mathcal L}
          \frac{\lambda(X)\overline{\lambda(Y)}}{f_\lambda}\right),
       &s_X&=K_{X,\mathbf1},\label{paper:source:reconstruction-gram}\\
 H_{XY}&=2K_{X,Y^*}-s_Xs_Y.
       &&\label{paper:source:reconstruction-pure-gram}
\end{align}
For actual fusion data, $K$ is integral of rank four, $H$ is integral
of rank three, and there is an alternating rational trilinear form
$\Omega$ on the space represented by $H$ such that
\begin{equation}\label{paper:source:quaternion-volume-square}
 \Omega(X,Y,Z)^2
       =-\det\bigl(H_{AB}\bigr)_{A,B\in(X,Y,Z)}.
\end{equation}
Choose one nonsingular triple. Its value of $\Omega$ has the two
possible signs of the square root in
\eqref{paper:source:quaternion-volume-square}; every other value follows by
rational change of basis. With one such global sign, the product is
\begin{align}
 N_{XY}^{Z}
  ={}&\sum_{\lambda\in\mathcal L}
       \frac{\lambda(X)\lambda(Y)\overline{\lambda(Z)}}{f_\lambda}
       \notag\\
    &+\frac1{4f}\bigl(
       s_Xs_Ys_{Z^*}+s_XH_{Y,Z^*}+s_YH_{X,Z^*}
                  +s_{Z^*}H_{XY}+\Omega(X,Y,Z^*)\bigr).
 \label{paper:source:one-matrix-fusion-formula}
\end{align}
The other sign gives $N_{YX}^{Z}$. Since basis duality is an
anti-isomorphism, the two resulting based rings are isomorphic after
reversing multiplication and relabelling by duality.
\end{theorem}

\begin{proof}
The unique degree-two character is fixed by Galois conjugation, hence
rational. Its rational simple component $B$ is a central simple algebra
of degree two, either split or quaternionic. Work in any complex
splitting representation and write $A_X$ for the image of $X$.
The positive involution can be realized as an adjoint, so
$A_{X^*}=A_X^*$. The coefficient-trace identity
$\tau(XY^*)=\delta_{X,Y}$ gives
\[
                  K_{XY}=\Tr(A_XA_Y^*),\qquad s_X=\Tr A_X.
\]
These are integers because they are values of the rational irreducible
character on ring elements. Put
$U_X=A_X-(s_X/2)I$. Then $U_X$ is traceless and
$H_{XY}=2\Tr(U_XU_Y)$. The source basis spans the whole simple algebra,
so $K$ has rank four and its pure part $H$ has rank three.

For traceless degree-two elements,
$UV+VU=\Tr(UV)I$. Therefore
$\Omega(U,V,W)=4\Tr(UVW)$ is alternating. In a quaternion presentation
$B=(a,b)_{\mathbb Q}$, choose pure basis $i,j,ij$ with
$i^2=a$, $j^2=b$, and $ij=-ji$. On this basis,
\[
 H=\operatorname{diag}(4a,4b,-4ab),\qquad
 \Omega(i,j,ij)=-8ab.
\]
Thus $\Omega^2=-\det H$ on the basis and, by change of coordinates,
on every triple. This calculation is equally valid for a split
degree-two algebra.

For a nonsingular triple, all coordinates of every other $U_X$ are
obtained by solving its rational Gram equations. The value of an
alternating three-form on that basis fixes it everywhere, so its
single sign is the only freedom. Expanding
$A_X=(s_X/2)I+U_X$ gives
\[
 \Tr(A_XA_YA_Z)=\frac14
 \bigl(s_Xs_Ys_Z+s_XH_{YZ}+s_YH_{XZ}+s_ZH_{XY}
                            +\Omega(X,Y,Z)\bigr).
\]
Finally $N_{XY}^{Z}=\tau(XYZ^*)$ gives
\eqref{paper:source:one-matrix-fusion-formula}. Exchanging $X,Y$ preserves
the symmetric terms and reverses the alternating term, proving the
opposite-ring assertion.
\end{proof}

\begin{corollary}[Integer square and coefficient tests before categorification]
\label{paper:source:one-matrix-volume-integrality}
In the preceding setting, for every source-basis triple the integer
$-\det(H_{AB})_{A,B\in(X,Y,Z)}$ is a nonnegative integer square.
Moreover every value $\Omega(X,Y,Z)$ is integral. A proposed full
linear table can therefore be rejected by failure of integrality
or the rank conditions on $K,H$, by one failed square test, or by
failure of a coefficient in \eqref{paper:source:one-matrix-fusion-formula}
to be a nonnegative integer.
\end{corollary}

\begin{proof}
The formula for the triple trace gives
\[
 \Omega(X,Y,Z)=4\chi_M(XYZ)-
       (s_Xs_Ys_Z+s_XH_{YZ}+s_YH_{XZ}+s_ZH_{XY}).
\]
Every term on the right is integral. Now apply
\eqref{paper:source:quaternion-volume-square}.
\end{proof}

\begin{theorem}[One-matrix volume divisibility and primitive lattice]
\label{paper:source:one-matrix-commutator-lattice-audit}
Suppose
$R_{\mathbb C}=\mathbb C^\ell\oplus M_2(\mathbb C)$, and let $f$
be the formal codegree of its degree-two character.  Use the matrices
$K,H$, the vector $s$, and the volume $\Omega$ of
Theorem~\ref{paper:source:one-matrix-reconstruction}.  Then
\begin{equation}\label{paper:source:volume-divisibility-audit}
 \Omega(X,Y,Z)=2f c_{XYZ}\in 2f\mathbb Z.
\end{equation}
Consequently, for every source triple,
\[
 \frac{-\det(H_{AB})_{A,B\in(X,Y,Z)}}{4f^2}
          =c_{XYZ}^{\,2}
\]
is an integer square.  Define
\[
                  G=2K-ss^{\mathsf T}.
\]
Then $G$ is an integral positive semidefinite matrix of rank three,
$G^2=2fG$, and
\begin{equation}\label{paper:source:one-matrix-volume-budget-audit}
 \det G_{I,I}=(2f)^2c_I^2,
 \qquad \sum_{|I|=3}c_I^2=2f.
\end{equation}

Let $W=\operatorname{im}_{\mathbb Q}G$, let
$\Lambda=W\cap\mathbb Z^r$, and let
$h=\gcd\{c_I:|I|=3\}>0$.  Thus $\Lambda$ is a primitive
rank-three lattice in the standard Euclidean lattice $\mathbb Z^r$.
Writing $\det\Lambda$ for its \emph{Gram determinant}, one has
\begin{equation}\label{paper:source:commutator-lattice-determinant-audit}
                        h^2\det\Lambda=2f.
\end{equation}
In particular $h^2\mid 2f$ and $\det\Lambda\mid 2f$.
\end{theorem}

\begin{proof}
All linear characters vanish on a commutator.  In the unique matrix
component, write $A_X=(s_X/2)I+U_X$.  Since the pure triple trace is
alternating,
\[
 c_{XYZ}
   =\frac1f\Tr([U_X,U_Y]U_Z)
   =\frac2f\Tr(U_XU_YU_Z)
   =\frac{\Omega(X,Y,Z)}{2f}.
\]
This proves the divisibility and the strengthened square test.

The orthogonal matrix-component projector is $K/f$, so $K^2=fK$.
Also $s=Ke_{\mathbf1}$, $K_{\mathbf1,\mathbf1}=2$, and hence
\[
 Ks=fs,\qquad s^{\mathsf T}s=2f.
\]
It follows that $G^2=2fG$.  On the four-dimensional image of $K$,
the line spanned by $s$ is its kernel, and the orthogonal complement
has eigenvalue $2f$.  This proves the rank and positivity claims.
Equivalently,
\[
                  G_{XY}=2\Tr(U_XU_Y^*).
\]

Let $\sigma$ denote the rational positive involution of the
degree-two simple algebra, restricted to its pure three-dimensional
part.  It preserves trace and reverses multiplication.  Therefore
\[
 \Omega(\sigma U,\sigma V,\sigma W)
    =\Omega(W,V,U)=-\Omega(U,V,W).
\]
As $\Omega$ is a nonzero volume form, $\det\sigma=-1$.  On any
three pure vectors the Gram matrices for $2\Tr(UV)$ and
$2\Tr(U\sigma(V))$ therefore have opposite determinants.  This
also holds for dependent triples, when both determinants vanish.
Thus
\[
             \det G_{I,I}=-\det H_{I,I}
                          =\Omega_I^2=(2f)^2c_I^2.
\]
The sum of the principal three-minors of $G$ is the product of its
three nonzero eigenvalues, namely $(2f)^3$.  Division by $(2f)^2$
gives the stated volume budget.  It also follows directly from
\eqref{paper:source:general-commutator-square-sum-audit}.

To justify the lattice assertion, let $T:\mathbb Q^r\longrightarrow
B^0$ send $e_X$ to $U_X$, where $B^0$ is the pure part of the rational
degree-two algebra.  The positive Gram formula gives
$\ker G=\ker T$.  The three-form $c$ is a nonzero scalar multiple
of the pullback under $T$ of the volume form on $B^0$.  Identify
covectors and vectors using the standard Euclidean form on
$\mathbb Q^r$.  Then $c$ is a decomposable rational three-vector
spanning $\bigwedge^3 W$.

Choose an integral basis $v_1,v_2,v_3$ of $\Lambda$.  Since
$\Lambda$ is primitive, the exterior vector
$p=v_1\wedge v_2\wedge v_3$ is primitive in
$\bigwedge^3\mathbb Z^r$: its coordinate gcd is one.  This follows,
for example, by extending a basis of the saturated submodule
$\Lambda$ to a basis of $\mathbb Z^r$.  As $c$ and $p$ span the
same rational line and all coefficients of $c$ are integers,
$c=\pm h p$.  Cauchy--Binet now gives
\[
 \sum_I c_I^2
     =h^2\sum_I p_I^2
     =h^2\det\bigl((v_i\cdot v_j)_{i,j=1}^3\bigr)
     =h^2\det\Lambda.
\]
Combine this with the volume budget.  The unique degree-two
character, and hence its formal codegree, are rational; formal
codegrees of fusion rings are algebraic integers.  Thus $f$ is a
positive integer, and the divisibility conclusions follow.
\end{proof}

\begin{corollary}[The quaternion class and the orientation test]
\label{paper:one-matrix-quaternion-determined}
The same full linear data determine the rational degree-two simple
algebra. Moreover, the two orientation choices have identical
integrality and nonnegativity outcomes for the proposed multiplication
coefficients, so only one global sign needs to be tested.
\end{corollary}

\begin{proof}
Diagonalize the nonsingular three-dimensional form represented by $H$
over $\mathbb Q$, obtaining $\langle u,v,w\rangle$. Write
$uvw=-q^2$ with $q\ne0$ rational. Rescaling the third vector by
$uv/(2q)$ gives
\[
 \langle u,v,-uv/4\rangle
     =\langle4a,4b,-4ab\rangle,\qquad a=u/4,\quad b=v/4.
\]
The rational simple algebra is therefore the quaternion algebra
$(a,b)_{\mathbb Q}$, including its split possibility. Its isomorphism
class is independent of the chosen diagonalization. The assertion
about testing the sign follows from
$N^{-}_{XY}{}^Z=N^{+}_{YX}{}^Z$.
\end{proof}

\begin{corollary}[The codegrees are recoverable from the linear table]
\label{paper:one-matrix-codegrees-recovered}
The formal codegrees in the reconstruction theorem need not be
specified independently. They are determined by the full linear table:
\[
 f_\lambda=\sum_X|\lambda(X)|^2,\qquad
 f=\frac{2}{\displaystyle1-\sum_{\lambda\in\mathcal L}1/f_\lambda}.
\]
Thus the complete aligned linear character table and source duality
alone determine a one-matrix fusion ring up to opposite, and hence up
to based-ring isomorphism.
\end{corollary}

\begin{proof}
The first formula is linear-character orthogonality. The coefficient
trace of the unit is
$1=\sum_{\lambda\in\mathcal L}1/f_\lambda+2/f$, which gives the
second formula. Apply Theorem~\ref{paper:source:one-matrix-reconstruction}.
\end{proof}

\begin{corollary}[A lower bound for the pure-component lattice]
\label{paper:commutator-lattice-minimum}
In Theorem~\ref{paper:source:one-matrix-commutator-lattice-audit},
\[
                   \det\Lambda\ge4,\qquad h^2\le f/2.
\]
If $f=6$, then $h=1$ and $\det\Lambda=12$.
\end{corollary}

\begin{proof}
The image of $G$ is the pure matrix component and is orthogonal to
the Perron component. Every distinguished basis element has strictly
positive Perron value. Hence no coordinate unit vector $e_X$ belongs
to this image, and the integral lattice $\Lambda$ has no vector of
squared length one.

Suppose its Gram determinant were at most three. Choose an integral
basis and place its three vectors as rows of a $3\times r$ matrix.
Cauchy--Binet expresses the Gram determinant as the sum of the squares
of its maximal minors. At least one minor must be $1$ or $-1$.
An integral unimodular row change and a column permutation therefore
put the basis matrix into the form $[I_3\mid A]$, with $A$ integral.
Since $AA^{\mathsf T}$ is positive semidefinite,
\[
 3\ge\det(I_3+AA^{\mathsf T})
       \ge1+\Tr(AA^{\mathsf T}).
\]
Thus $A$ has at most two nonzero integral entries. One of its three
rows vanishes, and the corresponding lattice basis vector is a
coordinate unit vector, a contradiction. This proves
$\det\Lambda\ge4$. The identity $h^2\det\Lambda=2f$ gives
$h^2\le f/2$. When $f=6$, the positive square $h^2$ divides twelve
and is at most three, so $h=1$.
\end{proof}

\subsection{Local coefficient tests from Cayley--Hamilton}
\label{paper:ch:sec:local-coefficients}

Several coefficients can be recovered before a complete character table
has been assembled.  This gives local necessary conditions on individual
rows, pairs of rows, and triples of rows.  Throughout this subsection,
$R$ is a fusion ring with distinguished basis $\{X_i\}_{i\in I}$,
$X_0=\mathbf1$, basis involution $i\mapsto i^*$, and
\[
 R\otimes_{\mathbb Z}\mathbb C\simeq
       \mathbb C^{\ell}\oplus M_2(\mathbb C).
\]
Let $\mathcal L$ denote all linear characters, including the positive
dimension character.  Write $f_\lambda$ for their formal codegrees and
$F$ for the formal codegree of the unique degree-two character.  With
$M_i$ the image of $X_i$ in a complex splitting of the latter component,
the coefficient-of-unit trace has decomposition
\begin{equation}\label{paper:ch:coefficient-trace}
 \tau(x)=\sum_{\lambda\in\mathcal L}\frac{\lambda(x)}{f_\lambda}
                    +\frac{\operatorname{tr}M(x)}{F}.
\end{equation}
Set $s_i=\operatorname{tr}M_i$, $\nu_i=\det M_i$, and
\[
 A_{ij}=\sum_{\lambda\in\mathcal L}
             \frac{\lambda(X_i)\lambda(X_j)}{f_\lambda},
 \qquad
 C_{ijk}=\sum_{\lambda\in\mathcal L}
             \frac{\lambda(X_i)\lambda(X_j)\lambda(X_{k^*})}
                  {f_\lambda}.
\]
The character ordering is common to every row under consideration.
The equalities $i=j$ and $i=j^*$ below are equalities of basis labels;
they cannot be inferred from equality of character values.

\begin{theorem}[Recovery of square and power coefficients]
\label{paper:ch:two-row-recovery}
The integers $s_i$ and $\nu_i$ are determined by one complete linear
character row and its specified basis duality:
\begin{equation}\label{paper:ch:trace-determinant-recovery}
 s_i=F(\delta_{i0}-A_{i0}),\qquad
 \nu_i=\frac{s_i^2-F(\delta_{i,i^*}-A_{ii})}{2}.
\end{equation}
For $q_i(T)=T^2-s_iT+\nu_i$, one has
\begin{align}
 N_{ii}^{j}
 &=\sum_{\lambda\in\mathcal L}
       \frac{q_i(\lambda(X_i))\lambda(X_{j^*})}{f_\lambda}
          +s_i\delta_{ij}-\nu_i\delta_{j0}
       \label{paper:ch:square-coefficient}\\
 &=C_{iij}+s_i(\delta_{ij}-A_{i,j^*})
                         -\frac{\nu_i s_j}{F}.
       \label{paper:ch:square-coefficient-projector}
\end{align}
More generally, let $p\in\mathbb Z[T]$ and divide it by $q_i$:
\[
              p(T)=h(T)q_i(T)+uT+v,
              \qquad h\in\mathbb Z[T],\quad u,v\in\mathbb Z.
\]
The coefficient of $X_j$ in $p(X_i)$ is
\begin{equation}\label{paper:ch:polynomial-coefficient}
 [p(X_i):X_j]=
   \sum_{\lambda\in\mathcal L}
       \frac{\bigl(p(\lambda(X_i))-u\lambda(X_i)-v\bigr)
                         \lambda(X_{j^*})}{f_\lambda}
            +u\delta_{ij}+v\delta_{j0}.
\end{equation}
Thus every power coefficient $[X_i^m:X_j]$ is determined by two
aligned linear rows, without choosing coordinates or an orientation
for the degree-two component.  The statements hold for both split
and nonsplit rational degree-two components.
\end{theorem}

\begin{proof}
Galois conjugation preserves irreducible representation degrees.  Since
there is only one degree-two character, it is rational.  Its values
on integral ring elements are rational algebraic integers, hence
integers.  The eigenvalues of $M_i$ are algebraic integers: they occur
among the eigenvalues of the integral matrix of multiplication by
$X_i$ in the regular representation.  Their product $\nu_i$ is
therefore an algebraic integer.  It is rational because
$2\nu_i=s_i^2-\operatorname{tr}(M_i^2)$, so $\nu_i\in\mathbb Z$.

Apply \eqref{paper:ch:coefficient-trace} to $X_i$ and $X_i^2$, using
$\tau(X_i)=\delta_{i0}$ and $\tau(X_i^2)=\delta_{i,i^*}$.
This gives \eqref{paper:ch:trace-determinant-recovery}.  Cayley--Hamilton
annihilates the matrix image of $q_i(X_i)$.  Consequently this element
belongs to the sum of the linear simple components.  Taking
$\tau(q_i(X_i)X_{j^*})$ gives \eqref{paper:ch:square-coefficient}.
The identity
\[
 \sum_{\lambda\in\mathcal L}
           \frac{\lambda(X_{j^*})}{f_\lambda}
                  =\delta_{j0}-\frac{s_j}{F}
\]
gives the second expression.  Here $s_{j^*}=s_j$, since the character
is rational and the positive involution acts as an adjoint in a
unitary realization.  Finally $p(X_i)-uX_i-v\mathbf1$ has zero matrix
image.  The same coefficient-trace argument proves
\eqref{paper:ch:polynomial-coefficient}.
\end{proof}

\begin{corollary}[Norm products in the symplectic case]
\label{paper:ch:symplectic-norm-coefficients}
Suppose the involution on the rational degree-two component is
symplectic.  If $t_i=\operatorname{tr}(M_iM_{i^*})$, then
$t_i=2\nu_i$ and
\begin{equation}\label{paper:ch:norm-coefficient}
 N_{i,i^*}^{j}=
  \sum_{\lambda\in\mathcal L}
   \frac{\bigl(\lambda(X_i)\lambda(X_{i^*})-\nu_i\bigr)
                   \lambda(X_{j^*})}{f_\lambda}
                  +\nu_i\delta_{j0}.
\end{equation}
The tests at $j=i$ and $j=i^*$ add no condition to the corresponding
square test: in every fusion ring,
\begin{equation}\label{paper:ch:unary-norm-redundancy}
                  N_{i,i^*}^{i}=N_{i,i^*}^{i^*}=N_{ii}^{i}.
\end{equation}
No analogous redundancy of all distinct-row norm-product tests is
asserted.
\end{corollary}

\begin{proof}
The symplectic involution on a degree-two algebra is reduced
conjugation.  Thus $M_{i^*}=s_iI-M_i$ and
$M_iM_{i^*}=\nu_iI$.  Subtract $\nu_i\mathbf1$ from $X_iX_{i^*}$
and use the coefficient trace, as above.  For the last assertion,
$X_iX_{i^*}$ is selfdual, so its coefficients at $i$ and $i^*$ agree;
cyclicity gives
$\tau(X_iX_{i^*}X_i)=\tau(X_i^2X_{i^*})$.
\end{proof}

\begin{corollary}[Local integrality, positivity, and dimension bounds]
\label{paper:ch:local-coefficient-bounds}
Let $d_i$ be the positive dimension of $X_i$.  Each value computed in
\eqref{paper:ch:square-coefficient} is a nonnegative integer and satisfies
\[
               N_{ii}^{j}\le d_j,\qquad
               d_jN_{ii}^{j}\le d_i^2.
\]
For any set $S$ of distinct nonunit basis labels,
\begin{equation}\label{paper:ch:partial-square-budget}
          \sum_{j\in S}d_jN_{ii}^{j}
                    \le d_i^2-\delta_{i,i^*}.
\end{equation}
The symplectic norm-product coefficients satisfy the same individual
bounds and the stronger unit-subtracted budget
\begin{equation}\label{paper:ch:partial-norm-budget}
          \sum_{j\in S}d_jN_{i,i^*}^{j}\le d_i^2-1.
\end{equation}
\end{corollary}

\begin{proof}
Nonnegative integrality is a based-ring axiom.  Apply the positive
dimension character to $X_i^2$ or $X_iX_{i^*}$, whose unit
coefficients are $\delta_{i,i^*}$ and $1$, respectively.  This proves
the dimension budgets.  Frobenius reciprocity gives
$N_{ii}^{j}=N_{i^*,j}^{i}$; comparison with the dimension of
$X_{i^*}X_j$ yields $d_iN_{ii}^{j}\le d_id_j$.
The norm-product inequality follows in the same way.
\end{proof}

\begin{proposition}[Three-row tests without an orientation choice]
\label{paper:ch:three-row-coefficient-pair}
Define
\begin{align}
 J_{ijk}={}&2C_{ijk}
       +s_i(\delta_{jk}-A_{j,k^*})
       +s_j(\delta_{ik}-A_{i,k^*})\notag\\
       &+s_k(\delta_{i,j^*}-A_{ij})-\frac{s_is_js_k}{F},
       \label{paper:ch:jordan-coefficient}\\
 H_{ab}={}&2F(\delta_{a,b^*}-A_{ab})-s_as_b,
       \qquad
 \Delta_{ijk}=-\frac{\det(H_{ab})_{a,b\in(i,j,k^*)}}{4F^2}.
       \label{paper:ch:local-commutator-square}
\end{align}
Then
\begin{equation}\label{paper:ch:coefficient-pair-data}
 J_{ijk}=N_{ij}^{k}+N_{ji}^{k},\qquad
 \Delta_{ijk}=(N_{ij}^{k}-N_{ji}^{k})^2.
\end{equation}
In particular, $\Delta_{ijk}$ is a nonnegative integer square.
Writing $q=\sqrt{\Delta_{ijk}}$, the two coefficients, without an
ordering, are
\begin{equation}\label{paper:ch:unordered-coefficient-pair}
             \left\{\frac{J_{ijk}+q}{2},
                          \frac{J_{ijk}-q}{2}\right\}.
\end{equation}
Consequently necessary local tests are
\[
 J_{ijk}\in\mathbb Z,\qquad
 J_{ijk}\ge q,\qquad J_{ijk}\equiv q\pmod2,
\]
and
\begin{equation}\label{paper:ch:three-row-dimension-bound}
 \frac{J_{ijk}+q}{2}
       \le \min\left\{\frac{d_id_j}{d_k},
                        \frac{d_id_k}{d_j},
                        \frac{d_jd_k}{d_i}\right\}.
\end{equation}
These quantities require only the three indicated aligned linear rows,
their dual labels, and the formal codegrees.
\end{proposition}

\begin{proof}
For degree-two matrices $A,B$,
\[
 AB+BA=(\operatorname{tr}A)B+(\operatorname{tr}B)A
       +\bigl(\operatorname{tr}(AB)
                      -\operatorname{tr}A\operatorname{tr}B\bigr)I.
\]
Apply this with $A=M_i$, $B=M_j$, multiply by $M_{k^*}$, and take
the trace.  The identity
\[
       \operatorname{tr}(M_aM_b)
                          =F(\delta_{a,b^*}-A_{ab})
\]
then proves the formula for $J_{ijk}$.

Put $U_a=M_a-(s_a/2)I$.  The matrices $U_a$ are traceless and
$H_{ab}=2\operatorname{tr}(U_aU_b)$.  For three traceless degree-two
matrices $U,V,W$,
\[
       \bigl(4\operatorname{tr}(UVW)\bigr)^2
                   =-\det\bigl(2\operatorname{tr}(AB)\bigr)_{A,B\in(U,V,W)}.
\]
Indeed $UV+VU=\operatorname{tr}(UV)I$ makes
$\operatorname{tr}(UVW)$ alternating.  Both sides of the displayed
identity change by the square of the determinant under a change of
the three vectors.  On the basis
\[
 \begin{pmatrix}1&0\\0&-1\end{pmatrix},\qquad
 \begin{pmatrix}0&1\\1&0\end{pmatrix},\qquad
 \begin{pmatrix}0&1\\-1&0\end{pmatrix}
\]
the Gram matrix is $\operatorname{diag}(4,4,-4)$ and the alternating
trace is $-8$, proving the identity, including dependent triples.
It is a polynomial identity in a splitting algebra and therefore
also applies to nonsplit rational forms.

Linear characters vanish on commutators.  Hence
\[
 N_{ij}^{k}-N_{ji}^{k}
 =\frac{\operatorname{tr}([U_i,U_j]U_{k^*})}{F}
 =\frac{2\operatorname{tr}(U_iU_jU_{k^*})}{F}.
\]
This proves the formula for $\Delta_{ijk}$ and the integrality,
parity, and positivity conditions.  Finally Frobenius reciprocity
identifies $N_{ij}^{k}$ with $N_{i^*,k}^{j}$ and $N_{k,j^*}^{i}$.
Applying the dimension character to these three products proves
\eqref{paper:ch:three-row-dimension-bound}; it applies equally to
$N_{ji}^{k}$.
\end{proof}

\paragraph{Exact use at dimension $5400$.}
Under the standing center hypotheses, put $n=5400$ and
$a_{\mathcal O}=n/f_{\mathcal O}$ for each non-Perron Galois orbit
of linear characters.  These are integers.  If $\alpha_i$ is the
value of a representative character in its full field $K_{\mathcal O}$,
set
\begin{align*}
 B_{ij}^{(2)}={}&(d_i^2-s_id_i+\nu_i)d_j\\
 &+\sum_{\mathcal O}a_{\mathcal O}
       \operatorname{Tr}_{K_{\mathcal O}/\mathbb Q}
        \bigl((\alpha_i^2-s_i\alpha_i+\nu_i)
                                   \overline{\alpha_j}\bigr).
\end{align*}
Then
\[
 N_{ii}^{j}=\frac{B_{ij}^{(2)}}{5400}
                         +s_i\delta_{ij}-\nu_i\delta_{j0}.
\]
All operations are exact integer field traces.  Thus divisibility
$5400\mid B_{ij}^{(2)}$ can be tested before positivity or the
dimension inequalities.  This divisibility concerns coefficient
numerators, not the source dimensions $d_i$.

Rows with a common fixed part may have several imaginary variants.
A class is eliminated by a unary test only when every variant fails;
an edge between two classes is eliminated only when every variant
pair fails a necessary test.  The two dual orientations and both
directions of a square-coefficient test must be included.  Conversely,
pairwise existence of compatible variants need not give a single
compatible choice for an entire table.  Local tests therefore precede,
and do not replace, the global character orthogonality, matrix-component
rank, common orientation, and based-ring verification conditions.

These identities are consequences of degree-two reconstruction, not
independent existence or rank theorems.  Their role is to enforce exact
multiplication constraints before a complete table is available.


\subsection{Two matrix components}
With two degree-two components the linear table leaves a two-dimensional
nonlinear center.  The integral commutator tensor determines its possible
attachments to the two simple Lie ideals.

\begin{theorem}[At most two two-matrix products for fixed linear and commutator data]
\label{paper:source:two-matrix-reconstruction}
Suppose a fusion ring satisfies
\[
 R_{\mathbb C}\simeq\mathbb C^\ell\oplus M_2(\mathbb C)^2,
\]
and its two nonlinear formal codegrees are rational. The following
data determine at most two multiplication tables on the given basis:
\begin{enumerate}
\item the full aligned table of all linear characters;
\item the unit and basis duality;
\item the integral commutator tensor $c_{XYZ}$.
\end{enumerate}
The two nonlinear codegrees are recovered from the nonzero spectrum
of $G$. The conclusion includes nonsplit rational quaternion
components and a quadratic full field for the nonlinear orbit.
It is conditional on the full specified data; it does not construct
those data or assert that either candidate has nonnegative integral
coefficients.
\end{theorem}

\begin{proof}
The linear formal codegrees and primitive central idempotents are
known from character orthogonality:
\[
 f_\lambda=\sum_X|\lambda(X)|^2,
 \qquad e_\lambda=\frac1{f_\lambda}
                       \sum_X\lambda(X^*)X.
\]
Put $e=\mathbf1-\sum_\lambda e_\lambda$. The commutator tensor and
basis duality recover the full Lie bracket,
\[
                    [X,Y]=\sum_Zc_{X,Y,Z^*}Z.
\]
They therefore recover its center $Z(R_{\mathbb C})$ and its
semisimple commutator Lie algebra
$L\simeq\mathfrak{sl}_2(\mathbb C)\oplus\mathfrak{sl}_2(\mathbb C)$.
The two simple Lie ideals are intrinsic and can be found from the
commutant of the adjoint representation. Their dimensions are three.

Inside the center, let
\[
 W=Z(R_{\mathbb C})\cap\bigcap_{\lambda\ \mathrm{linear}}\ker\lambda.
\]
This space has dimension two, is defined over $\mathbb Q$, and its
unit is $e$. The restriction of $\beta$ to it is nondegenerate.
Choose a nonzero rational $v\in W$ with $\tau(v)=0$ and write
\[
 g=\beta(v,v)\ne0,\qquad
 t_i=2/f_i,\qquad t=t_1+t_2=\tau(e)>0.
\]
There is such a $v$, since $W$ has dimension two and $\tau(e)>0$;
nondegeneracy and $\beta(e,v)=\tau(v)=0$ imply $g\ne0$.

The actual nonlinear primitive idempotents must have traces $t_1,t_2$,
sum $e$, and be mutually $\beta$-orthogonal. Hence they have exactly
the form
\begin{equation}\label{paper:binary-central-idempotents}
 e_1=\frac{t_1}{t}e+uv,\qquad
 e_2=\frac{t_2}{t}e-uv,\qquad
 u^2=\frac{t_1t_2}{tg}.
\end{equation}
Indeed $\beta(e_i,e_i)=\tau(e_i)=t_i$ gives this quadratic equation.
If $f_1\ne f_2$, there are at most two ordered central decompositions.
If $f_1=f_2$, the two signs merely exchange $e_1,e_2$, giving one
unordered pair of central idempotents.

It remains to attach each $e_i$ to one simple Lie ideal $L_i$.
The restriction of the Killing form of $L_i$ is
\[
                         \kappa_i=4f_i\,\beta|_{L_i}.
\]
This follows from $\kappa(A,B)=4\Tr(AB)$ on traceless two-by-two
matrices and $\beta=\Tr/f_i$ on the corresponding algebra block.
Thus distinct codegrees determine the attachment uniquely; with
equal codegrees there are at most two attachments.

Once the attachment is fixed, all products are forced. The central
idempotents multiply orthogonally, act as the identity on their
attached ideals and annihilate the other ideals. For $x,y\in L_i$,
\begin{equation}\label{paper:pure-product-reconstruction}
 xy=\frac{f_i}{2}\beta(x,y)e_i+\frac12[x,y].
\end{equation}
This is the identity $AB+BA=\Tr(AB)I$ for traceless two-by-two
matrices. Different simple ideals have zero associative product.
The linear components multiply through their known characters.
These rules determine every product by bilinearity.

For unequal codegrees there are at most two central decompositions
and one attachment for each. For equal codegrees there is one
unordered central decomposition and at most two attachments.
Thus there are at most two products in either case. The argument
works after complexification, and the imposed rational, integral,
positive and duality conditions can only remove candidate products.
It does not assume that either rational degree-two algebra is split.
\end{proof}

\begin{corollary}[Recovering the nonlinear full field]
\label{paper:source:two-matrix-center-field}
Use the notation of the preceding proof. If $f_1\ne f_2$, the
number $t_1t_2/(tg)$ must be a rational square. If $f_1=f_2=f$, the
nonlinear center is split over $\mathbb Q$ precisely when $fg$ is
a rational square; otherwise the two nonlinear characters form a
quadratic orbit with full field
\[
                              \mathbb Q(\sqrt{fg}).
\]
Thus the field and its real or imaginary nature are recoverable
from the linear table, duality and commutator tensor.
\end{corollary}

\begin{proof}
Distinct rational codegrees cannot be exchanged by Galois symmetry,
so their primitive idempotents are rational. This forces $u$ in
\eqref{paper:binary-central-idempotents} to be rational. For equal
codegrees the formula is
$e_{1,2}=e/2\pm v/\sqrt{fg}$. Their coefficient fields are either
rational or the displayed quadratic field. The identity
$\chi_i(X)=f\beta(e_i,X)$ identifies the same field with the full
character field. Rescaling the nonzero rational vector $v$ changes
$fg$ by a rational square and does not affect the conclusion.
\end{proof}

\begin{corollary}[An independent field check from the rational Lie centroid]
\label{paper:source:two-matrix-centroid-check}
Let $L_{\mathbb Q}=[R_{\mathbb Q},R_{\mathbb Q}]$. Its adjoint
commutant
\[
 \mathcal Z_L=\{T\in\operatorname{End}_{\mathbb Q}(L_{\mathbb Q}):
                  T\operatorname{ad}(x)=\operatorname{ad}(x)T
                  \text{ for every }x\in L_{\mathbb Q}\}
\]
is a two-dimensional commutative algebra canonically isomorphic to
the nonlinear associative center. It is therefore computable from
$c$ and duality alone. For equal nonlinear codegrees, its split or
quadratic field class must agree with the class determined by $fg$
in Corollary~\ref{paper:source:two-matrix-center-field}.
\end{corollary}

\begin{proof}
Multiplication by a nonlinear central element acts on $L_{\mathbb Q}$
and commutes with all adjoint maps. This action is faithful: a scalar
on either matrix component that annihilates its traceless part is
zero. After extension to $\mathbb C$, the two simple adjoint modules
are nonisomorphic modules for $\mathfrak{sl}_2\oplus\mathfrak{sl}_2$,
and their commutant consists exactly of the two independent scalar
actions. Thus both rational algebras have dimension two, so the
faithful map is an isomorphism. Their field classes must coincide.
\end{proof}

\begin{corollary}[A finite integral replacement for the two-matrix search]
\label{paper:source:finite-commutator-search}
For fixed canonical codegrees $f_1,f_2$, all two-matrix fusion rings
with a fixed aligned linear table and duality occur among the finite
integer alternating tensors satisfying
\[
                   \sum_{X<Y<Z}c_{XYZ}^2=2(f_1+f_2).
\]
Every such tensor must have the kernel, spectrum and duality properties
above; its induced bracket must satisfy Jacobi and have commutator
Lie algebra $\mathfrak{sl}_2\oplus\mathfrak{sl}_2$ after complexification.
For each retained tensor there are at most two candidate products.
In particular at most $2(f_1+f_2)$ unordered coefficients can be
nonzero, and each has absolute value at most
$\lfloor\sqrt{2(f_1+f_2)}\rfloor$.
\end{corollary}

\begin{proof}
The norm identity is
Theorem~\ref{paper:source:commutator-three-form-audit}. It bounds every integer
coordinate of a tensor in a finite-dimensional space. Every actual
ring has the necessary identities and Lie structure, and the preceding
theorem reconstructs its multiplication among at most two possibilities.
The support and coefficient bounds follow immediately from the sum
of integer squares.
\end{proof}

\begin{remark}[Sharpness for multiplication tables on a fixed labeled basis]
\label{paper:two-matrix-sharpness}
Let
\[
 G=\langle z,x,y\mid z^4=x^2=y^2=1,\ z\text{ central},\ yx=z^2xy\rangle
\]
be the Pauli group of order $16$. On its group basis, let $T$ exchange
only $z$ and $z^3$, and fix every other group element. The original
product and the transported product
\[
                         a\circ b=T\bigl(T(a)T(b)\bigr)
\]
are distinct: $z\circ x=z^3x$, whereas the original product is $zx$.
They have the same basis duality, full linear character table, and
commutator tensor. Indeed $T$ commutes with inversion, every linear
character takes the same value on $z$ and $z^3$, and $T$ fixes every
commutator. To see the last assertion, central labels commute and
the commutator of two noncentral labels is a difference of noncentral
labels whenever it is nonzero. The coefficient trace and the known
form $\beta$ are also preserved by $T$.

The group has eight linear characters and two degree-two characters,
so both products have the required two-matrix shape. They are
based-isomorphic through $T$, but are different tables on the fixed
labeled basis. Thus the upper bound two cannot in general be replaced
by one in Theorem~\ref{paper:source:two-matrix-reconstruction}; no claim of two
nonisomorphic fusion rings is made. The exact audit independently
checks all products, associativity, duality, linear characters and
commutator coefficients for both tables.
\end{remark}

\subsection{Symplectic duality}
Here basis duality is extended linearly to the complex algebra; it is
distinguished from the conjugate-linear positive involution.  On an
individually preserved degree-two component, symplectic linear duality
is $A\mapsto\operatorname{Tr}(A)I-A$.

\begin{theorem}[Individual symplectic codegree divisibility]
\label{paper:source:symplectic-individual-codegree-divisibility}
Suppose all nonlinear complex simple components of a fusion ring have
degree two, and basis duality preserves each of them individually with
symplectic restriction. Then $f_E/4$ is an algebraic integer for every
nonlinear component $E$. In particular, every rational such codegree
is an integer divisible by four. The number of components is arbitrary.
\end{theorem}

\begin{proof}
Use the integral alternating tensor $c_{XYZ}=\tau([X,Y]Z)$ and Gram
matrix of Theorem~\ref{paper:source:commutator-gram-spectrum}. Replacing any one
argument by its dual negates $c$: in each nonlinear component the
replacement is the negative of the original matrix plus a scalar,
and scalar matrices do not contribute. Linear components contribute
nothing. Hence $c$ vanishes when an argument is selfdual, or when two
arguments belong to the same source-duality pair.

Choose one representative from each nonselfdual pair and put
$d_{ijk}=c_{X_iX_jX_k}$. Let
\[
 H_{ij}=\sum_{u<v}d_{iuv}d_{juv}.
\]
This is an integral positive semidefinite matrix. In the ordered basis
of selfdual labels followed by representatives and their duals,
\[
 G=0\oplus4\begin{pmatrix}H&-H\\-H&H\end{pmatrix}.
\]
The factor four counts the choices of representatives in the two
contracted arguments. Thus its nonzero eigenvalues are eight times
those of $H$. By the cited Gram theorem, those same eigenvalues are
$2f_E$, each with multiplicity three. The eigenvalues $f_E/4$ of an
integral matrix are algebraic integers. Rational algebraic integers
are integers, proving the last statement.
\end{proof}

\begin{proposition}[The reduced tensor for one matrix component]
\label{paper:source:reduced-symplectic-tensor}
Suppose $R_{\mathbb C}\simeq\mathbb C^\ell\oplus M_2(\mathbb C)$ and
the matrix component is symplectic. Its character $s_X$ and formal
codegree $f$ are rational integers, $4\mid f$, and every $s_X$ is even.
The tensor $d$ above is a decomposable integral three-form satisfying
\[
 \sum_{i<j<k}d_{ijk}^2=f/4=:k,\qquad
 g_i:=\sum_{j<l}d_{ijl}^2,\qquad G_{X_iX_i}=4g_i.
\]
If $b_0$ is the number of nonzero slices, then $b_0\le k+2$.
Writing $h_X=s_X/2$, one also has
\[
 \sum_Xh_X^2=f/2,\qquad h_{\mathbf1}=1,\qquad
 \sum_Xd_Xh_X=0.
\]
\end{proposition}

\begin{proof}
The unique nonlinear character is Galois fixed, so its values and
codegree are rational algebraic integers. Its determinant $n_X$ is
also an integer, by the degree-two characteristic polynomial. On this
component $E(X^*)=s_XI-E(X)$, whence
\[
 G_{XX}=2\Tr(E(X)E(X)^*)-s_X^2=4n_X-s_X^2.
\]
The preceding block formula makes $G_{XX}$ divisible by four, and
therefore $s_X$ is even. The commutator form on the three-dimensional
traceless component is a volume form; its pullback to the source
basis, and then to the representatives, is decomposable. The full
three-form has squared norm $2f$ by its Gram spectrum, and each
nonzero reduced coordinate occurs eight times in that norm. This
gives the norm $f/4$.

The active coordinate columns represent a loopless rank-three
matroid. Fixing one basis, each other column permits a basis exchange
that includes that column and no other outside column. Hence there
are at least $1+(b_0-3)=b_0-2$ nonzero maximal minors. Each nonzero
integer coordinate contributes at least one to the squared norm.
Finally character orthogonality gives
$\sum_Xs_X^2=2f$ and $\sum_Xd_Xs_X=0$.
\end{proof}

\begin{corollary}[The first two symplectic codegrees]
\label{paper:source:symplectic-four-invertible}
In the preceding one-matrix setting, $f=4$ forces a nonunit invertible
source label. It is therefore impossible under the standing source
hypotheses. For $f=8$, precisely four source-duality pairs are active
in the commutator tensor, with reduced slice norms $(2,2,1,1)$.
\end{corollary}

\begin{proof}
For $f=4$, the nonunit contribution to $\sum h_X^2$ is one. There is
exactly one such label with value $\pm1$, necessarily selfdual; Perron
orthogonality forces its dimension to be one. For $f=8$, the reduced
norm is two. There are exactly two nonzero coordinates, each $\pm1$.
Decomposability forces their triples to share two indices, yielding
the asserted four slices. One nonzero coordinate of squared norm two
is impossible over the integers.
\end{proof}

\begin{corollary}[The trace support at codegree eight]
\label{paper:source:symplectic-eight-trace-support}
At $f=8$, the nonzero nonunit $h$-values occur either at three selfdual
labels, or at one nonselfdual pair and one selfdual label, and all
these values have absolute value one. In the latter case, if the pair
has dimension $p$ and the selfdual label has dimension $e$, then
$e=2p\pm1$. The pair with nonzero trace may be a scalar pair; it need
not belong to the four active commutator pairs.
\end{corollary}

\begin{proof}
The remaining squared trace budget is $f/2-1=3$. Duals have equal
values. The two listed supports exhaust its integral partitions, and
$1+2\varepsilon p+\eta e=0$ proves the dimension assertion.
\end{proof}

The individual-codegree theorem requires symplectic duality on every
nonlinear component and individual preservation of those components.
An orthogonal component or a pair exchanged by duality does not satisfy
the argument-by-argument sign identity.  The finite classification of decomposable reduced tensors on at most
seven active pairs is recorded in
Theorem~\ref{paper:symplectic-rank16-profiles} below.

\begin{proposition}[A limited parity extension]
\label{paper:symplectic-codegree-parity}
Under the individual-preservation hypotheses of
Theorem~\ref{paper:source:symplectic-individual-codegree-divisibility}, suppose
there are at most three nonlinear components and all their character
values are rational. Then each of those character values is even.
\end{proposition}

\begin{proof}
Each determinant is an integer. The diagonal Gram formula gives
$G_{XX}=\sum_E(4n_E(X)-s_E(X)^2)$. The block formula makes this
divisible by four. Among at most three squares, the sum can be zero
modulo four only when every square is zero modulo four. This argument
does not assert the same pointwise conclusion for four or more
components, although the individual codegree theorem remains valid.
\end{proof}

\subsection{Complete small symplectic slice catalogues}
\label{paper:sec:symplectic-profiles}

\begin{theorem}[Exact profiles on at most seven active pairs]
\label{paper:symplectic-rank16-profiles}
For decomposable integral three-forms of norm $k$ supported on at most
seven slots, the numbers of possible sorted positive slice-norm
profiles are as follows:
\[
\begin{array}{c|rrrrrrrrrrrrrrr}
k&1&2&3&4&5&6&7&8&9&10&11&12&13&14&15\\\hline
\#&1&1&2&4&3&4&4&7&7&8&6&14&10&9&11.
\end{array}
\]
For $k=18,25,27$ the respective numbers are $17,18,29$.
Zero slices may be added to reach the prescribed number of source
pairs. These catalogues apply to a one-matrix symplectic source of
rank sixteen, where there are at most seven nonselfdual pairs.
\end{theorem}

\begin{proof}
Here is a finite complete construction. The compact companion includes
\path{symplectic_profiles.py} and its independent exact replay
\path{verify_symplectic_profiles.py}; their comparison reproduces all
profile sets below, including the nonunit charts. If a coordinate is $\pm1$, change orientation and label that
coordinate $123$. The form is the maximal-minor vector of
$[I_3\mid A]$, with integral columns. Zero columns may be omitted;
there are at most four further columns. Orient each column so its
first nonzero coordinate is positive and sort the columns, allowing
repetitions. By Cauchy--Binet,
\[
 k=\det(I_3+AA^{\mathsf T})\ge1+\sum_j\|A_j\|^2.
\]
Thus every column and every finite multiset lies in an explicit
bounded search. The producer computes the determinant and the
contractions $A_j^{\mathsf T}\operatorname{adj}(I+AA^{\mathsf T})A_j$;
the independent replay computes all maximal minors directly, including
all omitted-pivot slices. The increasing sums of squared minors give
a valid pruning bound in the latter enumeration.

If no coordinate is a unit, every nonzero coordinate has absolute
value at least two. For the displayed norms the only square
partitions are
\[
4,\quad4+4,\quad9,\quad4+4+4,\quad9+4,\quad
9+9,\quad25,\quad16+9,\quad9+4+4+4+4,\quad9+9+9.
\]
The auditor enumerates every placement and sign of these coordinates
on the $\binom73=35$ triples, fixing one largest coordinate in the
first triple. It reconstructs that affine chart and tests every
maximal-minor identity, which is equivalent to decomposability.
This provides an independent exhaustive check of the nonunit cases.

For clarity, the only five-coordinate case is
$25=9+4+4+4+4$. With pivot three, the one-exterior numerators are four
entries $\pm2$ in a $3\times4$ matrix. Every two-by-two minor must
vanish, since a nonzero such minor is $\pm4$ or $\pm8$ and is not
divisible by three. The remaining matrix has rank one, and its
support is a $1\times4$ or $2\times2$ rectangle. Of $7,920$ signed
arrays, exactly $192$ pass: $48$ of the first shape and $144$ of the
second. They add the profiles
\[
 (25,25,9,4,4,4,4),\qquad(25,17,17,8,8).
\]
Combining these cases with the unit charts and removing duplicate
profiles gives precisely the table and the three further counts.
\end{proof}

\section{Codegree-six supports and cubic character configurations}

\subsection{Linear characters of formal codegree six}
\label{paper:sec:codegree-six}

A canonical degree-nine-hundred summand corresponds to source formal
codegree six. Its multiplicity need not be one. The following results
apply to multiplicity one, equivalently to a linear character of
$K_0(\mathcal C)$; the ring itself need not be commutative unless this
is explicitly stated.

\begin{theorem}[The full fields and value supports at codegree six]
\label{paper:linear-codegree-six-fields}
Let $h$ be a linear irreducible character of $K_0(\mathcal C)$ with
formal codegree $6$. Then its full character field is
\[
                    K_h=\mathbb Q\quad\text{or}\quad
                    \mathbb Q(\sqrt{-3}).
\]
If $K_h=\mathbb Q$, its nonzero off-unit values have precisely one
of these two patterns:
\begin{enumerate}
\item five values, each equal to $1$ or $-1$;
\item two values, one equal to $2\varepsilon$ and the other to
$-\varepsilon$, where $\varepsilon\in\{1,-1\}$.
\end{enumerate}
If $K_h=\mathbb Q(\sqrt{-3})$, exactly six simple labels have nonzero
values. They are the unit, one other self-dual simple with value
$1$ or $-1$, and two distinct dual pairs. On each dual pair the
values are a conjugate pair of primitive third or sixth roots.
Every other character value is zero.
\end{theorem}

\begin{proof}
Proposition~\ref{paper:source:canonical-character-bridge} gives
$K_h\subseteq\mathbb Q(\zeta_{12})$.
The proper quadratic subfields are
\[
 \mathbb Q(i),\qquad \mathbb Q(\sqrt3),\qquad
 \mathbb Q(\sqrt{-3}).
\]
We first record the pointwise bound
\begin{equation}\label{paper:six-pointwise-orbit-energy}
       \sum_{\sigma\in\operatorname{Gal}(K_h/\mathbb Q)}
                 |\sigma(h(T))|^2<6
       \qquad(T\in\operatorname{Irr}(\mathcal C)).
\end{equation}
Indeed, the symmetrizing trace decomposition gives
\[
 1=\tau(TT^*)=\sum_E\frac{\chi_E(TT^*)}{f_E}.
\]
Every term is nonnegative, since the representations preserve $*$.
The distinct Galois conjugates of $h$ contribute one sixth of the
left side of \eqref{paper:six-pointwise-orbit-energy}, while the
separate dimension representation contributes $d_T^2/5400>0$.

If $K_h=\mathbb Q(\sqrt3)$, a nonrational integral value
$a+b\sqrt3$ has conjugate-square sum $2a^2+6b^2\ge6$, contrary
to \eqref{paper:six-pointwise-orbit-energy}. Thus this field is
impossible.
If $K_h=\mathbb Q(i)$, orthogonality of $h$ and $\overline h$
and the norm identity give
\[
 \sum_T h(T)^2=0,\qquad \sum_T|h(T)|^2=6,
 \qquad\sum_T(\operatorname{Im}h(T))^2=3.
\]
The imaginary parts are integers. Their nonzero values occur
on distinct dual pairs with opposite signs, so their squared sum
is even. This excludes $\mathbb Q(i)$.

For the quartic field use the integral basis $1,i,\omega,i\omega$,
where $\omega^3=1$ and $\omega\ne1$. For
$\alpha=a+bi+c\omega+di\omega$ one has
\[
 \operatorname{Tr}_{\mathbb Q(\zeta_{12})/\mathbb Q}
                 (\alpha\overline\alpha)
 =4\bigl(a^2-ac+c^2+b^2-bd+d^2\bigr).
\]
The pointwise bound makes every nonzero value a twelfth root
of unity. The norm identity therefore leaves exactly six labels
in the support. Write their values as $\zeta_{12}^{a_j}$.
Orthogonality among the four conjugate characters, indexed by
$1,5,7,11$ modulo twelve, gives
\[
 \sum_j\zeta_{12}^{2a_j}=
 \sum_j\zeta_{12}^{4a_j}=
 \sum_j\zeta_{12}^{6a_j}=0.
\]
Together with the conjugate equations, these say that all
nontrivial Fourier coefficients of the counts of $a_j$ modulo
six vanish. Each residue modulo six consequently occurs once.
The unique label in residue three has value $i$ or $-i$.
Duality preserves that residue, forcing this label to be
self-dual, whereas its value is nonreal. This contradiction
excludes the quartic field.

For the rational field, integrality and
$\sum_{T\ne\mathbf1}|h(T)|^2=5$ give the two stated square
patterns. In the two-value pattern their signs are opposite:
otherwise $\sum_Td_Th(T)=0$ is impossible, since every nonunit
source degree is at least three.

Finally suppose $K_h=\mathbb Q(\sqrt{-3})$, and write
$h(T)=(a_T+b_T\sqrt{-3})/2$, with $a_T,b_T$ integers of the
same parity. Orthogonality with $\overline h$ gives
$\sum_T b_T^2=4$. Nonzero $b_T$ occur in opposite dual pairs,
so there are exactly two such pairs and all their $b_T$ equal
$1$ or $-1$. By \eqref{paper:six-pointwise-orbit-energy},
$a_T=1$ or $-1$ at these four labels. These are precisely
primitive third or sixth roots, each of absolute value one.
The remaining squared norm is two. One real value is the unit
value one; there is exactly one other real value, equal to
$1$ or $-1$. Its label is necessarily self-dual.
\end{proof}

\begin{proposition}[A universal adjoint constituent in the rational two-value case]
\label{paper:rational-six-two-values}
Suppose the second rational pattern in
Theorem~\ref{paper:linear-codegree-six-fields} occurs. Write
\[
 h(X)=2\varepsilon,\qquad h(Y)=-\varepsilon,
 \qquad d_X=u,\quad d_Y=v.
\]
Then $X,Y$ are self-dual and
\[
 v=2u+\varepsilon,\qquad
 \begin{cases}
 u\ge6,&\varepsilon=1,\\
 u\ge7,&\varepsilon=-1.
 \end{cases}
\]
In particular $v\ge13$. The simple $Y$ occurs in $TT^*$ for
every nonunit simple $T$ of $\mathcal C$. Consequently
\begin{equation}\label{paper:rational-six-type-bound}
 n_3=0,\qquad
 d_T\ge\lceil\sqrt{v+1}\rceil\quad(T\ne\mathbf1),\qquad
 \operatorname{rank}(\mathcal C)\le325.
\end{equation}
A degree-specific necessary bound is
\[
 \operatorname{rank}(\mathcal C)
 \le3+\left\lfloor
 \frac{5399-u^2-v^2}{\lceil\sqrt{v+1}\rceil^2}
 \right\rfloor.
\]

More explicitly, for $a=N_{XX}^X$ there is a positive
Grothendieck-ring element $Q$ supported off $\mathbf1,X,Y$
such that
\begin{align*}
 X^2&=\mathbf1+aX+(2a-3\varepsilon)Y+Q,\\
 XY&=(2a-3\varepsilon)X+(4a-4\varepsilon)Y+2Q,\\
 Y^2&=\mathbf1+(4a-4\varepsilon)X
                       +(8a-8\varepsilon)Y+4Q.
\end{align*}
Here $a\ge2$ for $\varepsilon=1$ and $a\ge0$ for
$\varepsilon=-1$. The element $Q$ is nonzero.
\end{proposition}

\begin{proof}
Duality preserves the real value of $h$, and both nonzero
magnitudes occur on unique off-unit labels, so $X,Y$ are
self-dual. Orthogonality to the dimension representation gives
$v=2u+\varepsilon$.
Let $\mathbf v=(h(T))_T$. The central primitive idempotent is
\[
 e_h=\frac{\mathbf1+2\varepsilon X-\varepsilon Y}{6}.
\]
Since $h$ is linear, left multiplication by $e_h$ is the rank-one
orthogonal projector $\mathbf v\mathbf v^{\mathsf t}/6$.
This assertion does not require the fusion ring to be
commutative. Consequently
\begin{equation}\label{paper:rational-six-matrix-identity}
 N_Y=2N_X+\varepsilon I-\varepsilon
                           \mathbf v\mathbf v^{\mathsf t}.
\end{equation}
For $T$ outside $\mathbf1,X,Y$, the diagonal entry
$N_{YT}^{T}=2N_{XT}^{T}+\varepsilon$ is an odd nonnegative
integer, hence positive.

Put $a=N_{XX}^X$, $b=N_{XX}^Y$.
The equality $h(X^2)=4$ gives $b=2a-3\varepsilon$.
Frobenius reciprocity and $h(XY)=-2$ give
$N_{YY}^X=4a-4\varepsilon$, and $h(Y^2)=1$ gives
$N_{YY}^Y=8a-8\varepsilon$.
Equation~\eqref{paper:rational-six-matrix-identity} gives the
remaining coefficients in the three displayed products.
Nonnegativity gives the stated bounds on $a$, makes $b>0$,
and makes $N_{YY}^Y\ge8$. Thus $Y$ occurs in every
nonunit adjoint, including those of $X$ and $Y$ themselves.
If $Q=0$, the three labels $\mathbf1,X,Y$ form a fusion
subcategory, contrary to simplicity and the stated rank
bound. Hence $Q\ne0$.

For $\varepsilon=1$, dimension in $X^2$ first gives
$u^2\ge1+2u+(2u+1)=4u+2$, so $u\ge5$.
For $\varepsilon=-1$ it gives
$u^2\ge1+3(2u-1)=6u-2$, so $u\ge6$.
At $u=6$ in the latter case, $a=0$ would leave $\dim Q=2$,
which is impossible; $a\ge1$ already exceeds $u^2$.
Thus $u\ge7$ in that case.
At this stage $v\ge11$. Universal occurrence in adjoints
gives $d_T^2\ge1+v$ for every nonunit $T$, excluding all
source dimensions three. In the positive case $u=5$ would
force $a=2$ and $\dim Q=3$, now impossible. Thus $u\ge6$
there, and $v\ge13$ in both cases.
The degree-specific rank bound follows by reserving the
squares $u^2,v^2$ and using the minimum degree for every
other nonunit.

For the uniform bound, if $\varepsilon=1$ and $u\ge7$,
the degree-specific bound is at most $323$; if
$\varepsilon=-1$ and $u\ge8$, it is at most $322$.
At $(u,v)=(6,13)$ it is at most $327$.
Ranks $327$ and $326$ would require, respectively, excess
$10$ and $26$ above the sum obtained by giving every other
nonunit degree four. Each positive excess from such an
object is $9,20,33,\ldots$, and neither total is possible.
At $(u,v)=(7,13)$ the preliminary bound is $326$; this
rank would similarly require excess $13$, also impossible.
This proves the uniform bound $325$ without asserting
existence at that rank.
\end{proof}

\begin{theorem}[Disjoint dual-pair supports at codegree six]
\label{paper:codegree-six-disjoint-supports}
Let $q_6$ be the number of nonreal conjugate pairs of linear
characters of formal codegree six. Their two-pair supports
from Theorem~\ref{paper:linear-codegree-six-fields} are pairwise
disjoint. Thus $\mathcal C$ has at least $2q_6$ distinct
non-self-dual simple pairs.

If $K_0(\mathcal C)$ is commutative, its character table has
at least $2q_6$ nonreal conjugate pairs in total. In particular,
every such collection of $q_6$ codegree-six pairs requires at
least $q_6$ further nonreal pairs of other formal codegrees.
The canonical dimension budget gives $q_6\le2$.

If $q_6=2$, the two characters' unique off-unit real labels
are the same self-dual simple $R$, with opposite values.
After naming the character with value $+1$ at $R$ first,
the dimensions of their respective two source dual pairs
satisfy
\[
 d_R+1=|d_T\pm d_U|,\qquad
 d_R-1=|d_V\pm d_W|.
\]
The four pairs represented by $T,U,V,W$ are distinct.
\end{theorem}

\begin{proof}
Choose one character from each nonreal conjugate pair and
one representative of each non-self-dual simple pair.
The scaled imaginary-part vectors
\[
       \left(\frac2{\sqrt3}\operatorname{Im}h(T)\right)_T
\]
have entries $0,1,-1$, exactly two nonzero entries each,
and are pairwise orthogonal. Hence two supports either
are disjoint or coincide.

Suppose two supports coincide, represented by $T,U$.
Orient the simple pairs and conjugate either character as
needed to write
\[
\begin{array}{c|cc}
 &T&U\\\hline
 h_1&(a_1+\sqrt{-3})/2&(b_1+\sqrt{-3})/2\\
 h_2&(a_2+\sqrt{-3})/2&(b_2-\sqrt{-3})/2,
\end{array}
\qquad a_i,b_i\in\{1,-1\}.
\]
Let $R_i$ be the unique off-unit real label for $h_i$,
with value $\varepsilon_i\in\{1,-1\}$.
Orthogonality of real parts gives
\begin{equation}\label{paper:six-overlap-real-orthogonality}
 1+\delta_{R_1,R_2}\varepsilon_1\varepsilon_2
             +\frac{a_1a_2+b_1b_2}{2}=0.
\end{equation}
If $R_1=R_2$, this forces $\varepsilon_2=-\varepsilon_1$
and exactly one of $a_2=a_1$, $b_2=b_1$.
Adding the two orthogonality equations against the dimension
representation then gives $2\pm2d_T=0$ or $2\pm2d_U=0$.
This would make a nonunit simple invertible, impossible here.
Consequently $R_1\ne R_2$, and
\eqref{paper:six-overlap-real-orthogonality} gives
$a_2=-a_1$, $b_2=-b_1$.

Let $e$ be the sum of the four primitive idempotents belonging
to $h_1,\overline h_1,h_2,\overline h_2$.
Their coefficient formula gives
\[
              3e=2\mathbf1+\varepsilon_1R_1
                                +\varepsilon_2R_2
                    \in K_0(\mathcal C).
\]
Thus $3L_e$ has integral entries. On the other hand, the
rank-one character projectors give
\[
 (L_e)_{T,U}
 =\frac16\sum_{h\in\{h_1,\overline h_1,h_2,\overline h_2\}}
                       h(U)\overline{h(T)}
 =\frac{a_1b_1}{6}.
\]
Indeed the two conjugate-pair contributions before division
by six are $(a_1b_1+3)/2$ and $(a_1b_1-3)/2$.
The resulting entry of $3L_e$ is $1/2$ or $-1/2$, a
contradiction. Coincident supports are therefore impossible.

When the ring is commutative, the linear duality operator
has the same trace in the simple basis and in the primitive
idempotent basis. The number of non-self-dual simple pairs
therefore equals the total number of nonreal character pairs.
This proves the counting assertion. Each codegree-six pair
contributes $2/6=1/3$ to $\sum_E m_E/f_E=1$, while the
dimension representation contributes positively; thus $q_6\le2$.

Finally, for two disjoint supports the real-part inner
product is simply
$1+\delta_{R_1,R_2}\varepsilon_1\varepsilon_2$.
Its vanishing forces the same real label with opposite signs.
Orthogonality of each character to the dimension
representation gives the two displayed signed-sum identities.
\end{proof}

\begin{proposition}[Five unit-sized values with one negative sign]
\label{paper:rational-six-one-negative}
Suppose the first rational pattern in
Theorem~\ref{paper:linear-codegree-six-fields} consists of four
positive values on $P_1,\ldots,P_4$ and one negative value
on $Y$. Then
\[
 d_Y=1+\sum_{i=1}^4d_{P_i}\ge13.
\]
For every nonunit simple $T$ with $h(T)=0$, the adjoint
$TT^*$ contains $Y$. Hence $n_3\le4$ and
\[
                    \operatorname{rank}(\mathcal C)\le328.
\]
\end{proposition}

\begin{proof}
The dimension identity gives the first assertion. Put
$P=\sum_iP_i$ and $\mathbf v=(h(T))_T$.
The same rank-one idempotent calculation gives
\[
                  N_Y=I+N_P-\mathbf v\mathbf v^{\mathsf t}.
\]
At a zero-character label its diagonal is at least one.
Frobenius reciprocity therefore puts $Y$ in $TT^*$ and
makes such a nonunit $T$ have degree at least four.
Only the four positive labels can have degree three.

Write $q$ for the number of zero-character simple labels,
so the rank is $6+q$. The dimension equation gives
\[
 5400\ge1+13^2+4\cdot3^2+16q=206+16q,
\]
and hence $q\le324$. If $q=324$ or $323$, the excess above
this lower bound is $10$ or $26$. Raising even one positive
supported degree from three to four raises the supported
square sum by at least $7+(14^2-13^2)=34$.
Thus all four positive degrees would have to equal three
and $d_Y=13$. The required excess would then be a sum of
$9,20,33,\ldots$, again impossible. Hence $q\le322$ and
the rank is at most $328$.
\end{proof}

\begin{corollary}[Two nonreal codegree-six pairs exclude ranks ten and eleven]
\label{paper:two-six-pairs-minimum-rank}
If $q_6=2$, then $\operatorname{rank}(\mathcal C)\ge12$.
In particular, a commutative rank-eleven field profile has
$q_6\le1$. If $q_6=1$ there, its total number $q$ of
nonreal character pairs satisfies $2\le q\le4$.
\end{corollary}

\begin{proof}
The two characters provide four distinct simple dual pairs,
the common self-dual extra label $R$, and the unit. Thus
there are at least ten labels. Write the four pair dimensions
as $a,b,c,d$ and put $r=d_R$.
The signed-sum identities in
Theorem~\ref{paper:codegree-six-disjoint-supports} imply
$a+b+c+d\equiv0\pmod2$, hence
\[
                   2(a^2+b^2+c^2+d^2)\equiv0\pmod4.
\]
At rank ten the dimension equation would give
$5400\equiv1+r^2\pmod4$, impossible. At rank eleven the
one remaining simple is self-dual; writing its dimension
as $s$ gives $5400\equiv1+r^2+s^2\pmod4$, also impossible.
For the last assertion, the disjoint-support theorem gives $q\ge2$.
The unit and the real extra label are two distinct self-dual
simples, so rank eleven permits at most four dual pairs.
\end{proof}

\subsection{An integral parity identity and the orthogonal codegree-six exclusion}
\label{paper:sec:orthogonal-six-exclusion}

\begin{lemma}[A quadratic difference of two simple objects]
\label{paper:quadratic-simple-difference}
Let $A,B$ be distinct self-dual simple labels in an integral
fusion ring with positive dimension function, and suppose
that $D=B-A$ is central and satisfies
\[
                        D^2=D+2\mathbf1,
 \qquad                 d_B-d_A=2.
\]
Then there is a third nonunit self-dual simple $Q$,
distinct from $A,B$, with $d_Q=d_A+1$ and
\begin{equation}\label{paper:quadratic-difference-four-label-block}
 DA=-\mathbf1+A+Q,\qquad
 DB=\mathbf1+B+Q,\qquad DQ=A+B.
\end{equation}
\end{lemma}

\begin{proof}
The integral matrix $N_D$ is symmetric and has spectrum
in $\{-1,2\}$. Write $x=(N_D)_{A,A}$.
The unit column has entries $-1$ at $A$, $1$ at $B$ and
zero elsewhere. From $N_D^2=N_D+2I$ in positions
$(\mathbf1,A)$ and $(\mathbf1,B)$ one obtains
\[
 (N_D)_{A,B}=x-1,\qquad (N_D)_{B,B}=x.
\]
A diagonal entry lies between $-1$ and $2$.
The row norm identity
$\sum_Z(N_D)_{A,Z}^2=x+2$, together with the nonzero
unit entry, excludes $x=-1,2$.
If $x=0$, the unit and $B$ entries already exhaust this
row's squared norm two. Its dimension identity would be
$-1-d_B=2d_A$, impossible. Hence $x=1$.

The unit and diagonal entries now account for two of
this row's squared norm three. There is exactly one
further nonzero entry, equal to $1$ or $-1$, at a label
$Q\notin\{\mathbf1,A,B\}$. The dimension identity
$-1+d_A\pm d_Q=2d_A$ forces the positive sign and
$d_Q=d_A+1$. The same argument at $B$ supplies a label
$Q'$ with positive coefficient and dimension $d_B-1$.
The $(A,B)$ entry of $N_D^2$ must be zero. The path
through the unit contributes $-1$, so the two extra
labels must coincide: $Q'=Q$.
The $(A,Q)$ equation gives $(N_D)_{Q,Q}=0$.
Its row norm is therefore two and is exhausted by its
entries one at $A,B$. This proves the three products.
Finally $DA$ is self-dual because $D$ is central and
$A,D$ are self-dual. Its expansion makes $Q$ self-dual.
\end{proof}

\begin{lemma}[A parity identity from two supported dual pairs]
\label{paper:orthogonal-six-pair-parity}
Retain $R_{\mathbb C}\simeq\mathbb C^9\oplus M_2(\mathbb C)$,
and assume orthogonal matrix-block duality, now with arbitrary
formal codegree $f=f_E$. Suppose that the nonreal linear
characters have one of the following two configurations:
\begin{enumerate}
\item exactly one conjugate pair, of formal codegree six;
\item exactly two conjugate pairs, one of formal codegree six
and the other of formal codegree eight with full field
$\mathbb Q(\sqrt{-2})$.
\end{enumerate}
Let $h,\overline h$ be the codegree-six pair and put
$e_F=e_h+e_{\overline h}$. Its two supported source dual
pairs can be represented by $T,U$ so that
\begin{equation}\label{paper:two-skew-product-idempotent}
 (T-T^*)(U-U^*)=\frac f2 e_E-3e_F.
\end{equation}
Consequently, for every source simple $X$,
\begin{equation}\label{paper:matrix-linear-six-parity}
       \chi_E(X)\equiv h(X)+\overline h(X)\pmod2.
\end{equation}
The matrix skew space has norm squareclass $f/2$.
\end{lemma}

\begin{proof}
Orient both supported pairs so that
$h(T-T^*)=h(U-U^*)=\sqrt{-3}$.
In the first configuration there are $b=q+1=2$ source
pairs. The codegree-six field contributes the projection
$\frac12\left(\begin{smallmatrix}1&1\\1&1\end{smallmatrix}\right)$
on their normalized differences. The orthogonal
matrix-block skew projection is its complement.
Thus the two matrix images are opposite, have squared
norm $f/2$, and their product is $(f/2)I$.
Their product has value $-3$ on each of $h,\overline h$
and vanishes on all other linear characters, proving
\eqref{paper:two-skew-product-idempotent}.

In the second configuration let $g$ be the codegree-eight
character. Its integral values have the form
$a+b\sqrt{-2}$. Orthogonality of $g,\overline g$ gives
$\sum_X(\operatorname{Im}g(X))^2=4$.
Each nonzero dual pair contributes $4b^2$ to this sum,
so exactly one source pair has nonzero imaginary part,
equal to $\pm\sqrt2$.
The imaginary-part vectors of $g$ and $h$ are orthogonal,
so this pair is disjoint from the two pairs supported by
$h$. There are exactly three source pairs in total.
The $g$-pair fills the entire skew projection on its
one supported pair, and the matrix image of that
pair's difference is zero. The preceding two-pair
argument therefore applies without change to $T,U$.

The coefficient of a simple $X$ in the right side of
\eqref{paper:two-skew-product-idempotent} is
\[
 \frac12\bigl(\chi_E(X^*)-h(X^*)-\overline h(X^*)\bigr).
\]
The product on the left belongs to the integral based
ring. The displayed coefficient is an integer,
proving \eqref{paper:matrix-linear-six-parity}.
The norm statement was established by the same projection.
\end{proof}

\begin{theorem}[The orthogonal codegree-six pair obstruction]
\label{paper:orthogonal-six-exclusion}
Neither configuration in
Lemma~\ref{paper:orthogonal-six-pair-parity} is possible
when the degree-two formal codegree is itself six.
In particular an orthogonal rank-thirteen field state
with $f_E=6$ and $q=q_6=1$ is excluded.
\end{theorem}

\begin{proof}
The parity identity makes the degree-two traces odd on
exactly the two source pairs supported by $h$, and even
on every other simple label. The character-square
budget is twelve, of which the unit uses four.
The supported pairs must each have trace $s_i=\pm1$,
using four more. The remaining budget four forces
precisely one off-unit self-dual label $R$ to have
trace $2\eta$, $\eta\in\{1,-1\}$; every other
trace is zero. An additional dual pair cannot use this
budget, since a nonzero even trace on it costs at least
eight.

Let $S$ denote the unique off-unit real label of $h$,
with value $\varepsilon\in\{1,-1\}$.
Write the real numerators at the two supported pairs
as $a_1,a_2\in\{1,-1\}$.
If $S=R$, orthogonality gives
\[
            2+2\eta\varepsilon+s_1a_1+s_2a_2=0.
\]
It forces $\varepsilon=-\eta$ and opposite signs for
$s_1a_1,s_2a_2$. The two dimension-orthogonality
identities then force one supported nonunit dimension
to belong to $\{1,-1\}$, exactly as in the first case of the proof
of Corollary~\ref{paper:symplectic-six-linear-six-threshold}.
Thus $S\ne R$.
Orthogonality now gives $a_i=-s_i$, and dimension
orthogonality gives
\begin{equation}\label{paper:orthogonal-six-difference-dimensions}
                   2+\eta d_R+\varepsilon d_S=0.
\end{equation}
The two signs are opposite and the two positive nonunit
dimensions differ by two.

Put $e=e_E+e_h+e_{\overline h}$.
The exact primitive-idempotent coefficient formulas,
including cancellation on both supported pairs, give
\[
                     3e=2\mathbf1+\eta R+\varepsilon S.
\]
Name the smaller of $R,S$ by $A$ and the larger by $B$.
Equation~\eqref{paper:orthogonal-six-difference-dimensions}
then makes $3e=2\mathbf1+A-B$.
Consequently $D=B-A=2\mathbf1-3e$ is central,
$D^2=D+2\mathbf1$, and $d_B-d_A=2$.
Lemma~\ref{paper:quadratic-simple-difference} provides a
distinct self-dual $Q$ with $DQ=A+B$.
On the degree-two representation $D$ acts by $-I$.
Taking its character in the last identity gives
\[
             \chi_E(Q)=-\chi_E(A+B)=-2\eta.
\]
This contributes another four to a character-square
budget already exhausted by the unit, $R$ and the
two supported pairs. The contradiction proves the theorem.
\end{proof}

\begin{proposition}[The residual codegree-four-hundred-fifty block]
\label{paper:codegree450-block-forms}
In either configuration of
Lemma~\ref{paper:orthogonal-six-pair-parity}, if $f_E=450$
then the rational degree-two algebra is either $M_2(\mathbb Q)$
or $B_{2,3}$. Its positive trace form is respectively
\[
      \frac1{225}\langle1,1,1,1\rangle
      \quad\text{or}\quad
      \frac1{225}\langle1,1,3,3\rangle.
\]
\end{proposition}

\begin{proof}
The skew norm is $f_E/2=225$, so the degree-two algebra
contains $\mathbb Q(i)$.
The same integral-idempotent and order-discriminant
argument of
Theorem~\ref{paper:rational-matrix-ramification} makes
it unramified outside $2,3,5$. It splits at infinity
because its involution is positive and orthogonal.
The prime five splits in $\mathbb Q(i)$, so the algebra
cannot ramify there: a division algebra cannot contain
$\mathbb Q_5\times\mathbb Q_5$.
Even cardinality of the ramification set therefore
leaves the split algebra and $B_{2,3}$.
With skew squareclass one these are represented by
$(-1,1)$ and $(-1,3)$, giving the stated forms after
multiplication by $2/f_E=1/225$.
\end{proof}

\subsection{Character-coordinate roots and a cubic source-type restriction}
\label{paper:sec:character-type-roots}

The integral coordinate covariance gives information about individual
source labels, beyond its short-vector exclusions.  The following
statements concern fusion rings and their multiplication tables.  A
based-ring automorphism obtained below is not asserted to lift to a
tensor autoequivalence.

\begin{theorem}[The codegree-nine real cubic support]
\label{paper:cubic-nine-type-support}
Suppose a fusion ring has a full non-Perron orbit of linear characters
with field $K=\mathbb Q(\zeta_9+\zeta_9^{-1})$ and codegree $9$.
Then there are five distinct selfdual nonunit source labels
$X_1,X_2,X_3,Y,Z$ such that
\[
 d_{X_1}=d_{X_2}=d_{X_3}=d,\qquad d_Y=e,\qquad d_Z=e+1.
\]
For a chosen character in this orbit, the values on $X_1,X_2,X_3$
are one common sign of the three conjugates of
$t=\zeta_9+\zeta_9^{-1}$, while
\[
             \chi(Y)=1,\qquad \chi(Z)=-1.
\]
Every other nonunit source label has value zero.  The dimension $d$
may equal $e$ or $e+1$; the five \emph{labels}, rather than their
dimension values, are distinct.

If the source rank is odd, there are at least six selfdual nonunit
labels.  If in addition the source is integral of even total dimension,
the sum of squared dimensions of the remaining selfdual nonunit labels
has parity $d$.  In particular, odd $d$ forces a further selfdual label
of odd dimension.
\end{theorem}

\begin{proof}
In the ambient integral power basis of $\mathbb Q(\zeta_9)$, the real
fixed lattice has integral basis
\[
 \beta_0=1,\quad
 \beta_1=\zeta_9-\zeta_9^2+\zeta_9^4=2-t^2,\quad
 \beta_2=-\zeta_9+\zeta_9^2+\zeta_9^5=-t.
\]
Indeed its free ambient coordinates are those of $1,\zeta_9^4,
\zeta_9^5$; conjugation fixes a vector exactly when its coordinates
are $(a,b-c,-b+c,0,b,c)$.  In this basis conjugation is the identity,
the unit coordinate is $(1,0,0)$, and
\[
 G=\begin{pmatrix}3&0&0\\0&6&-3\\0&-3&6\end{pmatrix},
 \qquad
 H=9G^{-1}=\begin{pmatrix}3&0&0\\0&2&1\\0&1&2\end{pmatrix}.
\]
These formulas also follow directly from $t^3-3t+1=0$.

Let $a,b,c$ be the three source coordinate columns.  The columns
$b,c$ have squared norm two and inner product one.  Their images
are roots sharing exactly one endpoint with the same sign.
Their union consists of three nonunit selfdual labels of equal
dimension.  Up to permutation and a common sign, the coordinate
pairs on them are
\[
                      (1,1),\quad(-1,0),\quad(0,-1).
\]
All other $b,c$ coordinates vanish.  The column $a$ has squared
norm three, value one at the unit, and is orthogonal to both roots.
It is therefore constant on their three-label support.  That
constant must be zero, since a nonzero integer value there would
already contribute three in addition to the unit's contribution.

Consequently $a$ has precisely two further nonunit entries, each
$\pm1$.  Perron orthogonality
$1+\varepsilon d_Y+\eta d_Z=0$ and $d_Y,d_Z\ge1$ force opposite
signs and consecutive dimensions.  Duality fixes these labels
individually because $a$ is duality-invariant, the unit is fixed,
and its two other nonzero entries have opposite signs.
Finally the three cubic values are
\[
 \beta_1+\beta_2=2-t-t^2,\qquad -\beta_1=t^2-2,
 \qquad -\beta_2=t,
\]
or their negatives; these are the three conjugates of $t$.

The number of nonunit selfdual labels has parity $r-1$, giving the
odd-rank conclusion.  For even integral total dimension, nonselfdual
pairs contribute even squared mass.  The contribution of the unit
and the displayed five labels has parity
$1+3d^2+e^2+(e+1)^2\equiv d\pmod2$, proving the last assertion.
\end{proof}

\begin{corollary}[A positive cubic fusion subtable]
\label{paper:cubic-nine-positive-subtable}
Choose the cyclic order so that the cubic values are
$\chi(X_i)=\varepsilon t_i$, where $\varepsilon\in\{1,-1\}$
and $t_i^2=2+t_{i+1}$, with indices modulo three.  Put
\[
 H_i=\begin{cases}X_{i+1},&\varepsilon=1,\\
                  S-X_{i+1},&\varepsilon=-1.
       \end{cases}
\]
There is a common nonnegative integral source-basis combination $T$
which is selfdual, fixed by the cyclic based-ring automorphism,
has unit coefficient zero, and has zero cubic projection, such that
\[
 X_i^2=1+Y+T+H_i,\qquad
 X_iX_j=Z+T+H_k\quad(\{i,j,k\}=\{1,2,3\}).
\]
In particular
\[
 \FPdim(T)=
 \begin{cases}d(d-1)-\FPdim(Z),&\varepsilon=1,\\
               d(d-2)-\FPdim(Z),&\varepsilon=-1.
 \end{cases}
\]
If every nonunit source dimension is at least three and $d=3$,
then necessarily $\varepsilon=1$, $T=0$, $\FPdim(Y)=5$, and
$\FPdim(Z)=6$.  There is moreover a further selfdual simple $W$
of dimension twelve, distinct from these five labels and the unit,
such that
\[
 YX_i=X_i+W,\qquad ZX_i=S-X_i+W,\qquad
 Y^2=1+Y+Z+P,\quad \FPdim(P)=13,
\]
where $P$ is nonnegative and has unit coefficient zero.
This specifies a forced local product pattern;
it is not an exclusion or an existence assertion for $d=3$.
\end{corollary}

\begin{proof}
The triple differences are central, so the triple elements commute.
Their images outside the cubic component agree.  The cubic equations
therefore give
\[
 X_i^2-X_j^2=\varepsilon(X_{i+1}-X_{j+1}),\qquad
 X_iX_j=X_k^2-(1+Y-Z).
\]
Thus $B=X_i^2-H_i$ is independent of $i$.  Every source label is
omitted by at least one $H_i$, so the corresponding expression for
$B$ proves each of its coefficients nonnegative.  Its unit
coefficient is one, since $X_i$ is selfdual.  Positivity of the
offdiagonal product gives $[B:Y]\ge1$.  Consequently
$T=B-1-Y$ is nonnegative with unit coefficient zero and gives the
two displayed products.  Selfduality and cyclic invariance follow
from the formulas.  The cubic projection of each $H_i$ is
$t_{i+1}$, so that of $B$ is two and that of $T$ is zero.
Taking dimensions gives the formulas for $\FPdim(T)$.

When $d=3$, one has $\FPdim(Z)=\FPdim(Y)+1\ge4$.
The negative sign would give $\FPdim(T)=3-\FPdim(Z)<0$.
For the positive sign, $0\le\FPdim(T)=6-\FPdim(Z)\le2$;
the unit coefficient is zero and every other dimension is at least
three, so $T=0$.  Hence $\FPdim(Z)=6$ and $\FPdim(Y)=5$.

For the further assertion, $V=YX_i-X_i$ is independent of $i$,
since $Y$ acts by one on the cubic ideal.  Frobenius reciprocity
and the positive subtable show that $V$ is nonnegative, has no
unit or triple support, and has dimension twelve.  The equation
$Y=X_i^2-1-X_{i+1}$ shows that $Y$ commutes with the triple;
therefore $V$ is selfdual.  Frobenius reciprocity and
$Y^2-YZ=1-Z$ give $[V:Y]=[V:Z]=h$.  The dimension bound gives
$h\le1$, and $h=1$ would leave residual dimension one without
a unit.  Thus $h=0$.

Put $t=[Y^2:Y]$.  The same identities give $[Y^2:Z]=t$ and
zero triple coefficients, so
$Y^2=1+t(Y+Z)+P$ with $P\ge0$ and
$\FPdim(P)=24-11t$.  Expanding with $X_i^2=1+Y+X_{i+1}$ gives
\[
 \tau(V^2)=\tau((Y-1)^2X_i^2)=t+h=t.
\]
As $V$ is nonzero and selfdual, $t\ge1$, whereas dimensions give
$t\le2$.  The case $t=2$ would leave $\FPdim(P)=2$ without a
unit and is impossible.  Hence $t=1$, and the positive integral
vector $V$ has squared norm one.  It is one selfdual simple $W$,
necessarily of dimension twelve.  The formula for $ZX_i$ follows
from $(Y-Z)X_i=2X_i-S$, and the remaining assertions follow.
\end{proof}

These restrictions are conditional on the specified character orbit.
They give source-type and multiplication constraints; they do not
assert that the displayed types are realized and do not add a new
enumerated endpoint exclusion count.

\begin{corollary}[A specified difference of fusion matrices]
\label{paper:cubic-nine-ring-identities}
Under the preceding hypotheses let $S=X_1+X_2+X_3$ and $A=Y-Z$.
The rational central idempotent of the cubic component is
\[
                    e_K=\frac{1+Y-Z}{3},
\]
and $A$ is central.  In the Grothendieck ring one has
\begin{align*}
 A^2&=A+2,& YZ&=ZY,\\
 Y^2-YZ&=1-Z,& Z^2-YZ&=1+Y,\\
 AX_i&=2X_i-S\quad(1\le i\le3),&
 AW&=-W.
\end{align*}
The last identity is restricted to nonunit labels $W$ outside
$\{X_1,X_2,X_3,Y,Z\}$.  These are identities of virtual ring
elements.  They specify the entire matrix $N_Y-N_Z$: on the ordered
labels $(1,Y,Z)$ its block is
\[
 \begin{pmatrix}0&1&-1\\1&0&-1\\-1&-1&0\end{pmatrix};
\]
on $(X_1,X_2,X_3)$ its block is $2I_3-J_3$, and it is $-I$ on
all remaining labels.  In particular its trace is $9-r$.

There is also a based-ring automorphism of order three that cycles
$X_1,X_2,X_3$ and fixes every other source label.
\end{corollary}

\begin{proof}
The primitive central idempotents of linear characters give
$e_K=\frac19\sum_X\Tr_{K/\mathbb Q}(\chi(X^*))X$.
The support theorem reduces this to the displayed expression.
Thus $A=3e_K-1$, proving centrality and its quadratic identity.
Since $e_KY=e_K$ and $e_KZ=-e_K$, multiplication by $A$ gives
$AY=1-Z$ and $AZ=-1-Y$; centrality also gives $YZ=ZY$.

Each coordinate coefficient vector lies in the rational cubic
ideal: it is the adjoint of a coordinate functional of the
projection onto that ideal, with respect to the coefficient-trace
inner product.  Hence $X_i-X_j$ lies in the cubic ideal.  The three
elements have equal images in every other simple component, while
their cubic images sum to zero.  Therefore
$e_KX_i=X_i-S/3$.  Labels outside the displayed support have
$e_KW=0$.  These facts prove the remaining multiplication identities
and the stated matrix blocks.

Extend a generator of $\operatorname{Gal}(K/\mathbb Q)$ by the
identity on all other rational simple components.  It cycles the
three distinct cubic values and fixes every rational value $0,1,-1$.
Because the three labels agree in all other components, this algebra
automorphism cycles exactly $X_1,X_2,X_3$ and fixes the other basis
elements.  It is consequently a based-ring automorphism, preserving
unit, duality, and Perron dimensions.
\end{proof}

\begin{corollary}[Two cubic orbits force a dimension chain]
\label{paper:two-cubic-type-chain}
There are at most two such codegree-nine real cubic orbits.  If
there are two, their cubic triples have disjoint supports, and
outside those six labels there are three distinct selfdual labels
of dimensions $e,e+1,e+2$.  Thus the source contains nine distinct
selfdual nonunit labels with dimension multiset
\[
                (d,d,d,d',d',d',e,e+1,e+2).
\]
For odd source rank there are at least ten selfdual nonunit labels.
The two order-three based-ring automorphisms commute and generate
a subgroup $C_3\times C_3$.
\end{corollary}

\begin{proof}
Different rational character components are orthogonal.  Their
root planes therefore use disjoint triples.  Each norm-three unit
column is orthogonal to the other root plane, and its unit value
one forces it to vanish on that triple as well.  Write the columns
as $q_i=e_{\mathbf1}+v_i$, where
$v_i=e_{Y_i}-e_{Z_i}$ and $d_{Z_i}=d_{Y_i}+1$.
Orthogonality gives $v_1\cdot v_2=-1$.  Two such difference roots
share exactly one opposite-sign endpoint; they therefore form a
chain of three distinct labels with consecutive dimensions.

For three orbits the nonunit columns would have Gram matrix
$3I_3-J_3$, giving $v_1+v_2+v_3=0$.  But their Perron inner
products are all $-1$, a contradiction.  The same bound also follows
from the unit trace: each orbit has mass $3/9=1/3$, and the Perron
component has additional positive mass.  Finally the two cyclic
automorphisms act on disjoint triples, so commute and act faithfully
as $C_3\times C_3$.  The parity statement follows as before.
\end{proof}

\begin{corollary}[Additional selfdual support at the next endpoint]
\label{paper:cubic-three-parity-consequences}
Assume the source is integral of even total dimension, every nonunit
simple has dimension at least three, and a full non-Perron real cubic
conductor-nine linear-character orbit of codegree nine has cubic
triple dimension $d=3$.  Then there is a further selfdual simple of
odd dimension beyond the six selfdual nonunit labels of dimensions
\[
                         (3,3,3,5,6,12)
\]
forced by Corollary~\ref{paper:cubic-nine-positive-subtable}.
Consequently there are at least seven selfdual nonunit labels, and
at least eight when the source rank is odd.  At source rank sixteen
there are at most four nonselfdual source pairs.

Moreover, two distinct real cubic conductor-nine codegree-nine
linear-character orbits cannot both have cubic triple dimension three.
\end{corollary}

\begin{proof}
The unit and the six displayed labels contribute squared mass
\[
              1+3\cdot3^2+5^2+6^2+12^2=233,
\]
which is odd.  Every remaining nonselfdual pair contributes even
squared mass, since its two integral dimensions agree.  To obtain
even total dimension, the remaining selfdual labels must therefore
include an odd-dimensional label.  This proves the first lower bound.
The number of nonunit selfdual labels has parity $r-1$, proving the
odd-rank refinement.  At rank sixteen, at least eight labels in
total are selfdual, including the unit, leaving at most eight labels
in nonselfdual pairs.

If both cubic triples had dimension three, each orbit would have
its positive rational endpoint at dimension five and its negative
rational endpoint at dimension six.  By
Corollary~\ref{paper:two-cubic-type-chain}, the two endpoint pairs
must share an opposite-sign endpoint.  Such a source label would
have to have dimensions five and six simultaneously, a contradiction.
\end{proof}

\subsection{The complete orthogonal codegree-six matrix support}
\label{paper:sec:orthogonal-six-catalogue}

Suppose that the complexified fusion ring has a unique nonlinear
component, isomorphic to $M_2(\mathbb C)$, of formal codegree six,
and that its positive involution is orthogonal.  Write $E$ for the
matrix representation, $s_X=\operatorname{Tr}E(X)$ and
$t_X=\operatorname{Tr}(E(X)E(X)^*)$.  Its rational form is a
quaternion algebra $B$ split at infinity.  The pure matrix images
have an integral Gram matrix $G$ satisfying $G^2=12G$ and
$\operatorname{rank}G=3$.  For a selfdual label put $J_X=G_{XX}$;
for one representative of a nonselfdual pair put
\[
 J_X=\frac{G_{XX}+G_{X,X^*}}2,\qquad
 K_X=\frac{G_{XX}-G_{X,X^*}}2.
\]
Thus
\[
 t_X=\frac{s_X^2+J_X+K_X}{2},\qquad
 \det E(X)=\frac{s_X^2-J_X+K_X}{4},
\]
with $K_X=0$ in a selfdual row.  A label is called active here
when its pure matrix image is nonzero.

\begin{theorem}[Four orthogonal supports]
\label{paper:orthogonal-six-patterns}
Up to permuting labels, the active rows have exactly one of the
following four necessary patterns.  A pair row describes two labels.
All displayed signs belong to $\{1,-1\}$.
\[
\begin{array}{c|c|c|c}
 &\text{selfdual rows }(J,s)&\text{pair rows }(J,K,s)&
 \text{additional nonunit scalar image}\\ \hline
A&(8,0),(8,0),(8,0)&(0,3,\epsilon),(0,3,\eta)&\lambda I\\
B&(5,s_1),(5,s_2),(5,s_3),(5,s_4)&(2,6,0)&\lambda I\\
D&(5,s_1),(5,s_2),(8,0)&(3,6,\epsilon)&\lambda I\\
E&(4,2\tau),(4,0)&(4,3,-\tau),(4,3,-\tau)&\text{none}
\end{array}
\]
In pattern D, if $s_1=s_2$, their common value is $\epsilon$.
All other nonunit matrix images are zero.  In every case the
nonlinear character is nonzero on precisely six basis labels:
the unit, one nonunit selfdual label of trace $\pm2$, and four
labels of trace $\pm1$.

After pair and equal-row permutations, the signed active row
signatures number $3,5,4,2$ respectively, a total of fourteen.
This is a necessary support classification; it does not assert
realizability of any listed support.
\end{theorem}

\begin{proof}
The commutator support bound gives at most seven active labels.
The integral decomposable alternating three-form of squared norm
twelve has no coloop and hence at least four nonzero maximal
minors.  If none were a unit, its squared norm would be at least
sixteen.  Choose a unit minor.  The pure support can consequently
be represented by an integral matrix $T=[I_3\mid A]$ with at most
four extra nonzero columns.  Cauchy--Binet gives
\[
 12=\det(TT^{\mathsf T})
   =1+\sum A_{ij}^2+\text{nonnegative further squared minors},
\]
so $\sum A_{ij}^2\le11$.  Orient each extra column by requiring
its first nonzero entry to be positive, and order the columns.
There are $85$ possible oriented columns and $12\,905$ multisets
of at most four columns with this squared-norm bound, including
the empty multiset.  Exactly $129$ have squared maximal-minor
norm twelve and no coloop.

For each matrix, restore independently every extra-column sign.
Its Gram matrix is
\[
 G=T^{\mathsf T}\operatorname{adj}(TT^{\mathsf T})T,
 \qquad
 G_{ij}=\sum_{a<b}c_{iab}c_{jab},
\]
where $c_{iab}$ is the alternating maximal minor.  A basis duality
$D$ must satisfy $DGD=G$ and $\operatorname{Tr}(GD)=12$;
the latter selects a fixed pure plane and a one-dimensional skew
space.  Enumerate every such involution.  For its integral traces
enumerate $s_i\in\{-2,-1,0,1,2\}$ subject to
\[
 Ds=s,\quad Gs=0,\quad
 s_i^2\equiv G_{i,i^*}\pmod4,\quad
 (s_i^2+G_{ii})/2\in\{0,1,2,3,4,5\}.
\]
The unit uses four of the total trace-square norm twelve.  An
inactive rational scalar image has even trace, so
$8-\sum s_i^2\in\{0,4,8\}$.  These conditions give $1752$
normalized trace states.  Character integrality imposes
$G_{ij}+s_is_j\equiv0\pmod2$.  If every trace is even, the
integral central element
$\tfrac12\sum_Xs_{X^*}X=3e_E$ further imposes this congruence
modulo four.  Exactly $1344$ states remain.

Finally impose a positive, duality-invariant dimension vector
$d$ in the kernel of $G$.  The finite calculation leaves $528$
normalized states, distributed as
\[
\begin{array}{c|rrrr|r}
\text{pattern}&A&B&D&E&\text{total}\\ \hline
\text{positive states}&72&216&96&144&528.
\end{array}
\]
Before positivity, one further unsigned pattern occurs: three
pair rows $(3,2,\epsilon_1),(3,2,\epsilon_2),(6,2,0)$ and no
active selfdual rows.  Its trace relation makes the first two
fixed vectors collinear, while the third is independent, which
prevents a strictly positive dimension relation.

For reproducibility, \texttt{orthogonal\_six\_catalogue.py}
regenerates every input and stores the actual frame, involution,
trace vector and a positive rational kernel witness.  The
independent \texttt{audit\_orthogonal\_six.py} imports no producer
routine: it enumerates the bounded multisets directly, evaluates
Leibniz determinants and alternating-minor contractions, enumerates
involutions as disjoint transpositions, and uses exact separating
lines in the fixed plane to audit positivity.  It checks equality
of the complete regenerated state sets, rather than only their
counts.  The initial $23$ signed integral signatures reduce to
$16$ after positivity.  The unit-character relation eliminates
two E signatures, by Theorem~\ref{paper:orthogonal-six-perron-relations}, leaving the
fourteen asserted signatures.

The active trace-square mass is four in A, B and D, forcing one
additional nonunit scalar image $\lambda I$.  It is selfdual,
since a nonselfdual scalar pair would use eight instead.  In E
the active trace-square mass is eight, leaving no additional
nonzero scalar image.  This also proves the assertion on the
six nonzero character values.
\end{proof}

\begin{corollary}[The three rational matrix forms]
\label{paper:orthogonal-six-three-forms}
Let $\delta$ be the positive square class of the skew line.  If
$B_{2,3}$ denotes the rational quaternion algebra ramified exactly
at $2$ and $3$, then the only possibilities are
\[
 (\delta,B)=
 \begin{cases}
 (3,B_{2,3}),&A,\\
 (6,B_{2,3}),&B,D,\\
 (3,M_2(\mathbb Q)),&E.
 \end{cases}
\]
In particular skew square classes one and two are impossible,
as is the split form with skew square class six.
\end{corollary}

\begin{proof}
The skew Gram diagonal is $K_X=\delta z_X^2$, and the catalogue
gives $\delta=3$ or six as displayed.  Any nonzero fixed pure
image gives $B=(-\delta,J_X)_{\mathbb Q}$.  Taking respectively
$J_X=8,2,3,4$ in A, B, D, E gives the stated rational forms by
the elementary Hilbert-symbol calculation.  For E the second
entry is a square, so the algebra is split.  In the other three
presentations the finite negative Hilbert symbols occur exactly
at two and three.
\end{proof}

\subsection{Positive dimensions in the orthogonal six-supports}
\label{paper:sec:orthogonal-six-refinement}

\begin{lemma}[The fixed-plane frame]
\label{paper:orthogonal-six-fixed-frame}
Represent the fixed pure part of a selfdual image by $u_X$ and
that of a pair image by $v_X=(u_X+u_{X^*})/2$.  Then
\[
 \sum_{X=X^*}u_X\otimes u_X+
 2\sum_{\{X,X^*\}}v_X\otimes v_X=12I_2.
\]
The squared norms of these vectors are the entries $J_X$ in the
preceding table.
\end{lemma}
\begin{proof}
Restrict the tight-frame equation $G^2=12G$ to the fixed pure
plane.  The skew parts of a dual pair cancel, leaving two copies
of the outer product of its fixed part.
\end{proof}

\begin{lemma}[Trace and dimension relations]
\label{paper:orthogonal-six-two-relations}
In the fixed plane one has
\[
 \sum_{X=X^*}s_Xu_X+2\sum_{\{X,X^*\}}s_Xv_X=0,
 \qquad
 \sum_{X=X^*}d_Xu_X+2\sum_{\{X,X^*\}}d_Xv_X=0.
\]
Moreover, orthogonality of the nonlinear character to the Perron
character gives $\sum_Xd_Xs_X=0$, with the unit included.
\end{lemma}
\begin{proof}
These are respectively $Gs=0$, $Gd=0$ and ordinary character
orthogonality.  Pair dimensions and pair traces agree under duality.
\end{proof}

\begin{theorem}[Perron dimension relations]
\label{paper:orthogonal-six-perron-relations}
The dimensions in the four supports satisfy the following laws.
\begin{enumerate}
\item In A the three active selfdual dimensions are equal to a
common value $a$.  If the two pair dimensions are $p,q$ and the
additional scalar image is $\lambda I$ on a label of dimension
$r$, then
\[
 1+\epsilon p+\eta q+\lambda r=0,
 \qquad r=|1+\epsilon p+\eta q|.
\]
\item In B the four signed selfdual dimensions can be partitioned
into two ordered pairs such that, for the pair dimension $p$,
\[
 p=|s_1d_1-s_2d_2|=|s_3d_3-s_4d_4|.
\]
If the additional scalar label has dimension $r$ and image
$\lambda I$, then $2+2\lambda r+\sum_i s_id_i=0$.
\item In D write $p$ for the pair dimension and $q$ for the
trace-zero selfdual dimension.  Then the other selfdual dimensions
are $p+q$ and $|p-q|$, and the additional scalar label satisfies
\[
 r=2p+\epsilon,\qquad\lambda=-\epsilon.
\]
Equivalently, with the original signs,
\[
 s_1d_1+s_2d_2=2\epsilon p,\qquad
 2q=|s_1d_1-s_2d_2|.
\]
\item In E the two pair dimensions agree, say $p$.  Their trace
is $-\tau$, where the nonzero selfdual trace is $2\tau$.  The
trace-$2\tau$ and trace-zero selfdual dimensions are
\[
 d_2=2p-\tau,\qquad d_0=4p-\tau.
\]
\end{enumerate}
\end{theorem}

\begin{proof}
For A, the three norm-eight fixed vectors form a tight frame with
nonzero Gram eigenvalues twelve and twelve.  Its one-dimensional
kernel therefore has coordinates of equal absolute value.  A
positive dimension vector in that kernel makes all three
dimensions equal.  The scalar character relation gives the
remaining assertion.

For B put $h_i=s_iu_i$.  They have equal squared norm five and
sum zero.  Four equal-length vectors in a plane with sum zero
split into two antipodal pairs: indeed a chord midpoint determines
its unordered endpoints on the circle, unless it is the center,
in which case the chord is already a diameter.  Write them as
$a,-a,b,-b$.  They cannot all be collinear, since their frame
contribution in that direction would have eigenvalue twenty,
larger than twelve.  If $v$ is the pair's fixed vector, the frame
equation becomes
\[
 a\otimes a+b\otimes b+v\otimes v=6I_2,
 \qquad (\|a\|^2,\|b\|^2,\|v\|^2)=(5,5,2).
\]
The Gram matrix is $6I-zz^{\mathsf T}$ with
$(z_1^2,z_2^2,z_3^2)=(1,1,4)$.  Thus
$2v=\kappa_1a+\kappa_2b$ for signs $\kappa_i$.  Comparing the
two independent coefficients in the dimension relation gives
the two absolute differences.  Character orthogonality supplies
the scalar relation.

For D put $a=s_1u_1$, $b=s_2u_2$ and $w=\epsilon v$.
The trace relation is $a+b+2w=0$ and their squared norms are
$5,5,3$.  Hence $a\cdot b=1$; $a+b$ and $a-b$ are orthogonal
with squared norms twelve and eight.  The frame equation forces
the third selfdual vector to equal $a-b$ up to sign.  Taking
coefficients in the dimension relation proves the two equations
in the statement.  Their positive solutions have dimensions
$p+q,|p-q|,q$.  The scalar character relation is
$2+4\epsilon p+2\lambda r=0$, so positivity gives
$r=2p+\epsilon$ and $\lambda=-\epsilon$.

For E, initially allow pair traces $\epsilon,\eta$.  If $b$ is
the trace-$2\tau$ fixed vector, put
$b'=\tau b$, $w_1=\epsilon v_1$, $w_2=\eta v_2$.
The trace relation gives $b'+w_1+w_2=0$; all three squared norms
are four.  The frame equation forces the trace-zero selfdual
vector to be $b'$ up to sign.  Project the dimension relation
onto $w_1-w_2$: it gives $\epsilon p_1=\eta p_2$.  Positivity
therefore gives equal signs and $p_1=p_2=p$.  The remaining
dimension relation is $d_0=|d_2-2\epsilon\tau p|$.
Finally $1+\tau d_2+2\epsilon p=0$ excludes
$\epsilon=\tau$, since $d_2,p\ge1$.  Thus $\epsilon=\eta=-\tau$,
$d_2=2p-\tau$, and $d_0=4p-\tau$.
\end{proof}

\begin{corollary}[Finite parameters at dimension 5400 and rank sixteen]
\label{paper:orthogonal-six-fp-parameters}
Assume all nonunit dimensions are integers at least three and
the source rank is sixteen.  In D,
\[
 p,q\ge3,\quad |p-q|\ge3,\quad
 8p^2+3q^2+4\epsilon p+83\le5400,
\]
so $p\le25$ and $q\le41$.  In E,
\[
 p\ge3,\qquad24p^2-12\tau p+84\le5400;
\]
there are only twenty-five parameter pairs $(p,\tau)$:
$3\le p\le15$ for $\tau=1$ and $3\le p\le14$ for
$\tau=-1$.  In A a necessary mass bound is
\[
 3a^2+2p^2+2q^2+r^2+64\le5400.
\]
\end{corollary}
\begin{proof}
The six nonunit D labels have squared mass
$8p^2+3q^2+4\epsilon p+1$.  Including the unit and nine other
nonunit labels contributes at least another $82$.  For E the
six displayed nonunit labels have squared mass
$24p^2-12\tau p+2$, and the same remaining-label estimate applies.
The stated ranges follow by direct integer inequalities.  A has
eight displayed nonunit labels, leaving seven further nonunit
labels and the unit, of total squared mass at least $64$.
\end{proof}

\begin{lemma}[The precise opposite-image restriction]
\label{paper:orthogonal-six-no-opposite-images}
Two active selfdual labels cannot have opposite matrix images.
Also, no active label has $E(X^*)=-E(X)$.
\end{lemma}
\begin{proof}
The latter equality would require both trace zero and fixed pure
norm $J=0$, which no pair row in the table has.  For two selfdual
rows, opposite images require opposite traces and opposite pure
vectors.  In A the positive relation among the three norm-eight
vectors excludes opposite vectors.  In B opposite images would
make two signed vectors $s_iu_i$ equal.  Their total sum zero
and equal norms would then make all four collinear, contrary to
the frame bound used above.  The explicit two-plane descriptions
in D and E likewise exclude the equality.  The independent
catalogue audit checks these precise assertions on every positive
state.  It also records a useful control: in seventy-two positive
normalized A states, images from the two different nonselfdual
pairs are opposite.  Thus no assertion about arbitrary active
labels is intended.
\end{proof}

\begin{corollary}[The orthogonal-six forms at the small endpoints]
\label{paper:orthogonal-six-finite-forms}
\label{paper:rank14-orthogonal-six-forms}
For a unique orthogonal degree-two component of codegree six, the
rank-independent list of Corollary~\ref{paper:orthogonal-six-three-forms}
contains every rational positive form.  Its nonzero skew-norm list is
\((6)\), or \((3,3)\).  In the rank-thirteen input, \(q=0\)
permits only the first, and \(q=1\) permits both.
\end{corollary}
\begin{proof}
Apply the complete integral Gram classification and use \(b=q+1\).
A proposed certificate may use a larger six-form domain, but the
three additional forms in a stored input cannot inherit a negative
certificate merely by enlargement. Such a certificate is valid only
if every required row and branch is regenerated and checked over
that entire larger domain, as required by
Lemma~\ref{paper:source:certificate-principle}.
The classification is independent of the endpoint exclusions.
\end{proof}

\label{paper:source:codegree-six}
Suppose $R_{\mathbb C}$ has exactly one nonlinear component, of degree
two and codegree six.  Its character is rational because the component
is unique.  The symplectic divisibility theorem excludes symplectic
duality.  In the orthogonal case write $E$ for the matrix representation,
$s_X=\operatorname{Tr}E(X)$, and
$t_X=\operatorname{Tr}(E(X)E(X)^*)$.

\subsection{Integral supports and the positive fusion cone}
The pure Gram matrix satisfies $G^2=12G$ and has rank three.  For a
selfdual row put $J_X=G_{XX}$; for a representative of a dual pair put
\[
 J_X=(G_{XX}+G_{X,X^*})/2,\qquad
 K_X=(G_{XX}-G_{X,X^*})/2.
\]
Then $t_X=(s_X^2+J_X+K_X)/2$ and
$\det E(X)=(s_X^2-J_X+K_X)/4$, with $K_X=0$ for a selfdual row.
An active label is one whose pure matrix image is nonzero.

The complete support theorem
\ref{paper:orthogonal-six-patterns} gives patterns A, B, D and E,
with their signed rows and rational forms.  The first three have
algebra \(B_{2,3}\); E is split.  We now exclude E directly from
the positive fusion cone, independently of the source rank or linear
fields.  This uses the complete fixed-plane and dimension relations
proved in the preceding support classification.

\begin{theorem}[Exclusion of pattern E at dimension 5400]
\label{paper:source:orthogonal-E-cone-exclusion}
Assume that the source is integral, has dimension $5400$, and every
nonunit simple has dimension at least three. Suppose that its unique
nonlinear complex component is $M_2(\mathbb C)$ of formal codegree six,
with orthogonal positive involution. Then its support cannot be pattern E.
In particular its rational matrix algebra is the nonsplit algebra
$B_{2,3}$, and the remaining orthogonal supports are A, B and D.
No center hypothesis or bound on the source rank is used in this theorem.
\end{theorem}

\begin{proof}
Suppose E occurs. Write $P,Q$ for its two selfdual nonunit labels and
$U,U^*,V,V^*$ for its two pairs. The positive Perron relations give
\[
 d_P=2p-\tau,\qquad d_Q=4p-\tau,\qquad
 d_U=d_V=p,\qquad \tau\in\{1,-1\},\quad p\ge3.
\]
Every other nonunit label has zero matrix image. The fixed-plane frame
and the skew norms in pattern E give the following normal form, after
choosing representatives of the two pairs:
\[
 E(P)=\tau\begin{pmatrix}2&0\\0&0\end{pmatrix},\quad
 E(Q)=\tau\begin{pmatrix}-1&0\\0&1\end{pmatrix},\quad
 E(U)=-\tau\begin{pmatrix}0&0\\3&1\end{pmatrix},\quad
 E(V)=-\tau\begin{pmatrix}0&-1\\0&1\end{pmatrix}.
\]
Here $A^*=H^{-1}A^{\mathsf T}H$ for $H=\operatorname{diag}(3,1)$.
For completeness, first diagonalize the trace-$2\tau$ selfdual image.
The two fixed pair vectors have equal norm and opposite off-diagonal
coordinates, while the trace-zero selfdual vector is forced by the
positive dimension relation. Each pair has skew squared norm three;
choosing its representative fixes the sign of that skew part. This gives
the displayed matrices after rescaling the first coordinate by
$\sqrt3$. Thus the normal form does not discard a signed E support.
In particular
\[
 E(UU^*)=\begin{pmatrix}0&0\\0&4\end{pmatrix},\qquad
 E(U^2)=-\tau E(U).
\]

We use three consequences of positivity. First, write the active part of
the selfdual product $UU^*$, including its unit, as
\[
 \mathbf1+aP+bQ+k(U+U^*)+k'(V+V^*).
\]
Comparing its four matrix entries with $E(UU^*)$ gives
\[
 k'=k,\qquad a=2k+\tau,\qquad b=4k+3\tau.
\]
If $\tau=1$, the least permissible value is $k=0$; if $\tau=-1$,
it is $k=1$. The least possible active dimension is therefore
\[
 B_+(p)=14p-3,\qquad B_-(p)=10p+3.
\]
Increasing $k$ by one increases this dimension by
$2\ell$, where
\[
 \ell=12p-3\tau.
\]
The total squared mass of the unit and the six displayed nonunit labels is
\[
 S=1+(2p-\tau)^2+(4p-\tau)^2+4p^2
   =24p^2-12\tau p+3.
\]
The inequality $S\le5400$ leaves $3\le p\le15$ for $\tau=1$
and $3\le p\le14$ for $\tau=-1$. Within these twenty-five pairs,
$p^2\ge B_\tau(p)$ leaves only
\[
 (\tau,p)=(1,14),(1,15),(-1,11),(-1,12),(-1,13),(-1,14).
\]
For each of the six, the remainder $p^2-B_\tau(p)$ is smaller than
$2\ell$. Consequently the least active coefficients are forced. We may write
\[
 UU^*=\begin{cases}
  \mathbf1+P+3Q+T_0,&\tau=1,\\
  \mathbf1+P+Q+U+U^*+V+V^*+T_0,&\tau=-1,
 \end{cases}
\]
where $T_0$ is a positive selfdual object supported on zero-image labels.
Put $r=d_{T_0}=p^2-B_\tau(p)$ and $M=5400-S$; thus $M$ is the total
squared dimension mass of all zero-image basis labels.

Second, every positive zero-image object with no unit has active part
\[
 c\bigl(P+2Q+U^*+V\bigr)
 +d\bigl(P+2Q+U+V^*\bigr),\qquad c,d\in\mathbb Z_{\ge0}.
\]
This follows by comparing matrix entries: if its active coefficients are
$(a,b,u,u',v,v')$, they satisfy
\[
 a=v+v',\quad b=2v+2v',\quad u=v',\quad u'=v.
\]
Each displayed positive kernel generator has dimension $\ell$.
If a zero-image simple $Y$ occurs in $T_0$ with multiplicity $m_Y$,
Frobenius reciprocity gives $[YU:U]=m_Y$. The product $YU$ has no unit
and has zero matrix image, so its active dimension is at least $m_Y\ell$.
It follows that
\[
 p\,d_Y\ge m_Y\ell,\qquad
 \|T_0\|^2=\sum_Ym_Y^2\le\frac{p}{\ell}\sum_Ym_Yd_Y
              =\frac{pr}{\ell}.
\]
Here $\|\ \|$ denotes the norm in which the simple basis is orthonormal.
Cauchy--Schwarz on the zero-image basis now yields the necessary condition
\[
 r^2\le M\|T_0\|^2\le M\left\lfloor\frac{pr}{\ell}\right\rfloor.
\]
The relevant exact values are
\[
\begin{array}{rr|rrrr|r}
 \tau&p&r&\ell&M&\lfloor pr/\ell\rfloor&r^2-M\lfloor pr/\ell\rfloor\\\hline
 1&14&3&165&861&0&9\\
 1&15&18&177&177&1&147\\
 -1&11&8&135&2361&0&64\\
 -1&12&21&147&1797&1&-1356\\
 -1&13&36&159&1185&2&-1074\\
 -1&14&53&171&525&4&709
\end{array}
\]
The four positive entries in the last column are contradictions.

Third, in the two remaining cases $\tau=-1$ and $p=12,13$, one has
$E(U^2)=E(U)$. The positive solutions of this matrix equation with no
unit have active part $U$ plus a nonnegative combination of the same
two positive kernel generators. Since $p^2<p+\ell$ in both cases,
\[
 U^2=U+T,
\]
where $T$ is positive and supported on zero-image labels, with
$d_T=p^2-p$. The based-ring norm decomposition gives
\[
 \|UU^*\|^2-\|U^2\|^2
 =\frac1{6}\left(
  \operatorname{Tr}(E(UU^*)^2)
  -\operatorname{Tr}(E(U^2)E(U^2)^*)\right)
 =\frac{16-4}{6}=2.
\]
Indeed the Perron terms cancel, and every other simple component is
linear, where both norm terms equal $|\chi(U)|^4$. The least active
formula above has squared norm seven, so
\[
 \|T\|^2=4+\|T_0\|^2
 \le4+\left\lfloor\frac{pr}{\ell}\right\rfloor.
\]
A second Cauchy--Schwarz inequality on the same zero-image basis would give
\[
 (p^2-p)^2\le M\left(4+\left\lfloor\frac{pr}{\ell}\right\rfloor\right).
\]
For $p=12$ this asserts $17424\le8985$, and for $p=13$ it asserts
$24336\le7110$. Both are impossible. This excludes pattern E.
The rational-form assertion follows from the complete support catalogue.
\end{proof}

\subsection{Local products for the dimension-three cubic triple}
\label{paper:sec:cubic-three-closure}

The codegree-nine cubic support becomes much more rigid when its
three equal source dimensions are three.  Throughout this subsection
the source is integral, every nonunit simple has dimension at least
three, and there is a full non-Perron linear-character orbit with field
$\mathbb Q(\zeta_9+\zeta_9^{-1})$, codegree nine, and cubic triple
dimension three.  We use the notation of
Corollary~\ref{paper:cubic-nine-positive-subtable}:
\[
 d_{X_i}=3,\quad d_Y=5,\quad d_Z=6,\quad d_W=12,
 \qquad S=X_1+X_2+X_3,
\]
where these six nonunit labels are distinct and selfdual.  With cyclic
indices the previously established products are
\begin{align}
 X_i^2&=1+Y+X_{i+1},&
 X_iX_{i+1}&=Z+X_i,\notag\\
 YX_i&=X_i+W,&
 ZX_i&=S-X_i+W,\label{paper:cubic-three-starting-products}\\
 Y^2&=1+Y+Z+P,& d_P&=13.\notag
\end{align}
Here $P$ is a positive integral combination of source labels.  These
elements commute with the $X_i$, and $P$ is selfdual.  Write
$\langle U,V\rangle=\tau(UV^*)$ and $\|U\|^2=\langle U,U\rangle$.
For a chosen cubic character $\chi$, the values on $X_1,X_2,X_3$ are
the three conjugates $t_1,t_2,t_3$ of $t^3-3t+1=0$; its values on
$1,Y,Z$ are $1,1,-1$, and all remaining values are zero.

\begin{lemma}[Zero-value products and the last dimension-three labels]
\label{paper:cubic-zero-product-coefficients}
More generally, suppose the cubic support has dimensions
$(d,d,d,e,e+1)$ as in Theorem~\ref{paper:cubic-nine-type-support}.
For a simple $U$ outside $\{1,X_1,X_2,X_3,Y,Z\}$, there are
nonnegative integers $a,c$ and a positive combination $R$, supported
outside these six labels, such that
\[
 UU^*=1+aY+(a+1)Z+cS+R,
 \qquad
 d_U^2=e+2+(2e+1)a+3dc+d_R.
\]
In the dimension-three setting of this subsection, no such $U$ can
have dimension three.  Consequently
\[
                             n_3=3.
\]
Every simple outside the six displayed labels therefore has dimension
at least four.
\end{lemma}

\begin{proof}
The coefficient of the unit in $UU^*$ is one.  The only rational
relation among $t_1,t_2,t_3$ is their sum zero, and their rational
span contains no nonzero rational number.  Applying $\chi$ to
$UU^*$, whose value is zero, therefore makes its three triple
coefficients equal and its $Z$ coefficient one more than its $Y$
coefficient.  This proves both identities.  The same argument works
with the common negative sign of the cubic values.

For $d_U=3,d=3,e=5$, the dimension identity becomes
$9=7+11a+9c+d_R$.  Thus $a=c=0$ and $d_R=2$, impossible for a
positive unit-free combination when every nonunit dimension is at
least three.  The other displayed nonunit labels have dimensions
five and six, proving the exact value of $n_3$.
\end{proof}

\begin{lemma}[The dimension-four neighbors of the cubic triple]
\label{paper:cubic-dimension-four-neighbors}
Let $U$ be any simple of dimension four.  Exactly one of the following
necessary alternatives holds, with the same product for all three
$X_i$:
\begin{enumerate}
\item $UU^*=1+Z+R$, where $R$ is positive of dimension nine, supported
outside $\{1,X_1,X_2,X_3,Y,Z\}$, and $X_iU$ is one simple of
dimension twelve;
\item $UU^*=1+Z+S$, and $X_iU=U+V$ for one simple $V$ of dimension
eight.
\end{enumerate}
Selfduality of $U$ is not required.
\end{lemma}

\begin{proof}
The preceding lemma gives
$16=7+11a+9c+d_R$, so $a=0$ and either $c=0,d_R=9$ or
$c=1,d_R=0$.  Frobenius reciprocity and
\eqref{paper:cubic-three-starting-products} give
\[
 \|X_iU\|^2
   =\langle X_i^2,UU^*\rangle=1+a+c,
 \qquad [X_iU:U]=c.
\]
When $c=0$, a positive integral vector of squared norm one is a
single basis element.  When $c=1$, the squared norm is two and the
coefficient of $U$ is one; the remaining term is a single basis
element of dimension eight.  Finally $X_i-X_j$ belongs to the cubic
ideal and annihilates $U$, whose cubic projection is zero.
\end{proof}

\begin{theorem}[Three possible dimension-thirteen remainders]
\label{paper:cubic-three-remainder-types}
The positive remainder $P$ in
\eqref{paper:cubic-three-starting-products} has one of exactly three
necessary decomposition types:
\[
 \boxed{\ P=U_{13},\qquad
         P=U_4+V_9,\qquad
         P=U_6+V_7.\ }
\]
Subscripts indicate dimensions.  Every displayed summand is simple
and selfdual, and is distinct from $1,X_1,X_2,X_3,Y,Z,W$.

Put $k=\|P\|^2\in\{1,2\}$.  There is a selfdual positive
combination $Q$, supported outside
$\{1,X_1,X_2,X_3,Y,Z,W\}$, such that
\begin{align}
 X_iW&=Y+Z+W+P,& X_iP&=kW+Q,\notag\\
 YW&=S+(k+1)W+Q,& ZW&=S+(k+2)W+Q,
 \label{paper:cubic-three-next-products}\\
 d_Q&=39-12k.&&\notag
\end{align}
All these identities concern actual positive products.  They do not
assert that any of the three remainder types is realized.
\end{theorem}

\begin{proof}
The proof of Corollary~\ref{paper:cubic-nine-positive-subtable} gives
zero coefficients of $1,X_i,Y,Z$ in $P$.  Its dimension thirteen
also rules out $W$: one copy would leave dimension one with no unit,
and two copies exceed the dimension.  Its cubic projection is zero.

Using $W=(Y-1)X_i$ and the preceding products gives
\[
 X_iW=(Y-1)X_i^2=Y^2+W-1=Y+Z+W+P.
\]
Let $u=[YW:W]$.  Taking the squared norm of this product in two ways
gives
\[
 3+\|P\|^2=\|X_iW\|^2
 =\langle X_i^2,W^2\rangle=2+u,
\]
since $[X_{i+1}W:W]=1$.  Thus $u=k+1$.
Frobenius reciprocity gives coefficient one on each $X_i$ in $YW$,
zero on $1,Y,Z$, and the already computed coefficient on $W$.
Subtracting these terms defines $Q\ge0$ with the asserted support.
Both $Y$ and $W$ commute, so $Q$ is selfdual.  Dimensions give
$d_Q=39-12k$, hence $1\le k\le3$.

The case $k=3$ would make $Q$ have dimension three, contradicting
Lemma~\ref{paper:cubic-zero-product-coefficients}.  Thus $k=1$ or two,
and $P$ is multiplicity-free.  Its summands have dimensions at least
four, so the preliminary possibilities are $13$, $(4,9)$,
$(5,8)$, and $(6,7)$.  For a summand $U$ of $P$, reciprocity gives
$[X_iU:W]=[X_iW:U]=1$.  If $d_U=5$, removing $W$ leaves dimension
three.  This remainder is supported outside the six cubic-support
labels: the coefficients on $1,X_j,Y,Z$ are zero by
\eqref{paper:cubic-three-starting-products}.  It is impossible by
Lemma~\ref{paper:cubic-zero-product-coefficients}.
The pair $(5,8)$ is therefore excluded.
Selfduality of $P$ fixes its unique summand, or fixes both summands
when their dimensions differ.

Finally, direct multiplication gives
\[
 X_iP=Y^2X_i-X_i-YX_i-ZX_i=YW-S-W=kW+Q.
\]
The identity $(Y-Z)W=-W$ gives the formula for $ZW$.
\end{proof}

\begin{corollary}[More of the two-summand multiplication table]
\label{paper:cubic-two-summand-products}
If $P=U_4+V_9$, then
\begin{align*}
 X_iU_4&=W,& X_iV_9&=W+Q,& d_Q&=15,\\
 YU_4&=Y+Z+V_9,&
 ZU_4&=Y+Z+U_4+V_9.
\end{align*}
Moreover, either
\[
 U_4^2=1+Z+V_9
 \quad\text{or}\quad
 U_4^2=1+Z+U_4+T_5,
\]
where $T_5$ in the second alternative is a further selfdual simple
of dimension five and has zero cubic value.

If $P=U_6+V_7$, there are selfdual simples $T_6,T_9$ such that
\[
 X_iU_6=W+T_6,\qquad X_iV_7=W+T_9,\qquad Q=T_6+T_9.
\]
Here $T_9$ is a further label, while $T_6$ may equal $U_6$.
In particular this alternative forces at least nine selfdual
nonunit labels of dimensions
\[
                         (3,3,3,5,6,12,6,7,9).
\]
The repeated sixes refer to distinct labels $Z$ and $U_6$.
\end{corollary}

\begin{proof}
In the first case, the coefficient of $W$ and dimensions force
$X_iU_4=W$, and subtraction from $X_iP=2W+Q$ gives the next
identity.  Taking duals shows $U_4$ commutes with the $X_i$; since
$P$ does as well, so does $V_9$.  Comparing $X_i^2U_4$ with
$X_iW$ yields the formula for $YU_4$.  The formula for $ZU_4$
follows from $(Y-Z)U_4=-U_4$.

The coefficient identities in
Lemma~\ref{paper:cubic-zero-product-coefficients} now give
$U_4^2=1+Z+R$ with $d_R=9$ and zero cubic projection.  Furthermore
\[
 WU_4=(Y-1)W=S+2W+Q
      =X_iU_4^2=S+W+X_iR,
\]
so $X_iR=W+Q=X_iV_9$.  Comparing the coefficient of $W$ gives
$\langle R,P\rangle=1$.  Since every summand of $R$ has dimension
at least four, either $R=V_9$ or $R=U_4+T_5$.  Selfduality gives
the stated property of $T_5$.

In the second case, removing the forced $W$ leaves positive
remainders of dimensions six and nine.  The first is one simple
$T_6$.  Both remainders are supported outside the initial seven
labels.  Their sum is the selfdual $Q$.  If $T_6$ were nonselfdual,
its dual would have to occur in the dimension-nine remainder,
leaving the impossible dimension three.  Hence $T_6$ and the
dimension-nine remainder are selfdual.  The latter is either one
simple of dimension nine or a sum of dimensions four and five.
The second possibility would make $X_i$ times a dimension-four
simple contain $V_7$, contrary to
Lemma~\ref{paper:cubic-dimension-four-neighbors}.  Thus it is one
selfdual simple $T_9$.  Its dimension and support make it distinct
from all the other displayed labels.
\end{proof}

\begin{lemma}[A selfduality budget for the remainders]
\label{paper:cubic-remainder-selfduality-budget}
If $P$ has two summands, the source has at least ten selfdual labels
including the unit.  If $P=U_{13}$ is simple, then
\[
             [Q:U_{13}]\le1,\qquad \|Q\|^2\le4,
\]
and no dimension-four label occurs in $Q$.
\end{lemma}

\begin{proof}
For two summands the initial labels and $P$ already give nine
selfdual labels.  The $(6,7)$ alternative supplies the additional
selfdual $T_9$.  In the $(4,9)$ alternative, suppose there were no
further selfdual label.  The coefficient of $U_4$ in $Q$ is zero:
$[X_iV_9:U_4]=[V_9:X_iU_4]=[V_9:W]=0$.
Consequently the only possible selfdual summand of $Q$ would be
$V_9$.  Its odd dimension fifteen forces one copy of $V_9$,
leaving dimension six for nonselfdual pairs.  Every such pair has
dimension at least eight, a contradiction.

Now let $P=U_{13}$, and put $p=[Q:P]$.  Two copies of $P$ would
leave dimension one in $Q$, so $p\le1$.  Since $X_iP=W+Q$,
\[
 1+\|Q\|^2=\|X_iP\|^2
   =1+[YP:P]+[X_{i+1}P:P]
   =1+[YP:P]+p.
\]
Reciprocity gives coefficients one on $Y,Z$ and $p$ on $W$ in
$YP$.  Counting these and its coefficient on $P$ gives
\[
 65\ge5+6+12p+13(\|Q\|^2-p)=11+13\|Q\|^2-p.
\]
Thus $\|Q\|^2\le4$.  Finally, a dimension-four summand of $Q$
would make $X_i$ times that label contain $P$, impossible because
its product dimension is twelve whereas $d_P=13$.
\end{proof}

\paragraph{Exact audit and scope.}
The standalone program \path{verify_cubic_three_closure.py} checks
these cubic coefficient relations, local polynomial identities, and
finite positive-remainder alternatives. The written arguments supply
the positivity and Frobenius-reciprocity steps. Its polynomial model
is an identity-checking device and is not asserted to be a fusion
ring. The subsequent fusion-matrix completion theorem gives the
stronger consequence at small source ranks.

\subsection{A second cubic orbit beside the dimension-three triple}
\label{paper:sec:two-cubic-coexistence}

Retain the hypotheses and notation of
Section~\ref{paper:sec:cubic-three-closure}, and suppose there is a
second full non-Perron linear-character orbit with field
$\mathbb Q(\zeta_9+\zeta_9^{-1})$ and codegree nine.
Denote its three source labels by $A_1,A_2,A_3$, their common
dimension by $h$, and its positive and negative rational endpoint
labels by $U,V$.  The two orbits are distinct character components;
their fields may, of course, be isomorphic.

\begin{theorem}[The only endpoint extension of the dimension-three orbit]
\label{paper:two-cubic-three-endpoints}
The second endpoints necessarily have dimensions four and five,
with $V=Y$ and $U=U_4$.  The first remainder and the square of the
new endpoint satisfy
\[
             P=U_4+B_9,\qquad U_4^2=1+Z+B_9,
\]
where $B_9$ is a selfdual simple of dimension nine whose values in
both cubic components are zero.  Moreover $h\ge5$, and the source
contains at least twelve distinct selfdual labels, including the
unit, of dimensions
\[
                    (1,3,3,3,4,5,6,9,12,h,h,h).
\]
If $h=5$, all nine products $X_iA_j$ are the same single simple
$K$ of dimension fifteen.  The right products are $A_jX_i=K^*$;
selfduality of $K$ is not assumed.
\end{theorem}

\begin{proof}
The chain theorem, Corollary~\ref{paper:two-cubic-type-chain},
leaves the endpoint pairs $(4,5)$ or $(6,7)$, sharing respectively
$Y$ or $Z$ with the first pair $(5,6)$.  In the second case
$U=Z$ and the second character $\psi$ gives
\[
                 \psi(P)=\psi(Y^2-1-Y-Z)=-2.
\]
The positive combination $P$ has no unit or $Z$ support.  Its
coefficients on the second triple must be equal, by the cubic
rational relation.  Hence its coefficient on the second negative
endpoint of dimension seven is two.  This contradicts $d_P=13$.

Thus $V=Y$, $d_U=4$, and $\psi(P)=1$.  The same coefficient
argument, using the absence of $Y$ and the unit from $P$, gives
$[P:U]=1$.  Theorem~\ref{paper:cubic-three-remainder-types}
therefore forces $P=U_4+B_9$, and $\psi(B_9)=0$.
Corollary~\ref{paper:cubic-two-summand-products} gives either the
asserted square or $U_4^2=1+Z+U_4+T_5$, with $T_5$ distinct
from $Y$.  Applying $\psi$ to the latter alternative would give
$\psi(T_5)=-1$.  The only source label with second character
value $-1$ is its negative endpoint $Y$, a contradiction.

The second triple is disjoint from the first support.  In addition
$\psi(W)=\psi((Y-1)X_i)=0$, and $\psi(B_9)=0$ as just proved.
Thus its three labels are distinct from all nine known labels.
They are selfdual by the cubic support theorem.  The exact equality
$n_3=3$ gives $h\ge4$.

Let $T'$ be the positive, unit-free remainder for the second
positive subtable, so that
\[
 A_j^2=1+U_4+T'+H'_j,\qquad
 d_{T'}=\begin{cases}h(h-1)-5,&\varepsilon'=1,\\
                      h(h-2)-5,&\varepsilon'=-1.
       \end{cases}
\]
Here $H'_j$ is supported on the second triple, and that triple
and $U_4$ have first character value zero.  Consequently the
first character satisfies $\chi(T')=-1$.  Its coefficient rule
gives
\begin{equation}
 T'=aY+(a+1)Z+cS+R,\qquad
 d_{T'}=6+11a+9c+d_R,
 \label{paper:second-cubic-remainder-coefficients}
\end{equation}
where $a,c\ge0$ and $R$ is a positive unit-free combination
outside the first cubic support.  For $h=4$, the two possible
dimensions of $T'$ are seven and three.  Three is below the
mandatory dimension six, and seven would leave residual dimension
one.  Both are impossible, proving $h\ge5$.

If $h=5$, the remainder dimensions are fifteen and ten.  Equation
\eqref{paper:second-cubic-remainder-coefficients} gives $a=0$ and
$c\in\{0,1\}$, with $c=1$ possible only when $T'=Z+S$.
The differences $X_i-X_k$ belong to the first cubic ideal, and
$A_j-A_l$ belong to the second.  Each annihilates the other
triple.  Thus $F=X_iA_j$ is independent of both indices.
Frobenius reciprocity gives
\[
        \|F\|^2=\langle X_i^2,A_j^2\rangle=1+c,
        \qquad [F:A_j]=[A_j^2:X_i]=c.
\]
If $c=1$, the three distinct labels $A_j$ all occur in $F$,
contradicting $\|F\|^2=2$.  Hence $c=0$ and $\|F\|^2=1$,
so $F$ is one simple $K$ of dimension $3\cdot5=15$.
Duality gives $A_jX_i=K^*$.  The character values of $K$ in
both cubic components are zero.
\end{proof}

\begin{corollary}[Two cubic orbits do not fit the two-matrix rank-sixteen shape]
\label{paper:two-cubic-two-matrix16}
In the $5400$ source problem, the shape
$\mathbb C^8\oplus M_2(\mathbb C)^2$ cannot contain two full
codegree-nine real cubic linear-character orbits.  This assertion
does not require that either cubic triple have dimension three.
\end{corollary}

\begin{proof}
The two cubic orbits use six linear characters, each at canonical
degree $5400/9=600$.  Besides the Perron character, precisely one
linear character remains; call its canonical degree $a$.  Write
$b,c$ for the two nonlinear canonical degrees, each with
multiplicity two.  The canonical mass identity is
\[
                         a+2b+2c=1799.
\]
All three degrees belong to the established palette
\[
\begin{gathered}
10,12,15,18,20,24,30,36,40,45,50,54,60,72,75,90,100,108,\\
120,135,150,180,200,216,225,270,300,360,450,540,600,675,900.
\end{gathered}
\]
The odd entries give $a\in\{15,45,75,135,225,675\}$, all
divisible by fifteen.  Hence $b+c\equiv1\pmod3$ and
$b+c\equiv2\pmod5$.  Every palette entry is divisible by
three or five.  The first congruence makes at least one of $b,c$
not divisible by three; every such entry is divisible by five.
The second congruence then makes the other entry not divisible by
five, so it is divisible by three.  The congruences restrict these
two entries, respectively, to
\[
                         \{10,40,100\},\qquad\{12,72\}.
\]
Their sum is at most $172$, whereas the mass identity and
$a\le675$ require $b+c\ge562$.  This is impossible.
\end{proof}

\paragraph{Independent exact audit and scope.}
The standalone program \path{verify_two_cubic_coexistence.py}
checks the endpoint character values, cubic coefficient rules,
positive-remainder dimensions, and the canonical palette argument.
The written proof supplies the positivity and reciprocity steps.
These assertions restrict the coexistence of the two specified
character components; they do not assert realization of the local
product pattern.

\subsection{Conditional fusion-matrix completion for the small cubic triple}
\label{paper:sec:cubic-tail-exclusion}

The local products can be tested against a single fusion matrix before
attempting to construct the whole ring. This gives a stronger small-rank
conclusion than separate selfduality and dimension-square arguments.

\begin{lemma}[Strict squared row and column bounds]
\label{paper:fusion-matrix-row-norm}
In a fusion-simple integral fusion category, let $X$ be a noninvertible
simple of dimension $d$. Every row and every column of its fusion
matrix has squared Euclidean norm at most $d^2-1$. In particular
every fusion coefficient in this matrix is at most $d-1$.
\end{lemma}

\begin{proof}
The object $XX^*$ contains a nonunit simple $Y$. Fusion simplicity
makes $N_Y$ irreducible, as in
Lemma~\ref{paper:regular-trace-upper-bound}. Since $N_{XX^*}$ contains
$N_Y$ entrywise, it is irreducible. The positive dimension vector
is an eigenvector with eigenvalue $d^2$. Every row has a positive
off-diagonal entry, so its diagonal entry is strictly less than
$d^2$. The equality $N_{XX^*}=N_XN_X^{\mathsf t}$ identifies these
diagonal entries with the squared row norms. The same argument
applied to $X^*X$ gives the column bounds. Integrality then gives
the asserted inequalities, including the entry bound.
\end{proof}

\begin{lemma}[The smaller norm in the four-plus-nine remainder]
\label{paper:cubic-four-nine-q-norm}
In the $d=3$ cubic branch, suppose $P=U_4+V_9$. Write
$X_iV_9=W+Q$, where $Q$ has dimension fifteen. Then
\[
                        \|Q\|^2\le3,
\]
and $Q$ has no dimension-four summand.
\end{lemma}

\begin{proof}
Put $p=[Q:V_9]\in\{0,1\}$ and $q=[YV_9:V_9]$.
The product $X_iU_4=W$ gives $[Q:U_4]=0$.
Using $X_iP=2W+Q$ and $X_i^2=1+Y+X_{i+1}$, reciprocity gives
\[
 4+\|Q\|^2
 =\langle X_i^2P,P\rangle
 =2+(2+q)+p,
\]
where $\langle YP,P\rangle=2+q$ follows from
$YU_4=Y+Z+V_9$. Thus $\|Q\|^2=q+p$.
The object $YV_9$ contains
$Y+Z+U_4+pW+qV_9$, so
\[
                         45\ge15+12p+9q.
\]
For $p=0$ this gives $q\le3$, and for $p=1$ it gives $q\le2$.
The asserted norm bound follows. If a dimension-four simple $T$
occurred in $Q$, reciprocity would put $V_9$ in $X_iT$.
The dimension-four dichotomy of
Lemma~\ref{paper:cubic-dimension-four-neighbors} permits only one
dimension-twelve summand or dimensions four and eight, a contradiction.
\end{proof}

\begin{proposition}[A small symmetric tail matrix is necessary]
\label{paper:cubic-tail-matrix-model}
Under the standing integral, fusion-simple dimension-$5400$ source
hypotheses, with minimum nonunit dimension three, assume the full
non-Perron real cubic linear-character orbit of
conductor nine and codegree nine has its forced triple of source
dimension three. Remove the seven known labels
\[
                     1,X_0,X_1,X_2,Y_5,Z_6,W_{12}.
\]
Let $d=(d_j)$ list the remaining source dimensions, and let
$p_j=[P:T_j]$. For any $X_i$, the submatrix $B=(b_{jk})$ of
$N_{X_i}$ on these remaining labels satisfies
\begin{align}
 &B=B^{\mathsf t},\qquad b_{jk}\in\{0,1,2\},
       \qquad (3I-B)d=12p,\label{paper:cubic-tail-equations}\\
 &\sum_k b_{jk}^2+p_j^2\le8.\label{paper:cubic-tail-norms}
\end{align}
Every connected component of the graph of the off-diagonal entries
of $B$ meets $\operatorname{supp}p$, and $3I-B$ is positive definite.
The following additional row restrictions hold:
\begin{enumerate}
\item If $P=U_{13}$, the $U_{13}$ row is the coefficient vector
of $Q_{27}$. Its squared norm is at most four, its diagonal entry
is at most one, and every entry at a dimension-four label is zero.
\item If $P=U_4+V_9$, the $U_4$ row is zero. The $V_9$ row has
squared norm at most three, and all its dimension-four entries
are zero.
\item If $P=U_6+V_7$, the $U_6$ row has one entry equal to one,
at a dimension-six label $T_6$, and all other entries zero.
The $V_7$ row has one entry equal to one at the forced
dimension-nine label $T_9$, and all other entries zero.
The label $T_6$ may equal $U_6$; otherwise it is another
dimension-six label. The label $T_9$ is outside the initial seven.
\end{enumerate}
All remaining dimensions are between four and seventy-two, and
their squared sum is $5167$.
\end{proposition}

\begin{proof}
Each $X_i$ is selfdual, so its fusion matrix is symmetric. The
known local products show that the only connection from a
remaining label $T_j$ to the seven-label part is its coefficient
$p_j$ at $W$. The dimension equation therefore gives
\eqref{paper:cubic-tail-equations}. Lemma~\ref{paper:fusion-matrix-row-norm} gives \eqref{paper:cubic-tail-norms}, including
the entry bound two.

Fusion simplicity makes $N_{X_i}$ irreducible. Its seven-label
part is connected, so every remaining graph component must
meet a label with $p_j\ne0$. As a proper principal submatrix
of an irreducible nonnegative symmetric matrix of Perron root
three, $B$ has spectral radius strictly below three. Consequently
$3I-B$ is positive definite. For example, its two-by-two
principal minors give the additional necessary inequalities
\[
                  (3-b_{jj})(3-b_{kk})>b_{jk}^2\qquad(j\ne k).
\]

The first row restrictions are
Lemma~\ref{paper:cubic-remainder-selfduality-budget}; its no-four
assertion also follows directly because $X_iT_4$ has dimension
twelve and cannot contain $U_{13}$. The second follows from
$X_iU_4=W$ and Lemma~\ref{paper:cubic-four-nine-q-norm}.
The third is Corollary~\ref{paper:cubic-two-summand-products}.
The initial seven labels have squared mass
$1+3\cdot9+25+36+144=233$. Thus the remaining squared sum is
$5400-233=5167$. Exactness of $n_3=3$ excludes another
dimension-three label, and the stated source bound is seventy-two.
\end{proof}

\begin{conditionaltheorem}[The dimension-three cubic triple requires rank at least seventeen]
\label{paper:cubic-three-rank-at-least-seventeen}
Assume verification hypothesis $\mathsf V$ of
Section~\ref{paper:verification-status} for the finite scopes and
prerequisite exclusions specified in the proof below.
Under the standing integral, fusion-simple dimension-$5400$ center
hypotheses, if a full real cubic linear-character orbit of conductor
nine and codegree nine has its forced triple of source dimension
three, then
\[
                       \operatorname{rank}(\mathcal C)\ge17.
\]
The hypothesis concerns the common dimension of that particular
triple, not every occurrence of a dimension-three source simple.
\end{conditionaltheorem}

\begin{proof}
Theorem~\ref{paper:commutative-rank-thirteen-excluded},
Theorem~\ref{paper:both-first-rank-endpoints-excluded} and the initial
bounds of \cite[Section~30]{BDPfamily} exclude ranks below fourteen.
Their certificates do not use the present fusion-matrix argument. Suppose $14\le r\le16$. The three possible remainders
from Theorem~\ref{paper:cubic-three-remainder-types} force the
following minimal dimension multisets:
\[
\begin{array}{c|l}
 P&\text{forced source dimensions}\\\hline
 U_{13}&(1,3,3,3,5,6,12,13)\\
 U_4+V_9&(1,3,3,3,4,5,6,9,12)\\
 U_6+V_7&(1,3,3,3,5,6,6,7,9,12).
\end{array}
\]
The last row includes the independently forced simple $T_9$.
All further dimensions lie in $\{4,\ldots,72\}$ and make the
total squared sum $5400$. Complete unordered-square enumeration
gives the following input counts:
\[
\begin{array}{c|r|r|r|r}
 r&P=U_{13}&P=U_4+V_9&P=U_6+V_7&\text{union of types}\\\hline
 14&6523&1112&114&7641\\
 15&33190&7023&635&39901\\
 16&144857&34277&4805&177234.
\end{array}
\]
The first three columns count a type together with a choice of
remainder branch. A dimension type may occur in several columns;
the union column counts it only once.

For every branch and type, the finite necessary system of
Proposition~\ref{paper:cubic-tail-matrix-model} has no solution.
The exact verification package checks all $232536$ branch/type
cases, or $224776$ distinct dimension types across the three
ranks. Two independent algorithms are used. The first assigns
symmetric entries while maintaining exact weighted row budgets,
squared row bounds and two-by-two principal-minor inequalities.
The second generates whole row domains independently and enforces
their mutual symmetry; it does not use the first algorithm's
mutable budgets or spectral-minor pruning. Both retain all
necessary choices and return no matrix.

Completeness of the dimension input is separately checked by an
unbounded-square coefficient table, a nondecreasing integer
recursion, explicit type generation and inclusion--exclusion.
Equal-dimensional labels outside the fixed core can be permuted,
so $T_6$ in the last branch has precisely the two cases $T_6=U_6$
and $T_6\ne U_6$ when another dimension-six label is available;
one representative dimension-nine label can be chosen for $T_9$.
This loses no possible tail matrix. The finite systems therefore
cover every source type and fusion matrix allowed by the hypotheses,
giving the contradiction.
\end{proof}

The exclusion uses only one necessary fusion matrix after the local
cubic structure has been proved. It does not construct a full
character table or multiplication table, and it does not apply to
cubic triples of larger source dimension. The source-duality
choices for the unspecified labels are left unrestricted, which
enlarges the tested universe.

\begin{corollary}[A regular-character upper bound for a norm word]
\label{paper:norm-word-regular-upper-bound}
Under the hypotheses of Lemma~\ref{paper:fusion-matrix-row-norm},
with $r$ the source rank and $\lambda$ its regular character,
\[
 \lambda(XX^*)\le1+(r-1)(d_X^2-1).
\]
Equivalently, if $t_g=\sum_{E\in\mathcal G_g}\chi_E(XX^*)$
for a partition of all non-Perron irreducible characters into
groups of common representation degree $m_g$, then
\[
                   \sum_gm_gt_g\le(r-2)(d_X^2-1).
\]
\end{corollary}

\begin{proof}
The regular character of $XX^*$ is
$\Tr(N_XN_X^{\mathsf t})$, the sum of the squared row norms.
The unit row has squared norm one, by Frobenius reciprocity.
Apply Lemma~\ref{paper:fusion-matrix-row-norm} to the other rows.
Subtracting the Perron contribution $d_X^2$ gives the second
formulation. This further inequality is recorded for subsequent
row searches and is not claimed as a filter in the certified
canonical-length or two-matrix endpoint computations.
\end{proof}

\subsection{Five local patterns for a cubic triple of dimension four}
\label{paper:sec:cubic-four-initial}

\begin{theorem}[The dimension-four local alternatives]
\label{paper:cubic-four-five-patterns}
Assume the source is integral and fusion-simple, every nonunit
dimension is at least three, and the real conductor-nine linear
orbit of codegree nine has triple dimension $d=4$. In the notation
of Corollary~\ref{paper:cubic-nine-positive-subtable}, its alternatives
are precisely among the following necessary local patterns:
\[
\begin{array}{c|c|c|c}
\varepsilon&d_Y&d_Z&T\\\hline
+1&11&12&0\\
+1&3&4&U_8\\
+1&4&5&U_7\\
+1&5&6&U_6\\
-1&7&8&0.
\end{array}
\]
Each $U_s$ denotes a single selfdual source label of dimension $s$,
distinct from the unit, the cubic triple, $Y$, and $Z$. This is a
list of local alternatives, with no existence assertion.
\end{theorem}

\begin{proof}
Put $S=X_1+X_2+X_3$ and $e=d_Y$. The positive subtable gives
$\dim T=11-e$ for the positive sign and $7-e$ for the negative
sign. Since $e\ge3$, this dimension is at most eight. Cyclic
invariance rules out triple support in $T$, whose cost would be
at least twelve. Its zero cubic projection requires equal $Y,Z$
coefficients. The only possible mass for such a pair is seven
when $e=3$, leaving either zero or at most one; the actual
$\dim T$ values are eight or four, so this is impossible. Thus
$T$ is supported outside all six initial labels.

If a label $U$ has coefficient $t_U$ in $T$, reciprocity shows that
each $X_iU$ contains $t_U S$. Hence $d_U\ge3t_U$. Equality would
make the three triple labels a proper invariant block for right
multiplication by $U$, contradicting the irreducibility in
Lemma~\ref{paper:regular-trace-upper-bound}, applied also to right
multiplication by duality. Therefore $d_U\ge3t_U+1$.
The mass at most eight now forces every coefficient in $T$ to be
one, every supporting dimension to be at least four, and $T$ to
be either zero, one label of dimension four through eight, or
two distinct dimension-four labels.

Here are the further exclusions, retaining positivity throughout.
Suppose first $T=U_s$, so $U$ is selfdual. Write
\[
 a=[UU^*:Y],\quad b=[UU^*:U],\quad c=[UU^*:X_i].
\]
The last coefficient is independent of $i$. The identity
$(Y-Z)U=-U$ gives $[UU^*:Z]=a+1$. For the positive sign,
\[
 \|X_iU\|^2=1+a+b+c,\qquad
 s^2=e+2+(2e+1)a+sb+12c+\dim R,
\]
with $R$ positive and supported outside the listed labels.
Also $X_iU=S+Q$, where $Q$ has positive dimension $4s-12$,
so its squared norm is at least four. For $s=4,e=7$ the dimension
equation forces $a=c=0,b\le1$, contradicting that norm. For
$s=5,e=6$, it gives $13a+5b+12c\le17$ while $a+b+c\ge3$.
Only $(a,b,c)=(0,3,0)$ remains, leaving a positive nonunit
remainder of dimension two, again impossible.

For the negative sign the only additional case is $e=3,s=4$.
The corresponding norm is $1+a+b+2c\ge4$, but
$7a+4b+12c\le11$ forces $c=0$ and $a+b\le2$.
Finally let $T=U_4+V_4$ in the positive $e=3$ case; the labels
may be dual to one another. Set
$b=[UU^*:U]+[UU^*:V]$. Now $X_iU=S+W_4$ for one source
label $W_4$, so its squared norm is exactly four and
$a+b+c=3$. The dimension identity gives
$11=7a+4b+12c+\dim R$. It first forces $c=0$, and then gives
$11=12+3a+\dim R$, a contradiction. The five displayed
alternatives are all that remain.
\end{proof}

\subsection{Coupling the dimension-four cubic pattern to codegree six}
\label{paper:sec:cubic-four-six-small}

Throughout the following subsections the source is an integral
fusion-simple fusion category of dimension $5400$, every nonunit
simple has dimension at least three, and its complexified
Grothendieck algebra has a unique nonlinear component
$M_2(\mathbb C)$ of formal codegree six.  There is also a full
non-Perron real conductor-nine cubic orbit of codegree nine,
whose three distinguished simples have dimension four.  No
Drinfeld-center hypothesis is used in these subsections.
Use the five local alternatives of
Theorem~\ref{paper:cubic-four-five-patterns}, and write
$S=X_1+X_2+X_3$, $d_Y=e$, $d_Z=e+1$.

\begin{lemma}[The matrix images of the cubic support]
\label{paper:cubic-four-six-coupling}
The degree-two component is orthogonal and has pattern A or B.
In every one of the five local alternatives,
\[
 E(X_i)=E(Z)=E(T)=0,\qquad E(Y)=-I.
\]
Here $E(T)=0$ includes the alternatives $T=0$.
Consequently $X_i,Y,Z$ and the nonzero $T$ are central elements
of the source algebra.
\end{lemma}

\begin{proof}
The symplectic possibility is excluded by
Theorem~\ref{paper:source:symplectic-individual-codegree-divisibility}:
a rational codegree of an individually preserved symplectic
degree-two component is divisible by four, whereas six is not.

The differences $X_i-X_j$ belong to the cubic component, so the
three matrix images agree.  The orthogonal support catalogue
does not admit three equal nonzero pure selfdual images.  In A
this contradicts its positive three-vector relation; in B three
equal signed norm-five vectors could not be canceled by the
fourth norm-five vector; D and E do not have three equal-norm
selfdual rows of the required trace.  Hence the common image
is scalar.  Three equal nonzero rational scalar images would
use at least twelve of the trace-square norm, in addition to
the unit's four, whereas the total is twelve.  Thus they vanish.

The cubic idempotent gives $E(Y)-E(Z)=-I$.  The positive local
subtable gives $E(Y)+E(T)=-I$, so $E(Z)=-E(T)$.  If $T=0$ the
claim follows immediately.  Otherwise $T$ is a single selfdual
simple.  Lemma~\ref{paper:orthogonal-six-no-opposite-images}
excludes opposite active selfdual images, so $E(Z)$ and $E(T)$
are scalar.  There cannot be two distinct nonunit labels with
opposite nonzero scalar images in the catalogue.  Both vanish,
and again $E(Y)=-I$.  This also excludes E, which has no
nonunit scalar image.

In D the scalar image $-I$ has dimension $2p+1$, where $p$ is
the active pair dimension.  This excludes $e=3,4,5$.  For
$e=7$ it would give $p=3$.  Any zero-cubic-valued simple of
dimension three then satisfies $UU^*=1+Z$ by the dimension
identity below, but the D pair has matrix norm trace five,
whereas $\operatorname{Tr}E(1+Z)=2$.  For $e=11$, every
zero-cubic-valued simple has dimension at least five: dimension
three is too small, and dimension four would leave an
impossible zero-cubic remainder of dimension three.  D would
have $p=5$.  For that pair the same identity gives
$a=0$, $c\in\{0,1\}$ and a remainder of dimension at most
twelve but matrix trace three.  Such a remainder costs at
least fifteen, since the only available nonzero traces are
$\pm1$ on labels of dimension at least five.  This is impossible.
Only A and B remain.  Their asserted scalar matrix images,
together with the scalar linear components, imply centrality.
\end{proof}

\paragraph{Two elementary identities used below.}
For any simple $U$ outside $\{1,X_1,X_2,X_3,Y,Z\}$, the cubic
character gives
\begin{equation}
 UU^*=1+aY+(a+1)Z+cS+R,
 \quad
 d_U^2=e+2+(2e+1)a+12c+d_R,
 \quad
 \operatorname{Tr}E(R)=t_U-2+2a,
 \label{paper:cubic-four-zero-row}
\end{equation}
where $a,c\ge0$ and $R$ is positive and supported outside those
six labels.  This is the general zero-value coefficient identity
from Lemma~\ref{paper:cubic-zero-product-coefficients}, with the
matrix character also applied.

For a positive-sign alternative put $V=(Y-1)X_i$, independent
of $i$.  It is positive, selfdual, has no unit or triple
coefficient, and has dimension $4(e-1)$.  Write
\[
 Y^2=1+a(Y+Z)+cS+P,\qquad b=[Y^2:T]
\]
when $T$ is a nonzero simple.  Reciprocity and expansion of
$(Y-1)^2X_i^2$ give
\begin{equation}
 [V:Y]=[V:Z]=c,\qquad \|V\|^2=a+b+c.
 \label{paper:cubic-four-V-norm}
\end{equation}
Every summand here is an actual positive basis combination.

\begin{lemma}[Homogeneous scalar-image costs]
\label{paper:cubic-four-scalar-cost}
Consider a positive selfdual combination $Q$ supported away
from the unit and $Y$, with $E(Q)=\rho I$, $\rho>0$.
In A, if the positive-trace pair has dimension $p$ and the
negative-trace pair dimension $q$, then
$d_Q\ge2p\rho$.  If both pair traces are positive, the lower
bound is $2\min(p,q)\rho$.

For B put $h=e-1$.  After choosing the two antipodal pairs in
the fixed plane, their signed selfdual dimensions are
\begin{equation}
 (x_1,x_2,x_3,x_4)=
 \left(\frac{z+h+p}{2},\frac{z+h-p}{2},
       \frac{h-z+p}{2},\frac{h-z-p}{2}\right),\qquad z\ge0,
 \label{paper:cubic-four-B-dimensions}
\end{equation}
where $p$ is the nonselfdual pair dimension.  In the sign regimes
$z\ge p+h+2m$ and $p\ge z+h+2m$, where every displayed absolute
dimension is at least $m$, one has respectively
\[
 d_Q\ge\rho(z+h),\qquad d_Q\ge\rho(3p+h).
\]
When $h=2m$, the remaining possible regime is $z=p\ge2m$,
with lower bound $\rho(p+h)$.
\end{lemma}

\begin{proof}
For A the selfdual active trace-zero images can cancel only with
equal coefficients on all three labels.  A selfdual pair occurs
with equal coefficients on its two labels.  If these coefficients
are $k,l$, the scalar equation is $k-l=\rho$ for opposite traces,
or $k+l=\rho$ when both are positive.  The bounds follow.

For B let $n_i$ be the four selfdual coefficients, $k$ the
coefficient on each member of its pair, $s_i=\operatorname{sign}(x_i)$,
and $f_i=s_in_i$.  Vanishing of the fixed pure part and the
specified trace give
\[
 f_1-f_2=f_3-f_4=k,\qquad \sum_i f_i=2\rho,
 \quad f_3=\rho-f_2,\quad f_4=\rho-f_1.
\]
In the first regime the signs are $(+,+,-,-)$, forcing
$f_1\ge f_2\ge\rho$; its least mass is obtained at
$f_1=f_2=\rho,k=0$ and is $\rho(x_1+x_2)=\rho(z+h)$.
In the second the signs are $(+,-,+,-)$, forcing
$f_1\ge\rho$, $f_2\le0$ and $k\ge\rho$; the least mass is
$\rho(x_1+x_3+2p)=\rho(3p+h)$.  In the remaining equal regime
the same equations give the bound $\rho(p+h)$.  Contributions
from zero-image labels only increase the mass.  Formula
\eqref{paper:cubic-four-B-dimensions} follows by combining the two
signed dimension differences of magnitude $p$ with
$\sum_i x_i=2h$, then interchanging the two pairs to make
$z\ge0$.
\end{proof}

\begin{proposition}[The positive $e=3$ alternative forces rank at least seventeen]
\label{paper:cubic-four-e3-six-excluded}
The pattern $(e,d_T)=(3,8)$ forces source rank at least seventeen
under the standing assumptions.
\end{proposition}
\begin{proof}
Dimension nine forces $Y^2=1+P_8$: a $Y+Z$ contribution would
leave dimension one, and a triple contribution already costs
twelve.  Here $V$ has dimension eight and
$\|V\|^2=[Y^2:T]$.  Since $T$ has dimension eight, positivity
forces $Y^2=1+T$ and $V=W_8$ simple.  First suppose $W=T$.
Reciprocity then gives
\[
 YT\ge Y+Z+S.
\]
The remaining positive selfdual combination has dimension five
and matrix image $I$.  In A the pair relation is $p-q=2$, so
the scalar-image cost is at least $2p\ge10$.  In B every
active dimension is at least three; the two possible regimes
have scalar-image cost at least thirteen.  Both contradict
dimension five.

It remains essential to allow $W\ne T$, despite their equal
dimensions.  These eight initial labels have squared mass $202$.
In A, rank at most sixteen allows at most one further zero-image
label, of dimension $d$ with $d=0$ for absence.  Since $p-q=2$,
the dimension equation
\[
 3a^2+2p^2+2q^2+d^2=5198
\]
would give $3a^2+d^2\equiv2\pmod4$, impossible.

In B, the product $YT$ still contains $Y+Z$, and its triple
coefficients are zero because $[YT:X_i]=[T:YX_i]=[T:X_i+W]=0$.
Its remainder has dimension seventeen and image $I$.
The scalar cost bound excludes the large-$p$ regime
and gives $3\le p\le7$, $p+8\le z\le15$, with equal parities.
The complete dimension equation is
\[
 z^2+3p^2+d^2+e^2=5194,
\]
where zeros permit missing extras and every positive extra is
at least three.  Its only possibilities are
\[
 (p,z,d,e)=(4,12,39,59),\qquad(4,12,49,51).
\]
In both, the signed active dimensions are $(9,5,-3,-7)$.
An image-$I$ positive selfdual combination must contain the
first two active labels, of total dimension fourteen.  Any
additional active cone vector costs at least twenty-four, while
every available zero-image label has dimension at least eight.
The remaining dimension three is impossible.  The full input
inventory and the direct matrix-cone exclusion are independently
regenerated in \texttt{cubic\_four\_e3.py} and its auditor.
\end{proof}

\begin{proposition}[The positive $e=4$ alternative is impossible at rank at most sixteen]
\label{paper:cubic-four-e4-six-excluded}
The pattern $(e,d_T)=(4,7)$ forces source rank at least seventeen.
\end{proposition}
\begin{proof}
The dimension-sixteen square of $Y$ has $a\le1,c\le1$ in
\eqref{paper:cubic-four-V-norm}.  If $c=1$, then $a=b=0$ and
$\|V\|^2=1$, although $V$ contains both $Y$ and $Z$, a
contradiction.  If $a=1$, its remaining dimension-six part
would have matrix image $I$; the scalar cost is at least twelve
in A and fifteen in B.  Thus $a=c=0$.

Now $V$ has dimension twelve and norm $b>0$.  The square $Y^2$
cannot contain $2T$, which would leave dimension one.  Hence
$b=1$, and there is a simple zero-image $W$ of dimension twelve
such that
\[
 YX_i=X_i+W,\qquad Y^2=1+T+R_8.
\]
The remainder $R_8$ is positive selfdual with zero matrix image.
No active contribution can have dimension eight: in A its least
possible selfdual zero-image active mass is at least nine; in B
the two cone generators have masses $2z,z+3p$ in the first sign
regime and $3p+z,3p-z$ in the second, all larger than eight.
Nor can $R_8$ contain $T$ or $W$.  At rank at most sixteen it
uses at most one extra label in A or two in B.

For a summand $D$ of $R_8$ with multiplicity $r$, reciprocity
gives $YD\ge r(Y+Z)$, and the remainder has image $rI$.
This rules out every dimension-three, four or five summand
allowed by a decomposition of eight; two copies of a dimension-four
label are also excluded directly by $16<2(4+5)$.  Consequently
$R_8=D_8$ is simple.  The product $YD-Y-Z$ has dimension
twenty-three and matrix image $I$.

In A the pair relation is $p=q+3$.  The exact dimension equation is
\[
 3a^2+2p^2+2q^2+8^2=5117.
\]
Its only solutions with all nonunit dimensions at least three
are $(a,p,q)=(21,32,29),(35,20,17)$.  Their scalar costs
$2p$ exceed twenty-three.

In B, permitting one further zero-image label of dimension $d$
(or $d=0$ for its absence), the equation is
\[
 z^2+3p^2+d^2=5044.
\]
There are twenty-one descriptions satisfying the signed dimension
bounds.  Twenty have scalar cost greater than twenty-three.
The last is
\[
 p=3,\quad z=16,\quad d=69,\quad
 (x_1,x_2,x_3,x_4)=(11,8,-5,-8).
\]
Every label except the dimension-sixty-nine label then has
dimension at most twelve.  Multiplication by $Y$ of dimension
four cannot leave that proper basis subset, because $48<69$.
It contains the unit, so the fusion subcategory generated by
$Y$ is proper, contradicting fusion simplicity.

The complete integer inventories and these exclusions are
regenerated by \texttt{cubic\_four\_e4.py}.  The independent
\texttt{audit\_cubic\_four\_e4.py} checks the full dimension
sets through a square table and enumerates the four active
coefficients directly to audit the scalar-image cone.  Keeping
the two $p=3$ descriptions is a deliberate relaxation; they
could also be removed by the dimension-three matrix-norm test.
\end{proof}

\begin{proposition}[The positive $e=5$ alternative is impossible at rank at most sixteen]
\label{paper:cubic-four-e5-six-excluded}
The pattern $(e,d_T)=(5,6)$ forces source rank at least seventeen.
\end{proposition}
\begin{proof}
Equation~\eqref{paper:cubic-four-zero-row} excludes every
zero-cubic-valued simple of dimension three, so these dimensions
are at least four.  The positive vector $V$ has dimension sixteen.
If $c=1$, subtracting $Y+Z$ leaves dimension five and matrix
image $I$.  This must be one selfdual simple, contrary to the
catalogue: the only nonunit scalar image is $E(Y)=-I$.
Larger $c$ already exceed the dimension.  Thus $c=0$.

If $a=1$, the remaining part of $Y^2$ has dimension thirteen
and image $I$.  In A the pair relation is $p-q=4$, with
$q\ge4$, giving scalar cost at least sixteen.  In B the two
signed dimensions $x_2,x_3$ sum to four and have absolute values
at least four.  They therefore have opposite signs, so
$|z-p|\ge12$, and the scalar cost is at least twenty.
Thus $a=1$ is impossible; $a\ge2$ leaves dimension at most two.
Consequently $a=0$ and $b=\|V\|^2\ge1$.

Reciprocity gives $YT\ge b(Y+Z)$.  Its remainder is positive
selfdual with matrix image $bI$, so
\[
                       30\ge b(11+L),
\]
where $L$ is the homogeneous scalar-image cost.  In B,
$L\ge20$, an immediate contradiction.  In A, $L\ge16$
forces $b=1$, so $V$ is one zero-image simple of dimension
sixteen.  There are then fifteen forced labels before any
further zero-image label.  Allow its dimension to be $d$, with
$d=0$ denoting its absence.  The dimension equation is
\[
 3a^2+2p^2+2q^2+d^2=4998,\qquad p-q=4.
\]
The pair contribution is divisible by four, leaving
$3a^2+d^2\equiv2\pmod4$, impossible.  This closes A as well.
\end{proof}

\subsection{The negative dimension-four cubic alternative}
\label{paper:sec:cubic-four-negative-six}

\begin{conditionaltheorem}[The negative alternative forces rank at least seventeen]
\label{paper:cubic-four-negative-six-rank17}
Assume verification hypothesis $\mathsf V$ of
Section~\ref{paper:verification-status} for the finite scopes and
prerequisite exclusions specified in the proof below.
Under the standing assumptions of
Subsection~\ref{paper:sec:cubic-four-six-small}, the negative cubic
alternative $e=7,T=0$ has source rank at least seventeen.
\end{conditionaltheorem}

\begin{proof}
Its products are $X_i^2=1+Y+S-X_{i+1}$.  Put
$V=(Y-1)X_i$, of dimension twenty-four, and write
\[
 Y^2=1+t(Y+Z)+hS+P.
\]
Then $[V:Y]=[V:Z]=h$, $\|V\|^2=t+2h$,
$d_P=48-15t-12h$, and $E(P)=tI$.  The bound on $d_V$ gives
$h\le1$.

If $h=1$, the positive selfdual vector $V-Y-Z$ has dimension
nine and image $I$.  In A, opposite pair traces have positive
pair dimension at least nine and scalar cost at least eighteen.
If both pair traces are positive, their dimension sum is six,
so both dimensions are three.  An image-$I$ vector of dimension
nine must then contain a complete pair of mass six and at
least one further label; its squared norm is at least three.
Thus $\|V\|^2\ge5$ and $t\ge3$, contradicting $d_P\ge0$.
In B, a zero-cubic label of dimension three would satisfy
$UU^*=1+Z$ and hence $t_U=2$, but no B row outside $Y$ has
that norm trace.  All its zero-cubic dimensions are therefore
at least four.  With $h_B=e-1=6$ in
Lemma~\ref{paper:cubic-four-scalar-cost}, its scalar cost is at
least twenty-four, again excluding dimension nine.  Consequently
the coefficient $h$ in the displayed square is zero.

The scalar cost inequality $48\ge t(15+L)$ now gives $t=1$
in B and in A with opposite pair traces.  The remaining A case
has both pair dimensions three and permits only $t=1,2$.

\paragraph{The same-sign A case.}
For any member $U$ of one of its dimension-three pairs,
\eqref{paper:cubic-four-zero-row} gives $UU^*=1+Z$, so
$\|X_iU\|^2=1$.  Thus $X_iU$ is a zero-image simple $D$ of
dimension twelve.  Reciprocity gives $[X_iD:U]>0$.
In the matrix trace of $X_iD$, the only negative contribution
can come from $Y$, whose coefficient is $[V:D]$; all active
pair traces are positive, and the active selfdual traces vanish.
Therefore $[V:D]>0$.  The possibility $t=1$ would make $V$
simple of dimension twenty-four, a contradiction.  For $t=2$,
$V$ must consist of two distinct dimension-twelve zero-image
labels.  At rank at most sixteen, the dimension equation becomes
\[
                         3a^2+g^2=4914,
\]
where $g=0$ permits absence of the final extra label.  Modulo
four, $3a^2+g^2$ never equals two.  This closes the same-sign case.

\paragraph{The opposite-sign A case.}
Now $V=W$ is a simple zero-image label of dimension twenty-four,
and
\begin{equation}
 Y^2=1+Y+Z+P_{33},\qquad
 X_iW=Y+Z+2W+P_{33}.
 \label{paper:cubic-negative-W24}
\end{equation}
The pair dimensions satisfy $p=q+6$, with $p\ge9$.  The
scalar-image cost of $P_{33}$ gives $p\le16$.  At rank at most
sixteen the dimension equation is
\[
 3a^2+2p^2+2q^2+d^2+e^2=4662.
\]
Reduction modulo four forces $a$ even and two additional
zero-image labels of odd dimensions $d,e$.  The full integer
inventory has fifty-five descriptions.  Every $P_{33}$ must be
\[
 (M+M^*)+h(A_1+A_2+A_3)+uD+vE,
 \qquad 2p+3ah+du+ev=33,
\]
where $M$ is the positive-trace pair.  There are six allocations
in the complete inventory.  Reciprocity fixes
$[X_iU:W]=[P_{33}:U]$.  For a member $N$ of the negative-trace
pair this coefficient is zero, and the dimension of $X_iN$
must belong to the semigroup generated by
$3a,p+q,d,e$.  Five allocations fail this necessary condition.
The last is
\[
 (a,p,q,d,e)=(36,13,7,7,17),\qquad P_{33}=M+M^*+D_7.
\]
Its only possible dimension allocation gives $X_iN=4D_7$.
The distinct extra dimensions make $D_7$ selfdual, so also
$X_iN^*=4D_7$.  Reciprocity would force
$X_iD_7\ge4(N+N^*)$, of dimension fifty-six, whereas its
dimension is twenty-eight.

\paragraph{The B case.}
Again \eqref{paper:cubic-negative-W24} holds.  The signed
dimensions are
\[
 (x_1,x_2,x_3,x_4)=
 \left(\frac{z+6+p}{2},\frac{z+6-p}{2},
       \frac{6-z+p}{2},\frac{6-z-p}{2}\right).
\]
All four absolute dimensions and $p$ are at least four.
The dimension-four pair case is also impossible: its square
would have a remainder of dimension seven and matrix trace
two, while two positive unit traces cost at least eight.
The same argument with trace minus two excludes a zero-image
extra label of dimension four.  Hence $p\ge5$ and all extra
zero-image dimensions are at least five.

The scalar-image bound for $P_{33}$ excludes the large-$p$
regime and gives
\[
 5\le p\le13,\qquad p+14\le z\le27,
 \qquad z\equiv p\pmod2.
\]
In this regime $P_{33}$ must contain $A_1+A_2$; adding a
nontrivial zero-image active cone vector costs more than its
remaining budget.  Thus
\[
 P_{33}=A_1+A_2+R_{27-z},
\]
where the last term is either zero or one extra zero-image
simple.  If the extra dimensions are $d_1,\ldots,d_k$, $0\le k\le3$,
the complete equation is
\[
 z^2+3p^2+\sum_jd_j^2=4626,
 \qquad d_j\ge5.
\]
It has exactly fifty-six descriptions satisfying the remainder
condition, all at rank sixteen.

For completeness, the final row test is explicit.  For each
required simple $U\ne W$ outside the initial cubic support,
enumerate a possible row of $X_iU$ on
$(W,A_1,A_2,A_3,A_4,M,M^*,D_1,\ldots,D_k)$:
\[
 (w,n_1,n_2,n_3,n_4,l,l,u_1,\ldots,u_k).
\]
Its dimension is $4d_U$, its $W$ coefficient is
$w=[P_{33}:U]$, and, putting $f_j=\operatorname{sign}(x_j)n_j$,
the zero matrix image requires
\[
 f_1-f_2=f_3-f_4=l,\qquad \sum_jf_j=0.
\]
Let $c$ be the row coefficient of $U$ and $N$ the squared
coefficient norm.  Then $a=N-1-2c\ge0$ by the negative
subtable.  Fusion simplicity gives $c\le3$: a diagonal
coefficient four would make the corresponding basis ray
invariant under the dimension-four generator.  The remainder
in \eqref{paper:cubic-four-zero-row} must have
\[
 d_R=d_U^2-9-15a-12c\ge0,
 \qquad \operatorname{Tr}E(R)=t_U-2+2a.
\]
Its dimension and trace must be a nonnegative integer sum of
\[
 (|x_j|,\operatorname{sign}(x_j)),\quad
 (2p,0),\quad(24,0),\quad(d_j,0).
\]
Allowing each zero-image extra independently only relaxes its
possible duality.  Every one of the fifty-six descriptions has
a required $X_iU$ row for which these necessary conditions fail.
The missing dimensions are four in eleven cases, five in
twenty-one, six in fourteen, and seven in ten.

The producer \texttt{cubic\_four\_negative.py} regenerates both
complete inventories and every certificate.  The independent
\texttt{audit\_cubic\_four\_negative.py} uses raw square
inventories, enumerates all four active row coefficients before
testing the matrix kernel, and uses literal sets of trace values
for the remainder semigroup.  It verifies equality of the full
input sets and all missing-row certificates.  Thus all negative
alternatives at rank at most sixteen are excluded.
\end{proof}

\subsection{Simplicity of the dimension-forty continuation}
\label{paper:sec:cubic-four-pre-w}

This step is needed before the final A/B enumeration: a positive
element of dimension forty cannot simply be treated as a simple
source object. The following proof establishes the required simplicity.
Assume the dimension-four positive pattern has $e=11,T=0$, the
unique nonlinear component has codegree six, and the source rank
is at most sixteen. Use the orthogonal coupling described below:
$E(X_i)=E(Z)=0$ and $E(Y)=-I$. The symplectic possibility is
excluded by Theorem~\ref{paper:source:symplectic-individual-codegree-divisibility};
the coupled D and E cases are separately excluded in the local
codegree-six argument. Only supports A and B enter this subsection.

\begin{lemma}[The individual product equations before simplicity]
\label{paper:cubic-four-pre-w-product-equations}
Put $V=(Y-1)X_i$. It is independent of $i$, positive selfdual,
has no unit or triple support, has dimension forty, and has zero
matrix image. Write
\[
 h=[V:Y]=[V:Z],\quad t=[Y^2:Y]=[Y^2:Z],\qquad
 Y^2=1+t(Y+Z)+hS+P.
\]
Then
\[
 \|V\|^2=t+h,\quad h\le1,\quad
 \dim P=120-23t-12h,\quad E(P)=tI.
\]
Every source label outside the six initial labels has dimension at
least five. For such a label $U$, put $v=[V:U]$ and let $R_U$ be
the nonnegative remainder in
\[
                     X_iU=v(Y+Z)+R_U.
\]
Its matrix image is $vI$ and its dimension is $4d_U-23v$.
If $c=[X_iU:U]$ and $a=\|X_iU\|^2-1-c$, then
\[
 0\le c\le3,\qquad a\ge0,\qquad
 UU^*=1+aY+(a+1)Z+cS+Q,
\]
where
\[
 \dim Q=d_U^2-13-23a-12c,\qquad
 \chi_E(Q)=\chi_E(UU^*)-2+2a.
\]
All remainders have zero cubic projection and the displayed
excluded support; $Q$ is selfdual.
\end{lemma}

\begin{proof}
The positive cubic subtable and reciprocity give the statements
about $V$ and the coefficients of $Y^2$. Direct expansion yields
$\tau((Y-1)^2X_i^2)=t+h$. The mass of $h(Y+Z)$ is $23h$,
so $h\le1$. Matrix evaluation gives $E(P)=tI$.
For any outside label, the identity $(Y-Z)U=-U$ forces the
$Y,Z$ coefficients of $UU^*$ to differ by one. Its triple
coefficients coincide, and the unit coefficient is one. Thus its
dimension is at least thirteen; dimension three is impossible.
At dimension four the remainder would have dimension three after
$1+Z$, with no possible outside support, so four is impossible too.
This proves the lower bound five.

Reciprocity gives the $v(Y+Z)$ terms in $X_iU$, and there is no
triple support since $T=0$. The norm identity is
$\|X_iU\|^2=\tau(X_i^2UU^*)=1+a+c$. The strict diagonal bound
for $N_{X_i}$ gives $c\le3$. The remaining formulas follow by
dimension and matrix trace. In these formulas the three equal
nonlinear images $E(X_i)=0$ make each $X_i$ central, so all
products used in taking duals have the asserted selfduality.
\end{proof}

\begin{conditionalproposition}[The dimension-forty element is simple]
\label{paper:cubic-four-W40}
Assume verification hypothesis $\mathsf V$ of
Section~\ref{paper:verification-status} for the finite scopes and
prerequisite exclusions specified in the proof below.
Under the hypotheses of this subsection,
\[
 V=W_{40}\quad\hbox{is a single selfdual source label},\qquad
 Y^2=1+Y+Z+P_{97}.
\]
\end{conditionalproposition}

\begin{proof}
The nonzero traces on $P$'s allowed support are $\pm1$, at labels
of dimensions at least five. Consequently $\dim P\ge10t$,
so $33t+12h\le120$ and $t\le3$. If $h=1$, the positive
selfdual element $V-Y-Z$ has image $I$, dimension seventeen,
and squared norm $t-1\le2$. A single scalar-$I$ label is
unavailable, since $Y$ is the unique nonunit scalar label and
has image $-I$. In A all active selfdual traces are zero, so
two such labels cannot sum to image $I$; a dual pair has even
dimension. These alternatives exhaust squared norm at most two.
In B the scalar-cone lower bound displayed below is at least
twenty, already greater than seventeen. This proves $h=0$.

It remains to rule out $t=2,3$. Then $V$ has $t$ distinct labels,
each with coefficient one. In B, any three active selfdual
traces have odd sum; a pair plus a selfdual active label also
has nonzero trace. A pair with a zero-image label cannot cancel
its nonzero fixed part. In A the only three-active-label
zero-image sum is the three equal-dimension selfdual labels,
whose total dimension is divisible by three and cannot be forty.
Lemma~\ref{paper:orthogonal-six-no-opposite-images} handles supports of size two. Thus
$V$ is supported on zero-image labels in both cases. Each such
label in $V$ has dimension at least seven, since $X_iU$ contains
$Y+Z$ and a positive remainder: dimension six would leave one.

For A, the pair dimensions satisfy either $p-q=10$, with
$p>q\ge5$, or $p=q=5$. Its active selfdual dimensions are
$a,a,a$. A positive selfdual scalar image $\rho I$ costs at
least $2p\rho$ in the opposite-sign case and $10\rho$ in
the same-positive case. For $t=2$, $P$ has dimension seventy-four
and image $2I$, forcing $15\le p\le18$ in the opposite-sign case.
The full dimension equation is
\[
 3a^2+2p^2+2q^2+\sum d_j^2=5086,
\]
with at most three zero-image labels; two of them sum to forty.
The exact bounded search gives twenty-seven descriptions. Testing
$P$'s complete scalar cone leaves just
\[
 (a,p,q)=(35,16,6),\quad(d_j)=(19,21,5),\quad
 P=2(M+M^*)+2D_5.
\]
For this description the only positive selfdual zero-image
product of dimension twenty on the allowed support is
$X_iD_5=4D_5$, contradicting the strict diagonal bound three.
For $t=3$, the opposite-sign case already exceeds the mass
fifty-one of $P$. With $p=q=5$, the dimension equation and
three zero-image dimensions summing to forty give only
$a=38$ and dimensions $(7,11,22)$. The product $X_iD_7$
would leave, after $Y+Z$, a positive selfdual element of
dimension five and image $I$, below the scalar-cone cost ten.

For B, write its four signed selfdual dimensions as
\[
 x=\left(\frac{z+10+p}{2},\frac{z+10-p}{2},
          \frac{10-z+p}{2},\frac{10-z-p}{2}\right),
 \qquad |x_i|\ge5,
\]
with $p\ge5,z\ge0$. Its homogeneous lower cost for image $I$
is $z+10$ when $z\ge p+20$, $3p+10$ when $p\ge z+20$,
and $p+10$ in the remaining case $z=p$. These costs are all
at least twenty. Thus $t=3$ would cost at least sixty inside
$P$ of dimension fifty-one. For $t=2$ the cost is at most
thirty-seven. The complete dimension equation is
\[
 z^2+3p^2+\sum d_j^2=4986,
\]
with at most four zero-image labels and two distinguished labels
summing to forty. It gives exactly four rank-fifteen and 155
rank-sixteen descriptions. Of these, 131 have no scalar-$2I$
remainder of dimension seventy-four. The other twenty-eight
fail an individual product condition from
Lemma~\ref{paper:cubic-four-pre-w-product-equations}.

Here are the exact finite tests, to make the computation reviewable.
For a positive element of matrix image $\rho I$, let $n_i$ be
the coefficients at the four selfdual active labels, let $k$ be
the common coefficient at the pair, and put $f_i=\operatorname{sgn}(x_i)n_i$.
The necessary and sufficient matrix equations are
\[
 f_1-f_2=f_3-f_4=k,\qquad\sum_i f_i=2\rho,
 \qquad n_i,k\ge0.
\]
All other coefficients belong to the zero-image labels. Enumerate
these equations at the required dimension. For every proposed
$X_iU$ row, the norm and diagonal formulas in the lemma determine
$a,c$ and the required dimension and trace of $Q$. Its possible
dimension/trace pairs lie in the additive semigroup generated by
\[
 (|x_i|,\operatorname{sgn}(x_i)),\quad (2p,0),
 \quad(d_j,0).
\]
Every one of the twenty-eight descriptions has a required product
with no eligible row. The independent replay uses the variables
$(k,n_2)$ instead of $(n_1,n_2)$, checks the literal matrix equations,
and computes semigroup sets independently. It also regenerates the
full 159-input set. Thus every $t=2,3$ case is excluded. The only
remaining norm is $t=1$, making $V$ a simple selfdual label.
\end{proof}

\subsection{The positive cubic pattern with endpoint dimension eleven}
\label{paper:sec:cubic-four-e11-complete}

Proposition~\ref{paper:cubic-four-W40} supplies, for the remaining
positive pattern $e=11,T=0$, a selfdual zero-image simple $W$
of dimension forty and the products
\begin{equation}
 X_iY=X_i+W,\qquad X_iZ=S-X_i+W,\qquad
 Y^2=1+Y+Z+P_{97},\qquad X_iW=Y+Z+W+P_{97}.
 \label{paper:cubic-four-forty-products}
\end{equation}
The last equation follows by multiplying $W=(Y-1)X_i$ and
using $X_i^2=1+Y+X_{i+1}$.  The positive selfdual remainder
$P_{97}$ has matrix image $I$.  Its support excludes the initial
six labels, but it may contain $W$; this possibility is retained
in every calculation below.

\begin{conditionalproposition}[The A continuation is impossible through rank sixteen]
\label{paper:cubic-four-six-forty-continuation}
Assume verification hypothesis $\mathsf V$ of
Section~\ref{paper:verification-status} for the finite scopes and
prerequisite exclusions specified in the proof below.
The A support cannot realize \eqref{paper:cubic-four-forty-products}
at source rank at most sixteen.
\end{conditionalproposition}

\begin{proof}
All zero-cubic-valued simples have dimensions at least five.
The three active selfdual dimensions are $a,a,a$, and the pair
dimensions satisfy either $p=q+10$ with opposite traces, or
$p=q=5$ with both traces positive.  The seven-label core
$(1,X_1,X_2,X_3,Y,Z,W)$ has squared mass $1914$.  Thus,
allowing at most two extra zero-image labels,
\[
 3a^2+2p^2+2q^2+d^2+e^2=3486.
\]
In either sign case $p,q$ have the same parity.  Modulo four,
the equation forces $a$ even and exactly two extras of odd
dimensions $d,e$; an absent extra is represented by zero and
cannot satisfy this congruence.

If both traces are positive, a pair member $U_5$ has, by
\eqref{paper:cubic-four-zero-row}, either $a_U=c_U=0$ or
$a_U=0,c_U=1$.  In the first case $\|X_iU_5\|^2=1$ would
give a zero-image simple of dimension twenty, absent from the
type.  In the second $X_iU_5=U_5+V_{15}$, so
$E(V_{15})=-E(U_5)$ would have trace minus one, also absent.
This excludes the same-positive case.

Hence $p=q+10\ge15$.  A zero-image extra of dimension five
is likewise impossible: its square remainder would have trace
minus two, excluding $c_U=1$ with zero remainder; the other
case again requires an absent zero-image simple of dimension
twenty.  Therefore $d,e\ge7$ are odd.  The complete equation
has seventy-four descriptions.

Write $M,M^*$ for the positive-trace pair of dimension $p$ and
$N,N^*$ for the negative-trace pair of dimension $q$.  Every
possible positive selfdual scalar-image-$I$ allocation has form
\[
 P_{97}=h(A_1+A_2+A_3)+(k+1)(M+M^*)+k(N+N^*)
           +wW+uD+vE,
\]
with
\[
 3ah+2p+2k(p+q)+40w+du+ev=97.
\]
Nineteen of the seventy-four descriptions admit such an
allocation, with twenty-three allocations in total.

Reciprocity in \eqref{paper:cubic-four-forty-products} fixes
$[X_iU:W]=[P_{97}:U]$.  The five corresponding mass inequalities
\[
 40h\le4a,\quad40(k+1)\le4p,\quad40k\le4q,\quad
 40u\le4d,\quad40v\le4e
\]
leave eight allocations.  For a member $N$ of the negative-trace
pair, its $W$ coefficient is exactly $k$.  Every other term of
the zero-image product has dimension in the semigroup generated
by $3a,p+q,d,e$.  Seven of the eight allocations cannot represent
$4q-40k$ in this semigroup.

The sole remaining description and allocation are
\[
 (a,p,q,d,e)=(24,15,5,23,27),\qquad
 P_{97}=M+M^*+W+E_{27}.
\]
For each of the three active selfdual labels $A_j$, the product
$X_iA_j$ is selfdual, has zero matrix image and zero $W$
coefficient.  Its dimension ninety-six has exactly one allocation
in the generators $3a,2(p+q),d,e$:
\[
                         X_iA_j=3D_{23}+E_{27}.
\]
Reciprocity then forces
$X_iD_{23}\ge3(A_1+A_2+A_3)$, of dimension $216$, although
$d_{X_iD_{23}}=92$.  This contradiction closes the last case.

The producer \texttt{cubic\_four\_e11\_A.py} generates all
seventy-four inputs, all twenty-three allocations and every
certificate.  Its independent auditor uses a separate two-square
table, Cartesian coefficient enumeration and finite semigroup
sets.  Both give the progression
\[
                    74\longrightarrow19\text{ types}/23\text{ allocations}
                    \longrightarrow8\longrightarrow0.
\]
\end{proof}

\begin{conditionalproposition}[The B continuation is impossible through rank sixteen]
\label{paper:e11-B-rank16-excluded}
Assume verification hypothesis $\mathsf V$ of
Section~\ref{paper:verification-status} for the finite scopes and
prerequisite exclusions specified in the proof below.
The B support cannot realize \eqref{paper:cubic-four-forty-products}
at source rank at most sixteen.
\end{conditionalproposition}

\begin{proof}
Use signed dimensions
\[
 x=\left(\frac{z+10+p}{2},\frac{z+10-p}{2},
          \frac{10-z+p}{2},\frac{10-z-p}{2}\right),
 \qquad p\ge5,\quad |x_i|\ge5,\quad z\ge0.
\]
The four active selfdual labels and the active pair, together
with the seven-label core, account for thirteen labels.  If
there are $k\le3$ additional zero-image labels of dimensions
$d_1\le\cdots\le d_k$, the exact dimension equation is
\begin{equation}
 z^2+3p^2+\sum_jd_j^2=3386,\qquad d_j\ge5.
 \label{paper:e11-B-dimension-equation}
\end{equation}
It bounds $p\le33$, $z\le58$ and requires $z\equiv p\pmod2$.
The scalar-image cost of $P_{97}$ is
\[
 \mu=
 \begin{cases}
 z+10,&z\ge p+20,\\
 3p+10,&p\ge z+20,\\
 p+10,&z=p\ge10.
 \end{cases}
\]
Thus $\mu\le97$.  Enumerating the complete bounded equation
with this necessary inequality gives no rank-thirteen or
rank-fourteen description, sixty-three at rank fifteen, and
846 at rank sixteen.  No unverified preselected type list is
used as input.

Here is the final necessary row system, which excludes all
909 descriptions.  For a simple $U\ne W$ outside the initial
six labels, $X_iU$ has no initial-support coefficient by
reciprocity.  On the ordered labels
$(W,A_1,A_2,A_3,A_4,M,M^*,D_1,\ldots,D_k)$ write its row as
\[
       (w,n_1,n_2,n_3,n_4,l,l,u_1,\ldots,u_k).
\]
All coefficients are nonnegative integers and their dimension
sum is $4d_U$.  Put $f_j=\operatorname{sign}(x_j)n_j$.
Its zero matrix image is equivalent to
\[
          f_1-f_2=f_3-f_4=l,\qquad \sum_j f_j=0.
\]
Let $N$ be the row's squared coefficient norm and $c$ its
coefficient at $U$.  Frobenius reciprocity gives
\[
 a=N-1-c\ge0,\qquad0\le c\le3.
\]
The strict bound on $c$ follows from fusion simplicity and the
dimension-four generator, or directly from
Lemma~\ref{paper:fusion-matrix-row-norm}.  The remaining positive
selfdual part of $UU^*$ must have
\[
 d_R=d_U^2-13-23a-12c\ge0,\qquad
 \operatorname{Tr}E(R)=t_U-2+2a,
\]
where $t_U=3$ for an active selfdual label, four for a pair
member, and zero for a zero-image extra.  A necessary condition
is that this dimension/trace pair lie in the semigroup generated by
\[
 (|x_j|,\operatorname{sign}(x_j)),\quad(2p,0),\quad(40,0),\quad(d_j,0).
\]
These conditions deliberately allow extra zero-image labels
independently even when source duality would couple them.
Consequently failure is a valid exclusion for every such duality.

Every input description lacks at least one required $X_iU$
row.  For the 846 rank-sixteen descriptions the dimensions of
the first certified missing rows are
\[
\begin{array}{c|rrrrrrrrrrrr}
 d_U&5&6&7&8&9&10&11&12&13&14&15&26\\ \hline
 \text{count}&358&213&130&66&50&2&14&8&2&1&1&1.
\end{array}
\]
All sixty-three rank-fifteen descriptions fail the same test.
The producer \texttt{cubic\_four\_e11\_B.py} regenerates the
complete square inventory, solves the matrix kernel in the
variables $(n_1,n_2)$ and uses exact trace bitsets.  The independent
auditor \texttt{audit\_cubic\_four\_B\_and\_pre\_w.py}
regenerates the full 909-input set, uses $(l,n_2)$ and the literal
matrix equations, and computes trace semigroups as ordinary sets.
It verifies every final missing-row certificate.
\end{proof}

\begin{conditionaltheorem}[The complete dimension-four, codegree-six cubic exclusion]
\label{paper:cubic-four-six-rank17-complete}
Assume verification hypothesis $\mathsf V$ of
Section~\ref{paper:verification-status} for the finite scopes and
prerequisite exclusions specified in the proof below.
Let $\mathcal C$ be an integral fusion-simple fusion category
of dimension $5400$, with every nonunit simple of dimension at
least three.  Suppose its complexified Grothendieck algebra has
one nonlinear simple component, isomorphic to $M_2(\mathbb C)$,
of formal codegree six.  If a full non-Perron real cubic
linear-character orbit of conductor nine and codegree nine has
its distinguished triple of source dimension four, then
\[
                        \operatorname{rank}(\mathcal C)\ge17.
\]
This theorem does not require that a modular category under
consideration be a Drinfeld center.
\end{conditionaltheorem}

\begin{proof}
The complete local reduction
Theorem~\ref{paper:cubic-four-five-patterns} has five alternatives.
The negative one is excluded through rank sixteen by
Theorem~\ref{paper:cubic-four-negative-six-rank17}.  The positive
$e=3,4,5$ alternatives are excluded by
Propositions~\ref{paper:cubic-four-e3-six-excluded},
\ref{paper:cubic-four-e4-six-excluded} and
\ref{paper:cubic-four-e5-six-excluded}.  In the last positive
alternative $e=11$, Proposition~\ref{paper:cubic-four-W40}
proves that the dimension-forty continuation is simple.
The two propositions of this subsection then exclude its
complete A and B continuations.  Lemma~\ref{paper:cubic-four-six-coupling} excludes every other matrix support.  This exhausts all cases.
\end{proof}

\subsection{A cubic orbit and the source-rank restriction}
The integral covariance identity has a particularly rigid consequence
for a real linear-character orbit of conductor nine and codegree nine.
Theorem~\ref{paper:cubic-nine-type-support} identifies three selfdual
source labels of one common dimension, called its distinguished triple.
The cubic field hypothesis concerns the full character orbit.

\begin{conditionaltheorem}[The small cubic triple]
\label{paper:source:small-cubic-triple}
Assume verification hypothesis $\mathsf V$ of
Section~\ref{paper:verification-status} for the finite scopes and
prerequisite exclusions specified in the proof below.
Under \eqref{paper:source:assumption}, the distinguished conductor-nine,
codegree-nine triple cannot have dimension three at source rank at
most sixteen.  If the unique nonlinear component has degree two and
codegree six, the triple cannot have dimension four at those ranks
either.  Consequently, in the latter case at rank sixteen, its common
dimension is at least five.
\end{conditionaltheorem}
\begin{proof}
The dimension-three assertion is
Theorem~\ref{paper:cubic-three-rank-at-least-seventeen}.  Under the
specified matrix hypothesis, the dimension-four assertion is
Theorem~\ref{paper:cubic-four-six-rank17-complete}.  The distinguished
triple has common integral dimension at least three by
Theorem~\ref{paper:cubic-nine-type-support}; hence its dimension is
at least five in the stated rank-sixteen case.
\end{proof}

\section{Unramified local corners and Frobenius constraints}
\label{paper:source:corners}

\label{paper:source:local-corners}
The local argument uses the symmetric bilinear trace $\tau(xy)$.
Its integral idempotents need not be self-adjoint or central.

\begin{lemma}[Perfect integral corners]
\label{paper:source:perfect-integral-corners}
Let $\mathcal O$ be a complete discrete valuation ring and let $A$ be a
finite free $\mathcal O$-algebra with a perfect symmetric associative
pairing $\beta(x,y)=\tau(xy)$.  For every idempotent $e\in A$, the
restriction of $\beta$ to $eAe$ is perfect.
\end{lemma}

\begin{proof}
The integral map $P(x)=exe$ is an idempotent projection.  Cyclicity of
$\tau$ gives
\[
 \beta(Px,y)=\tau(exey)=\tau(xeye)=\beta(x,Py).
\]
Thus $A=\operatorname{im}P\oplus\ker P$ is an orthogonal direct sum.
Both summands are free over $\mathcal O$.  Their Gram determinants
multiply to a unit, so both determinants are units.  Since
$\operatorname{im}P=eAe$, this proves the assertion.
\end{proof}

Assume there is one nonlinear irreducible representation, of degree
two and rational formal codegree $F$, and that all linear codegrees
$f_i$ are rational integers.  Embeddings of each full field are counted
separately.  Call a prime $p$ eligible if every linear character field
is unramified at $p$.  A finite unramified extension $K/\mathbb Q_p$
then contains all linear values.  On
$\mathcal O_K\otimes_{\mathbb Z}R$, the trace pairing is perfect
because its basis Gram matrix is the permutation matrix for duality.

\begin{theorem}[Trace and parity of an unramified local corner]
\label{paper:source:unramified-corner-parity}
Partition the linear characters by equality of their complete reductions
over the residue field of $K$.  For each class $S$ there is an integral
idempotent $e_S$, with the different $e_S$ pairwise orthogonal, whose
linear coordinates are one precisely on $S$.  Write $k_S\in\{0,1,2\}$
for its reduced matrix rank.  Then
\[
 \sum_S k_S\le2,\qquad
 \sum_{i\in S}\frac1{f_i}+\frac{k_S}{F}\in\mathcal O_K,
 \qquad
 \sum_{i\in S}v_p(f_i)+k_Sv_p(F)\equiv0\pmod2.
\]
If the matrix algebra remains a division quaternion algebra over $K$,
every $k_S$ belongs to $\{0,2\}$.  The statements include $p=2$.
\end{theorem}

\begin{proof}
Each reduction class is a one-dimensional simple module of the residue
algebra.  In its semisimple quotient, take the central idempotent for
that simple factor.  Idempotents lift first through the nilpotent
Jacobson radical and then through the $p$-adically complete algebra.
Here is an explicit lifting argument which does not divide by two.
For an approximate idempotent $x$, the element $2x-1$ is a unit because
\[
 (2x-1)^2=1+4(x^2-x).
\]
Replacing $x$ by
$x-(x^2-x)(2x-1)^{-1}$ squares the error up to a unit.  This terminates
modulo a nilpotent ideal and converges in the complete setting.  To lift
a collection orthogonally, repeat the construction in the complementary
corner $(1-e)A(1-e)$ after each lift.  All these assertions hold in the
commutative subalgebra generated by the one approximate idempotent;
centrality of its lift in $A$ is not needed.

For a linear character $\chi$, the value $\chi(e_S)$ is an idempotent
of the field $K$, hence zero or one.  Its reduction determines which
value occurs.  This gives exactly the stated set of linear coordinates.
The matrix images of the lifts are orthogonal idempotents, so their
reduced ranks have sum at most two.  A division algebra has only the
idempotents zero and one, of reduced ranks zero and two.

Over $K$, the coefficient-of-unit trace decomposes as
\[
 \tau=\sum_i\frac{\chi_i}{f_i}+\frac{\operatorname{Trd}}F.
\]
In particular $\tau(e_S)$ gives the asserted integrality condition.
The matrix coefficient is $k_S/F$, not $2k_S/F$.
The fraction-field corner is $K^{|S|}$ together with a matrix corner
of dimension $k_S^2$.  The unscaled trace pairing on a rank-one corner
has determinant one; on a split rank-two corner it has determinant
$-1$ in matrix units.  In a quaternion presentation
$i^2=a$, $j^2=b$, $ij=-ji$, the reduced-trace pairing in the basis
$(1,i,j,ij)$ has determinant
\[
 -16a^2b^2=-(4ab)^2.
\]
Consequently its valuation is even, also at the dyadic prime.  A change
of fraction-field basis changes a Gram determinant by a square.
Lemma~\ref{paper:source:perfect-integral-corners} therefore gives
\[
 \sum_{i\in S}v_p(f_i)+k_S^2v_p(F)\equiv0\pmod2.
\]
Unramifiedness identifies these valuations with the ordinary valuations
of the rational codegrees, and $k_S^2\equiv k_S\pmod2$.
\end{proof}

\begin{corollary}[The valuation-one resource rule]
\label{paper:source:corner-valuation-one}
At an eligible prime $p$, let $n_u$ count the linear characters for
which $v_p(f_i)=1$ and $f_i/p\equiv u\pmod p$.  Counts include each
embedding in a full Galois orbit.  If $v_p(F)\ne1$, these characters
must pair in opposite residue bins.  If $v_p(F)=1$, deleting zero,
one or two entries from the single bin
\[
 u\equiv-F/p\pmod p
\]
must leave opposite bins balanced.  Each deleted entry consumes one
matrix rank.  At $p=2$, balance means even cardinality.  If the matrix
algebra is still division over $K$, no deletion is permitted.
\end{corollary}

\begin{proof}
For two distinct congruent linear characters, column orthogonality
applied to their difference gives
\[
 p^2\mid f_i+f_j.
\]
For three distinct congruent characters, pairing the differences of
the first with the other two gives $p^2\mid f_i$, and similarly for
the others.  Thus a class containing a valuation-one character consists
of either an opposite-residue valuation-one pair or a singleton.
A character of unit codegree is isolated as well: its integral
primitive central idempotent distinguishes it from every other
reduction class.  These central idempotents can be chosen first in
the orthogonal collection, and their matrix ranks are zero.

For a singleton write $f_i=pu$ and $F=p^tv$ with $u,v$ units.  Its
rank is nonzero.  If $k=p^{v_p(k)}w$, trace integrality forces
\[
 t=1+v_p(k),\qquad v+wu\equiv0\pmod p.
\]
The rank-two possibilities have $(k,t)=(2,1)$ for odd $p$ and
$(k,t)=(2,2)$ for $p=2$.  They both violate corner parity:
$1+kt$ is odd.  Only $k=t=1$ remains, with $u\equiv-v\pmod p$.
At most two such singleton classes fit in the matrix component.
Two deletions mean two rank-one singleton corners, not one rank-two
singleton corner.  Ordinary pairs may be assigned rank zero in this
necessary relaxation.  The division assertion follows from the
absence of rank-one idempotents.
\end{proof}

\subsection{Full fields and Frobenius}
Eligibility and residue degrees are read from the actual field, not
merely its displayed ambient conductor.  For a field
$\mathbb Q(\zeta_c)^H$, write $c=p^am$ with $(m,p)=1$ and put
\[
 I_p=\{g\in(\mathbb Z/c\mathbb Z)^\times:g\equiv1\pmod m\}.
\]
It is unramified at $p$ exactly when $I_p\subseteq H$.  In that case
choose $g_p\equiv p\pmod m$ and $g_p\equiv1\pmod{p^a}$.
The residue degree is the order of $g_pH$ in the quotient group
\cite[Chapters~6--8]{MilneANT}.  Thus redundant factors in an ambient
conductor cause no loss of eligible fields.

Let $h$ be the least common multiple of the linear residue degrees.
A rational quaternion algebra ramified at $p$ remains division over
the common unramified degree-$h$ field precisely when $h$ is odd
\cite[Proposition~13.4.4]{VoightQuaternion}.  Each admissible rational
quaternion alternative must therefore be tested in its own local
context.  A state can be excluded only if every alternative fails at
some eligible prime.  A ramified linear factor makes this particular
test inapplicable at that prime.

\begin{lemma}[Descent of a singleton corner to its unramified field]
\label{paper:source:singleton-corner-descent}
Use the hypotheses and notation of
Theorem~\ref{paper:source:unramified-corner-parity}.  Let $K_h/\mathbb Q_p$ be a
common unramified extension containing every linear character value,
and let $\sigma$ be its Frobenius generator.  Suppose a valuation-one
linear character $\chi$ is a singleton in the partition by complete
reduced character columns.  If its Frobenius orbit has length $l$, the
rational matrix algebra splits over the unramified degree-$l$ field
$K_l$.  In particular, if the matrix algebra is ramified over
$\mathbb Q_p$, then $l$ is even.
\end{lemma}

\begin{proof}
The complete reduction $\lambda=\bar\chi$ has the same Frobenius
period $l$ as $\chi$.  Indeed, a smaller period would make a distinct
Frobenius conjugate of $\chi$ congruent to it, contradicting
singletonness.  Thus $l\mid h$, the character values lie in $K_l$,
and $\lambda$ takes values in its residue field $\mathbb F_{p^l}$.

Set $A_l=\mathcal O_{K_l}\otimes R$.  The character $\lambda$
defines a one-dimensional simple factor $\mathbb F_{p^l}$ of the
semisimple quotient of its residue algebra.  Lift the central
idempotent for that factor to an idempotent $e\in A_l$, using the
lifting construction in Theorem~\ref{paper:source:unramified-corner-parity}.
The lift need not be central.  After extension to $\mathcal O_{K_h}$,
its linear coordinate is one exactly on characters reducing to
$\lambda$.  To see this directly, every reduced linear character
annihilates the radical.  On the selected simple factor there is
only one $\mathbb F_{p^l}$-algebra homomorphism to
$\mathbb F_{p^h}$; on every other factor the selected idempotent
is zero.  Singletonness therefore leaves exactly the coordinate
of $\chi$.

The trace and parity calculation over $K_h$ forces the matrix image
of $e$ to have reduced rank one.  Scalar extension does not change
that rank.  The matrix algebra over $K_l$ therefore already contains
a nonzero proper idempotent and must be split.  A ramified local
quaternion algebra splits over a finite extension exactly when the
extension degree is even, giving the last assertion.
\end{proof}

\begin{theorem}[Frobenius-equivariant valuation-one matching]
\label{paper:source:frobenius-corner-matching}
At an eligible prime $p$, let $c_{u,l}$ denote the number of
Frobenius cycles of length $l$ on linear characters with
$v_p(f_i)=1$ and $f_i/p\equiv u\pmod p$.  A necessary condition is
the following finite matching rule.

If $v_p(F)\ne1$, no singleton cycle may be removed.  If
$v_p(F)=1$, singleton removals are restricted to the bin
$u\equiv-F/p\pmod p$ and consist of one of
\[
 \varnothing,\qquad
 \text{one cycle of length one},\qquad
 \text{two cycles of length one},\qquad
 \text{one cycle of length two}.
\]
For a quaternion alternative ramified over $\mathbb Q_p$, both
length-one possibilities are forbidden, even when the common
unramified splitting field has even degree.

After the removal, if $p$ is odd, one must have
\[
 c_{u,l}=c_{p-u,l}\qquad\text{for every }u,l.
\]
If $p=2$, one must instead have $c_{1,l}$ even for each odd $l$;
even-length cycles have no further parity requirement in this
matching relaxation.
\end{theorem}

\begin{proof}
On reduced columns, $\sigma$ acts by the $p$th-power Frobenius.
Hence equality of complete reduced columns is a
Frobenius-invariant equivalence relation.  By
Corollary~\ref{paper:source:corner-valuation-one}, each valuation-one class
is a singleton or an opposite-residue pair.  Its singleton set is
therefore a union of complete Frobenius cycles, and has at most two
vertices.  Each singleton consumes one matrix rank.  This gives the
four displayed removal possibilities and the prescribed residue
bin.  Lemma~\ref{paper:source:singleton-corner-descent} gives the additional
restriction for a ramified quaternion alternative.

The pairs define a fixed-point-free involution commuting with
Frobenius on the remaining characters.  For odd $p$, it sends the
bin $u$ bijectively to the distinct bin $-u$.  An equivariant
bijection pairs cycles of equal length, which gives the displayed
equalities.  For $p=2$ there is only one bin.  A cycle can pair with
another cycle only if the lengths agree.  Within one cycle of
length $l$, an equivariant pairing is a translation by some $a$;
the conditions are $a\not\equiv0$ and $2a\equiv0\pmod l$.
Thus an internal pairing exists precisely for even $l$, and is
the half-turn.  Odd cycles must pair with other cycles of the
same length, giving their even-multiplicity condition.

Conversely, those cycle conditions construct the indicated
abstract equivariant matching, by pairing equal-length cycles
and using half-turns on unpaired even cycles.  This converse
concerns only the matching relaxation.  It does not construct
character values or a fusion ring.
\end{proof}

For a full abelian field of degree $d$ and residue degree $l$, the
Frobenius permutation on its embeddings has $d/l$ cycles of length
$l$.  The full-field assignment therefore determines all counts in
the matching criterion.  The criterion concerns possible reduction
classes; it does not construct a common character table.  Once actual
columns are supplied, trace integrality and determinant parity must
hold in each of their fixed classes and the complementary matrix-only
corner as well.

\section{Complete finite certificates and the source bounds}

\label{paper:source:certified-restrictions}
The finite computations use the necessary conditions just proved and
the specialized support lemmas retained in the following results.
We describe their common proof mechanism before stating the resulting
bounds. The compact package contains the bounded checks listed in its
status map, not the complete large inputs and certificates described
below. Each Conditional statement assumes $\mathsf V$ for its own
screening scope and all prerequisite exclusions. Named legacy certificate
files in these proofs identify the records that would have to be
regenerated and verified; their names do not assert that they are supplied
or have been independently replayed in this release.

\begin{lemma}[Source-pair budgets for a Galois-stable group]
\label{paper:source:aggregate-duality-budget}
A group of \(c\) representations of common degree \(m\) contributes
to the source dual-pair count one of the values
\[
 \frac{cm(m-1)}2+mq\quad
 (m\text{ odd},\ 0\le q\le\lfloor c/2\rfloor),
\]
or one of
\[
 \frac{cm(m-1)}2+mj\quad
 (m\text{ even},\ 0\le j\le c).
\]
These are necessary supersets; full scalar-field data determine the
linear contribution exactly.  Their sums satisfy \(2b\le r-3\)
at odd source rank and \(2b\le r-2\) at even source rank.
\end{lemma}
\begin{proof}
A real orthogonal degree-\(m\) component has skew dimension
\(m(m-1)/2\).  A real symplectic component, possible only for
even \(m\), adds \(m\).  A conjugate pair of components has skew
dimension \(m^2\), adding \(m\) to its two orthogonal baselines.
These give the stated ranges.  At even rank the unit forces a second
selfdual label.  At odd rank there must be at least three: if the unit
were the only one, the source dimension would be odd.
\end{proof}

\subsection{Finite universes and negative certificates}
At fixed canonical length, let $c_{a,m}$ count representations of degree
$m$ and canonical dimension $a$.  Removing the Perron character gives
\begin{equation}\label{paper:source:canonical-universe}
 \sum_{a,m}m c_{a,m}=L-1,\qquad
 \sum_{a,m}am c_{a,m}=5399,\qquad
 r=1+\sum_{a,m}m^2c_{a,m},\qquad c_{a,m}\in\mathbb Z_{\ge0}.
\end{equation}
These equations define a finite set over the palette $\mathcal A$.
Equivalently, one chooses a nonincreasing list of $L-1$ weighted degrees
and partitions each repeated multiplicity into representation degrees.
The generating function
\[
 \prod_{a\in\mathcal A}\prod_{m=1}^{L-1}
                   (1-x^m y^{am})^{-1}
\]
provides an independent count.  At a specified matrix shape, repeated
weighted entries are allocated to its nonlinear constituents instead.

For each repeated $(a,m)$ group, all eligible Galois orbit decompositions
are retained.  Full fields are abelian subfields of the bounded
cyclotomic conductor, enumerated by subgroups of its Galois group and
tested at every actual multiplicity.  This includes degrees five, six,
and eight whenever the input permits them.  A representation unique
at $(a,m)$ is rational; odd representation degree alone does not imply
rationality.  Rational nonlinear components retain every Schur and
involution alternative admitted by the necessary form tests.

The exact row model gives selfdual and paired row vectors $v$ and a
global target $T$.  The following elementary principle is the common
certificate behind the moment eliminations.

\begin{lemma}[Exact certificates for a complete necessary row set]
\label{paper:source:certificate-principle}
\label{paper:finite-endpoint-certificates}
\label{paper:integer-dual-pair-branching}
Let \(V=\{v_j\}\subset\mathbb Z^N\) contain every admissible
row vector in a specified case, including its atom-count and duality
coordinates.  Every actual case must express its target as
\(T=\sum_j x_jv_j\), with \(x_j\in\mathbb Z_{\ge0}\).
Each of the following is an exact exclusion certificate.
\begin{enumerate}
\item An integer covector \(z\) is nonnegative on every \(v_j\)
and has \(z\cdot T<0\).  More generally, if the nonnegative
coordinates give bounds \(x_j\le U_j\), it suffices that
\[
 z\cdot T<\sum_jU_j\min(0,z\cdot v_j).
\]
One may take the minimum of
\(\lfloor T_\ell/(v_j)_\ell\rfloor\) over positive entries of
\(v_j\) in coordinates nonnegative on the entire row set.
\item An exhaustive integer recurrence shows that a required
count, dimension-square, or canonical-remainder target is absent
from its finite semigroup.
\item Exactly \(b\) source-pair atoms are required.  Enumerate all
multisets of \(b\) eligible pair atoms, allowing repetition.  For
each residual target, a negative coordinate that is nonnegative on
all selfdual atoms, or a certificate of the first kind on those
atoms, excludes that branch.  Excluding every branch excludes the case.
\end{enumerate}
A certificate verified on a complete necessary row set also excludes
every subset, with the same target and with any retained multiplicity
bounds still valid. Enlarging the set requires a fresh check on every
added row; in the bounded version the new negative contributions and
valid bounds must also be included. A semigroup or branching certificate
must likewise be rerun on the enlarged generator set. An incomplete
enumeration or a failed search for a certificate gives no exclusion.
\end{lemma}
\begin{proof}
Dot the required equality with \(z\).  In the bounded version,
\(x_j(z\cdot v_j)\ge U_j\min(0,z\cdot v_j)\) term by term.
For a semigroup certificate, every actual decomposition is among the
enumerated integer choices.  For pair branching, different source
pairs may have identical vectors, which is why repetition is allowed;
the residual is a nonnegative integer sum of selfdual atoms.
Each asserted negative condition is therefore impossible for an
actual case on the declared input. Restricting that input only removes
possible decompositions. Enlarging it may introduce a decomposition and
invalidates the conclusion unless the witness is rechecked.
\end{proof}

The direction of this restriction is essential. For example,
$V=\{(1,0)\}$, $T=(1,1)$ and $z=(0,-1)$ satisfy the first
certificate test, but adjoining $(1,1)$ makes the target representable.
The first coordinate already records one atom, so an atom-count
constraint does not prevent this failure.

A valid certificate must record its precise input, full field and matrix
alternatives, dimension ranges, row construction, target, and exact
witness. Its verifier must regenerate the complete necessary superset
and check the witness with integer and rational arithmetic. In particular,
newly retained fields, basis involutions, nonsplit quaternion alternatives,
repeated spectral roots and enlarged dimension ranges all require renewed
coverage and witness checks. A hash or an aggregate count is not a
substitute for these checks.  For
negative dimension or atom claims, every row asserted absent is
regenerated in its complete stated range.  A numerical optimizer may
propose a separating vector, but its numerical status is not a proof.

Coverage is checked against the independently regenerated input
universe.  A merged certificate must include every required identifier
exactly once and close every required branch.  For overlapping tests,
exclusions are combined by the exact identifier sets.  Aggregate
counts alone neither establish coverage nor justify subtracting one
calculation from another.

\subsection{Canonical length and the first possible source ranks}
\begin{conditionaltheorem}[Source rank and canonical length]
\label{paper:source:global-bounds}
Assume verification hypothesis $\mathsf V$ of
Section~\ref{paper:verification-status} for the finite scopes and
prerequisite exclusions specified in the proof below.
Under \eqref{paper:source:assumption},
\[
 L\ge13,\qquad
 \begin{cases}
 14\le r\le462,&R\text{ commutative},\\
 16\le r\le490,&R\text{ noncommutative}.
 \end{cases}
\]
At noncommutative rank sixteen,
\[
 R_{\mathbb C}\simeq\mathbb C^{12}\oplus M_2(\mathbb C),
 \qquad L=14.
\]
More generally, a degree-$m$ irreducible representation requires
$r\ge13+m(m-1)$.
\end{conditionaltheorem}
\begin{proof}
We give the dependency order explicitly, since several of the finite
calculations use smaller-rank exclusions.  The starting estimates in
\cite[Section~30]{BDPfamily} are $L\ge10$, commutative rank at least
eleven, and noncommutative rank at least thirteen.

The complete canonical universes and the endpoint calculations have
the following coverage.  The entries are different, sometimes
overlapping counting units; no sum of the table is used.
\begin{center}
\begin{tabular}{lll}
Scope&Initial canonical unit&Number\\\hline
$L=10$, all ranks&multiplicity assignments&208\\
Commutative $r=11$&weighted profiles&497\\
Noncommutative $r=13$&one-matrix placements&830\\
Commutative $r=12$&weighted profiles&6\,547\\
$r=14$, one matrix&placements&12\,921\\
$r=15$, two matrices&placements&747\\
$L=11$, all ranks&multiplicity assignments&4\,475\\
Commutative $r=13$&weighted profiles&62\,776\\
$r=15$, one matrix&placements&143\,657\\
$r=16$, two matrices&placements&12\,257\\
$L=12$, all ranks&multiplicity assignments&69\,312
\end{tabular}
\end{center}
Under $\mathsf V$ for each of these scopes, the exact certificate
checks specified below would exclude them completely. Only the bounded
length-ten scope is independently replayed by the delivered compact
checks; the table does not certify the other exclusions.

First, the $208$ length-ten assignments are excluded by the determinant,
canonical residue, and skew-form argument of
Theorem~\ref{paper:canonical-length-at-least-eleven}; it uses only the
starting rank estimates.  Independently, the full-field rows and
integer source-pair branching of
Theorem~\ref{paper:both-first-rank-endpoints-excluded} exclude commutative
rank eleven and noncommutative rank thirteen.  The resulting field and
positive-word calculations exclude commutative rank twelve and the
one-matrix shape at rank fourteen
(Theorems~\ref{paper:commutative-rank-twelve-excluded} and
\ref{paper:noncommutative-rank14-exclusion}).
At rank fourteen, all other nonlinear shapes have $L\le10$ and are
already impossible.

The two-matrix rank-fifteen universe has $747$ placements.  Its exact
preliminary conditions reduce to $153$ placements, of which $27$ have
direct aggregate-remainder certificates.  The remaining $126$ give
$484$ full matrix states after the form tests; all are excluded in
Theorem~\ref{paper:rank15-two-matrix-excluded}.  The length-eleven
preliminary tests leave $589$ assignments.  Their $153$ rank-fifteen
cases are matched to this two-matrix calculation entry by entry;
the other $436$ are the higher-rank cases handled by the moment,
polynomial, and remainder certificates of
Theorem~\ref{paper:canonical-length-at-least-twelve}.  Therefore $L\ge12$.

Commutative rank thirteen is excluded by the complete $40\,476$-state
field ledger of Theorem~\ref{paper:commutative-rank-thirteen-excluded}.
The one-matrix rank-fifteen calculation excludes all $125\,545$ of its
field-duality states
(Theorem~\ref{paper:noncommutative-rank15-exclusion}).  Together with the
two-matrix calculation, these prove the rank lower bounds without
using $L\ge13$.

Finally, the $69\,312$ length-twelve multiplicity assignments have
$12\,905$ survivors of their exact preliminary conditions.  Of these,
$9\,379$ have higher rank and are excluded by
Proposition~\ref{paper:canonical12-higher-rank-exclusion}.
The other $3\,526$ are checked to lie in the complete $12\,257$-placement
rank-sixteen two-matrix universe.  The latter has $9\,268$ full
field/involution states after its necessary form tests, all excluded
by Theorem~\ref{paper:rank16-two-matrix-exclusion}.  Thus no length-twelve
assignment remains, proving $L\ge13$.

At noncommutative rank sixteen, the defect
\eqref{paper:source:nonlinear-defect} is positive, even, and at most
three.  It is therefore two, produced by exactly one degree-two
representation, giving the stated algebra and $L=14$.  The general
degree bound follows from the contribution $m(m-1)$ to the same
defect.  The upper bounds were established in
Proposition~\ref{paper:source:palette-and-type}.
\end{proof}

\begin{conditionalcorollary}[The remaining small-rank complex algebra shapes]
\label{paper:small-rank-wedderburn-shapes}
Assume verification hypothesis $\mathsf V$ of
Section~\ref{paper:verification-status} for the finite scopes and
prerequisite exclusions specified in the proof below.
Write $M_j=M_j(\mathbb C)$ and suppose the source ring is
noncommutative. At ranks sixteen through eighteen the only shapes
allowed by the canonical-length bound are
\[
 \mathbb C^{r-4}\oplus M_2,
 \qquad
 \mathbb C^{r-8}\oplus M_2^{\oplus2}
       \quad\text{with the second shape allowed only for }r\ge17.
\]
At rank nineteen they are
\[
 \mathbb C^{15}\oplus M_2,\quad
 \mathbb C^{11}\oplus M_2^{\oplus2},\quad
 \mathbb C^7\oplus M_2^{\oplus3},\quad
 \mathbb C^{10}\oplus M_3.
\]
In particular, a degree-three irreducible representation, or three
degree-two representations, requires $r\ge19$.
These are necessary algebra shapes, not realized fusion rings.
\end{conditionalcorollary}

\begin{proof}
The nonnegative even defect satisfies
\[
                 \sum_E m_E(m_E-1)=r-L\le r-13.
\]
A degree-two component contributes two, and a degree-three component
contributes six. A component of larger degree contributes at least
twelve. Enumerating these possibilities for $16\le r\le19$ gives
the lists; the number of linear components is then
$r-\sum_{m_E>1}m_E^2$.
\end{proof}

\subsection{The first source-rank endpoints}
\label{paper:sec:endpoint-character-search}
The base bounds from \cite[Section~30]{BDPfamily} are \(L\ge10\),
commutative rank at least eleven, and noncommutative rank at least
thirteen.  The following calculation uses these bounds only.

\begin{lemma}[A rational column forced by an odd block]
\label{paper:odd-canonical-block-rational-character}
Suppose $\mathcal B$ is a Galois-stable collection of $k$
linear characters of rational formal codegree $f$. Put
$c=\gcd(2f,f^2)$, using the identification of the character
field with that of its canonical central simple.
If every odd divisor of $\varphi(c)$ between three and $k$
is absent and $k$ is odd, then $\mathcal B$ contains a
rational, hence integer-valued, character. Removing one
such character leaves a Galois-stable collection.

In particular this applies to odd linear-character blocks
at $f=6,8,10,12$. A unique two-dimensional irreducible
representation also has integer-valued character, independently
of its formal codegree.
\end{lemma}
\begin{proof}
The conductor bound puts the full character field in
$\mathbb Q(\zeta_c)$. Every character orbit has size
dividing $\varphi(c)$. If an odd collection had no singleton,
one of its orbits would have odd size between three and
$k$, contrary to the hypothesis. The last assertion follows
because Galois conjugation preserves the representation
dimension and therefore fixes the unique representation.
\end{proof}

\begin{conditionalproposition}[Complete canonical and field input]
\label{paper:exhaustive-endpoint-input}
\label{paper:exact-endpoint-reduction}
Assume verification hypothesis $\mathsf V$ of
Section~\ref{paper:verification-status} for the finite scopes and
prerequisite exclusions specified in the proof below.
At commutative rank eleven, the ten nonunit weighted canonical
dimensions give \(497\) nonincreasing lists from \(\mathcal A\)
with sum \(5399\).  At noncommutative rank thirteen the algebra is
\(\mathbb C^9\oplus M_2(\mathbb C)\); allocating its double
canonical occurrence gives \(830\) placements on those lists.
The complete preliminary field and moment tests retain respectively
\(59\) profiles and \(79\) placements.
\end{conditionalproposition}
\begin{proof}
For commutative rank eleven, \(L=r=11\).  At noncommutative rank
thirteen, \(r-L\) is positive and even, whereas \(L\ge10\);
thus \(L=11\) and the defect is two.  This forces exactly one
degree-two representation.  Removing the Perron occurrence gives
the asserted canonical sum.  To make the finite enumeration explicit,
order \(\mathcal A=\{b_0<\cdots<b_h\}\) and use
\[
 P(0,t,j)=\mathbf1_{t=0},\qquad
 P(\ell,t,j)=\sum_{i=0}^jP(\ell-1,t-b_i,i).
\]
Then \(P(10,5399,h)=497\).  Summing, over these lists, the number
of distinct entries occurring at least twice gives \(830\).

For a repeated linear block of codegree \(f\), enumerate all
orbits with full field in \(\mathbb Q(\zeta_c)\), where
\(c=\gcd(2f,f^2)\).  An orbit of size at most three is rational,
quadratic, or, when \(9\mid c\), the real cubic field of conductor
nine.  Every block of size at least four has
\(c\in\{9,12,16,20,24\}\); enumerating all subgroups of the
corresponding unit groups is therefore complete.  At codegree six
use \(\mathbb Q\) and \(\mathbb Q(\sqrt{-3})\), by
Theorem~\ref{paper:linear-codegree-six-fields}; canonical degree ten
has no rational orbit by Corollary~\ref{paper:no-rational-degree-ten}.
Compare the full discriminant squareclasses and apply the duality
and codegree-six support bounds.  A rational degree-two component
has square determinant contribution, by
Proposition~\ref{paper:rank-thirteen-discriminant}.

For the preliminary row test, isolate rational columns using
Lemma~\ref{paper:odd-canonical-block-rational-character} and isolate
the matrix character when present.  Enumerate the complete necessary
rows of Lemma~\ref{paper:source:row-moments} for \(3\le d\le72\),
using the integral cyclotomic norm lattice.  This preliminary test
allows field decompositions to vary between rows and does not impose
source-pair alignment.  It is therefore a valid enlargement of the
fixed-field problem.  The targets and separating-covector proof are
those of Lemma~\ref{paper:source:certificate-principle}.
The field tests exclude \(338\) and \(586\) inputs; the row
certificates exclude \(318\) and \(557\), with intersections
\(218\) and \(392\).  Exact identifier unions consequently leave
\(59\) and \(79\).  The scripts \path{endpoint_core.py} and
\path{field_discriminant_profiles.py}, with the integer certificate
verifier, regenerate the complete input and both unions.
\end{proof}

\subsection{Closing the first endpoints by fixed-field rows}
\label{paper:sec:exact-endpoints}
\begin{conditionaltheorem}[Exclusion of both first source-rank endpoints]
\label{paper:both-first-rank-endpoints-excluded}
Assume verification hypothesis $\mathsf V$ of
Section~\ref{paper:verification-status} for the finite scopes and
prerequisite exclusions specified in the proof below.
Under the center hypothesis, neither a commutative source of rank
eleven nor a noncommutative source of rank thirteen can occur.
\end{conditionaltheorem}
\begin{proof}
Fix each full character field across the whole table.  The complete
positive trace-form tests reduce the \(59\) and \(79\) preliminary
cases of Proposition~\ref{paper:exact-endpoint-reduction} to \(24\)
and \(61\), with respectively \(31\) and \(99\) field/duality
records.  Enumerate their integral rows, including entries in proper
subfields, and impose the moments of
Lemma~\ref{paper:source:row-moments}.  Integer separators leave
\(11\) and \(35\) cases.  The annihilator divisibility test of
Lemma~\ref{paper:row-annihilator-test} and the resulting moments
exclude every commutative case and leave five orthogonal matrix
assignments on four noncommutative placements.

The remaining pair branches are small enough to specify completely.
The profile index is the zero-based descending canonical enumeration;
\(q\) is the number of nonreal linear-character pairs, so \(b=q+1\).
\[
\begin{array}{c|r|r|r|r|r}
\text{profile}&a_M&q&\text{eligible pair atoms}&
\text{pair multisets}&\text{separated branches}\\\hline
177&900&1&34&595&595\\
177&12&1&19&190&190\\
281&900&1&24&300&300\\
328&12&2&23&2300&877\\
328&12&1&11&66&66
\end{array}
\]
In the fourth row, \(1423\) multisets already exceed a nonnegative
moment budget; all other branches have exact integer separators.
Lemma~\ref{paper:source:certificate-principle} excludes all five
assignments.  The fixed-field and pair certificates are in
\path{first_endpoints_certificates.json} and
\path{first_endpoints_pair_certificates.json};
\path{verify_first_endpoints.py} regenerates every required row,
target, and pair multiset and verifies complete coverage.
\end{proof}

\subsection{Canonical length ten is impossible}
\label{paper:sec:canonical-length-ten}

\begin{lemma}[The skew determinant of an odd-degree rational matrix block]
\label{paper:odd-matrix-skew-determinant}
Let a rational simple-algebra component have odd degree $m$, rational
character field, and a positive involution of the first kind.
Then it is $M_m(\mathbb Q)$ and its involution is orthogonal.
If its symmetrizing trace is the reduced trace divided by $f$, the
positive form on its skew subspace has dimension $m(m-1)/2$ and
determinant squareclass
\begin{equation}\label{paper:odd-matrix-skew-class}
                      (2/f)^{m(m-1)/2}.
\end{equation}
In particular at $m=3,f=6$ this squareclass is three.
\end{lemma}
\begin{proof}
An involution of the first kind identifies the algebra with its opposite,
so its Brauer class has order dividing two. Its order also divides its
degree $m$. Since $m$ is odd, the class is trivial and the algebra is
split. A symplectic involution requires even degree, so the involution
is orthogonal. Write it as $X\mapsto H^{-1}X^{\mathsf t}H$.
After rational diagonalization, positivity permits
$H=\operatorname{diag}(h_1,\ldots,h_m)$ with $h_i>0$.

For $i<j$ the skew vector
$E_{ij}-(h_i/h_j)E_{ji}$ has norm $2h_i/(fh_j)$.
These vectors form an orthogonal basis. The product of their ratios
$h_i/h_j$ is
\[
                  \prod_{i=1}^m h_i^{\,m+1-2i},
\]
a rational square because $m$ is odd. This proves
\eqref{paper:odd-matrix-skew-class}. At $m=3,f=6$ the determinant
is $(1/3)^3$ modulo squares, hence has squareclass three.
\end{proof}

\begin{theorem}[Exclusion of canonical length ten]
\label{paper:canonical-length-at-least-eleven}
Under the standing center hypotheses at dimension $5400$, the
canonical algebra $A=I(\mathbf1)$ satisfies
\[
                  \operatorname{length}(A)
                    =\sum_E\dim E\ge11.
\]
This statement applies to every source rank and every complex
representation-multiplicity pattern. In particular it excludes
all twenty-six possible weighted canonical degree lists of
length ten, not only the rank-fourteen two-matrix placements.
\end{theorem}
\begin{proof}
The initial bound is $L=\sum_E\dim E\ge10$
\cite[Section~30]{BDPfamily}.
Suppose $L=10$. The nine nonunit weighted canonical degrees sum
to $5399$ and form one of exactly twenty-six nonincreasing lists in the canonical palette;
the complete profile recursion regenerates this universe.
At each degree $a$ of weighted multiplicity $k_a$, choose an
integer partition of $k_a$ into the dimensions of the distinct
source-ring representations at that degree. The complete enumeration
has $208$ such multiplicity assignments.

Every $k_a$ in these lists is at most five. Thus any odd-dimensional
nonlinear representation, of dimension three or five, is the only
representation of its dimension at that canonical degree. Its full
character field is rational. By
Theorem~\ref{paper:source:even-degree-trace-determinant}, only the fields
of the linear representations can change the determinant squareclass.
Moreover the rational baseline squareclass is independent of the
chosen multiplicity partitions and equals
\[
                           \operatorname{sf}\left(\prod_a a^{k_a}\right),
\]
where $\operatorname{sf}$ denotes positive squarefree part.
Indeed the number of odd-dimensional complex representations is even,
since it has the parity of $L=10$; the power of $5400$ in the
formal-codegree denominators is therefore a square. At each $a$,
the number of odd-dimensional representations has the same parity
as $k_a$.

The conductor-only field calculation is finite and small. The possible
products of absolute field-discriminant squareclasses in a repeated
linear degree block are contained in the following sets:
\[
\begin{array}{c|c|c}
 a&f=5400/a&\text{possible products of squareclasses}\\\hline
 900&6&1,3\\
 675&8&1,2\\
 600&9&1,3\\
 450&12&1,2,3,6\\
 540&10&1,5\\
 100&54&1,3\\
 12&450&1,3,5,15.
\end{array}
\]
These are conservative lists: no orbit size or multiplication is inferred
from a passing squareclass. For blocks with at most two linear characters,
they follow by listing quadratic fundamental discriminants dividing $2f$.
The larger blocks lie in the cyclotomic fields of conductors $9,12,16$.
Their complete subfield lists give exactly the displayed squareclass sets:
odd-degree Galois fields have square discriminant, and the quartic
subfields contribute squareclass one or two as appropriate.
The verifier constructs these subfields by all subgroups of the unit
group and computes their discriminant characters from the signs of
the regular quotient permutations.

Applying the determinant condition leaves $37$ assignments.
The stated source-rank bounds remove $16$ of them: a length-ten
source has even rank, and the discarded ranks are ten in the
commutative case or twelve in the noncommutative case.
This uses only the initial commutative and noncommutative rank bounds
eleven and thirteen from \cite[Section~30]{BDPfamily}.
Of the remaining $21$ assignments, $20$ have the forbidden pattern
of Proposition~\ref{paper:forbidden-canonical-24-200-pattern}.
The remaining counts by profile are
\[
\begin{array}{c|rrrrr}
 \text{profile index}&10&12&20&23&25\\\hline
 \text{assignments after determinant and rank}&1&1&7&1&11\\
 \text{after the forbidden pattern}&1&0&0&0&0.
\end{array}
\]

The unique remaining assignment is
\[
 (900,900,900,900,900,540,200,135,24),
 \qquad 900:\ (3,1,1),
\]
with every other weighted entry a singleton. Its source algebra is
$\mathbb C^7\oplus M_3(\mathbb C)$, of rank sixteen.
The degree-three character is rational and has formal codegree six.
The other two degree-$900$ constituents are linear. The determinant
condition requires them to form a quadratic orbit of discriminant
squareclass three. The conductor bound alone leaves
$K=\mathbb Q(\sqrt{-3})$ or $\mathbb Q(\sqrt3)$.
It is harmless here to retain the latter even though the earlier
linear-codegree-six theorem already excludes it.

By Lemma~\ref{paper:odd-matrix-skew-determinant}, the rational
matrix block is split and its three-dimensional skew form has
squareclass three. If $K=\mathbb Q(\sqrt{-3})$, the linear
skew form is $\langle1\rangle$, so the entire skew space has
dimension four and determinant squareclass three. But four source
dual pairs require the form $2I_4$, whose determinant is a square.
If $K=\mathbb Q(\sqrt3)$, the linear skew space is zero; the
entire skew space has dimension three and squareclass three,
whereas $2I_3$ has squareclass two. Both alternatives are impossible.

The complete reduction is therefore
\[
                       208\longrightarrow37\longrightarrow21
                              \longrightarrow1\longrightarrow0.
\]
The standard-library verifier \path{verify_length10.py} enumerates all
multiplicity assignments and all relevant conductor subfields, records
every determinant choice and tests both final skew alternatives.
Together with $L\ge10$, this proves $L\ge11$.
\end{proof}

\subsection{Commutative rank twelve}
\label{paper:sec:commutative-rank-twelve}
\begin{conditionalproposition}[The complete rank-twelve field input]
\label{paper:comm12-field-input}
Assume verification hypothesis $\mathsf V$ of
Section~\ref{paper:verification-status} for the finite scopes and
prerequisite exclusions specified in the proof below.
If the source Grothendieck ring is commutative of rank twelve, its eleven
nonunit canonical degrees form exactly one of $6\,547$ nonincreasing lists
from the central palette, with sum $5399$.
There are $319$ lists containing exactly one degree-ten constituent, which
are impossible because that constituent would have rational full field.
For the remaining $6\,228$ lists, exhaustive full-field enumeration gives
$172\,519$ field assignments before the trace-form tests. Of these, $17\,158$
satisfy the preliminary discriminant and duality conditions, and $2\,516$
assignments on $1\,070$ profiles pass the full rational trace-form test.
\end{conditionalproposition}
\begin{proof}
Commutativity makes every irreducible ring representation one-dimensional,
so every canonical multiplicity is one. Apply the recurrence of
Proposition~\ref{paper:exhaustive-endpoint-input} with eleven nonunit entries
instead of ten. A unique constituent of a given dimension is Galois fixed;
the rational degree-ten exclusion applies to the stated $319$ profiles.

For codegree $f$, put $c=\gcd(2f,f^2)$, and let $k$ be its number of
linear characters. A full field is an abelian subfield of
$\mathbb Q(\zeta_c)$ of degree at most $k$. Quadratic fields are enumerated
by all fundamental discriminants dividing $c$. The only possible cubic
field is $\mathbb Q(\zeta_9)^+$. For a quartic field every character of its
Galois group has order dividing four. Consequently its conductor divides
$\gcd(c,240)$: the wild parts at three and five disappear, and the
$2$-part has conductor at most sixteen. Enumerate every subgroup of index
four in the unit group for this smaller conductor. The only applicable
sextic field is $\mathbb Q(\zeta_9)$.

There are no missing degrees five or seven. The complete list of repeated
blocks with $k\ge4$, recording the maximal multiplicity, is
\[
\begin{array}{c|rrrrrrrrrrrr}
f&6&8&9&10&12&15&18&20&24&25&27&108\\\hline
k_{\max}&5&7&7&7&6&5&4&4&4&4&4&4
\end{array}
\]
Their conductor bounds exclude degree five whenever the multiplicity could
reach five, and exclude degree seven throughout. Degree six is available
only in the stated conductor-nine block. The codegree-six field list is
further restricted to $\mathbb Q$ and $\mathbb Q(\sqrt{-3})$ by Theorem~\ref{paper:linear-codegree-six-fields}. No rational orbit is allowed at canonical degree ten.
An independent subgroup-closure enumeration checks all $25$ applicable
pairs $(f,k)$ with $k\ge4$ against this field list, including the full
conductor-$216$ unit group; the lists agree exactly.

Partition each equal-codegree block into full-field orbits once for the
whole table. Let $q_6$ count the nonreal codegree-six pairs and $q$ all
nonreal pairs. Impose $q_6\le2$, $q\ge2q_6$, and $q\le5$.
The first inequality is the canonical mass budget, the second is the
disjoint-support theorem, and the last follows because even dimension
requires an odd-dimensional selfdual nonunit in addition to the unit.
The whole positive trace form and its two duality eigenspaces must be
\[
 I_{12},\qquad I_{12-2q}\mathbin\perp2I_q,\qquad 2I_q,
\]
respectively, by Theorem~\ref{paper:source:positive-rational-trace-form}.
Compute determinants and all local Hilbert invariants exactly. The
finite counts in the statement are the resulting exhaustive enumeration.
\end{proof}

\begin{lemma}[Eligible source dimensions and canonical remainders]
\label{paper:comm12-remainder-and-type}
Fix any full-field assignment for a commutative source of rank twelve,
and let $\mathcal R$ be any complete necessary superset of its source
rows. This implication does not require a complete catalogue of
assignments. Let $E$, $E_+$ and $E_-$ be the dimensions
appearing in all rows, selfdual rows and nonselfdual rows, respectively.
Then every canonical remainder $a-1$ belongs to the additive semigroup
generated by $E$. Also
\[
 5399=\sum_{i=1}^{11-2q}d_i^2+2\sum_{j=1}^q e_j^2,
 \qquad d_i\in E_+,\quad e_j\in E_-.
\]
Repetitions in these lists are permitted. Either failure excludes the full
field assignment.
\end{lemma}
\begin{proof}
For a canonical simple $Z$ of dimension $a$, adjunction gives
$[FZ:\mathbf1]=1$. Hence $FZ-\mathbf1$ is supported on nonunit source
simples and has dimension $a-1$. Every such constituent must have dimension
in $E$, proving the semigroup condition. The dimension identity is the
sum of squared dimensions of the eleven nonunit simples, with the two
members of each of the $q$ dual pairs counted together. Both tests are
finite integer recurrences; neither requires a source dimension to divide
$5400$.
\end{proof}

\begin{conditionaltheorem}[No commutative source of rank twelve]
\label{paper:commutative-rank-twelve-excluded}
Assume verification hypothesis $\mathsf V$ of
Section~\ref{paper:verification-status} for the finite scopes and
prerequisite exclusions specified in the proof below.
Under the center hypothesis, a commutative source cannot have rank twelve.
\end{conditionaltheorem}
\begin{proof}
The complete field input above leaves \(2\,516\) records.  For each,
enumerate the common integral rows for \(3\le d\le72\), with
\(11-2q\) selfdual atoms and \(q\) pair atoms.  Apply the canonical
remainder and dimension-square tests of
Lemma~\ref{paper:comm12-remainder-and-type}, the annihilator condition,
and exact moment certificates.  For each compatible fixed choice of
entry values impose
\begin{equation}\label{paper:comm12-power-numerators}
 d^j+\sum_i a_i\operatorname{Tr}_{\mathcal B_i}(X^j)
     \in5400\mathbb Z_{\ge0}\qquad(0\le j\le12),
\end{equation}
with a positive return by exponent twelve.  These are necessary by
Theorem~\ref{paper:source:finite-power-trace-test} and
Lemma~\ref{paper:source:short-positive-power}; the same entry
choices are used for every exponent.

Rational and quadratic entry recurrences leave nine assignments on
seven profiles.  The complete higher-field enumeration of
Lemma~\ref{paper:comm12-integral-entry-powers} closes them all.
Only \(\mathbb Q(\zeta_9)^+\), \(\mathbb Q(\zeta_5)\), and
\(\mathbb Q(\zeta_{20})^+\) occur in this terminal step: the
eligible dimension sets are empty, \(\{72\}\), or \(\{45,55\}\),
whereas the required remainders are respectively \(53\), \(23\),
or \(53\).  None belongs to its indicated semigroup.  The full
ledger \path{rank12_certificates.json} covers all \(2\,516\) field
records; \path{verify_rank12.py} regenerates the \(6\,547\) profiles,
all full-field assignments and necessary rows, and checks their
integer certificates by Lemma~\ref{paper:source:certificate-principle}.
\end{proof}

\subsection{The one-matrix shape at rank fourteen}
\label{paper:sec:noncommutative-rank-fourteen}
\begin{conditionaltheorem}[The noncommutative rank-fourteen endpoint is excluded]
\label{paper:noncommutative-rank14-exclusion}
Assume verification hypothesis $\mathsf V$ of
Section~\ref{paper:verification-status} for the finite scopes and
prerequisite exclusions specified in the proof below.
Under the center hypothesis,
\(R_{\mathbb C}\simeq\mathbb C^{10}\oplus M_2(\mathbb C)\)
is impossible.  Consequently a noncommutative source has rank at least
fifteen.
\end{conditionaltheorem}
\begin{proof}
The unique matrix character is rational.  The \(6\,547\) length-twelve
weighted profiles give \(12\,921\) placements; all orthogonal and
symplectic quaternion alternatives of
Corollary~\ref{paper:finite-quaternion-trace-forms} are retained.
The full field, canonical-residue and trace-form tests leave
\(7\,483\) field/duality states on \(2\,572\) placements.  Elementary
row and type certificates close \(83\) placements; the remaining
\(6\,907\) states are the final row scope.
\label{paper:noncommutative-rank14-exhaustion}

Use the common exact row model for \(3\le d\le72\), retaining
full cubic, quartic and sextic fields and all entries in their proper
subfields.  If \(q\) counts nonreal linear pairs, the source has
\(b=q+1\) or \(q+3\) pairs according to matrix duality.  In
particular the symplectic case \(q=3\), with its permitted zero-image
source pair, is included.  No rank-thirteen pair bound is imposed.

Pure powers and positive return are tested through exponent fourteen;
mixed words use \(0\le\ell\le k\le11\) whenever the stored
values determine their traces.  A determinant is fixed by the square
identity when available.  Unknown higher-field word data are retained
conservatively.  Within each test a single choice of values and
determinant must work at every exponent; permitting different choices
between tests enlarges the necessary row set.  The four terminal
higher-field states also use their full characteristic polynomials,
norm powers through fourteen, and
Proposition~\ref{paper:unique-matrix-commutator-congruence}.

All \(6\,907\) states have certificates from
Lemma~\ref{paper:source:certificate-principle}: canonical remainders
\(a-m\), the exact \(13-2b\) selfdual and \(b\) paired
dimension-square budgets, or integer moment separators.  For clarity,
the four higher-field terminal dimension sets and missing remainders are
\[
 (\{72\},23),\quad(\{45,55\},53),\quad
 (\{26,34,72\},14),\quad(\{72\},23).
\]
Each is an immediate semigroup obstruction.  The exact checker
\path{verify_rank14.py} regenerates the complete canonical and
field input, all preliminary placement exclusions, and every final
row certificate.  Its field audit includes every applicable quartic
field/codegree pair and the sextic field; the quaternion list is
complete.  Thus the displayed shape is excluded.

Every other nonlinear rank-fourteen shape has \(L\le10\), so
Theorem~\ref{paper:canonical-length-at-least-eleven} excludes it.
The first endpoints and the base bounds exclude the smaller ranks.
\end{proof}

\subsection{The canonical input for two degree-two components at rank fifteen}
\label{paper:supp:source:two-matrix15-input}

\begin{conditionalproposition}[A preliminary reduction for the two-block rank-fifteen shape]
\label{paper:rank15-two-blocks-preliminary}
Assume verification hypothesis $\mathsf V$ of
Section~\ref{paper:verification-status} for the finite scopes and
prerequisite exclusions specified in the proof below.
For the shape $\mathbb C^7\oplus M_2(\mathbb C)\oplus M_2(\mathbb C)$,
the following necessary tests reduce the initial $747$ placements to
$153$ placements on $104$ weighted canonical profiles:
\[
\begin{array}{l|r|r}
 \text{stage}&\text{placements}&\text{profiles admitting them}\\\hline
 \text{initial input}&747&490\\
 \text{rational degree-ten exclusion}&727&478\\
 \text{linear full-field determinant}&176&115\\
 \text{degree-24/200 pattern}&161&105\\
 \text{rational canonical residues}&153&104.
\end{array}
\]
Among the retained placements, $80$ have equal nonlinear canonical
degrees. For these the two matrix characters may be individually rational
or may form a quadratic Galois orbit. The preliminary reduction does not
assert that any retained placement supports a source fusion ring.
\end{conditionalproposition}
\begin{proof}
A representation unique at its pair $(a,m)$ of canonical degree and
representation dimension is rational. If $a=10$, this contradicts the
global rational degree-ten exclusion. This removes twenty placements.

There are seven linear characters including the Perron character. The
even-component determinant theorem discards both degree-two components,
also when they are a single quadratic Galois orbit or have nontrivial
Schur index. If $a_i$ are the six nonunit linear canonical degrees, the
remaining determinant condition is
\[
       5400\left(\prod_i a_i\right)
              \left(\prod_{\mathcal O\text{ linear}}
                         |\operatorname{disc}K_{\mathcal O}|\right)
                       \in\mathbb Q^{\times2}.
\]
The factor $5400$ remains because the number of linear characters is odd.
Equivalently its rational baseline is the squarefree part of the product
of $5400$ and all ten weighted nonunit canonical degrees; the nonlinear
weighted contributions are squares.

The linear-field catalogue is complete. A repeated linear degree block
has size at most five. Blocks of size at least four occur only at
codegrees six and eight. At codegree six use the proved fields
$\mathbb Q$ and $\mathbb Q(\sqrt{-3})$; at codegree eight use every
subfield of $\mathbb Q(\zeta_{16})$ of eligible degree, including its
three quartic subfields. All other blocks have size at most three.
Their quadratic fields are exactly the fundamental discriminants dividing
$\gcd(2f,f^2)$, and the only possible cubic field is
$\mathbb Q(\zeta_9)^+$. There are no eligible degrees five or higher.
Every partition of a linear block into these full-field orbits is allowed.
The determinant condition leaves $176$ placements.

Fifteen have the forbidden degree-$24/200$ pattern and are excluded by
Proposition~\ref{paper:forbidden-canonical-24-200-pattern}.
For each remaining constituent unique at $(a,m)$, write
$F(Z)=m\mathbf1+P$ and $r=a-m$. The scalar of $P$ on every source
irreducible representation is an integer, and the strict bound gives
an integer $\mu$ with $-r<\mu<r$ on its own representation.
If $g$ is the gcd of the other nonunit weights $m_Ea_E$, the trace
identity requires
\[
                              r+ma\mu\equiv0\pmod g.
\]
This is precisely the rational canonical residue obstruction.
Exactly eight further placements fail. In every case the failing
constituent has $(a,m)=(12,2)$ and $g=25$; the congruence
$10+24\mu\equiv0\pmod{25}$ has no integer solution in $[-9,9]$.

The standalone standard-library verifier
\path{verify_rank15_two_blocks.py} enumerates all $747$ placements,
all permitted linear-field determinant squareclasses, and every tested
integer interval. It retains the possible quadratic matrix orbit whenever
the two nonlinear degrees agree. A separately arranged multiplicity-partition
enumeration agrees on all $747$ squareclasses and stage decisions.
No aggregate source row or fusion multiplication is used for this
preliminary input.
\end{proof}

\begin{lemma}[Safe aggregate source rows for the two-block shape]
\label{paper:rank15-two-block-aggregate-rows}
Fix a weighted canonical profile and a two-matrix placement. At canonical
degree $a$, let $k_a$ be its weighted multiplicity and $j_a$ the number
of degree-two source representations there. Put $f_a=5400/a$.
For a nonunit source simple $X$ of dimension $d$, define
\[
 s_a=\sum_{E:\,a_E=a}\chi_E(X),\qquad
 t_a=\sum_{E:\,a_E=a}\chi_E(XX^*).
\]
Then $s_a,t_a$ are integers and satisfy
\begin{gather*}
 \sum_a a s_a=-d,\qquad \sum_a a t_a=5400-d^2,\qquad
 |s_a|\le k_a d-1,\qquad s_a^2\le k_a t_a,\\
 0\le t_a\le\min\left\{
 f_a-1,\ (k_a+2j_a)f_a-k_a,\ k_a d^2,
                 \left\lfloor\frac{5400-d^2}{a}\right\rfloor\right\}.
\end{gather*}
If $k_a=1$, then $t_a=s_a^2$.
These are necessary conditions even when the two nonlinear characters
form a quadratic Galois orbit or have nontrivial Schur index.
At rank fifteen it suffices to consider $3\le d\le72$.
\end{lemma}
\begin{proof}
Each degree block is Galois stable, so both aggregate traces are rational
algebraic integers. The two unit-coefficient identities give the displayed
linear and norm equations. The sum of representation dimensions in the
block is $k_a$, giving Cauchy--Schwarz and the strict trace bound.
The separate positive Perron contribution makes $t_a<f_a$.
The sum of squared representation dimensions in this block is
$k_a+2j_a$. Summing its norm traces over all source simples gives
$(k_a+2j_a)f_a$; subtracting the unit contribution $k_a$ gives
the stated total off-unit norm bound. The remaining two norm bounds follow
from the operator norm and the global norm equation, respectively.

When $k_a=1$, the block is a single rational linear character, hence its
norm trace is the square of its integral value. At rank fifteen there
are fourteen nonunit source simples, each of dimension at least three.
A dimension at least $73$ would require
$73^2+13\cdot3^2=5446>5399$, proving the range.

For these deliberately weak aggregate rows the total off-unit bound is
redundant: $(k_a+2j_a)f_a-k_a\ge f_a-1$. Thus their eligibility sets
depend only on the weighted profile, even though the subsequent
canonical-remainder tests depend on the nonlinear placement.
\end{proof}

\begin{conditionalproposition}[Aggregate rows and remainders for the two-block rank-fifteen input]
\label{paper:rank15-two-blocks-126}
Assume verification hypothesis $\mathsf V$ of
Section~\ref{paper:verification-status} for the finite scopes and
prerequisite exclusions specified in the proof below.
Among the $153$ placements in
Proposition~\ref{paper:rank15-two-blocks-preliminary}, the aggregate
source-row and canonical-remainder tests retain $126$ placements on
$86$ weighted canonical profiles. The necessary equation for fourteen
nonunit source squares removes none of these further. Of the $126$
retained placements, $64$ have equal nonlinear canonical degrees and
continue to allow a quadratic orbit of degree-two characters.
\end{conditionalproposition}
\begin{proof}
For each of the $104$ profiles in the preliminary input, let
$\mathcal D$ contain the dimensions $3\le d\le72$ for which the
integer equations of Lemma~\ref{paper:rank15-two-block-aggregate-rows}
have a solution. Every actual nonunit source dimension belongs to
$\mathcal D$.

For every constituent of representation multiplicity $m$, its forgetful
remainder has dimension $a-m$ and is a genuine positive combination of
nonunit source simples. Consequently each actual linear occurrence at
degree $a$ requires $a-1$ to belong to the additive semigroup generated
by $\mathcal D$, and each actual nonlinear occurrence requires $a-2$.
These conditions apply to every known constituent, including both
members of a possible quadratic nonlinear orbit. Applying all such
remainder tests gives $153\to126$ placements and $104\to86$ profiles.

Finally test the necessary source dimension equation
\[
                   5399=d_1^2+\cdots+d_{14}^2,
                   \qquad d_i\in\mathcal D,
\]
allowing repetitions. Every one of the $126$ placements passes this
last relaxation, so it supplies no further exclusion at this stage.

The standalone checker \path{rank15_aggregate_extension/aggregate_screen.py}
reads the authenticated $153$-placement input and compares two independently
arranged exact feasibility algorithms on all $104\cdot70=7280$ source-degree
tests. One solves the two equations by a balanced split and exact sum sets;
the other enumerates the norm equation first and then solves the signed
trace equation. Their eligibility sets agree on every test.
The remainder semigroups and the fourteen-square equation use finite
integer recurrences. The separate output
\path{rank15_aggregate_extension/aggregate_results.json} records every
eligibility set, required remainder and retained placement, together with
the unchanged preliminary-input hash.
The input consists of necessary character data, without a claim that a
source type, character table, multiplication or category realizes it.
\end{proof}

\subsection{Two degree-two components at source rank fifteen}
\label{paper:sec:two-matrix15}

\begin{conditionaltheorem}[No two-matrix source ring at rank fifteen]
\label{paper:rank15-two-matrix-excluded}
Assume verification hypothesis $\mathsf V$ of
Section~\ref{paper:verification-status} for the finite scopes and
prerequisite exclusions specified in the proof below.
Under the standing hypotheses at dimension $5400$, the source
Grothendieck algebra cannot have the shape
\[
             K_0(\mathcal C)\otimes_{\mathbb Z}\mathbb C
                 \cong\mathbb C^7\oplus M_2(\mathbb C)
                                  \oplus M_2(\mathbb C).
\]
Equivalently, all $126$ placements retained in
Proposition~\ref{paper:rank15-two-blocks-126} are excluded.
The proof includes the two individually rational matrix characters,
their possible quadratic Galois orbit when the canonical degrees
coincide, and all Schur indices compatible with a positive involution.
\end{conditionaltheorem}

Under the stated $\mathsf V$ hypothesis, the two certificate stages
are as follows. Exact character-field
and positive-form tests leave $484$ field/involution states on $116$
placements. Source rows, global moment identities, and positive word
traces exclude every state. In particular, the quadratic-orbit case
is not silently replaced by two rational characters.

\subsubsection{Finite positive forms without a splitting assumption}

\begin{lemma}[Parity and the two source-pair budgets]
\label{paper:two-matrix-parity-budgets}
Let $q$ be the number of nonreal conjugate pairs among the seven
linear characters, and let $b$ be the number of nonselfdual pairs
of source simples. Then
\[
                              q\le2,\qquad b\le6.
\]
For individually rational matrix characters, write $O$ and $S$ for
their orthogonal and symplectic involution types. Their contributions
to $b$ are respectively one and three. For a quadratic matrix orbit,
the contribution is two in the real orthogonal case, six in the real
symplectic case, and four in the imaginary quadratic case.
\end{lemma}
\begin{proof}
The canonical mass is $5400$. Its unit contributes one, every
degree-two representation contributes the even mass $2a$, and every
nonreal conjugate pair of linear characters contributes the even
mass $2a$. If $q=3$, these account for all the nonunit characters,
making the total mass odd. Thus $q\le2$. Equivalently there is a
nonunit real linear character with odd canonical degree.

The source dimension is also even. If the unit were the only
selfdual source label, its contribution one and the even
contributions $2d^2$ of all dual pairs would give an odd dimension.
At odd rank fifteen there must therefore be at least three
selfdual labels, so $b\le6$.

The last assertions are the dimensions of the negative eigenspaces
of the involution. For a rational degree-two component these have
dimension one or three. A real quadratic orbit doubles this
dimension. An imaginary quadratic center is conjugated by the
involution and the two eigenspaces each have rational dimension four.
The linear contribution is $q$ in every case.
\end{proof}

\begin{lemma}[Finite orthogonal skew squareclasses]
\label{paper:two-matrix-integral-skew-line}
For any individually rational orthogonal degree-two character of
formal codegree $f$, its positive skew-line squareclass has a
representative $\delta$ satisfying
\[
 \delta>0\text{ squarefree},\qquad \delta\mid f,
 \qquad f/\delta=z_1^2+\cdots+z_b^2\quad(z_i\in\mathbb Z).
\]
This assertion does not require that the component be the only
nonlinear component.
\end{lemma}
\begin{proof}
Choose one label from each source dual pair and put
$a_i=X_i-X_i^*$. For the rational degree-two character $\chi$, set
\[
        K_{ij}=-\frac12\chi(a_i a_j)
              =\chi(X_iX_j^*)-\chi(X_iX_j).
\]
The second equality uses cyclicity and
$\chi(x^*)=\chi(x)$. Consequently $K$ is integral. The positive
trace form identifies $K/f$ with the orthogonal projection of the
source skew space onto the skew line of this component. Thus
\[
                     K\succeq0,\qquad K^2=fK,
                     \operatorname{rank}K=1.
\]
An integral positive matrix of rank one is $c vv^{\mathsf t}$
with $v$ primitive integral and $c$ a positive integer: take a
primitive integral vector spanning its rational image and use
B\'ezout on its coordinates to see that the proportionality
coefficient is integral. Write $c=\delta h^2$, where $\delta$ is
squarefree, and put $z=hv$. Then
$K=\delta zz^{\mathsf t}$. Taking the trace gives the claimed
identity. The vector $z$ is integral but need not be primitive.
\end{proof}

The rational quaternion ramification theorem
(Theorem~\ref{paper:rational-matrix-ramification}) leaves a finite
list at each $f$. In the orthogonal case, write the algebra as
$(-\delta,c)$ with $c>0$. The whole positive form, its positive
subspace, and its negative subspace have diagonal representatives
\[
 \frac2f\langle1,\delta,c,\delta c\rangle,\qquad
 \frac2f\langle1,c,\delta c\rangle,\qquad
 \frac2f\langle\delta\rangle.
\]
Its finite ramification set belongs to
\[
                 \varnothing,\quad\{2,3\},\quad
                 \{2,5\},\quad\{3,5\};
\]
all its primes divide $f$ and are nonsplit in
$\mathbb Q(\sqrt{-\delta})$.
For every permitted set the verifier finds a positive integral $c$
having precisely that ramification set, and asserts coverage of
the entire finite list. The search for a representative is not an
assumption that all parameters satisfy a numerical size bound.
Quaternion classification and Skolem--Noether identify any actual
component with one of these representatives, with its specified
skew squareclass.

In the symplectic case the four possible finite ramification sets
and the corresponding unscaled whole positive forms are
\[
\begin{array}{c|c}
 \{2\}&\langle1,1,1,1\rangle\\
 \{3\}&\langle1,1,3,3\rangle\\
 \{5\}&\langle1,2,5,10\rangle\\
 \{2,3,5\}&\langle1,3,10,30\rangle.
\end{array}
\]
Multiply these by $2/f$, take their first coordinate as the positive
line and their other three coordinates as the negative subspace,
and retain only ramification primes dividing $f$.

Let $W,L_+,L_-$ be the whole, positive, and negative forms from the
seven linear characters, including the Perron contribution
$\langle1/5400\rangle$. For rational matrix forms $B_1,B_2$,
the verifier imposes
\[
 W\perp B_1\perp B_2\cong I_{15},\qquad
 L_-\perp B_{1,-}\perp B_{2,-}\cong2I_b.
\]
It also verifies the corresponding positive-space isometry
with $I_{15-2b}\perp2I_b$, as follows by cancellation.
All comparisons use dimension, determinant squareclass, and exact
local Hilbert symbols.

For equal matrix degrees the full matrix field can instead be
quadratic of fundamental discriminant $D$ with
$|D|\mid\gcd(2f,f^2)$. No quaternion presentation or Schur-index
restriction over this field is imposed. The following weaker
necessary transfer conditions suffice. Put $d=D$ if $D$ is odd,
and $d=D/4$ otherwise.
\begin{itemize}
\item For a real symplectic orbit, the positive matrix subspace is
the center, with form $\langle4/f,4d/f\rangle$.
\item For an imaginary quadratic orbit, multiplication by
$\sqrt{-|d|}$ identifies its negative form with $|d|$ times its
positive form. Eliminating the unknown positive form gives
\[
 L_+\perp\frac1{|d|}(2I_b)
       \cong (I_{15-2b}\perp2I_b)\perp\frac1{|d|}L_-.
\]
\item For a real orthogonal orbit, the two-dimensional negative
form is the transfer of a positive one-dimensional form over
$\mathbb Q(\sqrt d)$. Hence its determinant divided by $d$
is a norm from this quadratic field. Its required determinant
squareclass is $2^b/\det L_-$, so the verifier tests
\[
       (d,\,2^b/(d\det L_-))_p=1
\]
at every relevant prime. No sufficiency assertion about these
quadratic-orbit tests is needed.
\end{itemize}

The full linear-field catalogue is the one already proved complete
in Proposition~\ref{paper:rank15-two-blocks-preliminary}. Each
degree block is partitioned into all permitted full-field orbits.
The determinant contribution of both matrix components is square
by Theorem~\ref{paper:source:even-degree-trace-determinant}, whether they
are rational or form a quadratic orbit. The rational canonical
residue obstruction is applied to every rational constituent.
The universal codegree-six support results are used only for
individually rational symplectic components: their coexistence
with a nonreal linear codegree-six pair requires rank at least
sixteen, and two such rational symplectic components require rank
at least eighteen. Neither result is applied to a quadratic matrix
orbit.

\subsubsection{Exact source rows for every remaining matrix field}

Write each remaining Galois-stable piece as $(a_i,k_i,j_i)$, where
$a_i=5400/f_i$, $k_i$ is the sum of its representation dimensions,
and $j_i$ is its number of degree-two representations. Thus
$j_i=0$ for a linear piece, $(k_i,j_i)=(2,1)$ for one rational
matrix character, and $(k_i,j_i)=(4,2)$ for a quadratic matrix orbit.
Some rational linear characters at a common degree may be combined;
this is a necessary relaxation, not a claimed common full field.
For a nonunit source simple of dimension $3\le d\le72$, set
\[
     s_i=\sum_{E\text{ in piece }i}\chi_E(X),\qquad
     t_i=\sum_{E\text{ in piece }i}\chi_E(XX^*).
\]
Every source row satisfies the exact equations
\[
                  \sum_i a_i s_i=-d,\qquad
                  \sum_i a_i t_i=5400-d^2.
\]
It also satisfies the stated strict spectral, field-integrality,
codegree-six support, and positive norm bounds. In particular
\[
 0\le t_i\le\min\{f_i-1,k_i d^2-1,
                   \lfloor(5400-d^2)/a_i\rfloor\},\qquad
 |s_i|<k_i d.
\]
The scalar entries are enumerated in the full rings of integers
of their permitted fields. An entry lying in a proper subfield is
retained, with the appropriate repeated conjugates.

For a rational matrix component put
\[
                    s=\operatorname{Tr}A,\qquad
                    t=\operatorname{Tr}(AA^*),\qquad
                    n=\det A.
\]
All three are integers. The determinant choices obey
$2|n|\le t$. An orthogonal component also obeys
$t+2n-s^2\ge0$; a symplectic component obeys $2n=t$.
For a selfdual row, $2n=s^2-t$. These statements concern reduced
traces and determinants and remain valid when the rational
quaternion algebra is nonsplit.

For a quadratic matrix orbit set
\[
 \alpha=\operatorname{Tr}A\in\mathcal O_K,\qquad
 \beta=\det A\in\mathcal O_K,\qquad
 \gamma=\operatorname{Tr}(AA^*)\in\mathcal O_K.
\]
Here $s=\operatorname{Tr}_{K/\mathbb Q}\alpha$ and
$t=\operatorname{Tr}_{K/\mathbb Q}\gamma<f$.
When $K$ is real, $\gamma$ is totally nonnegative and
\[
 2\gamma-\alpha^2\ge0,\qquad \gamma\pm2\beta\ge0
\]
at both embeddings. In the orthogonal case also
$\gamma+2\beta-\alpha^2\ge0$; in the symplectic case
$\gamma=2\beta$. A selfdual row has
$2\beta=\alpha^2-\gamma$, and in the real symplectic case
also $\alpha/2\in\mathcal O_K$ and
$\beta=(\alpha/2)^2$.
When $K$ is imaginary, $\gamma$ is a nonnegative rational integer,
\[
            2\gamma\ge|\alpha|^2,\qquad
            4|\beta|^2\le\gamma^2,
\]
and selfduality makes both $\alpha$ and $\beta$ rational.
The necessary nonnegative difference of two word traces is
\[
 G=\gamma^2-\gamma|\alpha|^2
          +2\operatorname{Re}(\alpha^2\overline\beta)
          -4|\beta|^2\ge0.
\]

All these variables are enumerated exactly in the integral basis
$1,(1+\sqrt D)/2$ when $D$ is odd and $1,\sqrt{D/4}$ otherwise.
Equivalently write $\alpha=(s+v\sqrt D)/2$ with the appropriate
integral parity condition. With row weight $w=1$ for a selfdual
label and $w=2$ for a dual-pair representative, character
orthogonality gives
\[
       ws^2\le4f-16,\qquad
       w\operatorname{Tr}_{K/\mathbb Q}(\alpha\overline\alpha)
                       \le4f-8.
\]
The conjugation in the latter expression is ordinary complex
conjugation, hence trivial for real $K$. In particular each
enumeration is finite. The determinant search is bounded by
\[
 \operatorname{Tr}_{K/\mathbb Q}(\beta^2)\le t^2/4
                  \quad(K\text{ real}),\qquad
 2|\beta|^2\le t^2/8\quad(K\text{ imaginary}).
\]
These are necessary bounds; the enumerator is allowed to retain
triples that have no actual representation matrix.

\subsubsection{The two-matrix moment and word conditions}
The common moment targets in Lemma~\ref{paper:source:row-moments}
apply with \(14-2b\) selfdual labels and \(b\) pair atoms.
A quadratic nonlinear orbit has the additional target
\[
 \sum_Xv(X)^2=4f/|D|.
\]
For a real quadratic orbit, the targets of \(d_Xv(X)\) and
\(s_i(X)v(X)\) are zero.  They are not imposed in the imaginary
case, where changing orientation changes the sign of \(v\).
These identities follow by adding and subtracting the orthogonality
relations for the two conjugate degree-two characters.

Theorem~\ref{paper:trace-norm-defect} gives \(|s_i|\le t_i\), with
strict inequality for at least one non-Perron piece.  With the
recurrences of Proposition~\ref{paper:matrix-power-word-traces}, test
\(X^k\) for \(0\le k\le15\), \((XX^*)^k\) for
\(1\le k\le15\), and \(X^k(X^*)^\ell\) for
\(0\le\ell\le k\le10\).  Their coefficient traces must be
nonnegative integers, with one consistent value and determinant choice
throughout each test.  A positive return is required by exponent
\(L=11\).  Pure powers of higher-field scalar entries use the full
field characteristic polynomial; a mixed or norm trace not determined
by the stored conjugate data is omitted conservatively.

The coupled commutator condition is
\begin{equation}\label{paper:two-matrix-coupled-commutator}
 5400\mid\sum_E a_EG(E(X)),
\end{equation}
where \(G\) is defined in Lemma~\ref{paper:matrix-commutator-traces}
and the sum includes both components, or their full field trace.
For a rational component it is
\(G=(t-2\det A)(t+2\det A-s^2)\).
No separate divisibility by the individual codegree is imposed.

\subsubsection{Completeness and the final two-matrix obstructions}
The \(126\) placements give \(777\) full linear-field assignments.
Parity, determinant and matrix-form tests retain \(484\) matrix
states on \(116\) placements.  Their complete nonlinear alternatives
are rational types \(OO,OS,SO\), imaginary quadratic discriminants
\(-4,-3\), and real orthogonal discriminants \(12,8\); no real
symplectic state passes the form tests.  Interchanging equal canonical
matrix degrees does not duplicate a state.

The universal symplectic-six theorem excludes \(54\) states with
fewer than five source pairs.  The others have exact remainder,
dimension-square, or moment certificates from
Lemma~\ref{paper:source:certificate-principle}.  After the stated
word refinement, \(36\) terminal states have a missing required
remainder.  The last two, with matrix canonical degrees
\((900,54)\) and \((900,135)\), force two source pairs of dimension
\(27\).  Their selfdual dimension sets are respectively
\[
 \{8,25,26,27,33,48\},\qquad\{25,26,48,54\}.
\]
Neither supplies ten squares with sum
\(5399-4\cdot27^2=2483\).  For the second set,
\(10\cdot25^2>2483\).  For the first, reduction modulo three
requires the number of dimensions from \(\{27,33,48\}\) to be
two modulo three; the budget permits only two.  Any of the other
eight dimensions being at least \(25\) would give a sum at least
\(2\cdot27^2+7\cdot8^2+25^2=2\,531\), so all eight are \(8\).
The last two cannot include \(48\), and their odd squares cannot
sum to \(2483-8\cdot8^2=1971\).

The exact checker \path{verify_two_matrix15.py} regenerates the full
field/form input and every necessary row used by a certificate, and
checks coverage of all \(484\) states.  The ten placements with no
admissible form are included in the input coverage.  Together these
certificates prove Theorem~\ref{paper:rank15-two-matrix-excluded}.

\subsection{Exclusion of canonical length eleven at every rank}
\label{paper:sec:canonical-length12}
\begin{conditionaltheorem}[Exclusion of canonical length eleven]
\label{paper:canonical-length-at-least-twelve}
Assume verification hypothesis $\mathsf V$ of
Section~\ref{paper:verification-status} for the finite scopes and
prerequisite exclusions specified in the proof below.
Under the center hypothesis, \(L\ge12\).  This includes all
representation multiplicities and nontrivial Galois orbits of odd
nonlinear degree.
\end{conditionaltheorem}
\begin{proof}
Theorem~\ref{paper:canonical-length-at-least-eleven} gives \(L\ge11\).
At \(L=11\), partitioning the repeated entries of all \(497\)
weighted profiles gives \(4\,475\) assignments, with ranks computed
from \(r=1+\sum m^2\) and no imposed upper cutoff.
The odd-degree determinant baseline is
\(\operatorname{sf}(5400\prod_a a^{k_a})\), by
Theorem~\ref{paper:source:even-degree-trace-determinant}.  A singleton
\((a,m)\) group is rational; odd \(m\) alone does not imply this.
Every repeated group retains all conductor-eligible field orbits.
For groups of size at least four, the conductor is in
\(\{9,12,16,20,24\}\), so enumeration of every unit-group subgroup
is complete.  Smaller groups use the full quadratic and real cubic
conductor-nine lists.

The determinant, rational degree-ten, proved rank, degree-\(24,200\)
and rational-residue conditions leave \(589\) assignments.  Exactly
\(153\) are the rank-fifteen two-matrix placements of
Proposition~\ref{paper:rank15-two-blocks-preliminary}; their identifiers
are matched entry by entry.  The \(27\) aggregate-remainder cases
of Proposition~\ref{paper:rank15-two-blocks-126} and the complete
\(126\)-placement calculation of
Theorem~\ref{paper:rank15-two-matrix-excluded} close this part.

The other \(436\) assignments have ranks seventeen through fifty-three.
Use Galois-stable groups \((a,m,c)\), the common rows and moments,
the jointly solved skew coordinates of \eqref{paper:aggregate-skew-row},
and the odd-rank pair budget of
Lemma~\ref{paper:source:aggregate-duality-budget}.  For scalar
groups, the codegree-six lists include both the rational and Eisenstein
possibilities, with the common real label and opposite signs permitted
for two nonreal pairs.  The commutative disjointness clause is not
imposed on nonlinear sources.  Outside codegree six, every scalar
block of size at least four has codegree eight, and all three quartic
subfields of \(\mathbb Q(\zeta_{16})\) are included.  Smaller blocks
use the complete quadratic and cubic lists.  This accounts for every
full scalar field at the actual multiplicities.

Two assignments fail Proposition~\ref{paper:integral-norm-rank-bound}:
a degree-seven character of codegree eight requires rank at least
\(56\), whereas its assignment has rank \(53\); the analogous
degree-six character requires \(42\), whereas its rank is \(41\).
Canonical remainders, exact dimension-square recurrences and integer
moment certificates close every other assignment except twelve.
For these, enumerate the complete integral polynomial superset of
Lemma~\ref{paper:aggregate-polynomial-superset}.  One consistent
polynomial choice must give nonnegative integral coefficient traces
through power fifteen, the correct square trace, a positive return
by \(L=11\), and \(\operatorname{Tr}L_{X^j}\ge r\tau(X^j)\).
The regular traces of \(X\) and \(XX^*\) are also imposed.
This uses no splitting assumption.  In each of the twelve terminal
cases a required positive canonical remainder is smaller than every
eligible source dimension, so it cannot be realized.

The exact checker \path{verify_canonical11.py} regenerates all
\(4\,475\) assignments, every initial field and residue decision,
the rank-fifteen identifier match, and all higher-rank rows and
certificates.  The complete report is
\path{canonical11_verification.json}.  Its coverage and the common
certificate principle exclude every length-eleven assignment.
\end{proof}

\subsection{The complete rank-thirteen commutative field input}
\label{paper:comm13-field-input-section}

At commutative rank thirteen the full field catalogue is larger
than at rank twelve: a
codegree-$10$ block of multiplicity eight can carry the full field
$\mathbb Q(\zeta_{20})$.

\begin{conditionalproposition}[Exhaustive canonical and full-field input at rank thirteen]
\label{paper:comm13-field-input}
Assume verification hypothesis $\mathsf V$ of
Section~\ref{paper:verification-status} for the finite scopes and
prerequisite exclusions specified in the proof below.
Under the standing center hypothesis, suppose the source Grothendieck
ring is commutative of rank thirteen. Its twelve nonunit canonical
degrees form one of exactly $62\,776$ nonincreasing profiles
\[
 (a_1,\ldots,a_{12}),\qquad \sum_i a_i=5399,
 \qquad a_i\in\mathcal A,
\]
where $\mathcal A$ is the stated divisor palette: the divisors of
$5400$ between $10$ and $900$ having at least two distinct prime factors.
The following exact preliminary reductions hold:
\[
\begin{array}{l|r|r}
 \text{test}&\text{profiles removed}&\text{profiles retained}\\\hline
 \text{initial canonical profile list}&0&62\,776\\
 \text{unique rational degree }10&4\,714&58\,062\\
 \text{two-character canonical modulus}&17&58\,045\\
 \text{full-field trace forms and duality}&45\,058&12\,987.
\end{array}
\]
On the $58\,045$ profiles entering field enumeration there are
$3\,571\,554$ full-field decompositions. After the support, duality and
whole-form determinant restrictions, $376\,486$ decompositions remain
for full rational-isometry checks. The whole form and both involution
eigenspaces retain exactly $40\,476$ full-field assignments on the
$12\,987$ profiles. Every later exclusion must address these assignments
or give a certificate excluding their whole profile.
\end{conditionalproposition}

\begin{proof}
The profile recursion chooses each next degree no larger than its
predecessor, retains the exact remaining sum and number of entries,
and stops only when both are zero. Its lower and upper sum bounds
prune impossible branches without changing the solution set. Thus
there is neither ordering duplication nor omission. The rational
canonical-degree-$10$ obstruction and the general two-character
modulus apply before any field choices are made.

For a formal-codegree block $(f,k)$, the stated conductor bound is
$c=\gcd(2f,f^2)$ and the full character field is an abelian subfield
of $\mathbb Q(\zeta_c)$ of degree at most $k$. The maximum actual
multiplicity is eight. The possible orbit degrees in this profile
universe are
\[
                         1,2,3,4,5,6,8.
\]
The quadratic fields are enumerated by their fundamental
 discriminants. The only cubic field is the real cubic field of
conductor nine. Every quartic field has conductor dividing
$\gcd(c,240)$: an abelian quotient of order four kills the wild
$3$- and $5$-parts and all $2$-primary characters of conductor greater
than sixteen. All index-four subgroups of the reduced unit group are
enumerated through two successive index-two kernels. The eligible
quartic ambient conductors include $60,80,120,240$.

For the remaining orbit degrees, the actual multiplicity/conductor
inventory leaves only the real cyclic quintic of conductor $25$;
the sextic fields in $\mathbb Q(\zeta_9)$ or
$\mathbb Q(\zeta_{36})$; and the full octic fields
$\mathbb Q(\zeta_{20})$ and $\mathbb Q(\zeta_{24})$.
At codegree six, the stated support theorem permits only
$\mathbb Q$ and $\mathbb Q(\sqrt{-3})$. At codegree $540$, the
rational field is removed by the degree-$10$ obstruction.
An independent enumeration by adjoining subgroup generators in the
\emph{full} conductor groups agrees with the proposed list on all
$55$ actual blocks of multiplicity at least four. The lower-degree
quadratic and cubic lists are explicit. Thus this catalogue includes
full fields, not just the fields generated by individual entries.

Let $q_6$ count the nonreal codegree-six pairs and $q$ all nonreal
character pairs. The canonical mass and disjoint-support restrictions
give $q_6\le2$ and $q\ge2q_6$. Also $q\le5$: otherwise all twelve
nonunit source simples would form six dual pairs, making the source
dimension odd. For each simultaneous full-field assignment, the
positive rational trace form and its two involution eigenspaces must
be isometric to
\[
 I_{13},\qquad I_{13-2q}\mathbin\perp2I_q,\qquad2I_q,
\]
respectively. Their determinants and finite local Hilbert invariants
are evaluated with exact rational arithmetic. The stated counts are
those of the complete recursion and these exact form tests.
\end{proof}

\subsection{Exclusion of commutative source rank thirteen}
\label{paper:comm13-completion-section}
\begin{conditionaltheorem}[Commutative source rank thirteen is excluded]
\label{paper:commutative-rank-thirteen-excluded}
Assume verification hypothesis $\mathsf V$ of
Section~\ref{paper:verification-status} for the finite scopes and
prerequisite exclusions specified in the proof below.
Under the center hypothesis, a commutative source has rank at least
fourteen.
\end{conditionaltheorem}
\begin{proof}
Proposition~\ref{paper:comm13-field-input} leaves \(40\,476\) complete
full-field assignments from \(62\,776\) weighted profiles.  Enumerate
the common integral rows for \(3\le d\le72\), using the full
cyclotomic fixed lattices and retaining repeated characteristic roots
for entries in proper subfields.  Apply the coefficient-power and
positive-return tests through thirteen and the regular-trace bounds.
Bounded canonical remainders search every dimension that can contribute
an eligible atom; the support threshold is \(\lambda(X)\ge m\),
with paired coins doubled for selfdual canonical constituents.
Rows of regular character zero remain eligible for the separate source
dimension equation.

Every assignment has an exact certificate: integral coordinate
covariance, a canonical remainder, a dimension-square recurrence,
an integer moment or support-cover separator, or the character-frame
obstruction.  The last is the explicit coloring proof of
Proposition~\ref{paper:rational-8370-frame}; it excludes twelve mutually
orthogonal nonunit character rows.  The other witnesses follow from
the coordinate and support theorems and
Lemma~\ref{paper:source:certificate-principle}.
The complete ledger \path{rank13_certificates.json} and report
\path{rank13_complete_verification.json} regenerate the initial
universe, all full-field trace-form decisions, and every necessary
row or bounded dimension range used by a certificate.  All \(40\,476\)
records are excluded.  The already proved lower endpoints then give
the asserted rank bound.
\end{proof}

\subsection{The one-matrix shape at rank fifteen}
\label{paper:sec:one-matrix-rank15}
Here \(R_{\mathbb C}\simeq\mathbb C^{11}\oplus M_2(\mathbb C)\)
and \(L=13\).  The unique matrix character is rational.  The
complete canonical recursion gives \(62\,776\) weighted profiles and
\(143\,657\) placements, and the full field/form enumeration leaves
\(125\,545\) field/duality states.  The field catalogue includes the
real quintic field of conductor \(25\), the sextic subfields of
conductor \(36\), and the sextic field of conductor \(9\).
No degree-eight linear orbit fits: every repeated linear block has
size at most seven.  Thus higher fields are retained whenever the
actual multiplicities permit them.

\begin{proposition}[Parity bounds at odd source rank]
\label{paper:rank15-one-matrix-parity}
Let $q$ count the nonreal pairs among the eleven linear characters.
At source rank fifteen,
\[
 q\le4\quad\hbox{for orthogonal matrix duality},\qquad
 q\le3\quad\hbox{for symplectic matrix duality}.
\]
The rank-independent orthogonal codegree-six catalogue of
Corollary~\ref{paper:orthogonal-six-finite-forms} applies at this rank.  For a symplectic codegree-six matrix block one has
$q\in\{2,3\}$ and $q_6=0$.
\end{proposition}

\begin{proof}
The Perron character is a real linear character.  If $q=5$, all ten
remaining linear characters occur in conjugate pairs with equal
canonical degrees.  The canonical mass equation would then read
$5400=1+2a_E+2\sum a$, an impossibility.  Thus $q\le4$, independently of
the matrix involution.

Let $b$ denote the number of source dual pairs.  At odd source rank the
number of self-dual labels is odd.  It cannot be one: if the unit were
the only self-dual label, the source dimension would be
$1+2\sum d_i^2$, which is odd.  Hence at least three labels are
self-dual and $b\le6$.  The trace of duality gives $b=q+1$ in the
orthogonal case and $b=q+3$ in the symplectic case.  This proves the
second bound.

In particular, orthogonal duality has $b\le5$.  The full
orthogonal-six classification is rank-independent, so its endpoint
corollary applies without an additional rank assumption.  The universal symplectic
support of Theorem~\ref{paper:universal-symplectic-six-support} requires
five source pairs, hence $q\ge2$.  Finally
Corollary~\ref{paper:symplectic-six-linear-six-threshold} excludes a
nonreal linear codegree-six pair below rank sixteen.
\end{proof}

The field/form test combines these parity bounds with
Proposition~\ref{paper:one-matrix-rank-independent-restrictions},
the complete quaternion classes, and the skew-lattice condition
\(\delta\mid f_E\), with \(f_E/\delta\) a sum of at most
\(q+1\) integer squares in the orthogonal case.  The imaginary
quadratic, real quadratic norm-one and norm-two, and integral
coordinate-covariance obstructions leave \(82\,964\) states, whose
largest linear-field degrees are one through four.  The sextic
fields were included and excluded by these proved restrictions.

\begin{lemma}[A bounded certificate for a canonical remainder]
\label{paper:bounded-one-matrix-remainder}
Fix a field--duality state and a canonical constituent of degree $a$
and multiplicity $m\in\{1,2\}$.  Put $R=a-m$.  Suppose $D_R$ is a set
containing every possible source-simple dimension $d\le\min(R,72)$
that can occur in that state.  If $R$ is not a nonnegative integer
combination of elements of $D_R$, that state is impossible.
\end{lemma}

\begin{proof}
The forgetful image of the constituent contains $m$ copies of the
unit.  The remaining positive object has dimension $R$, so every
nonunit constituent has dimension at most $R$.  At rank fifteen all
nonunit source dimensions are between three and seventy-two, by the
previous source-dimension bound.  Its dimension must therefore be a
nonnegative integer combination of dimensions in $D_R$.
\end{proof}

\begin{lemma}[Dimension certificates with source duality]
\label{paper:one-matrix-density-certificates}
In a rank-fifteen field--duality state, put
\[
 b=q+1\quad\text{or}\quad b=q+3
\]
according as the matrix involution is orthogonal or symplectic, and put
$s=14-2b$.  There are exactly $s$ nonunit self-dual labels and $b$
nonselfdual pairs.  If their eligible dimensions have lower bounds
$u$ and $v$, respectively, the state is impossible whenever
\[
                         su^2+2bv^2>5399.
\]
More precisely, let $u_o,u_e$ bound the odd and even self-dual
dimensions, respectively.  The state is impossible whenever
\[
 \min_{\substack{1\le j\le s-1\\j\text{ odd}}}
       \bigl(ju_o^2+(s-j)u_e^2+2bv^2\bigr)>5399.
\]
One may instead retain eligible sets $D_s,D_p$ and test exactly whether
$5399$ is a sum of $s$ squares from $D_s$ and twice $b$ squares from
$D_p$.  Absence of such a sum excludes the state.  Finally, absence of
any eligible row with $3\le d\le19$ already excludes a rank-fifteen
state, since $14\cdot20^2>5399$.
\end{lemma}

\begin{proof}
The trace of duality gives the stated counts.  Removing the unit from
the source dimension identity leaves
\[
 \sum_{i=1}^{s}d_i^2+2\sum_{j=1}^{b}e_j^2=5399.
\]
The first bound and the exact finite sum test follow immediately.
Here $s$ is positive and even.  Reducing the identity modulo two shows
that an odd number of the $d_i$ are odd.  This number is between one
and $s-1$, proving the stronger bound and ensuring that both parities
of self-dual dimension occur.  The last assertion follows from the
same dimension identity if all fourteen nonunit labels had dimension
at least twenty.  These tests impose no regular-character threshold:
rows of regular character zero remain eligible for the source basis.
\end{proof}

For moment certificates use Lemma~\ref{paper:source:row-moments}.
In particular the matrix norm total is \(4f_E-2\), while its
character-square total is \(2f_E-4\).  Selfdual rows contribute
once and pair rows twice, with separate atom-count coordinates.
The support-cover inequalities of Lemma~\ref{paper:canonical-support-cover}
may be appended using their nonnegative slack coordinates.
The quaternion row tests and central-atom theorem in
Section~\ref{paper:source:matrix-row-domains} apply without a
splitting assumption.

\begin{conditionaltheorem}[The noncommutative rank-fifteen endpoint is excluded]
\label{paper:noncommutative-rank15-exclusion}
Assume verification hypothesis $\mathsf V$ of
Section~\ref{paper:verification-status} for the finite scopes and
prerequisite exclusions specified in the proof below.
Under the center hypothesis,
\(R_{\mathbb C}\simeq\mathbb C^{11}\oplus M_2(\mathbb C)\)
is impossible.  Consequently a noncommutative source has rank at
least sixteen.
\end{conditionaltheorem}
\begin{proof}
Every one of the \(125\,545\) full field/duality states has an exact
certificate: a quadratic or coordinate-covariance obstruction, a
bounded or duality-sensitive canonical remainder, a dimension-density
or dimension-square obstruction, a forced central atom, or an integer
moment separator.  The common certificate theorem applies to the
complete necessary row sets, with \(3\le d\le72\) and every
allowed quaternion class.  Every claimed absent dimension or atom is
exhaustively regenerated in its stated range.  A row surviving a
weaker filter is retained as a necessary possibility; absence is never
inferred from a partial search.

The code/certificate group \path{verification/source/one_matrix15}
contains the complete input, ledger and replay.  Its independent trace
replayer regenerates all \(143\,657\) placements, including those
excluded before the final state list; the row merger requires exactly
one authenticated certificate for each of the \(125\,545\) states.
The supplied \path{replay_all.py} verifies the complete exclusion by
integer and rational arithmetic.  This excludes the one-matrix shape.

By Theorem~\ref{paper:noncommutative-rank14-exclusion}, the remaining
rank is at least fifteen.  At rank fifteen, \(L\ge11\) gives
\(0<\sum_E m_E(m_E-1)\le4\), so there are one or two
degree-two components.  The two-component case is excluded by
Theorem~\ref{paper:rank15-two-matrix-excluded}, including a possible
quadratic Galois orbit.  Rank fifteen is therefore impossible.
\end{proof}

\subsection{The two-matrix shape at rank sixteen}
\label{paper:sec:two-matrix16}
For \(\mathbb C^8\oplus M_2(\mathbb C)^{\oplus2}\), the
canonical length is twelve.  The common recursion gives \(12\,257\)
placements; an independent generating function agrees for each
unordered pair of nonlinear canonical degrees, as well as in total.

\begin{lemma}[The rank-sixteen duality budgets]
\label{paper:two-matrix16-duality-budgets}
Let $q$ count nonreal conjugate pairs of linear characters and let
$b$ count nonselfdual pairs of source labels.  Then $q\le3$ and
$b\le7$.  For two rational matrix characters,
\[
 b=q+2,\quad q+4,\quad q+6
\]
in types $OO$, $OS$ and $SS$, respectively.  A quadratic nonlinear
orbit instead contributes two, six or four to $b$, according as
it is real orthogonal, real symplectic or imaginary quadratic.
\end{lemma}
\begin{proof}
The Perron character is rational, leaving seven nonunit linear
characters, hence at most three nonreal pairs.  The source unit is
selfdual, so at even rank sixteen there are at most seven dual
pairs.  The remaining statements are the negative-eigenspace
dimensions of the positive involution, as in
Lemma~\ref{paper:two-matrix-parity-budgets}.  In particular, the
rank-fifteen bounds $q\le2$ and $b\le6$ are not used here.
\end{proof}

Use the complete rational quaternion lists and quadratic transfer
forms of the rank-fifteen two-matrix argument, with targets
\(I_{16}\), \(2I_b\), and \(I_{16-2b}\perp2I_b\).
Every orthogonal skew squareclass satisfies \(\delta\mid f\)
and the integral \(b\)-square representation condition.  Quadratic
nonlinear components retain every Schur index allowed by these
necessary forms.  Their codegree-six cases are excluded by
Corollary~\ref{paper:quadratic-six-orbit-exclusion}, whose all-nonlinear
orbit hypothesis holds here.  Rational symplectic-six components use
the universal five-pair and two-component rank-eighteen bounds.
Coexistence with a nonreal linear codegree-six pair is permitted at
rank sixteen; the stronger rank-fifteen exclusion is not imposed.

Linear multiplicities are at most five.  Multiplicity five occurs
only at codegrees \(6,8,10\), whose conductors have no degree-five
quotient.  Hence only degrees one through four occur.  Enumerating
all subgroups of the full conductor unit groups for the \(42\)
applicable blocks of multiplicity at least three independently
checks this catalogue.  Coordinate-covariance restrictions are applied
to the linear orbits, and the canonical residue and degree-\(24,200\)
conditions only with their stated hypotheses.

The common rows have \(3\le d\le72\), or \(3\le d\le51\)
for pair representatives, with \(15-2b\) selfdual atoms and \(b\)
pair atoms.  Their regular character satisfies
\(0\le\lambda(X)\le15(d_X-1)\), by
Lemma~\ref{paper:regular-trace-upper-bound}; rows of regular character
zero remain eligible for the source basis.  Canonical support and
central atoms use Lemma~\ref{paper:canonical-support-cover} and
Theorem~\ref{paper:central-canonical-atom}, including the quadratic
scalarity formulas of Section~\ref{paper:source:matrix-row-domains}.

Pure and norm powers are tested through sixteen, mixed words through
\(0\le\ell\le k\le10\), and positive return by \(L=12\).
In addition to nonnegative integral coefficient traces, every tested
positive word satisfies
\[
 \lambda(W)\ge16\tau(W).
\]
This follows by expanding \(W\) in the positive source basis: the
unit contributes \(16\tau(W)\) and every other regular-character
contribution is nonnegative.  Word traces not determined by stored
higher imaginary-field data are omitted conservatively.

One forced-atom example illustrates a restriction beyond a remainder
semigroup.  For the profile
\((900^5,675,180,12^2,10^2)\), put orthogonal matrix constituents
at \(12\) and \(900\), with linear fields
\((6,\mathbb Q)^3,(8,\mathbb Q),(30,\mathbb Q),
(540,\mathbb Q(\sqrt{-15}))\), where the first coordinate is the
codegree.  The degree-twelve remainder is ten.  Its eligible selfdual
dimensions are eight and nine, so it must be a single dimension-five
pair.  Every eligible pair row has matrix norm \(t=45\) in that
component, whereas centrality requires \(t=2\det A\), an even
integer.  This is an exact central-atom certificate.

\begin{conditionaltheorem}[Exclusion of the rank-sixteen two-matrix shape]
\label{paper:rank16-two-matrix-exclusion}
Assume verification hypothesis $\mathsf V$ of
Section~\ref{paper:verification-status} for the finite scopes and
prerequisite exclusions specified in the proof below.
Under the center hypothesis,
\(R_{\mathbb C}\simeq\mathbb C^8\oplus M_2(\mathbb C)^{\oplus2}\)
is impossible.  Both rational matrix characters and a quadratic
nonlinear orbit are covered, with all compatible Schur indices.
\end{conditionaltheorem}
\begin{proof}
The full field and form enumeration over the \(12\,257\) placements
retains \(9\,268\) states.  The retained nonlinear alternatives are
rational types \(OO,OS,SO,SS\), imaginary quadratic discriminants
\(-24,-20,-8,-4,-3\), real orthogonal discriminants \(5,8,12,24\),
and real symplectic discriminant \(8\).  The input includes states
with \(q=3\), with \(b=7\), and with rational symplectic-six and
nonreal linear-six coexistence.

Every state has a canonical degree-\(24,200\), symplectic-support,
remainder, central-atom, dimension-square or integer moment certificate.
The last four kinds regenerate the complete necessary rows and any
word refinement used in their certificate.  Their soundness is exactly
Lemma~\ref{paper:source:certificate-principle} and the scalarity
theorem, so every state is excluded.  The two-matrix-rank-sixteen
certificate group authenticates its canonical input, full trace-form
input and all state witnesses; the strict replay merger requires all
\(9\,268\) identifiers exactly once and no unresolved state.
\end{proof}

\subsection{Canonical length twelve at every source rank}
\label{paper:sec:canonical12-enumeration}
\begin{conditionalproposition}[The complete length-twelve universe]
\label{paper:canonical12-complete-universe}
Assume verification hypothesis $\mathsf V$ of
Section~\ref{paper:verification-status} for the finite scopes and
prerequisite exclusions specified in the proof below.
At \(L=12\), the \(6\,547\) weighted profiles give \(69\,312\)
representation-multiplicity assignments, of ranks twelve through
fifty-six.  The stated necessary tests retain \(12\,905\):
\(3\,526\) at rank sixteen and \(9\,379\) at higher rank.
No upper source-rank cutoff is imposed in this enumeration.
\end{conditionalproposition}
\begin{proof}
Partition the multiplicity of each weighted degree in the common
canonical recursion.  An independent count-vector enumeration gives
the same complete assignment set.  Its rank is even.  The determinant
baseline is now \(\operatorname{sf}(\prod_a a^{k_a})\), with no
extra factor \(5400\), since \(L\) is even.  Full fields are
enumerated at each \((a,m)\) multiplicity, and rationality is
forced only for singleton groups.

Every quadratic field has fundamental discriminant dividing
\(c=\gcd(2f,f^2)\); the cubic field is the real conductor-nine
field.  A quartic quotient kills odd wild inertia and has conductor
dividing \(\gcd(c,240)\), so all index-four subgroups there are
enumerated.  The actual repeated groups have size at most seven,
with no eligible degree-five or degree-seven orbit; every eligible
sextic orbit has conductor nine and is included.  Thus nonlinear
odd-degree Galois orbits are covered as well.

The determinant, rational degree-ten, proved rank, degree-\(24,200\),
rational-residue and aggregate norm-capacity conditions leave the
stated \(12\,905\) assignments.  The rank lower bounds used are
fourteen for a commutative source and sixteen otherwise.  The capacity
test is \((r-1)(f-1)\ge fm^2c_{a,m}-mc_{a,m}\), from
Proposition~\ref{paper:integral-norm-rank-bound}.
At rank sixteen the defect is four, so there are exactly two
degree-two components.  The \(3\,526\) retained inputs are matched
entry by entry to the complete \(12\,257\)-placement universe of
Theorem~\ref{paper:rank16-two-matrix-exclusion}.
\end{proof}

\paragraph{Complete fields and rows in the higher-rank part.}
The \(9\,379\) assignments have \(99\) scalar-block parameters,
with multiplicity at most five.  Size five occurs only at codegrees
\(6,8,9,10\), which permit no quintic field.  Every eligible
quadratic and cubic field and all index-four subfields are retained;
independent full-group subgroup closure checks all \(99\) catalogues.
The codegree-six lists use the rank-independent field and support
theorems, including the allowed common real label for two Eisenstein
orbits.  No scalar quintic or sextic is omitted.

Among the \(3\,239\) assignments with only odd nonlinear degrees,
\(3\,209\) have singleton nonlinear groups and hence rational
characters.  Only these are passed to
Lemma~\ref{paper:odd-rational-skew-determinant}; the other \(30\)
keep their nontrivial Galois orbits in aggregate rows.
Use the common moments and the even-rank pair budget of
Lemma~\ref{paper:source:aggregate-duality-budget}, with all
rank-dependent source dimensions and jointly solved skew coordinates.
The regular traces satisfy
\[
 0\le d+\sum_gm_gs_g\le(r-1)(d-1),\qquad
 d^2+\sum_gm_gt_g\ge r.
\]
The first is Lemma~\ref{paper:regular-trace-upper-bound}; the second
follows because \(XX^*\) contains the unit once and the remaining
regular characters are nonnegative.  Apply also the canonical-capacity
knapsack, cubic residue, and subset-capacity inequalities, with their
slack coordinates.  A singleton canonical group is rational and
selfdual, so its nonselfdual source atoms have dimension \(2d\).

When needed, the complete aggregate-polynomial superset must produce
nonnegative integral coefficient traces through power fifteen, the
correct square trace, regular-trace inequalities, and a positive
return by \(L=12\).  The same polynomial choices are used throughout.
This retains all Schur indices.  The common certificate theorem then
applies to exact moments, dimension squares and eligible remainders.

\paragraph{An explicit residual dimension-square contradiction.}
As an illustration, zero-based assignment $1987$ has rank twenty
and groups
\[
 (900,3,1),\ (900,1,2),\ (675,1,1),\ (60,2,1),
 (50,1,1),\ (36,1,1),\ (18,1,1).
\]
Its $198$ aggregate rows reduce to $37$ after the word tests.
The terminal branch has seven source dual pairs,
each of dimension sixteen, and five nonunit selfdual labels with
dimensions in
\[
       \{14,15,16,17,18,19,20,29,31,32,33,35,48\}.
\]
Their five squared dimensions would have to sum to
$5399-14\cdot16^2=1815$.

No dimension at least $33$ is possible, because
$33^2+4\cdot14^2>1815$, and at most one dimension can be at
least $29$. If all five dimensions are at most twenty, reduction
modulo four forces exactly three odd dimensions. The largest
possible squared sum is then
$3\cdot19^2+2\cdot20^2=1883$. Any change preserving these
parities decreases it by at least $72$, whereas the required
decrease is $68$.

If one dimension is $32$, the four others have squared sum
$791=4\cdot14^2+7$, impossible since the first positive
increment is $15^2-14^2=29$. If one is $31$, the corresponding
increment is $70$, which is not a sum of at most four elements
of $\{29,60\}$; every other positive increment is already
larger than $70$. Finally, a dimension $29$ leaves increment
$190$ above $4\cdot14^2$. Modulo four exactly two of the four
remaining dimensions are odd. Their two increments have sum
$58$, $122$, or $186$ below $190$, whereas the two even
increments have sum in
$\{0,60,120,128,188,256\}$. None is the needed complement
$132$, $68$, or $4$. This excludes the last branch without
an optimizer or an integer separating vector.

\begin{conditionalproposition}[Complete higher-rank exclusion at length twelve]
\label{paper:canonical12-higher-rank-exclusion}
Assume verification hypothesis $\mathsf V$ of
Section~\ref{paper:verification-status} for the finite scopes and
prerequisite exclusions specified in the proof below.
Every one of the \(9\,379\) higher-rank assignments in
Proposition~\ref{paper:canonical12-complete-universe} is impossible.
Without using the separate rank-sixteen exclusion, this proves that
\(L=12\) would force \(r=16\) and
\(R_{\mathbb C}\simeq\mathbb C^8\oplus M_2(\mathbb C)^{\oplus2}\).
\end{conditionalproposition}
\begin{proof}
The complete ledger closes \(2\,753\) assignments by rational odd-matrix
forms, \(4\,173\) by bounded or full canonical remainders, and
\(2\,399\) by moment or square certificates before full words.
The last \(54\) assignments use the stated word refinement:
\(53\) then have missing regular canonical remainders, and the
last is the explicit square contradiction above.
The checker \path{verify_canonical12.py} independently regenerates
the complete canonical and scalar-field input, the rank-sixteen match,
and every necessary higher-rank row used by a witness.  The exact
replay requires every one of the \(9\,379\) expected identifiers and
closes every branch.  All certificates are instances of the proved
conditions, so the higher-rank part is empty.  The preliminary tests
exclude smaller ranks, and \(r-L=4\) then gives the asserted
two-matrix algebra at rank sixteen.
\end{proof}

\section{Canonical-length thirteen and the rank-sixteen catalogue}

The catalogues and screening counts in this section are conditional on
$\mathsf V$ for their stated field, row and identifier scopes, including
the prerequisite source-rank exclusions. The structural lemmas and the
explicit thirteen-row exclusion below retain their stated independent
proofs. In particular the latter proves one specified assignment
impossible without assuming the completeness of a global inventory.

\subsection{The complete canonical-length-thirteen catalogue}
\label{paper:sec:canonical13-complete}
\begin{conditionaltheorem}[Rank and algebra at canonical length thirteen]
\label{paper:source:length-thirteen}
\label{paper:canonical13-high-tail}
\label{paper:canonical13-final-ranks-shapes}
Assume verification hypothesis $\mathsf V$ of
Section~\ref{paper:verification-status} for the finite scopes and
prerequisite exclusions specified in the proof below.
If \(L=13\), then
\[
 r\in\{17,19,21,23,25,27,29\},\qquad m_E\le4.
\]
The complete stated necessary tests retain \(146\,923\) multiplicity
assignments on \(25\,477\) weighted canonical profiles, distributed as
\[
\begin{array}{c|rrrrrrr}
r&17&19&21&23&25&27&29\\\hline
\text{assignments}&52\,481&21\,423&34\,731&10\,071&15\,450&10\,984&1\,783.
\end{array}
\]
Writing \(a,b,c\) for the numbers of matrix components of degrees
two, three and four, the possible complex algebras are exactly the
fifteen shapes specified by
\[
 2a+6b+12c=r-13,\qquad
 R_{\mathbb C}\simeq
 \mathbb C^{\,13-2a-3b-4c}\oplus M_2(\mathbb C)^a
 \oplus M_3(\mathbb C)^b\oplus M_4(\mathbb C)^c,
\]
with a positive scalar exponent, except for
\(\mathbb C\oplus M_2(\mathbb C)^6\).  These are necessary
assignments and algebra shapes; existence is not asserted.
\end{conditionaltheorem}
\begin{proof}
The direct weighted-entry recursion and an independent count-vector
enumeration both give \(62\,776\) profiles and \(807\,656\) assignments.
The rank, determinant, forbidden-pattern, rational-residue and capacity
conditions retain \(185\,421\).  The full-conductor catalogue includes
all allowed degrees five, six and eight; its \(322\) determinant and
\(154\) scalar catalogues are independently regenerated.

The complete tail of \(9\,438\) assignments of rank at least \(31\)
has exact certificates from bounded or full-row canonical remainders,
odd-matrix trace forms, central-atom geometry, complete integral central
spectra, moment separators, and positive words.  Each witness regenerates
its entire necessary row or spectral set; the positive-return cutoff is
the full length \(13\).  Thus \(r\le29\).  Since \(r\equiv L\pmod2\),
and the endpoint exclusions assumed in $\mathsf V$ remove \(13\) and \(15\), the
displayed ranks exhaust the possibilities.

Over all ranks, the rational odd-matrix gate excludes \(32\,180\)
assignments.  It is applied only where singleton nonlinear \((a,m)\)
groups force rationality; the \(528\) repeated-group inputs are
retained outside that test.  The exact overlap with the high-rank tail
is \(3\,120\), so identifier subtraction gives
\[
185\,421-(32\,180+9\,438-3\,120)=146\,923.
\]
The streamed subtraction yields the stated profile and rank counts;
the complete assignment and shape ledgers are in the certificate data.
Finally \(\sum_E m_E(m_E-1)=r-13\le16\) excludes \(m_E\ge5\).
Enumerating the three remaining matrix degrees gives the displayed
shape equation.  Its sole extra shape
\(\mathbb C\oplus M_2(\mathbb C)^6\) would imply
\(5400=1+2\sum_Ea_E\), contradicting parity.
\end{proof}

\subsection{Complete rank-sixteen one-matrix input and rows}
\label{paper:sec:one-matrix16-input}
Assuming the exclusions of Conditional
Theorem~\ref{paper:source:global-bounds}, at noncommutative rank sixteen,
\(R_{\mathbb C}\simeq\mathbb C^{12}\oplus M_2(\mathbb C)\)
and \(L=14\).  The unique matrix character is rational, with every
compatible quaternion alternative retained.  The thirteen nonunit
weighted entries sum to \(5399\), with the matrix canonical degree
occurring twice.  The common profile recursion agrees with the
coefficient of \(x^{13}y^{5399}\) in
\(\prod_{a\in\mathcal A}(1-xy^a)^{-1}\); placements at a doubled
degree \(a\) are independently checked by the coefficient of
\(x^{11}y^{5399-2a}\).

The full linear fields have degrees \(1,2,3,4,5,6,8\), including
the real quintic of conductor \(25\), the eligible sextic fields
(including those represented at conductor \(72\)), and the eligible
octic fields.  Integral covariance witnesses and positive rational
whole/fixed/skew forms use their complete full-field data.  Exact
Hasse-signature composition is checked against direct rational
diagonalization for every distinct encountered signature.  Its finite
support is verified field by field, without an assumption about
intermediate diagonalizations.

The row domain is the common model with \(15-2b\) selfdual atoms,
\(b=q+1\) or \(q+3\) pairs, and dimensions \(3\le d\le72\)
(\(3\le d\le51\) for pair representatives).  Retain every integral
field entry and test the rational matrix determinant against the square
identity, quaternion class and integral skew diagonal.  The regular
character has matrix multiplicity two and source rank sixteen; mixed
words are tested through exponent fourteen.  Undetermined higher-field
word traces are retained conservatively.  Cauchy--Schwarz on the
remaining blocks gives the optional exact bound
\[
 (d+as)^2\le(5399-ak)(5400-d^2-at)
\]
for a selected block option \((a,k;s,t)\).  The row-generation
algorithms are checked against the complete recursive enumeration.

Canonical remainders, regular support and forced central atoms use the
general results above.  A further dimension-density certificate is
particularly short: absence of every row with \(3\le d\le18\)
would give
\[
 5399=\sum_{X\ne\mathbf1}d_X^2\ge15\cdot19^2=5415.
\]
Every asserted absent dimension is exhaustively checked; a surviving
necessary row is only a non-exclusion.

\subsection{Symplectic tensor profiles as source-mass constraints}
Write $f=4k$ and $\chi(X)=2h_X$.  For the positive symplectic involution,
$h_X$ is an integer and every self-dual matrix image is scalar.  On the
$b=q+3$ source pairs the normalized pure commutator is an integral
decomposable alternating three-form $d_{ijk}$ with
\[
 \sum_{i<j<k}d_{ijk}^2=k,\qquad
 g_i=\sum_{j<\ell}d_{ij\ell}^2,\qquad \sum_i g_i=3k.
\]
A row with trace $2h$ has reduced norm $h^2+g$, and hence
\[
 (s,t)=(2h,2(h^2+g)).
\]
For a self-dual row $g=0$.  Character orthogonality, with the unit removed,
gives
\[
 \sum_{\mathrm{SD},\,X\ne\mathbf1}h_X^2
       +2\sum_{\text{source pairs}}h_X^2=2k-1.
\]
The strict single-row norm bound gives $h^2+g\le2k-1$.

For each category $(\mathrm{SD},|h|,g)$ the row engine proves a lower bound
on its source dimension by exhaustively testing all smaller dimensions.
For each complete slice profile it then minimizes the squared source mass
by an integer dynamic program with the preceding exact trace-norm budget.
A pair has weight two in both budgets.  If every slice profile has minimum
mass greater than $5399$, the full field-duality state is excluded.  The
replay regenerates the complete profile scope and every asserted lower
bound.  Categories not yet searched retain the elementary bound three.

For $b=3$ the only possible profile exists when $k$ is a square, and is
$(k,k,k)$.  For $b=4$, the Hodge dual is an arbitrary integral one-form:
if $k=z_1^2+z_2^2+z_3^2+z_4^2$, the slice profile is
$(k-z_1^2,\ldots,k-z_4^2)$, up to order.  Thus these two source-pair counts
can be handled for every formal codegree.  For $5\le b\le7$ the complete
finite tensor catalogue supplies the remaining stated scopes.  Each selected refinement covers its complete declared input scope;
overlaps are deduplicated by the full-field/duality identifier.

\subsection{The formal-codegree-eight Perron refinement}
For $f=8$ one has $k=2$ and the nonunit normalized trace-square budget is
three.  Hence the nonzero $h_X$ have one of precisely two duality patterns:
three self-dual labels with $|h|=1$, or one self-dual label and one dual pair
with $|h|=1$.  Perron orthogonality is
\[
 1+\sum_{X\ne\mathbf1}d_Xh_X=0.
\]
In the first pattern, if all three dimensions are at least $L\ge2$,
the least possible squared mass is
\[
 2L^2+(2L-1)^2.
\]
Indeed one positive dimension must be the sum of the other two, plus or
minus one; the smaller choice is obtained from two dimensions $L$ and a
third dimension $2L-1$.

In the second pattern let $p\ge P$ be the pair dimension and $r\ge L$ the
self-dual dimension.  Necessarily $r=2p+\delta$ with $\delta\in\{-1,1\}$.
Consequently its minimum mass is
\[
 \min_{\delta\in\{-1,1\}}
 \left((2p_\delta+\delta)^2+2p_\delta^2\right),\qquad
 p_\delta=\max\left(P,\left\lceil\frac{L-\delta}{2}\right\rceil\right).
\]
For every complete pure slice profile the remaining zero-trace labels
contribute their ordinary categorical lower bounds.  The nonzero-trace
pair is allowed to have $g=0$; it is not assumed noncentral.  Taking the
minimum over all profiles, both duality patterns, all positions of the
pair, and both signs supplies a valid lower bound on $5399$.

The bound for each individual row category records its actual searched
prefix.  A category whose search stopped early at a witness never receives
a later global search depth.  The verifier independently replays every
asserted missing lower row, checks the exact integer mass minimization,
and compares the closed formulas with direct finite enumeration of signs
and dimensions.

\subsection{The complete source-row screening scope}
\label{paper:sec:one-matrix16-symplectic-completion}
Every one of the \(67\,041\) symplectic full-field states is subjected
to the primary row tests.  The complete selected tensor refinements
cover every formal codegree for \(b=3,4\), and
\(f=36,40,60,72,100,108\) for \(5\le b\le7\); the Perron
refinement covers every retained state with \(f=8\).  The
dimension-three cubic-support theorem is checked on its exact
twenty-seven applicable inputs.  No additional numerical exclusion
is assigned to the dimension-four cubic theorem.

The exact row kernel checks every negative certificate and the
identifier union, leaving \(2\,065\) symplectic and \(974\,802\)
orthogonal states, or \(976\,867\) in total.  This is the complete
input to the local-corner and aligned-table arguments.  If a tensor
mass pattern uses only row categories whose minimum dimensions have
already been attained, it may safely establish a non-exclusion without
further search: that particular pattern's bound cannot rise.  Such a
pattern is not a source type or a multiplication certificate.
All negative lower-dimension claims exhaust the asserted finite ranges.

\subsection{The rank-sixteen one-matrix branch}
The rank-sixteen branch has $L=14$, so it is disjoint from the preceding
length-thirteen calculation.  Its canonical input has thirteen
nonunit weighted entries, with the matrix degree occurring twice.
A full state records this profile, the matrix placement, every full
linear character orbit, and the matrix duality type.  All admissible
rational quaternion alternatives remain attached to the state.

For symplectic duality there is a short common-table obstruction in
the fixed part of the source basis.

\begin{lemma}[Symplectic fixed-space pairing]
\label{paper:source:symplectic-fixed-pairing}
Suppose the unique nonlinear component is $M_2(\mathbb C)$, with
symplectic basis duality and rational formal codegree $F$. For a source
label $X$ put $s_X=\Tr E(X)$ and $X^+=(X+X^*)/2$. Let the linear Galois
orbits have fields $K_j$, codegrees $f_j$, and chosen character values
$\alpha_{j,X}\in K_j$. Then
\[
 \langle X^+,Y^+\rangle
 =\sum_j\frac{\Tr_{K_j/\mathbb Q}
  (\alpha_{j,X}\overline{\alpha_{j,Y}}
      +\alpha_{j,X}\alpha_{j,Y})}{2f_j}
   +\frac{s_Xs_Y}{2F}.
\]
The Perron character is included in the sum. Distinct source-duality
orbits give orthogonal nonzero vectors. A selfdual label has fixed
squared norm one, and a nonselfdual representative has fixed squared
norm one half.
\end{lemma}

\begin{proof}
On the matrix component, symplectic basis duality is
$A\mapsto\Tr(A)I-A$, so $E(X^+)=s_XI/2$. On a linear orbit it is complex
conjugation. Apply the coefficient-of-unit trace decomposition to
$X^+Y^+$ and use invariance of the field trace under conjugation. The
orthogonality and diagonal assertions also follow directly from the
orthonormal distinguished basis.
\end{proof}

In an exact row pool, identify rows with the same nonzero fixed vector
and join two classes when their fixed vectors are orthogonal.
Distinct source duality orbits must occupy distinct classes.  Thus
the required numbers of pair and selfdual classes must form one
orthogonal configuration; passing each row separately is insufficient.

\begin{conditionalproposition}[The proposed rank-sixteen catalogue]
\label{paper:source:rank-sixteen-catalogue}
Assume verification hypothesis $\mathsf V$ of
Section~\ref{paper:verification-status} for the finite scopes and
prerequisite exclusions specified in the proof below.
Under \eqref{paper:source:assumption}, every noncommutative rank-sixteen
source has one of the $600\,558$ full field-duality states retained by
the specified complete tests:
\[
 600\,558=599\,070\text{ orthogonal}+1\,488\text{ symplectic}.
\]
They lie on $96\,332$ weighted canonical profiles and $175\,512$ matrix
placements.  In the symplectic case the matrix codegree is neither
$60$ nor $120$.  The catalogue consists of necessary data only.
\end{conditionalproposition}
\begin{proof}
The complete specified field/form and source-row screens give the
\(976\,867\)-state set above.  The unramified corner criteria exclude
\(364\,918\), leaving \(610\,384\) orthogonal and \(1\,565\)
symplectic states.  Frobenius matching and singleton descent exclude
a further \(11\,387\), leaving \(599\,070+1\,492=600\,562\).
Eligibility is checked on each full field and every compatible
quaternion alternative.  Independent verifiers reconstruct the field
inertia and Frobenius actions and compare exact identifier sets.

Theorems~\ref{paper:symplectic-sixty-exclusion} and
\ref{paper:aligned-additional-exclusions} exclude four specified
symplectic assignments by complete fixed-space orthogonality:
states \(836343,47412,461210,1284815\), with matrix codegrees
\(60,120,120,40\), respectively.  Their full integral row domains
and graph obstructions are proved below.  All four identifiers belong
to the \(600\,562\)-state set and are distinct.  Exact streamed
subtraction therefore gives the asserted total, profile and placement
counts.  The complete input has only the displayed one codegree-sixty
and two codegree-120 symplectic assignments, so those branches are
closed.  The codegree-forty certificate excludes only its stated
assignment.
\end{proof}

At codegree six, the rational matrix factor is $B_{2,3}$ by the
rank-independent positive-cone theorem.  If the specified conductor-nine,
codegree-nine cubic orbit occurs, its distinguished triple has common dimension at least five under
$\mathsf V$ for Conditional Theorem~\ref{paper:source:small-cubic-triple}.

The certificate catalogue and the reconstruction theorem leave a clear
remaining problem.  A necessary field-duality state must first support
one common aligned character table.  Its reconstructed coefficients
must be nonnegative integers and satisfy the unit, duality, dimension,
reciprocity, and associativity identities, and the resulting based ring
must be simple.  An exact verifier checks these properties on a supplied
table.  Even such a table does not establish categorification or the
simple-center condition.  Alternatively, within a known modular category,
one must specify a connected Lagrangian algebra structure, not merely its
multiplicity vector; that algebra recovers the source category of modules
\cite[Section~4]{DMNO}.

The compact code package distinguishes freshly replayed bounded checks
from historical code and the large records required by $\mathsf V$.
An incomplete row pool, an unfinished alignment, or an unsupported
context receives no exclusion credit.

\subsection{Fixed-space orthogonality at symplectic codegree sixty}
\label{paper:sec:symplectic-sixty-fixed}

The local unramified tests do not apply to every state. The specified symplectic codegree-sixty state has, for each of $2,3,5$,
at least one linear field ramified at that prime. It is nevertheless
excluded by a small common-row obstruction.
This illustrates the distinction between individual moment rows and an
aligned character table.

\begin{theorem}[Exclusion of a specified symplectic codegree-sixty assignment]
\label{paper:symplectic-sixty-exclusion}
There is no integral source of dimension $5400$, with nonunit dimensions
at least three, whose unique nonlinear component is symplectic of
degree two and codegree $60$, and whose complete linear field/codegree
assignment is
\[
\begin{array}{c|rrrrrrr}
 f_j&5400&6&8&9&27&120&450\\\hline
 K_j&\mathbb Q&\mathbb Q,\mathbb Q&\mathbb Q(i)&
 \mathbb Q(\zeta_9+\zeta_9^{-1})&\mathbb Q&\mathbb Q&
 \mathbb Q(\sqrt{-15}).
\end{array}
\]
This excludes the displayed full-field assignment. The additional claim
that this is the only surviving symplectic codegree-sixty assignment at
rank sixteen requires $\mathsf V$ for the global inventory; it is not
part of this unconditional conclusion.
\end{theorem}
\begin{proof}
There are two conjugate pairs of linear characters, and the matrix
anti-fixed space has dimension three. Thus the source must have five
nonselfdual pairs. Let $X$ represent one such pair and write its
dimension as $d$. Necessarily $3\le d\le51$.

Write its two rational codegree-six values as $x,y$, its codegree-eight
value as $a+bi\in\mathbb Z[i]$, its real cubic value as $\beta$, and its
codegree-$27$ and codegree-$120$ values as $u,v$. Put
$w=\Tr_{\mathbb Q(\sqrt{-15})/\mathbb Q}\chi_{450}(X)$ and
$s=\Tr E(X)$. All displayed variables except $\beta$ are integers,
and $s$ is even by Proposition~\ref{paper:source:reduced-symplectic-tensor}.
Orthogonality to the unit and fixed squared norm one half give
\begin{align}
 d+900(x+y)+1350a+600\Tr(\beta)+200u+45v+12w+90s&=0,
 \label{paper:sixty-fixed-unit}\\
 2d^2+1800(x^2+y^2)+2700a^2+1200\Tr(\beta^2)
       +400u^2+90v^2+12w^2+90s^2&=5400.
 \label{paper:sixty-fixed-norm}
\end{align}

Put $\theta=\zeta_9+\zeta_9^{-1}$. The integral basis
$(1,\theta,\theta^2)$ has trace Gram matrix
\[
 \begin{pmatrix}3&0&6\\0&6&-3\\6&-3&18\end{pmatrix}.
\]
For $\beta=A+B\theta+C\theta^2$,
\[
 \Tr(\beta^2)=3(A+2C)^2+6(B-C/2)^2+\tfrac92C^2.
\]
Equation~\eqref{paper:sixty-fixed-norm} gives $\Tr(\beta^2)\le4$,
so $\beta\in\{0,1,-1\}$. The integral-basis statement can also be
checked directly in the cyclotomic power basis: the conjugation-fixed
space has rank three, and the displayed basis has determinant $\pm1$
on the free coordinates $1,\zeta_9^4,\zeta_9^5$.

For the remaining five variables set
\[
 L=d+200u+45v+12w+90s,\qquad
 Q=2d^2+400u^2+90v^2+12w^2+90s^2.
\]
Cauchy--Schwarz gives
\[
 L^2\le225Q,\qquad
 225=\tfrac12+100+\tfrac{45}2+12+90.
\]
Equality forces $d=L/450$.
If $\beta=\pm1$, positivity forces $x=y=a=0$, while $Q=1800$ and
$|L|=1800$, a contradiction. Hence $\beta=0$.
If $a=\pm1$, then either $x=y=0$, which would require
$1350^2\le225\cdot2700$, or exactly one of $x,y$ is nonzero.
The latter case is possible in Cauchy--Schwarz only with opposite
signs and equality $450^2=225\cdot900$, which forces $|d|=1$.
Thus $a=0$. Finally $x+y=\pm2$ contradicts Cauchy--Schwarz, and
$x+y=\pm1$ forces equality $900^2=225\cdot3600$ and $|d|=2$.
It follows that $y=-x$ and $x\in\{-1,0,1\}$.

The necessary fixed rows therefore satisfy
\[
 d+200u+45v+12w+90s=0,\qquad
 2d^2+400u^2+90v^2+12w^2+90s^2=R,
\]
where $R=5400-3600x^2$, $3\le d\le51$, and $s$ is even.
This small integer system has precisely the following thirteen solutions:
\[
\begin{array}{r|rrrrrrr}
 i&d&s&x&y&u&v&w\\\hline
 1&8&-2&0&0&-1&4&16\\
 2&8&2&0&0&-1&-4&16\\
 3&10&2&0&0&1&-6&-10\\
 4&17&0&0&0&-1&-1&19\\
 5&18&-2&-1&1&0&2&6\\
 6&18&-2&1&-1&0&2&6\\
 7&19&-4&0&0&1&5&-7\\
 8&34&2&0&0&-2&2&8\\
 9&35&0&0&0&-1&5&-5\\
 10&35&4&0&0&-1&-3&-5\\
 11&37&-4&0&0&1&3&-1\\
 12&37&0&0&0&1&-5&-1\\
 13&46&-2&0&0&1&-2&2
\end{array}
\]
For completeness, this enumeration has explicit tiny bounds. For each
$x$, take $|u|\le\lfloor\sqrt{R/400}\rfloor$,
$|v|\le\lfloor\sqrt{R/90}\rfloor$, and even
$|s|\le\lfloor\sqrt{R/90}\rfloor$. There are $1185$ triples in all.
Writing $A=200u+45v+90s$ and $B=400u^2+90v^2+90s^2$, eliminate $d$.
The possible integer $w$ are exactly
\[
 w=\frac{-48A\pm\sqrt{1200R-96A^2-1200B}}{600},\qquad d=-A-12w.
\]
Testing square discriminants, divisibility by $600$, and the displayed
dimension interval gives the table. A separate certificate enumerates
the bounded $d,u,v,s$ box directly and obtains the same thirteen rows.

By Lemma~\ref{paper:source:symplectic-fixed-pairing}, two distinct duality
pairs must have zero inner product for the positive diagonal form
\[
 \operatorname{diag}(2,90,1800,1800,400,90,12)
\]
on the listed coordinates. Every row has squared norm $5400$.
The only orthogonal pairs of distinct rows are
\[
 \{4,9\},\quad\{7,8\},\quad\{9,12\},\quad\{10,11\}.
\]
This graph has no triangle. In particular it cannot supply the five
mutually orthogonal nonzero fixed vectors required by the five source
duality pairs. Repeated rows are also impossible, since their squared
norm is positive. This proves the exclusion.

The exclusion just proved uses no global field inventory. If
$\mathsf V$ verifies that the rank-sixteen inventory has exactly this
one surviving symplectic codegree-sixty assignment (state $836343$),
then the global branch is also closed. This last implication is
conditional, and is not established by the thirteen-row calculation.
\end{proof}

The historical full integral-basis enumeration records a compatible
certificate, whose complete replay is not claimed for this compact release.
It lifts all $66$ necessary moment rows to $99$ exact algebraic rows:
$65$ selfdual and $34$ nonselfdual. The latter give eleven fixed
classes, contained in the thirteen-row superset above; the two additional
dimension-eight rows are isolated in the larger graph. An independent
Galois-embedding trace calculation reproduces the same four edges.
The short proof above does not depend on completeness of the
row-lifting implementation.

\subsection{The codegree-120 endpoints and a codegree-40 assignment}
\label{paper:sec:aligned-additional-exclusions}

The fixed-space obstruction also excludes the two remaining symplectic
codegree-$120$ assignments and one of the two remaining codegree-$40$
assignments. Every exclusion is checked by the same independent
integral-basis and Galois-embedding procedure.

Put $K_3=\mathbb Q(\sqrt{-3})$, $K_4=\mathbb Q(i)$,
$K_{15}=\mathbb Q(\sqrt{-15})$, and
$K_9^+=\mathbb Q(\zeta_9+\zeta_9^{-1})$.
In the following table, $(K,f)^m$ denotes $m$ distinct linear Galois
orbits with field $K$ and formal codegree $f$.  Every row also contains
the Perron orbit $(\mathbb Q,5400)$.
\begin{center}
\begin{tabular}{rcl}
State&$F$&Non-Perron linear orbits\\\hline
$47412$&$120$&
$(\mathbb Q,6)^3,(K_3,6),(\mathbb Q,24),(\mathbb Q,25),$\\
&&$(\mathbb Q,27),(\mathbb Q,50),(K_3,180)$\\[2pt]
$461210$&$120$&
$(\mathbb Q,6)^3,(\mathbb Q,8),(K_9^+,9),$\\
&&$(\mathbb Q,90),(\mathbb Q,108),(K_{15},450)$\\[2pt]
$1284815$&$40$&
$(\mathbb Q,8),(K_4,8),(\mathbb Q,9)^2,$\\
&&$(K_9^+,9),(K_3,135),(\mathbb Q,225)$
\end{tabular}
\end{center}

\begin{conditionaltheorem}[Further aligned-row exclusions]
\label{paper:aligned-additional-exclusions}
Assume verification hypothesis $\mathsf V$ of
Section~\ref{paper:verification-status} for the finite scopes and
prerequisite exclusions specified in the proof below.
Under the standing source hypotheses at rank sixteen, none of the
three full-field assignments above can occur with symplectic matrix
duality.  Consequently the symplectic codegree-$120$ branch is closed.
Together with Theorem~\ref{paper:symplectic-sixty-exclusion}, both the
codegree-$60$ and codegree-$120$ symplectic branches are closed.
\end{conditionaltheorem}
\begin{proof}
The construction and completeness checks are finite and exact.
The stated moment-row kernel gives the necessary dimension, trace,
and norm rows.  For each linear field, every algebraic integer with
the required trace and positive norm is then enumerated in a full
integral basis.  If $G$ is its positive trace Gram matrix and the norm
bound is $b$, the coordinate bound
$|a_i|^2\le b(G^{-1})_{ii}$ makes this enumeration exhaustive.
For a non-Perron orbit of codegree $f_j$, take $b=f_j-1$:
the positive Perron contribution makes the integral positive norm
strictly smaller than $f_j$.
The bases are $1$ over $\mathbb Q$, the usual quadratic integer bases,
and $(1,\theta,\theta^2)$ for $K_9^+$, whose integral-basis
verification is in the proof of Theorem~\ref{paper:symplectic-sixty-exclusion}.
All individual nonunit dimensions are included: $3\le d\le72$ for
selfdual labels and $3\le d\le51$ for nonselfdual representatives.
No condition $d\mid5400$ is imposed.

The complete numbers of necessary rows are
\[
\begin{array}{r|rrrrr}
\text{state}&\text{moment rows}&\text{exact rows}&
\text{selfdual}&\text{nonselfdual}&\text{pair fixed classes}\\\hline
47412&205&369&260&109&25\\
461210&158&340&123&217&44\\
1284815&87&145&78&67&12.
\end{array}
\]
Here a pair fixed class identifies precisely equal nonzero fixed
vectors, whose representatives differ only in their possible
anti-fixed coordinates.  Distinct source-duality pairs cannot use
the same fixed class, by positive definiteness.

Use Lemma~\ref{paper:source:symplectic-fixed-pairing} to form the orthogonality
graph on these pair fixed classes.  In state $47412$ this graph has
$25$ vertices, $26$ edges, eight triangles, and no four-clique,
whereas five mutually orthogonal pair vectors are required.
In state $1284815$ it has $12$ vertices and five edges, with no
triangle, whereas again five pair vectors are required.
Both cases are therefore excluded.

State $461210$ requires four pair vectors and seven nonunit selfdual
vectors.  Its pair graph has $44$ vertices, $126$ edges, $138$
triangles, and exactly $64$ four-cliques.  For $22$ of these cliques,
only six selfdual fixed classes are orthogonal to all four pair
vectors, which is insufficient.  Each of the other $42$ cliques has
eight common selfdual classes.  The $42\binom87=336$ possible
seven-element subsets are all checked; every one contains a pair
with nonzero inner product.  Thus no full orthogonal fixed table
exists in this case either.

The independent row audit uses the preserved original moment kernel
and separate integral-basis enumeration.  The independent graph
audit computes traces by summing all distinct Galois embeddings and
enumerates pair cliques first, followed by all required selfdual
subsets.  It uses neither the optimized search tree nor its
dimension, trace-square, or correlation bounds.  The complete
records include a nonorthogonality witness for every rejected
selfdual subset.

Finally, the authenticated input inventory has exactly two
symplectic codegree-$120$ states, namely $47412$ and $461210$.
This proves the asserted branch conclusion.  The codegree-$40$
inventory had two states, so the last exclusion closes only the
specified one; the other remains in the necessary catalogue.
\end{proof}

\begin{remark}[Exact scope of the aligned-row certificates]
The four assignments, including state $836343$, are full-field
necessary states, not source dimension types. The thirteen-row exclusion
of state $836343$ is independent; the other three exclusions and the
asserted disjointness from the local-corner and Frobenius stages require
$\mathsf V$, including exact membership in the corresponding remainder.  They do not imply a global rank-seventeen bound, and do
not construct a fusion ring or decide categorifiability of a
surviving state.  An interrupted or incomplete row search receives
no exclusion credit.
\end{remark}


\section{Remaining questions}\label{paper:questions}

The family theorem reduces the unresolved even problem to dimension
$5400^2$. A sufficient final step would show that every integral,
perfect modular category of that dimension has a nontrivial proper
fusion subcategory. The coprime splitting criterion gives such a
subcategory when the complementary principal intersection is the unit.
The general surviving intersection need not have that form, and the
homogeneity or isolated-fiber assumptions used in the stronger block
results have not been proved for every remaining configuration.

The odd frontier is different in nature. The parameter $(5,7,23)$
satisfies the stated uniform arithmetic conditions, but the compatible
cyclic mask is not modular data. Further progress requires either a
new obstruction for this parameter or additional categorical structure
that strengthens the support criteria. Dense difference envelopes can
be stable under the support operation without providing a fusion ring;
Proposition~\ref{paper:dense-envelope-barrier} gives a precise limitation.
Likewise, a positive spectral projection need not have nonnegative
coefficients in the distinguished basis. These distinctions explain why
local arithmetic consistency alone does not settle the categorical
problem.

Under the additional center hypothesis and the full verification
hypothesis $\mathsf V$, the proposed first open source ranks are
fourteen in the commutative case and sixteen in the noncommutative case.
Their status is conditional in this compact release. At the latter rank,
the resulting unique complex algebra shape allows
reconstruction from aligned linear characters and one integral
commutator tensor. The Cayley--Hamilton tests can be applied earlier,
to individual rows and their pairs or triples. A remaining field
assignment must still admit one common character table, compatible
basis duality, and a nonnegative integral multiplication satisfying
all fusion-ring axioms. The conditional finite catalogues record proposed necessary data and
exclusions; neither they nor the bounded checks construct such a ring.
A resulting fusion ring would in turn require categorification and
verification of the prescribed center properties.


\begin{thebibliography}{99}
\small
\setlength{\itemsep}{1pt}
\bibitem{BH}
M.~Beck and S.~Ho\c{s}ten, \href{https://doi.org/10.4310/mrl.2006.v13.n4.a10}{\emph{Cyclotomic polytopes and growth series of cyclotomic lattices}}, Math. Res. Lett. \textbf{13} (2006), 607--622.

\bibitem{BGNPWSuper}
P.~Bruillard, C.~Galindo, S.-H.~Ng, J.~Y.~Plavnik, E.~C.~Rowell and Z.~Wang,
\href{https://doi.org/10.1007/s10468-019-09873-9}{\emph{Classification of super-modular categories by rank}},
Algebr. Represent. Theory \textbf{23} (2020), 795--809.

\bibitem{BNRWFinite}
P.~Bruillard, S.-H.~Ng, E.~C.~Rowell, and Z.~Wang,
\href{https://doi.org/10.1090/jams/842}{\emph{Rank-finiteness for modular categories}},
J. Amer. Math. Soc. \textbf{29} (2016), no.~3, 857--881.

\bibitem{BurciuSubalgebras2025}
S.~Burciu,
\href{https://doi.org/10.1007/s10468-025-10314-z}{\emph{Subalgebras
of \'{E}tale Algebras in Braided Fusion Categories}},
Algebr. Represent. Theory \textbf{28} (2025), 211--238.
The primary preprint is
\href{https://arxiv.org/abs/2105.11202}{arXiv:2105.11202}.

\bibitem{BDPfamily}
S.~Burciu, J.~Dong and S.~Palcoux,
\emph{Prime factors of fusion-category dimensions: weak
group-theoreticity, rank and decomposition}, preprint, 2026.
\url{https://hexagonmath.org/2611.00233v3}.

\bibitem{BP}
S.~Burciu and S.~Palcoux,
\href{https://doi.org/10.1142/S0219199726500070}{\emph{Burnside-type results for fusion categories}},
Commun. Contemp. Math. (2026), article 2650007.

\bibitem{ChlSchur}
M.~Chlouveraki, \href{https://doi.org/10.2140/ant.2008.2.689}{\emph{Rouquier blocks of the cyclotomic Ariki--Koike algebras}}, Algebra Number Theory \textbf{2} (2008), no.~6, 689--720.

\bibitem{DMNO}
A.~Davydov, M.~M\"uger, D.~Nikshych, and V.~Ostrik,
\href{https://doi.org/10.1515/crelle.2012.014}{\emph{The Witt group of non-degenerate braided fusion categories}},
J. Reine Angew. Math. \textbf{677} (2013), 135--177.

\bibitem{DLNCongruence}
C.~Dong, X.~Lin, and S.-H.~Ng,
\href{https://doi.org/10.2140/ant.2015.9.2121}{\emph{Congruence property in conformal field theory}},
Algebra \& Number Theory \textbf{9} (2015), no.~9, 2121--2166.

\bibitem{DongWeakIntegral}
J.~Dong,
\href{https://doi.org/10.1007/s11401-026-0019-6}{\emph{Classification of certain weakly integral fusion categories}},
Chinese Ann. Math. Ser. B \textbf{47} (2026), no.~3,
555--572.

\bibitem{DNVSmall}
J.~Dong, S.~Natale and L.~Vendramin,
\href{https://arxiv.org/abs/1209.1726}{\emph{Frobenius property
for fusion categories of small integral dimension}},
J. Algebra Appl. \textbf{14} (2015), no.~2, 1550011.

\bibitem{DPLCM}
J.~Dong and S.~Palcoux, \emph{A new criterion for integral modular
categorification}, \href{https://arxiv.org/abs/2510.10311}{arXiv:2510.10311}.

\bibitem{DGNO}
V.~Drinfeld, S.~Gelaki, D.~Nikshych, and V.~Ostrik,
\href{https://doi.org/10.1007/s00029-010-0017-z}{\emph{On braided fusion categories I}},
Selecta Math. (N.S.) \textbf{16} (2010), 1--119.

\bibitem{EGNO}
P.~Etingof, S.~Gelaki, D.~Nikshych, and V.~Ostrik,
\href{https://math.mit.edu/~etingof/egnobookfinal.pdf}{\emph{Tensor Categories}},
Mathematical Surveys and Monographs, vol.~205,
American Mathematical Society, 2015.

\bibitem{ENOannals}
P.~Etingof, D.~Nikshych, and V.~Ostrik,
\href{https://doi.org/10.4007/annals.2005.162.581}{\emph{On fusion categories}},
Ann. of Math. \textbf{162} (2005), no.~2, 581--642.

\bibitem{ENOsolvable}
P.~Etingof, D.~Nikshych, and V.~Ostrik,
\href{https://doi.org/10.1016/j.aim.2010.06.009}{\emph{Weakly group-theoretical and solvable fusion categories}},
Adv. Math. \textbf{226} (2011), no.~1, 176--205.

\bibitem{FerroniMinimal}
L.~Ferroni, \emph{On the Ehrhart Polynomial of Minimal Matroids},
Discrete \& Computational Geometry \textbf{68} (2022), 255--273,
Theorem~2.1 (the minimal-base theorem of Dinolt and Murty).
\url{https://doi.org/10.1007/s00454-021-00313-4}.

\bibitem{FRSEvenFusion}
J.~Fuchs, I.~Runkel and C.~Schweigert,
\href{https://doi.org/10.1088/0305-4470/35/19/102}{\emph{A reason for fusion rules to be even}},
J. Phys. A: Math. Gen. \textbf{35} (2002), no.~19, L255--L260.

\bibitem{GeerPatureau2018}
N.~Geer and B.~Patureau-Mirand,
\emph{The trace on projective representations of quantum groups},
Letters in Mathematical Physics \textbf{108} (2018), 117--140.
\href{https://doi.org/10.1007/s11005-017-0993-4}{doi:10.1007/s11005-017-0993-4}.
See Section~1.1; the author preprint is
\href{https://arxiv.org/abs/1610.09129}{arXiv:1610.09129}.

\bibitem{GNNcenters}
S.~Gelaki, D.~Naidu and D.~Nikshych,
\href{https://msp.org/ant/2009/3-8/ant-v3-n8-p05-s.pdf}{\emph{Centers of graded fusion categories}},
Algebra Number Theory \textbf{3} (2009), 959--990.

\bibitem{GNnil}
S.~Gelaki and D.~Nikshych, \href{https://doi.org/10.1016/j.aim.2007.08.001}{\emph{Nilpotent fusion categories}}, Adv. Math. \textbf{217} (2008), 1053--1071.

\bibitem{JFR}
T.~Johnson-Freyd and D.~Reutter,
\href{https://doi.org/10.1090/jams/1023}{\emph{Minimal nondegenerate extensions}},
J. Amer. Math. Soc. \textbf{37} (2024), 81--150.

\bibitem{MilneANT}
J.~S.~Milne, \emph{Algebraic Number Theory}, lecture notes (2020).
\url{https://www.jmilne.org/math/CourseNotes/ANT.pdf}.

\bibitem{NSTV}
S.~Nasseh, S.~Sather-Wagstaff, R.~Takahashi and K.~VandeBogert,
\href{https://doi.org/10.1016/j.jpaa.2018.06.006}{\emph{Applications and homological properties of local rings with
decomposable maximal ideals}}, J. Pure Appl. Algebra \textbf{223}
(2019), 1272--1287.

\bibitem{NataleWGT}
S.~Natale,
\href{https://doi.org/10.4171/jncg/177}{\emph{On weakly group-theoretical non-degenerate braided fusion categories}},
J. Noncommut. Geom. \textbf{8} (2014), 1043--1060.

\bibitem{NataleCore}
S.~Natale,
\href{https://doi.org/10.1142/S0129167X1850012X}{\emph{The core of a weakly group-theoretical braided fusion category}},
Internat. J. Math. \textbf{29} (2018), no.~2, 1850012.

\bibitem{NSfs}
S.-H.~Ng and P.~Schauenburg, \href{https://doi.org/10.1016/j.aim.2006.07.017}{\emph{Frobenius--Schur indicators and exponents of spherical categories}}, Adv. Math. \textbf{211} (2007), 34--71.

\bibitem{NSind}
S.-H.~Ng and P.~Schauenburg, \href{https://doi.org/10.1007/s00220-010-1096-6}{\emph{Congruence subgroups and generalized Frobenius--Schur indicators}}, Comm. Math. Phys. \textbf{300} (2010), 1--46.

\bibitem{OstrikFormalCodegrees}
V.~Ostrik,
\href{https://doi.org/10.4310/MRL.2009.v16.n5.a11}{\emph{On formal codegrees of fusion categories}},
Math. Res. Lett. \textbf{16} (2009), no.~5, 895--901.

\bibitem{OstrikRankThree}
V.~Ostrik,
\href{https://doi.org/10.17323/1609-4514-2015-15-2-373-396}{\emph{Pivotal
fusion categories of rank 3}},
Moscow Math. J. \textbf{15} (2015), no.~2, 373--396.
The primary preprint is
\href{https://arxiv.org/abs/1309.4822}{arXiv:1309.4822}.

\bibitem{Sharifi}
R.~Sharifi,
\href{https://www.math.ucla.edu/~sharifi/algnum.pdf}{\emph{Algebraic Number Theory}}, lecture notes, Chapters~3, 6, 7, and~9.

\bibitem{VoightQuaternion}
J.~Voight,
\emph{Quaternion Algebras},
Graduate Texts in Mathematics, vol.~288, Springer, 2021.
\href{https://doi.org/10.1007/978-3-030-56694-4}{doi:10.1007/978-3-030-56694-4}.

\bibitem{YuSlightlyDegenerate}
Z.~Yu,
\href{https://doi.org/10.1016/j.jalgebra.2020.04.022}{\emph{On slightly degenerate fusion categories}},
J. Algebra \textbf{559} (2020), 408--431.

\end{thebibliography}
\end{document}
